% Auto-generated from data/segments.json + layout.json (+ transforms) — do not edit by hand.
% volume: 03Vin03
% mode: printing
% segments: 2834
% transform_rules: 126
% toc_marks: 582

% Volume layout defaults (layout.json ``layout``).
\csromanlayoutapply{1}{1.5}{21.6pt}{8pt}{8pt}{65pt}{2.5em}%

\pitaka{\textbf{วินยปิฏก}}
\csromanmatikaanchor{mtk.0}
\csromantocmarknopage{chapter}{มหาวคฺคปาฬิ}{mtk.0}
\bookmark[dest={mtk.0},level=0]{มหาวคฺคปาฬิ}
\gambhira{\textbf{มหาวคฺคปาฬิ}}
\namakkaram{นโม ตสฺส ภควโต อรหโต สมฺมา\-สมฺพุทฺธสฺส.}
\csromanmatikaanchor{mtk.1}
\csromantocmarknopage{chapter}{1. มหาขนฺธก}{mtk.1}
\bookmark[dest={mtk.1},level=0]{1. มหาขนฺธก}
\chapterhead{1. \textbf{มหาขนฺธก}}
\csromanmatikaanchor{mtk.2}
\csromantocmarkref{section}{1. โพธิกถา}{mtk.2}
\bookmark[dest={mtk.2},level=1]{1. โพธิกถา}
\csromanheader{1}{1. \textbf{โพธิกถา}}
\proseitem{1}{\csromanfolio{1}\csromansymbolfootnote{*}{ขุ 1. 77 ปิฏฺฐาทีสุ อุทาเนปิ.}เตน สมเยน พุทฺโธ ภควา อุรุเวลายํ วิหรติ นชฺชา เนรญฺชราย ตีเร โพธิรุกฺข\-มูเล ปฐมา\-ภิสมฺพุทฺโธ. อถ โข ภควา โพธิรุกฺข\-มูเล สตฺตาหํ เอกปลฺลงฺเกน นิสีทิ วิมุตฺติสุข\-ปฏิสํเวที\csromansharedfootnote{0083ed4ed8a320ba}{วิมุตฺติสุขํ ปฏิสํเวที (ก)}. อถ โข ภควา รตฺติยา ปฐมํ ยามํ ปฏิจฺจ\-สมุปฺปาทํ อนุโลมปฏิโลมํ มนสากาสิ “อวิชฺชาปจฺจยา สงฺขารา, สงฺขาร\-ปจฺจยา วิญฺญาณํ, วิญฺญาณปจฺจยา นามรูปํ, นามรูป\-ปจฺจยา สฬายตนํ, สฬายตน\-ปจฺจยา ผสฺโส, ผสฺสปจฺจยา เวทนา, เวทนาปจฺจยา ตณฺหา, ตณฺหาปจฺจยา อุปาทานํ, อุปาทานปจฺจยา ภโว, ภวปจฺจยา ชาติ, ชาติปจฺจยา ชรามรณํ โสก\-ปริเทว\-ทุกฺข\-โทมนสฺสุ\-ปายาสา สมฺภวนฺติ, เอวเมตสฺส เกวลสฺส ทุกฺข\-กฺขนฺธสฺส สมุทโย โหติ. อวิชฺชายเตฺวว อเสส\-วิราค\-นิโรธา สงฺขาร\-นิโรโธ, สงฺขาร\-นิโรธา วิญฺญาณนิโรโธ, วิญฺญาณนิโรธา นามรูป\-นิโรโธ, นามรูป\-นิโรธา สฬายตน\-นิโรโธ, สฬายตน\-นิโรธา ผสฺสนิโรโธ, ผสฺสนิโรธา เวทนานิโรโธ, เวทนานิโรธา ตณฺหานิโรโธ, ตณฺหานิโรธา อุปาทานนิโรโธ, อุปาทานนิโรธา ภวนิโรโธ, ภวนิโรธา ชาตินิโรโธ, ชาตินิโรธา ชรามรณํ โสก\-ปริเทว\-ทุกฺข\-โทมนสฺสุ\-ปายาสา นิรุชฺฌนฺติ, เอวเมตสฺส เกวลสฺส ทุกฺข\-กฺขนฺธสฺส นิโรโธ โหตี”ติ.}

\prose{\csromanfolio{2}อถ โข ภควา เอตมตฺถํ วิทิตฺวา ตายํ เวลายํ อิมํ อุทานํ อุทาเนสิ\csromandash{}}

\setlength{\csromangathapagewidth}{0pt}
\setlength{\csromangathawakwidth}{0pt}
\csromangathameasuregroup{}{“ยทา หเว ปาตุภวนฺติ ธมฺมา, \\ อาตาปิโน ฌายโต พฺราหฺมณสฺส. \\ อถสฺส กงฺขา วปยนฺติ สพฺพา, \\ ยโต ปชานาติ สเหตุธมฺมนฺ”ติ.}
\csromangathameasurewak{“ยทา หเว ปาตุภวนฺติ ธมฺมา, \\ อาตาปิโน ฌายโต พฺราหฺมณสฺส. \\ อถสฺส กงฺขา วปยนฺติ สพฺพา, \\ ยโต ปชานาติ สเหตุธมฺมนฺ”ติ.}
\setlength{\csromangathaleftcol}{0pt}
\csromangathagroup{\csromangathastanza{\csromangathawak{“ยทา หเว ปาตุภวนฺติ ธมฺมา, \\ อาตาปิโน ฌายโต พฺราหฺมณสฺส. \\ อถสฺส กงฺขา วปยนฺติ สพฺพา, \\ ยโต ปชานาติ สเหตุธมฺมนฺ”ติ.}}}

\proseitem{2}{\csromansymbolfootnote{*}{ขุ 1. 78 ปิฏฺเฐ อุทาเนปิ.}อถ โข ภควา รตฺติยา มชฺฌิมํ ยามํ ปฏิจฺจ\-สมุปฺปาทํ อนุโลมปฏิโลมํ มนสากาสิ “อวิชฺชาปจฺจยา สงฺขารา, สงฺขาร\-ปจฺจยา วิญฺญาณํ, วิญฺญาณปจฺจยา นามรูปํ ฯเปฯ เอวเมตสฺส เกวลสฺส ทุกฺข\-กฺขนฺธสฺส สมุทโย โหติ ฯเปฯ นิโรโธ โหตี”ติ.}

\prose{อถ โข ภควา เอตมตฺถํ วิทิตฺวา ตายํ เวลายํ อิมํ อุทานํ อุทาเนสิ\csromandash{}}

\setlength{\csromangathapagewidth}{0pt}
\setlength{\csromangathawakwidth}{0pt}
\csromangathameasuregroup{}{“ยทา หเว ปาตุภวนฺติ ธมฺมา, \\ อาตาปิโน ฌายโต พฺราหฺมณสฺส. \\ อถสฺส กงฺขา วปยนฺติ สพฺพา, \\ ยโต ขยํ ปจฺจยานํ อเวที”ติ.}
\csromangathameasurewak{“ยทา หเว ปาตุภวนฺติ ธมฺมา, \\ อาตาปิโน ฌายโต พฺราหฺมณสฺส. \\ อถสฺส กงฺขา วปยนฺติ สพฺพา, \\ ยโต ขยํ ปจฺจยานํ อเวที”ติ.}
\setlength{\csromangathaleftcol}{0pt}
\csromangathagroup{\csromangathastanza{\csromangathawak{“ยทา หเว ปาตุภวนฺติ ธมฺมา, \\ อาตาปิโน ฌายโต พฺราหฺมณสฺส. \\ อถสฺส กงฺขา วปยนฺติ สพฺพา, \\ ยโต ขยํ ปจฺจยานํ อเวที”ติ.}}}

\proseitem{3}{\csromansymbolfootnote{+}{ขุ 1. 78-79 ปิฏฺเฐสุปิ.}อถ โข ภควา รตฺติยา ปจฺฉิมํ ยามํ ปฏิจฺจ\-สมุปฺปาทํ อนุโลมปฏิโลมํ มนสากาสิ “อวิชฺชาปจฺจยา สงฺขารา, สงฺขาร\-ปจฺจยา วิญฺญาณํ, วิญฺญาณปจฺจยา นามรูปํ ฯเปฯ เอวเมตสฺส เกวลสฺส ทุกฺข\-กฺขนฺธสฺส สมุทโย โหติ ฯเปฯ นิโรโธ โหตี”ติ.}

\prose{อถ โข ภควา เอตมตฺถํ วิทิตฺวา ตายํ เวลายํ อิมํ อุทานํ อุทาเนสิ\csromandash{}}

\setlength{\csromangathapagewidth}{0pt}
\setlength{\csromangathawakwidth}{0pt}
\csromangathameasuregroup{}{“ยทา หเว ปาตุภวนฺติ ธมฺมา, \\ อาตาปิโน ฌายโต พฺราหฺมณสฺส. \\ วิธูปยํ ติฏฺฐติ มารเสนํ, \\ สูริโยว\textsuperscript{1} โอภาสยมนฺตลิกฺขนฺ”ติ.}
\csromangathameasurewak{“ยทา หเว ปาตุภวนฺติ ธมฺมา, \\ อาตาปิโน ฌายโต พฺราหฺมณสฺส. \\ วิธูปยํ ติฏฺฐติ มารเสนํ, \\ สูริโยว\textsuperscript{1} โอภาสยมนฺตลิกฺขนฺ”ติ.}
\setlength{\csromangathaleftcol}{0pt}
\csromangathagroupwithcloser{\csromangathastanza{\csromangathawak{“ยทา หเว ปาตุภวนฺติ ธมฺมา, \\ อาตาปิโน ฌายโต พฺราหฺมณสฺส. \\ วิธูปยํ ติฏฺฐติ มารเสนํ, \\ สูริโยว\csromansharedfootnote{87243025fe8e079f}{สุริโยว (สี, สฺยา, กํ)} โอภาสยมนฺตลิกฺขนฺ”ติ.}}}{โพธิกถา นิฏฺฐิตา.}
\csromanmatikaanchor{mtk.3}
\csromantocmarkref{section}{2. อชปาลกถา}{mtk.3}
\bookmark[dest={mtk.3},level=1]{2. อชปาลกถา}
\csromanheader{1}{2. \textbf{อชปาลกถา}}
\proseitem{4}{\csromanfolio{3}\csromansymbolfootnote{*}{ขุ 1. 80 ปิฏฺเฐ อุทาเนปิ.}อถ โข ภควา สตฺตาหสฺส อจฺจเยน ตมฺหา สมาธิมฺหา วุฏฺฐหิตฺวา โพธิรุกฺข\-มูลา เยน อชปาล\-นิโคฺรโธ เตนุปสงฺกมิ, อุปสงฺกมิตฺวา อชปาล\-นิโคฺรธมูเล สตฺตาหํ เอกปลฺลงฺเกน นิสีทิ วิมุตฺติสุข\-ปฏิสํเวที. อถ โข อญฺญตโร หุํหุงฺกชาติโก พฺราหฺมโณ เยน ภควา เตนุปสงฺกมิ, อุปสงฺกมิตฺวา ภควตา สทฺธิํ สมฺโมทิ, สมฺโมทนียํ กถํ สารณียํ วีติสาเรตฺวา เอกมนฺตํ อฏฺฐาสิ, เอกมนฺตํ ฐิโต โข โส พฺราหฺมโณ ภควนฺตํ เอตทโวจ “กิตฺตาวตา นุ โข โภ โคตม พฺราหฺมโณ โหติ, กตเม จ ปน พฺราหฺมณกรณา\csromansharedfootnote{b1fa353c589dc398}{พฺราหฺมณการกา (ก) พฺราหฺมณกราณา (?)} ธมฺมา”ติ. อถ โข ภควา เอตมตฺถํ วิทิตฺวา ตายํ เวลายํ อิมํ อุทานํ อุทาเนสิ\csromandash{}}

\setlength{\csromangathapagewidth}{0pt}
\setlength{\csromangathawakwidth}{0pt}
\csromangathameasuregroup{}{\textsuperscript{+}“โย พฺราหฺมโณ พาหิตปาปธมฺโม, \\ นิหุํหุงฺโก นิกฺกสาโว ยตตฺโต. \\ เวทนฺตคู วุสิตพฺรหฺมจริโย, \\ ธมฺเมน โส พฺรหฺมวาทํ วเทยฺย. \\ ยสฺสุสฺสทา นตฺถิ กุหิญฺจิ โลเก”ติ.}
\csromangathameasurewak{\textsuperscript{+}“โย พฺราหฺมโณ พาหิตปาปธมฺโม, \\ นิหุํหุงฺโก นิกฺกสาโว ยตตฺโต. \\ เวทนฺตคู วุสิตพฺรหฺมจริโย, \\ ธมฺเมน โส พฺรหฺมวาทํ วเทยฺย. \\ ยสฺสุสฺสทา นตฺถิ กุหิญฺจิ โลเก”ติ.}
\setlength{\csromangathaleftcol}{0pt}
\csromangathagroupwithcloserruled{\csromangathastanza{\csromangathawak{\csromansymbolfootnote{+}{ขุ 10. 129 ปิฏฺเฐ เนตฺติปาฬิยมฺปิ. ++ ขุ 1. 87, 88 ปิฏฺเฐสุ อุทาเนปิ.}“โย พฺราหฺมโณ พาหิต\-ปาปธมฺโม, \\ นิหุํหุงฺโก นิกฺกสาโว ยตตฺโต. \\ เวทนฺตคู วุสิต\-พฺรหฺมจริโย, \\ ธมฺเมน โส พฺรหฺมวาทํ วเทยฺย.}}\csromangathastanza{\csromangathawak{ยสฺสุสฺสทา นตฺถิ กุหิญฺจิ โลเก”ติ.}}}{อชปาลกถา นิฏฺฐิตา.}
\csromanmatikaanchor{mtk.4}
\csromantocmarkref{section}{3. มุจลินฺทกถา}{mtk.4}
\bookmark[dest={mtk.4},level=1]{3. มุจลินฺทกถา}
\csromanheader{1}{3. \textbf{มุจลินฺทกถา}}
\proseitem{5}{\csromansymbolmark{+}+ อถ โข ภควา สตฺตาหสฺส อจฺจเยน ตมฺหา สมาธิมฺหา วุฏฺฐหิตฺวา อชปาล\-นิโคฺรธมูลา เยน มุจลินฺโท เตนุปสงฺกมิ, อุปสงฺกมิตฺวา มุจลินฺทมูเล สตฺตาหํ เอกปลฺลงฺเกน นิสีทิ วิมุตฺติสุข\-ปฏิสํเวที. เตน โข ปน สมเยน มหา อกาลเมโฆ อุทปาทิ สตฺตาห\-วทฺทลิกา สีตวาต\-ทุทฺทินี. อถ โข มุจลินฺโท นาคราชา สกภวนา นิกฺขมิตฺวา ภควโต กายํ สตฺตกฺขตฺตุํ โภเคหิ ปริกฺขิปิตฺวา อุปริมุทฺธนิ มหนฺตํ ผณํ กริตฺวา อฏฺฐาสิ “มา ภควนฺตํ สีตํ, มา ภควนฺตํ อุณฺหํ, มา ภควนฺตํ ฑํสมกส\-วาตา\-ตป\-สรีสป\-สมฺผสฺโส”ติ\csromansharedfootnote{88a08bd42d02d72d}{สิริํสป... (สี, สฺยา, กํ)}. \csromanfolio{4}อถ โข มุจลินฺโท นาคราชา สตฺตาหสฺส อจฺจเยน วิทฺธํ วิคต\-วลาหกํ เทวํ วิทิตฺวา ภควโต กายา โภเค วินิเวเฐตฺวา สกวณฺณํ ปฏิสํหริตฺวา มาณวกวณฺณํ อภินิมฺมินิตฺวา ภควโต ปุรโต อฏฺฐาสิ ปญฺชลิโก ภควนฺตํ นมสฺสมาโน. อถ โข ภควา เอตมตฺถํ วิทิตฺวา ตายํ เวลายํ อิมํ อุทานํ อุทาเนสิ\csromandash{}}

\setlength{\csromangathapagewidth}{0pt}
\setlength{\csromangathawakwidth}{0pt}
\csromangathameasuregroup{\textsuperscript{*}“สุโข วิเวโก ตุฏฺฐสฺส, \\ อพฺยาปชฺชํ สุขํ โลเก, \\ \textsuperscript{*}สุขา วิราคตา โลเก, \\ อสฺมิมานสฺส โย วินโย,}{\csromangathabat{\textsuperscript{*}“สุโข วิเวโก ตุฏฺฐสฺส,}{สุตธมฺมสฺส ปสฺสโต.} \\ \csromangathabat{อพฺยาปชฺชํ สุขํ โลเก,}{ปาณภูเตสุ สํยโม.} \\ \csromangathabat{\textsuperscript{*}สุขา วิราคตา โลเก,}{กามานํ สมติกฺกโม.} \\ \csromangathabat{อสฺมิมานสฺส โย วินโย,}{เอตํ เว ปรมํ สุขนฺ”ติ.}}
\csromangathasetleft{\textsuperscript{*}“สุโข วิเวโก ตุฏฺฐสฺส, \\ อพฺยาปชฺชํ สุขํ โลเก, \\ \textsuperscript{*}สุขา วิราคตา โลเก, \\ อสฺมิมานสฺส โย วินโย,}
\csromangathagroupwithcloserruled{\csromangathastanza{\csromangathabat{\csromansymbolfootnote{*}{อภิ 4. 163 ปิฏฺเฐ กถาวตฺถุปาฬิยมฺปิ.}“สุโข วิเวโก ตุฏฺฐสฺส,}{สุตธมฺมสฺส ปสฺสโต.} \\ \csromangathabat{อพฺยาปชฺชํ สุขํ โลเก,}{ปาณภูเตสุ สํยโม.}}\csromangathastanza{\csromangathabat{\csromansymbolmark{*}สุขา วิราคตา โลเก,}{กามานํ สมติกฺกโม.} \\ \csromangathabat{อสฺมิมานสฺส โย วินโย,}{เอตํ เว ปรมํ สุขนฺ”ติ.}}}{มุจลินฺทกถา นิฏฺฐิตา.}
\csromanmatikaanchor{mtk.5}
\csromantocmarkref{section}{4. ราชายตนกถา}{mtk.5}
\bookmark[dest={mtk.5},level=1]{4. ราชายตนกถา}
\csromanheader{1}{4. \textbf{ราชายตน}\-\textbf{กถา}}
\proseitemwithcloserruled{6}{อถ โข ภควา สตฺตาหสฺส อจฺจเยน ตมฺหา สมาธิมฺหา วุฏฺฐหิตฺวา มุจลินฺทมูลา เยน ราชายตนํ เตนุปสงฺกมิ, อุปสงฺกมิตฺวา ราชายตนมูเล สตฺตาหํ เอกปลฺลงฺเกน นิสีทิ วิมุตฺติสุข\-ปฏิสํเวที. เตน โข ปน สมเยน ตปุสฺส\csromansharedfootnote{ba76d0f820698ab7}{ตปสฺสุ (สี)} ภลฺลิกา วาณิชา อุกฺกลา ตํ เทสํ อทฺธาน\-มคฺค\-ปฺปฏิปนฺนา โหนฺติ. อถ โข ตปุสฺส\-ภลฺลิกานํ วาณิชานํ ญาติสาโลหิตา เทวตา ตปุสฺสภลฺลิเก วาณิเช เอตทโวจ “อยํ มาริสา ภควา ราชายตนมูเล วิหรติ ปฐมา\-ภิสมฺพุทฺโธ, คจฺฉถ ตํ ภควนฺตํ มนฺเถน จ มธุปิณฺฑิกาย จ ปติมาเนถ, ตํ โว ภวิสฺสติ ทีฆรตฺตํ หิตาย สุขายา”ติ. อถ โข ตปุสฺส\-ภลฺลิกา วาณิชา มนฺถญฺจ มธุปิณฺฑิ\-กญฺจ อาทาย เยน ภควา เตนุ\-ปสงฺกมิํสุ, อุปสงฺกมิตฺวา ภควนฺตํ อภิวาเทตฺวา เอกมนฺตํ อฏฺฐํสุ, เอกมนฺตํ ฐิตา โข ตปุสฺส\-ภลฺลิกา วาณิชา ภควนฺตํ เอตทโวจุํ “ปฏิคฺคณฺหาตุ โน ภนฺเต ภควา มนฺถญฺจ มธุปิณฺฑิ\-กญฺจ, ยํ อมฺหากํ อสฺส ทีฆรตฺตํ หิตาย สุขายา”ติ. อถ โข ภควโต เอตทโหสิ “น โข ตถาคตา หตฺเถสุ ปฏิคฺคณฺหนฺติ, กิมฺหิ นุ โข อหํ ปฏิคฺคเณฺหยฺยํ มนฺถญฺจ มธุปิณฺฑิ\-กญฺจา”ติ. \csromanfolio{5}อถ โข จตฺตาโร มหาราชาโน ภควโต เจตสา เจโต\-ปริวิตกฺกม\-ญฺญาย จตุทฺทิสา จตฺตาโร เสลมเย ปตฺเต ภควโต อุปนาเมสุํ “อิธ ภนฺเต ภควา ปฏิคฺคณฺหาตุ มนฺถญฺจ มธุปิณฺฑิ\-กญฺจา”ติ. ปฏิคฺคเหสิ ภควา ปจฺจคฺเฆ เสลมเย ปตฺเต มนฺถญฺจ มธุปิณฺฑิ\-กญฺจ, ปฏิคฺคเหตฺวา ปริภุญฺชิ. อถ โข ตปุสฺส\-ภลฺลิกา วาณิชา ภควนฺตํ (\csromansymbolfoottext{( )}{สีสฺยาโปตฺถเกสุ นตฺถิ.}โอนีต\-ปตฺต\-ปาณิํ วิทิตฺวา ภควโต ปาเทสุ สิรสา นิปติตฺวา ภควนฺตํ) เอตทโวจุํ “เอเต มยํ ภนฺเต ภควนฺตํ สรณํ คจฺฉาม ธมฺมญฺจ, อุปาสเก โน ภควา ธาเรตุ อชฺชตคฺเค ปาณุเปเต สรณํ คเต”ติ. เต จ โลเก ปฐมํ อุปาสกา อเหสุํ เทฺววาจิกา.}{ราชายตน\-กถา นิฏฺฐิตา.}
\csromanmatikaanchor{mtk.6}
\csromantocmarkref{section}{5. พฺรหฺมยาจนกถา}{mtk.6}
\bookmark[dest={mtk.6},level=1]{5. พฺรหฺมยาจนกถา}
\csromanheader{1}{5. \textbf{พฺรหฺม}\-\textbf{ยาจน}\-\textbf{กถา}}
\proseitem{7}{\csromansymbolfootnote{*}{อยํ พฺรหฺมยาจนกถา ที 2. 30 ปิฏฺฐาทีสุ จ; ม 1. 224 ปิฏฺฐาทีสุ จ, ม 2. 291 ปิฏฺฐาทีสุ จ; สํ 1. 138 ปิฏฺฐาทีสุปิ จ อาคตา.}อถ โข ภควา สตฺตาหสฺส อจฺจเยน ตมฺหา สมาธิมฺหา วุฏฺฐหิตฺวา ราชายตนมูลา เยน อชปาล\-นิโคฺรโธ เตนุปสงฺกมิ. ตตฺร สุทํ ภควา อชปาล\-นิโคฺรธมูเล วิหรติ. อถ โข ภควโต รโหคตสฺส ปฏิสลฺลีนสฺส เอวํ เจตโส ปริวิตกฺโก อุทปาทิ “อธิคโต โข มฺยายํ ธมฺโม คมฺภีโร ทุทฺทโส ทุรนุโภโธ สนฺโต ปณีโต อตกฺกาวจโร นิปุโณ ปณฺฑิต\-เวทนีโย, อาลยรามา โข ปนายํ ปชา อาลยรตา อาลยสมฺมุทิตา, อาลยรามาย โข ปน ปชาย อาลยรตาย อาลย\-สมฺมุทิตาย ทุทฺทสํ อิทํ ฐานํ ยทิทํ อิทปฺปจฺจยตา\-ปฏิจฺจสมุปฺปาโท, อิทมฺปิ โข ฐานํ สุทุทฺทสํ ยทิทํ สพฺพ\-สงฺขาร\-สมโถ สพฺพู\-ปธิ\-ปฏินิสฺสคฺโค ตณฺหากฺขโย วิราโค นิโรโธ นิพฺพานํ. อหญฺเจว โข ปน ธมฺมํ เทเสยฺยํ, ปเร จ เม น อาชาเนยฺยุํ, โส มมสฺส กิลมโถ, สา มมสฺส วิเหสา”ติ. อปิสฺสุ ภควนฺตํ อิมา อนจฺฉริยา คาถาโย ปฏิภํสุ ปุพฺเพ อสฺสุตปุพฺพา\csromandash{}}

\setlength{\csromangathapagewidth}{0pt}
\setlength{\csromangathawakwidth}{0pt}
\csromangathameasuregroup{“กิจฺเฉน เม อธิคตํ, \\ ราคโทสปเรเตหิ, \\ ปฏิโสตคามิํ นิปุณํ, \\ ราครตฺตา น ทกฺขนฺติ,}{\csromangathabat{“กิจฺเฉน เม อธิคตํ,}{หลํ ทานิ ปกาสิตุํ.} \\ \csromangathabat{ราคโทสปเรเตหิ,}{นายํ ธมฺโม สุสมฺพุโธ.} \\ \csromangathabat{ปฏิโสตคามิํ นิปุณํ,}{คมฺภีรํ ทุทฺทสํ อณุํ.} \\ \csromangathabat{ราครตฺตา น ทกฺขนฺติ,}{ตโมขนฺเธน อาวุฏา\textsuperscript{1}”ติ.}}
\csromangathasetleft{“กิจฺเฉน เม อธิคตํ, \\ ราคโทสปเรเตหิ, \\ ปฏิโสตคามิํ นิปุณํ, \\ ราครตฺตา น ทกฺขนฺติ,}
\csromangathagroup{\csromangathastanza{\csromangathabat{“กิจฺเฉน เม อธิคตํ,}{หลํ ทานิ ปกาสิตุํ.} \\ \csromangathabat{ราคโทสปเรเตหิ,}{นายํ ธมฺโม สุสมฺพุโธ.}}\csromangathastanza{\csromanfolio{6}\csromangathabat{ปฏิโสตคามิํ นิปุณํ,}{คมฺภีรํ ทุทฺทสํ อณุํ.} \\ \csromangathabat{ราครตฺตา น ทกฺขนฺติ,}{ตโมขนฺเธน อาวุฏา\csromansharedfootnote{e54b7d22427c43de}{อาวฏา (สี)}”ติ.}}}

\prose{อิติห ภควโต ปฏิสญฺจิกฺขโต อปฺโปสฺสุกฺกตาย จิตฺตํ นมติ, โน ธมฺมเทสนาย.}

\proseitem{8}{อถ โข พฺรหฺมุโน สหมฺปติสฺส ภควโต เจตสา เจโต\-ปริวิตกฺกม\-ญฺญาย เอตทโหสิ “นสฺสติ วต โภ โลโก, วินสฺสติ วต โภ โลโก, ยตฺร หิ นาม ตถาคตสฺส อรหโต สมฺมา\-สมฺพุทฺธสฺส อปฺโปสฺสุกฺกตาย จิตฺตํ นมติ\csromansharedfootnote{fc1092df9bb1e965}{นมิสฺสติ (?)}, โน ธมฺมเทสนายา”ติ. อถ โข พฺรหฺมา สหมฺปติ เสยฺยถาปิ นาม พลวา ปุริโส สมิญฺชิตํ วา พาหํ ปสาเรยฺย, ปสาริตํ วา พาหํ สมิญฺเชยฺย, เอวเมว พฺรหฺมโลเก อนฺตรหิโต ภควโต ปุรโต ปาตุรโหสิ. อถ โข พฺรหฺมา สหมฺปติ เอกํสํ อุตฺตราสงฺคํ กริตฺวา ทกฺขิณ\-ชาณุมณฺฑลํ ปถวิยํ นิหนฺตฺวา เยน ภควา เตนญฺชลิํ ปณาเมตฺวา ภควนฺตํ เอตทโวจ “เทเสตุ ภนฺเต ภควา ธมฺมํ, เทเสตุ สุคโต ธมฺมํ, สนฺติ สตฺตา อปฺปรชกฺข\-ชาติกา, อสฺสวนตา ธมฺมสฺส ปริหายนฺติ, ภวิสฺสนฺติ ธมฺมสฺส อญฺญาตาโร”ติ. อิทมโวจ พฺรหฺมา สหมฺปติ, อิทํ วตฺวาน อถาปรํ เอตทโวจ\csromandash{}}

\setlength{\csromangathapagewidth}{0pt}
\setlength{\csromangathawakwidth}{0pt}
\csromangathameasuregroup{}{“ปาตุรโหสิ มคเธสุ ปุพฺเพ, \\ ธมฺโม อสุทฺโธ สมเลหิ จินฺติโต. \\ อปาปุเรตํ\textsuperscript{1} อมตสฺส ทฺวารํ, \\ สุณนฺตุ ธมฺมํ วิมเลนานุพุทฺธํ. \\ เสเล ยถา ปพฺพตมุทฺธนิฏฺฐิโต, \\ ยถาปิ ปสฺเส ชนตํ สมนฺตโต. \\ ตถูปมํ ธมฺมมยํ สุเมธ, \\ ปาสาทมารุยฺห สมนฺตจกฺขุ. \\ โสกาวติณฺณํ ชนตมเปตโสโก, \\ อเวกฺขสฺสุ ชาติชราภิภูตํ. \\ อุฏฺเฐหิ วีร วิชิตสงฺคาม, \\ สตฺถวาห อณณ\textsuperscript{1} วิจร โลเก. \\ เทสสฺสุ\textsuperscript{1} ภควา ธมฺมํ, \\ อญฺญาตาโร ภวิสฺสนฺตี”ติ.}
\csromangathameasurewak{“ปาตุรโหสิ มคเธสุ ปุพฺเพ, \\ ธมฺโม อสุทฺโธ สมเลหิ จินฺติโต. \\ อปาปุเรตํ\textsuperscript{1} อมตสฺส ทฺวารํ, \\ สุณนฺตุ ธมฺมํ วิมเลนานุพุทฺธํ. \\ เสเล ยถา ปพฺพตมุทฺธนิฏฺฐิโต, \\ ยถาปิ ปสฺเส ชนตํ สมนฺตโต. \\ ตถูปมํ ธมฺมมยํ สุเมธ, \\ ปาสาทมารุยฺห สมนฺตจกฺขุ. \\ โสกาวติณฺณํ ชนตมเปตโสโก, \\ อเวกฺขสฺสุ ชาติชราภิภูตํ. \\ อุฏฺเฐหิ วีร วิชิตสงฺคาม, \\ สตฺถวาห อณณ\textsuperscript{1} วิจร โลเก. \\ เทสสฺสุ\textsuperscript{1} ภควา ธมฺมํ, \\ อญฺญาตาโร ภวิสฺสนฺตี”ติ.}
\setlength{\csromangathaleftcol}{0pt}
\csromangathagroup{\csromangathastanza{\csromangathawak{“ปาตุรโหสิ มคเธสุ ปุพฺเพ, \\ ธมฺโม อสุทฺโธ สมเลหิ จินฺติโต. \\ อปาปุเรตํ\csromansharedfootnote{08efd866b43f7eb5}{อวาปุเรตํ (สี)} อมตสฺส ทฺวารํ, \\ สุณนฺตุ ธมฺมํ วิมเลนา\-นุพุทฺธํ.}}\csromangathastanza{\csromangathawak{เสเล ยถา ปพฺพตมุทฺธนิ\-ฏฺฐิโต, \\ ยถาปิ ปสฺเส ชนตํ สมนฺตโต. \\ ตถูปมํ ธมฺมมยํ สุเมธ, \\ ปาสาทมารุยฺห สมนฺตจกฺขุ.}}\csromangathastanza{\csromangathawak{โสกาวติณฺณํ ชนตม\-เปตโสโก, \\ อเวกฺขสฺสุ ชาติชรา\-ภิภูตํ. \\ อุฏฺเฐหิ วีร วิชิตสงฺคาม, \\ สตฺถวาห อณณ\csromansharedfootnote{fb6d53d1469cd44f}{อนณ (ก)} วิจร โลเก. \\ \csromanfolio{7}เทสสฺสุ\csromansharedfootnote{7a86865577559d98}{เทเสตุ (ก) [ ] สีสฺยาโปตฺถเกสุ นตฺถิ, มูลปณฺณาสเกปิ ปาสราสิสุตฺเต พฺรหฺมยาจนา สกิํ เยว อาคตา.} ภควา ธมฺมํ, \\ อญฺญาตาโร ภวิสฺสนฺตี”ติ.}}}

\prose{[เอวํ วุตฺเต ภควา พฺรหฺมานํ สหมฺปติํ เอตทโวจ “มยฺหมฺปิ โข พฺรหฺเม เอตทโหสิ ‘อธิคโต โข มฺยายํ ธมฺโม คมฺภีโร ทุทฺทโส ทุรนุโพโธ สนฺโต ปณีโต อตกฺกาวจโร นิปุโณ ปณฺฑิต\-เวทนีโย, อาลยรามา โข ปนายํ ปชา อาลยรตา อาลยสมฺมุทิตา, อาลยรามาย โข ปน ปชาย อาลยรตาย อาลย\-สมฺมุทิตาย ทุทฺทสํ อิทํ ฐานํ ยทิทํ อิทปฺปจฺจยตา\-ปฏิจฺจสมุปฺปาโท, อิทมฺปิ โข ฐานํ สุทุทฺทสํ ยทิทํ สพฺพ\-สงฺขาร\-สมโถ สพฺพู\-ปธิ\-ปฏินิสฺสคฺโค ตณฺหากฺขโย วิราโค นิโรโธ นิพฺพานํ. อหญฺเจว โข ปน ธมฺมํ เทเสยฺยํ, ปเร จ เม น อาชาเนยฺยุํ, โส มมสฺส กิลมโถ, สา มมสฺส วิเหสา’ติ. อปิสฺสุ มํ พฺรหฺเม อิมา อนจฺฉริยา คาถาโย ปฏิภํสุ ปุพฺเพ อสฺสุตปุพฺพา\csromandash{}}

\setlength{\csromangathapagewidth}{0pt}
\setlength{\csromangathawakwidth}{0pt}
\csromangathameasuregroup{“กิจฺเฉน เม อธิคตํ, \\ ราคโทสปเรเตหิ, \\ ปฏิโสตคามิํ นิปุณํ, \\ ราครตฺตา น ทกฺขนฺติ,}{\csromangathabat{“กิจฺเฉน เม อธิคตํ,}{หลํ ทานิ ปกาสิตุํ.} \\ \csromangathabat{ราคโทสปเรเตหิ,}{นายํ ธมฺโม สุสมฺพุโธ.} \\ \csromangathabat{ปฏิโสตคามิํ นิปุณํ,}{คมฺภีรํ ทุทฺทสํ อณุํ.} \\ \csromangathabat{ราครตฺตา น ทกฺขนฺติ,}{ตโมขนฺเธน อาวุฏา’ติ.}}
\csromangathasetleft{“กิจฺเฉน เม อธิคตํ, \\ ราคโทสปเรเตหิ, \\ ปฏิโสตคามิํ นิปุณํ, \\ ราครตฺตา น ทกฺขนฺติ,}
\csromangathagroup{\csromangathastanza{\csromangathabat{“กิจฺเฉน เม อธิคตํ,}{หลํ ทานิ ปกาสิตุํ.} \\ \csromangathabat{ราคโทสปเรเตหิ,}{นายํ ธมฺโม สุสมฺพุโธ.}}\csromangathastanza{\csromangathabat{ปฏิโสตคามิํ นิปุณํ,}{คมฺภีรํ ทุทฺทสํ อณุํ.} \\ \csromangathabat{ราครตฺตา น ทกฺขนฺติ,}{ตโมขนฺเธน อาวุฏา’ติ.}}}

\prose{อิติห เม พฺรหฺเม ปฏิสญฺจิกฺขโต อปฺโปสฺสุกฺกตาย จิตฺตํ นมติ, โน ธมฺมเทสนายา”ติ.}

\prose{ทุติยมฺปิ โข พฺรหฺมา สหมฺปติ ภควนฺตํ เอตทโวจ “เทเสตุ ภนฺเต ภควา ธมฺมํ, เทเสตุ สุคโต ธมฺมํ, สนฺติ สตฺตา อปฺปรชกฺข\-ชาติกา, อสฺสวนตา ธมฺมสฺส ปริหายนฺติ, ภวิสฺสนฺติ ธมฺมสฺส อญฺญาตาโร”ติ. อิทมโวจ พฺรหฺมา สหมฺปติ, อิทํ วตฺวาน อถาปรํ เอตทโวจ\csromandash{}}

\prose{\csromanfolio{8}“ปาตุรโหสิ มคเธสุ ปุพฺเพ.}

\prose{ธมฺโม อสุทฺโธ สมเลหิ จินฺติโต.}

\prose{อปาปุเรตํ อมตสฺส ทฺวารํ.}

\vspace{\tipitakaparskip}
\csromancenter{สุณนฺตุ ธมฺมํ วิมเลนา\-นุพุทฺธํ.}
\setlength{\csromangathapagewidth}{0pt}
\setlength{\csromangathawakwidth}{0pt}
\csromangathameasuregroup{}{เสเล ยถา ปพฺพตมุทฺธนิฏฺฐิโต, \\ ยถาปิ ปสฺเส ชนตํ สมนฺตโต. \\ ตถูปมํ ธมฺมมยํ สุเมธ, \\ ปาสาทมารุยฺห สมนฺตจกฺขุ. \\ โสกาวติณฺณํ ชนตมเปตโสโก, \\ อเวกฺขสฺสุ ชาติชราภิภูตํ. \\ อุฏฺเฐหิ วีร วิชิตสงฺคาม, \\ สตฺถวาห อณณ วิจร โลเก. \\ เทสสฺสุ ภควา ธมฺมํ, \\ อญฺญาตาโร ภวิสฺสนฺตี”ติ.}
\csromangathameasurewak{เสเล ยถา ปพฺพตมุทฺธนิฏฺฐิโต, \\ ยถาปิ ปสฺเส ชนตํ สมนฺตโต. \\ ตถูปมํ ธมฺมมยํ สุเมธ, \\ ปาสาทมารุยฺห สมนฺตจกฺขุ. \\ โสกาวติณฺณํ ชนตมเปตโสโก, \\ อเวกฺขสฺสุ ชาติชราภิภูตํ. \\ อุฏฺเฐหิ วีร วิชิตสงฺคาม, \\ สตฺถวาห อณณ วิจร โลเก. \\ เทสสฺสุ ภควา ธมฺมํ, \\ อญฺญาตาโร ภวิสฺสนฺตี”ติ.}
\setlength{\csromangathaleftcol}{0pt}
\csromangathagroup{\csromangathastanza{\csromangathawak{เสเล ยถา ปพฺพตมุทฺธนิ\-ฏฺฐิโต, \\ ยถาปิ ปสฺเส ชนตํ สมนฺตโต. \\ ตถูปมํ ธมฺมมยํ สุเมธ, \\ ปาสาทมารุยฺห สมนฺตจกฺขุ.}}\csromangathastanza{\csromangathawak{โสกาวติณฺณํ ชนตม\-เปตโสโก, \\ อเวกฺขสฺสุ ชาติชรา\-ภิภูตํ. \\ อุฏฺเฐหิ วีร วิชิตสงฺคาม, \\ สตฺถวาห อณณ วิจร โลเก. \\ เทสสฺสุ ภควา ธมฺมํ, \\ อญฺญาตาโร ภวิสฺสนฺตี”ติ.}}}

\prose{ทุติยมฺปิ โข ภควา พฺรหฺมานํ สหมฺปติํ เอตทโวจ “มยฺหมฺปิ โข พฺรหฺเม เอตทโหสิ ‘อธิคโต โข มฺยายํ ธมฺโม คมฺภีโร ทุทฺทโส ทุรนุโพโธ สนฺโต ปณีโต อตกฺกาวจโร นิปุโณ ปณฺฑิต\-เวทนีโย, อาลยรามา โข ปนายํ ปชา อาลยรตา อาลยสมฺมุทิตา, อาลยรามาย โข ปน ปชาย อาลยรตาย อาลย\-สมฺมุทิตาย ทุทฺทสํ อิทํ ฐานํ ยทิทํ อิทปฺปจฺจยตา\-ปฏิจฺจสมุปฺปาโท, อิทมฺปิ โข ฐานํ สุทุทฺทสํ ยทิทํ สพฺพ\-สงฺขาร\-สมโถ สพฺพู\-ปธิ\-ปฏินิสฺสคฺโค ตณฺหากฺขโย วิราโค นิโรโธ นิพฺพานํ. อหญฺเจว โข ปน ธมฺมํ เทเสยฺยํ ปเร จ เม น อาชาเนยฺยุํ, โส มมสฺส กิลมโถ, สา มมสฺส วิเหสา’ติ. อปิสฺสุ มํ พฺรหฺเม อิมา อนจฺฉริยา คาถาโย ปฏิภํสุ ปุพฺเพ อสฺสุตปุพฺพา\csromandash{}}

\setlength{\csromangathapagewidth}{0pt}
\setlength{\csromangathawakwidth}{0pt}
\csromangathameasuregroup{‘กิจฺเฉน เม อธิคตํ, \\ ราคโทสปเรเตหิ, \\ ปฏิโสตคามิํ นิปุณํ, \\ ราครตฺตา น ทกฺขนฺติ,}{\csromangathabat{‘กิจฺเฉน เม อธิคตํ,}{หลํ ทานิ ปกาสิตุํ.} \\ \csromangathabat{ราคโทสปเรเตหิ,}{นายํ ธมฺโม สุสมฺพุโธ.} \\ \csromangathabat{ปฏิโสตคามิํ นิปุณํ,}{คมฺภีรํ ทุทฺทสํ อณุํ.} \\ \csromangathabat{ราครตฺตา น ทกฺขนฺติ,}{ตโมขนฺเธน อาวุฏา’ติ.}}
\csromangathasetleft{‘กิจฺเฉน เม อธิคตํ, \\ ราคโทสปเรเตหิ, \\ ปฏิโสตคามิํ นิปุณํ, \\ ราครตฺตา น ทกฺขนฺติ,}
\csromangathagroup{\csromangathastanza{\csromangathabat{‘กิจฺเฉน เม อธิคตํ,}{หลํ ทานิ ปกาสิตุํ.} \\ \csromangathabat{ราคโทสปเรเตหิ,}{นายํ ธมฺโม สุสมฺพุโธ.}}\csromangathastanza{\csromangathabat{ปฏิโสตคามิํ นิปุณํ,}{คมฺภีรํ ทุทฺทสํ อณุํ.} \\ \csromangathabat{ราครตฺตา น ทกฺขนฺติ,}{ตโมขนฺเธน อาวุฏา’ติ.}}}

\prose{\csromanfolio{9}อิติห เม พฺรหฺเม ปฏิสญฺจิกฺขโต อปฺโปสฺสุกฺกตาย จิตฺตํ นมติ, โน ธมฺมเทสนายา”ติ.}

\prose{ตติยมฺปิ โข พฺรหฺมา สหมฺปติ ภควนฺตํ เอตทโวจ “เทเสตุ ภนฺเต ภควา ธมฺมํ, เทเสตุ สุคโต ธมฺมํ, สนฺติ สตฺตา อปฺปรชกฺข\-ชาติกา, อสฺสวนตา ธมฺมสฺส ปริหายนฺติ, ภวิสฺสนฺติ ธมฺมสฺส อญฺญาตาโร”ติ. อิทมโวจ พฺรหฺมา สหมฺปติ, อิทํ วตฺวาน อถาปรํ เอตทโวจ\csromandash{}}

\setlength{\csromangathapagewidth}{0pt}
\setlength{\csromangathawakwidth}{0pt}
\csromangathameasuregroup{“ปาตุรโหสิ มคเธสุ ปุพฺเพ,}{\csromangathabat{“ปาตุรโหสิ มคเธสุ ปุพฺเพ,}{ธมฺโม อสุทฺโธ สมเลหิ จินฺติโต.} \\ อปาปุเรตํ อมตสฺส ทฺวารํ. \\ สุณนฺตุ ธมฺมํ วิมเลนานุพุทฺธํ. \\ เสเล ยถา ปพฺพตมุทฺธนิฏฺฐิโต, \\ ยถาปิ ปสฺเส ชนตํ สมนฺตโต. \\ ตถูปมํ ธมฺมมยํ สุเมธ, \\ ปาสาทมารุยฺห สมนฺตจกฺขุ. \\ โสกาวติณฺณํ ชนตมเปตโสโก, \\ อเวกฺขสฺสุ ชาติชราภิภูตํ. \\ อุฏฺเฐหิ วีร วิชิตสงฺคาม, \\ สตฺถวาห อณณ วิจร โลเก. \\ เทสสฺสุ ภควา ธมฺมํ,}
\csromangathameasurewak{เสเล ยถา ปพฺพตมุทฺธนิฏฺฐิโต, \\ ยถาปิ ปสฺเส ชนตํ สมนฺตโต. \\ ตถูปมํ ธมฺมมยํ สุเมธ, \\ ปาสาทมารุยฺห สมนฺตจกฺขุ. \\ โสกาวติณฺณํ ชนตมเปตโสโก, \\ อเวกฺขสฺสุ ชาติชราภิภูตํ. \\ อุฏฺเฐหิ วีร วิชิตสงฺคาม, \\ สตฺถวาห อณณ วิจร โลเก. \\ เทสสฺสุ ภควา ธมฺมํ,}
\csromangathasetleft{“ปาตุรโหสิ มคเธสุ ปุพฺเพ,}
\csromangathagroup{\csromangathastanza{\csromangathabat{“ปาตุรโหสิ มคเธสุ ปุพฺเพ,}{ธมฺโม อสุทฺโธ สมเลหิ จินฺติโต.} \\ อปาปุเรตํ อมตสฺส ทฺวารํ. \\ สุณนฺตุ ธมฺมํ วิมเลนา\-นุพุทฺธํ.}\csromangathastanza{\csromangathawak{เสเล ยถา ปพฺพตมุทฺธนิ\-ฏฺฐิโต, \\ ยถาปิ ปสฺเส ชนตํ สมนฺตโต. \\ ตถูปมํ ธมฺมมยํ สุเมธ, \\ ปาสาทมารุยฺห สมนฺตจกฺขุ.}}\csromangathastanza{\csromangathawak{โสกาวติณฺณํ ชนตม\-เปตโสโก, \\ อเวกฺขสฺสุ ชาติชรา\-ภิภูตํ. \\ อุฏฺเฐหิ วีร วิชิตสงฺคาม, \\ สตฺถวาห อณณ วิจร โลเก.}}\csromangathastanza{\csromangathawak{เทสสฺสุ ภควา ธมฺมํ,}}}

\prose{อญฺญาตาโร ภวิสฺสนฺตี”ติ.]}

\proseitem{9}{อถ โข ภควา พฺรหฺมุโน จ อชฺเฌสนํ วิทิตฺวา สตฺเตสุ จ การุญฺญตํ ปฏิจฺจ พุทฺธจกฺขุนา โลกํ โวโลเกสิ, อทฺทสา โข ภควา พุทฺธจกฺขุนา โลกํ โวโลเกนฺโต สตฺเต อปฺปรชกฺเข มหารชกฺเข ติกฺขินฺทฺริเย มุทินฺทฺริเย สฺวากาเร ทฺวากาเร สุวิญฺญาปเย ทุวิญฺญาปเย อปฺเปกจฺเจ ปรโลก\-วชฺช\-ภยทสฺสาวิเน\csromansharedfootnote{7977f1ee055940f6}{ทสฺสาวิโน (สี, สฺยา, กํ)} วิหรนฺเต อปฺเปกจฺเจ น ปรโลก\-วชฺช\-ภยทสฺสาวิเน\csromansharedfootnote{7977f1ee055940f6}{ทสฺสาวิโน (สี, สฺยา, กํ)} วิหรนฺเต, เสยฺยถาปิ นาม อุปฺปลินิยํ วา ปทุมินิยํ วา ปุณฺฑรีกินิยํ วา อปฺเปกจฺจานิ อุปฺปลานิ วา ปทุมานิ วา ปุณฺฑรีกานิ วา อุทเก ชาตานิ อุทเก สํวฑฺฒานิ อุทกานุคฺคตานิ อนฺโต \csromanfolio{10}นิมุคฺคโปสีนิ, อปฺเปกจฺจานิ อุปฺปลานิ วา ปทุมานิ วา ปุณฺฑรีกานิ วา อุทเก ชาตานิ อุทเก สํวฑฺฒานิ สโมทกํ ฐิตานิ, อปฺเปกจฺจานิ อุปฺปลานิ วา ปทุมานิ วา ปุณฺฑรีกานิ วา อุทเก ชาตานิ อุทเก สํวฑฺฒานิ อุทกํ อจฺจุคฺคมฺม ฐิตานิ\csromansharedfootnote{dc091edcce453433}{ติฏฺฐนฺติ (สี, สฺยา)} อนุปลิตฺตานิ อุทเกน, เอวเมวํ ภควา พุทฺธจกฺขุนา โลกํ โวโลเกนฺโต อทฺทส สตฺเต อปฺปรชกฺเข มหารชกฺเข ติกฺขินฺทฺริเย มุทินฺทฺริเย สฺวากาเร ทฺวากาเร สุวิญฺญาปเย ทุวิญฺญาปเย อปฺเปกจฺเจ ปรโลก\-วชฺช\-ภยทสฺสาวิเน วิหรนฺเต อปฺเปกจฺเจ น ปรโลก\-วชฺช\-ภยทสฺสาวิเน วิหรนฺเต, ทิสฺวาน พฺรหฺมานํ สหมฺปติํ คาถาย ปจฺจภาสิ\csromandash{}}

\setlength{\csromangathapagewidth}{0pt}
\setlength{\csromangathawakwidth}{0pt}
\csromangathameasuregroup{}{“อปารุตา เตสํ อมตสฺส ทฺวารา, \\ เย โสตวนฺโต ปมุญฺจนฺตุ สทฺธํ. \\ วิหิํสสญฺญี ปคุณํ น ภาสิํ, \\ ธมฺมํ ปณีตํ มนุเชสุ พฺรหฺเม”ติ.}
\csromangathameasurewak{“อปารุตา เตสํ อมตสฺส ทฺวารา, \\ เย โสตวนฺโต ปมุญฺจนฺตุ สทฺธํ. \\ วิหิํสสญฺญี ปคุณํ น ภาสิํ, \\ ธมฺมํ ปณีตํ มนุเชสุ พฺรหฺเม”ติ.}
\setlength{\csromangathaleftcol}{0pt}
\csromangathagroup{\csromangathastanza{\csromangathawak{“อปารุตา เตสํ อมตสฺส ทฺวารา, \\ เย โสตวนฺโต ปมุญฺจนฺตุ สทฺธํ. \\ วิหิํสสญฺญี ปคุณํ น ภาสิํ, \\ ธมฺมํ ปณีตํ มนุเชสุ พฺรหฺเม”ติ.}}}

\prosewithcloserruled{อถ โข พฺรหฺมา สหมฺปติ “กตาวกาโส โขมฺหิ ภควตา ธมฺมเทสนายา”ติ ภควนฺตํ อภิวาเทตฺวา ปทกฺขิณํ กตฺวา ตตฺเถ\-ว\-นฺตรธายิ.}{พฺรหฺม\-ยาจน\-กถา นิฏฺฐิตา.}
\csromanmatikaanchor{mtk.7}
\csromantocmarkref{section}{6. ปญฺจวคฺคิยกถา}{mtk.7}
\bookmark[dest={mtk.7},level=1]{6. ปญฺจวคฺคิยกถา}
\csromanheader{1}{6. \textbf{ปญฺจวคฺคิย}\-\textbf{กถา}}
\proseitem{10}{\csromansymbolfootnote{*}{ม 1. 226 ปิฏฺฐาทีสุปิ; ม 2. 293 ปิฏฺฐาทีสุปิ จ.}อถ โข ภควโต เอตทโหสิ “กสฺส นุ โข อหํ ปฐมํ ธมฺมํ เทเสยฺยํ, โก อิมํ ธมฺมํ ขิปฺปเมว อาชานิสฺสตี”ติ. อถ โข ภควโต เอตทโหสิ “อยํ โข อาฬาโร กาลาโม ปณฺฑิโต พฺยตฺโต เมธาวี ทีฆรตฺตํ อปฺปรชกฺข\-ชาติโก, ยนฺนูนาหํ อาฬารสฺส กาลามสฺส ปฐมํ ธมฺมํ เทเสยฺยํ, โส อิมํ ธมฺมํ ขิปฺปเมว อาชานิสฺสตี”ติ. อถ โข อนฺตรหิตา เทวตา ภควโต อาโรเจสิ “สตฺตาห\-กาลงฺกโต ภนฺเต อาฬาโร กาลาโม”ติ. ภควโตปิ โข ญาณํ อุทปาทิ “สตฺตาห\-กาลงฺกโต \csromanfolio{11}อาฬาโร กาลาโม”ติ. อถ โข ภควโต เอตทโหสิ “มหาชานิโย โข อาฬาโร กาลาโม, สเจ หิ โส อิมํ ธมฺมํ สุเณยฺย, ขิปฺปเมว อาชาเนยฺยา”ติ. อถ โข ภควโต เอตทโหสิ “กสฺส นุ โข อหํ ปฐมํ ธมฺมํ เทเสยฺยํ, โก อิมํ ธมฺมํ ขิปฺปเมว อาชานิสฺสตี”ติ. อถ โข ภควโต เอตทโหสิ “อยํ โข อุทโก\csromansharedfootnote{7b9546b5c8280955}{อุทฺทโก (สี, สฺยา)} รามปุตฺโต ปณฺฑิโต พฺยตฺโต เมธาวี ทีฆรตฺตํ อปฺปรชกฺข\-ชาติโก, ยนฺนูนาหํ อุทกสฺส รามปุตฺตสฺส ปฐมํ ธมฺมํ เทเสยฺยํ, โส อิมํ ธมฺมํ ขิปฺปเมว อาชานิสฺสตี”ติ. อถ โข อนฺตรหิตา เทวตา ภควโต อาโรเจสิ “อภิโทส\-กาลงฺกโต ภนฺเต อุทโก รามปุตฺโต”ติ. ภควโตปิ โข ญาณํ อุทปาทิ “อภิโทส\-กาลงฺกโต อุทโก รามปุตฺโต”ติ. อถ โข ภควโต เอตทโหสิ “มหาชานิโย โข อุทโก รามปุตฺโต, สเจ หิ โส อิมํ ธมฺมํ สุเณยฺย, ขิปฺปเมว อาชาเนยฺยา”ติ.}

\prose{อถ โข ภควโต เอตทโหสิ “กสฺส นุ โข อหํ ปฐมํ ธมฺมํ เทเสยฺยํ, โก อิมํ ธมฺมํ ขิปฺปเมว อาชานิสฺสตี”ติ. อถ โข ภควโต เอตทโหสิ “พหุการา โข เม ปญฺจวคฺคิยา ภิกฺขู, เย มํ ปธาน\-ปหิตตฺตํ อุปฏฺฐหิํสุ, ยนฺนูนาหํ ปญฺจ\-วคฺคิยานํ ภิกฺขูนํ ปฐมํ ธมฺมํ เทเสยฺยนฺ”ติ. อถ โข ภควโต เอตทโหสิ “กหํ นุ โข เอตรหิ ปญฺจวคฺคิยา ภิกฺขู วิหรนฺตี”ติ. อทฺทสา โข ภควา ทิพฺเพน จกฺขุนา วิสุทฺเธน อติกฺกนฺต\-มานุสเกน ปญฺจวคฺคิเย ภิกฺขู พาราณสิยํ วิหรนฺเต อิสิปตเน มิคทาเย. อถ โข ภควา อุรุเวลายํ ยถาภิรนฺตํ วิหริตฺวา เยน พาราณสี เตน จาริกํ ปกฺกามิ.}

\proseitem{11}{อทฺทสา โข อุปโก อาชีวโก ภควนฺตํ อนฺตรา จ คยํ อนฺตรา จ โพธิํ อทฺธาน\-มคฺค\-ปฺปฏิปนฺนํ, ทิสฺวาน ภควนฺตํ เอตทโวจ “วิปฺปสนฺนานิ โข เต อาวุโส อินฺทฺริยานิ, ปริสุทฺโธ ฉวิวณฺโณ ปริโยทาโต, กํ’สิ ตฺวํ อาวุโส อุทฺทิสฺส ปพฺพชิโต, โก วา เต สตฺถา, กสฺส วา ตฺวํ ธมฺมํ โรเจสี”ติ. เอวํ วุตฺเต ภควา อุปกํ อาชีวกํ คาถาหิ อชฺฌภาสิ\csromandash{}}

\prose{\csromanfolio{12}\csromansymbolfootnote{+}{ขุ 1. 64 ปิฏฺเฐ ธมฺมปเทปิ.}* “สพฺพาภิภู สพฺพวิทูหมสฺมิ,}

\prose{สพฺเพสุ ธมฺเมสุ อนูปลิตฺโต.}

\setlength{\csromangathapagewidth}{0pt}
\setlength{\csromangathawakwidth}{0pt}
\csromangathameasuregroup{}{สพฺพญฺชโห ตณฺหากฺขเย วิมุตฺโต, \\ สยํ อภิญฺญาย ก’มุทฺทิเสยฺยํ.}
\csromangathameasurewak{สพฺพญฺชโห ตณฺหากฺขเย วิมุตฺโต, \\ สยํ อภิญฺญาย ก’มุทฺทิเสยฺยํ.}
\setlength{\csromangathaleftcol}{0pt}
\csromangathagroup{\csromangathastanza{\csromangathawak{สพฺพญฺชโห ตณฺหากฺขเย วิมุตฺโต, \\ สยํ อภิญฺญาย ก’มุทฺทิเสยฺยํ.}}}

\prose{++\csromansymbolfootnote{*}{อภิ 4. 218 ปิฏฺเฐ กถาวตฺถุปาฬิยมฺปิ. ++ ขุ 11. 228 ปิฏฺเฐ มิลินฺทปเญฺหปิ.}น เม อาจริโย อตฺถิ, สทิโส เม น วิชฺชติ.}

\setlength{\csromangathapagewidth}{0pt}
\setlength{\csromangathawakwidth}{0pt}
\csromangathameasuregroup{สเทวกสฺมิํ โลกสฺมิํ, \\ \textsuperscript{*}อหํ หิ อรหา โลเก, \\ เอโกมฺหิ สมฺมาสมฺพุทฺโธ, \\ \textsuperscript{*}ธมฺมจกฺกํ ปวตฺเตตุํ, \\ อนฺธีภูตสฺมิํ โลกสฺมิํ, \\ ยถา โข ตฺวํ อาวุโส ปฏิชานาสิ, \\ \textsuperscript{*}“มาทิสา เว ชินา โหนฺติ, \\ ชิตา เม ปาปกา ธมฺมา,}{\csromangathabat{สเทวกสฺมิํ โลกสฺมิํ,}{นตฺถิ เม ปฏิปุคฺคโล.} \\ \csromangathabat{\textsuperscript{*}อหํ หิ อรหา โลเก,}{อหํ สตฺถา อนุตฺตโร.} \\ \csromangathabat{เอโกมฺหิ สมฺมาสมฺพุทฺโธ,}{สีติภูโตสฺมิ นิพฺพุโต.} \\ \csromangathabat{\textsuperscript{*}ธมฺมจกฺกํ ปวตฺเตตุํ,}{คจฺฉามิ กาสินํ ปุรํ.} \\ \csromangathabat{อนฺธีภูตสฺมิํ โลกสฺมิํ,}{อาหญฺฉํ\textsuperscript{1} อมตทุนฺทุภินฺ”ติ.} \\ \csromangathabat{ยถา โข ตฺวํ อาวุโส ปฏิชานาสิ,}{อรหสิ อนนฺตชิโนติ.} \\ \csromangathabat{\textsuperscript{*}“มาทิสา เว ชินา โหนฺติ,}{เย ปตฺตา อาสวกฺขยํ.} \\ \csromangathabat{ชิตา เม ปาปกา ธมฺมา,}{ตสฺมาหมุปก\textsuperscript{1} ชิโน”ติ.}}
\csromangathasetleft{สเทวกสฺมิํ โลกสฺมิํ, \\ \textsuperscript{*}อหํ หิ อรหา โลเก, \\ เอโกมฺหิ สมฺมาสมฺพุทฺโธ, \\ \textsuperscript{*}ธมฺมจกฺกํ ปวตฺเตตุํ, \\ อนฺธีภูตสฺมิํ โลกสฺมิํ, \\ ยถา โข ตฺวํ อาวุโส ปฏิชานาสิ, \\ \textsuperscript{*}“มาทิสา เว ชินา โหนฺติ, \\ ชิตา เม ปาปกา ธมฺมา,}
\csromangathagroup{\csromangathastanza{\csromangathabat{สเทวกสฺมิํ โลกสฺมิํ,}{นตฺถิ เม ปฏิปุคฺคโล.} \\ \csromangathabat{\csromansymbolmark{*}อหํ หิ อรหา โลเก,}{อหํ สตฺถา อนุตฺตโร.}}\csromangathastanza{\csromangathabat{เอโกมฺหิ สมฺมาสมฺพุทฺโธ,}{สีติภูโตสฺมิ นิพฺพุโต.} \\ \csromangathabat{\csromansymbolmark{*}ธมฺมจกฺกํ ปวตฺเตตุํ,}{คจฺฉามิ กาสินํ ปุรํ.}}\csromangathastanza{\csromangathabat{อนฺธีภูตสฺมิํ โลกสฺมิํ,}{อาหญฺฉํ\csromansharedfootnote{2f32347d10dd9814}{อาหญฺญิํ (ก)} อมตทุนฺทุภินฺ”ติ.} \\ \csromangathabat{ยถา โข ตฺวํ อาวุโส ปฏิชานาสิ,}{อรหสิ อนนฺตชิโนติ.}}\csromangathastanza{\csromangathabat{\csromansymbolmark{*}“มาทิสา เว ชินา โหนฺติ,}{เย ปตฺตา อาสวกฺขยํ.} \\ \csromangathabat{ชิตา เม ปาปกา ธมฺมา,}{ตสฺมาหมุปก\csromansharedfootnote{8c81af929783f102}{ตสฺมาหมุปกา (สี)} ชิโน”ติ.}}}

\prose{เอวํ วุตฺเต อุปโก อาชีวโก “หุเปยฺยปาวุโส\csromansharedfootnote{aa43c499fbdedfb2}{หุเวยฺยปาวุโส (สี) หุเวยฺยาวุโส (สฺยา)}”ติ วตฺวา สีสํ โอกมฺเปตฺวา อุมฺมคฺคํ คเหตฺวา ปกฺกามิ.}

\proseitem{12}{อถ โข ภควา อนุปุพฺเพน จาริกํ จรมาโน เยน พาราณสี อิสิปตนํ มิคทาโย, เยน ปญฺจวคฺคิยา ภิกฺขู เตนุปสงฺกมิ. อทฺทสํสุ โข ปญฺจวคฺคิยา ภิกฺขู ภควนฺตํ ทูรโตว อาคจฺฉนฺตํ, ทิสฺวาน อญฺญมญฺญํ กติกํ\csromansharedfootnote{1572d5858345b76b}{อิทํ ปทํ เกสุจิ นตฺถิ.} สณฺฐเปสุํ “อยํ อาวุโส สมโณ โคตโม อาคจฺฉติ พาหุลฺลิโก ปธาน\-วิพฺภนฺโต อาวตฺโต พาหุลฺลาย, โส เนว อภิวาเทตพฺโพ, น ปจฺจุฏฺฐาตพฺโพ, นาสฺส ปตฺตจีวรํ ปฏิคฺคเหตพฺพํ, อปิ จ โข อาสนํ ฐเปตพฺพํ, สเจ โส อากงฺขิสฺสติ, นิสีทิสฺสตี”ติ. ยถา ยถา โข ภควา ปญฺจวคฺคิเย ภิกฺขู อุปสงฺกมติ ตถา ตถา\csromansharedfootnote{3d8bc2c05ee0c6bb}{ตถา ตถา เต (สี, สฺยา)} ปญฺจวคฺคิยา ภิกฺขู นาสกฺขิํสุ สกาย กติกาย \csromanfolio{13}สณฺฐาตุํ, อสณฺฐหนฺตา ภควนฺตํ ปจฺจุคฺคนฺตฺวา เอโก ภควโต ปตฺตจีวรํ ปฏิคฺคเหสิ, เอโก อาสนํ ปญฺญเปสิ, เอโก ปาโททกํ เอโก ปาทปีฐํ เอโก ปาทกฐลิกํ อุปนิกฺขิปิ. นิสีทิ ภควา ปญฺญตฺเต อาสเน, นิสชฺช โข ภควา ปาเท ปกฺขาเลสิ. อปิสฺสุ\csromansharedfootnote{287f6914660a9940}{อปิ จ โข (ปาสราสิสุตฺต)} ภควนฺตํ นาเมน จ อาวุโสวาเทน จ สมุทาจรนฺติ. เอวํ วุตฺเต ภควา ปญฺจวคฺคิเย ภิกฺขู เอตทโวจ “มา ภิกฺขเว ตถาคตํ นาเมน จ อาวุโสวาเทน จ สมุทาจรถ\csromansharedfootnote{c2aff3d834c6f10a}{สมุทาจริตฺถ (สี, สฺยา)}, อรหํ ภิกฺขเว ตถาคโต สมฺมาสมฺพุทฺโธ, โอทหถ ภิกฺขเว โสตํ, อมตม\-ธิคตํ, อหมนุสาสามิ, อหํ ธมฺมํ เทเสมิ, ยถานุสิฏฺฐํ ตถา ปฏิปชฺชมานา\csromansharedfootnote{86367d76c5c77ae8}{ยถานุสิฏฺฐํ ปฏิปชฺชมานา (สฺยา)} น จิรสฺเสว ยสฺสตฺถาย กุลปุตฺตา สมฺมเทว อคารสฺมา อนคาริยํ ปพฺพชนฺติ, ตทนุตฺตรํ พฺรหฺมจริย\-ปริโยสานํ ทิฏฺเฐวธมฺเม สยํ อภิญฺญา สจฺฉิกตฺวา อุปสมฺปชฺช วิหริสฺสถา”ติ. เอวํ วุตฺเต ปญฺจวคฺคิยา ภิกฺขู ภควนฺตํ เอตทโวจุํ “ตายปิ โข ตฺวํ อาวุโส โคตม อิริยาย\csromansharedfootnote{119087231df48186}{จริยาย (สฺยา)} ตาย ปฏิปทาย ตาย ทุกฺกร\-การิกาย เนวชฺฌคา อุตฺตริ มนุสฺสธมฺมา\csromansharedfootnote{e7239d21ee4bec47}{อุตฺตริมนุสฺสธมฺมํ (สฺยา, ก)} อลม\-ริย\-ญาณ\-ทสฺสน\-วิเสสํ, กิํ ปน ตฺวํ เอตรหิ พาหุลฺลิโก ปธาน\-วิพฺภนฺโต อาวตฺโต พาหุลฺลาย อธิคมิสฺสสิ อุตฺตริ มนุสฺสธมฺมา อลมริย\-ญาณทสฺสนวิเสสนฺ”ติ. เอวํ วุตฺเต ภควา ปญฺจวคฺคิเย ภิกฺขู เอตทโวจ “น ภิกฺขเว ตถาคโต พาหุลฺลิโก, น ปธาน\-วิพฺภนฺโต, น อาวตฺโต พาหุลฺลาย, อรหํ ภิกฺขเว ตถาคโต สมฺมาสมฺพุทฺโธ, โอทหถ ภิกฺขเว โสตํ, อมตม\-ธิคตํ, อหมนุสาสามิ, อหํ ธมฺมํ เทเสมิ, ยถานุสิฏฺฐํ ตถา ปฏิปชฺชมานา น จิรสฺเสว ยสฺสตฺถาย กุลปุตฺตา สมฺมเทว อคารสฺมา อนคาริยํ ปพฺพชนฺติ, ตทนุตฺตรํ พฺรหฺมจริย\-ปริโยสานํ ทิฏฺเฐวธมฺเม สยํ อภิญฺญา สจฺฉิกตฺวา อุปสมฺปชฺช วิหริสฺสถา”ติ. ทุติยมฺปิ โข ปญฺจวคฺคิยา ภิกฺขู ภควนฺตํ เอตทโวจุํ ฯเปฯ. ทุติยมฺปิ โข ภควา ปญฺจวคฺคิเย ภิกฺขู เอตทโวจ ฯเปฯ. ตติยมฺปิ โข ปญฺจวคฺคิยา ภิกฺขู ภควนฺตํ เอตทโวจุํ “ตาย ปิ โข ตฺวํ อาวุโส โคตม อิริยาย ตาย ปฏิปทาย ตาย ทุกฺกร\-การิกาย เนวชฺฌคา อุตฺตริ มนุสฺสธมฺมา อลม\-ริย\-ญาณ\-ทสฺสน\-วิเสสํ, กิํ ปน ตฺวํ เอตรหิ พาหุลฺลิโก \csromanfolio{14}ปธาน\-วิพฺภนฺโต อาวตฺโต พาหุลฺลาย อธิคมิสฺสสิ อุตฺตริ มนุสฺสธมฺมา อลมริย\-ญาณทสฺสนวิเสสนฺ”ติ. เอวํ วุตฺเต ภควา ปญฺจวคฺคิเย ภิกฺขู เอตทโวจ “อภิชานาถ เม โน ตุเมฺห ภิกฺขเว อิโต ปุพฺเพ เอวรูปํ ปภาวิตเมตนฺ”ติ\csromansharedfootnote{e7114b840f7d1f88}{ภาสิตเมตนฺติ (สี, สฺยา, ก) ฏีกาโย โอโลเกตพฺพา.}. โนเหหํ ภนฺเต. อรหํ ภิกฺขเว ตถาคโต สมฺมาสมฺพุทฺโธ, โอทหถ ภิกฺขเว โสตํ, อมตม\-ธิคตํ, อหมนุสาสามิ, อหํ ธมฺมํ เทเสมิ, ยถานุสิฏฺฐํ ตถา ปฏิปชฺชมานา น จิรสฺเสว ยสฺสตฺถาย กุลปุตฺตา สมฺมเทว อคารสฺมา อนคาริยํ ปพฺพชนฺติ, ตทนุตฺตรํ พฺรหฺมจริย\-ปริโยสานํ ทิฏฺเฐวธมฺเม สยํ อภิญฺญา สจฺฉิกตฺวา อุปสมฺปชฺช วิหริสฺสถาติ. อสกฺขิ โข ภควา ปญฺจวคฺคิเย ภิกฺขู สญฺญาเปตุํ. อถ โข ปญฺจวคฺคิยา ภิกฺขู ภควนฺตํ สุสฺสูสิํสุ, โสตํ โอทหิํสุ, อญฺญา จิตฺตํ อุปฏฺฐาเปสุํ.}

\proseitem{13}{อถ โข ภควา ปญฺจวคฺคิเย ภิกฺขู อามนฺเตสิ\csromandash{}}

\prose{\csromansymbolfootnote{*}{สํ 3. 368 ปิฏฺฐาทีสุปิ.}เทฺว’เม ภิกฺขเว อนฺตา ปพฺพชิเตน น เสวิตพฺพา. กตเม เทฺว\csromansharedfootnote{805459fff1190d69}{อิทํ ปททฺวยํ สีสฺยาโปตฺถเกสุ นตฺถิ.}, โย จายํ กาเมสุ กามสุขลฺลิกา\-นุโยโค หีโน คมฺโม โปถุชฺชนิโก อนริโย อนตฺถสํหิโต, โย จายํ อตฺต\-กิลมถา\-นุโยโค ทุกฺโข อนริโย อนตฺถสํหิโต. เอเต โข ภิกฺขเว อุโภ อนฺเต อนุปคมฺม มชฺฌิมา ปฏิปทา ตถาคเตน อภิสมฺพุทฺธา จกฺขุกรณี ญาณกรณี อุปสมาย อภิญฺญาย สมฺโพธาย นิพฺพานาย สํวตฺตติ. กตมา จ สา ภิกฺขเว มชฺฌิมา ปฏิปทา ตถาคเตน อภิสมฺพุทฺธา จกฺขุกรณี ญาณกรณี อุปสมาย อภิญฺญาย สมฺโพธาย นิพฺพานาย สํวตฺตติ, อยเมว อริโย อฏฺฐงฺคิโก มคฺโค. เสยฺยถิทํ, สมฺมาทิฏฺฐิ สมฺมาสงฺกปฺโป สมฺมาวาจา สมฺมากมฺมนฺโต สมฺมาอาชีโว สมฺมาวายาโม สมฺมาสติ สมฺมาสมาธิ. อยํ โข สา ภิกฺขเว มชฺฌิมา ปฏิปทา ตถาคเตน อภิสมฺพุทฺธา จกฺขุกรณี ญาณกรณี อุปสมาย อภิญฺญาย สมฺโพธาย นิพฺพานาย สํวตฺตติ.}

\proseitem{14}{อิทํ โข ปน ภิกฺขเว ทุกฺขํ อริยสจฺจํ, ชาติปิ ทุกฺขา, ชราปิ ทุกฺขา, พฺยาธิปิ ทุกฺโข, มรณมฺปิ ทุกฺขํ, อปฺปิเยหิ สมฺปโยโค ทุกฺโข, ปิเยหิ วิปฺปโยโค ทุกฺโข, ยมฺปิจฺฉํ น ลภติ ตมฺปิ ทุกฺขํ, สํขิตฺเตน \csromanfolio{15}ปญฺจุ\-ปาทาน\-กฺขนฺธา\csromansharedfootnote{feb51bbb692b3682}{ปญฺจุปาทานกฺขนฺธาปิ (ก)} ทุกฺขา. อิทํ โข ปน ภิกฺขเว ทุกฺข\-สมุทยํ\csromansharedfootnote{01e03502805da113}{เอตฺถ “อิทํ ทุกฺขํ อริยสจฺจนฺติ อาทีสุ ทุกฺขสมุทโย ทุกฺขนิโรโธติ วตฺตพฺเพ ทุกฺขสมุทยํ ทุกฺขนิโรธนฺติ ลิงฺควิปลฺลาโส กโต”ติ ปฏิสมฺภิทามคฺคฏฺฐกถายํ วุตฺตํ. วิสุทฺธิมคฺคฏีกายํ ปน อุปฺปาโท ภยนฺติปาฐวณฺณนายํ “สติปิ ทฺวินฺนํ ปทานํ สมานาธิกรณภาเว ลิงฺคเภโท คหิโต, ยถา ทุกฺขสมุทโย อริยสจฺจนฺ”ติ วุตฺตํ. เตสุ “ทุกฺขสมุทโย อริยสจฺจนฺ”ติ สกลิงฺคิกปาโฐ “ทุกฺขนิโรธคามินี ปฏิปทา อริยสจฺจนฺ”ติ ปาฬิยา สเมติ.} อริยสจฺจํ, ยายํ ตณฺหา โปโนพฺภวิกา\csromansharedfootnote{1e8468b810b12adc}{โปโนภวิกา (ก)} นนฺที\-ราค\-สหคตา\csromansharedfootnote{4fbe7f645803518e}{นนฺทิราคสหคตา (สี, สฺยา)} ตตฺร\-ตตฺรา\-ภินนฺทินี. เสยฺยถิทํ, กามตณฺหา ภวตณฺหา วิภวตณฺหา. อิทํ โข ปน ภิกฺขเว ทุกฺขนิโรธํ\csromansharedfootnote{01e03502805da113}{เอตฺถ “อิทํ ทุกฺขํ อริยสจฺจนฺติ อาทีสุ ทุกฺขสมุทโย ทุกฺขนิโรโธติ วตฺตพฺเพ ทุกฺขสมุทยํ ทุกฺขนิโรธนฺติ ลิงฺควิปลฺลาโส กโต”ติ ปฏิสมฺภิทามคฺคฏฺฐกถายํ วุตฺตํ. วิสุทฺธิมคฺคฏีกายํ ปน อุปฺปาโท ภยนฺติปาฐวณฺณนายํ “สติปิ ทฺวินฺนํ ปทานํ สมานาธิกรณภาเว ลิงฺคเภโท คหิโต, ยถา ทุกฺขสมุทโย อริยสจฺจนฺ”ติ วุตฺตํ. เตสุ “ทุกฺขสมุทโย อริยสจฺจนฺ”ติ สกลิงฺคิกปาโฐ “ทุกฺขนิโรธคามินี ปฏิปทา อริยสจฺจนฺ”ติ ปาฬิยา สเมติ.} อริยสจฺจํ, โย ตสฺสา เยว\csromansharedfootnote{19cf0a39ee311ce1}{ตสฺสาเยว สนฺธิ \csromandash{} ม.พ.ป.} ตณฺหาย อเสส\-วิราค\-นิโรโธ จาโค ปฏินิสฺสคฺโค มุตฺติ อนาลโย. อิทํ โข ปน ภิกฺขเว ทุกฺขนิโรธ\-คามินี ปฏิปทา อริยสจฺจํ, อยเมว อริโย อฏฺฐงฺคิโก มคฺโค. เสยฺยถิทํ, สมฺมาทิฏฺฐิ สมฺมาสงฺกปฺโป สมฺมาวาจา สมฺมากมฺมนฺโต สมฺมาอาชีโว สมฺมาวายาโม สมฺมาสติ สมฺมาสมาธิ.}

\proseitem{15}{“อิทํ ทุกฺขํ อริยสจฺจนฺ”ติ เม ภิกฺขเว ปุพฺเพ อนนุสฺสุเตสุ ธมฺเมสุ จกฺขุํ อุทปาทิ ญาณํ อุทปาทิ ปญฺญา อุทปาทิ วิชฺชา อุทปาทิ อาโลโก อุทปาทิ. “ตํ โข ปนิทํ ทุกฺขํ อริยสจฺจํ ปริญฺเญยฺยนฺ”ติ เม ภิกฺขเว ปุพฺเพ อนนุสฺสุเตสุ ธมฺเมสุ จกฺขุํ อุทปาทิ ญาณํ อุทปาทิ ปญฺญา อุทปาทิ วิชฺชา อุทปาทิ อาโลโก อุทปาทิ. “ตํ โข ปนิทํ ทุกฺขํ อริยสจฺจํ ปริญฺญาตนฺ”ติ เม ภิกฺขเว ปุพฺเพ อนนุสฺสุเตสุ ธมฺเมสุ จกฺขุํ อุทปาทิ ญาณํ อุทปาทิ ปญฺญา อุทปาทิ วิชฺชา อุทปาทิ อาโลโก อุทปาทิ.}

\prose{“อิทํ ทุกฺข\-สมุทยํ อริยสจฺจนฺ”ติ เม ภิกฺขเว ปุพฺเพ อนนุสฺสุเตสุ ธมฺเมสุ จกฺขุํ อุทปาทิ ญาณํ อุทปาทิ ปญฺญา อุทปาทิ วิชฺชา อุทปาทิ อาโลโก อุทปาทิ. “ตํ โข ปนิทํ ทุกฺข\-สมุทยํ อริยสจฺจํ ปหาตพฺพนฺ”ติ เม ภิกฺขเว ปุพฺเพ อนนุสฺสุเตสุ ธมฺเมสุ จกฺขุํ อุทปาทิ ญาณํ อุทปาทิ ปญฺญา อุทปาทิ วิชฺชา อุทปาทิ อาโลโก อุทปาทิ. “ตํ โข ปนิทํ ทุกฺข\-สมุทยํ อริยสจฺจํ ปหีนนฺ”ติ เม ภิกฺขเว ปุพฺเพ อนนุสฺสุเตสุ ธมฺเมสุ จกฺขุํ อุทปาทิ ญาณํ อุทปาทิ ปญฺญา อุทปาทิ วิชฺชา อุทปาทิ อาโลโก อุทปาทิ.}

\prose{“อิทํ ทุกฺขนิโรธํ อริยสจฺจนฺ”ติ เม ภิกฺขเว ปุพฺเพ อนนุสฺสุเตสุ ธมฺเมสุ จกฺขุํ อุทปาทิ ญาณํ อุทปาทิ ปญฺญา อุทปาทิ วิชฺชา อุทปาทิ \csromanfolio{16}อาโลโก อุทปาทิ. “ตํ โข ปนิทํ ทุกฺขนิโรธํ อริยสจฺจํ สจฺฉิกาตพฺพนฺ”ติ เม ภิกฺขเว ปุพฺเพ อนนุสฺสุเตสุ ธมฺเมสุ จกฺขุํ อุทปาทิ ญาณํ อุทปาทิ ปญฺญา อุทปาทิ วิชฺชา อุทปาทิ อาโลโก อุทปาทิ. “ตํ โข ปนิทํ ทุกฺขนิโรธํ อริยสจฺจํ สจฺฉิกตนฺ”ติ เม ภิกฺขเว ปุพฺเพ อนนุสฺสุเตสุ ธมฺเมสุ จกฺขุํ อุทปาทิ ญาณํ อุทปาทิ ปญฺญา อุทปาทิ วิชฺชา อุทปาทิ อาโลโก อุทปาทิ.}

\prose{“อิทํ ทุกฺขนิโรธ\-คามินี ปฏิปทา อริยสจฺจนฺ”ติ เม ภิกฺขเว ปุพฺเพ อนนุสฺสุเตสุ ธมฺเมสุ จกฺขุํ อุทปาทิ ญาณํ อุทปาทิ ปญฺญา อุทปาทิ วิชฺชา อุทปาทิ อาโลโก อุทปาทิ. “ตํ โข ปนิทํ ทุกฺขนิโรธ\-คามินี ปฏิปทา อริยสจฺจํ ภาเวตพฺพนฺ”ติ เม ภิกฺขเว ปุพฺเพ อนนุสฺสุเตสุ ธมฺเมสุ จกฺขุํ อุทปาทิ ญาณํ อุทปาทิ ปญฺญา อุทปาทิ วิชฺชา อุทปาทิ อาโลโก อุทปาทิ. “ตํ โข ปนิทํ ทุกฺขนิโรธ\-คามินี ปฏิปทา อริยสจฺจํ ภาวิตนฺ”ติ เม ภิกฺขเว ปุพฺเพ อนนุสฺสุเตสุ ธมฺเมสุ จกฺขุํ อุทปาทิ ญาณํ อุทปาทิ ปญฺญา อุทปาทิ วิชฺชา อุทปาทิ อาโลโก อุทปาทิ.}

\proseitem{16}{ยาวกีวญฺจ เม ภิกฺขเว อิเมสุ จตูสุ อริยสจฺเจสุ เอวํ ติปริวฏฺฏํ ทฺวาทสาการํ ยถาภูตํ ญาณทสฺสนํ น สุวิสุทฺธํ อโหสิ, เนว ตาวาหํ ภิกฺขเว สเทวเก โลเก สมารเก สพฺรหฺมเก สสฺสมณ\-พฺราหฺมณิยา ปชาย สเทว\-มนุสฺสาย “อนุตฺตรํ สมฺมาสมฺโพธิํ อภิสมฺพุทฺโธ”ติ\csromansharedfootnote{e01903969db9d80f}{อภิสมฺพุทฺโธ (สี, สฺยา)} ปจฺจญฺญาสิํ.}

\prose{ยโต จ โข เม ภิกฺขเว อิเมสุ จตูสุ อริยสจฺเจสุ เอวํ ติปริวฏฺฏํ ทฺวาทสาการํ ยถาภูตํ ญาณทสฺสนํ สุวิสุทฺธํ อโหสิ, อถาหํ ภิกฺขเว สเทวเก โลเก สมารเก สพฺรหฺมเก สสฺสมณ\-พฺราหฺมณิยา ปชาย สเทว\-มนุสฺสาย “อนุตฺตรํ สมฺมาสมฺโพธิํ อภิสมฺพุทฺโธ”ติ\csromansharedfootnote{61785d03c494c2ae}{[มิสฺสินฺคฺ โนเต 1]} ปจฺจญฺญาสิํ.}

\prose{ญาณญฺจ ปน เม ทสฺสนํ อุทปาทิ, อกุปฺปา เม วิมุตฺติ, อยมนฺติมา ชาติ, นตฺถิ ทานิ ปุนพฺภโวติ.}

\prose{อิทมโวจ ภควา, อตฺตมนา ปญฺจวคฺคิยา ภิกฺขู ภควโต ภาสิตํ อภินนฺทุนฺติ\csromansharedfootnote{a4737a24874c2654}{[มิสฺสินฺคฺ โนเต 2]}.}

\prose{อิมสฺมิํ จ ปน เวยฺยากรณสฺมิํ ภญฺญมาเน อายสฺมโต โกณฺฑญฺญสฺส วิรชํ วีตมลํ ธมฺมจกฺขุํ อุทปาทิ “ยํ กิญฺจิ สมุทย\-ธมฺมํ, สพฺพํ ตํ นิโรธธมฺมนฺ”ติ.}

\proseitem{17}{\csromanfolio{17}ปวตฺติเต จ ปน ภควตา ธมฺมจกฺเก ภุมฺมา เทวา สทฺทม\-นุสฺสาเวสุํ “เอตํ ภควตา พาราณสิยํ อิสิปตเน มิคทาเย อนุตฺตรํ ธมฺมจกฺกํ ปวตฺติตํ อปฺปฏิวตฺติยํ สมเณน วา พฺราหฺมเณน วา เทเวน วา มาเรน วา พฺรหฺมุนา วา เกนจิ วา โลกสฺมิํ”ติ. ภุมฺมานํ เทวานํ สทฺทํ สุตฺวา จาตุมหาราชิกา เทวา สทฺทม\-นุสฺสาเวสุํ ฯเปฯ. จาตุมหาราชิกานํ เทวานํ สทฺทํ สุตฺวา ตาวติํสา เทวา ฯเปฯ ยามา เทวา ฯเปฯ ตุสิตา เทวา ฯเปฯ นิมฺมานรตี เทวา ฯเปฯ ปรนิมฺมิต\-วส\-วตฺตี เทวา ฯเปฯ พฺรหฺมกายิกา เทวา สทฺทม\-นุสฺสาเวสุํ “เอตํ ภควตา พาราณสิยํ อิสิปตเน มิคทาเย อนุตฺตรํ ธมฺมจกฺกํ ปวตฺติตํ อปฺปฏิวตฺติยํ สมเณน วา พฺราหฺมเณน วา เทเวน วา มาเรน วา พฺรหฺมุนา วา เกนจิ วา โลกสฺมิํ”ติ. อิติห เตน ขเณน เตน ลเยน\csromansharedfootnote{0c626d1232646e3c}{เตน ลเยนาติ ปททฺวยํ สีสฺยาโปตฺถเกสุ นตฺถิ.} เตน มุหุตฺเตน ยาว พฺรหฺมโลกา สทฺโท อพฺภุคฺคจฺฉิ. อยญฺจ ทสสหสฺสิ\-โลกธาตุ สํกมฺปิ สมฺปกมฺปิ สมฺปเวธิ, อปฺปมาโณ จ อุฬาโร โอภาโส โลเก ปาตุรโหสิ อติกฺกมฺม เทวานํ เทวานุภาวํ. อถ โข ภควา อิมํ อุทานํ อุทาเนสิ “อญฺญาสิ วต โภ โกณฺฑญฺโญ, อญฺญาสิ วต โภ โกณฺฑญฺโญ”ติ. อิติ หิทํ อายสฺมโต โกณฺฑญฺญสฺส “อญฺญาสิโกณฺฑญฺโญ”เตฺวว นามํ อโหสิ.}

\proseitem{18}{อถ โข อายสฺมา อญฺญาสิโกณฺฑญฺโญ ทิฏฺฐธมฺโม ปตฺตธมฺโม วิทิตธมฺโม ปริโยคาฬฺห\-ธมฺโม ติณฺณ\-วิจิกิจฺโฉ วิคต\-กถํกโถ เวสารชฺช\-ปฺปตฺโต อปรปฺปจฺจโย สตฺถุสาสเน ภควนฺตํ เอตทโวจ “ลเภยฺยาหํ ภนฺเต ภควโต สนฺติเก ปพฺพชฺชํ, ลเภยฺยํ อุปสมฺปทนฺ”ติ. “เอหิ ภิกฺขู”ติ ภควา อโวจ “สฺวากฺขาโต ธมฺโม, จร พฺรหฺมจริยํ สมฺมา ทุกฺขสฺส อนฺตกิริยายา”ติ. สาว ตสฺส อายสฺมโต อุปสมฺปทา อโหสิ.}

\proseitem{19}{อถ โข ภควา ตทวเสเส ภิกฺขู ธมฺมิยา กถาย โอวทิ อนุสาสิ. อถ โข อายสฺมโต จ วปฺปสฺส อายสฺมโต จ ภทฺทิยสฺส ภควตา ธมฺมิยา กถาย โอวทิยมานานํ อนุสาสิย\-มานานํ วิรชํ วีตมลํ ธมฺมจกฺขุํ อุทปาทิ “ยํ กิญฺจิ สมุทย\-ธมฺมํ, สพฺพํ ตํ นิโรธธมฺมนฺ”ติ.}

\prose{\csromanfolio{18}เต ทิฏฺฐธมฺมา ปตฺตธมฺมา วิทิตธมฺมา ปริโยคาฬฺห\-ธมฺมา ติณฺณ\-วิจิกิจฺฉา วิคต\-กถํกถา เวสารชฺช\-ปฺปตฺตา อปรปฺปจฺจยา สตฺถุสาสเน ภควนฺตํ เอตทโวจุํ “ลเภยฺยาม มยํ ภนฺเต ภควโต สนฺติเก ปพฺพชฺชํ, ลเภยฺยาม อุปสมฺปทนฺ”ติ.}

\prose{“เอถ ภิกฺขโว”ติ ภควา อโวจ “สฺวากฺขาโต ธมฺโม, จรถ พฺรหฺมจริยํ สมฺมา ทุกฺขสฺส อนฺตกิริยายา”ติ.}

\setlength{\csromangathapagewidth}{0pt}
\setlength{\csromangathawakwidth}{0pt}
\csromangathameasuregroup{}{สาว เตสํ อายสฺมนฺตานํ อุปสมฺปทา อโหสิ.}
\csromangathameasurewak{สาว เตสํ อายสฺมนฺตานํ อุปสมฺปทา อโหสิ.}
\setlength{\csromangathaleftcol}{0pt}
\csromangathagroup{\csromangathastanza{\csromangathawak{สาว เตสํ อายสฺมนฺตานํ อุปสมฺปทา อโหสิ.}}}

\prose{อถ โข ภควา ตทวเสเส ภิกฺขู นีหารภตฺโต ธมฺมิยา กถาย โอวทิ อนุสาสิ. ยํ ตโย ภิกฺขู ปิณฺฑาย จริตฺวา อาหรนฺติ, เตน ฉพฺพคฺโค ยาเปติ. อถ โข อายสฺมโต จ มหานามสฺส อายสฺมโต จ อสฺสชิสฺส ภควตา ธมฺมิยา กถาย โอวทิยมานานํ อนุสาสิย\-มานานํ วิรชํ วีตมลํ ธมฺมจกฺขุํ อุทปาทิ “ยํ กิญฺจิ สมุทย\-ธมฺมํ, สพฺพํ ตํ นิโรธธมฺมนฺ”ติ. เต ทิฏฺฐธมฺมา ปตฺตธมฺมา วิทิตธมฺมา ปริโยคาฬฺห\-ธมฺมา ติณฺณ\-วิจิกิจฺฉา วิคต\-กถํกถา เวสารชฺช\-ปฺปตฺตา อปรปฺปจฺจยา สตฺถุสาสเน ภควนฺตํ เอตทโวจุํ “ลเภยฺยาม มยํ ภนฺเต ภควโต สนฺติเก ปพฺพชฺชํ, ลเภยฺยาม อุปสมฺปทนฺ”ติ. “เอถ ภิกฺขโว”ติ ภควา อโวจ “สฺวากฺขาโต ธมฺโม, จรถ พฺรหฺมจริยํ สมฺมา ทุกฺขสฺส อนฺตกิริยายา”ติ. สาว เตสํ อายสฺมนฺตานํ อุปสมฺปทา อโหสิ.}

\proseitem{20}{อถ โข ภควา ปญฺจวคฺคิเย ภิกฺขู อามนฺเตสิ\csromandash{}}

\prose{\csromansymbolfootnote{*}{สํ 2. 55 ปิฏฺฐาทีสุปิ.}รูปํ ภิกฺขเว อนตฺตา, รูปญฺจ หิท ภิกฺขเว อตฺตา อภวิสฺส, นยิทํ รูปํ อาพาธาย สํวตฺเตยฺย, ลพฺเภถ จ รูเป “เอวํ เม รูปํ โหตุ, เอวํ เม รูปํ มา อโหสี”ติ. ยสฺมา จ โข ภิกฺขเว รูปํ อนตฺตา, ตสฺมา รูปํ อาพาธาย สํวตฺตติ, น จ ลพฺภติ รูเป “เอวํ เม รูปํ โหตุ, เอวํ เม รูปํ มา อโหสี”ติ. เวทนา อนตฺตา, เวทนา จ หิทํ ภิกฺขเว อตฺตา อภวิสฺส, นยิทํ เวทนา อาพาธาย สํวตฺเตยฺย, ลพฺเภถ จ เวทนาย “เอวํ เม เวทนา โหตุ, เอวํ เม เวทนา มา อโหสี”ติ. ยสฺมา จ โข ภิกฺขเว เวทนา อนตฺตา, ตสฺมา เวทนา อาพาธาย สํวตฺตติ, น จ ลพฺภติ เวทนาย “เอวํ เม เวทนา โหตุ, เอวํ เม เวทนา มา อโหสี”ติ. สญฺญา อนตฺตา, สญฺญา จ หิทํ ภิกฺขเว อตฺตา อภวิสฺส, นยิทํ สญฺญา อาพาธาย สํวตฺเตยฺย, \csromanfolio{19}ลพฺเภถ จ สญฺญาย “เอวํ เม สญฺญา โหตุ, เอวํ เม สญฺญา มา อโหสี”ติ. ยสฺมา จ โข ภิกฺขเว สญฺญา อนตฺตา, ตสฺมา สญฺญา อาพาธาย สํวตฺตติ, น จ ลพฺภติ สญฺญาย “เอวํ เม สญฺญา โหตุ, เอวํ เม สญฺญา มา อโหสี”ติ. สงฺขารา อนตฺตา, สงฺขารา จ หิทํ ภิกฺขเว อตฺตา อภวิสฺสํสุ, นยิทํ\csromansharedfootnote{46472e5266cc95ef}{นยิเม (ก)} สงฺขารา อาพาธาย สํวตฺเตยฺยุํ, ลพฺเภถ จ สงฺขาเรสุ “เอวํ เม สงฺขารา โหนฺตุ, เอวํ เม สงฺขารา มา อเหสุนฺ”ติ. ยสฺมา จ โข ภิกฺขเว สงฺขารา อนตฺตา, ตสฺมา สงฺขารา อาพาธาย สํวตฺตนฺติ, น จ ลพฺภติ สงฺขาเรสุ “เอวํ เม สงฺขารา โหนฺตุ, เอวํ เม สงฺขารา มา อเหสุนฺ”ติ. วิญฺญาณํ อนตฺตา, วิญฺญาณญฺจ หิทํ ภิกฺขเว อตฺตา อภวิสฺส, นยิทํ วิญฺญาณํ อาพาธาย สํวตฺเตยฺย, ลพฺเภถ จ วิญฺญาเณ “เอวํ เม วิญฺญาณํ โหตุ, เอวํ เม วิญฺญาณํ มา อโหสี”ติ. ยสฺมา จ โข ภิกฺขเว วิญฺญาณํ อนตฺตา, ตสฺมา วิญฺญาณํ อาพาธาย สํวตฺตติ, น จ ลพฺภติ วิญฺญาเณ “เอวํ เม วิญฺญาณํ โหตุ, เอวํ เม วิญฺญาณํ มา อโหสี”ติ.}

\proseitem{21}{ตํ กิํ มญฺญถ ภิกฺขเว, รูปํ นิจฺจํ วา อนิจฺจํ วาติ. อนิจฺจํ ภนฺเต. ยํ ปนานิจฺจํ, ทุกฺขํ วา ตํ สุขํ วาติ. ทุกฺขํ ภนฺเต. ยํ ปนานิจฺจํ ทุกฺขํ วิปริณาม\-ธมฺมํ, กลฺลํ นุ ตํ สมนุปสฺสิตุํ “เอตํ มม, เอโสหมสฺมิ, เอโส เม อตฺตา”ติ. โน เหตํ ภนฺเต. เวทนา นิจฺจา วา อนิจฺจา วาติ. อนิจฺจา ภนฺเต. ยํ ปนานิจฺจํ, ทุกฺขํ วา ตํ สุขํ วาติ. ทุกฺขํ ภนฺเต. ยํ ปนานิจฺจํ ทุกฺขํ วิปริณาม\-ธมฺมํ, กลฺลํ นุ ตํ สมนุปสฺสิตุํ “เอตํ มม, เอโสหมสฺมิ, เอโส เม อตฺตา”ติ. โน เหตํ ภนฺเต. สญฺญา นิจฺจา วา อนิจฺจา วาติ. อนิจฺจา ภนฺเต. ยํ ปนานิจฺจํ ทุกฺขํ วา ตํ สุขํ วาติ. ทุกฺขํ ภนฺเต. ยํ ปนานิจฺจํ ทุกฺขํ วิปริณาม\-ธมฺมํ, กลฺลํ นุ ตํ สมนุปสฺสิตุํ “เอตํ มม, เอโสหมสฺมิ, เอโส เม อตฺตา”ติ. โน เหตํ ภนฺเต. สงฺขารา นิจฺจา วา อนิจฺจา วาติ. อนิจฺจา ภนฺเต. ยํ ปนานิจฺจํ, ทุกฺขํ วา ตํ สุขํ วาติ. ทุกฺขํ ภนฺเต. ยํ ปนานิจฺจํ ทุกฺขํ วิปริณาม\-ธมฺมํ, กลฺลํ นุ ตํ สมนุปสฺสิตุํ “เอตํ มม, เอโสหมสฺมิ, เอโส เม อตฺตา”ติ. โน เหตํ ภนฺเต. วิญฺญาณํ นิจฺจํ วา อนิจฺจํ วาติ. อนิจฺจํ ภนฺเต. ยํ ปนานิจฺจํ, ทุกฺขํ วา ตํ สุขํ วาติ. ทุกฺขํ ภนฺเต. ยํ ปนานิจฺจํ ทุกฺขํ วิปริณาม\-ธมฺมํ, กลฺลํ นุ ตํ สมนุปสฺสิตุํ “เอตํ มม, เอโสหมสฺมิ, เอโส เม อตฺตา”ติ. โน เหตํ ภนฺเต.}

\proseitem{22}{\csromanfolio{20}ตสฺมาติห ภิกฺขเว ยํ กิญฺจิ รูปํ อตีตา\-นาคต\-ปจฺจุปฺปนฺนํ อชฺฌตฺตํ วา พหิทฺธา วา โอฬาริกํ วา สุขุมํ วา หีนํ วา ปณีตํ วา, ยํ ทูเร\csromansharedfootnote{6afe1bc7fa1c5765}{ยํ ทูเร วา (สฺยา)} สนฺติเก วา, สพฺพํ รูปํ “เนตํ มม, เนโสหมสฺมิ, น เมโส อตฺตา”ติ เอวเมตํ ยถาภูตํ สมฺมปฺปญฺญาย ทฏฺฐพฺพํ. ยา กาจิ เวทนา อตีตา\-นาคต\-ปจฺจุปฺปนฺนา อชฺฌตฺตํ วา พหิทฺธา วา โอฬาริกา วา สุขุมา วา หีนา วา ปณีตา วา, ยา ทูเร สนฺติเก วา, สพฺพา เวทนา “เนตํ มม, เนโสหมสฺมิ, น เมโส อตฺตา”ติ เอวเมตํ ยถาภูตํ สมฺมปฺปญฺญาย ทฏฺฐพฺพํ. ยา กาจิ สญฺญา อตีตา\-นาคต\-ปจฺจุปฺปนฺนา อชฺฌตฺตํ วา พหิทฺธา วา โอฬาริกา วา สุขุมา วา หีนา วา ปณีตา วา, ยา ทูเร สนฺติเก วา, สพฺพา สญฺญา “เนตํ มม, เนโสหมสฺมิ, น เมโส อตฺตา”ติ เอวเมตํ ยถาภูตํ สมฺมปฺปญฺญาย ทฏฺฐพฺพํ. เย เกจิ สงฺขารา อตีตา\-นาคต\-ปจฺจุปฺปนฺนา อชฺฌตฺตํ วา พหิทฺธา วา โอฬาริกา วา สุขุมา วา หีนา วา ปณีตา วา, เย ทูเร สนฺติเก วา, สพฺเพ สงฺขารา “เนตํ มม, เนโสหมสฺมิ, น เมโส อตฺตา”ติ เอวเมตํ ยถาภูตํ สมฺมปฺปญฺญาย ทฏฺฐพฺพํ. ยํ กิญฺจิ วิญฺญาณํ อตีตา\-นาคต\-ปจฺจุปฺปนฺนํ อชฺฌตฺตํ วา พหิทฺธา วา โอฬาริกํ วา สุขุมํ วา หีนํ วา ปณีตํ วา, ยํ ทูเร สนฺติเก วา, สพฺพํ วิญฺญาณํ เนตํ มม, เนโสหมสฺมิ, น เมโส อตฺตา”ติ เอวเมตํ ยถาภูตํ สมฺมปฺปญฺญาย ทฏฺฐพฺพํ.}

\proseitem{23}{เอวํ ปสฺสํ ภิกฺขเว สุตวา อริยสาวโก รูปสฺมิมฺปิ นิพฺพินฺทติ, เวทนายปิ นิพฺพินฺทติ, สญฺญายปิ นิพฺพินฺทติ, สงฺขาเรสุปิ นิพฺพินฺทติ, วิญฺญาณสฺมิมฺปิ นิพฺพินฺทติ, นิพฺพินฺทํ วิรชฺชติ, วิราคา วิมุจฺจติ, วิมุตฺตสฺมิํ วิมุตฺตมิติ ญาณํ โหติ, “ขีณา ชาติ, วุสิตํ พฺรหฺมจริยํ, กตํ กรณียํ, นาปรํ อิตฺถตฺตายา”ติ ปชานาตีติ.}

\proseitemwithcloser{24}{อิทมโวจ ภควา, อตฺตมนา ปญฺจวคฺคิยา ภิกฺขู ภควโต ภาสิตํ อภินนฺทุนฺติ\csromansharedfootnote{3398327fac10016c}{อภินนฺทุํ (สฺยา)}. อิมสฺมิํ จ ปน เวยฺยากรณสฺมิํ ภญฺญมาเน ปญฺจ\-วคฺคิยานํ ภิกฺขูนํ อนุปาทาย อาสเวหิ จิตฺตานิ วิมุจฺจิํสุ. เตน โข ปน สมเยน ฉ โลเก อรหนฺโต โหนฺติ.}{ปญฺจวคฺคิย\-กถา นิฏฺฐิตา.}
\csromancenter{ปฐม\-ภาณวาโร.}
\csromanmatikaanchor{mtk.8}
\csromantocmarkref{section}{7. ปพฺพชฺชากถา}{mtk.8}
\bookmark[dest={mtk.8},level=1]{7. ปพฺพชฺชากถา}
\csromanheader{1}{7. \textbf{ปพฺพชฺชากถา}}
\proseitem{25}{\csromanfolio{21}เตน โข ปน สมเยน พาราณสิยํ ยโส นาม กุลปุตฺโต เสฏฺฐิปุตฺโต สุขุมาโล โหติ, ตสฺส ตโย ปาสาทา โหนฺติ เอโก เหมนฺติโก, เอโก คิมฺหิโก, เอโก วสฺสิโก, โส วสฺสิเก ปาสาเท จตฺตาโร มาเส\csromansharedfootnote{4959b46ba7ed28d4}{วสฺสิเก ปาสาเท วสฺสิเก จตฺตาโร มาเส (สี)} นิปฺปุริเสหิ ตูริเยหิ ปริจารยมาโน น เหฏฺฐาปาสาทํ โอโรหติ.}

\prose{อถ โข ยสสฺส กุลปุตฺตสฺส ปญฺจหิ กามคุเณหิ สมปฺปิตสฺส สมงฺคีภูตสฺส ปริจารย\-มานสฺส ปฏิกจฺเจว\csromansharedfootnote{61785d03c494c2ae}{[มิสฺสินฺคฺ โนเต 1]} นิทฺทา โอกฺกมิ, ปริชนสฺสปิ นิทฺทา โอกฺกมิ, สพฺพรตฺติโย จ เตลปทีโป ฌายติ.}

\prose{อถ โข ยโส กุลปุตฺโต ปฏิกจฺเจว ปพุชฺฌิตฺวา อทฺทส สกํ ปริชนํ สุปนฺตํ อญฺญิสฺสา กจฺเฉ วีณํ อญฺญิสฺสา กณฺเฐ มุทิงฺคํ อญฺญิสฺสา กจฺเฉ อาฬมฺพรํ อญฺญํ วิเกสิกํ อญฺญํ วิกฺเขฬิกํ อญฺญา วิปฺปลปนฺติโย หตฺถปฺปตฺตํ สุสานํ มญฺเญ, ทิสฺวานสฺส อาทีนโว ปาตุรโหสิ, นิพฺพิทาย จิตฺตํ สณฺฐาสิ.}

\prose{อถ โข ยโส กุลปุตฺโต อุทานํ อุทาเนสิ “อุปทฺทุตํ วต โภ, อุปสฺสฏฺฐํ วต โภ”ติ.}

\prose{อถ โข ยโส กุลปุตฺโต สุวณฺณ\-ปาทุกาโย อาโรหิตฺวา เยน นิเวสนทฺวารํ เตนุปสงฺกมิ, อมนุสฺสา ทฺวารํ วิวริํสุ “มา ยสสฺส กุลปุตฺตสฺส โกจิ อนฺตรายมกาสิ อคารสฺมา อนคาริยํ ปพฺพชฺชายา”ติ.}

\prose{อถ โข ยโส กุลปุตฺโต เยน นครทฺวารํ เตนุปสงฺกมิ, อมนุสฺสา ทฺวารํ วิวริํสุ “มา ยสสฺส กุลปุตฺตสฺส โกจิ อนฺตรายมกาสิ อคารสฺมา อนคาริยํ ปพฺพชฺชายา”ติ.}

\setlength{\csromangathapagewidth}{0pt}
\setlength{\csromangathawakwidth}{0pt}
\csromangathameasuregroup{}{อถ โข ยโส กุลปุตฺโต เยน อิสิปตนํ มิคทาโย เตนุปสงฺกมิ.}
\csromangathameasurewak{อถ โข ยโส กุลปุตฺโต เยน อิสิปตนํ มิคทาโย เตนุปสงฺกมิ.}
\setlength{\csromangathaleftcol}{0pt}
\csromangathagroup{\csromangathastanza{\csromangathawak{อถ โข ยโส กุลปุตฺโต เยน อิสิปตนํ มิคทาโย เตนุปสงฺกมิ.}}}

\proseitem{26}{เตน โข ปน สมเยน ภควา รตฺติยา ปจฺจูสสมยํ ปจฺจุฏฺฐาย อชฺโฌกาเส จงฺกมติ, อทฺทสา โข ภควา ยสํ กุลปุตฺตํ ทูรโตว อาคจฺฉนฺตํ, ทิสฺวาน จงฺกมา โอโรหิตฺวา ปญฺญตฺเต อาสเน นิสีทิ. อถ โข ยโส กุลปุตฺโต ภควโต อวิทูเร อุทานํ อุทาเนสิ “อุปทฺทุตํ วต โภ, อุปสฺสฏฺฐํ วต โภ”ติ. อถ โข ภควา ยสํ กุลปุตฺตํ เอตทโวจ “อิทํ โข ยส อนุปทฺทุตํ, อิทํ อนุปสฺสฏฺฐํ, เอหิ ยส นิสีท, ธมฺมํ เต เทเสสฺสามี”ติ. อถ โข ยโส กุลปุตฺโต “อิทํ กิร \csromanfolio{22}อนุปทฺทุตํ, อิทํ อนุปสฺสฏฺฐนฺ”ติ หฏฺโฐ อุทคฺโค สุวณฺณ\-ปาทุกาหิ โอโรหิตฺวา เยน ภควา เตนุปสงฺกมิ, อุปสงฺกมิตฺวา ภควนฺตํ อภิวาเทตฺวา เอกมนฺตํ นิสีทิ, เอกมนฺตํ นิสินฺนสฺส โข ยสสฺส กุลปุตฺตสฺส ภควา อนุปุพฺพิํ กถํ กเถสิ, เสยฺยถิทํ, ทานกถํ สีลกถํ สคฺคกถํ กามานํ อาทีนวํ โอการํ สํกิเลสํ เนกฺขมฺเม อานิสํสํ ปกาเสสิ. ยทา ภควา อญฺญาสิ ยสํ กุลปุตฺตํ กลฺลจิตฺตํ มุทุจิตฺตํ วินีวรณ\-จิตฺตํ อุทคฺคจิตฺตํ ปสนฺนจิตฺตํ, อถ ยา พุทฺธานํ สามุกฺกํสิกา ธมฺมเทสนา, ตํ ปกาเสสิ ทุกฺขํ สมุทยํ นิโรธํ มคฺคํ, เสยฺยถาปิ นาม สุทฺธํ วตฺถํ อปคตกาฬกํ สมฺมเทว รชนํ ปฏิคฺคเณฺหยฺย, เอวเมว ยสสฺส กุลปุตฺตสฺส ตสฺมิํเยว อาสเน วิรชํ วีตมลํ ธมฺมจกฺขุํ อุทปาทิ “ยํ กิญฺจิ สมุทย\-ธมฺมํ, สพฺพํ ตํ นิโรธธมฺมนฺ”ติ.}

\proseitem{27}{อถ โข ยสสฺส กุลปุตฺตสฺส มาตา ปาสาทํ อภิรุหิตฺวา ยสํ กุลปุตฺตํ อปสฺสนฺตี เยน เสฏฺฐิ คหปติ เตนุปสงฺกมิ, อุปสงฺกมิตฺวา เสฏฺฐิํ คหปติํ เอตทโวจ “ปุตฺโต เต คหปติ ยโส น ทิสฺสตี”ติ. อถ โข เสฏฺฐิ คหปติ จตุทฺทิสา อสฺสทูเต อุยฺโยเชตฺวา สามํเยว เยน อิสิปตนํ มิคทาโย เตนุปสงฺกมิ. อทฺทสา โข เสฏฺฐิ คหปติ สุวณฺณ\-ปาทุกานํ นิกฺเขปํ, ทิสฺวาน ตํเยว อนุคมาสิ\csromansharedfootnote{ed039d3b4f89a2a8}{อนุคมา (สี, สฺยา)}. อทฺทสา โข ภควา เสฏฺฐิํ คหปติํ ทูรโตว อาคจฺฉนฺตํ, ทิสฺวาน ภควโต เอตทโหสิ “ยนฺนูนาหํ ตถารูปํ อิทฺธา\-ภิสงฺขารํ อภิสงฺขาเรยฺยํ, ยถา เสฏฺฐิ คหปติ อิธ นิสินฺโน อิธ นิสินฺนํ ยสํ กุลปุตฺตํ น ปสฺเสยฺยา”ติ, อถ โข ภควา ตถารูปํ อิทฺธา\-ภิสงฺขารํ อภิสงฺขเรสิ. อถ โข เสฏฺฐิ คหปติ เยน ภควา เตนุปสงฺกมิ, อุปสงฺกมิตฺวา ภควนฺตํ เอตทโวจ “อปิ ภนฺเต ภควา ยสํ กุลปุตฺตํ ปสฺเสยฺยา”ติ. เตน หิ คหปติ นิสีท, อปฺเปว นาม อิธ นิสินฺโน อิธ นิสินฺนํ ยสํ กุลปุตฺตํ ปสฺเสยฺยาสีติ. อถ โข เสฏฺฐิ คหปติ “อิเธว กิราหํ นิสินฺโน อิธ นิสินฺนํ ยสํ กุลปุตฺตํ ปสฺสิสฺสามี”ติ หฏฺโฐ อุทคฺโค ภควนฺตํ อภิวาเทตฺวา เอกมนฺตํ นิสีทิ, เอกมนฺตํ นิสินฺนสฺส โข เสฏฺฐิสฺส คหปติสฺส ภควา อนุปุพฺพิํ กถํ กเถสิ, เสยฺยถิทํ, ทานกถํ สีลกถํ สคฺคกถํ กามานํ อาทีนวํ โอการํ สํกิเลสํ เนกฺขมฺเม \csromanfolio{23}อานิสํสํ ปกาเสสิ. ยทา ภควา อญฺญาสิ เสฏฺฐิํ คหปติํ กลฺลจิตฺตํ มุทุจิตฺตํ วินีวรณ\-จิตฺตํ อุทคฺคจิตฺตํ ปสนฺนจิตฺตํ, อถ ยา พุทฺธานํ สามุกฺกํสิกา ธมฺมเทสนา, ตํ ปกาเสสิ ทุกฺขํ สมุทยํ นิโรธํ มคฺคํ, เสยฺยถาปิ นาม สุทฺธํ วตฺถํ อปคตกาฬกํ สมฺมเทว รชนํ ปฏิคฺคเณฺหยฺย, เอวเมว เสฏฺฐิสฺส คหปติสฺส ตสฺมิํเยว อาสเน วิรชํ วีตมลํ ธมฺมจกฺขุํ อุทปาทิ “ยํ กิญฺจิ สมุทย\-ธมฺมํ, สพฺพํ ตํ นิโรธธมฺมนฺ”ติ. อถ โข เสฏฺฐิ คหปติ ทิฏฺฐธมฺโม ปตฺตธมฺโม วิทิตธมฺโม ปริโยคาฬฺห\-ธมฺโม ติณฺณ\-วิจิกิจฺโฉ วิคต\-กถํกโถ เวสารชฺช\-ปฺปตฺโต อปรปฺปจฺจโย สตฺถุสาสเน ภควนฺตํ เอตทโวจ “อภิกฺกนฺตํ ภนฺเต, อภิกฺกนฺตํ ภนฺเต, เสยฺยถาปิ ภนฺเต นิกฺกุชฺชิตํ\csromansharedfootnote{c55eb59ddc1dec73}{นิกุชฺชิตํ (ก)} วา อุกฺกุชฺเชยฺย, ปฏิจฺฉนฺนํ วา วิวเรยฺย, มูฬฺหสฺส วา มคฺคํ อาจิกฺเขยฺย, อนฺธกาเร วา เตลปชฺโชตํ ธาเรยฺย ‘จกฺขุมนฺโต รูปานิ ทกฺขนฺตี’ติ, เอวเมวํ ภควตา อเนก\-ปริยาเยน ธมฺโม ปกาสิโต. เอสาหํ ภนฺเต ภควนฺตํ สรณํ คจฺฉามิ ธมฺมญฺจ ภิกฺขุสํฆญฺจ, อุปาสกํ มํ ภควา ธาเรตุ อชฺชตคฺเค ปาณุเปตํ สรณํ คตนฺ”ติ. โสว โลเก ปฐมํ อุปาสโก อโหสิ เตวาจิโก.}

\proseitemwithcloserruled{28}{อถ โข ยสสฺส กุลปุตฺตสฺส ปิตุโน ธมฺเม เทสิยมาเน ยถาทิฏฺฐํ ยถาวิทิตํ ภูมิํ ปจฺจเวกฺขนฺตสฺส อนุปาทาย อาสเวหิ จิตฺตํ วิมุจฺจิ. อถ โข ภควโต เอตทโหสิ “ยสสฺส โข กุลปุตฺตสฺส ปิตุโน ธมฺเม เทสิยมาเน ยถาทิฏฺฐํ ยถาวิทิตํ ภูมิํ ปจฺจเวกฺขนฺตสฺส อนุปาทาย อาสเวหิ จิตฺตํ วิมุตฺตํ, อภพฺโพ โข ยโส กุลปุตฺโต หีนายาวตฺติตฺวา กาเม ปริภุญฺชิตุํ เสยฺยถาปิ ปุพฺเพ อคาริกภูโต, ยนฺนูนาหํ ตํ อิทฺธา\-ภิสงฺขารํ ปฏิปฺปสฺสมฺเภยฺยนฺ”ติ. อถ โข ภควา ตํ อิทฺธา\-ภิสงฺขารํ ปฏิปฺปสฺสมฺเภสิ. อทฺทสา โข เสฏฺฐิ คหปติ ยสํ กุลปุตฺตํ นิสินฺนํ, ทิสฺวาน ยสํ กุลปุตฺตํ เอตทโวจ “มาตา เต ตาต ยส ปริเทว\csromansharedfootnote{dbed57fe27ce5fce}{ปริเทวิ (ก)} โสกสมาปนฺนา, เทหิ มาตุยา ชีวิตนฺ”ติ. อถ โข ยโส กุลปุตฺโต ภควนฺตํ อุลฺโลเกสิ. อถ โข ภควา เสฏฺฐิํ คหปติํ เอตทโวจ “ตํ กิํ มญฺญสิ คหปติ, ยสฺส เสกฺเขน ญาเณน เสกฺเขน ทสฺสเนน ธมฺโม ทิฏฺโฐ วิทิโต เสยฺยถาปิ ตยา, ตสฺส ยถาทิฏฺฐํ ยถาวิทิตํ ภูมิํ ปจฺจเวกฺขนฺตสฺส \csromanfolio{24}อนุปาทาย อาสเวหิ จิตฺตํ วิมุตฺตํ, ภพฺโพ นุ โข โส คหปติ หีนายาวตฺติตฺวา กาเม ปริภุญฺชิตุํ เสยฺยถาปิ ปุพฺเพ อคาริกภูโต”ติ. โน เหตํ ภนฺเต. ยสสฺส โข คหปติ กุลปุตฺตสฺส เสกฺเขน ญาเณน เสกฺเขน ทสฺสเนน ธมฺโม ทิฏฺโฐ วิทิโต เสยฺยถาปิ ตยา, ตสฺส ยถาทิฏฺฐํ ยถาวิทิตํ ภูมิํ ปจฺจเวกฺขนฺตสฺส อนุปาทาย อาสเวหิ จิตฺตํ วิมุตฺตํ, อภพฺโพ โข คหปติ ยโส กุลปุตฺโต หีนายาวตฺติตฺวา กาเม ปริภุญฺชิตุํ เสยฺยถาปิ ปุพฺเพ อคาริกภูโตติ. ลาภา ภนฺเต ยสสฺส กุลปุตฺตสฺส, สุลทฺธํ ภนฺเต ยสสฺส กุลปุตฺตสฺส, ยถา ยสสฺส กุลปุตฺตสฺส อนุปาทาย อาสเวหิ จิตฺตํ วิมุตฺตํ, อธิวาเสตุ เม ภนฺเต ภควา อชฺชตนาย ภตฺตํ ยเสน กุลปุตฺเตน ปจฺฉาสมเณนาติ. อธิวาเสสิ ภควา ตุณฺหีภาเวน. อถ โข เสฏฺฐิ คหปติ ภควโต อธิวาสนํ วิทิตฺวา อุฏฺฐายาสนา ภควนฺตํ อภิวาเทตฺวา ปทกฺขิณํ กตฺวา ปกฺกามิ. อถ โข ยโส กุลปุตฺโต อจิรปกฺกนฺเต เสฏฺฐิมฺหิ คหปติมฺหิ ภควนฺตํ เอตทโวจ “ลเภยฺยาหํ ภนฺเต ภควโต สนฺติเก ปพฺพชฺชํ, ลเภยฺยํ อุปสมฺปทนฺ”ติ. “เอหิ ภิกฺขู”ติ ภควา อโวจ “สฺวากฺขาโต ธมฺโม, จร พฺรหฺมจริยํ สมฺมา ทุกฺขสฺส อนฺตกิริยายา”ติ. สาว ตสฺส อายสฺมโต อุปสมฺปทา อโหสิ. เตน โข ปน สมเยน สตฺต โลเก อรหนฺโต โหนฺติ.}{ยสสฺส ปพฺพชฺชา นิฏฺฐิตา.}
\proseitem{29}{อถ โข ภควา ปุพฺพณฺหสมยํ นิวาเสตฺวา ปตฺต\-จีวรมา\-ทาย อายสฺมตา ยเสน ปจฺฉาสมเณน เยน เสฏฺฐิสฺส คหปติสฺส นิเวสนํ เตนุปสงฺกมิ, อุปสงฺกมิตฺวา ปญฺญตฺเต อาสเน นิสีทิ. อถ โข อายสฺมโต ยสสฺส มาตา จ ปุราณทุติยิกา จ เยน ภควา เตนุ\-ปสงฺกมิํสุ, อุปสงฺกมิตฺวา ภควนฺตํ อภิวาเทตฺวา เอกมนฺตํ นิสีทิํสุ. ตาสํ ภควา อนุปุพฺพิํ กถํ กเถสิ, เสยฺยถิทํ, ทานกถํ สีลกถํ สคฺคกถํ กามานํ อาทีนวํ โอการํ สํกิเลสํ เนกฺขมฺเม อานิสํสํ ปกาเสสิ. ยทา ตา ภควา อญฺญาสิ กลฺลจิตฺตา มุทุจิตฺตา วินีวรณ\-จิตฺตา อุทคฺคจิตฺตา ปสนฺนจิตฺตา, อถ ยา พุทฺธานํ สามุกฺกํสิกา ธมฺมเทสนา, ตํ ปกาเสสิ ทุกฺขํ สมุทยํ นิโรธํ \csromanfolio{25}มคฺคํ, เสยฺยถาปิ นาม สุทฺธํ วตฺถํ อปคตกาฬกํ สมฺมเทว รชนํ ปฏิคฺคเณฺหยฺย, เอวเมว ตาสํ ตสฺมิํเยว อาสเน วิรชํ วีตมลํ ธมฺมจกฺขุํ อุทปาทิ “ยํ กิญฺจิ สมุทย\-ธมฺมํ, สพฺพํ ตํ นิโรธธมฺมนฺ”ติ. ตา ทิฏฺฐธมฺมา ปตฺตธมฺมา วิทิตธมฺมา ปริโยคาฬฺห\-ธมฺมา ติณฺณ\-วิจิกิจฺฉา วิคต\-กถํกถา เวสารชฺช\-ปฺปตฺตา อปรปฺปจฺจยา สตฺถุสาสเน ภควนฺตํ เอตทโวจุํ “อภิกฺกนฺตํ ภนฺเต, อภิกฺกนฺตํ ภนฺเต ฯเปฯ. เอตา มยํ ภนฺเต ภควนฺตํ สรณํ คจฺฉาม ธมฺมญฺจ ภิกฺขุสํฆญฺจ, อุปาสิกาโย โน ภควา ธาเรตุ อชฺชตคฺเค ปาณุเปตา สรณํ คตา”ติ. ตา จ โลเก ปฐมํ อุปาสิกา อเหสุํ เตวาจิกา.}

\prose{อถ โข อายสฺมโต ยสสฺส มาตา จ ปิตา จ ปุราณทุติยิกา จ ภควนฺตญฺจ อายสฺมนฺตญฺจ ยสํ ปณีเตน ขาทนีเยน โภชนีเยน สหตฺถา สนฺตปฺเปตฺวา สมฺปวาเรตฺวา ภควนฺตํ ภุตฺตาวิํ โอนีต\-ปตฺต\-ปาณิํ เอกมนฺตํ นิสีทิํสุ. อถ โข ภควา อายสฺมโต ยสสฺส มาตรญฺจ ปิตรญฺจ ปุราณทุติยิกญฺจ ธมฺมิยา กถาย สนฺทสฺเสตฺวา สมาทเปตฺวา สมุตฺเตเชตฺวา สมฺปหํเสตฺวา อุฏฺฐายาสนา ปกฺกามิ.}

\proseitemwithcloserruled{30}{อสฺโสสุํ โข อายสฺมโต ยสสฺส จตฺตาโร คิหิสหายกา พาราณสิยํ เสฏฺฐา\-นุเสฏฺฐีนํ กุลานํ ปุตฺตา วิมโล สุพาหุ ปุณฺณชิ ควมฺปติ “ยโส กิร กุลปุตฺโต เกสมสฺสุํ โอหาเรตฺวา กาสายานิ วตฺถานิ อจฺฉาเทตฺวา อคารสฺมา อนคาริยํ ปพฺพชิโต”ติ, สุตฺวาน เนสํ เอตทโหสิ “น หิ นูน โส โอรโก ธมฺมวินโย, น สา โอรกา ปพฺพชฺชา, ยตฺถ ยโส กุลปุตฺโต เกสมสฺสุํ โอหาเรตฺวา กาสายานิ วตฺถานิ อจฺฉาเทตฺวา อคารสฺมา อนคาริยํ ปพฺพชิโต”ติ. เต\csromansharedfootnote{9ef6a7f17dc0c601}{เต จตฺตาโร ชนา (ก)} เยนายสฺมา ยโส เตนุ\-ปสงฺกมิํสุ, อุปสงฺกมิตฺวา อายสฺมนฺตํ ยสํ อภิวาเทตฺวา เอกมนฺตํ อฏฺฐํสุ. อถ โข อายสฺมา ยโส เต จตฺตาโร คิหิสหายเก อาทาย เยน ภควา เตนุปสงฺกมิ, อุปสงฺกมิตฺวา ภควนฺตํ อภิวาเทตฺวา เอกมนฺตํ นิสีทิ, เอกมนฺตํ นิสินฺโน โข อายสฺมา ยโส ภควนฺตํ เอตทโวจ “อิเม เม ภนฺเต จตฺตาโร คิหิสหายกา พาราณสิยํ เสฏฺฐา\-นุเสฏฺฐีนํ กุลานํ ปุตฺตา วิมโล สุพาหุ ปุณฺณชิ ควมฺปติ, อิเม\csromansharedfootnote{3544c830a5a0ceea}{อิเม จตฺตาโร (ก)} ภควา โอวทตุ \csromanfolio{26}อนุสาสตู”ติ. เตสํ ภควา อนุปุพฺพิํ กถํ กเถสิ, เสยฺยถิทํ, ทานกถํ สีลกถํ สคฺคกถํ กามานํ อาทีนวํ โอการํ สํกิเลสํ เนกฺขมฺเม อานิสํสํ ปกาเสสิ, ยทา เต ภควา อญฺญาสิ กลฺลจิตฺเต มุทุจิตฺเต วินีวรณ\-จิตฺเต อุทคฺคจิตฺเต ปสนฺนจิตฺเต, อถ ยา พุทฺธานํ สามุกฺกํสิกา ธมฺมเทสนา, ตํ ปกาเสสิ ทุกฺขํ สมุทยํ นิโรธํ มคฺคํ, เสยฺยถาปิ นาม สุทฺธํ วตฺถํ อปคตกาฬกํ สมฺมเทว รชนํ ปฏิคฺคเณฺหยฺย, เอวเมว เตสํ ตสฺมิํเยว อาสเน วิรชํ วีตมลํ ธมฺมจกฺขุํ อุทปาทิ “ยํ กิญฺจิ สมุทย\-ธมฺมํ, สพฺพํ ตํ นิโรธธมฺมนฺ”ติ. เต ทิฏฺฐธมฺมา ปตฺตธมฺมา วิทิตธมฺมา ปริโยคาฬฺห\-ธมฺมา ติณฺณ\-วิจิกิจฺฉา วิคต\-กถํกถา เวสารชฺช\-ปฺปตฺตา อปรปฺปจฺจยา สตฺถุสาสเน ภควนฺตํ เอตทโวจุํ “ลเภยฺยาม มยํ ภนฺเต ภควโต สนฺติเก ปพฺพชฺชํ, ลเภยฺยาม อุปสมฺปทนฺ”ติ. “เอถ ภิกฺขโว”ติ ภควา อโวจ “สฺวากฺขาโต ธมฺโม, จรถ พฺรหฺมจริยํ สมฺมา ทุกฺขสฺส อนฺตกิริยายา”ติ. สาว เตสํ อายสฺมนฺตานํ อุปสมฺปทา อโหสิ. อถ โข ภควา เต ภิกฺขู ธมฺมิยา กถาย โอวทิ อนุสาสิ, เตสํ ภควตา ธมฺมิยา กถาย โอวทิยมานานํ อนุสาสิย\-มานานํ อนุปาทาย อาสเวหิ จิตฺตานิ วิมุจฺจิํสุ. เตน โข ปน สมเยน เอกาทส โลเก อรหนฺโต โหนฺติ.}{จตุคิหิ\-สหายก\-ปพฺพชฺชา นิฏฺฐิตา.}
\proseitemwithcloserruled{31}{อสฺโสสุํ โข อายสฺมโต ยสสฺส ปญฺญาสมตฺตา คิหิสหายกา ชานปทา ปุพฺพา\-นุปุพฺพ\-กานํ กุลานํ ปุตฺตา “ยโส กิร กุลปุตฺโต เกสมสฺสุํ โอหาเรตฺวา กาสายานิ วตฺถานิ อจฺฉาเทตฺวา อคารสฺมา อนคาริยํ ปพฺพชิโต”ติ. สุตฺวาน เนสํ เอตทโหสิ “น หิ นูน โส โอรโก ธมฺมวินโย, น สา โอรกา ปพฺพชฺชา, ยตฺถ ยโส กุลปุตฺโต เกสมสฺสุํ โอหาเรตฺวา กาสายานิ วตฺถานิ อจฺฉาเทตฺวา อคารสฺมา อนคาริยํ ปพฺพชิโต”ติ, เต เยนายสฺมา ยโส เตนุ\-ปสงฺกมิํสุ, อุปสงฺกมิตฺวา อายสฺมนฺตํ ยสํ อภิวาเทตฺวา เอกมนฺตํ อฏฺฐํสุ. อถ โข อายสฺมา ยโส เต ปญฺญาสมตฺเต คิหิสหายเก อาทาย เยน ภควา เตนุปสงฺกมิ, \csromanfolio{27}อุปสงฺกมิตฺวา ภควนฺตํ อภิวาเทตฺวา เอกมนฺตํ นิสีทิ, เอกมนฺตํ นิสินฺโน โข อายสฺมา ยโส ภควนฺตํ เอตทโวจ “อิเม เม ภนฺเต ปญฺญาสมตฺตา คิหิสหายกา ชานปทา ปุพฺพา\-นุปุพฺพ\-กานํ กุลานํ ปุตฺตา, อิเม ภควา โอวทตุ อนุสาสตู”ติ. เตสํ ภควา อนุปุพฺพิํ กถํ กเถสิ, เสยฺยถิทํ, ทานกถํ สีลกถํ สคฺคกถํ กามานํ อาทีนวํ โอการํ สํกิเลสํ เนกฺขมฺเม อานิสํสํ ปกาเสสิ. ยทา เต ภควา อญฺญาสิ กลฺลจิตฺเต มุทุจิตฺเต วินีวรณ\-จิตฺเต อุทคฺคจิตฺเต ปสนฺนจิตฺเต, อถ ยา พุทฺธานํ สามุกฺกํสิกา ธมฺมเทสนา, ตํ ปกาเสสิ ทุกฺขํ สมุทยํ นิโรธํ มคฺคํ, เสยฺยถาปิ นาม สุทฺธํ วตฺถํ อปคตกาฬกํ สมฺมเทว รชนํ ปฏิคฺคเณฺหยฺย, เอวเมว เตสํ ตสฺมิํเยว อาสเน วิรชํ วีตมลํ ธมฺมจกฺขุํ อุทปาทิ “ยํ กิญฺจิ สมุทย\-ธมฺมํ, สพฺพํ ตํ นิโรธธมฺมนฺ”ติ. เต ทิฏฺฐธมฺมา ปตฺตธมฺมา วิทิตธมฺมา ปริโยคาฬฺห\-ธมฺมา ติณฺณ\-วิจิกิจฺฉา วิคต\-กถํกถา เวสารชฺช\-ปฺปตฺตา อปรปฺปจฺจยา สตฺถุสาสเน ภควนฺตํ เอตทโวจุํ “ลเภยฺยาม มยํ ภนฺเต ภควโต สนฺติเก ปพฺพชฺชํ, ลเภยฺยาม อุปสมฺปทนฺ”ติ. “เอถ ภิกฺขโว”ติ ภควา อโวจ “สฺวากฺขาโต ธมฺโม, จรถ พฺรหฺมจริยํ สมฺมา ทุกฺขสฺส อนฺตกิริยายา”ติ. สาว เตสํ อายสฺมนฺตานํ อุปสมฺปทา อโหสิ. อถ โข ภควา เต ภิกฺขู ธมฺมิยา กถาย โอวทิ อนุสาสิ, เตสํ ภควตา ธมฺมิยา กถาย โอวทิยมานานํ อนุสาสิย\-มานานํ อนุปาทาย อาสเวหิ จิตฺตานิ วิมุจฺจิํสุ. เตน โข ปน สมเยน เอกสฏฺฐิ โลเก อรหนฺโต โหนฺติ.}{ปญฺญาส\-คิหิ\-สหายก\-ปพฺพชฺชา นิฏฺฐิตา.}
\csromanmatikaanchor{mtk.9}
\csromantocmarkref{section}{8. มารกถา}{mtk.9}
\bookmark[dest={mtk.9},level=1]{8. มารกถา}
\csromanheader{1}{8. \textbf{มารกถา}}
\proseitem{32}{อถ โข ภควา เต ภิกฺขู อามนฺเตสิ\csromansymbolfootnote{*}{สํ 1. 106-7 ปิฏฺเฐสุ มารสํยุตฺเตปิ.}“มุตฺตาหํ ภิกฺขเว สพฺพปาเสหิ เย ทิพฺพา เย จ มานุสา, ตุเมฺหปิ ภิกฺขเว มุตฺตา สพฺพปาเสหิ เย ทิพฺพา เย จ มานุสา. จรถ ภิกฺขเว จาริกํ พหุชนหิตาย พหุชน\-สุขาย โลกานุกมฺปาย อตฺถาย หิตาย สุขาย เทวมนุสฺสานํ, มา เอเกน เทฺว อคมิตฺถ, เทเสถ ภิกฺขเว ธมฺมํ อาทิกลฺยาณํ มชฺเฌกลฺยาณํ ปริโยสาน\-กลฺยาณํ สาตฺถํ สพฺยญฺชนํ เกวล\-ปริปุณฺณํ ปริสุทฺธํ พฺรหฺมจริยํ ปกาเสถ, สนฺติ สตฺตา \csromanfolio{28}อปฺปรชกฺข\-ชาติกา, อสฺสวนตา ธมฺมสฺส ปริหายนฺติ, ภวิสฺสนฺติ ธมฺมสฺส อญฺญาตาโร, อหมฺปิ ภิกฺขเว เยน อุรุเวลา เสนานิคโม เตนุปสงฺกมิสฺสามิ ธมฺมเทสนายา”ติ.}

\proseitem{33}{อถ โข มาโร ปาปิมา เยน ภควา เตนุปสงฺกมิ, อุปสงฺกมิตฺวา ภควนฺตํ คาถาย อชฺฌภาสิ\csromandash{}}

\setlength{\csromangathapagewidth}{0pt}
\setlength{\csromangathawakwidth}{0pt}
\csromangathameasuregroup{“พทฺโธสิ สพฺพปาเสหิ, \\ มหาพนฺธนพทฺโธสิ, \\ มุตฺตาหํ\textsuperscript{1} สพฺพปาเสหิ, \\ มหาพนฺธนมุตฺโตมฺหิ, \\ \textsuperscript{*}“อนฺตลิกฺขจโร ปาโส, \\ เตน ตํ พาธยิสฺสามิ, \\ \textsuperscript{*}รูปา สทฺทา รสา คนฺธา, \\ เอตฺถ เม วิคโต ฉนฺโท,}{\csromangathabat{“พทฺโธสิ สพฺพปาเสหิ,}{เย ทิพฺพา เย จ มานุสา.} \\ \csromangathabat{มหาพนฺธนพทฺโธสิ,}{น เม สมณ โมกฺขสี”ติ.} \\ \csromangathabat{มุตฺตาหํ\textsuperscript{1} สพฺพปาเสหิ,}{เย ทิพฺพา เย จ มานุสา.} \\ \csromangathabat{มหาพนฺธนมุตฺโตมฺหิ,}{นิหโต ตฺวมสิ อนฺตกาติ.} \\ \csromangathabat{\textsuperscript{*}“อนฺตลิกฺขจโร ปาโส,}{ยฺวายํ จรติ มานโส.} \\ \csromangathabat{เตน ตํ พาธยิสฺสามิ,}{น เม สมณ โมกฺขสี”ติ.} \\ \csromangathabat{\textsuperscript{*}รูปา สทฺทา รสา คนฺธา,}{โผฏฺฐพฺพา จ มโนรมา.} \\ \csromangathabat{เอตฺถ เม วิคโต ฉนฺโท,}{นิหโต ตฺวมสิ อนฺตกาติ.}}
\csromangathasetleft{“พทฺโธสิ สพฺพปาเสหิ, \\ มหาพนฺธนพทฺโธสิ, \\ มุตฺตาหํ\textsuperscript{1} สพฺพปาเสหิ, \\ มหาพนฺธนมุตฺโตมฺหิ, \\ \textsuperscript{*}“อนฺตลิกฺขจโร ปาโส, \\ เตน ตํ พาธยิสฺสามิ, \\ \textsuperscript{*}รูปา สทฺทา รสา คนฺธา, \\ เอตฺถ เม วิคโต ฉนฺโท,}
\csromangathagroup{\csromangathastanza{\csromangathabat{“พทฺโธสิ สพฺพปาเสหิ,}{เย ทิพฺพา เย จ มานุสา.} \\ \csromangathabat{มหาพนฺธนพทฺโธสิ,}{น เม สมณ โมกฺขสี”ติ.}}\csromangathastanza{\csromangathabat{มุตฺตาหํ\csromansharedfootnote{426b77b50687c65e}{มุตฺโตหํ (สี, สฺยา)} สพฺพปาเสหิ,}{เย ทิพฺพา เย จ มานุสา.} \\ \csromangathabat{มหาพนฺธนมุตฺโตมฺหิ,}{นิหโต ตฺวมสิ อนฺตกาติ.}}\csromangathastanza{\csromangathabat{\csromansymbolfootnote{*}{สํ 1. 113 ปิฏฺเฐ มารสํยุตฺเตปิ.}“อนฺตลิกฺขจโร ปาโส,}{ยฺวายํ จรติ มานโส.} \\ \csromangathabat{เตน ตํ พาธยิสฺสามิ,}{น เม สมณ โมกฺขสี”ติ.}}\csromangathastanza{\csromangathabat{\csromansymbolmark{*}รูปา สทฺทา รสา คนฺธา,}{โผฏฺฐพฺพา จ มโนรมา.} \\ \csromangathabat{เอตฺถ เม วิคโต ฉนฺโท,}{นิหโต ตฺวมสิ อนฺตกาติ.}}}

\prosewithcloserruled{อถ โข มาโร ปาปิมา “ชานาติ มํ ภควา, ชานาติ มํ สุคโต”ติ ทุกฺขี ทุมฺมโน ตตฺเถ\-ว\-นฺตรธายีติ.}{มารกถา นิฏฺฐิตา.}
\csromanmatikaanchor{mtk.10}
\csromantocmarkref{section}{9. ปพฺพชฺชูปสมฺปทากถา}{mtk.10}
\bookmark[dest={mtk.10},level=1]{9. ปพฺพชฺชูปสมฺปทากถา}
\csromanheader{1}{9. \textbf{ปพฺพชฺชู}\-\textbf{ปสมฺปทา}\-\textbf{กถา}}
\proseitem{34}{เตน โข ปน สมเยน ภิกฺขู นานาทิสา นานาชนปทา ปพฺพชฺชาเปกฺเข จ อุปสมฺปทา\-เปกฺเข จ อาเนนฺติ “ภควา เน ปพฺพาเชสฺสติ อุปสมฺปาเท\-สฺสตี”ติ, ตตฺถ ภิกฺขู เจว กิลมนฺติ ปพฺพชฺชาเปกฺขา จ อุปสมฺปทา\-เปกฺขา จ. อถ โข ภควโต รโหคตสฺส ปฏิสลฺลีนสฺส เอวํ เจตโส ปริวิตกฺโก อุทปาทิ “เอตรหิ โข ภิกฺขู นานาทิสา นานาชนปทา ปพฺพชฺชาเปกฺเข จ อุปสมฺปทา\-เปกฺเข จ อาเนนฺติ ‘ภควา เน \csromanfolio{29}ปพฺพาเชสฺสติ อุปสมฺปาเท\-สฺสตี’ติ, ตตฺถ ภิกฺขู เจว กิลมนฺติ ปพฺพชฺชาเปกฺขา จ อุปสมฺปทา\-เปกฺขา จ, ยนฺนูนาหํ ภิกฺขูนํ อนุชาเนยฺยํ ตุเมฺหว ทานิ ภิกฺขเว ตาสุ ตาสุ ทิสาสุ เตสุ เตสุ ชนปเทสุ ปพฺพาเชถ อุปสมฺปาเทถา”ติ. อถ โข ภควา สายนฺหสมยํ ปฏิสลฺลานา วุฏฺฐิโต เอตสฺมิํ นิทาเน เอตสฺมิํ ปกรเณ ธมฺมิํ กถํ กตฺวา ภิกฺขู อามนฺเตสิ “อิธ มยฺหํ ภิกฺขเว รโหคตสฺส ปฏิสลฺลีนสฺส เอวํ เจตโส ปริวิตกฺโก อุทปาทิ ‘เอตรหิ โข ภิกฺขู นานาทิสา นานาชนปทา ปพฺพชฺชาเปกฺเข จ อุปสมฺปทา\-เปกฺเข จ อาเนนฺติ ภควา เน ปพฺพาเชสฺสติ อุปสมฺปาเท\-สฺสตีติ, ตตฺถ ภิกฺขู เจว กิลมนฺติ ปพฺพชฺชาเปกฺขา จ อุปสมฺปทา\-เปกฺขา จ, ยนฺนูนาหํ ภิกฺขูนํ อนุชาเนยฺยํ ตุเมฺหว ทานิ ภิกฺขเว ตาสุ ตาสุ ทิสาสุ เตสุ เตสุ ชนปเทสุ ปพฺพาเชถ อุปสมฺปาเทถา’ติ, อนุชานามิ ภิกฺขเว ตุเมฺหว ทานิ ตาสุ ตาสุ ทิสาสุ เตสุ เตสุ ชนปเทสุ ปพฺพาเชถ อุปสมฺปาเทถ. เอวญฺจ ปน ภิกฺขเว ปพฺพาเชตพฺโพ อุปสมฺ\-ปาเท\-ตพฺโพ\csromandash{}}

\prosewithcloserruled{ปฐมํ เกสมสฺสุํ โอหาราเปตฺวา\csromansharedfootnote{ae9b3d7dc47df51e}{โอหาเรตฺวา (ก)} กาสายานิ วตฺถานิ อจฺฉาทาเปตฺวา เอกํสํ อุตฺตราสงฺคํ การาเปตฺวา ภิกฺขูนํ ปาเท วนฺทาเปตฺวา อุกฺกุฏิกํ นิสีทาเปตฺวา อญฺชลิํ ปคฺคณฺหาเปตฺวา ‘เอวํ วเทหี’ติ วตฺตพฺโพ ‘พุทฺธํ สรณํ คจฺฉามิ, ธมฺมํ สรณํ คจฺฉามิ, สํฆํ สรณํ คจฺฉามิ. ทุติยมฺปิ พุทฺธํ สรณํ คจฺฉามิ, ทุติยมฺปิ ธมฺมํ สรณํ คจฺฉามิ, ทุติยมฺปิ สํฆํ สรณํ คจฺฉามิ. ตติยมฺปิ พุทฺธํ สรณํ คจฺฉามิ, ตติยมฺปิ ธมฺมํ สรณํ คจฺฉามิ, ตติยมฺปิ สํฆํ สรณํ คจฺฉามี’ติ. อนุชานามิ ภิกฺขเว อิเมหิ ตีหิ สรณคมเนหิ ปพฺพชฺชํ อุปสมฺปทนฺ”ติ.}{ตีหิ สรณคมเนหิ อุปสมฺปทา\-กถา นิฏฺฐิตา.}
\csromanmatikaanchor{mtk.11}
\csromantocmarkref{section}{10. ทุติยมารกถา}{mtk.11}
\bookmark[dest={mtk.11},level=1]{10. ทุติยมารกถา}
\csromanheader{1}{10. \textbf{ทุติย}\-\textbf{มาร}\-\textbf{กถา}}
\proseitem{35}{อถ โข ภควา วสฺสํวุฏฺโฐ\csromansharedfootnote{6077527f0285e9f8}{วสฺสํวุตฺโถ (สี)} ภิกฺขู อามนฺเตสิ\csromansymbolfootnote{*}{สํ 1. 116 ปิฏฺเฐปิ.}“มยฺหํ โข ภิกฺขเว โยนิโส มนสิการา โยนิโส สมฺมปฺปธานา อนุตฺตรา วิมุตฺติ อนุปฺปตฺตา อนุตฺตรา วิมุตฺติ สจฺฉิกตา, ตุเมฺหปิ ภิกฺขเว โยนิโส \csromanfolio{30}มนสิการา โยนิโส สมฺมปฺปธานา อนุตฺตรํ วิมุตฺติํ อนุปาปุณาถ, อนุตฺตรํ วิมุตฺติํ สจฺฉิกโรถา”ติ. อถ โข มาโร ปาปิมา เยน ภควา เตนุปสงฺกมิ, อุปสงฺกมิตฺวา ภควนฺตํ คาถาย อชฺฌภาสิ\csromandash{}}

\setlength{\csromangathapagewidth}{0pt}
\setlength{\csromangathawakwidth}{0pt}
\csromangathameasuregroup{“พทฺโธสิ มารปาเสหิ, \\ มหาพนฺธนพทฺโธสิ\textsuperscript{1}, \\ มุตฺตาหํ มารปาเสหิ, \\ มหาพนฺธนมุตฺโตมฺหิ\textsuperscript{1},}{\csromangathabat{“พทฺโธสิ มารปาเสหิ,}{เย ทิพฺพา เย จ มานุสา.} \\ \csromangathabat{มหาพนฺธนพทฺโธสิ\textsuperscript{1},}{น เม สมเณ โมกฺขสี”ติ.} \\ \csromangathabat{มุตฺตาหํ มารปาเสหิ,}{เย ทิพฺพา เย จ มานุสา.} \\ \csromangathabat{มหาพนฺธนมุตฺโตมฺหิ\textsuperscript{1},}{นิหโต ตฺวมสิ อนฺตกาติ.}}
\csromangathasetleft{“พทฺโธสิ มารปาเสหิ, \\ มหาพนฺธนพทฺโธสิ\textsuperscript{1}, \\ มุตฺตาหํ มารปาเสหิ, \\ มหาพนฺธนมุตฺโตมฺหิ\textsuperscript{1},}
\csromangathagroup{\csromangathastanza{\csromangathabat{“พทฺโธสิ มารปาเสหิ,}{เย ทิพฺพา เย จ มานุสา.} \\ \csromangathabat{มหาพนฺธนพทฺโธสิ\csromansharedfootnote{95d848d8c7d7cc9f}{มารพนฺธนพทฺโธสิ (สี, สฺยา)},}{น เม สมเณ โมกฺขสี”ติ.}}\csromangathastanza{\csromangathabat{มุตฺตาหํ มารปาเสหิ,}{เย ทิพฺพา เย จ มานุสา.} \\ \csromangathabat{มหาพนฺธนมุตฺโตมฺหิ\csromansharedfootnote{f718d072d8caab1a}{มารพนฺธนมุตฺโตมฺหิ (สี, สฺยา)},}{นิหโต ตฺวมสิ อนฺตกาติ.}}}

\prosewithcloserruled{อถ โข มาโร ปาปิมา “ชานาติ มํ ภควา, ชานาติ มํ สุคโต”ติ ทุกฺขี ทุมฺมโน ตตฺเถ\-ว\-นฺตรธายิ.}{ทุติย\-มาร\-กถา นิฏฺฐิตา.}
\csromanmatikaanchor{mtk.12}
\csromantocmarkref{section}{11. ภทฺทวคฺคิยวตฺถุ}{mtk.12}
\bookmark[dest={mtk.12},level=1]{11. ภทฺทวคฺคิยวตฺถุ}
\csromanheader{1}{11. \textbf{ภทฺท}\-\textbf{วคฺคิย}\-\textbf{วตฺถุ}}
\proseitemwithcloser{36}{อถ โข ภควา พาราณสิยํ ยถาภิรนฺตํ วิหริตฺวา เยน อุรุเวลา เตน จาริกํ ปกฺกามิ. อถ โข ภควา มคฺคา โอกฺกมฺม เยน อญฺญตโร วนสณฺโฑ เตนุปสงฺกมิ, อุปสงฺกมิตฺวา ตํ วนสณฺฑํ อชฺโฌคาเหตฺวา อญฺญตรสฺมิํ รุกฺขมูเล นิสีทิ. เตน โข ปน สมเยน ติํสมตฺตา ภทฺทวคฺคิยา สหายกา สปชาปติกา ตสฺมิํ วนสณฺเฑ ปริจาเรนฺติ, เอกสฺส ปชาปติ นาโหสิ, ตสฺส อตฺถาย เวสี อานีตา อโหสิ. อถ โข สา เวสี เตสุ ปมตฺเตสุ ปริจาเรนฺเตสุ ภณฺฑํ อาทาย ปลายิตฺถ, อถ โข เต สหายกา สหายกสฺส เวยฺยาวจฺจํ กโรนฺตา ตํ อิตฺถิํ คเวสนฺตา ตํ วนสณฺฑํ อาหิณฺฑนฺตา อทฺทสํสุ ภควนฺตํ อญฺญตรสฺมิํ รุกฺขมูเล นิสินฺนํ, ทิสฺวาน เยน ภควา เตนุ\-ปสงฺกมิํสุ, อุปสงฺกมิตฺวา ภควนฺตํ เอตทโวจุํ “อปิ ภนฺเต ภควา เอกํ อิตฺถิํ ปสฺเสยฺยา”ติ. กิํ ปน โว กุมารา อิตฺถิยาติ. อิธ มยํ ภนฺเต ติํสมตฺตา ภทฺทวคฺคิยา สหายกา สปชาปติกา อิมสฺมิํ วนสณฺเฑ ปริจาริมฺหา, เอกสฺส ปชาปติ นาโหสิ, ตสฺส อตฺถาย เวสี อานีตา อโหสิ, อถ โข สา \csromanfolio{31}ภนฺเต เวสี อเมฺหสุ ปมตฺเตสุ ปริจาเรนฺเตสุ ภณฺฑํ อาทาย ปลายิตฺถ, เต มยํ ภนฺเต สหายกา สหายกสฺส เวยฺยาวจฺจํ กโรนฺตา ตํ อิตฺถิํ คเวสนฺตา อิมํ วนสณฺฑํ อาหิณฺฑามาติ. ตํ กิํ มญฺญถ โว กุมารา, กตมํ นุ โข ตุมฺหากํ วรํ, ยํ วา ตุเมฺห อิตฺถิํ คเวเสยฺยาถ, ยํ วา อตฺตานํ คเวเสยฺยาถาติ. เอตเทว ภนฺเต อมฺหากํ วรํ, ยํ มยํ อตฺตานํ คเวเสยฺยามาติ. เตน หิ โว กุมารา นิสีทถ ธมฺมํ โว เทเสสฺสามีติ, “เอวํ ภนฺเต”ติ โข เต ภทฺทวคฺคิยา สหายกา ภควนฺตํ อภิวาเทตฺวา เอกมนฺตํ นิสีทิํสุ. เตสํ ภควา อนุปุพฺพิํ กถํ กเถสิ, เสยฺยถิทํ, ทานกถํ สีลกถํ สคฺคกถํ กามานํ อาทีนวํ โอการํ สํกิเลสํ เนกฺขมฺเม อานิสํสํ ปกาเสสิ, ยทา เต ภควา อญฺญาสิ กลฺลจิตฺเต มุทุจิตฺเต วินีวรณ\-จิตฺเต อุทคฺคจิตฺเต ปสนฺนจิตฺเต, อถ ยา พุทฺธานํ สามุกฺกํสิกา ธมฺมเทสนา, ตํ ปกาเสสิ ทุกฺขํ สมุทยํ นิโรธํ มคฺคํ, เสยฺยถาปิ นาม สุทฺธํ วตฺถํ อปคตกาฬกํ สมฺมเทว รชนํ ปฏิคฺคเณฺหยฺย, เอวเมว เตสํ ตสฺมิํเยว อาสเน วิรชํ วีตมลํ ธมฺมจกฺขุํ อุทปาทิ “ยํ กิญฺจิ สมุทย\-ธมฺมํ, สพฺพํ ตํ นิโรธธมฺมนฺ”ติ. เต ทิฏฺฐธมฺมา ปตฺตธมฺมา วิทิตธมฺมา ปริโยคาฬฺห\-ธมฺมา ติณฺณ\-วิจิกิจฺฉา วิคต\-กถํกถา เวสารชฺช\-ปฺปตฺตา อปรปฺปจฺจยา สตฺถุสาสเน ภควนฺตํ เอตทโวจุํ “ลเภยฺยาม มยํ ภนฺเต ภควโต สนฺติเก ปพฺพชฺชํ, ลเภยฺยาม อุปสมฺปทนฺ”ติ. “เอถ ภิกฺขโว”ติ ภควา อโวจ “สฺวากฺขาโต ธมฺโม, จรถ พฺรหฺมจริยํ สมฺมา ทุกฺขสฺส อนฺตกิริยายา”ติ. สาว เตสํ อายสฺมนฺตานํ อุปสมฺปทา อโหสิ.}{ภทฺท\-วคฺคิย\-สหายกานํ วตฺถุ นิฏฺฐิตํ.}
\nitthitamruled{ทุติย\-ภาณวาโร.}
\csromanmatikaanchor{mtk.13}
\csromantocmarkref{section}{12. อุรุเวฬปาฏิหาริยกถา}{mtk.13}
\bookmark[dest={mtk.13},level=1]{12. อุรุเวฬปาฏิหาริยกถา}
\csromanheader{1}{12. \textbf{อุรุเวล}\-\textbf{ปาฏิหาริย}\-\textbf{กถา}}
\proseitem{37}{อถ โข ภควา อนุปุพฺเพน จาริกํ จรมาโน เยน อุรุเวลา ตทวสริ. เตน โข ปน สมเยน อุรุเวลายํ ตโย ชฏิลา ปฏิวสนฺติ อุรุเวลกสฺสโป นทีกสฺสโป คยากสฺสโปติ, เตสุ \csromanfolio{32}อุรุเวลกสฺสโป ชฏิโล ปญฺจนฺนํ ชฏิลสตานํ นายโก โหติ วินายโก อคฺโค ปมุโข ปาโมกฺโข, นทีกสฺสโป ชฏิโล ติณฺณํ ชฏิลสตานํ นายโก โหติ วินายโก อคฺโค ปมุโข ปาโมกฺโข, คยากสฺสโป ชฏิโล ทฺวินฺนํ ชฏิลสตานํ นายโก โหติ วินายโก อคฺโค ปมุโข ปาโมกฺโข. อถ โข ภควา เยน อุรุเวล\-กสฺสปสฺส ชฏิลสฺส อสฺสโม เตนุปสงฺกมิ, อุปสงฺกมิตฺวา อุรุเวล\-กสฺสปํ ชฏิลํ เอตทโวจ “สเจ เต กสฺสป อครุ, วเสยฺยาม เอกรตฺตํ อคฺยาคาเร”ติ. น โข เม มหาสมณ ครุ, จณฺเฑตฺถ นาคราชา อิทฺธิมา อาสิวิโส โฆรวิโส, โส ตํ มา วิเหเฐสีติ. ทุติยมฺปิ โข ภควา อุรุเวล\-กสฺสปํ ชฏิลํ เอตทโวจ “สเจ เต กสฺสป อครุ, วเสยฺยาม เอกรตฺตํ อคฺยาคาเร”ติ. น โข เม มหาสมณ ครุ, จณฺเฑตฺถ นาคราชา อิทฺธิมา อาสิวิโส โฆรวิโส, โส ตํ มา วิเหเฐสีติ. ตติยมฺปิ โข ภควา อุรุเวล\-กสฺสปํ ชฏิลํ เอตทโวจ “สเจ เต กสฺสป อครุ, วเสยฺยาม เอกรตฺตํ อคฺยาคาเร”ติ. น โข เม มหาสมณ ครุ, จณฺเฑตฺถ นาคราชา อิทฺธิมา อาสิวิโส โฆรวิโส, โส ตํ มา วิเหเฐสีติ. อปฺเปว มํ น วิเหเฐยฺย, อิงฺฆ ตฺวํ กสฺสป อนุชานาหิ อคฺยาคารนฺติ. วิหร มหาสมณ ยถาสุขนฺติ. อถ โข ภควา อคฺยาคารํ ปวิสิตฺวา ติณสนฺถารกํ ปญฺญเปตฺวา นิสีทิ ปลฺลงฺกํ อาภุชิตฺวา อุชุํ กายํ ปณิธาย ปริมุขํ สติํ อุปฏฺฐเปตฺวา.}

\proseitem{38}{อทฺทสา โข โส นาโค ภควนฺตํ ปวิฏฺฐํ, ทิสฺวาน ทุมฺมโน\csromansharedfootnote{42d33c83d94010f6}{ทุกฺขี ทุมฺมโน (สี, สฺยา)} ปธูปายิ\csromansharedfootnote{cacee89a8af2bb0c}{ปขูปาสิ (ก)}. อถ โข ภควโต เอตทโหสิ “ยนฺนูนาหํ อิมสฺส นาคสฺส อนุปหจฺจ ฉวิญฺจ จมฺมญฺจ มํสญฺจ นฺหารุญฺจ อฏฺฐิญฺจ อฏฺฐิมิญฺชญฺจ เตชสา เตชํ ปริยาทิเยยฺยนฺ”ติ. อถ โข ภควา ตถารูปํ อิทฺธา\-ภิสงฺขารํ อภิส\-งฺขริ\-ตฺวา ปธูปายิ. อถ โข โส นาโค มกฺขํ อสหมาโน ปชฺชลิ. ภควาปิ เตโชธาตุํ สมาปชฺชิตฺวา ปชฺชลิ. อุภินฺนํ สโชติภูตานํ อคฺยาคารํ อาทิตฺตํ วิย โหติ สมฺปชฺชลิตํ สโชติภูตํ. อถ โข เต ชฏิลา อคฺยาคารํ ปริวาเรตฺวา เอวมาหํสุ “อภิรูโป วต โภ มหาสมโณ นาเคน วิเหฐิยตี”ติ. อถ โข ภควา \csromanfolio{33}ตสฺสา รตฺติยา อจฺจเยน ตสฺส นาคสฺส อนุปหจฺจ ฉวิญฺจ จมฺมญฺจ มํสญฺจ นฺหารุญฺจ อฏฺฐิญฺจ อฏฺฐิมิญฺชญฺจ เตชสา เตชํ ปริยาทิยิตฺวา ปตฺเต ปกฺขิปิตฺวา อุรุเวล\-กสฺสปสฺส ชฏิลสฺส ทสฺเสสิ “อยํ เต กสฺสป นาโค ปริยาทินฺโน\csromansharedfootnote{b60b9afbd95c2529}{ปริยาทิณฺโณ (ก)} อสฺส เตชสา เตโช”ติ. อถ โข อุรุเวล\-กสฺสปสฺส ชฏิลสฺส เอตทโหสิ “มหิทฺธิโก โข มหาสมโณ มหานุภาโว, ยตฺร หิ นาม จณฺฑสฺส นาคราชสฺส อิทฺธิมโต อาสิวิสสฺส โฆรวิสสฺส เตชสา เตชํ ปริยาทิยิสฺสติ, นเตฺวว จ โข อรหา ยถา อหนฺ”ติ.}

\setlength{\csromangathapagewidth}{0pt}
\setlength{\csromangathawakwidth}{0pt}
\csromangathameasuregroup{เนรญฺชรายํ ภควา, \\ สเจ เต กสฺสป อครุ, \\ น โข เม มหาสมณ ครุ, \\ ทิสฺวา อิสิํ ปวิฏฺฐํ, \\ สุมนมนโส อธิมโน\textsuperscript{1}, \\ มกฺขญฺจ อสหมาโน, \\ เตโชธาตุสุกุสโล,}{\csromangathabat{เนรญฺชรายํ ภควา,}{อุรุเวลกสฺสปํ ชฏิลํ อโวจ.} \\ \csromangathabat{สเจ เต กสฺสป อครุ,}{วิหเรมุ อชฺชโณฺห อคฺคิสาลมฺหีติ\textsuperscript{1}.} \\ \csromangathabat{น โข เม มหาสมณ ครุ,}{ผาสุกาโมว ตํ นิวาเรมิ.} \\ จณฺเฑตฺถ นาคราชา. \\ อิทฺธิมา อาสิวิโส โฆรวิโส. \\ โส ตํ มา วิเหเฐสีติ. \\ อปฺเปว มํ น วิเหเฐยฺย, \\ อิงฺฆ ตฺวํ กสฺสป อนุชานาหิ อคฺยาคารนฺติ. \\ ทินฺนนฺติ นํ วิทิตฺวา, \\ อภีโต\textsuperscript{1} ปาวิสิ ภยมตีโต. \\ \csromangathabat{ทิสฺวา อิสิํ ปวิฏฺฐํ,}{อหินาโค ทุมฺมโน ปธูปายิ.} \\ \csromangathabat{สุมนมนโส อธิมโน\textsuperscript{1},}{มนุสฺสนาโคปิ ตตฺถ ปธูปายิ.} \\ \csromangathabat{มกฺขญฺจ อสหมาโน,}{อหินาโค ปาวโกว ปชฺชลิ.} \\ \csromangathabat{เตโชธาตุสุกุสโล,}{มนุสฺสนาโคปิ ตตฺถ ปชฺชลิ.} \\ อุภินฺนํ สโชติภูตานํ, \\ อคฺยาคารํ อาทิตฺตํ โหติ สมฺปชฺชลิตํ สโชติภูตํ. \\ อุทิจฺฉเร ชฏิลา, \\ อภิรูโป วต โภ มหาสมโณ. \\ นาเคน วิเหฐิยตีติ ภณนฺติ. \\ อถ ตสฺสา รตฺติยา\textsuperscript{1} อจฺจเยน, \\ หตา นาคสฺส อจฺจิโย โหนฺติ\textsuperscript{1}. \\ อิทฺธิมโต ปน ฐิตา\textsuperscript{1}, \\ อเนกวณฺณา อจฺจิโย โหนฺติ. \\ นีลา อถ โลหิติกา, \\ มญฺชิฏฺฐา ปีตกา ผลิกวณฺณาโย. \\ องฺคีรสสฺส กาเย, \\ อเนกวณฺณา อจฺจิโย โหนฺติ. \\ ปตฺตมฺหิ โอทหิตฺวา, \\ อหินาคํ พฺราหฺมณสฺส ทสฺเสสิ. \\ อยํ เต กสฺสป นาโค, \\ ปริยาทินฺโน อสฺส เตชสา เตโชติ.}
\csromangathameasurewak{โส ตํ มา วิเหเฐสีติ. \\ อปฺเปว มํ น วิเหเฐยฺย, \\ อิงฺฆ ตฺวํ กสฺสป อนุชานาหิ อคฺยาคารนฺติ. \\ ทินฺนนฺติ นํ วิทิตฺวา, \\ อภีโต\textsuperscript{1} ปาวิสิ ภยมตีโต. \\ อุภินฺนํ สโชติภูตานํ, \\ อคฺยาคารํ อาทิตฺตํ โหติ สมฺปชฺชลิตํ สโชติภูตํ. \\ อุทิจฺฉเร ชฏิลา, \\ อภิรูโป วต โภ มหาสมโณ. \\ นาเคน วิเหฐิยตีติ ภณนฺติ. \\ อถ ตสฺสา รตฺติยา\textsuperscript{1} อจฺจเยน, \\ หตา นาคสฺส อจฺจิโย โหนฺติ\textsuperscript{1}. \\ อิทฺธิมโต ปน ฐิตา\textsuperscript{1}, \\ อเนกวณฺณา อจฺจิโย โหนฺติ. \\ นีลา อถ โลหิติกา, \\ มญฺชิฏฺฐา ปีตกา ผลิกวณฺณาโย. \\ องฺคีรสสฺส กาเย, \\ อเนกวณฺณา อจฺจิโย โหนฺติ. \\ ปตฺตมฺหิ โอทหิตฺวา, \\ อหินาคํ พฺราหฺมณสฺส ทสฺเสสิ. \\ อยํ เต กสฺสป นาโค, \\ ปริยาทินฺโน อสฺส เตชสา เตโชติ.}
\csromangathasetleft{เนรญฺชรายํ ภควา, \\ สเจ เต กสฺสป อครุ, \\ น โข เม มหาสมณ ครุ, \\ ทิสฺวา อิสิํ ปวิฏฺฐํ, \\ สุมนมนโส อธิมโน\textsuperscript{1}, \\ มกฺขญฺจ อสหมาโน, \\ เตโชธาตุสุกุสโล,}
\csromangathagroup{\csromangathastanza{\csromangathaitembat{39}{เนรญฺชรายํ ภควา,}{อุรุเวล\-กสฺสปํ ชฏิลํ อโวจ.} \\ \csromangathacontbat{สเจ เต กสฺสป อครุ,}{วิหเรมุ อชฺชโณฺห อคฺคิสาลมฺหีติ\csromansharedfootnote{ac037f9d64c2b557}{อคฺคิสรณมฺหีติ (สี, สฺยา)}.}}\csromangathastanza{\csromangathabat{น โข เม มหาสมณ ครุ,}{ผาสุกาโมว ตํ นิวาเรมิ.} \\ จณฺเฑตฺถ นาคราชา. \\ อิทฺธิมา อาสิวิโส โฆรวิโส.}\csromangathastanza{\csromangathawak{โส ตํ มา วิเหเฐสีติ. \\ อปฺเปว มํ น วิเหเฐยฺย, \\ อิงฺฆ ตฺวํ กสฺสป อนุชานาหิ อคฺยาคารนฺติ. \\ ทินฺนนฺติ นํ วิทิตฺวา,}}\csromangathastanza{\csromangathawak{อภีโต\csromansharedfootnote{e05d505e3e470d74}{อสมฺภีโต (สี)} ปาวิสิ ภยมตีโต.}}\csromangathastanza{\csromangathabat{ทิสฺวา อิสิํ ปวิฏฺฐํ,}{อหินาโค ทุมฺมโน ปธูปายิ.} \\ \csromangathabat{สุมนมนโส อธิมโน\csromansharedfootnote{0bcd5540bbf2c989}{อวิมโน (กตฺถจิ), นวิมโน (สฺยา)},}{มนุสฺสนาโคปิ ตตฺถ ปธูปายิ.}}\csromangathastanza{\csromangathabat{มกฺขญฺจ อสหมาโน,}{อหินาโค ปาวโกว ปชฺชลิ.} \\ \csromangathabat{เตโชธาตุสุกุสโล,}{มนุสฺสนาโคปิ ตตฺถ ปชฺชลิ.}}\csromangathastanza{\csromangathawak{อุภินฺนํ สโชติภูตานํ, \\ อคฺยาคารํ อาทิตฺตํ โหติ สมฺปชฺชลิตํ สโชติภูตํ. \\ อุทิจฺฉเร ชฏิลา, \\ อภิรูโป วต โภ มหาสมโณ.}}\csromangathastanza{\csromangathawak{นาเคน วิเหฐิยตีติ ภณนฺติ. \\ อถ ตสฺสา รตฺติยา\csromansharedfootnote{bb0d214b1a0b45e2}{อถ รตฺติยา (สี, สฺยา)} อจฺจเยน, \\ หตา นาคสฺส อจฺจิโย โหนฺติ\csromansharedfootnote{02ca53d4746629ca}{อหินาคสฺส อจฺจิโย น โหนฺติ (สี, สฺยา)}. \\ อิทฺธิมโต ปน ฐิตา\csromansharedfootnote{a8a2bff8d8bec74b}{อิทฺธิมโต ปนุฏฺฐิตา (สี)},}}\csromangathastanza{\csromangathawak{\csromanfolio{34}อเนกวณฺณา อจฺจิโย โหนฺติ. \\ นีลา อถ โลหิติกา, \\ มญฺชิฏฺฐา ปีตกา ผลิกวณฺณาโย. \\ องฺคีรสสฺส กาเย,}}\csromangathastanza{\csromangathawak{อเนกวณฺณา อจฺจิโย โหนฺติ. \\ ปตฺตมฺหิ โอทหิตฺวา, \\ อหินาคํ พฺราหฺมณสฺส ทสฺเสสิ. \\ อยํ เต กสฺสป นาโค,}}\csromangathastanza{\csromangathawak{ปริยาทินฺโน อสฺส เตชสา เตโชติ.}}}

\prosewithcloserruled{อถ โข อุรุเวลกสฺสโป ชฏิโล ภควโต อิมินา อิทฺธิ\-ปาฏิหาริเยน อภิปฺปสนฺโน ภควนฺตํ เอตทโวจ “อิเธว มหาสมณ วิหร, อหํ เต\csromansharedfootnote{fb7243e263bd9073}{เต อุปฏฺฐามิ (อิติปิ)} ธุวภตฺเตนา”ติ.}{ปฐมํ ปาฏิหาริยํ.}
\proseitemwithcloserruled{40}{อถ โข ภควา อุรุเวล\-กสฺสปสฺส ชฏิลสฺส อสฺสมสฺส อวิทูเร อญฺญตรสฺมิํ วนสณฺเฑ วิหาสิ. อถ โข จตฺตาโร มหาราชาโน อภิกฺกนฺตาย รตฺติยา อภิกฺกนฺต\-วณฺณา เกวลกปฺปํ วนสณฺฑํ โอภาเสตฺวา เยน ภควา เตนุ\-ปสงฺกมิํสุ, อุปสงฺกมิตฺวา ภควนฺตํ อภิวาเทตฺวา จตุทฺทิสา อฏฺฐํสุ เสยฺยถาปิ มหนฺตา อคฺคิกฺขนฺธา. อถ โข อุรุเวลกสฺสโป ชฏิโล ตสฺสา รตฺติยา อจฺจเยน เยน ภควา เตนุปสงฺกมิ, อุปสงฺกมิตฺวา ภควนฺตํ เอตทโวจ “กาโล มหาสมณ นิฏฺฐิตํ ภตฺตํ, เก นุ โข เต มหาสมณ อภิกฺกนฺตาย รตฺติยา อภิกฺกนฺต\-วณฺณา เกวลกปฺปํ วนสณฺฑํ โอภาเสตฺวา เยน ตฺวํ \csromanfolio{35}เตนุ\-ปสงฺกมิํสุ, อุปสงฺกมิตฺวา ตํ อภิวาเทตฺวา จตุทฺทิสา อฏฺฐํสุ เสยฺยถาปิ มหนฺตา อคฺคิกฺขนฺธา”ติ. เอเต โข กสฺสป จตฺตาโร มหาราชาโน เยนาหํ เตนุ\-ปสงฺกมิํสุ ธมฺม\-สฺสวนายาติ. อถ โข อุรุเวล\-กสฺสปสฺส ชฏิลสฺส เอตทโหสิ “มหิทฺธิโก โข มหาสมโณ มหานุภาโว, ยตฺร หิ นาม จตฺตาโรปิ มหาราชาโน อุปสงฺกมิสฺสนฺติ ธมฺม\-สฺสวนาย, น เตฺวว จ โข อรหา ยถา อหนฺ”ติ. อถ โข ภควา อุรุเวล\-กสฺสปสฺส ชฏิลสฺส ภตฺตํ ภุญฺชิตฺวา ตสฺมิํเยว วนสณฺเฑ วิหาสิ.}{ทุติยํ ปาฏิหาริยํ.}
\proseitem{41}{อถ โข สกฺโก เทวานมินฺโท อภิกฺกนฺตาย รตฺติยา อภิกฺกนฺต\-วณฺโณ เกวลกปฺปํ วนสณฺฑํ โอภาเสตฺวา เยน ภควา เตนุปสงฺกมิ, อุปสงฺกมิตฺวา ภควนฺตํ อภิวาเทตฺวา เอกมนฺตํ อฏฺฐาสิ เสยฺยถาปิ มหา\-อคฺคิกฺขนฺโธ ปุริมาหิ วณฺณนิภาหิ อภิกฺกนฺตตโร จ ปณีตตโร จ. อถ โข อุรุเวลกสฺสโป ชฏิโล ตสฺสา รตฺติยา อจฺจเยน เยน ภควา เตนุปสงฺกมิ, อุปสงฺกมิตฺวา ภควนฺตํ เอตทโวจ “กาโล มหาสมณ นิฏฺฐิตํ ภตฺตํ, โก นุ โข โส มหาสมณ อภิกฺกนฺตาย รตฺติยา อภิกฺกนฺต\-วณฺโณ เกวลกปฺปํ วนสณฺฑํ โอภาเสตฺวา เยน ตฺวํ เตนุปสงฺกมิ, อุปสงฺกมิตฺวา ตํ อภิวาเทตฺวา เอกมนฺตํ อฏฺฐาสิ เสยฺยถาปิ มหา\-อคฺคิกฺขนฺโธ ปุริมาหิ วณฺณนิภาหิ อภิกฺกนฺตตโร จ ปณีตตโร จา”ติ. เอโส โข กสฺสป สกฺโก เทวานมินฺโท เยนาหํ เตนุปสงฺกมิ ธมฺม\-สฺสวนายาติ. อถ โข อุรุเวล\-กสฺสปสฺส ชฏิลสฺส เอตทโหสิ “มหิทฺธิโก โข มหาสมโณ มหานุภาโว, ยตฺร หิ นาม สกฺโกปิ เทวานมินฺโท อุปสงฺกมิสฺสติ ธมฺม\-สฺสวนาย, น เตฺวว จ โข อรหา ยถา อหนฺ”ติ. อถ โข ภควา อุรุเวล\-กสฺสปสฺส ชฏิลสฺส ภตฺตํ ภุญฺชิตฺวา ตสฺมิํเยว วนสณฺเฑ วิหาสิ.}

\vspace{\tipitakaparskip}
\csromancenter{ตติยํ ปาฏิหาริยํ.}
\proseitemwithcloserruled{42}{\csromanfolio{36}อถ โข พฺรหฺมา สหมฺปติ อภิกฺกนฺตาย รตฺติยา อภิกฺกนฺต\-วณฺโณ เกวลกปฺปํ วนสณฺฑํ โอภาเสตฺวา เยน ภควา เตนุปสงฺกมิ, อุปสงฺกมิตฺวา ภควนฺตํ อภิวาเทตฺวา เอกมนฺตํ อฏฺฐาสิ เสยฺยถาปิ มหา\-อคฺคิกฺขนฺโธ ปุริมาหิ วณฺณนิภาหิ อภิกฺกนฺตตโร จ ปณีตตโร จ. อถ โข อุรุเวลกสฺสโป ชฏิโล ตสฺสา รตฺติยา อจฺจเยน เยน ภควา เตนุปสงฺกมิ, อุปสงฺกมิตฺวา ภควนฺตํ เอตทโวจ “กาโล มหาสมณ นิฏฺฐิตํ ภตฺตํ, โก นุ โข โส มหาสมณ อภิกฺกนฺตาย รตฺติยา อภิกฺกนฺต\-วณฺโณ เกวลกปฺปํ วนสณฺฑํ โอภาเสตฺวา เยน ตฺวํ เตนุปสงฺกมิ, อุปสงฺกมิตฺวา ตํ อภิวาเทตฺวา เอกมนฺตํ อฏฺฐาสิ เสยฺยถาปิ มหา\-อคฺคิกฺขนฺโธ ปุริมาหิ วณฺณนิภาหิ อภิกฺกนฺตตโร จ ปณีตตโร จา”ติ. เอโส โข กสฺสป พฺรหฺมา สหมฺปติ เยนาหํ เตนุปสงฺกมิ ธมฺม\-สฺสวนายาติ. อถ โข อุรุเวล\-กสฺสปสฺส ชฏิลสฺส เอตทโหสิ “มหิทฺธิโก โข มหาสมโณ มหานุภาโว, ยตฺร หิ นาม พฺรหฺมาปิ สหมฺปติ อุปสงฺกมิสฺสติ ธมฺม\-สฺสวนาย, น เตฺวว จ โข อรหา ยถา อหนฺ”ติ. อถ โข ภควา อุรุเวล\-กสฺสปสฺส ชฏิลสฺส ภตฺตํ ภุญฺชิตฺวา ตสฺมิํเยว วนสณฺเฑ วิหาสิ.}{จตุตฺถํ ปาฏิหาริยํ.}
\proseitemwithcloserruled{43}{เตน โข ปน สมเยน อุรุเวล\-กสฺสปสฺส ชฏิลสฺส มหายญฺโญ ปจฺจุปฏฺฐิโต โหติ, เกวลกปฺปา จ องฺคมคธา ปหูตํ ขาทนียํ โภชนียํ อาทาย อภิกฺกมิตุ\-กามา โหนฺติ. อถ โข อุรุเวล\-กสฺสปสฺส ชฏิลสฺส เอตทโหสิ “เอตรหิ โข เม มหายญฺโญ ปจฺจุปฏฺฐิโต, เกวลกปฺปา จ องฺคมคธา ปหูตํ ขาทนียํ โภชนียํ อาทาย อภิกฺกมิสฺสนฺติ, สเจ มหาสมโณ มหาชนกาเย อิทฺธิ\-ปาฏิหาริยํ กริสฺสติ, มหาสมณสฺส ลาภสกฺกาโร อภิวฑฺฒิสฺสติ, มม ลาภสกฺกาโร ปริหายิสฺสติ, อโห นูน มหาสมโณ สฺวาตนาย นาคจฺเฉยฺยา”ติ. อถ โข ภควา อุรุเวล\-กสฺสปสฺส ชฏิลสฺส เจตสา เจโต\-ปริวิตกฺกม\-ญฺญาย อุตฺตรกุรุํ คนฺตฺวา ตโต ปิณฺฑปาตํ อาหริตฺวา อโนตตฺตทเห ปริภุญฺชิตฺวา ตตฺเถว ทิวาวิหารํ อกาสิ. อถ โข อุรุเวลกสฺสโป ชฏิโล \csromanfolio{37}ตสฺสา รตฺติยา อจฺจเยน เยน ภควา เตนุปสงฺกมิ, อุปสงฺกมิตฺวา ภควนฺตํ เอตทโวจ “กาโล มหาสมณ นิฏฺฐิตํ ภตฺตํ, กิํ นุ โข มหาสมณ หิยฺโย นาคมาสิ, อปิ จ มยํ ตํ สราม ‘กิํ นุ โข มหาสมโณ นาคจฺฉตี’ติ, ขาทนียสฺส จ โภชนียสฺส จ เต ปฏิวีโส\csromansharedfootnote{0965e11ddf274316}{ปฏิวิํโส (สี), ปฏิวิโส (สฺยา)} ฐปิโต”ติ. นนุ เต กสฺสป เอตทโหสิ “เอตรหิ โข เม มหายญฺโญ ปจฺจุปฏฺฐิโต, เกวลกปฺปา จ องฺคมคธา ปหูตํ ขาทนียํ โภชนียํ อาทาย อภิกฺกมิสฺสนฺติ, สเจ มหาสมโณ มหาชนกาเย อิทฺธิ\-ปาฏิหาริยํ กริสฺสติ, มหาสมณสฺส ลาภสกฺกาโร อภิวฑฺฒิสฺสติ, มม ลาภสกฺกาโร ปริหายิสฺสติ, อโห นูน มหาสมโณ สฺวาตนาย นาคจฺเฉยฺยา”ติ, โส โข อหํ กสฺสป ตว เจตสา เจโต\-ปริวิตกฺกม\-ญฺญาย อุตฺตรกุรุํ คนฺตฺวา ตโต ปิณฺฑปาตํ อาหริตฺวา อโนตตฺตทเห ปริภุญฺชิตฺวา ตตฺเถว ทิวาวิหารํ อกาสินฺติ. อถ โข อุรุเวล\-กสฺสปสฺส ชฏิลสฺส เอตทโหสิ “มหิทฺธิโก โข มหาสมโณ มหานุภาโว, ยตฺร หิ นาม เจตสาปิ จิตฺตํ ปชานิสฺสติ, น เตฺวว จ โข อรหา ยถา อหนฺ”ติ. อถ โข ภควา อุรุวลกสฺสปสฺส ชฏิลสฺส ภตฺตํ ภุญฺชิตฺวา ตสฺมิํเยว วนสณฺเฑ วิหาสิ.}{ปญฺจมํ ปาฏิหาริยํ.}
\proseitem{44}{เตน โข ปน สมเยน ภควโต ปํสุกูลํ อุปฺปนฺนํ โหติ. อถ โข ภควโต เอตทโหสิ “กตฺถ นุ โข อหํ ปํสุกูลํ โธเวยฺยนฺ”ติ. อถ โข สกฺโก เทวานมินฺโท ภควโต เจตสา เจโต\-ปริวิตกฺกม\-ญฺญาย ปาณินา โปกฺขรณิํ ขณิตฺวา ภควนฺตํ เอตทโวจ “อิธ ภนฺเต ภควา ปํสุกูลํ โธวตู”ติ. อถ โข ภควโต เอตทโหสิ “กิมฺหิ นุ โข อหํ ปํสุกูลํ ปริมทฺเทยฺยนฺ”ติ. อถ โข สกฺโก เทวานมินฺโท ภควโต เจตสา เจโต\-ปริวิตกฺกม\-ญฺญาย มหติํ สิลํ อุปนิกฺขิปิ “อิธ ภนฺเต ภควา ปํสุกูลํ ปริมทฺทตู”ติ. อถ โข ภควโต เอตทโหสิ “กิมฺหิ นุ โข อหํ\csromansharedfootnote{96f60582236d1c6f}{อหํ ปํสุกูลํ (ก)} อาลมฺพิตฺวา อุตฺตเรยฺยนฺ”ติ. อถ โข กกุเธ อธิวตฺถา เทวตา \csromanfolio{38}ภควโต เจตสา เจโต\-ปริวิตกฺกม\-ญฺญาย สาขํ โอนาเมสิ “อิธ ภนฺเต ภควา อาลมฺพิตฺวา อุตฺตรตู”ติ. อถ โข ภควโต เอตทโหสิ “กิมฺหิ นุ โข อหํ ปํสุกูลํ วิสฺสชฺเชยฺยนฺ”ติ. อถ โข สกฺโก เทวานมินฺโท ภควโต เจตสา เจโต\-ปริวิตกฺกม\-ญฺญาย มหติํ สิลํ อุปนิกฺขิปิ “อิธ ภนฺเต ภควา ปํสุกูลํ วิสฺสชฺเชตู”ติ. อถ โข อุรุเวลกสฺสโป ชฏิโล ตสฺสา รตฺติยา อจฺจเยน เยน ภควา เตนุปสงฺกมิ, อุปสงฺกมิตฺวา ภควนฺตํ เอตทโวจ “กาโล มหาสมณ นิฏฺฐิตํ ภตฺตํ, กิํ นุ โข มหาสมณ นายํ ปุพฺเพ อิธ โปกฺขรณี สายํ อิธ โปกฺขรณี, นยิมา สิลา ปุพฺเพ อุปนิกฺขิตฺตา, เกนิมา สิลา อุปนิกฺขิตฺตา, นยิมสฺส กกุธสฺส ปุพฺเพ สาขา โอนตา, สายํ สาขา โอนตา”ติ. อิธ เม กสฺสป ปํสุกูลํ อุปฺปนฺนํ อโหสิ, ตสฺส มยฺหํ กสฺสป เอตทโหสิ “กตฺถ นุ โข อหํ ปํสุกูลํ โธเวยฺยนฺ”ติ, อถ โข กสฺสป สกฺโก เทวานมินฺโท มม เจตสา เจโต\-ปริวิตกฺกม\-ญฺญาย ปาณินา โปกฺขรณิํ ขณิตฺวา มํ เอตทโวจ “อิธ ภนฺเต ภควา ปํสุกูลํ โธวตู”ติ, สา’ยํ กสฺสป อมนุสฺเสน ปาณินา ขณิตา โปกฺขรณี. ตสฺส มยฺหํ กสฺสป เอตทโหสิ “กิมฺหิ นุ โข อหํ ปํสุกูลํ ปริมทฺเทยฺยนฺ”ติ, อถ โข กสฺสป สกฺโก เทวานมินฺโท มม เจตสา เจโต\-ปริวิตกฺกม\-ญฺญาย มหติํ สิลํ อุปนิกฺขิปิ “อิธ ภนฺเต ภควา ปํสุกูลํ ปริมทฺทตู”ติ, สา’ยํ กสฺสป อมนุสฺเสน อุปนิกฺขิตฺตา สิลา. ตสฺส มยฺหํ กสฺสป เอตทโหสิ “กิมฺหิ นุ โข อหํ อาลมฺพิตฺวา อุตฺตเรยฺยนฺ”ติ, อถ โข กสฺสป กกุเธ อธิวตฺถา เทวตา มม เจตสา เจโต\-ปริวิตกฺกม\-ญฺญาย สาขํ โอนาเมสิ “อิธ ภนฺเต ภควา อาลมฺพิตฺวา อุตฺตรตู”ติ, สฺวายํ อาหรหตฺโถ กกุโธ. ตสฺส มยฺหํ กสฺสป เอตทโหสิ “กิมฺหิ นุ โข อหํ ปํสุกูลํ วิสฺสชฺเชยฺยนฺ”ติ, อถ โข กสฺสป สกฺโก เทวานมินฺโท มม เจตสา เจโต\-ปริวิตกฺกม\-ญฺญาย มหติํ สิลํ อุปนิกฺขิปิ “อิธ ภนฺเต ภควา ปํสุกูลํ วิสฺสชฺเชตู”ติ, สา’ยํ กสฺสป อมนุสฺเสน อุปนิกฺขิตฺตา สิลาติ. อถ โข อุรุเวล\-กสฺสปสฺส ชฏิลสฺส เอตทโหสิ “มหิทฺธิโก โข มหาสมโณ มหานุภาโว, ยตฺร หิ นาม สกฺโกปิ เทวานมินฺโท เวยฺยาวจฺจํ กริสฺสติ, น เตฺวว จ โข อรหา ยถา อหนฺ”ติ. อถ โข ภควา อุรุเวล\-กสฺสปสฺส ชฏิลสฺส ภตฺตํ ภุญฺชิตฺวา ตสฺมิํเยว วนสณฺเฑ วิหาสิ.}

\prose{\csromanfolio{39}อถ โข อุรุเวลกสฺสโป ชฏิโล ตสฺสา รตฺติยา อจฺจเยน เยน ภควา เตนุปสงฺกมิ, อุปสงฺกมิตฺวา ภควโต กาลํ อาโรเจสิ “กาโล มหาสมณ นิฏฺฐิตํ ภตฺตนฺ”ติ. “คจฺฉ ตฺวํ กสฺสป อายามหนฺ”ติ อุรุเวล\-กสฺสปํ ชฏิลํ อุยฺโยเชตฺวา ยาย ชมฺพุยา ชมฺพุทีโป ปญฺญายติ, ตโต ผลํ คเหตฺวา ปฐมตรํ อาคนฺตฺวา อคฺยาคาเร นิสีทิ. อทฺทสา โข อุรุเวลกสฺสโป ชฏิโล ภควนฺตํ อคฺยาคาเร นิสินฺนํ, ทิสฺวาน ภควนฺตํ เอตทโวจ “กตเมน ตฺวํ มหาสมณ มคฺเคน อาคโต, อหํ ตยา ปฐมตรํ ปกฺกนฺโต, โส ตฺวํ ปฐมตรํ อาคนฺตฺวา อคฺยาคาเร นิสินฺโน”ติ. อิธาหํ กสฺสป ตํ อุยฺโยเชตฺวา ยาย ชมฺพุยา ชมฺพุทีโป ปญฺญายติ, ตโต ผลํ คเหตฺวา ปฐมตรํ อาคนฺตฺวา อคฺยาคาเร นิสินฺโน, อิทํ โข กสฺสป ชมฺพุผลํ วณฺณสมฺปนฺนํ คนฺธ\-สมฺปนฺนํ รสสมฺปนฺนํ, สเจ อากงฺขสิ ปริภุญฺชาติ. อลํ มหาสมณ ตฺวํเยว ตํ อรหสิ, ตฺวํเยว ตํ\csromansharedfootnote{9e56975acf447aec}{ตฺวํเยเวตํ อาหรสิ, ตฺวํเยเวตํ (สี, สฺยา)} ปริภุญฺชาหีติ. อถ โข อุรุเวล\-กสฺสปสฺส ชฏิลสฺส เอตทโหสิ “มหิทฺธิโก โข มหาสมโณ มหานุภาโว, ยตฺร หิ นาม มํ ปฐมตรํ อุยฺโยเชตฺวา ยาย ชมฺพุยา ชมฺพุทีโป ปญฺญายติ, ตโต ผลํ คเหตฺวา ปฐมตรํ อาคนฺตฺวา อคฺยาคาเร นิสีทิสฺสติ, น เตฺวว จ โข อรหา ยถา อหนฺ”ติ. อถ โข ภควา อุรุเวล\-กสฺสปสฺส ชฏิลสฺส ภตฺตํ ภุญฺชิตฺวา ตสฺมิํเยว วนสณฺเฑ วิหาสิ.}

\proseitem{45}{อถ โข อุรุเวลกสฺสโป ชฏิโล ตสฺสา รตฺติยา อจฺจเยน เยน ภควา เตนุปสงฺกมิ, อุปสงฺกมิตฺวา ภควโต กาลํ อาโรเจสิ “กาโล มหาสมณ นิฏฺฐิตํ ภตฺตนฺ”ติ. “คจฺฉ ตฺวํ กสฺสป อายามหนฺ”ติ อุรุเวล\-กสฺสปํ ชฏิลํ อุยฺโยเชตฺวา ยาย ชมฺพุยา ชมฺพุทีโป ปญฺญายติ, ตสฺสา อวิทูเร อมฺโพ ฯเปฯ ตสฺสา อวิทูเร อามลกี ฯเปฯ ตสฺสา อวิทูเร หรีตกี ฯเปฯ ตาวติํสํ คนฺตฺวา ปาริจฺฉตฺตก\-ปุปฺผํ คเหตฺวา ปฐมตรํ อาคนฺตฺวา อคฺยาคาเร นิสีทิ. อทฺทสา โข อุรุเวลกสฺสโป ชฏิโล ภควนฺตํ อคฺยาคาเร นิสินฺนํ, ทิสฺวาน ภควนฺตํ เอตทโวจ “กตเมน ตฺวํ มหาสมณ มคฺเคน อาคโต, อหํ ตยา ปฐมตรํ ปกฺกนฺโต, โส ตฺวํ ปฐมตรํ อาคนฺตฺวา อคฺยาคาเร นิสินฺโน”ติ. \csromanfolio{40}อิธาหํ กสฺสป ตํ อุยฺโยเชตฺวา ตาวติํสํ คนฺตฺวา ปาริจฺฉตฺตก\-ปุปฺผํ คเหตฺวา ปฐมตรํ อาคนฺตฺวา อคฺยาคาเร นิสินฺโน, อิทํ โข กสฺสป ปาริจฺฉตฺตก\-ปุปฺผํ วณฺณสมฺปนฺนํ คนฺธ\-สมฺปนฺนํ\csromansharedfootnote{375b59000ea7a227}{สุคนฺธิกํ (ก)}, (สเจ อากงฺขสิ คณฺหาติ. อลํ มหาสมณ ตฺวํเยว ตํ อรหสิ, ตฺวํเยว ตํ คณฺหาติ.)\csromansharedfootnote{3fb4ac26f11a1729}{( ) สีสฺยาโปตฺถเกสุ นตฺถิ.} อถ โข อุรุเวล\-กสฺสปสฺส ชฏิลสฺส เอตทโหสิ “มหิทฺธิโก โข มหาสมโณ มหานุภาโว, ยตฺร หิ นาม มํ ปฐมตรํ อุยฺโยเชตฺวา ตาวติํสํ คนฺตฺวา ปาริจฺฉตฺตก\-ปุปฺผํ คเหตฺวา ปฐมตรํ อาคนฺตฺวา อคฺยาคาเร นิสีทิสฺสติ, น เตฺวว จ โข อรหา ยถา อหนฺ”ติ.}

\proseitem{46}{เตน โข ปน สมเยน เต ชฏิลา อคฺคิํ ปริจริตุกามา น สกฺโกนฺติ กฏฺฐานิ ผาเลตุํ. อถ โข เตสํ ชฏิลานํ เอตทโหสิ “นิสฺสํสยํ โข มหาสมณสฺส อิทฺธานุภาโว ยถา มยํ น สกฺโกม กฏฺฐานิ ผาเลตุนฺ”ติ. อถ โข ภควา อุรุเวล\-กสฺสปํ ชฏิลํ เอตทโวจ “ผาลิยนฺตุ กสฺสป กฏฺฐานี”ติ. ผาลิยนฺตุ มหาสมณาติ. สกิเทว ปญฺจ กฏฺฐสตานิ ผาลิยิํสุ. อถ โข อุรุเวล\-กสฺสปสฺส ชฏิลสฺส เอตทโหสิ “มหิทฺธิโก โข มหาสมโณ มหานุภาโว, ยตฺร หิ นาม กฏฺฐานิปิ ผาลิยิสฺสนฺติ, น เตฺวว จ โข อรหา ยถา อหนฺ”ติ.}

\proseitem{47}{เตน โข ปน สมเยน เต ชฏิลา อคฺคิํ ปริจริตุกามา น สกฺโกนฺติ อคฺคิํ อุชฺชเลตุํ\csromansharedfootnote{428126323d5499d1}{ชาเลตุํ (สี), อุชฺชลิตุํ (ก)}. อถ โข เตสํ ชฏิลานํ เอตทโหสิ “นิสฺสํสยํ โข มหาสมณสฺส อิทฺธานุภาโว ยถา มยํ น สกฺโกม อคฺคิํ อุชฺชเลตุนฺ”ติ. อถ โข ภควา อุรุเวล\-กสฺสปํ ชฏิลํ เอตทโวจ “อุชฺชลิยนฺตุ กสฺสป อคฺคี”ติ. อุชฺชลิยนฺตุ มหาสมณาติ. สกิเทว ปญฺจ อคฺคิสตานิ อุชฺชลิยิํสุ. อถ โข อุรุเวล\-กสฺสปสฺส ชฏิลสฺส เอตทโหสิ “มหิทฺธิโก โข มหาสมโณ มหานุภาโว, ยตฺร หิ นาม อคฺคีปิ อุชฺชลิยิสฺสนฺติ, น เตฺวว จ โข อรหา ยถา อหนฺ”ติ.}

\proseitem{48}{เตน โข ปน สมเยน เต ชฏิลา อคฺคิํ ปริจริตฺวา น สกฺโกนฺติ อคฺคิํ วิชฺฌาเปตุํ. อถ โข เตสํ ชฏิลานํ เอตทโหสิ “นิสฺสํสยํ โข มหาสมณสฺส อิทฺธานุภาโว ยถา มยํ น สกฺโกม อคฺคิํ วิชฺฌาเปตุนฺ”ติ. อถ โข ภควา อุรุเวล\-กสฺสปํ ชฏิลํ \csromanfolio{41}เอตทโวจ “วิชฺฌายนฺตุ กสฺสป อคฺคี”ติ. วิชฺฌายนฺตุ มหาสมณาติ. สกิเทว ปญฺจ อคฺคิสตานิ วิชฺฌายิํสุ. อถ โข อุรุเวล\-กสฺสปสฺส ชฏิลสฺส เอตทโหสิ “มหิทฺธิโก โข มหาสมโณ มหานุภาโว, ยตฺร หิ นาม อคฺคีปิ วิชฺฌายิสฺสนฺติ, น เตฺวว จ โข อรหา ยถา อหนฺ”ติ.}

\proseitem{49}{เตน โข ปน สมเยน เต ชฏิลา สีตาสุ เหมนฺติกาสุ รตฺตีสุ อนฺตรฏฺฐกาสุ หิมปาตสมเย นชฺชา เนรญฺชราย อุมฺมุชฺชนฺติปิ นิมุชฺชนฺติปิ อุมฺมุชฺชน\-นิมุชฺชนมฺปิ กโรนฺติ. อถ โข ภควา ปญฺจมตฺตานิ มนฺทามุขิ\-สตานิ อภินิมฺมินิ ยตฺถ เต ชฏิลา อุตฺตริตฺวา วิสิพฺเพสุํ. อถ โข เตสํ ชฏิลานํ เอตทโหสิ “นิสฺสํสยํ โข มหาสมณสฺส อิทฺธานุภาโว ยถยิมา มนฺทามุขิโย นิมฺมิตา”ติ. อถ โข อุรุเวล\-กสฺสปสฺส ชฏิลสฺส เอตทโหสิ “มหิทฺธิโก โข มหาสมโณ มหานุภาโว, ยตฺร หิ นาม ตาว พหู มนฺทามุขิโยปิ อภินิมฺมินิสฺสติ, น เตฺวว จ โข อรหา ยถา อหนฺ”ติ.}

\proseitem{50}{เตน โข ปน สมเยน มหา อกาลเมโฆ ปาวสฺสิ, มหา อุทกวาหโก สญฺชายิ, ยสฺมิํ ปเทเส ภควา วิหรติ โส ปเทโส อุทเกน น โอตฺถโฏ\csromansharedfootnote{d25e09052b123b5f}{อุทเกน โอตฺถโฏ (สี, สฺยา)} โหติ. อถ โข ภควโต เอตทโหสิ “ยนฺนูนาหํ สมนฺตา อุทกํ อุสฺสาเรตฺวา มชฺเฌ เรณุหตาย ภูมิยา จงฺกเมยฺยนฺ”ติ. อถ โข ภควา สมนฺตา อุทกํ อุสฺสาเรตฺวา มชฺเฌ เรณุหตาย ภูมิยา จงฺกมิ. อถ โข อุรุเวลกสฺสโป ชฏิโล “มาเหว โข มหาสมโณ อุทเกน วูโฬฺห อโหสี”ติ นาวาย สมฺพหุเลหิ ชฏิเลหิ สทฺธิํ ยสฺมิํ ปเทเส ภควา วิหรติ ตํ ปเทสํ อคมาสิ. อทฺทสา โข อุรุเวลกสฺสโป ชฏิโล ภควนฺตํ สมนฺตา อุทกํ อุสฺสาเรตฺวา มชฺเฌ เรณุหตาย ภูมิยา จงฺกมนฺตํ, ทิสฺวาน ภควนฺตํ เอตทโวจ “อิทํ นุ ตฺวํ มหาสมณา”ติ. “อยมหมสฺมิ\csromansharedfootnote{033f146f315ef5ac}{อาม อหมสฺมิ (สฺยา)} กสฺสปา”ติ ภควา เวหาสํ อพฺภุคฺคนฺตฺวา นาวาย ปจฺจุฏฺฐาสิ. อถ โข อุรุเวล\-กสฺสปสฺส ชฏิลสฺส เอตทโหสิ “มหิทฺธิโก โข มหาสมโณ มหานุภาโว, ยตฺร หิ นาม อุทกมฺปิ น ปวาหิสฺสติ\csromansharedfootnote{36ac9f30a3bf0207}{นปฺปสหิสฺสติ (สี)}, น เตฺวว จ โข อรหา ยถา อหนฺ”ติ.}

\proseitem{51}{\csromanfolio{42}อถ โข ภควโต เอตทโหสิ “จิรมฺปิ โข อิมสฺส โมฆปุริสสฺส เอวํ ภวิสฺสติ ‘มหิทฺธิโก โข มหาสมโณ มหานุภาโว, น เตฺวว จ โข อรหา ยถา อหนฺติ, ยนฺนูนาหํ อิมํ ชฏิลํ สํเวเชยฺยนฺ”ติ. อถ โข ภควา อุรุเวล\-กสฺสปํ ชฏิลํ เอตทโวจ ‘เนว จ โข ตฺวํ กสฺสป อรหา, นาปิ อรหตฺตมคฺค\-สมาปนฺโน, สาปิ เต ปฏิปทา นตฺถิ, ยาย ตฺวํ อรหา วา อสฺสสิ อรหตฺตมคฺคํ วา สมาปนฺโน”ติ. อถ โข อุรุเวลกสฺสโป ชฏิโล ภควโต ปาเทสุ สิรสา นิปติตฺวา ภควนฺตํ เอตทโวจ “ลเภยฺยาหํ ภนฺเต ภควโต สนฺติเก ปพฺพชฺชํ, ลเภยฺยํ อุปสมฺปทนฺ”ติ. ตฺวํ โขสิ กสฺสป ปญฺจนฺนํ ชฏิลสตานํ นายโก วินายโก อคฺโค ปมุโข ปาโมกฺโข, เตปิ ตาว อปโลเกหิ, ยถา เต มญฺญิสฺสนฺติ ตถา เต กริสฺสนฺตีติ. อถ โข อุรุเวลกสฺสโป ชฏิโล เยน เต ชฏิลา เตนุปสงฺกมิ, อุปสงฺกมิตฺวา เต ชฏิเล เอตทโวจ “อิจฺฉามหํ โภ มหาสมเณ พฺรหฺมจริยํ จริตุํ, ยถา ภวนฺโต มญฺญนฺติ ตถา กโรนฺตู”ติ. จิรปฏิกา มยํ โภ มหาสมเณ อภิปฺปสนฺนา, สเจ ภวํ มหาสมเณ พฺรหฺมจริยํ จริสฺสติ, สพฺเพว มยํ มหาสมเณ พฺรหฺมจริยํ จริสฺสามาติ. อถ โข เต ชฏิลา เกสมิสฺสํ ชฏามิสฺสํ ขาริกาช\-มิสฺสํ อคฺคิหุตมิสฺสํ อุทเก ปวาเหตฺวา เยน ภควา เตนุ\-ปสงฺกมิํสุ, อุปสงฺกมิตฺวา ภควโต ปาเทสุ สิรสา นิปติตฺวา ภควนฺตํ เอตทโวจุํ “ลเภยฺยาม มยํ ภนฺเต ภควโต สนฺติเก ปพฺพชฺชํ, ลเภยฺยาม อุปสมฺปทนฺ”ติ. “เอถ ภิกฺขโว”ติ ภควา อโวจ “สฺวากฺขาโต ธมฺโม, จรถ พฺรหฺมจริยํ สมฺมา ทุกฺขสฺส อนฺตกิริยายา”ติ. สาว เตสํ อายสฺมนฺตานํ อุปสมฺปทา อโหสิ.}

\proseitem{52}{อทฺทสา โข นทีกสฺสโป ชฏิโล เกสมิสฺสํ ชฏามิสฺสํ ขาริกาช\-มิสฺสํ อคฺคิหุตมิสฺสํ อุทเก วุยฺหมาเน, ทิสฺวานสฺส เอตทโหสิ “มาเหว เม ภาตุโน อุปสคฺโค อโหสี”ติ, ชฏิเล ปาเหสิ “คจฺฉถ เม ภาตรํ ชานาถา”ติ, สามญฺจ ตีหิ ชฏิลสเตหิ สทฺธิํ เยนายสฺมา อุรุเวลกสฺสโป เตนุปสงฺกมิ, อุปสงฺกมิตฺวา อายสฺมนฺตํ อุรุเวล\-กสฺสปํ เอตทโวจ “อิทํ นุ โข กสฺสป เสยฺโย”ติ. อามาวุโส อิทํ เสยฺโยติ. อถ โข เต ชฏิลา เกสมิสฺสํ ชฏามิสฺสํ ขาริกาช\-มิสฺสํ อคฺคิหุตมิสฺสํ อุทเก ปวาเหตฺวา เยน \csromanfolio{43}ภควา เตนุ\-ปสงฺกมิํสุ, อุปสงฺกมิตฺวา ภควโต ปาเทสุ สิรสา นิปติตฺวา ภควนฺตํ เอตทโวจุํ “ลเภยฺยาม มยํ ภนฺเต ภควโต สนฺติเก ปพฺพชฺชํ, ลเภยฺยาม อุปสมฺปทนฺ”ติ. “เอถ ภิกฺขโว”ติ ภควา อโวจ “สฺวากฺขาโต ธมฺโม, จรถ พฺรหฺมจริยํ สมฺมา ทุกฺขสฺส อนฺตกิริยายา”ติ. สาว เตสํ อายสฺมนฺตานํ อุปสมฺปทา อโหสิ.}

\proseitem{53}{อทฺทสา โข คยากสฺสโป ชฏิโล เกสมิสฺสํ ชฏามิสฺสํ ขาริกาช\-มิสฺสํ อคฺคิหุตมิสฺสํ อุทเก วุยฺหมาเน, ทิสฺวานสฺส เอตทโหสิ “มาเหว เม ภาตูนํ อุปสคฺโค อโหสี”ติ, ชฏิเล ปาเหสิ “คจฺฉถ เม ภาตโร ชานาถา”ติ, สามญฺจ ทฺวีหิ ชฏิลสเตหิ สทฺธิํ เยนายสฺมา อุรุเวลกสฺสโป เตนุปสงฺกมิ, อุปสงฺกมิตฺวา อายสฺมนฺตํ อุรุเวล\-กสฺสปํ เอตทโวจ “อิทํ นุ โข กสฺสป เสยฺโย”ติ. อามาวุโส อิทํ เสยฺโยติ. อถ โข เต ชฏิลา เกสมิสฺสํ ชฏามิสฺสํ ขาริกาช\-มิสฺสํ อคฺคิหุตมิสฺสํ อุทเก ปวาเหตฺวา เยน ภควา เตนุ\-ปสงฺกมิํสุ, อุปสงฺกมิตฺวา ภควโต ปาเทสุ สิรสา นิปติตฺวา ภควนฺตํ เอตทโวจุํ “ลเภยฺยาม มยํ ภนฺเต ภควโต สนฺติเก ปพฺพชฺชํ, ลเภยฺยาม อุปสมฺปทนฺ”ติ. “เอถ ภิกฺขโว”ติ ภควา อโวจ “สฺวากฺขาโต ธมฺโม, จรถ พฺรหฺมจริยํ สมฺมา ทุกฺขสฺส อนฺตกิริยายา”ติ. สาว เตสํ อายสฺมนฺตานํ อุปสมฺปทา อโหสิ. ภควโต อธิฏฺฐาเนน ปญฺจ กฏฺฐสตานิ น ผาลิยิํสุ, ผาลิยิํสุ, อคฺคี น อุชฺชลิยิํสุ, อุชฺชลิยิํสุ, น วิชฺฌายิํสุ, วิชฺฌายิํสุ, ปญฺจ\-มนฺทามุขิ\-สตานิ อภินิมฺมินิ, เอเตน นเยน อฑฺฒุฑฺฒ\-ปาฏิหาริย\-สหสฺสานิ โหนฺติ.}

\proseitem{54}{อถ โข ภควา อุรุเวลายํ ยถาภิรนฺตํ วิหริตฺวา เยน คยาสีสํ เตน ปกฺกามิ มหตา ภิกฺขุ\-สํเฆน สทฺธิํ ภิกฺขุ\-สหสฺเสน สพฺเพเหว ปุราณชฏิเลหิ. ตตฺร สุทํ ภควา คยายํ วิหรติ คยาสีเส สทฺธิํ ภิกฺขุ\-สหสฺเสน. ตตฺร โข ภควา ภิกฺขู อามนฺเตสิ\csromandash{}}

\prose{\csromanfolio{44}\csromansymbolfootnote{*}{สํ 2. 251, 252 ปิฏฺเฐสุปิ}“สพฺพํ ภิกฺขเว อาทิตฺตํ, กิญฺจ ภิกฺขเว สพฺพํ อาทิตฺตํ, จกฺขุ อาทิตฺตํ, รูปา อาทิตฺตา, จกฺขุวิญฺญาณํ อาทิตฺตํ, จกฺขุ\-สมฺผสฺโส อาทิตฺโต, ยมิทํ จกฺขุ\-สมฺผสฺส\-ปจฺจยา อุปฺปชฺชติ เวทยิตํ สุขํ วา ทุกฺขํ วา อทุกฺขมสุขํ วา ตมฺปิ อาทิตฺตํ.}

\prose{เกน อาทิตฺตํ, ราคคฺคินา โทสคฺคินา โมหคฺคินา อาทิตฺตํ, ชาติยา ชราย มรเณน โสเกหิ ปริเทเวหิ ทุกฺเขหิ โทมนสฺเสหิ อุปายาเสหิ อาทิตฺตนฺติ วทามิ.}

\prose{โสตํ อาทิตฺตํ, สทฺทา อาทิตฺตา, โสตวิญฺญาณํ อาทิตฺตํ, โสตสมฺผสฺโส อาทิตฺโต, ยมิทํ โสต\-สมฺผสฺส\-ปจฺจยา อุปฺปชฺชติ เวทยิตํ สุขํ วา ทุกฺขํ วา อทุกฺขมสุขํ วา ตมฺปิ อาทิตฺตํ.}

\prose{เกน อาทิตฺตํ, ราคคฺคินา โทสคฺคินา โมหคฺคินา อาทิตฺตํ, ชาติยา ชราย มรเณน โสเกหิ ปริเทเวหิ ทุกฺเขหิ โทมนสฺเสหิ อุปายาเสหิ อาทิตฺตนฺติ วทามิ.}

\prose{ฆานํ อาทิตฺตํ, คนฺธา อาทิตฺตา, ฆานวิญฺญาณํ อาทิตฺตํ, ฆานสมฺผสฺโส อาทิตฺโต, ยมิทํ ฆาน\-สมฺผสฺส\-ปจฺจยา อุปฺปชฺชติ เวทยิตํ สุขํ วา ทุกฺขํ วา อทุกฺขมสุขํ วา ตมฺปิ อาทิตฺตํ.}

\prose{เกน อาทิตฺตํ, ราคคฺคินา โทสคฺคินา โมหคฺคินา อาทิตฺตํ, ชาติยา ชราย มรเณน โสเกหิ ปริเทเวหิ ทุกฺเขหิ โทมนสฺเสหิ อุปายาเสหิ อาทิตฺตนฺติ วทามิ.}

\prose{ชิวฺหา อาทิตฺตา, รสา อาทิตฺตา, ชิวฺหาวิญฺญาณํ อาทิตฺตํ, ชิวฺหาสมฺผสฺโส อาทิตฺโต, ยมิทํ ชิวฺหา\-สมฺผสฺส\-ปจฺจยา อุปฺปชฺชติ เวทยิตํ สุขํ วา ทุกฺขํ วา อทุกฺขมสุขํ วา ตมฺปิ อาทิตฺตํ.}

\prose{เกน อาทิตฺตํ, ราคคฺคินา โทสคฺคินา โมหคฺคินา อาทิตฺตํ, ชาติยา ชราย มรเณน โสเกหิ ปริเทเวหิ ทุกฺเขหิ โทมนสฺเสหิ อุปายาเสหิ อาทิตฺตนฺติ วทามิ.}

\prose{กาโย อาทิตฺโต, โผฏฺฐพฺพา อาทิตฺตา, กายวิญฺญาณํ อาทิตฺตํ, กายสมฺผสฺโส อาทิตฺโต, ยมิทํ กาย\-สมฺผสฺส\-ปจฺจยา อุปฺปชฺชติ เวทยิตํ สุขํ วา ทุกฺขํ วา อทุกฺขมสุขํ วา ตมฺปิ อาทิตฺตํ.}

\prose{เกน อาทิตฺตํ, ราคคฺคินา โทสคฺคินา โมหคฺคินา อาทิตฺตํ, ชาติยา ชราย มรเณน โสเกหิ ปริเทเวหิ ทุกฺเขหิ โทมนสฺเสหิ อุปายาเสหิ อาทิตฺตนฺติ วทามิ.}

\prose{มโน อาทิตฺโต, ธมฺมา อาทิตฺตา, มโนวิญฺญาณํ อาทิตฺตํ, มโนสมฺผสฺโส อาทิตฺโต, ยมิทํ มโน\-สมฺผสฺส\-ปจฺจยา อุปฺปชฺชติ เวทยิตํ สุขํ วา ทุกฺขํ วา อทุกฺขมสุขํ วา ตมฺปิ อาทิตฺตํ.}

\prose{เกน อาทิตฺตํ, ราคคฺคินา โทสคฺคินา โมหคฺคินา อาทิตฺตํ, ชาติยา ชราย มรเณน โสเกหิ ปริเทเวหิ ทุกฺเขหิ โทมนสฺเสหิ อุปายาเสหิ อาทิตฺตนฺ”ติ วทามิ.}

\prose{\csromanfolio{45}เอวํ ปสฺสํ ภิกฺขเว สุตวา อริยสาวโก จกฺขุสฺมิมฺปิ นิพฺพินฺทติ, รูเปสุปิ นิพฺพินฺทติ, จกฺขุ\-วิญฺญาเณปิ นิพฺพินฺทติ, จกฺขุ\-สมฺผสฺเสปิ นิพฺพินฺทติ, ยมิทํ จกฺขุ\-สมฺผสฺส\-ปจฺจยา อุปฺปชฺชติ เวทยิตํ สุขํ วา ทุกฺขํวา อทุกฺขมสุขํ วา ตสฺมิมฺปิ นิพฺพินฺทติ.}

\setlength{\csromangathapagewidth}{0pt}
\setlength{\csromangathawakwidth}{0pt}
\csromangathameasuregroup{โสตสฺมิมฺปิ นิพฺพินฺทติ, \\ ฆานสฺมิมฺปิ นิพฺพินฺทติ, \\ ชิวฺหายปิ นิพฺพินฺทติ, \\ กายสฺมิมฺปิ นิพฺพินฺทติ,}{\csromangathabat{โสตสฺมิมฺปิ นิพฺพินฺทติ,}{สทฺเทสุปิ นิพฺพินฺทติ ฯเปฯ.} \\ \csromangathabat{ฆานสฺมิมฺปิ นิพฺพินฺทติ,}{คนฺเธสุปิ นิพฺพินฺทติ ฯเปฯ.} \\ \csromangathabat{ชิวฺหายปิ นิพฺพินฺทติ,}{รเสสุปิ นิพฺพินฺทติ ฯเปฯ.} \\ \csromangathabat{กายสฺมิมฺปิ นิพฺพินฺทติ,}{โผฏฺฐพฺเพสุปิ นิพฺพินฺทติ ฯเปฯ.}}
\csromangathasetleft{โสตสฺมิมฺปิ นิพฺพินฺทติ, \\ ฆานสฺมิมฺปิ นิพฺพินฺทติ, \\ ชิวฺหายปิ นิพฺพินฺทติ, \\ กายสฺมิมฺปิ นิพฺพินฺทติ,}
\csromangathagroup{\csromangathastanza{\csromangathabat{โสตสฺมิมฺปิ นิพฺพินฺทติ,}{สทฺเทสุปิ นิพฺพินฺทติ ฯเปฯ.} \\ \csromangathabat{ฆานสฺมิมฺปิ นิพฺพินฺทติ,}{คนฺเธสุปิ นิพฺพินฺทติ ฯเปฯ.}}\csromangathastanza{\csromangathabat{ชิวฺหายปิ นิพฺพินฺทติ,}{รเสสุปิ นิพฺพินฺทติ ฯเปฯ.} \\ \csromangathabat{กายสฺมิมฺปิ นิพฺพินฺทติ,}{โผฏฺฐพฺเพสุปิ นิพฺพินฺทติ ฯเปฯ.}}}

\prose{มนสฺมิมฺปิ นิพฺพินฺทติ, ธมฺเมสุปิ นิพฺพินฺทติ, มโนวิญฺญาเณปิ นิพฺพินฺทติ, มโน\-สมฺผสฺเสปิ นิพฺพินฺทติ, ยมิทํ มโน\-สมฺผสฺส\-ปจฺจยา อุปฺปชฺชติ เวทยิตํ สุขํ วา ทุกฺขํ วา อทุกฺขมสุขํ วา ตสฺมิมฺปิ นิพฺพินฺทติ.}

\prose{นิพฺพินฺทํ วิรชฺชติ, วิราคา วิมุจฺจติ, วิมุตฺตสฺมิํ วิมุตฺตมิติ ญาณํ โหติ, ขีณา ชาติ, วุสิตํ พฺรหฺมจริยํ, กตํ กรณียํ, นาปรํ อิตฺถตฺตายาติ ปชานาตีติ.}

\proseitemwithcloser{54}{อิมสฺมิํ จ ปน เวยฺยากรณสฺมิํ ภญฺญมาเน ตสฺส ภิกฺขุ\-สหสฺสสฺส อนุปาทาย อาสเวหิ จิตฺตานิ วิมุจฺจิํสุ.}{อาทิตฺตปริยาย\-สุตฺตํ นิฏฺฐิตํ.}
\nitthitamruled{อุรุเวล\-ปาฏิหาริยํ, ตติย\-ภาณวาโร นิฏฺฐิโต.}
\csromanmatikaanchor{mtk.14}
\csromantocmarkref{section}{13. พิมฺพิสารสมาคมกถา}{mtk.14}
\bookmark[dest={mtk.14},level=1]{13. พิมฺพิสารสมาคมกถา}
\csromanheader{1}{13. \textbf{พิมฺพิสาร}\-\textbf{สมาคม}\-\textbf{กถา}}
\proseitem{55}{อถ โข ภควา คยาสีเส ยถาภิรนฺตํ วิหริตฺวา เยน ราชคหํ เตน จาริกํ ปกฺกามิ มหตา ภิกฺขุ\-สํเฆน สทฺธิํ ภิกฺขุ\-สหสฺเสน สพฺเพเหว ปุราณชฏิเลหิ. อถ โข ภควา อนุปุพฺเพน จาริกํ จรมาโน เยน ราชคหํ ตทวสริ, ตตฺร สุทํ ภควา ราชคเห วิหรติ ลฏฺฐิวเน\csromansharedfootnote{d4899ccd7e0c9079}{ลฏฺฐิวนุยฺยาเน (สฺยา)} สุปฺปติฏฺเฐ เจติเย. อสฺโสสิ โข ราชา มาคโธ เสนิโย พิมฺพิสาโร “สมโณ ขลุ โภ โคตโม สกฺยปุตฺโต สกฺยกุลา ปพฺพชิโต ราชคหํ อนุปฺปตฺโต ราชคเห วิหรติ ลฏฺฐิวเน\csromansharedfootnote{d4899ccd7e0c9079}{ลฏฺฐิวนุยฺยาเน (สฺยา)} สุปฺปติฏฺเฐ เจติเย, ตํ โข ปน ภควนฺตํ\csromansharedfootnote{5ea8844b9bde44dc}{ภวนฺตํ (ก)} โคตมํ เอวํ กลฺยาโณ กิตฺติสทฺโท อพฺภุคฺคโต ‘อิติปิ โส ภควา อรหํ สมฺมาสมฺพุทฺโธ วิชฺชา\-จรณ\-สมฺปนฺโน สุคโต โลกวิทู อนุตฺตโร ปุริส\-ทมฺม\-สารถิ สตฺถา เทวมนุสฺสานํ พุทฺโธ ภควา’\csromansharedfootnote{475b0b93102c8b46}{ภควาติ (ก)}, โส อิมํ โลกํ สเทวกํ \csromanfolio{46}สมารกํ สพฺรหฺมกํ สสฺสมณ\-พฺราหฺมณิํ ปชํ สเทวมนุสฺสํ สยํ อภิญฺญา สจฺฉิกตฺวา ปเวเทติ, โส ธมฺมํ เทเสติ อาทิกลฺยาณํ มชฺเฌกลฺยาณํ ปริโยสาน\-กลฺยาณํ สาตฺถํ สพฺยญฺชนํ เกวล\-ปริปุณฺณํ ปริสุทฺธํ พฺรหฺมจริยํ ปกาเสติ, สาธุ โข ปน ตถารูปานํ อรหตํ ทสฺสนํ โหตี”ติ.}

\prose{อถ โข ราชา มาคโธ เสนิโย พิมฺพิสาโร ทฺวาทสนหุเตหิ\csromansharedfootnote{1e609960f99dca82}{ทฺวาทสนิยุเตหิ (โยชนา)} มาคธิเกหิ พฺราหฺมณคหปติเกหิ ปริวุโต เยน ภควา เตนุปสงฺกมิ, อุปสงฺกมิตฺวา ภควนฺตํ อภิวาเทตฺวา เอกมนฺตํ นิสีทิ. เตปิ โข ทฺวาทสนหุตา มาคธิกา พฺราหฺมณ\-คหปติกา อปฺเปกจฺเจ ภควนฺตํ อภิวาเทตฺวา เอกมนฺตํ นิสีทิํสุ, อปฺเปกจฺเจ ภควตา สทฺธิํ สมฺโมทิํสุ, สมฺโมทนียํ กถํ สารณียํ วีติสาเรตฺวา เอกมนฺตํ นิสีทิํสุ, อปฺเปกจฺเจ เยน ภควา เตนญฺชลิํ ปณาเมตฺวา เอกมนฺตํ นิสีทิํสุ, อปฺเปกจฺเจ ภควโต สนฺติเก นามโคตฺตํ สาเวตฺวา เอกมนฺตํ นิสีทิํสุ, อปฺเปกจฺเจ ตุณฺหีภูตา เอกมนฺตํ นิสีทิํสุ. อถ โข เตสํ ทฺวาทสนหุตานํ\csromansharedfootnote{08e2403044ca318b}{ทฺวาทสนิยุตานํ (โยชนา)} มาคธิกานํ พฺราหฺมณ\-คหปติกานํ เอตทโหสิ “กิํ นุ โข มหาสมโณ อุรุเวลกสฺสเป พฺรหฺมจริยํ จรติ, อุทาหุ อุรุเวลกสฺสโป มหาสมเณ พฺรหฺมจริยํ จรตี”ติ. อถ โข ภควา เตสํ ทฺวาทสนหุตานํ มาคธิกานํ พฺราหฺมณ\-คหปติกานํ เจตสา เจโต\-ปริวิตกฺกม\-ญฺญาย อายสฺมนฺตํ อุรุเวล\-กสฺสปํ คาถาย อชฺฌภาสิ\csromandash{}}

\setlength{\csromangathapagewidth}{0pt}
\setlength{\csromangathawakwidth}{0pt}
\csromangathameasuregroup{“กิเมว ทิสฺวา อุรุเวลวาสิ, \\ ปุจฺฉามิ ตํ กสฺสป เอตมตฺถํ,}{\csromangathabat{“กิเมว ทิสฺวา อุรุเวลวาสิ,}{ปหาสิ อคฺคิํ กิสโกวทาโน.} \\ \csromangathabat{ปุจฺฉามิ ตํ กสฺสป เอตมตฺถํ,}{กถํ ปหีนํ ตว อคฺคิหุตฺตนฺติ.} \\ รูเป จ สทฺเท จ อโถ รเส จ, \\ กามิตฺถิโย จาภิวทนฺติ ยญฺญา. \\ เอตํ มลนฺติ อุปธีสุ ญตฺวา, \\ ตสฺมา น ยิฏฺเฐ น หุเต อรญฺชินฺติ. \\ เอตฺเถว เต มโน น รมิตฺถ (กสฺสปาติ ภควา,) \\ รูเปสุ สทฺเทสุ อโถ รเสสุ. \\ อถ โก จรหิ เทวมนุสฺสโลเก, \\ รโต มโน กสฺสป พฺรูหิ เม ตนฺติ. \\ ทิสฺวา ปทํ สนฺตมนูปธีกํ, \\ อกิญฺจนํ กามภเว อสตฺตํ. \\ อนญฺญถาภาวิมนญฺญเนยฺยํ, \\ ตสฺมา น ยิฏฺเฐ น หุเต อรญฺชินฺ”ติ.}
\csromangathameasurewak{รูเป จ สทฺเท จ อโถ รเส จ, \\ กามิตฺถิโย จาภิวทนฺติ ยญฺญา. \\ เอตํ มลนฺติ อุปธีสุ ญตฺวา, \\ ตสฺมา น ยิฏฺเฐ น หุเต อรญฺชินฺติ. \\ เอตฺเถว เต มโน น รมิตฺถ (กสฺสปาติ ภควา,) \\ รูเปสุ สทฺเทสุ อโถ รเสสุ. \\ อถ โก จรหิ เทวมนุสฺสโลเก, \\ รโต มโน กสฺสป พฺรูหิ เม ตนฺติ. \\ ทิสฺวา ปทํ สนฺตมนูปธีกํ, \\ อกิญฺจนํ กามภเว อสตฺตํ. \\ อนญฺญถาภาวิมนญฺญเนยฺยํ, \\ ตสฺมา น ยิฏฺเฐ น หุเต อรญฺชินฺ”ติ.}
\csromangathasetleft{“กิเมว ทิสฺวา อุรุเวลวาสิ, \\ ปุจฺฉามิ ตํ กสฺสป เอตมตฺถํ,}
\csromangathagroup{\csromangathastanza{\csromangathaitembat{55}{“กิเมว ทิสฺวา อุรุเวลวาสิ,}{ปหาสิ อคฺคิํ กิสโกวทาโน.} \\ \csromangathacontbat{ปุจฺฉามิ ตํ กสฺสป เอตมตฺถํ,}{กถํ ปหีนํ ตว อคฺคิหุตฺตนฺติ.}}\csromangathastanza{\csromangathawak{รูเป จ สทฺเท จ อโถ รเส จ, \\ กามิตฺถิโย จาภิวทนฺติ ยญฺญา. \\ เอตํ มลนฺติ อุปธีสุ ญตฺวา, \\ ตสฺมา น ยิฏฺเฐ น หุเต อรญฺชินฺติ.}}\csromangathastanza{\csromangathawak{เอตฺเถว เต มโน น รมิตฺถ (กสฺสปาติ ภควา,) \\ รูเปสุ สทฺเทสุ อโถ รเสสุ. \\ อถ โก จรหิ เทว\-มนุสฺส\-โลเก, \\ รโต มโน กสฺสป พฺรูหิ เม ตนฺติ.}}\csromangathastanza{\csromangathawak{\csromanfolio{47}ทิสฺวา ปทํ สนฺตม\-นูปธีกํ, \\ อกิญฺจนํ กามภเว อสตฺตํ. \\ อนญฺญถาภาวิม\-นญฺญเนยฺยํ, \\ ตสฺมา น ยิฏฺเฐ น หุเต อรญฺชินฺ”ติ.}}}

\proseitem{56}{อถ โข อายสฺมา อุรุเวลกสฺสโป อุฏฺฐายาสนา เอกํสํ อุตฺตราสงฺคํ กริตฺวา ภควโต ปาเทสุ สิรสา นิปติตฺวา ภควนฺตํ เอตทโวจ “สตฺถา เม ภนฺเต ภควา, สาวโกหมสฺมิ, สตฺถา เม ภนฺเต ภควา, สาวโกหมสฺมี”ติ. อถ โข เตสํ ทฺวาทสนหุตานํ มาคธิกานํ พฺราหฺมณ\-คหปติกานํ เอตทโหสิ “อุรุเวลกสฺสโป มหาสมเณ พฺรหฺมจริยํ จรตี”ติ. อถ โข ภควา เตสํ ทฺวาทสนหุตานํ มาคธิกานํ พฺราหฺมณ\-คหปติกานํ เจตสา เจโต\-ปริวิตกฺกม\-ญฺญาย อนุปุพฺพิํ กถํ กเถสิ, เสยฺยถิทํ, ทานกถํ สีลกถํ สคฺคกถํ กามานํ อาทีนวํ โอการํ สํกิเลสํ เนกฺขมฺเม อานิสํสํ ปกาเสสิ. ยทา เต ภควา อญฺญาสิ กลฺลจิตฺเต มุทุจิตฺเต วินีวรณ\-จิตฺเต อุทคฺคจิตฺเต ปสนฺนจิตฺเต, อถ ยา พุทฺธานํ สามุกฺกํสิกา ธมฺมเทสนา, ตํ ปกาเสสิ ทุกฺขํ สมุทยํ นิโรธํ มคฺคํ, เสยฺยถาปิ นาม สุทฺธํ วตฺถํ อปคตกาฬกํ สมฺมเทว รชนํ ปฏิคฺคเณฺหยฺย, เอวเมว เอกาทส\-นหุตานํ มาคธิกานํ พฺราหฺมณ\-คหปติกานํ พิมฺพิสาร\-ปฺปมุขานํ ตสฺมิํ เยว อาสเน วิรชํ วีตมลํ ธมฺมจกฺขุํ อุทปาทิ “ยํ กิญฺจิ สมุทย\-ธมฺมํ สพฺพํ ตํ นิโรธธมฺมนฺ”ติ. เอกนหุตํ อุปาสกตฺตํ ปฏิเวเทสิ.}

\proseitem{57}{อถ โข ราชา มาคโธ เสนิโย พิมฺพิสาโร ทิฏฺฐธมฺโม ปตฺตธมฺโม วิทิตธมฺโม ปริโยคาฬฺห\-ธมฺโม ติณฺณ\-วิจิกิจฺโฉ วิคต\-กถํกโถ เวสารชฺช\-ปฺปตฺโต อปรปฺปจฺจโย สตฺถุสาสเน ภควนฺตํ เอตทโวจ “ปุพฺเพ เม ภนฺเต กุมารสฺส สโต ปญฺจ อสฺสาสกา อเหสุํ, เต เม เอตรหิ สมิทฺธา. ปุพฺเพ เม ภนฺเต กุมารสฺส สโต เอตทโหสิ ‘อโห วต มํ รชฺเช อภิสิญฺเจยฺยุนฺ’ติ, อยํ โข เม ภนฺเต ปฐโม อสฺสาสโก อโหสิ, โส เม เอตรหิ สมิทฺโธ. ‘ตสฺส จ เม วิชิตํ อรหํ สมฺมาสมฺพุทฺโธ โอกฺกเมยฺยา’ติ, อยํ โข เม ภนฺเต ทุติโย อสฺสาสโก อโหสิ, โส เม เอตรหิ สมิทฺโธ. ‘ตญฺจาหํ ภควนฺตํ ปยิรุปาเสยฺยนฺ’ติ, อยํ โข เม ภนฺเต \csromanfolio{48}ตติโย อสฺสาสโก อโหสิ, โส เม เอตรหิ สมิทฺโธ. ‘โส จ เม ภควา ธมฺมํ เทเสยฺยา’ติ, อยํ โข เม ภนฺเต จตุตฺโถ อสฺสาสโก อโหสิ, โส เม เอตรหิ สมิทฺโธ. ‘ตสฺส จาหํ ภควโต ธมฺมํ อาชาเนยฺยนฺ’ติ, อยํ โข เม ภนฺเต ปญฺจโม อสฺสาสโก อโหสิ, โส เม เอตรหิ สมิทฺโธ. ปุพฺเพ เม ภนฺเต กุมารสฺส สโต อิเม ปญฺจ อสฺสาสกา อเหสุํ, เต เม เอตรหิ สมิทฺธา. อภิกฺกนฺตํ ภนฺเต, อภิกฺกนฺตํ ภนฺเต, เสยฺยถาปิ ภนฺเต นิกฺกุชฺชิตํ วา อุกฺกุชฺเชยฺย, ปฏิจฺฉนฺนํ วา วิวเรยฺย, มูฬฺหสฺส วา มคฺคํ อาจิกฺเขยฺย, อนฺธกาเร วา เตลปชฺโชตํ ธาเรยฺย ‘จกฺขุมนฺโต รูปานิ ทกฺขนฺตี’ติ, เอวเมวํ ภควตา อเนก\-ปริยาเยน ธมฺโม ปกาสิโต. เอสาหํ ภนฺเต ภควนฺตํ สรณํ คจฺฉามิ ธมฺมญฺจ ภิกฺขุสํฆญฺจ, อุปาสกํ มํ\csromansharedfootnote{4e562e7a53a0b922}{มํ ภนฺเต (ก)} ภควา ธาเรตุ อชฺชตคฺเค ปาณุเปตํ สรณํ คตํ, อธิวาเสตุ จ เม ภนฺเต ภควา สฺวาตนาย ภตฺตํ สทฺธิํ ภิกฺขุสํเฆนา”ติ. อธิวาเสสิ ภควา ตุณฺหีภาเวน. อถ โข ราชา มาคโธ เสนิโย พิมฺพิสาโร ภควโต อธิวาสนํ วิทิตฺวา อุฏฺฐายาสนา ภควนฺตํ อภิวาเทตฺวา ปทกฺขิณํ กตฺวา ปกฺกามิ. อถ โข ราชา มาคโธ เสนิโย พิมฺพิสาโร ตสฺสา รตฺติยา อจฺจเยน ปณีตํ ขาทนียํ โภชนียํ ปฏิยาทาเปตฺวา ภควโต กาลํ อาโรจาเปสิ “กาโล ภนฺเต นิฏฺฐิตํ ภตฺตนฺ”ติ.}

\proseitem{58}{อถ โข ภควา ปุพฺพณฺหสมยํ นิวาเสตฺวา ปตฺต\-จีวรมา\-ทาย ราชคหํ ปาวิสิ มหตา ภิกฺขุ\-สํเฆน สทฺธิํ ภิกฺขุ\-สหสฺเสน สพฺเพเหว ปุราณชฏิเลหิ. เตน โข ปน สมเยน สกฺโก เทวานมินฺโท มาณวกวณฺณํ อภินิมฺมินิตฺวา พุทฺธ\-ปฺปมุขสฺส ภิกฺขุ\-สํฆสฺส ปุรโต ปุรโต คจฺฉติ อิมา คาถาโย คายมาโน\csromandash{}}

\setlength{\csromangathapagewidth}{0pt}
\setlength{\csromangathawakwidth}{0pt}
\csromangathameasuregroup{“ทนฺโต ทนฺเตหิ สห ปุราณชฏิเลหิ, \\ สิงฺคีนิกฺขสวณฺโณ, \\ มุตฺโต มุตฺเตหิ สห ปุราณชฏิเลหิ, \\ สิงฺคีนิกฺขสวณฺโณ, \\ ทสวาโส ทสพโล, \\ โส ทสสตปริวาโร\textsuperscript{1},}{\csromangathabat{“ทนฺโต ทนฺเตหิ สห ปุราณชฏิเลหิ,}{วิปฺปมุตฺโต วิปฺปมุตฺเตหิ.} \\ \csromangathabat{สิงฺคีนิกฺขสวณฺโณ,}{ราชคหํ ปาวิสิ ภควา.} \\ \csromangathabat{มุตฺโต มุตฺเตหิ สห ปุราณชฏิเลหิ,}{วิปฺปมุตฺโต วิปฺปมุตฺเตหิ.} \\ \csromangathabat{สิงฺคีนิกฺขสวณฺโณ,}{ราชคหํ ปาวิสิ ภควา.} \\ ติณฺโณ ติณฺเณหิ สห ปุราณชฏิเลหิ, \\ วิปฺปมุตฺโต วิปฺปมุตฺเตหิ. \\ สิงฺคีนิกฺขสวณฺโณ, \\ ราชคหํ ปาวิสิ ภควา. \\ สนฺโต สนฺเตหิ สห ปุราณชฏิเลหิ, \\ วิปฺปมุตฺโต วิปฺปมุตฺเตหิ. \\ สิงฺคีนิกฺขสวณฺโณ, \\ ราชคหํ ปาวิสิ ภควา. \\ \csromangathabat{ทสวาโส ทสพโล,}{ทสธมฺมวิทู ทสภิ จุเปโต.} \\ \csromangathabat{โส ทสสตปริวาโร\textsuperscript{1},}{ราชคหํ ปาวิสิ ภควา”ติ.}}
\csromangathameasurewak{ติณฺโณ ติณฺเณหิ สห ปุราณชฏิเลหิ, \\ วิปฺปมุตฺโต วิปฺปมุตฺเตหิ. \\ สิงฺคีนิกฺขสวณฺโณ, \\ ราชคหํ ปาวิสิ ภควา. \\ สนฺโต สนฺเตหิ สห ปุราณชฏิเลหิ, \\ วิปฺปมุตฺโต วิปฺปมุตฺเตหิ. \\ สิงฺคีนิกฺขสวณฺโณ, \\ ราชคหํ ปาวิสิ ภควา.}
\csromangathasetleft{“ทนฺโต ทนฺเตหิ สห ปุราณชฏิเลหิ, \\ สิงฺคีนิกฺขสวณฺโณ, \\ มุตฺโต มุตฺเตหิ สห ปุราณชฏิเลหิ, \\ สิงฺคีนิกฺขสวณฺโณ, \\ ทสวาโส ทสพโล, \\ โส ทสสตปริวาโร\textsuperscript{1},}
\csromangathagroup{\csromangathastanza{\csromangathabat{“ทนฺโต ทนฺเตหิ สห ปุราณชฏิเลหิ,}{วิปฺปมุตฺโต วิปฺปมุตฺเตหิ.} \\ \csromangathabat{สิงฺคีนิกฺขสวณฺโณ,}{ราชคหํ ปาวิสิ ภควา.}}\csromangathastanza{\csromangathabat{มุตฺโต มุตฺเตหิ สห ปุราณชฏิเลหิ,}{วิปฺปมุตฺโต วิปฺปมุตฺเตหิ.} \\ \csromangathabat{สิงฺคีนิกฺขสวณฺโณ,}{ราชคหํ ปาวิสิ ภควา.}}\csromangathastanza{\csromangathawak{\csromanfolio{49}ติณฺโณ ติณฺเณหิ สห ปุราณชฏิเลหิ, \\ วิปฺปมุตฺโต วิปฺปมุตฺเตหิ. \\ สิงฺคีนิกฺขสวณฺโณ, \\ ราชคหํ ปาวิสิ ภควา.}}\csromangathastanza{\csromangathawak{สนฺโต สนฺเตหิ สห ปุราณชฏิเลหิ, \\ วิปฺปมุตฺโต วิปฺปมุตฺเตหิ. \\ สิงฺคีนิกฺขสวณฺโณ, \\ ราชคหํ ปาวิสิ ภควา.}}\csromangathastanza{\csromangathabat{ทสวาโส ทสพโล,}{ทสธมฺมวิทู ทสภิ จุเปโต.} \\ \csromangathabat{โส ทสสตปริวาโร\csromansharedfootnote{b971cb163da20356}{ปริวารโก (ก)},}{ราชคหํ ปาวิสิ ภควา”ติ.}}}

\prose{มนุสฺสา สกฺกํ เทวานมินฺทํ ปสฺสิตฺวา เอวมาหํสุ “อภิรูโป วตายํ มาณวโก, ทสฺสนีโย วตายํ มาณวโก, ปาสาทิโก วตายํ มาณวโก, กสฺส นุ โข อยํ มาณวโก”ติ. เอวํ วุตฺเต สกฺโก เทวานมินฺโท เต มนุสฺเส คาถาย อชฺฌภาสิ\csromandash{}}

\setlength{\csromangathapagewidth}{0pt}
\setlength{\csromangathawakwidth}{0pt}
\csromangathameasuregroup{“โย ธีโร สพฺพธิ ทนฺโต, \\ อรหํ สุคโต โลเก,}{\csromangathabat{“โย ธีโร สพฺพธิ ทนฺโต,}{สุทฺโธ อปฺปฏิปุคฺคโล.} \\ \csromangathabat{อรหํ สุคโต โลเก,}{ตสฺสาหํ ปริจารโก”ติ.}}
\csromangathasetleft{“โย ธีโร สพฺพธิ ทนฺโต, \\ อรหํ สุคโต โลเก,}
\csromangathagroup{\csromangathastanza{\csromangathabat{“โย ธีโร สพฺพธิ ทนฺโต,}{สุทฺโธ อปฺปฏิปุคฺคโล.} \\ \csromangathabat{อรหํ สุคโต โลเก,}{ตสฺสาหํ ปริจารโก”ติ.}}}

\proseitemwithcloserruled{59}{อถ โข ภควา เยน รญฺโญ มาคธสฺส เสนิยสฺส พิมฺพิสารสฺส นิเวสนํ เตนุปสงฺกมิ, อุปสงฺกมิตฺวา ปญฺญตฺเต อาสเน นิสีทิ สทฺธิํ ภิกฺขุ\-สํเฆน. อถ โข ราชา มาคโธ เสนิโย พิมฺพิสาโร พุทฺธ\-ปฺปมุขํ ภิกฺขุสํฆํ ปณีเตน ขาทนีเยน โภชนีเยน สหตฺถา สนฺตปฺเปตฺวา สมฺปวาเรตฺวา ภควนฺตํ ภุตฺตาวิํ โอนีต\-ปตฺต\-ปาณิํ เอกมนฺตํ นิสีทิ, เอกมนฺตํ นิสินฺนสฺส โข รญฺโญ มาคธสฺส เสนิยสฺส พิมฺพิสารสฺส เอตทโหสิ\csromansymbolfootnote{*}{วิ 4. 306 ปิฏฺเฐปิ.}“กตฺถ นุ โข ภควา วิหเรยฺย, ยํ อสฺส คามโต เนว อติทูเร น อจฺจาสนฺเน คมนาคมน\-สมฺปนฺนํ อตฺถิกานํ อตฺถิกานํ มนุสฺสานํ อภิกฺกมนียํ ทิวา อปฺปากิณฺณํ\csromansharedfootnote{c0a91bb1f8ca16e1}{อปฺปกิณฺณํ (สี, สฺยา), อพฺโพกิณฺณํ (ก)} รตฺติํ อปฺปสทฺทํ อปฺปนิคฺโฆสํ วิชนวาตํ มนุสฺส\-ราหสฺเสยฺยกํ ปฏิสลฺลานสารุปฺปนฺ”ติ. อถ โข รญฺโญ มาคธสฺส เสนิยสฺส พิมฺพิสารสฺส เอตทโหสิ “อิทํ โข อมฺหากํ เวฬุวนํ อุยฺยานํ คามโต เนว อติทูเร น อจฺจาสนฺเน \csromanfolio{50}คมนาคมน\-สมฺปนฺนํ อตฺถิกานํ อตฺถิกานํ มนุสฺสานํ อภิกฺกมนียํ ทิวา อปฺปากิณฺณํ รตฺติํ อปฺปสทฺทํ อปฺปนิคฺโฆสํ วิชนวาตํ มนุสฺส\-ราหสฺเสยฺยกํ ปฏิสลฺลาน\-สารุปฺปํ, ยนฺนูนาหํ เวฬุวนํ อุยฺยานํ พุทฺธ\-ปฺปมุขสฺส ภิกฺขุ\-สํฆสฺส ทเทยฺยนฺ”ติ. อถ โข ราชา มาคโธ เสนิโย พิมฺพิสาโร โสวณฺณมยํ ภิงฺการํ คเหตฺวา ภควโต โอโณเชสิ “เอตาหํ ภนฺเต เวฬุวนํ อุยฺยานํ พุทฺธ\-ปฺปมุขสฺส ภิกฺขุ\-สํฆสฺส ทมฺมี”ติ. ปฏิคฺคเหสิ ภควา อารามํ. อถ โข ภควา ราชานํ มาคธํ เสนิยํ พิมฺพิสารํ ธมฺมิยา กถาย สนฺทสฺเสตฺวา สมาทเปตฺวา สมุตฺเตเชตฺวา สมฺปหํเสตฺวา อุฏฺฐายาสนา ปกฺกามิ. อถ โข ภควา เอตสฺมิํ นิทาเน เอตสฺมิํ ปกรเณ ธมฺมิํ กถํ กตฺวา ภิกฺขู อามนฺเตสิ “อนุชานามิ ภิกฺขเว อารามนฺ”ติ.}{พิมฺพิสาร\-สมาคม\-กถา นิฏฺฐิตา.}
\csromanmatikaanchor{mtk.15}
\csromantocmarkref{section}{14. สาริปุตฺตโมคฺคลฺลานปพฺพชฺชากถา}{mtk.15}
\bookmark[dest={mtk.15},level=1]{14. สาริปุตฺตโมคฺคลฺลานปพฺพชฺชากถา}
\csromanheader{1}{14. \textbf{สาริปุตฺตโมคฺคลฺลาน}\-\textbf{ปพฺพชฺชา}\-\textbf{กถา}}
\proseitem{60}{เตน โข ปน สมเยน สญฺจโย\csromansharedfootnote{099540057c8e0893}{สญฺชโย (สี, สฺยา)} ปริพฺพาชโก ราชคเห ปฏิวสติ มหติยา ปริพฺพาชก\-ปริสาย สทฺธิํ อฑฺฒเตยฺเยหิ ปริพฺพาชก\-สเตหิ. เตน โข ปน สมเยน สาริปุตฺต\-โมคฺคลฺลานา สญฺจเย ปริพฺพาชเก พฺรหฺมจริยํ จรนฺติ, เตหิ กติกา กตา โหติ “โย ปฐมํ อมตํ อธิคจฺฉติ, โส อิตรสฺส อาโรเจตู”ติ. อถ โข อายสฺมา อสฺสชิ ปุพฺพณฺหสมยํ นิวาเสตฺวา ปตฺต\-จีวรมา\-ทาย ราชคหํ ปิณฺฑาย ปาวิสิ ปาสาทิเกน อภิกฺกนฺเตน ปฏิกฺกนฺเตน อาโลกิเตน วิโลกิเตน สมิญฺชิเตน ปสาริเตน โอกฺขิตฺตจกฺขุ อิริยาปถ\-สมฺปนฺโน. อทฺทสา โข สาริปุตฺโต ปริพฺพาชโก อายสฺมนฺตํ อสฺสชิํ ราชคเห ปิณฺฑาย จรนฺตํ ปาสาทิเกน อภิกฺกนฺเตน ปฏิกฺกนฺเตน อาโลกิเตน วิโลกิเตน สมิญฺชิเตน ปสาริเตน โอกฺขิตฺต\-จกฺขุํ อิริยาปถ\-สมฺปนฺนํ, ทิสฺวานสฺส เอตทโหสิ “เย วต โลเก อรหนฺโต วา อรหตฺตมคฺคํ วา สมาปนฺนา, อยํ เตสํ ภิกฺขุ อญฺญตโร, ยนฺนูนาหํ อิมํ ภิกฺขุํ \csromanfolio{51}อุปสงฺกมิตฺวา ปุจฺเฉยฺยํ ‘กํสิ ตฺวํ อาวุโส อุทฺทิสฺส ปพฺพชิโต, โก วา เต สตฺถา, กสฺส วา ตฺวํ ธมฺมํ โรเจสี”ติ. อถ โข สาริปุตฺตสฺส ปริพฺพาชกสฺส เอตทโหสิ “อกาโล โข อิมํ ภิกฺขุํ ปุจฺฉิตุํ อนฺตรฆรํ ปวิฏฺโฐ ปิณฺฑาย จรติ, ยนฺนูนาหํ อิมํ ภิกฺขุํ ปิฏฺฐิโต ปิฏฺฐิโต อนุพนฺเธยฺยํ อตฺถิเกหิ อุปญฺญาตํ มคฺคนฺ”ติ. อถ โข อายสฺมา อสฺสชิ ราชคเห ปิณฺฑาย จริตฺวา ปิณฺฑปาตํ อาทาย ปฏิกฺกมิ. อถ โข สาริปุตฺโตปิ ปริพฺพาชโก เยนายสฺมา อสฺสชิ เตนุปสงฺกมิ, อุปสงฺกมิตฺวา อายสฺมตา อสฺสชินา สทฺธิํ สมฺโมทิ, สมฺโมทนียํ กถํ สารณียํ วีติสาเรตฺวา เอกมนฺตํ อฏฺฐาสิ, เอกมนฺตํ ฐิโต โข สาริปุตฺโต ปริพฺพาชโก อายสฺมนฺตํ อสฺสชิํ เอตทโวจ “วิปฺปสนฺนานิ โข เต อาวุโส อินฺทฺริยานิ, ปริสุทฺโธ ฉวิวณฺโณ ปริโยทาโต, กํสิ ตฺวํ อาวุโส อุทฺทิสฺส ปพฺพชิโต, โก วา เต สตฺถา, กสฺส วา ตฺวํ ธมฺมํ โรเจสี”ติ. อตฺถาวุโส มหาสมโณ สกฺยปุตฺโต สกฺยกุลา ปพฺพชิโต, ตาหํ ภควนฺตํ อุทฺทิสฺส ปพฺพชิโต, โส จ เม ภควา สตฺถา, ตสฺส จาหํ ภควโต ธมฺมํ โรเจมีติ. กิํวาที ปนายสฺมโต สตฺถา กิมกฺขายีติ. อหํ โข อาวุโส นโว อจิรปพฺพชิโต อธุนาคโต อิมํ ธมฺมวินยํ, น ตาหํ สกฺโกมิ วิตฺถาเรน ธมฺมํ เทเสตุํ, อปิ จ เต สํขิตฺเตน อตฺถํ วกฺขามีติ. อถ โข สาริปุตฺโต ปริพฺพาชโก อายสฺมนฺตํ อสฺสชิํ เอตทโวจ “โหตุ อาวุโส,}

\setlength{\csromangathapagewidth}{0pt}
\setlength{\csromangathawakwidth}{0pt}
\csromangathameasuregroup{อปฺปํ วา พหุํ วา ภาสสฺสุ, \\ อตฺเถเนว เม อตฺโถ,}{\csromangathabat{อปฺปํ วา พหุํ วา ภาสสฺสุ,}{อตฺถํเยว เม พฺรูหิ.} \\ \csromangathabat{อตฺเถเนว เม อตฺโถ,}{กิํ กาหสิ พฺยญฺชนํ พหุนฺ”ติ.}}
\csromangathasetleft{อปฺปํ วา พหุํ วา ภาสสฺสุ, \\ อตฺเถเนว เม อตฺโถ,}
\csromangathagroup{\csromangathastanza{\csromangathabat{อปฺปํ วา พหุํ วา ภาสสฺสุ,}{อตฺถํเยว เม พฺรูหิ.} \\ \csromangathabat{อตฺเถเนว เม อตฺโถ,}{กิํ กาหสิ พฺยญฺชนํ พหุนฺ”ติ.}}}

\prose{อถ โข อายสฺมา อสฺสชิ สาริปุตฺตสฺส ปริพฺพาชกสฺส อิมํ ธมฺม\-ปริยายํ อภาสิ\csromandash{}}

\setlength{\csromangathapagewidth}{0pt}
\setlength{\csromangathawakwidth}{0pt}
\csromangathameasuregroup{\textsuperscript{*}“เย ธมฺมา เหตุปฺปภวา, \\ เตสญฺจ โย นิโรโธ,}{\csromangathabat{\textsuperscript{*}“เย ธมฺมา เหตุปฺปภวา,}{เตสํ เหตุํ ตถาคโต อาห.} \\ \csromangathabat{เตสญฺจ โย นิโรโธ,}{เอวํวาที มหาสมโณ”ติ.}}
\csromangathasetleft{\textsuperscript{*}“เย ธมฺมา เหตุปฺปภวา, \\ เตสญฺจ โย นิโรโธ,}
\csromangathagroup{\csromangathastanza{\csromangathabat{\csromansymbolfootnote{*}{ขุ 3. 27 ปิฏฺเฐ เถราปทาเนปิ.}“เย ธมฺมา เหตุปฺปภวา,}{เตสํ เหตุํ ตถาคโต อาห.} \\ \csromangathabat{เตสญฺจ โย นิโรโธ,}{เอวํวาที มหาสมโณ”ติ.}}}

\prose{อถ โข สาริปุตฺตสฺส ปริพฺพาชกสฺส อิมํ ธมฺม\-ปริยายํ สุตฺวา วิรชํ วีตมลํ ธมฺมจกฺขุํ อุทปาทิ “ยํ กิญฺจิ สมุทย\-ธมฺมํ สพฺพํ ตํ นิโรธธมฺมนฺ”ติ.}

\setlength{\csromangathapagewidth}{0pt}
\setlength{\csromangathawakwidth}{0pt}
\csromangathameasuregroup{\textsuperscript{*}เอเสว ธมฺโม ยทิ ตาวเทว, \\ อทิฏฺฐํ อพฺภตีตํ,}{\csromangathabat{\textsuperscript{*}เอเสว ธมฺโม ยทิ ตาวเทว,}{ปจฺจพฺยตฺถ ปทมโสกํ.} \\ \csromangathabat{อทิฏฺฐํ อพฺภตีตํ,}{พหุเกหิ กปฺปนหุเตหีติ.}}
\csromangathasetleft{\textsuperscript{*}เอเสว ธมฺโม ยทิ ตาวเทว, \\ อทิฏฺฐํ อพฺภตีตํ,}
\csromangathagroup{\csromangathastanza{\csromanfolio{52}\csromangathabat{\csromansymbolfootnote{*}{ขุ 3. 28 ปิฏฺเฐ เถราปทาเนปิ.}เอเสว ธมฺโม ยทิ ตาวเทว,}{ปจฺจพฺยตฺถ ปทมโสกํ.} \\ \csromangathabat{อทิฏฺฐํ อพฺภตีตํ,}{พหุเกหิ กปฺปนหุเตหีติ.}}}

\proseitem{61}{อถ โข สาริปุตฺโต ปริพฺพาชโก เยน โมคฺคลฺลาโน ปริพฺพาชโก เตนุปสงฺกมิ. อทฺทสา โข โมคฺคลฺลาโน ปริพฺพาชโก สาริปุตฺตํ ปริพฺพาชกํ ทูรโตว อาคจฺฉนฺตํ, ทิสฺวาน สาริปุตฺตํ ปริพฺพาชกํ เอตทโวจ “วิปฺปสนฺนานิ โข เต อาวุโส อินฺทฺริยานิ, ปริสุทฺโธ ฉวิวณฺโณ ปริโยทาโต, กจฺจิ โน ตฺวํ อาวุโส อมตํ อธิคโต”ติ. อามาวุโส อมตํ อธิคโตติ. ยถากถํ ปน ตฺวํ อาวุโส อมตํ อธิคโตติ. อิธาหํ อาวุโส อทฺทสํ อสฺสชิํ ภิกฺขุํ ราชคเห ปิณฺฑาย จรนฺตํ ปาสาทิเกน อภิกฺกนฺเตน ปฏิกฺกนฺเตน อาโลกิเตน วิโลกิเตน สมิญฺชิเตน ปสาริเตน โอกฺขิตฺต\-จกฺขุํ อิริยาปถ\-สมฺปนฺนํ, ทิสฺวาน เม เอตทโหสิ “เย วต โลเก อรหนฺโต วา อรหตฺตมคฺคํ วา สมาปนฺนา, อยํ เตสํ ภิกฺขุ อญฺญตโร, ยนฺนูนาหํ อิมํ ภิกฺขุํ อุปสงฺกมิตฺวา ปุจฺเฉยฺยํ “กํสิ ตฺวํ อาวุโส อุทฺทิสฺส ปพฺพชิโต, โก วา เต สตฺถา, กสฺส วา ตฺวํ ธมฺมํ โรเจสี”ติ. ตสฺส มยฺหํ อาวุโส เอตทโหสิ “อกาโล โข อิมํ ภิกฺขุํ ปุจฺฉิตุํ อนฺตรฆรํ ปวิฏฺโฐ ปิณฺฑาย จรติ, ยนฺนูนาหํ อิมํ ภิกฺขุํ ปิฏฺฐิโต ปิฏฺฐิโต อนุพนฺเธยฺยํ อตฺถิเกหิ อุปญฺญาตํ มคฺคนฺ”ติ. อถ โข อาวุโส อสฺสชิ ภิกฺขุ ราชคเห ปิณฺฑาย จริตฺวา ปิณฺฑปาตํ อาทาย ปฏิกฺกมิ. อถ ขฺวาหํ อาวุโส เยน อสฺสชิ ภิกฺขุ เตนุปสงฺกมิํ, อุปสงฺกมิตฺวา อสฺสชินา ภิกฺขุนา สทฺธิํ สมฺโมทิํ, สมฺโมทนียํ กถํ สารณียํ วีติสาเรตฺวา เอกมนฺตํ อฏฺฐาสิํ, เอกมนฺตํ ฐิโต โข อหํ อาวุโส อสฺสชิํ ภิกฺขุํ เอตทโวจํ “วิปฺปสนฺนานิ โข เต อาวุโส อินฺทฺริยานิ, ปริสุทฺโธ ฉวิวณฺโณ ปริโยทาโต, กํสิ ตฺวํ อาวุโส อุทฺทิสฺส ปพฺพชิโต, โก วา เต สตฺถา, กสฺส วา ตฺวํ ธมฺมํ โรเจสีติ. อตฺถาวุโส มหาสมโณ สกฺยปุตฺโต สกฺยกุลา ปพฺพชิโต, ตาหํ ภควนฺตํ อุทฺทิสฺส ปพฺพชิโต, โส จ เม ภควา สตฺถา, ตสฺส จาหํ ภควโต ธมฺมํ โรเจมีติ. กิํ วาที ปนายสฺมโต สตฺถา กิมกฺขายีติ. อหํ โข อาวุโส นโว อจิรปพฺพชิโต อธุนาคโต อิมํ ธมฺมวินยํ, น ตาหํ สกฺโกมิ วิตฺถาเรน ธมฺมํ เทเสตุํ, อปิ จ เต สํขิตฺเตน อตฺถํ \csromanfolio{53}วกฺขามี”ติ. อถ ขฺวาหํ อาวุโส อสฺสชิํ ภิกฺขุํ เอตทโวจํ “โหตุ อาวุโส,}

\setlength{\csromangathapagewidth}{0pt}
\setlength{\csromangathawakwidth}{0pt}
\csromangathameasuregroup{อปฺปํ วา พหุํ วา ภาสสฺสุ, \\ อตฺเถเนว เม อตฺโถ,}{\csromangathabat{อปฺปํ วา พหุํ วา ภาสสฺสุ,}{อตฺถํเยว เม พฺรูหิ.} \\ \csromangathabat{อตฺเถเนว เม อตฺโถ,}{กิํ กาหสิ พฺยญฺชนํ พหุนฺ”ติ.}}
\csromangathasetleft{อปฺปํ วา พหุํ วา ภาสสฺสุ, \\ อตฺเถเนว เม อตฺโถ,}
\csromangathagroup{\csromangathastanza{\csromangathabat{อปฺปํ วา พหุํ วา ภาสสฺสุ,}{อตฺถํเยว เม พฺรูหิ.} \\ \csromangathabat{อตฺเถเนว เม อตฺโถ,}{กิํ กาหสิ พฺยญฺชนํ พหุนฺ”ติ.}}}

\prose{อถ โข อาวุโส อสฺสชิ ภิกฺขุ อิมํ ธมฺม\-ปริยายํ อภาสิ\csromandash{}}

\setlength{\csromangathapagewidth}{0pt}
\setlength{\csromangathawakwidth}{0pt}
\csromangathameasuregroup{“เย ธมฺมา เหตุปฺปภวา, \\ เตสญฺจ โย นิโรโธ,}{\csromangathabat{“เย ธมฺมา เหตุปฺปภวา,}{เตสํ เหตุํ ตถาคโต อาห.} \\ \csromangathabat{เตสญฺจ โย นิโรโธ,}{เอวํวาที มหาสมโณ”ติ.}}
\csromangathasetleft{“เย ธมฺมา เหตุปฺปภวา, \\ เตสญฺจ โย นิโรโธ,}
\csromangathagroup{\csromangathastanza{\csromangathabat{“เย ธมฺมา เหตุปฺปภวา,}{เตสํ เหตุํ ตถาคโต อาห.} \\ \csromangathabat{เตสญฺจ โย นิโรโธ,}{เอวํวาที มหาสมโณ”ติ.}}}

\prose{อถ โข โมคฺคลฺลานสฺส ปริพฺพาชกสฺส อิมํ ธมฺม\-ปริยายํ สุตฺวา วิรชํ วีตมลํ ธมฺมจกฺขุํ อุทปาทิ “ยํ กิญฺจิ สมุทย\-ธมฺมํ สพฺพํ ตํ นิโรธธมฺมนฺ”ติ.}

\setlength{\csromangathapagewidth}{0pt}
\setlength{\csromangathawakwidth}{0pt}
\csromangathameasuregroup{เอเสว ธมฺโม ยทิ ตาวเทว, \\ อทิฏฺฐํ อพฺภตีตํ,}{\csromangathabat{เอเสว ธมฺโม ยทิ ตาวเทว,}{ปจฺจพฺยตฺถ ปทมโสกํ.} \\ \csromangathabat{อทิฏฺฐํ อพฺภตีตํ,}{พหุเกหิ กปฺปนหุเตหีติ.}}
\csromangathasetleft{เอเสว ธมฺโม ยทิ ตาวเทว, \\ อทิฏฺฐํ อพฺภตีตํ,}
\csromangathagroup{\csromangathastanza{\csromangathabat{เอเสว ธมฺโม ยทิ ตาวเทว,}{ปจฺจพฺยตฺถ ปทมโสกํ.} \\ \csromangathabat{อทิฏฺฐํ อพฺภตีตํ,}{พหุเกหิ กปฺปนหุเตหีติ.}}}

\proseitem{62}{อถ โข โมคฺคลฺลาโน ปริพฺพาชโก สาริปุตฺตํ ปริพฺพาชกํ เอตทโวจ “คจฺฉาม มยํ อาวุโส ภควโต สนฺติเก, โส โน ภควา สตฺถา”ติ. อิมานิ โข อาวุโส อฑฺฒเตยฺยานิ ปริพฺพาชก\-สตานิ อเมฺห นิสฺสาย อเมฺห สมฺปสฺสนฺตา อิธ วิหรนฺติ เตปิ ตาว อปโลเกม\csromansharedfootnote{1a22a4fdc4e5fa44}{อปโลกาม (ก)}, ยถา เต มญฺญิสฺสนฺติ, ตถา เต กริสฺสนฺตีติ. อถ โข สาริปุตฺต\-โมคฺคลฺลานา เยน เต ปริพฺพาชกา เตนุ\-ปสงฺกมิํสุ, อุปสงฺกมิตฺวา เต ปริพฺพาชเก เอตทโวจุํ “คจฺฉาม มยํ อาวุโส ภควโต สนฺติเก, โส โน ภควา สตฺถา”ติ. มยํ อายสฺมนฺเต นิสฺสาย อายสฺมนฺเต สมฺปสฺสนฺตา อิธ วิหราม, สเจ อายสฺมนฺตา มหาสมเณ พฺรหฺมจริยํ จริสฺสนฺติ, สพฺเพว มยํ มหาสมเณ พฺรหฺมจริยํ จริสฺสามาติ. อถ โข สาริปุตฺต\-โมคฺคลฺลานา เยน สญฺจโย ปริพฺพาชโก เตนุ\-ปสงฺกมิํสุ, อุปสงฺกมิตฺวา สญฺจยํ ปริพฺพาชกํ เอตทโวจุํ “คจฺฉาม มยํ อาวุโส ภควโต สนฺติเก, โส โน ภควา สตฺถา”ติ. อลํ อาวุโส มา อคมิตฺถ, สพฺเพว ตโย อิมํ คณํ ปริหริสฺสามาติ. ทุติยมฺปิ โข ฯเปฯ. ตติยมฺปิ โข สาริปุตฺต\-โมคฺคลฺลานา \csromanfolio{54}สญฺจยํ ปริพฺพาชกํ เอตทโวจุํ “คจฺฉาม มยํ อาวุโส ภควโต สนฺติเก, โส โน ภควา สตฺถา”ติ. อลํ อาวุโส มา อคมิตฺถ, สพฺเพว ตโย อิมํ คณํ ปริหริสฺสามาติ. อถ โข สาริปุตฺต\-โมคฺคลฺลานา ตานิ อฑฺฒเตยฺยานิ ปริพฺพาชก\-สตานิ อาทาย เยน เวฬุวนํ เตนุ\-ปสงฺกมิํสุ. สญฺจยสฺส ปน ปริพฺพาชกสฺส ตตฺเถว อุณฺหํ โลหิตํ มุขโต อุคฺคญฺฉิ.}

\prose{อทฺทสา โข ภควา\csromansharedfootnote{82988d6d22501655}{ภควา เต (ก)} สาริปุตฺต\-โมคฺคลฺลาเน ทูรโตว อาคจฺฉนฺเต, ทิสฺวาน ภิกฺขู อามนฺเตสิ “เอเต ภิกฺขเว เทฺว สหายกา อาคจฺฉนฺติ โกลิโต อุปติสฺโส จ, เอตํ เม สาวกยุคํ ภวิสฺสติ อคฺคํ ภทฺทยุคนฺ”ติ.}

\setlength{\csromangathapagewidth}{0pt}
\setlength{\csromangathawakwidth}{0pt}
\csromangathameasuregroup{คมฺภีเร ญาณวิสเย, \\ วิมุตฺเต อปฺปตฺเต เวฬุวนํ, \\ เอเต เทฺว สหายกา, \\ เอตํ เม สาวกยุคํ,}{\csromangathabat{คมฺภีเร ญาณวิสเย,}{อนุตฺตเร อุปธิสงฺขเย.} \\ \csromangathabat{วิมุตฺเต อปฺปตฺเต เวฬุวนํ,}{อถ เน สตฺถา พฺยากาสิ.} \\ \csromangathabat{เอเต เทฺว สหายกา,}{อาคจฺฉนฺติ โกลิโต อุปติสฺโส จ.} \\ \csromangathabat{เอตํ เม สาวกยุคํ,}{ภวิสฺสติ อคฺคํ ภทฺทยุคนฺติ.}}
\csromangathasetleft{คมฺภีเร ญาณวิสเย, \\ วิมุตฺเต อปฺปตฺเต เวฬุวนํ, \\ เอเต เทฺว สหายกา, \\ เอตํ เม สาวกยุคํ,}
\csromangathagroup{\csromangathastanza{\csromangathabat{คมฺภีเร ญาณวิสเย,}{อนุตฺตเร อุปธิสงฺขเย.} \\ \csromangathabat{วิมุตฺเต อปฺปตฺเต เวฬุวนํ,}{อถ เน สตฺถา พฺยากาสิ.}}\csromangathastanza{\csromangathabat{เอเต เทฺว สหายกา,}{อาคจฺฉนฺติ โกลิโต อุปติสฺโส จ.} \\ \csromangathabat{เอตํ เม สาวกยุคํ,}{ภวิสฺสติ อคฺคํ ภทฺทยุคนฺติ.}}}

\prose{อถ โข สาริปุตฺต\-โมคฺคลฺลานา เยน ภควา เตนุ\-ปสงฺกมิํสุ, อุปสงฺกมิตฺวา ภควโต ปาเทสุ สิรสา นิปติตฺวา ภควนฺตํ เอตทโวจุํ “ลเภยฺยาม มยํ ภนฺเต ภควโต สนฺติเก ปพฺพชฺชํ, ลเภยฺยาม อุปสมฺปทนฺ”ติ.}

\prose{“เอถ ภิกฺขโว”ติ ภควา อโวจ “สฺวากฺขาโต ธมฺโม, จรถ พฺรหฺมจริยํ สมฺมา ทุกฺขสฺส อนฺตกิริยายา”ติ.}

\setlength{\csromangathapagewidth}{0pt}
\setlength{\csromangathawakwidth}{0pt}
\csromangathameasuregroup{}{สาว เตสํ อายสฺมนฺตานํ อุปสมฺปทา อโหสิ.}
\csromangathameasurewak{สาว เตสํ อายสฺมนฺตานํ อุปสมฺปทา อโหสิ.}
\setlength{\csromangathaleftcol}{0pt}
\csromangathagroup{\csromangathastanza{\csromangathawak{สาว เตสํ อายสฺมนฺตานํ อุปสมฺปทา อโหสิ.}}}

\csromanheader{2}{\textbf{อภิญฺญาตานํ ปพฺพชฺชา}}
\prose{เตน โข ปน สมเยน อภิญฺญาตา อภิญฺญาตา มาคธิกา กุลปุตฺตา ภควติ พฺรหฺมจริยํ จรนฺติ. มนุสฺสา อุชฺฌายนฺติ ขิยฺยนฺติ วิปาเจนฺติ “อปุตฺตกตาย ปฏิปนฺโน สมโณ โคตโม, เวธพฺยาย ปฏิปนฺโน สมโณ โคตโม, กุลุปจฺเฉทาย ปฏิปนฺโน สมโณ โคตโม, อิทานิ อเนน ชฏิลสหสฺสํ ปพฺพาชิตํ, อิมานิ จ อฑฺฒเตยฺยานิ ปริพฺพาชก\-สตานิ สญฺจยานิ\csromansharedfootnote{a2897d4c309ecf25}{สญฺเชยฺยานิ (สี), สญฺชยานิ (สฺยา)} ปพฺพาชิตานิ, อิเม จ อภิญฺญาตา อภิญฺญาตา}

\prose{\csromanfolio{55}มาคธิกา กุลปุตฺตา สมเณ โคตเม พฺรหฺมจริยํ จรนฺตี”ติ. อปิสฺสุ ภิกฺขู ทิสฺวา อิมาย คาถาย โจเทนฺติ\csromandash{}}

\setlength{\csromangathapagewidth}{0pt}
\setlength{\csromangathawakwidth}{0pt}
\csromangathameasuregroup{“อาคโต โข มหาสมโณ, \\ สพฺเพ สญฺจเย เนตฺวาน\textsuperscript{1},}{\csromangathabat{“อาคโต โข มหาสมโณ,}{มาคธานํ คริพฺพชํ.} \\ \csromangathabat{สพฺเพ สญฺจเย เนตฺวาน\textsuperscript{1},}{กํสุ ทานิ นยิสฺสตี”ติ.}}
\csromangathasetleft{“อาคโต โข มหาสมโณ, \\ สพฺเพ สญฺจเย เนตฺวาน\textsuperscript{1},}
\csromangathagroup{\csromangathastanza{\csromangathabat{“อาคโต โข มหาสมโณ,}{มาคธานํ คริพฺพชํ.} \\ \csromangathabat{สพฺเพ สญฺจเย เนตฺวาน\csromansharedfootnote{5540deef31421000}{สญฺเชยฺยเก เนตฺวา (สี)},}{กํสุ ทานิ นยิสฺสตี”ติ.}}}

\prose{อสฺโสสุํ โข ภิกฺขู เตสํ มนุสฺสานํ อุชฺฌายนฺตานํ ขิยฺยนฺตานํ วิปาเจนฺตานํ. อถ โข เต ภิกฺขู ภควโต เอตมตฺถํ อาโรเจสุํ ฯเปฯ. น ภิกฺขเว โส สทฺโท จิรํ ภวิสฺสติ, สตฺตาหเมว ภวิสฺสติ, สตฺตาหสฺส อจฺจเยน อนฺตรธายิสฺสติ, เตน หิ ภิกฺขเว เย ตุเมฺห อิมาย คาถาย โจเทนฺติ\csromandash{}}

\setlength{\csromangathapagewidth}{0pt}
\setlength{\csromangathawakwidth}{0pt}
\csromangathameasuregroup{“อาคโต โข มหาสมโณ, \\ สพฺเพ สญฺจเย เนตฺวาน,}{\csromangathabat{“อาคโต โข มหาสมโณ,}{มาคธานํ คิริพฺพชํ.} \\ \csromangathabat{สพฺเพ สญฺจเย เนตฺวาน,}{กํสุ ทานิ นยิสฺสตี”ติ.}}
\csromangathasetleft{“อาคโต โข มหาสมโณ, \\ สพฺเพ สญฺจเย เนตฺวาน,}
\csromangathagroup{\csromangathastanza{\csromangathabat{“อาคโต โข มหาสมโณ,}{มาคธานํ คิริพฺพชํ.} \\ \csromangathabat{สพฺเพ สญฺจเย เนตฺวาน,}{กํสุ ทานิ นยิสฺสตี”ติ.}}}

\prose{เต ตุเมฺห อิมาย คาถาย ปฏิโจเทถ\csromandash{}}

\setlength{\csromangathapagewidth}{0pt}
\setlength{\csromangathawakwidth}{0pt}
\csromangathameasuregroup{“นยนฺติ เว มหาวีรา, \\ ธมฺเมน นยมานานํ\textsuperscript{1},}{\csromangathabat{“นยนฺติ เว มหาวีรา,}{สทฺธมฺเมน ตถาคตา.} \\ \csromangathabat{ธมฺเมน นยมานานํ\textsuperscript{1},}{กา อุสูยา\textsuperscript{1} วิชานตนฺ”ติ.}}
\csromangathasetleft{“นยนฺติ เว มหาวีรา, \\ ธมฺเมน นยมานานํ\textsuperscript{1},}
\csromangathagroup{\csromangathastanza{\csromangathabat{“นยนฺติ เว มหาวีรา,}{สทฺธมฺเมน ตถาคตา.} \\ \csromangathabat{ธมฺเมน นยมานานํ\csromansharedfootnote{5e0b7d73af1aa8d9}{นียมานานํ (ก)},}{กา อุสูยา\csromansharedfootnote{668d1a11c38705d4}{อุสฺสุยา (ก)} วิชานตนฺ”ติ.}}}

\prose{เตน โข ปน สมเยน มนุสฺสา ภิกฺขู ทิสฺวา อิมาย คาถาย โจเทนฺติ\csromandash{}}

\setlength{\csromangathapagewidth}{0pt}
\setlength{\csromangathawakwidth}{0pt}
\csromangathameasuregroup{“อาคโต โข มหาสมโณ, \\ สพฺเพ สญฺจเย เนตฺวาน,}{\csromangathabat{“อาคโต โข มหาสมโณ,}{มาคธานํ คิริพฺพชํ.} \\ \csromangathabat{สพฺเพ สญฺจเย เนตฺวาน,}{กํสุ ทานิ นยิสฺสตี”ติ.}}
\csromangathasetleft{“อาคโต โข มหาสมโณ, \\ สพฺเพ สญฺจเย เนตฺวาน,}
\csromangathagroup{\csromangathastanza{\csromangathabat{“อาคโต โข มหาสมโณ,}{มาคธานํ คิริพฺพชํ.} \\ \csromangathabat{สพฺเพ สญฺจเย เนตฺวาน,}{กํสุ ทานิ นยิสฺสตี”ติ.}}}

\prose{ภิกฺขู เต มนุสฺเส อิมาย คาถาย ปฏิโจเทนฺติ\csromandash{}}

\setlength{\csromangathapagewidth}{0pt}
\setlength{\csromangathawakwidth}{0pt}
\csromangathameasuregroup{“นยนฺติ เว มหาวีรา, \\ ธมฺเมน นยมานานํ,}{\csromangathabat{“นยนฺติ เว มหาวีรา,}{สทฺธมฺเมน ตถาคตา.} \\ \csromangathabat{ธมฺเมน นยมานานํ,}{กา อุสูยา วิชานตนฺ”ติ.}}
\csromangathasetleft{“นยนฺติ เว มหาวีรา, \\ ธมฺเมน นยมานานํ,}
\csromangathagroup{\csromangathastanza{\csromangathabat{“นยนฺติ เว มหาวีรา,}{สทฺธมฺเมน ตถาคตา.} \\ \csromangathabat{ธมฺเมน นยมานานํ,}{กา อุสูยา วิชานตนฺ”ติ.}}}

\prosewithcloser{มนุสฺสา “ธมฺเมน กิร สมณา สกฺยปุตฺติยา เนนฺติ, โน อธมฺเมนา”ติ สตฺตาหเมว โส สทฺโท อโหสิ, สตฺตาหสฺส อจฺจเยน อนฺตรธายิ.}{สาริปุตฺตโมคฺคลฺลาน\-ปพฺพชฺชา\-กถา นิฏฺฐิตา.}
\nitthitam{จตุตฺถ\-ภาณวาโร นิฏฺฐิโต.}
\csromanmatikaanchor{mtk.16}
\csromantocmarkref{section}{15. อุปชฺฌายวตฺตกถา}{mtk.16}
\bookmark[dest={mtk.16},level=1]{15. อุปชฺฌายวตฺตกถา}
\csromanheader{1}{15. \textbf{อุปชฺฌายวตฺต}\-\textbf{กถา}}
\proseitem{64}{\csromanfolio{56}เตน โข ปน สมเยน ภิกฺขู อนุปชฺฌายกา อนาจริยกา\csromansharedfootnote{0aaa8df640f92eb8}{อิทํ ปทํ สีสฺยาโปตฺถเกสุ นตฺถิ.} อโนวทิยมานา อนนุสาสิยมานา ทุนฺนิวตฺถา ทุปฺปารุตา อนากปฺป\-สมฺปนฺนา ปิณฺฑาย จรนฺติ, มนุสฺสานํ\csromansharedfootnote{5ef96dc687adce99}{เต มนุสฺสานํ (ก)} ภุญฺชมานานํ อุปริโภชเนปิ อุตฺติฏฺฐปตฺตํ อุปนาเมนฺติ, อุปริ\-ขาทนีเยปิ อุตฺติฏฺฐปตฺตํ อุปนาเมนฺติ, อุปริ\-สายนีเยปิ อุตฺติฏฺฐปตฺตํ อุปนาเมนฺติ, อุปริปานีเยปิ อุตฺติฏฺฐปตฺตํ อุปนาเมนฺติ, สามํ สูปมฺปิ โอทนมฺปิ วิญฺญาเปตฺวา ภุญฺชนฺติ, ภตฺตคฺเคปิ อุจฺจาสทฺทา มหาสทฺทา วิหรนฺติ, มนุสฺสา อุชฺฌายนฺติ ขิยฺยนฺติ วิปาเจนฺติ “กถํ หิ นาม สมณา สกฺยปุตฺติยา ทุนฺนิวตฺถา ทุปฺปารุตา อนากปฺป\-สมฺปนฺนา ปิณฺฑาย จริสฺสนฺติ, มนุสฺสานํ ภุญฺชมานานํ อุปริโภชเนปิ อุตฺติฏฺฐปตฺตํ อุปนาเมสฺสนฺติ, อุปริ\-ขาทนีเยปิ อุตฺติฏฺฐปตฺตํ อุปนาเมสฺสนฺติ, อุปริ\-สายนีเยปิ อุตฺติฏฺฐปตฺตํ อุปนาเมสฺสนฺติ, อุปริปานีเยปิ อุตฺติฏฺฐปตฺตํ อุปนาเมสฺสนฺติ, สามํ สูปมฺปิ โอทนมฺปิ วิญฺญาเปตฺวา ภุญฺชิสฺสนฺติ, ภตฺตคฺเคปิ อุจฺจาสทฺทา มหาสทฺทา วิหริสฺสนฺติ เสยฺยถาปิ พฺราหฺมณา พฺราหฺมณ\-โภชเน”ติ.}

\prose{อสฺโสสุํ โข ภิกฺขู เตสํ มนุสฺสานํ อุชฺฌายนฺตานํ ขิยฺยนฺตานํ วิปาเจนฺตานํ. เย เต ภิกฺขู อปฺปิจฺฉา สนฺตุฏฺฐา ลชฺชิโน กุกฺกุจฺจกา สิกฺขากามา, เต อุชฺฌายนฺติ ขิยฺยนฺติ วิปาเจนฺติ “กถํ หิ นาม ภิกฺขู ทุนฺนิวตฺถา ทุปฺปารุตา อนากปฺป\-สมฺปนฺนา ปิณฺฑาย จริสฺสนฺติ, มนุสฺสานํ ภุญฺชมานานํ อุปริโภชเนปิ อุตฺติฏฺฐปตฺตํ อุปนาเมสฺสนฺติ, อุปริ\-ขาทนีเยปิ อุตฺติฏฺฐปตฺตํ อุปนาเมสฺสนฺติ, อุปริ\-สายนีเยปิ อุตฺติฏฺฐปตฺตํ อุปนาเมสฺสนฺติ, อุปริปานีเยปิ อุตฺติฏฺฐปตฺตํ อุปนาเมสฺสนฺติ, สามํ สูปมฺปิ โอทนมฺปิ วิญฺญาเปตฺวา ภุญฺชิสฺสนฺติ, ภตฺตคฺเคปิ อุจฺจาสทฺทา มหาสทฺทา วิหริสฺสนฺตี”ติ. อถ โข เต ภิกฺขู ฯเปฯ ภควโต เอตมตฺถํ อาโรเจสุํ.}

\proseitem{64}{อถ โข ภควา เอตสฺมิํ นิทาเน เอตสฺมิํ ปกรเณ ภิกฺขุสํฆํ สนฺนิปาตาเปตฺวา ภิกฺขู ปฏิปุจฺฉิ “สจฺจํ กิร ภิกฺขเว ภิกฺขู ทุนฺนิวตฺถา ทุปฺปารุตา อนากปฺป\-สมฺปนฺนา ปิณฺฑาย จรนฺติ, มนุสฺสานํ ภุญฺชมานานํ อุปริโภชเนปิ อุตฺติฏฺฐปตฺตํ อุปนาเมนฺติ, อุปริ\-ขาทนีเยปิ อุตฺติฏฺฐปตฺตํ อุปนาเมนฺติ, อุปริ\-สายนีเยปิ อุตฺติฏฺฐปตฺตํ อุปนาเมนฺติ, อุปริปานีเยปิ \csromanfolio{57}อุตฺติฏฺฐปตฺตํ อุปนาเมนฺติ, สามํ สูปมฺปิ โอทนมฺปิ วิญฺญาเปตฺวา ภุญฺชนฺติ, ภตฺตคฺเคปิ อุจฺจาสทฺทา มหาสทฺทา วิหรนฺตี”ติ. สจฺจํ ภควาติ. วิครหิ พุทฺโธ ภควา “อนนุจฺฉวิกํ ภิกฺขเว เตสํ โมฆปุริสานํ อนนุโลมิกํ อปฺปติรูปํ อสฺสามณกํ อกปฺปิยํ อกรณียํ, กถํ หิ นาม เต ภิกฺขเว โมฆปุริสา ทุนฺนิวตฺถา ทุปฺปารุตา อนากปฺป\-สมฺปนฺนา ปิณฺฑาย จริสฺสนฺติ, มนุสฺสานํ ภุญฺชมานานํ อุปริโภชเนปิ อุตฺติฏฺฐปตฺตํ อุปนาเมสฺสนฺติ, อุปริ\-ขาทนีเยปิ อุตฺติฏฺฐปตฺตํ อุปนาเมสฺสนฺติ, อุปริ\-สายนีเยปิ อุตฺติฏฺฐปตฺตํ อุปนาเมสฺสนฺติ, อุปริปานีเยปิ อุตฺติฏฺฐปตฺตํ อุปนาเมสฺสนฺติ, สามํ สูปมฺปิ โอทนมฺปิ วิญฺญาเปตฺวา ภุญฺชิสฺสนฺติ, ภตฺตคฺเคปิ อุจฺจาสทฺทา มหาสทฺทา วิหริสฺสนฺติ. เนตํ ภิกฺขเว อปฺปสนฺนานํ วา ปสาทาย ปสนฺนานํ วา ภิยฺโยภาวาย, อถ เขฺวตํ ภิกฺขเว อปฺปสนฺนาน\-ญฺเจว อปฺปสาทาย ปสนฺนานญฺจ เอกจฺจานํ อญฺญถตฺตายา”ติ. อถ โข ภควา เต ภิกฺขู อเนก\-ปริยาเยน วิครหิตฺวา ทุพฺภรตาย ทุปฺโปสตาย มหิจฺฉตาย อสนฺตุฏฺฐิตาย\csromansharedfootnote{7dfaecd6a678338b}{อสนฺตุฏฺฐิยา (สี), อสนฺตุฏฺฐตาย (สฺยา)} สงฺคณิกาย โกสชฺชสฺส อวณฺณํ ภาสิตฺวา อเนก\-ปริยาเยน สุภรตาย สุโปสตาย อปฺปิจฺฉสฺส สนฺตุฏฺฐสฺส สลฺเลขสฺส ธุตสฺส ปาสาทิกสฺส อปจยสฺส วีริยา\-รมฺภสฺส\csromansharedfootnote{9267ffdd082a3e79}{วิริยารมฺภสฺส (สี, สฺยา)} วณฺณํ ภาสิตฺวา ภิกฺขูนํ ตทนุจฺฉวิกํ ตทนุโลมิกํ ธมฺมิํ กถํ กตฺวา ภิกฺขู อามนฺเตสิ\csromandash{}}

\proseitem{65}{อนุชานามิ ภิกฺขเว อุปชฺฌายํ, อุปชฺฌาโย ภิกฺขเว สทฺธิ\-วิหาริกมฺหิ ปุตฺตจิตฺตํ อุปฏฺฐเปสฺสติ, สทฺธิวิหาริโก อุปชฺฌายมฺหิ ปิตุจิตฺตํ อุปฏฺฐเปสฺสติ, เอวํ เต อญฺญมญฺญํ สคารวา สปฺปติสฺสา สภาควุตฺติโน วิหรนฺตา อิมสฺมิํ ธมฺมวินเย วุฑฺฒิํ วิรุฬฺหิํ เวปุลฺลํ อาปชฺชิสฺสนฺติ. เอวญฺจ ปน ภิกฺขเว อุปชฺฌาโย คเหตพฺโพ, เอกํสํ อุตฺตราสงฺคํ กริตฺวา ปาเท วนฺทิตฺวา อุกฺกุฏิกํ นิสีทิตฺวา อญฺชลิํ ปคฺคเหตฺวา เอวมสฺส วจนีโย “อุปชฺฌาโย เม ภนฺเต โหหิ, อุปชฺฌาโย เม ภนฺเต โหหิ, อุปชฺฌาโย เม ภนฺเต โหหี”ติ. “สาหู”ติ วา “ลหู”ติ วา “โอปายิกนฺ”ติ วา “ปติรูปนฺ”ติ วา “ปาสาทิเกน สมฺปาเทหี”ติ วา กาเยน วิญฺญาเปติ, วาจาย วิญฺญาเปติ, กาเยน วาจาย วิญฺญาเปติ, คหิโต โหติ อุปชฺฌาโย. น กาเยน วิญฺญาเปติ, น วาจาย \csromanfolio{58}วิญฺญาเปติ, น กาเยน วาจาย\csromansharedfootnote{52a7bdd92b9bdedf}{น วาจาย (ก)} วิญฺญาเปติ, น คหิโต โหติ อุปชฺฌาโย.}

\proseitem{66}{\csromansymbolfootnote{*}{วิ 4. 393 ปิฏฺฐาทีสุปิ.}สทฺธิ\-วิหาริเกน ภิกฺขเว อุปชฺฌายมฺหิ สมฺมา วตฺติตพฺพํ. ตตฺรายํ สมฺมาวตฺตนา\csromandash{}}

\prose{กาลสฺเสว วุฏฺฐาย อุปาหนา โอมุญฺจิตฺวา เอกํสํ อุตฺตราสงฺคํ กริตฺวา ทนฺตกฏฺฐํ ทาตพฺพํ, มุโขทกํ ทาตพฺพํ, อาสนํ ปญฺญเปตพฺพํ.}

\prose{สเจ ยาคุ โหติ, ภาชนํ โธวิตฺวา ยาคุ อุปนาเมตพฺพา, ยาคุํ ปีตสฺส อุทกํ ทตฺวา ภาชนํ ปฏิคฺคเหตฺวา นีจํ กตฺวา สาธุกํ อปฺปฏิฆํสนฺเตน โธวิตฺวา ปฏิสาเมตพฺพํ, อุปชฺฌายมฺหิ วุฏฺฐิเต อาสนํ อุทฺธริตพฺพํ.}

\setlength{\csromangathapagewidth}{0pt}
\setlength{\csromangathawakwidth}{0pt}
\csromangathameasuregroup{สเจ โส เทโส อุกฺลาโป โหติ,}{\csromangathabat{สเจ โส เทโส อุกฺลาโป โหติ,}{โส เทโส สมฺมชฺชิตพฺโพ.}}
\csromangathasetleft{สเจ โส เทโส อุกฺลาโป โหติ,}
\csromangathagroup{\csromangathastanza{\csromangathaitembat{66}{สเจ โส เทโส อุกฺลาโป โหติ,}{โส เทโส สมฺมชฺชิตพฺโพ.}}}

\prose{สเจ อุปชฺฌาโย คามํ ปวิสิตุกาโม โหติ, นิวาสนํ ทาตพฺพํ ปฏินิวาสนํ ปฏิคฺคเหตพฺพํ, กายพนฺธนํ ทาตพฺพํ, สคุณํ กตฺวา สํฆาฏิโย ทาตพฺพา, โธวิตฺวา ปตฺโต โสทโก\csromansharedfootnote{6a149cae96967c2d}{สอุทโก (ก)} ทาตพฺโพ.}

\prose{สเจ อุปชฺฌาโย ปจฺฉาสมณํ อากงฺขติ, ติมณฺฑลํ ปฏิจฺฉาเทนฺเตน ปริมณฺฑลํ นิวาเสตฺวา กายพนฺธนํ พนฺธิตฺวา สคุณํ กตฺวา สํฆาฏิโย ปารุปิตฺวา คณฺฐิกํ ปฏิมุญฺจิตฺวา โธวิตฺวา ปตฺตํ คเหตฺวา อุปชฺฌายสฺส ปจฺฉาสมเณน โหตพฺพํ, นาติทูเร คนฺตพฺพํ, นาจฺจาสนฺเน คนฺตพฺพํ, ปตฺต\-ปริยาปนฺนํ ปฏิคฺคเหตพฺพํ, น อุปชฺฌายสฺส ภณมานสฺส อนฺตรนฺตรา กถา โอปาเตตพฺพา, อุปชฺฌาโย อาปตฺติสามนฺตา ภณมาโน นิวาเรตพฺโพ.}

\prose{นิวตฺตนฺเตน ปฐมตรํ อาคนฺตฺวา อาสนํ ปญฺญเปตพฺพํ, ปาโททกํ ปาทปีฐํ ปาทกถลิกํ อุปนิกฺขิปิตพฺพํ, ปจฺจุคฺคนฺตฺวา ปตฺตจีวรํ ปฏิคฺคเหตพฺพํ, ปฏินิวาสนํ ทาตพฺพํ, นิวาสนํ ปฏิคฺคเหตพฺพํ.}

\prose{สเจ จีวรํ สินฺนํ โหติ, มุหุตฺตํ อุเณฺห โอตาเปตพฺพํ, น จ อุเณฺห จีวรํ นิทหิตพฺพํ, จีวรํ สงฺฆริตพฺพํ, จีวรํ สงฺฆรนฺเตน จตุรงฺคุลํ กณฺณํ อุสฺสาเรตฺวา จีวรํ สงฺฆริตพฺพํ “มา มชฺเฌ ภงฺโค อโหสี”ติ, โอโภเค กายพนฺธนํ กาตพฺพํ.}

\proseitem{66}{สเจ ปิณฺฑปาโต โหติ อุปชฺฌาโย จ ภุญฺชิตุกาโม โหติ, อุทกํ ทตฺวา ปิณฺฑปาโต อุปนาเมตพฺโพ, อุปชฺฌาโย ปานีเยน ปุจฺฉิตพฺโพ, ภุตฺตาวิสฺส อุทกํ ทตฺวา ปตฺตํ ปฏิคฺคเหตฺวา นีจํ กตฺวา \csromanfolio{59}สาธุกํ อปฺปฏิฆํสนฺเตน โธวิตฺวา โวทกํ กตฺวา มุหุตฺตํ อุเณฺห โอตาเปตพฺโพ, น จ อุเณฺห ปตฺโต นิทหิตพฺโพ. ปตฺตจีวรํ นิกฺขิปิตพฺพํ, ปตฺตํ นิกฺขิปนฺเตน เอเกน หตฺเถน ปตฺตํ คเหตฺวา เอเกน หตฺเถน เหฏฺฐามญฺจํ วา เหฏฺฐาปีฐํ วา ปรามสิตฺวา ปตฺโต นิกฺขิปิตพฺโพ, น จ อนนฺตรหิตาย ภูมิยา ปตฺโต นิกฺขิปิตพฺโพ. จีวรํ นิกฺขิปนฺเตน เอเกน หตฺเถน จีวรํ คเหตฺวา เอเกน หตฺเถน จีวรวํสํ วา จีวรรชฺชุํ วา ปมชฺชิตฺวา ปารโต อนฺตํ โอรโต โภคํ กตฺวา จีวรํ นิกฺขิปิตพฺพํ. อุปชฺฌายมฺหิ วุฏฺฐิเต อาสนํ อุทฺธริตพฺพํ, ปาโททกํ ปาทปีฐํ ปาทกถลิกํ ปฏิสาเมตพฺพํ. สเจ โส เทโส อุกฺลาโป โหติ, โส เทโส สมฺมชฺชิตพฺโพ.}

\setlength{\csromangathapagewidth}{0pt}
\setlength{\csromangathawakwidth}{0pt}
\csromangathameasuregroup{สเจ อุปชฺฌาโย นหายิตุกาโม โหติ, \\ สเจ สีเตน อตฺโถ โหติ, \\ สเจ อุเณฺหน อตฺโถ โหติ,}{\csromangathabat{สเจ อุปชฺฌาโย นหายิตุกาโม โหติ,}{นหานํ ปฏิยาเทตพฺพํ.} \\ \csromangathabat{สเจ สีเตน อตฺโถ โหติ,}{สีตํ ปฏิยาเทตพฺพํ.} \\ \csromangathabat{สเจ อุเณฺหน อตฺโถ โหติ,}{อุณฺหํ ปฏิยาเทตพฺพํ.}}
\csromangathasetleft{สเจ อุปชฺฌาโย นหายิตุกาโม โหติ, \\ สเจ สีเตน อตฺโถ โหติ, \\ สเจ อุเณฺหน อตฺโถ โหติ,}
\csromangathagroup{\csromangathastanza{\csromangathaitembat{66}{สเจ อุปชฺฌาโย นหายิตุกาโม โหติ,}{นหานํ ปฏิยาเทตพฺพํ.} \\ \csromangathacontbat{สเจ สีเตน อตฺโถ โหติ,}{สีตํ ปฏิยาเทตพฺพํ.} \\ \csromangathacontbat{สเจ อุเณฺหน อตฺโถ โหติ,}{อุณฺหํ ปฏิยาเทตพฺพํ.}}}

\prose{สเจ อุปชฺฌาโย ชนฺตาฆรํ ปวิสิตุกาโม โหติ, จุณฺณํ สนฺเนตพฺพํ, มตฺติกา เตเมตพฺพา, ชนฺตาฆรปีฐํ อาทาย อุปชฺฌายสฺส ปิฏฺฐิโต ปิฏฺฐิโต คนฺตฺวา ชนฺตาฆรปีฐํ ทตฺวา จีวรํ ปฏิคฺคเหตฺวา เอกมนฺตํ นิกฺขิปิตพฺพํ, จุณฺณํ ทาตพฺพํ, มตฺติกา ทาตพฺพา.}

\prose{สเจ อุสฺสหติ, ชนฺตาฆรํ ปวิสิตพฺพํ, ชนฺตาฆรํ ปวิสนฺเตน มตฺติกาย มุขํ มกฺเขตฺวา ปุรโต จ ปจฺฉโต จ ปฏิจฺฉาเทตฺวา ชนฺตาฆรํ ปวิสิตพฺพํ, น เถเร ภิกฺขู อนุปขชฺช นิสีทิตพฺพํ, น นวา ภิกฺขู อาสเนน ปฏิพาหิตพฺพา, ชนฺตาฆเร อุปชฺฌายสฺส ปริกมฺมํ กาตพฺพํ, ชนฺตาฆรา นิกฺขมนฺเตน ชนฺตาฆรปีฐํ อาทาย ปุรโต จ ปจฺฉโต จ ปฏิจฺฉาเทตฺวา ชนฺตาฆรา นิกฺขมิตพฺพํ.}

\prose{อุทเกปิ อุปชฺฌายสฺส ปริกมฺมํ กาตพฺพํ, นหาเตน ปฐมตรํ อุตฺตริตฺวา อตฺตโน คตฺตํ โวทกํ กตฺวา นิวาเสตฺวา อุปชฺฌายสฺส คตฺตโต อุทกํ ปมชฺชิตพฺพํ, นิวาสนํ ทาตพฺพํ, สํฆาฏิ ทาตพฺพา, ชนฺตาฆรปีฐํ อาทาย ปฐมตรํ อาคนฺตฺวา อาสนํ ปญฺญเปตพฺพํ, ปาโททกํ ปาทปีฐํ ปาทกถลิกํ อุปนิกฺขิปิตพฺพํ, อุปชฺฌาโย ปานีเยน ปุจฺฉิตพฺโพ.}

\prose{สเจ อุทฺทิสาเปตุกาโม โหติ อุทฺทิสิตพฺโพ, สเจ ปริปุจฺฉิตุ\-กาโม โหติ ปริปุจฺฉิตพฺโพ.}

\proseitem{66}{ยสฺมิํ วิหาเร อุปชฺฌาโย วิหรติ สเจ โส วิหาโร อุกฺลาโป โหติ สเจ อุสฺสหติ, โสเธตพฺโพ. วิหารํ โสเมนฺเตน ปฐมํ \csromanfolio{60}ปตฺตจีวรํ นีหริตฺวา เอกมนฺตํ นิกฺขิปิตพฺพํ, นิสีทน\-ปจฺจตฺถรณํ นีหริตฺวา เอกมนฺตํ นิกฺขิปิตพฺพํ, ภิสิพิพฺโพหนํ\csromansharedfootnote{33b95556041a3969}{ภิสิพิมฺโพหนํ (สี, สฺยา)} นีหริตฺวา เอกมนฺตํ นิกฺขิปิตพฺพํ, มญฺโจ นีจํ กตฺวา สาธุกํ อปฺปฏิฆํสนฺเตน อสํฆฏฺเฏนฺเตน กวาฏปิฏฺฐํ นีหริตฺวา เอกมนฺตํ นิกฺขิปิตพฺโพ, ปีฐํ นีจํ กตฺวา สาธุกํ อปฺปฏิฆํสนฺเตน อสํฆฏฺเฏนฺเตน กวาฏปิฏฺฐํ นีหริตฺวา เอกมนฺตํ นิกฺขิปิตพฺพํ, มญฺจ\-ปฏิปาทกา นีหริตฺวา เอกมนฺตํ นิกฺขิปิตพฺพา. เขฬมลฺลโก นีหริตฺวา เอกมนฺตํ นิกฺขิปิตพฺโพ. อปสฺเสนผลกํ นีหริตฺวา เอกมนฺตํ นิกฺขิปิตพฺพํ. ภูมตฺถรณํ ยถาปญฺญตฺตํ สลฺลกฺเขตฺวา นีหริตฺวา เอกมนฺตํ นิกฺขิปิตพฺพํ. สเจ วิหาเร สนฺตานกํ โหติ, อุลฺโลกา ปฐมํ โอหาเรตพฺพํ. อาโลกสนฺธิ\-กณฺณภาคา ปมชฺชิตพฺพา. สเจ เครุก\-ปริกมฺม\-กตา ภิตฺติ กณฺณกิตา โหติ, โจฬกํ เตเมตฺวา ปีเฬตฺวา ปมชฺชิตพฺพา. สเจ กาฬวณฺณกตา ภูมิ กณฺณกิตา โหติ, โจฬกํ เตเมตฺวา ปีเฬตฺวา ปมชฺชิตพฺพา. สเจ อกตา โหติ ภูมิ, อุทเกน ปริปฺโผสิตฺวา สมฺมชฺชิตพฺพา “มา วิหาโร รเชน อุหญฺญี”ติ. สงฺการํ วิจินิตฺวา เอกมนฺตํ ฉฑฺเฑตพฺพํ.}

\prose{ภูมตฺถรณํ โอตาเปตฺวา โสเธตฺวา ปปฺโผเฏตฺวา อติหริตฺวา ยถาปญฺญตฺตํ ปญฺญเปตพฺพํ. มญฺจ\-ปฏิปาทกา โอตาเปตฺวา ปมชฺชิตฺวา อติหริตฺวา ยถาฐาเน ฐเปตพฺพา, มญฺโจ โอตาเปตฺวา โสเธตฺวา ปปฺโผเฏตฺวา นีจํ กตฺวา สาธุกํ อปฺปฏิฆํสนฺเตน อสํฆฏฺเฏนฺเตน กวาฏปิฏฺฐํ อติหริตฺวา ยถาปญฺญตฺตํ ปญฺญเปตพฺโพ. ปีฐํ โอตาเปตฺวา โสเธตฺวา ปปฺโผเฏตฺวา นีจํ กตฺวา สาธุกํ อปฺปฏิฆํสนฺเตน อสํฆฏฺเฏนฺเตน กวาฏปิฏฺฐํ อติหริตฺวา ยถาปญฺญตฺตํ ปญฺญเปตพฺพํ. ภิสิพิพฺโพหนํ โอตาเปตฺวา โสเธตฺวา ปปฺโผเฏตฺวา อติหริตฺวา ยถาปญฺญตฺตํ ปญฺญเปตพฺพํ. นิสีทน\-ปจฺจตฺถรณํ โอตาเปตฺวา โสเธตฺวา ปปฺโผเฏตฺวา อติหริตฺวา ยถาปญฺญตฺตํ ปญฺญเปตพฺพํ. เขฬมลฺลโก โอตาเปตฺวา ปมชฺชิตฺวา อติหริตฺวา ยถาฐาเน ฐเปตพฺโพ. อปสฺเสนผลกํ โอตาเปตฺวา ปมชฺชิตฺวา อติหริตฺวา ยถาฐาเน ฐเปตพฺพํ. ปตฺตจีวรํ นิกฺขิปิตพฺพํ. ปตฺตํ นิกฺขิปนฺเตน เอเกน หตฺเถน ปตฺตํ คเหตฺวา เอเกน หตฺเถน เหฏฺฐามญฺจํ วา เหฏฺฐาปีฐํ วา ปรามสิตฺวา ปตฺโต นิกฺขิปิตพฺโพ, น จ อนนฺตรหิตาย ภูมิยา ปตฺโต นิกฺขิปิตพฺโพ. จีวรํ นิกฺขิปนฺเตน เอเกน หตฺเถน จีวรํ \csromanfolio{61}คเหตฺวา เอเกน หตฺเถน จีวรวํสํ วา จีวรรชฺชุํ วา ปมชฺชิตฺวา ปารโต อนฺตํ โอรโต โภคํ กตฺวา จีวรํ นิกฺขิปิตพฺพํ.}

\setlength{\csromangathapagewidth}{0pt}
\setlength{\csromangathawakwidth}{0pt}
\csromangathameasuregroup{สเจ ปริเวณํ อุกฺลาปํ โหติ, \\ สเจ โกฏฺฐโก อุกฺลาโป โหติ, \\ สเจ อคฺคิสาลา อุกฺลาปา โหติ, \\ สเจ วจฺจกุฏิ อุกฺลาปา โหติ, \\ สเจ ปานียํ น โหติ, \\ สเจ ปริโภชนียํ น โหติ,}{สเจ ปุรตฺถิมา สรชา วาตา วายนฺติ, \\ ปุรตฺถิมา วาตปานา ถเกตพฺพา. \\ สเจ ปจฺฉิมา สรชา วาตา วายนฺติ, \\ ปจฺฉิมา วาตปานา ถเกตพฺพา. \\ สเจ อุตฺตรา สรชา วาตา วายนฺติ, \\ อุตฺตรา วาตปานา ถเกตพฺพา. \\ สเจ ทกฺขิณา สรชา วาตา วายนฺติ, \\ ทกฺขิณา วาตปานา ถเกตพฺพา. \\ สเจ สีตกาโล โหติ, \\ ทิวา วาตปานา วิวริตพฺพา, รตฺติํ ถเกตพฺพา. \\ สเจ อุณฺหกาโล โหติ, \\ ทิวา วาตปานา ถเกตพฺพา, รตฺติํ วิวริตพฺพา. \\ \csromangathabat{สเจ ปริเวณํ อุกฺลาปํ โหติ,}{ปริเวณํ สมฺมชฺชิตพฺพํ.} \\ \csromangathabat{สเจ โกฏฺฐโก อุกฺลาโป โหติ,}{โกฏฺฐโก สมฺมชฺชิตพฺโพ.} \\ สเจ อุปฏฺฐานสาลา อุกฺลาปา โหติ, \\ อุปฏฺฐานสาลา สมฺมชฺชิตพฺพา. \\ \csromangathabat{สเจ อคฺคิสาลา อุกฺลาปา โหติ,}{อคฺคิสาลา สมฺมชฺชิตพฺพา.} \\ \csromangathabat{สเจ วจฺจกุฏิ อุกฺลาปา โหติ,}{วจฺจกุฏิ สมฺมชฺชิตพฺพา.} \\ \csromangathabat{สเจ ปานียํ น โหติ,}{ปานียํ อุปฏฺฐาเปตพฺพํ.} \\ \csromangathabat{สเจ ปริโภชนียํ น โหติ,}{ปริโภชนียํ อุปฏฺฐาเปตพฺพํ.}}
\csromangathasetleft{สเจ ปริเวณํ อุกฺลาปํ โหติ, \\ สเจ โกฏฺฐโก อุกฺลาโป โหติ, \\ สเจ อคฺคิสาลา อุกฺลาปา โหติ, \\ สเจ วจฺจกุฏิ อุกฺลาปา โหติ, \\ สเจ ปานียํ น โหติ, \\ สเจ ปริโภชนียํ น โหติ,}
\csromangathagroup{\csromangathastanza{\csromangathaitemline{66}{สเจ ปุรตฺถิมา สรชา วาตา วายนฺติ,} \\ \csromangathacontline{ปุรตฺถิมา วาตปานา ถเกตพฺพา.} \\ \csromangathacontline{สเจ ปจฺฉิมา สรชา วาตา วายนฺติ,} \\ \csromangathacontline{ปจฺฉิมา วาตปานา ถเกตพฺพา.}}\csromangathastanza{\csromangathaitemline{66}{สเจ อุตฺตรา สรชา วาตา วายนฺติ,} \\ \csromangathacontline{อุตฺตรา วาตปานา ถเกตพฺพา.} \\ \csromangathacontline{สเจ ทกฺขิณา สรชา วาตา วายนฺติ,} \\ \csromangathacontline{ทกฺขิณา วาตปานา ถเกตพฺพา.}}\csromangathastanza{\csromangathaitemline{66}{สเจ สีตกาโล โหติ,} \\ \csromangathacontline{ทิวา วาตปานา วิวริตพฺพา, รตฺติํ ถเกตพฺพา.} \\ \csromangathacontline{สเจ อุณฺหกาโล โหติ,} \\ \csromangathacontline{ทิวา วาตปานา ถเกตพฺพา, รตฺติํ วิวริตพฺพา.}}\csromangathastanza{\csromangathabat{สเจ ปริเวณํ อุกฺลาปํ โหติ,}{ปริเวณํ สมฺมชฺชิตพฺพํ.} \\ \csromangathabat{สเจ โกฏฺฐโก อุกฺลาโป โหติ,}{โกฏฺฐโก สมฺมชฺชิตพฺโพ.}}\csromangathastanza{สเจ อุปฏฺฐานสาลา อุกฺลาปา โหติ, \\ อุปฏฺฐานสาลา สมฺมชฺชิตพฺพา. \\ \csromangathabat{สเจ อคฺคิสาลา อุกฺลาปา โหติ,}{อคฺคิสาลา สมฺมชฺชิตพฺพา.}}\csromangathastanza{\csromangathabat{สเจ วจฺจกุฏิ อุกฺลาปา โหติ,}{วจฺจกุฏิ สมฺมชฺชิตพฺพา.} \\ \csromangathabat{สเจ ปานียํ น โหติ,}{ปานียํ อุปฏฺฐาเป\-ตพฺพํ.} \\ \csromangathabat{สเจ ปริโภชนียํ น โหติ,}{ปริโภชนียํ อุปฏฺฐาเป\-ตพฺพํ.}}}

\prose{สเจ อาจมนกุมฺภิยา อุทกํ น โหติ, อาจมนกุมฺภิยา อุทกํ อาสิญฺจิตพฺพํ.}

\proseitem{66}{สเจ อุปชฺฌายสฺส อนภิรติ อุปฺปนฺนา โหติ, สทฺธิ\-วิหาริเกน วูปกาเสตพฺโพ วูปกาสาเป\-ตพฺโพ ธมฺมกถา วาสฺส กาตพฺพา. สเจ อุปชฺฌายสฺส กุกฺกุจฺจํ อุปฺปนฺนํ โหติ, สทฺธิ\-วิหาริเกน วิโนเทตพฺพํ วิโนทาเปตพฺพํ ธมฺมกถา วาสฺส กาตพฺพา. สเจ อุปชฺฌายสฺส ทิฏฺฐิคตํ อุปฺปนฺนํ โหติ, สทฺธิ\-วิหาริเกน วิเวเจตพฺพํ วิเวจาเปตพฺพํ ธมฺมกถา วาสฺส กาตพฺพา. สเจ อุปชฺฌาโย ครุธมฺมํ อชฺฌาปนฺโน โหติ ปริวาสารโห, สทฺธิ\-วิหาริเกน อุสฺสุกฺกํ กาตพฺพํ “กินฺติ นุ โข สํโฆ อุปชฺฌายสฺส ปริวาสํ ทเทยฺยา”ติ. สเจ อุปชฺฌาโย มูลาย ปฏิกสฺสนา\-รโห โหติ, สทฺธิ\-วิหาริเกน อุสฺสุกฺกํ กาตพฺพํ “กินฺติ นุ โข สํโฆ อุปชฺฌายํ มูลาย ปฏิกสฺเสยฺยา”ติ. สเจ อุปชฺฌาโย มานตฺตารโห โหติ, สทฺธิ\-วิหาริเกน อุสฺสุกฺกํ กาตพฺพํ “กินฺติ นุ โข สํโฆ อุปชฺฌายสฺส มานตฺตํ ทเทยฺยา”ติ. สเจ อุปชฺฌาโย อพฺภานารโห โหติ, สทฺธิ\-วิหาริเกน อุสฺสุกฺกํ กาตพฺพํ “กินฺติ นุ โข สํโฆ อุปชฺฌายํ อพฺเภยฺยา”ติ. สเจ สํโฆ \csromanfolio{62}อุปชฺฌายสฺส กมฺมํ กตฺตุกาโม โหติ ตชฺชนียํ วา นิยสฺสํ\csromansharedfootnote{52f2271a5c610e49}{นิยสํ (ก)} วา ปพฺพาชนียํ วา ปฏิสารณียํ วา อุกฺเขปนียํ วา, สทฺธิ\-วิหาริเกน อุสฺสุกฺกํ กาตพฺพํ “กินฺติ นุ โข สํโฆ อุปชฺฌายสฺส กมฺมํ น กเรยฺย ลหุกาย วา ปริณาเมยฺยา”ติ. กตํ วา ปนสฺส โหติ สํเฆน กมฺมํ ตชฺชนียํ วา นิยสฺสํ วา ปพฺพาชนียํ วา ปฏิสารณียํ วา อุกฺเขปนียํ วา, สทฺธิ\-วิหาริเกน อุสฺสุกฺกํ กาตพฺพํ “กินฺติ นุ โข อุปชฺฌาโย สมฺมา วตฺเตยฺย, โลมํ ปาเตยฺย, เนตฺถารํ วตฺเตยฺย, สํโฆ ตํ กมฺมํ ปฏิปฺปสฺสมฺเภยฺยา”ติ.}

\prose{สเจ อุปชฺฌายสฺส จีวรํ โธวิตพฺพํ โหติ, สทฺธิ\-วิหาริเกน โธวิตพฺพํ, อุสฺสุกฺกํ วา กาตพฺพํ “กินฺติ นุ โข อุปชฺฌายสฺส จีวรํ โธวิเยถา”ติ. สเจ อุปชฺฌายสฺส จีวรํ กาตพฺพํ โหติ, สทฺธิ\-วิหาริเกน กาตพฺพํ, อุสฺสุกฺกํ วา กาตพฺพํ “กินฺติ นุ โข อุปชฺฌายสฺส จีวรํ กริเยถา”ติ. สเจ อุปชฺฌายสฺส รชนํ ปจิตพฺพํ โหติ, สทฺธิ\-วิหาริเกน ปจิตพฺพํ, อุสฺสุกฺกํ วา กาตพฺพํ “กินฺติ นุ โข อุปชฺฌายสฺส รชนํ ปจิเยถา”ติ. สเจ อุปชฺฌายสฺส จีวรํ รชิตพฺพํ\csromansharedfootnote{796f8e6e787c55dd}{รเชตพฺพํ (สี, สฺยา)} โหติ, สทฺธิ\-วิหาริเกน รชิตพฺพํ, อุสฺสุกฺกํ วา กาตพฺพํ “กินฺติ นุ โข อุปชฺฌายสฺส จีวรํ รชิเยถา”ติ. จีวรํ รชนฺเตน\csromansharedfootnote{c683858abbe9a749}{รเชนฺเตน (สี, สฺยา)} สาธุกํ สมฺปริวตฺตกํ สมฺปริวตฺตกํ รชิตพฺพํ, น จ อจฺฉินฺเน เถเว ปกฺกมิตพฺพํ.}

\prosewithcloserruled{น อุปชฺฌายํ อนาปุจฺฉา เอกจฺจสฺส ปตฺโต ทาตพฺโพ, น เอกจฺจสฺส ปตฺโต ปฏิคฺคเหตพฺโพ, น เอกจฺจสฺส จีวรํ ทาตพฺพํ, น เอกจฺจสฺส จีวรํ ปฏิคฺคเหตพฺพํ, น เอกจฺจสฺส ปริกฺขาโร ทาตพฺโพ, น เอกจฺจสฺส ปริกฺขาโร ปฏิคฺคเหตพฺโพ, น เอกจฺจสฺส เกสา เฉเทตพฺพา\csromansharedfootnote{4c72805cf12f1b7a}{เฉตฺตพฺพา (สี), เฉทิตพฺพา (ก)}, น เอกจฺเจน เกสา เฉทาเปตพฺพา, น เอกจฺจสฺส ปริกมฺมํ กาตพฺพํ, น เอกจฺเจน ปริกมฺมํ การาเปตพฺพํ, น เอกจฺจสฺส เวยฺยาวจฺโจ\csromansharedfootnote{80fa427d5ad5abf1}{เวยฺยาวจฺจํ (กตฺถจิ)} กาตพฺโพ, น เอกจฺเจน เวยฺยาวจฺโจ การาเปตพฺโพ, น เอกจฺจสฺส ปจฺฉาสมเณน โหตพฺพํ, น เอกจฺโจ ปจฺฉาสมโณ อาทาตพฺโพ, น เอกจฺจสฺส ปิณฺฑปาโต นีหริตพฺโพ, น เอกจฺเจน ปิณฺฑปาโต นีหราเปตพฺโพ, น อุปชฺฌายํ \csromanfolio{63}อนาปุจฺฉา คาโม ปวิสิตพฺโพ, น สุสานํ คนฺตพฺพํ, น ทิสา ปกฺกมิตพฺพา. สเจ อุปชฺฌาโย คิลาโน โหติ, ยาวชีวํ อุปฏฺฐาตพฺโพ, วุฏฺฐานมสฺส อาคเมตพฺพนฺติ.}{อุปชฺฌาย\-วตฺตํ นิฏฺฐิตํ.}
\csromanmatikaanchor{mtk.17}
\csromantocmarkref{section}{16. สทฺธิวิหาริกวตฺตกถา}{mtk.17}
\bookmark[dest={mtk.17},level=1]{16. สทฺธิวิหาริกวตฺตกถา}
\csromanheader{1}{16. \textbf{สทฺธิวิหาริกวตฺต}\-\textbf{กถา}}
\proseitem{67}{\csromansymbolfootnote{*}{วิ 4. 398 ปิฏฺฐาทีสุปิ.}อุปชฺฌาเยน ภิกฺขเว สทฺธิ\-วิหาริกมฺหิ สมฺมา วตฺติตพฺพํ, ตตฺรายํ สมฺมาวตฺตนา\csromandash{}}

\prose{อุปชฺฌาเยน ภิกฺขเว สทฺธิวิหาริโก สงฺคเหตพฺโพ อนุคฺคเหตพฺโพ อุทฺเทเสน ปริปุจฺฉาย โอวาเทน อนุสาสนิยา. สเจ อุปชฺฌายสฺส ปตฺโต โหติ สทฺธิ\-วิหาริกสฺส ปตฺโต น โหติ, อุปชฺฌาเยน สทฺธิ\-วิหาริกสฺส ปตฺโต ทาตพฺโพ, อุสฺสุกฺกํ วา กาตพฺพํ “กินฺติ นุ โข สทฺธิ\-วิหาริกสฺส ปตฺโต อุปฺปชฺชิเยถา”ติ. สเจ อุปชฺฌายสฺส จีวรํ โหติ สทฺธิ\-วิหาริกสฺส จีวรํ น โหติ, อุปชฺฌาเยน สทฺธิ\-วิหาริกสฺส จีวรํ ทาตพฺพํ, อุสฺสุกฺกํ วา กาตพฺพํ “กินฺติ นุ โข สทฺธิ\-วิหาริกสฺส จีวรํ อุปฺปชฺชิเยถา”ติ. สเจ อุปชฺฌายสฺส ปริกฺขาโร โหติ สทฺธิ\-วิหาริกสฺส ปริกฺขาโร น โหติ, อุปชฺฌาเยน สทฺธิ\-วิหาริกสฺส ปริกฺขาโร ทาตพฺโพ, อุสฺสุกฺกํ วา กาตพฺพํ “กินฺติ นุ โข สทฺธิ\-วิหาริกสฺส ปริกฺขาโร อุปฺปชฺชิเยถา”ติ.}

\prose{สเจ สทฺธิวิหาริโก คิลาโน โหติ, กาลสฺเสว อุฏฺฐาย ทนฺตกฏฺฐํ ทาตพฺพํ, มุโขทกํ ทาตพฺพํ, อาสนํ ปญฺญเปตพฺพํ.}

\prose{สเจ ยาคุ โหติ, ภาชนํ โธวิตฺวา ยาคุ อุปนาเมตพฺพา, ยาคุํ ปีตสฺส อุทกํ ทตฺวา ภาชนํ ปฏิคฺคเหตฺวา นีจํ กตฺวา สาธุกํ อปฺปฏิฆํสนฺเตน โธวิตฺวา ปฏิสาเมตพฺพํ, สทฺธิ\-วิหาริกมฺหิ วุฏฺฐิเต อาสนํ อุทฺธริตพฺพํ.}

\setlength{\csromangathapagewidth}{0pt}
\setlength{\csromangathawakwidth}{0pt}
\csromangathameasuregroup{สเจ โส เทโส อุกฺลาโป โหติ,}{\csromangathabat{สเจ โส เทโส อุกฺลาโป โหติ,}{โส เทโส สมฺมชฺชิตพฺโพ.}}
\csromangathasetleft{สเจ โส เทโส อุกฺลาโป โหติ,}
\csromangathagroup{\csromangathastanza{\csromangathabat{สเจ โส เทโส อุกฺลาโป โหติ,}{โส เทโส สมฺมชฺชิตพฺโพ.}}}

\prose{สเจ สทฺธิวิหาริโก คามํ ปวิสิตุกาโม โหติ, นิวาสนํ ทาตพฺพํ, ปฏินิวาสนํ ปฏิคฺคเหตพฺพํ, กายพนฺธนํ ทาตพฺพํ, สคุณํ กตฺวา สํฆาฏิโย ทาตพฺพา, โธวิตฺวา ปตฺโต โสทโก ทาตพฺโพ.}

\proseitem{67}{“เอตฺตาวตฺตา นิวตฺติสฺสตี”ติ อาสนํ ปญฺญเปตพฺพํ, ปาโททกํ ปาทปีฐํ \csromanfolio{64}ปาทกถลิกํ อุปนิกฺขิปิตพฺพํ, ปจฺจุคฺคนฺตฺวา ปตฺตจีวรํ ปฏิคฺคเหตพฺพํ, ปฏินิวาสนํ ทาตพฺพํ, นิวาสนํ ปฏิคฺคเหตพฺพํ, สเจ จีวรํ สินฺนํ โหติ, มุหุตฺตํ อุเณฺห โอตาเปตพฺพํ, น จ อุเณฺห จีวรํ นิทหิตพฺพํ, จีวรํ สงฺฆริตพฺพํ, จีวรํ สงฺฆรนฺเตน จตุรงฺคุลํ กณฺณํ อุสฺสาเรตฺวา จีวรํ สงฺฆริตพฺพํ “มา มชฺเฌ ภงฺโค อโหสี”ติ. โอโภเค กายพนฺธนํ กาตพฺพํ.}

\prose{สเจ ปิณฺฑปาโต โหติ สทฺธิวิหาริโก จ ภุญฺชิตุกาโม โหติ อุทกํ ทตฺวา ปิณฺฑปาโต อุปนาเมตพฺโพ, สทฺธิวิหาริโก ปานีเยน ปุจฺฉิตพฺโพ, ภุตฺตาวิสฺส อุทกํ ทตฺวา ปตฺตํ ปฏิคฺคเหตฺวา นีจํ กตฺวา สาธุกํ อปฺปฏิฆํสนฺเตน โธวิตฺวา โวทกํ กตฺวา มุหุตฺตํ อุเณฺห โอตาเปตพฺโพ, น จ อุเณฺห ปตฺโต นิทหิตพฺโพ.}

\prose{ปตฺตจีวรํ นิกฺขิปิตพฺพํ. ปตฺตํ นิกฺขิปนฺเตน เอเกน หตฺเถน ปตฺตํ คเหตฺวา เอเกน หตฺเถน เหฏฺฐามญฺจํ วา เหฏฺฐาปีฐํ วา ปรามสิตฺวา ปตฺโต นิกฺขิปิตพฺโพ, น จ อนนฺตรหิตาย ภูมิยา ปตฺโต นิกฺขิปิตพฺโพ. จีวรํ นิกฺขิปนฺเตน เอเกน หตฺเถน จีวรํ คเหตฺวา เอเกน หตฺเถน จีวรวํสํ วา จีวรรชฺชุํ วา ปมชฺชิตฺวา ปารโต อนฺตํ โอรโต โภคํ กตฺวา จีวรํ นิกฺขิปิตพฺพํ. สทฺธิ\-วิหาริกมฺหิ วุฏฺฐิเต อาสนํ อุทฺธริตพฺพํ, ปาโททกํ ปาทปีฐํ ปาทกถลิกํ ปฏิสาเมตพฺพํ. สเจ โส เทโส อุกฺลาโป โหติ, โส เทโส สมฺมชฺชิตพฺโพ.}

\setlength{\csromangathapagewidth}{0pt}
\setlength{\csromangathawakwidth}{0pt}
\csromangathameasuregroup{สเจ สทฺธิวิหาริโก นหายิตุกาโม โหติ, \\ สเจ สีเตน อตฺโถ โหติ, \\ สเจ อุเณฺหน อตฺโถ โหติ,}{\csromangathabat{สเจ สทฺธิวิหาริโก นหายิตุกาโม โหติ,}{นหานํ ปฏิยาเทตพฺพํ.} \\ \csromangathabat{สเจ สีเตน อตฺโถ โหติ,}{สีตํ ปฏิยาเทตพฺพํ.} \\ \csromangathabat{สเจ อุเณฺหน อตฺโถ โหติ,}{อุณฺหํ ปฏิยาเทตพฺพํ.}}
\csromangathasetleft{สเจ สทฺธิวิหาริโก นหายิตุกาโม โหติ, \\ สเจ สีเตน อตฺโถ โหติ, \\ สเจ อุเณฺหน อตฺโถ โหติ,}
\csromangathagroup{\csromangathastanza{\csromangathabat{สเจ สทฺธิวิหาริโก นหายิตุกาโม โหติ,}{นหานํ ปฏิยาเทตพฺพํ.} \\ \csromangathabat{สเจ สีเตน อตฺโถ โหติ,}{สีตํ ปฏิยาเทตพฺพํ.} \\ \csromangathabat{สเจ อุเณฺหน อตฺโถ โหติ,}{อุณฺหํ ปฏิยาเทตพฺพํ.}}}

\prose{สเจ สทฺธิวิหาริโก ชนฺตาฆรํ ปวิสิตุกาโม โหติ, จุณฺณํ สนฺเนตพฺพํ, มตฺติกา เตเมตพฺพา, ชนฺตาฆรปีฐํ อาทาย คนฺตฺวา ชนฺตาฆรปีฐํ ทตฺวา จีวรํ ปฏิคฺคเหตฺวา เอกมนฺตํ นิกฺขิปิตพฺพํ, จุณฺณํ ทาตพฺพํ, มตฺติกา ทาตพฺพา.}

\prose{สเจ อุสฺสหติ ชนฺตาฆรํ ปวิสิตพฺพํ, ชนฺตาฆรํ ปวิสนฺเตน มตฺติกาย มุขํ มกฺเขตฺวา ปุรโต จ ปจฺฉโต จ ปฏิจฺฉาเทตฺวา ชนฺตาฆรํ ปวิสิตพฺพํ, น เถเร ภิกฺขู อนุปขชฺช นิสีทิตพฺพํ, น นวา ภิกฺขู อาสเนน ปฏิพาหิตพฺพา, ชนฺตาฆเร สทฺธิ\-วิหาริกสฺส ปริกมฺมํ กาตพฺพํ, ชนฺตาฆรา นิกฺขมนฺเตน ชนฺตาฆรปีฐํ อาทาย ปุรโต จ ปจฺฉโต จ ปฏิจฺฉาเทตฺวา ชนฺตาฆรา นิกฺขมิตพฺพํ.}

\prose{\csromanfolio{65}อุทเกปิ สทฺธิ\-วิหาริกสฺส ปริกมฺมํ กาตพฺพํ, นหาเตน ปฐมตรํ อุตฺตริตฺวา อตฺตโน คตฺตํ โวทกํ กตฺวา นิวาเสตฺวา สทฺธิ\-วิหาริกสฺส คตฺตโต อุทกํ ปมชฺชิตพฺพํ, นิวาสนํ ทาตพฺพํ, สํฆาฏิ ทาตพฺพา, ชนฺตาฆรปีฐํ อาทาย ปฐมตรํ อาคนฺตฺวา อาสนํ ปญฺญเปตพฺพํ.}

\prose{ปาโททกํ ปาทปีฐํ ปาทกถลิกํ อุปนิกฺขิปิตพฺพํ. สทฺธิวิหาริโก ปานีเยน ปุจฺฉิตพฺโพ.}

\prose{ยสฺมิํ วิหาเร สทฺธิวิหาริโก วิหรติ สเจ โส วิหาโร อุกฺลาโป โหติ สเจ อุสฺสหติ, โสเธตพฺโพ.}

\prose{วิหารํ โสเธนฺเตน ปฐมํ ปตฺตจีวรํ นีหริตฺวา เอกมนฺตํ นิกฺขิปิตพฺพํ, นิสีทน\-ปจฺจตฺถรณํ นีหริตฺวา เอกมนฺตํ นิกฺขิปิตพฺพํ, ภิสิพิพฺโพหนํ นีหริตฺวา เอกมนฺตํ นิกฺขิปิตพฺพํ, มญฺโจ นีจํ กตฺวา สาธุกํ อปฺปฏิฆํสนฺเตน อสํฆฏฺเฏนฺเตน กวาฏปิฏฺฐํ นีหริตฺวา เอกมนฺตํ นิกฺขิปิตพฺโพ, ปีฐํ นีจํ กตฺวา สาธุกํ อปฺปฏิฆํสนฺเตน อสํฆฏฺเฏนฺเตน กวาฏปิฏฺฐํ นีหริตฺวา เอกมนฺตํ นิกฺขิปิตพฺพํ, มญฺจ\-ปฏิปาทกา นีหริตฺวา เอกมนฺตํ นิกฺขิปิตพฺพา.}

\prose{เขฬมลฺลโก นีหริตฺวา เอกมนฺตํ นิกฺขิปิตพฺโพ, อปสฺเสนผลกํ นีหริตฺวา เอกมนฺตํ นิกฺขิปิตพฺพํ.}

\prose{ภูมตฺถรณํ ยถาปญฺญตฺตํ สลฺลกฺเขตฺวา นีหริตฺวา เอกมนฺตํ นิกฺขิปิตพฺพํ. สเจ วิหาเร สนฺตานกํ โหติ, อุลฺโลกา ปฐมํ โอหาเรตพฺพํ, อาโลกสนฺธิ\-กณฺณภาคา ปมชฺชิตพฺพา. สเจ เครุก\-ปริกมฺม\-กตา ภิตฺติ กณฺณกิตา โหติ, โจฬกํ เตเมตฺวา ปีเฬตฺวา ปมชฺชิตพฺพา. สเจ กาฬวณฺณกตา ภูมิ กณฺณกิตา โหติ, โจฬกํ เตเมตฺวา ปีเฬตฺวา ปมชฺชิตพฺพา. สเจ อกตา โหติ ภูมิ, อุทเกน ปริปฺโผสิตฺวา สมฺมชฺชิตพฺพา “มา วิหาโร รเชน อุหญฺญี”ติ, สงฺการํ วิจินิตฺวา เอกมนฺตํ ฉฑฺเฑตพฺพํ.}

\proseitem{67}{ภูมตฺถรณํ โอตาเปตฺวา โสเธตฺวา ปปฺโผเฏตฺวา อติหริตฺวา ยถาปญฺญตฺตํ ปญฺญเปตพฺพํ. มญฺจ\-ปฏิปาทกา โอตาเปตฺวา ปมชฺชิตฺวา อติหริตฺวา ยถาฐาเน ฐเปตพฺพา, มญฺโจ โอตาเปตฺวา โสเธตฺวา ปปฺโผเฏตฺวา นีจํ กตฺวา สาธุกํ อปฺปฏิฆํสนฺเตน อสํฆฏฺเฏนฺเตน กวาฏปิฏฺฐํ อติหริตฺวา ยถาปญฺญตฺตํ ปญฺญเปตพฺโพ. ปีฐํ โอตาเปตฺวา โสเธตฺวา ปปฺโผเฏตฺวา นีจํ กตฺวา สาธุกํ อปฺปฏิฆํสนฺเตน อสํฆฏฺเฏนฺเตน กวาฏปิฏฺฐํ อติหริตฺวา ยถาปญฺญตฺตํ ปญฺญเปตพฺพํ. ภิสิพิพฺโพหนํ โอตาเปตฺวา โสเธตฺวา ปปฺโผเฏตฺวา อติหริตฺวา \csromanfolio{66}ยถาปญฺญตฺตํ ปญฺญเปตพฺพํ. นิสีทน\-ปจฺจตฺถรณํ โอตาเปตฺวา โสเธตฺวา ปปฺโผเฏตฺวา อติหริตฺวา ยถาปญฺญตฺตํ ปญฺญเปตพฺพํ. เขฬมลฺลโก โอตาเปตฺวา ปมชฺชิตฺวา อติหริตฺวา ยถาฐาเน ฐเปตพฺโพ. อปสฺเสนผลกํ โอตาเปตฺวา ปมชฺชิตฺวา อติหริตฺวา ยถาฐาเน ฐเปตพฺพํ. ปตฺตจีวรํ นิกฺขิปิตพฺพํ. ปตฺตํ นิกฺขิปนฺเตน เอเกน หตฺเถน ปตฺตํ คเหตฺวา เอเกน หตฺเถน เหฏฺฐามญฺจํ วา เหฏฺฐาปีฐํ วา ปรามสิตฺวา ปตฺโต นิกฺขิปิตพฺโพ, น จ อนนฺตรหิตาย ภูมิยา ปตฺโต นิกฺขิปิตพฺโพ. จีวรํ นิกฺขิปนฺเตน เอเกน หตฺเถน จีวรํ คเหตฺวา เอเกน หตฺเถน จีวรวํสํ วา จีวรรชฺชุํ วา ปมชฺชิตฺวา ปารโต อนฺตํ โอรโต โภคํ กตฺวา จีวรํ นิกฺขิปิตพฺพํ.}

\setlength{\csromangathapagewidth}{0pt}
\setlength{\csromangathawakwidth}{0pt}
\csromangathameasuregroup{สเจ ปริเวณํ อุกฺลาปํ โหติ, \\ สเจ โกฏฺฐโก อุกฺลาโป โหติ, \\ สเจ อคฺคิสาลา อุกฺลาปา โหติ,}{สเจ ปุรตฺถิมา สรชา วาตา วายนฺติ, \\ ปุรตฺถิมา วาตปานา ถเกตพฺพา. \\ สเจ ปจฺฉิมา สรชา วาตา วายนฺติ, \\ ปจฺฉิมา วาตปานา ถเกตพฺพา. \\ สเจ อุตฺตรา สรชา วาตา วายนฺติ, \\ อุตฺตรา วาตปานา ถเกตพฺพา. \\ สเจ ทกฺขิณา สรชา วาตา วายนฺติ, \\ ทกฺขิณา วาตปานา ถเกตพฺพา. \\ สเจ สีตกาโล โหติ, \\ ทิวา วาตปานา วิวริตพฺพา, รตฺติํ ถเกตพฺพา. \\ สเจ อุณฺหกาโล โหติ, \\ ทิวา วาตปานา ถเกตพฺพา, รตฺติํ วิวริตพฺพา. \\ \csromangathabat{สเจ ปริเวณํ อุกฺลาปํ โหติ,}{ปริเวณํ สมฺมชฺชิตพฺพํ.} \\ \csromangathabat{สเจ โกฏฺฐโก อุกฺลาโป โหติ,}{โกฏฺฐโก สมฺมชฺชิตพฺโพ.} \\ สเจ อุปฏฺฐานสาลา อุกฺลาปา โหติ, \\ อุปฏฺฐานสาลา สมฺมชฺชิตพฺพา. \\ \csromangathabat{สเจ อคฺคิสาลา อุกฺลาปา โหติ,}{อคฺคิสาลา สมฺมชฺชิตพฺพา.}}
\csromangathasetleft{สเจ ปริเวณํ อุกฺลาปํ โหติ, \\ สเจ โกฏฺฐโก อุกฺลาโป โหติ, \\ สเจ อคฺคิสาลา อุกฺลาปา โหติ,}
\csromangathagroup{\csromangathastanza{\csromangathaitemline{67}{สเจ ปุรตฺถิมา สรชา วาตา วายนฺติ,} \\ \csromangathacontline{ปุรตฺถิมา วาตปานา ถเกตพฺพา.} \\ \csromangathacontline{สเจ ปจฺฉิมา สรชา วาตา วายนฺติ,} \\ \csromangathacontline{ปจฺฉิมา วาตปานา ถเกตพฺพา.}}\csromangathastanza{\csromangathaitemline{67}{สเจ อุตฺตรา สรชา วาตา วายนฺติ,} \\ \csromangathacontline{อุตฺตรา วาตปานา ถเกตพฺพา.} \\ \csromangathacontline{สเจ ทกฺขิณา สรชา วาตา วายนฺติ,} \\ \csromangathacontline{ทกฺขิณา วาตปานา ถเกตพฺพา.}}\csromangathastanza{\csromangathaitemline{67}{สเจ สีตกาโล โหติ,} \\ \csromangathacontline{ทิวา วาตปานา วิวริตพฺพา, รตฺติํ ถเกตพฺพา.} \\ \csromangathacontline{สเจ อุณฺหกาโล โหติ,} \\ \csromangathacontline{ทิวา วาตปานา ถเกตพฺพา, รตฺติํ วิวริตพฺพา.}}\csromangathastanza{\csromangathabat{สเจ ปริเวณํ อุกฺลาปํ โหติ,}{ปริเวณํ สมฺมชฺชิตพฺพํ.} \\ \csromangathabat{สเจ โกฏฺฐโก อุกฺลาโป โหติ,}{โกฏฺฐโก สมฺมชฺชิตพฺโพ.}}\csromangathastanza{สเจ อุปฏฺฐานสาลา อุกฺลาปา โหติ, \\ อุปฏฺฐานสาลา สมฺมชฺชิตพฺพา. \\ \csromangathabat{สเจ อคฺคิสาลา อุกฺลาปา โหติ,}{อคฺคิสาลา สมฺมชฺชิตพฺพา.}}}

\prose{สเจ วจฺจกุฏิ อุกฺลาปา โหติ วจฺจกุฏิ สมฺมชฺชิตพฺพา. สเจ ปานียํ น โหติ, ปานียํ อุปฏฺฐาเป\-ตพฺพํ. สเจ ปริโภชนียํ น โหติ, ปริโภชนียํ อุปฏฺฐาเป\-ตพฺพํ. สเจ อาจมนกุมฺภิยา อุทกํ น โหติ, อาจมนกุมฺภิยา อุทกํ อาสิญฺจิตพฺพํ.}

\proseitem{67}{สเจ สทฺธิ\-วิหาริกสฺส อนภิรติ อุปฺปนฺนา โหติ, อุปชฺฌาเยน วูปกาเสตพฺโพ วูปกาสาเป\-ตพฺโพ ธมฺมกถา วาสฺส กาตพฺพา. สเจ สทฺธิ\-วิหาริกสฺส กุกฺกุจฺจํ อุปฺปนฺนํ โหติ, อุปชฺฌาเยน วิโนเทตพฺพํ วิโนทาเปตพฺพํ ธมฺมกถา วาสฺส กาตพฺพา. สเจ สทฺธิ\-วิหาริกสฺส ทิฏฺฐิคตํ อุปฺปนฺนํ โหติ, อุปชฺฌาเยน วิเวเจตพฺพํ วิเวจาเปตพฺพํ ธมฺมกถา วาสฺส กาตพฺพา. สเจ สทฺธิวิหาริโก ครุธมฺมํ อชฺฌาปนฺโน โหติ ปริวาสารโห, อุปชฺฌาเยน อุสฺสุกฺกํ กาตพฺพํ “กินฺติ นุ โข \csromanfolio{67}สํโฆ สทฺธิ\-วิหาริกสฺส ปริวาสํ ทเทยฺยา”ติ. สเจ สทฺธิวิหาริโก มูลาย ปฏิกสฺสนา\-รโห โหติ, อุปชฺฌาเยน อุสฺสุกฺกํ กาตพฺพํ “กินฺติ นุ โข สํโฆ สทฺธิ\-วิหาริกํ มูลาย ปฏิกสฺเสยฺยา”ติ. สเจ สทฺธิวิหาริโก มานตฺตารโห โหติ, อุปชฺฌาเยน อุสฺสุกฺกํ กาตพฺพํ “กินฺติ นุ โข สํโฆ สทฺธิ\-วิหาริกสฺส มานตฺตํ ทเทยฺยา”ติ. สเจ สทฺธิวิหาริโก อพฺภานารโห โหติ, อุปชฺฌาเยน อุสฺสุกฺกํ กาตพฺพํ “กินฺติ นุ โข สํโฆ สทฺธิ\-วิหาริกํ อพฺเภยฺยา”ติ. สเจ สํโฆ สทฺธิ\-วิหาริกสฺส กมฺมํ กตฺตุกาโม โหติ ตชฺชนียํ วา นิยสฺสํ วา ปพฺพาชนียํ วา ปฏิสารณียํ วา อุกฺเขปนียํ วา, อุปชฺฌาเยน อุสฺสุกฺกํ กาตพฺพํ “กินฺติ นุ โข สํโฆ สทฺธิ\-วิหาริกสฺส กมฺมํ น กเรยฺย ลหุกาย วา ปริณาเมยฺยา”ติ. กตํ วา ปนสฺส โหติ สํเฆน กมฺมํ ตชฺชนียํ วา นิยสฺสํ วา ปพฺพาชนียํ วา ปฏิสารณียํ วา อุกฺเขปนียํ วา, อุปชฺฌาเยน อุสฺสุกฺกํ กาตพฺพํ “กินฺติ นุ โข สทฺธิวิหาริโก สมฺมา วตฺเตยฺย, โลมํ ปาเตยฺย, เนตฺถารํ วตฺเตยฺย, สํโฆ ตํ กมฺมํ ปฏิปฺปสฺสมฺเภยฺยา”ติ.}

\prosewithcloser{สเจ สทฺธิ\-วิหาริกสฺส จีวรํ โธวิตพฺพํ โหติ, อุปชฺฌาเยน อาจิกฺขิตพฺพํ “เอวํ โธเวยฺยาสี”ติ, อุสฺสุกฺกํ วา กาตพฺพํ “กินฺติ นุ โข สทฺธิ\-วิหาริกสฺส จีวรํ โธวิเยถา”ติ. สเจ สทฺธิ\-วิหาริกสฺส จีวรํ กาตพฺพํ โหติ, อุปชฺฌาเยน อาจิกฺขิตพฺพํ “เอวํ กเรยฺยาสี”ติ, อุสฺสุกฺกํ วา กาตพฺพํ “กินฺติ นุ โข สทฺธิ\-วิหาริกสฺส จีวรํ กริเยถา”ติ. สเจ สทฺธิ\-วิหาริกสฺส รชนํ ปจิตพฺพํ โหติ, อุปชฺฌาเยน อาจิกฺขิตพฺพํ “เอวํ ปเจยฺยาสี”ติ, อุสฺสุกฺกํ วา กาตพฺพํ “กินฺติ นุ โข สทฺธิ\-วิหาริกสฺส รชนํ ปจิเยถา”ติ. สเจ สทฺธิ\-วิหาริกสฺส จีวรํ รชิตพฺพํ โหติ, อุปชฺฌาเยน อาจิกฺขิตพฺพํ “เอวํ รเชยฺยาสี”ติ, อุสฺสุกฺกํ วา กาตพฺพํ “กินฺติ นุ โข สทฺธิ\-วิหาริกสฺส จีวรํ รชิเยถา”ติ. จีวรํ รชนฺเตน สาธุกํ สมฺปริวตฺตกํ สมฺปริวตฺตกํ รชิตพฺพํ, น จ อจฺฉินฺเน เถเว ปกฺกมิตพฺพํ. สเจ สทฺธิวิหาริโก คิลาโน โหติ, ยาวชีวํ อุปฏฺฐาตพฺโพ, วุฏฺฐานมสฺส อาคเมตพฺพนฺติ.}{สทฺธิวิหาริก\-วตฺตํ นิฏฺฐิตํ.}
\csromanmatikaanchor{mtk.18}
\csromantocmarkref{section}{17. ปณามิตกถา}{mtk.18}
\bookmark[dest={mtk.18},level=1]{17. ปณามิตกถา}
\csromanheader{1}{17. \textbf{ปณามิตกถา}}
\proseitem{68}{\csromanfolio{68}เตน โข ปน สมเยน สทฺธิวิหาริกา อุปชฺฌาเยสุ น สมฺมา วตฺตนฺติ. เย เต ภิกฺขู อปฺปิจฺฉา ฯเปฯ เต อุชฺฌายนฺติ ขิยฺยนฺติ วิปาเจนฺติ “กถํ หิ นาม สทฺธิวิหาริกา อุปชฺฌาเยสุ น สมฺมา วตฺติสฺสนฺตี”ติ. อถ โข เต ภิกฺขู ภควโต เอตมตฺถํ อาโรเจสุํ ฯเปฯ “สจฺจํ กิร ภิกฺขเว สทฺธิวิหาริกา อุปชฺฌาเยสุ น สมฺมา วตฺตนฺตี”ติ. ส จฺจํ ภควาติ. วิครหิ พุทฺโธ ภควา ฯเปฯ. กถํ หิ นาม ภิกฺขเว สทฺธิวิหาริกา อุปชฺฌาเยสุ น สมฺมา วตฺติสฺสนฺตีติ ฯเปฯ วิครหิตฺวา ฯเปฯ ธมฺมิํ กถํ กตฺวา ภิกฺขู อามนฺเตสิ “น ภิกฺขเว สทฺธิ\-วิหาริเกน อุปชฺฌายมฺหิ น สมฺมา วตฺติตพฺพํ, โย น สมฺมา วตฺเตยฺย, อาปตฺติ ทุกฺกฏสฺสา”ติ. เนว สมฺมา วตฺตนฺติ. ภควโต เอตมตฺถํ อาโรเจสุํ. อนุชานามิ ภิกฺขเว อสมฺมา\-วตฺตนฺตํ ปณาเมตุํ. เอวญฺจ ปน ภิกฺขเว ปณาเมตพฺโพ “ปณาเมมิ ตนฺ”ติ วา “มายิธ ปฏิกฺกมี”ติ วา “นีหร เต ปตฺตจีวรนฺ”ติ วา “นาหํ ตยา อุปฏฺฐาตพฺโพ”ติ วา กาเยน วิญฺญาเปติ, วาจาย วิญฺญาเปติ, กาเยน วาจาย วิญฺญาเปติ, ปณามิโต โหติ สทฺธิวิหาริโก. น กาเยน วิญฺญาเปติ, น วาจาย วิญฺญาเปติ, น กาเยน วาจาย วิญฺญาเปติ, น ปณามิโต โหติ สทฺธิ\-วิหาริโกติ.}

\prose{เตน โข ปน สมเยน สทฺธิวิหาริกา ปณามิตา น ขมาเปนฺติ. ภควโต เอตมตฺถํ อาโรเจสุํ. อนุชานามิ ภิกฺขเว ขมาเปตุนฺติ. เนว ขมาเปนฺติ. ภควโต เอตมตฺถํ อาโรเจสุํ. น ภิกฺขเว ปณามิเตน น ขมาเปตพฺโพ, โย น ขมาเปยฺย, อาปตฺติ ทุกฺกฏสฺสาติ.}

\prose{เตน โข ปน สมเยน อุปชฺฌายา ขมาปิยมานา น ขมนฺติ. ภควโต เอตมตฺถํ อาโรเจสุํ. อนุชานามิ ภิกฺขเว ขมิตุนฺติ. เนว ขมนฺติ. สทฺธิวิหาริกา ปกฺกมนฺติปิ วิพฺภมนฺติปิ ติตฺถิเยสุปิ สงฺกมนฺติ. ภควโต เอตมตฺถํ อาโรเจสุํ. น ภิกฺขเว ขมาปิย\-มาเนน น ขมิตพฺพํ, โย น ขเมยฺย, อาปตฺติ ทุกฺกฏสฺสาติ.}

\proseitem{68}{เตน โข ปน สมเยน อุปชฺฌายา สมฺมาวตฺตนฺตํ ปณาเมนฺติ. อสมฺมา\-วตฺตนฺตํ น ปณาเมนฺติ. ภควโต เอตมตฺถํ อาโรเจสุํ. น ภิกฺขเว สมฺมาวตฺตนฺโต ปณาเมตพฺโพ, โย ปณาเมยฺย, อาปตฺติ \csromanfolio{69}ทุกฺกฏสฺส. น จ ภิกฺขเว อสมฺมาวตฺตนฺโต น ปณาเมตพฺโพ, โย น ปณาเมยฺย, อาปตฺติ ทุกฺกฏสฺสาติ.}

\prose{ปญฺจหิ ภิกฺขเว องฺเคหิ สมนฺนาคโต สทฺธิวิหาริโก ปณาเมตพฺโพ. อุปชฺฌายมฺหิ นาธิมตฺตํ เปมํ โหติ, นาธิมตฺโต ปสาโท โหติ, นาธิมตฺตา หิรี โหติ, นาธิมตฺโต คารโว โหติ, นาธิมตฺตา ภาวนา โหติ. อิเมหิ โข ภิกฺขเว ปญฺจหงฺเคหิ สมนฺนาคโต สทฺธิวิหาริโก ปณาเมตพฺโพ.}

\prose{ปญฺจหิ ภิกฺขเว องฺเคหิ สมนฺนาคโต สทฺธิวิหาริโก น ปณาเมตพฺโพ. อุปชฺฌายมฺหิ อธิมตฺตํ เปมํ โหติ, อธิมตฺโต ปสาโท โหติ, อธิมตฺตา หิรี โหติ, อธิมตฺโต คารโว โหติ, อธิมตฺตา ภาวนา โหติ. อิเมหิ โข ภิกฺขเว ปญฺจหงฺเคหิ สมนฺนาคโต สทฺธิวิหาริโก น ปณาเมตพฺโพ.}

\prose{ปญฺจหิ ภิกฺขเว องฺเคหิ สมนฺนาคโต สทฺธิวิหาริโก อลํ ปณาเมตุํ. อุปชฺฌายมฺหิ นาธิมตฺตํ เปมํ โหติ, นาธิมตฺโต ปสาโท โหติ, นาธิมตฺตา หิรี โหติ, นาธิมตฺโต คารโว โหติ, นาธิมตฺตา ภาวนา โหติ. อิเมหิ โข ภิกฺขเว ปญฺจหงฺเคหิ สมนฺนาคโต สทฺธิวิหาริโก อลํ ปณาเมตุํ.}

\prose{ปญฺจหิ ภิกฺขเว องฺเคหิ สมนฺนาคโต สทฺธิวิหาริโก นาลํ ปณาเมตุํ. อุปชฺฌายมฺหิ อธิมตฺตํ เปมํ โหติ, อธิมตฺโต ปสาโท โหติ, อธิมตฺตา หิรี โหติ, อธิมตฺโต คารโว โหติ, อธิมตฺตา ภาวนา โหติ. อิเมหิ โข ภิกฺขเว ปญฺจหงฺเคหิ สมนฺนาคโต สทฺธิวิหาริโก นาลํ ปณาเมตุํ.}

\proseitem{68}{ปญฺจหิ ภิกฺขเว องฺเคหิ สมนฺนาคตํ สทฺธิ\-วิหาริกํ อปฺปณาเมนฺโต อุปชฺฌาโย สาติสาโร โหติ, ปณาเมนฺโต อนติสาโร โหติ. อุปชฺฌายมฺหิ นาธิมตฺตํ เปมํ โหติ, นาธิมตฺโต ปสาโท โหติ, นาธิมตฺตา หิรี โหติ, นาธิมตฺโต คารโว โหติ, นาธิมตฺตา ภาวนา โหติ. อิเมหิ โข ภิกฺขเว ปญฺจหงฺเคหิ \csromanfolio{70}สมนฺนาคตํ สทฺธิ\-วิหาริกํ อปฺปณาเมนฺโต อุปชฺฌาโย สาติสาโร โหติ, ปณาเมนฺโต อนติสาโร โหติ.}

\prose{ปญฺจหิ ภิกฺขเว องฺเคหิ สมนฺนาคตํ สทฺธิ\-วิหาริกํ ปณาเมนฺโต อุปชฺฌาโย สาติสาโร โหติ, อปฺปณาเมนฺโต อนติสาโร โหติ. อุปชฺฌายมฺหิ อธิมตฺตํ เปมํ โหติ, อธิมตฺโต ปสาโท โหติ, อธิมตฺตา หิรี โหติ, อธิมตฺโต คารโว โหติ, อธิมตฺตา ภาวนา โหติ. อิเมหิ โข ภิกฺขเว ปญฺจหงฺเคหิ สมนฺนาคตํ สทฺธิ\-วิหาริกํ ปณาเมนฺโต อุปชฺฌาโย สาติสาโร โหติ, อปฺปณาเมนฺโต อนติสาโร โหตีติ.}

\proseitem{69}{เตน โข ปน สมเยน อญฺญตโร พฺราหฺมโณ ภิกฺขู อุปสงฺกมิตฺวา ปพฺพชฺชํ ยาจิ, ตํ ภิกฺขู น อิจฺฉิํสุ ปพฺพาเชตุํ, โส ภิกฺขูสุ ปพฺพชฺชํ อลภมาโน กิโส อโหสิ ลูโข ทุพฺพณฺโณ อุปฺปณฺฑุ\-ปฺปณฺฑุก\-ชาโต ธมนิสนฺถตคตฺโถ. อทฺทสา โข ภควา ตํ พฺราหฺมณํ กิสํ ลูขํ ทุพฺพณฺณํ อุปฺปณฺฑุ\-ปฺปณฺฑุก\-ชาตํ ธมนิสนฺถต\-คตฺตํ. ทิสฺวาน ภิกฺขู อามนฺเตสิ “กิํ นุ โข โส ภิกฺขเว พฺราหฺมโณ กิโส ลูโข ทุพฺพณฺโณ อุปฺปณฺฑุ\-ปฺปณฺฑุก\-ชาโต ธมนิสนฺถต\-คตฺโต”ติ. เอโส ภนฺเต พฺราหฺมโณ ภิกฺขู อุปสงฺกมิตฺวา ปพฺพชฺชํ ยาจิ, ตํ ภิกฺขู น อิจฺฉิํสุ ปพฺพาเชตุํ, โส ภิกฺขูสุ ปพฺพชฺชํ อลภมาโน กิโส ลูโข ทุพฺพณฺโณ อุปฺปณฺฑุ\-ปฺปณฺฑุก\-ชาโต ธมนิสนฺถต\-คตฺโตติ. อถ โข ภควา ภิกฺขู อามนฺเตสิ “โก นุ โข ภิกฺขเว ตสฺส พฺราหฺมณสฺส อธิการํ สรสี”ติ. เอวํ วุตฺเต อายสฺมา สาริปุตฺโต ภควนฺตํ เอตทโวจ “อหํ โข ภนฺเต ตสฺส พฺราหฺมณสฺส อธิการํ สรามี”ติ. กิํ ปน ตฺวํ สาริปุตฺต ตสฺส พฺราหฺมณสฺส อธิการํ สรสีติ. อิธ เม ภนฺเต โส พฺราหฺมโณ ราชคเห ปิณฺฑาย จรนฺตสฺส กฏจฺฉุภิกฺขํ ทาเปสิ, อิมํ โข อหํ ภนฺเต ตสฺส พฺราหฺมณสฺส อธิการํ สรามีติ. สาธุ สาธุ สาริปุตฺต, กตญฺญุโน หิ สาริปุตฺต สปฺปุริสา กตเวทิโน, เตน หิ ตฺวํ สาริปุตฺต ตํ พฺราหฺมณํ ปพฺพาเชหิ อุปสมฺปาเทหีติ. กถาหํ ภนฺเต ตํ พฺราหฺมณํ ปพฺพาเชมิ อุปสมฺปาเทมีติ. อถ โข ภควา เอตสฺมิํ นิทาเน เอตสฺมิํ ปกรเณ ธมฺมิํ กถํ กตฺวา ภิกฺขู อามนฺเตสิ\csromandash{}ยา สา ภิกฺขเว มยา ตีหิ สรณคมเนหิ อุปสมฺปทา อนุญฺญาตา, ตํ อชฺชตคฺเค ปฏิกฺขิปามิ. อนุชานามิ ภิกฺขเว ญตฺติ\-จตุตฺเถน กมฺเมน \csromanfolio{71}อุปสมฺปาเทตุํ\csromansharedfootnote{418b29f465016665}{อุปสมฺปทํ (สี, สฺยา)}. เอวญฺจ ปน ภิกฺขเว อุปสมฺ\-ปาเท\-ตพฺโพ, พฺยตฺเตน ภิกฺขุนา ปฏิพเลน สํโฆ ญาเปตพฺโพ\csromandash{}}

\proseitem{70}{“สุณาตุ เม ภนฺเต สํโฆ, อยํ อิตฺถนฺนาโม อิตฺถนฺนามสฺส อายสฺมโต อุปสมฺปทา\-เปกฺโข, ยทิ สํฆสฺส ปตฺตกลฺลํ, สํโฆ อิตฺถนฺนามํ อุปสมฺปาเทยฺย อิตฺถนฺนาเมน อุปชฺฌาเยน, เอสา ญตฺติ.}

\prose{สุณาตุ เม ภนฺเต สํโฆ, อยํ อิตฺถนฺนาโม อิตฺถนฺนามสฺส อายสฺมโต อุปสมฺปทา\-เปกฺโข, สํโฆ อิตฺถนฺนามํ อุปสมฺปาเทติ อิตฺถนฺนาเมน อุปชฺฌาเยน, ยสฺสายสฺมโต ขมติ อิตฺถนฺนามสฺส อุปสมฺปทา อิตฺถนฺนาเมน อุปชฺฌาเยน, โส ตุณฺหสฺส, ยสฺส นกฺขมติ, โส ภาเสยฺย.}

\prose{ทุติยมฺปิ เอตมตฺถํ วาทามิ. สุณาตุ เม ภนฺเต สํโฆ, อยํ อิตฺถนฺนาโม อิตฺถนฺนามสฺส อายสฺมโต อุปสมฺปทา\-เปกฺโข, สํโฆ อิตฺถนฺนามํ อุปสมฺปาเทติ อิตฺถนฺนาเมน อุปชฺฌาเยน, ยสฺสายสฺมโต ขมติ อิตฺถนฺนามสฺส อุปสมฺปทา อิตฺถนฺนาเมน อุปชฺฌาเยน, โส ตุณฺหสฺส, ยสฺส นกฺขมติ, โส ภาเสยฺย.}

\prose{ตติยมฺปิ เอตมตฺถํ วทามิ. สุณาตุ เม ภนฺเต สํโฆ, อยํ อิตฺถนฺนาโม อิตฺถนฺนามสฺส อายสฺมโต อุปสมฺปทา\-เปกฺโข, สํโฆ อิตฺถนฺนามํ อุปสมฺปาเทติ อิตฺถนฺนาเมน อุปชฺฌาเยน, ยสฺสายสฺมโต ขมติ อิตฺถนฺนามสฺส อุปสมฺปทา อิตฺถนฺนาเมน อุปชฺฌาเยน, โส ตุณฺหสฺส, ยสฺส นกฺขมติ, โส ภาเสยฺย.}

\prose{อุปสมฺปนฺโน สํเฆน อิตฺถนฺนาโม อิตฺถนฺนาเมน อุปชฺฌาเยน, ขมติ สํฆสฺส, ตสฺมา ตุณฺหี, เอวเมตํ ธารยามี”ติ.}

\proseitem{71}{เตน โข ปน สมเยน อญฺญตโร ภิกฺขุ อุปสมฺปนฺน\-สมนนฺตรา อนาจารํ อาจรติ. ภิกฺขู เอวมาหํสุ “มาวุโส เอวรูปํ อกาสิ, เนตํ กปฺปตี”ติ. โส เอวมาห “เนวาหํ อายสฺมนฺเต ยาจิํ ‘อุปสมฺปาเทถ มนฺ’ติ, กิสฺส มํ ตุเมฺห อยาจิตา อุปสมฺปาทิตฺถา”ติ. ภควโต เอตมตฺถํ อาโรเจสุํ. น ภิกฺขเว อยาจิเตน \csromanfolio{72}อุปสมฺ\-ปาเท\-ตพฺโพ, โย อุปสมฺปาเทยฺย, อาปตฺติ ทุกฺกฏสฺส. อนุชานามิ ภิกฺขเว ยาจิเตน อุปสมฺปาเทตุํ. เอวญฺจ ปน ภิกฺขเว ยาจิตพฺโพ. เตน อุปสมฺปทา\-เปกฺเขน สํฆํ อุปสงฺกมิตฺวา เอกํสํ อุตฺตราสงฺคํ กริตฺวา ภิกฺขูนํ ปาเท วนฺทิตฺวา อุกฺกุฏิกํ นิสีทิตฺวา อญฺชลิํ ปคฺคเหตฺวา เอวมสฺส วจนีโย “สํฆํ ภนฺเต อุปสมฺปทํ ยาจามิ. อุลฺลุมฺปตุ มํ ภนฺเต สํโฆ อนุกมฺปํ อุปาทายา”ติ. ทุติยมฺปิ ยาจิตพฺโพ. ตติยมฺปิ ยาจิตพฺโพ. พฺยตฺเตน ภิกฺขุนา ปฏิพเลน สํโฆ ญาเปตพฺโพ\csromandash{}}

\proseitem{72}{“สุณาตุ เม ภนฺเต สํโฆ, อยํ อิตฺถนฺนาโม อิตฺถนฺนามสฺส อายสฺมโต อุปสมฺปทา\-เปกฺโข, อิตฺถนฺนาโม สํฆํ อุปสมฺปทํ ยาจติ อิตฺถนฺนาเมน อุปชฺฌาเยน, ยทิ สํฆสฺส ปตฺตกลฺลํ, สํโฆ อิตฺถนฺนามํ อุปสมฺปาเทยฺย อิตฺถนฺนาเมน อุปชฺฌาเยน, เอสา ญตฺติ.}

\prose{สุณาตุ เม ภนฺเต สํโฆ, อยํ อิตฺถนฺนาโม อิตฺถนฺนามสฺส อายสฺมโต อุปสมฺปทา\-เปกฺโข, อิตฺถนฺนาโม สํฆํ อุปสมฺปทํ ยาจติ อิตฺถนฺนาเมน อุปชฺฌาเยน, สํโฆ อิตฺถนฺนามํ อุปสมฺปาเทติ อิตฺถนฺนาเมน อุปชฺฌาเยน, ยสฺสายสฺมโต ขมติ อิตฺถนฺนามสฺส อุปสมฺปทา อิตฺถนฺนาเมน อุปชฺฌาเยน, โส ตุณฺหสฺส, ยสฺส นกฺขมติ, โส ภาเสยฺย.}

\setlength{\csromangathapagewidth}{0pt}
\setlength{\csromangathawakwidth}{0pt}
\csromangathameasuregroup{}{ทุติยมฺปิ เอตมตฺถํ วทามิ ฯเปฯ. ตติยมฺปิ เอตมตฺถํ วทามิ ฯเปฯ.}
\csromangathameasurewak{ทุติยมฺปิ เอตมตฺถํ วทามิ ฯเปฯ. ตติยมฺปิ เอตมตฺถํ วทามิ ฯเปฯ.}
\setlength{\csromangathaleftcol}{0pt}
\csromangathagroup{\csromangathastanza{\csromangathawak{ทุติยมฺปิ เอตมตฺถํ วทามิ ฯเปฯ. ตติยมฺปิ เอตมตฺถํ วทามิ ฯเปฯ.}}}

\prose{อุปสมฺปนฺโน สํเฆน อิตฺถนฺนาโม อิตฺถนฺนาเมน อุปชฺฌาเยน, ขมติ สํฆสฺส, ตสฺมา ตุณฺหี, เอวเมตํ ธารยามี”ติ.}

\proseitemwithcloserruled{73}{เตน โข ปน สมเยน ราชคเห ปณีตานํ ภตฺตานํ ภตฺตปฏิปาฏิ อฏฺฐิตา\csromansharedfootnote{831b03332f561769}{อธิฏฺฐิตา (ก)} โหติ. อถ โข อญฺญตรสฺส พฺราหฺมณสฺส เอตทโหสิ “อิเม โข สมณา สกฺยปุตฺติยา สุขสีลา สุขสมาจารา สุโภชนานิ ภุญฺชิตฺวา นิวาเตสุ สยเนสุ สยนฺติ, ยนฺนูนาหํ สมเณสุ สกฺยปุตฺติเยสุ ปพฺพเชยฺยนฺ”ติ. อถ โข โส พฺราหฺมโณ ภิกฺขู อุปสงฺกมิตฺวา ปพฺพชฺชํ ยาจิ. ตํ ภิกฺขู ปพฺพาเชสุํ อุปสมฺปาเทสุํ. ตสฺมิํ ปพฺพชิเต ภตฺตปฏิปาฏิ ขียิตฺถ. ภิกฺขู เอวมาหํสุ “เอหิ ทานิ อาวุโส ปิณฺฑาย จริสฺสามา”ติ. โส เอวมาห “นาหํ อาวุโส เอตํการณา ปพฺพชิโต ‘ปิณฺฑาย จริสฺสามี’ติ, สเจ เม ทสฺสถ \csromanfolio{73}ภุญฺชิสฺสามิ, โน เจ เม ทสฺสถ วิพฺภมิสฺสามี”ติ. กิํ ปน ตฺวํ อาวุโส อุทรสฺส การณา ปพฺพชิโตติ. เอวมาวุโสติ. เย เต ภิกฺขู อปฺปิจฺฉา ฯเปฯ เต อุชฺฌายนฺติ ขิยฺยนฺติ วิปาเจนฺติ “กถํ หิ นาม ภิกฺขุ เอวํ สฺวากฺขาเต ธมฺมวินเย อุทรสฺส การณา ปพฺพชิสฺสตี”ติ. เต ภิกฺขู ภควโต เอตมตฺถํ อาโรเจสุํ ฯเปฯ “สจฺจํ กิร ตฺวํ ภิกฺขุ อุทรสฺส การณา ปพฺพชิโต”ติ. สจฺจํ ภควาติ. วิครหิ พุทฺโธ ภควา ฯเปฯ. กถํ หิ นาม ตฺวํ โมฆปุริส เอวํ สฺวากฺขาเต ธมฺมวินเย อุทรสฺส การณา ปพฺพชิสฺสสิ. เนตํ โมฆปุริส อปฺปสนฺนานํ วา ปสาทาย ปสนฺนานํ วา ภิยฺโยภาวาย ฯเปฯ วิครหิตฺวา ฯเปฯ ธมฺมิํ กถํ กตฺวา ภิกฺขู อามนฺเตสิ “อนุชานามิ ภิกฺขเว อุปสมฺปาเทนฺเตน จตฺตาโร นิสฺสเย อาจิกฺขิตุํ. ปิณฺฑิยาโลป\-โภชนํ นิสฺสาย ปพฺพชฺชา, ตตฺถ เต ยาวชีวํ อุสฺสาโห กรณีโย, อติเรกลาโภ สํฆภตฺตํ อุทฺเทสภตฺตํ นิมนฺตนํ สลากภตฺตํ ปกฺขิกํ อุโปสถิกํ ปาฏิปทิกํ. ปํสุกูล\-จีวรํ นิสฺสาย ปพฺพชฺชา, ตตฺถ เต ยาวชีวํ อุสฺสาโห กรณีโย, อติเรกลาโภ โขมํ กปฺปาสิกํ โกเสยฺยํ กมฺพลํ สาณํ ภงฺคํ. รุกฺขมูล\-เสนาสนํ นิสฺสาย ปพฺพชฺชา, ตตฺถ เต ยาวชีวํ อุสฺสาโห กรณีโย, อติเรกลาโภ วิหาโร อฑฺฒโยโค ปาสาโท หมฺมิยํ คุหา. ปูติมุตฺต\-เภสชฺชํ นิสฺสาย ปพฺพชฺชา, ตตฺถ เต ยาวชีวํ อุสฺสาโห กรณีโย, อติเรกลาโภ สปฺปิ นวนีตํ เตลํ มธุ ผาณิตํ”ติ.}{อุปชฺฌาย\-วตฺต\-ภาณวาโร นิฏฺฐิโต ปญฺจโม.}
\csromanmatikaanchor{mtk.19}
\csromantocmarkref{section}{18. อาจริยวตฺตกถา}{mtk.19}
\bookmark[dest={mtk.19},level=1]{18. อาจริยวตฺตกถา}
\csromanheader{1}{18. \textbf{อาจริยวตฺต}\-\textbf{กถา}}
\proseitem{74}{เตน โข ปน สมเยน อญฺญตโร มาณวโก ภิกฺขู อุปสงฺกมิตฺวา ปพฺพชฺชํ ยาจิ. ตสฺส ภิกฺขู ปฏิกจฺเจว นิสฺสเย อาจิกฺขิํสุ. โส เอวมาห “สเจ เม ภนฺเต ปพฺพชิเต นิสฺสเย อาจิกฺเขยฺยาถ, อภิรเมยฺยามหํ\csromansharedfootnote{667d683e897febbc}{อภิรเมยฺยญฺจาหํ (สี), อภิรเมยฺยํ สฺวาหํ (ก)}, น ทานาหํ ภนฺเต ปพฺพชิสฺสามิ, เชคุจฺฉา เม \csromanfolio{74}นิสฺสยา ปฏิกูลา”ติ. ภควโต เอตมตฺถํ อาโรเจสุํ. น ภิกฺขเว ปฏิกจฺเจว นิสฺสยา อาจิกฺขิตพฺพา, โย อาจิกฺเขยฺย, อาปตฺติ ทุกฺกฏสฺส. อนุชานามิ ภิกฺขเว อุปสมฺปนฺน\-สมนนฺตรา นิสฺสเย อาจิกฺขิตุนฺติ.}

\prose{เตน โข ปน สมเยน ภิกฺขู ทุวคฺเคนปิ ติวคฺเคนปิ คเณน อุปสมฺปาเทนฺติ. ภควโต เอตมตฺถํ อาโรเจสุํ. น ภิกฺขเว อูนทสวคฺเคน คเณน อุปสมฺ\-ปาเท\-ตพฺโพ, โย อุปสมฺปาเทยฺย, อาปตฺติ ทุกฺกฏสฺส. อนุชานามิ ภิกฺขเว ทสวคฺเคน วา อติเรกทส\-วคฺเคน วา คเณน อุปสมฺปาเทตุนฺติ.}

\proseitem{75}{เตน โข ปน สมเยน ภิกฺขู เอกวสฺสาปิ ทุวสฺสาปิ สทฺธิ\-วิหาริกํ อุปสมฺปาเทนฺติ. อายสฺมาปิ อุปเสโน วงฺคนฺตปุตฺโต เอกวสฺโส สทฺธิ\-วิหาริกํ อุปสมฺปาเทติ, โส วสฺสํวุฏฺโฐ ทุวสฺโส เอกวสฺสํ สทฺธิ\-วิหาริกํ อาทาย เยน ภควา เตนุปสงฺกมิ, อุปสงฺกมิตฺวา ภควนฺตํ อภิวาเทตฺวา เอกมนฺตํ นิสีทิ. อาจิณฺณํ โข ปเนตํ พุทฺธานํ ภควนฺตานํ อาคนฺตุเกหิ ภิกฺขูหิ สทฺธิํ ปฏิสมฺโมทิตุํ. อถ โข ภควา อายสฺมนฺตํ อุปเสนํ วงฺคนฺตปุตฺตํ เอตทโวจ “กจฺจิ ภิกฺขุ ขมนียํ, กจฺจิ ยาปนียํ, กจฺจิ ตฺวํ อปฺป\-กิลมเถน อทฺธานํ อาคโต”ติ. ขมนียํ ภควา, ยาปนียํ ภควา, อปฺป\-กิลมเถน มยํ ภนฺเต อทฺธานํ อาคตาติ. ชานนฺตาปิ ตถาคตา ปุจฺฉนฺติ, ชานนฺตาปิ น ปุจฺฉนฺติ. กาลํ วิทิตฺวา ปุจฺฉนฺติ, กาลํ วิทิตฺวา น ปุจฺฉนฺติ. อตฺถสํหิตํ ตถาคตา ปุจฺฉนฺติ, โน อนตฺถ\-สํหิตํ. อนตฺถสํหิเต เสตุฆาโต ตถาคตานํ. ทฺวีหิ อากาเรหิ พุทฺธา ภควนฺโต ภิกฺขู ปฏิปุจฺฉนฺติ “ธมฺมํ วา เทเสสฺสาม, สาวกานํ วา สิกฺขาปทํ ปญฺญเปสฺสามา”ติ. อถ โข ภควา อายสฺมนฺตํ อุปเสนํ วงฺคนฺตปุตฺตํ เอตทโวจ “กติวสฺโสสิ ตฺวํ ภิกฺขู”ติ. ทุวสฺโสหํ ภควาติ. อยํ ปน ภิกฺขุ กติวสฺโสติ. เอกวสฺโส ภควาติ. กินฺตายํ ภิกฺขุ โหตีติ. สทฺธิวิหาริโก เม ภควาติ. วิครหิ พุทฺโธ ภควา อนนุจฺฉวิกํ โมฆปุริส อนนุโลมิกํ อปฺปติรูปํ อสฺสามณกํ อกปฺปิยํ อกรณียํ, กถํ หิ นาม ตฺวํ โมฆปุริส อญฺเญหิ โอวทิโย อนุสาสิโย อญฺญํ โอวทิตุํ อนุสาสิตุํ มญฺญิสฺสสิ. อติลหุํ โข ตฺวํ โมฆปุริส พาหุลฺลาย อาวตฺโต ยทิทํ คณพนฺธิกํ. เนตํ โมฆปุริส อปฺปสนฺนานํ วา ปสาทาย \csromanfolio{75}ปสนฺนานํ วา ภิยฺโยภาวาย ฯเปฯ วิครหิตฺวา ฯเปฯ ธมฺมิํ กถํ กตฺวา ภิกฺขู อามนฺเตสิ “น ภิกฺขเว อูนทสวสฺเสน อุปสมฺ\-ปาเท\-ตพฺโพ, โย อุปสมฺปาเทยฺย, อาปตฺติ ทุกฺกฏสฺส. อนุชานามิ ภิกฺขเว ทสวสฺเสน วา อติเรก\-ทสวสฺเสน วา อุปสมฺปาเทตุนฺ”ติ.}

\proseitem{76}{เตน โข ปน สมเยน ภิกฺขู “ทสวสฺสมฺหา ทสวสฺสมฺหา”ติ พาลา อพฺยตฺตา อุปสมฺปาเทนฺติ. ทิสฺสนฺติ อุปชฺฌายา พาลา สทฺธิวิหาริกา ปณฺฑิตา, ทิสฺสนฺติ อุปชฺฌายา อพฺยตฺตา สทฺธิวิหาริกา พฺยตฺตา, ทิสฺสนฺติ อุปชฺฌายา อปฺปสฺสุตา สทฺธิวิหาริกา พหุสฺสุตา, ทิสฺสนฺติ อุปชฺฌายา ทุปฺปญฺญา สทฺธิวิหาริกา ปญฺญวนฺโต. อญฺญตโรปิ อญฺญ\-ติตฺถิย\-ปุพฺโพ อุปชฺฌาเยน สหธมฺมิกํ วุจฺจมาโน อุปชฺฌายสฺส วาทํ อาโรเปตฺวา ตํเยว ติตฺถายตนํ สงฺกมิ. เย เต ภิกฺขู อปฺปิจฺฉา ฯเปฯ เต อุชฺฌายนฺติ ขิยฺยนฺติ วิปาเจนฺติ “กถํ หิ นาม ภิกฺขู ‘ทสวสฺสมฺหา ทสวสฺสมฺหา’ติ พาลา อพฺยตฺตา อุปสมฺปาเทสฺสนฺติ. ทิสฺสนฺติ อุปชฺฌายา พาลา สทฺธิวิหาริกา ปณฺฑิตา, ทิสฺสนฺติ อุปชฺฌายา อพฺยตฺตา สทฺธิวิหาริกา พฺยตฺตา, ทิสฺสนฺติ อุปชฺฌายา อปฺปสฺสุตา สทฺธิวิหาริกา พหุสฺสุตา, ทิสฺสนฺติ อุปชฺฌายา ทุปฺปญฺญา สทฺธิวิหาริกา ปญฺญวนฺโต”ติ. อถ โข เต ภิกฺขู ภควโต เอตมตฺถํ อาโรเจสุํ. สจฺจํ กิร ภิกฺขเว ภิกฺขู “ทสวสฺสมฺหา ทสวสฺสมฺหา”ติ พาลา อพฺยตฺตา อุปสมฺปาเทนฺติ. ทิสฺสนฺติ อุปชฺฌายา พาลา สทฺธิวิหาริกา ปณฺฑิตา, ทิสฺสนฺติ อุปชฺฌายา อพฺยตฺตา สทฺธิวิหาริกา พฺยตฺตา, ทิสฺสนฺติ อุปชฺฌายา อปฺปสฺสุตา สทฺธิวิหาริกา พหุสฺสุตา, ทิสฺสนฺติ อุปชฺฌายา ทุปฺปญฺญา สทฺธิวิหาริกา ปญฺญวนฺโตติ. สจฺจํ ภควาติ. วิครหิ พุทฺโธ ภควา ฯเปฯ. กถํ หิ นาม เต ภิกฺขเว โมฆปุริสา “ทสวสฺสมฺหา ทสวสฺสมฺหา”ติ พาลา อพฺยตฺตา อุปสมฺปาเทสฺสนฺติ. ทิสฺสนฺติ อุปชฺฌายา พาลา สทฺธิวิหาริกา ปณฺฑิตา, ทิสฺสนฺติ อุปชฺฌายา อพฺยตฺตา สทฺธิวิหาริกา พฺยตฺตา, ทิสฺสนฺติ อุปชฺฌายา อปฺปสฺสุตา สทฺธิวิหาริกา พหุสฺสุตา, ทิสฺสนฺติ อุปชฺฌายา ทุปฺปญฺญา สทฺธิวิหาริกา ปญฺญวนฺโต. เนตํ ภิกฺขเว อปฺปสนฺนานํ วา ปสาทาย ฯเปฯ วิครหิตฺวา ฯเปฯ ธมฺมิํ กถํ กตฺวา ภิกฺขู อามนฺเตสิ “น ภิกฺขเว พาเลน อพฺยตฺเตน อุปสมฺ\-ปาเท\-ตพฺโพ, โย อุปสมฺปาเทยฺย, อาปตฺติ ทุกฺกฏสฺส. อนุชานามิ ภิกฺขเว พฺยตฺเตน ภิกฺขุนา ปฏิพเลน ทสวสฺเสน วา อติเรก\-ทสวสฺเสน วา อุปสมฺปาเทตุนฺ”ติ.}

\proseitem{77}{\csromanfolio{76}เตน โข ปน สมเยน ภิกฺขู อุปชฺฌาเยสุ ปกฺกนฺเตสุปิ วิพฺภนฺเตสุปิ กาลงฺกเตสุปิ ปกฺข\-สงฺกนฺเตสุปิ อนาจริยกา อโนวทิยมานา อนนุสาสิยมานา ทุนฺนิวตฺถา ทุปฺปารุตา อนากปฺป\-สมฺปนฺนา ปิณฺฑาย จรนฺติ, มนุสฺสานํ ภุญฺชมานานํ อุปริโภชเนปิ อุตฺติฏฺฐปตฺตํ อุปนาเมนฺติ, อุปริ\-ขาทนีเยปิ อุปริ\-สายนีเยปิ อุปริปานีเยปิ อุตฺติฏฺฐปตฺตํ อุปนาเมนฺติ, สามํ สูปมฺปิ โอทนมฺปิ วิญฺญาเปตฺวา ภุญฺชนฺติ, ภตฺตคฺเคปิ อุจฺจาสทฺทา มหาสทฺทา วิหรนฺติ. มนุสฺสา อุชฺฌายนฺติ ขิยฺยนฺติ วิปาเจนฺติ “กถํ หิ นาม สมณา สกฺยปุตฺติยา ทุนฺนิวตฺถา ทุปฺปารุตา อนากปฺป\-สมฺปนฺนา ปิณฺฑาย จริสฺสนฺติ, มนุสฺสานํ ภุญฺชมานานํ อุปริโภชเนปิ อุตฺติฏฺฐปตฺตํ อุปนาเมสฺสนฺติ, อุปริ\-ขาทนีเยปิ อุปริ\-สายนีเยปิ อุปริปานีเยปิ อุตฺติฏฺฐปตฺตํ อุปนาเมสฺสนฺติ, สามํ สูปมฺปิ โอทนมฺปิ วิญฺญาเปตฺวา ภุญฺชิสฺสนฺติ, ภตฺตคฺเคปิ อุจฺจาสทฺทา มหาสทฺทา วิหริสฺสนฺติ เสยฺยถาปิ พฺราหฺมณา พฺราหฺมณ\-โภชเน”ติ. อสฺโสสุํ โข ภิกฺขู เตสํ มนุสฺสานํ อุชฺฌายนฺตานํ ขิยฺยนฺตานํ วิปาเจนฺตานํ ฯเปฯ. อถ โข เต ภิกฺขู ภควโต เอตมตฺถํ อาโรเจสุํ ฯเปฯ. สจฺจํ กิร ภิกฺขเว ฯเปฯ. สจฺจํ ภควาติ ฯเปฯ. วิครหิตฺวา ธมฺมิํ กถํ กตฺวา ภิกฺขู อามนฺเตสิ\csromandash{}}

\prose{อนุชานามิ ภิกฺขเว อาจริยํ, อาจริโย ภิกฺขเว อนฺเตวาสิกมฺหิ ปุตฺตจิตฺตํ อุปฏฺฐาเปสฺสติ, อนฺเตวาสิโก อาจริยมฺหิ ปิตุจิตฺตํ อุปฏฺฐาเปสฺสติ, เอวํ เต อญฺญมญฺญํ สคารวา สปฺปติสฺสา สภาควุตฺติโน วิหรนฺตา อิมสฺมิํ ธมฺมวินเย วุทฺธิํ วิรุฬฺหิํ เวปุลฺลํ อาปชฺชิสฺสนฺติ. อนุชานามิ ภิกฺขเว ทสวสฺสํ นิสฺสาย วตฺถุํ ทสวสฺเสน นิสฺสยํ ทาตุํ. เอวญฺจ ปน ภิกฺขเว อาจริโย คเหตพฺโพ, เอกํสํ อุตฺตราสงฺคํ กริตฺวา ปาเท วนฺทิตฺวา อุกฺกุฏิกํ นิสีทิตฺวา อญฺชลิํ ปคฺคเหตฺวา เอวมสฺส วจนีโย “อาจริโย เม ภนฺเต โหหิ, อายสฺมโต นิสฺสาย วจฺฉามิ. อาจริโย เม ภนฺเต โหหิ, อายสฺมโต นิสฺสาย วจฺฉามิ. อาจริโย เม ภนฺเต โหหิ, อายสฺมโต นิสฺสาย วจฺฉามี”ติ. “สาหู”ติ วา “ลหู”ติ วา “โอปายิกนฺ”ติ วา “ปติรูปนฺ”ติ วา “ปาสาทิเกน สมฺปาเทหี”ติ วา กาเยน วิญฺญาเปติ, วาจาย วิญฺญาเปติ, กาเยน วาจาย วิญฺญาเปติ, คหิโต โหติ อาจริโย. น กาเยน วิญฺญาเปติ, น วาจาย วิญฺญาเปติ, น กาเยน วาจาย วิญฺญาเปติ, น คหิโต โหติ อาจริโย.}

\proseitem{78}{\csromanfolio{77}\csromansymbolfootnote{*}{วิ 4. 402 ปิฏฺฐาทีสุปิ.}อนฺเตวาสิเกน ภิกฺขเว อาจริยมฺหิ สมฺมา วตฺติตพฺพํ. ตตฺรายํ สมฺมาวตฺตนา\csromandash{}}

\prose{กาลสฺเสว อุฏฺฐาย อุปาหนํ โอมุญฺจิตฺวา เอกํสํ อุตฺตราสงฺคํ กริตฺวา ทนฺตกฏฺฐํ ทาตพฺพํ, มุโขทกํ ทาตพฺพํ, อาสนํ ปญฺญเปตพฺพํ.}

\prose{สเจ ยาคุ โหติ, ภาชนํ โธวิตฺวา ยาคุ อุปนาเมตพฺพา, ยาคุํ ปีตสฺส อุทกํ ทตฺวา ภาชนํ ปฏิคฺคเหตฺวา นีจํ กตฺวา สาธุกํ อปฺปฏิฆํสนฺเตน โธวิตฺวา ปฏิสาเมตพฺพํ.}

\prose{อาจริยมฺหิ วุฏฺฐิเต อาสนํ อุทฺธริตพฺพํ. สเจ โส เทโส อุกฺลาโป โหติ, โส เทโส สมฺมชฺชิตพฺโพ.}

\prose{สเจ อาจริโย คามํ ปวิสิตุกาโม โหติ, นิวาสนํ ทาตพฺพํ, ปฏินิวาสนํ ปฏิคฺคเหตพฺพํ, กายพนฺธนํ ทาตพฺพํ, สคุณํ กตฺวา สํฆาฏิโย ทาตพฺพา, โธวิตฺวา ปตฺโต โสทโก ทาตพฺโพ.}

\prose{สเจ อาจริโย ปจฺฉาสมณํ อากงฺขติ, ติมณฺฑลํ ปฏิจฺฉาเทนฺเตน ปริมณฺฑลํ นิวาเสตฺวา กายพนฺธนํ พนฺธิตฺวา สคุณํ กตฺวา สํฆาฏิโย ปารุปิตฺวา คณฺฐิกํ ปฏิมุญฺจิตฺวา โธวิตฺวา ปตฺตํ คเหตฺวา อาจริยสฺส ปจฺฉาสมเณน โหตพฺพํ, นาติทูเร คนฺตพฺพํ, นาจฺจาสนฺเน คนฺตพฺพํ, ปตฺต\-ปริยาปนฺนํ ปฏิคฺคเหตพฺพํ.}

\prose{น อาจริยสฺส ภณมานสฺส อนฺตรนฺตรา กถา โอปาเตตพฺพา, อาจริโย อาปตฺติสามนฺตา ภณมาโน นิวาเรตพฺโพ.}

\prose{นิวตฺตนฺเตน ปฐมตรํ อาคนฺตฺวา อาสนํ ปญฺญเปตพฺพํ, ปาโททกํ ปาทปีฐํ ปาทกถลิกํ อุปนิกฺขิปิตพฺพํ, ปจฺจุคฺคนฺตฺวา ปตฺตจีวรํ ปฏิคฺคเหตพฺพํ, ปฏินิวาสนํ ทาตพฺพํ, นิวาสนํ ปฏิคฺคเหตพฺพํ.}

\prose{สเจ จีวรํ สินฺนํ โหติ, มุหุตฺตํ อุเณฺห โอตาเปตพฺพํ, น จ อุเณฺห จีวรํ นิทหิตพฺพํ.}

\prose{จีวรํ สงฺฆริตพฺพํ, จีวรํ สงฺฆรนฺเตน จตุรงฺคุลํ กณฺณํ อุสฺสาเรตฺวา จีวรํ สงฺฆริตพฺพํ “มา มชฺเฌ ภงฺโค อโหสี”ติ.}

\setlength{\csromangathapagewidth}{0pt}
\setlength{\csromangathawakwidth}{0pt}
\csromangathameasuregroup{}{โอโภเค กายพนฺธนํ กาตพฺพํ.}
\csromangathameasurewak{โอโภเค กายพนฺธนํ กาตพฺพํ.}
\setlength{\csromangathaleftcol}{0pt}
\csromangathagroup{\csromangathastanza{\csromangathawak{โอโภเค กายพนฺธนํ กาตพฺพํ.}}}

\proseitem{78}{สเจ ปิณฺฑปาโต โหติ, อาจริโย จ ภุญฺชิตุกาโม โหติ, อุทกํ ทตฺวา ปิณฺฑปาโต อุปนาเมตพฺโพ, อาจริโย ปานีเยน ปุจฺฉิตพฺโพ, ภุตฺตาวิสฺส อุทกํ ทตฺวา ปตฺตํ ปฏิคฺคเหตฺวา นีจํ กตฺวา สาธุกํ อปฺปฏิฆํสนฺเตน โธวิตฺวา โวทกํ กตฺวา มุหุตฺตํ อุเณฺห โอตาเปตพฺโพ, น จ อุเณฺห ปตฺโต นิทหิตพฺโพ. ปตฺตจีวรํ นิกฺขิปิตพฺพํ. ปตฺตํ นิกฺขิปนฺเตน เอเกน หตฺเถน ปตฺตํ คเหตฺวา เอเกน หตฺเถน เหฏฺฐามญฺจํ วา เหฏฺฐาปีฐํ วา ปรามสิตฺวา ปตฺโต นิกฺขิปิตพฺโพ, น จ อนนฺตรหิตาย ภูมิยา \csromanfolio{78}ปตฺโต นิกฺขิปิตพฺโพ. จีวรํ นิกฺขิปนฺเตน เอเกน หตฺเถน จีวรํ คเหตฺวา เอเกน หตฺเถน จีวรวํสํ วา จีวรรชฺชุํ วา ปมชฺชิตฺวา ปารโต อนฺตํ โอรโต โภคํ กตฺวา จีวรํ นิกฺขิปิตพฺพํ. อาจริยมฺหิ วุฏฺฐิเต อาสนํ อุทฺธริตพฺพํ. ปาโททกํ ปาทปีฐํ ปาทกถลิกํ ปฏิสาเมตพฺพํ. สเจ โส เทโส อุกฺลาโป โหติ, โส เทโส สมฺมชฺชิตพฺโพ.}

\setlength{\csromangathapagewidth}{0pt}
\setlength{\csromangathawakwidth}{0pt}
\csromangathameasuregroup{สเจ อาจริโย นหายิตุกาโม โหติ, \\ สเจ สีเตน อตฺโถ โหติ, \\ สเจ อุเณฺหน อตฺโถ โหติ,}{\csromangathabat{สเจ อาจริโย นหายิตุกาโม โหติ,}{นหานํ ปฏิยาเทตพฺพํ.} \\ \csromangathabat{สเจ สีเตน อตฺโถ โหติ,}{สีตํ ปฏิยาเทตพฺพํ.} \\ \csromangathabat{สเจ อุเณฺหน อตฺโถ โหติ,}{อุณฺหํ ปฏิยาเทตพฺพํ.}}
\csromangathasetleft{สเจ อาจริโย นหายิตุกาโม โหติ, \\ สเจ สีเตน อตฺโถ โหติ, \\ สเจ อุเณฺหน อตฺโถ โหติ,}
\csromangathagroup{\csromangathastanza{\csromangathaitembat{78}{สเจ อาจริโย นหายิตุกาโม โหติ,}{นหานํ ปฏิยาเทตพฺพํ.} \\ \csromangathacontbat{สเจ สีเตน อตฺโถ โหติ,}{สีตํ ปฏิยาเทตพฺพํ.} \\ \csromangathacontbat{สเจ อุเณฺหน อตฺโถ โหติ,}{อุณฺหํ ปฏิยาเทตพฺพํ.}}}

\prose{สเจ อาจริโย ชนฺตาฆรํ ปวิสิตุกาโม โหติ, จุณฺณํ สนฺเนตพฺพํ, มตฺติกา เตเมตพฺพา, ชนฺตาฆรปีฐํ อาทาย อาจริยสฺส ปิฏฺฐิโต ปิฏฺฐิโต คนฺตฺวา ชนฺตาฆรปีฐํ ทตฺวา จีวรํ ปฏิคฺคเหตฺวา เอกมนฺตํ นิกฺขิปิตพฺพํ, จุณฺณํ ทาตพฺพํ, มตฺติกา ทาตพฺพา.}

\prose{สเจ อุสฺสหติ ชนฺตาฆรํ ปวิสิตพฺพํ, ชนฺตาฆรํ ปวิสนฺเตน มตฺติกาย มุขํ มกฺเขตฺวา ปุรโต จ ปจฺฉโต จ ปฏิจฺฉาเทตฺวา ชนฺตาฆรํ ปวิสิตพฺพํ, น เถเร ภิกฺขู อนุปขชฺช นิสีทิตพฺพํ, น นวา ภิกฺขู อาสเนน ปฏิพาหิตพฺพา, ชนฺตาฆเร อาจริยสฺส ปริกมฺมํ กาตพฺพํ, ชนฺตาฆรา นิกฺขมนฺเตน ชนฺตาฆรปีฐํ อาทาย ปุรโต จ ปจฺฉโต จ ปฏิจฺฉาเทตฺวา ชนฺตาฆรา นิกฺขมิตพฺพํ.}

\prose{อุทเกปิ อาจริยสฺส ปริกมฺมํ กาตพฺพํ, นหาเตน ปฐมตรํ อุตฺตริตฺวา อตฺตโน คตฺตํ โวทกํ กตฺวา นิวาเสตฺวา อาจริยสฺส คตฺตโต อุทกํ ปมชฺชิตพฺพํ, นิวาสนํ ทาตพฺพํ, สํฆาฏิ ทาตพฺพา, ชนฺตาฆรปีฐํ อาทาย ปฐมตรํ อาคนฺตฺวา อาสนํ ปญฺญเปตพฺพํ.}

\prose{ปาโททกํ ปาทปีฐํ ปาทกถลิกํ อุปนิกฺขิปิตพฺพํ. อาจริโย ปานีเยน ปุจฺฉิตพฺโพ. สเจ อุทฺทิสาเปตุกาโม โหติ, อุทฺทิสาเปตพฺโพ. สเจ ปริปุจฺฉิตุ\-กาโม โหติ, ปริปุจฺฉิตพฺโพ.}

\proseitem{78}{ยสฺมิํ วิหาเร อาจริโย วิหรติ สเจ โส วิหาโร อุกฺลาโป โหติ สเจ อุสฺสหติ, โสเธตพฺโพ. วิหารํ โสเธนฺเตน ปฐมํ ปตฺตจีวรํ นีหริตฺวา เอกมนฺตํ นิกฺขิปิตพฺพํ, นิสีทน\-ปจฺจตฺถรณํ นีหริตฺวา เอกมนฺตํ นิกฺขิปิตพฺพํ, ภิสิพิพฺโพหนํ นีหริตฺวา เอกมนฺตํ นิกฺขิปิตพฺพํ, มญฺโจ นีจํ กตฺวา สาธุกํ อปฺปฏิฆํสนฺเตน อสํฆฏฺเฏนฺเตน กวาฏปิฏฺฐํ นีหริตฺวา เอกมนฺตํ นิกฺขิปิตพฺโพ, ปีฐํ นีจํ กตฺวา สาธุกํ อปฺปฏิฆํสนฺเตน \csromanfolio{79}อสํฆฏฺเฏนฺเตน กวาฏปิฏฺฐํ นีหริตฺวา เอกมนฺตํ นิกฺขิปิตพฺพํ, มญฺจ\-ปฏิปาทกา นีหริตฺวา เอกมนฺตํ นิกฺขิปิตพฺพา. เขฬมลฺลโก นีหริตฺวา เอกมนฺตํ นิกฺขิปิตพฺโพ, อปสฺเสนผลกํ นีหริตฺวา เอกมนฺตํ นิกฺขิปิตพฺพํ. ภูมตฺถรณํ ยถาปญฺญตฺตํ สลฺลกฺเขตฺวา นีหริตฺวา เอกมนฺตํ นิกฺขิปิตพฺพํ. สเจ วิหาเร สนฺตานกํ โหติ, อุลฺโลกา ปฐมํ โอหาเรตพฺพํ, อาโลกสนฺธิ\-กณฺณภาคา ปมชฺชิตพฺพา. สเจ เครุก\-ปริกมฺม\-กตา ภิตฺติ กณฺณกิตา โหติ, โจฬกํ เตเมตฺวา ปีเฬตฺวา ปมชฺชิตพฺพา. สเจ กาฬวณฺณกตา ภูมิ กณฺณกิตา โหติ, โจฬกํ เตเมตฺวา ปีเฬตฺวา ปมชฺชิตพฺพา. สเจ อกตา โหติ ภูมิ, อุทเกน ปริปฺโผสิตฺวา สมฺมชฺชิตพฺพา “มา วิหาโร รเชน อุหญฺญี”ติ, สงฺการํ วิจินิตฺวา เอกมนฺตํ ฉฑฺเฑตพฺพํ.}

\prose{ภูมตฺถรณํ โอตาเปตฺวา โสเธตฺวา ปปฺโผเฏตฺวา อติหริตฺวา ยถาปญฺญตฺตํ ปญฺญเปตพฺพํ. มญฺจ\-ปฏิปาทกา โอตาเปตฺวา ปมชฺชิตฺวา อติหริตฺวา ยถาฐาเน ฐเปตพฺพา, มญฺโจ โอตาเปตฺวา โสเธตฺวา ปปฺโผเฏตฺวา นีจํ กตฺวา สาธุกํ อปฺปฏิฆํสนฺเตน อสํฆฏฺเฏนฺเตน กวาฏปิฏฺฐํ อติหริตฺวา ยถาปญฺญตฺตํ ปญฺญเปตพฺโพ. ปีฐํ โอตาเปตฺวา โสเธตฺวา ปปฺโผเฏตฺวา นีจํ กตฺวา สาธุกํ อปฺปฏิฆํสนฺเตน อสํฆฏฺเฏนฺเตน กวาฏปิฏฺฐํ อติหริตฺวา ยถาปญฺญตฺตํ ปญฺญเปตพฺพํ. ภิสิพิพฺโพหนํ โอตาเปตฺวา โสเธตฺวา ปปฺโผเฏตฺวา อติหริตฺวา ยถาปญฺญตฺตํ ปญฺญเปตพฺพํ. นิสีทน\-ปจฺจตฺถรณํ โอตาเปตฺวา โสเธตฺวา ปปฺโผเฏตฺวา อติหริตฺวา ยถาปญฺญตฺตํ ปญฺญเปตพฺพํ. เขฬมลฺลโก โอตาเปตฺวา ปมชฺชิตฺวา อติหริตฺวา ยถาฐาเน ฐเปตพฺโพ. อปสฺเสนผลกํ โอตาเปตฺวา ปมชฺชิตฺวา อติหริตฺวา ยถาฐาเน ฐเปตพฺพํ. ปตฺตจีวรํ นิกฺขิปิตพฺพํ. ปตฺตํ นิกฺขิปนฺเตน เอเกน หตฺเถน ปตฺตํ คเหตฺวา เอเกน หตฺเถน เหฏฺฐามญฺจํ วา เหฏฺฐาปีฐํ วา ปรามสิตฺวา ปตฺโต นิกฺขิปิตพฺโพ, น จ อนนฺตรหิตาย ภูมิยา ปตฺโต นิกฺขิปิตพฺโพ. จีวรํ นิกฺขิปนฺเตน เอเกน หตฺเถน จีวรํ คเหตฺวา เอเกน หตฺเถน จีวรวํสํ วา จีวรรชฺชุํ วา ปมชฺชิตฺวา ปารโต อนฺตํ โอรโต โภคํ กตฺวา จีวรํ นิกฺขิปิตพฺพํ.}

\prose{สเจ ปุรตฺถิมา สรชา วาตา วายนฺติ, ปุรตฺถิมา วาตปานา ถเกตพฺพา. สเจ ปจฺฉิมา สรชา วาตา วายนฺติ, ปจฺฉิมา วาตปานา ถเกตพฺพา. สเจ อุตฺตรา สรชา วาตา วายนฺติ, อุตฺตรา วาตปานา \csromanfolio{80}ถเกตพฺพา. สเจ ทกฺขิณา สรชา วาตา วายนฺติ, ทกฺขิณา วาตปานา ถเกตพฺพา. สเจ สีตกาโล โหติ, ทิวา วาตปานา วิวริตพฺพา, รตฺติํ ถเกตพฺพา. สเจ อุณฺหกาโล โหติ, ทิวา วาตปานา ถเกตพฺพา, รตฺติํ วิวริตพฺพา.}

\prose{สเจ ปริเวณํ อุกฺลาปํ โหติ, ปริเวณํ สมฺมชฺชิตพฺพํ. สเจ โกฏฺฐโก อุกฺลาโป โหติ, โกฏฺฐโก สมฺมชฺชิตพฺโพ. สเจ อุปฏฺฐานสาลา อุกฺลาปา โหติ, อุปฏฺฐานสาลา สมฺมชฺชิตพฺพา. สเจ อคฺคิสาลา อุกฺลาปา โหติ, อคฺคิสาลา สมฺมชฺชิตพฺพา. สเจ วจฺจกุฏิ อุกฺลาปา โหติ, วจฺจกุฏิ สมฺมชฺชิตพฺพา. สเจ ปานียํ น โหติ, ปานียํ อุปฏฺฐาเป\-ตพฺพํ. สเจ ปริโภชนียํ น โหติ, ปริโภชนียํ อุปฏฺฐาเป\-ตพฺพํ. สเจ อาจมนกุมฺภิยํ อุทกํ น โหติ, อาจมนกุมฺภิยา อุทกํ อาสิญฺจิตพฺพํ.}

\prose{สเจ อาจริยสฺส อนภิรติ อุปฺปนฺนา โหติ, อนฺเตวาสิเกน วูปกาเสตพฺโพ วูปกาสาเป\-ตพฺโพ ธมฺมกถา วาสฺส กาตพฺพา. สเจ อาจริยสฺส กุกฺกุจฺจํ อุปฺปนฺนํ โหติ, อนฺเตวาสิเกน วิโนเทตพฺพํ วิโนทาเปตพฺพํ ธมฺมกถา วาสฺส กาตพฺพา. สเจ อาจริยสฺส ทิฏฺฐิคตํ อุปฺปนฺนํ โหติ, อนฺเตวาสิเกน วิเวเจตพฺพํ วิเวจาเปตพฺพํ ธมฺมกถา วาสฺส กาตพฺพา. สเจ อาจริโย ครุธมฺมํ อชฺฌาปนฺโน โหติ ปริวาสารโห, อนฺเตวาสิเกน อุสฺสุกฺกํ กาตพฺพํ “กินฺติ นุ โข สํโฆ อาจริยสฺส ปริวาสํ ทเทยฺยา”ติ. สเจ อาจริโย มูลาย ปฏิกสฺสนา\-รโห โหติ, อนฺเตวาสิเกน อุสฺสุกฺกํ กาตพฺพํ “กินฺติ นุ โข สํโฆ อาจริยํ มูลาย ปฏิกสฺเสยฺยา”ติ. สเจ อาจริโย มานตฺตารโห โหติ, อนฺเตวาสิเกน อุสฺสุกฺกํ กาตพฺพํ “กินฺติ นุ โข สํโฆ อาจริยสฺส มานตฺตํ ทเทยฺยา”ติ. สเจ อาจริโย อพฺภานารโห โหติ, อนฺเตวาสิเกน อุสฺสุกฺกํ กาตพฺพํ “กินฺติ นุ โข สํโฆ อาจริยํ อพฺเภยฺยา”ติ. สเจ สํโฆ อาจริยสฺส กมฺมํ กตฺตุกาโม โหติ ตชฺชนียํ วา นิยสฺสํ วา ปพฺพาชนียํ วา ปฏิสารณียํ วา อุกฺเขปนียํ วา, อนฺเตวาสิเกน อุสฺสุกฺกํ กาตพฺพํ “กินฺติ นุ โข สํโฆ อาจริยสฺส กมฺมํ น กเรยฺย, ลหุกาย วา ปริณาเมยฺยา”ติ. กตํ วา ปนสฺส โหติ สํเฆน กมฺมํ ตชฺชนียํ วา นิยสฺสํ วา ปพฺพาชนียํ วา \csromanfolio{81}ปฏิสารณียํ วา อุกฺเขปนียํ วา, อนฺเตวาสิเกน อุสฺสุกฺกํ กาตพฺพํ “กินฺติ นุ โข อาจริโย สมฺมา วตฺเตยฺย, โลมํ ปาเตยฺย, เนตฺถารํ วตฺเตยฺย, สํโฆ ตํ กมฺมํ ปฏิปฺปสฺสมฺเภยฺยา”ติ.}

\prose{สเจ อาจริยสฺส จีวรํ โธวิตพฺพํ โหติ, อนฺเตวาสิเกน โธวิตพฺพํ, อุสฺสุกฺกํ วา กาตพฺพํ “กินฺติ นุ โข อาจริยสฺส จีวรํ โธวิเยถา”ติ.}

\prose{สเจ อาจริยสฺส จีวรํ กาตพฺพํ โหติ, อนฺเตวาสิเกน กาตพฺพํ, อุสฺสุกฺกํ วา กาตพฺพํ “กินฺติ นุ โข อาจริยสฺส จีวรํ กริเยถา”ติ.}

\prose{สเจ อาจริยสฺส รชนํ ปจิตพฺพํ โหติ, อนฺเตวาสิเกน ปจิตพฺพํ, อุสฺสุกฺกํ วา กาตพฺพํ “กินฺติ นุ โข อาจริยสฺส รชนํ ปจิเยถา”ติ.}

\prose{สเจ อาจริยสฺส จีวรํ รชิตพฺพํ โหติ, อนฺเตวาสิเกน รชิตพฺพํ, อุสฺสุกฺกํ วา กาตพฺพํ “กินฺติ นุ โข อาจริยสฺส จีวรํ รชิเยถา”ติ.}

\prose{จีวรํ รชนฺเตน สาธุกํ สมฺปริวตฺตกํ สมฺปริวตฺตกํ รชิตพฺพํ, น จ อจฺฉินฺเน เถเว ปกฺกมิตพฺพํ.}

\prose{น อาจริยํ อนาปุจฺฉา เอกจฺจสฺส ปตฺโต ทาตพฺโพ, น เอกจฺจสฺส ปตฺโต ปฏิคฺคเหตพฺโพ, น เอกจฺจสฺส จีวรํ ทาตพฺพํ, น เอกจฺจสฺส จีวรํ ปฏิคฺคเหตพฺพํ, น เอกจฺจสฺส ปริกฺขาโร ทาตพฺโพ, น เอกจฺจสฺส ปริกฺขาโร ปฏิคฺคเหตพฺโพ, น เอกจฺจสฺส เกสา เฉเทตพฺพา, น เอกจฺเจน เกสา เฉทาเปตพฺพา, น เอกจฺจสฺส ปริกมฺมํ กาตพฺพํ, น เอกจฺเจน ปริกมฺมํ การาเปตพฺพํ, น เอกจฺจสฺส เวยฺยาวจฺโจ กาตพฺโพ, น เอกจฺเจน เวยฺยาวจฺโจ การาเปตพฺโพ, น เอกจฺจสฺส ปจฺฉาสมเณน โหตพฺพํ, น เอกจฺโจ ปจฺฉาสมโณ อาทาตพฺโพ, น เอกจฺจสฺส ปิณฺฑปาโต นีหริตพฺโพ, น เอกจฺเจน ปิณฺฑปาโต นีหราเปตพฺโพ, น อาจริยํ อนาปุจฺฉา คาโม ปวิสิตพฺโพ, น สุสานํ คนฺตพฺพํ, น ทิสา ปกฺกมิตพฺพา.}

\prosewithcloserruled{สเจ อาจริโย คิลาโน โหติ, ยาวชีวํ อุปฏฺฐาตพฺโพ, วุฏฺฐานมสฺส อาคเมตพฺพนฺติ.}{อาจริยวตฺตํ นิฏฺฐิตํ.}
\csromanmatikaanchor{mtk.20}
\csromantocmarkref{section}{19. อนฺเตวาสิกวตฺตกถา}{mtk.20}
\bookmark[dest={mtk.20},level=1]{19. อนฺเตวาสิกวตฺตกถา}
\csromanheader{1}{19. \textbf{อนฺเตวาสิกวตฺต}\-\textbf{กถา}}
\proseitem{79}{\csromansymbolfootnote{*}{วิ 4. 407 ปิฏฺฐาทีสุปิ.}อาจริเยน ภิกฺขเว อนฺเตวาสิกมฺหิ สมฺมา วตฺติตพฺพํ. ตตฺรายํ สมฺมาวตฺตนา\csromandash{}}

\prose{\csromanfolio{82}อาจริเยน ภิกฺขเว อนฺเตวาสิโก สงฺคเหตพฺโพ อนุคฺคเหตพฺโพ อุทฺเทเสน ปริปุจฺฉาย โอวาเทน อนุสาสนิยา. สเจ อาจริยสฺส ปตฺโต โหติ อนฺเตวาสิกสฺส ปตฺโต น โหติ, อาจริเยน อนฺเตวาสิกสฺส ปตฺโต ทาตพฺโพ, อุสฺสุกฺกํ วา กาตพฺพํ “กินฺติ นุ โข อนฺเตวาสิกสฺส ปตฺโต อุปฺปชฺชิเยถา”ติ. สเจ อาจริยสฺส จีวรํ โหติ อนฺเตวาสิกสฺส จีวรํ น โหติ, อาจริเยน อนฺเตวาสิกสฺส จีวรํ ทาตพฺพํ, อุสฺสุกฺกํ วา กาตพฺพํ “กินฺติ นุ โข อนฺเตวาสิกสฺส จีวรํ อุปฺปชฺชิเยถา”ติ. สเจ อาจริยสฺส ปริกฺขาโร โหติ อนฺเตวาสิกสฺส ปริกฺขาโร น โหติ, อาจริเยน อนฺเตวาสิกสฺส ปริกฺขาโร ทาตพฺโพ, อุสฺสุกฺกํ วา กาตพฺพํ “กินฺติ นุ โข อนฺเตวาสิกสฺส ปริกฺขาโร อุปฺปชฺชิเยถา”ติ.}

\prose{สเจ อนฺเตวาสิโก คิลาโน โหติ, กาลสฺเสว อุฏฺฐาย ทนฺตกฏฺฐํ ทาตพฺพํ, มุโขทกํ ทาตพฺพํ, อาสนํ ปญฺญเปตพฺพํ.}

\prose{สเจ ยาคุ โหติ, ภาชนํ โธวิตฺวา ยาคุ อุปนาเมตพฺพา, ยาคุํ ปีตสฺส อุทกํ ทตฺวา ภาชนํ ปฏิคฺคเหตฺวา นีจํ กตฺวา สาธุกํ อปฺปฏิฆํสนฺเตน โธวิตฺวา ปฏิสาเมตพฺพํ.}

\prose{อนฺเตวาสิกมฺหิ วุฏฺฐิเต อาสนํ อุทฺธริตพฺพํ. สเจ โส เทโส อุกฺลาโป โหติ, โส เทโส สมฺมชฺชิตพฺโพ.}

\prose{สเจ อนฺเตวาสิโก คามํ ปวิสิตุกาโม โหติ, นิวาสนํ ทาตพฺพํ, ปฏินิวาสนํ ปฏิคฺคเหตพฺพํ, กายพนฺธนํ ทาตพฺพํ, สคุณํ กตฺวา สํฆาฏิโย ทาตพฺพา, โธวิตฺวา ปตฺโต โสทโก ทาตพฺโพ.}

\prose{“เอตฺตาวตา นิวตฺติสฺสตี”ติ อาสนํ ปญฺญเปตพฺพํ, ปาโททกํ ปาทปีฐํ ปาทกถลิกํ อุปนิกฺขิปิตพฺพํ, ปจฺจุคฺคนฺตฺวา ปตฺตจีวรํ ปฏิคฺคเหตพฺพํ, ปฏินิวาสนํ ทาตพฺพํ, นิวาสนํ ปฏิคฺคเหตพฺพํ.}

\prose{สเจ จีวรํ สินฺนํ โหติ, มุหุตฺตํ อุเณฺห โอตาเปตพฺพํ, น จ อุเณฺห จีวรํ นิทหิตพฺพํ, จีวรํ สงฺฆริตพฺพํ, จีวรํ สงฺฆรนฺเตน จตุรงฺคุลํ กณฺณํ อุสฺสาเรตฺวา จีวรํ สงฺฆริตพฺพํ “มา มชฺเฌ ภงฺโค อโหสี”ติ.}

\setlength{\csromangathapagewidth}{0pt}
\setlength{\csromangathawakwidth}{0pt}
\csromangathameasuregroup{}{โอโภเค กายพนฺธนํ กาตพฺพํ.}
\csromangathameasurewak{โอโภเค กายพนฺธนํ กาตพฺพํ.}
\setlength{\csromangathaleftcol}{0pt}
\csromangathagroup{\csromangathastanza{\csromangathawak{โอโภเค กายพนฺธนํ กาตพฺพํ.}}}

\proseitem{79}{สเจ ปิณฺฑปาโต โหติ อนฺเตวาสิโก จ ภุญฺชิตุกาโม โหติ อุทกํ ทตฺวา ปิณฺฑปาโต อุปนาเมตพฺโพ, อนฺเตวาสิโก ปานีเยน ปุจฺฉิตพฺโพ, ภุตฺตาวิสฺส อุทกํ ทตฺวา ปตฺตํ ปฏิคฺคเหตฺวา นีจํ กตฺวา สาธุกํ อปฺปฏิฆํสนฺเตน โธวิตฺวา โวทกํ กตฺวา มุหุตฺตํ \csromanfolio{83}อุเณฺห โอตาเปตพฺโพ, น จ อุเณฺห ปตฺโต นิทหิตพฺโพ. ปตฺตจีวรํ นิกฺขิปิตพฺพํ. ปตฺตํ นิกฺขิปนฺเตน เอเกน หตฺเถน ปตฺตํ คเหตฺวา เอเกน หตฺเถน เหฏฺฐามญฺจํ วา เหฏฺฐาปีฐํ วา ปรามสิตฺวา ปตฺโต นิกฺขิปิตพฺโพ, น จ อนนฺตรหิตาย ภูมิยา ปตฺโต นิกฺขิปิตพฺโพ. จีวรํ นิกฺขิปนฺเตน เอเกน หตฺเถน จีวรํ คเหตฺวา เอเกน หตฺเถน จีวรวํสํ วา จีวรรชฺชุํ วา ปมชฺชิตฺวา ปารโต อนฺตํ โอรโต โภคํ กตฺวา จีวรํ นิกฺขิปิตพฺพํ. อนฺเตวาสิกมฺหิ วุฏฺฐิเต อาสนํ อุทฺธริตพฺพํ. ปาโททกํ ปาทปีฐํ ปาทกถลิกํ ปฏิสาเมตพฺพํ. สเจ โส เทโส อุกฺลาโป โหติ, โส เทโส สมฺมชฺชิตพฺโพ.}

\setlength{\csromangathapagewidth}{0pt}
\setlength{\csromangathawakwidth}{0pt}
\csromangathameasuregroup{สเจ อนฺเตวาสิโก นหายิตุกาโม โหติ, \\ สเจ สีเตน อตฺโถ โหติ, \\ สเจ อุเณฺหน อตฺโถ โหติ,}{\csromangathabat{สเจ อนฺเตวาสิโก นหายิตุกาโม โหติ,}{นหานํ ปฏิยาเทตพฺพํ.} \\ \csromangathabat{สเจ สีเตน อตฺโถ โหติ,}{สีตํ ปฏิยาเทตพฺพํ.} \\ \csromangathabat{สเจ อุเณฺหน อตฺโถ โหติ,}{อุณฺหํ ปฏิยาเทตพฺพํ.}}
\csromangathasetleft{สเจ อนฺเตวาสิโก นหายิตุกาโม โหติ, \\ สเจ สีเตน อตฺโถ โหติ, \\ สเจ อุเณฺหน อตฺโถ โหติ,}
\csromangathagroup{\csromangathastanza{\csromangathaitembat{79}{สเจ อนฺเตวาสิโก นหายิตุกาโม โหติ,}{นหานํ ปฏิยาเทตพฺพํ.} \\ \csromangathacontbat{สเจ สีเตน อตฺโถ โหติ,}{สีตํ ปฏิยาเทตพฺพํ.} \\ \csromangathacontbat{สเจ อุเณฺหน อตฺโถ โหติ,}{อุณฺหํ ปฏิยาเทตพฺพํ.}}}

\prose{สเจ อนฺเตวาสิโก ชนฺตาฆรํ ปวิสิตุกาโม โหติ, จุณฺณํ สนฺเนตพฺพํ, มตฺติกา เตเมตพฺพา, ชนฺตาฆรปีฐํ อาทาย คนฺตฺวา ชนฺตาฆรปีฐํ ทตฺวา จีวรํ ปฏิคฺคเหตฺวา เอกมนฺตํ นิกฺขิปิตพฺพํ, จุณฺณํ ทาตพฺพํ, มตฺติกา ทาตพฺพา.}

\prose{สเจ อุสฺสหติ ชนฺตาฆรํ ปวิสิตพฺพํ, ชนฺตาฆรํ ปวิสนฺเตน มตฺติกาย มุขํ มกฺเขตฺวา ปุรโต จ ปจฺฉโต จ ปฏิจฺฉาเทตฺวา ชนฺตาฆรํ ปวิสิตพฺพํ, น จ เถเร ภิกฺขู อนุปขชฺช นิสีทิตพฺพํ, น นวา ภิกฺขู อาสเนน ปฏิพาหิตพฺพา, ชนฺตาฆเร อนฺเตวาสิกสฺส ปริกมฺมํ กาตพฺพํ, ชนฺตาฆรา นิกฺขมนฺเตน ชนฺตาฆรปีฐํ อาทาย ปุรโต จ ปจฺฉโต จ ปฏิจฺฉาเทตฺวา ชนฺตาฆรา นิกฺขมิตพฺพํ.}

\prose{อุทเกปิ อนฺเตวาสิกสฺส ปริกมฺมํ กาตพฺพํ, นหาเตน ปฐมตรํ อุตฺตริตฺวา อตฺตโน คตฺตํ โวทกํ กตฺวา นิวาเสตฺวา อนฺเตวาสิกสฺส คตฺตโต อุทกํ ปมชฺชิตพฺพํ, นิวาสนํ ทาตพฺพํ, สํฆาฏิ ทาตพฺพา, ชนฺตาฆรปีฐํ อาทาย ปฐมตรํ อาคนฺตฺวา อาสนํ ปญฺญเปตพฺพํ.}

\prose{ปาโททกํ ปาทปีฐํ ปาทกถลิกํ อุปนิกฺขิปิตพฺพํ. อนฺเตวาสิโก ปานีเยน ปุจฺฉิตพฺโพ.}

\proseitem{79}{ยสฺมิํ วิหาเร อนฺเตวาสิโก วิหรติ สเจ โส วิหาโร อุกฺลาโป โหติ สเจ อุสฺสหติ, โสเธตพฺโพ. วิหารํ โสเธนฺเตน ปฐมํ ปตฺตจีวรํ นีหริตฺวา เอกมนฺตํ นิกฺขิปิตพฺพํ, นิสีทน\-ปจฺจตฺถรณํ นีหริตฺวา เอกมนฺตํ นิกฺขิปิตพฺพํ, ภิสิพิพฺโพหนํ นีหริตฺวา เอกมนฺตํ \csromanfolio{84}นิกฺขิปิตพฺพํ, มญฺโจ นีจํ กตฺวา สาธุกํ อปฺปฏิฆํสนฺเตน อสํฆฏฺเฏนฺเตน กวาฏปิฏฺฐํ นีหริตฺวา เอกมนฺตํ นิกฺขิปิตพฺโพ, ปีฐํ นีจํ กตฺวา สาธุกํ อปฺปฏิฆํสนฺเตน อสํฆฏฺเฏนฺเตน กวาฏปิฏฺฐํ นีหริตฺวา เอกมนฺตํ นิกฺขิปิตพฺพํ. มญฺจ\-ปฏิปาทกา นีหริตฺวา เอกมนฺตํ นิกฺขิปิตพฺพา. เขฬมลฺลโก นีหริตฺวา เอกมนฺตํ นิกฺขิปิตพฺโพ. อปสฺเสนผลกํ นีหริตฺวา เอกมนฺตํ นิกฺขิปิตพฺพํ. ภูมตฺถรณํ ยถาปญฺญตฺตํ สลฺลกฺเขตฺวา นีหริตฺวา เอกมนฺตํ นิกฺขิปิตพฺพํ. สเจ วิหาเร สนฺตานกํ โหติ, อุลฺโลกา ปฐมํ โอตาเรตพฺพํ, อาโลกสนฺธิ\-กณฺณภาคา ปมชฺชิตพฺพา. สเจ เครุก\-ปริกมฺม\-กตา ภิตฺติ กณฺณกิตา โหติ, โจฬกํ เตเมตฺวา ปีเฬตฺวา ปมชฺชิตพฺพา. สเจ กาฬวณฺณกตา ภูมิ กณฺณกิตา โหติ, โจฬกํ เตเมตฺวา ปีเฬตฺวา ปมชฺชิตพฺพา. สเจ อกตา โหติ ภูมิ, อุทเกน ปริปฺโผสิตฺวา สมฺมชฺชิตพฺพา “มา วิหาโร รเชน อุหญฺญี”ติ, สงฺการํ วิจินิตฺวา เอกมนฺตํ ฉฑฺเฑตพฺพํ.}

\prose{ภูมตฺถรณํ โอตาเปตฺวา โสเธตฺวา ปปฺโผเฏตฺวา อติหริตฺวา ยถาปญฺญตฺตํ ปญฺญเปตพฺพํ. มญฺจ\-ปฏิปาทกา โอตาเปตฺวา ปมชฺชิตฺวา อติหริตฺวา ยถาฐาเน ฐเปตพฺพา, มญฺโจ โอตาเปตฺวา โสเธตฺวา ปปฺโผเฏตฺวา นีจํ กตฺวา สาธุกํ อปฺปฏิฆํสนฺเตน อสํฆฏฺเฏนฺเตน กวาฏปิฏฺฐํ อติหริตฺวา ยถาปญฺญตฺตํ ปญฺญเปตพฺโพ. ปีฐํ โอตาเปตฺวา โสเธตฺวา ปปฺโผเฏตฺวา นีจํ กตฺวา สาธุกํ อปฺปฏิฆํสนฺเตน อสํฆฏฺเฏนฺเตน กวาฏปิฏฺฐํ อติหริตฺวา ยถาปญฺญตฺตํ ปญฺญเปตพฺพํ. ภิสิพิพฺโพหนํ โอตาเปตฺวา โสเธตฺวา ปปฺโผเฏตฺวา อติหริตฺวา ยถาปญฺญตฺตํ ปญฺญเปตพฺพํ. นิสีทน\-ปจฺจตฺถรณํ โอตาเปตฺวา โสเธตฺวา ปปฺโผเฏตฺวา อติหริตฺวา ยถาปญฺญตฺตํ ปญฺญเปตพฺพํ. เขฬมลฺลโก โอตาเปตฺวา ปมชฺชิตฺวา อติหริตฺวา ยถาฐาเน ฐเปตพฺโพ. อปสฺเสนผลกํ โอตาเปตฺวา ปมชฺชิตฺวา อติหริตฺวา ยถาฐาเน ฐเปตพฺพํ. ปตฺตจีวรํ นิกฺขิปิตพฺพํ. ปตฺตํ นิกฺขิปนฺเตน เอเกน หตฺเถน ปตฺตํ คเหตฺวา เอเกน หตฺเถน เหฏฺฐามญฺจํ วา เหฏฺฐาปีฐํ วา ปรามสิตฺวา ปตฺโต นิกฺขิปิตพฺโพ, น จ อนนฺตรหิตาย ภูมิยา ปตฺโต นิกฺขิปิตพฺโพ. จีวรํ นิกฺขิปนฺเตน เอเกน หตฺเถน จีวรํ คเหตฺวา เอเกน หตฺเถน จีวรวํสํ วา จีวรรชฺชุํ วา ปมชฺชิตฺวา ปารโต อนฺตํ โอรโต โภคํ กตฺวา จีวรํ นิกฺขิปิตพฺพํ.}

\setlength{\csromangathapagewidth}{0pt}
\setlength{\csromangathawakwidth}{0pt}
\csromangathameasuregroup{สเจ ปริเวณํ อุกฺลาปํ โหติ, \\ สเจ โกฏฺฐโก อุกฺลาโป โหติ, \\ สเจ อคฺคิสาลา อุกฺลาปา โหติ, \\ สเจ วจฺจกุฏิ อุกฺลาปา โหติ, \\ สเจ ปานียํ น โหติ, \\ สเจ ปริโภชนียํ น โหติ,}{สเจ ปุรตฺถิมา สรชา วาตา วายนฺติ, \\ ปุรตฺถิมา วาตปานา ถเกตพฺพา. \\ สเจ ปจฺฉิมา สรชา วาตา วายนฺติ, \\ ปจฺฉิมา วาตปานา ถเกตพฺพา. \\ สเจ อุตฺตรา สรชา วาตา วายนฺติ, \\ อุตฺตรา วาตปานา ถเกตพฺพา. \\ สเจ ทกฺขิณา สรชา วาตา วายนฺติ, \\ ทกฺขิณา วาตปานา ถเกตพฺพา. \\ สเจ สีตกาโล โหติ, \\ ทิวา วาตปานา วิวริตพฺพา, รตฺติํ ถเกตพฺพา. \\ สเจ อุณฺหกาโล โหติ, \\ ทิวา วาตปานา ถเกตพฺพา, รตฺติํ วิวริตพฺพา. \\ \csromangathabat{สเจ ปริเวณํ อุกฺลาปํ โหติ,}{ปริเวณํ สมฺมชฺชิตพฺพํ.} \\ \csromangathabat{สเจ โกฏฺฐโก อุกฺลาโป โหติ,}{โกฏฺฐโก สมฺมชฺชิตพฺโพ.} \\ สเจ อุปฏฺฐานสาลา อุกฺลาปา โหติ, \\ อุปฏฺฐานสาลา สมฺมชฺชิตพฺพา. \\ \csromangathabat{สเจ อคฺคิสาลา อุกฺลาปา โหติ,}{อคฺคิสาลา สมฺมชฺชิตพฺพา.} \\ \csromangathabat{สเจ วจฺจกุฏิ อุกฺลาปา โหติ,}{วจฺจกุฏิ สมฺมชฺชิตพฺพา.} \\ \csromangathabat{สเจ ปานียํ น โหติ,}{ปานียํ อุปฏฺฐาเปตพฺพํ.} \\ \csromangathabat{สเจ ปริโภชนียํ น โหติ,}{ปริโภชนียํ อุปฏฺฐาเปตพฺพํ.}}
\csromangathasetleft{สเจ ปริเวณํ อุกฺลาปํ โหติ, \\ สเจ โกฏฺฐโก อุกฺลาโป โหติ, \\ สเจ อคฺคิสาลา อุกฺลาปา โหติ, \\ สเจ วจฺจกุฏิ อุกฺลาปา โหติ, \\ สเจ ปานียํ น โหติ, \\ สเจ ปริโภชนียํ น โหติ,}
\csromangathagroup{\csromangathastanza{\csromanfolio{85}สเจ ปุรตฺถิมา สรชา วาตา วายนฺติ, \\ ปุรตฺถิมา วาตปานา ถเกตพฺพา. \\ สเจ ปจฺฉิมา สรชา วาตา วายนฺติ, \\ ปจฺฉิมา วาตปานา ถเกตพฺพา.}\csromangathastanza{สเจ อุตฺตรา สรชา วาตา วายนฺติ, \\ อุตฺตรา วาตปานา ถเกตพฺพา. \\ สเจ ทกฺขิณา สรชา วาตา วายนฺติ, \\ ทกฺขิณา วาตปานา ถเกตพฺพา.}\csromangathastanza{สเจ สีตกาโล โหติ, \\ ทิวา วาตปานา วิวริตพฺพา, รตฺติํ ถเกตพฺพา. \\ สเจ อุณฺหกาโล โหติ, \\ ทิวา วาตปานา ถเกตพฺพา, รตฺติํ วิวริตพฺพา.}\csromangathastanza{\csromangathabat{สเจ ปริเวณํ อุกฺลาปํ โหติ,}{ปริเวณํ สมฺมชฺชิตพฺพํ.} \\ \csromangathabat{สเจ โกฏฺฐโก อุกฺลาโป โหติ,}{โกฏฺฐโก สมฺมชฺชิตพฺโพ.}}\csromangathastanza{สเจ อุปฏฺฐานสาลา อุกฺลาปา โหติ, \\ อุปฏฺฐานสาลา สมฺมชฺชิตพฺพา. \\ \csromangathabat{สเจ อคฺคิสาลา อุกฺลาปา โหติ,}{อคฺคิสาลา สมฺมชฺชิตพฺพา.}}\csromangathastanza{\csromangathabat{สเจ วจฺจกุฏิ อุกฺลาปา โหติ,}{วจฺจกุฏิ สมฺมชฺชิตพฺพา.} \\ \csromangathabat{สเจ ปานียํ น โหติ,}{ปานียํ อุปฏฺฐาเป\-ตพฺพํ.} \\ \csromangathabat{สเจ ปริโภชนียํ น โหติ,}{ปริโภชนียํ อุปฏฺฐาเป\-ตพฺพํ.}}}

\prose{สเจ อาจมนกุมฺภิยา อุทกํ น โหติ, อาจมนกุมฺภิยา อุทกํ อาสิญฺจิตพฺพํ.}

\proseitem{79}{สเจ อนฺเตวาสิกสฺส อนภิรติ อุปฺปนฺนา โหติ, อาจริเยน วูปกาเสตพฺโพ วูปกาสาเป\-ตพฺโพ ธมฺมกถา วาสฺส กาตพฺพา. สเจ อนฺเตวาสิกสฺส กุกฺกุจฺจํ อุปฺปนฺนํ โหติ, อาจริเยน วิโนเทตพฺพํ วิโนทาเปตพฺพํ ธมฺมกถา วาสฺส กาตพฺพา. สเจ อนฺเตวาสิกสฺส ทิฏฺฐิคตํ อุปฺปนฺนํ โหติ, อาจริเยน วิเวเจตพฺพํ วิเวจาเปตพฺพํ ธมฺมกถา วาสฺส กาตพฺพา. สเจ อนฺเตวาสิโก ครุธมฺมํ อชฺฌาปนฺโน โหติ ปริวาสารโห, อาจริเยน อุสฺสุกฺกํ กาตพฺพํ “กินฺติ นุ โข สํโฆ อนฺเตวาสิกสฺส ปริวาสํ ทเทยฺยา”ติ. สเจ อนฺเตวาสิโก มูลาย ปฏิกสฺสนา\-รโห โหติ, อาจริเยน อุสฺสุกฺกํ กาตพฺพํ “กินฺติ นุ โข สํโฆ อนฺเตวาสิกํ มูลาย ปฏิกสฺเสยฺยา”ติ. สเจ อนฺเตวาสิโก มานตฺตารโห โหติ, อาจริเยน อุสฺสุกฺกํ กาตพฺพํ “กินฺติ นุ โข สํโฆ อนฺเตวาสิกสฺส มานตฺตํ ทเทยฺยา”ติ. สเจ อนฺเตวาสิโก อพฺภานารโห โหติ, อาจริเยน อุสฺสุกฺกํ กาตพฺพํ “กินฺติ นุ โข สํโฆ อนฺเตวาสิกํ อพฺเภยฺยา”ติ. สเจ สํโฆ อนฺเตวาสิกสฺส กมฺมํ กตฺตุกาโม โหติ ตชฺชนียํ วา นิยสฺสํ วา ปพฺพาชนียํ วา ปฏิสารณียํ วา \csromanfolio{86}อุกฺเขปนียํ วา, อาจริเยน อุสฺสุกฺกํ กาตพฺพํ “กินฺติ นุ โข สํโฆ อนฺเตวาสิกสฺส กมฺมํ น กเรยฺย, ลหุกาย วา ปริณาเมยฺยา”ติ. กตํ วา ปนสฺส โหติ สํเฆน กมฺมํ ตชฺชนียํ วา นิยสฺสํ วา ปพฺพาชนียํ วา ปฏิสารณียํ วา อุกฺเขปนียํ วา, อาจริเยน อุสฺสุกฺกํ กาตพฺพํ “กินฺติ นุ โข อนฺเตวาสิโก สมฺมา วตฺเตยฺย, โลมํ ปาเตยฺย, เนตฺถารํ วตฺเตยฺย, สํโฆ ตํ กมฺมํ ปฏิปฺปสฺสมฺเภยฺยา”ติ.}

\prosewithcloser{สเจ อนฺเตวาสิกสฺส จีวรํ โธวิตพฺพํ โหติ, อาจริเยน อาจิกฺขิตพฺพํ “เอวํ โธเวยฺยาสี”ติ, อุสฺสุกฺกํ วา กาตพฺพํ “กินฺติ นุ โข อนฺเตวาสิกสฺส จีวรํ โธวิเยถา”ติ. สเจ อนฺเตวาสิกสฺส จีวรํ กาตพฺพํ โหติ, อาจริเยน อาจิกฺขิตพฺพํ “เอวํ กเรยฺยาสี”ติ, อุสฺสุกฺกํ วา กาตพฺพํ “กินฺติ นุ โข อนฺเตวาสิกสฺส จีวรํ กริเยถา”ติ. สเจ อนฺเตวาสิกสฺส รชนํ ปจิตพฺพํ โหติ, อาจริเยน อาจิกฺขิตพฺพํ “เอวํ ปเจยฺยาสี”ติ, อุสฺสุกฺกํ วา กาตพฺพํ “กินฺติ นุ โข อนฺเตวาสิกสฺส รชนํ ปจิเยถา”ติ. สเจ อนฺเตวาสิกสฺส จีวรํ รชิตพฺพํ โหติ, อาจริเยน อาจิกฺขิตพฺพํ “เอวํ รเชยฺยาสี”ติ, อุสฺสุกฺกํ วา กาตพฺพํ “กินฺติ นุ โข อนฺเตวาสิกสฺส จีวรํ รชิเยถา”ติ. จีวรํ รชนฺเตน สาธุกํ สมฺปริวตฺตกํ สมฺปริวตฺตกํ รชิตพฺพํ, น จ อจฺฉินฺเน เถเว ปกฺกมิตพฺพํ. สเจ อนฺเตวาสิโก คิลาโน โหติ, ยาวชีวํ อุปฏฺฐาตพฺโพ, วุฏฺฐานมสฺส อาคเมตพฺพนฺติ.}{อนฺเตวาสิก\-วตฺตํ นิฏฺฐิตํ.}
\nitthitamruled{ฉฏฺโฐ ภาณวาโร.}
\csromanmatikaanchor{mtk.21}
\csromantocmarkref{section}{20. ปณามนา ขมาปนา}{mtk.21}
\bookmark[dest={mtk.21},level=1]{20. ปณามนา ขมาปนา}
\csromanheader{1}{20. \textbf{ปณามนา ขมาปนา}}
\proseitem{80}{เตน โข ปน สมเยน อนฺเตวาสิกา อาจริเยสุ น สมฺมา วตฺตนฺติ ฯเปฯ. ภควโต เอตมตฺถํ อาโรเจสุํ ฯเปฯ. น ภิกฺขเว อนฺเตวาสิเกน อาจริยมฺหิ น สมฺมา วตฺติตพฺพํ, โย น สมฺมา วตฺเตยฺย, อาปตฺติ ทุกฺกฏสฺสาติ. เนว สมฺมา วตฺตนฺติ. ภควโต เอตมตฺถํ อาโรเจสุํ ฯเปฯ. อนุชานามิ ภิกฺขเว อสมฺมา\-วตฺตนฺตํ ปณาเมตุํ. เอวญฺจ ปน ภิกฺขเว ปณาเมตพฺโพ “ปณาเมมิ ตนฺ”ติ วา “มายิธ \csromanfolio{87}ปฏิกฺกมี”ติ วา “นีหร เต ปตฺตจีวรนฺ”ติ วา “นาหํ ตยา อุปฏฺฐาตพฺโพ”ติ วา กาเยน วิญฺญาเปติ, วาจาย วิญฺญาเปติ, กาเยน วาจาย วิญฺญาเปติ, ปณามิโต โหติ อนฺเตวาสิโก. น กาเยน วิญฺญาเปติ, น วาจาย วิญฺญาเปติ, น กาเยน วาจาย วิญฺญาเปติ, น ปณามิโต โหติ อนฺเตวาสิโกติ.}

\prose{เตน โข ปน สมเยน อนฺเตวาสิกา ปณามิตา น ขมาเปนฺติ. ภควโต เอตมตฺถํ อาโรเจสุํ. อนุชานามิ ภิกฺขเว ขมาเปตุนฺติ. เนว ขมาเปนฺติ. ภควโต เอตมตฺถํ อาโรเจสุํ. น ภิกฺขเว ปณามิเตน น ขมาเปตพฺโพ, โย น ขมาเปยฺย, อาปตฺติ ทุกฺกฏสฺสาติ.}

\prose{เตน โข ปน สมเยน อาจริยา ขมาปิยมานา น ขมนฺติ. ภควโต เอตมตฺถํ อาโรเจสุํ. อนุชานามิ ภิกฺขเว ขมิตุนฺติ. เนว ขมนฺติ. อนฺเตวาสิกา ปกฺกมนฺติปิ วิพฺภมนฺติปิ ติตฺถิเยสุปิ สงฺกมนฺติ. ภควโต เอตมตฺถํ อาโรเจสุํ. น ภิกฺขเว ขมาปิย\-มาเนน น ขมิตพฺพํ, โย น ขเมยฺย, อาปตฺติ ทุกฺกฏสฺสาติ.}

\prose{เตน โข ปน สมเยน อาจริยา สมฺมาวตฺตนฺตํ ปณาเมนฺติ, อสมฺมา\-วตฺตนฺตํ น ปณาเมนฺติ.}

\prose{ภควโต เอตมตฺถํ อาโรเจสุํ. น ภิกฺขเว สมฺมาวตฺตนฺโต ปณาเมตพฺโพ, โย ปณาเมยฺย, อาปตฺติ ทุกฺกฏสฺส. น จ ภิกฺขเว อสมฺมาวตฺตนฺโต น ปณาเมตพฺโพ, โย น ปณาเมยฺย, อาปตฺติ ทุกฺกฏสฺส.}

\proseitem{81}{ปญฺจหิ ภิกฺขเว องฺเคหิ สมนฺนาคโต อนฺเตวาสิโก ปณาเมตพฺโพ. อาจริยมฺหิ นาธิมตฺตํ เปมํ โหติ, นาธิมตฺโต ปสาโท โหติ, นาธิมตฺตา หิรี โหติ, นาธิมตฺโต คารโว โหติ, นาธิมตฺตา ภาวนา โหติ. อิเมหิ โข ภิกฺขเว ปญฺจหงฺเคหิ สมนฺนาคโต อนฺเตวาสิโก ปณาเมตพฺโพ.}

\prose{ปญฺจหิ ภิกฺขเว องฺเคหิ สมนฺนาคโต อนฺเตวาสิโก น ปณาเมตพฺโพ. อาจริยมฺหิ อธิมตฺตํ เปมํ โหติ, อธิมตฺโต ปสาโท โหติ, อธิมตฺตา หิรี โหติ, อธิมตฺโต คารโว โหติ, อธิมตฺตา ภาวนา โหติ. อิเมหิ โข ภิกฺขเว ปญฺจหงฺเคหิ สมนฺนาคโต อนฺเตวาสิโก น ปณาเมตพฺโพ.}

\csromanmatikaanchor{mtk.22}
\csromantocmarkref{section}{21. พาลอพฺยตฺตวตฺถุ}{mtk.22}
\bookmark[dest={mtk.22},level=1]{21. พาลอพฺยตฺตวตฺถุ}
\prose{\csromanfolio{88}ปญฺจหิ ภิกฺขเว องฺเคหิ สมนฺนาคโต อนฺเตวาสิโก อลํ ปณาเมตุํ. อาจริยมฺหิ นาธิมตฺตํ เปมํ โหติ, นาธิมตฺโต ปสาโท โหติ, นาธิมตฺตา หิรี โหติ, นาธิมตฺโต คารโว โหติ, นาธิมตฺตา ภาวนา โหติ. อิเมหิ โข ภิกฺขเว ปญฺจหงฺเคหิ สมนฺนาคโต อนฺเตวาสิโก อลํ ปณาเมตุํ.}

\prose{ปญฺจหิ ภิกฺขเว องฺเคหิ สมนฺนาคโต อนฺเตวาสิโก นาลํ ปณาเมตุํ. อาจริยมฺหิ อธิมตฺตํ เปมํ โหติ, อธิมตฺโต ปสาโท โหติ, อธิมตฺตา หิรี โหติ, อธิมตฺโต คารโว โหติ, อธิมตฺตา ภาวนา โหติ. อิเมหิ ภิกฺขเว ปญฺจหงฺเคหิ สมนฺนาคโต อนฺเตวาสิโก นาลํ ปณาเมตุํ.}

\prose{ปญฺจหิ ภิกฺขเว องฺเคหิ สมนฺนาคตํ อนฺเตวาสิกํ อปฺปณาเมนฺโต อาจริโย สาติสาโร โหติ, ปณาเมนฺโต อนติสาโร โหติ.}

\prose{อาจริยมฺหิ นาธิมตฺตํ เปมํ โหติ, นาธิมตฺโต ปสาโท โหติ, นาธิมตฺตา หิรี โหติ, นาธิมตฺโต คารโว โหติ, นาธิมตฺตา ภาวนา โหติ.}

\prose{อิเมหิ โข ภิกฺขเว ปญฺจหงฺเคหิ สมนฺนาคตํ อนฺเตวาสิกํ อปฺปณาเมนฺโต อาจริโย สาติสาโร โหติ, ปณาเมนฺโต อนติสาโร โหติ.}

\prose{ปญฺจหิ ภิกฺขเว องฺเคหิ สมนฺนาคตํ อนฺเตวาสิกํ ปณาเมนฺโต อาจริโย สาติสาโร โหติ, อปฺปณาเมนฺโต อนติสาโร โหติ.}

\prose{อาจริยมฺหิ อธิมตฺตํ เปมํ โหติ, อธิมตฺโต ปสาโท โหติ, อธิมตฺตา หิรี โหติ, อธิมตฺโต คารโว โหติ, อธิมตฺตา ภาวนา โหติ.}

\prose{อิเมหิ โข ภิกฺขเว ปญฺจหงฺเคหิ สมนฺนาคตํ อนฺเตวาสิกํ ปณาเมนฺโต อาจริโย สาติสาโร โหติ, อปฺปณาเมนฺโต อนติสาโร โหตีติ.}

\csromanheader{2}{21. พาล\-อพฺยตฺต\-วตฺถุ}
\proseitem{82}{เตน โข ปน สมเยน ภิกฺขู “ทสวสฺสมฺหา ทสวสฺสมฺหา”ติ พาลา อพฺยตฺตา นิสฺสยํ เทนฺติ. ทิสฺสนฺติ อาจริยา พาลา อนฺเตวาสิกา \csromanfolio{89}ปณฺฑิตา, ทิสฺสนฺติ อาจริยา อพฺยตฺตา อนฺเตวาสิกา พฺยตฺตา, ทิสฺสนฺติ อาจริยา อปฺปสฺสุตา อนฺเตวาสิกา พหุสฺสุตา, ทิสฺสนฺติ อาจริยา ทุปฺปญฺญา อนฺเตวาสิกา ปญฺญวนฺโต. เย เต ภิกฺขู อปฺปิจฺฉา ฯเปฯ เต อุชฺฌายนฺติ ขิยฺยนฺติ วิปาเจนฺติ “กถํ หิ นาม ภิกฺขู ‘ทสวสฺสมฺหา ทสวสฺสมฺหา’ติ พาลา อพฺยตฺตา นิสฺสยํ ทสฺสนฺติ. ทิสฺสนฺติ อาจริยา พาลา อนฺเตวาสิกา ปณฺฑิตา, ทิสฺสนฺติ อาจริยา อพฺยตฺตา อนฺเตวาสิกา พฺยตฺตา, ทิสฺสนฺติ อาจริยา อปฺปสฺสุตา อนฺเตวาสิกา พหุสฺสุตา, ทิสฺสนฺติ อาจริยา ทุปฺปญฺญา อนฺเตวาสิกา ปญฺญวนฺโต”ติ. อถ โข เต ภิกฺขู ภควโต เอตมตฺถํ อาโรเจสุํ ฯเปฯ. สจฺจํ กิร ภิกฺขเว ภิกฺขู “ทสวสฺสมฺหา ทสวสฺสมฺหา”ติ พาลา อพฺยตฺตา นิสฺสยํ เทนฺติ ฯเปฯ. สจฺจํ ภควาติ. วิครหิ พุทฺโธ ภควา ฯเปฯ วิครหิตฺวา ฯเปฯ ธมฺมิํ กถํ กตฺวา ภิกฺขู อามนฺเตสิ. น ภิกฺขเว พาเลน อพฺยตฺเตน นิสฺสโย ทาตพฺโพ, โย ทเทยฺย, อาปตฺติ ทุกฺกฏสฺส. อนุชานามิ ภิกฺขเว พฺยตฺเตน ภิกฺขุนา ปฏิพเลน ทสวสฺเสน วา อติเรก\-ทสวสฺเสน วา นิสฺสยํ ทาตุนฺติ.}

\csromanmatikaanchor{mtk.23}
\csromantocmarkref{section}{22. นิสฺสยปฏิปฺปสฺสทฺธิกถา}{mtk.23}
\bookmark[dest={mtk.23},level=1]{22. นิสฺสยปฏิปฺปสฺสทฺธิกถา}
\csromanheader{1}{22. \textbf{นิสฺสยปฏิปฺปสฺสทฺธิ}\-\textbf{กถา}}
\proseitem{83}{เตน โข ปน สมเยน ภิกฺขู อาจริยุ\-ปชฺฌาเยสุ ปกฺกนฺเตสุปิ วิพฺภนฺเตสุปิ กาลงฺกเตสุปิ ปกฺข\-สงฺกนฺเตสุปิ นิสฺสย\-ปฏิปฺปสฺสทฺธิโย น ชานนฺติ. ภควโต เอตมตฺถํ อาโรเจสุํ.}

\prose{ปญฺจิมา ภิกฺขเว นิสฺสย\-ปฏิปฺปสฺสทฺธิโย อุปชฺฌายมฺหา. อุปชฺฌาโย ปกฺกนฺโต วา โหติ, วิพฺภนฺโต วา, กาลงฺกโต วา, ปกฺขสงฺกนฺโต วา, อาณตฺติเยว ปญฺจมี. อิมา โข ภิกฺขเว ปญฺจ นิสฺสย\-ปฏิปฺปสฺสทฺธิโย อุปชฺฌายมฺหา.}

\prose{ฉยิมา ภิกฺขเว นิสฺสย\-ปฏิปฺปสฺสทฺธิโย อาจริยมฺหา. อาจริโย ปกฺกนฺโต วา โหติ, วิพฺภนฺโต วา, กาลงฺกโต วา, ปกฺขสงฺกนฺโต วา, อาณตฺติเยว ปญฺจมี อุปชฺฌาเยน วา สโมธานคโต โหติ. อิมา โข ภิกฺขเว ฉ นิสฺสย\-ปฏิปฺปสฺสทฺธิโย อาจริยมฺหา.}

\csromanmatikaanchor{mtk.24}
\csromantocmarkref{section}{23. อุปสมฺปาเทตพฺพปญฺจก}{mtk.24}
\bookmark[dest={mtk.24},level=1]{23. อุปสมฺปาเทตพฺพปญฺจก}
\csromanheader{1}{23. \textbf{อุปสมฺปาเทตพฺพ}\-\textbf{ปญฺจก}}
\proseitem{84}{\csromanfolio{90}ปญฺจหิ ภิกฺขเว องฺเคหิ สมนฺนาคเตน ภิกฺขุนา น อุปสมฺ\-ปาเท\-ตพฺพํ, น นิสฺสโย ทาตพฺโพ, น สามเณโร อุปฏฺฐาเปตพฺโพ. น อเสกฺเขน\csromansharedfootnote{1d69a17a81509fe1}{น อเสเขน (ก)} สีลกฺขนฺเธน สมนฺนาคโต โหติ, น อเสกฺเขน สมาธิ\-กฺขนฺเธน สมนฺนาคโต โหติ, น อเสกฺเขน ปญฺญา\-กฺขนฺเธน สมนฺนาคโต โหติ, น อเสกฺเขน วิมุตฺติ\-กฺขนฺเธน สมนฺนาคโต โหติ, น อเสกฺเขน วิมุตฺติ\-ญาณ\-ทสฺสนกฺขนฺเธน สมนฺนาคโต โหติ. อิเมหิ โข ภิกฺขเว ปญฺจหงฺเคหิ สมนฺนาคเตน ภิกฺขุนา น อุปสมฺ\-ปาเท\-ตพฺพํ, น นิสฺสโย ทาตพฺโพ, น สามเณโร อุปฏฺฐาเปตพฺโพ. (1)}

\prose{ปญฺจหิ ภิกฺขเว องฺเคหิ สมนฺนาคเตน ภิกฺขุนา อุปสมฺ\-ปาเท\-ตพฺพํ, นิสฺสโย ทาตพฺโพ, สามเณโร อุปฏฺฐาเปตพฺโพ. อเสกฺเขน สีลกฺขนฺเธน สมนฺนาคโต โหติ, อเสกฺเขน สมาธิ\-กฺขนฺเธน สมนฺนาคโต โหติ, อเสกฺเขน ปญฺญา\-กฺขนฺเธน สมนฺนาคโต โหติ, อเสกฺเขน วิมุตฺติ\-กฺขนฺเธน สมนฺนาคโต โหติ, อเสกฺเขน วิมุตฺติ\-ญาณ\-ทสฺสนกฺขนฺเธน สมนฺนาคโต โหติ. อิเมหิ โข ภิกฺขเว ปญฺจหงฺเคหิ สมนฺนาคเตน ภิกฺขุนา อุปสมฺ\-ปาเท\-ตพฺพํ, นิสฺสโย ทาตพฺโพ, สามเณโร อุปฏฺฐาเปตพฺโพ. (2)}

\prose{อปเรหิปิ ภิกฺขเว ปญฺจหงฺเคหิ สมนฺนาคเตน ภิกฺขุนา น อุปสมฺ\-ปาเท\-ตพฺพํ, น นิสฺสโย ทาตพฺโพ, น สามเณโร อุปฏฺฐาเปตพฺโพ. อตฺตนา น อเสกฺเขน สีลกฺขนฺเธน สมนฺนาคโต โหติ, น ปรํ อเสกฺเข สีลกฺขนฺเธ สมาทเปตา. อตฺตนา น อเสกฺเขน สมาธิ\-กฺขนฺเธน สมนฺนาคโต โหติ, น ปรํ อเสกฺเข สมธิกฺขนฺเธ สมาทเปตา. อตฺตนา น อเสกฺเขน ปญฺญา\-กฺขนฺเธน สมนฺนาคโต โหติ, น ปรํ อเสกฺเข ปญฺญากฺขนฺเธ สมาทเปตา. อตฺตนา น อเสกฺเขน วิมุตฺติ\-กฺขนฺเธน สมนฺนาคโต โหติ, น ปรํ อเสกฺเข วิมุตฺติ\-กฺขนฺเธ สมาทเปตา. อตฺตนา น อเสกฺเขน วิมุตฺติ\-ญาณ\-ทสฺสนกฺขนฺเธน สมนฺนาคโต โหติ, น ปรํ อเสกฺเข วิมุตฺติ\-ญาณ\-ทสฺสนกฺขนฺเธ สมาทเปตา. อิเมหิ โข ภิกฺขเว ปญฺจหงฺเคหิ สมนฺนาคเตน ภิกฺขุนา น อุปสมฺ\-ปาเท\-ตพฺพํ, น นิสฺสโย ทาตพฺโพ, น สามเณโร อุปฏฺฐาเปตพฺโพ. (3)}

\prose{\csromanfolio{91}ปญฺจหิ ภิกฺขเว องฺเคหิ สมนฺนาคเตน ภิกฺขุนา อุปสมฺ\-ปาเท\-ตพฺพํ, นิสฺสโย ทาตพฺโพ, สามเณโร อุปฏฺฐาเปตพฺโพ. อตฺตนา อเสกฺเขน สีลกฺขนฺเธน สมนฺนาคโต โหติ, ปรํ อเสกฺเข สีลกฺขนฺเธ สมาทเปตา, อตฺตนา อเสกฺเขน สมาธิ\-กฺขนฺเธน สมนฺนาคโต โหติ, ปรํ อเสกฺเข สมาธิ\-กฺขนฺเธ สมาทเปตา. อตฺตนา อเสกฺเขน ปญฺญา\-กฺขนฺเธน สมนฺนาคโต โหติ, ปรํ อเสกฺเข ปญฺญากฺขนฺเธ สมาทเปตา. อตฺตนา อเสกฺเขน วิมุตฺติ\-กฺขนฺเธน สมนฺนาคโต โหติ, ปรํ อเสกฺเข วิมุตฺติ\-กฺขนฺเธ สมาทเปตา. อตฺตนา อเสกฺเขน วิมุตฺติ\-ญาณ\-ทสฺสนกฺขนฺเธน สมนฺนาคโต โหติ, ปรํ อเสกฺเข วิมุตฺติ\-ญาณ\-ทสฺสนกฺขนฺเธ สมาทเปตา. อิเมหิ โข ภิกฺขเว ปญฺจหงฺเคหิ สมนฺนาคเตน ภิกฺขุนา อุปสมฺ\-ปาเท\-ตพฺพํ, นิสฺสโย ทาตพฺโพ, สามเณโร อุปฏฺฐาเปตพฺโพ. (4)}

\prose{อปเรหิปิ ภิกฺขเว ปญฺจหงฺเคหิ สมนฺนาคเตน ภิกฺขุนา น อุปสมฺ\-ปาเท\-ตพฺพํ, น นิสฺสโย ทาตพฺโพ, น สามเณโร อุปฏฺฐาเปตพฺโพ. อสฺสทฺโธ โหติ, อหิริโก โหติ, อโนตฺตปฺปี โหติ, กุสีโต โหติ, มุฏฺฐสฺสติ โหติ. อิเมหิ โข ภิกฺขเว ปญฺจงฺเคหิ สมนฺนาคเตน ภิกฺขุนา น อุปสมฺ\-ปาเท\-ตพฺพํ, น นิสฺสโย ทาตพฺโพ, น สามเณโร อุปฏฺฐาเปตพฺโพ. (5)}

\prose{ปญฺจหิ ภิกฺขเว องฺเคหิ สมนฺนาคเตน ภิกฺขุนา อุปสมฺ\-ปาเท\-ตพฺพํ, นิสฺสโย ทาตพฺโพ, สามเณโร อุปฏฺฐาเปตพฺโพ. สทฺโธ โหติ, หิริมา โหติ, โอตฺตปฺปี โหติ, อารทฺธวีริโย โหติ, อุปฏฺฐิตสฺสติ โหติ. อิเมหิ โข ภิกฺขเว ปญฺจหงฺเคหิ สมนฺนาคเตน ภิกฺขุนา อุปสมฺ\-ปาเท\-ตพฺพํ, นิสฺสโย ทาตพฺโพ, สามเณโร อุปฏฺฐาเปตพฺโพ. (6)}

\prose{อปเรหิปิ ภิกฺขเว ปญฺจหงฺเคหิ สมนฺนาคเตน ภิกฺขุนา น อุปสมฺ\-ปาเท\-ตพฺพํ, น นิสฺสโย ทาตพฺโพ, น สามเณโร อุปฏฺฐาเปตพฺโพ. อธิสีเล สีลวิปนฺโน โหติ, อชฺฌาจาเร อาจารวิปนฺโน โหติ, อติทิฏฺฐิยา ทิฏฺฐิวิปนฺโน โหติ, อปฺปสฺสุโต โหติ, ทุปฺปญฺโญ โหติ. อิเมหิ โข ภิกฺขเว ปญฺจหงฺเคหิ สมนฺนาคเตน ภิกฺขุนา น อุปสมฺ\-ปาเท\-ตพฺพํ, น นิสฺสโย ทาตพฺโพ, น สามเณโร อุปฏฺฐาเปตพฺโพ. (7)}

\prose{\csromanfolio{92}ปญฺจหิ ภิกฺขเว องฺเคหิ สมนฺนาคเตน ภิกฺขุนา อุปสมฺ\-ปาเท\-ตพฺพํ, นิสฺสโย ทาตพฺโพ, สามเณโร อุปฏฺฐาเปตพฺโพ. น อธิสีเล สีลวิปนฺโน โหติ, น อชฺฌาจาเร อาจารวิปนฺโน โหติ, น อติทิฏฺฐิยา ทิฏฺฐิวิปนฺโน โหติ, พหุสฺสุโต โหติ, ปญฺญวา โหติ. อิเมหิ โข ภิกฺขเว ปญฺจหงฺเคหิ สมนฺนาคเตน ภิกฺขุนา อุปสมฺ\-ปาเท\-ตพฺพํ, นิสฺสโย ทาตพฺโพ, สามเณโร อุปฏฺฐาเปตพฺโพ. (8)}

\prose{อปเรหิปิ ภิกฺขเว ปญฺจหงฺเคหิ สมนฺนาคเตน ภิกฺขุนา น อุปสมฺ\-ปาเท\-ตพฺพํ, น นิสฺสโย ทาตพฺโพ, น สามเณโร อุปฏฺฐาเปตพฺโพ. น ปฏิพโล โหติ อนฺเตวาสิํ วา สทฺธิวิหาริํ วา คิลานํ อุปฏฺฐาตุํ วา อุปฏฺฐาเปตุํ วา, อนภิรตํ\csromansharedfootnote{3899dab7c7218692}{อนภิรติํ (สฺยา), อุปฺปนฺนํ อนภิรติํ (ก)} วูปกาเสตุํ วา วูปกาสาเปตุํ วา, อุปฺปนฺนํ กุกฺกุจฺจํ ธมฺมโต วิโนเทตุํ\csromansharedfootnote{df2201f6e8f610f7}{วิโนเทตุํ วา วิโนทาเปตุํ วา (สพฺพตฺถ, วิมติวิโนทนีฏีกา โอโลเกตพฺพา)}, อาปตฺติํ น ชานาติ, อาปตฺติยา วุฏฺฐานํ น ชานาติ. อิเมหิ โข ภิกฺขเว ปญฺจหงฺเคหิ สมนฺนาคเตน ภิกฺขุนา น อุปสมฺ\-ปาเท\-ตพฺพํ, น นิสฺสโย ทาตพฺโพ, น สามเณโร อุปฏฺฐาเปตพฺโพ. (9)}

\prose{ปญฺจหิ ภิกฺขเว องฺเคหิ สมนฺนาคเตน ภิกฺขุนา อุปสมฺ\-ปาเท\-ตพฺพํ, นิสฺสโย ทาตพฺโพ, สามเณโร อุปฏฺฐาเปตพฺโพ. ปฏิพโล โหติ อนฺเตวาสิํ วา สทฺธิวิหาริํ วา คิลานํ อุปฏฺฐาตุํ วา อุปฏฺฐาเปตุํ วา, อนภิรตํ วูปกาเสตุํ วา วูปกาสาเปตุํ วา, อุปฺปนฺนํ กุกฺกุจฺจํ ธมฺมโต วิโนเทตุํ, อาปตฺติํ ชานาติ, อาปตฺติยา วุฏฺฐานํ ชานาติ. อิเมหิ โข ภิกฺขเว ปญฺจหงฺเคหิ สมนฺนาคเตน ภิกฺขุนา อุปสมฺ\-ปาเท\-ตพฺพํ, นิสฺสโย ทาตพฺโพ, สามเณโร อุปฏฺฐาเปตพฺโพ. (10)}

\prose{อปเรหิปิ ภิกฺขเว ปญฺจหงฺเคหิ สมนฺนาคเตน ภิกฺขุนา น อุปสมฺ\-ปาเท\-ตพฺพํ, น นิสฺสโย ทาตพฺโพ, น สามเณโร อุปฏฺฐาเปตพฺโพ. น ปฏิพโล โหติ อนฺเตวาสิํ วา สทฺธิวิหาริํ วา อภิสมาจาริกาย สิกฺขาย สิกฺขาเปตุํ, อาทิ\-พฺรหฺม\-จริย\-กาย สิกฺขาย วิเนตุํ, อภิธมฺเม วิเนตุํ, อภิวินเย วิเนตุํ, อุปฺปนฺนํ ทิฏฺฐิคตํ ธมฺมโต วิเวเจตุํ. อิเมหิ โข ภิกฺขเว ปญฺจหงฺเคหิ สมนฺนาคเตน ภิกฺขุนา น อุปสมฺ\-ปาเท\-ตพฺพํ, น นิสฺสโย ทาตพฺโพ, น สามเณโร อุปฏฺฐาเปตพฺโพ. (11)}

\prose{\csromanfolio{93}ปญฺจหิ ภิกฺขเว องฺเคหิ สมนฺนาคเตน ภิกฺขุนา อุปสมฺ\-ปาเท\-ตพฺพํ, นิสฺสโย ทาตพฺโพ, สามเณโร อุปฏฺฐาเปตพฺโพ. ปฏิพโล โหติ อนฺเตวาสิํ วา สทฺธิวิหาริํ วา อภิสมาจาริกาย สิกฺขาย สิกฺขาเปตุํ, อาทิ\-พฺรหฺม\-จริย\-กาย สิกฺขาย วิเนตุํ, อภิธมฺเม วิเนตุํ, อภิวินเย วิเนตุํ, อุปฺปนฺนํ ทิฏฺฐิคตํ ธมฺมโต วิเวเจตุํ. อิเมหิ โข ภิกฺขเว ปญฺจหงฺเคหิ สมนฺนาคเตน ภิกฺขุนา อุปสมฺ\-ปาเท\-ตพฺพํ, นิสฺสโย ทาตพฺโพ, สามเณโร อุปฏฺฐาเปตพฺโพ. (12)}

\prose{อปเรหิปิ ภิกฺขเว ปญฺจหงฺเคหิ สมนฺนาคเตน ภิกฺขุนา น อุปสมฺ\-ปาเท\-ตพฺพํ, น นิสฺสโย ทาตพฺโพ, น สามเณโร อุปฏฺฐาเปตพฺโพ. อาปตฺติํ น ชานาติ, อนาปตฺติํ น ชานาติ, ลหุกํ อาปตฺติํ น ชานาติ, ครุกํ อาปตฺติํ น ชานาติ, อุภยานิ โข ปนสฺส ปาติโมกฺขานิ วิตฺถาเรน น สฺวาคตานิ โหนฺติ น สุวิภตฺตานิ น สุปฺปวตฺตีนิ น สุวินิจฺฉิตานิ สุตฺตโส อนุพฺยญฺชนโส. อิเมหิ โข ภิกฺขเว ปญฺจหงฺเคหิ สมนฺนาคเตน ภิกฺขุนา น อุปสมฺ\-ปาเท\-ตพฺพํ, น นิสฺสโย ทาตพฺโพ, น สามเณโร อุปฏฺฐาเปตพฺโพ. (13)}

\prose{ปญฺจหิ ภิกฺขเว องฺเคหิ สมนฺนาคเตน ภิกฺขุนา อุปสมฺ\-ปาเท\-ตพฺพํ, นิสฺสโย ทาตพฺโพ, สามเณโร อุปฏฺฐาเปตพฺโพ. อาปตฺติํ ชานาติ, อนาปตฺติํ ชานาติ, ลหุกํ อาปตฺติํ ชานาติ, ครุกํ อาปตฺติํ ชานาติ, อุภยานิ โข ปนสฺส ปาติโมกฺขานิ วิตฺถาเรน สฺวาคตานิ โหนฺติ สุวิภตฺตานิ สุปฺปวตฺตีนิ สุวินิจฺฉิตานิ สุตฺตโส อนุพฺยญฺชนโส. อิเมหิ โข ภิกฺขเว ปญฺจหงฺเคหิ สมนฺนาคเตน ภิกฺขุนา อุปสมฺ\-ปาเท\-ตพฺพํ, นิสฺสโย ทาตพฺโพ, สามเณโร อุปฏฺฐาเปตพฺโพ. (14)}

\prose{อปเรหิปิ ภิกฺขเว ปญฺจหงฺเคหิ สมนฺนาคเตน ภิกฺขุนา น อุปสมฺ\-ปาเท\-ตพฺพํ, น นิสฺสโย ทาตพฺโพ, น สามเณโร อุปฏฺฐาเปตพฺโพ. อาปตฺติํ น ชานาติ, อนาปตฺติํ น ชานาติ, ลหุกํ อาปตฺติํ น ชานาติ, ครุกํ อาปตฺติํ น ชานาติ, อูนทสวสฺโส โหติ. อิเมหิ โข ภิกฺขเว ปญฺจหงฺเคหิ สมนฺนาคเตน ภิกฺขุนา น อุปสมฺ\-ปาเท\-ตพฺพํ, น นิสฺสโย ทาตพฺโพ, น สามเณโร อุปฏฺฐาเปตพฺโพ. (15)}

\csromanmatikaanchor{mtk.25}
\csromantocmarkref{section}{24. อุปสมฺปาเทตพฺพฉกฺก}{mtk.25}
\bookmark[dest={mtk.25},level=1]{24. อุปสมฺปาเทตพฺพฉกฺก}
\prosewithcloserruled{ปญฺจหิ ภิกฺขเว องฺเคหิ สมนฺนาคเตน ภิกฺขุนา อุปสมฺ\-ปาเท\-ตพฺพํ, นิสฺสโย ทาตพฺโพ, สามเณโร อุปฏฺฐาเปตพฺโพ. อาปตฺติํ ชานาติ, \csromanfolio{94}อนาปตฺติํ ชานาติ, ลหุกํ อาปตฺติํ ชานาติ, ครุกํ อาปตฺติํ ชานาติ, ทสวสฺโส วา โหติ อติเรก\-ทสวสฺโส วา. อิเมหิ โข ภิกฺขเว ปญฺจหงฺเคหิ สมนฺนาคเตน ภิกฺขุนา อุปสมฺ\-ปาเท\-ตพฺพํ, นิสฺสโย ทาตพฺโพ, สามเณโร อุปฏฺฐาเป\-ตพฺโพติ. (16)}{อุปสมฺปาเทตพฺพ\-ปญฺจก\-โสฬส\-วาโร นิฏฺฐิโต.}
\csromancenter{\textbf{อุปสมฺปาเทตพฺพ}\-\textbf{ฉกฺก} 85. ฉหิ ภิกฺขเว องฺเคหิ สมนฺนาคเตน ภิกฺขุนา น อุปสมฺ\-ปาเท\-ตพฺพํ,น นิสฺสโย ทาตพฺโพ, น สามเณโร อุปฏฺฐาเปตพฺโพ. น อเสกฺเขน สีลกฺขนฺเธน สมนฺนาคโต โหติ, น อเสกฺเขน สมาธิ\-กฺขนฺเธน สมนฺนาคโต โหติ, น อเสกฺเขน ปญฺญา\-กฺขนฺเธน สมนฺนาคโต โหติ, น อเสกฺเขน วิมุตฺติ\-กฺขนฺเธน สมนฺนาคโต โหติ, น อเสกฺเขน วิมุตฺติ\-ญาณ\-ทสฺสนกฺขนฺเธน สมนฺนาคโต โหติ, อูนทสวสฺโส โหติ. อิเมหิ โข ภิกฺขเว ฉหงฺเคหิ สมนฺนาคเตน ภิกฺขุนา น อุปสมฺ\-ปาเท\-ตพฺพํ, น นิสฺสโย ทาตพฺโพ, น สามเณโร อุปฏฺฐาเปตพฺโพ. (1)}
\prose{ฉหิ ภิกฺขเว องฺเคหิ สมนฺนาคเตน ภิกฺขุนา อุปสมฺ\-ปาเท\-ตพฺพํ, นิสฺสโย ทาตพฺโพ, สามเณโร อุปฏฺฐาเปตพฺโพ. อเสกฺเขน สีลกฺขนฺเธน สมนฺนาคโต โหติ, อเสกฺเขน สมาธิ\-กฺขนฺเธน สมนฺนาคโต โหติ, อเสกฺเขน ปญฺญา\-กฺขนฺเธน สมนฺนาคโต โหติ, อเสกฺเขน วิมุตฺติ\-กฺขนฺเธน สมนฺนาคโต โหติ, อเสกฺเขน วิมุตฺติ\-ญาณ\-ทสฺสนกฺขนฺเธน สมนฺนาคโต โหติ, ทสวสฺโส วา โหติ อติเรก\-ทสวสฺโส วา. อิเมหิ โข ภิกฺขเว ฉหงฺเคหิ สมนฺนาคเตน ภิกฺขุนา อุปสมฺ\-ปาเท\-ตพฺพํ, นิสฺสโย ทาตพฺโพ, สามเณโร อุปฏฺฐาเปตพฺโพ. (2)}

\prose{อปเรหิปิ ภิกฺขเว ฉหงฺเคหิ สมนฺนาคเตน ภิกฺขุนา น อุปสมฺ\-ปาเท\-ตพฺพํ, น นิสฺสโย ทาตพฺโพ, น สามเณโร อุปฏฺฐาเปตพฺโพ. อตฺตนา น อเสกฺเขน สีลกฺขนฺเธน สมนฺนาคโต โหติ, น ปรํ อเสกฺเข สีลกฺขนฺเธ สมาทเปตา. อตฺตนา น อเสกฺเขน สมาธิ\-กฺขนฺเธน \csromanfolio{95}สมนฺนาคโต โหติ, น ปรํ อเสกฺเข สมาธิ\-กฺขนฺเธ สมาทเปตา. อตฺตนา น อเสกฺเขน ปญฺญา\-กฺขนฺเธน สมนฺนาคโต โหติ, น ปรํ อเสกฺเข ปญฺญากฺขนฺเธ สมาทเปตา. อตฺตนา น อเสกฺเขน วิมุตฺติ\-กฺขนฺเธน สมนฺนาคโต โหติ, น ปรํ อเสกฺเข วิมุตฺติ\-กฺขนฺเธ สมาทเปตา. อตฺตนา น อเสกฺเขน วิมุตฺติ\-ญาณ\-ทสฺสนกฺขนฺเธน สมนฺนาคโต โหติ, น ปรํ อเสกฺเข วิมุตฺติ\-ญาณ\-ทสฺสนกฺขนฺเธ สมาทเปตา. อูนทสวสฺโส โหติ. อิเมหิ โข ภิกฺขเว ฉหงฺเคหิ สมนฺนาคเตน ภิกฺขุนา น อุปสมฺ\-ปาเท\-ตพฺพํ, น นิสฺสโย ทาตพฺโพ, น สามเณโร อุปฏฺฐาเปตพฺโพ. (3)}

\prose{ฉหิ ภิกฺขเว องฺเคหิ สมนฺนาคเตน ภิกฺขุนา อุปสมฺ\-ปาเท\-ตพฺพํ, นิสฺสโย ทาตพฺโพ, สามเณโร อุปฏฺฐาเปตพฺโพ. อตฺตนา อเสกฺเขน สีลกฺขนฺเธน สมนฺนาคโต โหติ, ปรํ อเสกฺเข สีลกฺขนฺเธ สมาทเปตา. อตฺตนา อเสกฺเขน สมาธิ\-กฺขนฺเธน สมนฺนาคโต โหติ, ปรํ อเสกฺเข สมาธิ\-กฺขนฺเธ สมาทเปตา. อตฺตนา อเสกฺเขน ปญฺญา\-กฺขนฺเธน สมนฺนาคโต โหติ, ปรํ อเสกฺเข ปญฺญากฺขนฺเธ สมาทเปตา. อตฺตนา อเสกฺเขน วิมุตฺติ\-กฺขนฺเธน สมนฺนาคโต โหติ, ปรํ อเสกฺเข วิมุตฺติ\-กฺขนฺเธ สมาทเปตา. อตฺตนา อเสกฺเขน วิมุตฺติ\-ญาณ\-ทสฺสนกฺขนฺเธน สมนฺนาคโต โหติ, ปรํ อเสกฺเข วิมุตฺติ\-ญาณ\-ทสฺสนกฺขนฺเธ สมาทเปตา. ทสวสฺโส วา โหติ อติเรก\-ทสวสฺโส วา. อิเมหิ โข ภิกฺขเว ฉหงฺเคหิ สมนฺนาคเตน ภิกฺขุนา อุปสมฺ\-ปาเท\-ตพฺพํ, นิสฺสโย ทาตพฺโพ, สามเณโร อุปฏฺฐาเปตพฺโพ. (4)}

\prose{อปเรหิปิ ภิกฺขเว ฉหงฺเคหิ สมนฺนาคเตน ภิกฺขุนา น อุปสมฺ\-ปาเท\-ตพฺพํ, น นิสฺสโย ทาตพฺโพ, น สามเณโร อุปฏฺฐาเปตพฺโพ. อสฺสทฺโธ โหติ, อหิริโก โหติ, อโนตฺตปฺปี โหติ, กุสีโต โหติ, มุฏฺฐสฺสติ โหติ, อูนทสวสฺโส โหติ. อิเมหิ โข ภิกฺขเว ฉหงฺเคหิ สมนฺนาคเตน ภิกฺขุนา น อุปสมฺ\-ปาเท\-ตพฺพํ, น นิสฺสโย ทาตพฺโพ, น สามเณโร อุปฏฺฐาเปตพฺโพ. (5)}

\prose{ฉหิ ภิกฺขเว องฺเคหิ สมนฺนาคเตน ภิกฺขุนา อุปสมฺ\-ปาเท\-ตพฺพํ, นิสฺสโย ทาตพฺโพ, สามเณโร อุปฏฺฐาเปตพฺโพ. สทฺโธ โหติ, หิริมา โหติ, โอตฺตปฺปี โหติ, อารทฺธวีริโย โหติ, อุปฏฺฐิตสฺสติ โหติ, ทสวสฺโส วา โหติ อติเรก\-ทสวสฺโส วา. \csromanfolio{96}อิเมหิ โข ภิกฺขเว ฉหงฺเคหิ สมนฺนาคเตน ภิกฺขุนา อุปสมฺ\-ปาเท\-ตพฺพํ, นิสฺสโย ทาตพฺโพ, สามเณโร อุปฏฺฐาเปตพฺโพ. (6)}

\prose{อปเรหิปิ ภิกฺขเว ฉหงฺเคหิ สมนฺนาคเตน ภิกฺขุนา น อุปสมฺ\-ปาเท\-ตพฺพํ, น นิสฺสโย ทาตพฺโพ, น สามเณโร อุปฏฺฐาเปตพฺโพ. อธิสีเล สีลวิปนฺโน โหติ, อชฺฌาจาเร อาจารวิปนฺโน โหติ, อติทิฏฺฐิยา ทิฏฺฐิวิปนฺโน โหติ, อปฺปสฺสุโต โหติ, ทุปฺปญฺโญ โหติ, อูนทสวสฺโส โหติ. อิเมหิ โข ภิกฺขเว ฉหงฺเคหิ สมนฺนาคเตน ภิกฺขุนา น อุปสมฺ\-ปาเท\-ตพฺพํ, น นิสฺสโย ทาตพฺโพ, น สามเณโร อุปฏฺฐาเปตพฺโพ. (7)}

\prose{ฉหิ ภิกฺขเว องฺเคหิ สมนฺนาคเตน ภิกฺขุนา อุปสมฺ\-ปาเท\-ตพฺพํ, นิสฺสโย ทาตพฺโพ, สามเณโร อุปฏฺฐาเปตพฺโพ. น อธิสีเล สีลวิปนฺโน โหติ, น อชฺฌาจาเร อาจารวิปนฺโน โหติ, น อติทิฏฺฐิยา ทิฏฺฐิวิปนฺโน โหติ, พหุสฺสุโต โหติ, ปญฺญวา โหติ, ทสวสฺโส วา โหติ อติเรก\-ทสวสฺโส วา. อิเมหิ โข ภิกฺขเว ฉหงฺเคหิ สมนฺนาคเตน ภิกฺขุนา อุปสมฺ\-ปาเท\-ตพฺพํ, นิสฺสโย ทาตพฺโพ, สามเณโร อุปฏฺฐาเปตพฺโพ. (8)}

\prose{อปเรหิปิ ภิกฺขเว ฉหงฺเคหิ สมนฺนาคเตน ภิกฺขุนา น อุปสมฺ\-ปาเท\-ตพฺพํ, น นิสฺสโย ทาตพฺโพ, น สามเณโร อุปฏฺฐาเปตพฺโพ. น ปฏิพโล โหติ อนฺเตวาสิํ วา สทฺธิวิหาริํ วา คิลานํ อุปฏฺฐาตุํ วา อุปฏฺฐาเปตุํ วา, อนภิรตํ วูปกาเสตุํ วา วูปกาสาเปตุํ วา, อุปฺปนฺนํ กุกฺกุจฺจํ ธมฺมโต วิโนเทตุํ, อาปตฺติํ น ชานาติ, อาปตฺติยา วุฏฺฐานํ น ชานาติ, อูนทสวสฺโส โหติ. อิเมหิ โข ภิกฺขเว ฉหงฺเคหิ สมนฺนาคเตน ภิกฺขุนา น อุปสมฺ\-ปาเท\-ตพฺพํ, น นิสฺสโย ทาตพฺโพ, น สามเณโร อุปฏฺฐาเปตพฺโพ. (9)}

\prose{ฉหิ ภิกฺขเว องฺเคหิ สมนฺนาคเตน ภิกฺขุนา อุปสมฺ\-ปาเท\-ตพฺพํ, นิสฺสโย ทาตพฺโพ, สามเณโร อุปฏฺฐาเปตพฺโพ. ปฏิพโล โหติ อนฺเตวาสิํ วา สทฺธิวิหาริํ วา คิลานํ อุปฏฺฐาตุํ วา อุปฏฺฐาเปตุํ วา, อนภิรตํ วูปกาเสตุํ วา วูปกาสาเปตุํ วา, อุปฺปนฺนํ กุกฺกุจฺจํ ธมฺมโต วิโนเทตุํ, อาปตฺติํ ชานาติ, อาปตฺติยา วุฏฺฐานํ ชานาติ, \csromanfolio{97}ทสวสฺโส วา โหติ อติเรก\-ทสวสฺโส วา. อิเมหิ โข ภิกฺขเว ฉหงฺเคหิ สมนฺนาคเตน ภิกฺขุนา อุปสมฺ\-ปาเท\-ตพฺพํ, นิสฺสโย ทาตพฺโพ, สามเณโร อุปฏฺฐาเปตพฺโพ. (10)}

\prose{อปเรหิปิ ภิกฺขเว ฉหงฺเคหิ สมนฺนาคเตน ภิกฺขุนา น อุปสมฺ\-ปาเท\-ตพฺพํ, น นิสฺสโย ทาตพฺโพ, น สามเณโร อุปฏฺฐาเปตพฺโพ. น ปฏิพโล โหติ อนฺเตวาสิํ วา สทฺธิวิหาริํ วา อภิสมาจาริกาย สิกฺขาย สิกฺขาเปตุํ, อาทิ\-พฺรหฺม\-จริย\-กาย สิกฺขาย วิเนตุํ, อภิธมฺเม วิเนตุํ, อภิวินเย วิเนตุํ, อุปฺปนฺนํ ทิฏฺฐิคตํ ธมฺมโต วิเวเจตุํ, อูนทสวสฺโส โหติ. อิเมหิ โข ภิกฺขเว ฉหงฺเคหิ สมนฺนาคเตน ภิกฺขุนา น อุปสมฺ\-ปาเท\-ตพฺพํ, น นิสฺสโย ทาตพฺโพ, น สามเณโร อุปฏฺฐาเปตพฺโพ. (11)}

\prose{ฉหิ ภิกฺขเว องฺเคหิ สมนฺนาคเตน ภิกฺขุนา อุปสมฺ\-ปาเท\-ตพฺพํ, นิสฺสโย ทาตพฺโพ, สามเณโร อุปฏฺฐาเปตพฺโพ, ปฏิพโล โหติ อนฺเตวาสิํ วา สทฺธิวิหาริํ วา อภิสมาจาริกาย สิกฺขาย สิกฺขาเปตุํ, อาทิ\-พฺรหฺม\-จริย\-กาย สิกฺขาย วิเนตุํ, อภิธมฺเม วิเนตุํ, อภิวินเย วิเนตุํ, อุปฺปนฺนํ ทิฏฺฐิคตํ ธมฺมโต วิเวเจตุํ, ทสวสฺโส วา โหติ อติเรก\-ทสวสฺโส วา. อิเมหิ โข ภิกฺขเว ฉหงฺเคหิ สมนฺนาคเตน ภิกฺขุนา อุปสมฺ\-ปาเท\-ตพฺพํ, นิสฺสโย ทาตพฺโพ, สามเณโร อุปฏฺฐาเปตพฺโพ. (12)}

\prose{อปเรหิปิ ภิกฺขเว ฉหงฺเคหิ สมนฺนาคเตน ภิกฺขุนา น อุปสมฺ\-ปาเท\-ตพฺพํ, น นิสฺสโย ทาตพฺโพ, น สามเณโร อุปฏฺฐาเปตพฺโพ. อาปตฺติํ น ชานาติ, อนาปตฺติํ น ชานาติ, ลหุกํ อาปตฺติํ น ชานาติ, ครุกํ อาปตฺติํ น ชานาติ, อุภยานิ โข ปนสฺส ปาติโมกฺขานิ วิตฺถาเรน น สฺวาคตานิ โหนฺติ น สุวิภตฺตานิ น สุปฺปวตฺตีนิ น สุวินิจฺฉิตานิ สุตฺตโส อนุพฺยญฺชนโส, อูนทสวสฺโส โหติ. อิเมหิ โข ภิกฺขเว ฉหงฺเคหิ สมนฺนาคเตน ภิกฺขุนา น อุปสมฺ\-ปาเท\-ตพฺพํ, น นิสฺสโย ทาตพฺโพ, น สามเณโร อุปฏฺฐาเปตพฺโพ. (13)}

\prosewithcloserruled{ฉหิ ภิกฺขเว องฺเคหิ สมนฺนาคเตน ภิกฺขุนา อุปสมฺ\-ปาเท\-ตพฺพํ, นิสฺสโย ทาตพฺโพ, สามเณโร อุปฏฺฐาเปตพฺโพ. อาปตฺติํ \csromanfolio{98}ชานาติ, อนาปตฺติํ ชานาติ, ลหุกํ อาปตฺติํ ชานาติ, ครุกํ อาปตฺติํ ชานาติ, อุภยานิ โข ปนสฺส ปาติโมกฺขานิ วิตฺถาเรน สฺวาคตานิ โหนฺติ สุวิภตฺตานิ สุปฺปวตฺตีนิ สุวินิจฺฉิตานิ สุตฺตโส อนุพฺยญฺชนโส, ทสวสฺโส วา โหติ อติเรก\-ทสวสฺโส วา. อิเมหิ โข ภิกฺขเว ฉหงฺเคหิ สมนฺนาคเตน ภิกฺขุนา อุปสมฺ\-ปาเท\-ตพฺพํ, นิสฺสโย ทาตพฺโพ, สามเณโร อุปฏฺฐาเป\-ตพฺโพติ. (14)}{อุปสมฺปาเทตพฺพ\-ฉกฺก\-จุทฺทส\-วาโร นิฏฺฐิโต.}
\csromanmatikaanchor{mtk.26}
\csromantocmarkref{section}{25. อญฺญติตฺถิยปุพฺพกถา}{mtk.26}
\bookmark[dest={mtk.26},level=1]{25. อญฺญติตฺถิยปุพฺพกถา}
\csromanheader{1}{25. \textbf{อญฺญติตฺถิยปุพฺพ}\-\textbf{กถา}}
\proseitem{86}{เตน โข ปน สมเยน โย โส อญฺญ\-ติตฺถิย\-ปุพฺโพ\csromansharedfootnote{40c9bac9d325c54f}{โย โส ปสุรปริพฺพาชโก อญฺญติตฺถิยปุพฺโพ (ก)} อุปชฺฌาเยน สหธมฺมิกํ วุจฺจมาโน อุปชฺฌายสฺส วาทํ อาโรเปตฺวา ตํเยว ติตฺถายตนํ สงฺกมิ, โส ปุน ปจฺจาคนฺตฺวา ภิกฺขู อุปสมฺปทํ ยาจิ. ภิกฺขู ภควโต เอตมตฺถํ อาโรเจสุํ. โย โส ภิกฺขเว อญฺญ\-ติตฺถิย\-ปุพฺโพ อุปชฺฌาเยน สหธมฺมิกํ วุจฺจมาโน อุปชฺฌายสฺส วาทํ อาโรเปตฺวา ตํเยว ติตฺถายตนํ สงฺกนฺโต, โส อาคโต น อุปสมฺ\-ปาเท\-ตพฺโพ. โย โส ภิกฺขเว อญฺโญปิ อญฺญ\-ติตฺถิย\-ปุพฺโพ อิมสฺมิํ ธมฺมวินเย อากงฺขติ ปพฺพชฺชํ, อากงฺขติ อุปสมฺปทํ, ตสฺส จตฺตาโร มาเส ปริวาโส ทาตพฺโพ. เอวญฺจ ปน ภิกฺขเว ทาตพฺโพ\csromandash{}ปฐมํ เกสมสฺสุํ โอหาราเปตฺวา กาสายานิ วตฺถานิ อจฺฉาทาเปตฺวา เอกํสํ อุตฺตราสงฺคํ การาเปตฺวา ภิกฺขูนํ ปาเท วนฺทาเปตฺวา อุกฺกุฏิกํ นิสีทาเปตฺวา อญฺชลิํ ปคฺคณฺหาเปตฺวา “เอวํ วเทหี”ติ วตฺตพฺโพ “พุทฺธํ สรณํ คจฺฉามิ, ธมฺมํ สรณํ คจฺฉามิ, สํฆํ สรณํ คจฺฉามิ. ทุติยมฺปิ พุทฺธํ สรณํ คจฺฉามิ, ทุติยมฺปิ ธมฺมํ สรณํ คจฺฉามิ, ทุติยมฺปิ สํฆํ สรณํ คจฺฉามิ. ตติยมฺปิ พุทฺธํ สรณํ คจฺฉามิ, ตติยมฺปิ ธมฺมํ สรณํ คจฺฉามิ, ตติยมฺปิ สํฆํ สรณํ คจฺฉามี”ติ.}

\prose{เตน ภิกฺขเว อญฺญ\-ติตฺถิย\-ปุพฺเพน สํฆํ อุปสงฺกมิตฺวา เอกํสํ อุตฺตราสงฺคํ กริตฺวา ภิกฺขูนํ ปาเท วนฺทิตฺวา อุกฺกุฏิกํ นิสีทิตฺวา อญฺชลิํ ปคฺคเหตฺวา เอวมสฺส วจนีโย “อหํ ภนฺเต อญฺญ\-ติตฺถิย\-ปุพฺโพ \csromanfolio{99}อิมสฺมิํ ธมฺมวินเย อากงฺขามิ อุปสมฺปทํ, โสหํ ภนฺเต สํฆํ จตฺตาโร มาเส ปริวาสํ ยาจามี”ติ. ทุติยมฺปิ ยาจิตพฺโพ. ตติยมฺปิ ยาจิตพฺโพ. พฺยตฺเตน ภิกฺขุนา ปฏิพเลน สํโฆ ญาเปตพฺโพ\csromandash{}}

\prose{“สุณาตุ เม ภนฺเต สํโฆ, อยํ อิตฺถนฺนาโม อญฺญ\-ติตฺถิย\-ปุพฺโพ อิมสฺมิํ ธมฺมวินเย อากงฺขติ อุปสมฺปทํ, โส สํฆํ จตฺตาโร มาเส ปริวาสํ ยาจติ, ยทิ สํฆสฺส ปตฺตกลฺลํ, สํโฆ อิตฺถนฺนามสฺส อญฺญ\-ติตฺถิย\-ปุพฺพสฺส จตฺตาโร มาเส ปริวาสํ ทเทยฺย, เอสา ญตฺติ.}

\prose{สุณาตุ เม ภนฺเต สํโฆ, อยํ อิตฺถนฺนาโม อญฺญ\-ติตฺถิย\-ปุพฺโพ อิมสฺมิํ ธมฺมวินเย อากงฺขติ อุปสมฺปทํ, โส สํฆํ จตฺตาโร มาเส ปริวาสํ ยาจติ, สํโฆ อิตฺถนฺนามสฺส อญฺญ\-ติตฺถิย\-ปุพฺพสฺส จตฺตาโร มาเส ปริวาสํ เทติ, ยสฺสายสฺมโต ขมติ อิตฺถนฺนามสฺส อญฺญ\-ติตฺถิย\-ปุพฺพสฺส จตฺตาโร มาเส ปริวาสสฺส ทานํ, โส ตุณฺหสฺส, ยสฺส นกฺขมติ, โส ภาเสยฺย.}

\prose{ทินฺโน สํเฆน อิตฺถนฺนามสฺส อญฺญ\-ติตฺถิย\-ปุพฺพสฺส จตฺตาโร มาเส ปริวาโส, ขมติ สํฆสฺส, ตสฺมา ตุณฺหี, เอวเมตํ ธารยามี”ติ.}

\setlength{\csromangathapagewidth}{0pt}
\setlength{\csromangathawakwidth}{0pt}
\csromangathameasuregroup{}{เอวํ โข ภิกฺขเว อญฺญติตฺถิยปุพฺโพ อาราธโก โหติ, \\ เอวํ อนาราธโก.}
\setlength{\csromangathaleftcol}{0pt}
\csromangathagroup{\csromangathastanza{\csromangathaitemline{87}{เอวํ โข ภิกฺขเว อญฺญ\-ติตฺถิย\-ปุพฺโพ อาราธโก โหติ,} \\ \csromangathacontline{เอวํ อนาราธโก.}}}

\prose{กถญฺจ ภิกฺขเว อญฺญ\-ติตฺถิย\-ปุพฺโพ อนาราธโก โหติ. อิธ ภิกฺขเว อญฺญ\-ติตฺถิย\-ปุพฺโพ อติกาเลน คามํ ปวิสติ, อติทิวา ปฏิกฺกมติ. เอวมฺปิ ภิกฺขเว อญฺญ\-ติตฺถิย\-ปุพฺโพ อนาราธโก โหติ.}

\prose{ปุน จปรํ ภิกฺขเว อญฺญ\-ติตฺถิย\-ปุพฺโพ เวสิยาโคจโร วา โหติ, วิธวาโคจโร วา โหติ, ถุลฺลกุมาริกา\-โคจโร วา โหติ, ปณฺฑกโคจโร วา โหติ, ภิกฺขุนิ\-โคจโร วา โหติ.}

\setlength{\csromangathapagewidth}{0pt}
\setlength{\csromangathawakwidth}{0pt}
\csromangathameasuregroup{}{เอวมฺปิ ภิกฺขเว อญฺญติตฺถิยปุพฺโพ อนาราธโก โหติ.}
\csromangathameasurewak{เอวมฺปิ ภิกฺขเว อญฺญติตฺถิยปุพฺโพ อนาราธโก โหติ.}
\setlength{\csromangathaleftcol}{0pt}
\csromangathagroup{\csromangathastanza{\csromangathawak{เอวมฺปิ ภิกฺขเว อญฺญ\-ติตฺถิย\-ปุพฺโพ อนาราธโก โหติ.}}}

\prose{ปุน จปรํ ภิกฺขเว อญฺญ\-ติตฺถิย\-ปุพฺโพ ยานิ ตานิ สพฺรหฺมจารีนํ อุจฺจาวจานิ กรณียานิ, ตตฺถ น ทกฺโข โหติ น อนลโส น ตตฺรุปายาย วีมํสาย สมนฺนาคโต น อลํ กาตุํ น อลํ สํวิธาตุํ.}

\setlength{\csromangathapagewidth}{0pt}
\setlength{\csromangathawakwidth}{0pt}
\csromangathameasuregroup{}{เอวมฺปิ ภิกฺขเว อญฺญติตฺถิยปุพฺโพ อนาราธโก โหติ.}
\csromangathameasurewak{เอวมฺปิ ภิกฺขเว อญฺญติตฺถิยปุพฺโพ อนาราธโก โหติ.}
\setlength{\csromangathaleftcol}{0pt}
\csromangathagroup{\csromangathastanza{\csromangathawak{เอวมฺปิ ภิกฺขเว อญฺญ\-ติตฺถิย\-ปุพฺโพ อนาราธโก โหติ.}}}

\proseitem{87}{\csromanfolio{100}ปุน จปรํ ภิกฺขเว อญฺญ\-ติตฺถิย\-ปุพฺโพ น ติพฺพจฺฉนฺโท โหติ อุทฺเทเส ปริปุจฺฉาย อธิสีเล อธิจิตฺเต อธิปญฺญาย. เอวมฺปิ ภิกฺขเว อญฺญ\-ติตฺถิย\-ปุพฺโพ อนาราธโก โหติ.}

\prose{ปุน จปรํ ภิกฺขเว อญฺญ\-ติตฺถิย\-ปุพฺโพ ยสฺส ติตฺถายตนา สงฺกนฺโต โหติ, ตสฺส สตฺถุโน ตสฺส ทิฏฺฐิยา ตสฺส ขนฺติยา ตสฺส รุจิยา ตสฺส อาทายสฺส อวณฺเณ ภญฺญมาเน กุปิโต โหติ อนตฺตมโน อนภิรทฺโธ, พุทฺธสฺส วา ธมฺมสฺส วา สํฆสฺส วา อวณฺเณ ภญฺญมาเน อตฺตมโน โหติ อุทคฺโค อภิรทฺโธ.}

\prose{ยสฺส วา ปน ติตฺถายตนา สงฺกนฺโต โหติ, ตสฺส สตฺถุโน ตสฺส ทิฏฺฐิยา ตสฺส ขนฺติยา ตสฺส รุจิยา ตสฺส อาทายสฺส วณฺเณ ภญฺญมาเน อตฺตมโน โหติ อุทคฺโค อภิรทฺโธ, พุทฺธสฺส วา ธมฺมสฺส วา สํฆสฺส วา วณฺเณ ภญฺญมาเน กุปิโต โหติ อนตฺตมโน อนภิรทฺโธ.}

\prose{อิทํ ภิกฺขเว สงฺฆาตนิกํ อญฺญ\-ติตฺถิย\-ปุพฺพสฺส อนาราธนียสฺมิํ. เอวมฺปิ โข ภิกฺขเว อญฺญ\-ติตฺถิย\-ปุพฺโพ อนาราธโก โหติ. เอวํ อนาราธโก โข ภิกฺขเว อญฺญ\-ติตฺถิย\-ปุพฺโพ อาคโต น อุปสมฺ\-ปาเท\-ตพฺโพ.}

\prose{กถญฺจ ภิกฺขเว อญฺญ\-ติตฺถิย\-ปุพฺโพ อาราธโก โหติ. อิธ ภิกฺขเว อญฺญ\-ติตฺถิย\-ปุพฺโพ นาติกาเลน คามํ ปวิสติ, นาติทิวา ปฏิกฺกมติ. เอวมฺปิ ภิกฺขเว อญฺญ\-ติตฺถิย\-ปุพฺโพ อาราธโก โหติ.}

\prose{ปุน จปรํ ภิกฺขเว อญฺญ\-ติตฺถิย\-ปุพฺโพ น เวสิยาโคจโร โหติ, น วิธวาโคจโร โหติ, น ถุลฺลกุมาริกา\-โคจโร โหติ, น ปณฺฑกโคจโร โหติ, น ภิกฺขุนิ\-โคจโร โหติ.}

\setlength{\csromangathapagewidth}{0pt}
\setlength{\csromangathawakwidth}{0pt}
\csromangathameasuregroup{}{เอวมฺปิ ภิกฺขเว อญฺญติตฺถิยปุพฺโพ อาราธโก โหติ.}
\csromangathameasurewak{เอวมฺปิ ภิกฺขเว อญฺญติตฺถิยปุพฺโพ อาราธโก โหติ.}
\setlength{\csromangathaleftcol}{0pt}
\csromangathagroup{\csromangathastanza{\csromangathawak{เอวมฺปิ ภิกฺขเว อญฺญ\-ติตฺถิย\-ปุพฺโพ อาราธโก โหติ.}}}

\prose{ปุน จปรํ ภิกฺขเว อญฺญ\-ติตฺถิย\-ปุพฺโพ ยานิ ตานิ สพฺรหฺมจารีนํ อุจฺจาวจานิ กรณียานิ, ตตฺถ ทกฺโข โหติ อนลโส ตตฺรุปายาย วีมํสาย สมนฺนาคโต อลํ กาตุํ อลํ สํวิธาตุํ.}

\setlength{\csromangathapagewidth}{0pt}
\setlength{\csromangathawakwidth}{0pt}
\csromangathameasuregroup{}{เอวมฺปิ ภิกฺขเว อญฺญติตฺถิยปุพฺโพ อาราธโก โหติ.}
\csromangathameasurewak{เอวมฺปิ ภิกฺขเว อญฺญติตฺถิยปุพฺโพ อาราธโก โหติ.}
\setlength{\csromangathaleftcol}{0pt}
\csromangathagroup{\csromangathastanza{\csromangathawak{เอวมฺปิ ภิกฺขเว อญฺญ\-ติตฺถิย\-ปุพฺโพ อาราธโก โหติ.}}}

\proseitem{87}{ปุน จปรํ ภิกฺขเว อญฺญ\-ติตฺถิย\-ปุพฺโพ ติพฺพจฺฉนฺโท โหติ อุทฺเทเส ปริปุจฺฉาย อธิสีเล อธิจิตฺเต อธิปญฺญาย. เอวมฺปิ ภิกฺขเว อญฺญ\-ติตฺถิย\-ปุพฺโพ อาราธโก โหติ.}

\prose{\csromanfolio{101}ปุน จปรํ ภิกฺขเว อญฺญ\-ติตฺถิย\-ปุพฺโพ ยสฺส ติตฺถายตนา สงฺกนฺโต โหติ, ตสฺส สตฺถุโน ตสฺส ทิฏฺฐิยา ตสฺส ขนฺติยา ตสฺส รุจิยา ตสฺส อาทายสฺส อวณฺเณ ภญฺญมาเน อตฺตมโน โหติ อุทคฺโค อภิรทฺโธ, พุทฺธสฺส วา ธมฺมสฺส วา สํฆสฺส วา อวณฺเณ ภญฺญมาเน กุปิโต โหติ อนตฺตมโน อนภิรทฺโธ.}

\prose{ยสฺส วา ปน ติตฺถายตนา สงฺกนฺโต โหติ, ตสฺส สตฺถุโน ตสฺส ทิฏฺฐิยา ตสฺส ขนฺติยา ตสฺส รุจิยา ตสฺส อาทายสฺส วณฺเณ ภญฺญมาเน กุปิโต โหติ อนตฺตมโน อนภิรทฺโธ, พุทฺธสฺส วา ธมฺมสฺส วา สํฆสฺส วา วณฺเณ ภญฺญมาเน อตฺตมาโน โหติ อุทคฺโค อภิรทฺโธ.}

\prose{อิทํ ภิกฺขเว สงฺฆาตนิกํ อญฺญ\-ติตฺถิย\-ปุพฺพสฺส อาราธนียสฺมิํ. เอวมฺปิ โข ภิกฺขเว อญฺญ\-ติตฺถิย\-ปุพฺโพ อาราธโก โหติ. เอวํ อาราธโก โข ภิกฺขเว อญฺญ\-ติตฺถิย\-ปุพฺโพ อาคโต อุปสมฺ\-ปาเท\-ตพฺโพ.}

\prose{สเจ ภิกฺขเว อญฺญ\-ติตฺถิย\-ปุพฺโพ นคฺโค อาคจฺฉติ, อุปชฺฌายมูลกํ จีวรํ ปริเยสิตพฺพํ.}

\setlength{\csromangathapagewidth}{0pt}
\setlength{\csromangathawakwidth}{0pt}
\csromangathameasuregroup{สเจ อจฺฉินฺนเกโส อาคจฺฉติ,}{\csromangathabat{สเจ อจฺฉินฺนเกโส อาคจฺฉติ,}{สํโฆ อปโลเกตพฺโพ ภณฺฑุกมฺมาย.}}
\csromangathasetleft{สเจ อจฺฉินฺนเกโส อาคจฺฉติ,}
\csromangathagroup{\csromangathastanza{\csromangathabat{สเจ อจฺฉินฺนเกโส อาคจฺฉติ,}{สํโฆ อปโลเกตพฺโพ ภณฺฑุกมฺมาย.}}}

\prose{เย เต ภิกฺขเว อคฺคิกา ชฏิลกา, เต อาคตา อุปสมฺ\-ปาเท\-ตพฺพา, น เตสํ ปริวาโส ทาตพฺโพ.}

\setlength{\csromangathapagewidth}{0pt}
\setlength{\csromangathawakwidth}{0pt}
\csromangathameasuregroup{ตํ กิสฺสเหตุ,}{\csromangathabat{ตํ กิสฺสเหตุ,}{กมฺมวาทิโน เอเต ภิกฺขเว กิริยวาทิโน.}}
\csromangathasetleft{ตํ กิสฺสเหตุ,}
\csromangathagroup{\csromangathastanza{\csromangathabat{ตํ กิสฺสเหตุ,}{กมฺมวาทิโน เอเต ภิกฺขเว กิริยวาทิโน.}}}

\prosewithcloser{สเจ ภิกฺขเว ชาติยา สากิโย อญฺญ\-ติตฺถิย\-ปุพฺโพ อาคจฺฉติ, โส อาคโต อุปสมฺ\-ปาเท\-ตพฺโพ, น ตสฺส ปริวาโส ทาตพฺโพ, อิมาหํ ภิกฺขเว ญาตีนํ อาเวณิกํ ปริหารํ ทมฺมีติ.}{อญฺญติตฺถิยปุพฺพ\-กถา นิฏฺฐิตา.}
\nitthitamruled{สตฺตโม ภาณวาโร.}
\csromanmatikaanchor{mtk.27}
\csromantocmarkref{section}{26. ปญฺจาพาธวตฺถุ}{mtk.27}
\bookmark[dest={mtk.27},level=1]{26. ปญฺจาพาธวตฺถุ}
\csromanheader{1}{26. \textbf{ปญฺจา}\-\textbf{พาธ}\-\textbf{วตฺถุ}}
\proseitem{88}{เตน โข ปน สมเยน มคเธสุ ปญฺจ อาพาธา อุสฺสนฺนา โหนฺติ กุฏฺฐํ คณฺโฑ กิลาโส โสโส อปมาโร. มนุสฺสา ปญฺจหิ อาพาเธหิ ผุฏฺฐา ชีวกํ โกมารภจฺจํ อุปสงฺกมิตฺวา เอวํ วทนฺติ “สาธุ โน อาจริย ติกิจฺฉาหี”ติ. อหํ ขฺวยฺโย พหุกิจฺโจ พหุกรณีโย, ราชา จ เม มาคโธ เสนิโย พิมฺพิสาโร \csromanfolio{102}อุปฏฺฐาตพฺโพ อิตฺถาคารญฺจ พุทฺธปฺปมุโข จ ภิกฺขุสํโฆ, นาหํ สกฺโกมิ ติกิจฺฉิตุนฺติ. สพฺพํ สาปเตยฺยญฺจ เต อาจริย โหตุ, มยญฺจ เต ทาสา, สาธุ โน อาจริย ติกิจฺฉาหีติ. อหํ ขฺวยฺโย พหุกิจฺโจ พหุกรณีโย, ราชา จ เม มาคโธ เสนิโย พิมฺพิสาโร อุปฏฺฐาตพฺโพ อิตฺถาคารญฺจ พุทฺธปฺปมุโข จ ภิกฺขุสํโฆ, นาหํ สกฺโกมิ ติกิจฺฉิตุนฺติ. อถ โข เตสํ มนุสฺสานํ เอตทโหสิ “อิเม โข สมณา สกฺยปุตฺติยา สุขสีลา สุขสมาจารา สุโภชนานิ ภุญฺชิตฺวา นิวาเตสุ สยเนสุ สยนฺติ, ยนฺนูน มยํ สมเณสุ สกฺยปุตฺติเยสุ ปพฺพเชยฺยาม, ตตฺถ ภิกฺขู เจว อุปฏฺฐหิสฺสนฺติ, ชีวโก จ โกมารภจฺโจ ติกิจฺฉิสฺสตี”ติ. อถ โข เต มนุสฺสา ภิกฺขู อุปสงฺกมิตฺวา ปพฺพชฺชํ ยาจิํสุ. เต ภิกฺขู ปพฺพาเชสุํ อุปสมฺปาเทสุํ. เต ภิกฺขู เจว อุปฏฺฐหิํสุ, ชีวโก จ โกมารภจฺโจ ติกิจฺฉิ. เตน โข ปน สมเยน ภิกฺขู พหู คิลาเน ภิกฺขู อุปฏฺฐหนฺตา ยาจนพหุลา วิญฺญตฺติพหุลา วิหรนฺติ “คิลานภตฺตํ เทถ, คิลานุปฏฺฐาก\-ภตฺตํ เทถ, คิลาน\-เภสชฺชํ เทถา”ติ. ชีวโกปิ โกมารภจฺโจ พหู คิลาเน ภิกฺขู ติกิจฺฉนฺโต อญฺญตรํ ราชกิจฺจํ ปริหาเปสิ.}

\proseitem{89}{อญฺญตโรปิ ปุริโส ปญฺจหิ อาพาเธหิ ผุฏฺโฐ ชีวกํ โกมารภจฺจํ อุปสงฺกมิตฺวา เอตทโวจ “สาธุ มํ อาจริย ติกิจฺฉาหี”ติ. อหํ ขฺวยฺโย พหุกิจฺโจ พหุกรณีโย, ราชา จ เม มาคโธ เสนิโย พิมฺพิสาโร อุปฏฺฐาตพฺโพ อิตฺถาคารญฺจ พุทฺธปฺปมุโข จ ภิกฺขุสํโฆ, นาหํ สกฺโกมิ ติกิจฺฉิตุนฺติ. สพฺพํ สาปเตยฺยญฺจ เต อาจริย โหตุ, อหญฺจ เต ทาโส, สาธุ มํ อาจริย ติกิจฺฉาหีติ. อหํ ขฺวยฺโย พหุกิจฺโจ พหุกรณีโย, ราชา จ เม มาคโธ เสนิโย พิมฺพิสาโร อุปฏฺฐาตพฺโพ อิตฺถาคารญฺจ พุทฺธปฺปมุโข จ ภิกฺขุสํโฆ, นาหํ สกฺโกมิ ติกิจฺฉิตุนฺติ. อถ โข ตสฺส ปุริสสฺส เอตทโหสิ “อิเม โข สมณา สกฺยปุตฺติยา สุขสีลา สุขสมาจารา สุโภชนานิ ภุญฺชิตฺวา นิวาเตสุ สยเนสุ สยนฺติ, ยนฺนูนาหํ สมเณสุ สกฺยปุตฺติเยสุ ปพฺพเชยฺยํ, ตตฺถ ภิกฺขู เจว อุปฏฺฐหิสฺสนฺติ, ชีวโก จ โกมารภจฺโจ ติกิจฺฉิสฺสติ, โสมฺหิ\csromansharedfootnote{b4e1ff1ce6d0c3e2}{โสหํ (พหูสุ, วิมติวิโนทนีฏีกา โอโลเกตพฺพา)} อโรโค วิพฺภมิสฺสามี”ติ. อถ โข โส ปุริโส ภิกฺขู \csromanfolio{103}อุปสงฺกมิตฺวา ปพฺพชฺชํ ยาจิ, ตํ ภิกฺขู ปพฺพาเชสุํ, อุปสมฺปาเทสุํ, ตํ ภิกฺขู เจว อุปฏฺฐหิํสุ, ชีวโก จ โกมารภจฺโจ ติกิจฺฉิ, โส อโรโค วิพฺภมิ. อทฺทสา โข ชีวโก โกมารภจฺโจ ตํ ปุริสํ วิพฺภนฺตํ, ทิสฺวาน ตํ ปุริสํ เอตทโวจ “นนุ ตฺวํ อยฺโย ภิกฺขูสุ ปพฺพชิโต อโหสี”ติ. เอวํ อาจริยาติ. กิสฺส ปน ตฺวํ อยฺโย เอวรูปมกาสีติ. อถ โข โส ปุริโส ชีวกสฺส โกมาร\-ภจฺจสฺส เอตมตฺถํ อาโรเจสิ. ชีวโก โกมารภจฺโจ อุชฺฌายติ ขิยฺยติ วิปาเจติ “กถํ หิ นาม ภทนฺตา\csromansharedfootnote{ae3c6722d12f3ee3}{ภทฺทนฺตา (ก)} ปญฺจหิ อาพาเธหิ ผุฏฺฐํ ปพฺพาเชสฺสนฺตี”ติ. อถ โข ชีวโก โกมารภจฺโจ เยน ภควา เตนุปสงฺกมิ, อุปสงฺกมิตฺวา ภควนฺตํ อภิวาเทตฺวา เอกมนฺตํ นิสีทิ, เอกมนฺตํ นิสินฺโน โข ชีวโก โกมารภจฺโจ ภควนฺตํ เอตทโวจ “สาธุ ภนฺเต อยฺยา ปญฺจหิ อาพาเธหิ ผุฏฺฐํ น ปพฺพาเชยฺยุนฺ”ติ. อถ โข ภควา ชีวกํ โกมารภจฺจํ ธมฺมิยา กถาย สนฺทสฺเสสิ สมาทเปสิ สมุตฺเตเชสิ สมฺปหํเสสิ. อถ โข ชีวโก โกมารภจฺโจ ภควตา ธมฺมิยา กถาย สนฺทสฺสิโต สมาทปิโต สมุตฺเตชิโต สมฺปหํสิโต อุฏฺฐายาสนา ภควนฺตํ อภิวาเทตฺวา ปทกฺขิณํ กตฺวา ปกฺกามิ. อถ โข ภควา เอตสฺมิํ นิทาเน เอตสฺมิํ ปกรเณ ธมฺมิํ กถํ กตฺวา ภิกฺขู อามนฺเตสิ “น ภิกฺขเว ปญฺจหิ อาพาเธหิ ผุฏฺโฐ ปพฺพาเชตพฺโพ, โย ปพฺพาเชยฺย, อาปตฺติ ทุกฺกฏสฺสา”ติ.}

\csromanmatikaanchor{mtk.28}
\csromantocmarkref{section}{27. ราชภฏวตฺถุ}{mtk.28}
\bookmark[dest={mtk.28},level=1]{27. ราชภฏวตฺถุ}
\csromanheader{1}{27. \textbf{ราชภฏ}\-\textbf{วตฺถุ}}
\proseitem{90}{เตน โข ปน สมเยน รญฺโญ มาคธสฺส เสนิยสฺส พิมฺพิสารสฺส ปจฺจนฺโต กุปิโต โหติ. อถ โข ราชา มาคโธ เสนิโย พิมฺพิสาโร เสนานายเก มหามตฺเต อาณาเปสิ “คจฺฉถ ภเณ ปจฺจนฺตํ อุจฺจินถา”ติ. “เอวํ เทวา”ติ โข เสนานายกา มหามตฺตา รญฺโญ มาคธสฺส เสนิยสฺส พิมฺพิสารสฺส ปจฺจสฺโสสุํ. อถ โข อภิญฺญาตานํ อภิญฺญาตานํ โยธานํ เอตทโหสิ “มยํ โข ยุทฺธา\-ภินนฺทิโน คจฺฉนฺตา ปาปญฺจ กโรม, พหุญฺจ อปุญฺญํ ปสวาม, เกน นุ โข มยํ อุปาเยน ปาปา จ วิรเมยฺยาม, กลฺยาณญฺจ กเรยฺยามา”ติ. อถ โข เตสํ โยธานํ เอตทโหสิ “อิเม \csromanfolio{104}โข สมณา สกฺยปุตฺติยา ธมฺมจาริโน สมจาริโน พฺรหฺมจาริโน สจฺจวาทิโน สีลวนฺโต กลฺยาณธมฺมา, สเจ โข มยํ สมเณสุ สกฺยปุตฺติเยสุ ปพฺพเชยฺยาม, เอวํ มยํ ปาปา จ วิรเมยฺยาม, กลฺยาณญฺจ กเรยฺยามา”ติ. อถ โข เต โยธา ภิกฺขู อุปสงฺกมิตฺวา ปพฺพชฺชํ ยาจิํสุ, เต ภิกฺขู ปพฺพาเชสุํ อุปสมฺปาเทสุํ. เสนานายกา มหามตฺตา ราชภเฏ ปุจฺฉิํสุ “กิํ นุ โข ภเณ อิตฺถนฺนาโม จ อิตฺถนฺนาโม จ โยธา น ทิสฺสนฺตี”ติ. อิตฺถนฺนาโม จ อิตฺถนฺนาโม จ สามิ โยธา ภิกฺขูสุ ปพฺพชิตาติ. เสนานายกา มหามตฺตา อุชฺฌายนฺติ ขิยฺยนฺติ วิปาเจนฺติ “กถํ หิ นาม สมณา สกฺยปุตฺติยา ราชภฏํ ปพฺพาเชสฺสนฺตี”ติ. เสนานายกา มหามตฺตา รญฺโญ มาคธสฺส เสนิยสฺส พิมฺพิสารสฺส เอตมตฺถํ อาโรเจสุํ. อถ โข ราชา มาคโธ เสนิโย พิมฺพิสาโร โวหาริเก มหามตฺเต ปุจฺฉิ “โย ภเณ ราชภฏํ ปพฺพาเชติ, กิํ โส ปสวตี”ติ. อุปชฺฌายสฺส เทว สีสํ เฉตพฺพํ, อนุสฺสาวกสฺส\csromansharedfootnote{68594fe4c53a7b7f}{อนุสาวกสฺส (ก)} ชิวฺหา อุทฺธริตพฺพา, คณสฺส อุปฑฺฒผาสุกา ภญฺชิตพฺพาติ. อถ โข ราชา มาคโธ เสนิโย พิมฺพิสาโร เยน ภควา เตนุปสงฺกมิ, อุปสงฺกมิตฺวา ภควนฺตํ อภิวาเทตฺวา เอกมนฺตํ นิสีทิ, เอกมนฺตํ นิสินฺโน โข ราชา มาคโธ เสนิโย พิมฺพิสาโร ภควนฺตํ เอตทโวจ “สนฺติ ภนฺเต ราชาโน อสฺสทฺธา อปฺปสนฺนา, เต อปฺป\-มตฺต\-เกนปิ ภิกฺขู วิเหเฐยฺยุํ, สาธุ ภนฺเต อยฺยา ราชภฏํ น ปพฺพาเชยฺยุนฺ”ติ. อถ โข ภควา ราชานํ มาคธํ เสนิยํ พิมฺพิสารํ ธมฺมิยา กถาย สนฺทสฺเสสิ สมาทเปสิ สมุตฺเตเชสิ สมฺปหํเสสิ. อถ โข ราชา มาคโธ เสนิโย พิมฺพิสาโร ภควตา ธมฺมิยา กถาย สนฺทสฺสิโต สมาทปิโต สมุตฺเตชิโต สมฺปหํสิโต อุฏฺฐายาสนา ภควนฺตํ อภิวาเทตฺวา ปทกฺขิณํ กตฺวา ปกฺกามิ. อถ โข ภควา เอตสฺมิํ นิทาเน เอตสฺมิํ ปกรเณ ธมฺมิํ กถํ กตฺวา ภิกฺขู อามนฺเตสิ “น ภิกฺขเว ราชภโฏ ปพฺพาเชตพฺโพ, โย ปพฺพาเชยฺย, อาปตฺติ ทุกฺกฏสฺสา”ติ.}

\csromanmatikaanchor{mtk.29}
\csromantocmarkref{section}{28. องฺคุลิมาลโจรวตฺถุ}{mtk.29}
\bookmark[dest={mtk.29},level=1]{28. องฺคุลิมาลโจรวตฺถุ}
\csromanheader{1}{28. \textbf{องฺคุลิมาล}\-\textbf{โจร}\-\textbf{วตฺถุ}}
\proseitem{91}{เตน โข ปน สมเยน โจโร องฺคุลิมาโล ภิกฺขูสุ ปพฺพชิโต โหติ. มนุสฺสา ปสฺสิตฺวา อุพฺพิชฺชนฺติปิ อุตฺตสนฺติปิ ปลายนฺติปิ อญฺเญนปิ คจฺฉนฺติ อญฺเญนปิ มุขํ กโรนฺติ ทฺวารมฺปิ ถเกนฺติ. มนุสฺสา \csromanfolio{105}อุชฺฌายนฺติ ขิยฺยนฺติ วิปาเจนฺติ “กถํ หิ นาม สมณา สกฺยปุตฺติยา ธชพนฺธํ โจรํ ปพฺพาเชสฺสนฺตี”ติ. อสฺโสสุํ โข ภิกฺขู เตสํ มนุสฺสานํ อุชฺฌายนฺตานํ ขิยฺยนฺตานํ วิปาเจนฺตานํ. อถ โข เต ภิกฺขู ภควโต เอตมตฺถํ อาโรเจสุํ ฯเปฯ. น ภิกฺขเว ธชพนฺโธ โจโร ปพฺพาเชตพฺโพ, โย ปพฺพาเชยฺย, อาปตฺติ ทุกฺกฏสฺสาติ.}

\csromanmatikaanchor{mtk.30}
\csromantocmarkref{section}{29. การเภทกโจรวตฺถุ}{mtk.30}
\bookmark[dest={mtk.30},level=1]{29. การเภทกโจรวตฺถุ}
\csromanheader{1}{29. \textbf{การเภทก}\-\textbf{โจร}\-\textbf{วตฺถุ}}
\proseitem{92}{เตน โข ปน สมเยน รญฺญา มาคเธน เสนิเยน พิมฺพิสาเรน อนุญฺญาตํ โหติ “เย สมเณสุ สกฺยปุตฺติเยสุ ปพฺพชนฺติ, น เต ลพฺภา กิญฺจิ กาตุํ, สฺวากฺขาโต ธมฺโม, จรนฺตุ พฺรหฺมจริยํ สมฺมา ทุกฺขสฺส อนฺตกิริยายา”ติ. เตน โข ปน สมเยน อญฺญตโร ปุริโส โจริกํ กตฺวา การาย พทฺโธ โหติ, โส การํ ภินฺทิตฺวา ปลายิตฺวา ภิกฺขูสุ ปพฺพชิโต โหติ. มนุสฺสา ปสฺสิตฺวา เอวมาหํสุ “อยํ โส การเภทโก จาโร, หนฺท นํ เนมา”ติ. เอกจฺเจ เอวมาหํสุ มายฺโย เอวํ อวจุตฺถ, อนุญฺญาตํ รญฺญา มาคเธน เสนิเยน พิมฺพิสาเรน “เย สมเณสุ สกฺยปุตฺติเยสุ ปพฺพชนฺติ, น เต ลพฺภา กิญฺจิ กาตุํ, สฺวากฺขาโต ธมฺโม, จรนฺตุ พฺรหฺมจริยํ สมฺมา ทุกฺขสฺส อนฺตกิริยายา”ติ. มนุสฺสา อุชฺฌายนฺติ ขิยฺยนฺติ วิปาเจนฺติ “อภยูวรา อิเม สมณา สกฺยปุตฺติยา, นยิเม ลพฺภา กิญฺจิ กาตุํ, กถํ หิ นาม สมณา สกฺยปุตฺติยา การเภทกํ โจรํ ปพฺพาเชสฺสนฺตี”ติ. ภควโต เอตมตฺถํ อาโรเจสุํ. น ภิกฺขเว การเภทโก โจโร ปพฺพาเชตพฺโพ, โย ปพฺพาเชยฺย, อาปตฺติ ทุกฺกฏสฺสาติ.}

\csromanmatikaanchor{mtk.31}
\csromantocmarkref{section}{30. ลิขิตกโจรวตฺถุ}{mtk.31}
\bookmark[dest={mtk.31},level=1]{30. ลิขิตกโจรวตฺถุ}
\csromanheader{1}{30. \textbf{ลิขิตก}\-\textbf{โจร}\-\textbf{วตฺถุ}}
\proseitem{93}{เตน โข ปน สมเยน อญฺญตโร ปุริโส โจริกํ กตฺวา ปลายิตฺวา ภิกฺขูสุ ปพฺพชิโต โหติ, โส จ รญฺโญ อนฺเตปุเร ลิขิโต โหติ “ยตฺถ ปสฺสติ, ตตฺถ หนฺตพฺโพ”ติ. มนุสฺสา ปสฺสิตฺวา เอวมาหํสุ “อยํ โส ลิขิตโก โจโร, หนฺท นํ หนามา”ติ. เอกจฺเจ เอวมาหํสุ มายฺโย เอวํ อวจุตฺถ, อนุญฺญาตํ รญฺญา มาคเธน เสนิเยน พิมฺพิสาเรน “เย สมเณสุ สกฺยปุตฺติเยสุ ปพฺพชนฺติ, น เต ลพฺภา กิญฺจิ กาตุํ, สฺวากฺขาโต ธมฺโม, จรนฺตุ พฺรหฺมจริยํ สมฺมา ทุกฺขสฺส อนฺตกิริยายา”ติ. มนุสฺสา อุชฺฌายนฺติ ขิยฺยนฺติ วิปาเจนฺติ \csromanfolio{106}“อภยูวรา อิเม สมณา สกฺยปุตฺติยา, นยิเม ลพฺภา กิญฺจิ กาตุํ, กถํ หิ นาม สมณา สกฺยปุตฺติยา ลิขิตกํ โจรํ ปพฺพาเชสฺสนฺตี”ติ. ภควโต เอตมตฺถํ อาโรเจสุํ. น ภิกฺขเว ลิขิตโก โจโร ปพฺพาเชตพฺโพ, โย ปพฺพาเชยฺย, อาปตฺติ ทุกฺกฏสฺสาติ.}

\csromanmatikaanchor{mtk.32}
\csromantocmarkref{section}{31. กสาหตวตฺถุ}{mtk.32}
\bookmark[dest={mtk.32},level=1]{31. กสาหตวตฺถุ}
\csromanheader{1}{31. \textbf{กสาหตวตฺถุ}}
\proseitem{94}{เตน โข ปน สมเยน อญฺญตโร ปุริโส กสาหโต กตทณฺฑกมฺโม ภิกฺขูสุ ปพฺพชิโต โหติ. มนุสฺสา อุชฺฌายนฺติ ขิยฺยนฺติ วิปาเจนฺติ “กถํ หิ นาม สมณา สกฺยปุตฺติยา กสาหตํ กตทณฺฑกมฺมํ ปพฺพาเชสฺสนฺตี”ติ. ภควโต เอตมตฺถํ อาโรเจสุํ. น ภิกฺขเว กสาหโต กตทณฺฑกมฺโม ปพฺพาเชตพฺโพ, โย ปพฺพาเชยฺย, อาปตฺติ ทุกฺกฏสฺสาติ.}

\csromanmatikaanchor{mtk.33}
\csromantocmarkref{section}{32. ลกฺขณาหตวตฺถุ}{mtk.33}
\bookmark[dest={mtk.33},level=1]{32. ลกฺขณาหตวตฺถุ}
\csromanheader{1}{32. \textbf{ลกฺขณา}\-\textbf{หต}\-\textbf{วตฺถุ}}
\proseitem{95}{เตน โข ปน สมเยน อญฺญตโร ปุริโส ลกฺขณาหโต กตทณฺฑกมฺโม ภิกฺขูสุ ปพฺพชิโต โหติ. มนุสฺสา อุชฺฌายนฺติ ขิยฺยนฺติ วิปาเจนฺติ “กถํ หิ นาม สมณา สกฺยปุตฺติยา ลกฺขณาหตํ กตทณฺฑกมฺมํ ปพฺพาเชสฺสนฺตี”ติ. ภควโต เอตมตฺถํ อาโรเจสุํ. น ภิกฺขเว ลกฺขณาหโต กตทณฺฑกมฺโม ปพฺพาเชตพฺโพ, โย ปพฺพาเชยฺย, อาปตฺติ ทุกฺกฏสฺสาติ.}

\csromanmatikaanchor{mtk.34}
\csromantocmarkref{section}{33. อิณายิกวตฺถุ}{mtk.34}
\bookmark[dest={mtk.34},level=1]{33. อิณายิกวตฺถุ}
\csromanheader{1}{33. \textbf{อิณายิกวตฺถุ}}
\proseitem{96}{เตน โข ปน สมเยน อญฺญตโร ปุริโส อิณายิโก ปลายิตฺวา ภิกฺขูสุ ปพฺพชิโต โหติ. ธนิยา ปสฺสิตฺวา เอวมาหํสุ “อยํ โส อมฺหากํ อิณายิโก, หนฺท นํ เนมา”ติ. เอกจฺเจ เอวมาหํสุ มายฺโย เอวํ อวจุตฺถ, อนุญฺญาตํ รญฺญา มาคเธน เสนิเยน พิมฺพิสาเรน “เย สมเณสุ สกฺยปุตฺติเยสุ ปพฺพชนฺติ, น เต ลพฺภา กิญฺจิ กาตุํ, สฺวากฺขาโต ธมฺโม, จรนฺตุ พฺรหฺมจริยํ สมฺมา ทุกฺขสฺส อนฺตกิริยายา”ติ. มนุสฺสา อุชฺฌายนฺติ ขิยฺยนฺติ วิปาเจนฺติ “อภยูวรา อิเม สมณา สกฺยปุตฺติยา, นยิเม ลพฺภา กิญฺจิ กาตุํ, กถํ หิ นาม สมณา สกฺยปุตฺติยา อิณายิกํ ปพฺพาเชสฺสนฺตี”ติ. ภควโต เอตมตฺถํ \csromanfolio{107}อาโรเจสุํ. น ภิกฺขเว อิณายิโก ปพฺพเชตพฺโพ, โย ปพฺพาเชยฺย, อาปตฺติ ทุกฺกฏสฺสาติ.}

\csromanmatikaanchor{mtk.35}
\csromantocmarkref{section}{34. ทาสวตฺถุ}{mtk.35}
\bookmark[dest={mtk.35},level=1]{34. ทาสวตฺถุ}
\csromanheader{1}{34. \textbf{ทาสวตฺถุ}}
\proseitem{94}{เตน โข ปน สมเยน อญฺญตโร ทาโส ปลายิตฺวา ภิกฺขูสุ ปพฺพชิโต โหติ. อยฺยกา\csromansharedfootnote{2811d52d73fedc83}{อยฺยิกา (ก), อยิรกา (สี)} ปสฺสิตฺวา เอวมาหํสุ “อยํ โส อมฺหากํ ทาโส, หนฺท นํ เนมา”ติ. เอกจฺเจ เอวมาหํสุ มายฺโย เอวํ อวจุตฺถ, อนุญฺญาตํ รญฺญา มาคเธน เสนิเยน พิมฺพิสาเรน “เย สมเณสุ สกฺยปุตฺติเยสุ ปพฺพชนฺติ, น เต ลพฺภา กิญฺจิ กาตุํ, สฺวากฺขาโต ธมฺโม, จรนฺตุ พฺรหฺมจริยํ สมฺมา ทุกฺขสฺส อนฺตกิริยายา”ติ. มนุสฺสา อุชฺฌายนฺติ ขิยฺยนฺติ วิปาเจนฺติ “อภยูวรา อิเม สมณา สกฺยปุตฺติยา, นยิเม ลพฺภา กิญฺจิ กาตุํ, กถํ หิ นาม สมณา สกฺยปุตฺติยา ทาสํ ปพฺพาเชสฺสนฺตี”ติ. ภควโต เอตมตฺถํ อาโรเจสุํ. น ภิกฺขเว ทาโส ปพฺพาเชตพฺโพ, โย ปพฺพาเชยฺย, อาปตฺติ ทุกฺกฏสฺสาติ.}

\csromanmatikaanchor{mtk.36}
\csromantocmarkref{section}{35. กมฺมารภณฺฑุวตฺถุ}{mtk.36}
\bookmark[dest={mtk.36},level=1]{35. กมฺมารภณฺฑุวตฺถุ}
\csromanheader{1}{35. \textbf{กมฺมารภณฺฑุ}\-\textbf{วตฺถุ}}
\proseitem{98}{เตน โข ปน สมเยน อญฺญตโร กมฺมารภณฺฑุ มาตาปิตูหิ สทฺธิํ ภณฺฑิตฺวา อารามํ คนฺตฺวา ภิกฺขูสุ ปพฺพชิโต โหติ. อถ โข ตสฺส กมฺมาร\-ภณฺฑุสฺส มาตาปิตโร ตํ กมฺมารภณฺฑุํ วิจินนฺตา อารามํ คนฺตฺวา ภิกฺขู ปุจฺฉิํสุ “อปิ ภนฺเต เอวรูปํ ทารกํ ปสฺเสยฺยาถา”ติ. ภิกฺขู อชานํเยว อาหํสุ “น ชานามา”ติ, อปสฺสํเยว อาหํสุ “น ปสฺสามา”ติ. อถ โข ตสฺส กมฺมาร\-ภณฺฑุสฺส มาตาปิตโร ตํ กมฺมารภณฺฑุํ วิจินนฺตา ภิกฺขูสุ ปพฺพชิตํ ทิสฺวา อุชฺฌายนฺติ ขิยฺยนฺติ วิปาเจนฺติ “อลชฺชิโน อิเม สมณา สกฺยปุตฺติยา ทุสฺสีลา มุสาวาทิโน ชานํเยว อาหํสุ ‘น ชานามา’ติ, ปสฺสํเยว อาหํสุ ‘น ปสฺสามา’ติ, อยํ ทารโก ภิกฺขูสุ ปพฺพชิโต”ติ. อสฺโสสุํ โข ภิกฺขู ตสฺส กมฺมาร\-ภณฺฑุสฺส มาตาปิตูนํ อุชฺฌายนฺตานํ ขิยฺยนฺตานํ วิปาเจนฺตานํ. อถ โข เต ภิกฺขู ภควโต เอตมตฺถํ อาโรเจสุํ. อนุชานามิ ภิกฺขเว สํฆํ อปโลเกตุํ ภณฺฑุกมฺมายาติ.}

\csromanmatikaanchor{mtk.37}
\csromantocmarkref{section}{36. อุปาลิทารกวตฺถุ}{mtk.37}
\bookmark[dest={mtk.37},level=1]{36. อุปาลิทารกวตฺถุ}
\csromanheader{1}{36. \textbf{อุปาลิทารก}\-\textbf{วตฺถุ}}
\proseitem{99}{\csromanfolio{108}\csromansymbolfootnote{*}{อิทํ วตฺถุ วิ 2. 169 ปิฏฺฐาทีสุปิ อาคตํ.}เตน โข ปน สมเยน ราชคเห สตฺตรส\-วคฺคิยา ทารกา สหายกา โหนฺติ, อุปาลิทารโก เตสํ ปาโมกฺโข โหติ. อถ โข อุปาลิสฺส มาตาปิตูนํ เอตทโหสิ “เกน นุ โข อุปาเยน อุปาลิ อมฺหากํ อจฺจเยน สุขญฺจ ชีเวยฺย, น จ กิลเมยฺยา”ติ. อถ โข อุปาลิสฺส มาตาปิตูนํ เอตทโหสิ “สเจ โข อุปาลิ เลขํ สิกฺเขยฺย, เอวํ โข อุปาลิ อมฺหากํ อจฺจเยน สุขญฺจ ชีเวยฺย, น จ กิลเมยฺยา”ติ. อถ โข อุปาลิสฺส มาตาปิตูนํ เอตทโหสิ “สเจ โข อุปาลิ เลขํ สิกฺขิสฺสติ, องฺคุลิโย ทุกฺขา ภวิสฺสนฺติ, สเจ โข อุปาลิ คณนํ สิกฺเขยฺย, เอวํ โข อุปาลิ อมฺหากํ อจฺจเยน สุขญฺจ ชีเวยฺย, น จ กิลเมยฺยา”ติ. อถ โข อุปาลิสฺส มาตาปิตูนํ เอตทโหสิ “สเจ โข อุปาลิ คณนํ สิกฺขิสฺสติ, อุรสฺส ทุกฺโข ภวิสฺสติ, สเจ โข อุปาลิ รูปํ สิกฺเขยฺย, เอวํ โข อุปาลิ อมฺหากํ อจฺจเยน สุขญฺจ ชีเวยฺย, น จ กิลเมยฺยา”ติ. อถ โข อุปาลิสฺส มาตาปิตูนํ เอตทโหสิ “สเจ โข อุปาลิ รูปํ สิกฺขิสฺสติ, อกฺขีนิ ทุกฺขา ภวิสฺสนฺติ, อิเม โข สมณา สกฺยปุตฺติยา สุขสีลา สุขสมาจารา สุโภชนานิ ภุญฺชิตฺวา นิวาเตสุ สยเนสุ สยนฺติ. สเจ โข อุปาลิ สมเณสุ สกฺยปุตฺติเยสุ ปพฺพเชยฺย, เอวํ โข อุปาลิ อมฺหากํ อจฺจเยน สุขญฺจ ชีเวยฺย, น จ กิลเมยฺยา”ติ.}

\prose{อสฺโสสิ โข อุปาลิทารโก มาตาปิตูนํ อิมํ กถาสลฺลาปํ. อถ โข อุปาลิทารโก เยน เต ทารกา เตนุปสงฺกมิ, อุปสงฺกมิตฺวา เต ทารเก เอตทโวจ “เอถ มยํ อยฺยา สมเณสุ สกฺยปุตฺติเยสุ ปพฺพชิสฺสามา”ติ. สเจ โข ตฺวํ อยฺย ปพฺพชิสฺสสิ, เอวํ มยมฺปิ ปพฺพชิสฺสามาติ. อถ โข เต ทารกา เอกเมกสฺส มาตาปิตโร อุปสงฺกมิตฺวา เอตทโวจุํ “อนุชานาถ มํ อคารสฺมา อนาคาริยํ ปพฺพชฺชายา”ติ. อถ โข เตสํ ทารกานํ มาตาปิตโร “สพฺเพปิเม ทารกา สมานจฺฉนฺทา กลฺยาณา\-ธิปฺปายา”ติ อนุชานิํสุ. เต ภิกฺขู อุปสงฺกมิตฺวา ปพฺพชฺชํ ยาจิํสุ. เต ภิกฺขู ปพฺพาเชสุํ อุปสมฺปาเทสุํ. เต รตฺติยา ปจฺจูสสมยํ ปจฺจุฏฺฐาย โรทนฺติ “ยาคุํ เทถ ภตฺตํ เทถ ขาทนียํ เทถา”ติ. ภิกฺขู เอวมาหํสุ “อาคเมถ อาวุโส ยาว รตฺติ วิภายติ, สเจ ยาคุ \csromanfolio{109}ภวิสฺสติ ปิวิสฺสถ, สเจ ภตฺตํ ภวิสฺสติ ภุญฺชิสฺสถ, สเจ ขาทนียํ ภวิสฺสติ ขาทิสฺสถ, โน เจ ภวิสฺสติ ยาคุ วา ภตฺตํ วา ขาทนียํ วา ปิณฺฑาย จริตฺวา ภุญฺชิสฺสถา”ติ. เอวมฺปิ โข เต ภิกฺขู ภิกฺขูหิ วุจฺจมานา โรทนฺติเยว “ยาคุํ เทถ ภตฺตํ เทถ ขาทนียํ เทถา”ติ. เสนาสนํ อุหทนฺติปิ อุมฺมิหนฺติปิ.}

\prose{อสฺโสสิ โข ภควา รตฺติยา ปจฺจูสสมยํ ปจฺจุฏฺฐาย ทารกสทฺทํ, สุตฺวาน อายสฺมนฺตํ อานนฺทํ อามนฺเตสิ “กิํ นุ โข โส อานนฺท ทารกสทฺโท”ติ. อถ โข อายสฺมา อานนฺโท ภควโต เอตมตฺถํ อาโรเจสิ ฯเปฯ. สจฺจํ กิร ภิกฺขเว ภิกฺขู ชานํ อูนวีสติวสฺสํ ปุคฺคลํ อุปสมฺ\-ปาเท\-นฺตีติ. สจฺจํ ภควาติ. วิครหิ พุทฺโธ ภควา ฯเปฯ. กถํ หิ นาม เต ภิกฺขเว โมฆปุริสา ชานํ อูนวีสติวสฺสํ ปุคฺคลํ อุปสมฺปาเทสฺสนฺติ. อูนวีสติวสฺโส ภิกฺขเว ปุคฺคโล อกฺขโม โหติ สีตสฺส อุณฺหสฺส ชิฆจฺฉาย ปิปาสาย ฑํสมกส\-วาตา\-ตป\-สรีสป\-สมฺผสฺสานํ ทุรุตฺตานํ ทุราคตานํ วจนปถานํ, อุปฺปนฺนานํ สารีริกานํ เวทนานํ ทุกฺขานํ ติพฺพานํ ขรานํ กฏุกานํ อสาตานํ อมนาปานํ ปาณหรานํ อนธิวาสก\-ชาติโก โหติ. วีสติวสฺโสว โข ภิกฺขเว ปุคฺคโล ขโม โหติ สีตสฺส อุณฺหสฺส ชิฆจฺฉาย ปิปาสาย ฑํสมกส\-วาตา\-ตป\-สรีสป\-สมฺผสฺสานํ ทุรุตฺตานํ ทุราคตานํ วจนปถานํ, อุปฺปนฺนานํ สารีริกานํ เวทนานํ ทุกฺขานํ ติพฺพานํ ขรานํ กฏุกานํ อสาตานํ อมนาปานํ ปาณหรานํ อธิวาสก\-ชาติโก โหติ. เนตํ ภิกฺขเว อปฺปสนฺนานํ วา ปสาทาย ปสนฺนานํ วา ภิยฺโยภาวาย ฯเปฯ วิครหิตฺวา ฯเปฯ ธมฺมิํ กถํ กตฺวา ภิกฺขู อามนฺเตสิ “น ภิกฺขเว ชานํ อูนวีสติวสฺโส ปุคฺคโล อุปสมฺ\-ปาเท\-ตพฺโพ, โย อุปสมฺปาเทยฺย, ยถาธมฺโม กาเรตพฺโพ”ติ.}

\csromanmatikaanchor{mtk.38}
\csromantocmarkref{section}{37. อหิวาตกโรควตฺถุ}{mtk.38}
\bookmark[dest={mtk.38},level=1]{37. อหิวาตกโรควตฺถุ}
\csromanheader{1}{37. \textbf{อหิวาตกโรค}\-\textbf{วตฺถุ}}
\proseitem{100}{เตน โข ปน สมเยน อญฺญตรํ กุลํ อหิวาตก\-โรเคน กาลงฺกตํ โหติ, ตสฺส ปิตาปุตฺตกา เสสา โหนฺติ, เต ภิกฺขูสุ ปพฺพชิตฺวา เอกโตว ปิณฺฑาย จรนฺติ. อถ โข โส ทารโก ปิตุโน ภิกฺขาย ทินฺนาย อุปธาวิตฺวา เอตทโวจ “มยฺหมฺปิ ตาต เทหิ, มยฺหมฺปิ ตาต เทหี”ติ. มนุสฺสา อุชฺฌายนฺติ ขิยฺยนฺติ วิปาเจนฺติ “อพฺรหฺมจาริโน อิเม สมณา สกฺยปุตฺติยา, อยมฺปิ ทารโก ภิกฺขุนิยา ชาโต”ติ. อสฺโสสุํ โข ภิกฺขู เตสํ มนุสฺสานํ \csromanfolio{110}อุชฺฌายนฺตานํ ขิยฺยนฺตานํ วิปาเจนฺตานํ. อถ โข เต ภิกฺขู ภควโต เอตมตฺถํ อาโรเจสุํ. น ภิกฺขเว อูนปนฺนรส\-วสฺโส ทารโก ปพฺพาเชตพฺโพ, โย ปพฺพาเชยฺย, อาปตฺติ ทุกฺกฏสฺสาติ.}

\prose{เตน โข ปน สมเยน อายสฺมโต อานนฺทสฺส อุปฏฺฐากกุลํ สทฺธํ ปสนฺนํ อหิวาตก\-โรเคน กาลงฺกตํ โหติ, เทฺว จ ทารกา เสสา โหนฺติ, เต โปราณเกน อาจิณฺณกปฺเปน ภิกฺขู ปสฺสิตฺวา อุปธาวนฺติ, ภิกฺขู อปสาเทนฺติ, เต ภิกฺขูหิ อปสาทิยมานา โรทนฺติ. อถ โข อายสฺมโต อานนฺทสฺส เอตทโหสิ “ภควตา ปญฺญตฺตํ ‘น อูนปนฺนรส\-วสฺโส ทารโก ปพฺพาเชตพฺโพ’ติ, อิเม จ ทารกา อูนปนฺนรส\-วสฺสา, เกน นุ โข อุปาเยน อิเม ทารกา น วินสฺเสยฺยุนฺ”ติ. อถ โข อายสฺมา อานนฺโท ภควโต เอตมตฺถํ อาโรเจสิ. อุสฺสหนฺติ ปน เต อานนฺท ทารกา กาเก อุฑฺฑาเปตุนฺติ. อุสฺสหนฺติ ภควาติ. อถ โข ภควา เอตสฺมิํ นิทาเน เอตสฺมิํ ปกรเณ ธมฺมิํ กถํ กตฺวา ภิกฺขู อามนฺเตสิ “อนุชานามิ ภิกฺขเว อูนปนฺนรส\-วสฺสํ ทารกํ กากุฑฺเฑปกํ ปพฺพาเชตุนฺ”ติ.}

\csromanmatikaanchor{mtk.39}
\csromantocmarkref{section}{38. กณฺฏกวตฺถุ}{mtk.39}
\bookmark[dest={mtk.39},level=1]{38. กณฺฏกวตฺถุ}
\csromanheader{1}{38. \textbf{กณฺฏกวตฺถุ}}
\proseitem{101}{เตน โข ปน สมเยน อายสฺมโต อุปนนฺทสฺส สกฺยปุตฺตสฺส เทฺว สามเณรา โหนฺติ กณฺฏโก จ มหโก จ, เต อญฺญมญฺญํ ทูเสสุํ. ภิกฺขู อุชฺฌายนฺติ ขิยฺยนฺติ วิปาเจนฺติ “กถํ หิ นาม สามเณรา เอวรูปํ อนาจารํ อาจริสฺสนฺตี”ติ. ภควโต เอตมตฺถํ อาโรเจสุํ. น ภิกฺขเว เอเกน เทฺว สามเณรา อุปฏฺฐาเปตพฺพา, โย อุปฏฺฐาเปยฺย, อาปตฺติ ทุกฺกฏสฺสาติ.}

\csromanmatikaanchor{mtk.40}
\csromantocmarkref{section}{39. อาหุนฺทริกวตฺถุ}{mtk.40}
\bookmark[dest={mtk.40},level=1]{39. อาหุนฺทริกวตฺถุ}
\csromanheader{1}{39. \textbf{อาหุนฺทริก}\-\textbf{วตฺถุ}}
\proseitem{102}{เตน โข ปน สมเยน ภควา ตตฺเถว ราชคเห วสฺสํ วสิ ตตฺถ เหมนฺตํ ตตฺถ คิมฺหํ. มนุสฺสา อุชฺฌายนฺติ ขิยฺยนฺติ วิปาเจนฺติ “อาหุนฺทริกา สมณานํ สกฺย\-ปุตฺติ\-ยานํ ทิสา อนฺธการา, น อิเมสํ ทิสา ปกฺขายนฺตี”ติ. อสฺโสสุํ โข ภิกฺขู เตสํ มนุสฺสานํ อุชฺฌายนฺตานํ ขิยฺยนฺตานํ วิปาเจนฺตานํ. อถ โข เต ภิกฺขู ภควโต เอตมตฺถํ อาโรเจสุํ. อถ โข ภควา อายสฺมนฺตํ อานนฺทํ \csromanfolio{111}อามนฺเตสิ “คจฺฉานนฺท อวาปุรณํ\csromansharedfootnote{2fca3d04659ef030}{อปาปุรณํ (ก)} อาทาย อนุปริเวณิยํ ภิกฺขูนํ อาโรเจหิ อิจฺฉตาวุโส ภควา ทกฺขิณาคิริํ จาริกํ ปกฺกมิตุํ, ยสฺสายสฺมโต อตฺโถ, โส อาคจฺฉตู”ติ. “เอวํ ภนฺเต”ติ โข อายสฺมา อานนฺโท ภควโต ปฏิสฺสุณิตฺวา อวาปุรณํ อาทาย อนุปริเวณิยํ ภิกฺขูนํ อาโรเจสิ “อิจฺฉตาวุโส ภควา ทกฺขิณาคิริํ จาริกํ ปกฺกมิตุํ, ยสฺสายสฺมโต อตฺโถ, โส อาคจฺฉตู”ติ. ภิกฺขู เอวมาหํสุ “ภควตา อาวุโส อานนฺท ปญฺญตฺตํ ทสวสฺสานิ นิสฺสาย วตฺถุํ ทสวสฺเสน นิสฺสยํ ทาตุํ, ตตฺถ จ โน คนฺตพฺพํ ภวิสฺสติ, นิสฺสโย จ คเหตพฺโพ ภวิสฺสติ, อิตฺตโร จ วาโส ภวิสฺสติ, ปุน จ ปจฺจาคนฺตพฺพํ ภวิสฺสติ, ปุน จ นิสฺสโย คเหตพฺโพ ภวิสฺสติ. สเจ อมฺหากํ อาจริยุ\-ปชฺฌายา คมิสฺสนฺติ, มยมฺปิ คมิสฺสาม, โน เจ อมฺหากํ อาจริยุ\-ปชฺฌายา คมิสฺสนฺติ, มยมฺปิ น คมิสฺสาม, ลหุจิตฺตกตา โน อาวุโส อานนฺท ปญฺญายิสฺสตี”ติ. อถ โข ภควา โอคเณน ภิกฺขุ\-สํเฆน ทกฺขิณาคิริํ จาริกํ ปกฺกามิ.}

\csromanmatikaanchor{mtk.41}
\csromantocmarkref{section}{40. นิสฺสยมุจฺจนกกถา}{mtk.41}
\bookmark[dest={mtk.41},level=1]{40. นิสฺสยมุจฺจนกกถา}
\csromanheader{1}{40. \textbf{นิสฺสย}\-\textbf{มุจฺจนก}\-\textbf{กถา}}
\proseitem{103}{อถ โข ภควา ทกฺขิณา\-คิริสฺมิํ ยถาภิรนฺตํ วิหริตฺวา ปุนเทว ราชคหํ ปจฺจาคจฺฉิ. อถ โข ภควา อายสฺมนฺตํ อานนฺทํ อามนฺเตสิ “กิํ นุ โข อานนฺท ตถาคโต โอคเณน ภิกฺขุ\-สํเฆน ทกฺขิณาคิริํ จาริกํ ปกฺกนฺโต”ติ. อถ โข อายสฺมา อานนฺโท ภควโต เอตมตฺถํ อาโรเจสิ. อถ โข ภควา เอตสฺมิํ นิทาเน เอตสฺมิํ ปกรเณ ธมฺมิํ กถํ กตฺวา ภิกฺขู อามนฺเตสิ “อนุชานามิ ภิกฺขเว พฺยตฺเตน ภิกฺขุนา ปฏิพเลน ปญฺจวสฺสานิ นิสฺสาย วตฺถุํ, อพฺยตฺเตน ยาวชีวํ”.}

\prose{ปญฺจหิ ภิกฺขเว องฺเคหิ สมนฺนาคเตน ภิกฺขุนา น อนิสฺสิเตน วตฺถพฺพํ. น อเสกฺเขน สีลกฺขนฺเธน สมนฺนาคโต โหติ, น อเสกฺเขน สมาธิ\-กฺขนฺเธน. น อเสกฺเขน ปญฺญา\-กฺขนฺเธน. น อเสกฺเขน วิมุตฺติ\-กฺขนฺเธน. น อเสกฺเขน วิมุตฺติ\-ญาณ\-ทสฺสนกฺขนฺเธน สมนฺนาคโต โหติ. อิเมหิ โข ภิกฺขเว ปญฺจหงฺเคหิ สมนฺนาคเตน ภิกฺขุนา น อนิสฺสิเตน วตฺถพฺพํ. (1)}

\prose{\csromanfolio{112}ปญฺจหิ ภิกฺขเว องฺเคหิ สมนฺนาคเตน ภิกฺขุนา อนิสฺสิเตน วตฺถพฺพํ. อเสกฺเขน สีลกฺขนฺเธน สมนฺนาคโต โหติ. อเสกฺเขน สมาธิ\-กฺขนฺเธน. อเสกฺเขน ปญฺญา\-กฺขนฺเธน. อเสกฺเขน วิมุตฺติ\-กฺขนฺเธน. อเสกฺเขน วิมุตฺติ\-ญาณ\-ทสฺสนกฺขนฺเธน สมนฺนาคโต โหติ. อิเมหิ โข ภิกฺขเว ปญฺจหงฺเคหิ สมนฺนาคเตน ภิกฺขุนา อนิสฺสิเตน วตฺถพฺพํ. (2)}

\prose{อปเรหิปิ ภิกฺขเว ปญฺจงฺเคหิ สมนฺนาคเตน ภิกฺขุนา น อนิสฺสิเตน วตฺถพฺพํ. อสฺสทฺโธ โหติ, อหิริโก โหติ, อโนตฺตปฺปี โหติ, กุสีโต โหติ, มุฏฺฐสฺสติ โหติ. อิเมหิ โข ภิกฺขเว ปญฺจหงฺเคหิ สมนฺนาคเตน ภิกฺขุนา น อนิสฺสิเตน วตฺถพฺพํ. (3)}

\prose{ปญฺจหิ ภิกฺขเว องฺเคหิ สมนฺนาคเตน ภิกฺขุนา อนิสฺสิเตน วตฺถพฺพํ. สทฺโธ โหติ, หิริมา โหติ, โอตฺตปฺปี โหติ, อารทฺธวีริโย โหติ, อุปฏฺฐิตสฺสติ โหติ. อิเมหิ โข ภิกฺขเว ปญฺจหงฺเคหิ สมนฺนาคเตน ภิกฺขุนา อนิสฺสิเตน วตฺถพฺพํ. (4)}

\prose{อปเรหิปิ ภิกฺขเว ปญฺจหงฺเคหิ สมนฺนาคเตน ภิกฺขุนา น อนิสฺสิเตน วตฺถพฺพํ. อธิสีเล สีลวิปนฺโน โหติ, อชฺฌาจาเร อาจารวิปนฺโน โหติ, อติทิฏฺฐิยา ทิฏฺฐิวิปนฺโน โหติ, อปฺปสฺสุโต โหติ, ทุปฺปญฺโญ โหติ. อิเมหิ โข ภิกฺขเว ปญฺจหงฺเคหิ สมนฺนาคเตน ภิกฺขุนา น อนิสฺสิเตน วตฺถพฺพํ. (5)}

\prose{ปญฺจหิ ภิกฺขเว องฺเคหิ สมนฺนาคเตน ภิกฺขุนา อนิสฺสิเตน วตฺถพฺพํ. น อธิสีเล สีลวิปนฺโน โหติ, น อชฺฌาจาเร อาจารวิปนฺโน โหติ, น อติทิฏฺฐิยา ทิฏฺฐิวิปนฺโน โหติ, พหุสฺสุโต โหติ, ปญฺญวา โหติ. อิเมหิ โข ภิกฺขเว ปญฺจหงฺเคหิ สมนฺนาคเตน ภิกฺขุนา อนิสฺสิเตน วตฺถพฺพํ. (6)}

\prose{อปเรหิปิ ภิกฺขเว ปญฺจหงฺเคหิ สมนฺนาคเตน ภิกฺขุนา น อนิสฺสิเตน วตฺถพฺพํ. อาปตฺติํ น ชานาติ, อนาปตฺติํ น ชานาติ, ลหุกํ อาปตฺติํ น ชานาติ, ครุกํ อาปตฺติํ น ชานาติ, อุภยานิ โข ปนสฺส ปาติโมกฺขานิ วิตฺถาเรน น สฺวาคตานิ โหนฺติ น สุวิภตฺตานิ น สุปฺปวตฺตีนิ น สุวินิจฺฉิตานิ สุตฺตโส อนุพฺยญฺชนโส. อิเมหิ โข ภิกฺขเว ปญฺจหงฺเคหิ สมนฺนาคเตน ภิกฺขุนา น อนิสฺสิเตน วตฺถพฺพํ. (7)}

\prose{\csromanfolio{113}ปญฺจหิ ภิกฺขเว องฺเคหิ สมนฺนาคเตน ภิกฺขุนา อนิสฺสิเตน วตฺถพฺพํ. อาปตฺติํ ชานาติ, อนาปตฺติํ ชานาติ, ลหุกํ อาปตฺติํ ชานาติ, ครุกํ อาปตฺติํ ชานาติ, อุภยานิ โข ปนสฺส ปาติโมกฺขานิ วิตฺถาเรน สฺวาคตานิ โหนฺติ สุวิภตฺตานิ สุปฺปวตฺตีนิ สุวินิจฺฉิตานิ สุตฺตโส อนุพฺยญฺชนโส. อิเมหิ โข ภิกฺขเว ปญฺจหงฺเคหิ สมนฺนาคเตน ภิกฺขุนา อนิสฺสิเตน วตฺถพฺพํ. (8)}

\prose{อปเรหิปิ ภิกฺขเว ปญฺจหงฺเคหิ สมนฺนาคเตน ภิกฺขุนา น อนิสฺสิเตน วตฺถพฺพํ. อาปตฺติํ น ชานาติ, อนาปตฺติํ น ชานาติ, ลหุกํ อาปตฺติํ น ชานาติ, ครุกํ อาปตฺติํ น ชานาติ, อูนปญฺจวสฺโส โหติ. อิเมหิ โข ภิกฺขเว ปญฺจหงฺเคหิ สมนฺนาคเตน ภิกฺขุนา น อนิสฺสิเตน วตฺถพฺพํ. (9)}

\prosewithcloserruled{ปญฺจหิ ภิกฺขเว องฺเคหิ สมนฺนาคเตน ภิกฺขุนา อนิสฺสิเตน วตฺถพฺพํ. อาปตฺติํ ชานาติ, อนาปตฺติํ ชานาติ, ลหุกํ อาปตฺติํ ชานาติ, ครุกํ อาปตฺติํ ชานาติ, ปญฺจวสฺโส วา โหติ อติเรก\-ปญฺจวสฺโส วา. อิเมหิ โข ภิกฺขเว ปญฺจหงฺเคหิ สมนฺนาคเตน ภิกฺขุนา อนิสฺสิเตน วตฺถพฺพํ. (10)}{ปญฺจก\-ทส\-วาโร นิฏฺฐิโต.}
\proseitem{104}{ฉหิ ภิกฺขเว องฺเคหิ สมนฺนาคเตน ภิกฺขุนา น อนิสฺสิเตน วตฺถพฺพํ. น อเสกฺเขน สีลกฺขนฺเธน สมนฺนาคโต โหติ, น อเสกฺเขน สมาธิ\-กฺขนฺเธน. น อเสกฺเขน ปญฺญา\-กฺขนฺเธน. น อเสกฺเขน วิมุตฺติ\-กฺขนฺเธน. น อเสกฺเขน วิมุตฺติ\-ญาณ\-ทสฺสนกฺขนฺเธน สมนฺนาคโต โหติ, อูนปญฺจวสฺโส โหติ. อิเมหิ โข ภิกฺขเว ฉหงฺเคหิ สมนฺนาคเตน ภิกฺขุนา น อนิสฺสิเตน วตฺถพฺพํ. (1)}

\prose{ฉหิ ภิกฺขเว องฺเคหิ สมนฺนาคเตน ภิกฺขุนา อนิสฺสิเตน วตฺถพฺพํ. อเสกฺเขน สีลกฺขนฺเธน สมนฺนาคโต โหติ, อเสกฺเขน สมาธิ\-กฺขนฺเธน. อเสกฺเขน ปญฺญา\-กฺขนฺเธน. อเสกฺเขน วิมุตฺติ\-กฺขนฺเธน. อเสกฺเขน วิมุตฺติ\-ญาณ\-ทสฺสนกฺขนฺเธน สมนฺนาคโต โหติ, ปญฺจวสฺโส วา โหติ อติเรก\-ปญฺจวสฺโส วา. อิเมหิ โข ภิกฺขเว ฉหงฺเคหิ สมนฺนาคเตน ภิกฺขุนา อนิสฺสิเตน วตฺถพฺพํ. (2)}

\prose{\csromanfolio{114}อปเรหิปิ ภิกฺขเว ฉหงฺเคหิ สมนฺนาคเตน ภิกฺขุนา น อนิสฺสิเตน วตฺถพฺพํ. อสฺสทฺโธ โหติ, อหิริโก โหติ, อโนตฺตปฺปี โหติ, กุสีโต โหติ, มุฏฺฐสฺสติ โหติ, อูนปญฺจวสฺโส โหติ. อิเมหิ โข ภิกฺขเว ฉหงฺเคหิ สมนฺนาคเตน ภิกฺขุนา น อนิสฺสิเตน วตฺถพฺพํ. (3)}

\prose{ฉหิ ภิกฺขเว องฺเคหิ สมนฺนาคเตน ภิกฺขุนา อนิสฺสิเตน วตฺถพฺพํ. สทฺโธ โหติ, หิริมา โหติ, โอตฺตปฺปี โหติ, อารทฺธวีริโย โหติ, อุปฏฺฐิตสฺสติ โหติ, ปญฺจวสฺโส วา โหติ อติเรก\-ปญฺจวสฺโส วา. อิเมหิ โข ภิกฺขเว ฉหงฺเคหิ สมนฺนาคเตน ภิกฺขุนา อนิสฺสิเตน วตฺถพฺพํ. (4)}

\prose{อปเรหิปิ ภิกฺขเว ฉหงฺเคหิ สมนฺนาคเตน ภิกฺขุนา น อนิสฺสิเตน วตฺถพฺพํ. อธิสีเล สีลวิปนฺโน โหติ, อชฺฌาจาเร อาจารวิปนฺโน โหติ, อติทิฏฺฐิยา ทิฏฺฐิวิปนฺโน, โหติ, อปฺปสฺสุโต โหติ, ทุปฺปญฺโญ โหติ, อูนปญฺจวสฺโส โหติ. อิเมหิ โข ภิกฺขเว ฉหงฺเคหิ สมนฺนาคเตน ภิกฺขุนา น อนิสฺสิเตน วตฺถพฺพํ. (5)}

\prose{ฉหิ ภิกฺขเว องฺเคหิ สมนฺนาคเตน ภิกฺขุนา อนิสฺสิเตน วตฺถพฺพํ. น อธิสีเล สีลวิปนฺโน โหติ, น อชฺฌาจาเร อาจารวิปนฺโน โหติ, น อติทิฏฺฐิยา ทิฏฺฐิวิปนฺโน โหติ, พหุสฺสุโต โหติ, ปญฺญวา โหติ, ปญฺจวสฺโส วา โหติ อติเรก\-ปญฺจวสฺโส วา. อิเมหิ โข ภิกฺขเว ฉหงฺเคหิ สมนฺนาคเตน ภิกฺขุนา อนิสฺสิเตน วตฺถพฺพํ. (6)}

\prose{อปเรหิปิ ภิกฺขเว ฉหงฺเคหิ สมนฺนาคเตน ภิกฺขุนา น อนิสฺสิเตน วตฺถพฺพํ. อาปตฺติํ น ชานาติ, อนาปตฺติํ น ชานาติ, ลหุกํ อาปตฺติํ น ชานาติ, ครุกํ อาปตฺติํ น ชานาติ, อุภยานิ โข ปนสฺส ปาติโมกฺขานิ วิตฺถาเรน น สฺวาคตานิ โหนฺติ น สุวิภตฺตานิ น สุปฺปวตฺตีนิ น สุวินิจฺฉิตานิ สุตฺตโส อนุพฺยญฺชนโส, อูนปญฺจวสฺโส โหติ. อิเมหิ โข ภิกฺขเว ฉหงฺเคหิ สมนฺนาคเตน ภิกฺขุนา น อนิสฺสิเตน วตฺถพฺพํ. (7)}

\prosewithcloserruled{ฉหิ ภิกฺขเว องฺเคหิ สมนฺนาคเตน ภิกฺขุนา อนิสฺสิเตน วตฺถพฺพํ. อาปตฺติํ ชานาติ, อนาปตฺติํ ชานาติ, ลหุกํ อาปตฺติํ ชานาติ, ครุกํ อาปตฺติํ ชานาติ, อุภยานิ โข ปนสฺส ปาติโมกฺขานิ วิตฺถาเรน สฺวาคตานิ โหนฺติ สุวิภตฺตานิ สุปฺปวตฺตีนิ สุวินิจฺฉิตานิ \csromanfolio{115}สุตฺตโส อนุพฺยญฺชนโส, ปญฺจวสฺโส วา โหติ อติเรก\-ปญฺจวสฺโส วา. อิเมหิ โข ภิกฺขเว ฉหงฺเคหิ สมนฺนาคเตน ภิกฺขุนา อนิสฺสิเตน วตฺถพฺพนฺติ. (8)}{อภยูวร\-ภาณวาโร นิฏฺฐิโต อฏฺฐโม.}
\csromanmatikaanchor{mtk.42}
\csromantocmarkref{section}{41. ราหุลวตฺถุ}{mtk.42}
\bookmark[dest={mtk.42},level=1]{41. ราหุลวตฺถุ}
\csromanheader{1}{41. \textbf{ราหุลวตฺถุ}}
\proseitem{105}{อถ โข ภควา ราชคเห ยถาภิรนฺตํ วิหริตฺวา เยน กปิลวตฺถุ เตน จาริกํ ปกฺกามิ, อนุปุพฺเพน จาริกํ จรมาโน เยน กปิลวตฺถุ ตทวสริ, ตตฺร สุทํ ภควา สกฺเกสุ วิหรติ กปิล\-วตฺถุสฺมิํ นิโคฺรธาราเม. อถ โข ภควา ปุพฺพณฺหสมยํ นิวาเสตฺวา ปตฺต\-จีวรมา\-ทาย เยน สุทฺโธทนสฺส สกฺกสฺส นิเวสนํ เตนุปสงฺกมิ, อุปสงฺกมิตฺวา ปญฺญตฺเต อาสเน นิสีทิ. อถ โข ราหุลมาตา เทวี ราหุลํ กุมารํ เอตทโวจ “เอโส เต ราหุล ปิตา, คจฺฉสฺสุ\csromansharedfootnote{b0f778a8cca31ec9}{คจฺฉสฺส (สฺยา)} ทายชฺชํ ยาจาหี”ติ. อถ โข ราหุโล กุมาโร เยน ภควา เตนุปสงฺกมิ, อุปสงฺกมิตฺวา ภควโต ปุรโต อฏฺฐาสิ “สุขา เต สมณ ฉายา”ติ. อถ โข ภควา อุฏฺฐายาสนา ปกฺกามิ. อถ โข ราหุโล กุมาโร ภควนฺตํ ปิฏฺฐิโต ปิฏฺฐิโต อนุพนฺธิ “ทายชฺชํ เม สมณ เทหิ, ทายชฺชํ เม สมณ เทหี”ติ. อถ โข ภควา อายสฺมนฺตํ สาริปุตฺตํ อามนฺเตสิ “เตน หิ ตฺวํ สาริปุตฺต ราหุลํ กุมารํ ปพฺพาเชหี”ติ. กถาหํ ภนฺเต ราหุลํ กุมารํ ปพฺพาเชมีติ. อถ โข ภควา เอตสฺมิํ นิทาเน เอตสฺมิํ ปกรเณ ธมฺมิํ กถํ กตฺวา ภิกฺขู อามนฺเตสิ “อนุชานามิ ภิกฺขเว ตีหิ สรณคมเนหิ สามเณร\-ปพฺพชฺชํ. เอวญฺจ ปน ภิกฺขเว ปพฺพาเชตพฺโพ, ปฐมํ เกสมสฺสุํ โอหาราเปตฺวา กาสายานิ วตฺถานิ อจฺฉาทาเปตฺวา เอกํสํ อุตฺตราสงฺคํ การาเปตฺวา ภิกฺขูนํ ปาเท วนฺทาเปตฺวา อุกฺกุฏิกํ นิสีทาเปตฺวา อญฺชลิํ ปคฺคณฺหาเปตฺวา ‘เอวํ วเทหี’ติ วตฺตพฺโพ ‘พุทฺธํ สรณํ คจฺฉามิ, ธมฺมํ สรณํ คจฺฉามิ, สํฆํ สรณํ คจฺฉามิ. ทุติยมฺปิ พุทฺธํ สรณํ คจฺฉามิ, ทุติยมฺปิ ธมฺมํ สรณํ คจฺฉามิ, ทุติยมฺปิ สํฆํ สรณํ คจฺฉามิ. ตติยมฺปิ พุทฺธํ สรณํ คจฺฉามิ, \csromanfolio{116}ตติยมฺปิ ธมฺมํ สรณํ คจฺฉามิ, ตติยมฺปิ สํฆํ สรณํ คจฺฉามี’ติ. อนุชานามิ ภิกฺขเว อิเมหิ ตีหิ สรณคมเนหิ สามเณรปพฺพชฺชนฺ”ติ. อถ โข อายสฺมา สาริปุตฺโต ราหุลํ กุมารํ ปพฺพาเชสิ.}

\prose{อถ โข สุทฺโธทโน สกฺโก เยน ภควา เตนุปสงฺกมิ, อุปสงฺกมิตฺวา ภควนฺตํ อภิวาเทตฺวา เอกมนฺตํ นิสีทิ, เอกมนฺตํ นิสินฺโน โข สุทฺโธทโน สกฺโก ภควนฺตํ เอตทโวจ “เอกาหํ ภนฺเต ภควนฺตํ วรํ ยาจามี”ติ. อติกฺกนฺตวรา โข โคตม ตถาคตาติ. ยญฺจ ภนฺเต กปฺปติ, ยญฺจ อนวชฺชนฺติ. วเทหิ โคตมาติ. ภควติ เม ภนฺเต ปพฺพชิเต อนปฺปกํ ทุกฺขํ อโหสิ, ตถา นนฺเท, อธิมตฺตํ ราหุเล, ปุตฺตเปมํ ภนฺเต ฉวิํ ฉินฺทติ, ฉวิํ เฉตฺวา จมฺมํ ฉินฺทติ, จมฺมํ เฉตฺวา มํสํ ฉินฺทติ, มํสํ เฉตฺวา นฺหารุํ ฉินฺทติ, นฺหารุํ เฉตฺวา อฏฺฐิํ ฉินฺทติ, อฏฺฐิํ เฉตฺวา อฏฺฐิมิญฺชํ อาหจฺจ ติฏฺฐติ, สาธุ ภนฺเต อยฺยา อนนุญฺญาตํ มาตาปิตูหิ ปุตฺตํ น ปพฺพาเชยฺยุนฺติ. อถ โข ภควา สุทฺโธทนํ สกฺกํ ธมฺมิยา กถาย สนฺทสฺเสสิ สมาทเปสิ สมุตฺเตเชสิ สมฺปหํเสสิ. อถ โข สุทฺโธทโน สกฺโก ภควตา ธมฺมิยา กถาย สนฺทสฺสิโต สมาทปิโต สมุตฺเตชิโต สมฺปหํสิโต อุฏฺฐายาสนา ภควนฺตํ อภิวาเทตฺวา ปทกฺขิณํ กตฺวา ปกฺกามิ. อถ โข ภควา เอตสฺมิํ นิทาเน เอตสฺมิํ ปกรเณ ธมฺมิํ กถํ กตฺวา ภิกฺขู อามนฺเตสิ “น ภิกฺขเว อนนุญฺญาโต มาตาปิตูหิ ปุตฺโต ปพฺพาเชตพฺโพ, โย ปพฺพาเชยฺย, อาปตฺติ ทุกฺกฏสฺสา”ติ.}

\prose{อถ โข ภควา กปิล\-วตฺถุสฺมิํ ยถาภิรนฺตํ วิหริตฺวา เยน สาวตฺถิ เตน จาริกํ ปกฺกามิ, อนุปุพฺเพน จาริกํ จรมาโน เยน สาวตฺถิ ตทวสริ, ตตฺร สุทํ ภควา สาวตฺถิยํ วิหรติ เชตวเน อนาถ\-ปิณฺฑิกสฺส อาราเม. เตน โข ปน สมเยน อายสฺมโต สาริปุตฺตสฺส อุปฏฺฐากกุลํ อายสฺมโต สาริปุตฺตสฺส สนฺติเก ทารกํ ปาเหสิ “อิมํ ทารกํ เถโร ปพฺพาเชตู”ติ. อถ โข อายสฺมโต สาริปุตฺตสฺส เอตทโหสิ “ภควตา ปญฺญตฺตํ น เอเกน เทฺว สามเณรา อุปฏฺฐาเป\-ตพฺพาติ, อยญฺจ เม ราหุโล สามเณโร, กถํ นุ โข มยา ปฏิปชฺชิตพฺพนฺ”ติ. ภควโต เอตมตฺถํ อาโรเจสิ. อนุชานามิ ภิกฺขเว พฺยตฺเตน ภิกฺขุนา ปฏิพเลน เอเกน เทฺว สามเณเร อุปฏฺฐาเปตุํ, \csromanfolio{117}ยาวตเก วา ปน อุสฺสหติ โอวทิตุํ อนุสาสิตุํ, ตาวตเก อุปฏฺฐาเปตุนฺติ.}

\csromanmatikaanchor{mtk.43}
\csromantocmarkref{section}{42. สิกฺขาปทกถา}{mtk.43}
\bookmark[dest={mtk.43},level=1]{42. สิกฺขาปทกถา}
\csromanheader{1}{42. \textbf{สิกฺขาปท}\-\textbf{กถา}}
\proseitem{106}{อถ โข สามเณรานํ เอตทโหสิ “กติ นุ โข อมฺหากํ สิกฺขาปทานิ, กตฺถ จ อเมฺหหิ สิกฺขิตพฺพนฺ”ติ. ภควโต เอตมตฺถํ อาโรเจสุํ ฯเปฯ. อนุชานามิ ภิกฺขเว สามเณรานํ ทส สิกฺขาปทานิ, เตสุ จ สามเณเรหิ สิกฺขิตุํ. ปาณาติปาตา เวรมณี\csromansharedfootnote{610b1328d6460b0e}{เวรมณิ, เวรมณิํ (ก)}, อทินฺนาทานา เวรมณี, อพฺรหฺมจริยา เวรมณี, มุสาวาทา เวรมณี, สุรา\-เมรย\-มชฺช\-ปมาทฏฺฐานา เวรมณี, วิกาลโภชนา เวรมณี, นจฺจ\-คีต\-วาทิต\-วิสูก\-ทสฺสนา เวรมณี, มาลา\-คนฺธ\-วิเลปน\-ธารณ\-มณฺฑน\-วิภูสน\-ฏฺฐานา เวรมณี, อุจฺจาสยน\-มหาสยนา เวรมณี, ชาตรูป\-รชต\-ปฏิคฺคหณา เวรมณี. อนุชานามิ ภิกฺขเว สามเณรานํ อิมานิ ทส สิกฺขาปทานิ, อิเมสุ จ สามเณเรหิ สิกฺขิตุนฺติ.}

\csromanmatikaanchor{mtk.44}
\csromantocmarkref{section}{43. ทณฺฑกมฺมวตฺถุ}{mtk.44}
\bookmark[dest={mtk.44},level=1]{43. ทณฺฑกมฺมวตฺถุ}
\csromanheader{1}{43. \textbf{ทณฺฑกมฺม}\-\textbf{วตฺถุ}}
\proseitem{107}{เตน โข ปน สมเยน สามเณรา ภิกฺขูสุ อคารวา อปฺปติสฺสา อสภาค\-วุตฺติกา วิหรนฺติ. ภิกฺขู อุชฺฌายนฺติ ขิยฺยนฺติ วิปาเจนฺติ “กถํ หิ นาม สามเณรา ภิกฺขูสุ อคารวา อปฺปติสฺสา อสภาค\-วุตฺติกา วิหริสฺสนฺตี”ติ. ภควโต เอตมตฺถํ อาโรเจสุํ ฯเปฯ. อนุชานามิ ภิกฺขเว ปญฺจหงฺเคหิ สมนฺนาคตสฺส สามเณรสฺส ทณฺฑกมฺมํ กาตุํ. ภิกฺขูนํ อลาภาย ปริสกฺกติ, ภิกฺขูนํ อนตฺถาย ปริสกฺกติ, ภิกฺขูนํ อวาสาย ปริสกฺกติ, ภิกฺขู อกฺโกสติ ปริภาสติ, ภิกฺขู ภิกฺขูหิ เภเทติ. อนุชานามิ ภิกฺขเว อิเมหิ ปญฺจหงฺเคหิ สมนฺนาคตสฺส สามเณรสฺส ทณฺฑกมฺมํ กาตุนฺติ.}

\prose{อถ โข ภิกฺขูนํ เอตทโหสิ “กิํ นุ โข ทณฺฑกมฺมํ กาตพฺพนฺ”ติ. ภควโต เอตมตฺถํ อาโรเจสุํ. อนุชานามิ ภิกฺขเว อาวรณํ กาตุนฺติ.}

\prose{เตน โข ปน สมเยน ภิกฺขู สามเณรานํ สพฺพํ สํฆารามํ อาวรณํ กโรนฺติ, สามเณรา อารามํ ปวิสิตุํ อลภมานา \csromanfolio{118}ปกฺกมนฺติปิ วิพฺภมนฺติปิ ติตฺถิเยสุปิ สงฺกมนฺติ. ภควโต เอตมตฺถํ อาโรเจสุํ. น ภิกฺขเว สพฺโพ สํฆาราโม อาวรณํ กาตพฺโพ, โย กเรยฺย, อาปตฺติ ทุกฺกฏสฺส. อนุชานามิ ภิกฺขเว ยตฺถ วา วสติ, ยตฺถ วา ปฏิกฺกมติ, ตตฺถ อาวรณํ กาตุนฺติ.}

\prosewithcloserruled{เตน โข ปน สมเยน ภิกฺขู สามเณรานํ มุขทฺวาริกํ อาหารํ อาวรณํ กโรนฺติ. มนุสฺสา ยาคุปานมฺปิ สํฆภตฺตมฺปิ กโรนฺตา สามเณเร เอวํ วเทนฺติ “เอถ ภนฺเต ยาคุํ ปิวถ, เอถ ภนฺเต ภตฺตํ ภุญฺชถา”ติ. สามเณรา เอวํ วเทนฺติ “นาวุโส ลพฺภา ภิกฺขูหิ อาวรณํ กตนฺ”ติ. มนุสฺสา อุชฺฌายนฺติ ขิยฺยนฺติ วิปาเจนฺติ “กถํ หิ นาม ภทนฺตา สามเณรานํ มุขทฺวาริกํ อาหารํ อาวรณํ กริสฺสนฺตี”ติ. ภควโต เอตมตฺถํ อาโรเจสุํ. น ภิกฺขเว มุขทฺวาริโก อาหาโร อาวรณํ กาตพฺโพ, โย กเรยฺย, อาปตฺติ ทุกฺกฏสฺสาติ.}{ทณฺฑกมฺม\-วตฺถุ นิฏฺฐิตํ.}
\csromanmatikaanchor{mtk.45}
\csromantocmarkref{section}{44. อนาปุจฺฉาวรณวตฺถุ}{mtk.45}
\bookmark[dest={mtk.45},level=1]{44. อนาปุจฺฉาวรณวตฺถุ}
\csromanheader{1}{44. \textbf{อนาปุจฺฉา}\-\textbf{วรณ}\-\textbf{วตฺถุ}}
\proseitem{108}{เตน โข ปน สมเยน ฉพฺพคฺคิยา ภิกฺขู อุปชฺฌาเย อนาปุจฺฉา สามเณรานํ อาวรณํ กโรนฺติ. อุปชฺฌายา คเวสนฺติ กถํ\csromansharedfootnote{164deec8497f84f4}{กหํ (ก)} นุ โข อมฺหากํ สามเณรา น ทิสฺสนฺตีติ. ภิกฺขู เอวมาหํสุ “ฉพฺพคฺคิเยหิ อาวุโส ภิกฺขูหิ อาวรณํ กตนฺ”ติ. อุปชฺฌายา อุชฺฌายนฺติ ขิยฺยนฺติ วิปาเจนฺติ “กถํ หิ นาม ฉพฺพคฺคิยา ภิกฺขู อเมฺห อนาปุจฺฉา อมฺหากํ สามเณรานํ อาวรณํ กริสฺสนฺตี”ติ. ภควโต เอตมตฺถํ อาโรเจสุํ. น ภิกฺขเว อุปชฺฌาเย อนาปุจฺฉา อาวรณํ กาตพฺพํ, โย กเรยฺย, อาปตฺติ ทุกฺกฏสฺสาติ.}

\csromanmatikaanchor{mtk.46}
\csromantocmarkref{section}{45. อปลาฬนวตฺถุ}{mtk.46}
\bookmark[dest={mtk.46},level=1]{45. อปลาฬนวตฺถุ}
\csromanheader{1}{45. \textbf{อปลาฬน}\-\textbf{วตฺถุ}}
\prose{เตน โข ปน สมเยน ฉพฺพคฺคิยา ภิกฺขู เถรานํ ภิกฺขูนํ สามเณเร อปลาเฬนฺติ. เถรา สามํ ทนฺตกฏฺฐมฺปิ มุโขทกมฺปิ คณฺหนฺตา}

\prose{\csromanfolio{119}กิลมนฺติ. ภควโต เอตมตฺถํ อาโรเจสุํ. น ภิกฺขเว อญฺญสฺส ปริสา อปลาเฬตพฺพา, โย อปลาเฬยฺย, อาปตฺติ ทุกฺกฏสฺสาติ.}

\csromanmatikaanchor{mtk.47}
\csromantocmarkref{section}{46. กณฺฏกสามเณรวตฺถุ}{mtk.47}
\bookmark[dest={mtk.47},level=1]{46. กณฺฏกสามเณรวตฺถุ}
\csromanheader{1}{46. \textbf{กณฺฏก}\-\textbf{สามเณร}\-\textbf{วตฺถุ}}
\prose{เตน โข ปน สมเยน อายสฺมโต อุปนนฺทสฺส สกฺยปุตฺตสฺส กณฺฏโก นาม สามเณโร กณฺฏกิํ นาม ภิกฺขุนิํ ทูเสสิ. ภิกฺขู อุชฺฌายนฺติ ขิยฺยนฺติ วิปาเจนฺติ “กถํ หิ นาม สามเณโร เอวรูปํ อนาจารํ อาจริสฺสตี”ติ. ภควโต เอตมตฺถํ อาโรเจสุํ. อนุชานามิ ภิกฺขเว ทสหงฺเคหิ สมนฺนาคตํ สามเณรํ นาเสตุํ. ปาณาติปาตี โหติ, อทินฺนาทายี โหติ, อพฺรหฺมจารี โหติ, มุสาวาที โหติ, มชฺชปายี โหติ, พุทฺธสฺส อวณฺณํ ภาสติ, ธมฺมสฺส อวณฺณํ ภาสติ, สํฆสฺส อวณฺณํ ภาสติ, มิจฺฉาทิฏฺฐิโก โหติ, ภิกฺขุนิ\-ทูสโก โหติ. อนุชานามิ ภิกฺขเว อิเมหิ ทสหงฺเคหิ สมนฺนาคตํ สามเณรํ นาเสตุนฺติ.}

\csromanmatikaanchor{mtk.48}
\csromantocmarkref{section}{47. ปณฺฑกวตฺถุ}{mtk.48}
\bookmark[dest={mtk.48},level=1]{47. ปณฺฑกวตฺถุ}
\csromanheader{1}{47. \textbf{ปณฺฑกวตฺถุ}}
\proseitem{109}{เตน โข ปน สมเยน อญฺญตโร ปณฺฑโก ภิกฺขูสุ ปพฺพชิโต โหติ, โส ทหเร ทหเร ภิกฺขู อุปสงฺกมิตฺวา เอวํ วเทติ “เอถ มํ อายสฺมนฺโต ทูเสถา”ติ. ภิกฺขู อปสาเทนฺติ “นสฺส ปณฺฑก, วินสฺส ปณฺฑก, โก ตยา อตฺโถ”ติ. โส ภิกฺขูหิ อปสาทิโต มหนฺเต มหนฺเต โมฬิคลฺเล สามเณเร อุปสงฺกมิตฺวา เอวํ วเทติ “เอถ มํ อาวุโส ทูเสถา”ติ. สามเณรา อปสาเทนฺติ “นสฺส ปณฺฑก, วินสฺส ปณฺฑก, โก ตยา อตฺโถ”ติ. โส สามเณเรหิ อปสาทิโต หตฺถิภณฺเฑ อสฺสภณฺเฑ อุปสงฺกมิตฺวา เอวํ วเทติ “เอถ มํ อาวุโส ทูเสถา”ติ. หตฺถิภณฺฑา อสฺสภณฺฑา ทูเสสุํ, เต อุชฺฌายนฺติ ขิยฺยนฺติ วิปาเจนฺติ “ปณฺฑกา อิเม สมณา สกฺยปุตฺติยา, เยปิ อิเมสํ น ปณฺฑกา, เตปิ อิเม ปณฺฑเก ทูเสนฺติ, เอวํ อิเม สพฺเพว อพฺรหฺมจาริโน”ติ. อสฺโสสุํ โข ภิกฺขู เตสํ หตฺถิ\-ภณฺฑานํ อสฺสภณฺฑานํ อุชฺฌายนฺตานํ ขิยฺยนฺตานํ วิปาเจนฺตานํ. อถ โข เต ภิกฺขู ภควโต เอตมตฺถํ อาโรเจสุํ. ปณฺฑโก ภิกฺขเว อนุปสมฺปนฺโน น อุปสมฺ\-ปาเท\-ตพฺโพ, อุปสมฺปนฺโน นาเสตพฺโพติ. (1)}

\csromanmatikaanchor{mtk.49}
\csromantocmarkref{section}{48. เถยฺยสํวาสกวตฺถุ}{mtk.49}
\bookmark[dest={mtk.49},level=1]{48. เถยฺยสํวาสกวตฺถุ}
\csromanheader{1}{48. \textbf{เถยฺยสํวาสก}\-\textbf{วตฺถุ}}
\proseitem{110}{\csromanfolio{120}เตน โข ปน สมเยน อญฺญตโร ปุราณ\-กุลปุตฺโต ขีณโกลญฺโญ สุขุมาโล โหติ. อถ โข ตสฺส ปุราณ\-กุลปุตฺตสฺส ขีณโกลญฺญสฺส เอตทโหสิ “อหํ โข สุขุมาโล, น ปฏิพโล อนธิคตํ วา โภคํ อธิคนฺตุํ, อธิคตํ วา โภคํ ผาติํ กาตุํ, เกน นุ โข อหํ อุปาเยน สุขญฺจ ชีเวยฺยํ, น จ กิลเมยฺยนฺ”ติ. อถ โข ตสฺส ปุราณ\-กุลปุตฺตสฺส ขีณโกลญฺญสฺส เอตทโหสิ “อิเม โข สมณา สกฺยปุตฺติยา สุขสีลา สุขสมาจารา สุโภชนานิ ภุญฺชิตฺวา นิวาเตสุ สยเนสุ สยนฺติ, ยนฺนูนาหํ สามํ ปตฺตจีวรํ ปฏิยาเทตฺวา เกสมสฺสุํ โอหาเรตฺวา กาสายานิ วตฺถานิ อจฺฉาเทตฺวา อารามํ คนฺตฺวา ภิกฺขูหิ สทฺธิํ สํวเสยฺยนฺ”ติ. อถ โข โส ปุราณ\-กุลปุตฺโต ขีณโกลญฺโญ สามํ ปตฺตจีวรํ ปฏิยาเทตฺวา เกสมสฺสุํ โอหาเรตฺวา กาสายานิ วตฺถานิ อจฺฉาเทตฺวา อารามํ คนฺตฺวา ภิกฺขู อภิวาเทติ. ภิกฺขู เอวมาหํสุ “กติวสฺโสสิ ตฺวํ อาวุโส”ติ. กิํ เอตํ อาวุโส กติวสฺโส นามาติ. โก ปน เต อาวุโส อุปชฺฌาโยติ. กิํ เอตํ อาวุโส อุปชฺฌาโย นามาติ. ภิกฺขู อายสฺมนฺตํ อุปาลิํ เอตทโวจุํ “อิงฺฆาวุโส อุปาลิ อิมํ ปพฺพชิตํ อนุยุญฺชาหี”ติ. อถ โข โส ปุราณ\-กุลปุตฺโต ขีณโกลญฺโญ อายสฺมตา อุปาลินา อนุยุญฺชียมาโน เอตมตฺถํ อาโรเจสิ. อายสฺมา อุปาลิ ภิกฺขูนํ เอตมตฺถํ อาโรเจสิ. ภิกฺขู ภควโต เอตมตฺถํ อาโรเจสุํ. เถยฺย\-สํวาสโก ภิกฺขเว อนุปสมฺปนฺโน น อุปสมฺ\-ปาเท\-ตพฺโพ, อุปสมฺปนฺโน นาเสตพฺโพติ. ติตฺถิย\-ปกฺกนฺตโก ภิกฺขเว อนุปสมฺปนฺโน น อุปสมฺ\-ปาเท\-ตพฺโพ, อุปสมฺปนฺโน นาเสตพฺโพติ. \mbox{(2-3)}}

\csromanmatikaanchor{mtk.50}
\csromantocmarkref{section}{49. ติรจฺฉานคตวตฺถุ}{mtk.50}
\bookmark[dest={mtk.50},level=1]{49. ติรจฺฉานคตวตฺถุ}
\csromanheader{1}{49. \textbf{ติรจฺฉานคต}\-\textbf{วตฺถุ}}
\proseitem{111}{เตน โข ปน สมเยน อญฺญตโร นาโค นาคโยนิยา อฏฺฏียติ หรายติ ชิคุจฺฉติ. อถ โข ตสฺส นาคสฺส เอตทโหสิ “เกน นุ โข อหํ อุปาเยน นาคโยนิยา จ ปริมุจฺเฉยฺยํ, ขิปฺปญฺจ มนุสฺสตฺตํ ปฏิลเภยฺยนฺ”ติ. อถ โข ตสฺส นาคสฺส เอตทโหสิ “อิเม โข สมณา สกฺยปุตฺติยา ธมฺมจาริโน สมจาริโน พฺรหฺมจาริโน \csromanfolio{121}สจฺจวาทิโน สีลวนฺโต กลฺยาณธมฺมา, สเจ โข อหํ สมเณสุ สกฺยปุตฺติเยสุ ปพฺพเชยฺยํ, เอวาหํ นาคโยนิยา จ ปริมุจฺเจยฺยํ, ขิปฺปญฺจ มนุสฺสตฺตํ ปฏิลเภยฺยนฺ”ติ. อถ โข โส นาโค มาณวก\-วณฺเณน ภิกฺขู อุปสงฺกมิตฺวา ปพฺพชฺชํ ยาจิ, ตํ ภิกฺขู ปพฺพาเชสุํ อุปสมฺปาเทสุํ. เตน โข ปน สมเยน โส นาโค อญฺญตเรน ภิกฺขุนา สทฺธิํ ปจฺจนฺติเม วิหาเร ปฏิวสติ. อถ โข โส ภิกฺขุ รตฺติยา ปจฺจูสสมยํ ปจฺจุฏฺฐาย อชฺโฌกาเส จงฺกมติ. อถ โข โส นาโค ตสฺส ภิกฺขุโน นิกฺขนฺเต วิสฺสฏฺโฐ นิทฺทํ โอกฺกมิ, สพฺโพ วิหาโร อหินา ปุณฺโณ, วาตปาเนหิ โภคา นิกฺขนฺตา โหนฺติ. อถ โข โส ภิกฺขุ “วิหารํ ปวิสิสฺสามี”ติ กวาฏํ ปณาเมนฺโต อทฺทส สพฺพํ วิหารํ อหินา ปุณฺณํ วาตปาเนหิ โภเค นิกฺขนฺเต, ทิสฺวาน ภีโต วิสฺสรมกาสิ. ภิกฺขู อุปธาวิตฺวา ตํ ภิกฺขุํ เอตทโวจุํ “กิสฺส ตฺวํ อาวุโส วิสฺสรมกาสี”ติ. อยํ อาวุโส สพฺโพ วิหาโร อหินา ปุณฺโณ, วาตปาเนหิ โภคา นิกฺขนฺตาติ. อถ โข โส นาโค เตน สทฺเทน ปฏิพุชฺฌิตฺวา สเก อาสเน นิสีทิ. ภิกฺขู เอวมาหํสุ “โกสิ ตฺวํ อาวุโส”ติ. อหํ ภนฺเต นาโคติ. กิสฺส ปน ตฺวํ อาวุโส เอวรูปํ อกาสีติ. อถ โข โส นาโค ภิกฺขูนํ เอตมตฺถํ อาโรเจสิ. ภิกฺขู ภควโต เอตมตฺถํ อาโรเจสุํ. อถ โข ภควา เอตสฺมิํ นิทาเน เอตสฺมิํ ปกรเณ ภิกฺขุสํฆํ สนฺนิปาตาเปตฺวา ตํ นาคํ เอตทโวจ “ตุเมฺห โขตฺถ นาคา อวิรุฬฺหิธมฺมา อิมสฺมิํ ธมฺมวินเย, คจฺฉ ตฺวํ นาค, ตตฺเถว จาตุทฺทเส ปนฺนรเส อฏฺฐมิยา จ ปกฺขสฺส อุโปสถํ อุปวส, เอวํ ตฺวํ นาคโยนิยา จ ปริมุจฺจิสฺสสิ, ขิปฺปญฺจ มนุสฺสตฺตํ ปฏิลภิสฺสสี”ติ. อถ โข โส นาโค “อวิรุฬฺหิธมฺโม กิราหํ อิมสฺมิํ ธมฺมวินเย”ติ ทุกฺขี ทุมฺมโน อสฺสูนิ ปวตฺตยมาโน วิสฺสรํ กตฺวา ปกฺกามิ. อถ โข ภควา ภิกฺขู อามนฺเตสิ “เทฺวเม ภิกฺขเว ปจฺจยา นาคสฺส สภาว\-ปาตุกมฺมาย, ยทา จ สชาติยา เมถุนํ ธมฺมํ ปฏิเสวติ, ยทา จ วิสฺสฏฺโฐ นิทฺทํ โอกฺกมติ, อิเม โข ภิกฺขเว เทฺว ปจฺจยา นาคสฺส สภาว\-ปาตุกมฺมาย. ติรจฺฉานคโต ภิกฺขเว อนุปสมฺปนฺโน น อุปสมฺ\-ปาเท\-ตพฺโพ, อุปสมฺปนฺโน นาเสตพฺโพ”ติ. (4)}

\csromanmatikaanchor{mtk.51}
\csromantocmarkref{section}{50. มาตุฆาตกวตฺถุ}{mtk.51}
\bookmark[dest={mtk.51},level=1]{50. มาตุฆาตกวตฺถุ}
\csromanheader{1}{50. \textbf{มาตุฆาตก}\-\textbf{วตฺถุ}}
\proseitem{112}{เตน โข ปน สมเยน อญฺญตโร มาณวโก มาตรํ ชีวิตา โวโรเปสิ, โส เตน ปาปเกน กมฺเมน อฏฺฏียติ หรายติ \csromanfolio{122}ชิคุจฺฉติ. อถ โข ตสฺส มาณวกสฺส เอตทโหสิ “เกน นุ โข อหํ อุปาเยน อิมสฺส ปาปกสฺส กมฺมสฺส นิกฺขนฺติํ กเรยฺยนฺ”ติ. อถ โข ตสฺส มาณวกสฺส เอตทโหสิ “อิเม โข สมณา สกฺยปุตฺติยา ธมฺมจาริโน สมจาริโน พฺรหฺมจาริโน สจฺจวาทิโน สีลวนฺโต กลฺยาณธมฺมา, สเจ โข อหํ สมเณสุ สกฺยปุตฺติเยสุ ปพฺพเชยฺยํ, เอวาหํ อิมสฺส ปาปกสฺส กมฺมสฺส นิกฺขนฺติํ กเรยฺยนฺ”ติ. อถ โข โส มาณวโก ภิกฺขู อุปสงฺกมิตฺวา ปพฺพชฺชํ ยาจิ. ภิกฺขู อายสฺมนฺตํ อุปาลิํ เอตทโวจุํ “ปุพฺเพปิ โข อาวุโส อุปาลิ นาโค มาณวก\-วณฺเณน ภิกฺขูสุ ปพฺพชิโต, อิงฺฆาวุโส อุปาลิ อิมํ มาณวกํ อนุยุญฺชาหี”ติ. อถ โข โส มาณวโก อายสฺมตา อุปาลินา อนุยุญฺชียมาโน เอตมตฺถํ อาโรเจสิ. อายสฺมา อุปาลิ ภิกฺขูนํ เอตมตฺถํ อาโรเจสิ. ภิกฺขู ภควโต เอตมตฺถํ อาโรเจสุํ ฯเปฯ. มาตุฆาตโก ภิกฺขเว อนุปสมฺปนฺโน น อุปสมฺ\-ปาเท\-ตพฺโพ, อุปสมฺปนฺโน นาเสตพฺโพติ. (5)}

\csromanmatikaanchor{mtk.52}
\csromantocmarkref{section}{51. ปิตุฆาตกวตฺถุ}{mtk.52}
\bookmark[dest={mtk.52},level=1]{51. ปิตุฆาตกวตฺถุ}
\csromanheader{1}{51. \textbf{ปิตุฆาตก}\-\textbf{วตฺถุ}}
\proseitem{113}{เตน โข ปน สมเยน อญฺญตโร มาณวโก ปิตรํ ชีวิตา โวโรเปสิ, โส เตน ปาปเกน กมฺเมน อฏฺฏียติ หรายติ ชิคุจฺฉติ. อถ โข ตสฺส มาณวกสฺส เอตทโหสิ “เกน นุ โข อหํ อุปาเยน อิมสฺส ปาปกสฺส กมฺมสฺส นิกฺขนฺติํ กเรยฺยนฺ”ติ. อถ โข ตสฺส มาณวกสฺส เอตทโหสิ “อิเม โข สมณา สกฺยปุตฺติยา ธมฺมจาริโน สมจาริโน พฺรหฺมจาริโน สจฺจวาทิโน สีลวนฺโต กลฺยาณธมฺมา, สเจ โข อหํ สมเณสุ สกฺยปุตฺติเยสุ ปพฺพเชยฺยํ, เอวาหํ อิมสฺส ปาปกสฺส กมฺมสฺส นิกฺขนฺติํ กเรยฺยนฺ”ติ. อถ โข โส มาณวโก ภิกฺขู อุปสงฺกมิตฺวา ปพฺพชฺชํ ยาจิ. ภิกฺขู อายสฺมนฺตํ อุปาลิํ เอตทโวจุํ “ปุพฺเพปิ โข อาวุโส อุปาลิ นาโค มาณวก\-วณฺเณน ภิกฺขูสุ ปพฺพชิโต, อิงฺฆาวุโส อุปาลิ อิมํ มาณวกํ อนุยุญฺชาหี”ติ. อถ โข โส มาณวโก อายสฺมตา อุปาลินา อนุยุญฺชียมาโน เอตมตฺถํ อาโรเจสิ. อายสฺมา อุปาลิ ภิกฺขูนํ เอตมตฺถํ อาโรเจสิ. ภิกฺขู ภควโต เอตมตฺถํ อาโรเจสุํ. ปิตุฆาตโก ภิกฺขเว อนุปสมฺปนฺโน น อุปสมฺ\-ปาเท\-ตพฺโพ, อุปสมฺปนฺโน นาเสตพฺโพติ. (6)}

\csromanmatikaanchor{mtk.53}
\csromantocmarkref{section}{52. อรหนฺตฆาตกวตฺถุ}{mtk.53}
\bookmark[dest={mtk.53},level=1]{52. อรหนฺตฆาตกวตฺถุ}
\csromanheader{1}{52. \textbf{อรหนฺตฆาตก}\-\textbf{วตฺถุ}}
\proseitem{114}{\csromanfolio{123}เตน โข ปน สมเยน สมฺพหุลา ภิกฺขู สาเกตา สาวตฺถิํ อทฺธาน\-มคฺค\-ปฺปฏิปนฺนา โหนฺติ. อนฺตรามคฺเค โจรา นิกฺขมิตฺวา เอกจฺเจ ภิกฺขู อจฺฉินฺทิํสุ, เอกจฺเจ ภิกฺขู หนิํสุ. สาวตฺถิยา ราชภฏา นิกฺขมิตฺวา เอกจฺเจ โจเร อคฺคเหสุํ, เอกจฺเจ โจรา ปลายิํสุ. เย เต ปลายิํสุ, เต ภิกฺขูสุ ปพฺพชิํสุ. เย เต คหิตา, เต วธาย โอนิยฺยนฺติ. อทฺทสํสุ โข เต ปลายิตฺวา ปพฺพชิตา เต โจเร วธาย โอนิยฺยมาเน, ทิสฺวาน เอวมาหํสุ “สาธุ โข มยํ ปลายิมฺหา, สจา จ\csromansharedfootnote{0367d20c65d468b8}{สเจ จ, สจชฺช (อฏฺฐกถายํ ปาฐนฺตรา)} มยํ คเยฺหยฺยาม\csromansharedfootnote{6293d7182e70f87e}{คเณฺหยฺยาม (ก)}, มยมฺปิ เอวเมว หญฺเญยฺยามา”ติ. ภิกฺขู เอวมาหํสุ “กิํ ปน ตุเมฺห อาวุโส อกตฺถา”ติ. อถ โข เต ปพฺพชิตา ภิกฺขูนํ เอตมตฺถํ อาโรเจสุํ. ภิกฺขู ภควโต เอตมตฺถํ อาโรเจสุํ. อรหนฺโต เอเต ภิกฺขเว ภิกฺขู, อรหนฺต\-ฆาตโก ภิกฺขเว อนุปสมฺปนฺโน น อุปสมฺ\-ปาเท\-ตพฺโพ, อุปสมฺปนฺโน นาเสตพฺโพติ. (7)}

\csromanmatikaanchor{mtk.54}
\csromantocmarkref{section}{53. ภิกฺขุนีทูสกวตฺถุ}{mtk.54}
\bookmark[dest={mtk.54},level=1]{53. ภิกฺขุนีทูสกวตฺถุ}
\csromanheader{1}{53. \textbf{ภิกฺขุนี}\-\textbf{ทูสก}\-\textbf{วตฺถุ}}
\proseitem{115}{เตน โข ปน สมเยน สมฺพหุลา ภิกฺขุนิโย สาเกตา สาวตฺถิํ อทฺธาน\-มคฺค\-ปฺปฏิปนฺนา โหนฺติ. อนฺตรามคฺเค โจรา นิกฺขมิตฺวา เอกจฺจา ภิกฺขุนิโย อจฺฉินฺทิํสุ, เอกจฺจา ภิกฺขุนิโย ทูเสสุํ. สาวตฺถิยา ราชภฏา นิกฺขมิตฺวา เอกจฺเจ โจเร อคฺคเหสุํ, เอกจฺเจ โจรา ปลายิํสุ. เย เต ปลายิํสุ, เต ภิกฺขูสุ ปพฺพชิํสุ. เย เต คหิตา, เต วธาย โอนิยฺยนฺติ. อทฺทสํสุ โข เต ปลายิตฺวา ปพฺพชิตา เต โจเร วธาย โอนิยฺยมาเน, ทิสฺวาน เอวมาหํสุ “สาธุ โข มยํ ปลายิมฺหา, สจา จ มยํ คเยฺหยฺยาม, มยมฺปิ เอวเมว หญฺเญยฺยามา”ติ. ภิกฺขู เอวมาหํสุ “กิํ ปน ตุเมฺห อาวุโส อกตฺถา”ติ. อถ โข เต ปพฺพชิตา ภิกฺขูนํ เอตมตฺถํ อาโรเจสุํ. ภิกฺขู ภควโต เอตมตฺถํ อาโรเจสุํ. ภิกฺขุนิ\-ทูสโก ภิกฺขเว อนุปสมฺปนฺโน น อุปสมฺ\-ปาเท\-ตพฺโพ, อุปสมฺปนฺโน นาเสตพฺโพติ. สํฆเภทโก ภิกฺขเว อนุปสมฺปนฺโน น อุปสมฺ\-ปาเท\-ตพฺโพ, อุปสมฺปนฺโน นาเสตพฺโพติ. โลหิตุปฺปาทโก ภิกฺขเว อนุปสมฺปนฺโน น อุปสมฺ\-ปาเท\-ตพฺโพ, อุปสมฺปนฺโน นาเสตพฺโพติ. (8-9-10)}

\csromanmatikaanchor{mtk.55}
\csromantocmarkref{section}{54. อุภโตพฺยญฺชนกวตฺถุ}{mtk.55}
\bookmark[dest={mtk.55},level=1]{54. อุภโตพฺยญฺชนกวตฺถุ}
\csromanheader{1}{54. \textbf{อุภโตพฺยญฺชนก}\-\textbf{วตฺถุ}}
\proseitem{116}{\csromanfolio{124}เตน โข ปน สมเยน อญฺญตโร อุภโต\-พฺยญฺชนโก ภิกฺขูสุ ปพฺพชิโต โหติ, โส กโรติปิ การาเปติปิ. ภควโต เอตมตฺถํ อาโรเจสุํ. อุภโต\-พฺยญฺชนโก ภิกฺขเว อนุปสมฺปนฺโน น อุปสมฺ\-ปาเท\-ตพฺโพ, อุปสมฺปนฺโน นาเสตพฺโพติ. (11)}

\csromanmatikaanchor{mtk.56}
\csromantocmarkref{section}{55. อนุปชฺฌายกาทิวตฺถูนิ}{mtk.56}
\bookmark[dest={mtk.56},level=1]{55. อนุปชฺฌายกาทิวตฺถูนิ}
\csromanheader{1}{55. \textbf{อนุปชฺฌายกา}\-\textbf{ทิ}\-\textbf{วตฺถูนิ}}
\proseitem{117}{เตน โข ปน สมเยน ภิกฺขู อนุปชฺฌายกํ อุปสมฺปาเทนฺติ. ภควโต เอตมตฺถํ อาโรเจสุํ. น ภิกฺขเว อนุปชฺฌายโก อุปสมฺ\-ปาเท\-ตพฺโพ, โย อุปสมฺปาเทยฺย, อาปตฺติ ทุกฺกฏสฺสาติ. (12)}

\prose{เตน โข ปน สมเยน ภิกฺขู สํเฆน อุปชฺฌาเยน อุปสมฺปาเทนฺติ. ภควโต เอตมตฺถํ อาโรเจสุํ. น ภิกฺขเว สํเฆน อุปชฺฌาเยน อุปสมฺ\-ปาเท\-ตพฺโพ, โย อุปสมฺปาเทยฺย, อาปตฺติ ทุกฺกฏสฺสาติ. (13)}

\prose{เตน โข ปน สมเยน ภิกฺขู คเณน อุปชฺฌาเยน อุปสมฺปาเทนฺติ. ภควโต เอตมตฺถํ อาโรเจสุํ. น ภิกฺขเว คเณน อุปชฺฌาเยน อุปสมฺ\-ปาเท\-ตพฺโพ, โย อุปสมฺปาเทยฺย, อาปตฺติ ทุกฺกฏสฺสาติ. (14)}

\prose{เตน โข ปน สมเยน ภิกฺขู ปณฺฑกุ\-ปชฺฌาเยน อุปสมฺปาเทนฺติ ฯเปฯ เถยฺยสํวาสกุ\-ปชฺฌาเยน อุปสมฺปาเทนฺติ ฯเปฯ ติตฺถิยปกฺกนฺตกุ\-ปชฺฌาเยน อุปสมฺปาเทนฺติ ฯเปฯ ติรจฺฉานคตุ\-ปชฺฌาเยน อุปสมฺปาเทนฺติ ฯเปฯ มาตุฆาตกุ\-ปชฺฌาเยน อุปสมฺปาเทนฺติ ฯเปฯ ปิตุฆาตกุ\-ปชฺฌาเยน อุปสมฺปาเทนฺติ ฯเปฯ อรหนฺตฆาตกุ\-ปชฺฌาเยน อุปสมฺปาเทนฺติ ฯเปฯ ภิกฺขุนิทูสกุ\-ปชฺฌาเยน อุปสมฺปาเทนฺติ ฯเปฯ สํฆเภทกุปชฺฌาเยน อุปสมฺปาเทนฺติ ฯเปฯ โลหิตุปฺปาทกุ\-ปชฺฌาเยน อุปสมฺปาเทนฺติ ฯเปฯ อุภโตพฺยญฺชนกุ\-ปชฺฌาเยน อุปสมฺปาเทนฺติ. ภควโต เอตมตฺถํ อาโรเจสุํ. น ภิกฺขเว ปณฺฑกุ\-ปชฺฌาเยน อุปสมฺ\-ปาเท\-ตพฺโพ ฯเปฯ. น ภิกฺขเว เถยฺยสํวาสกุ\-ปชฺฌาเยน อุปสมฺ\-ปาเท\-ตพฺโพ ฯเปฯ. น ภิกฺขเว ติตฺถิยปกฺกนฺตกุ\-ปชฺฌาเยน อุปสมฺ\-ปาเท\-ตพฺโพ ฯเปฯ. น ภิกฺขเว ติรจฺฉานคตุ\-ปชฺฌาเยน อุปสมฺ\-ปาเท\-ตพฺโพ ฯเปฯ. น ภิกฺขเว มาตุฆาตกุ\-ปชฺฌาเยน อุปสมฺ\-ปาเท\-ตพฺโพ ฯเปฯ. น ภิกฺขเว ปิตุฆาตกุ\-ปชฺฌาเยน อุปสมฺ\-ปาเท\-ตพฺโพ ฯเปฯ. น ภิกฺขเว อรหนฺตฆาตกุ\-ปชฺฌาเยน อุปสมฺ\-ปาเท\-ตพฺโพ ฯเปฯ. น ภิกฺขเว ภิกฺขุนิทูสกุ\-ปชฺฌาเยน อุปสมฺ\-ปาเท\-ตพฺโพ \csromanfolio{125}ฯเปฯ. น ภิกฺขเว สํฆเภทกุปชฺฌาเยน อุปสมฺ\-ปาเท\-ตพฺโพ ฯเปฯ. น ภิกฺขเว โลหิตุปฺปาทกุ\-ปชฺฌาเยน อุปสมฺ\-ปาเท\-ตพฺโพ ฯเปฯ. น ภิกฺขเว อุภโตพฺยญฺชนกุ\-ปชฺฌาเยน อุปสมฺ\-ปาเท\-ตพฺโพ, โย อุปสมฺปาเทยฺย, อาปตฺติ ทุกฺกฏสฺสาติ. (15)}

\csromanmatikaanchor{mtk.57}
\csromantocmarkref{section}{56. อปตฺตกาทิวตฺถุ}{mtk.57}
\bookmark[dest={mtk.57},level=1]{56. อปตฺตกาทิวตฺถุ}
\csromanheader{1}{56. \textbf{อปตฺตกา}\-\textbf{ทิ}\-\textbf{วตฺถุ}}
\proseitem{118}{เตน โข ปน สมเยน ภิกฺขู อปตฺตกํ อุปสมฺปาเทนฺติ, หตฺเถสุ ปิณฺฑาย จรนฺติ. มนุสฺสา อุชฺฌายนฺติ ขิยฺยนฺติ วิปาเจนฺติ “เสยฺยถาปิ ติตฺถิยา”ติ. ภควโต เอตมตฺถํ อาโรเจสุํ. น ภิกฺขเว อปตฺตโก อุปสมฺ\-ปาเท\-ตพฺโพ, โย อุปสมฺปาเทยฺย, อาปตฺติ ทุกฺกฏสฺสาติ. (16)}

\prose{เตน โข ปน สมเยน ภิกฺขู อจีวรกํ อุปสมฺปาเทนฺติ, นคฺคา ปิณฺฑาย จรนฺติ. มนุสฺสา อุชฺฌายนฺติ ขิยฺยนฺติ วิปาเจนฺติ “เสยฺยถาปิ ติตฺถิยา”ติ. ภควโต เอตมตฺถํ อาโรเจสุํ. น ภิกฺขเว อจีวรโก อุปสมฺ\-ปาเท\-ตพฺโพ, โย อุปสมฺปาเทยฺย, อาปตฺติ ทุกฺกฏสฺสาติ. (17)}

\prose{เตน โข ปน สมเยน ภิกฺขู อปตฺตจีวรกํ อุปสมฺปาเทนฺติ, นคฺคา หตฺเถสุ ปิณฺฑาย จรนฺติ. มนุสฺสา อุชฺฌายนฺติ ขิยฺยนฺติ วิปาเจนฺติ “เสยฺยถาปิ ติตฺถิยา”ติ. ภควโต เอตมตฺถํ อาโรเจสุํ. น ภิกฺขเว อปตฺตจีวรโก อุปสมฺ\-ปาเท\-ตพฺโพ, โย อุปสมฺปาเทยฺย, อาปตฺติ ทุกฺกฏสฺสาติ. (18)}

\prose{เตน โข ปน สมเยน ภิกฺขู ยาจิตเกน ปตฺเตน อุปสมฺปาเทนฺติ, อุปสมฺปนฺเน ปตฺตํ ปฏิหรนฺติ, หตฺเถสุ ปิณฺฑาย จรนฺติ. มนุสฺสา อุชฺฌายนฺติ ขิยฺยนฺติ วิปาเจนฺติ “เสยฺยถาปิ ติตฺถิยา”ติ. ภควโต เอตมตฺถํ อาโรเจสุํ. น ภิกฺขเว ยาจิตเกน ปตฺเตน อุปสมฺ\-ปาเท\-ตพฺโพ, โย อุปสมฺปาเทยฺย, อาปตฺติ ทุกฺกฏสฺสาติ. (19)}

\prose{เตน โข ปน สมเยน ภิกฺขู ยาจิตเกน จีวเรน อุปสมฺปาเทนฺติ, อุปสมฺปนฺเน จีวรํ ปฏิหรนฺติ, นคฺคา ปิณฺฑาย จรนฺติ. มนุสฺสา อุชฺฌายนฺติ ขิยฺยนฺติ วิปาเจนฺติ “เสยฺยถาปิ ติตฺถิยา”ติ. ภควโต เอตมตฺถํ อาโรเจสุํ. น ภิกฺขเว ยาจิตเกน จีวเรน อุปสมฺ\-ปาเท\-ตพฺโพ, โย อุปสมฺปาเทยฺย, อาปตฺติ ทุกฺกฏสฺสาติ. (20)}

\prosewithcloserruled{\csromanfolio{126}เตน โข ปน สมเยน ภิกฺขู ยาจิตเกน ปตฺตจีวเรน อุปสมฺปาเทนฺติ, อุปสมฺปนฺเน ปตฺตจีวรํ ปฏิหรนฺติ, นคฺคา หตฺเถสุ ปิณฺฑาย จรนฺติ. มนุสฺสา อุชฺฌายนฺติ ขิยฺยนฺติ วิปาเจนฺติ “เสยฺยถาปิ ติตฺถิยา”ติ. ภควโต เอตมตฺถํ อาโรเจสุํ. น ภิกฺขเว ยาจิตเกน ปตฺตจีวเรน อุปสมฺ\-ปาเท\-ตพฺโพ, โย อุปสมฺปาเทยฺย, อาปตฺติ ทุกฺกฏสฺสาติ. (21)}{นอุปสมฺปาเทตพฺเพกวีสติวาโร นิฏฺฐิโต.}
\csromanmatikaanchor{mtk.58}
\csromantocmarkref{section}{57. นปพฺพาเชตพฺพทฺวตฺติํสวาร}{mtk.58}
\bookmark[dest={mtk.58},level=1]{57. นปพฺพาเชตพฺพทฺวตฺติํสวาร}
\csromanheader{1}{57. \textbf{นปพฺพาเชตพฺพ}\-\textbf{ทฺวตฺติํส}\-\textbf{วาร}}
\proseitemwithcloser{119}{เตน โข ปน สมเยน ภิกฺขู หตฺถจฺฉินฺนํ ปพฺพาเชนฺติ ฯเปฯ ปาทจฺฉินฺนํ ปพฺพาเชนฺติ ฯเปฯ หตฺถ\-ปาท\-จฺฉินฺนํ ปพฺพาเชนฺติ ฯเปฯ กณฺณจฺฉินฺนํ ปพฺพาเชนฺติ ฯเปฯ นาสจฺฉินฺนํ ปพฺพาเชนฺติ ฯเปฯ กณฺณ\-นาส\-จฺฉินฺนํ ปพฺพาเชนฺติ ฯเปฯ องฺคุลิจฺฉินฺนํ ปพฺพาเชนฺติ ฯเปฯ อฬจฺฉินฺนํ ปพฺพาเชนฺติ ฯเปฯ กณฺฑร\-จฺฉินฺนํ ปพฺพาเชนฺติ ฯเปฯ ผณหตฺถกํ ปพฺพาเชนฺติ ฯเปฯ ขุชฺชํ ปพฺพาเชนฺติ ฯเปฯ วามนํ ปพฺพาเชนฺติ ฯเปฯ คลคณฺฑิํ ปพฺพาเชนฺติ ฯเปฯ ลกฺขณาหตํ ปพฺพาเชนฺติ ฯเปฯ กสาหตํ ปพฺพาเชนฺติ ฯเปฯ ลิขิตกํ ปพฺพาเชนฺติ ฯเปฯ สีปทิํ ปพฺพาเชนฺติ ฯเปฯ ปาปโรคิํ ปพฺพาเชนฺติ ฯเปฯ ปริสทูสกํ ปพฺพาเชนฺติ ฯเปฯ กาณํ ปพฺพาเชนฺติ ฯเปฯ กุณิํ ปพฺพาเชนฺติ ฯเปฯ ขญฺชํ ปพฺพาเชนฺติ ฯเปฯ ปกฺขหตํ ปพฺพาเชนฺติ ฯเปฯ ฉินฺนิริยาปถํ ปพฺพาเชนฺติ ฯเปฯ ชราทุพฺพลํ ปพฺพาเชนฺติ ฯเปฯ อนฺธํ ปพฺพาเชนฺติ ฯเปฯ มูคํ ปพฺพาเชนฺติ ฯเปฯ พธิรํ ปพฺพาเชนฺติ ฯเปฯ อนฺธมูคํ ปพฺพาเชนฺติ ฯเปฯ อนฺธพธิรํ ปพฺพาเชนฺติ ฯเปฯ มูคพธิรํ ปพฺพาเชนฺติ ฯเปฯ อนฺธ\-มูค\-พธิรํ ปพฺพาเชนฺติ. ภควโต เอตมตฺถํ อาโรเจสุํ ฯเปฯ. น ภิกฺขเว หตฺถจฺฉินฺโน ปพฺพาเชตพฺโพ ฯเปฯ. น ภิกฺขเว ปาทจฺฉินฺโน ปพฺพาเชตพฺโพ ฯเปฯ. น ภิกฺขเว หตฺถ\-ปาท\-จฺฉินฺโน ปพฺพาเชตพฺโพ ฯเปฯ. น ภิกฺขเว กณฺณจฺฉินฺโน ปพฺพาเชตพฺโพ ฯเปฯ. น ภิกฺขเว นาสจฺฉินฺโน ปพฺพาเชตพฺโพ ฯเปฯ. น ภิกฺขเว กณฺณ\-นาส\-จฺฉินฺโน ปพฺพาเชตพฺโพ ฯเปฯ. น ภิกฺขเว องฺคุลิจฺฉินฺโน ปพฺพาเชตพฺโพ ฯเปฯ. น ภิกฺขเว อฬจฺฉินฺโน ปพฺพาเชตพฺโพ ฯเปฯ. น ภิกฺขเว กณฺฑรจฺฉินฺโน ปพฺพาเชตพฺโพ ฯเปฯ. น ภิกฺขเว ผณหตฺถโก ปพฺพาเชตพฺโพ ฯเปฯ. น ภิกฺขเว ขุชฺโช ปพฺพาเชตพฺโพ ฯเปฯ. น ภิกฺขเว วามโน ปพฺพาเชตพฺโพ ฯเปฯ. น ภิกฺขเว คลคณฺฑี ปพฺพาเชตพฺโพ ฯเปฯ. น ภิกฺขเว \csromanfolio{127}ลกฺขณาหโต ปพฺพาเชตพฺโพ ฯเปฯ. น ภิกฺขเว กสาหโต ปพฺพาเชตพฺโพ ฯเปฯ. น ภิกฺขเว ลิขิตโก ปพฺพาเชตพฺโพ ฯเปฯ. น ภิกฺขเว สีปที ปพฺพาเชตพฺโพ ฯเปฯ. น ภิกฺขเว ปาปโรคี ปพฺพาเชตพฺโพ ฯเปฯ. น ภิกฺขเว ปริสทูสโก ปพฺพาเชตพฺโพ ฯเปฯ. น ภิกฺขเว กาโณ ปพฺพาเชตพฺโพ ฯเปฯ. น ภิกฺขเว กุณี ปพฺพาเชตพฺโพ ฯเปฯ. น ภิกฺขเว ขญฺโช ปพฺพาเชตพฺโพ ฯเปฯ. น ภิกฺขเว ปกฺขหโต ปพฺพาเชตพฺโพ ฯเปฯ. น ภิกฺขเว ฉินฺนิริยาปโถ ปพฺพาเชตพฺโพ ฯเปฯ. น ภิกฺขเว ชราทุพฺพโล ปพฺพาเชตพฺโพ ฯเปฯ. น ภิกฺขเว อนฺโธ ปพฺพาเชตพฺโพ ฯเปฯ. น ภิกฺขเว มูโค ปพฺพาเชตพฺโพ ฯเปฯ. น ภิกฺขเว พธิโร ปพฺพาเชตพฺโพ ฯเปฯ. น ภิกฺขเว อนฺธมูโค ปพฺพาเชตพฺโพ ฯเปฯ. น ภิกฺขเว อนฺธพธิโร ปพฺพาเชตพฺโพ ฯเปฯ. น ภิกฺขเว มูคพธิโร ปพฺพาเชตพฺโพ ฯเปฯ. น ภิกฺขเว อนฺธ\-มูค\-พธิโร ปพฺพาเชตพฺโพ, โย ปพฺพาเชยฺย, อาปตฺติ ทุกฺกฏสฺสาติ.}{นปพฺพาเชตพฺพ\-ทฺวตฺติํส\-วาโร นิฏฺฐิโต.}
\nitthitamruled{ทายชฺช\-ภาณวาโร นิฏฺฐิโต นวโม.}
\csromanmatikaanchor{mtk.59}
\csromantocmarkref{section}{58. อลชฺชีนิสฺสยวตฺถูนิ}{mtk.59}
\bookmark[dest={mtk.59},level=1]{58. อลชฺชีนิสฺสยวตฺถูนิ}
\csromanheader{1}{58. \textbf{อลชฺชี}\-\textbf{นิสฺสย}\-\textbf{วตฺถูนิ}}
\proseitem{120}{เตน โข ปน สมเยน ฉพฺพคฺคิยา ภิกฺขู อลชฺชีนํ นิสฺสยํ เทนฺติ. ภควโต เอตมตฺถํ อาโรเจสุํ. น ภิกฺขเว อลชฺชีนํ นิสฺสโย ทาตพฺโพ, โย ทเทยฺย, อาปตฺติ ทุกฺกฏสฺสาติ.}

\prose{เตน โข ปน สมเยน ภิกฺขู อลชฺชีนํ นิสฺสาย วสนฺติ, เตปิ นจิรสฺเสว อลชฺชิโน โหนฺติ ปาปกา.}

\prose{ภิกฺขู ภควโต เอตมตฺถํ อาโรเจสุํ. น ภิกฺขเว อลชฺชีนํ นิสฺสาย วตฺถพฺพํ, โย วเสยฺย, อาปตฺติ ทุกฺกฏสฺสาติ.}

\prose{อถ โข ภิกฺขูนํ เอตทโหสิ “ภควตา ปญฺญตฺตํ ‘น อลชฺชีนํ นิสฺสโย ทาตพฺโพ, น อลชฺชีนํ นิสฺสาย วตฺถพฺพนฺ’ติ, กถํ นุ โข มยํ ชาเนยฺยาม ลชฺชิํ วา อลชฺชิํ วา”ติ.}

\prose{ภควโต เอตมตฺถํ อาโรเจสุํ. อนุชานามิ ภิกฺขเว จตูหปญฺจาหํ อาคเมตุํ ยาว ภิกฺขุ\-สภาคตํ ชานามีติ.}

\csromanmatikaanchor{mtk.60}
\csromantocmarkref{section}{59. คมิกาทินิสฺสยวตฺถูนิ}{mtk.60}
\bookmark[dest={mtk.60},level=1]{59. คมิกาทินิสฺสยวตฺถูนิ}
\csromanheader{1}{59. \textbf{คมิกา}\-\textbf{ทิ}\-\textbf{นิสฺสย}\-\textbf{วตฺถูนิ}}
\proseitem{121}{\csromanfolio{128}เตน โข ปน สมเยน อญฺญตโร ภิกฺขุ โกสเลสุ ชนปเท อทฺธาน\-มคฺค\-ปฺปฏิปนฺโน โหติ. อถ โข ตสฺส ภิกฺขุโน เอตทโหสิ “ภควตา ปญฺญตฺตํ ‘น อนิสฺสิเตน วตฺถพฺพนฺ’ติ, อหญฺจมฺหิ นิสฺสยกรณีโย อทฺธาน\-มคฺค\-ปฺปฏิปนฺโน, กถํ นุ โข มยา ปฏิปชฺชิตพฺพนฺ”ติ. ภควโต เอตมตฺถํ อาโรเจสุํ. อนุชานามิ ภิกฺขเว อทฺธาน\-มคฺค\-ปฺปฏิปนฺเนน ภิกฺขุนา นิสฺสยํ อลภมาเนน อนิสฺสิเตน วตฺถุนฺติ.}

\prose{เตน โข ปน สมเยน เทฺว ภิกฺขู โกสเลสุ ชนปเท อทฺธาน\-มคฺค\-ปฺปฏิปนฺนา โหนฺติ, เต อญฺญตรํ อาวาสํ อุปคจฺฉิํสุ.}

\prose{ตตฺถ เอโก ภิกฺขุ คิลาโน โหติ. อถ โข ตสฺส คิลานสฺส ภิกฺขุโน เอตทโหสิ “ภควตา ปญฺญตฺตํ ‘น อนิสฺสิเตน วตฺถพฺพนฺ’ติ, อหญฺจมฺหิ นิสฺสยกรณีโย คิลาโน, กถํ นุ โข มยา ปฏิปชฺชิตพฺพนฺ”ติ. ภควโต เอตมตฺถํ อาโรเจสุํ. อนุชานามิ ภิกฺขเว คิลาเนน ภิกฺขุนา นิสฺสยํ อลภมาเนน อนิสฺสิเตน วตฺถุนฺติ.}

\prose{อถ โข ตสฺส คิลานุ\-ปฏฺฐากสฺส ภิกฺขุโน เอตทโหสิ “ภควตา ปญฺญตฺตํ ‘น อนิสฺสิเตน วตฺถพฺพนฺ’ติ, อหญฺจมฺหิ นิสฺสยกรณีโย, อยญฺจ ภิกฺขุ คิลาโน, กถํ นุ โข มยา ปฏิปชฺชิตพฺพนฺ”ติ.}

\prose{ภควโต เอตมตฺถํ อาโรเจสุํ. อนุชานามิ ภิกฺขเว คิลานุ\-ปฏฺฐาเกน ภิกฺขุนา นิสฺสยํ อลภมาเนน ยาจิยมาเนน อนิสฺสิเตน วตฺถุนฺติ.}

\prose{เตน โข ปน สมเยน อญฺญตโร ภิกฺขุ อรญฺเญ วิหรติ, ตสฺส จ ตสฺมิํ เสนาสเน ผาสุ โหติ.}

\prose{อถ โข ตสฺส ภิกฺขุโน เอตทโหสิ “ภควตา ปญฺญตฺตํ ‘น อนิสฺสิเตน วตฺถพฺพนฺ’ติ, อหญฺจมฺหิ นิสฺสยกรณีโย อรญฺเญ วิหรามิ, มยฺหญฺจ อิมสฺมิํ เสนาสเน ผาสุ โหติ, กถํ นุ โข มยา ปฏิปชฺชิตพฺพนฺ”ติ.}

\prose{ภควโต เอตมตฺถํ อาโรเจสุํ. อนุชานามิ ภิกฺขเว อารญฺญิเกน ภิกฺขุนา ผาสุวิหารํ สลฺลกฺเขนฺเตน นิสฺสยํ อลภมาเนน อนิสฺสิเตน วตฺถุํ “ยทา ปติรูโป นิสฺสยทายโก อาคจฺฉิสฺสติ, ตทา ตสฺส นิสฺสาย วสิสฺสามี”ติ.}

\csromanmatikaanchor{mtk.61}
\csromantocmarkref{section}{60. โคตฺเตน อนุสฺสาวนานุชานนา}{mtk.61}
\bookmark[dest={mtk.61},level=1]{60. โคตฺเตน อนุสฺสาวนานุชานนา}
\csromanheader{1}{60. \textbf{โคตฺเตน อนุสฺสาวนา}\-\textbf{นุชานนา}}
\proseitem{122}{\csromanfolio{129}เตน โข ปน สมเยน อายสฺมโต มหากสฺสปสฺส อุปสมฺปทา\-เปกฺโข โหติ. อถ โข อายสฺมา มหากสฺสโป อายสฺมโต อานนฺทสฺส สนฺติเก ทูตํ ปาเหสิ “อาคจฺฉตุ อานนฺโท อิมํ อนุสฺสาเวสฺสตู”ติ\csromansharedfootnote{9b99c250d07c1c2b}{อนุสฺสาเวสฺสตีติ (สฺยา)}. อายสฺมา อานนฺโท เอวมาห “นาหํ อุสฺสหามิ เถรสฺส นามํ คเหตุํ, ครุ เม เถโร”ติ. ภควโต เอตมตฺถํ อาโรเจสุํ. อนุชานามิ ภิกฺขเว โคตฺเตนปิ อนุสฺสาเวตุนฺติ.}

\csromanmatikaanchor{mtk.62}
\csromantocmarkref{section}{61. เทฺวอุปสมฺปทาเปกฺขาทิวตฺถุ}{mtk.62}
\bookmark[dest={mtk.62},level=1]{61. เทฺวอุปสมฺปทาเปกฺขาทิวตฺถุ}
\csromanheader{1}{61. \textbf{เทฺว อุปสมฺปทาเปกฺขาทิวตฺถุ}}
\proseitem{123}{เตน โข ปน สมเยน อายสฺมโต มหากสฺสปสฺส เทฺว อุปสมฺปทา\-เปกฺขา โหนฺติ. เต วิวทนฺติ “อหํ ปฐมํ อุปสมฺปชฺชิสฺสามิ, อหํ ปฐมํ อุปสมฺปชฺชิสฺสามี”ติ. ภควโต เอตมตฺถํ อาโรเจสุํ. อนุชานามิ ภิกฺขเว เทฺว เอกานุสฺสาวเน กาตุนฺติ.}

\prose{เตน โข ปน สมเยน สมฺพหุลานํ เถรานํ อุปสมฺปทา\-เปกฺขา โหนฺติ. เต วิวทนฺติ “อหํ ปฐมํ อุปสมฺปชฺชิสฺสามิ, อหํ ปฐมํ อุปสมฺปชฺชิสฺสามี”ติ. เถรา เอวมาหํสุ “หนฺท มยํ อาวุโส สพฺเพว เอกานุสฺสาวเน กโรมา”ติ. ภควโต เอตมตฺถํ อาโรเจสุํ. อนุชานามิ ภิกฺขเว เทฺว ตโย เอกานุสฺสาวเน กาตุํ, ตญฺจ โข เอเกน อุปชฺฌาเยน, น เตฺวว นานุปชฺฌาเยนาติ.}

\csromanmatikaanchor{mtk.63}
\csromantocmarkref{section}{62. คพฺภวีสูปสมฺปทานุชานนา}{mtk.63}
\bookmark[dest={mtk.63},level=1]{62. คพฺภวีสูปสมฺปทานุชานนา}
\csromanheader{1}{62. \textbf{คพฺภ}\-\textbf{วีสู}\-\textbf{ปสมฺปทา}\-\textbf{นุชานนา}}
\proseitem{124}{เตน โข ปน สมเยน อายสฺมา กุมารกสฺสโป คพฺภวีโส อุปสมฺปนฺโน อโหสิ. อถ โข อายสฺมโต กุมาร\-กสฺสปสฺส เอตทโหสิ “ภควตา ปญฺญตฺตํ ‘น อูนวีสติวสฺโส ปุคฺคโล อุปสมฺ\-ปาเท\-ตพฺโพ’ติ, อหญฺจมฺหิ คพฺภวีโส อุปสมฺปนฺโน, อุปสมฺปนฺโน นุ โขมฺหิ นนุ โข อุปสมฺปนฺโน”ติ. ภควโต เอตมตฺถํ อาโรเจสุํ. ยํ ภิกฺขเว มาตุกุจฺฉิสฺมิํ ปฐมํ จิตฺตํ อุปฺปนฺนํ, ปฐมํ วิญฺญาณํ ปาตุภูตํ, ตทุปาทาย สาวสฺส ชาติ. อนุชานามิ ภิกฺขเว คพฺภวีสํ อุปสมฺปาเทตุนฺติ.}

\csromanmatikaanchor{mtk.64}
\csromantocmarkref{section}{63. อุปสมฺปทาวิธิ}{mtk.64}
\bookmark[dest={mtk.64},level=1]{63. อุปสมฺปทาวิธิ}
\csromanheader{1}{63. \textbf{อุปสมฺปทาวิธิ}}
\proseitem{125}{\csromanfolio{130}เตน โข ปน สมเยน อุปสมฺปนฺนา ทิสฺสนฺติ กุฏฺฐิกาปิ คณฺฑิกาปิ กิลาสิกาปิ โสสิกาปิ อปมาริกาปิ. ภควโต เอตมตฺถํ อาโรเจสุํ. อนุชานามิ ภิกฺขเว อุปสมฺปาเทนฺเตน เตรส\csromansharedfootnote{f4c06bebf72a231c}{ตสฺส (ก)} อนฺตรายิเก ธมฺเม ปุจฺฉิตุํ. เอวญฺจ ปน ภิกฺขเว ปุจฺฉิตพฺโพ “สนฺติ เต เอวรูปา อาพาธา กุฏฺฐํ คณฺโฑ กิลาโส โสโส อปมาโร, มนุสฺโสสิ, ปุริโสสิ, ภุชิสฺโสสิ, อณโณสิ, นสิ ราชภโฏ, อนุญฺญาโตสิ มาตาปิตูหิ, ปริปุณฺณวีสติวสฺโสสิ, ปริปุณฺณํ เต ปตฺตจีวรํ, กิํนาโมสิ, โกนาโม เต อุปชฺฌาโย”ติ.}

\prose{เตน โข ปน สมเยน ภิกฺขู อนนุสิฏฺเฐ อุปสมฺปทา\-เปกฺเข อนฺตรายิเก ธมฺเม ปุจฺฉนฺติ. อุปสมฺปทา\-เปกฺขา วิตฺถายนฺติ, มงฺกู โหนฺติ, น สกฺโกนฺติ วิสฺสชฺเชตุํ. ภควโต เอตมตฺถํ อาโรเจสุํ. อนุชานามิ ภิกฺขเว ปฐมํ อนุสาสิตฺวา ปจฺฉา อนฺตรายิเก ธมฺเม ปุจฺฉิตุนฺติ.}

\proseitem{125}{ตตฺเถว สํฆมชฺเฌ อนุสาสนฺติ, อุปสมฺปทา\-เปกฺขา ตเถว วิตฺถายนฺติ, มงฺกู โหนฺติ, น สกฺโกนฺติ วิสฺสชฺเชตุํ. ภควโต เอตมตฺถํ อาโรเจสุํ. อนุชานามิ ภิกฺขเว เอกมนฺตํ อนุสาสิตฺวา สํฆมชฺเฌ อนฺตรายิเก ธมฺเม ปุจฺฉิตุํ. เอวญฺจ ปน ภิกฺขเว อนุสาสิตพฺโพ\csromandash{}}

\proseitem{126}{ปฐมํ อุปชฺฌํ คาหาเปตพฺโพ, อุปชฺฌํ คาหาเปตฺวา ปตฺตจีวรํ อาจิกฺขิตพฺพํ, อยํ เต ปตฺโต, อยํ สํฆาฏิ, อยํ อุตฺตราสงฺโค, อยํ อนฺตรวาสโก.}

\setlength{\csromangathapagewidth}{0pt}
\setlength{\csromangathawakwidth}{0pt}
\csromangathameasuregroup{}{คจฺฉ อมุมฺหิ โอกาเส ติฏฺฐาหีติ.}
\csromangathameasurewak{คจฺฉ อมุมฺหิ โอกาเส ติฏฺฐาหีติ.}
\setlength{\csromangathaleftcol}{0pt}
\csromangathagroup{\csromangathastanza{\csromangathawak{\csromangathaitemline{126}{คจฺฉ อมุมฺหิ โอกาเส ติฏฺฐาหีติ.}}}}

\prose{พาลา อพฺยตฺตา อนุสาสนฺติ, ทุรนุสิฏฺฐา อุปสมฺปทา\-เปกฺขา วิตฺถายนฺติ, มงฺกู โหนฺติ, น สกฺโกนฺติ วิสฺสชฺเชตุํ.}

\prose{ภควโต เอตมตฺถํ อาโรเจสุํ. น ภิกฺขเว พาเลน อพฺยตฺเตน อนุสาสิตพฺโพ, โย อนุสาเสยฺย, อาปตฺติ ทุกฺกฏสฺส. อนุชานามิ ภิกฺขเว พฺยตฺเตน ภิกฺขุนา ปฏิพเลน อนุสาสิตุนฺติ.}

\prose{อสมฺมตา อนุสาสนฺติ. ภควโต เอตมตฺถํ อาโรเจสุํ. น ภิกฺขเว อสมฺมเตน อนุสาสิตพฺโพ, โย อนุสาเสยฺย, อาปตฺติ ทุกฺกฏสฺส. \csromanfolio{131}อนุชานามิ ภิกฺขเว สมฺมเตน อนุสาสิตุํ. เอวญฺจ ปน ภิกฺขเว สมฺมนฺนิตพฺโพ\csromansharedfootnote{7f044bd59346eed7}{สมฺมนิตพฺโพ (ก)}, อตฺตนา วา\csromansharedfootnote{3f4250394bbf1e60}{อตฺตนาว (สฺยา)} อตฺตานํ สมฺมนฺนิตพฺพํ, ปเรน วา ปโร สมฺมนฺนิตพฺโพ.}

\proseitem{126}{กถญฺจ อตฺตนาว อตฺตานํ สมฺมนฺนิตพฺพํ, พฺยตฺเตน ภิกฺขุนา ปฏิพเลน สํโฆ ญาเปตพฺโพ “สุณาตุ เม ภนฺเต สํโฆ, อิตฺถนฺนาโม อิตฺถนฺนามสฺส อายสฺมโต อุปสมฺปทา\-เปกฺโข, ยทิ สํฆสฺส ปตฺตกลฺลํ, อหํ อิตฺถนฺนามํ อนุสาเสยฺยนฺ”ติ, เอวํ อตฺตนาว อตฺตานํ สมฺมนฺนิตพฺพํ.}

\prose{กถญฺจ ปน ปเรน ปโร สมฺมนฺนิตพฺโพ, พฺยตฺเตน ภิกฺขุนา ปฏิพเลน สํโฆ ญาเปตพฺโพ “สุณาตุ เม ภนฺเต สํโฆ, อิตฺถนฺนาโม อิตฺถนฺนามสฺส อายสฺมโต อุปสมฺปทา\-เปกฺโข, ยทิ สํฆสฺส ปตฺตกลฺลํ, อิตฺถนฺนาโม อิตฺถนฺนามํ อนุสาเสยฺยา”ติ, เอวํ ปเรน ปโร สมฺมนฺนิตพฺโพ.}

\prose{เตน สมฺมเตน ภิกฺขุนา อุปสมฺปทา\-เปกฺโข อุปสงฺกมิตฺวา เอวมสฺส วจนีโย “สุณสิ อิตฺถนฺนาม อยํ เต สจฺจกาโล ภูตกาโล, ยํ ชาตํ, ตํ สํฆมชฺเฌ ปุจฺฉนฺเต สนฺตํ ‘อตฺถี’ติ วตฺตพฺพํ, อสนฺตํ ‘นตฺถี’ติ วตฺตพฺพํ, มา โข วิตฺถายิ, มา โข มงฺกุ อโหสิ, เอวํ ตํ ปุจฺฉิสฺสนฺติ\csromandash{}สนฺติ เต เอวรูปา อาพาธา กุฏฺฐํ คณฺโฑ กิลาโส โสโส อปมาโร, มนุสฺโสสิ, ปุริโสสิ, ภุชิสฺโสสิ, อณโณสิ, นสิ ราชภโฏ, อนุญฺญาโตสิ มาตาปิตูหิ, ปริปุณฺณวีสติวสฺโสสิ, ปริปุณฺณํ เต ปตฺตจีวรํ, กิํนาโมสิ, โกนาโม เต อุปชฺฌาโย”ติ.}

\prose{เอกโต อาคจฺฉนฺติ. น ภิกฺขเว เอกโต อาคนฺตพฺพํ, อนุสาสเกน ปฐมตรํ อาคนฺตฺวา สํโฆ ญาเปตพฺโพ “สุณาตุ เม ภนฺเต สํโฆ, อิตฺถนฺนาโม อิตฺถนฺนามสฺส อายสฺมโต อุปสมฺปทา\-เปกฺโข, อนุสิฏฺโฐ โส มยา, ยทิ สํฆสฺส ปตฺตกลฺลํ, อิตฺถนฺนาโม อาคจฺเฉยฺยา”ติ. “อาคจฺฉาหี”ติ วตฺตพฺโพ.}

\proseitem{126}{เอกํสํ อุตฺตราสงฺค การาเปตฺวา ภิกฺขูนํ ปาเท วนฺทาเปตฺวา อุกฺกุฏิกํ นิสีทาเปตฺวา อญฺชลิํ ปคฺคณฺหาเปตฺวา อุปสมฺปทํ ยาจาเปตพฺโพ “สํฆํ ภนฺเต อุปสมฺปทํ ยาจามิ, อุลฺลุมฺปตุ มํ ภนฺเต สํโฆ อนุกมฺปํ อุปาทาย. ทุติยมฺปิ ภนฺเต สํฆํ อุปสมฺปทํ ยาจามิ, อุลฺลุมฺปตุ มํ ภนฺเต สํโฆ \csromanfolio{132}อนุกมฺปํ อุปาทาย. ตติยมฺปิ ภนฺเต สํฆํ อุปสมฺปทํ ยาจามิ, อุลฺลุมฺปตุ มํ ภนฺเต สํโฆ อนุกมฺปํ อุปาทายา”ติ. พฺยตฺเตน ภิกฺขุนา ปฏิพเลน สํโฆ ญาเปตพฺโพ\csromandash{}}

\prose{“สุณาตุ เม ภนฺเต สํโฆ, อยํ อิตฺถนฺนาโม อิตฺถนฺนามสฺส อายสฺมโต อุปสมฺปทา\-เปกฺโข, ยทิ สํฆสฺส ปตฺตกลฺลํ, อหํ อิตฺถนฺนามํ อนฺตรายิเก ธมฺเม ปุจฺเฉยฺยนฺ”ติ. สุณสิ อิตฺถนฺนาม อยํ เต สจฺจกาโล ภูตกาโล, ยํ ชาตํ, ตํ ปุจฺฉามิ, สนฺตํ “อตฺถี”ติ วตฺตพฺพํ, อสนฺตํ “นตฺถี”ติ วตฺตพฺพํ, สนฺติ เต เอวรูปา อาพาธา กุฏฺฐํ คณฺโฑ กิลาโส โสโส อปมาโร, มนุสฺโสสิ, ปุริโสสิ, ภุชิสฺโสสิ, อณโณสิ, นสิ ราชภโฏ, อนุญฺญาโตสิ มาตาปิตูหิ, ปริปุณฺณวีสติวสฺโสสิ, ปริปุณฺณํ เต ปตฺตจีวรํ, กิํนาโมสิ, โกนาโม เต อุปชฺฌาโยติ. พฺยตฺเตน ภิกฺขุนา ปฏิพเลน สํโฆ ญาเปตพฺโพ\csromandash{}}

\proseitem{127}{“สุณาตุ เม ภนฺเต สํโฆ, อยํ อิตฺถนฺนาโม อิตฺถนฺนามสฺส อายสฺมโต อุปสมฺปทา\-เปกฺโข, ปริสุทฺโธ อนฺตรายิเกหิ ธมฺเมหิ, ปริปุณฺณ’สฺส ปตฺตจีวรํ, อิตฺถนฺนาโม สํฆํ อุปสมฺปทํ ยาจติ อิตฺถนฺนาเมน อุปชฺฌาเยน, ยทิ สํฆสฺส ปตฺตกลฺลํ, สํโฆ อิตฺถนฺนามํ อุปสมฺปาเทยฺย อิตฺถนฺนาเมน อุปชฺฌาเยน, เอสา ญตฺติ.}

\prose{สุณาตุ เม ภนฺเต สํโฆ, อยํ อิตฺถนฺนาโม อิตฺถนฺนามสฺส อายสฺมโต อุปสมฺปทา\-เปกฺโข, ปริสุทฺโธ อนฺตรายิเกหิ ธมฺเมหิ, ปริปุณฺณ’สฺส ปตฺตจีวรํ, อิตฺถนฺนาโม สํฆํ อุปสมฺปทํ ยาจติ อิตฺถนฺนาเมน อุปชฺฌาเยน, สํโฆ อิตฺถนฺนามํ อุปสมฺปาเทติ อิตฺถนฺนาเมน อุปชฺฌาเยน, ยสฺสายสฺมโต ขมติ อิตฺถนฺนามสฺส อุปสมฺปทา อิตฺถนฺนาเมน อุปชฺฌาเยน, โส ตุณฺหสฺส, ยสฺส นกฺขมติ, โส ภาเสยฺย.}

\prose{ทุติยมฺปิ เอตมตฺถํ วทามิ. สุณาตุ เม ภนฺเต สํโฆ, อยํ อิตฺถนฺนาโม อิตฺถนฺนามสฺส อายสฺมโต อุปสมฺปทา\-เปกฺโข, ปริสุทฺโธ อนฺตรายิเกหิ ธมฺเมหิ, ปริปุณฺณ’สฺส ปตฺตจีวรํ, อิตฺถนฺนาโม สํฆํ อุปสมฺปทํ ยาจติ อิตฺถนฺนาเมน อุปชฺฌาเยน, สํโฆ อิตฺถนฺนามํ อุปสมฺปาเทติ อิตฺถานฺนาเมน อุปชฺฌาเยน, ยสฺสายสฺมโต ขมติ อิตฺถนฺนามสฺส อุปสมฺปทา อิตฺถนฺนาเมน อุปชฺฌาเยน, โส ตุณฺหสฺส, ยสฺส นกฺขมติ, โส ภาเสยฺย.}

\prose{\csromanfolio{133}ตติยมฺปิ เอตมตฺถํ วทามิ. สุณาตุ เม ภนฺเต สํโฆ, อยํ อิตฺถนฺนาโม อิตฺถนฺนามสฺส อายสฺมโต อุปสมฺปทา\-เปกฺโข, ปริสุทฺโธ อนฺตรายิเกหิ ธมฺเมหิ, ปริปุณฺณ’สฺส ปตฺตจีวรํ, อิตฺถนฺนาโม สํฆํ อุปสมฺปทํ ยาจติ อิตฺถนฺนาเมน อุปชฺฌาเยน, สํโฆ อิตฺถนฺนามํ อุปสมฺปาเทติ อิตฺถนฺนาเมน อุปชฺฌาเยน, ยสฺสายสฺมโต ขมติ อิตฺถนฺนามสฺส อุปสมฺปทา อิตฺถนฺนาเมน อุปชฺฌาเยน, โส ตุณฺห’สฺส, ยสฺส นกฺขมติ, โส ภาเสยฺย.}

\prosewithcloserruled{อุปสมฺปนฺโน สํเฆน อิตฺถนฺนาโม อิตฺถนฺนาเมน อุปชฺฌาเยน, ขมติ สํฆสฺส, ตสฺมา ตุณฺหี, เอวเมตํ ธารยามี”ติ.}{อุปสมฺปทา\-กมฺมํ นิฏฺฐิตํ.}
\csromanmatikaanchor{mtk.65}
\csromantocmarkref{section}{64. จตฺตาโร นิสฺสยา}{mtk.65}
\bookmark[dest={mtk.65},level=1]{64. จตฺตาโร นิสฺสยา}
\csromanheader{1}{64. \textbf{จตฺตาโร นิสฺสยา}}
\proseitem{128}{ตาวเทว ฉายา เมตพฺพา, อุตุปฺปมาณํ อาจิกฺขิตพฺพํ, ทิวสภาโค อาจิกฺขิตพฺโพ, สงฺคีติ อาจิกฺขิตพฺพา, จตฺตาโร นิสฺสยา อาจิกฺขิตพฺพา\csromansharedfootnote{8bc4b023fcf6df51}{อาจิกฺขิตพฺพา, จตฺตาริ อกรณียานิ อาจิกฺขิตพฺพานิ. (ก)}\csromandash{}}

\prose{ปิณฺฑิยาโลป\-โภชนํ นิสฺสาย ปพฺพชฺชา, ตตฺถ เต ยาวชีวํ อุสฺสาโห กรณีโย, อติเรกลาโภ สํฆภตฺตํ อุทฺเทสภตฺตํ นิมนฺตนํ สลากภตฺตํ ปกฺขิกํ อุโปสถิกํ ปาฏิปทิกํ.}

\prose{ปํสุกูล\-จีวรํ นิสฺสาย ปพฺพชฺชา, ตตฺถ เต ยาวชีวํ อุสฺสาโห กรณีโย, อติเรกลาโภ โขมํ กปฺปาสิกํ โกเสยฺยํ กมฺพลํ สาณํ ภงฺคํ.}

\prose{รุกฺขมูล\-เสนาสนํ นิสฺสาย ปพฺพชฺชา, ตตฺถ เต ยาวชีวํ อุสฺสาโห กรณีโย, อติเรกลาโภ วิหาโร อฑฺฒโยโค ปาสาโท หมฺมิยํ คุหา.}

\prosewithcloser{ปูติมุตฺต\-เภสชฺชํ นิสฺสาย ปพฺพชฺชา, ตตฺถ เต ยาวชีวํ อุสฺสาโห กรณีโย, อติเรกลาโภ สปฺปิ นวนีตํ เตลํ มธุ ผาณิตนฺติ.}{จตฺตาโร นิสฺสยา นิฏฺฐิตา.}
\csromanmatikaanchor{mtk.66}
\csromantocmarkref{section}{65. จตฺตาริ อกรณียานิ}{mtk.66}
\bookmark[dest={mtk.66},level=1]{65. จตฺตาริ อกรณียานิ}
\csromanheader{1}{65. \textbf{จตฺตาริ อกรณียานิ}}
\proseitem{129}{\csromanfolio{134}เตน โข ปน สมเยน ภิกฺขู อญฺญตรํ ภิกฺขุํ อุปสมฺปาเทตฺวา เอกกํ โอหาย ปกฺกมิํสุ. โส ปจฺฉา เอกโกว อาคจฺฉนฺโต อนฺตรามคฺเค ปุราณ\-ทุติยิกาย สมาคญฺฉิ. สา เอวมาห “กิํทานิ ปพฺพชิโตสี”ติ. อาม ปพฺพชิโตมฺหีติ. ทุลฺลโภ โข ปพฺพชิตานํ เมถุโน ธมฺโม, เอหิ เมถุนํ ธมฺมํ ปฏิเสวาติ. โส ตสฺสา เมถุนํ ธมฺมํ ปฏิเสวิตฺวา จิเรน อคมาสิ. ภิกฺขู เอวมาหํสุ “กิสฺส ตฺวํ อาวุโส เอวํ จิรํ อกาสี”ติ. อถ โข โส ภิกฺขุ ภิกฺขูนํ เอตมตฺถํ อาโรเจสิ. ภิกฺขู ภควโต เอตมตฺถํ อาโรเจสุํ. อนุชานามิ ภิกฺขเว อุปสมฺปาเทตฺวา ทุติยํ ทาตุํ, จตฺตาริ จ อกรณียานิ อาจิกฺขิตุํ\csromandash{}}

\prose{อุปสมฺปนฺเนน ภิกฺขุนา เมถุโน ธมฺโม น ปฏิเสวิตพฺโพ อนฺตมโส ติรจฺฉาน\-คตายปิ. โย ภิกฺขุ เมถุนํ ธมฺมํ ปฏิเสวติ, อสฺสมโณ โหติ อสกฺยปุตฺติโย. เสยฺยถาปิ นาม ปุริโส สีสจฺฉินฺโน อภพฺโพ เตน สรีร\-พนฺธเนน ชีวิตุํ, เอวเมว ภิกฺขุ เมถุนํ ธมฺมํ ปฏิเสวิตฺวา อสฺสมโณ โหติ อสกฺยปุตฺติโย, ตํ เต ยาวชีวํ อกรณียํ.}

\prose{อุปสมฺปนฺเนน ภิกฺขุนา อทินฺนํ เถยฺยสงฺขาตํ น อาทาตพฺพํ อนฺตมโส ติณสลากํ อุปาทาย. โย ภิกฺขุ ปาทํ วา ปาทารหํ วา อติเรกปาทํ วา อทินฺนํ เถยฺยสงฺขาตํ อาทิยติ, อสฺสมโณ โหติ อสกฺยปุตฺติโย. เสยฺยถาปิ นาม ปณฺฑุปลาโส พนฺธนา ปวุตฺโต อภพฺโพ หริตตฺถาย, เอวเมว ภิกฺขุ ปาทํ วา ปาทารหํ วา อติเรกปาทํ วา อทินฺนํ เถยฺยสงฺขาตํ อาทิยิตฺวา อสฺสมโณ โหติ อสกฺยปุตฺติโย ตํ เต ยาวชีวํ อกรณียํ.}

\prose{อุปสมฺปนฺเนน ภิกฺขุนา สญฺจิจฺจ ปาโณ ชีวิตา น โวโรเปตพฺโพ อนฺตมโส กุนฺถ\-กิปิลฺลิกํ อุปาทาย. โย ภิกฺขุ สญฺจิจฺจ มนุสฺส\-วิคฺคหํ ชีวิตา โวโรเปติ อนฺตมโส คพฺภปาตนํ อุปาทาย, อสฺสมโณ โหติ อสกฺยปุตฺติโย. เสยฺยถาปิ นาม ปุถุสิลา เทฺวธา ภินฺนา อปฺปฏิสนฺธิกา โหติ, เอวเมว ภิกฺขุ สญฺจิจฺจ มนุสฺส\-วิคฺคหํ ชีวิตา โวโรเปตฺวา อสฺสมโณ โหติ อสกฺยปุตฺติโย, ตํ เต ยาวชีวํ อกรณียํ.}

\prosewithcloserruled{\csromanfolio{135}อุปสมฺปนฺเนน ภิกฺขุนา อุตฺตริ\-มนุสฺสธมฺโม น อุลฺลปิตพฺโพ อนฺตมโส “สุญฺญาคาเร อภิรมามี”ติ. โย ภิกฺขุ ปาปิจฺโฉ อิจฺฉาปกโต อสนฺตํ อภูตํ อุตฺตริ\-มนุสฺสธมฺมํ อุลฺลปติ ฌานํ วา วิโมกฺขํ วา สมาธิํ วา สมาปตฺติํ วา มคฺคํ วา ผลํ วา, อสฺสมโณ โหติ อสกฺยปุตฺติโย. เสยฺยถาปิ นาม ตาโล มตฺถก\-จฺฉินฺโน อภพฺโพ ปุน วิรุฬฺหิยา, เอวเมว ภิกฺขุ ปาปิจฺโฉ อิจฺฉาปกโต อสนฺตํ อภูตํ อุตฺตริ\-มนุสฺสธมฺมํ อุลฺลปิตฺวา อสฺสมโณ โหติ อสกฺยปุตฺติโย, ตํ เต ยาวชีวํ อกรณียนฺติ.}{จตฺตาริ อกรณียานิ นิฏฺฐิตานิ.}
\csromanmatikaanchor{mtk.67}
\csromantocmarkref{section}{66. อาปตฺติยา อทสฺสเน อุกฺขิตฺตกวตฺถูนิ}{mtk.67}
\bookmark[dest={mtk.67},level=1]{66. อาปตฺติยา อทสฺสเน อุกฺขิตฺตกวตฺถูนิ}
\csromanheader{1}{66. \textbf{อาปตฺติยา อทสฺสเน อุกฺขิตฺตก}\-\textbf{วตฺถูนิ}}
\proseitem{130}{เตน โข ปน สมเยน อญฺญตโร ภิกฺขุ อาปตฺติยา อทสฺสเน อุกฺขิตฺตโก วิพฺภมิ, โส ปุน ปจฺจาคนฺตฺวา ภิกฺขู อุปสมฺปทํ ยาจิ.}

\setlength{\csromangathapagewidth}{0pt}
\setlength{\csromangathawakwidth}{0pt}
\csromangathameasuregroup{}{ภควโต เอตมตฺถํ อาโรเจสุํ.}
\csromangathameasurewak{ภควโต เอตมตฺถํ อาโรเจสุํ.}
\setlength{\csromangathaleftcol}{0pt}
\csromangathagroup{\csromangathastanza{\csromangathawak{\csromangathaitemline{130}{ภควโต เอตมตฺถํ อาโรเจสุํ.}}}}

\prose{อิธ ปน ภิกฺขเว ภิกฺขุ อาปตฺติยา อทสฺสเน อุกฺขิตฺตโก วิพฺภมติ, โส ปุน ปจฺจาคนฺตฺวา ภิกฺขู อุปสมฺปทํ ยาจติ, โส เอวมสฺส วจนีโย “ปสฺสิสฺสสิ ตํ อาปตฺตินฺ”ติ.}

\setlength{\csromangathapagewidth}{0pt}
\setlength{\csromangathawakwidth}{0pt}
\csromangathameasuregroup{สจาหํ ปสฺสิสฺสามีติ, \\ สจาหํ น ปสฺสิสฺสามีติ,}{\csromangathabat{สจาหํ ปสฺสิสฺสามีติ,}{ปพฺพาเชตพฺโพ.} \\ \csromangathabat{สจาหํ น ปสฺสิสฺสามีติ,}{น ปพฺพาเชตพฺโพ.}}
\csromangathasetleft{สจาหํ ปสฺสิสฺสามีติ, \\ สจาหํ น ปสฺสิสฺสามีติ,}
\csromangathagroup{\csromangathastanza{\csromangathabat{สจาหํ ปสฺสิสฺสามีติ,}{ปพฺพาเชตพฺโพ.} \\ \csromangathabat{สจาหํ น ปสฺสิสฺสามีติ,}{น ปพฺพาเชตพฺโพ.}}}

\prose{ปพฺพาเชตฺวา วตฺตพฺโพ “ปสฺสิสฺสสิ ตํ อาปตฺตินฺ”ติ. สจาหํ ปสฺสิสฺสามีติ, อุปสมฺ\-ปาเท\-ตพฺโพ. สจาหํ น ปสฺสิสฺสามีติ, น อุปสมฺ\-ปาเท\-ตพฺโพ. อุปสมฺปาเทตฺวา วตฺตพฺโพ “ปสฺสิสฺสสิ ตํ อาปตฺตินฺ”ติ. สจาหํ ปสฺสิสฺสามีติ, โอสาเรตพฺโพ. สจาหํ น ปสฺสิสฺสามีติ, น โอสาเรตพฺโพ. โอสาเรตฺวา วตฺตพฺโพ “ปสฺสสิ\csromansharedfootnote{338dd9d64fd7ee1e}{[มิสฺสินฺคฺ โนเต 0]} ตํ อาปตฺตินฺ”ติ. สเจ ปสฺสติ, อิจฺเจตํ กุสลํ. โน เจ ปสฺสติ, ลพฺภมานาย สามคฺคิยา ปุน อุกฺขิปิตพฺโพ. อลพฺภมานาย สามคฺคิยา อนาปตฺติ สมฺโภเค สํวาเส.}

\csromanmatikaanchor{mtk.68}
\csromantocmarkref{section}{67. อุทฺทานคาถา}{mtk.68}
\bookmark[dest={mtk.68},level=1]{67. อุทฺทานคาถา}
\proseitem{130}{อิธ ปน ภิกฺขเว ภิกฺขุ อาปตฺติยา อปฺปฏิกมฺเม อุกฺขิตฺตโก วิพฺภมติ, โส ปุน ปจฺจาคนฺตฺวา ภิกฺขู อุปสมฺปทํ ยาจติ, โส เอวมสฺส วจนีโย “ปฏิกริสฺสสิ ตํ อาปตฺตินฺ”ติ. สจาหํ ปฏิกริสฺสามีติ, ปพฺพาเชตพฺโพ. สจาหํ น ปฏิกริสฺสามีติ, น ปพฺพาเชตพฺโพ. ปพฺพาเชตฺวา วตฺตพฺโพ \csromanfolio{136}“ปฏิกริสฺสสิ ตํ อาปตฺตินฺ”ติ. สจาหํ ปฏิกริสฺสามีติ, อุปสมฺ\-ปาเท\-ตพฺโพ. สจาหํ น ปฏิกริสฺสามีติ, น อุปสมฺ\-ปาเท\-ตพฺโพ. อุปสมฺปาเทตฺวา วตฺตพฺโพ “ปฏิกริสฺสสิ ตํ อาปตฺตินฺ”ติ. สจาหํ ปฏิกริสฺสามีติ, โอสาเรตพฺโพ. สจาหํ น ปฏิกริสฺสามีติ, น โอสาเรตพฺโพ. โอสาเรตฺวา วตฺตพฺโพ “ปฏิกโรหิ ตํ อาปตฺตินฺ”ติ. สเจ ปฏิกโรติ, อิจฺเจตํ กุสลํ. โน เจ ปฏิกโรติ, ลพฺภมานาย สามคฺคิยา ปุน อุกฺขิปิตพฺโพ. อลพฺภมานาย สามคฺคิยา อนาปตฺติ สมฺโภเค สํวาเส.}

\prose{อิธ ปน ภิกฺขเว ภิกฺขุ ปาปิกาย ทิฏฺฐิยา อปฺปฏินิสฺสคฺเค อุกฺขิตฺตโก วิพฺภมติ, โส ปุน ปจฺจาคนฺตฺวา ภิกฺขู อุปสมฺปทํ ยาจติ, โส เอวมสฺส วจนีโย “ปฏินิสฺสชฺชิสฺสสิ ตํ ปาปิกํ ทิฏฺฐินฺ”ติ. สจาหํ ปฏินิสฺสชฺชิสฺสามีติ, ปพฺพาเชตพฺโพ. สจาหํ น ปฏินิสฺสชฺชิสฺสามีติ, น ปพฺพาเชตพฺโพ. ปพฺพาเชตฺวา วตฺตพฺโพ “ปฏินิสฺสชฺชิสฺสสิ ตํ ปาปิกํ ทิฏฺฐินฺ”ติ. สจาหํ ปฏินิสฺสชฺชิสฺสามีติ, อุปสมฺ\-ปาเท\-ตพฺโพ. สจาหํ น ปฏินิสฺสชฺชิสฺสามีติ, น อุปสมฺ\-ปาเท\-ตพฺโพ. อุปสมฺปาเทตฺวา วตฺตพฺโพ “ปฏินิสฺสชฺชิสฺสสิ ตํ ปาปิกํ ทิฏฺฐินฺ”ติ. สจาหํ ปฏินิสฺสชฺชิสฺสามีติ, โอสาเรตพฺโพ. สจาหํ น ปฏินิสฺสชฺชิสฺสามีติ, น โอสาเรตพฺโพ. โอสาเรตฺวา วตฺตพฺโพ “ปฏินิสฺสชฺเชหิ ตํ ปาปิกํ ทิฏฺฐินฺ”ติ. สเจ ปฏินิสฺสชฺชติ, อิจฺเจตํ กุสลํ. โน เจ ปฏินิสฺสชฺชติ, ลพฺภมานาย สามคฺคิยา ปุน อุกฺขิปิตพฺโพ. อลพฺภมานาย สามคฺคิยา อนาปตฺติ สมฺโภเค สํวาเสติ.}

\vspace{\tipitakaparskip}
\csromancenter{มหาขนฺธโก ปฐโม.}
\csromancenter{67. \textbf{ตสฺสุทฺทานํ}}
\setlength{\csromangathapagewidth}{0pt}
\setlength{\csromangathawakwidth}{0pt}
\csromangathameasuregroup{วินยมฺหิ มหตฺเถสุ, \\ นิคฺคหานญฺจ ปาปิจฺเฉ, \\ สาสนาธารเณ เจว, \\ อนญฺญวิสเย เขเม, \\ ขนฺธเก วินเย เจว, \\ ยถาตฺถการี กุสโล, \\ โย ควํ น วิชานาติ, \\ เอวํ สีลํ อชานนฺโต, \\ ปมุฏฺฐมฺหิ จ สุตฺตนฺเต, \\ วินเย อวินฏฺฐมฺหิ, \\ ตสฺมา สงฺคาหณาเหตุํ\textsuperscript{1}, \\ ปวกฺขามิ ยถาญายํ, \\ วตฺถุ นิทานํ อาปตฺติ, \\ ทุกฺกรํ ตํ อเสเสตุํ, \\ โพธิ ราชายตนญฺจ, \\ พฺรหฺมา อาฬาโร อุทโก, \\ โกณฺฑญฺโญ วปฺโป ภทฺทิโย, \\ ยโส จตฺตาโร ปญฺญาส, \\ วตฺถุ มาเรหิ ติํสา จ, \\ อคฺยาคารํ มหาราชา, \\ ปํสุกูลํ โปกฺขรณี, \\ ชมฺพุ อมฺโพ จ อามโล, \\ ผาลิยนฺตุ อุชฺชลนฺตุ, \\ นิมุชฺชนฺติ มุขี เมโฆ, \\ อุปติสฺโส โกลิโต จ, \\ ทุนฺนิวตฺถา ปณามนา, \\ อนาจารํ อาจรติ, \\ วสฺสํ พาเลหิ ปกฺกนฺโต, \\ น วตฺตนฺติ ปณาเมตุํ, \\ โย โส อญฺโญ จ นคฺโค จ, \\ มคเธสุ ปญฺจาพาธา, \\ มาคโธ จ อนุญฺญาสิ, \\ ลกฺขณา อิณา ทาโส จ, \\ สทฺธํ กุลํ กณฺฏโก จ, \\ วตฺถุมฺหิ ทารโก สิกฺขา, \\ สพฺพํ มุขํ อุปชฺฌาเย, \\ ปณฺฑโก เถยฺยปกฺกนฺโต, \\ อรหนฺตภิกฺขุนีเภทา, \\ อนุปชฺฌาย สํเฆน, \\ อจีวรํ ตทุภยํ, \\ หตฺถา ปาทา หตฺถปาทา, \\ องฺคุลิอฬกณฺฑรํ, \\ คลคณฺฑี ลกฺขณา เจว, \\ ปาปปริสทูสี จ, \\ ขญฺชํ ปกฺขหตญฺเจว, \\ ชรานฺธมูคพธิรํ, \\ อนฺธพธิรํ ยํ วุตฺตํ, \\ อนฺธมูคพธิรญฺจ, \\ วตฺถพฺพญฺจ ตถาทฺธานํ, \\ อาคจฺฉตุ วิวทนฺติ, \\ ทิสฺสนฺติ อุปสมฺปนฺนา, \\ อนนุสิฏฺฐา วิตฺเถนฺติ, \\ สํเฆปิ จ อโถ พาลา, \\ อุลฺลุมฺปตุปสมฺปทา,}{\csromangathabat{วินยมฺหิ มหตฺเถสุ,}{เปสลานํ สุขาวเห.} \\ \csromangathabat{นิคฺคหานญฺจ ปาปิจฺเฉ,}{ลชฺชีนํ ปคฺคเหสุ จ.} \\ \csromangathabat{สาสนาธารเณ เจว,}{สพฺพญฺญุชินโคจเร.} \\ \csromangathabat{อนญฺญวิสเย เขเม,}{สุปญฺญตฺเต อสํสเย.} \\ \csromangathabat{ขนฺธเก วินเย เจว,}{ปริวาเร จ มาติเก.} \\ \csromangathabat{ยถาตฺถการี กุสโล,}{ปฏิปชฺชติ โยนิโส.} \\ \csromangathabat{โย ควํ น วิชานาติ,}{น โส รกฺขติ โคคณํ.} \\ \csromangathabat{เอวํ สีลํ อชานนฺโต,}{กิํ โส รกฺเขยฺย สํวรํ.} \\ \csromangathabat{ปมุฏฺฐมฺหิ จ สุตฺตนฺเต,}{อภิธมฺเม จ ตาวเท.} \\ \csromangathabat{วินเย อวินฏฺฐมฺหิ,}{ปุน ติฏฺฐติ สาสนํ.} \\ \csromangathabat{ตสฺมา สงฺคาหณาเหตุํ\textsuperscript{1},}{อุทฺทานํ อนุปุพฺพโส.} \\ \csromangathabat{ปวกฺขามิ ยถาญายํ,}{สุณาถ มม ภาสโต.} \\ \csromangathabat{วตฺถุ นิทานํ อาปตฺติ,}{นยา เปยฺยาลเมว จ.} \\ \csromangathabat{ทุกฺกรํ ตํ อเสเสตุํ,}{นยโต ตํ วิชานถาติ.} \\ \csromangathabat{โพธิ ราชายตนญฺจ,}{อชปาโล สหมฺปติ.} \\ \csromangathabat{พฺรหฺมา อาฬาโร อุทโก,}{ภิกฺขุ จ อุปโก อิสิ.} \\ \csromangathabat{โกณฺฑญฺโญ วปฺโป ภทฺทิโย,}{มหานาโม จ อสฺสชิ.} \\ \csromangathabat{ยโส จตฺตาโร ปญฺญาส,}{สพฺเพ เปเสสิ โส ทิสา.} \\ \csromangathabat{วตฺถุ มาเรหิ ติํสา จ,}{อุรุเวลํ ตโย ชฏี.} \\ \csromangathabat{อคฺยาคารํ มหาราชา,}{สกฺโก พฺรหฺมา จ เกวลา.} \\ \csromangathabat{ปํสุกูลํ โปกฺขรณี,}{สิลา จ กกุโธ สิลา.} \\ \csromangathabat{ชมฺพุ อมฺโพ จ อามโล,}{ปาริปุปฺผญฺจ อาหริ.} \\ \csromangathabat{ผาลิยนฺตุ อุชฺชลนฺตุ,}{วิชฺฌายนฺตุ จ กสฺสป.} \\ \csromangathabat{นิมุชฺชนฺติ มุขี เมโฆ,}{คยา ลฏฺฐิ จ มาคโธ.} \\ \csromangathabat{อุปติสฺโส โกลิโต จ,}{อภิญฺญาตา จ ปพฺพชุํ.} \\ \csromangathabat{ทุนฺนิวตฺถา ปณามนา,}{กิโส ลูโข จ พฺราหฺมโณ.} \\ \csromangathabat{อนาจารํ อาจรติ,}{อุทรํ มาณโว คโณ.} \\ \csromangathabat{วสฺสํ พาเลหิ ปกฺกนฺโต,}{ทส วสฺสานิ นิสฺสโย.} \\ \csromangathabat{น วตฺตนฺติ ปณาเมตุํ,}{พาลา ปสฺสทฺธิ ปญฺจ ฉ.} \\ \csromangathabat{โย โส อญฺโญ จ นคฺโค จ,}{อจฺฉินฺนชฏิลสากิโย.} \\ \csromangathabat{มคเธสุ ปญฺจาพาธา,}{เอโก ราชา\textsuperscript{1} จ องฺคุลิ.} \\ \csromangathabat{มาคโธ จ อนุญฺญาสิ,}{การา ลิขิ กสาหโต.} \\ \csromangathabat{ลกฺขณา อิณา ทาโส จ,}{ภณฺฑุโก อุปาลิ อหิ.} \\ \csromangathabat{สทฺธํ กุลํ กณฺฏโก จ,}{อาหุนฺทริกเมว จ.} \\ \csromangathabat{วตฺถุมฺหิ ทารโก สิกฺขา,}{วิหรนฺติ จ กิํ นุ โข.} \\ \csromangathabat{สพฺพํ มุขํ อุปชฺฌาเย,}{อปลาฬน กณฺฏโก.} \\ \csromangathabat{ปณฺฑโก เถยฺยปกฺกนฺโต,}{อหิ จ มาตรี ปิตา.} \\ \csromangathabat{อรหนฺตภิกฺขุนีเภทา,}{รุหิเรน จ พฺยญฺชนํ.} \\ \csromangathabat{อนุปชฺฌาย สํเฆน,}{คณปณฺฑกปตฺตโก.} \\ \csromangathabat{อจีวรํ ตทุภยํ,}{ยาจิเตนปิ เย ตโย.} \\ \csromangathabat{หตฺถา ปาทา หตฺถปาทา,}{กณฺณา นาสา ตทูภยํ.} \\ \csromangathabat{องฺคุลิอฬกณฺฑรํ,}{ผณํ ขุชฺชญฺจ วามนํ.} \\ \csromangathabat{คลคณฺฑี ลกฺขณา เจว,}{กสา ลิขิตสีปที.} \\ \csromangathabat{ปาปปริสทูสี จ,}{กาณํ กุณิ ตเถว จ.} \\ \csromangathabat{ขญฺชํ ปกฺขหตญฺเจว,}{สจฺฉินฺนอิริยาปถํ.} \\ \csromangathabat{ชรานฺธมูคพธิรํ,}{อนฺธมูคญฺจ ยํ ตหิํ.} \\ \csromangathabat{อนฺธพธิรํ ยํ วุตฺตํ,}{มูคพธิรเมว จ.} \\ \csromangathabat{อนฺธมูคพธิรญฺจ,}{อลชฺชีนญฺจ นิสฺสยํ.} \\ \csromangathabat{วตฺถพฺพญฺจ ตถาทฺธานํ,}{ยาจมาเนน ลกฺขณา\textsuperscript{1}.} \\ \csromangathabat{อาคจฺฉตุ วิวทนฺติ,}{เอกุปชฺฌาเยน กสฺสโป.} \\ \csromangathabat{ทิสฺสนฺติ อุปสมฺปนฺนา,}{อาพาเธหิ จ ปีฬิตา.} \\ \csromangathabat{อนนุสิฏฺฐา วิตฺเถนฺติ,}{ตตฺเถว อนุสาสนา.} \\ \csromangathabat{สํเฆปิ จ อโถ พาลา,}{อสมฺมตา จ เอกโต.} \\ \csromangathabat{อุลฺลุมฺปตุปสมฺปทา,}{นิสฺสโย เอกโก ตโยติ.} \\ อิมมฺหิ ขนฺธเก วตฺถูนิ เอกสตญฺจ ทฺวาสตฺตติ.}
\csromangathameasurewak{อิมมฺหิ ขนฺธเก วตฺถูนิ เอกสตญฺจ ทฺวาสตฺตติ.}
\csromangathasetleft{วินยมฺหิ มหตฺเถสุ, \\ นิคฺคหานญฺจ ปาปิจฺเฉ, \\ สาสนาธารเณ เจว, \\ อนญฺญวิสเย เขเม, \\ ขนฺธเก วินเย เจว, \\ ยถาตฺถการี กุสโล, \\ โย ควํ น วิชานาติ, \\ เอวํ สีลํ อชานนฺโต, \\ ปมุฏฺฐมฺหิ จ สุตฺตนฺเต, \\ วินเย อวินฏฺฐมฺหิ, \\ ตสฺมา สงฺคาหณาเหตุํ\textsuperscript{1}, \\ ปวกฺขามิ ยถาญายํ, \\ วตฺถุ นิทานํ อาปตฺติ, \\ ทุกฺกรํ ตํ อเสเสตุํ, \\ โพธิ ราชายตนญฺจ, \\ พฺรหฺมา อาฬาโร อุทโก, \\ โกณฺฑญฺโญ วปฺโป ภทฺทิโย, \\ ยโส จตฺตาโร ปญฺญาส, \\ วตฺถุ มาเรหิ ติํสา จ, \\ อคฺยาคารํ มหาราชา, \\ ปํสุกูลํ โปกฺขรณี, \\ ชมฺพุ อมฺโพ จ อามโล, \\ ผาลิยนฺตุ อุชฺชลนฺตุ, \\ นิมุชฺชนฺติ มุขี เมโฆ, \\ อุปติสฺโส โกลิโต จ, \\ ทุนฺนิวตฺถา ปณามนา, \\ อนาจารํ อาจรติ, \\ วสฺสํ พาเลหิ ปกฺกนฺโต, \\ น วตฺตนฺติ ปณาเมตุํ, \\ โย โส อญฺโญ จ นคฺโค จ, \\ มคเธสุ ปญฺจาพาธา, \\ มาคโธ จ อนุญฺญาสิ, \\ ลกฺขณา อิณา ทาโส จ, \\ สทฺธํ กุลํ กณฺฏโก จ, \\ วตฺถุมฺหิ ทารโก สิกฺขา, \\ สพฺพํ มุขํ อุปชฺฌาเย, \\ ปณฺฑโก เถยฺยปกฺกนฺโต, \\ อรหนฺตภิกฺขุนีเภทา, \\ อนุปชฺฌาย สํเฆน, \\ อจีวรํ ตทุภยํ, \\ หตฺถา ปาทา หตฺถปาทา, \\ องฺคุลิอฬกณฺฑรํ, \\ คลคณฺฑี ลกฺขณา เจว, \\ ปาปปริสทูสี จ, \\ ขญฺชํ ปกฺขหตญฺเจว, \\ ชรานฺธมูคพธิรํ, \\ อนฺธพธิรํ ยํ วุตฺตํ, \\ อนฺธมูคพธิรญฺจ, \\ วตฺถพฺพญฺจ ตถาทฺธานํ, \\ อาคจฺฉตุ วิวทนฺติ, \\ ทิสฺสนฺติ อุปสมฺปนฺนา, \\ อนนุสิฏฺฐา วิตฺเถนฺติ, \\ สํเฆปิ จ อโถ พาลา, \\ อุลฺลุมฺปตุปสมฺปทา,}
\csromangathagroupwithclosermajor{\csromangathastanza{\csromangathaitembat{131}{วินยมฺหิ มหตฺเถสุ,}{เปสลานํ สุขาวเห.} \\ \csromangathacontbat{นิคฺคหานญฺจ ปาปิจฺเฉ,}{ลชฺชีนํ ปคฺคเหสุ จ.}}\csromangathastanza{\csromangathabat{สาสนาธารเณ เจว,}{สพฺพญฺญุชินโคจเร.} \\ \csromangathabat{อนญฺญวิสเย เขเม,}{สุปญฺญตฺเต อสํสเย.}}\csromangathastanza{\csromangathabat{ขนฺธเก วินเย เจว,}{ปริวาเร จ มาติเก.} \\ \csromangathabat{ยถาตฺถการี กุสโล,}{ปฏิปชฺชติ โยนิโส.}}\csromangathastanza{\csromanfolio{137}\csromangathabat{โย ควํ น วิชานาติ,}{น โส รกฺขติ โคคณํ.} \\ \csromangathabat{เอวํ สีลํ อชานนฺโต,}{กิํ โส รกฺเขยฺย สํวรํ.}}\csromangathastanza{\csromangathabat{ปมุฏฺฐมฺหิ จ สุตฺตนฺเต,}{อภิธมฺเม จ ตาวเท.} \\ \csromangathabat{วินเย อวินฏฺฐมฺหิ,}{ปุน ติฏฺฐติ สาสนํ.}}\csromangathastanza{\csromangathabat{ตสฺมา สงฺคาหณาเหตุํ\csromansharedfootnote{2a91c2d931bcce61}{สงฺคาหนาเหตุํ (ก)},}{อุทฺทานํ อนุปุพฺพโส.} \\ \csromangathabat{ปวกฺขามิ ยถาญายํ,}{สุณาถ มม ภาสโต.}}\csromangathastanza{\csromangathabat{วตฺถุ นิทานํ อาปตฺติ,}{นยา เปยฺยาลเมว จ.} \\ \csromangathabat{ทุกฺกรํ ตํ อเสเสตุํ,}{นยโต ตํ วิชานถาติ.}}\csromangathastanza{\csromangathabat{โพธิ ราชายตนญฺจ,}{อชปาโล สหมฺปติ.} \\ \csromangathabat{พฺรหฺมา อาฬาโร อุทโก,}{ภิกฺขุ จ อุปโก อิสิ.}}\csromangathastanza{\csromangathabat{โกณฺฑญฺโญ วปฺโป ภทฺทิโย,}{มหานาโม จ อสฺสชิ.} \\ \csromangathabat{ยโส จตฺตาโร ปญฺญาส,}{สพฺเพ เปเสสิ โส ทิสา.}}\csromangathastanza{\csromangathabat{วตฺถุ มาเรหิ ติํสา จ,}{อุรุเวลํ ตโย ชฏี.} \\ \csromangathabat{อคฺยาคารํ มหาราชา,}{สกฺโก พฺรหฺมา จ เกวลา.}}\csromangathastanza{\csromangathabat{ปํสุกูลํ โปกฺขรณี,}{สิลา จ กกุโธ สิลา.} \\ \csromangathabat{ชมฺพุ อมฺโพ จ อามโล,}{ปาริปุปฺผญฺจ อาหริ.}}\csromangathastanza{\csromangathabat{ผาลิยนฺตุ อุชฺชลนฺตุ,}{วิชฺฌายนฺตุ จ กสฺสป.} \\ \csromangathabat{นิมุชฺชนฺติ มุขี เมโฆ,}{คยา ลฏฺฐิ จ มาคโธ.}}\csromangathastanza{\csromangathabat{อุปติสฺโส โกลิโต จ,}{อภิญฺญาตา จ ปพฺพชุํ.} \\ \csromangathabat{ทุนฺนิวตฺถา ปณามนา,}{กิโส ลูโข จ พฺราหฺมโณ.}}\csromangathastanza{\csromangathabat{อนาจารํ อาจรติ,}{อุทรํ มาณโว คโณ.} \\ \csromangathabat{วสฺสํ พาเลหิ ปกฺกนฺโต,}{ทส วสฺสานิ นิสฺสโย.}}\csromangathastanza{\csromangathabat{น วตฺตนฺติ ปณาเมตุํ,}{พาลา ปสฺสทฺธิ ปญฺจ ฉ.} \\ \csromangathabat{โย โส อญฺโญ จ นคฺโค จ,}{อจฺฉินฺนชฏิลสากิโย.}}\csromangathastanza{\csromangathabat{มคเธสุ ปญฺจาพาธา,}{เอโก ราชา\csromansharedfootnote{e7e43b1452dc5939}{ภโฏ โจโร (สฺยา)} จ องฺคุลิ.} \\ \csromangathabat{มาคโธ จ อนุญฺญาสิ,}{การา ลิขิ กสาหโต.}}\csromangathastanza{\csromanfolio{138}\csromangathabat{ลกฺขณา อิณา ทาโส จ,}{ภณฺฑุโก อุปาลิ อหิ.} \\ \csromangathabat{สทฺธํ กุลํ กณฺฏโก จ,}{อาหุนฺทริกเมว จ.}}\csromangathastanza{\csromangathabat{วตฺถุมฺหิ ทารโก สิกฺขา,}{วิหรนฺติ จ กิํ นุ โข.} \\ \csromangathabat{สพฺพํ มุขํ อุปชฺฌาเย,}{อปลาฬน กณฺฏโก.}}\csromangathastanza{\csromangathabat{ปณฺฑโก เถยฺยปกฺกนฺโต,}{อหิ จ มาตรี ปิตา.} \\ \csromangathabat{อรหนฺตภิกฺขุนีเภทา,}{รุหิเรน จ พฺยญฺชนํ.}}\csromangathastanza{\csromangathabat{อนุปชฺฌาย สํเฆน,}{คณปณฺฑกปตฺตโก.} \\ \csromangathabat{อจีวรํ ตทุภยํ,}{ยาจิเตนปิ เย ตโย.}}\csromangathastanza{\csromangathabat{หตฺถา ปาทา หตฺถปาทา,}{กณฺณา นาสา ตทูภยํ.} \\ \csromangathabat{องฺคุลิอฬกณฺฑรํ,}{ผณํ ขุชฺชญฺจ วามนํ.}}\csromangathastanza{\csromangathabat{คลคณฺฑี ลกฺขณา เจว,}{กสา ลิขิตสีปที.} \\ \csromangathabat{ปาปปริสทูสี จ,}{กาณํ กุณิ ตเถว จ.}}\csromangathastanza{\csromangathabat{ขญฺชํ ปกฺขหตญฺเจว,}{สจฺฉินฺนอิริยาปถํ.} \\ \csromangathabat{ชรานฺธมูคพธิรํ,}{อนฺธมูคญฺจ ยํ ตหิํ.}}\csromangathastanza{\csromangathabat{อนฺธพธิรํ ยํ วุตฺตํ,}{มูคพธิรเมว จ.} \\ \csromangathabat{อนฺธมูคพธิรญฺจ,}{อลชฺชีนญฺจ นิสฺสยํ.}}\csromangathastanza{\csromangathabat{วตฺถพฺพญฺจ ตถาทฺธานํ,}{ยาจมาเนน ลกฺขณา\csromansharedfootnote{4dd984373e76b51f}{เปกฺขนา (สพฺพตฺถ)}.} \\ \csromangathabat{อาคจฺฉตุ วิวทนฺติ,}{เอกุปชฺฌาเยน กสฺสโป.}}\csromangathastanza{\csromangathabat{ทิสฺสนฺติ อุปสมฺปนฺนา,}{อาพาเธหิ จ ปีฬิตา.} \\ \csromangathabat{อนนุสิฏฺฐา วิตฺเถนฺติ,}{ตตฺเถว อนุสาสนา.}}\csromangathastanza{\csromangathabat{สํเฆปิ จ อโถ พาลา,}{อสมฺมตา จ เอกโต.} \\ \csromangathabat{อุลฺลุมฺปตุปสมฺปทา,}{นิสฺสโย เอกโก ตโยติ.}}\csromangathastanza{\csromangathawak{อิมมฺหิ ขนฺธเก วตฺถูนิ เอกสตญฺจ ทฺวาสตฺตติ.}}}{\textbf{มหาขนฺธโก นิฏฺฐิโต.}}
\csromanmatikaanchor{mtk.69}
\csromantocmarknopage{chapter}{2. อุโปสถกฺขนฺธก}{mtk.69}
\bookmark[dest={mtk.69},level=0]{2. อุโปสถกฺขนฺธก}
\chapterheadpageread{2. \textbf{อุโปสถ}\-\textbf{กฺขนฺธก}}
\csromanmatikaanchor{mtk.70}
\csromantocmarkref{section}{68. สนฺนิปาตานุชานนา}{mtk.70}
\bookmark[dest={mtk.70},level=1]{68. สนฺนิปาตานุชานนา}
\csromanheader{1}{68. \textbf{สนฺนิปาตา}\-\textbf{นุชานนา}}
\proseitem{132}{\csromanfolio{139}เตน สมเยน พุทฺโธ ภควา ราชคเห วิหรติ คิชฺฌกูเฏ ปพฺพเต. เตน โข ปน สมเยน อญฺญติตฺถิยา ปริพฺพาชกา จาตุทฺทเส ปนฺนรเส อฏฺฐมิยา จ ปกฺขสฺส สนฺนิปติตฺวา ธมฺมํ ภาสนฺติ. เต มนุสฺสา อุปสงฺกมนฺติ ธมฺม\-สฺสวนาย, เต ลภนฺติ อญฺญติตฺถิเยสุ ปริพฺพาชเกสุ เปมํ, ลภนฺติ ปสาทํ, ลภนฺติ อญฺญติตฺถิยา ปริพฺพาชกา ปกฺขํ. อถ โข รญฺโญ มาคธสฺส เสนิยสฺส พิมฺพิสารสฺส รโหคตสฺส ปฏิสลฺลีนสฺส เอวํ เจตโส ปริวิตกฺโก อุทปาทิ “เอตรหิ โข อญฺญติตฺถิยา ปริพฺพาชกา จาตุทฺทเส ปนฺนรเส อฏฺฐมิยา จ ปกฺขสฺส สนฺนิปติตฺวา ธมฺมํ ภาสนฺติ. เต มนุสฺสา อุปสงฺกมนฺติ ธมฺม\-สฺสวนาย, เต ลภนฺติ อญฺญติตฺถิเยสุ ปริพฺพาชเกสุ เปมํ, ลภนฺติ ปสาทํ, ลภนฺติ อญฺญติตฺถิยา ปริพฺพาชกา ปกฺขํ, ยนฺนูน อยฺยาปิ จาตุทฺทเส ปนฺนรเส อฏฺฐมิยา จ ปกฺขสฺส สนฺนิปเตยฺยุนฺ”ติ. อถ โข ราชา มาคโธ เสนิโย พิมฺพิสาโร เยน ภควา เตนุปสงฺกมิ, อุปสงฺกมิตฺวา ภควนฺตํ อภิวาเทตฺวา เอกมนฺตํ นิสีทิ, เอกมนฺตํ นิสินฺโน โข ราชา มาคโธ เสนิโย พิมฺพิสาโร ภควนฺตํ เอตทโวจ “อิธ มยฺหํ ภนฺเต รโหคตสฺส ปฏิสลฺลีนสฺส เอวํ เจตโส ปริวิตกฺโก อุทปาทิ ‘เอตรหิ โข อญฺญติตฺถิยา ปริพฺพาชกา จาตุทฺทเส ปนฺนรเส อฏฺฐมิยา จ ปกฺขสฺส สนฺนิปติตฺวา ธมฺมํ ภาสนฺติ. เต มนุสฺสา อุปสงฺกมนฺติ ธมฺม\-สฺสวนาย, เต ลภนฺติ อญฺญติตฺถิเยสุ ปริพฺพาชเกสุ เปมํ, ลภนฺติ ปสาทํ, ลภนฺติ อญฺญติตฺถิยา ปริพฺพาชกา ปกฺขํ, ยนฺนูน อยฺยาปิ จาตุทฺทเส ปนฺนรเส อฏฺฐมิยา จ ปกฺขสฺส สนฺนิปเตยฺยุนฺ’ติ. สาธุ ภนฺเต อยฺยาปิ จาตุทฺทเส ปนฺนรเส อฏฺฐมิยา จ ปกฺขสฺส สนฺนิปเตยฺยุนฺ”ติ. อถ โข ภควา ราชานํ มาคธํ เสนิยํ พิมฺพิสารํ ธมฺมิยา กถาย สนฺทสฺเสสิ สมาทเปสิ สมุตฺเตเชสิ สมฺปหํเสสิ. อถ โข ราชา มาคโธ เสนิโย พิมฺพิสาโร ภควตา ธมฺมิยา กถาย สนฺทสฺสิโต สมาทปิโต สมุตฺเตชิโต สมฺปหํสิโต อุฏฺฐายาสนา ภควนฺตํ อภิวาเทตฺวา ปทกฺขิณํ กตฺวา ปกฺกามิ. อถ โข ภควา เอตสฺมิํ นิทาเน เอตสฺมิํ ปกรเณ ธมฺมิํ กถํ กตฺวา ภิกฺขู \csromanfolio{140}อามนฺเตสิ “อนุชานามิ ภิกฺขเว จาตุทฺทเส ปนฺนรเส อฏฺฐมิยา จ ปกฺขสฺส สนฺนิปติตุนฺ”ติ.}

\prose{เตน โข ปน สมเยน ภิกฺขู “ภควตา อนุญฺญาตา จาตุทฺทเส ปนฺนรเส อฏฺฐมิยา จ ปกฺขสฺส สนฺนิปติตุนฺ”ติ จาตุทฺทเส ปนฺนรเส อฏฺฐมิยา จ ปกฺขสฺส สนฺนิปติตฺวา ตุณฺหี นิสีทนฺติ. เต มนุสฺสา อุปสงฺกมนฺติ ธมฺม\-สฺสวนาย. เต อุชฺฌายนฺติ ขิยฺยนฺติ วิปาเจนฺติ “กถํ หิ นาม สมณา สกฺยปุตฺติยา จาตุทฺทเส ปนฺนรเส อฏฺฐมิยา จ ปกฺขสฺส สนฺนิปติตฺวา ตุณฺหี นิสีทิสฺสนฺติ เสยฺยถาปิ มูคสูกรา, นนุ นาม สนฺนิปติเตหิ ธมฺโม ภาสิตพฺโพ”ติ. อสฺโสสุํ โข ภิกฺขู เตสํ มนุสฺสานํ อุชฺฌายนฺตานํ ขิยฺยนฺตานํ วิปาเจนฺตานํ. อถ โข เต ภิกฺขู ภควโต เอตมตฺถํ อาโรเจสุํ ฯเปฯ. อถ โข ภควา เอตสฺมิํ นิทาเน เอตสฺมิํ ปกรเณ ธมฺมิํ กถํ กตฺวา ภิกฺขู อามนฺเตสิ “อนุชานามิ ภิกฺขเว จาตุทฺทเส ปนฺนรเส อฏฺฐมิยา จ ปกฺขสฺส สนฺนิปติตฺวา ธมฺมํ ภาสิตุนฺ”ติ.}

\csromanmatikaanchor{mtk.71}
\csromantocmarkref{section}{69. ปาติโมกฺขุทฺเทสานุชานนา}{mtk.71}
\bookmark[dest={mtk.71},level=1]{69. ปาติโมกฺขุทฺเทสานุชานนา}
\csromanheader{1}{96. \textbf{ปาติโมกฺขุ}\-\textbf{ทฺเทสา}\-\textbf{นุชานนา}}
\proseitem{133}{อถ โข ภควโต รโหคตสฺส ปฏิสลฺลีนสฺส เอวํ เจตโส ปริวิตกฺโก อุทปาทิ “ยนฺนูนาหํ ยานิ มยา ภิกฺขูนํ ปญฺญตฺตานิ สิกฺขาปทานิ, ตานิ เนสํ ปาติโมกฺขุ\-ทฺเทสํ อนุชาเนยฺยํ, โส เนสํ ภวิสฺสติ อุโปสถกมฺมนฺ”ติ. อถ โข ภควา สายนฺหสมยํ ปฏิสลฺลานา วุฏฺฐิโต เอตสฺมิํ นิทาเน เอตสฺมิํ ปกรเณ ธมฺมิํ กถํ กตฺวา ภิกฺขู อามนฺเตสิ “อิธ มยฺหํ ภิกฺขเว รโหคตสฺส ปฏิสลฺลีนสฺส เอวํ เจตโส ปริวิตกฺโก อุทปาทิ ยนฺนูนาหํ ยานิ มยา ภิกฺขูนํ ปญฺญตฺตานิ สิกฺขาปทานิ, ตานิ เนสํ ปาติโมกฺขุ\-ทฺเทสํ อนุชาเนยฺยํ, โส เนสํ ภวิสฺสติ อุโปสถกมฺมนฺ”ติ. อนุชานามิ ภิกฺขเว ปาติโมกฺขํ อุทฺทิสิตุํ. เอวญฺจ ปน ภิกฺขเว อุทฺทิสิตพฺพํ, พฺยตฺเตน ภิกฺขุนา ปฏิพเลน สํโฆ ญาเปตพฺโพ\csromandash{}}

\proseitem{134}{“สุณาตุ เม ภนฺเต สํโฆ, ยทิ สํฆสฺส ปตฺตกลฺลํ, สํโฆ อุโปสถํ กเรยฺย, ปาติโมกฺขํ อุทฺทิเสยฺย. กิํ สํฆสฺส ปุพฺพกิจฺจํ, ปาริสุทฺธิํ อายสฺมนฺโต อาโรเจถ, ปาติโมกฺขํ อุทฺทิสิสฺสามิ, ตํ สพฺเพว สนฺตา สาธุกํ สุโณม มนสิ กโรม. ยสฺส สิยา \csromanfolio{141}อาปตฺติ, โส อาวิกเรยฺย, อสนฺติยา \textbf{อาปตฺติ}ยา ตุณฺหี ภวิตพฺพํ, ตุณฺหีภาเวน โข ปนายสฺมนฺเต ปริสุทฺธาติ เวทิสฺสามิ. ยถา โข ปน ปจฺเจก\-ปุฏฺฐสฺส เวยฺยากรณํ โหติ, เอวเมวํ\csromansharedfootnote{66733781a5984088}{เอวเมว (ก)} เอวรูปาย ปริสาย ยาวตติยํ อนุสฺสาวิตํ โหติ. โย ปน ภิกฺขุ ยาวตติยํ อนุสฺสาวิยมาเน สรมาโน \textbf{สนฺติํ อาปตฺติํ นาวิกเรยฺย, สมฺปชาน}\-\textbf{มุสาวาทสฺส โหติ.} สมฺปชาน\-มุสาวาโท โข ปนายสฺมนฺโต อนฺตรายิโก ธมฺโม วุตฺโต ภควตา, ตสฺมา สรมาเนน ภิกฺขุนา อาปนฺเนน วิสุทฺธา\-เปกฺเขน สนฺตี อาปตฺติ อาวิกาตพฺพา, อาวิกตา หิสฺส ผาสุ โหตี”ติ.}

\proseitem{135}{\textbf{ปาติโมกฺข}นฺติ อาทิเมตํ มุขเมตํ ปมุขเมตํ กุสลานํ ธมฺมานํ, เตน วุจฺจติ “ปาติโมกฺขนฺ”ติ. \textbf{อายสฺมนฺโต}ติ ปิยวจนเมตํ, ครุวจนเมตํ, สคารวสปฺปติสฺสาธิวจนเมตํ “อายสฺมนฺโต”ติ. \textbf{อุทฺทิสิสฺสามี}ติ อาจิกฺขิสฺสามิ เทเสสฺสามิ ปญฺญเปสฺสามิ ปฏฺฐเปสฺสามิ วิวริสฺสามิ วิภชิสฺสามิ อุตฺตานิํ กริสฺสามิ\csromansharedfootnote{9d0818a1bc7c0c2a}{อุตฺตานี กริสฺสามิ (สี, สฺยา)} ปกาเสสฺสามิ. \textbf{ต}นฺติ ปาติโมกฺขํ วุจฺจติ. \textbf{สพฺเพว สนฺตา}ติ ยาวติกา ตสฺสา ปริสาย เถรา จ นวา จ มชฺฌิมา จ, เอเต วุจฺจนฺติ “สพฺเพว สนฺตา”ติ. \textbf{สาธุกํ สุโณมา}ติ อฏฺฐิํ กตฺวา มนสิ กตฺวา สพฺพเจตสา\csromansharedfootnote{a43900b3058e4d83}{สพฺพํ เจตสา (สฺยา, ก)} สมนฺนาหราม. \textbf{มนสิ กโรมา}ติ เอกคฺคจิตฺตา อวิกฺขิตฺตจิตฺตา อวิสาหฏจิตฺตา นิสาเมม. \textbf{ยสฺส สิยา อาปตฺตี}ติ เถรสฺส วา นวสฺส วา มชฺฌิมสฺส วา ปญฺจนฺนํ วา อาปตฺติ\-กฺขนฺธานํ อญฺญตรา อาปตฺติ, สตฺตนฺนํ วา อาปตฺติ\-กฺขนฺธานํ อญฺญตรา อาปตฺติ. \textbf{โส อาวิกเรยฺยา}ติ โส เทเสยฺย, โส วิวเรยฺย, โส อุตฺตานิํ กเรยฺย, โส ปกาเสยฺย สํฆมชฺเฌ วา คณมชฺเฌ วา เอกปุคฺคเล วา. \textbf{อสนฺตี} นาม \textbf{อาปตฺติ} อนชฺฌาปนฺนา วา โหติ, อาปชฺชิตฺวา วา วุฏฺฐิตา. \textbf{ตุณฺหี ภวิตพฺพ}นฺติ อธิวาเสตพฺพํ, น พฺยาหริตพฺพํ. \textbf{ปริสุทฺธาติ} \textbf{เวทิสฺสามี}ติ ชานิสฺสามิ ธาเรสฺสามิ. \textbf{ยถา โข ปน ปจฺเจก}\-\textbf{ปุฏฺฐสฺส เวยฺยากรณํ โหตี}ติ ยถา เอเกน เอโก ปุฏฺโฐ พฺยากเรยฺย, เอวเมว ตสฺสา ปริสาย ชานิตพฺพํ “มํ ปุจฺฉตี”ติ. \textbf{เอวรูปา} นาม \textbf{ปริสา} ภิกฺขุปริสา วุจฺจติ. \textbf{ยาวตติยํ อนุสฺสาวิตํ โหตี}ติ สกิมฺปิ อนุสฺสาวิตํ โหติ, ทุติยมฺปิ อนุสฺสาวิตํ โหติ, ตติยมฺปิ \csromanfolio{142}อนุสฺสาวิตํ โหติ. \textbf{สรมาโน}ติ ชานมาโน สญฺชานมาโน. \textbf{สนฺตี} นาม \textbf{อาปตฺติ} อชฺฌาปนฺนา วา โหติ, อาปชฺชิตฺวา วา อวุฏฺฐิตา. \textbf{นาวิกเรยฺยา}ติ น เทเสยฺย, น วิวเรยฺย, น อุตฺตานิํ กเรยฺย, น ปกาเสยฺย สํฆมชฺเฌ วา คณมชฺเฌ วา เอกปุคฺคเล วา. สฺ\textbf{อมฺปชานมุสาวาทสฺส โหตี}ติ สมฺปชาน\-มุสาวาเท กิํ โหติ, ทุกฺกฏํ โหติ. \textbf{อนฺตรายิโก ธมฺโม วุตฺโต ภควตา}ติ กิสฺส อนฺตรายิโก, ปฐมสฺส ฌานสฺส อธิคมาย อนฺตรายิโก, ทุติยสฺส ฌานสฺส อธิคมาย อนฺตรายิโก, ตติยสฺส ฌานสฺส อธิคมาย อนฺตรายิโก, จตุตฺถสฺส ฌานสฺส อธิคมาย อนฺตรายิโก, ฌานานํ วิโมกฺขานํ สมาธีนํ สมาปตฺตีนํ เนกฺขมฺมานํ นิสฺสรณานํ ปวิเวกานํ กุสลานํ ธมฺมานํ อธิคมาย อนฺตรายิโก. \textbf{ตสฺมา}ติ ตํการณา. \textbf{สรมาเนนา}ติ ชานมาเนน สญฺชานมาเนน. \textbf{วิสุทฺธาเปกฺเขนา}ติ วุฏฺฐาตุกาเมน วิสุชฺฌิตุ\-กาเมน. \textbf{สนฺตี} นาม \textbf{อาปตฺติ} อชฺฌาปนฺนา วา โหติ, อาปชฺชิตฺวา วา อวุฏฺฐิตา. \textbf{อาวิกาตพฺพา}ติ อาวิกาตพฺพา สํฆมชฺเฌ วา คณมชฺเฌ วา เอกปุคฺคเล วา. \textbf{อาวิกตา หิสฺส ผาสุ โหตี}ติ กิสฺส ผาสุ โหติ, ปฐมสฺส ฌานสฺส อธิคมาย ผาสุ โหติ, ทุติยสฺส ฌานสฺส อธิคมาย ผาสุ โหติ, ตติยสฺส ฌานสฺส อธิคมาย ผาสุ โหติ, จตุตฺถสฺส ฌานสฺส อธิคมาย ผาสุ โหติ, ฌานานํ วิโมกฺขานํ สมาธีนํ สมาปตฺตีนํ เนกฺขมฺมานํ นิสฺสรณานํ ปวิเวกานํ กุสลานํ ธมฺมานํ อธิคมาย ผาสุ โหตีติ.}

\proseitem{136}{เตน โข ปน สมเยน ภิกฺขู “ภควตา ปาติโมกฺขุ\-ทฺเทโส อนุญฺญาโต”ติ เทวสิกํ ปาติโมกฺขํ อุทฺทิสนฺติ. ภควโต เอตมตฺถํ อาโรเจสุํ. น ภิกฺขเว เทวสิกํ ปาติโมกฺขํ อุทฺทิสิตพฺพํ, โย อุทฺทิเสยฺย, อาปตฺติ ทุกฺกฏสฺส. อนุชานามิ ภิกฺขเว อุโปสเถ ปาติโมกฺขํ อุทฺทิสิตุนฺติ.}

\prose{เตน โข ปน สมเยน ภิกฺขู “ภควตา อุโปสเถ ปาติโมกฺขุ\-ทฺเทโส อนุญฺญาโต”ติ ปกฺขสฺส ติกฺขตฺตุํ ปาติโมกฺขํ อุทฺทิสนฺติ จาตุทฺทเส ปนฺนรเส อฏฺฐมิยา จ ปกฺขสฺส. ภควโต เอตมตฺถํ อาโรเจสุํ. น ภิกฺขเว ปกฺขสฺส ติกฺขตฺตุํ ปาติโมกฺขํ อุทฺทิสิตพฺพํ, โย อุทฺทิเสยฺย, อาปตฺติ ทุกฺกฏสฺส. อนุชานามิ ภิกฺขเว สกิํ ปกฺขสฺส จาตุทฺทเส วา ปนฺนรเส วา ปาติโมกฺขํ อุทฺทิสิตุนฺติ.}

\prose{\csromanfolio{143}เตน โข ปน สมเยน ฉพฺพคฺคิยา ภิกฺขู ยถาปริสาย ปาติโมกฺขํ อุทฺทิสนฺติ สกาย สกาย ปริสาย. ภควโต เอตมตฺถํ อาโรเจสุํ. น ภิกฺขเว ยถาปริสาย ปาติโมกฺขํ อุทฺทิสิตพฺพํ สกาย สกาย ปริสาย, โย อุทฺทิเสยฺย, อาปตฺติ ทุกฺกฏสฺส. อนุชานามิ ภิกฺขเว สมคฺคานํ อุโปสถ\-กมฺมนฺติ.}

\prose{อถ โข ภิกฺขูนํ เอตทโหสิ “ภควตา ปญฺญตฺตํ ‘สมคฺคานํ อุโปสถกมฺมนฺ’ติ, กิตฺตาวตา นุ โข สามคฺคี โหติ, ยาวตา เอกาวาโส, อุทาหุ สพฺพา ปถวี”ติ.}

\prose{ภควโต เอตมตฺถํ อาโรเจสุํ. อนุชานามิ ภิกฺขเว เอตฺตาวตา สามคฺคี, ยาวตา เอกาวาโสติ.}

\csromanmatikaanchor{mtk.72}
\csromantocmarkref{section}{70. มหากปฺปินวตฺถุ}{mtk.72}
\bookmark[dest={mtk.72},level=1]{70. มหากปฺปินวตฺถุ}
\csromanheader{1}{70. \textbf{มหากปฺปิน}\-\textbf{วตฺถุ}}
\proseitem{137}{เตน โข ปน สมเยน อายสฺมา มหากปฺปิโน ราชคเห วิหรติ มทฺทกุจฺฉิมฺหิ มิคทาเย. อถ โข อายสฺมโต มหากปฺปินสฺส รโหคตสฺส ปฏิสลฺลีนสฺส เอวํ เจตโส ปริวิตกฺโก อุทปาทิ “คจฺเฉยฺยํ วาหํ อุโปสถํ น วา คจฺเฉยฺยํ, คจฺเฉยฺยํ วาหํ สํฆกมฺมํ น วา คจฺเฉยฺยํ, อถ ขฺวาหํ วิสุทฺโธ ปรมาย วิสุทฺธิยา”ติ. อถ โข ภควา อายสฺมโต มหากปฺปินสฺส เจตสา เจโต\-ปริวิตกฺกม\-ญฺญาย เสยฺยถาปิ นาม พลวา ปุริโส สมิญฺชิตํ วา พาหํ ปสาเรยฺย, ปสาริตํ วา พาหํ สมิญฺเชยฺย, เอวเมว คิชฺฌกูเฏ ปพฺพเต อนฺตรหิโต มทฺทกุจฺฉิมฺหิ มิคทาเย อายสฺมโต มหากปฺปินสฺส สมฺมุเข ปาตุรโหสิ, นิสีทิ ภควา ปญฺญตฺเต อาสเน. อายสฺมาปิ โข มหากปฺปิโน ภควนฺตํ อภิวาเทตฺวา เอกมนฺตํ นิสีทิ, เอกมนฺตํ นิสินฺนํ โข อายสฺมนฺตํ มหากปฺปินํ ภควา เอตทโวจ “นนุ เต กปฺปิน รโหคตสฺส ปฏิสลฺลีนสฺส เอวํ เจตโส ปริวิตกฺโก อุทปาทิ คจฺเฉยฺยํ วาหํ อุโปสถํ น วา คจฺเฉยฺยํ, คจฺเฉยฺยํ วาหํ สํฆกมฺมํ น วา คจฺเฉยฺยํ, อถ ขฺวาหํ วิสุทฺโธ ปรมาย วิสุทฺธิยา”ติ. เอวํ ภนฺเต. ตุเมฺห เจ พฺราหฺมณา อุโปสถํ น สกฺกริสฺสถ น ครุกริสฺสถ\csromansharedfootnote{d25b2759f3400ad1}{น ครุํ กริสฺสถ (ก)} น มาเนสฺสถ น ปูเชสฺสถ, อถ โก จรหิ อุโปสถํ สกฺกริสฺสติ ครุกริสฺสติ มาเนสฺสติ ปูเชสฺสติ. คจฺฉ ตฺวํ พฺราหฺมณ อุโปสถํ, มา โน อคมาสิ. คจฺฉ ตฺวํ สํฆกมฺมํ, มา โน อคมาสีติ. “เอวํ ภนฺเต”ติ โข อายสฺมา \csromanfolio{144}มหากปฺปิโน ภควโต ปจฺจสฺโสสิ. อถ โข ภควา อายสฺมนฺตํ มหากปฺปินํ ธมฺมิยา กถาย สนฺทสฺเสตฺวา สมาทเปตฺวา สมุตฺเตเชตฺวา สมฺปหํเสตฺวา เสยฺยถาปิ นาม พลวา ปุริโส สมิญฺชิตํ วา พาหํ ปสาเรยฺย, ปสาริตํ วา พาหํ สมิญฺเชยฺย, เอวเมว มทฺทกุจฺฉิมฺหิ มิคทาเย อายสฺมโต มหากปฺปินสฺส สมฺมุเข อนฺตรหิโต คิชฺฌกูเฏ ปพฺพเต ปาตุรโหสิ.}

\csromanmatikaanchor{mtk.73}
\csromantocmarkref{section}{71. สีมานุชานนา}{mtk.73}
\bookmark[dest={mtk.73},level=1]{71. สีมานุชานนา}
\csromanheader{1}{71. \textbf{สีมานุชานนา}}
\proseitem{138}{อถ โข ภิกฺขูนํ เอตทโหสิ “ภควตา ปญฺญตฺตํ ‘เอตฺตาวตา สามคฺคี, ยาวตา เอกาวาโส’ติ, กิตฺตาวตา นุ โข เอกาวาโส โหตี”ติ. ภควโต เอตมตฺถํ อาโรเจสุํ. อนุชานามิ ภิกฺขเว สีมํ สมฺมนฺนิตุํ. เอวญฺจ ปน ภิกฺขเว สมฺมนฺนิตพฺพา, ปฐมํ นิมิตฺตา กิตฺเตตพฺพา ปพฺพต\-นิมิตฺตํ ปาสาณนิมิตฺตํ วนนิมิตฺตํ รุกฺขนิมิตฺตํ มคฺคนิมิตฺตํ วมฺมิก\-นิมิตฺตํ นทีนิมิตฺตํ อุทกนิมิตฺตํ. นิมิตฺเต กิตฺเตตฺวา พฺยตฺเตน ภิกฺขุนา ปฏิพเลน สํโฆ ญาเปตพฺโพ\csromandash{}}

\proseitem{139}{“สุณาตุ เม ภนฺเต สํโฆ, ยาวตา สมนฺตา นิมิตฺตา กิตฺติตา, ยทิ สํฆสฺส ปตฺตกลฺลํ, สํโฆ เอเตหิ นิมิตฺเตหิ สีมํ สมฺมนฺเนยฺย สมานสํวาสํ เอกุโปสถํ\csromansharedfootnote{460399214dd81796}{เอกูโปสถํ (ก)}, เอสา ญตฺติ.}

\prose{สุณาตุ เม ภนฺเต สํโฆ, ยาวตา สมนฺตา นิมิตฺตา กิตฺติตา, สํโฆ เอเตหิ นิมิตฺเตหิ สีมํ สมฺมนฺนติ สมานสํวาสํ เอกุโปสถํ, ยสฺสายสฺมโต ขมติ เอเตหิ นิมิตฺเตหิ สีมาย สมฺมุติ\csromansharedfootnote{5c6f5c0884ad928c}{สมฺมติ (สฺยา)} สมาน\-สํวาสาย เอกุโปสถาย, โส ตุณฺหสฺส, ยสฺส นกฺขมติ, โส ภาเสยฺย.}

\prose{สมฺมตา สีมา สํเฆน เอเตหิ นิมิตฺเตหิ สมานสํวาสา เอกุโปสถา, ขมติ สํฆสฺส, ตสฺมา ตุณฺหี, เอวเมตํ ธารยามี”ติ.}

\proseitem{140}{เตน โข ปน สมเยน ฉพฺพคฺคิยา ภิกฺขู “ภควตา สีมาสมฺมุติ อนุญฺญาตา”ติ อติมหติโย สีมาโย สมฺมนฺนนฺติ จตุโยชนิกาปิ ปญฺจ\-โยชนิ\-กาปิ ฉโยชนิกาปิ. ภิกฺขู อุโปสถํ อาคจฺฉนฺตา อุทฺทิสฺสมาเนปิ ปาติโมกฺเข อาคจฺฉนฺติ, อุทฺทิฏฺฐ\-มตฺเตปิ อาคจฺฉนฺติ, \csromanfolio{145}อนฺตราปิ ปริวสนฺติ. ภควโต เอตมตฺถํ อาโรเจสุํ. น ภิกฺขเว อติมหตี สีมา สมฺมนฺนิตพฺพา จตุโยชนิกา วา ปญฺจโยชนิกา วา ฉโยชนิกา วา, โย สมฺมนฺเนยฺย, อาปตฺติ ทุกฺกฏสฺส. อนุชานามิ ภิกฺขเว ติโยชนปรมํ สีมํ สมฺมนฺนิตุนฺติ.}

\prose{เตน โข ปน สมเยน ฉพฺพคฺคิยา ภิกฺขู นทีปารสีมํ\csromansharedfootnote{f5ff6cd8d628f8ed}{นทีปารํ สีมํ (สี, สฺยา)} สมฺมนฺนนฺติ. อุโปสถํ อาคจฺฉนฺตา ภิกฺขูปิ วุยฺหนฺติ, ปตฺตาปิ วุยฺหนฺติ, จีวรานิปิ วุยฺหนฺติ, ภควโต เอตมตฺถํ อาโรเจสุํ. น ภิกฺขเว นทีปารสีมา สมฺมนฺนิตพฺพา, โย สมฺมนฺเนยฺย, อาปตฺติ ทุกฺกฏสฺส. อนุชานามิ ภิกฺขเว ยตฺถสฺส ธุวนาวา วา ธุวเสตุ วา, เอวรูปํ นทีปารสีมํ สมฺมนฺนิตุนฺติ.}

\csromanmatikaanchor{mtk.74}
\csromantocmarkref{section}{72. อุโปสถาคารกถา}{mtk.74}
\bookmark[dest={mtk.74},level=1]{72. อุโปสถาคารกถา}
\csromanheader{1}{72. \textbf{อุโปสถา}\-\textbf{คาร}\-\textbf{กถา}}
\proseitem{141}{เตน โข ปน สมเยน ภิกฺขู อนุปริเวณิยํ ปาติโมกฺขํ อุทฺทิสนฺติ อสงฺเกเตน. อาคนฺตุกา ภิกฺขู น ชานนฺติ “กตฺถ วา อชฺชุโปสโถ กรียิสฺสตี”ติ. ภควโต เอตมตฺถํ อาโรเจสุํ. น ภิกฺขเว อนุปริเวณิยํ ปาติโมกฺขํ อุทฺทิสิตพฺพํ อสงฺเกเตน, โย อุทฺทิเสยฺย, อาปตฺติ ทุกฺกฏสฺส. อนุชานามิ ภิกฺขเว อุโปสถาคารํ สมฺมนฺนิตฺวา อุโปสถํ กาตุํ, ยํ สํโฆ อากงฺขติ วิหารํ วา อฑฺฒโยคํ วา ปาสาทํ วา หมฺมิยํ วา คุหํ วา. เอวญฺจ ปน ภิกฺขเว สมฺมนฺนิตพฺพํ, พฺยตฺเตน ภิกฺขุนา ปฏิพเลน สํโฆ ญาเปตพฺโพ\csromandash{}}

\prose{“สุณาตุ เม ภนฺเต สํโฆ, ยทิ สํฆสฺส ปตฺตกลฺลํ, สํโฆ อิตฺถนฺนามํ วิหารํ อุโปสถาคารํ สมฺมนฺเนยฺย, เอสา ญตฺติ.}

\prose{สุณาตุ เม ภนฺเต สํโฆ, สํโฆ อิตฺถนฺนามํ วิหารํ อุโปสถาคารํ สมฺมนฺนติ, ยสฺสายสฺมโต ขมติ อิตฺถนฺนามสฺส วิหารสฺส อุโปสถา\-คารสฺส สมฺมุติ, โส ตุณฺหสฺส, ยสฺส นกฺขมติ, โส ภาเสยฺย.}

\prose{สมฺมโต สํเฆน อิตฺถนฺนาโม วิหาโร อุโปสถาคารํ, ขมติ สํฆสฺส, ตสฺมา ตุณฺหี, เอวเมตํ ธารยามี”ติ.}

\proseitem{141}{เตน โข ปน สมเยน อญฺญตรสฺมิํ อาวาเส เทฺว อุโปสถาคารานิ สมฺมตานิ โหนฺติ. ภิกฺขู อุภยตฺถ สนฺนิปตนฺติ “อิธ อุโปสโถ กรียิสฺสติ, อิธ อุโปสโถ กรียิสฺสตี”ติ. ภควโต เอตมตฺถํ \csromanfolio{146}อาโรเจสุํ. น ภิกฺขเว เอกสฺมิํ อาวาเส เทฺว อุโปสถาคารานิ สมฺมนฺนิ\-ตพฺพานิ, โย สมฺมนฺเนยฺย, อาปตฺติ ทุกฺกฏสฺส. อนุชานามิ ภิกฺขเว เอกํ สมูหนิตฺวา\csromansharedfootnote{e2960f8143f68c9f}{สมุหนิตฺวา (ก)} เอกตฺถ อุโปสถํ กาตุํ. เอวญฺจ ปน ภิกฺขเว สมูหนฺตพฺพํ, พฺยตฺเตน ภิกฺขุนา ปฏิพเลน สํโฆ ญาเปตพฺโพ\csromandash{}}

\prose{“สุณาตุ เม ภนฺเต สํโฆ, ยทิ สํฆสฺส ปตฺตกลฺลํ, สํโฆ อิตฺถนฺนามํ อุโปสถาคารํ สมูหเนยฺย,\csromansharedfootnote{c12b95b661606aee}{สมุหเนยฺย (ก)} เอสา ญตฺติ.}

\prose{สุณาตุ เม ภนฺเต สํโฆ, สํโฆ อิตฺถนฺนามํ อุโปสถาคารํ สมูหนติ, ยสฺสายสฺมโต ขมติ อิตฺถนฺนามสฺส อุโปสถา\-คารสฺส สมุคฺฆาโต, โส ตุณฺหสฺส, ยสฺส นกฺขมติ, โส ภาเสยฺย.}

\prose{สมูหตํ สํเฆน อิตฺถนฺนามํ อุโปสถาคารํ, ขมติ สํฆสฺส, ตสฺมา ตุณฺหี, เอวเมตํ ธารยามี”ติ.}

\csromanmatikaanchor{mtk.75}
\csromantocmarkref{section}{73. อุโปสถปฺปมุขานุชานนา}{mtk.75}
\bookmark[dest={mtk.75},level=1]{73. อุโปสถปฺปมุขานุชานนา}
\csromanheader{1}{73. \textbf{อุโปสถปฺปมุขา}\-\textbf{นุชานนา}}
\proseitem{142}{เตน โข ปน สมเยน อญฺญตรสฺมิํ อาวาเส อติขุทฺทกํ อุโปสถาคารํ สมฺมตํ โหติ ตทหุโปสเถ มหา\-ภิกฺขุสํโฆ สนฺนิปติโต โหติ. ภิกฺขู อสมฺมตาย ภูมิยา นิสินฺนา ปาติโมกฺขํ อสฺโสสุํ. อถ โข เตสํ ภิกฺขุนํ เอตทโหสิ “ภควตา ปญฺญตฺตํ ‘อุโปสถาคารํ สมฺมนฺนิตฺวา อุโปสโถ กาตพฺโพ’ติ, มยญฺจมฺหา อสมฺมตาย ภูมิยา นิสินฺนา ปาติโมกฺขํ อสฺสุมฺหา, กโต นุ โข อมฺหากํ อุโปสโถ, อกโต นุ โข”ติ. ภควโต เอตมตฺถํ อาโรเจสุํ. สมฺมตาย วา ภิกฺขเว ภูมิยา นิสินฺนา อสมฺมตาย วา ยโต ปาติโมกฺขํ สุณาติ, กโตวสฺส อุโปสโถ. เตน หิ ภิกฺขเว สํโฆ ยาว มหนฺตํ อุโปสถปฺปมุขํ\csromansharedfootnote{c0460c10abe3e281}{อุโปสถมุขํ (สฺยา)} อากงฺขติ, ตาว มหนฺตํ อุโปสถปฺปมุขํ สมฺมนฺนตุ. เอวญฺจ ปน ภิกฺขเว สมฺมนฺนิตพฺพํ, ปฐมํ นิมิตฺตา กิตฺเตตพฺพา, นิมิตฺเต กิตฺเตตฺวา พฺยตฺเตน ภิกฺขุนา ปฏิพเลน สํโฆ ญาเปตพฺโพ\csromandash{}}

\prose{“สุณาตุ เม ภนฺเต สํโฆ, ยาวตา สมนฺตา นิมิตฺตา กิตฺติตา, ยทิ สํฆสฺส ปตฺตกลฺลํ, สํโฆ เอเตหิ นิมิตฺเตหิ อุโปสถปฺปมุขํ สมฺมนฺเนยฺย, เอสา ญตฺติ.}

\prose{สุณาตุ เม ภนฺเต สํโฆ, ยาวตา สมนฺตา นิมิตฺตา กิตฺติตา, สํโฆ เอเตหิ นิมิตฺเตหิ อุโปสถปฺปมุขํ สมฺมนฺนติ, \csromanfolio{147}ยสฺสายสฺมโต ขมติ เอเตหิ นิมิตฺเตหิ อุโปสถปฺปมุขสฺส สมฺมุติ, โส ตุณฺหสฺส, ยสฺส นกฺขมติ, โส ภาเสยฺย. สมฺมตํ สํเฆน เอเตหิ นิมิตฺเตหิ อุโปสถปฺปมุขํ, ขมติ สํฆสฺส, ตสฺมา ตุณฺหี, เอวเมตํ ธารยามี”ติ.}

\prose{เตน โข ปน สมเยน อญฺญตรสฺมิํ อาวาเส ตทหุโปสเถ นวกา ภิกฺขู ปฐมตรํ สนฺนิปติตฺวา “น ตาว เถรา อาคจฺฉนฺตี”ติ ปกฺกมิํสุ. อุโปสโถ วิกาเล อโหสิ. ภควโต เอตมตฺถํ อาโรเจสุํ. อนุชานามิ ภิกฺขเว ตทหุโปสเถ เถเรหิ ภิกฺขูหิ ปฐมตรํ สนฺนิปติตุนฺติ.}

\prose{เตน โข ปน สมเยน ราชคเห สมฺพหุลา อาวาสา สมานสีมา โหนฺติ. ตตฺถ ภิกฺขู “วิวทนฺติ อมฺหากํ อาวาเส อุโปสโถ กรียตุ, อมฺหากํ อาวาเส อุโปสโถ กรียตู”ติ. ภควโต เอตมตฺถํ อาโรเจสุํ. อิธ ปน ภิกฺขเว สมฺพหุลา อาวาสา สมานสีมา โหนฺติ. ตตฺถ ภิกฺขู วิวทนฺติ “อมฺหากํ อาวาเส อุโปสโถ กรียตุ, อมฺหากํ อาวาเส อุโปสโถ กรียตู”ติ. เตหิ ภิกฺขเว ภิกฺขูหิ สพฺเพเหว เอกชฺฌํ สนฺนิปติตฺวา อุโปสโถ กาตพฺโพ. ยตฺถ วา ปน เถโร ภิกฺขุ วิหรติ, ตตฺถ สนฺนิปติตฺวา อุโปสโถ กาตพฺโพ, น เตฺวว วคฺเคน สํเฆน อุโปสโถ กาตพฺโพ, โย กเรยฺย, อาวตฺติ ทุกฺกฏสฺสาติ.}

\csromanmatikaanchor{mtk.76}
\csromantocmarkref{section}{74. อวิปฺปวาสสีมานุชานนา}{mtk.76}
\bookmark[dest={mtk.76},level=1]{74. อวิปฺปวาสสีมานุชานนา}
\csromanheader{1}{74. \textbf{อวิปฺปวาส}\-\textbf{สีมา}\-\textbf{นุชานนา}}
\proseitem{143}{เตน โข ปน สมเยน อายสฺมา มหากสฺสโป อนฺธกวินฺทา ราชคหํ อุโปสถํ อาคจฺฉนฺโต อนฺตรามคฺเค นทิํ ตรนฺโต มนํ วูโฬฺห อโหสิ, จีวรานิสฺส\csromansharedfootnote{30244f06d61f6652}{เตน จีวรานิสฺส (ก)} อลฺลานิ. ภิกฺขู อายสฺมนฺตํ มหากสฺสปํ เอตทโวจุํ “กิสฺส เต อาวุโส จีวรานิ อลฺลานี”ติ. อิธาหํ อาวุโส อนฺธกวินฺทา ราชคหํ อุโปสถํ อาคจฺฉนฺโต อนฺตรามคฺเค นทิํ ตรนฺโต มนมฺหิ วูโฬฺห, เตน เม จีวรานิ อลฺลานีติ. ภควโต เอตมตฺถํ อาโรเจสุํ. ยา สา ภิกฺขเว สํเฆน สีมา สมฺมตา สมานสํวาสา เอกุโปสถา, สํโฆ ตํ สีมํ ติจีวเรน อวิปฺปวาสํ สมฺมนฺนตุ. เอวญฺจ ปน ภิกฺขเว สมฺมนฺนิตพฺพา, พฺยตฺเตน ภิกฺขุนา ปฏิพเลน สํโฆ ญาเปตพฺโพ\csromandash{}}

\prose{\csromanfolio{148}“สุณาตุ เม ภนฺเต สํโฆ, ยา สา สํเฆน สีมา สมฺมตา สมานสํวาสา เอกุโปสถา, ยทิ สํฆสฺส ปตฺตกลฺลํ, สํโฆ ตํ สีมํ ติจีวเรน อวิปฺปวาสํ สมฺมนฺเนยฺย, เอสา ญตฺติ.}

\prose{สุณาตุ เม ภนฺเต สํโฆ, ยา สา สํเฆน สีมา สมฺมตา สมานสํวาสา เอกุโปสถา, สํโฆ ตํ สีมํ ติจีวเรน อวิปฺปวาสํ สมฺมนฺนติ, ยสฺสายสฺมโต ขมติ เอติสฺสา สีมาย ติจีวเรน อวิปฺปวาสาย\csromansharedfootnote{d424aef0904b1d84}{อวิปฺปวาสสฺส (สฺยา)} สมฺมุติ, โส ตุณฺหสฺส, ยสฺส นกฺขมติ, โส ภาเสยฺย.}

\prose{สมฺมตา สา สีมา สํเฆน ติจีวเรน อวิปฺปวาสา\csromansharedfootnote{61785d03c494c2ae}{[มิสฺสินฺคฺ โนเต 1]}, ขมติ สํฆสฺส, ตสฺมา ตุณฺหี, เอวเมตํ ธารยามี”ติ.}

\proseitem{143}{เตน โข ปน สมเยน ภิกฺขู “ภควตา ติจีวเรน อวิปฺปวาส\-สมฺมุติ อนุญฺญาตา”ติ อนฺตรฆเร จีวรานิ นิกฺขิปนฺติ, ตานิ จีวรานิ นสฺสนฺติปิ ฑยฺหนฺติปิ อุนฺทูเรหิปิ ขชฺชนฺติ. ภิกฺขู ทุจฺโจฬา โหนฺติ ลูขจีวรา. ภิกฺขู เอวมาหํสุ “กิสฺส ตุเมฺห อาวุโส ทุจฺโจฬา ลูขจีวรา”ติ. อิธ มยํ อาวุโส ภควตา “ติจีวเรน อวิปฺปวาส\-สมฺมุติ อนุญฺญาตา”ติ อนฺตรฆเร จีวรานิ นิกฺขิปิมฺหา, ตานิ จีวรานิ นฏฺฐานิปิ ทฑฺฒานิปิ อุนฺทูเรหิปิ ขายิตานิ, เตน มยํ ทุจฺโจฬา ลูขจีวราติ. ภควโต เอตมตฺถํ อาโรเจสุํ. ยา สา ภิกฺขเว สํเฆน สีมา สมฺมตา สมานสํวาสา เอกุโปสถา, สํโฆ ตํ สีมํ ติจีวเรน อวิปฺปวาสํ สมฺมนฺนตุ ฐเปตฺวา คามญฺจ คามูปจารญฺจ, เอวญฺจ ปน ภิกฺขเว สมฺมนฺนิตพฺพา, พฺยตฺเตน ภิกฺขุนา ปฏิพเลน สํโฆ ญาเปตพฺโพ\csromandash{}}

\proseitem{144}{“สุณาตุ เม ภนฺเต สํโฆ, ยา สา สํเฆน สีมา สมฺมตา สมานสํวาสา เอกุโปสถา, ยทิ สํฆสฺส ปตฺตกลฺลํ, สํโฆ ตํ สีมํ ติจีวเรน อวิปฺปวาสํ สมฺมนฺเนยฺย ฐเปตฺวา คามญฺจ คามูปจารญฺจ, เอสา ญตฺติ.}

\prose{สุณาตุ เม ภนฺเต สํโฆ, ยา สา สํเฆน สีมา สมฺมตา สมานสํวาสา เอกุโปสถา, สํโฆ ตํ สีมํ ติจีวเรน อวิปฺปวาสํ สมฺมนฺนติ ฐเปตฺวา คามญฺจ คามูปจารญฺจ, ยสฺสายสฺมโต ขมติ เอติสฺสา สีมาย ติจีวเรน อวิปฺปวาสาย\csromansharedfootnote{d424aef0904b1d84}{อวิปฺปวาสสฺส (สฺยา)} สมฺมุติ ฐเปตฺวา คามญฺจ คามูปจารญฺจ, โส ตุณฺหสฺส, ยสฺส นกฺขมติ, โส ภาเสยฺย. สมฺมตา \csromanfolio{149}สา สีมา สํเฆน ติจีวเรน อวิปฺปวาสา\csromansharedfootnote{776935c27fe35206}{อวิปฺปวาโส (สฺยา)} ฐเปตฺวา คามญฺจ คามูปจารญฺจ, ขมติ สํฆสฺส, ตสฺมา ตุณฺหี, เอวเมตํ ธารยามี”ติ.}

\csromanmatikaanchor{mtk.77}
\csromantocmarkref{section}{75. สีมาสมูหนน}{mtk.77}
\bookmark[dest={mtk.77},level=1]{75. สีมาสมูหนน}
\csromanheader{1}{75. \textbf{สีมาสมูหนน}}
\prose{สีมํ ภิกฺขเว สมฺมนฺนนฺเตน ปฐมํ สมาน\-สํวาส\-สีมา\csromansharedfootnote{22e6a39f05a9b881}{สมานสํวาสา สีมา (สฺยา)} สมฺมนฺนิตพฺพา, ปจฺฉา ติจีวเรน อวิปฺปวาโส สมฺมนฺนิตพฺโพ. สีมํ ภิกฺขเว สมูหนนฺเตน ปฐมํ ติจีวเรน อวิปฺปวาโส สมูหนฺตพฺโพ, ปจฺฉา สมาน\-สํวาส\-สีมา สมูหนฺตพฺพา. เอวญฺจ ปน ภิกฺขเว ติจีวเรน อวิปฺปวาโส สมูหนฺตพฺโพ, พฺยตฺเตน ภิกฺขุนา ปฏิพเลน สํโฆ ญาเปตพฺโพ\csromandash{}}

\proseitem{145}{“สุณาตุ เม ภนฺเต สํโฆ, โย โส สํเฆน ติจีวเรน อวิปฺปวาโส สมฺมโต, ยทิ สํฆสฺส ปตฺตกลฺลํ, สํโฆ ตํ ติจีวเรน อวิปฺปวาสํ สมูหเนยฺย, เอสา ญตฺติ.}

\prose{สุณาตุ เม ภนฺเต สํโฆ, โย โส สํเฆน ติจีวเรน อวิปฺปวาโส สมฺมโต, สํโฆ ตํ ติจีวเรน อวิปฺปวาสํ สมูหนติ, ยสฺสายสฺมโต ขมติ เอตสฺส ติจีวเรน อวิปฺปวาสสฺส สมุคฺฆาโต, โส ตุณฺหสฺส, ยสฺส นกฺขมติ, โส ภาเสยฺย.}

\prose{สมูหโต โส สํเฆน ติจีวเรน อวิปฺปวาโส, ขมติ สํฆสฺส, ตสฺมา ตุณฺหี, เอวเมตํ ธารยามี”ติ.}

\proseitem{145}{เอวญฺจ ปน ภิกฺขเว สีมา\csromansharedfootnote{22e6a39f05a9b881}{สมานสํวาสา สีมา (สฺยา)} สมูหนฺตพฺพา, พฺยตฺเตน ภิกฺขุนา ปฏิพเลน สํโฆ ญาเปตพฺโพ\csromandash{}}

\proseitem{146}{“สุณาตุ เม ภนฺเต สํโฆ, ยา สา สํเฆน สีมา สมฺมตา สมานสํวาสา เอกุโปสถา, ยทิ สํฆสฺส ปตฺตกลฺลํ, สํโฆ ตํ สีมํ สมูหเนยฺย สมานสํวาสํ เอกุโปสถํ, เอสา ญตฺติ.}

\prose{สุณาตุ เม ภนฺเต สํโฆ, ยา สา สํเฆน สีมา สมฺมตา สมานสํวาสา เอกุโปสถา, สํโฆ ตํ สีมํ สมูหนติ สมานสํวาสํ เอกุโปสถํ, ยสฺสายสฺมโต ขมติ เอติสฺสา สีมาย สมาน\-สํวาสาย เอกุโปสถาย สมุคฺฆาโต, โส ตุณฺหสฺส, ยสฺส นกฺขมติ, โส ภาเสยฺย.}

\prose{สมูหตา สา สีมา สํเฆน สมานสํวาสา เอกุโปสถา, ขมติ สํฆสฺส, ตสฺมา ตุณฺหี, เอวเมตํ ธารยามี”ติ.}

\csromanmatikaanchor{mtk.78}
\csromantocmarkref{section}{76. คามสีมาทิ}{mtk.78}
\bookmark[dest={mtk.78},level=1]{76. คามสีมาทิ}
\csromanheader{1}{76. \textbf{คามสีมาทิ}}
\proseitem{147}{\csromanfolio{150}อสมฺมตาย ภิกฺขเว สีมาย อฏฺฐปิตาย ยํ คามํ วา นิคมํ วา อุปนิสฺสาย วิหรติ. ยา ตสฺส วา คามสฺส คามสีมา, นิคมสฺส วา นิคมสีมา, อยํ ตตฺถ สมานสํวาสา เอกุโปสถา. อคามเก เจ ภิกฺขเว อรญฺเญ สมนฺตา สตฺตพฺภนฺตรา, อยํ ตตฺถ สมานสํวาสา เอกุโปสถา. สพฺพา ภิกฺขเว นที อสีมา, สพฺโพ สมุทฺโท อสีโม, สพฺโพ ชาตสฺสโร อสีโม. นทิยา วา ภิกฺขเว สมุทฺเท วา ชาตสฺสเร วา ยํ มชฺฌิมสฺส ปุริสสฺส สมนฺตา อุทกุกฺเขปา, อยํ ตตฺถ สมานสํวาสา เอกุโปสถาติ.}

\proseitem{148}{เตน โข ปน สมเยน ฉพฺพคฺคิยา ภิกฺขู สีมาย สีมํ สมฺภินฺทนฺติ. ภควโต เอตมตฺถํ อาโรเจสุํ. เยสํ ภิกฺขเว สีมา ปฐมํ สมฺมตา, เตสํ ตํ กมฺมํ ธมฺมิกํ อกุปฺปํ ฐานารหํ. เยสํ ภิกฺขเว สีมา ปจฺฉา สมฺมตา, เตสํ ตํ กมฺมํ อธมฺมิกํ กุปฺปํ อฏฺฐานารหํ. น ภิกฺขเว สีมาย สีมา สมฺภินฺทิตพฺพา, โย สมฺภินฺเทยฺย, อาปตฺติ ทุกฺกฏสฺสาติ.}

\prose{เตน โข ปน สมเยน ฉพฺพคฺคิยา ภิกฺขู สีมาย สีมํ อชฺโฌตฺถรนฺติ. ภควโต เอตมตฺถํ อาโรเจสุํ. เยสํ ภิกฺขเว สีมา ปฐมํ สมฺมตา, เตสํ ตํ กมฺมํ ธมฺมิกํ อกุปฺปํ ฐานารหํ. เยสํ ภิกฺขเว สีมา ปจฺฉา สมฺมตา, เตสํ ตํ กมฺมํ อธมฺมิกํ กุปฺปํ อฏฺฐานารหํ. น ภิกฺขเว สีมาย สีมา อชฺโฌตฺถริตพฺพา, โย อชฺโฌตฺถเรยฺย, อาปตฺติ ทุกฺกฏสฺสาติ. อนุชานามิ ภิกฺขเว สีมํ สมฺมนฺนนฺเตน สีมนฺตริกํ ฐเปตฺวา สีมํ สมฺมนฺนิตุนฺติ.}

\csromanmatikaanchor{mtk.79}
\csromantocmarkref{section}{77. อุโปสถเภทาทิ}{mtk.79}
\bookmark[dest={mtk.79},level=1]{77. อุโปสถเภทาทิ}
\csromanheader{1}{77. \textbf{อุโปสถ}\-\textbf{เภทาทิ}}
\proseitem{149}{อถ โข ภิกฺขูนํ เอตทโหสิ “กติ นุ โข อุโปสถา”ติ. ภควโต เอตมตฺถํ อาโรเจสุํ. เทฺวเม ภิกฺขเว อุโปสถา จาตุทฺทสิโก จ ปนฺนรสิโก จ. อิเม โข ภิกฺขเว เทฺว อุโปสถาติ.}

\prose{อถ โข ภิกฺขูนํ เอตทโหสิ “กติ นุ โข อุโปสถกมฺมานี”ติ. ภควโต เอตมตฺถํ อาโรเจสุํ. จตฺตาริมานิ ภิกฺขเว \csromanfolio{151}อุโปสถ\-กมฺมานิ, อธมฺเมน วคฺคํ อุโปสถกมฺมํ, อธมฺเมน สมคฺคํ อุโปสถกมฺมํ, ธมฺเมน วคฺคํ อุโปสถกมฺมํ, ธมฺเมน สมคฺคํ อุโปสถ\-กมฺมนฺติ. ตตฺร ภิกฺขเว ยทิทํ อธมฺเมน วคฺคํ อุโปสถกมฺมํ, น ภิกฺขเว เอวรูปํ อุโปสถกมฺมํ กาตพฺพํ, น จ มยา เอวรูปํ อุโปสถกมฺมํ อนุญฺญาตํ. ตตฺร ภิกฺขเว ยทิทํ อธมฺเมน สมคฺคํ อุโปสถกมฺมํ, น ภิกฺขเว เอวรูปํ อุโปสถกมฺมํ กาตพฺพํ, น จ มยา เอวรูปํ อุโปสถกมฺมํ อนุญฺญาตํ. ตตฺร ภิกฺขเว ยทิทํ ธมฺเมน วคฺคํ อุโปสถกมฺมํ, น ภิกฺขเว เอวรูปํ อุโปสถกมฺมํ กาตพฺพํ, น จ มยา เอวรูปํ อุโปสถกมฺมํ อนุญฺญาตํ. ตตฺร ภิกฺขเว ยทิทํ ธมฺเมน สมคฺคํ อุโปสถกมฺมํ, เอวรูปํ ภิกฺขเว อุโปสถกมฺมํ กาตพฺพํ, เอวรูปญฺจ มยา อุโปสถกมฺมํ อนุญฺญาตํ. ตสฺมาติห ภิกฺขเว เอวรูปํ อุโปสถกมฺมํ กริสฺสาม ยทิทํ ธมฺเมน สมคฺคนฺติ. เอวญฺหิ โว ภิกฺขเว สิกฺขิตพฺพนฺติ.}

\csromanmatikaanchor{mtk.80}
\csromantocmarkref{section}{78. สํขิตฺเตน ปาติโมกฺขุทฺเทสาทิ}{mtk.80}
\bookmark[dest={mtk.80},level=1]{78. สํขิตฺเตน ปาติโมกฺขุทฺเทสาทิ}
\csromanheader{1}{78. \textbf{สํขิตฺเตน ปาติโมกฺขุ}\-\textbf{ทฺเทสาทิ}}
\proseitem{150}{อถ โข ภิกฺขูนํ เอตทโหสิ “กติ นุ โข ปาติโมกฺขุ\-ทฺเทสา”ติ. ภควโต เอตมตฺถํ อาโรเจสุํ. ปญฺจิเม ภิกฺขเว ปาติโมกฺขุ\-ทฺเทสา, นิทานํ อุทฺทิสิตฺวา อวเสสํ สุเตน สาเวตพฺพํ, อยํ ปฐโม ปาติโมกฺขุ\-ทฺเทโส. นิทานํ อุทฺทิสิตฺวา จตฺตาริ ปาราชิกานิ อุทฺทิสิตฺวา อวเสสํ สุเตน สาเวตพฺพํ, อยํ ทุติโย ปาติโมกฺขุ\-ทฺเทโส. นิทานํ อุทฺทิสิตฺวา จตฺตาริ ปาราชิกานิ อุทฺทิสิตฺวา เตรส สํฆาทิเสเส อุทฺทิสิตฺวา อวเสสํ สุเตน สาเวตพฺพํ, อยํ ตติโย ปาติโมกฺขุ\-ทฺเทโส. นิทานํ อุทฺทิสิตฺวา จตฺตาริ ปาราชิกานิ อุทฺทิสิตฺวา เตรส สํฆาทิเสเส อุทฺทิสิตฺวา เทฺว อนิยเต อุทฺทิสิตฺวา อวเสสํ สุเตน สาเวตพฺพํ, อยํ จตุตฺโถ ปาติโมกฺขุ\-ทฺเทโส. วิตฺถาเรเนว ปญฺจโม. อิเม โข ภิกฺขเว ปญฺจ ปาติโมกฺขุ\-ทฺเทสาติ.}

\prose{เตน โข ปน สมเยน ภิกฺขู “ภควตา สํขิตฺเตน ปาติโมกฺขุ\-ทฺเทโส อนุญฺญาโต”ติ สพฺพกาลํ สํขิตฺเตน ปาติโมกฺขํ อุทฺทิสนฺติ. ภควโต เอตมตฺถํ อาโรเจสุํ. น ภิกฺขเว สํขิตฺเตน ปาติโมกฺขํ อุทฺทิสิตพฺพํ, โย อุทฺทิเสยฺย, อาปตฺติ ทุกฺกฏสฺสาติ.}

\proseitem{150}{\csromanfolio{152}เตน โข ปน สมเยน โกสเลสุ ชนปเท อญฺญตรสฺมิํ อาวาเส ตทหุโปสเถ สวรภยํ\csromansharedfootnote{dcfee9be1d982bfe}{สํจรภยํ (สฺยา)} อโหสิ. ภิกฺขู นาสกฺขิํสุ วิตฺถาเรน ปาติโมกฺขํ อุทฺทิสิตุํ. ภควโต เอตมตฺถํ อาโรเจสุํ. อนุชานามิ ภิกฺขเว สติ อนฺตราเย สํขิตฺเตน ปาติโมกฺขํ อุทฺทิสิตุนฺติ.}

\prose{เตน โข ปน สมเยน ฉพฺพคฺคิยา ภิกฺขู อสติปิ อนฺตราเย สํขิตฺเตน ปาติโมกฺขํ อุทฺทิสนฺติ. ภควโต เอตมตฺถํ อาโรเจสุํ. น ภิกฺขเว อสติ อนฺตราเย สํขิตฺเตน ปาติโมกฺขํ อุทฺทิสิตพฺพํ, โย อุทฺทิเสยฺย, อาปตฺติ ทุกฺกฏสฺส. อนุชานามิ ภิกฺขเว สติ อนฺตราเย สํขิตฺเตน ปาติโมกฺขํ อุทฺทิสิตุํ. ตตฺริเม อนฺตรายา, ราชนฺตราโย โจรนฺตราโย อคฺยนฺตราโย อุทกนฺตราโย มนุสฺสนฺตราโย อมนุสฺสนฺตราโย วาฬนฺตราโย สรีสปนฺตราโย ชีวิตนฺตราโย พฺรหฺมจริย\-นฺตราโยติ. อนุชานามิ ภิกฺขเว เอวรูเปสุ อนฺตราเยสุ สํขิตฺเตน ปาติโมกฺขํ อุทฺทิสิตุํ, อสติ อนฺตราเย วิตฺถาเรนาติ.}

\prose{เตน โข ปน สมเยน ฉพฺพคฺคิยา ภิกฺขู สํฆมชฺเฌ อนชฺฌิฏฺฐา ธมฺมํ ภาสนฺติ. ภควโต เอตมตฺถํ อาโรเจสุํ. น ภิกฺขเว สํฆมชฺเฌ อนชฺฌิฏฺเฐน ธมฺโม ภาสิตพฺโพ, โย ภาเสยฺย, อาปตฺติ ทุกฺกฏสฺส. อนุชานามิ ภิกฺขเว เถเรน ภิกฺขุนา สามํ วา ธมฺมํ ภาสิตุํ ปรํ วา อชฺเฌสิตุนฺติ.}

\csromanmatikaanchor{mtk.81}
\csromantocmarkref{section}{79. วินยปุจฺฉนกถา}{mtk.81}
\bookmark[dest={mtk.81},level=1]{79. วินยปุจฺฉนกถา}
\csromanheader{1}{79. \textbf{วินย}\-\textbf{ปุจฺฉน}\-\textbf{กถา}}
\proseitem{151}{เตน โข ปน สมเยน ฉพฺพคฺคิยา ภิกฺขู สํฆมชฺเฌ อสมฺมตา วินยํ ปุจฺฉนฺติ. ภควโต เอตมตฺถํ อาโรเจสุํ. น ภิกฺขเว สํฆมชฺเฌ อสมฺมเตน วินโย ปุจฺฉิตพฺโพ, โย ปุจฺเฉยฺย, อาปตฺติ ทุกฺกฏสฺส. อนุชานามิ ภิกฺขเว สํฆมชฺเฌ สมฺมเตน วินยํ ปุจฺฉิตุํ. เอวญฺจ ปน ภิกฺขเว สมฺมนฺนิตพฺโพ, อตฺตนา วา\csromansharedfootnote{3f4250394bbf1e60}{อตฺตนาว (สฺยา)} อตฺตานํ สมฺมนฺนิตพฺพํ, ปเรน วา ปโร สมฺมนฺนิตพฺโพ, กถญฺจ อตฺตนาว อตฺตานํ สมฺมนฺนิตพฺพํ, พฺยตฺเตน ภิกฺขุนา ปฏิพเลน สํโฆ ญาเปตพฺโพ\csromandash{}}

\prose{\csromanfolio{153}“สุณาตุ เม ภนฺเต สํโฆ, ยทิ สํฆสฺส ปตฺตกลฺลํ, อหํ อิตฺถานฺนามํ วินยํ ปุจฺเฉยฺยนฺ”ติ. เอวํ อตฺตนาว อตฺตานํ สมฺมนฺนิตพฺพํ.}

\prose{กถญฺจ ปเรน ปโร สมฺมนฺนิตพฺโพ, พฺยตฺเตน ภิกฺขุนา ปฏิพเลน สํโฆ ญาเปตพฺโพ\csromandash{}}

\prose{“สุณาตุ เม ภนฺเต สํโฆ, ยทิ สํฆสฺส ปตฺตกลฺลํ, อิตฺถนฺนาโม อิตฺถนฺนามํ วินยํ ปุจฺเฉยฺยา”ติ.}

\setlength{\csromangathapagewidth}{0pt}
\setlength{\csromangathawakwidth}{0pt}
\csromangathameasuregroup{}{เอวํ ปเรน ปโร สมฺมนฺนิตพฺโพติ.}
\csromangathameasurewak{เอวํ ปเรน ปโร สมฺมนฺนิตพฺโพติ.}
\setlength{\csromangathaleftcol}{0pt}
\csromangathagroup{\csromangathastanza{\csromangathawak{\csromangathaitemline{151}{เอวํ ปเรน ปโร สมฺมนฺนิ\-ตพฺโพติ.}}}}

\prose{เตน โข ปน สมเยน เปสลา ภิกฺขู สํฆมชฺเฌ สมฺมตา วินยํ ปุจฺฉนฺติ. ฉพฺพคฺคิยา ภิกฺขู ลภนฺติ อาฆาตํ, ลภนฺติ อปฺปจฺจยํ, วเธน ตชฺเชนฺติ. ภควโต เอตมตฺถํ อาโรเจสุํ. อนุชานามิ ภิกฺขเว สํฆมชฺเฌ สมฺมเตนปิ ปริสํ โอโลเกตฺวา ปุคฺคลํ ตุลยิตฺวา วินยํ ปุจฺฉิตุนฺติ.}

\csromanmatikaanchor{mtk.82}
\csromantocmarkref{section}{80. วินยวิสฺสชฺชนกถา}{mtk.82}
\bookmark[dest={mtk.82},level=1]{80. วินยวิสฺสชฺชนกถา}
\csromanheader{1}{80. \textbf{วินย}\-\textbf{วิสฺสชฺชน}\-\textbf{กถา}}
\proseitem{152}{เตน โข ปน สมเยน ฉพฺพคฺคิยา ภิกฺขู สํฆมชฺเฌ อสมฺมตา วินยํ วิสฺสชฺเชนฺติ\csromansharedfootnote{e61cdab786bfe9e7}{วิสฺสชฺชนฺติ (ก)}. ภควโต เอตมตฺถํ อาโรเจสุํ. น ภิกฺขเว สํฆมชฺเฌ อสมฺมเตน วินโย วิสฺสชฺเชตพฺโพ, โย วิสฺสชฺเชยฺย, อาปตฺติ ทุกฺกฏสฺส. อนุชานามิ ภิกฺขเว สํฆมชฺเฌ สมฺมเตน วินยํ วิสฺสชฺเชตุํ. เอวญฺจ ปน ภิกฺขเว สมฺมนฺนิตพฺพํ, อตฺตนา วา\csromansharedfootnote{3f4250394bbf1e60}{อตฺตนาว (สฺยา)} อตฺตานํ สมฺมนฺนิตพฺพํ, ปเรน วา ปโร สมฺมนฺนิตพฺโพ, กถญฺจ อตฺตนาว อตฺตานํ สมฺมนฺนิตพฺพํ, พฺยตฺเตน ภิกฺขุนา ปฏิพเลน สํโฆ ญาเปตพฺโพ\csromandash{}}

\prose{“สุณาตุ เม ภนฺเต สํโฆ, ยทิ สํฆสฺส ปตฺตกลฺลํ, อหํ อิตฺถนฺนาเมน วินยํ ปุฏฺโฐ วิสฺสชฺเชยฺยนฺ”ติ. เอวํ อตฺตนาว อตฺตานํ สมฺมนฺนิตพฺพํ.}

\prose{กถญฺจ ปเรน ปโร สมฺมนฺนิตพฺโพ, พฺยตฺเตน ภิกฺขุนา ปฏิพเลน สํโฆ ญาเปตพฺโพ\csromandash{}}

\prose{“สุณาตุ เม ภนฺเต สํโฆ, ยทิ สํฆสฺส ปตฺตกลฺลํ, อิตฺถนฺนาโม อิตฺถนฺนาเมน วินยํ ปุฏฺโฐ วิสฺสชฺเชยฺยา”ติ.}

\setlength{\csromangathapagewidth}{0pt}
\setlength{\csromangathawakwidth}{0pt}
\csromangathameasuregroup{}{เอวํ ปเรน ปโร สมฺมนฺนิตพฺโพ”ติ.}
\csromangathameasurewak{เอวํ ปเรน ปโร สมฺมนฺนิตพฺโพ”ติ.}
\setlength{\csromangathaleftcol}{0pt}
\csromangathagroup{\csromangathastanza{\csromangathawak{\csromangathaitemline{152}{เอวํ ปเรน ปโร สมฺมนฺนิตพฺโพ”ติ.}}}}

\prose{\csromanfolio{154}เตน โข ปน สมเยน เปสลา ภิกฺขู สํฆมชฺเฌ สมฺมตา วินยํ วิสฺสชฺเชนฺติ. ฉพฺพคฺคิยา ภิกฺขู ลภนฺติ อาฆาตํ, ลภนฺติ อปฺปจฺจยํ, วเธน ตชฺเชนฺติ. ภควโต เอตมตฺถํ อาโรเจสุํ. อนุชานามิ ภิกฺขเว สํฆมชฺเฌ สมฺมเตนปิ ปริสํ โอโลเกตฺวา ปุคฺคลํ ตุลยิตฺวา วินยํ วิสฺสชฺเชตุนฺติ.}

\csromanmatikaanchor{mtk.83}
\csromantocmarkref{section}{81. โจทนากถา}{mtk.83}
\bookmark[dest={mtk.83},level=1]{81. โจทนากถา}
\csromanheader{1}{81. \textbf{โจทนากถา}}
\proseitem{153}{เตน โข ปน สมเยน ฉพฺพคฺคิยา ภิกฺขู อโนกาสกตํ ภิกฺขุํ อาปตฺติยา โจเทนฺติ. ภควโต เอตมตฺถํ อาโรเจสุํ. น ภิกฺขเว อโนกาสกโต ภิกฺขุ อาปตฺติยา โจเทตพฺโพ, โย โจเทยฺย, อาปตฺติ ทุกฺกฏสฺส. อนุชานามิ ภิกฺขเว โอกาสํ การาเปตฺวา อาปตฺติยา โจเทตุํ “กโรตุ อายสฺมา โอกาสํ อหํ ตํ วตฺตุกาโม”ติ.}

\prose{เตน โข ปน สมเยน เปสลา ภิกฺขู ฉพฺพคฺคิเย ภิกฺขู โอกาสํ การาเปตฺวา อาปตฺติยา โจเทนฺติ. ฉพฺพคฺคิยา ภิกฺขู ลภนฺติ อาฆาตํ, ลภนฺติ อปฺปจฺจยํ, วเธน ตชฺเชนฺติ. ภควโต เอตมตฺถํ อาโรเจสุํ. อนุชานามิ ภิกฺขเว กเตปิ โอกาเส ปุคฺคลํ ตุลยิตฺวา อาปตฺติยา โจเทตุนฺติ.}

\prose{เตน โข ปน สมเยน ฉพฺพคฺคิยา ภิกฺขู “ปุรมฺหากํ เปสลา ภิกฺขู โอกาสํ การาเปนฺตี”ติ ปฏิกจฺเจว สุทฺธานํ ภิกฺขูนํ อนาปตฺติกานํ อวตฺถุสฺมิํ อการเณ โอกาสํ การาเปนฺติ. ภควโต เอตมตฺถํ อาโรเจสุํ. น ภิกฺขเว สุทฺธานํ ภิกฺขูนํ อนาปตฺติกานํ อวตฺถุสฺมิํ อการเณ โอกาโส การาเปตพฺโพ, โย การาเปยฺย, อาปตฺติ ทุกฺกฏสฺส. อนุชานามิ ภิกฺขเว ปุคฺคลํ ตุลยิตฺวา โอกาสํ กาตุนฺ\csromansharedfootnote{7afb863e16614a73}{การาเปตุํ (สฺยา)} ติ.}

\csromanmatikaanchor{mtk.84}
\csromantocmarkref{section}{82. อธมฺมกมฺมปฏิกฺโกสนาทิ}{mtk.84}
\bookmark[dest={mtk.84},level=1]{82. อธมฺมกมฺมปฏิกฺโกสนาทิ}
\csromanheader{1}{82. \textbf{อธมฺมกมฺม}\-\textbf{ปฏิกฺโกสนาทิ}}
\proseitem{154}{เตน โข ปน สมเยน ฉพฺพคฺคิยา ภิกฺขู สํฆมชฺเฌ อธมฺมกมฺมํ กโรนฺติ. ภควโต เอตมตฺถํ อาโรเจสุํ. น ภิกฺขเว อธมฺมกมฺมํ กาตพฺพํ, โย กเรยฺย, อาปตฺติ ทุกฺกฏสฺสาติ. กโรนฺติเยว อธมฺมกมฺมํ. ภควโต เอตมตฺถํ อาโรเจสุํ. อนุชานามิ ภิกฺขเว อธมฺมกมฺเม กยิรมาเน ปฏิกฺโกสิตุนฺติ.}

\prose{\csromanfolio{155}เตน โข ปน สมเยน เปสลา ภิกฺขู ฉพฺพคฺคิเยหิ ภิกฺขูหิ อธมฺมกมฺเม กยิรมาเน ปฏิกฺโกสนฺติ. ฉพฺพคฺคิยา ภิกฺขู ลภนฺติ อาฆาตํ, ลภนฺติ อปฺปจฺจยํ, วเธน ตชฺเชนฺติ. ภควโต เอตมตฺถํ อาโรเจสุํ. อนุชานามิ ภิกฺขเว ทิฏฺฐิมฺปิ อาวิกาตุนฺติ. เตสํเยว สนฺติเก ทิฏฺฐิํ อาวิกโรนฺติ. ฉพฺพคฺคิยา ภิกฺขู ลภนฺติ อาฆาตํ, ลภนฺติ อปฺปจฺจยํ, วเธน ตชฺเชนฺติ. ภควโต เอตมตฺถํ อาโรเจสุํ. อนุชานามิ ภิกฺขเว จตูหิ ปญฺจหิ ปฏิกฺโกสิตุํ, ทฺวีหิ ตีหิ ทิฏฺฐิํ อาวิกาตุํ, เอเกน อธิฏฺฐาตุํ “น เมตํ ขมตี”ติ.}

\prose{เตน โข ปน สมเยน ฉพฺพคฺคิยา ภิกฺขู สํฆมชฺเฌ ปาติโมกฺขํ อุทฺทิสมานา สญฺจิจฺจ น สาเวนฺติ. ภควโต เอตมตฺถํ อาโรเจสุํ. น ภิกฺขเว ปาติโมกฺขุ\-ทฺเทสเกน สญฺจิจฺจ น สาเวตพฺพํ, โย น สาเวยฺย, อาปตฺติ ทุกฺกฏสฺสาติ.}

\prose{เตน โข ปน สมเยน อายสฺมา อุทายี สํฆสฺส ปาติโมกฺขุ\-ทฺเทสโก โหติ กากสฺสรโก. อถ โข อายสฺมโต อุทายิสฺส เอตทโหสิ “ภควตา ปญฺญตฺตํ ‘ปาติโมกฺขุ\-ทฺเทสเกน สาเวตพฺพนฺ’ติ, อหญฺจมฺหิ กากสฺสรโก, กถํ นุ โข มยา ปฏิปชฺชิตพฺพนฺ”ติ. ภควโต เอตมตฺถํ อาโรเจสุํ. อนุชานามิ ภิกฺขเว ปาติโมกฺขุ\-ทฺเทสเกน วายมิตุํ “กถํ สาเวยฺยนฺ”ติ. วายมนฺตสฺส อนาปตฺตีติ.}

\prose{เตน โข ปน สมเยน เทวทตฺโต สคหฏฺฐาย ปริสาย ปาติโมกฺขํ อุทฺทิสติ. ภควโต เอตมตฺถํ อาโรเจสุํ. น ภิกฺขเว สคหฏฺฐาย ปริสาย ปาติโมกฺขํ อุทฺทิสิตพฺพํ, โย อุทฺทิเสยฺย, อาปตฺติ ทุกฺกฏสฺสาติ.}

\prose{เตน โข ปน สมเยน ฉพฺพคฺคิยา ภิกฺขู สํฆมชฺเฌ อนชฺฌิฏฺฐา ปาติโมกฺขํ อุทฺทิสนฺติ. ภควโต เอตมตฺถํ อาโรเจสุํ. น ภิกฺขเว สํฆมชฺเฌ อนชฺฌิฏฺเฐน ปาติโมกฺขํ อุทฺทิสิตพฺพํ, โย อุทฺทิเสยฺย, อาปตฺติ ทุกฺกฏสฺส. อนุชานามิ ภิกฺขเว เถราธิกํ\csromansharedfootnote{0e4b8cb81b73907d}{เถราเธยฺยํ (อฏฺฐกถายํ ปาฐนฺตรํ)} ปาติโมกฺขนฺติ.}

\setlength{\csromangathapagewidth}{0pt}
\setlength{\csromangathawakwidth}{0pt}
\csromangathameasuregroup{}{อญฺญติตฺถิยภาณวาโร นิฏฺฐิโต ปฐโม\textsuperscript{1}.}
\csromangathameasurewak{อญฺญติตฺถิยภาณวาโร นิฏฺฐิโต ปฐโม\textsuperscript{1}.}
\setlength{\csromangathaleftcol}{0pt}
\csromangathagroup{\csromangathastanza{\csromangathawak{อญฺญติตฺถิย\-ภาณวาโร นิฏฺฐิโต ปฐโม\csromansharedfootnote{c9fe24a2ca2c6370}{เอกาทสโม (ก)}.}}}

\csromanmatikaanchor{mtk.85}
\csromantocmarkref{section}{83. ปาติโมกฺขุทฺเทสกอชฺเฌสนาทิ}{mtk.85}
\bookmark[dest={mtk.85},level=1]{83. ปาติโมกฺขุทฺเทสกอชฺเฌสนาทิ}
\csromanheader{1}{83. ปาติโมกฺขุ\-ทฺเทสก\-อชฺเฌสนาทิ}
\proseitem{155}{\csromanfolio{156}อถ โข ภควา ราชคเห ยถาภิรนฺตํ วิหริตฺวา เยน โจทนาวตฺถุ เตน จาริกํ ปกฺกามิ, อนุปุพฺเพน จาริกํ จรมาโน เยน โจทนาวตฺถุ ตทวสริ.}

\prose{เตน โข ปน สมเยน อญฺญตรสฺมิํ อาวาเส สมฺพหุลา ภิกฺขู วิหรนฺติ, ตตฺถ เถโร ภิกฺขุ พาโล โหติ อพฺยตฺโต, โส น ชานาติ อุโปสถํ วา อุโปสถกมฺมํ วา ปาติโมกฺขํ วา ปาติโมกฺขุ\-ทฺเทสํ วา.}

\prose{อถ โข เตสํ ภิกฺขูนํ เอตทโหสิ “ภควตา ปญฺญตฺตํ ‘เถราธิกํ ปาติโมกฺขนฺ’ติ, อยญฺจ อมฺหากํ เถโร พาโล อพฺยตฺโต, น ชานาติ อุโปสถํ วา อุโปสถกมฺมํ วา ปาติโมกฺขํ วา ปาติโมกฺขุ\-ทฺเทสํ วา, กถํ นุ โข อเมฺหหิ ปฏิปชฺชิตพฺพนฺ”ติ.}

\prose{ภควโต เอตมตฺถํ อาโรเจสุํ. อนุชานามิ ภิกฺขเว โย ตตฺถ ภิกฺขุ พฺยตฺโต ปฏิพโล, ตสฺสาเธยฺยํ ปาติโมกฺขนฺติ.}

\prose{เตน โข ปน สมเยน อญฺญตรสฺมิํ อาวาเส ตทหุโปสเถ สมฺพหุลา ภิกฺขู วิหรนฺติ พาลา อพฺยตฺตา, เต น ชานนฺติ อุโปสถํ วา อุโปสถกมฺมํ วา ปาติโมกฺขํ วา ปาติโมกฺขุ\-ทฺเทสํ วา, เต เถรํ อชฺเฌสิํสุ “อุทฺทิสตุ ภนฺเต เถโร ปาติโมกฺขนฺ”ติ.}

\prose{โส เอวมาห “น เม อาวุโส วตฺตตี”ติ. ทุติยํ เถรํ อชฺเฌสิํสุ “อุทฺทิสตุ ภนฺเต เถโร ปาติโมกฺขนฺ”ติ. โสปิ เอวมาห “น เม อาวุโส วตฺตตี”ติ. ตติยํ เถรํ อชฺเฌสิํสุ “อุทฺทิสตุ ภนฺเต เถโร ปาติโมกฺขนฺ”ติ. โสปิ เอวมาห “น เม อาวุโส วตฺตตี”ติ. เอเตเนว อุปาเยน ยาว สํฆนวกํ อชฺเฌสิํสุ “อุทฺทิสตุ อายสฺมา ปาติโมกฺขนฺ”ติ. โสปิ เอวมาห “น เม ภนฺเต วตฺตตี”ติ. ภควโต เอตมตฺถํ อาโรเจสุํ.}

\prose{อิธ ปน ภิกฺขเว อญฺญตรสฺมิํ อาวาเส ตทหุโปสเถ สมฺพหุลา ภิกฺขู วิหรนฺติ พาลา อพฺยตฺตา, เต น ชานนฺติ อุโปสถํ วา อุโปสถกมฺมํ วา ปาติโมกฺขํ วา ปาติโมกฺขุ\-ทฺเทสํ วา, เต เถรํ อชฺเฌสนฺติ “อุทฺทิสตุ ภนฺเต เถโร ปาติโมกฺขนฺ”ติ.}

\prose{โส เอวํ วเทติ “น เม อาวุโส วตฺตตี”ติ. ทุติยํ เถรํ อชฺเฌสนฺติ “อุทฺทิสตุ ภนฺเต เถโร ปาติโมกฺขนฺ”ติ. โสปิ เอวํ วเทติ “น เม อาวุโส วตฺตตี”ติ. \csromanfolio{157}ตติยํ เถรํ อชฺเฌสนฺติ “อุทฺทิสตุ ภนฺเต เถโร ปาติโมกฺขนฺ”ติ. โสปิ เอวํ วเทติ “น เม อาวุโส วตฺตตี”ติ. เอเตเนว อุปาเยน ยาว สํฆนวกํ อชฺเฌสนฺติ “อุทฺทิสตุ อายสฺมา ปาติโมกฺขนฺ”ติ. โสปิ เอวํ วเทติ “น เม ภนฺเต วตฺตตี”ติ. เตหิ ภิกฺขเว ภิกฺขูหิ เอโก ภิกฺขุ สามนฺตา อาวาสา สชฺชุกํ ปาเหตพฺโพ “คจฺฉาวุโส สํขิตฺเตน วา วิตฺถาเรน วา ปาติโมกฺขํ ปริยาปุณิตฺวาน อาคจฺฉาหี”ติ.}

\proseitem{155}{อถ โข ภิกฺขูนํ เอตทโหสิ “เกน นุ โข ปาเหตพฺโพ”ติ. ภควโต เอตมตฺถํ อาโรเจสุํ. อนุชานามิ ภิกฺขเว เถเรน ภิกฺขุนา นวํ ภิกฺขุํ อาณาเปตุนฺติ. เถเรน อาณตฺตา นวา ภิกฺขู น คจฺฉนฺติ. ภควโต เอตมตฺถํ อาโรเจสุํ. น ภิกฺขเว เถเรน อาณตฺเตน อคิลาเนน น คนฺตพฺพํ, โย น คจฺเฉยฺย, อาปตฺติ ทุกฺกฏสฺสาติ.}

\csromanmatikaanchor{mtk.86}
\csromantocmarkref{section}{84. ปกฺขคณนาทิอุคฺคหณานุชานนา}{mtk.86}
\bookmark[dest={mtk.86},level=1]{84. ปกฺขคณนาทิอุคฺคหณานุชานนา}
\csromanheader{1}{84. ปกฺข\-คณนา\-ทิ\-อุคฺคหณา\-นุชานนา}
\proseitem{156}{\textbf{อถ โข ภควา โจทนา}\-\textbf{วตฺถุสฺมิํ ยถาภิรนฺตํ วิหริตฺวา} ปุนเทว ราชคหํ ปจฺจาคญฺฉิ.}

\prose{เตน โข ปน สมเยน มนุสฺสา ภิกฺขู ปิณฺฑาย จรนฺเต ปุจฺฉนฺติ “กติมี ภนฺเต ปกฺขสฺสา”ติ. ภิกฺขู เอวมาหํสุ “น โข มยํ อาวุโส ชานามา”ติ. มนุสฺสา อุชฺฌายนฺติ ขิยฺยนฺติ วิปาเจนฺติ “ปกฺขคณนมตฺตมฺปิเม สมณา สกฺยปุตฺติยา น ชานนฺติ, กิํ ปนิเม อญฺญํ กิญฺจิ กลฺยาณํ ชานิสฺสนฺตี”ติ. ภควโต เอตมตฺถํ อาโรเจสุํ. อนุชานามิ ภิกฺขเว ปกฺขคณนํ อุคฺคเหตุนฺติ. อถ โข ภิกฺขูนํ เอตทโหสิ “เกน นุ โข ปกฺขคณนา อุคฺคเหตพฺพา”ติ. ภควโต เอตมตฺถํ อาโรเจสุํ. อนุชานามิ ภิกฺขเว สพฺเพเหว ปกฺขคณนํ อุคฺคเหตุนฺติ.}

\proseitem{157}{เตน โข ปน สมเยน มนุสฺสา ภิกฺขู ปิณฺฑาย จรนฺเต ปุจฺฉนฺติ “กีวติกา ภนฺเต ภิกฺขู”ติ. ภิกฺขู เอวมาหํสุ “น โข มยํ อาวุโส ชานามา”ติ. มนุสฺสา อุชฺฌายนฺติ ขิยฺยนฺติ วิปาเจนฺติ “อญฺญมญฺญมฺปิเม สมณา สกฺยปุตฺติยา น ชานนฺติ, กิํ ปนิเม อญฺญํ กิญฺจิ กลฺยาณํ ชานิสฺสนฺตี”ติ. ภควโต เอตมตฺถํ อาโรเจสุํ. อนุชานามิ ภิกฺขเว ภิกฺขู คเณตุนฺติ.}

\prose{\csromanfolio{158}อถ โข ภิกฺขูนํ เอตทโหสิ “กทา นุ โข ภิกฺขู คเณตพฺพา”ติ. ภควโต เอตมตฺถํ อาโรเจสุํ. อนุชานามิ ภิกฺขเว ตทหุโปสเถ นามคฺเคน\csromansharedfootnote{c72a3f2595e4a644}{นามมตฺเตน (สฺยา), คณมคฺเคน (ก)} วา คเณตุํ, สลากํ วา คาเหตุนฺติ.}

\proseitem{158}{เตน โข ปน สมเยน ภิกฺขู อชานนฺตา “อชฺชุโปสโถ”ติ ทูรํ คามํ ปิณฺฑาย จรนฺติ. เต อุทฺทิสฺสมาเนปิ ปาติโมกฺเข อาคจฺฉนฺติ, อุทฺทิฏฺฐ\-มตฺเตปิ อาคจฺฉนฺติ. ภควโต เอตมตฺถํ อาโรเจสุํ. อนุชานามิ ภิกฺขเว อาโรเจสุํ อชฺชุโปสโถติ.}

\prose{อถ โข ภิกฺขูนํ เอตทโหสิ “เกน นุ โข อาโรเจตพฺโพ”ติ. ภควโต เอตมตฺถํ อาโรเจสุํ. อนุชานามิ ภิกฺขเว เถเรน ภิกฺขุนา กาลวโต อาโรเจตุนฺติ.}

\prose{เตน โข ปน สมเยน อญฺญตโร เถโร กาลวโต นสฺสรติ. ภควโต เอตมตฺถํ อโรเจสุํ. อนุชานามิ ภิกฺขเว ภตฺตกาเลปิ อาโรเจตุนฺติ.}

\prose{ภตฺตกาเลปิ นสฺสรติ. ภควโต เอตมตฺถํ อาโรเจสุํ. อนุชานามิ ภิกฺขเว ยํ กาลํ สรติ, ตํ กาลํ อาโรเจตุนฺติ.}

\csromanmatikaanchor{mtk.87}
\csromantocmarkref{section}{85. ปุพฺพกรณานุชานนา}{mtk.87}
\bookmark[dest={mtk.87},level=1]{85. ปุพฺพกรณานุชานนา}
\csromanheader{1}{85. \textbf{ปุพฺพกรณา}\-\textbf{นุชานนา}}
\proseitem{159}{เตน โข ปน สมเยน อญฺญตรสฺมิํ อาวาเส อุโปสถาคารํ อุกฺลาปํ โหติ. อาคนฺตุกา ภิกฺขู อุชฺฌายนฺติ ขิยฺยนฺติ วิปาเจนฺติ “กถํ หิ นาม อาวาสิกา ภิกฺขู อุโปสถาคารํ น สมฺมชฺชิสฺสนฺตี”ติ. ภควโต เอตมตฺถํ อาโรเจสุํ. อนุชานามิ ภิกฺขเว อุโปสถาคารํ สมฺมชฺชิตุนฺติ.}

\prose{อถ โข ภิกฺขูนํ เอตทโหสิ “เกน นุ โข อุโปสถาคารํ สมฺมชฺชิตพฺพนฺ”ติ. ภควโต เอตมตฺถํ อาโรเจสุํ. อนุชานามิ ภิกฺขเว เถเรน ภิกฺขุนา นวํ ภิกฺขุํ อาณาเปตุนฺติ.}

\prose{เถเรน อาณตฺตา นวา ภิกฺขู น สมฺมชฺชนฺติ. ภควโต เอตมตฺถํ อาโรเจสุํ. น ภิกฺขเว เถเรน อาณตฺเตน อคิลาเนน น สมฺมชฺชิตพฺพํ, โย น สมฺมชฺเชยฺย, อาปตฺติ ทุกฺกฏสฺสาติ.}

\proseitem{160}{\csromanfolio{159}เตน โข ปน สมเยน อุโปสถาคาเร อาสนํ อปญฺญตฺตํ โหติ. ภิกฺขู ฉมายํ นิสีทนฺติ, คตฺตานิปิ จีวรานิปิ ปํสุกิตานิ โหนฺติ. ภควโต เอตมตฺถํ อาโรเจสุํ. อนุชานามิ ภิกฺขเว อุโปสถาคาเร อาสนํ ปญฺญเปตุนฺติ.}

\prose{อถ โข ภิกฺขูนํ เอตทโหสิ “เกน นุ โข อุโปสถาคาเร อาสนํ ปญฺญเปตพฺพนฺ”ติ. ภควโต เอตมตฺถํ อาโรเจสุํ. อนุชานามิ ภิกฺขเว เถเรน ภิกฺขุนา นวํ ภิกฺขุํ อาณาเปตุนฺติ.}

\prose{เถเรน อาณตฺตา นวา ภิกฺขู น ปญฺญเปนฺติ. ภควโต เอตมตฺถํ อาโรเจสุํ. น ภิกฺขเว เถเรน อาณตฺเตน อคิลาเนน น ปญฺญเปตพฺพํ, โย น ปญฺญเปยฺย, อาปตฺติ ทุกฺกฏสฺสาติ.}

\proseitem{161}{เตน โข ปน สมเยน อุโปสถาคาเร ปทีโป น โหติ, ภิกฺขู อนฺธกาเร กายมฺปิ จีวรมฺปิ อกฺกมนฺติ. ภควโต เอตมตฺถํ อาโรเจสุํ. อนุชานามิ ภิกฺขเว อุโปสถาคาเร ปทีปํ กาตุนฺติ.}

\prose{อถ โข ภิกฺขูนํ เอตทโหสิ “เกน นุ โข อุโปสถาคาเร ปทีโป กาตพฺโพ”ติ. ภควโต เอตมตฺถํ อาโรเจสุํ. อนุชานามิ ภิกฺขเว เถเรน ภิกฺขุนา นวํ ภิกฺขุํ อาณาเปตุนฺติ.}

\prose{เถเรน อาณตฺตา นวา ภิกฺขู น ปทีเปนฺติ. ภควโต เอตมตฺถํ อาโรเจสุํ. น ภิกฺขเว เถเรน อาณตฺเตน อคิลาเนน น ปทีเปตพฺโพ, โย น ปทีเปยฺย, อาปตฺติ ทุกฺกฏสฺสาติ.}

\proseitem{162}{เตน โข ปน สมเยน อญฺญตรสฺมิํ อาวาเส อาวาสิกา ภิกฺขู เนว ปานียํ อุปฏฺฐาเปนฺติ, น ปริโภชนียํ อุปฏฺฐาเปนฺติ. อาคนฺตุกา ภิกฺขู อุชฺฌายนฺติ ขิยฺยนฺติ วิปาเจนฺติ “กถํ หิ นาม อาวาสิกา ภิกฺขู เนว ปานียํ อุปฏฺฐาเปสฺสนฺติ, น ปริโภชนียํ อุปฏฺฐาเปสฺสนฺตี”ติ. ภควโต เอตมตฺถํ อาโรเจสุํ. อนุชานามิ ภิกฺขเว ปานียํ ปริโภชนียํ อุปฏฺฐาเปตุนฺติ.}

\prose{อถ โข ภิกฺขูนํ เอตทโหสิ “เกน นุ โข ปานียํ ปริโภชนียํ อุปฏฺฐาเปตพฺพนฺ”ติ. ภควโต เอตมตฺถํ อาโรเจสุํ. อนุชานามิ ภิกฺขเว เถเรน ภิกฺขุนา นวํ ภิกฺขุํ อาณาเปตุนฺติ.}

\prose{\csromanfolio{160}เถเรน อาณตฺตา นวา ภิกฺขู น อุปฏฺฐาเปนฺติ. ภควโต เอตมตฺถํ อาโรเจสุํ. น ภิกฺขเว เถเรน อาณตฺเตน อคิลาเนน น อุปฏฺฐาเป\-ตพฺพํ, โย น อุปฏฺฐาเปยฺย, อาปตฺติ ทุกฺกฏสฺสาติ.}

\csromanmatikaanchor{mtk.88}
\csromantocmarkref{section}{86. ทิสํคมิกาทิวตฺถุ}{mtk.88}
\bookmark[dest={mtk.88},level=1]{86. ทิสํคมิกาทิวตฺถุ}
\csromanheader{1}{86. \textbf{ทิสํคมิกา}\-\textbf{ทิ}\-\textbf{วตฺถุ}}
\proseitem{163}{เตน โข ปน สมเยน สมฺพหุลา ภิกฺขู พาลา อพฺยตฺตา ทิสํคมิกา อาจริยุ\-ปชฺฌาเย อาปุจฺฉิํสุ\csromansharedfootnote{c495272ad2034724}{น อาปุจฺฉิํสุ (ก)}. ภควโต เอตมตฺถํ อาโรเจสุํ.}

\prose{อิธ ปน ภิกฺขเว สมฺพหุลา ภิกฺขู พาลา อพฺยตฺตา ทิสํคมิกา อาจริยุ\-ปชฺฌาเย อาปุจฺฉนฺติ\csromansharedfootnote{376496cb09a7f634}{น อาปุจฺฉนฺติ (ก)}. เต\csromansharedfootnote{055d6961c69e588d}{เตหิ (ก)} ภิกฺขเว อาจริยุ\-ปชฺฌาเยหิ ปุจฺฉิตพฺพา “กหํ คมิสฺสถ เกน สทฺธิํ คมิสฺสถา”ติ. เต เจ ภิกฺขเว พาลา อพฺยตฺตา อญฺเญ พาเล อพฺยตฺเต อปทิเสยฺยุํ. น ภิกฺขเว อาจริยุ\-ปชฺฌาเยหิ อนุชานิตพฺพา, อนุชาเนยฺยุํ เจ, อาปตฺติ ทุกฺกฏสฺส. เต จ ภิกฺขเว พาลา อพฺยตฺตา อนนุญฺญาตา อาจริยุ\-ปชฺฌาเยหิ คจฺเฉยฺยุํ เจ, อาปตฺติ ทุกฺกฏสฺส.}

\prose{อิธ ปน ภิกฺขเว อญฺญตรสฺมิํ อาวาเส สมฺพหุลา ภิกฺขู วิหรนฺติ พาลา อพฺยตฺตา, เต น ชานนฺติ อุโปสถํ วา อุโปสถกมฺมํ วา ปาติโมกฺขํ วา ปาติโมกฺขุ\-ทฺเทสํ วา.}

\prose{ตตฺถ อญฺโญ ภิกฺขุ อาคจฺฉติ พหุสฺสุโต อาคตาคโม ธมฺมธโร วินยธโร มาติกาธโร ปณฺฑิโต พฺยตฺโต เมธาวี ลชฺชี กุกฺกุจฺจโก สิกฺขากาโม. เตหิ ภิกฺขเว ภิกฺขูหิ โส ภิกฺขุ สงฺคเหตพฺโพ อนุคฺคเหตพฺโพ อุปลาเปตพฺโพ อุปฏฺฐาเปตพฺโพ จุณฺเณน มตฺติกาย ทนฺตกฏฺเฐน มุโขทเกน, โน เจ สงฺคเณฺหยฺยุํ อนุคฺคเณฺหยฺยุํ อุปลาเปยฺยุํ อุปฏฺฐาเปยฺยุํ จุณฺเณน มตฺติกาย ทนฺตกฏฺเฐน มุโขทเกน, อาปตฺติ ทุกฺกฏสฺส.}

\proseitem{163}{อิธ ปน ภิกฺขเว อญฺญตรสฺมิํ อาวาเส ตทหุโปสเถ สมฺพหุลา ภิกฺขู วิหรนฺติ พาลา อพฺยตฺตา, เต น ชานนฺติ อุโปสถํ วา อุโปสถกมฺมํ วา ปาติโมกฺขํ วา ปาติโมกฺขุ\-ทฺเทสํ วา. เตหิ ภิกฺขเว ภิกฺขูหิ เอโก ภิกฺขุ สามนฺตา อาวาสา สชฺชุกํ ปาเหตพฺโพ “คจฺฉาวุโส สํขิตฺเตน วา วิตฺถาเรน วา ปาติโมกฺขํ ปริยาปุณิตฺวา อาคจฺฉา”ติ. เอวญฺเจตํ ลเภถ, อิจฺเจตํ กุสลํ. โน เจ ลเภถ, เตหิ ภิกฺขเว \csromanfolio{161}ภิกฺขูหิ สพฺเพเหว ยตฺถ ชานนฺติ อุโปสถํ วา อุโปสถกมฺมํ วา ปาติโมกฺขํ วา ปาติโมกฺขุ\-ทฺเทสํ วา, โส อาวาโส คนฺตพฺโพ, โน เจ คจฺเฉยฺยุํ, อาปตฺติ ทุกฺกฏสฺส.}

\prose{อิธ ปน ภิกฺขเว อญฺญตรสฺมิํ อาวาเส สมฺพหุลา ภิกฺขู วสฺสํ วสนฺติ พาลา อพฺยตฺตา, เต น ชานนฺติ อุโปสถํ วา อุโปสถกมฺมํ วา ปาติโมกฺขํ วา ปาติโมกฺขุ\-ทฺเทสํ วา. เตหิ ภิกฺขเว ภิกฺขูหิ เอโก ภิกฺขุ สามนฺตา อาวาสา สชฺชุกํ ปาเหตพฺโพ “คจฺฉาวุโส สํขิตฺเตน วา วิตฺถาเรน วา ปาติโมกฺขํ ปริยาปุณิตฺวา อาคจฺฉา”ติ. เอวญฺเจตํ ลเภถ, อิจฺเจตํ กุสลํ. โน เจ ลเภถ, เอโก ภิกฺขุ สตฺถาหกาลิกํ ปาเหตพฺโพ “คจฺฉาวุโส สํขิตฺเตน วา วิตฺถาเรน วา ปาติโมกฺขํ ปริยาปุณิตฺวา อาคจฺฉา”ติ. เอวญฺเจตํ ลเภถ, อิจฺเจตํ กุสลํ. โน เจ ลเภถ, น ภิกฺขเว เตหิ ภิกฺขูหิ ตสฺมิํ อาวาเส วสฺสํ วสิตพฺพํ, วเสยฺยุํ เจ, อาปตฺติ ทุกฺกฏสฺสาติ.}

\csromanmatikaanchor{mtk.89}
\csromantocmarkref{section}{87. ปาริสุทฺธิทานกถา}{mtk.89}
\bookmark[dest={mtk.89},level=1]{87. ปาริสุทฺธิทานกถา}
\csromanheader{1}{87. \textbf{ปาริสุทฺธิ}\-\textbf{ทาน}\-\textbf{กถา}}
\proseitem{164}{อถ โข ภควา ภิกฺขู อามนฺเตสิ “สนฺนิปตถ ภิกฺขเว สํโฆ อุโปสถํ กริสฺสตี”ติ. เอวํ วุตฺเต อญฺญตโร ภิกฺขุ ภควนฺตํ เอตทโวจ “อตฺถิ ภนฺเต ภิกฺขุ คิลาโน, โส อนาคโต”ติ. อนุชานามิ ภิกฺขเว คิลาเนน ภิกฺขุนา ปาริสุทฺธิํ ทาตุํ. เอวญฺจ ปน ภิกฺขเว ทาตพฺพา. เตน คิลาเนน ภิกฺขุนา เอกํ ภิกฺขุํ อุปสงฺกมิตฺวา เอกํสํ อุตฺตราสงฺคํ กริตฺวา อุกฺกุฏิกํ นิสีทิตฺวา อญฺชลิํ ปคฺคเหตฺวา เอวมสฺส วจนีโย “ปาริสุทฺธิํ ทมฺมิ, ปาริสุทฺธิํ เม หร, ปาริสุทฺธิํ เม อาโรเจหี”ติ. กาเยน วิญฺญาเปติ, วาจาย วิญฺญาเปติ, กาเยน วาจาย วิญฺญาเปติ, ทินฺนา โหติ ปาริสุทฺธิ. น กาเยน วิญฺญาเปติ, น วาจาย วิญฺญาเปติ, น กาเยน วาจาย วิญฺญาเปติ, น ทินฺนา โหติ ปาริสุทฺธิ. เอวญฺเจตํ ลเภถ, อิจฺเจตํ กุสลํ. โน เจ ลเภถ, โส ภิกฺขเว คิลาโน ภิกฺขุ มญฺเจน วา ปีเฐน วา สํฆมชฺเฌ อาเนตฺวา อุโปสโถ กาตพฺโพ. สเจ ภิกฺขเว คิลานุ\-ปฏฺฐากานํ ภิกฺขูนํ เอวํ โหติ “สเจ โข มยํ คิลานํ ฐานา จาเวสฺสาม, อาพาโธ วา อภิวฑฺฒิสฺสติ, กาลํกิริยา วา ภวิสฺสตี”ติ. น ภิกฺขเว คิลาโน ภิกฺขุ ฐานา จาเวตพฺโพ, สํเฆน ตตฺถ คนฺตฺวา \csromanfolio{162}อุโปสโถ กาตพฺโพ, น เตฺวว วคฺเคน สํเฆน อุโปสโถ กาตพฺโพ, กเรยฺย เจ, อาปตฺติ ทุกฺกฏสฺส.}

\prose{ปาริสุทฺธิ\-หารโก เจ ภิกฺขเว ทินฺนาย ปาริสุทฺธิยา ตตฺเถว ปกฺกมติ, อญฺญสฺส ทาตพฺพา ปาริสุทฺธิ.}

\prose{ปาริสุทฺธิ\-หารโก เจ ภิกฺขเว ทินฺนาย ปาริสุทฺธิยา ตตฺเถว วิพฺภมติ ฯเปฯ กาลํ กโรติ. สามเณโร ปฏิชานาติ. สิกฺขํ ปจฺจกฺขาตโก ปฏิชานาติ. อนฺติมวตฺถุํ อชฺฌาปนฺนโก ปฏิชานาติ. อุมฺมตฺตโก ปฏิชานาติ. ขิตฺตจิตฺโต ปฏิชานาติ. เวทนาฏฺโฏ ปฏิชานาติ. อาปตฺติยา อทสฺสเน อุกฺขิตฺตโก ปฏิชานาติ. อาปตฺติยา อปฺปฏิกมฺเม อุกฺขิตฺตโก ปฏิชานาติ. ปาปิกาย ทิฏฺฐิยา อปฺปฏินิสฺสคฺเค อุกฺขิตฺตโก ปฏิชานาติ. ปณฺฑโก ปฏิชานาติ. เถยฺย\-สํวาสโก ปฏิชานาติ. ติตฺถิย\-ปกฺกนฺตโก ปฏิชานาติ. ติรจฺฉานคโต ปฏิชานาติ. มาตุฆาตโก ปฏิชานาติ. ปิตุฆาตโก ปฏิชานาติ. อรหนฺต\-ฆาตโก ปฏิชานาติ. ภิกฺขุนิ\-ทูสโก ปฏิชานาติ. สํฆเภทโก ปฏิชานาติ. โลหิตุปฺปาทโก ปฏิชานาติ. อุภโต\-พฺยญฺชนโก ปฏิชานาติ, อญฺญสฺส ทาตพฺพา ปาริสุทฺธิ.}

\prose{ปาริสุทฺธิ\-หารโก เจ ภิกฺขเว ทินฺนาย ปาริสุทฺธิยา อนฺตรามคฺเค ปกฺกมติ, อนาหฏา โหติ ปาริสุทฺธิ.}

\prose{ปาริสุทฺธิ\-หารโก เจ ภิกฺขเว ทินฺนาย ปาริสุทฺธิยา อนฺตรามคฺเค วิพฺภมติ ฯเปฯ. อุภโต\-พฺยญฺชนโก ปฏิชานาติ, อนาหฏา โหติ ปาริสุทฺธิ.}

\prose{ปาริสุทฺธิ\-หารโก เจ ภิกฺขเว ทินฺนาย ปาริสุทฺธิยา สํฆปฺปตฺโต ปกฺกมติ, อาหฏา โหติ ปาริสุทฺธิ.}

\prose{ปาริสุทฺธิ\-หารโก เจ ภิกฺขเว ทินฺนาย ปาริสุทฺธิยา สํฆปฺปตฺโต วิพฺภมติ ฯเปฯ. อุภโต\-พฺยญฺชนโก ปฏิชานาติ, อาหฏา โหติ ปาริสุทฺธิ.}

\prose{ปาริสุทฺธิ\-หารโก เจ ภิกฺขเว ทินฺนาย ปาริสุทฺธิยา สํฆปฺปตฺโต สุตฺโต น อาโรเจติ. ปมตฺโต น อาโรเจติ. สมาปนฺโน น อาโรเจติ, อาหฏา โหติ ปาริสุทฺธิ, ปาริสุทฺธิ\-หารกสฺส อนาปตฺติ.}

\prose{ปาริสุทฺธิ\-หารโก เจ ภิกฺขเว ทินฺนาย ปาริสุทฺธิยา สํฆปฺปตฺโต สญฺจิจฺจ น อาโรเจติ, อาหฏา โหติ ปาริสุทฺธิ, ปาริสุทฺธิ\-หารกสฺส อาปตฺติ ทุกฺกฏสฺสาติ.}

\csromanmatikaanchor{mtk.90}
\csromantocmarkref{section}{88. ฉนฺททานกถา}{mtk.90}
\bookmark[dest={mtk.90},level=1]{88. ฉนฺททานกถา}
\csromanheader{1}{88. \textbf{ฉนฺททาน}\-\textbf{กถา}}
\proseitem{165}{\csromanfolio{163}อถ โข ภควา ภิกฺขู อามนฺเตสิ “สนฺนิปตถ ภิกฺขเว สํโฆ กมฺมํ กริสฺสตี”ติ. เอวํ วุตฺเต อญฺญตโร ภิกฺขุ ภควนฺตํ เอตทโวจ “อตฺถิ ภนฺเต ภิกฺขุ คิลาโน, โส อนาคโต”ติ. อนุชานามิ ภิกฺขเว คิลาเนน ภิกฺขุนา ฉนฺทํ ทาตุํ. เอวญฺจ ปน ภิกฺขเว ทาตพฺโพ, เตน คิลาเนน ภิกฺขุนา เอกํ ภิกฺขุํ อุปสงฺกมิตฺวา เอกํสํ อุตฺตราสงฺคํ กริตฺวา อุกฺกุฏิกํ นิสีทิตฺวา อญฺชลิํ ปคฺคเหตฺวา เอวมสฺส วจนีโย “ฉนฺทํ ทมฺมิ, ฉนฺทํ เม หร, ฉนฺทํ เม อาโรเจหี”ติ. กาเยน วิญฺญาเปติ, วาจาย วิญฺญาเปติ, กาเยน วาจาย วิญฺญาเปติ, ทินฺโน โหติ ฉนฺโท. น กาเยน วิญฺญาเปติ, น วาจาย วิญฺญาเปติ, น กาเยน วาจาย วิญฺญาเปติ, น ทินฺโน โหติ ฉนฺโท. เอวญฺเจตํ ลเภถ, อิจฺเจตํ กุสลํ. โน เจ ลเภถ, โส ภิกฺขเว คิลาโน ภิกฺขุ มญฺเจน วา ปีเฐน วา สํฆมชฺเฌ อาเนตฺวา กมฺมํ กาตพฺพํ. สเจ ภิกฺขเว คิลานุ\-ปฏฺฐากานํ ภิกฺขูนํ เอวํ โหติ “สเจ โข มยํ คิลานํ ฐานา จาเวสฺสาม, อาพาโธ วา อภิวฑฺฒิสฺสติ, กาลํกิริยา วา ภวิสฺสตี”ติ. น ภิกฺขเว คิลาโน ภิกฺขุ ฐานา จาเวตพฺโพ, สํเฆน ตตฺถ คนฺตฺวา กมฺมํ กาตพฺพํ, น เตฺวว วคฺเคน สํเฆน กมฺมํ กาตพฺพํ, กเรยฺย เจ, อาปตฺติ ทุกฺกฏสฺส.}

\prose{ฉนฺทหารโก เจ ภิกฺขเว ทินฺเน ฉนฺเท ตตฺเถว ปกฺกมติ, อญฺญสฺส ทาตพฺโพ ฉนฺโท.}

\prose{ฉนฺทหารโก เจ ภิกฺขเว ทินฺเน ฉนฺเท ตตฺเถว วิพฺภมติ ฯเปฯ กาลํ กโรติ. สามเณโร ปฏิชานาติ. สิกฺขํ ปจฺจกฺขาตโก ปฏิชานาติ. อนฺติมวตฺถุํ อชฺฌาปนฺนโก ปฏิชานาติ. อุมฺมตฺตโก ปฏิชานาติ. ขิตฺตจิตฺโต ปฏิชานาติ. เวทนาฏฺโฏ ปฏิชานาติ. อาปตฺติยา อทสฺสเน อุกฺขิตฺตโก ปฏิชานาติ. อาปตฺติยา อปฺปฏิกมฺเม อุกฺขิตฺตโก ปฏิชานาติ. ปาปิกาย ทิฏฺฐิยา อปฺปฏินิสฺสคฺเค อุกฺขิตฺตโก ปฏิชานาติ. ปณฺฑโก ปฏิชานาติ. เถยฺย\-สํวาสโก ปฏิชานาติ. ติตฺถิย\-ปกฺกนฺตโก ปฏิชานาติ. ติรจฺฉานคโต ปฏิชานาติ. มาตุฆาตโก ปฏิชานาติ. ปิตุฆาตโก ปฏิชานาติ. อรหนฺต\-ฆาตโก ปฏิชานาติ. ภิกฺขุนิ\-ทูสโก ปฏิชานาติ. สํฆเภทโก ปฏิชานาติ. โลหิตุปฺปาทโก ปฏิชานาติ. อุภโต\-พฺยญฺชนโก ปฏิชานาติ, อญฺญสฺส ทาตพฺโพ ฉนฺโท.}

\prose{\csromanfolio{164}ฉนฺทหารโก เจ ภิกฺขเว ทินฺเน ฉนฺเท อนฺตรามคฺเค ปกฺกมติ, อนาหโฏ โหติ ฉนฺโท.}

\prose{ฉนฺทหารโก เจ ภิกฺขเว ทินฺเน ฉนฺเท อนฺตรามคฺเค วิพฺภมติ ฯเปฯ. อุภโต\-พฺยญฺชนโก ปฏิชานาติ, อนาหโฏ โหติ ฉนฺโท.}

\prose{ฉนฺทหารโก เจ ภิกฺขเว ทินฺเน ฉนฺเท สํฆปฺปตฺโต ปกฺกมติ, อาหโฏ โหติ ฉนฺโท.}

\prose{ฉนฺทหารโก เจ ภิกฺขเว ทินฺเน ฉนฺเท สํฆปฺปตฺโต วิพฺภมติ ฯเปฯ. อุภโต\-พฺยญฺชนโก ปฏิชานาติ, อาหโฏ โหติ ฉนฺโท.}

\prose{ฉนฺทหารโก เจ ภิกฺขเว ทินฺเน ฉนฺเท สํฆปฺปตฺโต สุตฺโต น อาโรเจติ. ปมตฺโต น อาโรเจติ. สมาปนฺโน น อาโรเจติ, อาหโฏ โหติ ฉนฺโท, ฉนฺท\-หารกสฺส อนาปตฺติ.}

\prose{ฉนฺทหารโก เจ ภิกฺขเว ทินฺเน ฉนฺเท สํฆปฺปตฺโต สญฺจิจฺจ น อาโรเจติ, อาหโฏ โหติ ฉนฺโท, ฉนฺท\-หารกสฺส อาปตฺติ ทุกฺกฏสฺส.}

\prose{อนุชานามิ ภิกฺขเว ตทหุโปสเถ ปาริสุทฺธิํ เทนฺเตน ฉนฺทมฺปิ ทาตุํ, สนฺติ สํฆสฺส กรณียนฺติ.}

\csromanmatikaanchor{mtk.91}
\csromantocmarkref{section}{89. ญาตกาทิคฺคหณกถา}{mtk.91}
\bookmark[dest={mtk.91},level=1]{89. ญาตกาทิคฺคหณกถา}
\csromanheader{1}{89. \textbf{ญาตกา}\-\textbf{ทิ}\-\textbf{คฺคหณ}\-\textbf{กถา}}
\proseitem{166}{เตน โข ปน สมเยน อญฺญตรํ ภิกฺขุํ ตทหุโปสเถ ญาตกา คณฺหิํสุ. ภควโต เอตมตฺถํ อาโรเจสุํ.}

\prose{อิธ ปน ภิกฺขเว ภิกฺขุํ ตทหุโปสเถ ญาตกา คณฺหนฺติ, เต ญาตกา ภิกฺขูหิ เอวมสฺสุ วจนียา “อิงฺฆ ตุเมฺห อายสฺมนฺโต อิมํ ภิกฺขุํ มุหุตฺตํ มุญฺจถ, ยาวายํ ภิกฺขุ อุโปสถํ กโรตี”ติ.}

\setlength{\csromangathapagewidth}{0pt}
\setlength{\csromangathawakwidth}{0pt}
\csromangathameasuregroup{เอวญฺเจตํ ลเภถ,}{\csromangathabat{เอวญฺเจตํ ลเภถ,}{อิจฺเจตํ กุสลํ.}}
\csromangathasetleft{เอวญฺเจตํ ลเภถ,}
\csromangathagroup{\csromangathastanza{\csromangathaitembat{166}{เอวญฺเจตํ ลเภถ,}{อิจฺเจตํ กุสลํ.}}}

\prose{โน เจ ลเภถ, เต ญาตกา ภิกฺขูหิ เอวมสฺสุ วจนียา “อิงฺฆ ตุเมฺห อายสฺมนฺโต มุหุตฺตํ เอกมนฺตํ โหถ, ยาวายํ ภิกฺขุ ปาริสุทฺธิํ เทตี”ติ.}

\setlength{\csromangathapagewidth}{0pt}
\setlength{\csromangathawakwidth}{0pt}
\csromangathameasuregroup{เอวญฺเจตํ ลเภถ,}{\csromangathabat{เอวญฺเจตํ ลเภถ,}{อิจฺเจตํ กุสลํ.}}
\csromangathasetleft{เอวญฺเจตํ ลเภถ,}
\csromangathagroup{\csromangathastanza{\csromangathaitembat{166}{เอวญฺเจตํ ลเภถ,}{อิจฺเจตํ กุสลํ.}}}

\prose{โน เจ ลเภถ, เต ญาตกา ภิกฺขูหิ เอวมสฺสุ วจนียา “อิงฺฆ ตุเมฺห อายสฺมนฺโต อิมํ ภิกฺขุํ มุหุตฺตํ นิสฺสีมํ เนถ, ยาว สํโฆ อุโปสถํ กโรตี”ติ.}

\setlength{\csromangathapagewidth}{0pt}
\setlength{\csromangathawakwidth}{0pt}
\csromangathameasuregroup{เอวญฺเจตํ ลเภถ,}{\csromangathabat{เอวญฺเจตํ ลเภถ,}{อิจฺเจตํ กุสลํ.}}
\csromangathasetleft{เอวญฺเจตํ ลเภถ,}
\csromangathagroup{\csromangathastanza{\csromangathaitembat{166}{เอวญฺเจตํ ลเภถ,}{อิจฺเจตํ กุสลํ.}}}

\prose{โน เจ ลเภถ, น เตฺวว วคฺเคน สํเฆน อุโปสโถ กาตพฺโพ, กเรยฺย เจ, อาปตฺติ ทุกฺกฏสฺส.}

\prose{\csromanfolio{165}อิธ ปน ภิกฺขเว ภิกฺขุํ ตทหุโปสเถ ราชาโน คณฺหนฺติ ฯเปฯ โจรา คณฺหนฺติ. ธุตฺตา คณฺหนฺติ. ภิกฺขุ\-ปจฺจตฺถิกา คณฺหนฺติ, เต ภิกฺขุ\-ปจฺจตฺถิกา ภิกฺขูหิ เอวมสฺสุ วจนียา “อิงฺฆ ตุเมฺห อายสฺมนฺโต อิมํ ภิกฺขุํ มุหุตฺตํ มุญฺจถ, ยาวายํ ภิกฺขุ อุโปสถํ กโรตี”ติ. เอวญฺเจตํ ลเภถ, อิจฺเจตํ กุสลํ. โน เจ ลเภถ, เต ภิกฺขุ\-ปจฺจตฺถิกา ภิกฺขูหิ เอวมสฺสุ วจนียา “อิงฺฆ ตุเมฺห อายสฺมนฺโต มุหุตฺตํ เอกมนฺตํ โหถ, ยาวายํ ภิกฺขุ ปาริสุทฺธิํ เทตี”ติ. เอวญฺเจตํ ลเภถ, อิจฺเจตํ กุสลํ. โน เจ ลเภถ, เต ภิกฺขุ\-ปจฺจตฺถิกา ภิกฺขูหิ เอวมสฺสุ วจนียา “อิงฺฆ ตุเมฺห อายสฺมนฺโต อิมํ ภิกฺขุํ มุหุตฺตํ นิสฺสีมํ เนถ, ยาว สํโฆ อุโปสถํ กโรตี”ติ. เอวญฺเจตํ ลเภถ, อิจฺเจตํ กุสลํ. โน เจ ลเภถ, น เตฺวว วคฺเคน สํเฆน อุโปสโถ กาตพฺโพ, กเรยฺย เจ, อาปตฺติ ทุกฺกฏสฺสาติ.}

\csromanmatikaanchor{mtk.92}
\csromantocmarkref{section}{90. อุมฺมตฺตกสมฺมุติ}{mtk.92}
\bookmark[dest={mtk.92},level=1]{90. อุมฺมตฺตกสมฺมุติ}
\csromanheader{1}{90. \textbf{อุมฺมตฺตก}\-\textbf{สมฺมุติ}}
\proseitem{167}{อถ โข ภควา ภิกฺขู อามนฺเตสิ “สนฺนิปตถ ภิกฺขเว อตฺถิ สํฆสฺส กรณียนฺ”ติ. เอวํ วุตฺเต อญฺญตโร ภิกฺขุ ภควนฺตํ เอตทโวจ “อตฺถิ ภนฺเต คคฺโค นาม ภิกฺขุ อุมฺมตฺตโก, โส อนาคโต”ติ.}

\prose{เทฺวเม ภิกฺขเว อุมฺมตฺตกา. อตฺถิ ภิกฺขเว ภิกฺขุ อุมฺมตฺตโก สรติปิ อุโปสถํ นปิ สรติ, สรติปิ สํฆกมฺมํ นปิ สรติ, อตฺถิ เนว สรติ. อาคจฺฉติปิ อุโปสถํ นปิ อาคจฺฉติ, อาคจฺฉติปิ สํฆกมฺมํ นปิ อาคจฺฉติ, อตฺถิ เนว อาคจฺฉติ. ตตฺร ภิกฺขเว ยฺวายํ อุมฺมตฺตโก สรติปิ อุโปสถํ นปิ สรติ, สรติปิ สํฆกมฺมํ นปิ สรติ, อาคจฺฉติปิ อุโปสถํ นปิ อาคจฺฉติ, อาคจฺฉติปิ สํฆกมฺมํ นปิ อาคจฺฉติ, อนุชานามิ ภิกฺขเว เอวรูปสฺส อุมฺมตฺตกสฺส อุมฺมตฺตก\-สมฺมุติํ ทาตุํ. เอวญฺจ ปน ภิกฺขเว ทาตพฺพา, พฺยตฺเตน ภิกฺขุนา ปฏิพเลน สํโฆ ญาเปตพฺโพ\csromandash{}}

\prose{“สุณาตุ เม ภนฺเต สํโฆ, คคฺโค ภิกฺขุ อุมฺมตฺตโก สรติปิ อุโปสถํ นปิ สรติ, สรติปิ สํฆกมฺมํ นปิ สรติ, อาคจฺฉติปิ อุโปสถํ นปิ อาคจฺฉติ, อาคจฺฉติปิ สํฆกมฺมํ นปิ อาคจฺฉติ, ยทิ สํฆสฺส ปตฺตกลฺลํ, สํโฆ คคฺคสฺส ภิกฺขุโน อุมฺมตฺตกสฺส \csromanfolio{166}อุมฺมตฺตก\-สมฺมุติํ ทเทยฺย, สเรยฺย วา คคฺโค ภิกฺขุ อุโปสถํ น วา สเรยฺย, สเรยฺย วา สํฆกมฺมํ น วา สเรยฺย, อาคจฺเฉยฺย วา อุโปสถํ น วา อาคจฺเฉยฺย, อาคจฺเฉยฺย วา สํฆกมฺมํ น วา อาคจฺเฉยฺย. สํโฆ สห วา คคฺเคน วินา วา คคฺเคน อุโปสถํ กเรยฺย, สํฆกมฺมํ กเรยฺย, เอสา ญตฺติ.}

\prose{สุณาตุ เม ภนฺเต สํโฆ, คคฺโค ภิกฺขุ อุมฺมตฺตโก สรติปิ อุโปถํ นปิ สรติ, สรติปิ สํฆกมฺมํ นปิ สรติ, อาคจฺฉติปิ อุโปสถํ นปิ อาคจฺฉติ, อาคจฺฉติปิ สํฆกมฺมํ นปิ อาคจฺฉติ.}

\prose{สํโฆ คคฺคสฺส ภิกฺขุโน อุมฺมตฺตกสฺส อุมฺมตฺตก\-สมฺมุติํ เทติ, สเรยฺย วา คคฺโค ภิกฺขุ อุโปสถํ น วา สเรยฺย, สเรยฺย วา สํฆกมฺมํ น วา สเรยฺย, อาคจฺเฉยฺย วา อุโปสถํ น วา อาคจฺเฉยฺย, อาคจฺเฉยฺย วา สํฆกมฺมํ น วา อาคจฺเฉยฺย, สํโฆ สห วา คคฺเคน วินา วา คคฺเคน อุโปสถํ กริสฺสติ, สํฆกมฺมํ กริสฺสติ.}

\prose{ยสฺสายสฺมโต ขมติ คคฺคสฺส ภิกฺขุโน อุมฺมตฺตกสฺส อุมฺมตฺตก\-สมฺมุติยา ทานํ, สเรยฺย วา คคฺโค ภิกฺขุ อุโปสถํ น วา สเรยฺย, สเรยฺย วา สํฆกมฺมํ น วา สเรยฺย, อาคจฺเฉยฺย วา อุโปสถํ น วา อาคจฺเฉยฺย, อาคจฺเฉยฺย วา สํฆกมฺมํ น วา อาคจฺเฉยฺย, สํโฆ สห วา คคฺเคน วินา วา คคฺเคน อุโปสถํ กริสฺสติ, สํฆกมฺมํ กริสฺสติ, โส ตุณฺหสฺส, ยสฺส นกฺขมติ, โส ภาเสยฺย.}

\prose{ทินฺนา สํเฆน คคฺคสฺส ภิกฺขุโน อุมฺมตฺตกสฺส อุมฺมตฺตก\-สมฺมุติ, สเรยฺย วา คคฺโค ภิกฺขุ อุโปสถํ น วา สเรยฺย, สเรยฺย วา สํฆกมฺมํ น วา สเรยฺย, อาคจฺเฉยฺย วา อุโปสถํ น วา อาคจฺเฉยฺย, อาคจฺเฉยฺย วา สํฆกมฺมํ น วา อาคจฺเฉยฺย, สํโฆ สห วา คคฺเคน วินา วา คคฺเคน อุโปสถํ กริสฺสติ, สํฆกมฺมํ กริสฺสติ, ขมติ สํฆสฺส, ตสฺมา ตุณฺหี, เอวเมตํ ธารยามี”ติ.}

\csromanmatikaanchor{mtk.93}
\csromantocmarkref{section}{91. สํฆุโปสถาทิปฺปเภท}{mtk.93}
\bookmark[dest={mtk.93},level=1]{91. สํฆุโปสถาทิปฺปเภท}
\csromanheader{1}{91. \textbf{สํฆุโปสถาทิปฺปเภท}}
\proseitem{168}{เตน โข ปน สมเยน อญฺญตรสฺมิํ อาวาเส ตทหุโปสเถ จตฺตาโร ภิกฺขู วิหรนฺติ. อถ โข เตสํ ภิกฺขูนํ เอตทโหสิ “ภควตา ปญฺญตฺตํ ‘อุโปสโถ กาตพฺโพ’ติ, มยญฺจมฺหา จตฺตาโร ชนา, กถํ นุ โข อเมฺหหิ อุโปสโถ กาตพฺโพ”ติ. ภควโต \csromanfolio{167}เอตมตฺถํ อาโรเจสุํ. อนุชานามิ ภิกฺขเว จตุนฺนํ ปาติโมกฺขํ อุทฺทิสิตุนฺติ.}

\prose{เตน โข ปน สมเยน อญฺญตรสฺมิํ อาวาเส ตทหุโปสเถ ตโย ภิกฺขู วิหรนฺติ. อถ โข เตสํ ภิกฺขูนํ เอตทโหสิ “ภควตา อนุญฺญาตํ จตุนฺนํ ปาติโมกฺขํ อุทฺทิสิตุํ, มยญฺจมฺหา ตโย ชนา, กถํ นุ โข อเมฺหหิ อุโปสโถ กาตพฺโพ”ติ. ภควโต เอตมตฺถํ อาโรเจสุํ. อนุชานามิ ภิกฺขเว ติณฺณํ ปาริสุทฺธิ\-อุโปสถํ กาตุํ. เอวญฺจ ปน ภิกฺขเว กาตพฺโพ, พฺยตฺเตน ภิกฺขุนา ปฏิพเลน เต ภิกฺขู ญาเปตพฺพา\csromandash{}}

\prose{“สุณนฺตุ เม อายสฺมนฺตา, อชฺชุโปสโถ ปนฺนรโส, ยทา\-ยสฺมนฺตานํ ปตฺตกลฺลํ, มยํ อญฺญมญฺญํ ปาริสุทฺธิ\-อุโปสถํ กเรยฺยามา”ติ.}

\prose{เถเรน ภิกฺขุนา เอกํสํ อุตฺตราสงฺคํ กริตฺวา อุกฺกุฏิกํ นิสีทิตฺวา อญฺชลิํ ปคฺคเหตฺวา เต ภิกฺขู เอวมสฺสุ วจนียา “ปริสุทฺโธ อหํ อาวุโส ปริสุทฺโธติ มํ ธาเรถ, ปริสุทฺโธ อหํ อาวุโส ปริสุทฺโธติ มํ ธาเรถ, ปริสุทฺโธ อหํ อาวุโส ปริสุทฺโธติ มํ ธาเรถา”ติ.}

\prose{นวเกน ภิกฺขุนา เอกํสํ อุตฺตราสงฺคํ กริตฺวา อุกฺกุฏิกํ นิสีทิตฺวา อญฺชลิํ ปคฺคเหตฺวา เต ภิกฺขู เอวมสฺสุ วจนียา “ปริสุทฺโธ อหํ ภนฺเต ปริสุทฺโธติ มํ ธาเรถ, ปริสุทฺโธ อหํ ภนฺเต ปริสุทฺโธติ มํ ธาเรถ, ปริสุทฺโธ อหํ ภนฺเต ปริสุทฺโธติ มํ ธาเรถา”ติ.}

\prose{เตน โข ปน สมเยน อญฺญตรสฺมิํ อาวาเส ตทหุโปสเถ เทฺว ภิกฺขู วิหรนฺติ. อถ โข เตสํ ภิกฺขูนํ เอตทโหสิ “ภควโต อนุญฺญาตํ จตุนฺนํ ปาติโมกฺขํ อุทฺทิสิตุํ, ติณฺณนฺนํ ปาริสุทฺธิ\-อุโปสถํ กาตุํ, มยญฺจมฺหา เทฺว ชนา, กถํ นุ โข อเมฺหหิ อุโปสโถ กาตพฺโพ”ติ. ภควโต เอตมตฺถํ อาโรเจสุํ. อนุชานามิ ภิกฺขเว ทฺวินฺนํ ปาริสุทฺธิ\-อุโปสถํ กาตุํ. เอวญฺจ ปน ภิกฺขเว กาตพฺโพ, เถเรน ภิกฺขุนา เอกํสํ อุตฺตราสงฺคํ กริตฺวา อุกฺกุฏิกํ นิสีทิตฺวา อญฺชลิํ ปคฺคเหตฺวา นโว ภิกฺขุ เอวมสฺส วจนีโย “ปริสุทฺโธ อหํ อาวุโส ปริสุทฺโธติ มํ ธาเรหิ, ปริสุทฺโธ อหํ อาวุโส ปริสุทฺโธติ มํ ธาเรหิ, ปริสุทฺโธ อหํ อาวุโส ปริสุทฺโธติ มํ ธาเรหี”ติ.}

\prose{\csromanfolio{168}นวเกน ภิกฺขุนา เอกํสํ อุตฺตราสงฺคํ กริตฺวา อุกฺกุฏิกํ นิสีทิตฺวา อญฺชลิํ ปคฺคเหตฺวา เถโร ภิกฺขุ เอวมสฺส วจนีโย “ปริสุทฺโธ อหํ ภนฺเต ปริสุทฺโธติ มํ ธาเรถ, ปริสุทฺโธ อหํ ภนฺเต ปริสุทฺโธติ มํ ธาเรถ, ปริสุทฺโธ อหํ ภนฺเต ปริสุทฺโธติ มํ ธาเรถา”ติ.}

\prose{เตน โข ปน สมเยน อญฺญตรสฺมิํ อาวาเส ตทหุโปสเถ เอโก ภิกฺขุ วิหรติ. อถ โข ตสฺส ภิกฺขุโน เอตทโหสิ “ภควตา อนุญฺญาตํ จตุนฺนํ ปาติโมกฺขํ อุทฺทิสิตุํ, ติณฺณนฺนํ ปาริสุทฺธิ\-อุโปสถํ กาตุํ, ทฺวินฺนํ ปาริสุทฺธิ\-อุโปสถํ กาตุํ, อหญฺจมฺหิ เอกโก, กถํ นุ โข มยา อุโปสโถ กาตพฺโพ”ติ. ภควโต เอตมตฺถํ อาโรเจสุํ. อิธ ปน ภิกฺขเว อญฺญตรสฺมิํ อาวาเส ตทหุโปสเถ เอโก ภิกฺขุ วิหรติ. เตน ภิกฺขเว ภิกฺขุนา ยตฺถ ภิกฺขู ปฏิกฺกมนฺติ อุปฏฺฐาน\-สาลาย วา มณฺฑเป วา รุกฺขมูเล วา, โส เทโส สมฺมชฺชิตฺวา ปานียํ ปริโภชนียํ อุปฏฺฐาเปตฺวา อาสนํ ปญฺญเปตฺวา ปทีปํ กตฺวา นิสีทิตพฺพํ. สเจ อญฺเญ ภิกฺขู อาคจฺฉนฺติ, เตหิ สทฺธิํ อุโปสโถ กาตพฺโพ. โน เจ อาคจฺฉนฺติ, “อชฺช เม อุโปสโถ”ติ อธิฏฺฐาตพฺโพ, โน เจ อธิฏฺฐเหยฺย, อาปตฺติ ทุกฺกฏสฺส.}

\prose{ตตฺร ภิกฺขเว ยตฺถ จตฺตาโร ภิกฺขู วิหรนฺติ, น เอกสฺส ปาริสุทฺธิํ อาหริตฺวา ตีหิ ปาติโมกฺขํ อุทฺทิสิตพฺพํ, อุทฺทิเสยฺยุํ เจ, อาปตฺติ ทุกฺกฏสฺส.}

\prose{ตตฺร ภิกฺขเว ยตฺถ ตโย ภิกฺขู วิหรนฺติ, น เอกสฺส ปาริสุทฺธิํ อาหริตฺวา ทฺวีหิ ปาริสุทฺธิ\-อุโปสโถ กาตพฺโพ, กเรยฺยุํ เจ, อาปตฺติ ทุกฺกฏสฺส.}

\prose{ตตฺร ภิกฺขเว ยตฺถ เทฺว ภิกฺขู วิหรนฺติ, น เอกสฺส ปาริสุทฺธิํ อาหริตฺวา เอเกน อธิฏฺฐาตพฺโพ, อธิฏฺฐเหยฺย เจ, อาปตฺติ ทุกฺกฏสฺสาติ.}

\csromanmatikaanchor{mtk.94}
\csromantocmarkref{section}{92. อาปตฺติปฏิกมฺมวิธิ}{mtk.94}
\bookmark[dest={mtk.94},level=1]{92. อาปตฺติปฏิกมฺมวิธิ}
\csromanheader{1}{92. \textbf{อาปตฺติปฏิกมฺมวิธิ}}
\proseitem{169}{เตน โข ปน สมเยน อญฺญตโร ภิกฺขุ ตทหุโปสเถ อาปตฺติํ อาปนฺโน โหติ. อถ โข ตสฺส ภิกฺขุโน เอตทโหสิ “ภควตา ปญฺญตฺตํ ‘น สาปตฺติเกน อุโปสโถ กาตพฺโพ’ติ, อหญฺจมฺหิ อาปตฺติํ อาปนฺโน, กถํ นุ โข มยา ปฏิปชฺชิตพฺพนฺ”ติ. ภควโต เอตมตฺถํ อาโรเจสุํ. อิธ ปน ภิกฺขเว ภิกฺขุ ตทหุโปสเถ อาปตฺติํ อาปนฺโน โหติ, เตน ภิกฺขเว ภิกฺขุนา เอกํ ภิกฺขุํ อุปสงฺกมิตฺวา เอกํสํ อุตฺตราสงฺคํ กริตฺวา อุกฺกุฏิกํ นิสีทิตฺวา อญฺชลิํ \csromanfolio{169}ปคฺคเหตฺวา เอวมสฺส วจนีโย “อหํ อาวุโส อิตฺถนฺนามํ อาปตฺติํ อาปนฺโน, ตํ ปฏิเทเสมี”ติ. เตน วตฺตพฺโพ “ปสฺสสี”ติ. “อาม ปสฺสามี”ติ. “อายติํ สํวเรยฺยาสี”ติ.}

\prose{อิธ ปน ภิกฺขเว ภิกฺขุ ตทหุโปสเถ อาปตฺติยา เวมติโก โหติ, เตน ภิกฺขเว ภิกฺขุนา เอกํ ภิกฺขุํ อุปสงฺกมิตฺวา เอกํสํ อุตฺตราสงฺคํ กริตฺวา อุกฺกุฏิกํ นิสีทิตฺวา อญฺชลิํ ปคฺคเหตฺวา เอวมสฺส วจนีโย “อหํ อาวุโส อิตฺถนฺนามาย อาปตฺติยา เวมติโก, ยทา นิพฺเพมติโก ภวิสฺสามิ, ตทา ตํ อาปตฺติํ ปฏิกริสฺสามี”ติ วตฺวา อุโปสโถ กาตพฺโพ, ปาติโมกฺขํ โสตพฺพํ, น เตฺวว ตปฺปจฺจยา อุโปสถสฺส อนฺตราโย กาตพฺโพติ.}

\prose{เตน โข ปน สมเยน ฉพฺพคฺคิยา ภิกฺขู สภาคํ อาปตฺติํ เทเสนฺติ. ภควโต เอตมตฺถํ อาโรเจสุํ. น ภิกฺขเว สภาคา อาปตฺติ เทเสตพฺพา, โย เทเสยฺย, อาปตฺติ ทุกฺกฏสฺสาติ.}

\prose{เตน โข ปน สมเยน ฉพฺพคฺคิยา ภิกฺขู สภาคํ อาปตฺติํ ปฏิคฺคณฺหนฺติ. ภควโต เอตมตฺถํ อาโรเจสุํ. น ภิกฺขเว สภาคา อาปตฺติ ปฏิคฺคเหตพฺพา, โย ปฏิคฺคเณฺหยฺย, อาปตฺติ ทุกฺกฏสฺสาติ.}

\csromanmatikaanchor{mtk.95}
\csromantocmarkref{section}{93. อาปตฺติอาวิกรณวิธิ}{mtk.95}
\bookmark[dest={mtk.95},level=1]{93. อาปตฺติอาวิกรณวิธิ}
\csromanheader{1}{93. อาปตฺติ\-อาวิกรณวิธิ}
\proseitem{170}{เตน โข ปน สมเยน อญฺญตโร ภิกฺขุ ปาติโมกฺเข อุทฺทิสฺสมาเน อาปตฺติํ สรติ, อถ โข ตสฺส ภิกฺขุโน เอตทโหสิ “ภควตา ปญฺญตฺตํ ‘น สาปตฺติเกน อุโปสโถ กาตพฺโพ’ติ, อหญฺจมฺหิ อาปตฺติํ อาปนฺโน, กถํ นุ โข มยา ปฏิปชฺชิตพฺพนฺ”ติ.}

\setlength{\csromangathapagewidth}{0pt}
\setlength{\csromangathawakwidth}{0pt}
\csromangathameasuregroup{}{ภควโต เอตมตฺถํ อาโรเจสุํ.}
\csromangathameasurewak{ภควโต เอตมตฺถํ อาโรเจสุํ.}
\setlength{\csromangathaleftcol}{0pt}
\csromangathagroup{\csromangathastanza{\csromangathawak{\csromangathaitemline{170}{ภควโต เอตมตฺถํ อาโรเจสุํ.}}}}

\prose{อิธ ปน ภกฺขเว ภิกฺขุ ปาติโมกฺเข อุทฺทิสฺสมาเน อาปตฺติํ สรติ, เตน ภิกฺขเว ภิกฺขุนา สามนฺโต ภิกฺขุ เอวมสฺส วจนีโย “อหํ อาวุโส อิตฺถนฺนามํ อาปตฺติํ อาปนฺโน, อิโต วุฏฺฐหิตฺวา ตํ อาปตฺติํ ปาฏิกริสฺสามี”ติ วตฺวา อุโปสโถ กาตพฺโพ, ปาติโมกฺขํ โสตพฺพํ, น เตฺวว ตปฺปจฺจยา อุโปสถสฺส อนฺตราโย กาตพฺโพ.}

\proseitem{170}{อิธ ปน ภิกฺขเว ภิกฺขุ ปาติโมกฺเข อุทฺทิสฺสมาเน อาปตฺติยา เวมติโก โหติ, เตน ภิกฺขเว ภิกฺขุนา สามนฺโต ภิกฺขุ เอวมสฺส \csromanfolio{170}วจนีโย “อหํ อาวุโส อิตฺถนฺนามาย อาปตฺติยา เวมติโก, ยทา นิพฺเพมติโก ภวิสฺสามิ, ตทา ตํ อาปตฺติํ ปฏิกริสฺสามี”ติ วตฺวา อุโปสโถ กาตพฺโพ, ปาติโมกฺขํ โสตพฺพํ, น เตฺวว ตปฺปจฺจยา อุโปสถสฺส อนฺตราโย กาตพฺโพติ.}

\csromanmatikaanchor{mtk.96}
\csromantocmarkref{section}{94. สภาคาปตฺติปฏิกมฺมวิธิ}{mtk.96}
\bookmark[dest={mtk.96},level=1]{94. สภาคาปตฺติปฏิกมฺมวิธิ}
\csromanheader{1}{94. \textbf{สภาคาปตฺติ}\-\textbf{ปฏิกมฺมวิธิ}}
\proseitem{171}{เตน โข ปน สมเยน อญฺญตรสฺมิํ อาวาเส ตทหุโปสเถ สพฺโพ สํโฆ สภาคํ อาปตฺติํ อาปนฺโน โหติ. อถ โข เตสํ ภิกฺขูนํ เอตทโหสิ “ภควตา ปญฺญตฺตํ ‘น สภาคา อาปตฺติ เทเสตพฺพา, น สภาคา อาปตฺติ ปฏิคฺคเหตพฺพา’ติ, อยญฺจ สพฺโพ สํโฆ สภาคํ อาปตฺติํ อาปนฺโน, กถํ นุ โข อเมฺหหิ ปฏิปชฺชิตพฺพนฺ”ติ. ภควโต เอตมตฺถํ อาโรเจสุํ.}

\prose{อิธ ปน ภิกฺขเว อญฺญตรสฺมิํ อาวาเส ตทหุโปสเถ สพฺโพ สํโฆ สภาคํ อาปตฺติํ อาปนฺโน โหติ, เตหิ ภิกฺขเว ภิกฺขูหิ เอโก ภิกฺขุ สามนฺตา อาวาสา สชฺชุกํ ปาเหตพฺโพ “คจฺฉาวุโส ตํ อาปตฺติํ ปฏิกริตฺวา อาคจฺฉ, มยํ เต สนฺติเก อาปตฺติํ ปฏิกริสฺสามา”ติ. เอวญฺเจตํ ลเภถ, อิจฺเจตํ กุสลํ. โน เจ ลเภถ, พฺยตฺเตน ภิกฺขุนา ปฏิพเลน สํโฆ ญาเปตพฺโพ\csromandash{}}

\prose{“สุณาตุ เม ภนฺเต สํโฆ, อยํ สพฺโพ สํโฆ สภาคํ อาปตฺติํ อาปนฺโน, ยทา อญฺญํ ภิกฺขุํ สุทฺธํ อนาปตฺติกํ ปสฺสิสฺสติ, ตทา ตสฺส สนฺติเก ตํ อาปตฺติํ ปฏิกริสฺสตี”ติ วตฺวา อุโปสโถ กาตพฺโพ, ปาติโมกฺขํ อุทฺทิสิตพฺพํ, น เตฺวว ตปฺปจฺจยา อุโปสถสฺส อนฺตราโย กาตพฺโพ.}

\prose{อิธ ปน ภิกฺขเว อญฺญตรสฺมิํ อาวาเส ตทหุโปสเถ สพฺโพ สํโฆ สภาคาย อาปตฺติยา เวมติโก โหติ, พฺยตฺเตน ภิกฺขุนา ปฏิพเลน สํโฆ ญาเปตพฺโพ\csromandash{}}

\prose{“สุณาตุ เม ภนฺเต สํโฆ, อยํ สพฺโพ สํโฆ สภาคาย อาปตฺติยา เวมติโก, ยทา นิพฺเพมติโก ภวิสฺสติ, ตทา ตํ อาปตฺติํ ปฏิกริสฺสตี”ติ วตฺวา อุโปสโถ กาตพฺโพ, ปาติโมกฺขํ อุทฺทิสิตพฺพํ, น เตฺวว ตปฺปจฺจยา อุโปสถสฺส อนฺตราโย กาตพฺโพ.}

\prose{\csromanfolio{171}อิธ ปน ภิกฺขเว อญฺญตรสฺมิํ อาวาเส วสฺสูปคโต สํโฆ สภาคํ อาปตฺติํ อาปนฺโน โหติ, เตหิ ภิกฺขเว ภิกฺขูหิ เอโก ภิกฺขุ สามนฺตา อาวาสา สชฺชุกํ ปาเหตพฺโพ “คจฺฉาวุโส ตํ อาปตฺติํ ปฏิกริตฺวา อาคจฺฉ, มยํ เต สนฺติเก ตํ อาปตฺติํ ปฏิกริสฺสามา”ติ.}

\setlength{\csromangathapagewidth}{0pt}
\setlength{\csromangathawakwidth}{0pt}
\csromangathameasuregroup{เอวญฺเจตํ ลเภถ,}{\csromangathabat{เอวญฺเจตํ ลเภถ,}{อิจฺเจตํ กุสลํ.}}
\csromangathasetleft{เอวญฺเจตํ ลเภถ,}
\csromangathagroup{\csromangathastanza{\csromangathabat{เอวญฺเจตํ ลเภถ,}{อิจฺเจตํ กุสลํ.}}}

\prose{โน เจ ลเภถ, เอโก ภิกฺขุ สตฺตาหกาลิกํ ปาเหตพฺโพ “คจฺฉาวุโส ตํ อาปตฺติํ ปฏิกริตฺวา อาคจฺฉ, มยํ เต สนฺติเก ตํ อาปตฺติํ ปฏิกริสฺสามา”ติ.}

\prose{เตน โข ปน สมเยน อญฺญตรสฺมิํ อาวาเส สพฺโพ สํโฆ สภาคํ อาปตฺติํ อาปนฺโน โหติ, โส น ชานาติ ตสฺสา อาปตฺติยา นามโคตฺตํ.}

\prose{ตตฺถ อญฺโญ ภิกฺขุ อาคจฺฉติ พหุสฺสุโต อาคตาคโม ธมฺมธโร วินยธโร มาติกาธโร ปณฺฑิโต พฺยตฺโต เมธาวี ลชฺชี กุกฺกุจฺจโก สิกฺขากาโม, ตเมนํ อญฺญตโร ภิกฺขุ เยน โส ภิกฺขุ เตนุปสงฺกมิ, อุปสงฺกมิตฺวา ตํ ภิกฺขุํ เอตทโวจ “โย นุ โข อาวุโส เอวญฺเจวญฺจ กโรติ, กิํ นาม โส อาปตฺติํ อาปชฺชตี”ติ.}

\prose{โส เอวมาห “โย โข อาวุโส เอวญฺเจวญฺจ กโรติ, อิมํ นาม โส อาปตฺติํ อาปชฺชติ, อิมํ นาม ตฺวํ อาวุโส อาปตฺติํ อาปนฺโน, ปฏิกโรหิ ตํ อาปตฺตินฺ”ติ.}

\prose{โส เอวมาห “น โข อหํ อาวุโส เอโกว อิมํ อาปตฺติํ อาปนฺโน, อยํ สพฺโพ สํโฆ อิมํ อาปตฺติํ อาปนฺโน”ติ.}

\prose{โส เอวมาห “กิํ เต อาวุโส กริสฺสติ ปโร อาปนฺโน วา อนาปนฺโน วา, อิงฺฆ ตฺวํ อาวุโส สกาย อาปตฺติยา วุฏฺฐาหี”ติ.}

\prose{อถ โข โส ภิกฺขุ ตสฺส ภิกฺขุโน วจเนน ตํ อาปตฺติํ ปฏิกริตฺวา เยน เต ภิกฺขู เตนุปสงฺกมิ, อุปสงฺกมิตฺวา เต ภิกฺขู เอตทโวจ “โย กิร อาวุโส เอวญฺเจวญฺจ กโรติ, อิมํ นาม โส อาปตฺติํ อาปชฺชติ, อิมํ นาม ตุเมฺห อาวุโส อาปตฺติํ อาปนฺนา, ปฏิกโรถ ตํ อาปตฺตินฺ”ติ.}

\prose{อถ โข เต ภิกฺขู น อิจฺฉิํสุ ตสฺส ภิกฺขุโน วจเนน ตํ อาปตฺติํ ปฏิกาตุํ. ภควโต เอตมตฺถํ อาโรเจสุํ.}

\proseitemwithcloserruled{171}{อิธ ปน ภิกฺขเว อญฺญตรสฺมิํ อาวาเส สพฺโพ สํโฆ สภาคํ อาปตฺติํ อาปนฺโน โหติ, โส น ชานาติ ตสฺสา อาปตฺติยา นามโคตฺตํ. ตตฺถ อญฺโญ ภิกฺขุ อาคจฺฉติ พหุสฺสุโต อาคตาคโม \csromanfolio{172}ธมฺมธโร วินยธโร มาติกาธโร ปณฺฑิโต พฺยตฺโต เมธาวี ลชฺชี กุกฺกุจฺจโก สิกฺขากาโม, ตเมนํ อญฺญตโร ภิกฺขุ เยน โส ภิกฺขุ เตนุปสงฺกมิ, อุปสงฺกมิตฺวา ตํ ภิกฺขุํ เอวํ วเทติ “โย นุ โข อาวุโส เอวญฺเจวญฺจ กโรติ, กิํ นาม โส อาปตฺติํ อาปชฺชตี”ติ. โส เอวํ วเทติ “โย โข อาวุโส เอวญฺเจวญฺจ กโรติ, อิมํ นาม โส อาปตฺติํ อาปชฺชติ, อิมํ นาม ตฺวํ อาวุโส อาปตฺติํ อาปนฺโน, ปฏิกโรหิ ตํ อาปตฺตินฺ”ติ. โส เอวํ วเทติ “น โข อหํ อาวุโส เอโกว อิมํ อาปตฺติํ อาปนฺโน, อยํ สพฺโพ สํโฆ อิมํ อาปตฺติํ อาปนฺโน”ติ. โส เอวํ วเทติ “กิํ เต อาวุโส กริสฺสติ ปโร อาปนฺโน วา อนาปนฺโน วา, อิงฺฆ ตฺวํ อาวุโส สกาย อาปตฺติยา วุฏฺฐาหี”ติ. โส เจ ภิกฺขเว ภิกฺขุ ตสฺส ภิกฺขุโน วจเนน ตํ อาปตฺติํ ปฏิกริตฺวา เยน เต ภิกฺขู เตนุปสงฺกมิ, อุปสงฺกมิตฺวา เต ภิกฺขู เอวํ วเทติ “โย กิร อาวุโส เอวญฺเจวญฺจ กโรติ, อิมํ นาม โส อาปตฺติํ อาปชฺชติ, อิมํ นาม ตุเมฺห อาวุโส อาปตฺติํ อาปนฺนา, ปฏิกโรถ ตํ อาปตฺตินฺ”ติ. เต เจ ภิกฺขเว ภิกฺขู ตสฺส ภิกฺขุโน วจเนน ตํ อาปตฺติํ ปฏิกเรยฺยุํ, อิจฺเจตํ กุสลํ. โน เจ ปฏิกเรยฺยุํ, น เต ภิกฺขเว ภิกฺขู เตน ภิกฺขุนา อกามา วจนียาติ.}{โจทนาวตฺถุ\-ภาณวาโร นิฏฺฐิโต ทุติโย.}
\csromanmatikaanchor{mtk.97}
\csromantocmarkref{section}{95. อนาปตฺติปนฺนรสก}{mtk.97}
\bookmark[dest={mtk.97},level=1]{95. อนาปตฺติปนฺนรสก}
\csromanheader{1}{95. \textbf{อนาปตฺติ}\-\textbf{ปนฺนรสก}}
\proseitem{172}{เตน โข ปน สมเยน อญฺญตรสฺมิํ อาวาเส ตทหุโปสเถ สมฺพหุลา อาวาสิกา ภิกฺขู สนฺนิปติํสุ จตฺตาโร วา อติเรกา วา, เต น ชานิํสุ “อตฺถญฺเญ อาวาสิกา ภิกฺขู อนาคตา”ติ, เต ธมฺมสญฺญิโน วินยสญฺญิโน วคฺคา สมคฺคสญฺญิโน อุโปสถํ อกํสุ ปาติโมกฺขํ อุทฺทิสิํสุ, เตหิ อุทฺทิสฺสมาเน ปาติโมกฺเข อถญฺเญ อาวาสิกา ภิกฺขู อาคจฺฉิํสุ พหุตรา. ภควโต เอตมตฺถํ อาโรเจสุํ.}

\prose{อิธ ปน ภิกฺขเว อญฺญตรสฺมิํ อาวาเส ตทหุโปสเถ สมฺพหุลา อาวาสิกา ภิกฺขู สนฺนิปตนฺติ จตฺตาโร วา อติเรกา วา, เต น ชานนฺติ \csromanfolio{173}“อตฺถญฺเญ อาวาสิกา ภิกฺขู อนาคตา”ติ, เต ธมฺมสญฺญิโน วินยสญฺญิโน วคฺคา สมคฺคสญฺญิโน อุโปสถํ กโรนฺติ ปาติโมกฺขํ อุทฺทิสนฺติ, เตหิ อุทฺทิสฺสมาเน ปาติโมกฺเข อถญฺเญ อาวาสิกา ภิกฺขู อาคจฺฉนฺติ พหุตรา, เตหิ ภิกฺขเว ภิกฺขูหิ ปุน ปาติโมกฺขํ อุทฺทิสิตพฺพํ, อุทฺเทสกานํ อนาปตฺติ. (1)}

\prose{อิธ ปน ภิกฺขเว อญฺญตรสฺมิํ อาวาเส ตทหุโปสเถ สมฺพหุลา อาวาสิกา ภิกฺขู สนฺนิปตนฺติ จตฺตาโร วา อติเรกา วา, เต น ชานนฺติ “อตฺถญฺเญ อาวาสิกา ภิกฺขู อนาคตา”ติ, เต ธมฺมสญฺญิโน วินยสญฺญิโน วคฺคา สมคฺคสญฺญิโน อุโปสถํ กโรนฺติ ปาติโมกฺขํ อุทฺทิสนฺติ, เตหิ อุทฺทิสฺสมาเน ปาติโมกฺเข อถญฺเญ อาวาสิกา ภิกฺขู อาคจฺฉนฺติ สมสมา, อุทฺทิฏฺฐํ สุอุทฺทิฏฺฐํ, อวเสสํ โสตพฺพํ. อุทฺเทสกานํ อนาปตฺติ. (2)}

\prose{อิธ ปน ภิกฺขเว อญฺญตรสฺมิํ อาวาเส ตทหุโปสเถ สมฺพหุลา อาวาสิกา ภิกฺขู สนฺนิปตนฺติ จตฺตาโร วา อติเรกา วา, เต น ชานนฺติ “อตฺถญฺเญ อาวาสิกา ภิกฺขู อนาคตา”ติ, เต ธมฺมสญฺญิโน วินยสญฺญิโน วคฺคา สมคฺคสญฺญิโน อุโปสถํ กโรนฺติ ปาติโมกฺขํ อุทฺทิสนฺติ, เตหิ อุทฺทิสฺสมาเน ปาติโมกฺเข อถญฺเญ อาวาสิกา ภิกฺขู อาคจฺฉนฺติ โถกตรา, อุทฺทิฏฺฐํ สุอุทฺทิฏฺฐํ, อวเสสํ โสตพฺพํ. อุทฺเทสกานํ อนาปตฺติ. (3)}

\prose{อิธ ปน ภิกฺขเว อญฺญตรสฺมิํ อาวาเส ตทหุโปสเถ สมฺพหุลา อาวาสิกา ภิกฺขู สนฺนิปตนฺติ จตฺตาโร วา อติเรกา วา, เต น ชานนฺติ “อตฺถญฺเญ อาวาสิกา ภิกฺขู อนาคตา”ติ, เต ธมฺมสญฺญิโน วินยสญฺญิโน วคฺคา สมคฺคสญฺญิโน อุโปสถํ กโรนฺติ ปาติโมกฺขํ อุทฺทิสนฺติ, เตหิ อุทฺทิฏฺฐมตฺเต ปาติโมกฺเข อถญฺเญ อาวาสิกา ภิกฺขู อาคจฺฉนฺติ พหุตรา, เตหิ ภิกฺขเว ภิกฺขูหิ ปุน ปาติโมกฺขํ อุทฺทิสิตพฺพํ. อุทฺเทสกานํ อนาปตฺติ. (4)}

\prose{อิธ ปน ภิกฺขเว อญฺญตรสฺมิํ อาวาเส ตทหุโปสเถ สมฺพหุลา อาวาสิกา ภิกฺขู สนฺนิปตนฺติ จตฺตาโร วา อติเรกา วา, เต น ชานนฺติ “อตฺถญฺเญ อาวาสิกา ภิกฺขู อนาคตา”ติ, เต ธมฺมสญฺญิโน วินยสญฺญิโน วคฺคา สมคฺคสญฺญิโน อุโปสถํ กโรนฺติ ปาติโมกฺขํ อุทฺทิสนฺติ, \csromanfolio{174}เตหิ อุทฺทิฏฺฐมตฺเต ปาติโมกฺเข อถญฺเญ อาวาสิกา ภิกฺขู อาคจฺฉนฺติ สมสมา, อุทฺทิฏฺฐํ สุอุทฺทิฏฺฐํ, เตสํ สนฺติเก ปาริสุทฺธิ อาโรเจตพฺพา. อุทฺเทสกานํ อนาปตฺติ. (5)}

\prose{อิธ ปน ภิกฺขเว อญฺญตรสฺมิํ อาวาเส ตทหุโปสเถ สมฺพหุลา อาวาสิกา ภิกฺขู สนฺนิปตนฺติ จตฺตาโร วา อติเรกา วา, เต น ชานนฺติ “อตฺถญฺเญ อาวาสิกา ภิกฺขู อนาคตา”ติ, เต ธมฺมสญฺญิโน วินยสญฺญิโน วคฺคา สมคฺคสญฺญิโน อุโปสถํ กโรนฺติ ปาติโมกฺขํ อุทฺทิสนฺติ, เตหิ อุทฺทิฏฺฐมตฺเต ปาติโมกฺเข อถญฺเญ อาวาสิกา ภิกฺขู อาคจฺฉนฺติ โถกตรา, อุทฺทิฏฺฐํ สุอุทฺทิฏฺฐํ, เตสํ สนฺติเก ปาริสุทฺธิ อาโรเจตพฺพา. อุทฺเทสกานํ อนาปตฺติ. (6)}

\prose{อิธ ปน ภิกฺขเว อญฺญตรสฺมิํ อาวาเส ตทหุโปสเถ สมฺพหุลา อาวาสิกา ภิกฺขู สนฺนิปตนฺติ จตฺตาโร วา อติเรกา วา, เต น ชานนฺติ “อตฺถญฺเญ อาวาสิกา ภิกฺขู อนาคตา”ติ, เต ธมฺมสญฺญิโน วินยสญฺญิโน วคฺคา สมคฺคสญฺญิโน อุโปสถํ กโรนฺติ ปาติโมกฺขํ อุทฺทิสนฺติ, เตหิ อุทฺทิฏฺฐมตฺเต ปาติโมกฺเข อวุฏฺฐิตาย ปริสาย อถญฺเญ อาวาสิกา ภิกฺขู อาคจฺฉนฺติ พหุตรา, เตหิ ภิกฺขเว ภิกฺขูหิ ปุน ปาติโมกฺขํ อุทฺทิสิตพฺพํ. อุทฺเทสกานํ อนาปตฺติ. (7)}

\prose{อิธ ปน ภิกฺขเว อญฺญตรสฺมิํ อาวาเส ตทหุโปสเถ สมฺพหุลา อาวาสิกา ภิกฺขู สนฺนิปตนฺติ จตฺตาโร วา อติเรกา วา, เต น ชานนฺติ “อตฺถญฺเญ อาวาสิกา ภิกฺขู อนาคตา”ติ, เต ธมฺมสญฺญิโน วินยสญฺญิโน วคฺคา สมคฺคสญฺญิโน อุโปสถํ กโรนฺติ ปาติโมกฺขํ อุทฺทิสนฺติ, เตหิ อุทฺทิฏฺฐมตฺเต ปาติโมกฺเข อวุฏฺฐิตาย ปริสาย อถญฺเญ อาวาสิกา ภิกฺขู อาคจฺฉนฺติ สมสมา, อุทฺทิฏฺฐํ สุอุทฺทิฏฺฐํ, เตสํ สนฺติเก ปาริสุทฺธิ อาโรเจตพฺพา. อุทฺเทสกานํ อนาปตฺติ. (8)}

\prose{อิธ ปน ภิกฺขเว อญฺญตรสฺมิํ อาวาเส ตทหุโปสเถ สมฺพหุลา อาวาสิกา ภิกฺขู สนฺนิปตนฺติ จตฺตาโร วา อติเรกา วา, เต น ชานนฺติ “อตฺถญฺเญ อาวาสิกา ภิกฺขู อนาคตา”ติ, เต ธมฺมสญฺญิโน วินยสญฺญิโน วคฺคา สมคฺคสญฺญิโน อุโปสถํ กโรนฺติ ปาติโมกฺขํ อุทฺทิสนฺติ, เตหิ อุทฺทิฏฺฐมตฺเต ปาติโมกฺเข อวุฏฺฐิตาย ปริสาย อถญฺเญ \csromanfolio{175}อาวาสิกา ภิกฺขู อาคจฺฉนฺติ โถกตรา, อุทฺทิฏฺฐํ สุอุทฺทิฏฺฐํ, เตสํ สนฺติเก ปาริสุทฺธิ อาโรเจตพฺพา. อุทฺเทสกานํ อนาปตฺติ. (9)}

\prose{อิธ ปน ภิกฺขเว อญฺญตรสฺมิํ อาวาเส ตทหุโปสเถ สมฺพหุลา อาวาสิกา ภิกฺขู สนฺนิปตนฺติ จตฺตาโร วา อติเรกา วา, เต น ชานนฺติ “อตฺถญฺเญ อาวาสิกา ภิกฺขู อนาคตา”ติ, เต ธมฺมสญฺญิโน วินยสญฺญิโน วคฺคา สมคฺคสญฺญิโน อุโปสถํ กโรนฺติ ปาติโมกฺขํ อุทฺทิสนฺติ, เตหิ อุทฺทิฏฺฐมตฺเต ปาติโมกฺเข เอกจฺจาย วุฏฺฐิตาย ปริสาย อถญฺเญ อาวาสิกา ภิกฺขู อาคจฺฉนฺติ พหุตรา, เตหิ ภิกฺขเว ภิกฺขูหิ ปุน ปาติโมกฺขํ อุทฺทิสิตพฺพํ. อุทฺเทสกานํ อนาปตฺติ. (10)}

\prose{อิธ ปน ภิกฺขเว อญฺญตรสฺมิํ อาวาเส ตทหุโปสเถ สมฺพหุลา อาวาสิกา ภิกฺขู สนฺนิปตนฺติ จตฺตาโร วา อติเรกา วา, เต น ชานนฺติ “อตฺถญฺเญ อาวาสิกา ภิกฺขู อนาคตา”ติ, เต ธมฺมสญฺญิโน วินยสญฺญิโน วคฺคา สมคฺคสญฺญิโน อุโปสถํ กโรนฺติ ปาติโมกฺขํ อุทฺทิสนฺติ, เตหิ อุทฺทิฏฺฐมตฺเต ปาติโมกฺเข เอกจฺจาย วุฏฺฐิตาย ปริสาย อถญฺเญ อาวาสิกา ภิกฺขู อาคจฺฉนฺติ สมสมา, อุทฺทิฏฺฐํ สุอุทฺทิฏฺฐํ, เตสํ สนฺติเก ปาริสุทฺธิ อาโรเจตพฺพา. อุทฺเทสกานํ อนาปตฺติ. (11)}

\prose{อิธ ปน ภิกฺขเว อญฺญตรสฺมิํ อาวาเส ตทหุโปสเถ สมฺพหุลา อาวาสิกา ภิกฺขู สนฺนิปตนฺติ จตฺตาโร วา อติเรกา วา, เต น ชานนฺติ “อตฺถญฺเญ อาวาสิกา ภิกฺขู อนาคตา”ติ, เต ธมฺมสญฺญิโน วินยสญฺญิโน วคฺคา สมคฺคสญฺญิโน อุโปสถํ กโรนฺติ ปาติโมกฺขํ อุทฺทิสนฺติ, เตหิ อุทฺทิฏฺฐมตฺเต ปาติโมกฺเข เอกจฺจาย วุฏฺฐิตาย ปริสาย อถญฺเญ อาวาสิกา ภิกฺขู อาคจฺฉนฺติ โถกตรา, อุทฺทิฏฺฐํ สุอุทฺทิฏฺฐํ, เตสํ สนฺติเก ปาริสุทฺธิ อาโรเจตพฺพา. อุทฺเทสกานํ อนาปตฺติ. (12)}

\prose{อิธ ปน ภิกฺขเว อญฺญตรสฺมิํ อาวาเส ตทหุโปสเถ สมฺพหุลา อาวาสิกา ภิกฺขู สนฺนิปตนฺติ จตฺตาโร วา อติเรกา วา, เต น ชานนฺติ “อตฺถญฺเญ อาวาสิกา ภิกฺขู อนาคตา”ติ, เต ธมฺมสญฺญิโน วินยสญฺญิโน วคฺคา สมคฺคสญฺญิโน อุโปสถํ กโรนฺติ ปาติโมกฺขํ อุทฺทิสนฺติ, เตหิ อุทฺทิฏฺฐมตฺเต ปาติโมกฺเข สพฺพาย วุฏฺฐิตาย ปริสาย อถญฺเญ อาวาสิกา ภิกฺขู อาคจฺฉนฺติ พหุตรา, เตหิ ภิกฺขเว ภิกฺขูหิ ปุน ปาติโมกฺขํ อุทฺทิสิตพฺพํ. อุทฺเทสกานํ อนาปตฺติ. (13)}

\prose{\csromanfolio{176}อิธ ปน ภิกฺขเว อญฺญตรสฺมิํ อาวาเส ตทหุโปสเถ สมฺพหุลา อาวาสิกา ภิกฺขู สนฺนิปตนฺติ จตฺตาโร วา อติเรกา วา, เต น ชานนฺติ “อตฺถญฺเญ อาวาสิกา ภิกฺขู อนาคตา”ติ, เต ธมฺมสญฺญิโน วินยสญฺญิโน วคฺคา สมคฺคสญฺญิโน อุโปสถํ กโรนฺติ ปาติโมกฺขํ อุทฺทิสนฺติ, เตหิ อุทฺทิฏฺฐมตฺเต ปาติโมกฺเข สพฺพาย วุฏฺฐิตาย ปริสาย อถญฺเญ อาวาสิกา ภิกฺขู อาคจฺฉนฺติ สมสมา, อุทฺทิฏฺฐํ สุอุทฺทิฏฺฐํ, เตสํ สนฺติเก ปาริสุทฺธิ อาโรเจตพฺพา. อุทฺเทสกานํ อนาปตฺติ. (14)}

\prosewithcloserruled{อิธ ปน ภิกฺขเว อญฺญตรสฺมิํ อาวาเส ตทหุโปสเถ สมฺพหุลา อาวาสิกา ภิกฺขู สนฺนิปตนฺติ จตฺตาโร วา อติเรกา วา, เต น ชานนฺติ “อตฺถญฺเญ อาวาสิกา ภิกฺขู อนาคตา”ติ, เต ธมฺมสญฺญิโน วินยสญฺญิโน วคฺคา สมคฺคสญฺญิโน อุโปสถํ กโรนฺติ ปาติโมกฺขํ อุทฺทิสนฺติ, เตหิ อุทฺทิฏฺฐมตฺเต ปาติโมกฺเข สพฺพาย วุฏฺฐิตาย ปริสาย อถญฺเญ อาวาสิกา ภิกฺขู อาคจฺฉนฺติ โถกตรา, อุทฺทิฏฺฐํ สุอุทฺทิฏฺฐํ, เตสํ สนฺติเก ปาริสุทฺธิ อาโรเจตพฺพา. อุทฺเทสกานํ อนาปตฺติ. (15)}{อนาปตฺติ\-ปนฺนรสกํ นิฏฺฐิตํ.}
\csromanmatikaanchor{mtk.98}
\csromantocmarkref{section}{96. วคฺคาวคฺคสญฺญีปนฺนรสก}{mtk.98}
\bookmark[dest={mtk.98},level=1]{96. วคฺคาวคฺคสญฺญีปนฺนรสก}
\csromanheader{1}{96. \textbf{วคฺคา}\-\textbf{วคฺคสญฺญี}\-\textbf{ปนฺนรสก}}
\proseitem{173}{อิธ ปน ภิกฺขเว อญฺญตรสฺมิํ อาวาเส ตทหุโปสเถ สมฺพหุลา อาวาสิกา ภิกฺขู สนฺนิปตนฺติ จตฺตาโร วา อติเรกา วา, เต ชานนฺติ “อตฺถญฺเญ อาวาสิกา ภิกฺขู อนาคตา”ติ, เต ธมฺมสญฺญิโน วินยสญฺญิโน วคฺคา วคฺคสญฺญิโน อุโปสถํ กโรนฺติ ปาติโมกฺขํ อุทฺทิสนฺติ, เตหิ อุทฺทิสฺสมาเน ปาติโมกฺเข อถญฺเญ อาวาสิกา ภิกฺขู อาคจฺฉนฺติ พหุตรา, เตหิ ภิกฺขเว ภิกฺขูหิ ปุน ปาติโมกฺขํ อุทฺทิสิตพฺพํ. อุทฺเทสกานํ อาปตฺติ ทุกฺกฏสฺส. (1)}

\prose{อิธ ปน ภิกฺขเว อญฺญตรสฺมิํ อาวาเส ตทหุโปสเถ สมฺพหุลา อาวาสิกา ภิกฺขู สนฺนิปตนฺติ จตฺตาโร วา อติเรกา วา, เต ชานนฺติ “อตฺถญฺเญ อาวาสิกา ภิกฺขู อนาคตา”ติ, เต ธมฺมสญฺญิโน วินยสญฺญิโน วคฺคา วคฺคสญฺญิโน อุโปสถํ กโรนฺติ ปาติโมกฺขํ อุทฺทิสนฺติ, เตหิ อุทฺทิสฺสมาเน ปาติโมกฺเข อถญฺเญ อาวาสิกา ภิกฺขู อาคจฺฉนฺติ \csromanfolio{177}สมสมา, อุทฺทิฏฺฐํ สุอุทฺทิฏฺฐํ, อวเสสํ โสตพฺพํ. อุทฺเทสกานํ อาปตฺติ ทุกฺกฏสฺส. (2)}

\prose{อิธ ปน ภิกฺขเว อญฺญตรสฺมิํ อาวาเส ตทหุโปสเถ สมฺพหุลา อาวาสิกา ภิกฺขู สนฺนิปตนฺติ จตฺตาโร วา อติเรกา วา, เต ชานนฺติ “อตฺถญฺเญ อาวาสิกา ภิกฺขู อนาคตา”ติ, เต ธมฺมสญฺญิโน วินยสญฺญิโน วคฺคา วคฺคสญฺญิโน อุโปสถํ กโรนฺติ ปาติโมกฺขํ อุทฺทิสนฺติ, เตหิ อุทฺทิสฺสมาเน ปาติโมกฺเข อถญฺเญ อาวาสิกา ภิกฺขู อาคจฺฉนฺติ โถกตรา, อุทฺทิฏฺฐํ สุอุทฺทิฏฺฐํ, อวเสสํ โสตพฺพํ. อุทฺเทสกานํ อาปตฺติ ทุกฺกฏสฺส. (3)}

\prosewithcloserruled{อิธ ปน ภิกฺขเว อญฺญตรสฺมิํ อาวาเส ตทหุโปสเถ สมฺพหุลา อาวาสิกา ภิกฺขู สนฺนิปตนฺติ จตฺตาโร วา อติเรกา วา, เต ชานนฺติ “อตฺถญฺเญ อาวาสิกา ภิกฺขู อนาคตา”ติ, เต ธมฺมสญฺญิโน วินยสญฺญิโน วคฺคา วคฺคสญฺญิโน อุโปสถํ กโรนฺติ ปาติโมกฺขํ อุทฺทิสนฺติ, เตหิ อุทฺทิฏฺฐมตฺเต ปาติโมกฺเข ฯเปฯ อวุฏฺฐิตาย ปริสาย ฯเปฯ เอกจฺจาย วุฏฺฐิตาย ปริสาย ฯเปฯ สพฺพาย วุฏฺฐิตาย ปริสาย อถญฺเญ อาวาสิกา ภิกฺขู อาคจฺฉนฺติ พหุตรา ฯเปฯ สมสมา ฯเปฯ โถกตรา, อุทฺทิฏฺฐํ สุอุทฺทิฏฺฐํ, เตสํ สนฺติเก ปาริสุทฺธิ อาโรเจตพฺพา. อุทฺเทสกานํ อาปตฺติ ทุกฺกฏสฺส. \mbox{(4-15)}}{วคฺคา\-วคฺค\-สญฺญิ\-ปนฺนรสกํ นิฏฺฐิตํ.}
\csromanmatikaanchor{mtk.99}
\csromantocmarkref{section}{97. เวมติกปนฺนรสก}{mtk.99}
\bookmark[dest={mtk.99},level=1]{97. เวมติกปนฺนรสก}
\csromanheader{1}{97. \textbf{เวมติก}\-\textbf{ปนฺนรสก}}
\proseitem{174}{อิธ ปน ภิกฺขเว อญฺญตรสฺมิํ อาวาเส ตทหุโปสเถ สมฺพหุลา อาวาสิกา ภิกฺขู สนฺนิปตนฺติ จตฺตาโร วา อติเรกา วา, เต ชานนฺติ “อตฺถญฺเญ อาวาสิกา ภิกฺขู อนาคตา”ติ, เต “กปฺปติ นุ โข อมฺหากํ อุโปสโถ กาตุํ, น นุ โข กปฺปตี”ติ เวมติกา อุโปสถํ กโรนฺติ ปาติโมกฺขํ อุทฺทิสนฺติ, เตหิ อุทฺทิสฺสมาเน ปาติโมกฺเข อถญฺเญ อาวาสิกา ภิกฺขู อาคจฺฉนฺติ พหุตรา, เตหิ ภิกฺขเว ภิกฺขูหิ ปุน ปาติโมกฺขํ อุทฺทิสิตพฺพํ. อุทฺเทสกานํ อาปตฺติ ทุกฺกฏสฺส. (1)}

\prose{\csromanfolio{178}อิธ ปน ภิกฺขเว อญฺญตรสฺมิํ อาวาเส ตทหุโปสเถ สมฺพหุลา อาวาสิกา ภิกฺขู สนฺนิปตนฺติ จตฺตาโร วา อติเรกา วา, เต ชานนฺติ “อตฺถญฺเญ อาวาสิกา ภิกฺขู อนาคตา”ติ, เต “กปฺปติ นุ โข อมฺหากํ อุโปสโถ กาตุํ, น นุ โข กปฺปตี”ติ เวมติกา อุโปสถํ กโรนฺติ ปาติโมกฺขํ อุทฺทิสนฺติ, เตหิ อุทฺทิสฺสมาเน ปาติโมกฺเข อถญฺเญ อาวาสิกา ภิกฺขู อาคจฺฉนฺติ สมสมา, อุทฺทิฏฺฐํ สุอุทฺทิฏฺฐํ, อวเสสํ โสตพฺพํ. อุทฺเทสกานํ อาปตฺติ ทุกฺกฏสฺส. (2)}

\prose{อิธ ปน ภิกฺขเว อญฺญตรสฺมิํ อาวาเส ตทหุโปสเถ สมฺพหุลา อาวาสิกา ภิกฺขู สนฺนิปตนฺติ จตฺตาโร วา อติเรกา วา, เต ชานนฺติ “อตฺถญฺเญ อาวาสิกา ภิกฺขู อนาคตา”ติ, เต “กปฺปติ นุ โข อมฺหากํ อุโปสโถ กาตุํ, น นุ โข กปฺปตี”ติ เวมติกา อุโปสถํ กโรนฺติ ปาติโมกฺขํ อุทฺทิสนฺติ, เตหิ อุทฺทิสฺสมาเน ปาติโมกฺเข อถญฺเญ อาวาสิกา ภิกฺขู อาคจฺฉนฺติ โถกตรา, อุทฺทิฏฺฐํ สุอุทฺทิฏฺฐํ, อวเสสํ โสตพฺพํ. อุทฺเทสกานํ อาปตฺติ ทุกฺกฏสฺส. (3)}

\prosewithcloser{อิธ ปน ภิกฺขเว อญฺญตรสฺมิํ อาวาเส ตทหุโปสเถ สมฺพหุลา อาวาสิกา ภิกฺขู สนฺนิปตนฺติ จตฺตาโร วา อติเรกา วา, เต ชานนฺติ “อตฺถญฺเญ อาวาสิกา ภิกฺขู อนาคตา”ติ, เต “กปฺปติ นุ โข อมฺหากํ อุโปสโถ กาตุํ, น นุ โข กปฺปตี”ติ เวมติกา อุโปสถํ กโรนฺติ ปาติโมกฺขํ อุทฺทิสนฺติ. เตหิ อุทฺทิฏฺฐมตฺเต ปาติโมกฺเข ฯเปฯ อวุฏฺฐิตาย ปริสาย ฯเปฯ เอกจฺจาย วุฏฺฐิตาย ปริสาย ฯเปฯ สพฺพาย วุฏฺฐิตาย ปริสาย อถญฺเญ อาวาสิกา ภิกฺขู อาคจฺฉนฺติ พหุตรา ฯเปฯ สมสมา ฯเปฯ โถกตรา, อุทฺทิฏฺฐํ สุอุทฺทิฏฺฐํ, เตสํ สนฺติเก ปาริสุทฺธิ อาโรเจตพฺพา. อุทฺเทสกานํ อาปตฺติ ทุกฺกฏสฺส.}{\mbox{(4-15)} เวมติก\-ปนฺนรสกํ นิฏฺฐิตํ.}
\csromanmatikaanchor{mtk.100}
\csromantocmarkref{section}{98. กุกฺกุจฺจปกตปนฺนรสก}{mtk.100}
\bookmark[dest={mtk.100},level=1]{98. กุกฺกุจฺจปกตปนฺนรสก}
\csromanheader{1}{98. \textbf{กุกฺกุจฺจ}\-\textbf{ปกต}\-\textbf{ปนฺนรสก}}
\proseitem{175}{อิธ ปน ภิกฺขเว อญฺญตรสฺมิํ อาวาเส ตทหุโปสเถ สมฺพหุลา อาวาสิกา ภิกฺขู สนฺนิปตนฺติ จตฺตาโร วา อติเรกา วา, เต ชานนฺติ “อตฺถญฺเญ อาวาสิกา ภิกฺขู อนาคตา”ติ, เต “กปฺปเตว \csromanfolio{179}อมฺหากํ อุโปสโถ กาตุํ, นามฺหากํ น กปฺปตี”ติ กุกฺกุจฺจปกตา อุโปสถํ กโรนฺติ ปาติโมกฺขํ อุทฺทิสนฺติ, เตหิ อุทฺทิสฺสมาเน ปาติโมกฺเข อถญฺเญ อาวาสิกา ภิกฺขู อาคจฺฉนฺติ พหุตรา, เตหิ ภิกฺขเว ภิกฺขูหิ ปุน ปาติโมกฺขํ อุทฺทิสิตพฺพํ. อุทฺเทสกานํ อาปตฺติ ทุกฺกฏสฺส. (1)}

\prose{อิธ ปน ภิกฺขเว อญฺญตรสฺมิํ อาวาเส ตทหุโปสเถ สมฺพหุลา อาวาสิกา ภิกฺขู สนฺนิปตนฺติ จตฺตาโร วา อติเรกา วา, เต ชานนฺติ “อตฺถญฺเญ อาวาสิกา ภิกฺขู อนาคตา”ติ, เต “กปฺปเตว อมฺหากํ อุโปสโถ กาตุํ, นามฺหากํ น กปฺปตี”ติ กุกฺกุจฺจปกตา อุโปสถํ กโรนฺติ ปาติโมกฺขํ อุทฺทิสนฺติ, เตหิ อุทฺทิสฺสมาเน ปาติโมกฺเข อถญฺเญ อาวาสิกา ภิกฺขู อาคจฺฉนฺติ สมสมา, อุทฺทิฏฺฐํ สุอุทฺทิฏฺฐํ, อวเสสํ โสตพฺพํ. อุทฺเทสกานํ อาปตฺติ ทุกฺกฏสฺส. (2)}

\prose{อิธ ปน ภิกฺขเว อญฺญตรสฺมิํ อาวาเส ตทหุโปสเถ สมฺพหุลา อาวาสิกา ภิกฺขู สนฺนิปตนฺติ จตฺตาโร วา อติเรกา วา, เต ชานนฺติ “อตฺถญฺเญ อาวาสิกา ภิกฺขู อนาคตา”ติ, เต “กปฺปเตว อมฺหากํ อุโปสโถ กาตุํ, นามฺหากํ น กปฺปตี”ติ กุกฺกุจฺจปกตา อุโปสถํ กโรนฺติ ปาติโมกฺขํ อุทฺทิสนฺติ, เตหิ อุทฺทิสฺสมาเน ปาติโมกฺเข อถญฺเญ อาวาสิกา ภิกฺขู อาคจฺฉนฺติ โถกตรา, อุทฺทิฏฺฐํ สุอุทฺทิฏฺฐํ, อวเสสํ โสตพฺพํ. อุทฺเทสกานํ อาปตฺติ ทุกฺกฏสฺส. (3)}

\prosewithcloser{อิธ ปน ภิกฺขเว อญฺญตรสฺมิํ อาวาเส ตทหุโปสเถ สมฺพหุลา อาวาสิกา ภิกฺขู สนฺนิปตนฺติ จตฺตาโร วา อติเรกา วา, เต ชานนฺติ “อตฺถญฺเญ อาวาสิกา ภิกฺขู อนาคตา”ติ, เต “กปฺปเตว อมฺหากํ อุโปสโถ กาตุํ, นามฺหากํ น กปฺปตี”ติ กุกฺกุจฺจปกตา อุโปสถํ กโรนฺติ ปาติโมกฺขํ อุทฺทิสนฺติ, เตหิ อุทฺทิฏฺฐมตฺเต ปาติโมกฺเข ฯเปฯ อวุฏฺฐิตาย ปริสาย ฯเปฯ เอกจฺจาย วุฏฺฐิตาย ปริสาย ฯเปฯ สพฺพาย วุฏฺฐิตาย ปริสาย อถญฺเญ อาวาสิกา ภิกฺขู อาคจฺฉนฺติ พหุตรา ฯเปฯ สมสมา ฯเปฯ โถกตรา, อุทฺทิฏฺฐํ สุอุทฺทิฏฺฐํ, เตสํ สนฺติเก ปาริสุทฺธิ อาโรเจตพฺพา. อุทฺเทสกานํ อาปตฺติ ทุกฺกฏสฺส.}{\mbox{(4-15)} กุกฺกุจฺจ\-ปกต\-ปนฺนรสกํ นิฏฺฐิตํ.}
\csromanmatikaanchor{mtk.101}
\csromantocmarkref{section}{99. เภทปุเรกฺขารปนฺนรสก}{mtk.101}
\bookmark[dest={mtk.101},level=1]{99. เภทปุเรกฺขารปนฺนรสก}
\csromanheader{1}{99. \textbf{เภท}\-\textbf{ปุเรกฺขาร}\-\textbf{ปนฺนรสก}}
\proseitem{176}{\csromanfolio{180}อิธ ปน ภิกฺขเว อญฺญตรสฺมิํ อาวาเส ตทหุโปสเถ สมฺพหุลา อาวาสิกา ภิกฺขู สนฺนิปตนฺติ จตฺตาโร วา อติเรกา วา, เต ชานนฺติ “อตฺถญฺเญ อาวาสิกา ภิกฺขู อนาคตา”ติ, เต “นสฺสนฺเตเต, วินสฺสนฺเตเต, โก เตหิ อตฺโถ”ติ เภท\-ปุเรกฺขารา อุโปสถํ กโรนฺติ ปาติโมกฺขํ อุทฺทิสนฺติ, เตหิ อุทฺทิสฺสมาเน ปาติโมกฺเข อถญฺเญ อาวาสิกา ภิกฺขู อาคจฺฉนฺติ พหุตรา, เตหิ ภิกฺขเว ภิกฺขูหิ ปุน ปาติโมกฺขํ อุทฺทิสิตพฺพํ. อุทฺเทสกานํ อาปตฺติ ถุลฺลจฺจยสฺส. (1)}

\prose{อิธ ปน ภิกฺขเว อญฺญตรสฺมิํ อาวาเส ตทหุโปสเถ สมฺพหุลา อาวาสิกา ภิกฺขู สนฺนิปตนฺติ จตฺตาโร วา อติเรกา วา, เต ชานนฺติ “อตฺถญฺเญ อาวาสิกา ภิกฺขู อนาคตา”ติ, เต “นสฺสนฺเตเต, วินสฺสนฺเตเต, โก เตหิ อตฺโถ”ติ เภท\-ปุเรกฺขารา อุโปสถํ กโรนฺติ ปาติโมกฺขํ อุทฺทิสนฺติ, เตหิ อุทฺทิสฺสมาเน ปาติโมกฺเข อถญฺเญ อาวาสิกา ภิกฺขู อาคจฺฉนฺติ สมสมา, อุทฺทิฏฺฐํ สุอุทฺทิฏฺฐํ, อวเสสํ โสตพฺพํ. อุทฺเทสกานํ อาปตฺติ ถุลฺลจฺจยสฺส. (2)}

\prose{อิธ ปน ภิกฺขเว อญฺญตรสฺมิํ อาวาเส ตทหุโปสเถ สมฺพหุลา อาวาสิกา ภิกฺขู สนฺนิปตนฺติ จตฺตาโร วา อติเรกา วา, เต ชานนฺติ “อตฺถญฺเญ อาวาสิกา ภิกฺขู อนาคตา”ติ, เต “นสฺสนฺเตเต, วินสฺสนฺเตเต, โก เตหิ อตฺโถ”ติ เภท\-ปุเรกฺขารา อุโปสถํ กโรนฺติ ปาติโมกฺขํ อุทฺทิสนฺติ, เตหิ อุทฺทิสฺสมาเน ปาติโมกฺเข อถญฺเญ อาวาสิกา ภิกฺขู อาคจฺฉนฺติ โถกตรา, อุทฺทิฏฺฐํ สุอุทฺทิฏฺฐํ, อวเสสํ โสตพฺพํ. อุทฺเทสกานํ อาปตฺติ ถุลฺลจฺจยสฺส. (3)}

\prose{อิธ ปน ภิกฺขเว อญฺญตรสฺมิํ อาวาเส ตทหุโปสเถ สมฺพหุลา อาวาสิกา ภิกฺขู สนฺนิปตนฺติ จตฺตาโร วา อติเรกา วา, เต ชานนฺติ “อตฺถญฺเญ อาวาสิกา ภิกฺขู อนาคตา”ติ, เต “นสฺสนฺเตเต, วินสฺสนฺเตเต, โก เตหิ อตฺโถ”ติ เภท\-ปุเรกฺขารา อุโปสถํ กโรนฺติ ปาติโมกฺขํ อุทฺทิสนฺติ, เตหิ อุทฺทิฏฺฐมตฺเต ปาติโมกฺเข อถญฺเญ อาวาสิกา ภิกฺขู อาคจฺฉนฺติ พหุตรา, เตหิ ภิกฺขเว ภิกฺขูหิ ปุน ปาติโมกฺขํ อุทฺทิสิตพฺพํ. อุทฺเทสกานํ อาปตฺติ ถุลฺลจฺจยสฺส. (4)}

\prose{\csromanfolio{181}อิธ ปน ภิกฺขเว อญฺญตรสฺมิํ อาวาเส ตทหุโปสเถ สมฺพหุลา อาวาสิกา ภิกฺขู สนฺนิปตนฺติ จตฺตาโร วา อติเรกา วา, เต ชานนฺติ “อตฺถญฺเญ อาวาสิกา ภิกฺขู อนาคตา”ติ, เต “นสฺสนฺเตเต, วินสฺสนฺเตเต, โก เตหิ อตฺโถ”ติ เภท\-ปุเรกฺขารา อุโปสถํ กโรนฺติ ปาติโมกฺขํ อุทฺทิสนฺติ, เตหิ อุทฺทิฏฺฐมตฺเต ปาติโมกฺเข อถญฺเญ อาวาสิกา ภิกฺขู อาคจฺฉนฺติ สมสมา, อุทฺทิฏฺฐํ สุอุทฺทิฏฺฐํ, เตสํ สนฺติเก ปาริสุทฺธิ อาโรเจตพฺพา. อุทฺเทสกานํ อาปตฺติ ถุลฺลจฺจยสฺส. (5)}

\prose{อิธ ปน ภิกฺขเว อญฺญตรสฺมิํ อาวาเส ตทหุโปสเถ สมฺพหุลา อาวาสิกา ภิกฺขู สนฺนิปตนฺติ จตฺตาโร วา อติเรกา วา, เต ชานนฺติ “อตฺถญฺเญ อาวาสิกา ภิกฺขู อนาคตา”ติ, เต “นสฺสนฺเตเต, วินสฺสนฺเตเต, โก เตหิ อตฺโถ”ติ เภท\-ปุเรกฺขารา อุโปสถํ กโรนฺติ ปาติโมกฺขํ อุทฺทิสนฺติ, เตหิ อุทฺทิฏฺฐมตฺเต ปาติโมกฺเข อถญฺเญ อาวาสิกา ภิกฺขู อาคจฺฉนฺติ โถกตรา, อุทฺทิฏฺฐํ สุอุทฺทิฏฺฐํ, เตสํ สนฺติเก ปาริสุทฺธิ อาโรเจตพฺพา. อุทฺเทสกานํ อาปตฺติ ถุลฺลจฺจยสฺส. (6)}

\prose{อิธ ปน ภิกฺขเว อญฺญตรสฺมิํ อาวาเส ตทหุโปสเถ สมฺพหุลา อาวาสิกา ภิกฺขู สนฺนิปตนฺติ จตฺตาโร วา อติเรกา วา, เต ชานนฺติ “อตฺถญฺเญ อาวาสิกา ภิกฺขู อนาคตา”ติ, เต “นสฺสนฺเตเต, วินสฺสนฺเตเต, โก เตหิ อตฺโถ”ติ เภท\-ปุเรกฺขารา อุโปสถํ กโรนฺติ ปาติโมกฺขํ อุทฺทิสนฺติ, เตหิ อุทฺทิฏฺฐมตฺเต ปาติโมกฺเข อวุฏฺฐิตาย ปริสาย อถญฺเญ อาวาสิกา ภิกฺขู อาคจฺฉนฺติ พหุตรา, เตหิ ภิกฺขเว ภิกฺขูหิ ปุน ปาติโมกฺขํ อุทฺทิสิตพฺพํ. อุทฺเทสกานํ อาปตฺติ ถุลฺลจฺจยสฺส. (7)}

\prose{อิธ ปน ภิกฺขเว อญฺญตรสฺมิํ อาวาเส ตทหุโปสเถ สมฺพหุลา อาวาสิกา ภิกฺขู สนฺนิปตนฺติ จตฺตาโร วา อติเรกา วา, เต ชานนฺติ “อตฺถญฺเญ อาวาสิกา ภิกฺขู อนาคตา”ติ, เต “นสฺสนฺเตเต, วินสฺสนฺเตเต, โก เตหิ อตฺโถ”ติ เภท\-ปุเรกฺขารา อุโปสถํ กโรนฺติ ปาติโมกฺขํ อุทฺทิสนฺติ, เตหิ อุทฺทิฏฺฐมตฺเต ปาติโมกฺเข อวุฏฺฐิตาย ปริสาย อถญฺเญ อาวาสิกา ภิกฺขู อาคจฺฉนฺติ สมสมา, อุทฺทิฏฺฐํ สุอุทฺทิฏฺฐํ, เตสํ สนฺติเก ปาริสุทฺธิ อาโรเจตพฺพา. อุทฺเทสกานํ อาปตฺติ ถุลฺลจฺจยสฺส. (8)}

\prose{\csromanfolio{182}อิธ ปน ภิกฺขเว อญฺญตรสฺมิํ อาวาเส ตทหุโปสเถ สมฺพหุลา อาวาสิกา ภิกฺขู สนฺนิปตนฺติ จตฺตาโร วา อติเรกา วา, เต ชานนฺติ “อตฺถญฺเญ อาวาสิกา ภิกฺขู อนาคตา”ติ, เต “นสฺสนฺเตเต, วินสฺสนฺเตเต, โก เตหิ อตฺโถ”ติ เภท\-ปุเรกฺขารา อุโปสถํ กโรนฺติ ปาติโมกฺขํ อุทฺทิสนฺติ, เตหิ อุทฺทิฏฺฐมตฺเต ปาติโมกฺเข อวุฏฺฐิตาย ปริสาย อถญฺเญ อาวาสิกา ภิกฺขู อาคจฺฉนฺติ โถกตรา, อุทฺทิฏฺฐํ สุอุทฺทิฏฺฐํ, เตสํ สนฺติเก ปาริสุทฺธิ อาโรเจตพฺพา. อุทฺเทสกานํ อาปตฺติ ถุลฺลจฺจยสฺส. (9)}

\prose{อิธ ปน ภิกฺขเว อญฺญตรสฺมิํ อาวาเส ตทหุโปสเถ สมฺพหุลา อาวาสิกา ภิกฺขู สนฺนิปตนฺติ จตฺตาโร วา อติเรกา วา, เต ชานนฺติ “อตฺถญฺเญ อาวาสิกา ภิกฺขู อนาคตา”ติ, เต “นสฺสนฺเตเต, วินสฺสนฺเตเต, โก เตหิ อตฺโถ”ติ เภท\-ปุเรกฺขารา อุโปสถํ กโรนฺติ ปาติโมกฺขํ อุทฺทิสนฺติ, เตหิ อุทฺทิฏฺฐมตฺเต ปาติโมกฺเข เอกจฺจาย วุฏฺฐิตาย ปริสาย อถญฺเญ อาวาสิกา ภิกฺขู อาคจฺฉนฺติ พหุตรา, เตหิ ภิกฺขเว ภิกฺขูหิ ปุน ปาติโมกฺขํ อุทฺทิสิตพฺพํ. อุทฺเทสกานํ อาปตฺติ ถุลฺลจฺจยสฺส. (10)}

\prose{อิธ ปน ภิกฺขเว อญฺญตรสฺมิํ อาวาเส ตทหุโปสเถ สมฺพหุลา อาวาสิกา ภิกฺขู สนฺนิปตนฺติ จตฺตาโร วา อติเรกา วา, เต ชานนฺติ “อตฺถญฺเญ อาวาสิกา ภิกฺขู อนาคตา”ติ, เต “นสฺสนฺเตเต, วินสฺสนฺเตเต, โก เตหิ อตฺโถ”ติ เภท\-ปุเรกฺขารา อุโปสถํ กโรนฺติ ปาติโมกฺขํ อุทฺทิสนฺติ, เตหิ อุทฺทิฏฺฐมตฺเต ปาติโมกฺเข เอกจฺจาย วุฏฺฐิตาย ปาริสาย อถญฺเญ อาวาสิกา ภิกฺขู อาคจฺฉนฺติ สมสมา, อุทฺทิฏฺฐํ สุอุทฺทิฏฺฐํ, เตสํ สนฺติเก ปาริสุทฺธิ อาโรเจตพฺพา. อุทฺเทสกานํ อาปตฺติ ถุลฺลจฺจยสฺส. (11)}

\prose{อิธ ปน ภิกฺขเว อญฺญตรสฺมิํ อาวาเส ตทหุโปสเถ สมฺพหุลา อาวาสิกา ภิกฺขู สนฺนิปตนฺติ จตฺตาโร วา อติเรกา วา, เต ชานนฺติ “อตฺถญฺเญ อาวาสิกา ภิกฺขู อนาคตา”ติ, เต “นสฺสนฺเตเต, วินสฺสนฺเตเต, โก เตหิ อตฺโถ”ติ เภท\-ปุเรกฺขารา อุโปสถํ กโรนฺติ ปาติโมกฺขํ อุทฺทิสนฺติ, เตหิ อุทฺทิฏฺฐมตฺเต ปาติโมกฺเข เอกจฺจาย วุฏฺฐิตาย ปริสาย อถญฺเญ อาวาสิกา ภิกฺขู อาคจฺฉนฺติ โถกตรา, \csromanfolio{183}อุทฺทิฏฺฐํ สุอุทฺทิฏฺฐํ, เตสํ สนฺติเก ปาริสุทฺธิ อาโรเจตพฺพา. อุทฺเทสกานํ อาปตฺติ ถุลฺลจฺจยสฺส. (12)}

\prose{อิธ ปน ภิกฺขเว อญฺญตรสฺมิํ อาวาเส ตทหุโปสเถ สมฺพหุลา อาวาสิกา ภิกฺขู สนฺนิปตนฺติ จตฺตาโร วา อติเรกา วา, เต ชานนฺติ “อตฺถญฺเญ อาวาสิกา ภิกฺขู อนาคตา”ติ, เต “นสฺสนฺเตเต, วินสฺสนฺเตเต, โก เตหิ อตฺโถ”ติ เภท\-ปุเรกฺขารา อุโปสถํ กโรนฺติ ปาติโมกฺขํ อุทฺทิสนฺติ, เตหิ อุทฺทิฏฺฐมตฺเต ปาติโมกฺเข สพฺพาย วุฏฺฐิตาย ปริสาย อถญฺเญ อาวาสิกา ภิกฺขู อาคจฺฉนฺติ พหุตรา, เตหิ ภิกฺขเว ภิกฺขูหิ ปุน ปาติโมกฺขํ อุทฺทิสิตพฺพํ. อุทฺเทสกานํ อาปตฺติ ถุลฺลจฺจยสฺส. (13)}

\prose{อิธ ปน ภิกฺขเว อญฺญตรสฺมิํ อาวาเส ตทหุโปสเถ สมฺพหุลา อาวาสิกา ภิกฺขู สนฺนิปตนฺติ จตฺตาโร วา อติเรกา วา, เต ชานนฺติ “อตฺถญฺเญ อาวาสิกา ภิกฺขู อนาคตา”ติ, เต “นสฺสนฺเตเต, วินสฺสนฺเตเต, โก เตหิ อตฺโถ”ติ เภท\-ปุเรกฺขารา อุโปสถํ กโรนฺติ ปาติโมกฺขํ อุทฺทิสนฺติ, เตหิ อุทฺทิฏฺฐมตฺเต ปาติโมกฺเข สพฺพาย วุฏฺฐิตาย ปริสาย อถญฺเญ อาวาสิกา ภิกฺขู อาคจฺฉนฺติ สมสมา, อุทฺทิฏฺฐํ สุอุทฺทิฏฺฐํ, เตสํ สนฺติเก ปาริสุทฺธิ อาโรเจตพฺพา. อุทฺเทสกานํ อาปตฺติ ถุลฺลจฺจยสฺส. (14)}

\prosewithcloser{อิธ ปน ภิกฺขเว อญฺญตรสฺมิํ อาวาเส ตทหุโปสเถ สมฺพหุลา อาวาสิกา ภิกฺขู สนฺนิปตนฺติ จตฺตาโร วา อติเรกา วา, เต ชานนฺติ “อตฺถญฺเญ อาวาสิกา ภิกฺขู อนาคตา”ติ, เต “นสฺสนฺเตเต, วินสฺสนฺเตเต, โก เตหิ อตฺโถ”ติ เภท\-ปุเรกฺขารา อุโปสถํ กโรนฺติ ปาติโมกฺขํ อุทฺทิสนฺติ, เตหิ อุทฺทิฏฺฐมตฺเต ปาติโมกฺเข สพฺพาย วุฏฺฐิตาย ปริสาย อถญฺเญ อาวาสิกา ภิกฺขู อาคจฺฉนฺติ โถกตรา, อุทฺทิฏฺฐํ สุอุทฺทิฏฺฐํ, เตสํ สนฺติเก ปาริสุทฺธิ อาโรเจตพฺพา. อุทฺเทสกานํ อาปตฺติ ถุลฺลจฺจยสฺส. (15)}{เภท\-ปุเรกฺขาร\-ปนฺนรสกํ นิฏฺฐิตํ.}
\nitthitam{ปญฺจวีสติกา นิฏฺฐิตา.}
\csromanmatikaanchor{mtk.102}
\csromantocmarkref{section}{100. สีโมกฺกนฺติกเปยฺยาล}{mtk.102}
\bookmark[dest={mtk.102},level=1]{100. สีโมกฺกนฺติกเปยฺยาล}
\csromanheader{1}{100. \textbf{สีโม}\-\textbf{กฺกนฺติก}\-\textbf{เปยฺยาล}}
\proseitem{177}{\csromanfolio{184}อิธ ปน ภิกฺขเว อญฺญตรสฺมิํ อาวาเส ตทหุโปสเถ สมฺพหุลา อาวาสิกา ภิกฺขู สนฺนิปตนฺติ จตฺตาโร วา อติเรกา วา, เต น ชานนฺติ “อญฺเญ อาวาสิกา ภิกฺขู อนฺโตสีมํ โอกฺกมนฺตี”ติ ฯเปฯ เต น ชานนฺติ “อญฺเญ อาวาสิกา ภิกฺขู อนฺโตสีมํ โอกฺกนฺตา”ติ ฯเปฯ เต น ปสฺสนฺติ อญฺเญ อาวาสิเก ภิกฺขู อนฺโตสีมํ โอกฺกมนฺเต ฯเปฯ เต น ปสฺสนฺติ อญฺเญ อาวาสิเก ภิกฺขู อนฺโตสีมํ โอกฺกนฺเต ฯเปฯ เต น สุณนฺติ “อญฺเญ อาวาสิกา ภิกฺขู อนฺโตสีมํ โอกฺกมนฺตี”ติ ฯเปฯ เต น สุณนฺติ “อญฺเญ อาวาสิกา ภิกฺขู อนฺโตสีมํ โอกฺกนฺตา”ติ ฯเปฯ. อาวาสิเกน อาวาสิกา เอกสต\-ปญฺจสตฺตติ ติกนยโต, อาวาสิเกน อาคนฺตุกา, อาคนฺตุเกน อาวาสิกา, อาคนฺตุเกน อาคนฺตุกา, เปยฺยาลมุเขน สตฺต ติกสตานิ โหนฺติ.}

\proseitem{178}{อิธ ปน ภิกฺขเว อาวาสิกานํ ภิกฺขูนํ จาตุทฺทโส โหติ อาคนฺตุกานํ ปนฺนรโส, สเจ อาวาสิกา พหุตรา โหนฺติ, อาคนฺตุเกหิ อาวาสิกานํ อนุวตฺติตพฺพํ.}

\setlength{\csromangathapagewidth}{0pt}
\setlength{\csromangathawakwidth}{0pt}
\csromangathameasuregroup{}{สเจ สมสมา โหนฺติ, \\ อาคนฺตุเกหิ อาวาสิกานํ อนุวตฺติตพฺพํ. \\ สเจ อาคนฺตุกา พหุตรา โหนฺติ, \\ อาวาสิเกหิ อาคนฺตุกานํ อนุวตฺติตพฺพํ.}
\setlength{\csromangathaleftcol}{0pt}
\csromangathagroup{\csromangathastanza{\csromangathaitemline{178}{สเจ สมสมา โหนฺติ,} \\ \csromangathacontline{อาคนฺตุเกหิ อาวาสิกานํ อนุวตฺติตพฺพํ.} \\ \csromangathacontline{สเจ อาคนฺตุกา พหุตรา โหนฺติ,} \\ \csromangathacontline{อาวาสิเกหิ อาคนฺตุกานํ อนุวตฺติตพฺพํ.}}}

\prose{อิธ ปน ภิกฺขเว อาวาสิกานํ ภิกฺขูนํ ปนฺนรโส โหติ อาคนฺตุกานํ จาตุทฺทโส, สเจ อาวาสิกา พหุตรา โหนฺติ, อาคนฺตุเกหิ อาวาสิกานํ อนุวตฺติตพฺพํ.}

\setlength{\csromangathapagewidth}{0pt}
\setlength{\csromangathawakwidth}{0pt}
\csromangathameasuregroup{}{สเจ สมสมา โหนฺติ, \\ อาคนฺตุเกหิ อาวาสิกานํ อนุวตฺติตพฺพํ. \\ สเจ อาคนฺตุกา พหุตรา โหนฺติ, \\ อาวาสิเกหิ อาคนฺตุกานํ อนุวตฺติตพฺพํ.}
\setlength{\csromangathaleftcol}{0pt}
\csromangathagroup{\csromangathastanza{สเจ สมสมา โหนฺติ, \\ อาคนฺตุเกหิ อาวาสิกานํ อนุวตฺติตพฺพํ. \\ สเจ อาคนฺตุกา พหุตรา โหนฺติ, \\ อาวาสิเกหิ อาคนฺตุกานํ อนุวตฺติตพฺพํ.}}

\proseitem{178}{อิธ ปน ภิกฺขเว อาวาสิกานํ ภิกฺขูนํ ปาฏิปโท โหติ อาคนฺตุกานํ ปนฺนรโส, สเจ อาวาสิกา พหุตรา โหนฺติ, อาวาสิเกหิ อาคนฺตุกานํ นากามา ทาตพฺพา สามคฺคี, อาคนฺตุเกหิ นิสฺสีมํ คนฺตฺวา อุโปสโถ กาตพฺโพ. สเจ สมสมา โหนฺติ, อาวาสิเกหิ อาคนฺตุกานํ นากามา ทาตพฺพา สามคฺคี, อาคนฺตุเกหิ \csromanfolio{185}นิสฺสีมํ คนฺตฺวา อุโปสโถ กาตพฺโพ. สเจ อาคนฺตุกา พหุตรา โหนฺติ, อาวาสิเกหิ อาคนฺตุกานํ สามคฺคี วา ทาตพฺพา นิสฺสีมํ วา คนฺตพฺพํ.}

\prose{อิธ ปน ภิกฺขเว อาวาสิกานํ ภิกฺขูนํ ปนฺนรโส โหติ อาคนฺตุกานํ ปาฏิปโท, สเจ อาวาสิกา พหุตรา โหนฺติ, อาคนฺตุเกหิ อาวาสิกานํ สามคฺคี วา ทาตพฺพา นิสฺสีมํ วา คนฺตพฺพํ. สเจ สมสมา โหนฺติ, อาคนฺตุเกหิ อาวาสิกานํ สามคฺคี วา ทาตพฺพา นิสฺสีมํ วา คนฺตพฺพํ. สเจ อาคนฺตุกา พหุตรา โหนฺติ, อาคนฺตุเกหิ อาวาสิกานํ นากามา ทาตพฺพา สามคฺคี, อาวาสิเกหิ นิสฺสีมํ คนฺตฺวา อุโปสโถ กาตพฺโพ.}

\csromanmatikaanchor{mtk.103}
\csromantocmarkref{section}{101. ลิงฺคาทิทสฺสน}{mtk.103}
\bookmark[dest={mtk.103},level=1]{101. ลิงฺคาทิทสฺสน}
\csromanheader{1}{101. \textbf{ลิงฺคาทิทสฺสน}}
\proseitem{179}{อิธ ปน ภิกฺขเว อาคนฺตุกา ภิกฺขู ปสฺสนฺติ อาวาสิกานํ ภิกฺขูนํ อาวาสิกาการํ อาวาสิกลิงฺคํ อาวาสิก\-นิมิตฺตํ อาวาสิกุทฺเทสํ สุปญฺญตฺตํ มญฺจปีฐํ ภิสิพิพฺโพหนํ ปานียํ ปริโภชนียํ สูปฏฺฐิตํ ปริเวณํ สุสมฺมฏฺฐํ, ปสฺสิตฺวา เวมติกา โหนฺติ “อตฺถิ นุ โข อาวาสิกา ภิกฺขู นตฺถิ นุ โข”ติ, เต เวมติกา น วิจินนฺติ, อวิจินิตฺวา อุโปสถํ กโรนฺติ, อาปตฺติ ทุกฺกฏสฺส. เต เวมติกา วิจินนฺติ, วิจินิตฺวา น ปสฺสนฺติ, อปสฺสิตฺวา อุโปสถํ กโรนฺติ, อนาปตฺติ. เต เวมติกา วิจินนฺติ วิจินิตฺวา ปสฺสนฺติ, ปสฺสิตฺวา เอกโต อุโปสถํ กโรนฺติ, อนาปตฺติ. เต เวมติกา วิจินนฺติ, วิจินิตฺวา ปสฺสนฺติ, ปสฺสิตฺวา ปาเฏกฺกํ อุโปสถํ กโรนฺติ, อาปตฺติ ทุกฺกฏสฺส. เต เวมติกา วิจินนฺติ, วิจินิตฺวา ปสฺสนฺติ, ปสฺสิตฺวา “นสฺสนฺเตเต, วินสฺสนฺเตเต, โก เตหิ อตฺโถ”ติ เภท\-ปุเรกฺขารา อุโปสถํ กโรนฺติ, อาปตฺติ ถุลฺลจฺจยสฺส.}

\prose{อิธ ปน ภิกฺขเว อาคนฺตุกา ภิกฺขู สุณนฺติ อาวาสิกานํ ภิกฺขูนํ อาวาสิกาการํ อาวาสิกลิงฺคํ อาวาสิก\-นิมิตฺตํ อาวาสิกุทฺเทสํ จงฺกมนฺตานํ ปทสทฺทํ สชฺฌายสทฺทํ อุกฺกาสิตสทฺทํ ขิปิตสทฺทํ, สุตฺวา เวมติกา โหนฺติ “อตฺถิ นุ โข อาวาสิกา ภิกฺขู นตฺถิ นุ โข”ติ, เต เวมติกา น วิจินนฺติ, อวิจินิตฺวา อุโปสถํ กโรนฺติ, อาปตฺติ ทุกฺกฏสฺส. เต เวมติกา วิจินนฺติ, วิจินิตฺวา น ปสฺสนฺติ, อปสฺสิตฺวา \csromanfolio{186}\textbf{อุโปสถํ กโรนฺติ, อนาปตฺติ. เต เวมติกา วิจินนฺติ, วิจินิตฺวา ปสฺสนฺติ,} ปสฺสิตฺวา เอกโต อุโปสถํ กโรนฺติ, อนาปตฺติ. เต เวมติกา วิจินนฺติ, วิจินิตฺวา ปสฺสนฺติ, ปสฺสิตฺวา ปาเฏกฺกํ อุโปสถํ กโรนฺติ, อาปตฺติ ทุกฺกฏสฺส. เต เวมติกา วิจินนฺติ, วิจินิตฺวา ปสฺสนฺติ, ปสฺสิตฺวา “นสฺสนฺเตเต, วินสฺสนฺเตเต, โก เตหิ อตฺโถ”ติ เภท\-ปุเรกฺขารา อุโปสถํ กโรนฺติ, อาปตฺติ ถุลฺลจฺจยสฺส.}

\prose{อิธ ปน ภิกฺขเว อาวาสิกา ภิกฺขู ปสฺสนฺติ อาคนฺตุกานํ ภิกฺขูนํ อาคนฺตุกาการํ อาคนฺตุกลิงฺคํ อาคนฺตุก\-นิมิตฺตํ อาคนฺตุกุทฺเทสํ อญฺญาตกํ ปตฺตํ อญฺญาตกํ จีวรํ อญฺญาตกํ นิสีทนํ ปาทานํ โธตํ อุทกนิสฺเสกํ, ปสฺสิตฺวา เวมติกา โหนฺติ “อตฺถิ นุ โข อาคนฺตุกา ภิกฺขู นตฺถิ นุ โข”ติ, เต เวมติกา น วิจินนฺติ, อวิจินิตฺวา อุโปสถํ กโรนฺติ, อาปตฺติ ทุกฺกฏสฺส.}

\prose{เต เวมติกา วิจินนฺติ, วิจินิตฺวา น ปสฺสนฺติ, อปสฺสิตฺวา อุโปสถํ กโรนฺติ, อนาปตฺติ.}

\prose{เต เวมติกา วิจินนฺติ, วิจินิตฺวา ปสฺสนฺติ, ปสฺสิตฺวา เอกโต อุโปสถํ กโรนฺติ, อนาปตฺติ.}

\prose{เต เวมติกา วิจินนฺติ, วิจินิตฺวา ปสฺสนฺติ, ปสฺสิตฺวา ปาเฏกฺกํ อุโปสถํ กโรนฺติ, อาปตฺติ ทุกฺกฏสฺส.}

\prose{เต เวมติกา วิจินนฺติ, วิจินิตฺวา ปสฺสนฺติ, ปสฺสิตฺวา “นสฺสนฺเตเต, วินสฺสนฺเตเต, โก เตหิ อตฺโถ”ติ เภท\-ปุเรกฺขารา อุโปสถํ กโรนฺติ, อาปตฺติ ถุลฺลจฺจยสฺส.}

\prose{อิธ ปน ภิกฺขเว อาวาสิกา ภิกฺขู สุณนฺติ อาคนฺตุกานํ ภิกฺขูนํ อาคนฺตุกาการํ อาคนฺตุกลิงฺคํ อาคนฺตุก\-นิมิตฺตํ อาคนฺตุกุทฺเทสํ อาคจฺฉนฺตานํ ปทสทฺทํ อุปาหน\-ปปฺโผฏน\-สทฺทํ อุกฺกาสิตสทฺทํ ขิปิตสทฺทํ, สุตฺวา เวมติกา โหนฺติ “อตฺถิ นุ โข อาคนฺตุกา ภิกฺขู นตฺถิ นุ โข”ติ, เต เวมติกา น วิจินนฺติ, อวิจินิตฺวา อุโปสถํ กโรนฺติ, อาปตฺติ ทุกฺกฏสฺส.}

\prose{เต เวมติกา วิจินนฺติ, วิจินิตฺวา น ปสฺสนฺติ, อปสฺสิตฺวา อุโปสถํ กโรนฺติ, อนาปตฺติ.}

\prose{เต เวมติกา วิจินนฺติ, วิจินิตฺวา ปสฺสนฺติ, ปสฺสิตฺวา เอกโต อุโปสถํ กโรนฺติ, อนาปตฺติ.}

\prose{เต เวมติกา วิจินนฺติ, วิจินิตฺวา ปสฺสนฺติ, ปสฺสิตฺวา ปาเฏกฺกํ อุโปสถํ กโรนฺติ, อาปตฺติ ทุกฺกฏสฺส.}

\prose{เต เวมติกา วิจินนฺติ, วิจินิตฺวา ปสฺสนฺติ, ปสฺสิตฺวา “นสฺสนฺเตเต, วินสฺสนฺเตเต, โก เตหิ อตฺโถ”ติ เภท\-ปุเรกฺขารา อุโปสถํ กโรนฺติ, อาปตฺติ ถุลฺลจฺจยสฺส.}

\csromanmatikaanchor{mtk.104}
\csromantocmarkref{section}{102. นานาสํวาสกาทีหิ อุโปสถกรณํ}{mtk.104}
\bookmark[dest={mtk.104},level=1]{102. นานาสํวาสกาทีหิ อุโปสถกรณํ}
\csromanheader{1}{102. \textbf{นานาสํวาสกาทีหิ อุโปสถ}\-\textbf{กรณํ}}
\csromanmatikaanchor{mtk.105}
\csromantocmarkref{section}{103. น คนฺตพฺพวาร}{mtk.105}
\bookmark[dest={mtk.105},level=1]{103. น คนฺตพฺพวาร}
\proseitem{180}{\csromanfolio{187}อิธ ปน ภิกฺขเว อาคนฺตุกา ภิกฺขู ปสฺสนฺติ อาวาสิเก ภิกฺขู นานาสํวาสเก, เต สมานสํวาสก\-ทิฏฺฐิํ ปฏิลภนฺติ, สมานสํวาสก\-ทิฏฺฐิํ ปฏิลภิตฺวา น ปุจฺฉนฺติ, อปุจฺฉิตฺวา เอกโต อุโปสถํ กโรนฺติ, อนาปตฺติ.}

\prose{เต ปุจฺฉนฺติ, ปุจฺฉิตฺวา นาภิวิตรนฺติ, อนภิวิตริตฺวา เอกโต อุโปสถํ กโรนฺติ, อาปตฺติ ทุกฺกฏสฺส.}

\prose{เต ปุจฺฉนฺติ, ปุจฺฉิตฺวา นาภิวิตรนฺติ, อนภิวิตริตฺวา ปาเฏกฺกํ อุโปสถํ กโรนฺติ, อนาปตฺติ.}

\prose{อิธ ปน ภิกฺขเว อาคนฺตุกา ภิกฺขู ปสฺสนฺติ อาวาสิเก ภิกฺขู สมาน\-สํวาสเก, เต นานา\-สํวาสก\-ทิฏฺฐิํ ปฏิลภนฺติ, นานา\-สํวาสก\-ทิฏฺฐิํ ปฏิลภิตฺวา น ปุจฺฉนฺติ, อปุจฺฉิตฺวา เอกโต อุโปสถํ กโรนฺติ, อาปตฺติ ทุกฺกฏสฺส.}

\prose{เต ปุจฺฉนฺติ, ปุจฺฉิตฺวา อภิวิตรนฺติ, อภิวิตริตฺวา ปาเฏกฺกํ อุโปสถํ กโรนฺติ, อาปตฺติ ทุกฺกฏสฺส.}

\prose{เต ปุจฺฉนฺติ, ปุจฺฉิตฺวา อภิวิตรนฺติ, อภิวิตริตฺวา เอกโต อุโปสถํ กโรนฺติ, อนาปตฺติ.}

\prose{อิธ ปน ภิกฺขเว อาวาสิกา ภิกฺขู ปสฺสนฺติ อาคนฺตุเก ภิกฺขู นานาสํวาสเก, เต สมานสํวาสก\-ทิฏฺฐิํ ปฏิลภนฺติ, สมานสํวาสก\-ทิฏฺฐิํ ปฏิลภิตฺวา น ปุจฺฉนฺติ, อปุจฺฉิตฺวา เอกโต อุโปสถํ กโรนฺติ, อนาปตฺติ.}

\prose{เต ปุจฺฉนฺติ, ปุจฺฉิตฺวา นาภิวิตรนฺติ, อนภิวิตริตฺวา เอกโต อุโปสถํ กโรนฺติ, อาปตฺติ ทุกฺกฏสฺส.}

\prose{เต ปุจฺฉนฺติ, ปุจฺฉิตฺวา นาภิวิตรนฺติ, อนภิวิตริตฺวา ปาเฏกฺกํ อุโปสถํ กโรนฺติ, อนาปตฺติ.}

\prose{อิธ ปน ภิกฺขเว อาวาสิกา ภิกฺขู ปสฺสนฺติ อาคนฺตุเก ภิกฺขู สมาน\-สํวาสเก, เต นานา\-สํวาสก\-ทิฏฺฐิํ ปฏิลภนฺติ, นานา\-สํวาสก\-ทิฏฺฐิํ ปฏิลภิตฺวา น ปุจฺฉนฺติ, อปุจฺฉิตฺวา เอกโต อุโปสถํ กโรนฺติ, อาปตฺติ ทุกฺกฏสฺส.}

\prose{เต ปุจฺฉนฺติ, ปุจฺฉิตฺวา อภิวิตรนฺติ, อภิวิตริตฺวา ปาเฏกฺกํ อุโปสถํ กโรนฺติ, อาปตฺติ ทุกฺกฏสฺส.}

\prose{เต ปุจฺฉนฺติ, ปุจฺฉิตฺวา อภิวิตรนฺติ, อภิวิตริตฺวา เอกโต อุโปสถํ กโรนฺติ, อนาปตฺติ.}

\csromanheader{2}{103. \textbf{นคนฺตพฺพวาร}}
\proseitem{181}{น ภิกฺขเว ตทหุโปสเถ สภิกฺขุกา อาวาสา อภิกฺขุโก อาวาโส คนฺตพฺโพ อญฺญตฺร สํเฆน อญฺญตฺร อนฺตรายา. น \csromanfolio{188}ภิกฺขเว ตทหุโปสเถ สภิกฺขุกา อาวาสา อภิกฺขุโก อนาวาโส คนฺตพฺโพ อญฺญตฺร สํเฆน อญฺญตฺร อนฺตรายา. น ภิกฺขเว ตทหุโปสเถ สภิกฺขุกา อาวาสา อภิกฺขุโก อาวาโส วา อนาวาโส วา คนฺตพฺโพ อญฺญตฺร สํเฆน อญฺญตฺร อนฺตรายา.}

\prose{น ภิกฺขเว ตทหุโปสเถ สภิกฺขุกา อนาวาสา อภิกฺขุโก อาวาโส คนฺตพฺโพ อญฺญตฺร สํเฆน อญฺญตฺร อนฺตรายา. น ภิกฺขเว ตทหุโปสเถ สภิกฺขุกา อนาวาสา อภิกฺขุโก อนาวาโส คนฺตพฺโพ อญฺญตฺร สํเฆน อญฺญตฺร อนฺตรายา. น ภิกฺขเว ตทหุโปสเถ สภิกฺขุกา อนาวาสา อภิกฺขุโก อาวาโส วา อนาวาโส วา คนฺตพฺโพ อญฺญตฺร สํเฆน อญฺญตฺร อนฺตรายา.}

\prose{น ภิกฺขเว ตทหุโปสเถ สภิกฺขุกา อาวาสา วา อนาวาสา วา อภิกฺขุโก อาวาโส คนฺตพฺโพ อญฺญตฺร สํเฆน อญฺญตฺร อนฺตรายา. น ภิกฺขเว ตทหุโปสเถ สภิกฺขุกา อาวาสา วา อนาวาสา วา อภิกฺขุโก อนาวาโส คนฺตพฺโพ อญฺญตฺร สํเฆน อญฺญตฺร อนฺตรายา. น ภิกฺขเว ตทหุโปสเถ สภิกฺขุกา อาวาสา วา อนาวาสา วา อภิกฺขุโก อาวาโส วา อนาวาโส วา คนฺตพฺโพ อญฺญตฺร สํเฆน อญฺญตฺร อนฺตรายา.}

\prose{น ภิกฺขเว ตทหุโปสเถ สภิกฺขุกา อาวาสา สภิกฺขุโก อาวาโส คนฺตพฺโพ ยตฺถสฺสุ ภิกฺขู นานาสํวาสกา อญฺญตฺร สํเฆน อญฺญตฺร อนฺตรายา. น ภิกฺขเว ตทหุโปสเถ สภิกฺขุกา อาวาสา สภิกฺขุโก อนาวาโส คนฺตพฺโพ ยตฺถสฺสุ ภิกฺขู นานาสํวาสกา อญฺญตฺร สํเฆน อญฺญตฺร อนฺตรายา. น ภิกฺขเว ตทหุโปสเถ สภิกฺขุกา อาวาสา สภิกฺขุโก อาวาโส วา อนาวาโส วา คนฺตพฺโพ ยตฺถสฺสุ ภิกฺขู นานาสํวาสกา อญฺญตฺร สํเฆน อญฺญตฺร อนฺตรายา.}

\prose{น ภิกฺขเว ตทหุโปสเถ สภิกฺขุกา อนาวาสา สภิกฺขุโก อาวาโส คนฺตพฺโพ ยตฺถสฺสุ ภิกฺขู นานาสํวาสกา อญฺญตฺร สํเฆน อญฺญตฺร อนฺตรายา. น ภิกฺขเว ตทหุโปสเถ สภิกฺขุกา อนาวาสา สภิกฺขุโก อนาวาโส คนฺตพฺโพ ยตฺถสฺสุ ภิกฺขู นานาสํวาสกา อญฺญตฺร สํเฆน อญฺญตฺร อนฺตรายา. น ภิกฺขเว \csromanfolio{189}ตทหุโปสเถ สภิกฺขุกา อนาวาสา สภิกฺขุโก อาวาโส วา อนาวาโส วา คนฺตพฺโพ ยตฺถสฺสุ ภิกฺขู นานาสํวาสกา อญฺญตฺร สํเฆน อญฺญตฺร อนฺตรายา.}

\prose{น ภิกฺขเว ตทหุโปสเถ สภิกฺขุกา อาวาสา วา อนาวาสา วา สภิกฺขุโก อาวาโส คนฺตพฺโพ ยตฺถสฺสุ ภิกฺขู นานาสํวาสกา อญฺญตฺร สํเฆน อญฺญตฺร อนฺตรายา. น ภิกฺขเว ตทหุโปสเถ สภิกฺขุกา อาวาสา วา อนาวาสา วา สภิกฺขุโก อนาวาโส คนฺตพฺโพ ยตฺถสฺสุ ภิกฺขู นานาสํวาสกา อญฺญตฺร สํเฆน อญฺญตฺร อนฺตรายา. น ภิกฺขเว ตทหุโปสเถ สภิกฺขุกา อาวาสา วา อนาวาสา วา สภิกฺขุโก อาวาโส วา อนาวาโส วา คนฺตพฺโพ ยตฺถสฺสุ ภิกฺขู นานาสํวาสกา อญฺญตฺร สํเฆน อญฺญตฺร อนฺตรายา.}

\csromanmatikaanchor{mtk.106}
\csromantocmarkref{section}{104. คนฺตพฺพวาร}{mtk.106}
\bookmark[dest={mtk.106},level=1]{104. คนฺตพฺพวาร}
\csromanheader{1}{104. \textbf{คนฺตพฺพวาร}}
\proseitem{182}{คนฺตพฺโพ ภิกฺขเว ตทหุโปสเถ สภิกฺขุกา อาวาสา สภิกฺขุโก อาวาโส ยตฺถสฺสุ ภิกฺขู สมาน\-สํวาสกา ยํ ชญฺญา “สกฺโกมิ อชฺเชว คนฺตุนฺ”ติ. คนฺตพฺโพ ภิกฺขเว ตทหุโปสเถ สภิกฺขุกา อาวาสา สภิกฺขุโก อนาวาโส ฯเปฯ สภิกฺขุโก อาวาโส วา อนาวาโส วา ยตฺถสฺสุ ภิกฺขู สมาน\-สํวาสกา ยํ ชญฺญา “สกฺโกมิ อชฺเชว คนฺตุนฺ”ติ.}

\prose{คนฺตพฺโพ ภิกฺขเว ตทหุโปสเถ สภิกฺขุกา อนาวาสา สภิกฺขุโก อาวาโส ฯเปฯ สภิกฺขุโก อนาวาโส ฯเปฯ สภิกฺขุโก อาวาโส วา อนาวาโส วา ยตฺถสฺสุ ภิกฺขู สมาน\-สํวาสกา ยํ ชญฺญา “สกฺโกมิ อชฺเชว คนฺตุนฺ”ติ.}

\prose{คนฺตพฺโพ ภิกฺขเว ตทหุโปสเถ สภิกฺขุกา อาวาสา วา อนาวาสา วา สภิกฺขุโก อาวาโส ฯเปฯ สภิกฺขุโก อนาวาโส ฯเปฯ สภิกฺขุโก อาวาโส วา อนาวาโส วา ยตฺถสฺสุ ภิกฺขู สมาน\-สํวาสกา ยํ ชญฺญา “สกฺโกมิ อชฺเชว คนฺตุนฺ”ติ.}

\csromanmatikaanchor{mtk.107}
\csromantocmarkref{section}{105. วชฺชนียปุคฺคลสนฺทสฺสนา}{mtk.107}
\bookmark[dest={mtk.107},level=1]{105. วชฺชนียปุคฺคลสนฺทสฺสนา}
\csromanheader{1}{105. \textbf{วชฺชนีย}\-\textbf{ปุคฺคล}\-\textbf{สนฺทสฺสนา}}
\csromanmatikaanchor{mtk.108}
\csromantocmarkref{section}{106. อุทฺทานคาถา}{mtk.108}
\bookmark[dest={mtk.108},level=1]{106. อุทฺทานคาถา}
\proseitem{183}{น ภิกฺขเว ภิกฺขุนิยา นิสินฺน\-ปริสาย ปาติโมกฺขํ อุทฺทิสิตพฺพํ, โย อุทฺทิเสยฺย, อาปตฺติ ทุกฺกฏสฺส. น สิกฺขมานาย ฯเปฯ. น \csromanfolio{190}สามเณรสฺส ฯเปฯ. น สามเณริยา ฯเปฯ. น สิกฺขา\-ปจฺจกฺขาตกสฺส ฯเปฯ. น อนฺติมวตฺถุํ อชฺฌาปนฺนกสฺส นิสินฺน\-ปริสาย ปาติโมกฺขํ อุทฺทิสิตพฺพํ, โย อุทฺทิเสยฺย, อาปตฺติ ทุกฺกฏสฺส.}

\prose{น อาปตฺติยา อทสฺสเน อุกฺขิตฺตกสฺส นิสินฺน\-ปริสาย ปาติโมกฺขํ อุทฺทิสิตพฺพํ, โย อุทฺทิเสยฺย, ยถาธมฺโม กาเรตพฺโพ.}

\prose{น อาปตฺติยา อปฺปฏิกมฺเม อุกฺขิตฺตกสฺส นิสินฺน\-ปริสาย ฯเปฯ. น ปาปิกาย ทิฏฺฐิยา อปฺปฏินิสฺสคฺเค อุกฺขิตฺตกสฺส นิสินฺน\-ปริสาย ปาติโมกฺขํ อุทฺทิสิตพฺพํ, โย อุทฺทิเสยฺย, ยถาธมฺโม กาเรตพฺโพ.}

\prose{น ปณฺฑกสฺส นิสินฺน\-ปริสาย ปาติโมกฺขํ อุทฺทิสิตพฺพํ, โย อุทฺทิเสยฺย, อาปตฺติ ทุกฺกฏสฺส.}

\prose{น เถยฺย\-สํวาสกสฺส ฯเปฯ. น ติตฺถิย\-ปกฺกนฺตกสฺส ฯเปฯ. น ติรจฺฉาน\-คตสฺส ฯเปฯ. น มาตุฆาตกสฺส ฯเปฯ. น ปิตุฆาตกสฺส ฯเปฯ. น อรหนฺต\-ฆาตกสฺส ฯเปฯ. น ภิกฺขุนิ\-ทูสกสฺส ฯเปฯ. น สํฆเภทกสฺส ฯเปฯ. น โลหิตุ\-ปฺปาทกสฺส ฯเปฯ. น อุภโต\-พฺยญฺชนกสฺส นิสินฺน\-ปริสาย ปาติโมกฺขํ อุทฺทิสิตพฺพํ, โย อุทฺทิเสยฺย, อาปตฺติ ทุกฺกฏสฺส.}

\proseitemwithcloser{183}{น ภิกฺขเว ปาริวาสิก\-ปาริสุทฺธิ\-ทาเนน อุโปสโถ กาตพฺโพ อญฺญตฺร อวุฏฺฐิตาย ปริสาย. น จ ภิกฺขเว อนุโปสเถ อุโปสโถ กาตพฺโพ อญฺญตฺร สํฆสามคฺคิยาติ.}{ตติย\-ภาณวาโร นิฏฺฐิโต.}
\csromancenter{อุโปสถ\-กฺขนฺธโก ทุติโย.}
\csromancenter{106. \textbf{ตสฺสุทฺทานํ}}
\setlength{\csromangathapagewidth}{0pt}
\setlength{\csromangathawakwidth}{0pt}
\csromangathameasuregroup{ติตฺถิยา พิมฺพิสาโร จ, \\ ธมฺมํ รโห ปาติโมกฺขํ, \\ ยถาปริสา สมคฺคํ, \\ สีมา มหตี นทิยา, \\ นวา ราชคเห เจว, \\ สมฺมนฺเน\textsuperscript{1} ปฐมํ สีมํ, \\ อสมฺมตา คามสีมา, \\ อุทกุกฺเขโป ภินฺทนฺติ, \\ กติ กมฺมานิ อุทฺเทโส, \\ ธมฺมํ วินยํ ตชฺเชนฺติ, \\ โจทนา กเต โอกาเส, \\ จตุปญฺจปรา อาวิ, \\ สคหฏฺฐา อนชฺฌิฏฺฐา, \\ สมฺพหุลา น ชานนฺติ, \\ กติมี กีวติกา ทูเร, \\ อุกฺลาปํ อาสนํ ทีโป, \\ สชฺชุกํ วสฺสุโปสโถ\textsuperscript{1}, \\ คคฺโค จตุตโย เทฺวโก, \\ สพฺโพ สํโฆ เวมติโก, \\ พหู สมสมา โถกา, \\ เอกจฺจา วุฏฺฐิตา สพฺพา, \\ กปฺปเตวาติ กุกฺกุจฺจา, \\ อาวาสิเกน อาคนฺตุ, \\ ปาฏิปโท ปนฺนรโส, \\ ปาริวาสานุโปสโถ, \\ เอเต วิภตฺตา อุทฺทานา,}{\csromangathabat{ติตฺถิยา พิมฺพิสาโร จ,}{สนฺนิปติตุํ ตุณฺหิกา.} \\ \csromangathabat{ธมฺมํ รโห ปาติโมกฺขํ,}{เทวสิกํ ตทา สกิํ.} \\ \csromangathabat{ยถาปริสา สมคฺคํ,}{สามคฺคี มทฺทกุจฺฉิ จ.} \\ \csromangathabat{สีมา มหตี นทิยา,}{อนุ เทฺว ขุทฺทกานิ จ.} \\ \csromangathabat{นวา ราชคเห เจว,}{สีมา อวิปฺปวาสนา.} \\ \csromangathabat{สมฺมนฺเน\textsuperscript{1} ปฐมํ สีมํ,}{ปจฺฉา สีมํ สมูหเน.} \\ \csromangathabat{อสมฺมตา คามสีมา,}{นทิยา สมุทฺเท สเร.} \\ \csromangathabat{อุทกุกฺเขโป ภินฺทนฺติ,}{ตเถวชฺโฌตฺถรนฺติ จ.} \\ \csromangathabat{กติ กมฺมานิ อุทฺเทโส,}{สวรา อสตีปิ จ.} \\ \csromangathabat{ธมฺมํ วินยํ ตชฺเชนฺติ,}{ปุน วินยตชฺชนา.} \\ \csromangathabat{โจทนา กเต โอกาเส,}{อธมฺมปฺปฏิกฺโกสนา.} \\ \csromangathabat{จตุปญฺจปรา อาวิ,}{สญฺจิจฺจ เจปิ วายเม.} \\ \csromangathabat{สคหฏฺฐา อนชฺฌิฏฺฐา,}{โจทนมฺหิ น ชานติ.} \\ \csromangathabat{สมฺพหุลา น ชานนฺติ,}{สชฺชุกํ น จ คจฺฉเร.} \\ \csromangathabat{กติมี กีวติกา ทูเร,}{อาโรเจตุญฺจ นสฺสริ.} \\ \csromangathabat{อุกฺลาปํ อาสนํ ทีโป,}{ทิสา อญฺโญ พหุสฺสุโต.} \\ \csromangathabat{สชฺชุกํ วสฺสุโปสโถ\textsuperscript{1},}{สุทฺธิกมฺมญฺจ ญาตกา.} \\ \csromangathabat{คคฺโค จตุตโย เทฺวโก,}{อาปตฺติสภาคา สริ.} \\ \csromangathabat{สพฺโพ สํโฆ เวมติโก,}{น ชานนฺติ พหุสฺสุโต.} \\ \csromangathabat{พหู สมสมา โถกา,}{ปริสา อวุฏฺฐิตาย จ.} \\ \csromangathabat{เอกจฺจา วุฏฺฐิตา สพฺพา,}{ชานนฺติ จ เวมติกา.} \\ \csromangathabat{กปฺปเตวาติ กุกฺกุจฺจา,}{ชานํ ปสฺสํ สุณนฺติ จ.} \\ \csromangathabat{อาวาสิเกน อาคนฺตุ,}{จาตุปนฺนรโส ปุน.} \\ \csromangathabat{ปาฏิปโท ปนฺนรโส,}{ลิงฺคสํวาสกา อุโภ.} \\ \csromangathabat{ปาริวาสานุโปสโถ,}{อญฺญตฺร สํฆสามคฺคิยา.} \\ \csromangathabat{เอเต วิภตฺตา อุทฺทานา,}{วตฺถุวิภูตการณาติ.} \\ อิมสฺมิํ ขนฺธเก วตฺถูนิ ฉอสีติ.}
\csromangathameasurewak{อิมสฺมิํ ขนฺธเก วตฺถูนิ ฉอสีติ.}
\csromangathasetleft{ติตฺถิยา พิมฺพิสาโร จ, \\ ธมฺมํ รโห ปาติโมกฺขํ, \\ ยถาปริสา สมคฺคํ, \\ สีมา มหตี นทิยา, \\ นวา ราชคเห เจว, \\ สมฺมนฺเน\textsuperscript{1} ปฐมํ สีมํ, \\ อสมฺมตา คามสีมา, \\ อุทกุกฺเขโป ภินฺทนฺติ, \\ กติ กมฺมานิ อุทฺเทโส, \\ ธมฺมํ วินยํ ตชฺเชนฺติ, \\ โจทนา กเต โอกาเส, \\ จตุปญฺจปรา อาวิ, \\ สคหฏฺฐา อนชฺฌิฏฺฐา, \\ สมฺพหุลา น ชานนฺติ, \\ กติมี กีวติกา ทูเร, \\ อุกฺลาปํ อาสนํ ทีโป, \\ สชฺชุกํ วสฺสุโปสโถ\textsuperscript{1}, \\ คคฺโค จตุตโย เทฺวโก, \\ สพฺโพ สํโฆ เวมติโก, \\ พหู สมสมา โถกา, \\ เอกจฺจา วุฏฺฐิตา สพฺพา, \\ กปฺปเตวาติ กุกฺกุจฺจา, \\ อาวาสิเกน อาคนฺตุ, \\ ปาฏิปโท ปนฺนรโส, \\ ปาริวาสานุโปสโถ, \\ เอเต วิภตฺตา อุทฺทานา,}
\csromangathagroupwithclosermajor{\csromangathastanza{\csromangathabat{ติตฺถิยา พิมฺพิสาโร จ,}{สนฺนิปติตุํ ตุณฺหิกา.} \\ \csromangathabat{ธมฺมํ รโห ปาติโมกฺขํ,}{เทวสิกํ ตทา สกิํ.}}\csromangathastanza{\csromangathabat{ยถาปริสา สมคฺคํ,}{สามคฺคี มทฺทกุจฺฉิ จ.} \\ \csromangathabat{สีมา มหตี นทิยา,}{อนุ เทฺว ขุทฺทกานิ จ.}}\csromangathastanza{\csromangathabat{นวา ราชคเห เจว,}{สีมา อวิปฺปวาสนา.} \\ \csromangathabat{สมฺมนฺเน\csromansharedfootnote{b997a02c20a49a1e}{สมฺมเน (ก)} ปฐมํ สีมํ,}{ปจฺฉา สีมํ สมูหเน.}}\csromangathastanza{\csromanfolio{191}\csromangathabat{อสมฺมตา คามสีมา,}{นทิยา สมุทฺเท สเร.} \\ \csromangathabat{อุทกุกฺเขโป ภินฺทนฺติ,}{ตเถวชฺโฌตฺถรนฺติ จ.}}\csromangathastanza{\csromangathabat{กติ กมฺมานิ อุทฺเทโส,}{สวรา อสตีปิ จ.} \\ \csromangathabat{ธมฺมํ วินยํ ตชฺเชนฺติ,}{ปุน วินยตชฺชนา.}}\csromangathastanza{\csromangathabat{โจทนา กเต โอกาเส,}{อธมฺมปฺปฏิกฺโกสนา.} \\ \csromangathabat{จตุปญฺจปรา อาวิ,}{สญฺจิจฺจ เจปิ วายเม.}}\csromangathastanza{\csromangathabat{สคหฏฺฐา อนชฺฌิฏฺฐา,}{โจทนมฺหิ น ชานติ.} \\ \csromangathabat{สมฺพหุลา น ชานนฺติ,}{สชฺชุกํ น จ คจฺฉเร.}}\csromangathastanza{\csromangathabat{กติมี กีวติกา ทูเร,}{อาโรเจตุญฺจ นสฺสริ.} \\ \csromangathabat{อุกฺลาปํ อาสนํ ทีโป,}{ทิสา อญฺโญ พหุสฺสุโต.}}\csromangathastanza{\csromangathabat{สชฺชุกํ วสฺสุโปสโถ\csromansharedfootnote{b799086ff7f3b541}{สชฺชุวสฺสรูโปสโถ (ก)},}{สุทฺธิกมฺมญฺจ ญาตกา.} \\ \csromangathabat{คคฺโค จตุตโย เทฺวโก,}{อาปตฺติสภาคา สริ.}}\csromangathastanza{\csromangathabat{สพฺโพ สํโฆ เวมติโก,}{น ชานนฺติ พหุสฺสุโต.} \\ \csromangathabat{พหู สมสมา โถกา,}{ปริสา อวุฏฺฐิตาย จ.}}\csromangathastanza{\csromangathabat{เอกจฺจา วุฏฺฐิตา สพฺพา,}{ชานนฺติ จ เวมติกา.} \\ \csromangathabat{กปฺปเตวาติ กุกฺกุจฺจา,}{ชานํ ปสฺสํ สุณนฺติ จ.}}\csromangathastanza{\csromangathabat{อาวาสิเกน อาคนฺตุ,}{จาตุปนฺนรโส ปุน.} \\ \csromangathabat{ปาฏิปโท ปนฺนรโส,}{ลิงฺคสํวาสกา อุโภ.}}\csromangathastanza{\csromangathabat{ปาริวาสานุโปสโถ,}{อญฺญตฺร สํฆสามคฺคิยา.} \\ \csromangathabat{เอเต วิภตฺตา อุทฺทานา,}{วตฺถุวิภูตการณาติ.}}\csromangathastanza{\csromangathawak{อิมสฺมิํ ขนฺธเก วตฺถูนิ ฉอสีติ.}}}{\textbf{อุโปสถ}\-\textbf{กฺขนฺธโก นิฏฺฐิโต.}}
\csromanmatikaanchor{mtk.109}
\csromantocmarknopage{chapter}{3. วสฺสูปนายิกกฺขนฺธก}{mtk.109}
\bookmark[dest={mtk.109},level=0]{3. วสฺสูปนายิกกฺขนฺธก}
\chapterheadpageread{3. \textbf{วสฺสู}\-\textbf{ปนายิก}\-\textbf{กฺขนฺธก}}
\csromanmatikaanchor{mtk.110}
\csromantocmarkref{section}{107. วสฺสูปนายิกานุชานนา}{mtk.110}
\bookmark[dest={mtk.110},level=1]{107. วสฺสูปนายิกานุชานนา}
\csromanheader{1}{107. \textbf{วสฺสูปนายิกา}\-\textbf{นุชานนา}}
\proseitem{184}{\csromanfolio{192}เตน สมเยน พุทฺโธ ภควา ราชคเห วิหรติ เวฬุวเน กลนฺทก\-นิวาเป. เตน โข ปน สมเยน ภควตา ภิกฺขูนํ วสฺสาวาโส อปญฺญตฺโต โหติ. เต’ธ ภิกฺขู เหมนฺตมฺปิ คิมฺหมฺปิ วสฺสมฺปิ จาริกํ จรนฺติ. มนุสฺสา อุชฺฌายนฺติ ขิยฺยนฺติ วิปาเจนฺติ “กถํ หิ นาม สมณา สกฺยปุตฺติยา เหมนฺตมฺปิ คิมฺหมฺปิ วสฺสมฺปิ จาริกํ จริสฺสนฺติ หริตานิ ติณานิ สมฺมทฺทนฺตา เอกินฺทฺริยํ ชีวํ วิเหเฐนฺตา พหู ขุทฺทเก ปาเณ สงฺฆาตํ อาปาเทนฺตา, อิเม หิ นาม อญฺญติตฺถิยา ทุรกฺขาต\-ธมฺมา วสฺสาวาสํ อลฺลียิสฺสนฺติ สงฺกสายิสฺสนฺติ\csromansharedfootnote{3780c91887cd92ab}{สงฺกาสยิสฺสนฺติ (สี, สฺยา)}, อิเม หิ นาม สกุนฺตกา รุกฺขคฺเคสุ กุลาวกานิ กริตฺวา วสฺสาวาสํ อลฺลียิสฺสนฺติ สงฺกสายิสฺสนฺติ, อิเม ปน สมณา สกฺยปุตฺติยา เหมนฺตมฺปิ คิมฺหมฺปิ วสฺสมฺปิ จาริกํ จรนฺติ หริตานิ ติณานิ สมฺมทฺทนฺตา เอกินฺทฺริยํ ชีวํ วิเหเฐนฺตา พหู ขุทฺทเก ปาเณ สงฺฆาตํ อาปาเทนฺตา”ติ. อสฺโสสุํ โข ภิกฺขู เตสํ มนุสฺสานํ อุชฺฌายนฺตานํ ขิยฺยนฺตานํ วิปาเจนฺตานํ. อถ โข เต ภิกฺขู ภควโต เอตมตฺถํ อาโรเจสุํ. อถ โข ภควา เอตสฺมิํ นิทาเน เอตสฺมิํ ปกรเณ ธมฺมิํ กถํ กตฺวา ภิกฺขู อามนฺเตสิ “อนุชานามิ ภิกฺขเว วสฺสํ อุปคนฺตุนฺ”ติ. อถ โข ภิกฺขูนํ เอตทโหสิ “กทา นุ โข วสฺสํ อุปคนฺตพฺพนฺ”ติ. ภควโต เอตมตฺถํ อาโรเจสุํ. อนุชานามิ ภิกฺขเว วสฺสาเน วสฺสํ อุปคนฺตุนฺติ.}

\prose{อถ โข ภิกฺขูนํ เอตทโหสิ “กติ นุ โข วสฺสูปนายิกา”ติ. ภควโต เอตมตฺถํ อาโรเจสุํ. เทฺวมา ภิกฺขเว วสฺสูปนายิกา ปุริมิกา ปจฺฉิมิกา, อปรชฺชุคตาย อาสาฬฺหิยา ปุริมิกา อุปคนฺตพฺพา, มาสคตาย อาสาฬฺหิยา ปจฺฉิมิกา อุปคนฺตพฺพา, อิมา โข ภิกฺขเว เทฺว วสฺสู\-ปนายิกาติ.}

\csromanmatikaanchor{mtk.111}
\csromantocmarkref{section}{108. วสฺสาเน จาริกาปฏิกฺเขปาทิ}{mtk.111}
\bookmark[dest={mtk.111},level=1]{108. วสฺสาเน จาริกาปฏิกฺเขปาทิ}
\csromanheader{1}{108. \textbf{วสฺสาเน จาริกา}\-\textbf{ปฏิกฺเขปาทิ}}
\proseitem{185}{เตน โข ปน สมเยน ฉพฺพคฺคิยา ภิกฺขู วสฺสํ อุปคนฺตฺวา อนฺตราวสฺสํ จาริกํ จรนฺติ. มนุสฺสา ตเถว อุชฺฌายนฺติ ขิยฺยนฺติ วิปาเจนฺติ \csromanfolio{193}“กถํ หิ นาม สมณา สกฺยปุตฺติยา เหมนฺตมฺปิ คิมฺหมฺปิ วสฺสมฺปิ จาริกํ จริสฺสนฺติ หริตานิ ติณานิ สมฺมทฺทนฺตา เอกินฺทฺริยํ ชีวํ วิเหเฐนฺตา พหู ขุทฺทเก ปาเณ สงฺฆาตํ อาปาเทนฺตา, อิเม หิ นาม อญฺญติตฺถิยา ทุรกฺขาต\-ธมฺมา วสฺสาวาสํ อลฺลียิสฺสนฺติ สงฺกสายิสฺสนฺติ, อิเม หิ นาม สกุนฺตกา รุกฺขคฺเคสุ กุลาวกานิ กริตฺวา วสฺสาวาสํ อลฺลียิสฺสนฺติ สงฺกสายิสฺสนฺติ, อิเม ปน สมณา สกฺยปุตฺติยา เหมนฺตมฺปิ คิมฺหมฺปิ วสฺสมฺปิ จาริกํ จรนฺติ หริตานิ ติณานิ สมฺมทฺทนฺตา เอกินฺทฺริยํ ชีวํ วิเหเฐนฺตา พหู ขุทฺทเก ปาเณ สงฺฆาตํ อาปาเทนฺตา”ติ. อสฺโสสุํ โข ภิกฺขู เตสํ มนุสฺสานํ อุชฺฌายนฺตานํ ขิยฺยนฺตานํ วิปาเจนฺตานํ. เย เต ภิกฺขู อปฺปิจฺฉา ฯเปฯ เต อุชฺฌายนฺติ ขิยฺยนฺติ วิปาเจนฺติ “กถํ หิ นาม ฉพฺพคฺคิยา ภิกฺขู วสฺสํ อุปคนฺตฺวา อนฺตราวสฺสํ จาริกํ จริสฺสนฺตี”ติ. อถ โข เต ภิกฺขู ภควโต เอตมตฺถํ อาโรเจสุํ ฯเปฯ. อถ โข ภควา เอตสฺมิํ นิทาเน เอตสฺมิํ ปกรเณ ธมฺมิํ กถํ กตฺวา ภิกฺขู อามนฺเตสิ “น ภิกฺขเว วสฺสํ อุปคนฺตฺวา ปุริมํ วา เตมาสํ ปจฺฉิมํ วา เตมาสํ อวสิตฺวา จาริกา ปกฺกมิตพฺพา, โย ปกฺกเมยฺย, อาปตฺติ ทุกฺกฏสฺสา”ติ.}

\proseitem{186}{เตน โข ปน สมเยน ฉพฺพคฺคิยา ภิกฺขู น อิจฺฉนฺติ วสฺสํ อุปคนฺตุํ. ภควโต เอตมตฺถํ อาโรเจสุํ. น ภิกฺขเว วสฺสํ น อุปคนฺตพฺพํ, โย น อุปคจฺเฉยฺย, อาปตฺติ ทุกฺกฏสฺสาติ.}

\prose{เตน โข ปน สมเยน ฉพฺพคฺคิยา ภิกฺขู ตทหุ วสฺสู\-ปนายิกาย วสฺสํ อนุปคนฺตุกามา สญฺจิจฺจ อาวาสํ อติกฺกมนฺติ. ภควโต เอตมตฺถํ อาโรเจสุํ. น ภิกฺขเว ตทหุ วสฺสู\-ปนายิกาย วสฺสํ อนุปคนฺตุ\-กาเมน สญฺจิจฺจ อาวาโส อติกฺกมิตพฺโพ, โย อติกฺกเมยฺย, อาปตฺติ ทุกฺกฏสฺสาติ.}

\proseitem{186}{เตน โข ปน สมเยน ราชา มาคโธ เสนิโย พิมฺพิสาโร วสฺสํ อุกฺกฑฺฒิตุกาโม ภิกฺขูนํ สนฺติเก ทูตํ ปาเหสิ “ยทิ ปนายฺยา อาคเม ชุเณฺห วสฺสํ อุปคจฺเฉยฺยุนฺ”ติ. ภควโต เอตมตฺถํ อาโรเจสุํ. อนุชานามิ ภิกฺขเว ราชูนํ อนุวตฺติตุนฺติ.}

\csromanmatikaanchor{mtk.112}
\csromantocmarkref{section}{109. สตฺตาหกรณียานุชานนา}{mtk.112}
\bookmark[dest={mtk.112},level=1]{109. สตฺตาหกรณียานุชานนา}
\csromanheader{1}{109. \textbf{สตฺตาหกรณียา}\-\textbf{นุชานนา}}
\proseitem{187}{อถ โข ภควา ราชคเห ยถาภิรนฺตํ วิหริตฺวา เยน สาวตฺถิ เตน จาริกํ ปกฺกามิ, อนุปุพฺเพน จาริกํ จรมาโน เยน \csromanfolio{194}สาวตฺถิ ตทวสริ, ตตฺร สุทํ ภควา สาวตฺถิยํ วิหรติ เชตวเน อนาถ\-ปิณฺฑิกสฺส อาราเม. เตน โข ปน สมเยน โกสเลสุ ชนปเท อุเทเนน อุปาสเกน สํฆํ อุทฺทิสฺส วิหาโร การาปิโต โหติ. โส ภิกฺขูนํ สนฺติเก ทูตํ ปาเหสิ “อาคจฺฉนฺตุ ภทนฺตา, อิจฺฉามิ ทานญฺจ ทาตุํ ธมฺมญฺจ โสตุํ ภิกฺขู จ ปสฺสิตุนฺ”ติ. ภิกฺขู เอวมาหํสุ “ภควตา อาวุโส ปญฺญตฺตํ ‘น วสฺสํ อุปคนฺตฺวา ปุริมํ วา เตมาสํ ปจฺฉิมํ วา เตมาสํ อวสิตฺวา จาริกา ปกฺกมิตพฺพา’ติ, อาคเมตุ อุเทโน อุปาสโก ยาว ภิกฺขู วสฺสํ วสนฺติ, วสฺสํวุฏฺฐา อาคมิสฺสนฺติ, สเจ ปนสฺส อจฺจายิกํ กรณียํ, ตตฺเถว อาวาสิกานํ ภิกฺขูนํ สนฺติเก วิหารํ ปติฏฺฐาเปตู”ติ. อุเทโน อุปาสโก อุชฺฌายติ ขิยฺยติ วิปาเจติ “กถํ หิ นาม ภทนฺตา มยา ปหิเต น อาคจฺฉิสฺสนฺติ, อหํ หิ ทายโก การโก สํฆุปฏฺฐาโก”ติ. อสฺโสสุํ โข ภิกฺขู อุเทนสฺส อุปาสกสฺส อุชฺฌายนฺตสฺส ขิยฺยนฺตสฺส วิปาเจนฺตสฺส. อถ โข เต ภิกฺขู ภควโต เอตมตฺถํ อาโรเจสุํ. อถ โข ภควา เอตสฺมิํ นิทาเน เอตสฺมิํ ปกรเณ ธมฺมิํ กถํ กตฺวา ภิกฺขู อามนฺเตสิ, อนุชานามิ ภิกฺขเว สตฺตนฺนํ สตฺตาห\-กรณีเยน ปหิเต คนฺตุํ, น เตฺวว อปฺปหิเต, ภิกฺขุสฺส ภิกฺขุนิยา สิกฺขมานาย สามเณรสฺส สามเณริยา อุปาสกสฺส อุปาสิกาย. อนุชานามิ ภิกฺขเว อิเมสํ สตฺตนฺนํ สตฺตาห\-กรณีเยน ปหิเต คนฺตุํ, น เตฺวว อปฺปหิเต, สตฺตาหํ สนฺนิวตฺโต กาตพฺโพ.}

\proseitem{188}{อิธ ปน ภิกฺขเว อุปาสเกน สํฆํ อุทฺทิสฺส วิหาโร การาปิโต โหติ, โส เจ ภิกฺขูนํ สนฺติเก ทูตํ ปหิเณยฺย “อาคจฺฉนฺตุ ภทนฺตา, อิจฺฉามิ ทานญฺจ ทาตุํ ธมฺมญฺจ โสตุํ ภิกฺขู จ ปสฺสิตุนฺ”ติ. คนฺตพฺพํ ภิกฺขเว สตฺตาห\-กรณีเยน ปหิเต, น เตฺวว อปฺปหิเต, สตฺตาหํ สนฺนิวตฺโต กาตพฺโพ.}

\prose{อิธ ปน ภิกฺขเว อุปาสเกน สํฆํ อุทฺทิสฺส อฑฺฒโยโค การาปิโต โหติ ฯเปฯ ปาสาโท การาปิโต โหติ. หมฺมิยํ การาปิตํ โหติ. คุหา การาปิตา โหติ. ปริเวณํ การาปิตํ โหติ. โกฏฺฐโก การาปิโต โหติ. อุปฏฺฐานสาลา การาปิตา โหติ. อคฺคิสาลา การาปิตา โหติ. กปฺปิยกุฏิ การาปิตา โหติ. วจฺจกุฏิ การาปิตา โหติ. จงฺกโม การาปิโต \csromanfolio{195}โหติ. จงฺกมนสาลา การาปิตา โหติ. อุทปาโน การาปิโต โหติ. อุทปานสาลา การาปิตา โหติ. ชนฺตาฆรํ การาปิตํ โหติ. ชนฺตาฆรสาลา การาปิตา โหติ. โปกฺขรณี การาปิตา โหติ. มณฺฑโป การาปิโต โหติ. อาราโม การาปิโต โหติ. อารามวตฺถุ การาปิตํ โหติ, โส เจ ภิกฺขูนํ สนฺติเก ทูตํ ปหิเณยฺย “อาคจฺฉนฺตุ ภทนฺตา, อิจฺฉามิ ทานญฺจ ทาตุํ ธมฺมญฺจ โสตุํ ภิกฺขู จ ปสฺสิตุนฺ”ติ. คนฺตพฺพํ ภิกฺขเว สตฺตาห\-กรณีเยน ปหิเต, น เตฺวว อปฺปหิเต, สตฺตาหํ สนฺนิวตฺโต กาตพฺโพ.}

\prose{อิธ ปน ภิกฺขเว อุปาสเกน สมฺพหุเล ภิกฺขู อุทฺทิสฺส ฯเปฯ เอกํ ภิกฺขุํ อุทฺทิสฺส วิหาโร การาปิโต โหติ ฯเปฯ อฑฺฒโยโค การาปิโต โหติ. ปาสาโท การาปิโต โหติ. หมฺมิยํ การาปิตํ โหติ. คุหา การาปิตา โหติ. ปริเวณํ การาปิตํ โหติ. โกฏฺฐโก การาปิโต โหติ. อุปฏฺฐานสาลา การาปิตา โหติ. อคฺคิสาลา การาปิตา โหติ. กปฺปิยกุฏิ การาปิตา โหติ. วจฺจกุฏิ การาปิตา โหติ. จงฺกโม การาปิโต โหติ. จงฺกมนสาลา การาปิตา โหติ. อุทปาโน การาปิโต โหติ. อุทปานสาลา การาปิตา โหติ. ชนฺตาฆรํ การาปิตํ โหติ. ชนฺตาฆรสาลา การาปิตา โหติ. โปกฺขรณี การาปิตา โหติ. มณฺฑโป การาปิโต โหติ. อาราโม การาปิโต โหติ. อารามวตฺถุ การาปิตํ โหติ, โส เจ ภิกฺขูนํ สนฺติเก ทูตํ ปหิเณยฺย “อาคจฺฉนฺตุ ภทนฺตา, อิจฺฉามิ ทานญฺจ ทาตุํ ธมฺมญฺจ โสตุํ ภิกฺขู จ ปสฺสิตุนฺ”ติ. คนฺตพฺพํ ภิกฺขเว สตฺตาห\-กรณีเยน ปหิเต, น เตฺวว อปฺปหิเต, สตฺตาหํ สนฺนิวตฺโต กาตพฺโพ.}

\prose{อิธ ปน ภิกฺขเว อุปาสเกน ภิกฺขุนิสํฆํ อุทฺทิสฺส ฯเปฯ สมฺพหุลา ภิกฺขุนิโย อุทฺทิสฺส ฯเปฯ เอกํ ภิกฺขุนิํ อุทฺทิสฺส ฯเปฯ สมฺพหุลา สิกฺขมานาโย อุทฺทิสฺส ฯเปฯ เอกํ สิกฺขมานํ อุทฺทิสฺส ฯเปฯ สมฺพหุเล สามเณเร อุทฺทิสฺส ฯเปฯ เอกํ สามเณรํ อุทฺทิสฺส ฯเปฯ สมฺพหุลา สามเณริโย อุทฺทิสฺส ฯเปฯ เอกํ สามเณริํ อุทฺทิสฺส วิหาโร การาปิโต โหติ ฯเปฯ อฑฺฒโยโค การาปิโต โหติ. ปาสาโท การาปิโต โหติ. \csromanfolio{196}หมฺมิยํ การาปิตํ โหติ. คุหา การาปิตา โหติ. ปริเวณํ การาปิตํ โหติ. โกฏฺฐโก การาปิโต โหติ. อุปฏฺฐานสาลา การาปิตา โหติ. อคฺคิสาลา การาปิตา โหติ. กปฺปิยกุฏิ การาปิตา โหติ. จงฺกโม การาปิโต โหติ. จงฺกมนสาลา การาปิตา โหติ. อุทปาโน การาปิโต โหติ. อุทปานสาลา การาปิตา โหติ. โปกฺขรณี การาปิตา โหติ. มณฺฑโป การาปิโต โหติ. อาราโม การาปิโต โหติ. อารามวตฺถุ การาปิตํ โหติ, โส เจ ภิกฺขูนํ สนฺติเก ทูตํ ปหิเณยฺย “อาคจฺฉนฺตุ ภทนฺตา, อิจฺฉามิ ทานญฺจ ทาตุํ ธมฺมญฺจ โสตุํ ภิกฺขู จ ปสฺสิตุนฺ”ติ. คนฺตพฺพํ ภิกฺขเว สตฺตาห\-กรณีเยน ปหิเต, น เตฺวว อปฺปหิเต, สตฺตาหํ สนฺนิวตฺโต กาตพฺโพ.}

\proseitem{198}{อิธ ปน ภิกฺขเว อุปาสเกน อตฺตโน อตฺถาย นิเวสนํ การาปิตํ โหติ ฯเปฯ สยนิฆรํ การาปิตํ โหติ. อุโทสิโต การาปิโต โหติ. อฏฺโฏ การาปิโต โหติ. มาโฬ การาปิโต โหติ. อาปโณ การาปิโต โหติ. อาปณสาลา การาปิตา โหติ. ปาสาโท การาปิโต โหติ. หมฺมิยํ การาปิตํ โหติ. คุหา การาปิตา โหติ. ปริเวณํ การาปิตํ โหติ. โกฏฺฐโก การาปิโต โหติ. อุปฏฺฐานสาลา การาปิตา โหติ. อคฺคิสาลา การาปิตา โหติ. รสวตี การาปิตา โหติ. จงฺกโม การาปิโต โหติ. จงฺกมนสาลา การาปิตา โหติ. อุทปาโน การาปิโต โหติ. อุทปานสาลา การาปิตา โหติ. โปกฺขรณี การาปิตา โหติ. มณฺฑโป การาปิโต โหติ. อาราโม การาปิโต โหติ. อารามวตฺถุ การาปิตํ โหติ. ปุตฺตสฺส วา วาเรยฺยํ โหติ. ธีตุยา วา วาเรยฺยํ โหติ. คิลาโน วา โหติ. อภิญฺญาตํ วา สุตฺตนฺตํ ภณติ, โส เจ ภิกฺขูนํ สนฺติเก ทูตํ ปหิเณยฺย “อาคจฺฉนฺตุ ภทนฺตา, อิมํ สุตฺตนฺตํ ปริยาปุณิสฺสนฺติ, ปุรายํ สุตฺตนฺโต ปลุชฺชตี”ติ. อญฺญตรํ วา ปนสฺส กิจฺจํ โหติ กรณียํ วา, โส เจ ภิกฺขูนํ สนฺติเก ทูตํ ปหิเณยฺย “อาคจฺฉนฺตุ ภทนฺตา, อิจฺฉามิ ทานญฺจ ทาตุํ ธมฺมญฺจ โสตุํ ภิกฺขู จ ปสฺสิตุนฺ”ติ. คนฺตพฺพํ ภิกฺขเว สตฺตาห\-กรณีเยน ปหิเต, น เตฺวว อปฺปหิเต, สตฺตาหํ สนฺนิวตฺโต กาตพฺโพ.}

\proseitem{190}{\csromanfolio{197}อิธ ปน ภิกฺขเว อุปาสิกาย สํฆํ อุทฺทิสฺส วิหาโร การาปิโต โหติ, สา เจ ภิกฺขูนํ สนฺติเก ทูตํ ปหิเณยฺย “อาคจฺฉนฺตุ อยฺยา, อิจฺฉามิ ทานญฺจ ทาตุํ ธมฺมญฺจ โสตุํ ภิกฺขู จ ปสฺสิตุนฺ”ติ.}

\prose{คนฺตพฺพํ ภิกฺขเว สตฺตาห\-กรณีเยน ปหิเต, น เตฺวว อปฺปหิเต, สตฺตาหํ สนฺนิวตฺโต กาตพฺโพ.}

\prose{อิธ ปน ภิกฺขเว อุปาสิกาย สํฆํ อุทฺทิสฺส อฑฺฒโยโค การาปิโต โหติ ฯเปฯ ปาสาโท การาปิโต โหติ. หมฺมิยํ การาปิตํ โหติ. คุหา การาปิตา โหติ. ปริเวณํ การาปิตํ โหติ. โกฏฺฐโก การาปิโต โหติ. อุปฏฺฐานสาลา การาปิตา โหติ. อคฺคิสาลา การาปิตา โหติ. กปฺปิยกุฏิ การาปิตา โหติ. วจฺจกุฏิ การาปิตา โหติ. จงฺกโม การาปิโต โหติ. จงฺกมนสาลา การาปิตา โหติ. อุทปาโน การาปิโต โหติ. อุทปานสาลา การาปิตา โหติ. ชนฺตาฆรํ การาปิตํ โหติ. ชนฺตาฆรสาลา การาปิตา โหติ. โปกฺขรณี การาปิตา โหติ. มณฺฑโป การาปิโต โหติ. อาราโม การาปิโต โหติ. อารามวตฺถุ การาปิตํ โหติ, สา เจ ภิกฺขูนํ สนฺติเก ทูตํ ปหิเณยฺย “อาคจฺฉนฺตุ อยฺยา, อิจฺฉามิ ทานญฺจ ทาตุํ ธมฺมญฺจ โสตุํ ภิกฺขู จ ปสฺสิตุนฺ”ติ. คนฺตพฺพํ ภิกฺขเว สตฺตาห\-กรณีเยน ปหิเต, น เตฺวว อปฺปหิเต, สตฺตาหํ สนฺนิวตฺโต กาตพฺโพ.}

\proseitem{190}{อิธ ปน ภิกฺขเว อุปาสิกาย สมฺพหุเล ภิกฺขู อุทฺทิสฺส ฯเปฯ เอกํ ภิกฺขุํ อุทฺทิสฺส ฯเปฯ ภิกฺขุนิสํฆํ อุทฺทิสฺส ฯเปฯ สมฺพหุลา ภิกฺขุนิโย อุทฺทิสฺส ฯเปฯ เอกํ ภิกฺขุนิํ อุทฺทิสฺส ฯเปฯ สมฺพหุลา สิกฺขมานาโย อุทฺทิสฺส ฯเปฯ เอกํ สิกฺขมานํ อุทฺทิสฺส ฯเปฯ สมฺพหุเล สามเณเร อุทฺทิสฺส ฯเปฯ เอกํ สามเณรํ อุทฺทิสฺส ฯเปฯ สมฺพหุลา สามเณริโย อุทฺทิสฺส ฯเปฯ เอกํ สามเณริํ อุทฺทิสฺส ฯเปฯ.}

\proseitem{191}{อิธ ปน ภิกฺขเว อุปาสิกาย อตฺตโน อตฺถาย นิเวสนํ การาปิตํ โหติ ฯเปฯ สยนิฆรํ การาปิตํ โหติ. อุโทสิโต การาปิโต โหติ. อฏฺโฏ การาปิโต โหติ. มาโฬ การาปิโต โหติ. อาปโณ การาปิโต โหติ. อาปณสาลา การาปิตา โหติ. ปาสาโท การาปิโต โหติ. \csromanfolio{198}หมฺมิยํ การาปิตํ โหติ. คุหา การาปิตา โหติ. ปริเวณํ การาปิตํ โหติ. โกฏฺฐโก การาปิโต โหติ. อุปฏฺฐานสาลา การาปิตา โหติ. อคฺคิสาลา การาปิตา โหติ. รสวตี การาปิตา โหติ. จงฺกโม การาปิโต โหติ. จงฺกมนสาลา การาปิตา โหติ. อุทปาโน การาปิโต โหติ. อุทปานสาลา การาปิตา โหติ. โปกฺขรณี การาปิตา โหติ. \textbf{มณฺฑโป การาปิโต โหติ. อาราโม การาปิโต โหติ. อารามวตฺถุ การาปิตํ โหติ.} ปุตฺตสฺส วา วาเรยฺยํ โหติ. ธีตุยา วา วาเรยฺยํ โหติ. คิลานา วา โหติ. อภิญฺญาตํ วา สุตฺตนฺตํ ภณติ, สา เจ ภิกฺขูนํ สนฺติเก ทูตํ ปหิเณยฺย “อาคจฺฉนฺตุ อยฺยา, อิมํ สุตฺตนฺตํ ปริยาปุณิสฺสนฺติ, ปุรายํ สุตฺตนฺโต ปลุชฺชตี”ติ. อญฺญตรํ วา ปนสฺสา กิจฺจํ โหติ กรณียํ วา, สา เจ ภิกฺขูนํ สนฺติเก ทูตํ ปหิเณยฺย “อาคจฺฉนฺตุ อยฺยา, อิจฺฉามิ ทานญฺจ ทาตุํ ธมฺมญฺจ โสตุํ ภิกฺขู จ ปสฺสิตุนฺ”ติ. คนฺตพฺพํ ภิกฺขเว สตฺตาห\-กรณีเยน ปหิเต, น เตฺวว อปฺปหิเต, สตฺตาหํ สนฺนิวตฺโต กาตพฺโพ.}

\proseitem{192}{อิธ ปน ภิกฺขเว ภิกฺขุนา สํฆํ อุทฺทิสฺส ฯเปฯ ภิกฺขุนิยา สํฆํ อุทฺทิสฺส. สิกฺขมานาย สํฆํ อุทฺทิสฺส. สามเณเรน สํฆํ อุทฺทิสฺส. สามเณริยา สํฆํ อุทฺทิสฺส. สมฺพหุเล ภิกฺขู อุทฺทิสฺส. เอกํ ภิกฺขุํ อุทฺทิสฺส. ภิกฺขุนิสํฆํ อุทฺทิสฺส. สมฺพหุลา ภิกฺขุนิโย อุทฺทิสฺส. เอกํ ภิกฺขุนิํ อุทฺทิสฺส. สมฺพหุลา สิกฺขมานาโย อุทฺทิสฺส. เอกํ สิกฺขมานํ อุทฺทิสฺส. สมฺพหุเล สามเณเร อุทฺทิสฺส. เอกํ สามเณรํ อุทฺทิสฺส. สมฺพหุลา สามเณริโย อุทฺทิสฺส. เอกํ สามเณริํ อุทฺทิสฺส. อตฺตโน อตฺถาย วิหาโร การาปิโต โหติ ฯเปฯ อฑฺฒโยโค การาปิโต โหติ. ปาสาโท การาปิโต โหติ. หมฺมิยํ การาปิตํ โหติ. คุหา การาปิตา โหติ. ปริเวณํ การาปิตํ โหติ. โกฏฺฐโก การาปิโต โหติ. อุปฏฺฐานสาลา การาปิตา โหติ. อคฺคิสาลา การาปิตา โหติ. กปฺปิยกุฏิ การาปิตา โหติ. จงฺกโม การาปิโต โหติ. จงฺกมนสาลา การาปิตา โหติ. อุทปาโน การาปิโต โหติ. อุทปานสาลา การาปิตา โหติ. โปกฺขรณี การาปิตา โหติ. มณฺฑโป การาปิโต โหติ. อาราโม การาปิโต โหติ. อารามวตฺถุ การาปิตํ โหติ, สา เจ \csromanfolio{199}ภิกฺขูนํ สนฺติเก ทูตํ ปหิเณยฺย “อาคจฺฉนฺตุ อยฺยา, อิจฺฉามิ ทานญฺจ ทาตุํ ธมฺมญฺจ โสตุํ ภิกฺขู จ ปสฺสิตุนฺ”ติ. คนฺตพฺพํ ภิกฺขเว สตฺตาห\-กรณีเยน ปหิเต, น เตฺวว อปฺปหิเต, สตฺตาหํ สนฺนิวตฺโต กาตพฺโพติ.}

\csromanmatikaanchor{mtk.113}
\csromantocmarkref{section}{110. ปญฺจนฺนํ อปฺปหิเตปิ อนุชานนา}{mtk.113}
\bookmark[dest={mtk.113},level=1]{110. ปญฺจนฺนํ อปฺปหิเตปิ อนุชานนา}
\csromanheader{1}{110. \textbf{ปญฺจนฺนํ อปฺปหิเตปิ อนุชานนา}}
\proseitem{193}{เตน โข ปน สมเยน อญฺญตโร ภิกฺขุ คิลาโน โหติ, โส ภิกฺขูนํ สนฺติเก ทูตํ ปาเหสิ “อหํ หิ คิลาโน, อาคจฺฉนฺตุ ภิกฺขู, อิจฺฉามิ ภิกฺขูนํ อาคตนฺ”ติ.}

\prose{ภควโต เอตมตฺถํ อาโรเจสุํ. อนุชานามิ ภิกฺขเว ปญฺจนฺนํ สตฺตาห\-กรณีเยน อปฺปหิเตปิ คนฺตุํ, ปเคว ปหิเต, ภิกฺขุสฺส ภิกฺขุนิยา สิกฺขมานาย สามเณรสฺส สามเณริยา, อนุชานามิ ภิกฺขเว อิเมสํ ปญฺจนฺนํ สตฺตาห\-กรณีเยน อปฺปหิเตปิ คนฺตุํ, ปเคว คหิเต, สตฺตาหํ สนฺนิวตฺโต กาตพฺโพ.}

\prose{อิธ ปน ภิกฺขเว ภิกฺขุ คิลาโน โหติ, โส เจ ภิกฺขูนํ สนฺติเก ทูตํ ปหิเณยฺย “อหํ หิ คิลาโน อาคจฺฉนฺตุ ภิกฺขู, อิจฺฉามิ ภิกฺขูนํ อาคตนฺ”ติ.}

\prose{คนฺตพฺพํ ภิกฺขเว สตฺตาห\-กรณีเยน อปฺปหิเตปิ, ปเคว ปหิเต “คิลานภตฺตํ วา ปริเยสิสฺสามิ, คิลานุปฏฺฐาก\-ภตฺตํ วา ปริเยสิสฺสามิ, คิลาน\-เภสชฺชํ วา ปริเยสิสฺสามิ, ปุจฺฉิสฺสามิ วา อุปฏฺฐหิสฺสามิ วา”ติ, สตฺตาหํ สนฺนิวตฺโต กาตพฺโพ.}

\prose{อิธ ปน ภิกฺขเว ภิกฺขุสฺส อนภิรติ อุปฺปนฺนา โหติ, โส เจ ภิกฺขูนํ สนฺติเก ทูตํ ปหิเณยฺย “อนภิรติ เม อุปฺปนฺนา, อาคจฺฉนฺตุ ภิกฺขู, อิจฺฉามิ ภิกฺขูนํ อาคตนฺ”ติ.}

\prose{คนฺตพฺพํ ภิกฺขเว สตฺตาห\-กรณีเยน อปฺปหิเตปิ, ปเคว ปหิเต “อนภิรตํ วูปกาเสสฺสามิ วา วูปกาสาเปสฺสามิ วา ธมฺมกถํ วาสฺส กริสฺสามี”ติ, สตฺตาหํ สนฺนิวตฺโต กาตพฺโพ.}

\prose{อิธ ปน ภิกฺขเว ภิกฺขุสฺส กุกฺกุจฺจํ อุปฺปนฺนํ โหติ, โส เจ ภิกฺขูนํ สนฺติเก ทูตํ ปหิเณยฺย “กุกฺกุจฺจํ เม อุปฺปนฺนํ, อาคจฺฉนฺตุ ภิกฺขู, อิจฺฉามิ ภิกฺขูนํ อาคตนฺ”ติ.}

\prose{คนฺตพฺพํ ภิกฺขเว สตฺตาห\-กรณีเยน อปฺปหิเตปิ, ปเคว ปหิเต “กุกฺกุจฺจํ วิโนเทสฺสามิ วา วิโนทาเปสฺสามิ วา ธมฺมกถํ วาสฺส กริสฺสามี”ติ, สตฺตาหํ สนฺนิวตฺโต กาตพฺโพ.}

\prose{\csromanfolio{200}อิธ ปน ภิกฺขเว ภิกฺขุสฺส ทิฏฺฐิคตํ อุปฺปนฺนํ โหติ, โส เจ ภิกฺขูนํ สนฺติเก ทูตํ ปหิเณยฺย “ทิฏฺฐิคตํ เม อุปฺปนฺนํ, อาคจฺฉนฺตุ ภิกฺขู, อิจฺฉามิ ภิกฺขูนํ อาคตนฺ”ติ.}

\prose{คนฺตพฺพํ ภิกฺขเว สตฺตาห\-กรณีเยน อปฺปหิเตปิ, ปเคว ปหิเต “ทิฏฺฐิคตํ วิเวเจสฺสามิ วา วิเวจาเปสฺสามิ วา ธมฺมกถํ วาสฺส กริสฺสามี”ติ, สตฺตาหํ สนฺนิวตฺโต กาตพฺโพ.}

\prose{อิธ ปน ภิกฺขเว ภิกฺขุ ครุธมฺมํ อชฺฌาปนฺโน โหติ ปริวาสารโห, โส เจ ภิกฺขูนํ สนฺติเก ทูตํ ปหิเณยฺย “อหํ หิ ครุธมฺมํ อชฺฌาปนฺโน ปริวาสารโห, อาคจฺฉนฺตุ ภิกฺขู, อิจฺฉามิ ภิกฺขูนํ อาคตนฺ”ติ.}

\prose{คนฺตพฺพํ ภิกฺขเว สตฺตาห\-กรณีเยน อปฺปหิเตปิ, ปเคว ปหิเต “ปริวาสทานํ อุสฺสุกฺกํ กริสฺสามิ วา อนุสฺสาเวสฺสามิ วา คณปูรโก วา ภวิสฺสามี”ติ, สตฺตาหํ สนฺนิวตฺโต กาตพฺโพ.}

\prose{อิธ ปน ภิกฺขเว ภิกฺขุ มูลาย ปฏิกสฺสนา\-รโห โหติ, โส เจ ภิกฺขูนํ สนฺติเก ทูตํ ปหิเณยฺย “อหํ หิ มูลาย ปฏิกสฺสนา\-รโห, อาคจฺฉนฺตุ ภิกฺขู, อิจฺฉามิ ภิกฺขูนํ อาคตนฺ”ติ.}

\prose{คนฺตพฺพํ ภิกฺขเว สตฺตาห\-กรณีเยน อปฺปหิเตปิ, ปเคว ปหิเต “มูลาย ปฏิกสฺสนํ อุสฺสุกฺกํ กริสฺสามิ วา อนุสฺสาเวสฺสามิ วา คณปูรโก วา ภวิสฺสามี”ติ, สตฺตาหํ สนฺนิวตฺโต กาตพฺโพ.}

\prose{อิธ ปน ภิกฺขเว ภิกฺขุ มานตฺตารโห โหติ, โส เจ ภิกฺขูนํ สนฺติเก ทูตํ ปหิเณยฺย “อหํ หิ มานตฺตารโห, อาคจฺฉนฺตุ ภิกฺขู, อิจฺฉามิ ภิกฺขูนํ อาคตนฺ”ติ.}

\prose{คนฺตพฺพํ ภิกฺขเว สตฺตาห\-กรณีเยน อปฺปหิเตปิ, ปเคว ปหิเต “มานตฺตทานํ อุสฺสุกฺกํ กริสฺสามิ วา อนุสฺสาเวสฺสามิ วา คณปูรโก วา ภวิสฺสามี”ติ, สตฺตาหํ สนฺนิวตฺโต กาตพฺโพ.}

\prose{อิธ ปน ภิกฺขเว ภิกฺขุ อพฺภานารโห โหติ, โส เจ ภิกฺขูนํ สนฺติเก ทูตํ ปหิเณยฺย “อหํ หิ อพฺภานารโห, อาคจฺฉนฺตุ ภิกฺขู, อิจฺฉามิ ภิกฺขูนํ อาคตนฺ”ติ.}

\prose{คนฺตพฺพํ ภิกฺขเว สตฺตาห\-กรณีเยน อปฺปหิเตปิ, ปเคว ปหิเต “อพฺภานํ อุสฺสุกฺกํ กริสฺสามิ วา อนุสฺสาเวสฺสามิ วา คณปูรโก วา ภวิสฺสามี”ติ, สตฺตาหํ สนฺนิวตฺโต กาตพฺโพ.}

\prose{\csromanfolio{201}อิธ ปน ภิกฺขเว ภิกฺขุสฺส สํโฆ กมฺมํ กตฺตุกาโม โหติ ตชฺชนียํ วา นิยสฺสํ วา ปพฺพาชนียํ วา ปฏิสารณียํ วา อุกฺเขปนียํ วา, โส เจ ภิกฺขูนํ สนฺติเก ทูตํ ปหิเณยฺย “สํโฆ เม กมฺมํ กตฺตุกาโม, อาคจฺฉนฺตุ ภิกฺขู, อิจฺฉามิ ภิกฺขูนํ อาคตนฺ”ติ.}

\prose{คนฺตพฺพํ ภิกฺขเว สตฺตาห\-กรณีเยน อปฺปหิเตปิ, ปเคว ปหิเต “กินฺติ นุ โข สํโฆ กมฺมํ น กเรยฺย, ลหุกาย วา ปริณาเมยฺยา”ติ, สตฺตาหํ สนฺนิวตฺโต กาตพฺโพ.}

\prose{กตํ วา ปนสฺส โหติ สํเฆน กมฺมํ ตชฺชนียํ วา นิยสฺสํ วา ปพฺพาชนียํ วา ปฏิสารณียํ วา อุกฺเขปนียํ วา, โส เจ ภิกฺขูนํ สนฺติเก ทูตํ ปหิเณยฺย “สํโฆ เม กมฺมํ อกาสิ, อาคจฺฉนฺตุ ภิกฺขู, อิจฺฉามิ ภิกฺขูนํ อาคตนฺ”ติ.}

\prose{คนฺตพฺพํ ภิกฺขเว สตฺตาห\-กรณีเยน อปฺปหิเตปิ, ปเคว ปหิเต “กินฺติ นุ โข สมฺมา วตฺเตยฺย โลมํ ปาเตยฺย เนตฺถารํ วตฺเตยฺย สํโฆ ตํ กมฺมํ ปฏิปฺปสฺสมฺเภยฺยา”ติ, สตฺตาหํ สนฺนิวตฺโต กาตพฺโพ.}

\proseitem{194}{อิธ ปน ภิกฺขเว ภิกฺขุนี คิลานา โหติ, สา เจ ภิกฺขูนํ สนฺติเก ทูตํ ปหิเณยฺย “อหํ หิ คิลานา, อาคจฺฉนฺตุ อยฺยา, อิจฺฉามิ อยฺยานํ อาคตนฺ”ติ.}

\prose{คนฺตพฺพํ ภิกฺขเว สตฺตาห\-กรณีเยน อปฺปหิเตปิ, ปเคว ปหิเต “คิลานภตฺตํ วา ปริเยสิสฺสามิ คิลานุปฏฺฐาก\-ภตฺตํ วา ปริเยสิสฺสามิ คิลาน\-เภสชฺชํ วา ปริเยสิสฺสามิ ปุจฺฉิสฺสามิ วา อุปฏฺฐหิสฺสามิ วา”ติ, สตฺตาหํ สนฺนิวตฺโต กาตพฺโพ.}

\prose{อิธ ปน ภิกฺขเว ภิกฺขุนิยา อนภิรติ อุปฺปนฺนา โหติ, สา เจ ภิกฺขูนํ สนฺติเก ทูตํ ปหิเณยฺย “อนภิรติ เม อุปฺปนฺนา, อาคจฺฉนฺตุ อยฺยา, อิจฺฉามิ อยฺยานํ อาคตนฺ”ติ.}

\prose{คนฺตพฺพํ ภิกฺขเว สตฺตาห\-กรณีเยน อปฺปหิเตปิ, ปเคว ปหิเต “อนภิรตํ วูปกาเสสฺสามิ วา วูปกาสาเปสฺสามิ วา ธมฺมกถํ วาสฺสา กริสฺสามี”ติ, สตฺตาหํ สนฺนิวตฺโต กาตพฺโพ.}

\proseitem{194}{อิธ ปน ภิกฺขเว ภิกฺขุนิยา กุกฺกุจฺจํ อุปฺปนฺนํ โหติ, สา เจ ภิกฺขูนํ สนฺติเก ทูตํ ปหิเณยฺย “กุกฺกุจฺจํ เม อุปฺปนฺนํ, อาคจฺฉนฺตุ อยฺยา, อิจฺฉามิ อยฺยานํ อาคตนฺ”ติ. คนฺตพฺพํ ภิกฺขเว สตฺตาห\-กรณีเยน อปฺปหิเตปิ, \csromanfolio{202}ปเคว ปหิเต “กุกฺกุจฺจํ วิโนเทสฺสามิ วา วิโนทาเปสฺสามิ วา ธมฺมกถํ วาสฺสา กริสฺสามี”ติ, สตฺตาหํ สนฺนิวตฺโต กาตพฺโพ.}

\prose{อิธ ปน ภิกฺขเว ภิกฺขุนิยา ทิฏฺฐิคตํ อุปฺปนฺนํ โหติ, สา เจ ภิกฺขูนํ สนฺติเก ทูตํ ปหิเณยฺย “ทิฏฺฐิคตํ เม อุปฺปนฺนํ, อาคจฺฉนฺตุ อยฺยา, อิจฺฉามิ อยฺยานํ อาคตนฺ”ติ.}

\prose{คนฺตพฺพํ ภิกฺขเว สตฺตาห\-กรณีเยน อปฺปหิเตปิ, ปเคว ปหิเต “ทิฏฺฐิคตํ วิเวเจสฺสามิ วา วิเวจาเปสฺสามิ วา ธมฺมกถํ วาสฺสา กริสฺสามี”ติ, สตฺตาหํ สนฺนิวตฺโต กาตพฺโพ.}

\prose{อิธ ปน ภิกฺขเว ภิกฺขุนี ครุธมฺมํ อชฺฌาปนฺนา โหติ มานตฺตารหา, สา เจ ภิกฺขูนํ สนฺติเก ทูตํ ปหิเณยฺย “อหํ หิ ครุธมฺมํ อชฺฌาปนฺนา มานตฺตารหา, อาคจฺฉนฺตุ อยฺยา, อิจฺฉามิ อยฺยานํ อาคตนฺ”ติ.}

\prose{คนฺตพฺพํ ภิกฺขเว สตฺตาห\-กรณีเยน อปฺปหิเตปิ, ปเคว ปหิเต “มานตฺตทานํ อุสฺสุกฺกํ กริสฺสามี”ติ, สตฺตาหํ สนฺนิวตฺโต กาตพฺโพ.}

\prose{อิธ ปน ภิกฺขเว ภิกฺขุนี มูลาย ปฏิกสฺสนา\-รหา โหติ, สา เจ ภิกฺขูนํ สนฺติเก ทูตํ ปหิเณยฺย “อหํ หิ มูลาย ปฏิกสฺสนา\-รหา, อาคจฺฉนฺตุ อยฺยา, อิจฺฉามิ อยฺยานํ อาคตนฺ”ติ.}

\prose{คนฺตพฺพํ ภิกฺขเว สตฺตาหกรณีเยนฺ อปฺปหิเตปิ, ปเคว ปหิเต “มูลาย ปฏิกสฺสนํ อุสฺสุกฺกํ กริสฺสามี”ติ, สตฺตาหํ สนฺนิวตฺโต กาตพฺโพ.}

\prose{อิธ ปน ภิกฺขเว ภิกฺขุนี อพฺภานารหา โหติ, สา เจ ภิกฺขูนํ สนฺติเก ทูตํ ปหิเณยฺย “อหํ หิ อพฺภานารหา, อาคจฺฉนฺตุ อยฺยา, อิจฺฉามิ อยฺยานํ อาคตนฺ”ติ.}

\prose{คนฺตพฺพํ ภิกฺขเว สตฺตาห\-กรณีเยน อปฺปหิเตปิ, ปเคว ปหิเต “อพฺภานํ อุสฺสุกฺกํ กริสฺสามี”ติ, สตฺตาหํ สนฺนิวตฺโต กาตพฺโพ.}

\prose{อิธ ปน ภิกฺขเว ภิกฺขุนิยา สํโฆ กมฺมํ กตฺตุกาโม โหติ ตชฺชนียํ วา นิยสฺสํ วา ปพฺพาชนียํ วา ปฏิสารณียํ วา อุกฺเขปนียํ วา, สา เจ ภิกฺขูนํ สนฺติเก ทูตํ ปหิเณยฺย “สํโฆ เม กมฺมํ กตฺตุกาโม, อาคจฺฉนฺตุ อยฺยา, อิจฺฉามิ อยฺยานํ อาคตนฺ”ติ.}

\prose{คนฺตพฺพํ ภิกฺขเว สตฺตาห\-กรณีเยน อปฺปหิเตปิ, ปเคว ปหิเต “กินฺติ นุ โข สํโฆ กมฺมํ น กเรยฺย ลหุกาย วา ปริณาเมยฺยา”ติ, สตฺตาหํ สนฺนิวตฺโต กาตพฺโพ.}

\prose{\csromanfolio{203}กตํ วา ปนสฺสา โหติ สํเฆน กมฺมํ ตชฺชนียํ วา นิยสฺสํ วา ปพฺพาชนียํ วา ปฏิสารณียํ วา อุกฺเขปนียํ วา, สา เจ ภิกฺขูนํ สนฺติเก ทูตํ ปหิเณยฺย “สํโฆ เม กมฺมํ อกาสิ, อาคจฺฉนฺตุ อยฺยา, อิจฺฉามิ อยฺยานํ อาคตนฺ”ติ.}

\prose{คนฺตพฺพํ ภิกฺขเว สตฺตาห\-กรณีเยน อปฺปหิเตปิ, ปเคว ปหิเต “กินฺติ นุ โข สมฺมา วตฺเตยฺย โลมํ ปาเตยฺย เนตฺถารํ วตฺเตยฺย สํโฆ ตํ กมฺมํ ปฏิปฺปสฺสมฺเภยฺยา”ติ, สตฺตาหํ สนฺนิวตฺโต กาตพฺโพ.}

\proseitem{195}{อิธ ปน ภิกฺขเว สิกฺขมานา คิลานา โหติ, สา เจ ภิกฺขูนํ สนฺติเก ทูตํ ปหิเณยฺย “อหํ หิ คิลานา, อาคจฺฉนฺตุ อยฺยา, อิจฺฉามิ อยฺยานํ อาคตนฺ”ติ.}

\prose{คนฺตพฺพํ ภิกฺขเว สตฺตาห\-กรณีเยน อปฺปหิเตปิ, ปเคว ปหิเต “คิลานภตฺตํ วา ปริเยสิสฺสามิ คิลานุปฏฺฐาก\-ภตฺตํ วา ปริเยสิสฺสามิ คิลาน\-เภสชฺชํ วา ปริเยสิสฺสามิ ปุจฺฉิสฺสามิ วา อุปฏฺฐหิสฺสามิ วา”ติ, สตฺตาหํ สนฺนิวตฺโต กาตพฺโพ.}

\prose{อิธ ปน ภิกฺขเว สิกฺขมานาย อนภิรติ อุปฺปนฺนา โหติ ฯเปฯ สิกฺขมานาย กุกฺกุจฺจํ อุปฺปนฺนํ โหติ. สิกฺขมานาย ทิฏฺฐิคตํ อุปฺปนฺนํ โหติ. สิกฺขมานาย สิกฺขา กุปิตา โหติ, สา เจ ภิกฺขูนํ สนฺติเก ทูตํ ปหิเณยฺย “สิกฺขา เม กุปิตา, อาคจฺฉนฺตุ อยฺยา, อิจฺฉามิ อยฺยานํ อาคตนฺ”ติ. คนฺตพฺพํ ภิกฺขเว สตฺตาห\-กรณีเยน อปฺปหิเตปิ, ปเคว ปหิเต “สิกฺขา\-สมาทานํ อุสฺสุกฺกํ กริสฺสามี”ติ, สตฺตาหํ สนฺนิวตฺโต กาตพฺโพ.}

\prose{อิธ ปน ภิกฺขเว สิกฺขมานา อุปสมฺปชฺชิตุ\-กามา โหติ, สา เจ ภิกฺขูนํ สนฺติเก ทูตํ ปหิเณยฺย “อหํ หิ อุปสมฺปชฺชิตุ\-กามา, อาคจฺฉนฺตุ อยฺยา, อิจฺฉามิ อยฺยานํ อาคตนฺ”ติ.}

\prose{คนฺตพฺพํ ภิกฺขเว สตฺตาห\-กรณีเยน อปฺปหิเตปิ, ปเคว ปหิเต “อุปสมฺปทํ อุสฺสุกฺกํ กริสฺสามิ วา อนุสฺสาเวสฺสามิ วา คณปูรโก วา ภวิสฺสามี”ติ, สตฺตาหํ สนฺนิวตฺโต กาตพฺโพ.}

\proseitem{196}{อิธ ปน ภิกฺขเว สามเณโร คิลาโน โหติ, โส เจ ภิกฺขูนํ สนฺติเก ทูตํ ปหิเณยฺย “อหํ หิ คิลาโน, อาคจฺฉนฺตุ \csromanfolio{204}ภิกฺขู, อิจฺฉามิ ภิกฺขูนํ อาคตนฺ”ติ. คนฺตพฺพํ ภิกฺขเว สตฺตาห\-กรณีเยน อปฺปหิเตปิ, ปเคว ปหิเต “คิลานภตฺตํ วา ปริเยสิสฺสามิ คิลานุปฏฺฐาก\-ภตฺตํ วา ปริเยสิสฺสามิ คิลาน\-เภสชฺชํ วา ปริเยสิสฺสามิ ปุจฺฉิสฺสามิ วา อุปฏฺฐหิสฺสามิ วา”ติ, สตฺตาหํ สนฺนิวตฺโต กาตพฺโพ.}

\prose{อิธ ปน ภิกฺขเว สามเณรสฺส อนภิรติ อุปฺปนฺนา โหติ ฯเปฯ สามเณรสฺส กุกฺกุจฺจํ อุปฺปนฺนํ โหติ. สามเณรสฺส ทิฏฺฐิคตํ อุปฺปนฺนํ โหติ. สามเณโร วสฺสํ ปุจฺฉิตุกาโม โหติ, โส เจ ภิกฺขูนํ สนฺติเก ทูตํ ปหิเณยฺย “อหํ หิ วสฺสํ ปุจฺฉิตุกาโม, อาคจฺฉนฺตุ ภิกฺขู, อิจฺฉามิ ภิกฺขูนํ อาคตนฺ”ติ. คนฺตพฺพํ ภิกฺขเว สตฺตาห\-กรณีเยน อปฺปหิเตปิ, ปเคว ปหิเต “ปุจฺฉิสฺสามิ วา อาจิกฺขิสฺสามิ วา”ติ, สตฺตาหํ สนฺนิวตฺโต กาตพฺโพ.}

\prose{อิธ ปน ภิกฺขเว สามเณโร อุปสมฺปชฺชิตุ\-กาโม โหติ, โส เจ ภิกฺขูนํ สนฺติเก ทูตํ ปหิเณยฺย “อหํ หิ อุปสมฺปชฺชิตุ\-กาโม, อาคจฺฉนฺตุ ภิกฺขู, อิจฺฉามิ ภิกฺขูนํ อาคตนฺ”ติ.}

\prose{คนฺตพฺพํ ภิกฺขเว สตฺตาห\-กรณีเยน อปฺปหิเตปิ, ปเคว ปหิเต “อุปสมฺปทํ อุสฺสุกฺกํ กริสฺสามิ วา อนุสฺสาเวสฺสามิ วา คณปูรโก วา ภวิสฺสามี”ติ, สตฺตาหํ สนฺนิวตฺโต กาตพฺโพ.}

\proseitem{197}{อิธ ปน ภิกฺขเว สามเณรี คิลานา โหติ, สา เจ ภิกฺขูนํ สนฺติเก ทูตํ ปหิเณยฺย “อหํ หิ คิลานา, อาคจฺฉนฺตุ อยฺยา, อิจฺฉามิ อยฺยานํ อาคตนฺ”ติ.}

\prose{คนฺตพฺพํ ภิกฺขเว สตฺตาหกรณียน อปฺปหิเตปิ, ปเคว ปหิเต “คิลานภตฺตํ วา ปริเยสิสฺสามิ คิลานุปฏฺฐาก\-ภตฺตํ วา ปริเยสิสฺสามิ คิลาน\-เภสชฺชํ วา ปริเยสิสฺสามิ ปุจฺฉิสฺสามิ วา อุปฏฺฐหิสฺสามิ วา”ติ, สตฺตาหํ สนฺนิวตฺโต กาตพฺโพ.}

\prose{อิธ ปน ภิกฺขเว สามเณริยา อนภิรติ อุปฺปนฺนา โหติ ฯเปฯ สามเณริยา กุกฺกุจฺจํ อุปฺปนฺนํ โหติ. สามเณริยา ทิฏฺฐิคตํ อุปฺปนฺนํ โหติ. สามเณรี วสฺสํ ปุจฺฉิตุกามา โหติ, สา เจ ภิกฺขูนํ สนฺติเก ทูตํ ปหิเณยฺย “อหํ หิ วสฺสํ ปุจฺฉิตุกามา, อาคจฺฉนฺตุ อยฺยา, อิจฺฉามิ อยฺยานํ อาคตนฺ”ติ. คนฺตพฺพํ ภิกฺขเว สตฺตาห\-กรณีเยน อปฺปหิเตปิ, ปเคว ปหิเต “ปุจฺฉิสฺสามิ วา อาจิกฺขิสฺสามิ วา”ติ, สตฺตาหํ สนฺนิวตฺโต กาตพฺโพ.}

\prose{\csromanfolio{205}อิธ ปน ภิกฺขเว สามเณรี สิกฺขํ สมาทิยิตุกามา โหติ, สา เจ ภิกฺขูนํ สนฺติเก ทูตํ ปหิเณยฺย “อหํ หิ สิกฺขํ สมาทิยิตุกามา, อาคจฺฉนฺตุ อยฺยา, อิจฺฉามิ อยฺยานํ อาคตนฺ”ติ.}

\prose{คนฺตพฺพํ ภิกฺขเว สตฺตาห\-กรณีเยน อปฺปหิเตปิ, ปเคว ปหิเต “สิกฺขา\-สมาทานํ อุสฺสุกฺกํ กริสฺสามี”ติ, สตฺตาหํ สนฺนิวตฺโต กาตพฺโพติ.}

\csromanmatikaanchor{mtk.114}
\csromantocmarkref{section}{111. สตฺตนฺนํ อปฺปหิเตปิ อนุชานนา}{mtk.114}
\bookmark[dest={mtk.114},level=1]{111. สตฺตนฺนํ อปฺปหิเตปิ อนุชานนา}
\csromanheader{1}{111. \textbf{สตฺตนฺนํ อปฺปหิเตปิ อนุชานนา}}
\proseitem{198}{เตน โข ปน สมเยน อญฺญตรสฺส ภิกฺขุโน มาตา คิลานา โหติ, สา ปุตฺตสฺส สนฺติเก ทูตํ ปาเหสิ “อหํ หิ คิลานา, อาคจฺฉตุ เม ปุตฺโต, อิจฺฉามิ ปุตฺตสฺส อาคตนฺ”ติ.}

\prose{อถ โข ตสฺส ภิกฺขุโน เอตทโหสิ “ภควตา ปญฺญตฺตํ “สตฺตนฺนํ สตฺตาห\-กรณีเยน ปหิเต คนฺตุํ, น เตฺวว อปฺปหิเต, ปญฺจนฺนํ สตฺตาห\-กรณีเยน อปฺปหิเตปิ คนฺตุํ, ปเคว ปหิเต’ติ, อยญฺจ เม มาตา คิลานา, สา จ อนุปาสิกา, กถํ นุ โข มยา ปฏิปชฺชิตพฺพนฺ”ติ. ภควโต เอตมตฺถํ อาโรเจสุํ. อนุชานามิ ภิกฺขเว สตฺตนฺนํ สตฺตาห\-กรณีเยน อปฺปหิเตปิ คนฺตุํ, ปเคว ปหิเต, ภิกฺขุสฺส ภิกฺขุนิยา สิกฺขมานาย สามเณรสฺส สามเณริยา มาตุยา จ ปิตุสฺส จ. อนุชานามิ ภิกฺขเว อิเมสํ สตฺตนฺนํ สตฺตาห\-กรณีเยน อปฺปหิเตปิ คนฺตุํ, ปเคว ปหิเต, สตฺตาหํ สนฺนิวตฺโต กาตพฺโพ.}

\prose{อิธ ปน ภิกฺขเว ภิกฺขุสฺส มาตา คิลานา โหติ, สา เจ ปุตฺตสฺส สนฺติเก ทูตํ ปหิเณยฺย “อหํ หิ คิลานา, อาคจฺฉตุ เม ปุตฺโต, อิจฺฉามิ ปุตฺตสฺส อาคตนฺ”ติ.}

\prose{คนฺตพฺพํ ภิกฺขเว สตฺตาห\-กรณีเยน อปฺปหิเตปิ, ปเคว ปหิเต “คิลานภตฺตํ วา ปริเยสิสฺสามิ คิลานุปฏฺฐาก\-ภตฺตํ วา ปริเยสิสฺสามิ คิลาน\-เภสชฺชํ วา ปริเยสิสฺสามิ ปุจฺฉิสฺสามิ วา อุปฏฺฐหิสฺสามิ วา”ติ, สตฺตาหํ สนฺนิวตฺโต กาตพฺโพ.}

\prose{อิธ ปน ภิกฺขเว ภิกฺขุสฺส ปิตา คิลาโน โหติ, โส เจ ปุตฺตสฺส สนฺติเก ทูตํ ปหิเณยฺย “อหํ หิ คิลาโน, อาคจฺฉตุ เม ปุตฺโต, อิจฺฉามิ ปุตฺตสฺส อาคตนฺ”ติ.}

\prose{คนฺตพฺพํ ภิกฺขเว สตฺตาห\-กรณีเยน อปฺปหิเตปิ, ปเคว ปหิเต “คิลานภตฺตํ วา ปริเยสิสฺสามิ คิลานุปฏฺฐาก\-ภตฺตํ วา ปริเยสิสฺสามิ คิลาน\-เภสชฺชํ วา ปริเยสิสฺสามิ ปุจฺฉิสฺสามิ วา อุปฏฺฐหิสฺสามิ วา”ติ, สตฺตาหํ สนฺนิวตฺโต กาตพฺโพ.}

\csromanmatikaanchor{mtk.115}
\csromantocmarkref{section}{112. ปหิเตเยว อนุชานนา}{mtk.115}
\bookmark[dest={mtk.115},level=1]{112. ปหิเตเยว อนุชานนา}
\csromanheader{1}{112. \textbf{ปหิเตเยว อนุชานนา}}
\proseitem{199}{\csromanfolio{206}อิธ ปน ภิกฺขเว ภิกฺขุสฺส ภาตา คิลาโน โหติ, โส เจ ภาตุโน สนฺติเก ทูตํ ปหิเณยฺย “อหํ หิ คิลาโน, อาคจฺฉตุ เม ภาตา, อิจฺฉามิ ภาตุโน อาคตนฺ”ติ.}

\prose{คนฺตพฺพํ ภิกฺขเว สตฺตาห\-กรณีเยน ปหิเต, น เตฺวว อปฺปหิเต, สตฺตาหํ สนฺนิวตฺโต กาตพฺโพ.}

\prose{อิธ ปน ภิกฺขเว ภิกฺขุสฺส ภคินี คิลานา โหติ, สา เจ ภาตุโน สนฺติเก ทูตํ ปหิเณยฺย “อหํ หิ คิลานา, อาคจฺฉตุ เม ภาตา, อิจฺฉามิ ภาตุโน อาคตนฺ”ติ.}

\prose{คนฺตพฺพํ ภิกฺขเว สตฺตาห\-กรณีเยน ปหิเต, น เตฺวว อปฺปหิเต, สตฺตาหํ สนฺนิวตฺโต กาตพฺโพ.}

\prose{อิธ ปน ภิกฺขเว ภิกฺขุสฺส ญาตโก คิลาโน โหติ, โส เจ ภิกฺขุสฺส สนฺติเก ทูตํ ปหิเณยฺย “อหํ หิ คิลาโน, อาคจฺฉตุ ภทนฺโต, อิจฺฉามิ ภทนฺตสฺส อาคตนฺ”ติ คนฺตพฺพํ ภิกฺขเว สตฺตาห\-กรณีเยน ปหิเต, น เตฺวว อปฺปหิเต, สตฺตาหํ สนฺนิวตฺโต กาตพฺโพ.}

\prose{อิธ ปน ภิกฺขเว ภิกฺขุคติโก คิลาโน โหติ, โส เจ ภิกฺขูนํ สนฺติเก ทูตํ ปหิเณยฺย “อหํ หิ คิลาโน, อาคจฺฉนฺตุ ภทนฺตา, อิจฺฉามิ ภทนฺตานํ อาคตนฺ”ติ.}

\prose{คนฺตพฺพํ ภิกฺขเว สตฺตาห\-กรณีเยน ปหิเต, น เตฺวว อปฺปหิเต, สตฺตาหํ สนฺนิวตฺโต กาตพฺโพ.}

\prose{เตน โข ปน สมเยน สํฆสฺส วิหาโร อุนฺทฺริยติ, อญฺญตเรน อุปาสเกน อรญฺเญ ภณฺฑํ เฉทาปิตํ โหติ.}

\prose{โส ภิกฺขูนํ สนฺติเก ทูตํ ปาเหสิ “สเจ ภทนฺตา ตํ ภณฺฑํ อาวหาเปยฺยุํ, ทชฺชาหํ ตํ ภณฺฑนฺ”ติ.}

\prosewithcloserruled{ภควโต เอตมตฺถํ อาโรเจสุํ. อนุชานามิ ภิกฺขเว สํฆกรณีเยน คนฺตุํ, สตฺตาหํ สนฺนิวตฺโต กาตพฺโพติ.}{วสฺสาวาส\-ภาณวาโร นิฏฺฐิโต.}
\csromanmatikaanchor{mtk.116}
\csromantocmarkref{section}{113. อนฺตราเย อนาปตฺติวสฺสจฺเฉทวาร}{mtk.116}
\bookmark[dest={mtk.116},level=1]{113. อนฺตราเย อนาปตฺติวสฺสจฺเฉทวาร}
\csromanheader{1}{113. \textbf{อนฺตราเย อนาปตฺติ}\-\textbf{วสฺสจฺเฉท}\-\textbf{วาร}}
\proseitem{200}{เตน โข ปน สมเยน โกสเลสุ ชนปเท อญฺญตรสฺมิํ อาวาเส วสฺสูปคตา ภิกฺขู วาเฬหิ อุพฺพาฬฺหา โหนฺติ, คณฺหิํสุปิ ปริปาติํสุปิ.}

\setlength{\csromangathapagewidth}{0pt}
\setlength{\csromangathawakwidth}{0pt}
\csromangathameasuregroup{}{ภควโต เอตมตฺถํ อาโรเจสุํ.}
\csromangathameasurewak{ภควโต เอตมตฺถํ อาโรเจสุํ.}
\setlength{\csromangathaleftcol}{0pt}
\csromangathagroup{\csromangathastanza{\csromangathawak{\csromangathaitemline{200}{ภควโต เอตมตฺถํ อาโรเจสุํ.}}}}

\prose{\csromanfolio{207}อิธ ปน ภิกฺขเว วสฺสูปคตา ภิกฺขู วาเฬหิ อุพฺพาฬฺหา โหนฺติ, คณฺหนฺติปิ ปริปาเตนฺติปิ.}

\setlength{\csromangathapagewidth}{0pt}
\setlength{\csromangathawakwidth}{0pt}
\csromangathameasuregroup{เอเสว อนฺตราโยติ ปกฺกมิตพฺพํ,}{\csromangathabat{เอเสว อนฺตราโยติ ปกฺกมิตพฺพํ,}{อนาปตฺติ วสฺสจฺเฉทสฺส.}}
\csromangathasetleft{เอเสว อนฺตราโยติ ปกฺกมิตพฺพํ,}
\csromangathagroup{\csromangathastanza{\csromangathabat{เอเสว อนฺตราโยติ ปกฺกมิตพฺพํ,}{อนาปตฺติ วสฺสจฺเฉทสฺส.}}}

\prose{อิธ ปน ภิกฺขเว วสฺสูปคตา ภิกฺขู สรีสเปหิ อุพฺพาฬฺหา โหนฺติ, ฑํสนฺติปิ ปริปาเตนฺติปิ.}

\setlength{\csromangathapagewidth}{0pt}
\setlength{\csromangathawakwidth}{0pt}
\csromangathameasuregroup{เอเสว อนฺตราโยติ ปกฺกมิตพฺพํ,}{\csromangathabat{เอเสว อนฺตราโยติ ปกฺกมิตพฺพํ,}{อนาปตฺติ วสฺสจฺเฉทสฺส.}}
\csromangathasetleft{เอเสว อนฺตราโยติ ปกฺกมิตพฺพํ,}
\csromangathagroup{\csromangathastanza{\csromangathabat{เอเสว อนฺตราโยติ ปกฺกมิตพฺพํ,}{อนาปตฺติ วสฺสจฺเฉทสฺส.}}}

\prose{อิธ ปน ภิกฺขเว วสฺสูปคตา ภิกฺขู โจเรหิ อุพฺพาฬฺหา โหนฺติ, วิลุมฺปนฺติปิ อาโกเฏนฺติปิ.}

\setlength{\csromangathapagewidth}{0pt}
\setlength{\csromangathawakwidth}{0pt}
\csromangathameasuregroup{เอเสว อนฺตราโยติ ปกฺกมิตพฺพํ,}{\csromangathabat{เอเสว อนฺตราโยติ ปกฺกมิตพฺพํ,}{อนาปตฺติ วสฺสจฺเฉทสฺส.}}
\csromangathasetleft{เอเสว อนฺตราโยติ ปกฺกมิตพฺพํ,}
\csromangathagroup{\csromangathastanza{\csromangathabat{เอเสว อนฺตราโยติ ปกฺกมิตพฺพํ,}{อนาปตฺติ วสฺสจฺเฉทสฺส.}}}

\prose{อิธ ปน ภิกฺขเว วสฺสูปคตา ภิกฺขู ปิสาเจหิ อุพฺพาฬฺหา โหนฺติ, อาวิสนฺติปิ หนนฺติปิ\csromansharedfootnote{a3a70d86b797656e}{โอชมฺปิ หรนฺติ (สี), หรนฺติปิ (สฺยา)}.}

\setlength{\csromangathapagewidth}{0pt}
\setlength{\csromangathawakwidth}{0pt}
\csromangathameasuregroup{เอเสว อนฺตราโยติ ปกฺกมิตพฺพํ,}{\csromangathabat{เอเสว อนฺตราโยติ ปกฺกมิตพฺพํ,}{อนาปตฺติ วสฺสจฺเฉทสฺส.}}
\csromangathasetleft{เอเสว อนฺตราโยติ ปกฺกมิตพฺพํ,}
\csromangathagroup{\csromangathastanza{\csromangathabat{เอเสว อนฺตราโยติ ปกฺกมิตพฺพํ,}{อนาปตฺติ วสฺสจฺเฉทสฺส.}}}

\prose{อิธ ปน ภิกฺขเว วสฺสูปคตานํ ภิกฺขูนํ คาโม อคฺคินา ทฑฺโฒ โหติ, ภิกฺขู ปิณฺฑเกน กิลมนฺติ.}

\setlength{\csromangathapagewidth}{0pt}
\setlength{\csromangathawakwidth}{0pt}
\csromangathameasuregroup{เอเสว อนฺตราโยติ ปกฺกมิตพฺพํ,}{\csromangathabat{เอเสว อนฺตราโยติ ปกฺกมิตพฺพํ,}{อนาปตฺติ วสฺสจฺเฉทสฺส.}}
\csromangathasetleft{เอเสว อนฺตราโยติ ปกฺกมิตพฺพํ,}
\csromangathagroup{\csromangathastanza{\csromangathabat{เอเสว อนฺตราโยติ ปกฺกมิตพฺพํ,}{อนาปตฺติ วสฺสจฺเฉทสฺส.}}}

\prose{อิธ ปน ภิกฺขเว วสฺสูปคตานํ ภิกฺขูนํ เสนาสนํ อคฺคินา ทฑฺฒํ โหติ, ภิกฺขู เสนาสเนน กิลมนฺติ.}

\setlength{\csromangathapagewidth}{0pt}
\setlength{\csromangathawakwidth}{0pt}
\csromangathameasuregroup{เอเสว อนฺตราโยติ ปกฺกมิตพฺพํ,}{\csromangathabat{เอเสว อนฺตราโยติ ปกฺกมิตพฺพํ,}{อนาปตฺติ วสฺสจฺเฉทสฺส.}}
\csromangathasetleft{เอเสว อนฺตราโยติ ปกฺกมิตพฺพํ,}
\csromangathagroup{\csromangathastanza{\csromangathabat{เอเสว อนฺตราโยติ ปกฺกมิตพฺพํ,}{อนาปตฺติ วสฺสจฺเฉทสฺส.}}}

\prose{อิธ ปน ภิกฺขเว วสฺสูปคตานํ ภิกฺขูนํ คาโม อุทเกน วูโฬฺห โหติ, ภิกฺขู ปิณฺฑเกน กิลมนฺติ.}

\setlength{\csromangathapagewidth}{0pt}
\setlength{\csromangathawakwidth}{0pt}
\csromangathameasuregroup{เอเสว อนฺตราโยติ ปกฺกมิตพฺพํ,}{\csromangathabat{เอเสว อนฺตราโยติ ปกฺกมิตพฺพํ,}{อนาปตฺติ วสฺสจฺเฉทสฺส.}}
\csromangathasetleft{เอเสว อนฺตราโยติ ปกฺกมิตพฺพํ,}
\csromangathagroup{\csromangathastanza{\csromangathabat{เอเสว อนฺตราโยติ ปกฺกมิตพฺพํ,}{อนาปตฺติ วสฺสจฺเฉทสฺส.}}}

\prose{อิธ ปน ภิกฺขเว วสฺสูปคตานํ ภิกฺขูนํ เสนาสนํ อุทเกน วูฬฺหํ โหติ, ภิกฺขู เสนาสเนน กิลมนฺติ.}

\setlength{\csromangathapagewidth}{0pt}
\setlength{\csromangathawakwidth}{0pt}
\csromangathameasuregroup{เอเสว อนฺตราโยติ ปกฺกมิตพฺพํ,}{\csromangathabat{เอเสว อนฺตราโยติ ปกฺกมิตพฺพํ,}{อนาปตฺติ วสฺสจฺเฉทสฺสาติ.}}
\csromangathasetleft{เอเสว อนฺตราโยติ ปกฺกมิตพฺพํ,}
\csromangathagroup{\csromangathastanza{\csromangathabat{เอเสว อนฺตราโยติ ปกฺกมิตพฺพํ,}{อนาปตฺติ วสฺสจฺเฉทสฺสาติ.}}}

\proseitem{201}{เตน โข ปน สมเยน อญฺญตรสฺมิํ อาวาเส วสฺสูปคตานํ ภิกฺขูนํ คาโม โจเรหิ วุฏฺฐาสิ. ภควโต เอตมตฺถํ อาโรเจสุํ. อนุชานามิ ภิกฺขเว เยน คาโม เตน คนฺตุนฺติ.}

\prose{\csromanfolio{208}คาโม เทฺวธา ภิชฺชิตฺถ. ภควโต เอตมตฺถํ อาโรเจสุํ. อนุชานามิ ภิกฺขเว เยน พหุตรา เตน คนฺตุนฺติ.}

\prose{พหุตรา อสฺสทฺธา โหนฺติ อปฺปสนฺนา. ภควโต เอตมตฺถํ อาโรเจสุํ. อนุชานามิ ภิกฺขเว เยน สทฺธา ปสนฺนา เตน คนฺตุนฺติ.}

\prose{เตน โข ปน สมเยน โกสเลสุ ชนปเท อญฺญตรสฺมิํ อาวาเส วสฺสูปคตา ภิกฺขู น ลภิํสุ ลูขสฺส วา ปณีตสฺส วา โภชนสฺส ยาวทตฺถํ ปาริปูริํ. ภควโต เอตมตฺถํ อาโรเจสุํ.}

\prose{อิธ ปน ภิกฺขเว วสฺสูปคตา ภิกฺขู น ลภนฺติ ลูขสฺส วา ปณีตสฺส วา โภชนสฺส ยาวทตฺถํ ปาริปูริํ. เอเสว อนฺตราโยติ ปกฺกมิตพฺพํ, อนาปตฺติ วสฺสจฺเฉทสฺส.}

\prose{อิธ ปน ภิกฺขเว วสฺสูปคตา ภิกฺขู ลภนฺติ ลูขสฺส วา ปณีตสฺส วา โภชนสฺส ยาวทตฺถํ ปาริปูริํ, น ลภนฺติ สปฺปายานิ โภชนานิ.}

\setlength{\csromangathapagewidth}{0pt}
\setlength{\csromangathawakwidth}{0pt}
\csromangathameasuregroup{เอเสว อนฺตราโยติ ปกฺกมิตพฺพํ,}{\csromangathabat{เอเสว อนฺตราโยติ ปกฺกมิตพฺพํ,}{อนาปตฺติ วสฺสจฺเฉทสฺส.}}
\csromangathasetleft{เอเสว อนฺตราโยติ ปกฺกมิตพฺพํ,}
\csromangathagroup{\csromangathastanza{\csromangathabat{เอเสว อนฺตราโยติ ปกฺกมิตพฺพํ,}{อนาปตฺติ วสฺสจฺเฉทสฺส.}}}

\prose{อิธ ปน ภิกฺขเว วสฺสูปคตา ภิกฺขู ลภนฺติ ลูขสฺส วา ปณีตสฺส วา โภชนสฺส ยาวทตฺถํ ปาริปูริํ, ลภนฺติ สปฺปายานิ โภชนานิ, น ลภนฺติ สปฺปายานิ เภสชฺชานิ.}

\setlength{\csromangathapagewidth}{0pt}
\setlength{\csromangathawakwidth}{0pt}
\csromangathameasuregroup{เอเสว อนฺตราโยติ ปกฺกมิตพฺพํ,}{\csromangathabat{เอเสว อนฺตราโยติ ปกฺกมิตพฺพํ,}{อนาปตฺติ วสฺสจฺเฉทสฺส.}}
\csromangathasetleft{เอเสว อนฺตราโยติ ปกฺกมิตพฺพํ,}
\csromangathagroup{\csromangathastanza{\csromangathabat{เอเสว อนฺตราโยติ ปกฺกมิตพฺพํ,}{อนาปตฺติ วสฺสจฺเฉทสฺส.}}}

\prose{อิธ ปน ภิกฺขเว วสฺสูปคตา ภิกฺขู ลภนฺติ ลูขสฺส วา ปณีตสฺส วา โภชนสฺส ยาวทตฺถํ ปาริปูริํ, ลภนฺติ สปฺปายานิ โภชนานิ, ลภนฺติ สปฺปายานิ เภสชฺชานิ, น ลภนฺติ ปติรูปํ อุปฏฺฐากํ.}

\setlength{\csromangathapagewidth}{0pt}
\setlength{\csromangathawakwidth}{0pt}
\csromangathameasuregroup{เอเสว อนฺตราโยติ ปกฺกมิตพฺพํ,}{\csromangathabat{เอเสว อนฺตราโยติ ปกฺกมิตพฺพํ,}{อนาปตฺติ วสฺสจฺเฉทสฺส.}}
\csromangathasetleft{เอเสว อนฺตราโยติ ปกฺกมิตพฺพํ,}
\csromangathagroup{\csromangathastanza{\csromangathabat{เอเสว อนฺตราโยติ ปกฺกมิตพฺพํ,}{อนาปตฺติ วสฺสจฺเฉทสฺส.}}}

\prose{อิธ ปน ภิกฺขเว วสฺสูปคตํ ภิกฺขุํ อิตฺถี นิมนฺเตติ “เอหิ ภนฺเต หิรญฺญํ วา เต เทมิ, สุวณฺณํ วา เต เทมิ, เขตฺตํ วา เต เทมิ, วตฺถุํ วา เต เทมิ, คาวุํ วา เต เทมิ, คาวิํ วา เต เทมิ, ทาสํ วา เต เทมิ, ทาสิํ วา เต เทมิ, ธีตรํ วา เต เทมิ ภริยตฺถาย, อหํ วา เต ภริยา โหมิ, อญฺญํ วา เต ภริยํ อาเนมี”ติ.}

\prose{ตตฺร เจ ภิกฺขุโน เอวํ โหติ “ลหุปริวตฺตํ โข จิตฺตํ วุตฺตํ ภควตา, สิยาปิ เม พฺรหฺม\-จริยสฺส อนฺตราโย”ติ ปกฺกมิตพฺพํ, อนาปตฺติ วสฺสจฺเฉทสฺส.}

\prose{\csromanfolio{209}อิธ ปน ภิกฺขเว วสฺสูปคตํ ภิกฺขุํ เวสี นิมนฺเตติ ฯเปฯ ถุลฺลกุมารี นิมนฺเตติ. ปณฺฑโก นิมนฺเตติ. ญาตกา นิมนฺเตนฺติ. ราชาโน นิมนฺเตนฺติ. โจรา นิมนฺเตนฺติ. ธุตฺตา นิมนฺเตนฺติ “เอหิ ภนฺเต หิรญฺญํ วา เต เทม, สุวณฺณํ วา เต เทม, เขตฺตํ วา เต เทม, วตฺถุํ วา เต เทม, คาวุํ วา เต เทม, คาวิํ วา เต เทม, ทาสํ วา เต เทม, ทาสิํ วา เต เทม, ธีตรํ วา เต เทม ภริยตฺถาย, อญฺญํ วา เต ภริยํ อาเนมา”ติ. ตตฺร เจ ภิกฺขุโน เอวํ โหติ “ลหุปริวตฺตํ โข จิตฺตํ วุตฺตํ ภควตา, สิยาปิ เม พฺรหฺม\-จริยสฺส อนฺตราโย”ติ ปกฺกมิตพฺพํ, อนาปตฺติ วสฺสจฺเฉทสฺส.}

\prose{อิธ ปน ภิกฺขเว วสฺสูปคโต ภิกฺขุ อสฺสามิกํ นิธิํ ปสฺสติ. ตตฺร เจ ภิกฺขุโน เอวํ โหติ “ลหุปริวตฺตํ โข จิตฺตํ วุตฺตํ ภควตา, สิยาปิ เม พฺรหฺม\-จริยสฺส อนฺตราโย”ติ ปกฺกมิตพฺพํ, อนาปตฺติ วสฺสจฺเฉทสฺส.}

\csromanmatikaanchor{mtk.117}
\csromantocmarkref{section}{114. สํฆเภเท อนาปตฺติวสฺสจฺเฉทวาร}{mtk.117}
\bookmark[dest={mtk.117},level=1]{114. สํฆเภเท อนาปตฺติวสฺสจฺเฉทวาร}
\csromanheader{1}{114. \textbf{สํฆเภเท อนาปตฺติ}\-\textbf{วสฺสจฺเฉท}\-\textbf{วาร}}
\proseitem{202}{อิธ ปน ภิกฺขเว วสฺสูปคโต ภิกฺขุ ปสฺสติ สมฺพหุเล ภิกฺขู สํฆเภทาย ปรกฺกมนฺเต. ตตฺร เจ ภิกฺขุโน เอวํ โหติ “ครุโก โข สํฆเภโท วุตฺโต ภควตา, มา มยิ สมฺมุขีภูเต สํโฆ ภิชฺชี”ติ ปกฺกมิตพฺพํ, อนาปตฺติ วสฺสจฺเฉทสฺส.}

\prose{อิธ ปน ภิกฺขเว วสฺสูปคโต ภิกฺขุ สุณาติ “อสุกสฺมิํ กิร อาวาเส สมฺพหุลา ภิกฺขุ สํฆเภทาย ปรกฺกมนฺตี”ติ. ตตฺร เจ ภิกฺขุโน เอวํ โหติ “ครุโก โข สํฆเหโท วุตฺโต ภควตา, มา มยิ สมฺมุขีภูเต สํโฆ ภิชฺชี”ติ ปกฺกมิตพฺพํ, อนาปตฺติ วสฺสจฺเฉทสฺส.}

\prose{อิธ ปน ภิกฺขเว วสฺสูปคโต ภิกฺขุ สุณาติ “อสุกสฺมิํ กิร อาวาเส สมฺพหุลา ภิกฺขู สํฆเภทาย ปรกฺกมนฺตี”ติ. ตตฺร เจ ภิกฺขุโน เอวํ โหติ “เต โข เม ภิกฺขู มิตฺตา, ตฺยาหํ วกฺขามิ ‘ครุโก โข อาวุโส สํฆเภโท วุตฺโต ภควตา, มายสฺมนฺตานํ สํฆเภโท รุจฺจิตฺถา’ติ, กริสฺสนฺติ เม วจนํ สุสฺสูสิสฺสนฺติ, โสตํ โอทหิสฺสนฺตี”ติ ปกฺกมิตพฺพํ, อนาปตฺติ วสฺสจฺเฉทสฺส.}

\prose{\csromanfolio{210}อิธ ปน ภิกฺขเว วสฺสูปคโต ภิกฺขุ สุณาติ “อสุกสฺมิํ กิร อาวาเส สมฺพหุลา ภิกฺขู สํฆเภทาย ปรกฺกมนฺตี”ติ. ตตฺร เจ ภิกฺขุโน เอวํ โหติ “เต โข เม ภิกฺขู น มิตฺตา, อปิ จ เย เตสํ มิตฺตา, เต เม มิตฺตา, ตฺยาหํ วกฺขามิ, เต วุตฺตา เต วกฺขนฺติ ‘ครุโก โข อาวุโส สํฆเภโท วุตฺโต ภควตา, มายสฺมนฺตานํ สํฆเภโท รุจฺจิตฺถา’ติ, กริสฺสนฺติ เตสํ วจนํ สุสฺสูสิสฺสนฺติ, โสตํ โอทหิสฺสนฺตี”ติ ปกฺกมิตพฺพํ, อนาปตฺติ วสฺสจฺเฉทสฺส.}

\prose{อิธ ปน ภิกฺขเว วสฺสูปคโต ภิกฺขุ สุณาติ “อสุกสฺมิํ กิร อาวาเส สมฺพหุเลหิ ภิกฺขูหิ สํโฆ ภินฺโน”ติ, ตตฺร เจ ภิกฺขุโน เอวํ โหติ “เต โข เม ภิกฺขู มิตฺตา, ตฺยาหํ วกฺขามิ ‘ครุโก โข อาวุโส สํฆเภโท วุตฺโต ภควตา, มายสฺมนฺตานํ สํฆเภโท รุจฺจิตฺถา’ติ, กริสฺสนฺติ เม วจนํ สุสฺสูสิสฺสนฺติ, โสตํ โอทหิสฺสนฺตี”ติ ปกฺกมิตพฺพํ, อนาปตฺติ วสฺสจฺเฉทสฺส.}

\prose{อิธ ปน ภิกฺขเว วสฺสูปคโต ภิกฺขุ สุณาติ “อสุกสฺมิํ กิร อาวาเส สมฺพหุเลหิ ภิกฺขูหิ สํโฆ ภินฺโน”ติ. ตตฺร เจ ภิกฺขุโน เอวํ โหติ “เต โข เม ภิกฺขู น มิตฺตา, อปิ จ เย เตสํ มิตฺตา, เต เม มิตฺตา, ตฺยาหํ วกฺขามิ, เต วุตฺตา เต วกฺขนฺติ ‘ครุโก โข อาวุโส สํฆเภโท วุตฺโต ภควตา, มายสฺมนฺตานํ สํฆเภโท รุจฺจิตฺถา’ติ, กริสฺสนฺติ เตสํ วจนํ สุสฺสูสิสฺสนฺติ, โสตํ โอทหิสฺสนฺตี”ติ ปกฺกมิตพฺพํ, อนาปตฺติ วสฺสจฺเฉทสฺส.}

\prose{อิธ ปน ภิกฺขเว วสฺสูปคโต ภิกฺขุ สุณาติ “อมุกสฺมิํ กิร อาวาเส สมฺพหุลา ภิกฺขุนิโย สํฆเภทาย ปรกฺกมนฺตี”ติ. ตตฺร เจ ภิกฺขุโน เอวํ โหติ “ตา โข เม ภิกฺขุนิโย มิตฺตา, ตาหํ วกฺขามิ ‘ครุโก โข ภคินิโย สํฆเภโท วุตฺโต ภควตา, มา ภคินีนํ สํฆเภโท รุจฺจิตฺถา’ติ, กริสฺสนฺติ เม วจนํ สุสฺสูสิสฺสนฺติ, โสตํ โอทหิสฺสนฺตี”ติ ปกฺกมิตพฺพํ, อนาปตฺติ วสฺสจฺเฉทสฺส.}

\proseitem{202}{อิธ ปน ภิกฺขเว วสฺสูปคโต ภิกฺขุ สุณาติ “อมุกสฺมิํ กิร อาวาเส สมฺพหุลา ภิกฺขุนิโย สํฆเภทาย ปรกฺกมนฺตี”ติ. ตตฺร เจ ภิกฺขุโน เอวํ โหติ “ตา โข เม ภิกฺขุนิโย น มิตฺตา, อปิ จ ยา ตาสํ มิตฺตา, ตา เม มิตฺตา, ตาหํ วกฺขามิ, ตา วุตฺตา ตา วกฺขนฺติ \csromanfolio{211}‘ครุโก โข ภคินิโย\csromansharedfootnote{7d892a5c6fb4306a}{อยฺยาโย (สี)} สํฆเภโท วุตฺโต ภควตา, มา ภคินีนํ\csromansharedfootnote{8bac68ae77ff4c76}{อยฺยานํ (สี)} สํฆเภโท รุจฺจิตฺถา’ติ, กริสฺสนฺติ ตาสํ วจนํ สุสฺสูสิสฺสนฺติ, โสตํ โอทหิสฺสนฺตี”ติ ปกฺกมิตพฺพํ, อนาปตฺติ วสฺสจฺเฉทสฺส.}

\prose{อิธ ปน ภิกฺขเว วสฺสูปคโต ภิกฺขุ สุณาติ “อมุกสฺมิํ กิร อาวาเส สมฺพหุลาหิ ภิกฺขุนีหิ สํโฆ ภินฺโน”ติ. ตตฺร เจ ภิกฺขุโน เอวํ โหติ “ตา โข เม ภิกฺขุนิโย มิตฺตา, ตาหํ วกฺขามิ ‘ครุโก โข ภคินิโย สํฆเภโท วุตฺโต ภควตา, มา ภคินีนํ สํฆเภโท รุจฺจิตฺถา’ติ, กริสฺสนฺติ เม วจนํ สุสฺสูสิสฺสนฺติ, โสตํ โอทหิสฺสนฺตี”ติ ปกฺกมิตพฺพํ, อนาปตฺติ วสฺสจฺเฉทสฺส.}

\prose{อิธ ปน ภิกฺขเว วสฺสูปคโต ภิกฺขุ สุณาติ “อมุกสฺมิํ กิร อาวาเส สมฺพหุลาหิ ภิกฺขุนีหิ สํโฆ ภินฺโน”ติ. ตตฺร เจ ภิกฺขุโน เอวํ โหติ “ตา โข เม ภิกฺขุนิโย น มิตฺตา, อปิ จ ยา ตาสํ มิตฺตา, ตา เม มิตฺตา, ตาหํ วกฺขามิ, ตา วุตฺตา ตา วกฺขนฺติ ‘ครุโก โข ภคินิโย\csromansharedfootnote{7d892a5c6fb4306a}{อยฺยาโย (สี)} สํฆเภโท วุตฺโต ภควตา, มา ภคินีนํ\csromansharedfootnote{8bac68ae77ff4c76}{อยฺยานํ (สี)} สํฆเภโท รุจฺจิตฺถา’ติ, กริสฺสนฺติ ตาสํ วจนํ สุสฺสูสิสฺสนฺติ, โสตํ โอทหิสฺสนฺตี”ติ ปกฺกมิตพฺพํ, อนาปตฺติ วสฺสจฺเฉทสฺสาติ.}

\csromanmatikaanchor{mtk.118}
\csromantocmarkref{section}{115. วชาทีสุ วสฺสูปคมนํ}{mtk.118}
\bookmark[dest={mtk.118},level=1]{115. วชาทีสุ วสฺสูปคมนํ}
\csromanheader{1}{115. \textbf{วชาทีสุ วสฺสูปคมนํ}}
\proseitem{203}{เตน โข ปน สมเยน อญฺญตโร ภิกฺขุ วเช วสฺสํ อุปคนฺตุกาโม โหติ. ภควโต เอตมตฺถํ อาโรเจสุํ. อนุชานามิ ภิกฺขเว วเช วสฺสํ อุปคนฺตุนฺติ. วโช วุฏฺฐาสิ. ภควโต เอตมตฺถํ อาโรเจสุํ. อนุชานามิ ภิกฺขเว เยน วโช เตน คนฺตุนฺติ.}

\prose{เตน โข ปน สมเยน อญฺญตโร ภิกฺขุ อุปกฏฺฐาย วสฺสู\-ปนายิกาย สตฺเถน คนฺตุกาโม โหติ. ภควโต เอตมตฺถํ อาโรเจสุํ. อนุชานามิ ภิกฺขเว สตฺเถ วสฺสํ อุปคนฺตุนฺติ.}

\prose{เตน โข ปน สมเยน อญฺญตโร ภิกฺขุ อุปกฏฺฐาย วสฺสู\-ปนายิกาย นาวาย คนฺตุกาโม โหติ. ภควโต เอตมตฺถํ อาโรเจสุํ. อนุชานามิ ภิกฺขเว นาวาย วสฺสํ อุปคนฺตุนฺติ.}

\csromanheader{2}{116. \textbf{วสฺสํ อนุปคนฺตพฺพฏฺฐานานิ}}
\csromanmatikaanchor{mtk.119}
\csromantocmarkref{section}{116. วสฺสํอนุปคนฺตพฺพฏฺฐานานิ}{mtk.119}
\bookmark[dest={mtk.119},level=1]{116. วสฺสํอนุปคนฺตพฺพฏฺฐานานิ}
\proseitem{204}{\csromanfolio{212}เตน โข ปน สมเยน ภิกฺขู รุกฺขสุสิเร วสฺสํ อุปคจฺฉนฺติ. มนุสฺสา อุชฺฌายนฺติ ขิยฺยนฺติ วิปาเจนฺติ “เสยฺยถาปิ ปิสาจิลฺลิกา”ติ. ภควโต เอตมตฺถํ อาโรเจสุํ. น ภิกฺขเว รุกฺขสุสิเร วสฺสํ อุปคนฺตพฺพํ, โย อุปคจฺเฉยฺย, อาปตฺติ ทุกฺกฏสฺสาติ.}

\prose{เตน โข ปน สมเยน ภิกฺขู รุกฺข\-วิฏภิยา วสฺสํ อุปคจฺฉนฺติ. มนุสฺสา อุชฺฌายนฺติ ขิยฺยนฺติ วิปาเจนฺติ “เสยฺยถาปิ มิคลุทฺทกา”ติ. ภควโต เอตมตฺถํ อาโรเจสุํ. น ภิกฺขเว รุกฺข\-วิฏภิยา วสฺสํ อุปคนฺตพฺพํ, โย อุปคจฺเฉยฺย, อาปตฺติ ทุกฺกฏสฺสาติ.}

\prose{เตน โข ปน สมเยน ภิกฺขู อชฺโฌกาเส วสฺสํ อุปคจฺฉนฺติ. เทเว วสฺสนฺเต รุกฺขมูลมฺปิ นิพฺพโกสมฺปิ อุปธาวนฺติ. ภควโต เอตมตฺถํ อาโรเจสุํ. น ภิกฺขเว อชฺโฌกาเส วสฺสํ อุปคนฺตพฺพํ, โย อุปคจฺเฉยฺย, อาปตฺติ ทุกฺกฏสฺสาติ.}

\prose{เตน โข ปน สมเยน ภิกฺขู อเสนาสนิกา วสฺสํ อุปคจฺฉนฺติ. สีเตนปิ กิลมนฺติ, อุเณฺหนปิ กิลมนฺติ. ภควโต เอตมตฺถํ อาโรเจสุํ. น ภิกฺขเว อเสนาสนิเกน วสฺสํ อุปคนฺตพฺพํ, โย อุปคจฺเฉยฺย, อาปตฺติ ทุกฺกฏสฺสาติ.}

\prose{เตน โข ปน สมเยน ภิกฺขู ฉวกุฏิกาย วสฺสํ อุปคจฺฉนฺติ. มนุสฺสา อุชฺฌายนฺติ ขิยฺยนฺติ วิปาเจนฺติ “เสยฺยถาปิ ฉวฑาหกา”ติ. ภควโต เอตมตฺถํ อาโรเจสุํ. น ภิกฺขเว ฉวกุฏิกาย วสฺสํ อุปคนฺตพฺพํ, โย อุปคจฺเฉยฺย, อาปตฺติ ทุกฺกฏสฺสาติ.}

\prose{เตน โข ปน สมเยน ภิกฺขู ฉตฺเต วสฺสํ อุปคจฺฉนฺติ. มนุสฺสา อุชฺฌายนฺติ ขิยฺยนฺติ วิปาเจนฺติ “เสยฺยถาปิ โคปาลกา”ติ. ภควโต เอตมตฺถํ อาโรเจสุํ. น ภิกฺขเว ฉตฺเต วสฺสํ อุปคนฺตพฺพํ, โย อุปคจฺเฉยฺย, อาปตฺติ ทุกฺกฏสฺสาติ.}

\prose{เตน โข ปน สมเยน ภิกฺขู จาฏิยา วสฺสํ อุปคจฺฉนฺติ. มนุสฺสา อุชฺฌายนฺติ ขิยฺยนฺติ วิปาเจนฺติ “เสยฺยถาปิ ติตฺถิยา”ติ. ภควโต เอตมตฺถํ อาโรเจสุํ. น ภิกฺขเว จาฏิยา วสฺสํ อุปคนฺตพฺพํ, โย อุปคจฺเฉยฺย, อาปตฺติ ทุกฺกฏสฺสาติ.}

\csromanmatikaanchor{mtk.120}
\csromantocmarkref{section}{117. อธมฺมิกกติกา}{mtk.120}
\bookmark[dest={mtk.120},level=1]{117. อธมฺมิกกติกา}
\csromanheader{1}{117. \textbf{อธมฺมิก}\-\textbf{กติกา}}
\proseitem{205}{\csromanfolio{213}เตน โข ปน สมเยน สาวตฺถิยา สํเฆน เอวรูปา กติกา กตา โหติ “อนฺตราวสฺสํ น ปพฺพาเชตพฺพนฺ”ติ. วิสาขาย มิคารมาตุยา นตฺตา ภิกฺขู อุปสงฺกมิตฺวา ปพฺพชฺชํ ยาจิ. ภิกฺขู เอวมาหํสุ “สํเฆน โข อาวุโส เอวรูปา กติกา กตา ‘อนฺตราวสฺสํ น ปพฺพาเชตพฺพนฺ’ติ. อาคเมหิ อาวุโส ยาว ภิกฺขู วสฺสํ วสนฺติ, วสฺสํวุฏฺฐา ปพฺพาเชสฺสนฺตี”ติ. อถ โข เต ภิกฺขู วสฺสํวุฏฺฐา วิสาขาย มิคารมาตุยา นตฺตารํ เอตทโวจุํ “เอหิ ทานิ อาวุโส ปพฺพาชาหี”ติ. โส เอวมาห “สจาหํ ภนฺเต ปพฺพชิโต อสฺสํ, อภิรเมยฺยามหํ\csromansharedfootnote{5179915982564cb4}{อภิรเมยฺยํ จาหํ (สี)}, น ทานาหํ ภนฺเต ปพฺพชิสฺสามี”ติ. วิสาขา มิคารมาตา อุชฺฌายติ ขิยฺยติ วิปาเจติ “กถํ หิ นาม อยฺยา เอวรูปํ กติกํ กริสฺสนฺติ ‘น อนฺตราวสฺสํ ปพฺพาเชตพฺพนฺ’ติ, กํ กาลํ ธมฺโม น จริตพฺโพ”ติ. อสฺโสสุํ โข ภิกฺขู วิสาขาย มิคารมาตุยา อุชฺฌายนฺติยา ขิยฺยนฺติยา วิปาเจนฺติยา. อถ โข เต ภิกฺขู ภควโต เอตมตฺถํ อาโรเจสุํ. น ภิกฺขเว เอวรูปา กติกา กาตพฺพา “น อนฺตราวสฺสํ ปพฺพาเชตพฺพนฺ”ติ, โย กเรยฺย, อาปตฺติ ทุกฺกฏสฺสาติ.}

\csromanmatikaanchor{mtk.121}
\csromantocmarkref{section}{118. ปฏิสฺสวทุกฺกฏาปตฺติ}{mtk.121}
\bookmark[dest={mtk.121},level=1]{118. ปฏิสฺสวทุกฺกฏาปตฺติ}
\csromanheader{1}{118. \textbf{ปฏิสฺสว}\-\textbf{ทุกฺกฏา}\-\textbf{ปตฺติ}}
\proseitem{206}{เตน โข ปน สมเยน อายสฺมตา อุปนนฺเทน สกฺยปุตฺเตน รญฺโญ ปเสนทิสฺส โกสลสฺส วสฺสาวาโส ปฏิสฺสุโต โหติ ปุริมิกาย. โส ตํ อาวาสํ คจฺฉนฺโต อทฺทส อนฺตรามคฺเค เทฺว อาวาเส พหุจีวรเก, ตสฺส เอตทโหสิ “ยนฺนูนาหํ อิเมสุ ทฺวีสุ อาวาเสสุ วสฺสํ วเสยฺยํ, เอวํ เม พหุํ จีวรํ\csromansharedfootnote{327d57a3f72c9e1b}{พหุจีวรํ (ก)} อุปฺปชฺชิสฺสตี”ติ. โส เตสุ ทฺวีสุ อาวาเสสุ วสฺสํ วสิ. ราชา ปเสนทิ โกสโล อุชฺฌายติ ขิยฺยติ วิปาเจติ “กถํ หิ นาม อยฺโย อุปนนฺโท สกฺยปุตฺโต อมฺหากํ วสฺสาวาสํ ปฏิสฺสุณิตฺวา วิสํวาเทสฺสติ, นนุ ภควตา อเนก\-ปริยาเยน มุสาวาโท ครหิโต, มุสาวาทา เวรมณี ปสตฺถา”ติ. อสฺโสสุํ โข ภิกฺขู รญฺโญ ปเสนทิสฺส โกสลสฺส อุชฺฌายนฺตสฺส ขิยฺยนฺตสฺส วิปาเจนฺตสฺส. เย เต ภิกฺขู อปฺปิจฺฉา ฯเปฯ เต อุชฺฌายนฺติ ขิยฺยนฺติ วิปาเจนฺติ “กถํ หิ นาม อายสฺมา อุปนนฺโท \csromanfolio{214}สกฺยปุตฺโต รญฺโญ ปเสนทิสฺส โกสลสฺส วสฺสาวาสํ ปฏิสฺสุณิตฺวา วิสํวาเทสฺสติ, นนุ ภควตา อเนก\-ปริยาเยน มุสาวาโท ครหิโต, มุสาวาทา เวรมณี ปสตฺถา”ติ. อถ โข เต ภิกฺขู ภควโต เอตมตฺถํ อาโรเจสุํ ฯเปฯ. อถ โข ภควา เอตสฺมิํ นิทาเน เอตสฺมิํ ปกรเณ ภิกฺขุสํฆํ สนฺนิปาตาเปตฺวา อายสฺมนฺตํ อุปนนฺทํ สกฺยปุตฺตํ ปฏิปุจฺฉิ “สจฺจํ กิร ตฺวํ อุปนนฺท รญฺโญ ปเสนทิสฺส โกสลสฺส วสฺสาวาสํ ปฏิสฺสุณิตฺวา วิสํวาเทสี”ติ. สจฺจํ ภควาติ. วิครหิ พุทฺโธ ภควา ฯเปฯ. กถํ หิ นาม ตฺวํ โมฆปุริส รญฺโญ ปเสนทิสฺส โกสลสฺส วสฺสาวาสํ ปฏิสฺสุณิตฺวา วิสํวาเทสฺสสิ. นนุ มยา โมฆปุริส อเนก\-ปริยาเยน มุสาวาโท ครหิโต, มุสาวาทา เวรมณี ปสตฺถา. เนตํ โมฆปุริส อปฺปสนฺนานํ วา ปสาทาย ฯเปฯ วิครหิตฺวา ฯเปฯ ธมฺมิํ กถํ กตฺวา ภิกฺขู อามนฺเตสิ\csromandash{}}

\proseitem{207}{อิธ ปน ภิกฺขเว ภิกฺขุนา วสฺสาวาโส ปฏิสฺสุโต โหติ ปุริมิกาย, โส ตํ อาวาสํ คจฺฉนฺโต ปสฺสติ อนฺตรามคฺเค เทฺว อาวาเส พหุจีวรเก, ตสฺส เอวํ โหติ “ยนฺนูนาหํ อิเมสุ ทฺวีสุ อาวาเสสุ วสฺสํ วเสยฺยํ, เอวํ เม พหุํ จีวรํ อุปฺปชฺชิสฺสตี”ติ, โส เตสุ ทฺวีสุ อาวาเสสุ วสฺสํ วสติ, ตสฺส ภิกฺขเว ภิกฺขุโน ปุริมิกา จ น ปญฺญายติ, ปฏิสฺสเว จ อาปตฺติ ทุกฺกฏสฺส.}

\prose{อิธ ปน ภิกฺขเว ภิกฺขุนา วสฺสาวาโส ปฏิสฺสุโต โหติ ปุริมิกาย, โส ตํ อาวาสํ คจฺฉนฺโต พหิทฺธา อุโปสถํ กโรติ, ปาฏิปเท\csromansharedfootnote{30e5ea1971107f66}{ปาฏิปเทน (ก)} วิหารํ อุเปติ, เสนาสนํ ปญฺญเปติ, ปานียํ ปริโภชนียํ อุปฏฺฐาเปติ, ปริเวณํ สมฺมชฺชติ, โส ตทเหว อกรณีโย ปกฺกมติ, ตสฺส ภิกฺขเว ภิกฺขุโน ปุริมิกา จ น ปญฺญายติ, ปฏิสฺสเว จ อาปตฺติ ทุกฺกฏสฺส.}

\proseitem{207}{อิธ ปน ภิกฺขเว ภิกฺขุนา วสฺสาวาโส ปฏิสฺสุโต โหติ ปุริมิกาย, โส ตํ อาวาสํ คจฺฉนฺโต พหิทฺธา อุโปสถํ กโรติ, ปาฏิปเท วิหารํ อุเปติ, เสนาสนํ ปญฺญเปติ, ปานียํ ปริโภชนียํ อุปฏฺฐาเปติ, ปริเวณํ สมฺมชฺชติ, โส ตทเหว สกรณีโย ปกฺกมติ, \csromanfolio{215}ตสฺส ภิกฺขเว ภิกฺขุโน ปุริมิกา จ น ปญฺญายติ, ปฏิสฺสเว จ อาปตฺติ ทุกฺกฏสฺส.}

\prose{อิธ ปน ภิกฺขเว ภิกฺขุนา วสฺสาวาโส ปฏิสฺสุโต โหติ ปุริมิกาย, โส ตํ อาวาสํ คจฺฉนฺโต พหิทฺธา อุโปสถํ กโรติ, ปาฏิปเท วิหารํ อุเปติ, เสนาสนํ ปญฺญเปติ, ปานียํ ปริโภชนียํ อุปฏฺฐาเปติ, ปริเวณํ สมฺมชฺชติ, โส ทฺวีหตีหํ วสิตฺวา อกรณีโย ปกฺกมติ, ตสฺส ภิกฺขเว ภิกฺขุโน ปุริมิกา จ น ปญฺญายติ, ปฏิสฺสเว จ อาปตฺติ ทุกฺกฏสฺส.}

\prose{อิธ ปน ภิกฺขเว ภิกฺขุนา วสฺสาวาโส ปฏิสฺสุโต โหติ ปุริมิกาย, โส ตํ อาวาสํ คจฺฉนฺโต พหิทฺธา อุโปสถํ กโรติ, ปาฏิปเท วิหารํ อุเปติ, เสนาสนํ ปญฺญเปติ, ปานียํ ปริโภชนียํ อุปฏฺฐาเปติ, ปริเวณํ สมฺมชฺชติ, โส ทฺวีหตีหํ วสิตฺวา สกรณีโย ปกฺกมติ, ตสฺส ภิกฺขเว ภิกฺขุโน ปุริมิกา จ น ปญฺญายติ, ปฏิสฺสเว จ อาปตฺติ ทุกฺกฏสฺส.}

\prose{อิธ ปน ภิกฺขเว ภิกฺขุนา วสฺสาวาโส ปฏิสฺสุโต โหติ ปุริมิกาย, โส ตํ อาวาสํ คจฺฉนฺโต พหิทฺธา อุโปสถํ กโรติ, ปาฏิปเท วิหารํ อุเปติ, เสนาสนํ ปญฺญเปติ, ปานียํ ปริโภชนียํ อุปฏฺฐาเปติ, ปริเวณํ สมฺมชฺชติ, โส ทฺวีหตีหํ วสิตฺวา สตฺตาห\-กรณีเยน ปกฺกมติ, โส ตํ สตฺตาหํ พหิทฺธา วีตินาเมติ, ตสฺส ภิกฺขเว ภิกฺขุโน ปุริมิกา จ น ปญฺญายติ, ปฏิสฺสเว จ อาปตฺติ ทุกฺกฏสฺส.}

\prose{อิธ ปน ภิกฺขเว ภิกฺขุนา วสฺสาวาโส ปฏิสฺสุโต โหติ ปุริมิกาย, โส ตํ อาวาสํ คจฺฉนฺโต พหิทฺธา อุโปสถํ กโรติ, ปาฏิปเท วิหารํ อุเปติ, เสนาสนํ ปญฺญเปติ, ปานียํ ปริโภชนียํ อุปฏฺฐาเปติ, ปริเวณํ สมฺมชฺชติ, โส ทฺวีหตีหํ วสิตฺวา สตฺตาห\-กรณีเยน ปกฺกมติ, โส ตํ สตฺตาหํ อนฺโต สนฺนิวตฺตํ กโรติ, ตสฺส ภิกฺขเว ภิกฺขุโน ปุริมิกา จ ปญฺญายติ, ปฏิสฺสเว จ อนาปตฺติ.}

\proseitem{207}{อิธ ปน ภิกฺขเว ภิกฺขุนา วสฺสาวาโส ปฏิสฺสุโต โหติ ปุริมิกาย, โส ตํ อาวาสํ คจฺฉนฺโต พหิทฺธา อุโปสถํ กโรติ, ปาฏิปเท วิหารํ อุเปติ, เสนาสนํ ปญฺญเปติ, ปานียํ ปริโภชนียํ \csromanfolio{216}อุปฏฺฐาเปติ, ปริเวณํ สมฺมชฺชติ, โส สตฺตาหํ อนาคตาย ปวารณาย สกรณีโย ปกฺกมติ, อาคจฺเฉยฺย วา โส ภิกฺขเว ภิกฺขุ ตํ อาวาสํ น วา อาคจฺเฉยฺย, ตสฺส ภิกฺขเว ภิกฺขุโน ปุริมิกา จ ปญฺญายติ, ปฏิสฺสเว จ อนาปตฺติ.}

\prose{อิธ ปน ภิกฺขเว ภิกฺขุนา วสฺสาวาโส ปฏิสฺสุโต โหติ ปุริมิกาย, โส ตํ อาวาสํ คนฺตฺวา อุโปสถํ กโรติ, ปาฏิปเท วิหารํ อุเปติ, เสนาสนํ ปญฺญเปติ, ปานียํ ปริโภชนียํ อุปฏฺฐาเปติ, ปริเวณํ สมฺมชฺชติ, โส ตทเหว อกรณีโย ปกฺกมติ, ตสฺส ภิกฺขเว ภิกฺขุโน ปุริมิกา จ น ปญฺญายติ, ปฏิสฺสเว จ อาปตฺติ ทุกฺกฏสฺส.}

\prose{อิธ ปน ภิกฺขเว ภิกฺขุนา วสฺสาวาโส ปฏิสฺสุโต โหติ ปุริมิกาย, โส ตํ อาวาสํ คนฺตฺวา อุโปสถํ กโรติ, ปาฏิปเท วิหารํ อุเปติ, เสนาสนํ ปญฺญเปติ, ปานียํ ปริโภชนียํ อุปฏฺฐาเปติ, ปริเวณํ สมฺมชฺชติ, โส ตทเหว สกรณีโย ปกฺกมติ ฯเปฯ โส ทฺวีหตีหํ วสิตฺวา อกรณีโย ปกฺกมติ ฯเปฯ โส ทฺวีหตีหํ วสิตฺวา สกรณีโย ปกฺกมติ ฯเปฯ โส ทฺวีหตีหํ วสิตฺวา สตฺตาห\-กรณีเยน ปกฺกมติ, โส ตํ สตฺตาหํ พหิทฺธา วีตินาเมติ, ตสฺส ภิกฺขเว ภิกฺขุโน ปุริมิกา จ น ปญฺญายติ, ปฏิสฺสเว จ อาปตฺติ ทุกฺกฏสฺส ฯเปฯ โส ทฺวีหตีหํ วสิตฺวา สตฺตาห\-กรณีเยน ปกฺกมติ, โส ตํ สตฺตาหํ อนฺโต สนฺนิวตฺตํ กโรติ, ตสฺส ภิกฺขเว ภิกฺขุโน ปุริมิกา จ ปญฺญายติ, ปฏิสฺสเว จ อนาปตฺติ ฯเปฯ โส สตฺตาหํ อนาคตาย ปวารณาย สกรณีโย ปกฺกมติ, อาคจฺเฉยฺย วา โส ภิกฺขเว ภิกฺขุ ตํ อาวาสํ น วา อาคจฺเฉยฺย, ตสฺส ภิกฺขเว ภิกฺขุโน ปุริมิกา จ ปญฺญายติ, ปฏิสฺสเว จ อนาปตฺติ.}

\proseitem{208}{อิธ ปน ภิกฺขเว ภิกฺขุนา วสฺสาวาโส ปฏิสฺสุโต โหติ ปจฺฉิมิกาย, โส ตํ อาวาสํ คจฺฉนฺโต พหิทฺธา อุโปสถํ กโรติ, ปาฏิปเท วิหารํ อุเปติ, เสนาสนํ ปญฺญเปติ, ปานียํ ปริโภชนียํ อุปฏฺฐาเปติ, ปริเวณํ สมฺมชฺชติ, โส ตทเหว อกรณีโย ปกฺกมติ, ตสฺส ภิกฺขเว ภิกฺขุโน ปจฺฉิมิกา จ น ปญฺญายติ, ปฏิสฺสเว จ อาปตฺติ ทุกฺกฏสฺส.}

\prose{\csromanfolio{217}อิธ ปน ภิกฺขเว ภิกฺขุนา วสฺสาวาโส ปฏิสฺสุโต โหติ ปจฺฉิมิกาย, โส ตํ อาวาสํ คจฺฉนฺโต พหิทฺธา อุโปสถํ กโรติ, ปาฏิปเท วิหารํ อุเปติ, เสนาสนํ ปญฺญเปติ, ปานียํ ปริโภชนียํ อุปฏฺฐาเปติ, ปริเวณํ สมฺมชฺชติ, โส ตทเหว สกรณีโย ปกฺกมติ, ตสฺส ภิกฺขเว ภิกฺขุโน ปจฺฉิมิกา จ น ปญฺญายติ, ปฏิสฺสเว จ อาปตฺติ ทุกฺกฏสฺส.}

\prose{อิธ ปน ภิกฺขเว ภิกฺขุนา วสฺสาวาโส ปฏิสฺสุโต โหติ ปจฺฉิมิกาย, โส ตํ อาวาสํ คจฺฉนฺโต พหิทฺธา อุโปสถํ กโรติ, ปาฏิปเท วิหารํ อุเปติ, เสนาสนํ ปญฺญเปติ, ปานียํ ปริโภชนียํ อุปฏฺฐาเปติ, ปริเวณํ สมฺมชฺชติ, โส ทฺวีหตีหํ วสิตฺวา อกรณีโย ปกฺกมติ, ตสฺส ภิกฺขเว ภิกฺขุโน ปจฺฉิมิกา จ น ปญฺญายติ, ปฏิสฺสเว จ อาปตฺติ ทุกฺกฏสฺส.}

\prose{อิธ ปน ภิกฺขเว ภิกฺขุนา วสฺสาวาโส ปฏิสฺสุโต โหติ ปจฺฉิมิกาย, โส ตํ อาวาสํ คจฺฉนฺโต พหิทฺธา อุโปสถํ กโรติ, ปาฏิปเท วิหารํ อุเปติ, เสนาสนํ ปญฺญเปติ, ปานียํ ปริโภชนียํ อุปฏฺฐาเปติ, ปริเวณํ สมฺมชฺชติ, โส ทฺวีหตีหํ วสิตฺวา สกรณีโย ปกฺกมติ, ตสฺส ภิกฺขเว ภิกฺขุโน ปจฺฉิมิกา จ น ปญฺญายติ, ปฏิสฺสเว จ อาปตฺติ ทุกฺกฏสฺส.}

\prose{อิธ ปน ภิกฺขเว ภิกฺขุนา วสฺสาวาโส ปฏิสฺสุโต โหติ ปจฺฉิมิกาย, โส ตํ อาวาสํ คจฺฉนฺโต พหิทฺธา อุโปสถํ กโรติ, ปาฏิปเท วิหารํ อุเปติ, เสนาสนํ ปญฺญเปติ, ปานียํ ปริโภชนียํ อุปฏฺฐาเปติ, ปริเวณํ สมฺมชฺชติ, โส ทฺวีหตีหํ วสิตฺวา สตฺตาห\-กรณีเยน ปกฺกมติ, โส ตํ สตฺตาหํ พหิทฺธา วีตินาเมติ, ตสฺส ภิกฺขเว ภิกฺขุโน ปจฺฉิมิกา จ น ปญฺญายติ, ปฏิสฺสเว จ อาปตฺติ ทุกฺกฏสฺส.}

\prose{อิธ ปน ภิกฺขเว ภิกฺขุนา วสฺสาวาโส ปฏิสฺสุโต โหติ ปจฺฉิมิกาย, โส ตํ อาวาสํ คจฺฉนฺโต พหิทฺธา อุโปสถํ กโรติ, ปาฏิปเท วิหารํ อุเปติ, เสนาสนํ ปญฺญเปติ, ปานียํ ปริโภชนียํ อุปฏฺฐาเปติ, ปริเวณํ สมฺมชฺชติ, โส ทฺวีหตีหํ วสิตฺวา สตฺตาห\-กรณีเยน ปกฺกมติ, โส ตํ สตฺตาหํ อนฺโต สนฺนิวตฺตํ กโรติ, ตสฺส ภิกฺขเว ภิกฺขุโน ปจฺฉิมิกา จ ปญฺญยติ, ปฏิสฺสเว จ อนาปตฺติ.}

\prose{\csromanfolio{218}อิธ ปน ภิกฺขเว ภิกฺขุนา วสฺสาวาโส ปฏิสฺสุโต โหติ ปจฺฉิมิกาย, โส ตํ อาวาสํ คจฺฉนฺโต พหิทฺธา อุโปสถํ กโรติ, ปาฏิปเท วิหารํ อุเปติ, เสนาสนํ ปญฺญเปติ, ปานียํ ปริโภชนียํ อุปฏฺฐาเปติ, ปริเวณํ สมฺมชฺชติ, โส สตฺตาหํ อนาคตาย โกมุทิยา จาตุมาสินิยา สกรณีโย ปกฺกมติ, อาคจฺเฉยฺย วา โส ภิกฺขเว ภิกฺขุ ตํ อาวาสํ น วา อาคจฺเฉยฺย, ตสฺส ภิกฺขเว ภิกฺขุโน ปจฺฉิมิกา จ ปญฺญายติ, ปฏิสฺสเว จ อนาปตฺติ.}

\prose{อิธ ปน ภิกฺขเว ภิกฺขุนา วสฺสาวาโส ปฏิสฺสุโต โหติ ปจฺฉิมิกาย, โส ตํ อาวาสํ คนฺตฺวา อุโปสถํ กโรติ, ปาฏิปเท วิหารํ อุเปติ, เสนาสนํ ปญฺญเปติ, ปานียํ ปริโภชนียํ อุปฏฺฐาเปติ, ปริเวณํ สมฺมชฺชติ, โส ตทเหว อกรณีโย ปกฺกมติ, ตสฺส ภิกฺขเว ภิกฺขุโน ปจฺฉิมิกา จ น ปญฺญายติ, ปฏิสฺสเว จ อาปตฺติ ทุกฺกฏสฺส.}

\prose{อิธ ปน ภิกฺขเว ภิกฺขุนา วสฺสาวาโส ปฏิสฺสุโต โหติ ปจฺฉิมิกาย, โส ตํ อาวาสํ คนฺตฺวา อุโปสถํ กโรติ, ปาฏิปเท วิหารํ อุเปติ, เสนาสนํ ปญฺญเปติ, ปานียํ ปริโภชนียํ อุปฏฺฐาเปติ, ปริเวณํ สมฺมชฺชติ, โส ตทเหว สกรณีโย ปกฺกมติ ฯเปฯ โส ทฺวีหตีหํ วสิตฺวา อกรณีโย ปกฺกมติ ฯเปฯ โส ทฺวีหตีหํ วสิตฺวา สกรณีโย ปกฺกมติ ฯเปฯ โส ทฺวีหตีหํ วสิตฺวา สตฺตาห\-กรณีเยน ปกฺกมติ, โส ตํ สตฺตาหํ พหิทฺธา วีตินาเมติ, ตสฺส ภิกฺขเว ภิกฺขุโน ปจฺฉิมิกา จ น ปญฺญายติ, ปฏิสฺสเว จ อาปตฺติ ทุกฺกฏสฺส ฯเปฯ โส ทฺวีหตีหํ วสิตฺวา สตฺตาห\-กรณีเยน ปกฺกมติ, โส ตํ สตฺตาหํ อนฺโต สนฺนิวตฺตํ กโรติ, ตสฺส ภิกฺขเว ภิกฺขุโน ปจฺฉิมิกา จ ปญฺญายติ, ปฏิสฺสเว จ อนาปตฺติ.}

\prose{อิธ ปน ภิกฺขเว ภิกฺขุนา วสฺสาวาโส ปฏิสฺสุโต โหติ ปจฺฉิมิกาย, โส ตํ อาวาสํ คนฺตฺวา อุโปสถํ กโรติ, ปาฏิปเท วิหารํ อุเปติ, เสนาสนํ ปญฺญเปติ, ปานียํ ปริโภชนียํ อุปฏฺฐาเปติ, ปริเวณํ สมฺมชฺชติ, โส สตฺตาหํ อนาคตาย โกมุทิยา จาตุมาสินิยา สกรณีโย ปกฺกมติ, อาคจฺเฉยฺย วา โส ภิกฺขเว ภิกฺขุ ตํ อาวาสํ น วา อาคจฺเฉยฺย, ตสฺส ภิกฺขเว ภิกฺขุโน ปจฺฉิมิกา จ ปญฺญายติ, ปฏิสฺสเว จ อนาปตฺตีติ.}

\vspace{\tipitakaparskip}
\csromancenter{วสฺสู\-ปนายิก\-กฺขนฺธโก ตติโย.}
\csromancenter{119. \textbf{ตสฺสุทฺทานํ}}
\csromanmatikaanchor{mtk.122}
\csromantocmarkref{section}{119. อุทฺทานคาถา}{mtk.122}
\bookmark[dest={mtk.122},level=1]{119. อุทฺทานคาถา}
\setlength{\csromangathapagewidth}{0pt}
\setlength{\csromangathawakwidth}{0pt}
\csromangathameasuregroup{อุปคนฺตุํ กทา เจว, \\ น อิจฺฉนฺติ จ สญฺจิจฺจ, \\ คิลาโน มาตา จ ปิตา, \\ ภิกฺขุคติโก วิหาโร, \\ โจโร เจว ปิสาจา จ, \\ วูโฬฺหทเกน วุฏฺฐาสิ, \\ ลูขปฺปณีตสปฺปาย, \\ อิตฺถี เวสี กุมารี จ, \\ ราชา โจรา ธุตฺตา นิธิ, \\ วชสตฺถา จ นาวา จ, \\ อชฺโฌกาเส วสฺสาวาโส, \\ ฉวกุฏิกา ฉตฺเต จ, \\ กติกา ปฏิสฺสุณิตฺวา, \\ ปุริมิกา ปจฺฉิมิกา, \\ อกรณี ปกฺกมติ, \\ ทฺวีหตีหา จ ปุน จ\textsuperscript{1}, \\ สตฺตาหนาคตา เจว, \\ วตฺถุทฺทาเน อนฺตริกา,}{\csromangathabat{อุปคนฺตุํ กทา เจว,}{กติ อนฺตราวสฺส จ.} \\ \csromangathabat{น อิจฺฉนฺติ จ สญฺจิจฺจ,}{อุกฺกฑฺฒิตุํ อุปาสโก.} \\ \csromangathabat{คิลาโน มาตา จ ปิตา,}{ภาตา จ อถ ญาตโก.} \\ \csromangathabat{ภิกฺขุคติโก วิหาโร,}{วาฬา จาปิ สรีสปา.} \\ \csromangathabat{โจโร เจว ปิสาจา จ,}{ทฑฺฒา ตทุภเยน จ.} \\ \csromangathabat{วูโฬฺหทเกน วุฏฺฐาสิ,}{พหุตรา จ ทายกา.} \\ \csromangathabat{ลูขปฺปณีตสปฺปาย,}{เภสชฺชุปฏฺฐเกน จ.} \\ \csromangathabat{อิตฺถี เวสี กุมารี จ,}{ปณฺฑโก ญาตเกน จ.} \\ \csromangathabat{ราชา โจรา ธุตฺตา นิธิ,}{เภทอฏฺฐวิเธน\textsuperscript{1} จ.} \\ \csromangathabat{วชสตฺถา จ นาวา จ,}{สุสิเร วิฏภิยา จ.} \\ \csromangathabat{อชฺโฌกาเส วสฺสาวาโส,}{อเสนาสนิเกน จ.} \\ \csromangathabat{ฉวกุฏิกา ฉตฺเต จ,}{จาฏิยา จ อุเปนฺติ เต.} \\ \csromangathabat{กติกา ปฏิสฺสุณิตฺวา,}{พหิทฺธา จ อุโปสถา.} \\ \csromangathabat{ปุริมิกา ปจฺฉิมิกา,}{ยถาญาเยน โยชเย.} \\ \csromangathabat{อกรณี ปกฺกมติ,}{สกรณี ตเถว จ.} \\ \csromangathabat{ทฺวีหตีหา จ ปุน จ\textsuperscript{1},}{สตฺตาหกรณีเยน จ.} \\ \csromangathabat{สตฺตาหนาคตา เจว,}{อาคจฺเฉยฺย น เอยฺย วา.} \\ \csromangathabat{วตฺถุทฺทาเน อนฺตริกา,}{ตนฺติมคฺคํ นิสามเยติ.} \\ อิมมฺหิ ขนฺธเก วตฺถูนิ เทฺวปณฺณาส.}
\csromangathameasurewak{อิมมฺหิ ขนฺธเก วตฺถูนิ เทฺวปณฺณาส.}
\csromangathasetleft{อุปคนฺตุํ กทา เจว, \\ น อิจฺฉนฺติ จ สญฺจิจฺจ, \\ คิลาโน มาตา จ ปิตา, \\ ภิกฺขุคติโก วิหาโร, \\ โจโร เจว ปิสาจา จ, \\ วูโฬฺหทเกน วุฏฺฐาสิ, \\ ลูขปฺปณีตสปฺปาย, \\ อิตฺถี เวสี กุมารี จ, \\ ราชา โจรา ธุตฺตา นิธิ, \\ วชสตฺถา จ นาวา จ, \\ อชฺโฌกาเส วสฺสาวาโส, \\ ฉวกุฏิกา ฉตฺเต จ, \\ กติกา ปฏิสฺสุณิตฺวา, \\ ปุริมิกา ปจฺฉิมิกา, \\ อกรณี ปกฺกมติ, \\ ทฺวีหตีหา จ ปุน จ\textsuperscript{1}, \\ สตฺตาหนาคตา เจว, \\ วตฺถุทฺทาเน อนฺตริกา,}
\csromangathagroupwithclosermajor{\csromangathastanza{\csromanfolio{219}\csromangathabat{อุปคนฺตุํ กทา เจว,}{กติ อนฺตราวสฺส จ.} \\ \csromangathabat{น อิจฺฉนฺติ จ สญฺจิจฺจ,}{อุกฺกฑฺฒิตุํ อุปาสโก.}}\csromangathastanza{\csromangathabat{คิลาโน มาตา จ ปิตา,}{ภาตา จ อถ ญาตโก.} \\ \csromangathabat{ภิกฺขุคติโก วิหาโร,}{วาฬา จาปิ สรีสปา.}}\csromangathastanza{\csromangathabat{โจโร เจว ปิสาจา จ,}{ทฑฺฒา ตทุภเยน จ.} \\ \csromangathabat{วูโฬฺหทเกน วุฏฺฐาสิ,}{พหุตรา จ ทายกา.}}\csromangathastanza{\csromangathabat{ลูขปฺปณีตสปฺปาย,}{เภสชฺชุปฏฺฐเกน จ.} \\ \csromangathabat{อิตฺถี เวสี กุมารี จ,}{ปณฺฑโก ญาตเกน จ.}}\csromangathastanza{\csromangathabat{ราชา โจรา ธุตฺตา นิธิ,}{เภทอฏฺฐวิเธน\csromansharedfootnote{e7584b7ade187572}{เภทา อฏฺฐวิเธน (สี, สฺยา)} จ.} \\ \csromangathabat{วชสตฺถา จ นาวา จ,}{สุสิเร วิฏภิยา จ.}}\csromangathastanza{\csromangathabat{อชฺโฌกาเส วสฺสาวาโส,}{อเสนาสนิเกน จ.} \\ \csromangathabat{ฉวกุฏิกา ฉตฺเต จ,}{จาฏิยา จ อุเปนฺติ เต.}}\csromangathastanza{\csromangathabat{กติกา ปฏิสฺสุณิตฺวา,}{พหิทฺธา จ อุโปสถา.} \\ \csromangathabat{ปุริมิกา ปจฺฉิมิกา,}{ยถาญาเยน โยชเย.}}\csromangathastanza{\csromangathabat{อกรณี ปกฺกมติ,}{สกรณี ตเถว จ.} \\ \csromangathabat{ทฺวีหตีหา จ ปุน จ\csromansharedfootnote{16fdf2fde37367bc}{ทฺวีหตีหํ วสิตฺวาน (สี)},}{สตฺตาห\-กรณีเยน จ.}}\csromangathastanza{\csromangathabat{สตฺตาหนาคตา เจว,}{อาคจฺเฉยฺย น เอยฺย วา.} \\ \csromangathabat{วตฺถุทฺทาเน อนฺตริกา,}{ตนฺติมคฺคํ นิสามเยติ.}}\csromangathastanza{\csromangathawak{อิมมฺหิ ขนฺธเก วตฺถูนิ เทฺวปณฺณาส.}}}{\textbf{วสฺสู}\-\textbf{ปนายิก}\-\textbf{กฺขนฺธโก นิฏฺฐิโต.}}
\csromanmatikaanchor{mtk.123}
\csromantocmarknopage{chapter}{4. ปวารณากฺขนฺธก}{mtk.123}
\bookmark[dest={mtk.123},level=0]{4. ปวารณากฺขนฺธก}
\chapterheadpageread{4. \textbf{ปวารณา}\-\textbf{กฺขนฺธก}}
\csromanmatikaanchor{mtk.124}
\csromantocmarkref{section}{120. อผาสุกวิหาร}{mtk.124}
\bookmark[dest={mtk.124},level=1]{120. อผาสุกวิหาร}
\csromanheader{1}{120. \textbf{อผาสุกวิหาร}}
\proseitem{209}{\csromanfolio{220}เตน สมเยน พุทฺโธ ภควา สาวตฺถิยํ วิหรติ เชตวเน อนาถ\-ปิณฺฑิกสฺส อาราเม. เตน โข ปน สมเยน สมฺพหุลา สนฺทิฏฺฐา สมฺภตฺตา ภิกฺขู โกสเลสุ ชนปเท อญฺญตรสฺมิํ อาวาเส วสฺสํ อุปคจฺฉิํสุ. อถ โข เตสํ ภิกฺขูนํ เอตทโหสิ “เกน นุ โข มยํ อุปาเยน สมคฺคา สมฺโมทมานา อวิวทมานา ผาสุกํ วสฺสํ วเสยฺยาม, น จ ปิณฺฑเกน กิลเมยฺยามา”ติ. อถ โข เตสํ ภิกฺขูนํ เอตทโหสิ “สเจ โข มยํ อญฺญมญฺญํ เนว อาลเปยฺยาม น สลฺลเปยฺยาม, โย ปฐมํ คามโต ปิณฺฑาย ปฏิกฺกเมยฺย, โส อาสนํ ปญฺญเปยฺย, ปาโททกํ ปาทปีฐํ ปาทกถลิกํ อุปนิกฺขิเปยฺย, อวกฺการปาติํ โธวิตฺวา อุปฏฺฐาเปยฺย, ปานียํ ปริโภชนียํ อุปฏฺฐาเปยฺย, โย ปจฺฉา คามโต ปิณฺฑาย ปฏิกฺกเมยฺย, สจสฺส ภุตฺตาวเสโส, สเจ อากงฺเขยฺย ภุญฺเชยฺย, โน เจ อากงฺเขยฺย อปฺปหริเต วา ฉฑฺเฑยฺย, อปฺปาณเก วา อุทเก โอปิลาเปยฺย, โส อาสนํ อุทฺธเรยฺย, ปาโททกํ ปาทปีฐํ ปาทกถลิกํ ปฏิสาเมยฺย, อวกฺการปาติํ โธวิตฺวา ปฏิสาเมยฺย, ปานียํ ปริโภชนียํ ปฏิสาเมยฺย, ภตฺตคฺคํ สมฺมชฺเชยฺย, โย ปสฺเสยฺย ปานียฆฏํ วา ปริโภชนียฆฏํ วา วจฺจฆฏํ วา ริตฺตํ ตุจฺฉํ, โส อุปฏฺฐาเปยฺย, สจสฺส โหติ อวิสยฺหํ, หตฺถวิกาเรน ทุติยํ อามนฺเตตฺวา หตฺถ\-วิลงฺฆเกน อุปฏฺฐาเปยฺย, น เตฺวว ตปฺปจฺจยา วาจํ ภินฺเทยฺย, เอวํ โข มยํ สมคฺคา สมฺโมทมานา อวิวทมานา ผาสุกํ วสฺสํ วเสยฺยาม, น จ ปิณฺฑเกน กิลเมยฺยามา”ติ. อถ โข เต ภิกฺขู อญฺญมญฺญํ เนว อาลปิํสุ, น สลฺลปิํสุ. โย ปฐมํ คามโต ปิณฺฑาย ปฏิกฺกมติ, โส อาสนํ ปญฺญเปติ, ปาโททกํ ปาทปีฐํ ปาทกถลิกํ อุปนิกฺขิปติ, อวกฺการปาติํ โธวิตฺวา อุปฏฺฐาเปติ, ปานียํ ปริโภชนียํ อุปฏฺฐาเปติ, โย ปจฺฉา คามโต ปิณฺฑาย ปฏิกฺกมติ, สเจ โหติ ภุตฺตาวเสโส, สเจ อากงฺขติ ภุญฺชติ, โน เจ อากงฺขติ อปฺปหริเต วา ฉฑฺเฑติ, อปฺปาณเก วา อุทเก โอปิลาเปติ, โส อาสนํ อุทฺธรติ, ปาโททกํ ปาทปีฐํ ปาทกถลิกํ ปฏิสาเมติ, อวกฺการปาติํ โธวิตฺวา ปฏิสาเมติ, ปานียํ ปริโภชนียํ ปฏิสาเมติ, \csromanfolio{221}ภตฺตคฺคํ สมฺมชฺชติ, โย ปสฺสติ ปานียฆฏํ วา ปริโภชนียฆฏํ วา วจฺจฆฏํ วา ริตฺตํ ตุจฺฉํ, โส อุปฏฺฐาเปติ, สจสฺส โหติ อวิสยฺหํ, หตฺถวิกาเรน ทุติยํ อามนฺเตตฺวา หตฺถ\-วิลงฺฆเกน อุปฏฺฐาเปติ, น เตฺวว ตปฺปจฺจยา วาจํ ภินฺทติ.}

\prose{อาจิณฺณํ โข ปเนตํ วสฺสํ\-วุฏฺฐานํ ภิกฺขูนํ ภควนฺตํ ทสฺสนาย อุปสงฺกมิตุํ, อถ โข เต ภิกฺขู วสฺสํวุฏฺฐา เตมาสจฺจเยน เสนาสนํ สํสาเมตฺวา ปตฺต\-จีวรมา\-ทาย เยน สาวตฺถิ เตน ปกฺกมิํสุ. อนุปุพฺเพน เยน สาวตฺถิ เชตวนํ อนาถ\-ปิณฺฑิกสฺส อาราโม เยน ภควา เตนุ\-ปสงฺกมิํสุ, อุปสงฺกมิตฺวา ภควนฺตํ อภิวาเทตฺวา เอกมนฺตํ นิสีทิํสุ. อาจิณฺณํ โข ปเนตํ พุทฺธานํ ภควนฺตานํ อาคนฺตุเกหิ ภิกฺขูหิ สทฺธิํ ปฏิสมฺโมทิตุํ. อถ โข ภควา เต ภิกฺขู เอตทโวจ “กจฺจิ ภิกฺขเว ขมนียํ, กจฺจิ ยาปนียํ, กจฺจิ สมคฺคา สมฺโมทมานา อวิวทมานา ผาสุกํ วสฺสํ วสิตฺถ, น จ ปิณฺฑเกน กิลมิตฺถา”ติ. ขมนียํ ภควา, ยาปนียํ ภควา, สมคฺคา จ มยํ ภนฺเต สมฺโมทมานา อวิวทมานา ผาสุกํ วสฺสํ วสิมฺหา, น จ ปิณฺฑเกน กิลมิมฺหาติ. ชานนฺตาปิ ตถาคตา ปุจฺฉนฺติ, ชานนฺตาปิ น ปุจฺฉนฺติ. กาลํ วิทิตฺวา ปุจฺฉนฺติ, กาลํ วิทิตฺวา น ปุจฺฉนฺติ. อตฺถสํหิตํ ตถาคตา ปุจฺฉนฺติ, โน อนตฺถ\-สํหิตํ, อนตฺถสํหิเต เสตุฆาโต ตถาคตานํ. ทฺวีหากาเรหิ พุทฺธา ภควนฺโต ภิกฺขู ปฏิปุจฺฉนฺติ “ธมฺมํ วา เทเสสฺสาม, สาวกานํ วา สิกฺขาปทํ ปญฺญเปสฺสามา”ติ. อถ โข ภควา เต ภิกฺขู เอตทโวจ “ยถากถํ ปน ตุเมฺห ภิกฺขเว สมคฺคา สมฺโมทมานา อวิวทมานา ผาสุกํ วสฺสํ วสิตฺถ, น จ ปิณฺฑเกน กิลมิตฺถา”ติ.}

\prose{อิธ มยํ ภนฺเต สมฺพหุลา สนฺทิฏฺฐา สมฺภตฺตา ภิกฺขู โกสเลสุ ชนปเท อญฺญตรสฺมิํ อาวาเส วสฺสํ อุปคจฺฉิมฺหา. เตสํ โน ภนฺเต อมฺหากํ เอตทโหสิ “เกน นุ โข มยํ อุปาเยน สมคฺคา สมฺโมทมานา อวิวทมานา ผาสุกํ วสฺสํ วเสยฺยาม, น จ ปิณฺฑเกน กิลเมยฺยามา”ติ. เตสํ โน ภนฺเต อมฺหากํ เอตทโหสิ “สเจ โข มยํ อญฺญมญฺญํ เนว อาลเปยฺยาม น สลฺลเปยฺยาม, โย ปฐมํ คามโต ปิณฺฑาย ปฏิกฺกเมยฺย, โส อาสนํ ปญฺญเปยฺย, ปาโททกํ ปาทปีฐํ ปาทกถลิกํ อุปนิกฺขิเปยฺย, อวกฺการปาติํ โธวิตฺวา อุปฏฺฐาเปยฺย, ปานียํ ปริโภชนียํ อุปฏฺฐาเปยฺย, โย ปจฺฉา คามโต ปิณฺฑาย ปฏิกฺกเมยฺย, สจสฺส \csromanfolio{222}ภุตฺตาวเสโส, สเจ อากงฺเขยฺย ภุญฺเชยฺย, โน เจ อากงฺเขยฺย อปฺปหริเต วา ฉฑฺเฑยฺย, อปฺปาณเก วา อุทเก โอปิลาเปยฺย, โส อาสนํ อุทฺธเรยฺย, ปาโททกํ ปาทปีฐํ ปาทกถลิกํ ปฏิสาเมยฺย, อวกฺการปาติํ โธวิตฺวา ปฏิสาเมยฺย, ปานียํ ปริโภชนียํ ปฏิสาเมยฺย, ภตฺตคฺคํ สมฺมชฺเชยฺย, โย ปสฺเสยฺย ปานียฆฏํ วา ปริโภชนียฆฏํ วา วจฺจฆฏํ วา ริตฺตํ ตุจฺฉํ, โส อุปฏฺฐาเปยฺย, สจสฺส โหติ อวิสยฺหํ, หตฺถวิกาเรน ทุติยํ อามนฺเตตฺวา หตฺถ\-วิลงฺฆเกน อุปฏฺฐาเปยฺย, น เตฺวว ตปฺปจฺจยา วาจํ ภินฺเทยฺย, เอวํ โข มยํ สมคฺคา สมฺโมทมานา อวิวทมานา ผาสุกํ วสฺสํ วเสยฺยาม, น จ ปิณฺฑเกน กิลเมยฺยามา”ติ. อถ โข มยํ ภนฺเต อญฺญมญฺญํ เนว อาลปิมฺหา น สลฺลปิมฺหา, โย ปฐมํ คามโต ปิณฺฑาย ปฏิกฺกมติ, โส อาสนํ ปญฺญเปติ, ปาโททกํ ปาทปีฐํ ปาทกถลิกํ อุปนิกฺขิปติ, อวกฺการปาติํ โธวิตฺวา อุปฏฺฐาเปติ, ปานียํ ปริโภชนียํ อุปฏฺฐาเปติ, โย ปจฺฉา คามโต ปิณฺฑาย ปฏิกฺกมติ, สเจ โหติ ภุตฺตาวเสโส, สเจ อากงฺขติ ภุญฺชติ, โน เจ อากงฺขติ อปฺปหริเต วา ฉฑฺเฑติ, อปฺปาณเก วา อุทเก โอปิลาเปติ, โส อาสนํ อุทฺธรติ, ปาโททกํ ปาทปีฐํ ปาทกถลิกํ ปฏิสาเมติ, อวกฺการปาติํ โธวิตฺวา ปฏิสาเมติ, ปานียํ ปริโภชนียํ ปฏิสาเมติ, ภตฺตคฺคํ สมฺมชฺชติ, โย ปสฺสติ ปานียฆฏํ วา ปริโภชนียฆฏํ วา วจฺจฆฏํ วา ริตฺตํ ตุจฺฉํ, โส อุปฏฺฐาเปติ, สจสฺส โหติ อวิสยฺหํ, หตฺถวิกาเรน ทุติยํ อามนฺเตตฺวา หตฺถ\-วิลงฺฆเกน อุปฏฺฐาเปติ, น เตฺวว ตปฺปจฺจยา วาจํ ภินฺทติ, เอวํ โข มยํ ภนฺเต สมคฺคา สมฺโมทมานา อวิวทมานา ผาสุกํ วสฺสํ วสิมฺหา, น จ ปิณฺฑเกน กิลมิมฺหาติ.}

\prose{อถ โข ภควา ภิกฺขู อามนฺเตสิ อผาสุญฺเญว\csromansharedfootnote{666db47f545d6640}{อผาสุกญฺเญว (สี)} กิร’เม\csromansharedfootnote{3f4bf6f29f0bf6b8}{กิริเม (ก)} ภิกฺขเว โมฆปุริสา วุฏฺฐา\csromansharedfootnote{29a58de4756de854}{วุตฺถา (ก)} สมานา “ผาสุมฺหา\csromansharedfootnote{d2e793567890c531}{ผาสุกมฺหา (สี)} วุฏฺฐา”ติ ปฏิชานนฺติ. ปสุสํวาสญฺเญว กิร’เม ภิกฺขเว โมฆปุริสา วุฏฺฐา สมานา “ผาสุมฺหา วุฏฺฐา”ติ ปฏิชานนฺติ. เอฬก\-สํวาสญฺเญว กิร’เม ภิกฺขเว โมฆปุริสา วุฏฺฐา สมานา “ผาสุมฺหา วุฏฺฐา”ติ ปฏิชานนฺติ. สปตฺต\-สํวาสญฺเญว กิร’เม ภิกฺขเว โมฆปุริสา วุฏฺฐา สมานา “ผาสุมฺหา วุฏฺฐา”ติ ปฏิชานนฺติ, กถํ หิ นามิเม ภิกฺขเว โมฆปุริสา มูคพฺพตํ ติตฺถิย\-สมาทานํ สมาทิยิสฺสนฺติ. เนตํ ภิกฺขเว อปฺปสนฺนานํ วา ปสาทาย ฯเปฯ วิครหิตฺวา \csromanfolio{223}ธมฺมิํ กถํ กตฺวา ภิกฺขู อามนฺเตสิ น ภิกฺขเว มูคพฺพตํ ติตฺถิย\-สมาทานํ สมาทิยิตพฺพํ, โย สมาทิเยยฺย, อาปตฺติ ทุกฺกฏสฺส. อนุชานามิ ภิกฺขเว วสฺสํ\-วุฏฺฐานํ ภิกฺขูนํ ตีหิ ฐาเนหิ ปวาเรตุํ ทิฏฺเฐน วา สุเตน วา ปริสงฺกาย วา, สา โว ภวิสฺสติ อญฺญมญฺญานุโลมตา อาปตฺติวุฏฺฐานตา วินยปุเรกฺขารตา. เอวญฺจ ปน ภิกฺขเว ปวาเรตพฺพํ, พฺยตฺเตน ภิกฺขุนา ปฏิพเลน สํโฆ ญาเปตพฺโพ\csromandash{}}

\proseitem{210}{“สุณาตุ เม ภนฺเต สํโฆ, อชฺช ปวารณา, ยทิ สํฆสฺส ปตฺตกลฺลํ, สํโฆ ปวาเรยฺยา”ติ.}

\prose{เถเรน ภิกฺขุนา เอกํสํ อุตฺตราสงฺคํ กริตฺวา อุกฺกุฏิกํ นิสีทิตฺวา อญฺชลิํ ปคฺคเหตฺวา เอวมสฺส วจนีโย “สํฆํ อาวุโส ปวาเรมิ ทิฏฺเฐน วา สุเตน วา ปริสงฺกาย วา, วทนฺตุ มํ อายสฺมนฺโต อนุกมฺปํ อุปาทาย, ปสฺสนฺโต ปฏิกริสฺสามิ.}

\prose{ทุติยมฺปิ อาวุโส สํฆํ ปวาเรมิ ทิฏฺเฐน วา สุเตน วา ปริสงฺกาย วา, วทนฺตุ มํ อายสฺมนฺโต อนุกมฺปํ อุปาทาย, ปสฺสนฺโต ปฏิกริสฺสามิ.}

\prose{ตติยมฺปิ อาวุโส สํฆํ ปวาเรมิ ทิฏฺเฐน วา สุเตน วา ปริสงฺกาย วา, วทนฺตุ มํ อายสฺมนฺโต อนุกมฺปํ อุปาทาย, ปสฺสนฺโต ปฏิกริสฺสามี”ติ.}

\prose{นวเกน ภิกฺขุนา เอกํสํ อุตฺตราสงฺคํ กริตฺวา อุกฺกุฏิกํ นิสีทิตฺวา อญฺชลิํ ปคฺคเหตฺวา เอวมสฺส วจนีโย “สํฆํ ภนฺเต ปวาเรมิ ทิฏฺเฐน วา สุเตน วา ปริสงฺกาย วา, วทนฺตุ มํ อายสฺมนฺโต อนุกมฺปํ อุปาทาย, ปสฺสนฺโต ปฏิกริสฺสามิ.}

\prose{ทุติยมฺปิ ภนฺเต สํฆํ ฯเปฯ. ตติยมฺปิ ภนฺเต สํฆํ ปวาเรมิ ทิฏฺเฐน วา สุเตน วา ปริสงฺกาย วา, วทนฺตุ มํ อายสฺมนฺโต อนุกมฺปํ อุปาทาย, ปสฺสนฺโต ปฏิกริสฺสามี”ติ.}

\proseitem{211}{เตน โข ปน สมเยน ฉพฺพคฺคิยา ภิกฺขู เถเรสุ ภิกฺขูสุ อุกฺกุฏิกํ นิสินฺเนสุ ปวารยมาเนสุ อาสเนสุ อจฺฉนฺติ. เย เต ภิกฺขู อปฺปิจฺฉา ฯเปฯ เต อุชฺฌายนฺติ ขิยฺยนฺติ วิปาเจนฺติ “กถํ หิ นาม ฉพฺพคฺคิยา ภิกฺขู เถเรสุ ภิกฺขูสุ อุกฺกุฏิกํ นิสินฺเนสุ ปวารยมาเนสุ อาสเนสุ อจฺฉิสฺสนฺตี”ติ. อถ โข เต ภิกฺขู ภควโต เอตมตฺถํ อาโรเจสุํ ฯเปฯ. สจฺจํ กิร ภิกฺขเว ฉพฺพคฺคิยา ภิกฺขู เถเรสุ ภิกฺขูสุ อุกฺกุฏิกํ นิสินฺเนสุ ปวารยมาเนสุ อาสเนสุ อจฺฉนฺตีติ. สจฺจํ ภควาติ. \csromanfolio{224}วิครหิ พุทฺโธ ภควา ฯเปฯ. กถํ หิ นาม เต ภิกฺขเว โมฆปุริสา เถเรสุ ภิกฺขูสุ อุกฺกุฏิกํ นิสินฺเนสุ ปวารยมาเนสุ อาสเนสุ อจฺฉิสฺสนฺติ. เนตํ ภิกฺขเว อปฺปสนฺนานํ วา ปสาทาย, ปสนฺนานํ วา ภิยฺโยภาวาย ฯเปฯ วิครหิตฺวา ฯเปฯ ธมฺมิํ กถํ กตฺวา ภิกฺขู อามนฺเตสิ “น ภิกฺขเว เถเรสุ ภิกฺขูสุ อุกฺกุฏิกํ นิสินฺเนสุ ปวารยมาเนสุ อาสเนสุ อจฺฉิตพฺพํ, โย อจฺเฉยฺย, อาปตฺติ ทุกฺกฏสฺส. อนุชานามิ ภิกฺขเว สพฺเพเหว อุกฺกุฏิกํ นิสินฺเนหิ ปวาเรตุนฺ”ติ.}

\prose{เตน โข ปน สมเยน อญฺญตโร เถโร ชราทุพฺพโล “ยาว สพฺเพ ปวาเรนฺตี”ติ\csromansharedfootnote{0bf4098958cfb0d5}{ยาว สพฺเพ ปวาเรนฺติ (สฺยา)} อุกฺกุฏิกํ นิสินฺโน อาคมยมาโน มุจฺฉิโต ปปติ. ภควโต เอตมตฺถํ อาโรเจสุํ. อนุชานามิ ภิกฺขเว ตทมนฺตรา อุกฺกุฏิกํ นิสีทิตุํ, ยาว ปวาเรติ, ปวาเรตฺวา อาสเน นิสีทิตุนฺติ.}

\csromanmatikaanchor{mtk.125}
\csromantocmarkref{section}{121. ปวารณาเภท}{mtk.125}
\bookmark[dest={mtk.125},level=1]{121. ปวารณาเภท}
\csromanheader{1}{121. \textbf{ปวารณาเภท}}
\proseitem{212}{อถ โข ภิกฺขูนํ เอตทโหสิ “กติ นุ โข ปวารณา”ติ, ภควโต เอตมตฺถํ อาโรเจสุํ.}

\prose{เทฺวมา ภิกฺขเว ปวารณา จาตุทฺทสิกา จ ปนฺนรสิกา จ, อิมา โข ภิกฺขเว เทฺว ปวารณาติ.}

\prose{อถ โข ภิกฺขูนํ เอตทโหสิ “กติ นุ โข ปวารณกมฺมานี”ติ\csromansharedfootnote{bd1ca77eef399e0f}{ปวารณากมฺมานีติ (สฺยา)}. ภควโต เอตมตฺถํ อาโรเจสุํ. จตฺตาริมานิ ภิกฺขเว ปวารณ\-กมฺมานิ, อธมฺเมน วคฺคํ ปวารณกมฺมํ, อธมฺเมน สมคฺคํ ปวารณกมฺมํ, ธมฺเมน วคฺคํ ปวารณกมฺมํ, ธมฺเมน สมคฺคํ ปวารณกมฺมํ. ตตฺร ภิกฺขเว ยทิทํ อธมฺเมน วคฺคํ ปวารณกมฺมํ, น ภิกฺขเว เอวรูปํ ปวารณกมฺมํ กาตพฺพํ, น จ มยา เอวรูปํ ปวารณกมฺมํ อนุญฺญาตํ. ตตฺร ภิกฺขเว ยทิทํ อธมฺเมน สมคฺคํ ปวารณกมฺมํ, น ภิกฺขเว เอวรูปํ ปวารณกมฺมํ กาตพฺพํ, น จ มยา เอวรูปํ ปวารณกมฺมํ อนุญฺญาตํ. ตตฺร ภิกฺขเว ยทิทํ ธมฺเมน วคฺคํ ปวารณกมฺมํ, น ภิกฺขเว เอวรูปํ ปวารณกมฺมํ กาตพฺพํ, น จ มยา เอวรูปํ ปวารณกมฺมํ อนุญฺญาตํ. ตตฺร ภิกฺขเว ยทิทํ ธมฺเมน สมคฺคํ ปวารณกมฺมํ, เอวรูปํ ภิกฺขเว ปวารณกมฺมํ กาตพฺพํ, เอวรูปญฺจ มยา ปวารณกมฺมํ อนุญฺญาตํ. ตสฺมาติห ภิกฺขเว เอวรูปํ ปวารณกมฺมํ กริสฺสาม ยทิทํ ธมฺเมน สมคฺคนฺติ, เอวญฺหิ โว ภิกฺขเว สิกฺขิตพฺพนฺติ.}

\csromanmatikaanchor{mtk.126}
\csromantocmarkref{section}{122. ปวารณาทานานุชานนา}{mtk.126}
\bookmark[dest={mtk.126},level=1]{122. ปวารณาทานานุชานนา}
\csromanheader{1}{122. \textbf{ปวารณา}\-\textbf{ทานา}\-\textbf{นุชานนา}}
\proseitem{213}{\csromanfolio{225}อถ โข ภควา ภิกฺขู อามนฺเตสิ “สนฺนิปตถ ภิกฺขเว สํโฆ ปวาเรสฺสตี”ติ. เอวํ วุตฺเต อญฺญตโร ภิกฺขุ ภควนฺตํ เอตทโวจ “อตฺถิ ภนฺเต ภิกฺขุ คิลาโน, โส อนาคโต”ติ. อนุชานามิ ภิกฺขเว คิลาเนน ภิกฺขุนา ปวารณํ ทาตุํ. เอวญฺจ ปน ภิกฺขเว ทาตพฺพา, เตน คิลาเนน ภิกฺขุนา เอกํ ภิกฺขุํ อุปสงฺกมิตฺวา เอกํสํ อุตฺตราสงฺคํ กริตฺวา อุกฺกุฏิกํ นิสีทิตฺวา อญฺชลิํ ปคฺคเหตฺวา เอวมสฺส วจนีโย “ปวารณํ ทมฺมิ, ปวารณํ เม หร, ปวารณํ เม อาโรเจหิ, มมตฺถาย ปวาเรหี”ติ, กาเยน วิญฺญาเปติ, วาจาย วิญฺญาเปติ, กาเยน วาจาย วิญฺญาเปติ, ทินฺนา โหติ ปวารณา. น กาเยน วิญฺญาเปติ, น วาจาย วิญฺญาเปติ, น กาเยน วาจาย วิญฺญาเปติ, น ทินฺนา โหติ ปวารณา. เอวญฺเจตํ ลเภถ, อิจฺเจตํ กุสลํ. โน เจ ลเภถ, โส ภิกฺขเว คิลาโน ภิกฺขุ มญฺเจน วา ปีเฐน วา สํฆมชฺเฌ อาเนตฺวา ปวาเรตพฺพํ. สเจ ภิกฺขเว คิลานุ\-ปฏฺฐากานํ ภิกฺขูนํ เอตทโหสิ “สเจ โข มยํ คิลานํ ฐานา จาเวสฺสาม, อาพาโธ วา อภิวฑฺฒิสฺสติ, กาลํกิริยา วา ภวิสฺสตี”ติ. น ภิกฺขเว คิลาโน ภิกฺขุ ฐานา จาเวตพฺโพ, สํเฆน ตตฺถ คนฺตฺวา ปวาเรตพฺพํ. น เตฺวว วคฺเคน สํเฆน ปวาเรตพฺพํ, ปวาเรยฺย เจ, อาปตฺติ ทุกฺกฏสฺส.}

\prose{ปวารณหารโก\csromansharedfootnote{af24b03b7f356b09}{ปวารณาหารโก (สฺยา)} เจ ภิกฺขเว ทินฺนาย ปวารณาย ตตฺเถว ปกฺกมติ, อญฺญสฺส ทาตพฺพา ปวารณา. ปวารณหารโก เจ ภิกฺขเว ทินฺนาย ปวารณาย ตตฺเถว วิพฺภมติ ฯเปฯ กาลํ กโรติ. สามเณโร ปฏิชานาติ. สิกฺขํ ปจฺจกฺขาตโก ปฏิชานาติ. อนฺติมวตฺถุํ อชฺฌาปนฺนโก ปฏิชานาติ. อุมฺมตฺตโก ปฏิชานาติ. ขิตฺตจิตฺโต ปฏิชานาติ. เวทนาฏฺโฏ ปฏิชานาติ. อาปตฺติยา อทสฺสเน อุกฺขิตฺตโก ปฏิชานาติ. อาปตฺติยา อปฺปฏิกมฺเม อุกฺขิตฺตโก ปฏิชานาติ. ปาปิกาย ทิฏฺฐิยา อปฺปฏินิสฺสคฺเค อุกฺขิตฺตโก ปฏิชานาติ. ปณฺฑโก ปฏิชานาติ. เถยฺย\-สํวาสโก ปฏิชานาติ. ติตฺถิย\-ปกฺกนฺตโก ปฏิชานาติ. ติรจฺฉานคโต ปติชานาติ. มาตุฆาตโก ปฏิชานาติ. ปิตุฆาตโก ปฏิชานาติ. อรหนฺต\-ฆาตโก ปฏิชานาติ. ภิกฺขุนิ\-ทูสโก ปฏิชานาติ. สํฆเภทโก \csromanfolio{226}ปฏิชานาติ. โลหิตุปฺปาทโก ปฏิชานาติ. อุภโต\-พฺยญฺชนโก ปฏิชานาติ, อญฺญสฺส ทาตพฺพา ปวารณา.}

\prose{ปวารณหารโก เจ ภิกฺขเว ทินฺนาย ปวารณาย อนฺตรามคฺเค ปกฺกมติ, อนาหฏา โหติ ปวารณา.}

\prose{ปวารณหารโก เจ ภิกฺขเว ทินฺนาย ปวารณาย อนฺตรามคฺเค วิพฺภมติ ฯเปฯ กาลํ กโรติ. สามเณโร ปฏิชานาติ. สิกฺขํ ปจฺจกฺขาตโก ปฏิชานาติ. อนฺติมวตฺถุํ อชฺฌาปนฺนโก ปฏิชานาติ. อุมฺมตฺตโก ปฏิชานาติ. ขิตฺตจิตฺโต ปฏิชานาติ. เวทนาฏฺโฏ ปฏิชานาติ. อาปตฺติยา อทสฺสเน อุกฺขิตฺตโก ปฏิชานาติ. อาปตฺติยา อปฺปฏิกมฺเม อุกฺขิตฺตโก ปฏิชานาติ. ปาปิกาย ทิฏฺฐิยา อปฺปฏินิสฺสคฺเค อุกฺขิตฺตโก ปฏิชานาติ. ปณฺฑโก ปฏิชานาติ. เถยฺย\-สํวาสโก ปฏิชานาติ. ติตฺถิย\-ปกฺกนฺตโก ปฏิชานาติ. ติรจฺฉานคโต ปฏิชานาติ. มาตุฆาตโก ปฏิชานาติ. ปิตุฆาตโก ปฏิชานาติ. อรหนฺต\-ฆาตโก ปฏิชานาติ. ภิกฺขุนิ\-ทูสโก ปฏิชานาติ. สํฆเภทโก ปฏิชานาติ. โลหิตุปฺปาทโก ปฏิชานาติ. อุภโต\-พฺยญฺชนโก ปฏิชานาติ, อนาหฏา โหติ ปวารณา.}

\prose{ปวารณหารโก เจ ภิกฺขเว ทินฺนาย ปวารณาย สํฆปฺปตฺโต ปกฺกมติ, อาหฏา โหติ ปวารณา.}

\prose{ปวารณหารโก เจ ภิกฺขเว ทินฺนาย ปวารณาย สํฆปฺปตฺโต วิพฺภมติ ฯเปฯ กาลํ กโรติ. สามเณโร ปฏิชานาติ. สิกฺขํ ปจฺจกฺขาตโก ปฏิชานาติ. อนฺติมวตฺถุํ อชฺฌาปนฺนโก ปฏิชานาติ. อุมฺมตฺตโก ปฏิชานาติ. ขิตฺตจิตฺโต ปฏิชานาติ. เวทนาฏฺโฏ ปฏิชานาติ. อาปตฺติยา อทสฺสเน อุกฺขิตฺตโก ปฏิชานาติ. อาปตฺติยา อปฺปฏิกมฺเม อุกฺขิตฺตโก ปฏิชานาติ. ปาปิกาย ทิฏฺฐิยา อปฺปฏินิสฺสคฺเค อุกฺขิตฺตโก ปฏิชานาติ. ปณฺฑโก ปฏิชานาติ. เถยฺย\-สํวาสโก ปฏิชานาติ. ติตฺถิย\-ปกฺกนฺตโก ปฏิชานาติ. ติรจฺฉานคโต ปฏิชานาติ. มาตุฆาตโก ปฏิชานาติ. ปิตุฆาตโก ปฏิชานาติ. อรหนฺต\-ฆาตโก ปฏิชานาติ. ภิกฺขุนิ\-ทูสโก ปฏิชานาติ. สํฆเภทโก ปฏิชานาติ. โลหิตุปฺปาทโก ปฏิชานาติ. อุภโต\-พฺยญฺชนโก ปฏิชานาติ, อาหฏา โหติ ปวารณา.}

\proseitem{213}{ปวารณหารโก เจ ภิกฺขเว ทินฺนาย ปวารณาย สํฆปฺปตฺโต สุตฺโต นาโรเจติ, อาหฏา โหติ ปวารณา, ปวารณ\-หารกสฺส \csromanfolio{227}อนาปตฺติ. ปวารณหารโก เจ ภิกฺขเว ทินฺนาย ปวารณาย สํฆปฺปตฺโต ปมตฺโต นาโรเจติ ฯเปฯ สมาปนฺโน นาโรเจติ, อาหฏา โหติ ปวารณา, ปวารณ\-หารกสฺส อนาปตฺติ.}

\prose{ปวารณหารโก เจ ภิกฺขเว ทินฺนาย ปวารณาย สํฆปฺปตฺโต สญฺจิจฺจ นาโรเจติ, อาหฏา โหติ ปวารณา, ปวารณ\-หารกสฺส อาปตฺติ ทุกฺกฏสฺส. อนุชานามิ ภิกฺขเว ตทหุ ปวารณาย ปวารณํ เทนฺเตน ฉนฺทมฺปิ ทาตุํ, สนฺติ สํฆสฺส กรณียนฺติ.}

\csromanmatikaanchor{mtk.127}
\csromantocmarkref{section}{123. ญาตกาทิคฺคหณกถา}{mtk.127}
\bookmark[dest={mtk.127},level=1]{123. ญาตกาทิคฺคหณกถา}
\csromanheader{1}{123. \textbf{ญาตกา}\-\textbf{ทิ}\-\textbf{คฺคหณ}\-\textbf{กถา}}
\proseitem{214}{เตน โข ปน สมเยน อญฺญตรํ ภิกฺขุํ ตทหุ ปวารณาย ญาตกา คณฺหิํสุ. ภควโต เอตมตฺถํ อาโรเจสุํ. อิธ ปน ภิกฺขเว ภิกฺขุํ ตทหุ ปวารณาย ญาตกา คณฺหนฺติ, เต ญาตกา ภิกฺขูหิ เอวมสฺสุ วจนียา “อิงฺฆ ตุเมฺห อายสฺมนฺโต อิมํ ภิกฺขุํ มุหุตฺตํ มุญฺจถ, ยาวายํ ภิกฺขุ ปวาเรตี”ติ. เอวญฺเจตํ ลเภถ, อิจฺเจตํ กุสลํ. โน เจ ลเภถ, เต ญาตกา ภิกฺขูหิ เอวมสฺสุ วจนียา “อิงฺฆ ตุเมฺห อายสฺมนฺโต มุหุตฺตํ เอกมนฺตํ โหถ, ยาวายํ ภิกฺขุ ปวารณํ เทตี”ติ. เอวญฺเจตํ ลเภถ, อิจฺเจตํ กุสลํ. โน เจ ลเภถ, เต ญาตกา ภิกฺขูหิ เอวมสฺสุ วจนียา “อิงฺฆ ตุเมฺห อายสฺมนฺโต อิมํ ภิกฺขุํ มุหุตฺตํ นิสฺสีมํ เนถ, ยาว สํโฆ ปวาเรตี”ติ. เอวญฺเจตํ ลเภถ, อิจฺเจตํ กุสลํ. โน เจ ลเภถ, น เตฺวว วคฺเคน สํเฆน ปวาเรตพฺพํ, ปวาเรยฺย เจ, อาปตฺติ ทุกฺกฏสฺส.}

\prose{อิธ ปน ภิกฺขเว ภิกฺขุํ ตทหุ ปวารณาย ราชาโน คณฺหนฺติ ฯเปฯ โจรา คณฺหนฺติ. ธุตฺตา คณฺหนฺติ. ภิกฺขุ\-ปจฺจตฺถิกา คณฺหนฺติ, เต ภิกฺขุ\-ปจฺจตฺถิกา ภิกฺขูหิ เอวมสฺสุ วจนียา “อิงฺฆ ตุเมฺห อายสฺมนฺโต อิมํ ภิกฺขุํ มุหุตฺตํ มุญฺจถ, ยาวายํ ภิกฺขุ ปวาเรตี”ติ. เอวญฺเจตํ ลเภถ, อิจฺเจตํ กุสลํ. โน เจ ลเภถ, เต ภิกฺขุ\-ปจฺจตฺถิกา ภิกฺขูหิ เอวมสฺสุ วจนียา “อิงฺฆ ตุเมฺห อายสฺมนฺโต มุหุตฺตํ เอกมนฺตํ โหถ, ยาวายํ ภิกฺขุ ปวารณํ เทตี”ติ. เอวญฺเจตํ ลเภถ, อิจฺเจตํ กุสลํ. โน เจ ลเภถ, เต ภิกฺขุ\-ปจฺจตฺถิกา ภิกฺขูหิ เอวมสฺสุ วจนียา “อิงฺฆ ตุเมฺห อายสฺมนฺโต อิมํ ภิกฺขุํ มุหุตฺตํ นิสฺสีมํ เนถ, ยาว สํโฆ ปวาเรตี”ติ. เอวญฺเจตํ ลเภถ, อิจฺเจตํ กุสลํ. โน เจ ลเภถ, \csromanfolio{228}น เตฺวว วคฺเคน สํเฆน ปวาเรตพฺพํ, ปวาเรยฺย เจ, อาปตฺติ ทุกฺกฏสฺสาติ.}

\csromanmatikaanchor{mtk.128}
\csromantocmarkref{section}{124. สํฆปวารณาทิปฺปเภท}{mtk.128}
\bookmark[dest={mtk.128},level=1]{124. สํฆปวารณาทิปฺปเภท}
\csromanheader{1}{124. \textbf{สํฆปวารณาทิปฺปเภท}}
\proseitem{215}{เตน โข ปน สมเยน อญฺญตรสฺมิํ อาวาเส ตทหุ ปวารณาย ปญฺจ ภิกฺขู วิหรนฺติ. อถ โข เตสํ ภิกฺขูนํ เอตทโหสิ “ภควตา ปญฺญตฺตํ ‘สํเฆน ปวาเรตพฺพนฺ’ติ, มยญฺจมฺหา ปญฺจ ชนา, กถํ นุ โข อเมฺหหิ ปวาเรตพฺพนฺ”ติ. ภควโต เอตมตฺถํ อาโรเจสุํ. อนุชานามิ ภิกฺขเว ปญฺจนฺนํ สํเฆ ปวาเรตุนฺติ.}

\proseitem{216}{เตน โข ปน สมเยน อญฺญตรสฺมิํ อาวาเส ตทหุ ปวารณาย จตฺตาโร ภิกฺขู วิหรนฺติ. อถ โข เตสํ ภิกฺขูนํ เอตทโหสิ “ภควตา อนุญฺญาตํ ‘ปญฺจนฺนํ สํเฆ ปวาเรตุนฺ’ติ, มยญฺจมฺหา จตฺตาโร ชนา, กถํ นุ โข อเมฺหหิ ปวาเรตพฺพนฺ”ติ. ภควโต เอตมตฺถํ อาโรเจสุํ. อนุชานามิ ภิกฺขเว จตุนฺนํ อญฺญมญฺญํ ปวาเรตุํ. เอวญฺจ ปน ภิกฺขเว ปวาเรตพฺพํ, พฺยตฺเตน ภิกฺขุนา ปฏิพเลน เต ภิกฺขู ญาเปตพฺพา\csromandash{}}

\prose{“สุณนฺตุ เม อายสฺมนฺโต, อชฺช ปวารณา, ยทา\-ยสฺมนฺตานํ ปตฺตกลฺลํ, มยํ อญฺญมญฺญํ ปวาเรยฺยามา”ติ.}

\prose{เถเรน ภิกฺขุนา เอกํสํ อุตฺตราสงฺคํ กริตฺวา อุกฺกุฏิกํ นิสีทิตฺวา อญฺชลิํ ปคฺคเหตฺวา เต ภิกฺขู เอวมสฺสุ วจนียา “อหํ อาวุโส อายสฺมนฺเต ปวาเรมิ ทิฏฺเฐน วา สุเตน วา ปริสงฺกาย วา, วทนฺตุ มํ อายสฺมนฺโต อนุกมฺปํ อุปาทาย, ปสฺสนฺโต ปฏิกริสฺสามิ.}

\prose{ทุติยมฺปิ ฯเปฯ. ตติยมฺปิ อหํ อาวุโส อายสฺมนฺเต ปวาเรมิ ทิฏฺเฐน วา สุเตน วา ปริสงฺกาย วา, วทนฺตุ มํ อายสฺมนฺโต อนุกมฺปํ อุปาทาย, ปสฺสนฺโต ปฏิกริสฺสามี”ติ.}

\proseitem{216}{นวเกน ภิกฺขุนา เอกํสํ อุตฺตราสงฺคํ กริตฺวา อุกฺกุฏิกํ นิสีทิตฺวา อญฺชลิํ ปคฺคเหตฺวา เต ภิกฺขู เอวมสฺสุ วจนียา “อหํ ภนฺเต อายสฺมนฺเต ปวาเรมิ ทิฏฺเฐน วา สุเตน วา ปริสงฺกาย วา, วทนฺตุ มํ อายสฺมนฺโต อนุกมฺปํ อุปาทาย, ปสฺสนฺโต ปฏิกริสฺสามิ. ทุติยมฺปิ ฯเปฯ. ตติยมฺปิ อหํ ภนฺเต อายสฺมนฺเต ปวาเรมิ ทิฏฺเฐน วา สุเตน วา \csromanfolio{229}ปริสงฺกาย วา, วทนฺตุ มํ อายสฺมนฺโต อนุกมฺปํ อุปาทาย, ปสฺสนฺโต ปฏิกริสฺสามี”ติ.}

\prose{เตน โข ปน สมเยน อญฺญตรสฺมิํ อาวาเส ตทหุ ปวารณาย ตโย ภิกฺขู วิหรนฺติ. อถ โข เตสํ ภิกฺขูนํ เอตทโหสิ “ภควตา อนุญฺญาตํ ปญฺจนฺนํ สํเฆ ปวาเรตุํ, จตุนฺนํ อญฺญมญฺญํ ปวาเรตุํ, มยญฺจมฺหา ตโย ชนา, กถํ นุ โข อเมฺหหิ ปวาเรตพฺพนฺ”ติ. ภควโต เอตมตฺถํ อาโรเจสุํ. อนุชานามิ ภิกฺขเว ติณฺณํ อญฺญมญฺญํ ปวาเรตุํ. เอวญฺจ ปน ภิกฺขเว ปวาเรตพฺพํ, พฺยตฺเตน ภิกฺขุนา ปฏิพเลน เต ภิกฺขู ญาเปตพฺพา\csromandash{}}

\prose{“สุณนฺตุ เม อายสฺมนฺตา, อชฺช ปวารณา, ยทา\-ยสฺมนฺตานํ ปตฺตกลฺลํ, มยํ อญฺญมญฺญํ ปวาเรยฺยามา”ติ.}

\prose{เถเรน ภิกฺขุนา เอกํสํ อุตฺตราสงฺคํ กริตฺวา อุกฺกุฏิกํ นิสีทิตฺวา อญฺชลิํ ปคฺคเหตฺวา เต ภิกฺขู เอวมสฺสุ วจนียา “อหํ อาวุโส อายสฺมนฺเต ปวาเรมิ ทิฏฺเฐน วา สุเตน วา ปริสงฺกาย วา, วทนฺตุ มํ อายสฺมนฺตา อนุกมฺปํ อุปาทาย, ปสฺสนฺโต ปฏิกริสฺสามิ.}

\prose{ทุติยมฺปิ ฯเปฯ. ตติยมฺปิ อหํ อาวุโส อายสฺมนฺเต ปวาเรมิ ทิฏฺเฐน วา สุเตน วา ปริสงฺกาย วา, วทนฺตุ มํ อายสฺมนฺตา อนุกมฺปํ อุปาทาย, ปสฺสนฺโต ปฏิกริสฺสามี”ติ.}

\prose{นวเกน ภิกฺขุนา เอกํสํ อุตฺตราสงฺคํ กริตฺวา อุกฺกุฏิกํ นิสีทิตฺวา อญฺชลิํ ปคฺคเหตฺวา เต ภิกฺขู เอวมสฺสุ วจนียา “อหํ ภนฺเต อายสฺมนฺเต ปวาเรมิ ทิฏฺเฐน วา สุเตน วา ปริสงฺกาย วา, วทนฺตุ มํ อายสฺมนฺตา อนุกมฺปํ อุปาทาย, ปสฺสนฺโต ปฏิกริสฺสามิ.}

\prose{ทุติยมฺปิ ฯเปฯ. ตติยมฺปิ อหํ ภนฺเต อายสฺมนฺเต ปวาเรมิ ทิฏฺเฐน วา สุเตน วา ปริสงฺกาย วา, วทนฺตุ มํ อายสฺมนฺตา อนุกมฺปํ อุปาทาย, ปสฺสนฺโต ปฏิกริสฺสามี”ติ.}

\proseitem{217}{เตน โข ปน สมเยน อญฺญตรสฺมิํ อาวาเส ตทหุ ปวารณาย เทฺว ภิกฺขู วิหรนฺติ. อถ โข เตสํ ภิกฺขูนํ เอตทโหสิ “ภควตา อนุญฺญาตํ ปญฺจนฺนํ สํเฆ ปวาเรตุํ, จตุนฺนํ อญฺญมญฺญํ ปวาเรตุํ, \csromanfolio{230}ติณฺณํ อญฺญมญฺญํ ปวาเรตุํ, มยญฺจมฺหา เทฺว ชนา, กถํ นุ โข อเมฺหหิ ปวาเรตพฺพนฺ”ติ. ภควโต เอตมตฺถํ อาโรเจสุํ. อนุชานามิ ภิกฺขเว ทฺวินฺนํ อญฺญมญฺญํ ปวาเรตุํ. เอวญฺจ ปน ภิกฺขเว ปวาเรตพฺพํ, เถเรน ภิกฺขุนา เอกํสํ อุตฺตราสงฺคํ กริตฺวา อุกฺกุฏิกํ นิสีทิตฺวา อญฺชลิํ ปคฺคเหตฺวา นโว ภิกฺขุ เอวมสฺส วจนีโย “อหํ อาวุโส อายสฺมนฺตํ ปวาเรมิ ทิฏฺเฐน วา สุเตน วา ปริสงฺกาย วา, วทตุ มํ อายสฺมา อนุกมฺปํ อุปาทาย, ปสฺสนฺโต ปฏิกริสฺสามิ. ทุติยมฺปิ ฯเปฯ. ตติยมฺปิ อหํ อาวุโส อายสฺมนฺตํ ปวาเรมิ ทิฏฺเฐน วา สุเตน วา ปริสงฺกาย วา, วทตุ มํ อายสฺมา อนุกมฺปํ อุปาทาย, ปสฺสนฺโต ปฏิกริสฺสามี”ติ.}

\prose{นวเกน ภิกฺขุนา เอกํสํ อุตฺตราสงฺคํ กริตฺวา อุกฺกุฏิกํ นิสีทิตฺวา อญฺชลิํ ปคฺคเหตฺวา เถโร ภิกฺขุ เอวมสฺส วจนีโย “อหํ ภนฺเต อายสฺมนฺตํ ปวาเรมิ ทิฏฺเฐน วา สุเตน วา ปริสงฺกาย วา, วทตุ มํ อายสฺมา อนุกมฺปํ อุปาทาย, ปสฺสนฺโต ปฏิกริสฺสามิ. ทุติยมฺปิ ฯเปฯ. ตติยมฺปิ อหํ ภนฺเต อายสฺมนฺตํ ปวาเรมิ ทิฏฺเฐน วา สุเตน วา ปริสงฺกาย วา, วทตุ มํ อายสฺมา อนุกมฺปํ อุปาทาย, ปสฺสนฺโต ปฏิกริสฺสามี”ติ.}

\proseitem{218}{เตน โข ปน สมเยน อญฺญตรสฺมิํ อาวาเส ตทหุ ปวารณาย เอโก ภิกฺขุ วิหรติ. อถ โข ตสฺส ภิกฺขุโน เอตทโหสิ “ภควตา อนุญฺญาตํ ปญฺจนฺนํ สํเฆ ปวาเรตุํ, จตุนฺนํ อญฺญมญฺญํ ปวาเรตุํ, ติณฺณํ อญฺญมญฺญํ ปวาเรตุํ, ทฺวินฺนํ อญฺญมญฺญํ ปวาเรตุํ, อหญฺจมฺหิ เอกโก, กถํ นุ โข มยา ปวาเรตพฺพนฺ”ติ. ภควโต เอตมตฺถํ อาโรเจสุํ.}

\prose{อิธ ปน ภิกฺขเว อญฺญตรสฺมิํ อาวาเส ตทหุ ปวารณาย เอโก ภิกฺขุ วิหรติ, เตน ภิกฺขเว ภิกฺขุนา ยตฺถ ภิกฺขู ปฏิกฺกมนฺติ อุปฏฺฐาน\-สาลาย วา มณฺฑเป วา รุกฺขมูเล วา, โส เทโส สมฺมชฺชิตฺวา ปานียํ ปริโภชนียํ อุปฏฺฐาเปตฺวา อาสนํ ปญฺญเปตฺวา ปทีปํ กตฺวา นิสีทิตพฺพํ.}

\setlength{\csromangathapagewidth}{0pt}
\setlength{\csromangathawakwidth}{0pt}
\csromangathameasuregroup{สเจ อญฺเญ ภิกฺขู อาคจฺฉนฺติ,}{\csromangathabat{สเจ อญฺเญ ภิกฺขู อาคจฺฉนฺติ,}{เตหิ สทฺธิํ ปวาเรตพฺพํ.}}
\csromangathasetleft{สเจ อญฺเญ ภิกฺขู อาคจฺฉนฺติ,}
\csromangathagroup{\csromangathastanza{\csromangathabat{สเจ อญฺเญ ภิกฺขู อาคจฺฉนฺติ,}{เตหิ สทฺธิํ ปวาเรตพฺพํ.}}}

\prose{โน เจ อาคจฺฉนฺติ, “อชฺช เม ปวารณา”ติ อธิฏฺฐาตพฺพํ.}

\setlength{\csromangathapagewidth}{0pt}
\setlength{\csromangathawakwidth}{0pt}
\csromangathameasuregroup{โน เจ อธิฏฺเฐยฺย,}{\csromangathabat{โน เจ อธิฏฺเฐยฺย,}{อาปตฺติ ทุกฺกฏสฺส.}}
\csromangathasetleft{โน เจ อธิฏฺเฐยฺย,}
\csromangathagroup{\csromangathastanza{\csromangathabat{โน เจ อธิฏฺเฐยฺย,}{อาปตฺติ ทุกฺกฏสฺส.}}}

\prose{\csromanfolio{231}ตตฺร ภิกฺขเว ยตฺถ ปญฺจ ภิกฺขู วิหรนฺติ, น เอกสฺส ปวารณํ อาหริตฺวา จตูหิ สํเฆ ปวาเรตพฺพํ, ปวาเรยฺยุํ เจ, อาปตฺติ ทุกฺกฏสฺส.}

\prose{ตตฺร ภิกฺขเว ยตฺถ จตฺตาโร ภิกฺขู วิหรนฺติ, น เอกสฺส ปวารณํ อาหริตฺวา ตีหิ อญฺญมญฺญํ ปวาเรตพฺพํ, ปวาเรยฺยุํ เจ, อาปตฺติ ทุกฺกฏสฺส.}

\prose{ตตฺร ภิกฺขเว ยตฺถ ตโย ภิกฺขู วิหรนฺติ, น เอกสฺส ปวารณํ อาหริตฺวา ทฺวีหิ อญฺญมญฺญํ ปวาเรตพฺพํ, ปวาเรยฺยุํ เจ, อาปตฺติ ทุกฺกฏสฺส.}

\prose{ตตฺร ภิกฺขเว ยตฺถ เทฺว ภิกฺขู วิหรนฺติ, น เอกสฺส ปวารณํ อาหริตฺวา เอเกน อธิฏฺฐาตพฺพํ, อธิฏฺเฐยฺย เจ, อาปตฺติ ทุกฺกฏสฺสาติ.}

\csromanmatikaanchor{mtk.129}
\csromantocmarkref{section}{125. อาปตฺติปฏิกมฺมวิธิ}{mtk.129}
\bookmark[dest={mtk.129},level=1]{125. อาปตฺติปฏิกมฺมวิธิ}
\csromanheader{1}{125. \textbf{อาปตฺติปฏิกมฺมวิธิ}}
\proseitem{219}{เตน โข ปน สมเยน อญฺญตโร ภิกฺขุ ตทหุ ปวารณาย อาปตฺติํ อาปนฺโน โหติ. อถ โข ตสฺส ภิกฺขุโน เอตทโหสิ “ภควตา ปญฺญตฺตํ ‘น สาปตฺติเกน ปวาเรตพฺพนฺ’ติ, อหญฺจมฺหิ อาปตฺติํ อาปนฺโน, กถํ นุ โข มยา ปฏิปชฺชิตพฺพนฺ”ติ. ภควโต เอตมตฺถํ อาโรเจสุํ\csromansharedfootnote{bb7547e98636a836}{อาโรเจสิ (ก)}.}

\prose{อิธ ปน ภิกฺขเว ภิกฺขุ ตทหุ ปวารณาย อาปตฺติํ อาปนฺโน โหติ, เตน ภิกฺขเว ภิกฺขุนา เอกํ ภิกฺขุํ อุปสงฺกมิตฺวา เอกํสํ อุตฺตราสงฺคํ กริตฺวา อุกฺกุฏิกํ นิสีทิตฺวา อญฺชลิํ ปคฺคเหตฺวา เอวมสฺส วจนีโย “อหํ อาวุโส อิตฺถนฺนามํ อาปตฺติํ อาปนฺโน, ตํ ปฏิเทเสมี”ติ.}

\prose{เตน วตฺตพฺโพ “ปสฺสสี”ติ, “อาม ปสฺสามี”ติ, “อายติํ สํวเรยฺยาสี”ติ.}

\prose{อิธ ปน ภิกฺขเว ภิกฺขุ ตทหุ ปวารณาย อาปตฺติยา เวมติโก โหติ, เตน ภิกฺขเว ภิกฺขุนา เอกํ ภิกฺขุํ อุปสงฺกมิตฺวา เอกํสํ อุตฺตราสงฺคํ กริตฺวา อุกฺกุฏิกํ นิสีทิตฺวา อญฺชลิํ ปคฺคเหตฺวา เอวมสฺส วจนีโย “อหํ อาวุโส อิตฺถนฺนามาย อาปตฺติยา เวมติโก, ยทา นิพฺเพมติโก ภวิสฺสามิ, ตทา ตํ อาปตฺติํ ปฏิกริสฺสามี”ติ วตฺวา ปวาเรตพฺพํ, น เตฺวว ตปฺปจฺจยา ปวารณาย อนฺตราโย กาตพฺโพติ.}

\csromanmatikaanchor{mtk.130}
\csromantocmarkref{section}{126. อาปตฺติอาวิกรณวิธิ}{mtk.130}
\bookmark[dest={mtk.130},level=1]{126. อาปตฺติอาวิกรณวิธิ}
\csromanheader{1}{126. อาปตฺติ\-อาวิกรณวิธิ}
\proseitem{220}{\csromanfolio{232}เตน โข ปน สมเยน อญฺญตโร ภิกฺขุ ปวารยมาโน อาปตฺติํ สรติ. อถ โข ตสฺส ภิกฺขุโน เอตทโหสิ “ภควตา ปญฺญตฺตํ ‘น สาปตฺติเกน ปวาเรตพฺพนฺ’ติ, อหญฺจมฺหิ อาปตฺติํ อาปนฺโน, กถํ นุ โข มยา ปฏิปชฺชิตพฺพนฺ”ติ. ภควโต เอตมตฺถํ อาโรเจสุํ.}

\prose{อิธ ปน ภิกฺขเว ภิกฺขุ ปวารยมาโน อาปตฺติํ สรติ, เตน ภิกฺขเว ภิกฺขุนา สามนฺโต ภิกฺขุ เอวมสฺส วจนีโย “อหํ อาวุโส อิตฺถนฺนามํ อาปตฺติํ อาปนฺโน, อิโต วุฏฺฐหิตฺวา ตํ อาปตฺติํ ปฏิกริสฺสามี”ติ วตฺวา ปวาเรตพฺพํ, น เตฺวว ตปฺปจฺจยา ปวารณาย อนฺตราโย กาตพฺโพ.}

\prose{อิธ ปน ภิกฺขเว ภิกฺขุ ปวารยมาโน อาปตฺติยา เวมติโก โหติ, เตน ภิกฺขเว ภิกฺขุนา สามนฺโต ภิกฺขุ เอวมสฺส วจนีโย “อหํ อาวุโส อิตฺถนฺนามาย อาปตฺติยา เวมติโก, ยทา นิพฺเพมติโก ภวิสฺสามิ, ตทา ตํ อาปตฺติํ ปฏิกริสฺสามี”ติ วตฺวา ปวาเรตพฺพํ, น เตฺวว ตปฺปจฺจยา ปวารณาย อนฺตราโย กาตพฺโพติ.}

\csromanmatikaanchor{mtk.131}
\csromantocmarkref{section}{127. สภาคาปตฺติปฏิกมฺมวิธิ}{mtk.131}
\bookmark[dest={mtk.131},level=1]{127. สภาคาปตฺติปฏิกมฺมวิธิ}
\csromanheader{1}{127. \textbf{สภาคาปตฺติ}\-\textbf{ปฏิกมฺมวิธิ}}
\proseitem{221}{เตน โข ปน สมเยน อญฺญตรสฺมิํ อาวาเส ตทหุ ปวารณาย สพฺโพ สํโฆ สภาคํ อาปตฺติํ อาปนฺโน โหติ. อถ โข เตสํ ภิกฺขูนํ เอตทโหสิ “ภควตา ปญฺญตฺตํ ‘น สภาคา อาปตฺติ เทเสตพฺพา, น สภาคา อาปตฺติ ปฏิคฺคเหตพฺพา’ติ, อยญฺจ สพฺโพ สํโฆ สภาคํ อาปตฺติํ อาปนฺโน, กถํ นุ โข อเมฺหหิ ปฏิปชฺชิตพฺพนฺ”ติ. ภควโต เอตมตฺถํ อาโรเจสุํ.}

\prose{อิธ ปน ภิกฺขเว อญฺญตรสฺมิํ อาวาเส ตทหุ ปวารณาย สพฺโพ สํโฆ สภาคํ อาปตฺติํ อาปนฺโน โหติ, เตหิ ภิกฺขเว ภิกฺขูหิ เอโก ภิกฺขุ สามนฺตา อาวาสา สชฺชุกํ ปาเหตพฺโพ “คจฺฉาวุโส ตํ อาปตฺติํ ปฏิกริตฺวา อาคจฺฉ, มยํ เต สนฺติเก ตํ อาปตฺติํ ปฏิกริสฺสามา”ติ. เอวญฺเจตํ ลเภถ, อิจฺเจตํ กุสลํ. โน เจ ลเภถ, พฺยตฺเตน ภิกฺขุนา ปฏิพเลน สํโฆ ญาเปตพฺโพ\csromandash{}}

\prose{\csromanfolio{233}“สุณาตุ เม ภนฺเต สํโฆ, อยํ สพฺโพ สํโฆ สภาคํ อาปตฺติํ อาปนฺโน, ยทา อญฺญํ ภิกฺขุํ สุทฺธํ อนาปตฺติกํ ปสฺสิสฺสติ, ตทา ตสฺส สนฺติเก ตํ อาปตฺติํ ปฏิกริสฺสตี”ติ วตฺวา ปวาเรตพฺพํ, น เตฺวว ตปฺปจฺจยา ปวารณาย อนฺตราโย กาตพฺโพ.}

\prose{อิธ ปน ภิกฺขเว อญฺญตรสฺมิํ อาวาเส ตทหุ ปวารณาย สพฺโพ สํโฆ สภาคาย อาปตฺติยา เวมติโก โหติ, พฺยตฺเตน ภิกฺขุนา ปฏิพเลน สํโฆ ญาเปตพฺโพ\csromandash{}}

\prosewithcloserruled{“สุณาตุ เม ภนฺเต สํโฆ, อยํ สพฺโพ สํโฆ สภาคาย อาปตฺติยา เวมติโก, ยทา นิพฺเพมติโก ภวิสฺสติ, ตทา ตํ อาปตฺติํ ปฏิกริสฺสตี”ติ วตฺวา ปวาเรตพฺพํ, น เตฺวว ตปฺปจฺจยา ปวารณาย อนฺตราโย กาตพฺโพติ.}{ปฐม\-ภาณวาโร นิฏฺฐิโต.}
\csromanmatikaanchor{mtk.132}
\csromantocmarkref{section}{128. อนาปตฺติปนฺนรสก}{mtk.132}
\bookmark[dest={mtk.132},level=1]{128. อนาปตฺติปนฺนรสก}
\csromanheader{1}{128. \textbf{อนาปตฺติ}\-\textbf{ปนฺนรสก}}
\proseitem{222}{เตน โข ปน สมเยน อญฺญตรสฺมิํ อาวาเส ตทหุ ปวารณาย สมฺพหุลา อาวาสิกา ภิกฺขู สนฺนิปติํสุ ปญฺจ วา อติเรกา วา, เต น ชานิํสุ “อตฺถญฺเญ อาวาสิกา ภิกฺขู อนาคตา”ติ, เต ธมฺมสญฺญิโน วินยสญฺญิโน วคฺคา สมคฺคสญฺญิโน ปวาเรสุํ, เตหิ ปวาริยมาเน อถญฺเญ อาวาสิกา ภิกฺขู อาคจฺฉิํสุ พหุตรา. ภควโต เอตมตฺถํ อาโรเจสุํ.}

\prose{อิธ ปน ภิกฺขเว อญฺญตรสฺมิํ อาวาเส ตทหุ ปวารณาย สมฺพหุลา อาวาสิกา ภิกฺขู สนฺนิปตนฺติ ปญฺจ วา อติเรกา วา, เต น ชานนฺติ “อตฺถญฺเญ อาวาสิกา ภิกฺขู อนาคตา”ติ, เต ธมฺมสญฺญิโน วินยสญฺญิโน วคฺคา สมคฺคสญฺญิโน ปวาเรนฺติ, เตหิ ปวาริยมาเน อถญฺเญ อาวาสิกา ภิกฺขู อาคจฺฉนฺติ พหุตรา, เตหิ ภิกฺขเว ภิกฺขูหิ ปุน ปวาเรตพฺพํ, ปวาริตานํ อนาปตฺติ. (1)}

\prose{อิธ ปน ภิกฺขเว อญฺญตรสฺมิํ อาวาเส ตทหุ ปวารณาย สมฺพหุลา อาวาสิกา ภิกฺขู สนฺนิปตนฺติ ปญฺจ วา อติเรกา วา, เต น ชานนฺติ “อตฺถญฺเญ อาวาสิกา ภิกฺขู อนาคตา”ติ, เต ธมฺมสญฺญิโน \csromanfolio{234}วินยสญฺญิโน วคฺคา สมคฺคสญฺญิโน ปวาเรนฺติ, เตหิ ปวาริยมาเน อถญฺเญ อาวาสิกา ภิกฺขู อาคจฺฉนฺติ สมสมา, ปวาริตา สุปฺปวาริตา, อวเสเสหิ ปวาเรตพฺพํ, ปวาริตานํ อนาปตฺติ. (2)}

\prose{อิธ ปน ภิกฺขเว อญฺญตรสฺมิํ อาวาเส ตทหุ ปวารณาย สมฺพหุลา อาวาสิกา ภิกฺขู สนฺนิปตนฺติ ปญฺจ วา อติเรกา วา, เต น ชานนฺติ “อตฺถญฺเญ อาวาสิกา ภิกฺขู อนาคตา”ติ, เต ธมฺมสญฺญิโน วินยสญฺญิโน วคฺคา สมคฺคสญฺญิโน ปวาเรนฺติ, เตหิ ปวาริยมาเน อถญฺเญ อาวาสิกา ภิกฺขู อาคจฺฉนฺติ โถกตรา, ปวาริตา สุปฺปวาริตา, อวเสเสหิ ปวาเรตพฺพํ, ปวาริตานํ อนาปตฺติ. (3)}

\prose{อิธ ปน ภิกฺขเว อญฺญตรสฺมิํ อาวาเส ตทหุ ปวารณาย สมฺพหุลา อาวาสิกา ภิกฺขู สนฺนิปตนฺติ ปญฺจ วา อติเรกา วา, เต น ชานนฺติ “อตฺถญฺเญ อาวาสิกา ภิกฺขู อนาคตา”ติ, เต ธมฺมสญฺญิโน วินยสญฺญิโน วคฺคา สมคฺคสญฺญิโน ปวาเรนฺติ, เตหิ ปวาริตมตฺเต อถญฺเญ อาวาสิกา ภิกฺขู อาคจฺฉนฺติ พหุตรา, เตหิ ภิกฺขเว ภิกฺขูหิ ปุน ปวาเรตพฺพํ, ปวาริตานํ อนาปตฺติ. (4)}

\prose{อิธ ปน ภิกฺขเว อญฺญตรสฺมิํ อาวาเส ตทหุ ปวารณาย สมฺพหุลา อาวาสิกา ภิกฺขู สนฺนิปตนฺติ ปญฺจ วา อติเรกา วา, เต น ชานนฺติ “อตฺถญฺเญ อาวาสิกา ภิกฺขู อนาคตา”ติ, เต ธมฺมสญฺญิโน วินยสญฺญิโน วคฺคา สมคฺคสญฺญิโน ปวาเรนฺติ, เตหิ ปวาริตมตฺเต อถญฺเญ อาวาสิกา ภิกฺขู อาคจฺฉนฺติ สมสมา, ปวาริตา สุปฺปวาริตา, เตสํ สนฺติเก ปวาเรตพฺพํ, ปวาริตานํ อนาปตฺติ. (5)}

\prose{อิธ ปน ภิกฺขเว อญฺญตรสฺมิํ อาวาเส ตทหุ ปวารณาย สมฺพหุลา อาวาสิกา ภิกฺขู สนฺนิปตนฺติ ปญฺจ วา อติเรกา วา, เต น ชานนฺติ “อตฺถญฺเญ อาวาสิกา ภิกฺขู อนาคตา”ติ, เต ธมฺมสญฺญิโน วินยสญฺญิโน วคฺคา สมคฺคสญฺญิโน ปวาเรนฺติ, เตหิ ปวาริตมตฺเต อถญฺเญ อาวาสิกา ภิกฺขู อาคจฺฉนฺติ โถกตรา, ปวาริตา สุปฺปวาริตา, เตสํ สนฺติเก ปวาเรตพฺพํ, ปวาริตานํ อนาปตฺติ. (6)}

\prose{อิธ ปน ภิกฺขเว อญฺญตรสฺมิํ อาวาเส ตทหุ ปวารณาย สมฺพหุลา อาวาสิกา ภิกฺขู สนฺนิปตนฺติ ปญฺจ วา อติเรกา วา, เต น ชานนฺติ “อตฺถญฺเญ อาวาสิกา ภิกฺขู อนาคตา”ติ, เต ธมฺมสญฺญิโน \csromanfolio{235}วินยสญฺญิโน วคฺคา สมคฺคสญฺญิโน ปวาเรนฺติ, เตหิ ปวาริตมตฺเต อวุฏฺฐิตาย ปริสาย อถญฺเญ อาวาสิกา ภิกฺขู อาคจฺฉนฺติ พหุตรา, เตหิ ภิกฺขเว ภิกฺขูหิ ปุน ปวาเรตพฺพํ, ปวาริตานํ อนาปตฺติ. (7)}

\prose{อิธ ปน ภิกฺขเว อญฺญตรสฺมิํ อาวาเส ตทหุ ปวารณาย สมฺพหุลา อาวาสิกา ภิกฺขู สนฺนิปตนฺติ ปญฺจ วา อติเรกา วา, เต น ชานนฺติ “อตฺถญฺเญ อาวาสิกา ภิกฺขู อนาคตา”ติ, เต ธมฺมสญฺญิโน วินยสญฺญิโน วคฺคา สมคฺคสญฺญิโน ปวาเรนฺติ, เตหิ ปวาริตมตฺเต อวุฏฺฐิตาย ปริสาย อถญฺเญ อาวาสิกา ภิกฺขู อาคจฺฉนฺติ สมสมา, ปวาริตา สุปฺปวาริตา, เตสํ สนฺติเก ปวาเรตพฺพํ, ปวาริตานํ อนาปตฺติ. (8)}

\prose{อิธ ปน ภิกฺขเว อญฺญตรสฺมิํ อาวาเส ตทหุ ปวารณาย สมฺพหุลา อาวาสิกา ภิกฺขู สนฺนิปตนฺติ ปญฺจ วา อติเรกา วา, เต น ชานนฺติ “อตฺถญฺเญ อาวาสิกา ภิกฺขู อนาคตา”ติ, เต ธมฺมสญฺญิโน วินยสญฺญิโน วคฺคา สมคฺคสญฺญิโน ปวาเรนฺติ, เตหิ ปวาริตมตฺเต อวุฏฺฐิตาย ปริสาย อถญฺเญ อาวาสิกา ภิกฺขู อาคจฺฉนฺติ โถกตรา, ปวาริตา สุปฺปวาริตา, เตสํ สนฺติเก ปวาเรตพฺพํ, ปวาริตานํ อนาปตฺติ. (9)}

\prose{อิธ ปน ภิกฺขเว อญฺญตรสฺมิํ อาวาเส ตทหุ ปวารณาย สมฺพหุลา อาวาสิกา ภิกฺขู สนฺนิปตนฺติ ปญฺจ วา อติเรกา วา, เต น ชานนฺติ “อตฺถญฺเญ อาวาสิกา ภิกฺขู อนาคตา”ติ, เต ธมฺมสญฺญิโน วินยสญฺญิโน วคฺคา สมคฺคสญฺญิโน ปวาเรนฺติ, เตหิ ปวาริตมตฺเต เอกจฺจาย วุฏฺฐิตาย ปริสาย อถญฺเญ อาวาสิกา ภิกฺขู อาคจฺฉนฺติ พหุตรา, เตหิ ภิกฺขเว ภิกฺขูหิ ปุน ปวาเรตพฺพํ, ปวาริตานํ อนาปตฺติ. (10)}

\prose{อิธ ปน ภิกฺขเว อญฺญตรสฺมิํ อาวาเส ตทหุ ปวารณาย สมฺพหุลา อาวาสิกา ภิกฺขู สนฺนิปตนฺติ ปญฺจ วา อติเรกา วา, เต น ชานนฺติ “อตฺถญฺเญ อาวาสิกา ภิกฺขู อนาคตา”ติ, เต ธมฺมสญฺญิโน วินยสญฺญิโน วคฺคา สมคฺคสญฺญิโน ปวาเรนฺติ, เตหิ ปวาริตมตฺเต เอกจฺจาย วุฏฺฐิตาย ปริสาย อถญฺเญ อาวาสิกา ภิกฺขู อาคจฺฉนฺติ สมสมา, ปวาริตา สุปฺปวาริตา, เตสํ สนฺติเก ปวาเรตพฺพํ, ปวาริตานํ อนาปตฺติ. (11)}

\prose{\csromanfolio{236}อิธ ปน ภิกฺขเว อญฺญตรสฺมิํ อาวาเส ตทหุ ปวารณาย สมฺพหุลา อาวาสิกา ภิกฺขู สนฺนิปตนฺติ ปญฺจ วา อติเรกา วา, เต น ชานนฺติ “อตฺถญฺเญ อาวาสิกา ภิกฺขู อนาคตา”ติ, เต ธมฺมสญฺญิโน วินยสญฺญิโน วคฺคา สมคฺคสญฺญิโน ปวาเรนฺติ, เตหิ ปวาริตมตฺเต เอกจฺจาย วุฏฺฐิตาย ปริสาย อถญฺเญ อาวาสิกา ภิกฺขู อาคจฺฉนฺติ โถกตรา, ปวาริตา สุปฺปวาริตา, เตสํ สนฺติเก ปวาเรตพฺพํ, ปวาริตานํ อนาปตฺติ. (12)}

\prose{อิธ ปน ภิกฺขเว อญฺญตรสฺมิํ อาวาเส ตทหุ ปวารณาย สมฺพหุลา อาวาสิกา ภิกฺขู สนฺนิปตนฺติ ปญฺจ วา อติเรกา วา, เต น ชานนฺติ “อตฺถญฺเญ อาวาสิกา ภิกฺขู อนาคตา”ติ, เต ธมฺมสญฺญิโน วินยสญฺญิโน วคฺคา สมคฺคสญฺญิโน ปวาเรนฺติ, เตหิ ปวาริตมตฺเต สพฺพาย วุฏฺฐิตาย ปริสาย อถญฺเญ อาวาสิกา ภิกฺขู อาคจฺฉนฺติ พหุตรา, เตหิ ภิกฺขเว ภิกฺขูหิ ปุน ปวาเรตพฺพํ, ปวาริตานํ อนาปตฺติ. (13)}

\prose{อิธ ปน ภิกฺขเว อญฺญตรสฺมิํ อาวาเส ตทหุ ปวารณาย สมฺพหุลา อาวาสิกา ภิกฺขู สนฺนิปตนฺติ ปญฺจ วา อติเรกา วา, เต น ชานนฺติ “อตฺถญฺเญ อาวาสิกา ภิกฺขู อนาคตา”ติ, เต ธมฺมสญฺญิโน วินยสญฺญิโน วคฺคา สมคฺคสญฺญิโน ปวาเรนฺติ, เตหิ ปวาริตมตฺเต สพฺพาย วุฏฺฐิตาย ปริสาย อถญฺเญ อาวาสิกา ภิกฺขู อาคจฺฉนฺติ สมสมา, ปวาริตา สุปฺปวาริตา, เตสํ สนฺติเก ปวาเรตพฺพํ, ปวาริตานํ อนาปตฺติ. (14)}

\prosewithcloser{อิธ ปน ภิกฺขเว อญฺญตรสฺมิํ อาวาเส ตทหุ ปวารณาย สมฺพหุลา อาวาสิกา ภิกฺขู สนฺนิปตนฺติ ปญฺจ วา อติเรกา วา, เต น ชานนฺติ “อตฺถญฺเญ อาวาสิกา ภิกฺขู อนาคตา”ติ, เต ธมฺมสญฺญิโน วินยสญฺญิโน วคฺคา สมคฺคสญฺญิโน ปวาเรนฺติ, เตหิ ปวาริตมตฺเต สพฺพาย วุฏฺฐิตาย ปริสาย อถญฺเญ อาวาสิกา ภิกฺขู อาคจฺฉนฺติ โถกตรา, ปวาริตา สุปฺปวาริตา, เตสํ สนฺติเก ปวาเรตพฺพํ, ปวาริตานํ อนาปตฺติ. (15)}{อนาปตฺติ\-ปนฺนรสกํ นิฏฺฐิตํ.}
\csromanmatikaanchor{mtk.133}
\csromantocmarkref{section}{129. วคฺคาวคฺคสญฺญีปนฺนรสก}{mtk.133}
\bookmark[dest={mtk.133},level=1]{129. วคฺคาวคฺคสญฺญีปนฺนรสก}
\csromanheader{1}{129. \textbf{วคฺคา}\-\textbf{วคฺคสญฺญี}\-\textbf{ปนฺนรสก}}
\proseitem{223}{\csromanfolio{237}อิธ ปน ภิกฺขเว อญฺญตรสฺมิํ อาวาเส ตทหุ ปวารณาย สมฺพหุลา อาวาสิกา ภิกฺขู สนฺนิปตนฺติ ปญฺจ วา อติเรกา วา, เต ชานนฺติ “อตฺถญฺเญ อาวาสิกา ภิกฺขู อนาคตา”ติ, เต ธมฺมสญฺญิโน วินยสญฺญิโน วคฺคา วคฺคสญฺญิโน ปวาเรนฺติ, เตหิ ปวาริยมาเน อถญฺเญ อาวาสิกา ภิกฺขู อาคจฺฉนฺติ พหุตรา, เตหิ ภิกฺขเว ภิกฺขูหิ ปุน ปวาเรตพฺพํ, ปวาริตานํ อาปตฺติ ทุกฺกฏสฺส. (1)}

\prose{อิธ ปน ภิกฺขเว อญฺญตรสฺมิํ อาวาเส ตทหุ ปวารณาย สมฺพหุลา อาวาสิกา ภิกฺขู สนฺนิปตนฺติ ปญฺจ วา อติเรกา วา, เต ชานนฺติ “อตฺถญฺเญ อาวาสิกา ภิกฺขู อนาคตา”ติ, เต ธมฺมสญฺญิโน วินยสญฺญิโน วคฺคา วคฺคสญฺญิโน ปวาเรนฺติ, เตหิ ปวาริยมาเน อถญฺเญ อาวาสิกา ภิกฺขู อาคจฺฉนฺติ สมสมา, ปวาริตา สุปฺปวาริตา, อวเสเสหิ ปวาเรตพฺพํ, ปวาริตานํ อาปตฺติ ทุกฺกฏสฺส. (2)}

\prose{อิธ ปน ภิกฺขเว อญฺญตรสฺมิํ อาวาเส ตทหุ ปวารณาย สมฺพหุลา อาวาสิกา ภิกฺขู สนฺนิปตนฺติ ปญฺจ วา อติเรกา วา, เต ชานนฺติ “อตฺถญฺเญ อาวาสิกา ภิกฺขู อนาคตา”ติ, เต ธมฺมสญฺญิโน วินยสญฺญิโน วคฺคา วคฺคสญฺญิโน ปวาเรนฺติ, เตหิ ปวาริยมาเน อถญฺเญ อาวาสิกา ภิกฺขู อาคจฺฉนฺติ โถกตรา, ปวาริตา สุปฺปวาริตา, อวเสเสหิ ปวาเรตพฺพํ, ปวาริตานํ อาปตฺติ ทุกฺกฏสฺส. (3)}

\prosewithcloser{อิธ ปน ภิกฺขเว อญฺญตรสฺมิํ อาวาเส ตทหุ ปวารณาย สมฺพหุลา อาวาสิกา ภิกฺขู สนฺนิปตนฺติ ปญฺจ วา อติเรกา วา, เต ชานนฺติ “อตฺถญฺเญ อาวาสิกา ภิกฺขู อนาคตา”ติ, เต ธมฺมสญฺญิโน วินยสญฺญิโน วคฺคา วคฺคสญฺญิโน ปวาเรนฺติ, เตหิ ปวาริตมตฺเต ฯเปฯ อวุฏฺฐิตาย ปริสาย ฯเปฯ เอกจฺจาย วุฏฺฐิตาย ปริสาย ฯเปฯ สพฺพาย วุฏฺฐิตาย ปริสาย อถญฺเญ อาวาสิกา ภิกฺขู อาคจฺฉนฺติ พหุตรา ฯเปฯ สมสมา ฯเปฯ โถกตรา, ปวาริตา สุปฺปวาริตา, เตสํ สนฺติเก ปวาเรตพฺพํ, ปวาริตานํ อาปตฺติ ทุกฺกฏสฺส. \mbox{(4-15)}}{วคฺคา\-วคฺค\-สญฺญิ\-ปนฺนรสกํ นิฏฺฐิตํ.}
\csromanmatikaanchor{mtk.134}
\csromantocmarkref{section}{130. เวมติกปนฺนรสก}{mtk.134}
\bookmark[dest={mtk.134},level=1]{130. เวมติกปนฺนรสก}
\csromanheader{1}{130. \textbf{เวมติก}\-\textbf{ปนฺนรสก}}
\proseitem{224}{\csromanfolio{238}อิธ ปน ภิกฺขเว อญฺญตรสฺมิํ อาวาเส ตทหุ ปวารณาย สมฺพหุลา อาวาสิกา ภิกฺขู สนฺนิปตนฺติ ปญฺจ วา อติเรกา วา, เต ชานนฺติ “อตฺถญฺเญ อาวาสิกา ภิกฺขู อนาคตา”ติ, เต “กปฺปติ นุ โข อมฺหากํ ปวาเรตุํ, น นุ โข กปฺปตี”ติ เวมติกา ปวาเรนฺติ, เตหิ ปวาริยมาเน อถญฺเญ อาวาสิกา ภิกฺขู อาคจฺฉนฺติ พหุตรา, เตหิ ภิกฺขเว ภิกฺขูหิ ปุน ปวาเรตพฺพํ, ปวาริตานํ อาปตฺติ ทุกฺกฏสฺส. (1)}

\prose{อิธ ปน ภิกฺขเว อญฺญตรสฺมิํ อาวาเส ตทหุ ปวารณาย สมฺพหุลา อาวาสิกา ภิกฺขู สนฺนิปตนฺติ ปญฺจ วา อติเรกา วา, เต ชานนฺติ “อตฺถญฺเญ อาวาสิกา ภิกฺขู อนาคตา”ติ, เต “กปฺปติ นุ โข อมฺหากํ ปวาเรตุํ, น นุ โข กปฺปตี”ติ เวมติกา ปวาเรนฺติ, เตหิ ปวาริยมาเน อถญฺเญ อาวาสิกา ภิกฺขู อาคจฺฉนฺติ สมสมา, ปวาริตา สุปฺปวาริตา, อวเสเสหิ ปวาเรตพฺพํ, ปวาริตานํ อาปตฺติ ทุกฺกฏสฺส. (2)}

\prose{อิธ ปน ภิกฺขเว อญฺญตรสฺมิํ อาวาเส ตทหุ ปวารณาย สมฺพหุลา อาวาสิกา ภิกฺขู สนฺนิปตนฺติ ปญฺจ วา อติเรกา วา, เต ชานนฺติ “อตฺถญฺเญ อาวาสิกา ภิกฺขู อนาคตา”ติ, เต “กปฺปติ นุ โข อมฺหากํ ปวาเรตุํ, น นุ โข กปฺปตี”ติ เวมติกา ปวาเรนฺติ, เตหิ ปวาริยมาเน อถญฺเญ อาวาสิกา ภิกฺขู อาคจฺฉนฺติ, โถกตรา, ปวาริตา สุปฺปวาริตา, อวเสเสหิ ปวาเรตพฺพํ, ปวาริตานํ อาปตฺติ ทุกฺกฏสฺส. (3)}

\prosewithcloserruled{อิธ ปน ภิกฺขเว อญฺญตรสฺมิํ อาวาเส ตทหุ ปวารณาย สมฺพหุลา อาวาสิกา ภิกฺขู สนฺนิปตนฺติ ปญฺจ วา อติเรกา วา, เต ชานนฺติ “อตฺถญฺเญ อาวาสิกา ภิกฺขู อนาคตา”ติ, เต “กปฺปติ นุ โข อมฺหากํ ปวาเรตุํ, น นุ โข กปฺปตี”ติ เวมติกา ปวาเรนฺติ, เตหิ ปวาริตมตฺเต ฯเปฯ อวุฏฺฐิตาย ปริสาย ฯเปฯ เอกจฺจาย วุฏฺฐิตาย ปริสาย ฯเปฯ สพฺพาย วุฏฺฐิตาย ปริสาย อถญฺเญ \csromanfolio{239}อาวาสิกา ภิกฺขู อาคจฺฉนฺติ พหุตรา ฯเปฯ สมสมา ฯเปฯ โถกตรา, ปวาริตา สุปฺปวาริตา, เตสํ สนฺติเก ปวาเรตพฺพํ, ปวาริตานํ อาปตฺติ ทุกฺกฏสฺส. \mbox{(4-15)}}{วิมติกปนฺนรสกํ นิฏฺฐิตํ.}
\csromanmatikaanchor{mtk.135}
\csromantocmarkref{section}{131. กุกฺกุจฺจปกตปนฺนรสก}{mtk.135}
\bookmark[dest={mtk.135},level=1]{131. กุกฺกุจฺจปกตปนฺนรสก}
\csromanheader{1}{131. \textbf{กุกฺกุจฺจ}\-\textbf{ปกต}\-\textbf{ปนฺนรสก}}
\proseitem{225}{อิธ ปน ภิกฺขเว อญฺญตรสฺมิํ อาวาเส ตทหุ ปวารณาย สมฺพหุลา อาวาสิกา ภิกฺขู สนฺนิปตนฺติ ปญฺจ วา อติเรกา วา, เต ชานนฺติ “อตฺถญฺเญ อาวาสิกา ภิกฺขู อนาคตา”ติ, เต “กปฺปเตว อมฺหากํ ปวาเรตุํ, นามฺหากํ น กปฺปตี”ติ กุกฺกุจฺจปกตา ปวาเรนฺติ, เตหิ ปวาริยมาเน อถญฺเญ อาวาสิกา ภิกฺขู อาคจฺฉนฺติ พหุตรา, เตหิ ภิกฺขเว ภิกฺขูหิ ปุน ปวาเรตพฺพํ, ปวาริตานํ อาปตฺติ ทุกฺกฏสฺส. (1)}

\prose{อิธ ปน ภิกฺขเว อญฺญตรสฺมิํ อาวาเส ตทหุ ปวารณาย สมฺพหุลา อาวาสิกา ภิกฺขู สนฺนิปตนฺติ ปญฺจ วา อติเรกา วา, เต ชานนฺติ “อตฺถญฺเญ อาวาสิกา ภิกฺขู อนาคตา”ติ, เต “กปฺปเตว อมฺหากํ ปวาเรตุํ, นามฺหากํ น กปฺปตี”ติ กุกฺกุจฺจปกตา ปวาเรนฺติ, เตหิ ปวาริยมาเน อถญฺเญ อาวาสิกา ภิกฺขู อาคจฺฉนฺติ สมสมา, ปวาริตา สุปฺปวาริตา, อวเสเสหิ ปวาเรตพฺพํ, ปวาริตานํ อาปตฺติ ทุกฺกฏสฺส. (2)}

\prose{อิธ ปน ภิกฺขเว อญฺญตรสฺมิํ อาวาเส ตทหุ ปวารณาย สมฺพหุลา อาวาสิกา ภิกฺขู สนฺนิปตนฺติ ปญฺจ วา อติเรกา วา, เต ชานนฺติ “อตฺถญฺเญ อาวาสิกา ภิกฺขู อนาคตา”ติ, เต “กปฺปเตว อมฺหากํ ปวาเรตุํ, นามฺหากํ น กปฺปตี”ติ กุกฺกุจฺจปกตา ปวาเรนฺติ, เตหิ ปวาริยมาเน อถญฺเญ อาวาสิกา ภิกฺขู อาคจฺฉนฺติ โถกตรา, ปวาริตา สุปฺปวาริตา, อวเสเสหิ\csromansharedfootnote{fe37d12218f286e4}{อวเสเสหิ เตสํ สนฺติเก (ก)} ปวาเรตพฺพํ, ปวาริตานํ อาปตฺติ ทุกฺกฏสฺส. (3)}

\prosewithcloserruled{\csromanfolio{240}อิธ ปน ภิกฺขเว อญฺญตรสฺมิํ อาวาเส ตทหุ ปวารณาย สมฺพหุลา อาวาสิกา ภิกฺขู สนฺนิปตนฺติ ปญฺจ วา อติเรกา วา, เต ชานนฺติ “อตฺถญฺเญ อาวาสิกา ภิกฺขู อนาคตา”ติ, เต “กปฺปเตว อมฺหากํ ปวาเรตุํ, นามฺหากํ น กปฺปตี”ติ กุกฺกุจฺจปกตา ปวาเรนฺติ, เตหิ ปวาริตมตฺเต ฯเปฯ อวุฏฺฐิตาย ปริสาย ฯเปฯ เอกจฺจาย วุฏฺฐิตาย ปริสาย ฯเปฯ สพฺพาย วุฏฺฐิตาย ปริสาย อถญฺเญ อาวาสิกา ภิกฺขู อาคจฺฉนฺติ พหุตรา ฯเปฯ สมสมา ฯเปฯ โถกตรา, ปวาริตา สุปฺปวาริตา, เตสํ สนฺติเก ปวาเรตพฺพํ, ปวาริตานํ อาปตฺติ ทุกฺกฏสฺส. \mbox{(4-15)}}{กุกฺกุจฺจ\-ปกต\-ปนฺนรสกํ นิฏฺฐิตํ.}
\csromanmatikaanchor{mtk.136}
\csromantocmarkref{section}{132. เภทปุเรกฺขารปนฺนรสก}{mtk.136}
\bookmark[dest={mtk.136},level=1]{132. เภทปุเรกฺขารปนฺนรสก}
\csromanheader{1}{132. \textbf{เภท}\-\textbf{ปุเรกฺขาร}\-\textbf{ปนฺนรสก}}
\proseitem{226}{อิธ ปน ภิกฺขเว อญฺญตรสฺมิํ อาวาเส ตทหุ ปวารณาย สมฺพหุลา อาวาสิกา ภิกฺขู สนฺนิปตนฺติ ปญฺจ วา อติเรกา วา, เต ชานนฺติ “อตฺถญฺเญ อาวาสิกา ภิกฺขู อนาคตา”ติ, เต “นสฺสนฺเตเต, วินสฺสนฺเตเต, โก เตหิ อตฺโถ”ติ เภท\-ปุเรกฺขารา ปวาเรนฺติ, เตหิ ปวาริยมาเน อถญฺเญ อาวาสิกา ภิกฺขู อาคจฺฉนฺติ พหุตรา, เตหิ ภิกฺขเว ภิกฺขูหิ ปุน ปวาเรตพฺพํ, ปวาริตานํ อาปตฺติ ถุลฺลจฺจยสฺส. (1)}

\prose{อิธ ปน ภิกฺขเว อญฺญตรสฺมิํ อาวาเส ตทหุ ปวารณาย สมฺพหุลา อาวาสิกา ภิกฺขู สนฺนิปตนฺติ ปญฺจ วา อติเรกา วา, เต ชานนฺติ “อตฺถญฺเญ อาวาสิกา ภิกฺขู อนาคตา”ติ, เต “นสฺสนฺเตเต, วินสฺสนฺเตเต, โก เตหิ อตฺโถ”ติ เภท\-ปุเรกฺขารา ปวาเรนฺติ, เตหิ ปวาริยมาเน อถญฺเญ อาวาสิกา ภิกฺขู อาคจฺฉนฺติ สมสมา, ปวาริตา สุปฺปวาริตา, อวเสเสหิ ปวาเรตพฺพํ, ปวาริตานํ อาปตฺติ ถุลฺลจฺจยสฺส. (2)}

\prose{อิธ ปน ภิกฺขเว อญฺญตรสฺมิํ อาวาเส ตทหุ ปวารณาย สมฺพหุลา อาวาสิกา ภิกฺขู สนฺนิปตนฺติ ปญฺจ วา อติเรกา วา, เต ชานนฺติ “อตฺถญฺเญ อาวาสิกา ภิกฺขู อนาคตา”ติ, เต “นสฺสนฺเตเต, \csromanfolio{241}วินสฺสนฺเตเต, โก เตหิ อตฺโถ”ติ เภท\-ปุเรกฺขารา ปวาเรนฺติ, เตหิ ปวาริยมาเน อถญฺเญ อาวาสิกา ภิกฺขู อาคจฺฉนฺติ โถกตรา, ปวาริตา สุปฺปวาริตา, อวเสเสหิ ปวาเรตพฺพํ, ปวาริตานํ อาปตฺติ ถุลฺลจฺจยสฺส. (3)}

\prose{อิธ ปน ภิกฺขเว อญฺญตรสฺมิํ อาวาเส ตทหุ ปวารณาย สมฺพหุลา อาวาสิกา ภิกฺขู สนฺนิปตนฺติ ปญฺจ วา อติเรกา วา, เต ชานนฺติ “อตฺถญฺเญ อาวาสิกา ภิกฺขู อนาคตา”ติ, เต “นสฺสนฺเตเต, วินสฺสนฺเตเต, โก เตหิ อตฺโถ”ติ เภท\-ปุเรกฺขารา ปวาเรนฺติ, เตหิ ปวาริตมตฺเต อถญฺเญ อาวาสิกา ภิกฺขู อาคจฺฉนฺติ พหุตรา, เตหิ ภิกฺขเว ภิกฺขูหิ ปุน ปวาเรตพฺพํ, ปวาริตานํ อาปตฺติ ถุลฺลจฺจยสฺส. (4)}

\prose{อิธ ปน ภิกฺขเว อญฺญตรสฺมิํ อาวาเส ตทหุ ปวารณาย สมฺพหุลา อาวาสิกา ภิกฺขู สนฺนิปตนฺติ ปญฺจ วา อติเรกา วา, เต ชานนฺติ “อตฺถญฺเญ อาวาสิกา ภิกฺขู อนาคตา”ติ, เต “นสฺสนฺเตเต, วินสฺสนฺเตเต, โก เตหิ อตฺโถ”ติ เภท\-ปุเรกฺขารา ปวาเรนฺติ, เตหิ ปวาริตมตฺเต อถญฺเญ อาวาสิกา ภิกฺขู อาคจฺฉนฺติ สมสมา, ปวาริตา สุปฺปวาริตา, เตสํ สนฺติเก ปวาเรตพฺพํ, ปวาริตานํ อาปตฺติ ถุลฺลจฺจยสฺส. (5)}

\prose{อิธ ปน ภิกฺขเว อญฺญตรสฺมิํ อาวาเส ตทหุ ปวารณาย สมฺพหุลา อาวาสิกา ภิกฺขู สนฺนิปตนฺติ ปญฺจ วา อติเรกา วา, เต ชานนฺติ “อตฺถญฺเญ อาวาสิกา ภิกฺขู อนาคตา”ติ, เต “นสฺสนฺเตเต, วินสฺสนฺเตเต, โก เตหิ อตฺโถ”ติ เภท\-ปุเรกฺขารา ปวาเรนฺติ, เตหิ ปวาริตมตฺเต อถญฺเญ อาวาสิกา ภิกฺขู อาคจฺฉนฺติ โถกตรา, ปวาริตา สุปฺปวาริตา, เตสํ สนฺติเก ปวาเรตพฺพํ, ปวาริตานํ อาปตฺติ ถุลฺลจฺจยสฺส. (6)}

\prose{อิธ ปน ภิกฺขเว อญฺญตรสฺมิํ อาวาเส ตทหุ ปวารณาย สมฺพหุลา อาวาสิกา ภิกฺขู สนฺนิปตนฺติ ปญฺจ วา อติเรกา วา, เต ชานนฺติ “อตฺถญฺเญ อาวาสิกา ภิกฺขู อนาคตา”ติ, “นสฺสนฺเตเต, วินสฺสนฺเตเต, โก เตหิ อตฺโถ”ติ เภท\-ปุเรกฺขารา ปวาเรนฺติ, เตหิ ปวาริตมตฺเต อวุฏฺฐิตาย ปริสาย อถญฺเญ อาวาสิกา \csromanfolio{242}ภิกฺขู อาคจฺฉนฺติ พหุตรา, เตหิ ภิกฺขเว ภิกฺขูหิ ปุน ปวาเรตพฺพํ, ปวาริตานํ อาปตฺติ ถุลฺลจฺจยสฺส. (7)}

\prose{อิธ ปน ภิกฺขเว อญฺญตรสฺมิํ อาวาเส ตทหุ ปวารณาย สมฺพหุลา อาวาสิกา ภิกฺขู สนฺนิปตนฺติ ปญฺจ วา อติเรกา วา, เต ชานนฺติ “อตฺถญฺเญ อาวาสิกา ภิกฺขู อนาคตา”ติ, เต “นสฺสนฺเตเต, วินสฺสนฺเตเต, โก เตหิ อตฺโถ”ติ เภท\-ปุเรกฺขารา ปวาเรนฺติ, เตหิ ปวาริตมตฺเต อวุฏฺฐิตาย ปริสาย อถญฺเญ อาวาสิกา ภิกฺขู อาคจฺฉนฺติ สมสมา, ปวาริตา สุปฺปวาริตา, เตสํ สนฺติเก ปวาเรตพฺพํ, ปวาริตานํ อาปตฺติ ถุลฺลจฺจยสฺส. (8)}

\prose{อิธ ปน ภิกฺขเว อญฺญตรสฺมิํ อาวาเส ตทหุ ปวารณาย สมฺพหุลา อาวาสิกา ภิกฺขู สนฺนิปตนฺติ ปญฺจ วา อติเรกา วา, เต ชานนฺติ “อตฺถญฺเญ อาวาสิกา ภิกฺขู อนาคตา”ติ, เต “นสฺสนฺเตเต, วินสฺสนฺเตเต, โก เตหิ อตฺโถ”ติ เภท\-ปุเรกฺขารา ปวาเรนฺติ, เตหิ ปวาริตมตฺเต อวุฏฺฐิตาย ปริสาย อถญฺเญ อาวาสิกา ภิกฺขู อาคจฺฉนฺติ โถกตรา, ปวาริตา สุปฺปวาริตา, เตสํ สนฺติเก ปวาเรตพฺพํ, ปวาริตานํ อาปตฺติ ถุลฺลจฺจยสฺส. (9)}

\prose{อิธ ปน ภิกฺขเว อญฺญตรสฺมิํ อาวาเส ตทหุ ปวารณาย สมฺพหุลา อาวาสิกา ภิกฺขู สนฺนิปตนฺติ ปญฺจ วา อติเรกา วา, เต ชานนฺติ “อตฺถญฺเญ อาวาสิกา ภิกฺขู อนาคตา”ติ, เต “นสฺสนฺเตเต, วินสฺสนฺเตเต, โก เตหิ อตฺโถ”ติ เภท\-ปุเรกฺขารา ปวาเรนฺติ, เตหิ ปวาริตมตฺเต เอกจฺจาย วุฏฺฐิตาย ปริสาย อถญฺเญ อาวาสิกา ภิกฺขู อาคจฺฉนฺติ พหุตรา, เตหิ ภิกฺขเว ภิกฺขูหิ ปุน ปวาเรตพฺพํ, ปวาริตานํ อาปตฺติ ถุลฺลจฺจยสฺส. (10)}

\prose{อิธ ปน ภิกฺขเว อญฺญตรสฺมิํ อาวาเส ตทหุ ปวารณาย สมฺพหุลา อาวาสิกา ภิกฺขู สนฺนิปตนฺติ ปญฺจ วา อติเรกา วา, เต ชานนฺติ “อตฺถญฺเญ อาวาสิกา ภิกฺขู อนาคตา”ติ, เต “นสฺสนฺเตเต, วินสฺสนฺเตเต, โก เตหิ อตฺโถ”ติ เภท\-ปุเรกฺขารา ปวาเรนฺติ, เตหิ ปวาริตมตฺเต เอกจฺจาย วุฏฺฐิตาย ปริสาย อถญฺเญ อาวาสิกา ภิกฺขู อาคจฺฉนฺติ สมสมา, ปวาริตา สุปฺปวาริตา, เตสํ สนฺติเก ปวาเรตพฺพํ, ปวาริตานํ อาปตฺติ ถุลฺลจฺจยสฺส. (11)}

\prose{\csromanfolio{243}อิธ ปน ภิกฺขเว อญฺญตรสฺมิํ อาวาเส ตทหุ ปวารณาย สมฺพหุลา อาวาสิกา ภิกฺขู สนฺนิปตนฺติ ปญฺจ วา อติเรกา วา, เต ชานนฺติ “อตฺถญฺเญ อาวาสิกา ภิกฺขู อนาคตา”ติ, เต “นสฺสนฺเตเต, วินสฺสนฺเตเต, โก เตหิ อตฺโถ”ติ เภท\-ปุเรกฺขารา ปวาเรนฺติ, เตหิ ปวาริตมตฺเต เอกจฺจาย วุฏฺฐิตาย ปริสาย อถญฺเญ อาวาสิกา ภิกฺขู อาคจฺฉนฺติ โถกตรา, ปวาริตา สุปฺปวาริตา, เตสํ สนฺติเก ปวาเรตพฺพํ, ปวาริตานํ อาปตฺติ ถุลฺลจฺจยสฺส. (12)}

\prose{อิธ ปน ภิกฺขเว อญฺญตรสฺมิํ อาวาเส ตทหุ ปวารณาย สมฺพหุลา อาวาสิกา ภิกฺขู สนฺนิปตนฺติ ปญฺจ วา อติเรกา วา, เต ชานนฺติ “อตฺถญฺเญ อาวาสิกา ภิกฺขู อนาคตา”ติ, เต “นสฺสนฺเตเต, วินสฺสนฺเตเต, โก เตหิ อตฺโถ”ติ เภท\-ปุเรกฺขารา ปวาเรนฺติ, เตหิ ปวาริตมตฺเต สพฺพาย วุฏฺฐิตาย ปริสาย อถญฺเญ อาวาสิกา ภิกฺขู อาคจฺฉนฺติ พหุตรา, เตหิ ภิกฺขเว ภิกฺขูหิ ปุน ปวาเรตพฺพํ, ปวาริตานํ อาปตฺติ ถุลฺลจฺจยสฺส. (13)}

\prose{อิธ ปน ภิกฺขเว อญฺญตรสฺมิํ อาวาเส ตทหุ ปวารณาย สมฺพหุลา อาวาสิกา ภิกฺขู สนฺนิปตนฺติ ปญฺจ วา อติเรกา วา, เต ชานนฺติ “อตฺถญฺเญ อาวาสิกา ภิกฺขู อนาคตา”ติ, เต “นสฺสนฺเตเต, วินสฺสนฺเตเต, โก เตหิ อตฺโถ”ติ เภท\-ปุเรกฺขารา ปวาเรนฺติ, เตหิ ปวาริตมตฺเต สพฺพาย วุฏฺฐิตาย ปริสาย อถญฺเญ อาวาสิกา ภิกฺขู อาคจฺฉนฺติ สมสมา, ปวาริตา สุปฺปวาริตา, เตสํ สนฺติเก ปวาเรตพฺพํ, ปวาริตานํ อาปตฺติ ถุลฺลจฺจยสฺส. (14)}

\prosewithcloser{อิธ ปน ภิกฺขเว อญฺญตรสฺมิํ อาวาเส ตทหุ ปวารณาย สมฺพหุลา อาวาสิกา ภิกฺขู สนฺนิปตนฺติ ปญฺจ วา อติเรกา วา, เต ชานนฺติ “อตฺถญฺเญ อาวาสิกา ภิกฺขู อนาคตา”ติ, เต “นสฺสนฺเตเต, วินสฺสนฺเตเต, โก เตหิ อตฺโถ”ติ เภท\-ปุเรกฺขารา ปวาเรนฺติ, เตหิ ปวาริตมตฺเต สพฺพาย วุฏฺฐิตาย ปริสาย อถญฺเญ อาวาสิกา ภิกฺขู อาคจฺฉนฺติ โถกตรา, ปวาริตา สุปฺปวาริตา, เตสํ สนฺติเก ปวาเรตพฺพํ, ปวาริตานํ อาปตฺติ ถุลฺลจฺจยสฺส. (15)}{เภท\-ปุเรกฺขาร\-ปนฺนรสกํ นิฏฺฐิตํ.}
\nitthitam{ปญฺจวีสตฺติกา นิฏฺฐิตา.}
\csromanmatikaanchor{mtk.137}
\csromantocmarkref{section}{133. สีโมกฺกนฺติกเปยฺยาล}{mtk.137}
\bookmark[dest={mtk.137},level=1]{133. สีโมกฺกนฺติกเปยฺยาล}
\csromanheader{1}{133. \textbf{สีโม}\-\textbf{กฺกนฺติก}\-\textbf{เปยฺยาล}}
\csromanmatikaanchor{mtk.138}
\csromantocmarkref{section}{134. ทิวสนานตฺถ}{mtk.138}
\bookmark[dest={mtk.138},level=1]{134. ทิวสนานตฺถ}
\proseitem{227}{\csromanfolio{244}อิธ ปน ภิกฺขเว อญฺญตรสฺมิํ อาวาเส ตทหุ ปวารณาย สมฺพหุลา อาวาสิกา ภิกฺขู สนฺนิปตนฺติ ปญฺจ วา อติเรกา วา, เต น ชานนฺติ “อญฺเญ อาวาสิกา ภิกฺขู อนฺโตสีมํ โอกฺกมนฺตี”ติ ฯเปฯ เต น ชานนฺติ “อญฺเญ อาวาสิกา ภิกฺขู อนฺโตสีมํ โอกฺกนฺตา”ติ ฯเปฯ เต น ปสฺสนฺติ อญฺเญ อาวาสิเก ภิกฺขู อนฺโตสีมํ โอกฺกมนฺเต ฯเปฯ เต น ปสฺสนฺติ อญฺเญ อาวาสิเก ภิกฺขู อนฺโตสีมํ โอกฺกนฺเต ฯเปฯ เต น สุณนฺติ “อญฺเญ อาวาสิกา ภิกฺขู อนฺโตสีมํ โอกฺกมนฺตี”ติ ฯเปฯ เต น สุณนฺติ “อญฺเญ อาวาสิกา ภิกฺขู อนฺโตสีมํ โอกฺกนฺตา”ติ ฯเปฯ.}

\prose{อาวาสิเกน อาวาสิกา เอกสต\-ปญฺจสตฺตติ ติกนยโต, อาวาสิเกน อาคนฺตุกา, อาคนฺตุเกน อาวาสิกา, อาคนฺตุเกน อาคนฺตุกา, เปยฺยาลมุเขน สตฺต ติกสตานิ โหนฺติ.}

\csromanheader{2}{134. \textbf{ทิวสนานตฺต}}
\proseitem{228}{อิธ ปน ภิกฺขเว อาวาสิกานํ ภิกฺขูนํ จาตุทฺทโส โหติ อาคนฺตุกานํ ปนฺนรโส, สเจ อาวาสิกา พหุตรา โหนฺติ, อาคนฺตุเกหิ อาวาสิกานํ อนุวตฺติตพฺพํ.}

\setlength{\csromangathapagewidth}{0pt}
\setlength{\csromangathawakwidth}{0pt}
\csromangathameasuregroup{}{สเจ สมสมา โหนฺติ, \\ อาคนฺตุเกหิ อาวาสิกานํ อนุวตฺติตพฺพํ. \\ สเจ อาคนฺตุกา พหุตรา โหนฺติ, \\ อาวาสิเกหิ อาคนฺตุกานํ อนุวตฺติตพฺพํ.}
\setlength{\csromangathaleftcol}{0pt}
\csromangathagroup{\csromangathastanza{\csromangathaitemline{228}{สเจ สมสมา โหนฺติ,} \\ \csromangathacontline{อาคนฺตุเกหิ อาวาสิกานํ อนุวตฺติตพฺพํ.} \\ \csromangathacontline{สเจ อาคนฺตุกา พหุตรา โหนฺติ,} \\ \csromangathacontline{อาวาสิเกหิ อาคนฺตุกานํ อนุวตฺติตพฺพํ.}}}

\prose{อิธ ปน ภิกฺขเว อาวาสิกานํ ภิกฺขูนํ ปนฺนรโส โหติ อาคนฺตุกานํ จาตุทฺทโส, สเจ อาวาสิกา พหุตรา โหนฺติ, อาคนฺตุเกหิ อาวาสิกานํ อนุวตฺติตพฺพํ.}

\setlength{\csromangathapagewidth}{0pt}
\setlength{\csromangathawakwidth}{0pt}
\csromangathameasuregroup{}{สเจ สมสมา โหนฺติ, \\ อาคนฺตุเกหิ อาวาสิกานํ อนุวตฺติตพฺพํ. \\ สเจ อาคนฺตุกา พหุตรา โหนฺติ, \\ อาวาสิเกหิ อาคนฺตุกานํ อนุวตฺติตพฺพํ.}
\setlength{\csromangathaleftcol}{0pt}
\csromangathagroup{\csromangathastanza{สเจ สมสมา โหนฺติ, \\ อาคนฺตุเกหิ อาวาสิกานํ อนุวตฺติตพฺพํ. \\ สเจ อาคนฺตุกา พหุตรา โหนฺติ, \\ อาวาสิเกหิ อาคนฺตุกานํ อนุวตฺติตพฺพํ.}}

\proseitem{228}{อิธ ปน ภิกฺขเว อาวาสิกานํ ภิกฺขูนํ ปาฏิปโท โหติ อาคนฺตุกานํ ปนฺนรโส, สเจ อาวาสิกา พหุตรา โหนฺติ, อาวาสิเกหิ อาคนฺตุกานํ นากามา ทาตพฺพา สามคฺคี, อาคนฺตุเกหิ นิสฺสีมํ คนฺตฺวา ปวาเรตพฺพํ. สเจ สมสมา โหนฺติ, อาวาสิเกหิ อาคนฺตุกานํ นากามา ทาตพฺพา สามคฺคี, อาคนฺตุเกหิ นิสฺสีมํ คนฺตฺวา \csromanfolio{245}ปวาเรตพฺพํ. สเจ อาคนฺตุกา พหุตรา โหนฺติ, อาวาสิเกหิ อาคนฺตุกานํ สามคฺคี วา ทาตพฺพา, นิสฺสีมํ วา คนฺตพฺพํ.}

\prose{อิธ ปน ภิกฺขเว อาวาสิกานํ ภิกฺขูนํ ปนฺนรโส โหติ, อาคนฺตุกานํ ปาฏิปโท, สเจ อาวาสิกา พหุตรา โหนฺติ, อาคนฺตุเกหิ อาวาสิกานํ สามคฺคี วา ทาตพฺพา, นิสฺสีมํ วา คนฺตพฺพํ. สเจ สมสมา โหนฺติ, อาคนฺตุเกหิ อาวาสิกานํ สามคฺคี วา ทาตพฺพา, นิสฺสีมํ วา คนฺตพฺพํ. สเจ อาคนฺตุกา พหุตรา โหนฺติ, อาคนฺตุเกหิ อาวาสิกานํ นากามา ทาตพฺพา สามคฺคี, อาวาสิเกหิ นิสฺสีมํ คนฺตฺวา ปวาเรตพฺพํ.}

\csromanmatikaanchor{mtk.139}
\csromantocmarkref{section}{135. ลิงฺคาทิทสฺสน}{mtk.139}
\bookmark[dest={mtk.139},level=1]{135. ลิงฺคาทิทสฺสน}
\csromanheader{1}{135. \textbf{ลิงฺคาทิทสฺสน}}
\proseitem{229}{อิธ ปน ภิกฺขเว อาคนฺตุกา ภิกฺขู ปสฺสนฺติ อาวาสิกานํ ภิกฺขูนํ อาวาสิกาการํ อาวาสิกลิงฺคํ อาวาสิก\-นิมิตฺตํ อาวาสิกุทฺเทสํ สุปฺปญฺญตฺตํ มญฺจปีฐํ ภิสิพิพฺโพหนํ ปานียํ ปริโภชนียํ สูปฏฺฐิตํ ปริเวณํ สุสมฺมฏฺฐํ, ปสฺสิตฺวา เวมติกา โหนฺติ “อตฺถิ นุ โข อาวาสิกา ภิกฺขู, นตฺถิ นุ โข”ติ. เต เวมติกา น วิจินนฺติ, อวิจินิตฺวา ปวาเรนฺติ, อาปตฺติ ทุกฺกฏสฺส ฯเปฯ. เต เวมติกา วิจินนฺติ, วิจินิตฺวา น ปสฺสนฺติ, อปสฺสิตฺวา ปวาเรนฺติ, อนาปตฺติ. เต เวมติกา วิจินนฺติ, วิจินิตฺวา ปสฺสนฺติ, ปสฺสิตฺวา เอกโต ปวาเรนฺติ, อนาปตฺติ. เต เวมติกา วิจินนฺติ, วิจินิตฺวา ปสฺสนฺติ, ปสฺสิตฺวา ปาเฏกฺกํ ปวาเรนฺติ, อาปตฺติ ทุกฺกฏสฺส. เต เวมติกา วิจินนฺติ, วิจินิตฺวา ปสฺสนฺติ, ปสฺสิตฺวา “นสฺสนฺเตเต, วินสฺสนฺเตเต, โก เตหิ อตฺโถ”ติ เภท\-ปุเรกฺขารา ปวาเรนฺติ, อาปตฺติ ถุลฺลจฺจยสฺส.}

\prose{อิธ ปน ภิกฺขเว อาคนฺตุกา ภิกฺขู สุณนฺติ อาวาสิกานํ ภิกฺขูนํ อาวาสิกาการํ อาวาสิกลิงฺคํ อาวาสิก\-นิมิตฺตํ อาวาสิกุทฺเทสํ จงฺกมนฺตานํ ปทสทฺทํ สชฺฌายสทฺทํ อุกฺกาสิตสทฺทํ ขิปิตสทฺทํ, สุตฺวา เวมติกา โหนฺติ “อตฺถิ นุ โข อาวาสิกา ภิกฺขู, นตฺถิ นุ โข”ติ. เต เวมติกา น วิจินนฺติ, อวิจินิตฺวา ปวาเรนฺติ, อาปตฺติ ทุกฺกฏสฺส. เต เวมติกา วิจินนฺติ, วิจินิตฺวา น ปสฺสนฺติ, อปสฺสิตฺวา ปวาเรนฺติ, อนาปตฺติ. เต เวมติกา วิจินนฺติ, วิจินิตฺวา ปสฺสนฺติ, ปสฺสิตฺวา เอกโต ปวาเรนฺติ, อนาปตฺติ. เต เวมติกา วิจินนฺติ, วิจินิตฺวา ปสฺสนฺติ, \csromanfolio{246}ปสฺสิตฺวา ปาเฏกฺกํ ปวาเรนฺติ, อาปตฺติ ทุกฺกฏสฺส. เต เวมติกา วิจินนฺติ วิจินิตฺวา ปสฺสนฺติ, ปสฺสิตฺวา “นสฺสนฺเตเต, วินสฺสนฺเตเต, โก เตหิ อตฺโถ”ติ เภท\-ปุเรกฺขารา ปวาเรนฺติ, อาปตฺติ ถุลฺลจฺจยสฺส.}

\prose{อิธ ปน ภิกฺขเว อาวาสิกา ภิกฺขู ปสฺสนฺติ อาคนฺตุกานํ ภิกฺขูนํ อาคนฺตุกาการํ อาคนฺตุกลิงฺคํ อาคนฺตุก\-นิมิตฺตํ อาคนฺตุกุทฺเทสํ อญฺญาตกํ ปตฺตํ อญฺญาตกํ จีวรํ อญฺญาตกํ นิสีทนํ ปาทานํ โธตํ อุทกนิสฺเสกํ, ปสฺสิตฺวา เวมติกา โหนฺติ “อตฺถิ นุ โข อาคนฺตุกา ภิกฺขู, นตฺถิ นุ โข”ติ.}

\setlength{\csromangathapagewidth}{0pt}
\setlength{\csromangathawakwidth}{0pt}
\csromangathameasuregroup{เต เวมติกา น วิจินนฺติ,}{\csromangathabat{เต เวมติกา น วิจินนฺติ,}{อวิจินิตฺวา ปวาเรนฺติ, อาปตฺติ ทุกฺกฏสฺส.}}
\csromangathasetleft{เต เวมติกา น วิจินนฺติ,}
\csromangathagroup{\csromangathastanza{\csromangathaitembat{229}{เต เวมติกา น วิจินนฺติ,}{อวิจินิตฺวา ปวาเรนฺติ, อาปตฺติ ทุกฺกฏสฺส.}}}

\prosecont{เต เวมติกา วิจินนฺติ, วิจินิตฺวา น ปสฺสนฺติ, อปสฺสิตฺวา ปวาเรนฺติ, อนาปตฺติ. เต เวมติกา วิจินนฺติ, วิจินิตฺวา ปสฺสนฺติ, ปสฺสิตฺวา เอกโต ปวาเรนฺติ, อนาปตฺติ. เต เวมติกา วิจินนฺติ, วิจินิตฺวา ปสฺสนฺติ, ปสฺสิตฺวา ปาเฏกฺกํ ปวาเรนฺติ, อาปตฺติ ทุกฺกฏสฺส. เต เวมติกา วิจินนฺติ, วิจินิตฺวา ปสฺสนฺติ, ปสฺสิตฺวา “นสฺสนฺเตเต, วินสฺสนฺเตเต, โก เตหิ อตฺโถ”ติ เภท\-ปุเรกฺขารา ปวาเรนฺติ, อาปตฺติ ถุลฺลจฺจยสฺส.}

\prose{อิธ ปน ภิกฺขเว อาวาสิกา ภิกฺขู สุณนฺติ อาคนฺตุกานํ ภิกฺขูนํ อาคนฺตุกาการํ อาคนฺตุกลิงฺคํ อาคนฺตุก\-นิมิตฺตํ อาคนฺตุกุทฺเทสํ อาคจฺฉนฺตานํ ปทสทฺทํ อุปาหน\-ปปฺโผฏน\-สทฺทํ อุกฺกาสิตสทฺทํ ขิปิตสทฺทํ, สุตฺวา เวมติกา โหนฺติ “อตฺถิ นุ โข อาคนฺตุกา ภิกฺขู, นตฺถิ นุ โข”ติ.}

\setlength{\csromangathapagewidth}{0pt}
\setlength{\csromangathawakwidth}{0pt}
\csromangathameasuregroup{เต เวมติกา น วิจินนฺติ,}{\csromangathabat{เต เวมติกา น วิจินนฺติ,}{อวิจินิตฺวา ปวาเรนฺติ, อาปตฺติ ทุกฺกฏสฺส.}}
\csromangathasetleft{เต เวมติกา น วิจินนฺติ,}
\csromangathagroup{\csromangathastanza{\csromangathabat{เต เวมติกา น วิจินนฺติ,}{อวิจินิตฺวา ปวาเรนฺติ, อาปตฺติ ทุกฺกฏสฺส.}}}

\prose{เต เวมติกา วิจินนฺติ, วิจินิตฺวา น ปสฺสนฺติ, อปสฺสิตฺวา ปวาเรนฺติ, อนาปตฺติ.}

\prose{เต เวมติกา วิจินนฺติ, วิจินิตฺวา ปสฺสนฺติ, ปสฺสิตฺวา เอกโต ปวาเรนฺติ, อนาปตฺติ.}

\prose{เต เวมติกา วิจินนฺติ, วิจินิตฺวา ปสฺสนฺติ, ปสฺสิตฺวา ปาเฏกฺกํ ปวาเรนฺติ, อาปตฺติ ทุกฺกฏสฺส.}

\prose{เต เวมติกา วิจินนฺติ, วิจินิตฺวา ปสฺสนฺติ, ปสฺสิตฺวา “นสฺสนฺเตเต, วินสฺสนฺเตเต, โก เตหิ อตฺโถ”ติ เภท\-ปุเรกฺขารา ปวาเรนฺติ, อาปตฺติ ถุลฺลจฺจยสฺส.}

\csromanmatikaanchor{mtk.140}
\csromantocmarkref{section}{136. นานาสํวาสกาทีหิ ปวารณา}{mtk.140}
\bookmark[dest={mtk.140},level=1]{136. นานาสํวาสกาทีหิ ปวารณา}
\csromanheader{1}{136. \textbf{นานาสํวาสกาทีหิ ปวารณา}}
\csromanmatikaanchor{mtk.141}
\csromantocmarkref{section}{137. นคนฺตพฺพวาร}{mtk.141}
\bookmark[dest={mtk.141},level=1]{137. นคนฺตพฺพวาร}
\proseitem{230}{อิธ ปน ภิกฺขเว อาคนฺตุกา ภิกฺขู ปสฺสนฺติ อาวาสิเก ภิกฺขู นานาสํวาสเก, เต สมานสํวาสก\-ทิฏฺฐิํ ปฏิลภนฺติ, \csromanfolio{247}สมานสํวาสก\-ทิฏฺฐิํ ปฏิลภิตฺวา น ปุจฺฉนฺติ, อปุจฺฉิตฺวา เอกโต ปวาเรนฺติ อนาปตฺติ. เต ปุจฺฉนฺติ, ปุจฺฉิตฺวา นาภิวิตรนฺติ, อนภิวิตริตฺวา เอกโต ปวาเรนฺติ, อาปตฺติ ทุกฺกฏสฺส. เต ปุจฺฉนฺติ, ปุจฺฉิตฺวา นาภิวิตรนฺติ, อนภิวิตริตฺวา ปาเฏกฺกํ ปวาเรนฺติ, อนาปตฺติ.}

\prose{อิธ ปน ภิกฺขเว อาคนฺตุกา ภิกฺขู ปสฺสนฺติ อาวาสิเก ภิกฺขู สมาน\-สํวาสเก, เต นานา\-สํวาสก\-ทิฏฺฐิํ ปฏิลภนฺติ, นานา\-สํวาสก\-ทิฏฺฐิํ ปฏิลภิตฺวา น ปุจฺฉนฺติ, อปุจฺฉิตฺวา เอกโต ปวาเรนฺติ, อาปตฺติ ทุกฺกฏสฺส.}

\prose{เต ปุจฺฉนฺติ, ปุจฺฉิตฺวา อภิวิตรนฺติ, อภิวิตริตฺวา ปาเฏกฺกํ ปวาเรนฺติ, อาปตฺติ ทุกฺกฏสฺส.}

\prose{เต ปุจฺฉนฺติ, ปุจฺฉิตฺวา อภิวิตรนฺติ, อภิวิตริตฺวา เอกโต ปวาเรนฺติ, อนาปตฺติ.}

\prose{อิธ ปน ภิกฺขเว อาวาสิกา ภิกฺขู ปสฺสนฺติ อาคนฺตุเก ภิกฺขู นานาสํวาสเก, เต สมานสํวาสก\-ทิฏฺฐิํ ปฏิลภนฺติ, สมานสํวาสก\-ทิฏฺฐิํ ปฏิลภิตฺวา น ปุจฺฉนฺติ, อปุจฺฉิตฺวา เอกโต ปวาเรนฺติ, อนาปตฺติ.}

\prose{เต ปุจฺฉนฺติ, ปุจฺฉิตฺวา นาภิวิตรนฺติ, อนภิวิตริตฺวา เอกโต ปวาเรนฺติ, อาปตฺติ ทุกฺกฏสฺส.}

\prose{เต ปุจฺฉนฺติ, ปุจฺฉิตฺวา นาภิวิตรนฺติ, อนภิวิตริตฺวา ปาเฏกฺกํ ปวาเรนฺติ, อนาปตฺติ.}

\prose{อิธ ปน ภิกฺขเว อาวาสิกา ภิกฺขู ปสฺสนฺติ อาคนฺตุเก ภิกฺขู สมาน\-สํวาสเก, เต นานา\-สํวาสก\-ทิฏฺฐิํ ปฏิลภนฺติ, นานา\-สํวาสก\-ทิฏฺฐิํ ปฏิลภิตฺวา น ปุจฺฉนฺติ, อปุจฺฉิตฺวา เอกโต ปวาเรนฺติ, อาปตฺติ ทุกฺกฏสฺส.}

\prose{เต ปุจฺฉนฺติ, ปุจฺฉิตฺวา อภิวิตรนฺติ, อภิวิตริตฺวา ปาเฏกฺกํ ปวาเรนฺติ, อาปตฺติ ทุกฺกฏสฺส.}

\prose{เต ปุจฺฉนฺติ, ปุจฺฉิตฺวา อภิวิตรนฺติ, อภิวิตริตฺวา เอกโต ปวาเรนฺติ, อนาปตฺติ.}

\csromanheader{2}{137. \textbf{น คนฺตพฺพวาร}}
\proseitem{231}{น ภิกฺขเว ตทหุ ปวารณาย สภิกฺขุกา อาวาสา อภิกฺขุโก อาวาโส คนฺตพฺโพ อญฺญตฺร สํเฆน อญฺญตฺร อนฺตรายา. น ภิกฺขเว ตทหุ ปวารณาย สภิกฺขุกา อาวาสา อภิกฺขุโก อนาวาโส คนฺตพฺโพ อญฺญตฺร สํเฆน อญฺญตฺร อนฺตรายา. น ภิกฺขเว ตทหุ ปวารณาย สภิกฺขุกา อาวาสา อภิกฺขุโก อาวาโส วา อนาวาโส วา คนฺตพฺโพ อญฺญตฺร สํเฆน อญฺญตฺร อนฺตรายา.}

\prose{\csromanfolio{248}น ภิกฺขเว ตทหุ ปวารณาย สภิกฺขุกา อนาวาสา อภิกฺขุโก อาวาโส คนฺตพฺโพ อญฺญตฺร สํเฆน อญฺญตฺร อนฺตรายา. น ภิกฺขเว ตทหุ ปวารณาย สภิกฺขุกา อนาวาสา อภิกฺขุโก อนาวาโส คนฺตพฺโพ อญฺญตฺร สํเฆน อญฺญตฺร อนฺตรายา. น ภิกฺขเว ตทหุ ปวารณาย สภิกฺขุกา อนาวาสา อภิกฺขุโก อาวาโส วา อนาวาโส วา คนฺตพฺโพ อญฺญตฺร สํเฆน อญฺญตฺร อนฺตรายา.}

\prose{น ภิกฺขเว ตทหุ ปวารณาย สภิกฺขุกา อาวาสา วา อนาวาสา วา อภิกฺขุโก อาวาโส คนฺตพฺโพ อญฺญตฺร สํเฆน อญฺญตฺร อนฺตรายา. น ภิกฺขเว ตทหุ ปวารณาย สภิกฺขุกา อาวาสา วา อนาวาสา วา อภิกฺขุโก อนาวาโส คนฺตพฺโพ อญฺญตฺร สํเฆน อญฺญตฺร อนฺตรายา. น ภิกฺขเว ตทหุ ปวารณาย สภิกฺขุกา อาวาสา วา อนาวาสา วา อภิกฺขุโก อาวาโส วา อนาวาโส วา คนฺตพฺโพ อญฺญตฺร สํเฆน อญฺญตฺร อนฺตรายา.}

\prose{น ภิกฺขเว ตทหุ ปวารณาย สภิกฺขุกา อาวาสา สภิกฺขุโก อาวาโส คนฺตพฺโพ ยตฺถสฺสุ ภิกฺขู นานาสํวาสกา อญฺญตฺร สํเฆน อญฺญตฺร อนฺตรายา. น ภิกฺขเว ตทหุ ปวารณาย สภิกฺขุกา อาวาสา สภิกฺขุโก อนาวาโส คนฺตพฺโพ ยตฺถสฺสุ ภิกฺขู นานาสํวาสกา อญฺญตฺร สํเฆน อญฺญตฺร อนฺตรายา. น ภิกฺขเว ตทหุ ปวารณาย สภิกฺขุกา อาวาสา สภิกฺขุโก อาวาโส วา อนาวาโส วา คนฺตพฺโพ ยตฺถสฺสุ ภิกฺขู นานาสํวาสกา อญฺญตฺร สํเฆน อญฺญตฺร อนฺตรายา.}

\prose{น ภิกฺขเว ตทหุ ปวารณาย สภิกฺขุกา อนาวาสา สภิกฺขุโก อาวาโส คนฺตพฺโพ ยตฺถสฺสุ ภิกฺขู นานาสํวาสกา อญฺญตฺร สํเฆน อญฺญตฺร อนฺตรายา. น ภิกฺขเว ตทหุ ปวารณาย สภิกฺขุกา อนาวาสา สภิกฺขุโก อนาวาโส คนฺตพฺโพ ยตฺถสฺสุ ภิกฺขู นานาสํวาสกา อญฺญตฺร สํเฆน อญฺญตฺร อนฺตรายา. น ภิกฺขเว ตทหุ ปวารณาย สภิกฺขุกา อนาวาสา สภิกฺขุโก อาวาโส วา อนาวาโส วา คนฺตพฺโพ ยตฺถสฺสุ ภิกฺขู นานาสํวาสกา อญฺญตฺร สํเฆน อญฺญตฺร อนฺตรายา.}

\prose{\csromanfolio{249}น ภิกฺขเว ตทหุ ปวารณาย สภิกฺขุกา อาวาสา วา อนาวาสา วา สภิกฺขุโก อาวาโส คนฺตพฺโพ ยตฺถสฺสุ ภิกฺขู นานาสํวาสกา อญฺญตฺร สํเฆน อญฺญตฺร อนฺตรายา. น ภิกฺขเว ตทหุ ปวารณาย สภิกฺขุกา อาวาสา วา อนาวาสา วา สภิกฺขุโก อนาวาโส คนฺตพฺโพ ยตฺถสฺสุ ภิกฺขู นานาสํวาสกา อญฺญตฺร สํเฆน อญฺญตฺร อนฺตรายา. น ภิกฺขเว ตทหุ ปวารณาย สภิกฺขุกา อาวาสา วา อนาวาสา วา สภิกฺขุโก อาวาโส วา อนาวาโส วา คนฺตพฺโพ ยตฺถสฺสุ ภิกฺขู นานาสํวาสกา อญฺญตฺร สํเฆน อญฺญตฺร อนฺตรายา.}

\csromanmatikaanchor{mtk.142}
\csromantocmarkref{section}{138. คนฺตพฺพวาร}{mtk.142}
\bookmark[dest={mtk.142},level=1]{138. คนฺตพฺพวาร}
\csromanheader{1}{138. \textbf{คนฺตพฺพวาร}}
\proseitem{232}{คนฺตพฺโพ ภิกฺขเว ตทหุ ปวารณาย สภิกฺขุกา อาวาสา สภิกฺขุโก อาวาโส ยตฺถสฺสุ ภิกฺขู สมาน\-สํวาสกา ยํ ชญฺญา “สกฺโกมิ อชฺเชว คนฺตุนฺ”ติ. คนฺตพฺโพ ภิกฺขเว ตทหุ ปวารณาย สภิกฺขุกา อาวาสา สภิกฺขุโก อนาวาโส ฯเปฯ สภิกฺขุโก อาวาโส วา อนาวาโส วา ยตฺถสฺสุ ภิกฺขู สมาน\-สํวาสกา ยํ ชญฺญา “สกฺโกมิ อชฺเชว คนฺตุนฺ”ติ.}

\prose{คนฺตพฺโพ ภิกฺขเว ตทหุ ปวารณาย สภิกฺขุกา อนาวาสา สภิกฺขุโก อาวาโส ฯเปฯ สภิกฺขุโก อนาวาโส ฯเปฯ สภิกฺขุโก อาวาโส วา อนาวาโส วา ยตฺถสฺสุ ภิกฺขู สมาน\-สํวาสกา ยํ ชญฺญา “สกฺโกมิ อชฺเชว คนฺตุนฺ”ติ.}

\prose{คนฺตพฺโพ ภิกฺขเว ตทหุ ปวารณาย สภิกฺขุกา อาวาสา วา อนาวาสา วา สภิกฺขุโก อาวาโส ฯเปฯ สภิกฺขุโก อนาวาโส ฯเปฯ สภิกฺขุโก อาวาโส วา อนาวาโส วา ยตฺถสฺสุ ภิกฺขู สมาน\-สํวาสกา ยํ ชญฺญา “สกฺโกมิ อชฺเชว คนฺตุนฺ”ติ.}

\csromanmatikaanchor{mtk.143}
\csromantocmarkref{section}{139. วชฺชนียปุคฺคลสนฺทสฺสนา}{mtk.143}
\bookmark[dest={mtk.143},level=1]{139. วชฺชนียปุคฺคลสนฺทสฺสนา}
\csromanheader{1}{139. \textbf{วชฺชนีย}\-\textbf{ปุคฺคล}\-\textbf{สนฺทสฺสนา}}
\proseitem{233}{น ภิกฺขเว ภิกฺขุนิยา นิสินฺน\-ปริสาย ปวาเรตพฺพํ, โย ปวาเรยฺย, อาปตฺติ ทุกฺกฏสฺส.}

\prose{น ภิกฺขเว สิกฺขมานาย ฯเปฯ. น สามเณรสฺส ฯเปฯ. น สามเณริยา ฯเปฯ. น สิกฺขํ ปจฺจกฺขาตกสฺส ฯเปฯ. น อนฺติมวตฺถุํ อชฺฌาปนฺนกสฺส นิสินฺน\-ปริสาย ปวาเรตพฺพํ, โย ปวาเรยฺย, อาปตฺติ ทุกฺกฏสฺส.}

\prose{\csromanfolio{250}น อาปตฺติยา อทสฺสเน อุกฺขิตฺตกสฺส นิสินฺน\-ปริสาย ปวาเรตพฺพํ, โย ปวาเรยฺย, ยถาธมฺโม กาเรตพฺโพ.}

\prose{น อาปตฺติยา อปฺปฏิกมฺเม อุกฺขิตฺตกสฺส ฯเปฯ. น ปาปิกาย ทิฏฺฐิยา อปฺปฏินิสฺสคฺเค อุกฺขิตฺตกสฺส นิสินฺน\-ปริสาย ปวาเรตพฺพํ, โย ปวาเรยฺย, ยถาธมฺโม กาเรตพฺโพ.}

\prose{น ปณฺฑกสฺส นิสินฺน\-ปริสาย ปวาเรตพฺพํ, โย ปวาเรยฺย, อาปตฺติ ทุกฺกฏสฺส.}

\prose{น เถยฺย\-สํวาสกสฺส ฯเปฯ. น ติตฺถิย\-ปกฺกนฺตกสฺส ฯเปฯ. น ติรจฺฉาน\-คตสฺส ฯเปฯ. น มาตุฆาตกสฺส ฯเปฯ. น ปิตุฆาตกสฺส ฯเปฯ. น อรหนฺต\-ฆาตกสฺส ฯเปฯ. น ภิกฺขุนิ\-ทูสกสฺส ฯเปฯ. น สํฆเภทกสฺส ฯเปฯ. น โลหิตุ\-ปฺปาทกสฺส ฯเปฯ. น อุภโต\-พฺยญฺชนกสฺส นิสินฺน\-ปริสาย ปวาเรตพฺพํ, โย ปวาเรยฺย, อาปตฺติ ทุกฺกฏสฺส.}

\proseitemwithcloserruled{233}{น ภิกฺขเว ปาริวาสิก\-ปวารณา\-ทาเนน ปวาเรตพฺพํ อญฺญตฺร อวุฏฺฐิตาย ปริสาย. น จ ภิกฺขเว อปฺปวารณาย ปวาเรตพฺพํ อญฺญตฺร สํฆสามคฺคิยาติ.}{ทุติย\-ภาณวาโร นิฏฺฐิโต.}
\csromanmatikaanchor{mtk.144}
\csromantocmarkref{section}{140. เทฺววาจิกาทิปวารณา}{mtk.144}
\bookmark[dest={mtk.144},level=1]{140. เทฺววาจิกาทิปวารณา}
\csromanheader{1}{140. \textbf{เทฺววาจิกา}\-\textbf{ทิ}\-\textbf{ปวารณา}}
\proseitem{234}{เตน โข ปน สมเยน โกสเลสุ ชนปเท อญฺญตรสฺมิํ อาวาเส ตทหุ ปวารณาย สวรภยํ อโหสิ. ภิกฺขู นาสกฺขิํสุ เตวาจิกํ ปวาเรตุํ. ภควโต เอตมตฺถํ อาโรเจสุํ. อนุชานามิ ภิกฺขเว เทฺววาจิกํ ปวาเรตุนฺติ.}

\prose{พาฬฺหตรํ สวรภยํ อโหสิ. ภิกฺขู นาสกฺขิํสุ เทฺววาจิกํ ปวาเรตุํ. ภควโต เอตมตฺถํ อาโรเจสุํ. อนุชานามิ ภิกฺขเว เอกวาจิกํ ปวาเรตุนฺติ.}

\prose{พาฬฺหตรํ สวรภยํ อโหสิ. ภิกฺขู นาสกฺขิํสุ เอกวาจิกํ ปวาเรตุํ. ภควโต เอตมตฺถํ อาโรเจสุํ. อนุชานามิ ภิกฺขเว สมานวสฺสิกํ ปวาเรตุนฺติ.}

\prose{เตน โข ปน สมเยน อญฺญตรสฺมิํ อาวาเส ตทหุ ปวารณาย มนุสฺเสหิ ทานํ เทนฺเตหิ เยภุยฺเยน รตฺติ เขปิตา โหติ. อถ โข \csromanfolio{251}เตสํ ภิกฺขูนํ เอตทโหสิ “มนุสฺเสหิ ทานํ เทนฺเตหิ เยภุยฺเยน รตฺติ เขปิตา, สเจ สํโฆ เตวาจิกํ ปวาเรสฺสติ, อปฺปวาริโตว สํโฆ ภวิสฺสติ, อถายํ รตฺติ วิภายิสฺสติ, กถํ นุ โข อเมฺหหิ ปฏิปชฺชิตพฺพนฺ”ติ. ภควโต เอตมตฺถํ อาโรเจสุํ. อิธ ปน ภิกฺขเว อญฺญตรสฺมิํ อาวาเส ตทหุ ปวารณาย มนุสฺเสหิ ทานํ เทนฺเตหิ เยภุยฺเยน รตฺติ เขปิตา โหติ, ตตฺร เจ ภิกฺขเว ภิกฺขูนํ เอวํ โหติ “มนุสฺเสหิ ทานํ เทนฺเตหิ เยภุยฺเยน รตฺติ เขปิตา, สเจ สํโฆ เตวาจิกํ ปวาเรสฺสติ, อปฺปวาริโตว สํโฆ ภวิสฺสติ, อถายํ รตฺติ วิภายิสฺสตี”ติ. พฺยตฺเตน ภิกฺขุนา ปฏิพเลน สํโฆ ญาเปตพฺโพ\csromandash{}}

\prose{“สุณาตุ เม ภนฺเต สํโฆ, มนุสฺเสหิ ทานํ เทนฺเตหิ เยภุยฺเยน รตฺติ เขปิตา, สเจ สํโฆ เตวาจิกํ ปวาเรสฺสติ, อปฺปวาริโตว สํโฆ ภวิสฺสติ, อถายํ รตฺติ วิภายิสฺสติ, ยทิ สํฆสฺส ปตฺตกลฺลํ, สํโฆ เทฺววาจิกํ เอกวาจิกํ สมานวสฺสิกํ ปวาเรยฺยา”ติ.}

\prose{อิธ ปน ภิกฺขเว อญฺญตรสฺมิํ อาวาเส ตทหุ ปวารณาย ภิกฺขูหิ ธมฺมํ ภณนฺเตหิ ฯเปฯ สุตฺตนฺติเกหิ สุตฺตนฺตํ สงฺคายนฺเตหิ. วินยธเรหิ วินยํ วินิจฺฉิน\-นฺเตหิ. ธมฺม\-กถิเกหิ ธมฺมํ สากจฺฉนฺเตหิ. ภิกฺขูหิ กลหํ กโรนฺเตหิ เยภุยฺเยน รตฺติ เขปิตา โหติ, ตตฺร เจ ภิกฺขูนํ เอวํ โหติ “ภิกฺขูหิ กลหํ กโรนฺเตหิ เยภุยฺเยน รตฺติ เขปิตา, สเจ สํโฆ เตวาจิกํ ปวาเรสฺสติ, อปฺปวาริโตว สํโฆ ภวิสฺสติ, อถายํ รตฺติ วิภายิสฺสตี”ติ. พฺยตฺเตน ภิกฺขุนา ปฏิพเลน สํโฆ ญาเปตพฺโพ\csromandash{}}

\prose{“สุณาตุ เม ภนฺเต สํโฆ, ภิกฺขูหิ กลหํ กโรนฺเตหิ เยภุยฺเยน รตฺติ เขปิตา, สเจ สํโฆ เตวาจิกํ ปวาเรสฺสติ, อปฺปวาริโตว สํโฆ ภวิสฺสติ, อถายํ รตฺติ วิภายิสฺสติ, ยทิ สํฆสฺส ปตฺตกลฺลํ, สํโฆ เทฺววาจิกํ เอกวาจิกํ สมานวสฺสิกํ ปวาเรยฺยา”ติ.}

\prose{เตน โข ปน สมเยน โกสเลสุ ชนปเท อญฺญตรสฺมิํ อาวาเส ตทหุ ปวารณาย มหา\-ภิกฺขุสํโฆ สนฺนิปติโต \csromanfolio{252}โหติ, ปริตฺตญฺจ อโนวสฺสิกํ\csromansharedfootnote{b7e3e9f0a8b635d7}{อโนวสฺสกํ (ก)} โหติ, มหา จ เมโฆ อุคฺคโต โหติ. อถ โข เตสํ ภิกฺขูนํ เอตทโหสิ “อยํ โข มหา\-ภิกฺขุสํโฆ สนฺนิปติโต, ปริตฺตญฺจ อโนวสฺสิกํ, มหา จ เมโฆ อุคฺคโต, สเจ สํโฆ เตวาจิกํ ปวาเรสฺสติ, อปฺปวาริโตว สํโฆ ภวิสฺสติ, อถายํ เมโฆ ปวสฺสิสฺสติ, กถํ นุ โข อเมฺหหิ ปฏิปชฺชิตพฺพนฺ”ติ. ภควโต เอตมตฺถํ อาโรเจสุํ. อิธ ปน ภิกฺขเว อญฺญตรสฺมิํ อาวาเส ตทหุ ปวารณาย มหา\-ภิกฺขุสํโฆ สนฺนิปติโต โหติ, ปริตฺตญฺจ อโนวสฺสิกํ โหติ, มหา จ เมโฆ อุคฺคโต โหติ, ตตฺร เจ ภิกฺขูนํ เอวํ โหติ “อยํ โข มหา\-ภิกฺขุสํโฆ สนฺนิปติโต, ปริตฺตญฺจ อโนวสฺสิกํ, มหา จ เมโฆ อุคฺคโต, สเจ สํโฆ เตวาจิกํ ปวาเรสฺสติ, อปฺปวาริโตว สํโฆ ภวิสฺสติ, อถายํ เมโฆ ปวสฺสิสฺสตี”ติ. พฺยตฺเตน ภิกฺขุนา ปฏิพเลน สํโฆ ญาเปตพฺโพ\csromandash{}}

\prose{“สุณาตุ เม ภนฺเต สํโฆ, อยํ มหา\-ภิกฺขุสํโฆ สนฺนิปติโต, ปริตฺตญฺจ อโนวสฺสิกํ, มหา จ เมโฆ อุคฺคโต, สเจ สํโฆ เตวาจิกํ ปวาเรสฺสติ, อปฺปวาริโตว สํโฆ ภวิสฺสติ, อถายํ เมโฆ ปวสฺสิสฺสติ, ยทิ สํฆสฺส ปตฺตกลฺลํ, สํโฆ เทฺววาจิกํ เอกวาจิกํ สมานวสฺสิกํ ปวาเรยฺยา”ติ. อิธ ปน ภิกฺขเว อญฺญตรสฺมิํ อาวาเส ตทหุ ปวารณาย ราชนฺตราโย โหติ ฯเปฯ โจรนฺตราโย โหติ. อคฺยนฺตราโย โหติ. อุทกนฺตราโย โหติ. มนุสฺสนฺตราโย โหติ. อมนุสฺสนฺตราโย โหติ. วาฬนฺตราโย โหติ. สรีสปนฺตราโย โหติ. ชีวิตนฺตราโย โหติ. พฺรหฺมจริย\-นฺตราโย โหติ, ตตฺร เจ ภิกฺขูนํ เอวํ โหติ “อยํ โข พฺรหฺมจริย\-นฺตราโย, สเจ สํโฆ เตวาจิกํ ปวาเรสฺสติ, อปฺปวาริโตว สํโฆ ภวิสฺสติ, อถายํ พฺรหฺมจริย\-นฺตราโย ภวิสฺสตี”ติ. พฺยตฺเตน ภิกฺขุนา ปฏิพเลน สํโฆ ญาเปตพฺโพ\csromandash{}}

\prose{“สุณาตุ เม ภนฺเต สํโฆ, อยํ พฺรหฺมจริย\-นฺตราโย, สเจ สํโฆ เตวาจิกํ ปวาเรสฺสติ, อปฺปวาริโตว สํโฆ ภวิสฺสติ, \csromanfolio{253}อถายํ พฺรหฺมจริย\-นฺตราโย ภวิสฺสติ, ยทิ สํฆสฺส ปตฺตกลฺลํ, สํโฆ เทฺววาจิกํ เอกวาจิกํ สมานวสฺสิกํ ปวาเรยฺยา”ติ.}

\csromanmatikaanchor{mtk.145}
\csromantocmarkref{section}{141. ปวารณาฐปน}{mtk.145}
\bookmark[dest={mtk.145},level=1]{141. ปวารณาฐปน}
\csromanheader{1}{141. \textbf{ปวารณา}\-\textbf{ฐปน}}
\proseitem{235}{เตน โข ปน สมเยน ฉพฺพคฺคิยา ภิกฺขู สาปตฺติกา ปวาเรนฺติ. ภควโต เอตมตฺถํ อาโรเจสุํ. น ภิกฺขเว สาปตฺติเกน ปวาเรตพฺพํ, โย ปวาเรยฺย, อาปตฺติ ทุกฺกฏสฺส. อนุชานามิ ภิกฺขเว โย สาปตฺติโก ปวาเรติ, ตสฺส โอกาสํ การาเปตฺวา อาปตฺติยา โจเทตุนฺติ.}

\prose{เตน โข ปน สมเยน ฉพฺพคฺคิยา ภิกฺขู โอกาสํ การาปิยมานา น อิจฺฉนฺติ โอกาสํ กาตุํ. ภควโต เอตมตฺถํ อาโรเจสุํ. อนุชานามิ ภิกฺขเว โอกาสํ อกโรนฺตสฺส ปวารณํ ฐเปตุํ. เอวญฺจ ปน ภิกฺขเว ฐเปตพฺพา, ตทหุ ปวารณาย จาตุทฺทเส วา ปนฺนรเส วา ตสฺมิํ ปุคฺคเล สมฺมุขีภูเต สํฆมชฺเฌ อุทาหริตพฺพํ “สุณาตุ เม ภนฺเต สํโฆ, อิตฺถนฺนาโม ปุคฺคโล สาปตฺติโก, ตสฺส ปวารณํ ฐเปมิ, น ตสฺมิํ สมฺมุขีภูเต ปวาเรตพฺพนฺ”ติ, ฐปิตา โหติ ปวารณาติ.}

\prose{เตน โข ปน สมเยน ฉพฺพคฺคิยา ภิกฺขู “ปุรมฺหากํ เปสลา ภิกฺขู ปวารณํ ฐเปนฺตี”ติ ปฏิกจฺเจว สุทฺธานํ ภิกฺขูนํ อนาปตฺติกานํ อวตฺถุสฺมิํ อการเณ ปวารณํ ฐเปนฺติ, ปวาริตานมฺปิ ปวารณํ ฐเปนฺติ.}

\prose{ภควโต เอตมตฺถํ อาโรเจสุํ. น ภิกฺขเว สุทฺธานํ ภิกฺขูนํ อนาปตฺติกานํ อวตฺถุสฺมิํ อการเณ ปวารณา ฐเปตพฺพา, โย ฐเปยฺย, อาปตฺติ ทุกฺกฏสฺส. น ภิกฺขเว ปวาริตานมฺปิ ปวารณา ฐเปตพฺพา, โย ฐเปยฺย, อาปตฺติ ทุกฺกฏสฺส.}

\proseitem{236}{เอวํ โข ภิกฺขเว ฐปิตา โหติ ปวารณา, เอวํ อฏฺฐปิตา. กถญฺจ ภิกฺขเว อฏฺฐปิตา โหติ ปวารณา, เตวาจิกาย เจ ภิกฺขเว ปวารณาย ภาสิตาย ลปิตาย ปริโยสิตาย ปวารณํ ฐเปติ, อฏฺฐปิตา โหติ ปวารณา. เทฺววาจิกาย เจ ภิกฺขเว. เอกวาจิกาย เจ ภิกฺขเว. สมาน\-วสฺสิกาย เจ ภิกฺขเว ปวารณาย ภาสิตาย ลปิตาย ปริโยสิตาย ปวารณํ \csromanfolio{254}ฐเปติ, อฏฺฐปิตา โหติ ปวารณา. เอวํ โข ภิกฺขเว อฏฺฐปิตา โหติ ปวารณา.}

\prose{กถญฺจ ภิกฺขเว ฐปิตา โหติ ปวารณา, เตวาจิกาย เจ ภิกฺขเว ปวารณาย ภาสิตาย ลปิตาย อปริโยสิตาย ปวารณํ ฐเปติ, ฐปิตา โหติ ปวารณา. เทฺววาจิกาย เจ ภิกฺขเว. เอกวาจิกาย เจ ภิกฺขเว. สมาน\-วสฺสิกาย เจ ภิกฺขเว ปวารณาย ภาสิตาย ลปิตาย อปริโยสิตาย ปวารณํ ฐเปติ, ฐปิตา โหติ ปวารณา. เอวํ โข ภิกฺขเว ฐปิตา โหติ ปวารณา.}

\proseitem{237}{อิธ ปน ภิกฺขเว ตทหุ ปวารณาย ภิกฺขุ ภิกฺขุสฺส ปวารณํ ฐเปติ, ตํ เจ ภิกฺขุํ อญฺเญ ภิกฺขู ชานนฺติ “อยํ โข อายสฺมา อปริสุทฺธ\-กายสมาจาโร อปริสุทฺธ\-วจีสมาจาโร อปริสุทฺธา\-ชีโว พาโล อพฺยตฺโต น ปฏิพโล อนุยุญฺชียมาโน อนุโยคํ ทาตุนฺ”ติ, “อลํ ภิกฺขุ มา ภณฺฑนํ มา กลหํ มา วิคฺคหํ มา วิวาทนฺ”ติ โอมทฺทิตฺวา สํเฆน ปวาเรตพฺพํ.}

\prose{อิธ ปน ภิกฺขเว ตทหุ ปวารณาย ภิกฺขุ ภิกฺขุสฺส ปวารณํ ฐเปติ, ตํ เจ ภิกฺขุํ อญฺเญ ภิกฺขู ชานนฺติ “อยํ โข อายสฺมา ปริสุทฺธ\-กาย\-สมาจาโร อปริสุทฺธ\-วจีสมาจาโร อปริสุทฺธา\-ชีโว พาโล อพฺยตฺโต น ปฏิพโล อนุยุญฺชียมาโน อนุโยคํ ทาตุนฺ”ติ, “อลํ ภิกฺขุ มา ภณฺฑนํ มา กลหํ มา วิคฺคหํ มา วิวาทนฺ”ติ โอมทฺทิตฺวา สํเฆน ปวาเรตพฺพํ.}

\prose{อิธ ปน ภิกฺขเว ตทหุ ปวารณาย ภิกฺขุ ภิกฺขุสฺส ปวารณํ ฐเปติ, ตํ เจ ภิกฺขุํ อญฺเญ ภิกฺขู ชานนฺติ “อยํ โข อายสฺมา ปริสุทฺธ\-กาย\-สมาจาโร ปริสุทฺธ\-วจี\-สมาจาโร อปริสุทฺธา\-ชีโว พาโล อพฺยตฺโต น ปฏิพโล อนุยุญฺชียมาโน อนุโยคํ ทาตุนฺ”ติ, “อลํ ภิกฺขุ มา ภณฺฑนํ มา กลหํ มา วิคฺคหํ มา วิวาทนฺ”ติ โอมทฺทิตฺวา สํเฆน ปวาเรตพฺพํ.}

\proseitem{237}{อิธ ปน ภิกฺขเว ตทหุ ปวารณาย ภิกฺขุ ภิกฺขุสฺส ปวารณํ ฐเปติ, ตํ เจ ภิกฺขุํ อญฺเญ ภิกฺขู ชานนฺติ “อยํ โข อายสฺมา ปริสุทฺธ\-กาย\-สมาจาโร ปริสุทฺธ\-วจี\-สมาจาโร ปริสุทฺธาชีโว พาโล \csromanfolio{255}อพฺยตฺโต น ปฏิพโล อนุยุญฺชียมาโน อนุโยคํ ทาตุนฺ”ติ, “อลํ ภิกฺขุ มา ภณฺฑนํ มา กลหํ มา วิคฺคหํ มา วิวาทนฺ”ติ โอมทฺทิตฺวา สํเฆน ปวาเรตพฺพํ.}

\prose{อิธ ปน ภิกฺขเว ตทหุ ปวารณาย ภิกฺขุ ภิกฺขุสฺส ปวารณํ ฐเปติ, ตํ เจ ภิกฺขุํ อญฺเญ ภิกฺขู ชานนฺติ “อยํ โข อายสฺมา ปริสุทฺธ\-กาย\-สมาจาโร ปริสุทฺธ\-วจี\-สมาจาโร ปริสุทฺธาชีโว ปณฺฑิโต พฺยตฺโต ปฏิพโล อนุยุญฺชียมาโน อนุโยคํ ทาตุนฺ”ติ, โส เอวมสฺส วจนีโย “ยํ โข ตฺวํ อาวุโส อิมสฺส ภิกฺขุโน ปวารณํ ฐเปสิ, กิมฺหิ นํ ฐเปสิ, สีลวิปตฺติยา วา ฐเปสิ, อาจารวิปตฺติยา วา ฐเปสิ, ทิฏฺฐิ\-วิปตฺติยา วา ฐเปสี”ติ. โส เจ เอวํ วเทยฺย “สีลวิปตฺติยา วา ฐเปมิ, อาจารวิปตฺติยา วา ฐเปมิ, ทิฏฺฐิ\-วิปตฺติยา วา ฐเปมี”ติ. โส เอวมสฺส วจนีโย “ชานาสิ ปนายสฺมา สีลวิปตฺติํ, ชานาสิ อาจารวิปตฺติํ, ชานาสิ ทิฏฺฐิวิปตฺตินฺ”ติ. โส เจ เอวํ วเทยฺย “ชานามิ โข อหํ อาวุโส สีลวิปตฺติํ, ชานามิ อาจารวิปตฺติํ, ชานามิ ทิฏฺฐิวิปตฺตินฺ”ติ. โส เอวมสฺส วจนีโย “กตมา ปนาวุโส สีลวิปตฺติ, กตมา อาจารวิปตฺติ, กตมา ทิฏฺฐิวิปตฺตี”ติ. โส เจ เอวํ วเทยฺย “จตฺตาริ ปาราชิกานิ เตรส สํฆาทิเสสา, อยํ สีลวิปตฺติ. ถุลฺลจฺจยํ ปาจิตฺติยํ ปาฏิเทสนียํ ทุกฺกฏํ ทุพฺภาสิตํ, อยํ อาจารวิปตฺติ. มิจฺฉาทิฏฺฐิ อนฺตคฺคาหิกา\-ทิฏฺฐิ, อยํ ทิฏฺฐิวิปตฺตี”ติ. โส เอวมสฺส วจนีโย “ยํ โข ตฺวํ อาวุโส อิมสฺส ภิกฺขุโน ปวารณํ ฐเปสิ, ทิฏฺเฐน วา ฐเปสิ, สุเตน วา ฐเปสิ, ปริสงฺกาย วา ฐเปสี”ติ. โส เจ เอวํ วเทยฺย “ทิฏฺเฐน วา ฐเปมิ, สุเตน วา ฐเปมิ, ปริสงฺกาย วา ฐเปมี”ติ. โส เอวมสฺส วจนีโย “ยํ โข ตฺวํ อาวุโส อิมสฺส ภิกฺขุโน ทิฏฺเฐน ปวารณํ ฐเปสิ, กิํ เต ทิฏฺฐํ, กินฺติ เต ทิฏฺฐํ, กทา เต ทิฏฺฐํ, กตฺถ เต ทิฏฺฐํ, ปาราชิกํ อชฺฌาปชฺชนฺโต ทิฏฺโฐ, สํฆาทิเสสํ อชฺฌาปชฺชนฺโต ทิฏฺโฐ, ถุลฺลจฺจยํ. ปาจิตฺติยํ. ปาฏิเทสนียํ. ทุกฺกฏํ. ทุพฺภาสิตํ อชฺฌาปชฺชนฺโต ทิฏฺโฐ, กตฺถ จ ตฺวํ อโหสิ, กตฺถ จายํ ภิกฺขุ อโหสิ, กิญฺจ ตฺวํ กโรสิ, กิญฺจายํ ภิกฺขุ กโรตี”ติ. โส เจ เอวํ วเทยฺย “น โข อหํ อาวุโส อิมสฺส ภิกฺขุโน ทิฏฺเฐน ปวารณํ ฐเปมิ, อปิ จ สุเตน ปวารณํ ฐเปมี”ติ. โส เอวมสฺส วจนีโย “ยํ โข ตฺวํ อาวุโส อิมสฺส \csromanfolio{256}ภิกฺขุโน สุเตน ปวารณํ ฐเปสิ, กิํ เต สุตํ, กินฺติ เต สุตํ, กทา เต สุตํ, กตฺถ เต สุตํ, ‘ปาราชิกํ อชฺฌาปนฺโน’ติ สุตํ, ‘สํฆาทิเสสํ อชฺฌาปนฺโน’ติ สุตํ, ถุลฺลจฺจยํ. ปาจิตฺติยํ. ปาฏิเทสนียํ. ทุกฺกฏํ. ‘ทุพฺภาสิตํ อชฺฌาปนฺโน’ติ สุตํ, ภิกฺขุสฺส สุตํ, ภิกฺขุนิยา สุตํ, สิกฺขมานาย สุตํ, สามเณรสฺส สุตํ, สามเณริยา สุตํ, อุปาสกสฺส สุตํ, อุปาสิกาย สุตํ, ราชูนํ สุตํ, ราช\-มหามตฺตานํ สุตํ, ติตฺถิยานํ สุตํ, ติตฺถิย\-สาวกานํ สุตนฺ”ติ. โส เจ เอวํ วเทยฺย “น โข อหํ อาวุโส อิมสฺส ภิกฺขุโน สุเตน ปวารณํ ฐเปมิ, อปิ จ ปริสงฺกาย ปวารณํ ฐเปมี”ติ. โส เอวมสฺส วจนีโย “ยํ โข ตฺวํ อาวุโส อิมสฺส ภิกฺขุโน ปริสงฺกาย ปวารณํ ฐเปสิ, กิํ ปริสงฺกสิ, กินฺติ ปริสงฺกสิ, กทา ปริสงฺกสิ, กตฺถ ปริสงฺกสิ, ‘ปาราชิกํ อชฺฌาปนฺโน’ติ ปริสงฺกสิ, ‘สํฆาทิเสสํ อชฺฌาปนฺโน’ติ ปริสงฺกสิ, ถุลฺลจฺจยํ. ปาจิตฺติยํ. ปาฏิเทสนียํ. ทุกฺกฏํ. ‘ทุพฺภาสิตํ อชฺฌาปนฺโน’ติ ปริสงฺกสิ, ภิกฺขุสฺส สุตฺวา ปริสงฺกสิ, ภิกฺขุนิยา สุตฺวา ปริสงฺกสิ, สิกฺขมานาย สุตฺวา ปริสงฺกสิ, สามเณรสฺส สุตฺวา ปริสงฺกสิ, สามเณริยา สุตฺวา ปริสงฺกสิ, อุปาสกสฺส สุตฺวา ปริสงฺกสิ, อุปาสิกาย สุตฺวา ปริสงฺกสิ, ราชูนํ สุตฺวา ปริสงฺกสิ, ราช\-มหามตฺตานํ สุตฺวา ปริสงฺกสิ, ติตฺถิยานํ สุตฺวา ปริสงฺกสิ, ติตฺถิย\-สาวกานํ สุตฺวา ปริสงฺกสี”ติ. โส เจ เอวํ วเทยฺย “น โข อหํ อาวุโส อิมสฺส ภิกฺขุโน ปริสงฺกาย ปวารณํ ฐเปมิ, อปิ จ อหมฺปิ น ชานามิ, เกน ปนาหํ อิมสฺส ภิกฺขุโน ปวารณํ ฐเปมี”ติ. โส เจ ภิกฺขเว โจทโก ภิกฺขุ อนุโยเคน วิญฺญูนํ สพฺรหฺมจารีนํ จิตฺตํ น อาราเธติ, “อนนุวาโท จุทิโต ภิกฺขู”ติ อลํ วจนาย. โส เจ ภิกฺขเว โจทโก ภิกฺขุ อนุโยเคน วิญฺญูนํ สพฺรหฺมจารีนํ จิตฺตํ อาราเธติ, “สานุวาโท จุทิโต ภิกฺขู”ติ อลํ วจนาย. โส เจ ภิกฺขเว โจทโก ภิกฺขุ อมูลเกน ปาราชิเกน อนุทฺธํสิตํ ปฏิชานาติ, สํฆาทิเสสํ อาโรเปตฺวา สํเฆน ปวาเรตพฺพํ. โส เจ ภิกฺขเว โจทโก ภิกฺขุ อมูลเกน สํฆาทิเสเสน อนุทฺธํสิตํ ปฏิชานาติ, ยถาธมฺมํ การาเปตฺวา สํเฆน ปวาเรตพฺพํ. โส เจ ภิกฺขเว โจทโก ภิกฺขุ อมูลเกน ถุลฺลจฺจเยน. ปาจิตฺติเยน. ปาฏิเทสนีเยน. ทุกฺกเฏน. ทุพฺภาสิเตน อนุทฺธํสิตํ ปฏิชานาติ, ยถาธมฺมํ การาเปตฺวา สํเฆน ปวาเรตพฺพํ. \csromanfolio{257}โส เจ ภิกฺขเว จุทิโต ภิกฺขุ “ปาราชิกํ อชฺฌาปนฺโน”ติ ปฏิชานาติ, นาเสตฺวา สํเฆน ปวาเรตพฺพํ. โส เจ ภิกฺขเว จุทิโต ภิกฺขุ “สํฆาทิเสสํ อชฺฌาปนฺโน”ติ ปฏิชานาติ, สํฆาทิเสสํ อาโรเปตฺวา สํเฆน ปวาเรตพฺพํ. โส เจ ภิกฺขเว จุทิโต ภิกฺขุ ถุลฺลจฺจยํ. ปาจิตฺติยํ. ปาฏิเทสนียํ. ทุกฺกฏํ. “ทุพฺภาสิตํ อชฺฌาปนฺโน”ติ ปฏิชานาติ, ยถาธมฺมํ การาเปตฺวา สํเฆน ปวาเรตพฺพํ.}

\csromanmatikaanchor{mtk.146}
\csromantocmarkref{section}{142. ถุลฺลจฺจยวตฺถุกาทิ}{mtk.146}
\bookmark[dest={mtk.146},level=1]{142. ถุลฺลจฺจยวตฺถุกาทิ}
\csromanheader{1}{142. \textbf{ถุลฺลจฺจย}\-\textbf{วตฺถุกาทิ}}
\proseitem{238}{อิธ ปน ภิกฺขเว ภิกฺขุ ตทหุ ปวารณาย ถุลฺลจฺจยํ อชฺฌาปนฺโน โหติ, เอกจฺเจ ภิกฺขู ถุลฺลจฺจย\-ทิฏฺฐิโน โหนฺติ, เอกจฺเจ ภิกฺขู สํฆาทิเสสทิฏฺฐิโน โหนฺติ.}

\prose{เย เต ภิกฺขเว ภิกฺขู ถุลฺลจฺจย\-ทิฏฺฐิโน, เตหิ โส ภิกฺขเว ภิกฺขุ เอกมนฺตํ อปเนตฺวา ยถาธมฺมํ การาเปตฺวา สํฆํ อุปสงฺกมิตฺวา เอวมสฺส วจนีโย “ยํ โข โส อาวุโส ภิกฺขุ อาปตฺติํ อาปนฺโน, สาสฺส ยถาธมฺมํ ปฏิกตา, ยทิ สํฆสฺส ปตฺตกลฺลํ, สํโฆ ปวาเรยฺยา”ติ.}

\prose{อิธ ปน ภิกฺขเว ภิกฺขุ ตทหุ ปวารณาย ถุลฺลจฺจยํ อชฺฌาปนฺโน โหติ, เอกจฺเจ ภิกฺขู ถุลฺลจฺจย\-ทิฏฺฐิโน โหนฺติ, เอกจฺเจ ภิกฺขู ปาจิตฺติย\-ทิฏฺฐิโน โหนฺติ ฯเปฯ เอกจฺเจ ภิกฺขู ถุลฺลจฺจย\-ทิฏฺฐิโน โหนฺติ, เอกจฺเจ ภิกฺขู ปาฏิเทสนีย\-ทิฏฺฐิโน โหนฺติ.}

\prose{เอกจฺเจ ภิกฺขู ถุลฺลจฺจย\-ทิฏฺฐิโน โหนฺติ, เอกจฺเจ ภิกฺขู ทุกฺกฏ\-ทิฏฺฐิโน โหนฺติ.}

\prose{เอกจฺเจ ภิกฺขู ถุลฺลจฺจย\-ทิฏฺฐิโน โหนฺติ, เอกจฺเจ ภิกฺขู ทุพฺภาสิต\-ทิฏฺฐิโน โหนฺติ.}

\prose{เย เต ภิกฺขเว ภิกฺขู ถุลฺลจฺจย\-ทิฏฺฐิโน, เตหิ โส ภิกฺขเว ภิกฺขุ เอกมนฺตํ อปเนตฺวา ยถาธมฺมํ การาเปตฺวา สํฆํ อุปสงฺกมิตฺวา เอวมสฺส วจนีโย “ยํ โข โส อาวุโส ภิกฺขุ อาปตฺติํ อาปนฺโน, สาสฺส ยถาธมฺมํ ปฏิกตา, ยทิ สํฆสฺส ปตฺตกลฺลํ, สํโฆ ปวาเรยฺยา”ติ.}

\proseitem{238}{อิธ ปน ภิกฺขเว ภิกฺขุ ตทหุ ปวารณาย ปาจิตฺติยํ อชฺฌาปนฺโน โหติ ฯเปฯ ปาฏิเทสนียํ อชฺฌาปนฺโน โหติ. ทุกฺกฏํ อชฺฌาปนฺโน โหติ. ทุพฺภาสิตํ อชฺฌาปนฺโน โหติ, เอกจฺเจ ภิกฺขู ทุพฺภาสิต\-ทิฏฺฐิโน โหนฺติ, เอกจฺเจ ภิกฺขู สํฆาทิเสสทิฏฺฐิโน โหนฺติ. เย เต ภิกฺขเว ภิกฺขู ทุพฺภาสิต\-ทิฏฺฐิโน, เตหิ โส ภิกฺขเว ภิกฺขุ \csromanfolio{258}เอกมนฺตํ อปเนตฺวา ยถาธมฺมํ การาเปตฺวา สํฆํ อุปสงฺกมิตฺวา เอวมสฺส วจนีโย “ยํ โข โส อาวุโส ภิกฺขุ อาปตฺติํ อาปนฺโน, สาสฺส ยถาธมฺมํ ปฏิกตา, ยทิ สํฆสฺส ปตฺตกลฺลํ, สํโฆ ปวาเรยฺยา”ติ.}

\prose{อิธ ปน ภิกฺขเว ภิกฺขุ ตทหุ ปวารณาย ทุพฺภาสิตํ อชฺฌาปนฺโน โหติ, เอกจฺเจ ภิกฺขู ทุพฺภาสิต\-ทิฏฺฐิโน โหนฺติ, เอกจฺเจ ภิกฺขู ถุลฺลจฺจย\-ทิฏฺฐิโน โหนฺติ ฯเปฯ เอกจฺเจ ภิกฺขู ทุพฺภาสิต\-ทิฏฺฐิโน โหนฺติ, เอกจฺเจ ภิกฺขู ปาจิตฺติย\-ทิฏฺฐิโน โหนฺติ. เอกจฺเจ ภิกฺขู ทุพฺภาสิต\-ทิฏฺฐิโน โหนฺติ, เอกจฺเจ ภิกฺขู ปาฏิเทสนีย\-ทิฏฺฐิโน โหนฺติ. เอกจฺเจ ภิกฺขู ทุพฺภาสิต\-ทิฏฺฐิโน โหนฺติ, เอกจฺเจ ภิกฺขู ทุกฺกฏ\-ทิฏฺฐิโน โหนฺติ. เย เต ภิกฺขเว ภิกฺขู ทุพฺภาสิต\-ทิฏฺฐิโน, เตหิ โส ภิกฺขเว ภิกฺขุ เอกมนฺตํ อปเนตฺวา ยถาธมฺมํ การาเปตฺวา สํฆํ อุปสงฺกมิตฺวา เอวมสฺส วจนีโย “ยํ โข โส อาวุโส ภิกฺขุ อาปตฺติํ อาปนฺโน, สาสฺส ยาถาธมฺมํ ปฏิกตา, ยทิ สํฆสฺส ปตฺตกลฺลํ, สํโฆ ปวาเรยฺยา”ติ.}

\csromanmatikaanchor{mtk.147}
\csromantocmarkref{section}{143. วตฺถุฐปนาทิ}{mtk.147}
\bookmark[dest={mtk.147},level=1]{143. วตฺถุฐปนาทิ}
\csromanheader{1}{143. \textbf{วตฺถุ}\-\textbf{ฐปนาทิ}}
\proseitem{239}{อิธ ปน ภิกฺขเว ภิกฺขุ ตทหุ ปวารณาย สํฆมชฺเฌ อุทาหเรยฺย “สุณาตุ เม ภนฺเต สํโฆ, อิทํ วตฺถุ ปญฺญายติ, น ปุคฺคโล, ยทิ สํฆสฺส ปตฺตกลฺลํ, วตฺถุํ ฐเปตฺวา สํโฆ ปวาเรยฺยา”ติ.}

\prose{โส เอวมสฺส วจนีโย “ภควตา โข อาวุโส วิสุทฺธานํ ปวารณา ปญฺญตฺตา, สเจ วตฺถุ ปญฺญายติ, น ปุคฺคโล, อิทาเนว นํ วเทหี”ติ.}

\prose{อิธ ปน ภิกฺขเว ภิกฺขุ ตทหุ ปวารณาย สํฆมชฺเฌ อุทาหเรยฺย “สุณาตุ เม ภนฺเต สํโฆ, อยํ ปุคฺคโล ปญฺญายติ, น วตฺถุ, ยทิ สํฆสฺส ปตฺตกลฺลํ, ปุคฺคลํ ฐเปตฺวา สํโฆ ปวาเรยฺยา”ติ.}

\prose{โส เอวมสฺส วจนีโย “ภควตา โข อาวุโส สมคฺคานํ ปวารณา ปญฺญตฺตา, สเจ ปุคฺคโล ปญฺญายติ, น วตฺถุ, อิทาเนว นํ วเทหี”ติ.}

\proseitem{239}{อิธ ปน ภิกฺขเว ภิกฺขุ ตทหุ ปวารณาย สํฆมชฺเฌ อุทาหเรยฺย “สุณาตุ เม ภนฺเต สํโฆ, อิทํ วตฺถุ จ ปุคฺคโล จ ปญฺญายติ, \csromanfolio{259}ยทิ สํฆสฺส ปตฺตกลฺลํ, วตฺถุญฺจ ปุคฺคลญฺจ ฐเปตฺวา สํโฆ ปวาเรยฺยา”ติ. โส เอวมสฺส วจนีโย “ภควตา โข อาวุโส วิสุทฺธานญฺจ สมคฺคานญฺจ ปวารณา ปญฺญตฺตา, สเจ วตฺถุ จ ปุคฺคโล จ ปญฺญายติ, อิทาเนว นํ วเทหี”ติ.}

\prose{ปุพฺเพ เจ ภิกฺขเว ปวารณาย วตฺถุ ปญฺญายติ, ปจฺฉา ปุคฺคโล, กลฺลํ วจนาย. ปุพฺเพ เจ ภิกฺขเว ปวารณาย ปุคฺคโล ปญฺญายติ, ปจฺฉา วตฺถุ, กลฺลํ วจนาย. ปุพฺเพ เจ ภิกฺขเว ปวารณาย วตฺถุ จ ปุคฺคโล จ ปญฺญายติ, ตํ เจ กตาย ปวารณาย อุกฺโกเฏติ, อุกฺโกฏนกํ ปาจิตฺติยนฺติ.}

\csromanmatikaanchor{mtk.148}
\csromantocmarkref{section}{144. ภณฺฑนการกวตฺถุ}{mtk.148}
\bookmark[dest={mtk.148},level=1]{144. ภณฺฑนการกวตฺถุ}
\csromanheader{1}{144. \textbf{ภณฺฑน}\-\textbf{การก}\-\textbf{วตฺถุ}}
\proseitem{240}{เตน โข ปน สมเยน สมฺพหุลา สนฺทิฏฺฐา สมฺภตฺตา ภิกฺขู โกสเลสุ ชนปเท อญฺญตรสฺมิํ อาวาเส วสฺสํ อุปคจฺฉิํสุ. เตสํ สามนฺตา อญฺเญ ภิกฺขู ภณฺฑนการกา กลหการกา วิวาทการกา ภสฺสการกา สํเฆ อธิกรณ\-การกา วสฺสํ อุปคจฺฉิํสุ “มยํ เตสํ ภิกฺขูนํ วสฺสํ\-วุฏฺฐานํ ปวารณาย ปวารณํ ฐเปสฺสามา”ติ. อสฺโสสุํ โข เต ภิกฺขู “อมฺหากํ กิร สามนฺตา อญฺเญ ภิกฺขู ภณฺฑนการกา กลหการกา วิวาทการกา ภสฺสการกา สํเฆ อธิกรณ\-การกา วสฺสํ อุปคตา มยํ เตสํ ภิกฺขูนํ วสฺสํ\-วุฏฺฐานํ ปวารณาย ปวารณํ ฐเปสฺสามา”ติ. กถํ นุ โข อเมฺหหิ ปฏิปชฺชิตพฺพนฺติ. ภควโต เอตมตฺถํ อาโรเจสุํ.}

\prose{อิธ ปน ภิกฺขเว สมฺพหุลา สนฺทิฏฺฐา สมฺภตฺตา ภิกฺขู อญฺญตรสฺมิํ อาวาเส วสฺสํ อุปคจฺฉนฺติ. เตสํ สามนฺตา อญฺเญ ภิกฺขู ภณฺฑนการกา กลหการกา วิวาทการกา ภสฺสการกา สํเฆ อธิกรณ\-การกา วสฺสํ อุปคจฺฉนฺติ “มยํ เตสํ ภิกฺขูนํ วสฺสํ\-วุฏฺฐานํ ปวารณาย ปวารณํ ฐเปสฺสามา”ติ. อนุชานามิ ภิกฺขเว เตหิ ภิกฺขูหิ เทฺว ตโย อุโปสเถ จาตุทฺทสิเก กาตุํ “กถํ มยํ เตหิ ภิกฺขูหิ ปฐมตรํ ปวาเรยฺยามา”ติ. เต เจ ภิกฺขเว ภิกฺขู ภณฺฑนการกา กลหการกา วิวาทการกา ภสฺสการกา สํเฆ อธิกรณ\-การกา ตํ อาวาสํ อาคจฺฉนฺติ, เตหิ ภิกฺขเว อาวาสิเกหิ ภิกฺขูหิ ลหุํ ลหุํ สนฺนิปติตฺวา ปวาเรตพฺพํ, ปวาเรตฺวา วตฺตพฺพา \csromanfolio{260}“ปวาริตา โข มยํ อาวุโส, ยถายสฺมนฺตา มญฺญนฺติ, ตถา กโรนฺตู”ติ. เต เจ ภิกฺขเว ภิกฺขู ภณฺฑนการกา กลหการกา วิวาทการกา ภสฺสการกา สํเฆ อธิกรณ\-การกา อสํวิหิตา ตํ อาวาสํ อาคจฺฉนฺติ, เตหิ ภิกฺขเว อาวาสิเกหิ ภิกฺขูหิ อาสนํ ปญฺญเปตพฺพํ, ปาโททกํ ปาทปีฐํ ปาทกถลิกํ อุปนิกฺขิปิตพฺพํ, ปจฺจุคฺคนฺตฺวา ปตฺตจีวรํ ปฏิคฺคเหตพฺพํ, ปานีเยน ปริปุจฺฉิตพฺพา, เตสํ วิกฺขิตฺวา\csromansharedfootnote{8a1b681b329577dc}{วิกฺขิปาเปตฺวา (ปฏิวิโสธกานํ มติ), อาจิกฺขิตฺวา (ก)} นิสฺสีมํ คนฺตฺวา ปวาเรตพฺพํ, ปวาเรตฺวา วตฺตพฺพา “ปวาริตา โข มยํ อาวุโส, ยถายสฺมนฺตา มญฺญนฺติ, ตถา กโรนฺตู”ติ. เอวญฺเจตํ ลเภถ, อิจฺเจตํ กุสลํ. โน เจ ลเภถ, อาวาสิเกน ภิกฺขุนา พฺยตฺเตน ปฏิพเลน อาวาสิกา ภิกฺขู ญาเปตพฺโพ\csromandash{}}

\prose{“สุณนฺตุ เม อายสฺมนฺโต\csromansharedfootnote{72ff93293253703f}{อายสฺมนฺตา (ก)} อาวาสิกา, ยทา\-ยสฺมนฺตานํ ปตฺตกลฺลํ, อิทานิ อุโปสถํ กเรยฺยาม, ปาติโมกฺขํ อุทฺทิเสยฺยาม, อาคเม กาเฬ ปวาเรยฺยามา”ติ. เต เจ ภิกฺขเว ภิกฺขู ภณฺฑนการกา กลหการกา วิวาทการกา ภสฺสการกา สํเฆ อธิกรณ\-การกา เต ภิกฺขู เอวํ วเทยฺยุํ “สาธาวุโส อิทาเนว โน ปวาเรถา”ติ. เต เอวมสฺสุ วจนียา “อนิสฺสรา โข ตุเมฺห อาวุโส อมฺหากํ ปวารณาย, น ตาว มยํ ปวาเรยฺยามา”ติ. เต เจ ภิกฺขเว ภิกฺขู ภณฺฑนการกา กลหการกา วิวาทการกา ภสฺสการกา สํเฆ อธิกรณ\-การกา ตํ กาฬํ อนุวเสยฺยุํ. อาวาสิเกน ภิกฺขเว ภิกฺขุนา พฺยตฺเตน ปฏิพเลน อาวาสิกา ภิกฺขู ญาเปตพฺพา\csromandash{}}

\prose{“สุณนฺตุ เม อายสฺมนฺโต อาวาสิกา, ยทา\-ยสฺมนฺตานํ ปตฺตกลฺลํ, อิทานิ อุโปสถํ กเรยฺยาม, ปาติโมกฺขํ อุทฺทิเสยฺยาม, อาคเม ชุเณฺห ปวาเรยฺยามา”ติ. เต เจ ภิกฺขเว ภิกฺขู ภณฺฑนการกา กลหการกา วิวาทการกา ภสฺสการกา สํเฆ อธิกรณ\-การกา เต ภิกฺขู เอวํ วเทยฺยุํ “สาธาวุโส อิทาเนว โน ปวาเรยฺยาถา”ติ. เต เอวมสฺสุ วจนียา “อนิสฺสรา โข ตุเมฺห อาวุโส อมฺหากํ ปวารณาย, น ตาว มยํ ปวาเรยฺยามา”ติ. เต เจ ภิกฺขเว ภิกฺขู ภณฺฑนการกา กลหการกา วิวาทการกา ภสฺสการกา สํเฆ \csromanfolio{261}อธิกรณ\-การกา ตมฺปิ ชุณฺหํ อนุวเสยฺยุํ. เตหิ ภิกฺขเว ภิกฺขูหิ สพฺเพเหว อาคเม ชุเณฺห โกมุทิยา จาตุมาสินิยา อกามา ปวาเรตพฺพํ.}

\prose{เตหิ เจ ภิกฺขเว ภิกฺขูหิ ปวาริยมาเน คิลาโน อคิลานสฺส ปวารณํ ฐเปติ, โส เอวมสฺส วจนีโย “อายสฺมา โข คิลาโน, คิลาโน จ อนนุโยคกฺขโม วุตฺโต ภควตา, อาคเมหิ อาวุโส ยาว อโรโค โหสิ, อโรโค อากงฺขมาโน โจเทสฺสสี”ติ. เอวญฺเจ วุจฺจมาโน โจเทติ, อนาทริเย ปาจิตฺติยํ. เตหิ เจ ภิกฺขเว ภิกฺขูหิ ปวาริยมาเน อคิลาโน คิลานสฺส ปวารณํ ฐเปติ, โส เอวมสฺส วจนีโย “อยํ โข อาวุโส ภิกฺขุ คิลาโน, คิลาโน จ อนนุโยคกฺขโม วุตฺโต ภควตา, อาคเมหิ อาวุโส ยาวายํ ภิกฺขุ อโรโค โหติ, อโรคํ อากงฺขมาโน โจเทสฺสสี”ติ. เอวญฺเจ วุจฺจมาโน โจเทติ, อนาทริเย ปาจิตฺติยํ. เตหิ เจ ภิกฺขเว ภิกฺขูหิ ปวาริยมาเน คิลาโน คิลานสฺส ปวารณํ ฐเปติ, โส เอวมสฺส วจนีโย “อายสฺมนฺตา โข คิลานา, คิลาโน จ อนนุโยคกฺขมวุตฺโต ภควตา, อาคเมหิ อาวุโส ยาว อโรคา โหถ, อโรโค อโรคํ อากงฺขมาโน โจเทสฺสสี”ติ\csromansharedfootnote{4a32fb69be5b0f5b}{ยาว อโรโค โหติ, อโรคํ อากงฺขมาโน โจเทสฺสสีติ (ก)}. เอวญฺเจ วุจฺจมาโน โจเทติ, อนาทริเย ปาจิตฺติยํ. เตหิ เจ ภิกฺขเว ภิกฺขูหิ ปวาริยมาเน อคิลาโน อคิลานสฺส ปวารณํ ฐเปติ, อุโภ สํเฆน สมนุยุญฺชิตฺวา สมนุคาหิตฺวา\csromansharedfootnote{033a486eb8053751}{สมนุภาสิตฺวา (สี)} ยถาธมฺมํ การาเปตฺวา สํเฆน ปวาเรตพฺพนฺติ.}

\csromanmatikaanchor{mtk.149}
\csromantocmarkref{section}{145. ปวารณาสงฺคห}{mtk.149}
\bookmark[dest={mtk.149},level=1]{145. ปวารณาสงฺคห}
\csromanheader{1}{145. \textbf{ปวารณา}\-\textbf{สงฺคห}}
\proseitem{241}{เตน โข ปน สมเยน สมฺพหุลา สนฺทิฏฺฐา สมฺภตฺตา ภิกฺขู โกสเลสุ ชนปเท อญฺญตรสฺมิํ อาวาเส วสฺสํ อุปคจฺฉิํสุ, เตสํ สมคฺคานํ สมฺโมทมานานํ อวิวทมานานํ วิหรตํ อญฺญตโร ผาสุวิหาโร อธิคโต โหติ. อถ โข เตสํ ภิกฺขูนํ เอตทโหสิ “อมฺหากํ โข สมคฺคานํ สมฺโมทมานานํ อวิวทมานานํ วิหรตํ อญฺญตโร ผาสุวิหาโร อธิคโต, สเจ มยํ อิทานิ \csromanfolio{262}ปวาเรสฺสาม, สิยาปิ ภิกฺขู ปวาเรตฺวา จาริกํ ปกฺกเมยฺยุํ, เอวํ มยํ อิมมฺหา ผาสุวิหารา ปริพาหิรา ภวิสฺสาม, กถํ นุ โข อเมฺหหิ ปฏิปชฺชิตพฺพนฺ”ติ. ภควโต เอตมตฺถํ อาโรเจสุํ.}

\prose{อิธ ปน ภิกฺขเว สมฺพหุลา สนฺทิฏฺฐา สมฺภตฺตา ภิกฺขู อญฺญตรสฺมิํ อาวาเส วสฺสํ อุปคจฺฉนฺติ, เตสํ สมคฺคานํ สมฺโมทมานานํ อวิวทมานานํ วิหรตํ อญฺญตโร ผาสุวิหาโร อธิคโต โหติ. ตตฺร เจ ภิกฺขูนํ เอวํ โหติ “อมฺหากํ โข สมคฺคานํ สมฺโมทมานานํ อวิวทมานานํ วิหรตํ อญฺญตโร ผาสุวิหาโร อธิคโต, สเจ มยํ อิทานิ ปวาเรสฺสาม, สิยาปิ ภิกฺขู ปวาเรตฺวา จาริกํ ปกฺกเมยฺยุํ, เอวํ มยํ อิมมฺหา ผาสุวิหารา ปริพาหิรา ภวิสฺสามา”ติ. อนุชานามิ ภิกฺขเว เตหิ ภิกฺขูหิ ปวารณา\-สงฺคหํ กาตุํ. เอวญฺจ ปน ภิกฺขเว กาตพฺโพ, สพฺเพเหว เอกชฺฌํ สนฺนิปติตพฺพํ, สนฺนิปติตฺวา พฺยตฺเตน ภิกฺขุนา ปฏิพเลน สํโฆ ญาเปตพฺโพ\csromandash{}}

\prose{“สุณาตุ เม ภนฺเต สํโฆ, อมฺหากํ สมคฺคานํ สมฺโมทมานานํ อวิวทมานานํ วิหรตํ อญฺญตโร ผาสุวิหาโร อธิคโต, สเจ มยํ อิทานิ ปวาเรสฺสาม, สิยาปิ ภิกฺขู ปวาเรตฺวา จาริกํ ปกฺกเมยฺยุํ, เอวํ มยํ อิมมฺหา ผาสุวิหารา ปริพาหิรา ภวิสฺสาม, ยทิ สํฆสฺส ปตฺตกลฺลํ, สํโฆ ปวารณา\-สงฺคหํ กเรยฺย, อิทานิ อุโปสถํ กเรยฺย, ปาติโมกฺขํ อุทฺทิเสยฺย, อาคเม ชุเณฺห โกมุทิยา จาตุมาสินิยา ปวาเรยฺย, เอสา ญตฺติ.}

\prose{สุณาตุ เม ภนฺเต สํโฆ, อมฺหากํ สมคฺคานํ สมฺโมทมานานํ อวิวทมานานํ วิหรตํ อญฺญตโร ผาสุวิหาโร อธิคโต, สเจ มยํ อิทานิ ปวาเรสฺสาม, สิยาปิ ภิกฺขู ปวาเรตฺวา จาริกํ ปกฺเกเมยฺยุํ, เอวํ มยํ อิมมฺหา ผาสุวิหารา ปริพาหิรา ภวิสฺสาม, สํโฆ ปวารณา\-สงฺคหํ กโรติ, อิทานิ อุโปสถํ กริสฺสติ, ปาติโมกฺขํ อุทฺทิสิสฺสติ, อาคเม ชุเณฺห โกมุทิยา จาตุมาสินิยา ปวาเรสฺสติ, ยสฺสายสฺมโต ขมติ ปวารณา\-สงฺคหสฺส กรณํ, อิทานิ อุโปสถํ กริสฺสติ, ปาติโมกฺขํ อุทฺทิสิสฺสติ, อาคเม ชุเณฺห โกมุทิยา จาตุมาสินิยา ปวาเรสฺสติ, โส ตุณฺหสฺส, ยสฺส นกฺขมติ, โส ภาเสยฺย.}

\csromanmatikaanchor{mtk.150}
\csromantocmarkref{section}{146. อุทฺทานคาถา}{mtk.150}
\bookmark[dest={mtk.150},level=1]{146. อุทฺทานคาถา}
\prose{\csromanfolio{263}กโต สํเฆน ปวารณา\-สงฺคโห, อิทานิ อุโปสถํ กริสฺสติ, ปาติโมกฺขํ อุทฺทิสิสฺสติ, อาคเม ชุเณฺห โกมุทิยา จาตุมาสินิยา ปวาเรสฺสติ, ขมติ สํฆสฺส, ตสฺมา ตุณฺหี, เอวเมตํ ธารยามี”ติ.}

\prose{เตหิ เจ ภิกฺขเว ภิกฺขูหิ กเต ปวารณา\-สงฺคเห อญฺญตโร ภิกฺขุ เอวํ วเทยฺย “อิจฺฉามหํ อาวุโส ชนปท\-จาริกํ ปกฺกมิตุํ, อตฺถิ เม ชนปเท กรณียนฺ”ติ.}

\prose{โส เอวมสฺส วจนีโย “สาธาวุโส ปวาเรตฺวา คจฺฉาหี”ติ. โส เจ ภิกฺขเว ภิกฺขุ ปวารยมาโน อญฺญตรสฺส ภิกฺขุโน ปวารณํ ฐเปติ, โส เอวมสฺส วจนีโย “อนิสฺสโร โข เม ตฺวํ อาวุโส ปวารณาย, น ตาวาหํ ปวาเรสฺสามี”ติ. ตสฺส เจ ภิกฺขเว ภิกฺขุโน ปวารย\-มานสฺส อญฺญตโร ภิกฺขุ ตสฺส ภิกฺขุโน ปวารณํ ฐเปติ, อุโภ สํเฆน สมนุยุญฺชิตฺวา สมนุคาหิตฺวา ยถาธมฺมํ การาเปตพฺพา. โส เจ ภิกฺขเว ภิกฺขุ ชนปเท ตํ กรณียํ ตีเรตฺวา ปุนเทว อนฺโต โกมุทิยา จาตุมาสินิยา ตํ อาวาสํ อาคจฺฉติ, เตหิ เจ ภิกฺขเว ภิกฺขูหิ ปวาริยมาเน อญฺญตโร ภิกฺขุ ตสฺส ภิกฺขุโน ปวารณํ ฐเปติ, โส เอวมสฺส วจนีโย “อนิสฺสโร โข เม ตฺวํ อาวุโส ปวารณาย, ปวาริโต อหนฺ”ติ. เตหิ เจ ภิกฺขเว ภิกฺขูหิ ปวาริยมาเน โส ภิกฺขุ อญฺญตรสฺส ภิกฺขุโน ปวารณํ ฐเปติ, อุโภ สํเฆน สมนุยุญฺชิตฺวา สมนุคาหิตฺวา ยถาธมฺมํ การาเปตฺวา สํเฆน ปวาเรตพฺพนฺติ.}

\vspace{\tipitakaparskip}
\csromancenter{ปวารณา\-กฺขนฺธโก จตุตฺโถ.}
\csromancenter{146. \textbf{ตสฺสุทฺทานํ}}
\setlength{\csromangathapagewidth}{0pt}
\setlength{\csromangathawakwidth}{0pt}
\csromangathameasuregroup{วสฺสํวุฏฺฐา โกสเลสุ, \\ อผาสุํ ปสุสํวาสํ, \\ ปวาเรนฺตา ปณามญฺจ\textsuperscript{1}, \\ ราชา โจรา จ ธุตฺตา จ, \\ ปญฺจ จตุตโย เทฺวโก, \\ สพฺโพ สํโฆ เวมติโก, \\ อาวาสิกา จาตุทฺทส, \\ คนฺตพฺพํ น นิสินฺนาย, \\ สวเรหิ เขปิตา เมโฆ, \\ น อิจฺฉนฺติ ปุรมฺหากํ, \\ กิมฺหิ วาติ กตมญฺจ, \\ โจทโก จุทิตโก จ, \\ ปวารณาสงฺคโห จ,}{\csromangathabat{วสฺสํวุฏฺฐา โกสเลสุ,}{อคมุํ สตฺถุ ทสฺสนํ.} \\ \csromangathabat{อผาสุํ ปสุสํวาสํ,}{อญฺญมญฺญานุโลมตา.} \\ \csromangathabat{ปวาเรนฺตา ปณามญฺจ\textsuperscript{1},}{กมฺมํ คิลานญาตกา.} \\ \csromangathabat{ราชา โจรา จ ธุตฺตา จ,}{ภิกฺขุปจฺจตฺถิกา ตถา.} \\ \csromangathabat{ปญฺจ จตุตโย เทฺวโก,}{อาปนฺโน เวมตี สริ.} \\ \csromangathabat{สพฺโพ สํโฆ เวมติโก,}{พหู สมา จ โถกิกา.} \\ \csromangathabat{อาวาสิกา จาตุทฺทส,}{ลิงฺคสํวาสกา อุโภ.} \\ \csromangathabat{คนฺตพฺพํ น นิสินฺนาย,}{ฉนฺททาเน ปวารณา\textsuperscript{1}.} \\ \csromangathabat{สวเรหิ เขปิตา เมโฆ,}{อนฺตรา จ ปวารณา.} \\ \csromangathabat{น อิจฺฉนฺติ ปุรมฺหากํ,}{อฏฺฐปิตา จ ภิกฺขุโน.} \\ \csromangathabat{กิมฺหิ วาติ กตมญฺจ,}{ทิฏฺเฐน สุตสงฺกาย.} \\ \csromangathabat{โจทโก จุทิตโก จ,}{ถุลฺลจฺจยํ วตฺถุ ภณฺฑนํ.} \\ \csromangathabat{ปวารณาสงฺคโห จ,}{อนิสฺสโร ปวารเยติ.} \\ อิมมฺหิ ขนฺธเก วตฺถูนิ ฉจตฺตารีสาติ.}
\csromangathasetleft{วสฺสํวุฏฺฐา โกสเลสุ, \\ อผาสุํ ปสุสํวาสํ, \\ ปวาเรนฺตา ปณามญฺจ\textsuperscript{1}, \\ ราชา โจรา จ ธุตฺตา จ, \\ ปญฺจ จตุตโย เทฺวโก, \\ สพฺโพ สํโฆ เวมติโก, \\ อาวาสิกา จาตุทฺทส, \\ คนฺตพฺพํ น นิสินฺนาย, \\ สวเรหิ เขปิตา เมโฆ, \\ น อิจฺฉนฺติ ปุรมฺหากํ, \\ กิมฺหิ วาติ กตมญฺจ, \\ โจทโก จุทิตโก จ, \\ ปวารณาสงฺคโห จ,}
\csromangathagroupwithclosermajor{\csromangathastanza{\csromangathabat{วสฺสํวุฏฺฐา โกสเลสุ,}{อคมุํ สตฺถุ ทสฺสนํ.} \\ \csromangathabat{อผาสุํ ปสุสํวาสํ,}{อญฺญมญฺญานุโลมตา.}}\csromangathastanza{\csromangathabat{ปวาเรนฺตา ปณามญฺจ\csromansharedfootnote{b0db9f5d61c351d3}{ปวาเรนฺตาสเน เทฺว จ (สี, สฺยา)},}{กมฺมํ คิลานญาตกา.} \\ \csromangathabat{ราชา โจรา จ ธุตฺตา จ,}{ภิกฺขุ\-ปจฺจตฺถิกา ตถา.}}\csromangathastanza{\csromangathabat{ปญฺจ จตุตโย เทฺวโก,}{อาปนฺโน เวมตี สริ.} \\ \csromangathabat{สพฺโพ สํโฆ เวมติโก,}{พหู สมา จ โถกิกา.}}\csromangathastanza{\csromanfolio{264}\csromangathabat{อาวาสิกา จาตุทฺทส,}{ลิงฺคสํวาสกา อุโภ.} \\ \csromangathabat{คนฺตพฺพํ น นิสินฺนาย,}{ฉนฺททาเน ปวารณา\csromansharedfootnote{c8bc4722a6bb590b}{ฉนฺททานปวารณา (ก)}.}}\csromangathastanza{\csromangathabat{สวเรหิ เขปิตา เมโฆ,}{อนฺตรา จ ปวารณา.} \\ \csromangathabat{น อิจฺฉนฺติ ปุรมฺหากํ,}{อฏฺฐปิตา จ ภิกฺขุโน.}}\csromangathastanza{\csromangathabat{กิมฺหิ วาติ กตมญฺจ,}{ทิฏฺเฐน สุตสงฺกาย.} \\ \csromangathabat{โจทโก จุทิตโก จ,}{ถุลฺลจฺจยํ วตฺถุ ภณฺฑนํ.}}\csromangathastanza{\csromangathabat{ปวารณา\-สงฺคโห จ,}{อนิสฺสโร ปวารเยติ.} \\ อิมมฺหิ ขนฺธเก วตฺถูนิ ฉจตฺตารีสาติ.}}{\textbf{ปวารณา}\-\textbf{กฺขนฺธโก นิฏฺฐิโต.}}
\csromanmatikaanchor{mtk.151}
\csromantocmarknopage{chapter}{5. จมฺมกฺขนฺธก}{mtk.151}
\bookmark[dest={mtk.151},level=0]{5. จมฺมกฺขนฺธก}
\chapterheadpageread{5. \textbf{จมฺม}\-\textbf{กฺขนฺธก}}
\csromanmatikaanchor{mtk.152}
\csromantocmarkref{section}{147. โสณโกฬิวิสวตฺถุ}{mtk.152}
\bookmark[dest={mtk.152},level=1]{147. โสณโกฬิวิสวตฺถุ}
\csromanheader{1}{147. \textbf{โสณ}\-\textbf{โกฬิวิส}\-\textbf{วตฺถุ}}
\proseitem{242}{\csromanfolio{265}เตน สมเยน พุทฺโธ ภควา ราชคเห วิหรติ คิชฺฌกูเฏ ปพฺพเต. เตน โข ปน สมเยน ราชา มาคโธ เสนิโย พิมฺพิสาโร อสีติยา คามสหสฺเสสุ อิสฺสริยาธิปจฺจํ รชฺชํ กาเรติ. เตน โข ปน สมเยน จมฺปายํ โสโณ นาม โกฬิวิโส\csromansharedfootnote{efaf8a9e4ccb8422}{โกฬิวีโส (สี)} เสฏฺฐิปุตฺโต สุขุมาโล โหติ, ตสฺส ปาทตเลสุ โลมานิ ชาตานิ โหนฺติ. อถ โข ราชา มาคโธ เสนิโย พิมฺพิสาโร ตานิ อสีติ คามิก\-สหสฺสานิ สนฺนิปาตาเปตฺวา เกนจิเทว กรณีเยน โสณสฺส โกฬิวิสสฺส สนฺติเก ทูตํ ปาเหสิ “อาคจฺฉตุ โสโณ, อิจฺฉามิ โสณสฺส อาคตนฺ”ติ. อถ โข โสณสฺส โกฬิวิสสฺส มาตาปิตโร โสณํ โกฬิวิสํ เอตทโวจุํ “ราชา เต ตาต โสณ ปาเท ทกฺขิตุกาโม, มา โข ตฺวํ ตาต โสณ เยน ราชา เตน ปาเท อภิปฺปสาเรยฺยาสิ, รญฺโญ ปุรโต ปลฺลงฺเกน นิสีท, นิสินฺนสฺส เต ราชา ปาเท ทกฺขิสฺสตี”ติ. อถ โข โสณํ โกฬิวิสํ สิวิกาย อาเนสุํ. อถ โข โสโณ โกฬิวิโส เยน ราชา มาคโธ เสนิโย พิมฺพิสาโร เตนุปสงฺกมิ, อุปสงฺกมิตฺวา ราชานํ มาคธํ เสนิยํ พิมฺพิสารํ อภิวาเทตฺวา รญฺโญ ปุรโต ปลฺลงฺเกน นิสีทิ. อทฺทสา โข ราชา มาคโธ เสนิโย พิมฺพิสาโร โสณสฺส โกฬิวิสสฺส ปาทตเลสุ โลมานิ ชาตานิ. อถ โข ราชา มาคโธ เสนิโย พิมฺพิสาโร ตานิ อสีติ คามิก\-สหสฺสานิ ทิฏฺฐธมฺมิเก อตฺเถ อนุสาสิตฺวา อุยฺโยเชสิ “ตุเมฺห ขฺวตฺถ ภเณ มยา ทิฏฺฐธมฺมิเก อตฺเถ อนุสาสิตา, คจฺฉถ ตํ ภควนฺตํ ปยิรุปาสถ, โส โน ภควา สมฺปรายิเก อตฺเถ อนุสาสิสฺสตี”ติ.}

\prose{อถ โข ตานิ อสีติ คามิก\-สหสฺสานิ เยน คิชฺฌกูโฏ ปพฺพโต เตนุ\-ปสงฺกมิํสุ. เตน โข ปน สมเยน อายสฺมา สาคโต ภควโต อุปฏฺฐาโก โหติ. อถ โข ตานิ อสีติ คามิก\-สหสฺสานิ เยนายสฺมา สาคโต เตนุ\-ปสงฺกมิํสุ, อุปสงฺกมิตฺวา \csromanfolio{266}อายสฺมนฺตํ สาคตํ เอตทโวจุํ “อิมานิ ภนฺเต อสีติ คามิก\-สหสฺสานิ อิธูปสงฺกนฺตานิ ภควนฺตํ ทสฺสนาย, สาธุ มยํ ภนฺเต ลเภยฺยาม ภควนฺตํ ทสฺสนายา”ติ. เตน หิ ตุเมฺห อายสฺมนฺโต มุหุตฺตํ อิเธว ตาว โหถ, ยาวาหํ ภควนฺตํ ปฏิเวเทมีติ. อถ โข อายสฺมา สาคโต เตสํ อสีติยา คามิก\-สหสฺสานํ ปุรโต เปกฺขมานานํ ปาฏิกาย นิมุชฺชิตฺวา ภควโต ปุรโต อุมฺมุชฺชิตฺวา ภควนฺตํ เอตทโวจ “อิมานิ ภนฺเต อสีติ คามิก\-สหสฺสานิ อิธูปสงฺกนฺตานิ ภควนฺตํ ทสฺสนาย, ยสฺส ทานิ ภนฺเต ภควา กาลํ มญฺญตี”ติ. เตน หิ ตฺวํ สาคต วิหาร\-ปจฺฉายายํ อาสนํ ปญฺญเปหีติ. “เอวํ ภนฺเต”ติ โข อายสฺมา สาคโต ภควโต ปฏิสฺสุณิตฺวา ปีฐํ คเหตฺวา ภควโต ปุรโต นิมุชฺชิตฺวา เตสํ อสีติยา คามิก\-สหสฺสานํ ปุรโต เปกฺขมานานํ ปาฏิกาย อุมฺมุชฺชิตฺวา วิหาร\-ปจฺฉายายํ อาสนํ ปญฺญเปติ. อถ โข ภควา วิหารา นิกฺขมิตฺวา วิหาร\-ปจฺฉายายํ ปญฺญตฺเต อาสเน นิสีทิ. อถ โข ตานิ อสีติ คามิก\-สหสฺสานิ เยน ภควา เตนุ\-ปสงฺกมิํสุ, อุปสงฺกมิตฺวา ภควนฺตํ อภิวาเทตฺวา เอกมนฺตํ นิสีทิํสุ. อถ โข ตานิ อสีติ คามิก\-สหสฺสานิ อายสฺมนฺตํเยว สาคตํ สมนฺนาหรนฺติ, โน ตถา ภควนฺตํ. อถ โข ภควา เตสํ อสีติยา คามิก\-สหสฺสานํ เจตสา เจโต\-ปริวิตกฺกม\-ญฺญาย อายสฺมนฺตํ สาคตํ อามนฺเตสิ “เตน หิ ตฺวํ สาคต ภิยฺโยโส\-มตฺตาย\csromansharedfootnote{ab67f8da3c3dc98c}{ภิยฺโยโส มตฺตาย เทฺว ปทานิ \csromandash{} ม.พ.ป.} อุตฺตริ\-มนุสฺสธมฺมํ อิทฺธิ\-ปาฏิหาริยํ ทสฺเสหี”ติ. “เอวํ ภนฺเต”ติ โข อายสฺมา สาคโต ภควโต ปฏิสฺสุณิตฺวา เวหาสํ อพฺภุคฺคนฺตฺวา อากาเส อนฺตลิกฺเข จงฺกมติปิ ติฏฺฐติปิ นิสีทติปิ เสยฺยมฺปิ กปฺเปติ ธูมายติปิ\csromansharedfootnote{8a9d3cc4e4689b18}{ธูปายติปิ (สี), ปธูปายติปิ (สฺยา)} ปชฺชลติปิ อนฺตร\-ธายติปิ. อถ โข อายสฺมา สาคโต อากาเส อนฺตลิกฺเข อเนกวิหิตํ อุตฺตริ\-มนุสฺสธมฺมํ อิทฺธิ\-ปาฏิหาริยํ ทสฺเสตฺวา ภควโต ปาเทสุ สิรสา นิปติตฺวา ภควนฺตํ เอตทโวจ “สตฺถา เม ภนฺเต ภควา, สาวโกหมสฺมิ, สตฺถา เม ภนฺเต ภควา, สาวโกหมสฺมี”ติ. อถ โข ตานิ อสีติ คามิก\-สหสฺสานิ “อจฺฉริยํ วต โภ, อพฺภุตํ วต โภ, สาวโกปิ นาม เอวํ มหิทฺธิโก ภวิสฺสติ เอวํ มหานุภาโว, อโห นูน สตฺถา”ติ ภควนฺตํเยว สมนฺนาหรนฺติ, โน ตถา อายสฺมนฺตํ สาคตํ.}

\prose{\csromanfolio{267}อถ โข ภควา เตสํ อสีติยา คามิก\-สหสฺสานํ เจตสา เจโต\-ปริวิตกฺกม\-ญฺญาย อนุปุพฺพิํ กถํ กเถสิ, เสยฺยถิทํ, ทานกถํ สีลกถํ สคฺคกถํ กามานํ อาทีนวํ โอการํ สํกิเลสํ เนกฺขมฺเม อานิสํสํ ปกาเสสิ. ยทา เต ภควา อญฺญาสิ กลฺลจิตฺเต มุทุจิตฺเต วินีวรณ\-จิตฺเต อุทคฺคจิตฺเต ปสนฺนจิตฺเต, อถ ยา พุทฺธานํ สามุกฺกํสิกา ธมฺมเทสนา, ตํ ปกาเสสิ ทุกฺขํ สมุทยํ นิโรธํ มคฺคํ, เสยฺยถาปิ นาม สุทฺธํ วตฺถํ อปคตกาฬกํ สมฺมเทว รชนํ ปฏิคฺคเณฺหยฺย, เอวเมวํ เตสํ อสีติยา คามิก\-สหสฺสานํ ตสฺมิํเยว อาสเน วิรชํ วีตมลํ ธมฺมจกฺขุํ อุทปาทิ “ยํกิญฺจิ สมุทย\-ธมฺมํ, สพฺพํ ตํ นิโรธธมฺมนฺ”ติ. เต ทิฏฺฐธมฺมา ปตฺตธมฺมา วิทิตธมฺมา ปริโยคาฬฺห\-ธมฺมา ติณฺณ\-วิจิกิจฺฉา วิคต\-กถํกถา เวสารชฺช\-ปฺปตฺตา อปรปฺปจฺจยา สตฺถุสาสเน ภควนฺตํ เอตทโวจุํ “อภิกฺกนฺตํ ภนฺเต อภิกฺกนฺตํ ภนฺเต, เสยฺยถาปิ ภนฺเต นิกฺกุชฺชิตํ วา อุกฺกุชฺเชยฺย, ปฏิจฺฉนฺนํ วา วิวเรยฺย, มูฬฺหสฺส วา มคฺคํ อาจิกฺเขยฺย, อนฺธกาเร วา เตลปชฺโชตํ ธาเรยฺย ‘จกฺขุมนฺโต รูปานิ ทกฺขนฺตี’ติ, เอวเมวํ ภควตา อเนก\-ปริยาเยน ธมฺโม ปกาสิโต, เอเต มยํ ภนฺเต ภควนฺตํ สรณํ คจฺฉาม ธมฺมญฺจ ภิกฺขุสํฆญฺจ, อุปาสเก โน ภควา ธาเรตุ อชฺชตคฺเค ปาณุเปเต สรณํ คเต”ติ.}

\csromanheader{2}{\textbf{โสณสฺส ปพฺพชฺชา}}
\proseitem{243}{อถ โข โสณสฺส โกฬิวิสสฺส เอตทโหสิ “ยถา ยถา โข อหํ ภควตา ธมฺมํ เทสิตํ อาชานามิ, นยิทํ สุกรํ อคารํ อชฺฌาวสตา เอกนฺต\-ปริปุณฺณํ เอกนฺต\-ปริสุทฺธํ สงฺขลิขิตํ พฺรหฺมจริยํ จริตุํ, ยนฺนูนาหํ เกสมสฺสุํ โอหาเรตฺวา กาสายานิ วตฺถานิ อจฺฉาเทตฺวา อคารสฺมา อนคาริยํ ปพฺพเชยฺยนฺ”ติ. อถ โข ตานิ อสีติ คามิก\-สหสฺสานิ ภควโต ภาสิตํ อภินนฺทิตฺวา อนุโมทิตฺวา อุฏฺฐายาสนา ภควนฺตํ อภิวาเทตฺวา ปทกฺขิณํ กตฺวา ปกฺกมิํสุ. อถ โข โสโณ โกฬิวิโส อจิร\-ปกฺกนฺเตสุ เตสุ อสีติยา คามิก\-สหสฺเสสุ เยน ภควา เตนุปสงฺกมิ, อุปสงฺกมิตฺวา ภควนฺตํ อภิวาเทตฺวา เอกมนฺตํ นิสีทิ, เอกมนฺตํ นิสินฺโน โข โสโณ โกฬิวิโส ภควนฺตํ เอตทโวจ “ยถา ยถาหํ ภนฺเต ภควตา ธมฺมํ เทสิตํ อาชานามิ, นยิทํ สุกรํ อคารํ อชฺฌาวสตา เอกนฺต\-ปริปุณฺณํ เอกนฺต\-ปริสุทฺธํ สงฺขลิขิตํ พฺรหฺมจริยํ จริตุํ, อิจฺฉามหํ ภนฺเต เกสมสฺสุํ \csromanfolio{268}โอหาเรตฺวา กาสายานิ วตฺถานิ อจฺฉาเทตฺวา อคารสฺมา อนคาริยํ ปพฺพชฺชิตุํ, ปพฺพาเชตุ มํ ภนฺเต ภควา”ติ. อลตฺถ โข โสโณ โกฬิวิโส ภควโต สนฺติเก ปพฺพชฺชํ, อลตฺถ อุปสมฺปทํ. อจิรุ\-ปสมฺปนฺโน จ ปนายสฺมา โสโณ สีตวเน วิหรติ, ตสฺส อจฺจารทฺธ\-วีริยสฺส จงฺกมโต ปาทา ภิชฺชิํสุ, จงฺกโม โลหิเตน ผุโฏ โหติ เสยฺยถาปิ ควาฆาตนํ.\csromansymbolfootnote{*}{อิโต ปรํ ยาว อิมสฺส วตฺถุสฺส อวสานํ ตาว ปาโฐ อํ 2. 329 ปิฏฺฐาทีสุปิ อาคโต.}อถ โข อายสฺมโต โสณสฺส รโหคตสฺส ปฏิสลฺลีนสฺส เอวํ เจตโส ปริวิตกฺโก อุทปาทิ “เย โข เกจิ ภควโต สาวกา อารทฺธวีริยา วิหรนฺติ, อหํ เตสํ อญฺญตโร, อถ จ ปน เม นานุปาทาย อาสเวหิ จิตฺตํ วิมุจฺจติ, สํวิชฺชนฺติ โข ปน เม กุเล โภคา, สกฺกา โภเค จ ภุญฺชิตุํ ปุญฺญานิ จ กาตุํ, ยนฺนูนาหํ หีนายาวตฺติตฺวา โภเค จ ภุญฺเชยฺยํ, ปุญฺญานิ จ กเรยฺยนฺ”ติ. อถ โข ภควา อายสฺมโต โสณสฺส เจตสา เจโต\-ปริวิตกฺกม\-ญฺญาย เสยฺยถาปิ นาม พลวา ปุริโส สมิญฺชิตํ วา พาหํ ปสาเรยฺย, ปสาริตํ วา พาหํ สมิญฺเชยฺย, เอวเมว คิชฺฌกูเฏ ปพฺพเต อนฺตรหิโต สีตวเน ปาตุรโหสิ. อถ โข ภควา สมฺพหุเลหิ ภิกฺขูหิ สทฺธิํ เสนาสนจาริกํ อาหิณฺฑนฺโต เยนายสฺมโต โสณสฺส จงฺกโม เตนุปสงฺกมิ. อทฺทสา โข ภควา อายสฺมโต โสณสฺส จงฺกมํ โลหิเตน ผุฏํ, ทิสฺวาน ภิกฺขู อามนฺเตสิ “กสฺส นฺวายํ ภิกฺขเว จงฺกโม โลหิเตน ผุโฏ เสยฺยถาปิ ควาฆาตนนฺ”ติ. อายาสฺมโต ภนฺเต โสณสฺส อจฺจารทฺธ\-วีริยสฺส จงฺกมโต ปาทา ภิชฺชิํสุ, ตสฺสายํ จงฺกโม โลหิเตน ผุโฏ เสยฺยถาปิ ควาฆาตนนฺติ.}

\prose{อถ โข ภควา เยนายสฺมโต โสณสฺส วิหาโร เตนุปสงฺกมิ, อุปสงฺกมิตฺวา ปญฺญตฺเต อาสเน นิสีทิ. อายสฺมาปิ โข โสโณ ภควนฺตํ อภิวาเทตฺวา เอกมนฺตํ นิสีทิ, เอกมนฺตํ นิสินฺนํ โข อายสฺมนฺตํ โสณํ ภควา เอตทโวจ “นนุ เต โสณ รโหคตสฺส ปฏิสลฺลีนสฺส เอวํ เจตโส ปริวิตกฺโก อุทปาทิ ‘เย โข เกจิ ภควโต สาวกา อารทฺธวีริยา วิหรนฺติ, อหํ เตสํ อญฺญตโร, อถ จ ปน เม นานุปาทาย อาสเวหิ จิตฺตํ วิมุจฺจติ, สํวิชฺชนฺติ โข ปน เม กุเล โภคา, สกฺกา โภเค จ ภุญฺชิตุํ, ปุญฺญานิ จ กาตุํ, ยนฺนูนาหํ หีนายาวตฺติตฺวา โภเค จ ภุญฺเชยฺยํ, ปุญฺญานิ จ กเรยฺยนฺ’ติ”. เอวํ \csromanfolio{269}ภนฺเตติ. ตํ กิํ มญฺญสิ โสณ กุสโล ตฺวํ ปุพฺเพ อคาริกภูโต วีณาย ตนฺติสฺสเรติ. เอวํ ภนฺเตติ. ตํ กิํ มญฺญสิ โสณ ยทา เต วีณาย ตนฺติโย อจฺจายตา โหนฺติ, อปิ นุ เต วีณา ตสฺมิํ สมเย สรวตี วา โหติ กมฺมญฺญา วาติ. โนเหตํ ภนฺเตติ. ตํ กิํ มญฺญสิ โสณ ยทา เต วีณาย ตนฺติโย อติสิถิลา โหนฺติ, อปิ นุ เต วีณา ตสฺมิํ สมเย สรวตี วา โหติ กมฺมญฺญา วาติ. โนเหตํ ภนฺเตติ. ตํ กิํ มญฺญสิ โสณ ยทา เต วีณาย ตนฺติโย เนว อจฺจายตา โหนฺติ นาติสิถิลา สเม คุเณ ปติฏฺฐิตา, อปิ นุ เต วีณา ตสฺมิํ สมเย สรวตี วา โหติ กมฺมญฺญา วาติ. เอวํ ภนฺเตติ. เอวเมว โข โสณ อจฺจารทฺธ\-วีริยํ อุทฺธจฺจาย สํวตฺตติ, อติลีนวีริยํ โกสชฺชาย สํวตฺตติ, ตสฺมาติห ตฺวํ โสณ วีริยสมตํ อธิฏฺฐห, อินฺทฺริยานญฺจ สมตํ ปฏิวิชฺฌ, ตตฺถ จ นิมิตฺตํ คณฺหาหีติ. เอวํ ภนฺเตติ โข อายสฺมา โสโณ ภควโต ปจฺจสฺโสสิ. อถ โข ภควา อายสฺมนฺตํ โสณํ อิมินา โอวาเทน โอวทิตฺวา เสยฺยถาปิ นาม พลวา ปุริโส สมิญฺชิตํ วา พาหํ ปสาเรยฺย, ปสาริตํ วา พาหํ สมิญฺเชยฺย, เอวเมว สีตวเน อายสฺมโต โสณสฺส สมฺมุเข อนฺตรหิโต คิชฺฌกูเฏ ปพฺพเต ปาตุรโหสิ. อถ โข อายสฺมา โสโณ อปเรน สมเยน วีริยสมตํ อธิฏฺฐาสิ, อินฺทฺริยานญฺจ สมตํ ปฏิวิชฺฌิ, ตตฺถ จ นิมิตฺตํ อคฺคเหสิ. อถ โข อายสฺมา โสโณ เอโก วูปกฏฺโฐ อปฺปมตฺโต อาตาปี ปหิตตฺโต วิหรนฺโต น จิรสฺเสว ‘ยสฺสตฺถาย กุลปุตฺตา สมฺมเทว อคารสฺมา อนคาริยํ ปพฺพชนฺติ’ ตทนุตฺตรํ พฺรหฺมจริย\-ปริโยสานํ ทิฏฺเฐว ธมฺเม สยํ อภิญฺญา สจฺฉิกตฺวา อุปสมฺปชฺช วิหาสิ, “ขีณา ชาติ วุสิตํ พฺรหฺมจริยํ กตํ กรณียํ, นาปรํ อิตฺถตฺตายา”ติ อภิญฺญาสิ, อญฺญตโร จ ปนายสฺมา โสโณ อรหตํ อโหสิ.}

\proseitem{244}{อถ โข อายสฺมโต โสณสฺส อรหตฺต\-ปฺปตฺตสฺส เอตทโหสิ “ยนฺนูนาหํ ภควโต สนฺติเก อญฺญํ พฺยากเรยฺยนฺ”ติ. อถ โข อายสฺมา โสโณ เยน ภควา เตนุปสงฺกมิ, อุปสงฺกมิตฺวา ภควนฺตํ อภิวาเทตฺวา เอกมนฺตํ นิสีทิ, เอกมนฺตํ นิสินฺโน โข อายสฺมา โสโณ ภควนฺตํ เอตทโวจ “โย โส ภนฺเต ภิกฺขุ อรหํ ขีณาสโว วุสิตวา กตกรณีโย โอหิตภาโร อนุปฺปตฺต\-สทตฺโถ \csromanfolio{270}ปริกฺขีณ\-ภว\-สญฺโญชโน สมฺมทญฺญา วิมุตฺโต, โส ฉฏฺฐานานิ อธิมุตฺโต โหติ, เนกฺขมฺมา\-ธิมุตฺโต โหติ, ปวิเวกา\-ธิมุตฺโต โหติ, อพฺยาปชฺชา\-ธิมุตฺโต โหติ, อุปาทานกฺขยา\-ธิมุตฺโต โหติ, ตณฺห\-กฺขยา\-ธิมุตฺโต โหติ, อสมฺโมหา\-ธิมุตฺโต โหติ.}

\prose{สิยา โข ปน ภนฺเต อิเธกจฺจสฺส อายสฺมโต เอวมสฺส ‘เกวลํ สทฺธามตฺตกํ นูน อยมายสฺมา นิสฺสาย เนกฺขมฺมา\-ธิมุตฺโต’ติ, น โข ปเนตํ ภนฺเต เอวํ ทฏฺฐพฺพํ, ขีณาสโว ภนฺเต ภิกฺขุ วุสิตวา กตกรณีโย กรณีย\-มตฺตานํ\csromansharedfootnote{84c3f78508d1c9dd}{กรณียํ อตฺตโน (องฺคุตฺตรปาฬิยํ)} อสมนุปสฺสนฺโต กตสฺส วา ปฏิจยํ ขยา ราคสฺส วีตราคตฺตา เนกฺขมฺมา\-ธิมุตฺโต โหติ, ขยา โทสสฺส วีตโทสตฺตา เนกฺขมฺมา\-ธิมุตฺโต โหติ, ขยา โมหสฺส วีตโมหตฺตา เนกฺขมฺมา\-ธิมุตฺโต โหติ.}

\prose{สิยา โข ปน ภนฺเต อิเธกจฺจสฺส อายสฺมโต เอวมสฺส ‘ลาภ\-สกฺการ\-สิโลกํ นูน อยมายสฺมา นิกามยมาโน ปวิเวกา\-ธิมุตฺโต’ติ, น โข ปเนตํ ภนฺเต เอวํ ทฏฺฐพฺพํ, ขีณาสโว ภนฺเต ภิกฺขุ วุสิตวา กตกรณีโย กรณีย\-มตฺตานํ อสมนุปสฺสนฺโต กตสฺส วา ปฏิจยํ ขยา ราคสฺส วีตราคตฺตา ปวิเวกา\-ธิมุตฺโต โหติ, ขยา โทสสฺส วีตโทสตฺตา ปวิเวกา\-ธิมุตฺโต โหติ, ขยา โมหสฺส วีตโมหตฺตา ปวิเวกา\-ธิมุตฺโต โหติ.}

\prose{สิยา โข ปน ภนฺเต อิเธกจฺจสฺส อายสฺมโต เอวมสฺส ‘สีลพฺพต\-ปรามาสํ นูน อยมายสฺมา สารโต ปจฺจาคจฺฉนฺโต อพฺยาปชฺชา\-ธิมุตฺโต’ติ, น โข ปเนตํ ภนฺเต เอวํ ทฏฺฐพฺพํ, ขีณาสโว ภนฺเต ภิกฺขุ วุสิตวา กตกรณีโย กรณีย\-มตฺตานํ อสมนุปสฺสนฺโต กตสฺส วา ปฏิจยํ ขยา ราคสฺส วีตราคตฺตา อพฺยาปชฺชา\-ธิมุตฺโต โหติ, ขยา โทสสฺส วีตโทสตฺตา อพฺยาปชฺชา\-ธิมุตฺโต โหติ, ขยา โมหสฺส วีตโมหตฺตา อพฺยาปชฺชา\-ธิมุตฺโต โหติ.}

\prose{ขยา ราคสฺส วีตราคตฺตา อุปาทานกฺขยา\-ธิมุตฺโต โหติ, ขยา โทสสฺส วีตโทสตฺตา อุปาทานกฺขยา\-ธิมุตฺโต โหติ, ขยา โมหสฺส วีตโมหตฺตา อุปาทานกฺขยา\-ธิมุตฺโต โหติ.}

\prose{\csromanfolio{271}ขยา ราคสฺส วีตราคตฺตา ตณฺห\-กฺขยา\-ธิมุตฺโต โหติ, ขยา โทสสฺส วีตโทสตฺตา ตณฺห\-กฺขยา\-ธิมุตฺโต โหติ, ขยา โมหสฺส วีตโมหตฺตา ตณฺห\-กฺขยา\-ธิมุตฺโต โหติ.}

\prose{ขยา ราคสฺส วีตราคตฺตา อสมฺโมหา\-ธิมุตฺโต โหติ, ขยา โทสสฺส วีตโทสตฺตา อสมฺโมหา\-ธิมุตฺโต โหติ, ขยา โมหสฺส วีตโมหตฺตา อสมฺโมหา\-ธิมุตฺโต โหติ.}

\prose{เอวํ สมฺมา วิมุตฺต\-จิตฺตสฺส ภนฺเต ภิกฺขุโน ภุสา เจปิ จกฺขุวิญฺเญยฺยา รูปา จกฺขุสฺส อาปาถํ อาคจฺฉนฺติ, เนวสฺส จิตฺตํ ปริยาทิยนฺติ, อมิสฺสีกตเม\-วสฺส จิตฺตํ โหติ, ฐิตํ อาเนญฺชปฺปตฺตํ, วยญฺจ\-สฺสา\-นุปสฺสติ.}

\prose{ภุสา เจปิ โสตวิญฺเญยฺยา สทฺทา ฯเปฯ ฆานวิญฺเญยฺยา คนฺธา. ชิวฺหาวิญฺเญยฺยา รสา. กายวิญฺเญยฺยา โผฏฺฐพฺพา. มโนวิญฺเญยฺยา ธมฺมา มนสฺส อาปาถํ อาคจฺฉนฺติ, เนวสฺส จิตฺตํ ปริยาทิยนฺติ, อมิสฺสีกตเม\-วสฺส จิตฺตํ โหติ, ฐิตํ อาเนญฺชปฺปตฺตํ, วยญฺจ\-สฺสา\-นุปสฺสติ, เสยฺยถาปิ ภนฺเต เสโล ปพฺพโต อจฺฉิทฺโท อสุสิโร เอกคฺฆโน ปุรตฺถิมาย เจปิ ทิสาย อาคจฺเฉยฺย ภุสา วาตวุฏฺฐิ, เนว นํ สงฺกมฺเปยฺย น สมฺปกมฺเปยฺย น สมฺปเวเธยฺย. ปจฺฉิมาย เจปิ ทิสาย อาคจฺเฉยฺย ภุสา วาตวุฏฺฐิ ฯเปฯ. อุตฺตราย เจปิ ทิสาย ฯเปฯ. ทกฺขิณาย เจปิ ทิสาย อาคจฺเฉยฺย ภุสา วาตวุฏฺฐิ, เนว นํ สงฺกมฺเปยฺย น สมฺปกมฺเปยฺย น สมฺปเวเธยฺย. เอวเมว โข ภนฺเต เอวํ สมฺมา วิมุตฺต\-จิตฺตสฺส ภิกฺขุโน ภุสา เจปิ จกฺขุวิญฺเญยฺยา รูปา จกฺขุสฺส อาปาถํ อาคจฺฉนฺติ, เนวสฺส จิตฺตํ ปริยาทิยนฺติ, อมิสฺสีกตเม\-วสฺส จิตฺตํ โหติ, ฐิตํ อาเนญฺชปฺปตฺตํ, วยญฺจ\-สฺสา\-นุปสฺสติ. ภุสา เจปิ โสตวิญฺเญยฺยา สทฺทา ฯเปฯ ฆานวิญฺเญยฺยา คนฺธา. ชิวฺหาวิญฺเญยฺยา รสา. กายวิญฺเญยฺยา โผฏฺฐพฺพา. มโนวิญฺเญยฺยา ธมฺมา มนสฺส อาปาถํ อาคจฺฉนฺติ, เนวสฺส จิตฺตํ ปริยาทิยนฺติ, อมิสฺสีกตเม\-วสฺส จิตฺตํ โหติ, ฐิตํ อาเนญฺชปฺปตฺตํ, วยญฺจสฺสานุปสฺสตี”ติ.}

\setlength{\csromangathapagewidth}{0pt}
\setlength{\csromangathawakwidth}{0pt}
\csromangathameasuregroup{เนกฺขมฺมํ อธิมุตฺตสฺส, \\ อพฺยาปชฺชาธิมุตฺตสฺส, \\ ตณฺหกฺขยาธิมุตฺตสฺส, \\ ทิสฺวา อายตนุปฺปาทํ, \\ ตสฺส สมฺมา วิมุตฺตสฺส, \\ กตสฺส ปฏิจโย นตฺถิ, \\ เสโล ยถา เอกคฺฆโน, \\ เอวํ รูปา รสา สทฺทา, \\ อิฏฺฐา ธมฺมา อนิฏฺฐา จ, \\ ฐิตํ จิตฺตํ วิปฺปมุตฺตํ,}{\csromangathabat{เนกฺขมฺมํ อธิมุตฺตสฺส,}{ปวิเวกญฺจ เจตโส.} \\ \csromangathabat{อพฺยาปชฺชาธิมุตฺตสฺส,}{อุปาทานกฺขยสฺส จ.} \\ \csromangathabat{ตณฺหกฺขยาธิมุตฺตสฺส,}{อสมฺโมหญฺจ เจตโส.} \\ \csromangathabat{ทิสฺวา อายตนุปฺปาทํ,}{สมฺมา จิตฺตํ วิมุจฺจติ.} \\ \csromangathabat{ตสฺส สมฺมา วิมุตฺตสฺส,}{สนฺตจิตฺตสฺส ภิกฺขุโน.} \\ \csromangathabat{กตสฺส ปฏิจโย นตฺถิ,}{กรณียํ น วิชฺชติ.} \\ \csromangathabat{เสโล ยถา เอกคฺฆโน,}{วาเตน น สมีรติ.} \\ \csromangathabat{เอวํ รูปา รสา สทฺทา,}{คนฺธา ผสฺสา จ เกวลา.} \\ \csromangathabat{อิฏฺฐา ธมฺมา อนิฏฺฐา จ,}{น ปเวเธนฺติ ตาทิโน.} \\ \csromangathabat{ฐิตํ จิตฺตํ วิปฺปมุตฺตํ,}{วยญฺจสฺสานุปสฺสตีติ.}}
\csromangathasetleft{เนกฺขมฺมํ อธิมุตฺตสฺส, \\ อพฺยาปชฺชาธิมุตฺตสฺส, \\ ตณฺหกฺขยาธิมุตฺตสฺส, \\ ทิสฺวา อายตนุปฺปาทํ, \\ ตสฺส สมฺมา วิมุตฺตสฺส, \\ กตสฺส ปฏิจโย นตฺถิ, \\ เสโล ยถา เอกคฺฆโน, \\ เอวํ รูปา รสา สทฺทา, \\ อิฏฺฐา ธมฺมา อนิฏฺฐา จ, \\ ฐิตํ จิตฺตํ วิปฺปมุตฺตํ,}
\csromangathagroup{\csromangathastanza{\csromangathabat{เนกฺขมฺมํ อธิมุตฺตสฺส,}{ปวิเวกญฺจ เจตโส.} \\ \csromangathabat{อพฺยาปชฺชาธิมุตฺตสฺส,}{อุปาทานกฺขยสฺส จ.}}\csromangathastanza{\csromangathabat{ตณฺห\-กฺขยา\-ธิมุตฺตสฺส,}{อสมฺโมหญฺจ เจตโส.} \\ \csromangathabat{ทิสฺวา อายตนุปฺปาทํ,}{สมฺมา จิตฺตํ วิมุจฺจติ.}}\csromangathastanza{\csromanfolio{272}\csromangathabat{ตสฺส สมฺมา วิมุตฺตสฺส,}{สนฺตจิตฺตสฺส ภิกฺขุโน.} \\ \csromangathabat{กตสฺส ปฏิจโย นตฺถิ,}{กรณียํ น วิชฺชติ.}}\csromangathastanza{\csromangathabat{เสโล ยถา เอกคฺฆโน,}{วาเตน น สมีรติ.} \\ \csromangathabat{เอวํ รูปา รสา สทฺทา,}{คนฺธา ผสฺสา จ เกวลา.}}\csromangathastanza{\csromangathabat{อิฏฺฐา ธมฺมา อนิฏฺฐา จ,}{น ปเวเธนฺติ ตาทิโน.} \\ \csromangathabat{ฐิตํ จิตฺตํ วิปฺปมุตฺตํ,}{วยญฺจ\-สฺสา\-นุปสฺสตีติ.}}}

\csromanmatikaanchor{mtk.153}
\csromantocmarkref{section}{148. ทิคุณาทิอุปาหนปฏิกฺเขป}{mtk.153}
\bookmark[dest={mtk.153},level=1]{148. ทิคุณาทิอุปาหนปฏิกฺเขป}
\csromanheader{1}{148. ทิคุณา\-ทิ\-อุปาหน\-ปฏิกฺเขป}
\proseitem{245}{อถ โข ภควา ภิกฺขู อามนฺเตสิ “เอวํ โข ภิกฺขเว กุลปุตฺตา อญฺญํ พฺยากโรนฺติ, อตฺโถ จ วุตฺโต อตฺตา จ อนุปนีโต, อถ จ ปนิเธกจฺเจ โมฆปุริสา หสมานกํ มญฺเญ อญฺญํ พฺยากโรนฺติ, เต ปจฺฉา วิฆาตํ อาปชฺชนฺตี”ติ. อถ โข ภควา อายสฺมนฺตํ โสณํ อามนฺเตสิ “ตฺวํ โขสิ โสณ สุขุมาโล, อนุชานามิ เต โสณ เอกปลาสิกํ อุปาหนนฺ”ติ. อหํ โข ภนฺเต อสีติ\-สกฏวาเห หิรญฺญํ โอหาย อคารสฺมา อนคาริยํ ปพฺพชิโต สตฺตหตฺถิกญฺจ อนีกํ, อถาหํ ภนฺเต เอกปลาสิกํ เจ อุปาหนํ ปริหริสฺสามิ, ตสฺส เม ภวิสฺสนฺติ วตฺตาโร “โสโณ โกฬิวิโส อสีติ\-สกฏวาเห หิรญฺญํ โอหาย อคารสฺมา อนคาริยํ ปพฺพชิโต สตฺตหตฺถิกญฺจ อนีกํ, โส ทานายํ เอกปลาสิกาสุ อุปาหนาสุ สตฺโต”ติ, สเจ ภควา ภิกฺขุ\-สํฆสฺส อนุชานิสฺสติ อหมฺปิ ปริภุญฺชิสฺสามิ, โน เจ ภควา ภิกฺขุ\-สํฆสฺส อนุชานิสฺสติ อหมฺปิ น ปริภุญฺชิสฺสามีติ. อถ โข ภควา เอตสฺมิํ นิทาเน เอตสฺมิํ ปกรเณ ธมฺมิํ กถํ กตฺวา ภิกฺขู อามนฺเตสิ “อนุชานามิ ภิกฺขเว เอกปลาสิกํ อุปาหนํ, น ภิกฺขเว ทิคุณา อุปาหนา ธาเรตพฺพา, น ติคุณา อุปาหนา ธาเรตพฺพา, น คุณงฺคุณู\-ปาหนา\csromansharedfootnote{19ca8b72d3659684}{คณงฺคณูปาหนา (พหูสุ)} ธาเรตพฺพา, โย ธาเรยฺย, อาปตฺติ ทุกฺกฏสฺสา”ติ.}

\csromanmatikaanchor{mtk.154}
\csromantocmarkref{section}{149. สพฺพนีลิกาทิปฏิกฺเขป}{mtk.154}
\bookmark[dest={mtk.154},level=1]{149. สพฺพนีลิกาทิปฏิกฺเขป}
\csromanheader{1}{149. \textbf{สพฺพนีลิกา}\-\textbf{ทิ}\-\textbf{ปฏิกฺเขป}}
\proseitem{246}{เตน โข ปน สมเยน ฉพฺพคฺคิยา ภิกฺขู สพฺพนีลิกา อุปาหนาโย ธาเรนฺติ ฯเปฯ สพฺพปีติกา อุปาหนาโย ธาเรนฺติ. สพฺพโลหิติกา อุปาหนาโย ธาเรนฺติ. สพฺพ\-มญฺชิฏฺฐิกา\csromansharedfootnote{507b73c713746bd9}{สพฺพมญฺเชฏฺฐิกา (ก)} อุปาหนาโย \csromanfolio{273}ธาเรนฺติ. สพฺพกณฺหา อุปาหนาโย ธาเรนฺติ. สพฺพ\-มหารงฺค\-รตฺตา อุปาหานาโย ธาเรนฺติ. สพฺพ\-มหา\-นาม\-รตฺตา อุปาหนาโย ธาเรนฺติ. มนุสฺสา อุชฺฌายนฺติ ขิยฺยนฺติ วิปาเจนฺติ “เสยฺยถาปิ คิหี กามโภคิโน”ติ. ภควโต เอตมตฺถํ อาโรเจสุํ. น ภิกฺขเว สพฺพนีลิกา อุปาหนา ธาเรตพฺพา ฯเปฯ น สพฺพปีติกา อุปาหนา ธาเรตพฺพา. น สพฺพโลหิติกา อุปาหนา ธาเรตพฺพา. น สพฺพ\-มญฺชิฏฺฐิกา อุปาหนา ธาเรตพฺพา. น สพฺพกณฺหา อุปาหนา ธาเรตพฺพา. น สพฺพ\-มหารงฺค\-รตฺตา อุปาหนา ธาเรตพฺพา. น สพฺพ\-มหา\-นาม\-รตฺตา อุปาหนา ธาเรตพฺพา, โย ธาเรยฺย, อาปตฺติ ทุกฺกฏสฺสาติ.}

\prose{เตน โข ปน สมเยน ฉพฺพคฺคิยา ภิกฺขูนี ลกวทฺธิกา\csromansharedfootnote{4eaf8c587a061574}{วฏฺฏิกา (สี)} อุปาหนาโย ธาเรนฺติ ฯเปฯ ปีตกวทฺธิกา อุปาหนาโย ธาเรนฺติ. โลหิตก\-วทฺธิกา อุปาหนาโย ธาเรนฺติ. มญฺชิฏฺฐิกวทฺธิกา อุปาหนาโย ธาเรนฺติ. กณฺหวทฺธิกา อุปาหนาโย ธาเรนฺติ. มหารงฺค\-รตฺต\-วทฺธิกา อุปาหนาโย ธาเรนฺติ. มหานาม\-รตฺต\-วทฺธิกา อุปาหนาโย ธาเรนฺติ. มนุสฺสา อุชฺฌายนฺติ ขิยฺยนฺติ วิปาเจนฺติ “เสยฺยถาปิ คิหี กามโภคิโน”ติ. ภควโต เอตมตฺถํ อาโรเจสุํ. น ภิกฺขเว นีลกวทฺธิกา อุปาหนา ธาเรตพฺพา ฯเปฯ น ปีตกวทฺธิกา อุปาหนา ธาเรตพฺพา. น โลหิตก\-วทฺธิกา อุปาหนา ธาเรตพฺพา. น มญฺชิฏฺฐิกวทฺธิกา อุปาหนา ธาเรตพฺพา. น กณฺหวทฺธิกา อุปาหนา ธาเรตพฺพา. น มหารงฺค\-รตฺต\-วทฺธิกา อุปาหนา ธาเรตพฺพา. น มหานาม\-รตฺต\-วทฺธิกา อุปาหนา ธาเรตพฺพา, โย ธาเรยฺย, อาปตฺติ ทุกฺกฏสฺสาติ.}

\prose{เตน โข ปน สมเยน ฉพฺพคฺคิยา ภิกฺขู ขลฺลกพทฺธา\csromansharedfootnote{ade68ce9b47ec360}{...พนฺธา (ก)} อุปาหนาโย ธาเรนฺติ ฯเปฯ ปุฏพทฺธา อุปาหนาโย ธาเรนฺติ. ปาลิคุณฺฐิมา อุปาหนาโย ธาเรนฺติ. ตูลปุณฺณิกา อุปาหนาโย ธาเรนฺติ. ติตฺติรปตฺติกา อุปาหนาโย ธาเรนฺติ. เมณฺฑ\-วิสาณ\-วทฺธิกา อุปาหนาโย ธาเรนฺติ. อชวิสาณวทฺธิกา อุปาหนาโย ธาเรนฺติ. วิจฺฉิกาฬิกา อุปาหนาโย ธาเรนฺติ. โมรปิญฺฉ\csromansharedfootnote{47a3c03834b724e2}{โมรปิญฺช (สี, สฺยา)} ปริสิพฺพิตา อุปาหนาโย ธาเรนฺติ. จิตฺรา อุปาหนาโย ธาเรนฺติ. มนุสฺสา อุชฺฌายนฺติ ขิยฺยนฺติ วิปาเจนฺติ “เสยฺยถาปิ คิหี กามโภคิโน”ติ. ภควโต เอตมตฺถํ อาโรเจสุํ. น \csromanfolio{274}ภิกฺขเว ขลฺลกพทฺธา อุปาหนา ธาเรตพฺพา ฯเปฯ น ปุฏพทฺธา อุปาหนา ธาเรตพฺพา. น ปาลิคุณฺฐิมา อุปาหนา ธาเรตพฺพา. น ตูลปุณฺณิกา อุปาหนา ธาเรตพฺพา. น ติตฺติรปตฺติกา อุปาหนา ธาเรตพฺพา. น เมณฺฑ\-วิสาณ\-วทฺธิกา อุปาหนา ธาเรตพฺพา. น อชวิสาณวทฺธิกา อุปาหนา ธาเรตพฺพา. น วิจฺฉิกาฬิกา อุปาหนา ธาเรตพฺพา. น โมร\-ปิญฺฉ\-ปริสิพฺพิตา อุปาหนา ธาเรตพฺพา. น จิตฺรา อุปาหนา ธาเรตพฺพา, โย ธาเรยฺย, อาปตฺติ ทุกฺกฏสฺสาติ.}

\prose{เตน โข ปน สมเยน ฉพฺพคฺคิยา ภิกฺขู สีหจมฺม\-ปริกฺขฏา อุปาหนาโย ธาเรนฺติ ฯเปฯ พฺยคฺฆ\-จมฺม\-ปริกฺขฏา อุปาหนาโย ธาเรนฺติ. ทีปิจมฺม\-ปริกฺขฏา อุปาหนาโย ธาเรนฺติ. อชินจมฺม\-ปริกฺขฏา อุปาหนาโย ธาเรนฺติ. อุทฺทจมฺม\-ปริกฺขฏา อุปาหนาโย ธาเรนฺติ. มชฺชารจมฺม\-ปริกฺขฏา อุปาหนาโย ธาเรนฺติ. กาฬกจมฺม\-ปริกฺขฏา อุปาหนาโย ธาเรนฺติ. ลุวกจมฺม\-ปริกฺขฏา\csromansharedfootnote{929e240b12c2ad6e}{อุลูกจมฺมปริกฺขฏา (โยชนา)} อุปาหนาโย ธาเรนฺติ. มนุสฺสา อุชฺฌายนฺติ ขิยฺยนฺติ วิปาเจนฺติ “เสยฺยถาปิ คิหี กามโภคิโน”ติ. ภควโต เอตมตฺถํ อาโรเจสุํ. น ภิกฺขเว สีหจมฺม\-ปริกฺขฏา อุปาหนา ธาเรตพฺพา ฯเปฯ น พฺยคฺฆ\-จมฺม\-ปริกฺขฏา อุปาหนา ธาเรตพฺพา. น ทีปิจมฺม\-ปริกฺขฏา อุปาหนา ธาเรตพฺพา. น อชินจมฺม\-ปริกฺขฏา อุปาหนา ธาเรตพฺพา. น อุทฺทจมฺม\-ปริกฺขฏา อุปาหนา ธาเรตพฺพา. น มชฺชารจมฺม\-ปริกฺขฏา อุปาหนา ธาเรตพฺพา. น กาฬกจมฺม\-ปริกฺขฏา อุปาหนา ธาเรตพฺพา. น ลุวกจมฺม\-ปริกฺขฏา อุปาหนา ธาเรตพฺพา, โย ธาเรยฺย, อาปตฺติ ทุกฺกฏสฺสาติ.}

\csromanmatikaanchor{mtk.155}
\csromantocmarkref{section}{150. โอมุกฺกคุณงฺคุณูปาหนานุชานน}{mtk.155}
\bookmark[dest={mtk.155},level=1]{150. โอมุกฺกคุณงฺคุณูปาหนานุชานน}
\csromanheader{1}{150. \textbf{โอมุกฺกคุณงฺคุณูปาหนานุชานน}}
\proseitem{247}{อถ โข ภควา ปุพฺพณฺหสมยํ นิวาเสตฺวา ปตฺต\-จีวรมา\-ทาย ราชคหํ ปิณฺฑาย ปาวิสิ อญฺญตเรน ภิกฺขุนา ปจฺฉาสมเณน. อถ โข โส ภิกฺขุ ขญฺชมาโน ภควนฺตํ ปิฏฺฐิโต ปิฏฺฐิโต อนุพนฺธิ. อทฺทสา โข อญฺญตโร อุปาสโก คุณงฺคุณู\-ปาหนา อาโรหิตฺวา ภควนฺตํ ทูรโตว อาคจฺฉนฺตํ, ทิสฺวา อุปาหนา โอโรหิตฺวา เยน ภควา เตนุปสงฺกมิ, อุปสงฺกมิตฺวา ภควนฺตํ อภิวาเทตฺวา เยน โส ภิกฺขุ เตนุปสงฺกมิ, อุปสงฺกมิตฺวา ตํ ภิกฺขุํ อภิวาเทตฺวา เอตทโวจ “กิสฺส ภนฺเต อยฺโย \csromanfolio{275}ขญฺชตี”ติ. ปาทา เม อาวุโส ผลิตาติ\csromansharedfootnote{eb0f240424da97dd}{ผาลิตาติ (ก)}. หนฺท ภนฺเต อุปาหนาโยติ. อลํ อาวุโส ปฏิกฺขิตฺตา ภควตา คุณงฺคุณู\-ปาหนาติ. คณฺหาเหตา ภิกฺขุ อุปาหนาโยติ. อถ โข ภควา เอตสฺมิํ นิทาเน เอตสฺมิํ ปกรเณ ธมฺมิํ กถํ กตฺวา ภิกฺขู อามนฺเตสิ “อนุชานามิ ภิกฺขเว โอมุกฺกํ คุณงฺคุณู\-ปาหนํ, น ภิกฺขเว นวา คุณงฺคุณู\-ปาหนา ธาเรตพฺพา, โย ธาเรยฺย, อาปตฺติ ทุกฺกฏสฺสา”ติ.}

\csromanmatikaanchor{mtk.156}
\csromantocmarkref{section}{151. อชฺฌาราเม อุปาหนปฏิกฺเขป}{mtk.156}
\bookmark[dest={mtk.156},level=1]{151. อชฺฌาราเม อุปาหนปฏิกฺเขป}
\csromanheader{1}{151. \textbf{อชฺฌาราเม อุปาหน}\-\textbf{ปฏิกฺเขป}}
\proseitem{248}{เตน โข ปน สมเยน ภควา อชฺโฌกาเส อนุปาหโน จงฺกมติ, สตฺถา อนุปาหโน จงฺกมตีติ เถราปิ ภิกฺขู อนุปาหนา จงฺกมนฺติ, ฉพฺพคฺคิยา ภิกฺขู สตฺถริ อนุปาหเน จงฺกมมาเน เถเรสุปิ ภิกฺขูสุ อนุปาหเนสุ จงฺกมมาเนสุ สอุปาหนา จงฺกมนฺติ. เย เต ภิกฺขู อปฺปิจฺฉา ฯเปฯ เต อุชฺฌายนฺติ ขิยฺยนฺติ วิปาเจนฺติ “กถํ หิ นาม ฉพฺพคฺคิยา ภิกฺขู สตฺถริ อนุปาหเน จงฺกมมาเน เถเรสุปิ ภิกฺขูสุ อนุปาหเนสุ จงฺกมมาเนสุ สอุปาหนา จงฺกมิสฺสนฺตี”ติ. อถ โข เต ภิกฺขู ภควโต เอตมตฺถํ อาโรเจสุํ ฯเปฯ. สจฺจํ กิร ภิกฺขเว ฉพฺพคฺคิยา ภิกฺขู สตฺถริ อนุปาหเน จงฺกมมาเน เถเรสุปิ ภิกฺขูสุ อนุปาหเนสุ จงฺกมมาเนสุ สอุปาหนา จงฺกมนฺตีติ. สจฺจํ ภควาติ. วิครหิ พุทฺโธ ภควา ฯเปฯ. กถํ หิ นาม เต ภิกฺขเว โมฆปุริสา สตฺถริ อนุปาหเน จงฺกมมาเน เถเรสุปิ ภิกฺขูสุ อนุปาหเนสุ จงฺกมมาเนสุ สอุปาหนา จงฺกมิสฺสนฺติ, อิเม หิ นาม ภิกฺขเว คิหี โอทาต\-วตฺถ\-วสนกา อภิชีวนิกสฺส สิปฺปสฺส การณา อาจริเยสุ สคารวา สปฺปติสฺสา สภาควุตฺติกา วิหริสฺสนฺติ, อิธ โข ตํ ภิกฺขเว โสเภถ, ยํ ตุเมฺห เอวํ สฺวากฺขาเต ธมฺมวินเย ปพฺพชิตา สมานา อาจาริเยสุ อาจริยมตฺเตสุ อุปชฺฌาเยสุ อุปชฺฌายมตฺเตสุ อคารวา อปฺปติสฺสา อสภาค\-วุตฺติกา\csromansharedfootnote{dadf278cd13073d7}{สคารวา สปฺปติสฺสา สภาควุตฺติกา (ก)} วิหเรยฺยาถ? เนตํ ภิกฺขเว อปฺปสนฺนานํ วา ปสาทาย ฯเปฯ วิครหิตฺวา ฯเปฯ ธมฺมิํ กถํ กตฺวา ภิกฺขู อามนฺเตสิ “น ภิกฺขเว อาจริเยสุ อาจริยมตฺเตสุ อุปชฺฌาเยสุ อุปชฺฌายมตฺเตสุ อนุปาหเนสุ จงฺกมมาเนสุ สอุปาหเนน จงฺกมิตพฺพํ, โย จงฺกเมยฺย, อาปตฺติ ทุกฺกฏสฺส. น จ ภิกฺขเว อชฺฌาราเม อุปาหนา ธาเรตพฺพา, โย ธาเรยฺย, อาปตฺติ ทุกฺกฏสฺสา”ติ.}

\proseitem{249}{\csromanfolio{276}เตน โข ปน สมเยน อญฺญตรสฺส ภิกฺขุโน ปาทขิลาพาโธ โหติ, ตํ ภิกฺขู ปริคฺคเหตฺวา อุจฺจารมฺปิ ปสฺสาวมฺปิ นิกฺขาเมนฺติ.}

\prose{อทฺทสา โข ภควา เสนาสนจาริกํ อาหิณฺฑนฺโต เต ภิกฺขู ตํ ภิกฺขุํ ปริคฺคเหตฺวา อุจฺจารมฺปิ ปสฺสาวมฺปิ นิกฺขาเมนฺเต, ทิสฺวาน เยน เต ภิกฺขู เตนุปสงฺกมิ, อุปสงฺกมิตฺวา เต ภิกฺขู เอตทโวจ “กิํ อิมสฺส ภิกฺขเว ภิกฺขุโน อาพาโธ”ติ.}

\prose{อิมสฺส ภนฺเต อายสฺมโต ปาทขิลาพาโธ, อิมํ มยํ ปริคฺคเหตฺวา อุจฺจารมฺปิ ปสฺสาวมฺปิ นิกฺขาเมมาติ.}

\prose{อถ โข ภควา เอตสฺมิํ นิทาเน เอตสฺมิํ ปกรเณ ธมฺมิํ กถํ กตฺวา ภิกฺขู อามนฺเตสิ “อนุชานามิ ภิกฺขเว ยสฺส ปาทา วา ทุกฺขา ปาทา วา ผลิตา ปาทขิโล วา อาพาโธ\csromansharedfootnote{338dd9d64fd7ee1e}{[มิสฺสินฺคฺ โนเต 0]} อุปาหนํ ธาเรตุนฺ”ติ.}

\prose{เตน โข ปน สมเยน ภิกฺขู อโธเตหิ ปาเทหิ มญฺจมฺปิ ปีฐมฺปิ อภิรุหนฺติ, จีวรมฺปิ เสนาสนมฺปิ ทุสฺสติ.}

\prose{ภควโต เอตมตฺถํ อาโรเจสุํ. อนุชานามิ ภิกฺขเว “อิทานิ มญฺจํ วา ปีฐํ วา อภิรุหิสฺสามี”ติ อุปาหนํ ธาเรตุนฺติ.}

\prose{เตน โข ปน สมเยน ภิกฺขู รตฺติยา อุโปสถคฺคมฺปิ สนฺนิสชฺชมฺปิ คจฺฉนฺตา อนฺธกาเร ขาณุมฺปิ กณฺฏกมฺปิ อกฺกมนฺติ, ปาทา ทุกฺขา โหนฺติ.}

\prose{ภควโต เอตมตฺถํ อาโรเจสุํ. อนุชานามิ ภิกฺขเว อชฺฌาราเม อุปาหนํ ธาเรตุํ อุกฺกํ ปทีปํ กตฺตร\-ทณฺฑนฺติ.}

\csromanmatikaanchor{mtk.157}
\csromantocmarkref{section}{152. กฏฺฐปาทุกาทิปฏิกฺเขป}{mtk.157}
\bookmark[dest={mtk.157},level=1]{152. กฏฺฐปาทุกาทิปฏิกฺเขป}
\csromanheader{1}{152. \textbf{กฏฺฐปาทุกา}\-\textbf{ทิ}\-\textbf{ปฏิกฺเขป}}
\proseitem{250}{เตน โข ปน สมเยน ฉพฺพคฺคิยา ภิกฺขู รตฺติยา ปจฺจูสสมยํ ปจฺจุฏฺฐาย กฏฺฐปาทุกาโย อภิรุหิตฺวา อชฺโฌกาเส จงฺกมนฺติ อุจฺจาสทฺทา มหาสทฺทา ขฏขฏสทฺทา อเนกวิหิตํ ติรจฺฉานกถํ กเถนฺตา, เสยฺยถิทํ,\csromansymbolfootnote{*}{อิมา ติรจฺฉานกถาโย วิ 2. 213; ที 1. 7; ม 2. 181; สํ 3. 368; อํ 3. 358 ปิฏฺฐาทีสุปิ อาคตา.}ราชกถํ โจรกถํ มหามตฺตกถํ เสนากถํ ภยกถํ ยุทฺธกถํ อนฺนกถํ ปานกถํ วตฺถกถํ สยนกถํ มาลากถํ คนฺธกถํ ญาติกถํ ยานกถํ คามกถํ นิคมกถํ นครกถํ ชนปทกถํ อิตฺถิกถํ\csromansharedfootnote{3974d88881fe2ec3}{อิตฺถิกถํ ปุริสกถํ (ก)} สูรกถํ วิสิขากถํ กุมฺภ\-ฏฺฐานกถํ ปุพฺพเปตกถํ นานตฺตกถํ โลกกฺขายิกํ สมุทฺท\-กฺขายิกํ อิติภวา\-ภวกถํ อิติ วา, กีฏกมฺปิ อกฺกมิตฺวา มาเรนฺติ, ภิกฺขูปิ สมาธิมฺหา จาเวนฺติ. \csromanfolio{277}เย เต ภิกฺขู อปฺปิจฺฉา ฯเปฯ เต อุชฺฌายนฺติ ขิยฺยนฺติ วิปาเจนฺติ “กถํ หิ นาม ฉพฺพคฺคิยา ภิกฺขู รตฺติยา ปจฺจูสสมยํ ปจฺจุฏฺฐาย กฏฺฐปาทุกาโย อภิรุหิตฺวา อชฺโฌกาเส จงฺกมิสฺสนฺติ อุจฺจาสทฺทา มหาสทฺทา ขฏขฏสทฺทา อเนกวิหิตํ ติรจฺฉานกถํ กเถนฺตา, เสยฺยถิทํ, ราชกถํ โจรกถํ ฯเปฯ อิติภวา\-ภวกถํ อิติ วา, กีฏกมฺปิ อกฺกมิตฺวา มาเรสฺสนฺติ, ภิกฺขูปิ สมาธิมฺหา จาเวสฺสนฺตี”ติ. อถ โข เต ภิกฺขู ภควโต เอตมตฺถํ อาโรเจสุํ ฯเปฯ. สจฺจํ กิร ภิกฺขเว ฉพฺพคฺคิยา ภิกฺขู รตฺติยา ปจฺจูสสมยํ ปจฺจุฏฺฐาย กฏฺฐปาทุกาโย อภิรุหิตฺวา อชฺโฌกาเส จงฺกมนฺติ อุจฺจาสทฺทา มหาสทฺทา ขฏขฏสทฺทา อเนกวิหิตํ ติรจฺฉานกถํ กเถนฺตา, เสยฺยถิทํ, ราชกถํ โจรกถํ ฯเปฯ อิติภวา\-ภวกถํ อิติ วา, กีฏกมฺปิ อกฺกมิตฺวา มาเรนฺติ, ภิกฺขูปิ สมาธิมฺหา จาเวนฺตีติ. สจฺจํ ภควาติ ฯเปฯ วิครหิตฺวา ฯเปฯ ธมฺมิํ กถํ กตฺวา ภิกฺขู อามนฺเตสิ “น ภิกฺขเว กฏฺฐปาทุกา ธาเรตพฺพา, โย ธาเรยฺย, อาปตฺติ ทุกฺกฏสฺสา”ติ.}

\prose{อถ โข ภควา ราชคเห ยถาภิรนฺตํ วิหริตฺวา เยน พาราณสี เตน จาริกํ ปกฺกามิ, อนุปุพฺเพน จาริกํ จรมาโน เยน พาราณสี ตทวสริ, ตตฺร สุทํ ภควา พาราณสิยํ วิหรติ อิสิปตเน มิคทาเย. เตน โข ปน สมเยน ฉพฺพคฺคิยา ภิกฺขู “ภควตา กฏฺฐปาทุกา ปฏิกฺขิตฺตา”ติ ตาลตรุเณ เฉทาเปตฺวา ตาลปตฺต\-ปาทุกาโย ธาเรนฺติ, ตานิ ตาลตรุณานิ ฉินฺนานิ มิลายนฺติ. มนุสฺสา อุชฺฌายนฺติ ขิยฺยนฺติ วิปาเจนฺติ “กถํ หิ นาม สมณา สกฺยปุตฺติยา ตาลตรุเณ เฉทาเปตฺวา ตาลปตฺต\-ปาทุกาโย ธาเรสฺสนฺติ, ตานิ ตาลตรุณานิ ฉินฺนานิ มิลายนฺติ, เอกินฺทฺริยํ สมณา สกฺยปุตฺติยา ชีวํ วิเหเฐนฺตี”ติ. อสฺโสสุํ โข ภิกฺขู เตสํ มนุสฺสานํ อุชฺฌายนฺตานํ ขิยฺยนฺตานํ วิปาเจนฺตานํ. อถ โข เต ภิกฺขู ภควโต เอตมตฺถํ อาโรเจสุํ ฯเปฯ. สจฺจํ กิร ภิกฺขเว ฉพฺพคฺคิยา ภิกฺขู ตาลตรุเณ เฉทาเปตฺวา ตาลปตฺต\-ปาทุกาโย ธาเรนฺติ, ตานิ ตาลตรุณานิ ฉินฺนานิ มิลายนฺตีติ. สจฺจํ ภควาติ. วิครหิ พุทฺโธ ภควา ฯเปฯ. กถํ หิ นาม เต ภิกฺขเว โมฆปุริสา ตาลตรุเณ เฉทาเปตฺวา ตาลปตฺต\-ปาทุกาโย ธาเรสฺสนฺติ, ตานิ ตาลตรุณานิ ฉินฺนานิ มิลายนฺติ, ชีวสญฺญิโน หิ ภิกฺขเว มนุสฺสา รุกฺขสฺมิํ. เนตํ ภิกฺขเว \csromanfolio{278}อปฺปสนฺนานํ วา ปสาทาย ฯเปฯ วิครหิตฺวา ฯเปฯ ธมฺมิํ กถํ กตฺวา ภิกฺขู อามนฺเตสิ “น ภิกฺขเว ตาลปตฺต\-ปาทุกา ธาเรตพฺพา, โย ธาเรยฺย, อาปตฺติ ทุกฺกฏสฺสา”ติ.}

\prose{เตน โข ปน สมเยน ฉพฺพคฺคิยา ภิกฺขู “ภควตา ตาลปตฺต\-ปาทุกา ปฏิกฺขิตฺตา”ติ เวฬุตรุเณ เฉทาเปตฺวา เวฬุปตฺต\-ปาทุกาโย ธาเรนฺติ, ตานิ เวฬุตรุณานิ ฉินฺนานิ มิลายนฺติ. มนุสฺสา อุชฺฌายนฺติ ขิยฺยนฺติ วิปาเจนฺติ “กถํ หิ นาม สมณา สกฺยปุตฺติยา เวฬุตรุเณ เฉทาเปตฺวา เวฬุปตฺต\-ปาทุกาโย ธาเรสฺสนฺติ, ตานิ เวฬุตรุณานิ ฉินฺนานิ มิลายนฺติ, เอกินฺทฺริยํ สมณา สกฺยปุตฺติยา ชีวํ วิเหเฐนฺตี”ติ. อสฺโสสุํ โข ภิกฺขู เตสํ มนุสฺสานํ อุชฺฌายนฺตานํ ขิยฺยนฺตานํ วิปาเจนฺตานํ. อถ โข เต ภิกฺขู ภควโต เอตมตฺถํ อาโรเจสุํ ฯเปฯ. ชีวสญฺญิโน หิ ภิกฺขเว มนุสฺสา รุกฺขสฺมิํ ฯเปฯ. น ภิกฺขเว เวฬุปตฺต\-ปาทุกา ธาเรตพฺพา, โย ธาเรยฺย, อาปตฺติ ทุกฺกฏสฺสาติ.}

\proseitem{251}{อถ โข ภควา พาราณสิยํ ยถาภิรนฺตํ วิหริตฺวา เยน ภทฺทิยํ เตน จาริกํ ปกฺกามิ, อนุปุพฺเพน จาริกํ จรมาโน เยน ภทฺทิยํ ตทวสริ, ตตฺร สุทํ ภควา ภทฺทิเย วิหรติ ชาติยา วเน. เตน โข ปน สมเยน ภทฺทิยา ภิกฺขู อเนกวิหิตํ ปาทุกมณฺฑนานุโยคมนุยุตฺตา วิหรนฺติ, ติณปาทุกํ กโรนฺติปิ การาเปนฺติปิ, มุญฺชปาทุกํ กโรนฺติปิ การาเปนฺติปิ, ปพฺพชปาทุกํ กโรนฺติปิ การาเปนฺติปิ, หินฺตาลปาทุกํ กโรนฺติปิ การาเปนฺติปิ, กมลปาทุกํ กโรนฺติปิ การาเปนฺติปิ, กมฺพลปาทุกํ กโรนฺติปิ การาเปนฺติปิ, ริญฺจนฺติ อุทฺเทสํ ปริปุจฺฉํ อธิสีลํ อธิจิตฺตํ อธิปญฺญํ. เย เต ภิกฺขู อปฺปิจฺฉา ฯเปฯ เต อุชฺฌายนฺติ ขิยฺยนฺติ วิปาเจนฺติ “กถํ หิ นาม ภทฺทิยา ภิกฺขู อเนกวิหิตํ ปาทุกมณฺฑนานุโยคมนุยุตฺตา วิหริสฺสนฺติ, ติณปาทุกํ กริสฺสนฺติปิ การาเปสฺสนฺติปิ, มุญฺชปาทุกํ กริสฺสนฺติปิ การาเปสฺสนฺติปิ, ปพฺพชปาทุกํ กริสฺสนฺติปิ การาเปสฺสนฺติปิ, หินฺตาลปาทุกํ กริสฺสนฺติปิ การาเปสฺสนฺติปิ, กมลปาทุกํ กริสฺสนฺติปิ การาเปสฺสนฺติปิ, กมฺพลปาทุกํ กริสฺสนฺติปิ การาเปสฺสนฺติปิ, ริญฺชิสฺสนฺติ อุทฺเทสํ ปริปุจฺฉํ อธิสีลํ อธิจิตฺตํ อธิปญฺญนฺ”ติ. อถ โข เต ภิกฺขู ภควโต เอตมตฺถํ อาโรเจสุํ ฯเปฯ. สจฺจํ กิร ภิกฺขเว ภทฺทิยา ภิกฺขู อเนกวิหิตํ ปาทุกมณฺฑนานุโยคมนุยุตฺตา วิหรนฺติ, ติณปาทุกํ กโรนฺติปิ การาเปนฺติปิ ฯเปฯ ริญฺจนฺติ อุทฺเทสํ ปริปุจฺฉํ อธิสีลํ \csromanfolio{279}อธิจิตฺตํ อธิปญฺญนฺ”ติ. สจฺจํ ภควาติ. วิครหิ พุทฺโธ ภควา ฯเปฯ . กถํ หิ นาม เต ภิกฺขเว โมฆปุริสา อเนกวิหิตํ ปาทุกมณฺฑนานุโยคมนุยุตฺตา วิหริสฺสนฺติ, ติณปาทุกํ กริสฺสนฺติปิ การาเปสฺสนฺติปิ ฯเปฯ ริญฺจิสฺสนฺติ อุทฺเทสํ ปริปุจฺฉํ อธิสีลํ อธิจิตฺตํ อธิปญฺญํ. เนตํ ภิกฺขเว อปฺปสนฺนานํ วา ปสาทาย ฯเปฯ วิครหิตฺวา ฯเปฯ ธมฺมิํ กถํ กตฺวา ภิกฺขู อามนฺเตสิ “น ภิกฺขเว ติณปาทุกา ธาเรตพฺพา, น มุญฺชปาทุกา ธาเรตพฺพา, น ปพฺพชปาทุกา ธาเรตพฺพา, น หินฺตาลปาทุกา ธาเรตพฺพา, น กมลปาทุกา ธาเรตพฺพา, น กมฺพลปาทุกา ธาเรตพฺพา, น โสวณฺณมยา ปาทุกา ธาเรตพฺพา, น รูปิยมยา ปาทุกา ธาเรตพฺพา, น มณิมยา ปาทุกา ธาเรตพฺพา, น เวฬุริยมยา ปาทุกา ธาเรตพฺพา, น ผลิกมยา ปาทุกา ธาเรตพฺพา, น กํสมยา ปาทุกา ธาเรตพฺพา, น กาจมยา ปาทุกา ธาเรตพฺพา, น ติปุมยา ปาทุกา ธาเรตพฺพา, น สีสมยา ปาทุกา ธาเรตพฺพา, น ตมฺพโลหมยา ปาทุกา ธาเรตพฺพา, โย ธาเรยฺย, อาปตฺติ ทุกฺกฏสฺส. น จ ภิกฺขเว กาจิ สงฺกมนิยา ปาทุกา ธาเรตพฺพา, โย ธาเรยฺย, อาปตฺติ ทุกฺกฏสฺส. อนุชานามิ ภิกฺขเว ติสฺโส ปาทุกา ธุวฏฺฐานิยา อสงฺกมนิยาโย, วจฺจปาทุกํ ปสฺสาวปาทุกํ อาจมนปาทุกนฺติ.}

\proseitem{252}{อถ โข ภควา ภทฺทิเย ยถาภิรนฺตํ วิหริตฺวา เยน สาวตฺถิ เตน จาริกํ ปกฺกามิ, อนุปุพฺเพน จาริกํ จรมาโน เยน สาวตฺถิ ตทวสริ, ตตฺร สุทํ ภควา สาวตฺถิยํ วิหรติ เชตวเน อนาถ\-ปิณฺฑิกสฺส อาราเม. เตน โข ปน สมเยน ฉพฺพคฺคิยา ภิกฺขู อจิรวติยา นทิยา คาวีนํ ตรนฺตีนํ วิสาเณสุปิ คณฺหนฺติ, กณฺเณสุปิ คณฺหนฺติ, คีวายปิ คณฺหนฺติ, เฉปฺปาปิ คณฺหนฺติ, ปิฏฺฐิมฺปิ อภิรุหนฺติ, รตฺตจิตฺตาปิ องฺคชาตํ ฉุปนฺติ, วจฺฉตริมฺปิ โอคาเหตฺวา มาเรนฺติ. มนุสฺสา อุชฺฌายนฺติ ขิยฺยนฺติ วิปาเจนฺติ “กถํ หิ นาม สมณา สกฺยปุตฺติยา คาวีนํ ตรนฺตีนํ วิสาเณสุปิ คเหสฺสนฺติ ฯเปฯ เสยฺยถาปิ คิหี กามโภคิโน”ติ. อสฺโสสุํ โข ภิกฺขู เตสํ มนุสฺสานํ อุชฺฌายนฺตานํ ขิยฺยนฺตานํ วิปาเจนฺตานํ. อถ โข เต ภิกฺขู ภควโต เอตมตฺถํ อาโรเจสุํ ฯเปฯ. สจฺจํ กิร ภิกฺขเว ฯเปฯ. สจฺจํ ภควาติ ฯเปฯ วิครหิตฺวา ฯเปฯ ธมฺมิํ กถํ กตฺวา ภิกฺขู อามนฺเตสิ “น ภิกฺขเว คาวีนํ วิสาเณสุ คเหตพฺพํ, น กณฺเณสุ คเหตพฺพํ, น คีวาย คเหตพฺพํ, น เฉปฺปาย คเหตพฺพํ, น ปิฏฺฐิ \csromanfolio{280}อภิรุหิตพฺพา, โย อภิรุเหยฺย, อาปตฺติ ทุกฺกฏสฺส. น จ ภิกฺขเว รตฺตจิตฺเตน องฺคชาตํ ฉุปิตพฺพํ, โย ฉุเปยฺย, อาปตฺติ ถุลฺลจฺจยสฺส. น วจฺฉตรี มาเรตพฺพา, โย มาเรยฺย, ยถาธมฺโม กาเรตพฺโพ”ติ.}

\csromanmatikaanchor{mtk.158}
\csromantocmarkref{section}{153. ยานาทิปฏิกฺเขป}{mtk.158}
\bookmark[dest={mtk.158},level=1]{153. ยานาทิปฏิกฺเขป}
\csromanheader{1}{153. \textbf{ยานา}\-\textbf{ทิ}\-\textbf{ปฏิกฺเขป}}
\proseitem{253}{เตน โข ปน สมเยน ฉพฺพคฺคิยา ภิกฺขู ยาเนน ยายนฺติ อิตฺถิยุตฺเตนปิ ปุริสนฺตเรน, ปุริสยุตฺเตนปิ อิตฺถนฺตเรน. มนุสฺสา อุชฺฌายนฺติ ขิยฺยนฺติ วิปาเจนฺติ “เสยฺยถาปิ คงฺคามหิยายา”ติ. ภควโต เอตมตฺถํ อาโรเจสุํ. น ภิกฺขเว ยาเนน ยายิตพฺพํ, โย ยาเยยฺย, อาปตฺติ ทุกฺกฏสฺสาติ. เตน โข ปน สมเยน อญฺญตโร ภิกฺขุ โกสเลสุ ชนปเท สาวตฺถิํ คจฺฉนฺโต ภควนฺตํ ทสฺสนาย อนฺตรามคฺเค คิลาโน โหติ. อถ โข โส ภิกฺขุ มคฺคา โอกฺกมฺม อญฺญตรสฺมิํ รุกฺขมูเล นิสีทิ. มนุสฺสา ตํ ภิกฺขุํ ทิสฺวา เอตทโวจุํ “กหํ ภนฺเต อยฺโย คมิสฺสตี”ติ. สาวตฺถิํ โข อหํ อาวุโส คมิสฺสามิ ภควนฺตํ ทสฺสนายาติ. เอหิ ภนฺเต คมิสฺสามาติ. นาหํ อาวุโส สกฺโกมิ คิลาโนมฺหีติ. เอหิ ภนฺเต ยานํ อภิรุหาติ. อลํ อาวุโส ปฏิกฺขิตฺตํ ภควตา ยานนฺติ กุกฺกุจฺจายนฺโต ยานํ นาภิรุหิ. อถ โข โส ภิกฺขุ สาวตฺถิํ คนฺตฺวา ภิกฺขูนํ เอตมตฺถํ อาโรเจสิ. ภิกฺขู ภควโต เอตมตฺถํ อาโรเจสุํ. อนุชานามิ ภิกฺขเว คิลานสฺส ยานนฺติ. อถ โข ภิกฺขูนํ เอตทโหสิ “อิตฺถิยุตฺตํ นุ โข ปุริสยุตฺตํ นุ โข”ติ. ภควโต เอตมตฺถํ อาโรเจสุํ. อนุชานามิ ภิกฺขเว ปุริสยุตฺตํ หตฺถวฏฺฏกนฺติ.}

\prose{เตน โข ปน สมเยน อญฺญตรสฺส ภิกฺขุโน ยานุคฺฆาเตน พาฬฺหตรํ อผาสุ อโหสิ. ภควโต เอตมตฺถํ อาโรเจสุํ. อนุชานามิ ภิกฺขเว สิวิกํ ปาฏงฺกินฺติ.}

\csromanmatikaanchor{mtk.159}
\csromantocmarkref{section}{154. อุจฺจาสยนมหาสยนปฏิกฺเขป}{mtk.159}
\bookmark[dest={mtk.159},level=1]{154. อุจฺจาสยนมหาสยนปฏิกฺเขป}
\csromanheader{1}{154. \textbf{อุจฺจาสยน}\-\textbf{มหาสยน}\-\textbf{ปฏิกฺเขป}}
\proseitem{254}{เตน โข ปน สมเยน ฉพฺพคฺคิยา ภิกฺขู อุจฺจาสยน\-มหาสยนานิ ธาเรนฺติ, เสยฺยถิทํ, อาสนฺทิํ ปลฺลงฺกํ โคนกํ จิตฺตกํ ปฏิกํ ปฏลิกํ ตูลิกํ วิกติกํ อุทฺธโลมิํ\csromansharedfootnote{4d4df5aa69054b96}{อุนฺทโลมิํ (ก), อุทฺทโลมิํ (ก)} เอกนฺตโลมิํ กฏฺฏิสฺสํ \csromanfolio{281}โกเสยฺยํ กุตฺตกํ หตฺถตฺถรํ อสฺสตฺถรํ รถตฺถรํ อชินปเวณิํ กทลิ\-มิค\-ปวร\-ปจฺจตฺถรณํ สอุตฺตรจฺฉทํ อุภโต\-โลหิตกู\-ปธานนฺติ. มนุสฺสา วิหารจาริกํ อาหิณฺฑนฺตา ปสฺสิตฺวา อุชฺฌายนฺติ ขิยฺยนฺติ วิปาเจนฺติ “เสยฺยถาปิ คิหี กามโภคิโน”ติ. ภควโต เอตมตฺถํ อาโรเจสุํ. น ภิกฺขเว อุจฺจาสยน\-มหาสยนานิ ธาเรตพฺพานิ, เสยฺยถิทํ, อาสนฺทิ ปลฺลงฺโก โคนโก จิตฺตโก ปฏิกา ปฏลิกา ตูลิกา วิกติกา อุทฺธโลมิ เอกนฺตโลมิ กฏฺฏิสฺสํ โกเสยฺยํ กุตฺตกํ หตฺถตฺถรํ อสฺสตฺถรํ รถตฺถรํ อชินปเวณิ กทลิ\-มิค\-ปวร\-ปจฺจตฺถรณํ สอุตฺตรจฺฉทํ อุภโต\-โลหิตกู\-ปธานํ, โย ธาเรยฺย, อาปตฺติ ทุกฺกฏสฺสาติ.}

\csromanmatikaanchor{mtk.160}
\csromantocmarkref{section}{155. สพฺพจมฺมปฏิกฺเขป}{mtk.160}
\bookmark[dest={mtk.160},level=1]{155. สพฺพจมฺมปฏิกฺเขป}
\csromanheader{1}{155. \textbf{สพฺพ}\-\textbf{จมฺม}\-\textbf{ปฏิกฺเขป}}
\proseitem{255}{เตน โข ปน สมเยน ฉพฺพคฺคิยา ภิกฺขู “ภควตา อุจฺจาสยน\-มหาสยนานิ ปฏิกฺขิตฺตานี”ติ มหาจมฺมานิ ธาเรนฺติ สีหจมฺมํ พฺยคฺฆจมฺมํ ทีปิจมฺมํ. ตานิ มญฺจปฺปมาเณนปิ ฉินฺนานิ โหนฺติ, ปีฐปฺปมาเณนปิ ฉินฺนานิ โหนฺติ, อนฺโตปิ มญฺเจ ปญฺญตฺตานิ โหนฺติ, พหิปิ มญฺเจ ปญฺญตฺตานิ โหนฺติ, อนฺโตปิ ปีเฐ ปญฺญตฺตานิ โหนฺติ, พหิปิ ปีเฐ ปญฺญตฺตานิ โหนฺติ. มนุสฺสา วิหารจาริกํ อาหิณฺฑนฺตา ปสฺสิตฺวา อุชฺฌายนฺติ ขิยฺยนฺติ วิปาเจนฺติ “เสยฺยถาปิ คิหี กามโภคิโน”ติ. ภควโต เอตมตฺถํ อาโรเจสุํ. น ภิกฺขเว มหาจมฺมานิ ธาเรตพฺพานิ, สีหจมฺมํ พฺยคฺฆจมฺมํ ทีปิจมฺมํ, โย ธาเรยฺย, อาปตฺติ ทุกฺกฏสฺสาติ.}

\prose{เตน โข ปน สมเยน ฉพฺพคฺคิยา ภิกฺขู “ภควตา มหาจมฺมานิ ปฏิกฺขิตฺตานี”ติ โคจมฺมานิ ธาเรนฺติ. ตานิ มญฺจปฺปมาเณนปิ ฉินฺนานิ โหนฺติ, ปีฐปฺปมาเณนปิ ฉินฺนานิ โหนฺติ, อนฺโตปิ มญฺเจ ปญฺญตฺตานิ โหนฺติ, พหิปิ มญฺเจ ปญฺญตฺตานิ โหนฺติ, อนฺโตปิ ปีเฐ ปญฺญตฺตานิ โหนฺติ, พหิปิ ปีเฐ ปญฺญตฺตานิ โหนฺติ. อญฺญตโรปิ ปาปภิกฺขุ อญฺญตรสฺส ปาปุปาสกสฺส กุลูปโก โหติ. อถ โข โส ปาปภิกฺขุ ปุพฺพณฺหสมยํ นิวาเสตฺวา ปตฺต\-จีวรมา\-ทาย เยน ตสฺส ปาปุปาสกสฺส นิเวสนํ เตนุปสงฺกมิ, อุปสงฺกมิตฺวา ปญฺญตฺเต อาสเน นิสีทิ. อถ โข โส ปาปุปาสโก เยน โส ปาปภิกฺขุ เตนุปสงฺกมิ, อุปสงฺกมิตฺวา ตํ ปาปภิกฺขุํ อภิวาเทตฺวา เอกมนฺตํ นิสีทิ. เตน โข ปน สมเยน ตสฺส ปาปุปาสกสฺส วจฺฉโก โหติ ตรุณโก อภิรูโป \csromanfolio{282}ทสฺสนีโย ปาสาทิโก จิโตฺร, เสยฺยถาปิ ทีปิจฺฉาโป. อถ โข โส ปาปภิกฺขุ ตํ วจฺฉกํ สกฺกจฺจํ อุปนิชฺฌายติ. อถ โข โส ปาปุปาสโก ตํ ปาปภิกฺขุํ เอตทโวจ “กิสฺส ภนฺเต อยฺโย อิมํ วจฺฉกํ สกฺกจฺจํ อุปนิชฺฌายตี”ติ. อตฺโถ เม อาวุโส อิมสฺส วจฺฉกสฺส จมฺเมนาติ. อถ โข โส ปาปุปาสโก ตํ วจฺฉกํ วธิตฺวา จมฺมํ วิธุนิตฺวา ตสฺส ปาปภิกฺขุโน ปาทาสิ. อถ โข โส ปาปภิกฺขุ ตํ จมฺมํ สงฺฆาฏิยา ปฏิจฺฉาเทตฺวา อคมาสิ. อถ โข สา คาวี วจฺฉคิทฺธินี ตํ ปาปภิกฺขุํ ปิฏฺฐิโต ปิฏฺฐิโต อนุพนฺธิ. ภิกฺขู เอวมาหํสุ “กิสฺส ตฺยายํ อาวุโส คาวี ปิฏฺฐิโต ปิฏฺฐิโต อนุพนฺธี”ติ. อหมฺปิ โข อาวุโส น ชานามิ “เกน\csromansharedfootnote{2f8d05c08c4b8f43}{เกนจิ (ก)} มฺยายํ คาวี ปิฏฺฐิโต ปิฏฺฐิโต อนุพนฺธี”ติ. เตน โข ปน สมเยน ตสฺส ปาปภิกฺขุโน สงฺฆาฏิ โลหิเตน มกฺขิตา โหติ. ภิกฺขู เอวมาหํสุ “อยํ ปน เต อาวุโส สํฆาฏิ กิํ กตา”ติ. อถ โข โส ปาปภิกฺขุ ภิกฺขูนํ เอตมตฺถํ อาโรเจสิ. กิํ ปน ตฺวํ อาวุโส ปาณาติปาเต สมาทเปสีติ. เอวมาวุโสติ. เย เต ภิกฺขู อปฺปิจฺฉา, เต อุชฺฌายนฺติ ขิยฺยนฺติ วิปาเจนฺติ “กถํ หิ นาม ภิกฺขุ ปาณาติปาเต สมาทเปสฺสติ, นนุ ภควตา อเนก\-ปริยาเยน ปาณาติปาโต ครหิโต, ปาณาติปาตา เวรมณี ปสตฺถา”ติ. อถ โข เต ภิกฺขู ภควโต เอตมตฺถํ อาโรเจสุํ.}

\prose{อถ โข ภควา เอตสฺมิํ นิทาเน เอตสฺมิํ ปกรเณ ภิกฺขุสํฆํ สนฺนิปาตาเปตฺวา ตํ ปาปภิกฺขุํ ปฏิปุจฺฉิ “สจฺจํ กิร ตฺวํ ภิกฺขุ ปาณาติปาเต สมาทเปสี”ติ. สจฺจํ ภควาติ ฯเปฯ. กถํ หิ นาม ตฺวํ โมฆปุริส ปาณาติปาเต สมาทเปสฺสสิ, นนุ มยา โมฆปุริส อเนก\-ปริยาเยน ปาณาติปาโต ครหิโต, ปาณาติปาตา เวรมณี ปสตฺถา. เนตํ โมฆปุริส อปฺปสนฺนานํ วา ปสาทาย ฯเปฯ วิครหิตฺวา ธมฺมิํ กถํ กตฺวา ภิกฺขู อามนฺเตสิ “น ภิกฺขเว ปาณาติปาเต สมาทเปตพฺพํ, โย สมาทเปยฺย, ยถาธมฺโม กาเรตพฺโพ. น ภิกฺขเว โคจมฺมํ ธาเรตพฺพํ, โย ธาเรยฺย, อาปตฺติ ทุกฺกฏสฺส. น จ ภิกฺขเว กิญฺจิ จมฺมํ ธาเรตพฺพํ, โย ธาเรยฺย, อาปตฺติ ทุกฺกฏสฺสา”ติ.}

\csromanmatikaanchor{mtk.161}
\csromantocmarkref{section}{156. คิหิวิกตานุญฺญาตาทิ}{mtk.161}
\bookmark[dest={mtk.161},level=1]{156. คิหิวิกตานุญฺญาตาทิ}
\csromanheader{1}{156. \textbf{คิหิวิกตา}\-\textbf{นุญฺญาตาทิ}}
\proseitem{256}{\csromanfolio{283}เตน โข ปน สมเยน มนุสฺสานํ มญฺจมฺปิ ปีฐมฺปิ จมฺโมนทฺธานิ โหนฺติ จมฺมวินทฺธานิ. ภิกฺขู กุกฺกุจฺจายนฺตา นาภินิสีทนฺติ. ภควโต เอตมตฺถํ อาโรเจสุํ. อนุชานามิ ภิกฺขเว คิหิวิกตํ อภินิสีทิตุํ, น เตฺวว อภินิ\-ปชฺชิ\-ตุนฺติ.}

\prose{เตน โข ปน สมเยน วิหารา จมฺมวทฺเธหิ โอคุมฺผิยนฺติ\csromansharedfootnote{727c1e871b5897ac}{โอคุมฺภิยนฺติ (ก)}. ภิกฺขู กุกฺกุจฺจายนฺตา นาภินิสีทนฺติ. ภควโต เอตมตฺถํ อาโรเจสุํ. อนุชานามิ ภิกฺขเว พนฺธนมตฺตํ อภินิ\-สีทิ\-ตุนฺติ.}

\prose{เตน โข ปน สมเยน ฉพฺพคฺคิยา ภิกฺขู สอุปาหนา คามํ ปวิสนฺติ. มนุสฺสา อุชฺฌายนฺติ ขิยฺยนฺติ วิปาเจนฺติ “เสยฺยถาปิ คิหี กามโภคิโน”ติ. ภควโต เอตมตฺถํ อาโรเจสุํ. น ภิกฺขเว สอุปาหเนน คาโม ปวิสิตพฺโพ, โย ปวิเสยฺย, อาปตฺติ ทุกฺกฏสฺสาติ.}

\prose{เตน โข ปน สมเยน อญฺญตโร ภิกฺขุ คิลาโน โหติ, น สกฺโกติ วินา อุปาหเนน คามํ ปวิสิตุํ.}

\prose{ภควโต เอตมตฺถํ อาโรเจสุํ. อนุชานามิ ภิกฺขเว คิลาเนน ภิกฺขุนา สอุปาหเนน คามํ ปวิสิตุนฺติ.}

\csromanmatikaanchor{mtk.162}
\csromantocmarkref{section}{157. โสณกุฏิกณฺณวตฺถุ}{mtk.162}
\bookmark[dest={mtk.162},level=1]{157. โสณกุฏิกณฺณวตฺถุ}
\csromanheader{1}{157. \textbf{โสณ}\-\textbf{กุฏิกณฺณ}\-\textbf{วตฺถุ}}
\proseitem{257}{\csromansymbolfootnote{*}{ขุ 1. 144 ปิฏฺเฐ อุทาเน โสณสุตฺเตน สํสนฺทิตฺวา ปสฺสิตพฺพํ.}เตน โข ปน สมเยน อายสฺมา มหากจฺจาโน อวนฺตีสุ วิหรติ กุรรฆเร\csromansharedfootnote{f6c9971650692bf0}{กุรุรฆเร (ก)} ปปตเก\csromansharedfootnote{814603cd4b043482}{ปปาเต (สี, สฺยา) ปวตฺเต (ขุ 1. 144)} ปพฺพเต. เตน โข ปน สมเยน โสโณ อุปาสโก กุฏิกณฺโณ อายสฺมโต มหากจฺจานสฺส อุปฏฺฐาโก โหติ. อถ โข โสโณ อุปาสโก กุฏิกณฺโณ เยนายสฺมา มหากจฺจาโน เตนุปสงฺกมิ, อุปสงฺกมิตฺวา อายสฺมนฺตํ มหากจฺจานํ อภิวาเทตฺวา เอกมนฺตํ นิสีทิ, เอกมนฺตํ นิสินฺโน โข โสโณ อุปาสโก กุฏิกณฺโณ อายสฺมนฺตํ มหากจฺจานํ เอตทโวจ “ยถา ยถาหํ ภนฺเต อยฺเยน มหากจฺจาเนน ธมฺมํ เทสิตํ อาชานามิ, นยิทํ สุกรํ อคารํ อชฺฌาวสตา เอกนฺต\-ปริปุณฺณํ เอกนฺต\-ปริสุทฺธํ สงฺขลิขิตํ พฺรหฺมจริยํ จริตุํ, อิจฺฉามหํ ภนฺเต เกสมสฺสุํ โอหาเรตฺวา กาสายานิ วตฺถานิ อจฺฉาเทตฺวา อคารสฺมา อนคาริยํ \csromanfolio{284}ปพฺพชิตุํ, ปพฺพาเชตุ มํ ภนฺเต อยฺโย มหากจฺจาโน”ติ. ( ) ทุกฺกรํ โข โสณ ยาวชีวํ เอกเสยฺยํ เอกภตฺตํ พฺรหฺมจริยํ จริตุํ, อิงฺฆ ตฺวํ โสณ ตตฺเถว อคาริกภูโต พุทฺธานํ สาสนํ อนุยุญฺช กาลยุตฺตํ เอกเสยฺยํ เอกภตฺตํ พฺรหฺม\-จริยนฺติ. อถ โข โสณสฺส อุปาสกสฺส กุฏิกณฺณสฺส โย อโหสิ ปพฺพชฺชา\-ภิสงฺขาโร, โส ปฏิปฺปสฺสมฺภิ. ทุติยมฺปิ โข โสโณ อุปาสโก กุฏิกณฺโณ ฯเปฯ. ตติยมฺปิ โข โสโณ อุปาสโก กุฏิกณฺโณ เยนายสฺมา มหากจฺจาโน เตนุปสงฺกมิ, อุปสงฺกมิตฺวา อายสฺมนฺตํ มหากจฺจานํ อภิวาเทตฺวา เอกมนฺตํ นิสีทิ, เอกมนฺตํ นิสินฺโน โข โสโณ อุปาสโก กุฏิกณฺโณ อายสฺมนฺตํ มหากจฺจานํ เอตทโวจ “ยถา ยถาหํ ภนฺเต อยฺเยน มหากจฺจาเนน ธมฺมํ เทสิตํ อาชานามิ, นยิทํ สุกรํ อคารํ อชฺฌาวสตา เอกนฺต\-ปริปุณฺณํ เอกนฺต\-ปริสุทฺธํ สงฺขลิขิตํ พฺรหฺมจริยํ จริตุํ, อิจฺฉามหํ ภนฺเต เกสมสฺสุํ โอหาเรตฺวา กาสายานิ วตฺถานิ อจฺฉาเทตฺวา อคารสฺมา อนคาริยํ ปพฺพาชิตุํ, ปพฺพาเชตุ มํ ภนฺเต อยฺโย มหากจฺจาโน”ติ. อถ โข อายสฺมา มหากจฺจาโน โสณํ อุปาสกํ กุฏิกณฺณํ ปพฺพาเชสิ. เตน โข ปน สมเยน อวนฺติ\-ทกฺขิณาปโถ อปฺปภิกฺขุโก โหติ. อถ โข อายสฺมา มหากจฺจาโน ติณฺณํ วสฺสานํ อจฺจเยน กิจฺเฉน กสิเรน ตโต ตโต ทสวคฺคํ ภิกฺขุสํฆํ สนฺนิปาตาเปตฺวา อายสฺมนฺตํ โสณํ อุปสมฺปาเทสิ.}

\csromanmatikaanchor{mtk.163}
\csromantocmarkref{section}{158. มหากจฺจานสฺส ปญฺจวรปริทสฺสนา}{mtk.163}
\bookmark[dest={mtk.163},level=1]{158. มหากจฺจานสฺส ปญฺจวรปริทสฺสนา}
\csromanheader{1}{158. \textbf{มหากจฺจานสฺส ปญฺจ}\-\textbf{วร}\-\textbf{ปริทสฺสนา}}
\prose{อถ โข อายสฺมโต โสณสฺส วสฺสํ\-วุฏฺฐสฺส รโหคตสฺส ปฏิสลฺลีนสฺส เอวํ เจตโส ปริวิตกฺโก อุทปาทิ “สุโตเยว โข เม ‘โส ภควา เอทิโส จ เอทิโส จา’ติ, น จ มยา สมฺมุขา ทิฏฺโฐ, คจฺเฉยฺยาหํ ตํ ภควนฺตํ ทสฺสนาย อรหนฺตํ สมฺมา\-สมฺพุทฺธํ, สเจ มํ อุปชฺฌาโย อนุชาเนยฺยา”ติ. อถ โข อายสฺมา โสโณ สายนฺหสมยํ ปฏิสลฺลานา วุฏฺฐิโต เยนายสฺมา มหากจฺจาโน เตนุปสงฺกมิ, อุปสงฺกมิตฺวา อายสฺมนฺตํ มหากจฺจานํ อภิวาเทตฺวา เอกมนฺตํ นิสีทิ, เอกมนฺตํ นิสินฺโน โข อายสฺมา โสโณ อายสฺมนฺตํ มหากจฺจานํ เอตทโวจ “อิธ มยฺหํ ภนฺเต}

\prose{\csromanfolio{285}รโหคตสฺส ปฏิสลฺลีนสฺส เอวํ เจตโส ปริวิตกฺโก อุทปาทิ ‘สุโต เยว โข เม ‘โส ภควา เอทิโส จ เอทิโส จา’ติ, น จ มยา สมฺมุขา ทิฏฺโฐ, คจฺเฉยฺยาหํ ตํ ภควนฺตํ ทสฺสนาย อรหนฺตํ สมฺมา\-สมฺพุทฺธํ, สเจ มํ อุปชฺฌาโย อนุชาเนยฺยา’ติ, คจฺเฉยฺยาหํ ภนฺเต ตํ ภควนฺตํ ทสฺสนาย อรหนฺตํ สมฺมา\-สมฺพุทฺธํ, สเจ มํ อุปชฺฌาโย อนุชานาตี”ติ. สาธุ สาธุ โสณ, คจฺฉ ตฺวํ โสณ ตํ ภควนฺตํ ทสฺสนาย รหนฺตํ สมฺมา\-สมฺพุทฺธํ, ทกฺขิสฺสสิ ตฺวํ โสณ ตํ ภควนฺตํ ปาสาทิกํ ปสาทนียํ สนฺตินฺทฺริยํ สนฺตมานสํ อุตฺตม\-ทมถ\-สมถํ อนุปฺปตฺตํ ทนฺตํ คุตฺตํ ยตินฺทฺริยํ นาคํ, เตน หิ ตฺวํ โสณ มม วจเนน ภควโต ปาเท สิรสา วนฺท “อุปชฺฌาโย เม ภนฺเต อายสฺมา มหากจฺจาโน ภควโต ปาเท สิรสา วนฺทตี”ติ, เอวญฺจ วเทหิ “อวนฺติ\-ทกฺขิณาปโถ ภนฺเต อปฺปภิกฺขุโก, ติณฺณํ เม วสฺสานํ อจฺจเยน กิจฺเฉน กสิเรน ตโต ตโต ทสวคฺคํ ภิกฺขุสํฆํ สนฺนิปาตาเปตฺวา อุปสมฺปทํ อลตฺถํ, อปฺเปว นาม ภควา อวนฺติ\-ทกฺขิณาปเถ อปฺปตเรน คเณน อุปสมฺปทํ อนุชาเนยฺย. อวนฺติ\-ทกฺขิณาปเถ ภนฺเต กณฺหุตฺตรา ภูมิ ขรา โคกณฺฏกหตา, อปฺเปว นาม ภควา อวนฺติ\-ทกฺขิณาปเถ คุณงฺคุณู\-ปาหนํ อนุชาเนยฺย. อวนฺติ\-ทกฺขิณาปเถ ภนฺเต นหานครุกา มนุสฺสา อุทกสุทฺธิกา, อปฺเปว นาม ภควา อวนฺติ\-ทกฺขิณาปเถ ธุวนหานํ อนุชาเนยฺย. อวนฺติ\-ทกฺขิณาปเถ ภนฺเต จมฺมานิ อตฺถรณานิ เอฬกจมฺมํ อชจมฺมํ มิคจมฺมํ, เสยฺยถาปิ ภนฺเต มชฺฌิเมสุ ชนปเทสุ เอรคู โมรคู มชฺชารู\csromansharedfootnote{4be5609d69e0c7cc}{มชฺฌารู (ก)} ชนฺตู, เอวเมว โข ภนฺเต อวนฺติ\-ทกฺขิณาปเถ จมฺมานิ อตฺถรณานิ เอฬกจมฺมํ อชจมฺมํ มิคจมฺมํ, อปฺเปว นาม ภควา อวนฺติ\-ทกฺขิณาปเถ จมฺมานิ อตฺถรณานิ อนุชาเนยฺย เอฬกจมฺมํ อชจมฺมํ มิคจมฺมํ. เอตรหิ ภนฺเต มนุสฺสา นิสฺสีมคตานํ ภิกฺขูนํ จีวรํ เทนฺติ ‘อิมํ จีวรํ อิตฺถนฺนามสฺส เทมา’ติ, เต อาคนฺตฺวา อาโรเจนฺติ ‘อิตฺถนฺนาเมหิ เต อาวุโส มนุสฺเสหิ จีวรํ ทินฺนนฺ’ติ, เต กุกฺกุจฺจายนฺตา น สาทิยนฺติ ‘มา โน นิสฺสคฺคิยํ อโหสี’ติ, อปฺเปว นาม ภควา จีวเร ปริยายํ อาจิกฺเขยฺยา”ติ. “เอวํ ภนฺเต”ติ โข อายสฺมา โสโณ อายสฺมโต มหากจฺจานสฺส ปฏิสฺสุตฺวา อุฏฺฐายาสนา อายสฺมนฺตํ มหากจฺจานํ อภิวาเทตฺวา ปทกฺขิณํ กตฺวา เสนาสนํ สํสาเมตฺวา}

\prose{\csromanfolio{286}ปตฺต\-จีวรมา\-ทาย เยน สาวตฺถิ เตน ปกฺกามิ, อนุปุพฺเพน เยน สาวตฺถิ เชตวนํ อนาถ\-ปิณฺฑิกสฺส อาราโม เยน ภควา เตนุปสงฺกมิ, อุปสงฺกมิตฺวา ภควนฺตํ อภิวาเทตฺวา เอกมนฺตํ นิสีทิ. อถ โข ภควา อายสฺมนฺตํ อานนฺทํ อามนฺเตสิ “อิมสฺส อานนฺท อาคนฺตุกสฺส ภิกฺขุโน เสนาสนํ ปญฺญาเปหี”ติ. อถ โข อายสฺมา อานนฺโท “ยสฺส โข มํ ภควา อาณาเปติ ‘อิมสฺส อานนฺท อาคนฺตุกสฺส ภิกฺขุโน เสนาสนํ ปญฺญาเปหี’ติ, อิจฺฉติ ภควา เตน ภิกฺขุนา สทฺธิํ เอกวิหาเร วตฺถุํ, อิจฺฉติ ภควา อายสฺมตา โสเณน สทฺธิํ เอกวิหาเร วตฺถุนฺ”ติ ยสฺมิํ วิหาเร ภควา วิหรติ, ตสฺมิํ วิหาเร อายสฺมโต โสณสฺส เสนาสนํ ปญฺญาเปสิ.}

\proseitem{258}{อถ โข ภควา พหุเทว รตฺติํ อชฺโฌกาเส วีตินาเมตฺวา วิหารํ ปาวิสิ, อายสฺมาปิ โข โสโณ พหุเทว รตฺติํ อชฺโฌกาเส วีตินาเมตฺวา วิหารํ ปาวิสิ. อถ โข ภควา รตฺติยา ปจฺจูสสมยํ ปจฺจุฏฺฐาย อายสฺมนฺตํ โสณํ อชฺเฌสิ “ปฏิภาตุ ตํ ภิกฺขุ ธมฺโม ภาสิตุนฺ”ติ. “เอวํ ภนฺเต”ติ โข อายสฺมา โสโณ ภควโต ปฏิสฺสุณิตฺวา สพฺพาเนว อฏฺฐก\-วคฺคิ\-กานิ สเรน อภาสิ. อถ โข ภควา อายสฺมโต โสณสฺส สรภญฺญ\-ปริโยสาเน อพฺภานุโมทิ, สาธุ สาธุ ภิกฺขุ, สุคฺคหิตานิ โข เต ภิกฺขุ อฏฺฐก\-วคฺคิ\-กานิ สุมนสิกตานิ สูปธาริตานิ, กลฺยาณิยาปิ วาจาย สมนฺนาคโต วิสฺสฏฺฐาย อเนลคลาย\csromansharedfootnote{fef8bdd933223291}{อเนฬคลาย (ก)} อตฺถสฺส วิญฺญาปนิยา, กติวสฺโสสิ ตฺวํ ภิกฺขูติ. เอกวสฺโสหํ ภควาติ. กิสฺส ปน ตฺวํ ภิกฺขุ เอวํ จิรํ อกาสีติ. จิรํ ทิฏฺโฐ เม ภนฺเต กาเมสุ อาทีนโว, อปิ จ สมฺพาธา ฆราวาสา พหุกิจฺจา พหุกรณียาติ. อถ โข ภควา เอตมตฺถํ วิทิตฺวา ตายํ เวลายํ อิมํ อุทานํ อุทาเนสิ\csromandash{}}

\setlength{\csromangathapagewidth}{0pt}
\setlength{\csromangathawakwidth}{0pt}
\csromangathameasuregroup{\textsuperscript{*}“ทิสฺวา อาทีนวํ โลเก, \\ อริโย น รมตี ปาเป,}{\csromangathabat{\textsuperscript{*}“ทิสฺวา อาทีนวํ โลเก,}{ญตฺวา ธมฺมํ นิรูปธิํ.} \\ \csromangathabat{อริโย น รมตี ปาเป,}{ปาเป น รมตี สุจี”ติ.}}
\csromangathasetleft{\textsuperscript{*}“ทิสฺวา อาทีนวํ โลเก, \\ อริโย น รมตี ปาเป,}
\csromangathagroup{\csromangathastanza{\csromangathabat{\csromansymbolfootnote{*}{ขุ 1. 148 ปิฏฺเฐ อุทาเนปิ.}“ทิสฺวา อาทีนวํ โลเก,}{ญตฺวา ธมฺมํ นิรูปธิํ.} \\ \csromangathabat{อริโย น รมตี ปาเป,}{ปาเป น รมตี สุจี”ติ.}}}

\prose{อถ โข อายสฺมา โสโณ “ปฏิสมฺโมทติ โข มํ ภควา, อยํ ขฺวสฺส กาโล, ยํ เม อุปชฺฌาโย ปริทสฺสี”ติ อุฏฺฐายาสนา เอกํสํ อุตฺตราสงฺคํ กริตฺวา ภควโต ปาเทสุ สิรสา นิปติตฺวา \csromanfolio{287}ภควนฺตํ เอตทโวจ “อุปชฺฌาโย เม ภนฺเต อายสฺมา มหากจฺจาโน ภควโต ปาเท สิรสา วนฺทติ, เอวญฺจ วเทติ ‘อวนฺติ\-ทกฺขิณาปโถ ภนฺเต อปฺปภิกฺขุโก, ติณฺณํ เม วสฺสานํ อจฺจเยน กิจฺเฉน กสิเรน ตโต ตโต ทสวคฺคํ ภิกฺขุสํฆํ สนฺนิปาตาเปตฺวา อุปสมฺปทํ อลตฺถํ, อปฺเปว นาม ภควา อวนฺติ\-ทกฺขิณาปเถ อปฺปตเรน คเณน อุปสมฺปทํ อนุชาเนยฺย. อวนฺติ\-ทกฺขิณาปเถ ภนฺเต กณฺหุตฺตรา ภูมิ ขรา โคกณฺฏกหตา, อปฺเปว นาม ภควา อวนฺติ\-ทกฺขิณาปเถ คุณงฺคุณู\-ปาหนํ อนุชาเนยฺย. อวนฺติ\-ทกฺขิณาปเถ ภนฺเต นหานครุกา มนุสฺสา อุทกสุทฺธิกา, อปฺเปว นาม ภควา อวนฺติ\-ทกฺขิณาปเถ ธุวนหานํ อนุชาเนยฺย. อวนฺติ\-ทกฺขิณาปเถ ภนฺเต จมฺมานิ อตฺถรณานิ เอฬกจมฺมํ อชจมฺมํ มิคจมฺมํ, เสยฺยถาปิ ภนฺเต มชฺฌิเมสุ ชนปเทสุ เอรคู โมรคู มชฺชารู ชนฺตู, เอวเมว โข ภนฺเต อวนฺติ\-ทกฺขิณาปเถ จมฺมานิ อตฺถรณานิ เอฬกจมฺมํ อชจมฺมํ มิคจมฺมํ, อปฺเปว นาม ภควา อวนฺติ\-ทกฺขิณาปเถ จมฺมานิ อตฺถรณานิ อนุชาเนยฺย เอฬกจมฺมํ อชจมฺมํ มิคจมฺมํ. เอตรหิ ภนฺเต มนุสฺสา นิสฺสีมคตานํ ภิกฺขูนํ จีวรํ เทนฺติ ‘อิมํ จีวรํ อิตฺถนฺนามสฺส เทมา’ติ, เต อาคนฺตฺวา อาโรเจนฺติ ‘อิตฺถนฺนาเมหิ เต อาวุโส มนุสฺเสหิ จีวรํ ทินฺนนฺ’ติ, เต กุกฺกุจฺจายนฺตา น สาทิยนฺติ ‘มา โน นิสฺสคฺคิยํ อโหสี’ติ, อปฺเปว นาม ภควา จีวเร ปริยายํ อาจิกฺเขยฺยา”ติ.}

\csromanmatikaanchor{mtk.164}
\csromantocmarkref{section}{159. อุทฺทานคาถา}{mtk.164}
\bookmark[dest={mtk.164},level=1]{159. อุทฺทานคาถา}
\proseitem{259}{อถ โข ภควา เอตสฺมิํ นิทาเน เอตสฺมิํ ปกรเณ ธมฺมิํ กถํ กตฺวา ภิกฺขู อามนฺเตสิ “อวนฺติ\-ทกฺขิณาปโถ ภิกฺขเว อปฺปภิกฺขุโก, อนุชานามิ ภิกฺขเว สพฺพ\-ปจฺจนฺติเมสุ ชนปเทสุ วินยธร\-ปญฺจเมน คเณน อุปสมฺปทํ. ตตฺริเม ปจฺจนฺติมา ชนปทา, ปุรตฺถิมาย ทิสาย คชงฺคลํ\csromansharedfootnote{ade80d65126ec54f}{กชงฺคลํ (สี, สฺยา)} นาม นิคโม, ตสฺส ปเรน มหาสาลา, ตโต ปรา ปจฺจนฺติมา ชนปทา, โอรโต มชฺเฌ. ปุรตฺถิม\-ทกฺขิณาย ทิสาย สลฺลวตี\csromansharedfootnote{17fac1ff7f3f642a}{สลลวตี (สี)} นาม นที, ตโต ปรา ปจฺจนฺติมา ชนปทา, โอรโต มชฺเฌ. ทกฺขิณาย ทิสาย เสตกณฺณิกํ นาม นิคโม, ตโต ปรา ปจฺจนฺติมา ชนปทา, โอรโต มชฺเฌ. ปจฺฉิมาย ทิสาย ถูณํ นาม พฺราหฺมณคาโม, ตโต ปรา ปจฺจนฺติมา ชนปทา, โอรโต มชฺเฌ. \csromanfolio{288}อุตฺตราย ทิสาย อุสีรทฺธโช นาม ปพฺพโต, ตโต ปรา ปจฺจนฺติมา ชนปทา, โอรโต มชฺเฌ. อนุชานามิ ภิกฺขเว เอวรูเปสุ ปจฺจนฺติเมสุ ชนปเทสุ วินยธร\-ปญฺจเมน คเณน อุปสมฺปทํ. อวนฺติ\-ทกฺขิณาปเถ ภิกฺขเว กณฺหุตฺตรา ภูมิ ขรา โคกณฺฏกหตา, อนุชานามิ ภิกฺขเว สพฺพ\-ปจฺจนฺติเมสุ ชนปเทสุ คุณงฺคุณู\-ปาหนํ. อวนฺติ\-ทกฺขิณาปเถ ภิกฺขเว นหานครุกา มนุสฺสา อุทกสุทฺธิกา, อนุชานามิ ภิกฺขเว สพฺพ\-ปจฺจนฺติเมสุ ชนปเทสุ ธุวนหานํ. อวนฺติ\-ทกฺขิณาปเถ ภิกฺขเว จมฺมานิ อตฺถรณานิ เอฬกจมฺมํ อชจมฺมํ มิคจมฺมํ, เสยฺยถาปิ ภิกฺขเว มชฺฌิเมสุ ชนปเทสุ เอรคู โมรคู มชฺชารู ชนฺตู, เอวเมว โข ภิกฺขเว อวนฺติ\-ทกฺขิณาปเถ จมฺมานิ อตฺถรณานิ เอฬกจมฺมํ อชจมฺมํ มิคจมฺมํ, อนุชานามิ ภิกฺขเว สพฺพ\-ปจฺจนฺติเมสุ ชนปเทสุ จมฺมานิ อตฺถรณานิ เอฬกจมฺมํ อชจมฺมํ มิคจมฺมํ. อิธ ปน ภิกฺขเว มนุสฺสา นิสฺสีมคตานํ \textbf{ภิกฺขูนํ จีวรํ เทนฺติ ‘อิมํ จีวรํ อิตฺถนฺนามสฺส เทมา’ติ,} \textbf{อนุชานามิ ภิกฺขเว สาทิตุํ, น ตาว ตํ คณนูปคํ ยาว น หตฺถํ} คจฺฉตี”ติ.}

\vspace{\tipitakaparskip}
\csromancenter{จมฺม\-กฺขนฺธโก ปญฺจโม.}
\csromancenter{159. \textbf{ตสฺสุทฺทานํ}}
\setlength{\csromangathapagewidth}{0pt}
\setlength{\csromangathawakwidth}{0pt}
\csromangathameasuregroup{ราชา จ มาคโธ โสโณ, \\ สาคโต คิชฺฌกูฏสฺมิํ, \\ ปพฺพชฺชารทฺธ ภิชฺชิํสุ, \\ นีลา ปีตา โลหิติกา, \\ มหารงฺคมหานามา, \\ ขลฺลกา ปุฏปาลิ จ, \\ วิจฺฉิกา โมรจิตฺรา จ, \\ อชินุทฺทา มชฺชารี จ, \\ ผลิตุปาหนา ขิลา, \\ ตาลเวฬุติณํ เจว, \\ กมลกมฺพลโสวณฺณา, \\ ผลิกา กํสกาจา จ, \\ คาวี ยานํ คิลาโน จ, \\ สยนานิ มหาจมฺมา, \\ คิหีนํ จมฺมวทฺเธหิ, \\ มหากจฺจายโน โสโณ, \\ อุปสมฺปทํ ปญฺจหิ, \\ จมฺมตฺถรณานุญฺญาสิ, \\ อทาสิ เม วเร ปญฺจ,}{\csromangathabat{ราชา จ มาคโธ โสโณ,}{อสีติสหสฺสิสฺสโร.} \\ \csromangathabat{สาคโต คิชฺฌกูฏสฺมิํ,}{พหุํ ทสฺเสติ อุตฺตริํ.} \\ \csromangathabat{ปพฺพชฺชารทฺธ ภิชฺชิํสุ,}{วีณํ เอกปลาสิกํ.} \\ \csromangathabat{นีลา ปีตา โลหิติกา,}{มญฺชิฏฺฐา กณฺหเมว จ.} \\ \csromangathabat{มหารงฺคมหานามา,}{วทฺธิกา จ ปฏิกฺขิปิ.} \\ \csromangathabat{ขลฺลกา ปุฏปาลิ จ,}{ตูลติตฺติรเมณฺฑชา.} \\ \csromangathabat{วิจฺฉิกา โมรจิตฺรา จ,}{สีหพฺยคฺฆา จ ทีปิกา.} \\ \csromangathabat{อชินุทฺทา มชฺชารี จ,}{กาฬลุวกปริกฺขฏา.} \\ \csromangathabat{ผลิตุปาหนา ขิลา,}{โธตขาณุขฏขฏา.} \\ \csromangathabat{ตาลเวฬุติณํ เจว,}{มุญฺชปพฺพชหินฺตาลา.} \\ \csromangathabat{กมลกมฺพลโสวณฺณา,}{รูปิกา มณิเวฬุริยา.} \\ \csromangathabat{ผลิกา กํสกาจา จ,}{ติปุสีสญฺจ ตมฺพกา.} \\ \csromangathabat{คาวี ยานํ คิลาโน จ,}{ปุริสายุตฺตสิวิกา.} \\ \csromangathabat{สยนานิ มหาจมฺมา,}{โคจมฺเมหิ จ ปาปโก.} \\ \csromangathabat{คิหีนํ จมฺมวทฺเธหิ,}{ปวิสนฺติ คิลายโน.} \\ \csromangathabat{มหากจฺจายโน โสโณ,}{สเรน อฏฺฐกวคฺคิกํ.} \\ \csromangathabat{อุปสมฺปทํ ปญฺจหิ,}{คุณงฺคุณา ธุวสินา.} \\ \csromangathabat{จมฺมตฺถรณานุญฺญาสิ,}{น ตาว คณนูปคํ.} \\ \csromangathabat{อทาสิ เม วเร ปญฺจ,}{โสณตฺเถรสฺส นายโกติ.} \\ อิมมฺหิ ขนฺธเก วตฺถูนิ เตสฏฺฐิ.}
\csromangathasetleft{ราชา จ มาคโธ โสโณ, \\ สาคโต คิชฺฌกูฏสฺมิํ, \\ ปพฺพชฺชารทฺธ ภิชฺชิํสุ, \\ นีลา ปีตา โลหิติกา, \\ มหารงฺคมหานามา, \\ ขลฺลกา ปุฏปาลิ จ, \\ วิจฺฉิกา โมรจิตฺรา จ, \\ อชินุทฺทา มชฺชารี จ, \\ ผลิตุปาหนา ขิลา, \\ ตาลเวฬุติณํ เจว, \\ กมลกมฺพลโสวณฺณา, \\ ผลิกา กํสกาจา จ, \\ คาวี ยานํ คิลาโน จ, \\ สยนานิ มหาจมฺมา, \\ คิหีนํ จมฺมวทฺเธหิ, \\ มหากจฺจายโน โสโณ, \\ อุปสมฺปทํ ปญฺจหิ, \\ จมฺมตฺถรณานุญฺญาสิ, \\ อทาสิ เม วเร ปญฺจ,}
\csromangathagroupwithclosermajor{\csromangathastanza{\csromangathabat{ราชา จ มาคโธ โสโณ,}{อสีติสหสฺสิสฺสโร.} \\ \csromangathabat{สาคโต คิชฺฌกูฏสฺมิํ,}{พหุํ ทสฺเสติ อุตฺตริํ.}}\csromangathastanza{\csromangathabat{ปพฺพชฺชารทฺธ ภิชฺชิํสุ,}{วีณํ เอกปลาสิกํ.} \\ \csromangathabat{นีลา ปีตา โลหิติกา,}{มญฺชิฏฺฐา กณฺหเมว จ.}}\csromangathastanza{\csromangathabat{มหารงฺคมหานามา,}{วทฺธิกา จ ปฏิกฺขิปิ.} \\ \csromangathabat{ขลฺลกา ปุฏปาลิ จ,}{ตูลติตฺติรเมณฺฑชา.}}\csromangathastanza{\csromangathabat{วิจฺฉิกา โมรจิตฺรา จ,}{สีหพฺยคฺฆา จ ทีปิกา.} \\ \csromangathabat{อชินุทฺทา มชฺชารี จ,}{กาฬลุวกปริกฺขฏา.}}\csromangathastanza{\csromangathabat{ผลิตุปาหนา ขิลา,}{โธตขาณุขฏขฏา.} \\ \csromangathabat{ตาลเวฬุติณํ เจว,}{มุญฺชปพฺพชหินฺตาลา.}}\csromangathastanza{\csromanfolio{289}\csromangathabat{กมลกมฺพลโสวณฺณา,}{รูปิกา มณิเวฬุริยา.} \\ \csromangathabat{ผลิกา กํสกาจา จ,}{ติปุสีสญฺจ ตมฺพกา.}}\csromangathastanza{\csromangathabat{คาวี ยานํ คิลาโน จ,}{ปุริสายุตฺตสิวิกา.} \\ \csromangathabat{สยนานิ มหาจมฺมา,}{โคจมฺเมหิ จ ปาปโก.}}\csromangathastanza{\csromangathabat{คิหีนํ จมฺมวทฺเธหิ,}{ปวิสนฺติ คิลายโน.} \\ \csromangathabat{มหากจฺจายโน โสโณ,}{สเรน อฏฺฐกวคฺคิกํ.}}\csromangathastanza{\csromangathabat{อุปสมฺปทํ ปญฺจหิ,}{คุณงฺคุณา ธุวสินา.} \\ \csromangathabat{จมฺมตฺถรณานุญฺญาสิ,}{น ตาว คณนูปคํ.}}\csromangathastanza{\csromangathabat{อทาสิ เม วเร ปญฺจ,}{โสณตฺเถรสฺส นายโกติ.} \\ อิมมฺหิ ขนฺธเก วตฺถูนิ เตสฏฺฐิ.}}{\textbf{จมฺม}\-\textbf{กฺขนฺธโก นิฏฺฐิโต.}}
\csromanmatikaanchor{mtk.165}
\csromantocmarknopage{chapter}{6. เภสชฺชกฺขนฺธก}{mtk.165}
\bookmark[dest={mtk.165},level=0]{6. เภสชฺชกฺขนฺธก}
\chapterheadpageread{6. \textbf{เภสชฺช}\-\textbf{กฺขนฺธก}}
\csromanmatikaanchor{mtk.166}
\csromantocmarkref{section}{160. ปญฺจเภสชฺชกถา}{mtk.166}
\bookmark[dest={mtk.166},level=1]{160. ปญฺจเภสชฺชกถา}
\csromanheader{1}{160. \textbf{ปญฺจ}\-\textbf{เภสชฺช}\-\textbf{กถา}}
\proseitem{260}{\csromanfolio{290}เตน สมเยน พุทฺโธ ภควา สาวตฺถิยํ วิหรติ เชตวเน อนาถ\-ปิณฺฑิกสฺส อาราเม. เตน โข ปน สมเยน ภิกฺขูนํ สารทิเกน อาพาเธน ผุฏฺฐานํ ยาคุปิ ปีตา อุคฺคจฺฉติ, ภตฺตมฺปิ ภุตฺตํ อุคฺคจฺฉติ, เต เตน กิสา โหนฺติ ลูขา ทุพฺพณฺณา อุปฺปณฺฑุ\-ปฺปณฺฑุก\-ชาตา ธมนิสนฺถต\-คตฺตา. อทฺทสา โข ภควา เต ภิกฺขู กิเส ลูเข ทุพฺพณฺเณ อุปฺปณฺฑุ\-ปฺปณฺฑุก\-ชาเต ธมนิสนฺถต\-คตฺเต, ทิสฺวาน อายสฺมนฺตํ อานนฺทํ อามนฺเตสิ “กิํ นุ โข อานนฺท เอตรหิ ภิกฺขู กิสา ลูขา ทุพฺพณฺณา อุปฺปณฺฑุ\-ปฺปณฺฑุก\-ชาตา ธมนิสนฺถต\-คตฺตา”ติ. เอตรหิ ภนฺเต ภิกฺขูนํ สารทิเกน อาพาเธน ผุฏฺฐานํ ยาคุปิ ปีตา อุคฺคจฺฉติ, ภตฺตมฺปิ ภุตฺตํ อุคฺคจฺฉติ, เต เตน กิสา โหนฺติ ลูขา ทุพฺพณฺณา อุปฺปณฺฑุ\-ปฺปณฺฑุก\-ชาตา ธมนิสนฺถต\-คตฺตาติ. อถ โข ภควโต รโหคตสฺส ปฏิสลฺลีนสฺส เอวํ เจตโส ปริวิตกฺโก อุทปาทิ “เอตรหิ โข ภิกฺขูนํ สารทิเกน อาพาเธน ผุฏฺฐานํ ยาคุปิ ปีตา อุคฺคจฺฉติ, ภตฺตมฺปิ ภุตฺตํ อุคฺคจฺฉติ, เต เตน กิสา โหนฺติ ลูขา ทุพฺพณฺณา อุปฺปณฺฑุ\-ปฺปณฺฑุก\-ชาตา ธมนิสนฺถต\-คตฺตา, กิํ นุ โข อหํ ภิกฺขูนํ เภสชฺชํ อนุชาเนยฺยํ, ยํ เภสชฺชญฺเจว อสฺส เภสชฺชสมฺมตญฺจ โลกสฺส, อาหารตฺถญฺจ ผเรยฺย, น จ โอฬาริโก อาหาโร ปญฺญาเยยฺยา”ติ. อถ โข ภควโต เอตทโหสิ “อิมานิ โข ปญฺจ เภสชฺชานิ, เสยฺยถิทํ, สปฺปิ นวนีตํ เตลํ มธุ ผาณิตํ, เภสชฺชานิ เจว เภสชฺช\-สมฺมตานิ จ โลกสฺส, อาหารตฺถญฺจ ผรนฺติ, น จ โอฬาริโก อาหาโร ปญฺญายติ, ยนฺนูนาหํ ภิกฺขูนํ อิมานิ ปญฺจ เภสชฺชานิ อนุชาเนยฺยํ กาเล ปฏิคฺคเหตฺวา กาเล ปริภุญฺชิตุนฺ”ติ. อถ โข ภควา สายนฺหสมยํ ปฏิสลฺลาน วุฏฺฐิโต เอตสฺมิํ นิทาเน เอตสฺมิํ ปกรเณ ธมฺมิํ กถํ กตฺวา ภิกฺขู อามนฺเตสิ “อิธ มยฺหํ ภิกฺขเว รโหคตสฺส ปฏิสลฺลีนสฺส เอวํ เจตโส ปริวิตกฺโก อุทปาทิ ‘เอตรหิ โข ภิกฺขูนํ สารทิเกน อาพาเธน ผุฏฺฐานํ ยาคุปิ ปีตา อุคฺคจฺฉติ, ภตฺตมฺปิ ภุตฺตํ อุคฺคจฺฉติ, เต เตน กิสา โหนฺติ ลูขา ทุพฺพณฺณา อุปฺปุณฺฑุปฺปณฺฑุกชาตา ธมนิสนฺถต\-คตฺตา, กิํ นุ โข อหํ ภิกฺขูนํ เภสชฺชํ อนุชาเนยฺยํ, ยํ เภสชฺชญฺเจว อสฺส เภสชฺชสมฺมตญฺจ โลกสฺส, \csromanfolio{291}อาหารตฺถญฺจ ผเรยฺย, น จ โอฬาริโก อาหาโร ปญฺญาเยยฺยา’ติ. ตสฺส มยฺหํ ภิกฺขเว เอตทโหสิ ‘อิมานิ โข ปญฺจ เภสชฺชานิ, เสยฺยถิทํ, สปฺปิ นวนีตํ เตลํ มธุ ผาณิตํ, เภสชฺชานิ เจว เภสชฺช\-สมฺมตานิ จ โลกสฺส, อาหารตฺถญฺจ ผรนฺติ, น จ โอฬาริโก อาหาโร ปญฺญายติ, ยนฺนูนาหํ ภิกฺขูนํ อิมานิ ปญฺจ เภสชฺชานิ อนุชาเนยฺยํ กาเล ปฏิคฺคเหตฺวา กาเล ปริภุญฺชิตุนฺ’ติ, อนุชานามิ ภิกฺขเว ตานิ ปญฺจ เภสชฺชานิ กาเล ปฏิคฺคเหตฺวา กาเล ปริภุญฺชิตุนฺ”ติ.}

\proseitem{261}{เตน โข ปน สมเยน ภิกฺขู ตานิ ปญฺจ เภสชฺชานิ กาเล ปฏิคฺคเหตฺวา กาเล ปริภุญฺชนฺติ. เตสํ ยานิปิ ตานิ ปากติกานิ ลูขานิ โภชนานิ, ตานิปิ นจฺฉาเทนฺติ, ปเคว เสเนสิตานิ\csromansharedfootnote{31e78f251afbea5a}{เสเนสิกานิ (สี, สฺยา), เสเนหิกานิ (โยชนา)}. เต เตน เจว สารทิเกน อาพาเธน ผุฏฺฐา อิมินา จ ภตฺตา\-จฺฉาทเกน\csromansharedfootnote{c282fd4d13dba400}{ภตฺตาจฺฉนฺนเกน (ก)} ตทุภเยน ภิยฺโยโส มตฺตาย กิสา โหนฺติ ลูขา ทุพฺพณฺณา อุปฺปณฺฑุ\-ปฺปณฺฑุก\-ชาตา ธมนิสนฺถต\-คตฺตา. อทฺทสา โข ภควา เต ภิกฺขู ภิยฺโยโส มตฺตาย กิเส ลูเข ทุพฺพณฺเณ อุปฺปณฺฑุ\-ปฺปณฺฑุก\-ชาเต ธมนิสนฺถต\-คตฺเต, ทิสฺวาน อายสฺมนฺตํ อานนฺทํ อามนฺเตสิ “กิํ นุ โข อานนฺท เอตรหิ ภิกฺขู ภิยฺโยโส มตฺตาย กิสา ลูขา ทุพฺพณฺณา อุปฺปณฺฑุปฺปณฺฑุชาตา ธมนิสนฺถต\-คตฺตา”ติ. เอตรหิ ภนฺเต ภิกฺขู ตานิ จ ปญฺจ เภสชฺชานิ กาเล ปฏิคฺคเหตฺวา กาเล ปริภุญฺชนฺติ. เตสํ ยานิปิ ตานิ ปากติกานิ ลูขานิ โภชนานิ, ตานิปิ นจฺฉาเทนฺติ, ปเคว เสเนสิตานิ. เต เตน เจว สารทิเกน อาพาเธน ผุฏฺฐา อิมินา จ ภตฺตา\-จฺฉาทเกน ตทุภเยน ภิยฺโยโส มตฺตาย กิสา ลูขา ทุพฺพณฺณา อุปฺปณฺฑุ\-ปฺปณฺฑุก\-ชาตา ธมนิสนฺถต\-คตฺตาติ. อถ โข ภควา เอตสฺมิํ นิทาเน เอตสฺมิํ ปกรเณ ธมฺมิํ กถํ กตฺวา ภิกฺขู อามนฺเตสิ “อนุชานามิ ภิกฺขเว ตานิ ปญฺจ เภสชฺชานิ ปฏิคฺคเหตฺวา กาเลปิ วิกาเลปิ ปริภุญฺชิตุนฺ”ติ.}

\proseitem{262}{เตน โข ปน สมเยน คิลานานํ ภิกฺขูนํ วเสหิ เภสชฺเชหิ อตฺโถ โหติ. ภควโต เอตมตฺถํ อาโรเจสุํ. อนุชานามิ ภิกฺขเว วสานิ เภสชฺชานิ อจฺฉวสํ มจฺฉวสํ \csromanfolio{292}สุสุกาวสํ สูกรวสํ คทฺรภวสํ กาเล ปฏิคฺคหิตํ กาเล นิปฺปกฺกํ\csromansharedfootnote{fdce9f6abd8864d4}{นิปกฺกํ (ก)} กาเล สํสฏฺฐํ เตล\-ปริโภเคน ปริภุญฺชิตุํ. วิกาเล เจ ภิกฺขเว ปฏิคฺคหิตํ วิกาเล นิปฺปกฺกํ วิกาเล สํสฏฺฐํ, ตํ เจ ปริภุญฺเชยฺย, อาปตฺติ ติณฺณํ ทุกฺกฏานํ. กาเล เจ ภิกฺขเว ปฏิคฺคหิตํ วิกาเล นิปฺปกฺกํ วิกาเล สํสฏฺฐํ, ตํ เจ ปริภุญฺเชยฺย, อาปตฺติ ทฺวินฺนํ ทุกฺกฏานํ. กาเล เจ ภิกฺขเว ปฏิคฺคหิตํ กาเล นิปฺปกฺกํ วิกาเล สํสฏฺฐํ, ตํ เจ ปริภุญฺเชยฺย, อาปตฺติ ทุกฺกฏสฺส. กาเล เจ ภิกฺขเว ปฏิคฺคหิตํ กาเล นิปฺปกฺกํ กาเล สํสฏฺฐํ, ตํ เจ ปริภุญฺเชยฺย, อนาปตฺตีติ.}

\csromanmatikaanchor{mtk.167}
\csromantocmarkref{section}{161. มูลาทิเภสชฺชกถา}{mtk.167}
\bookmark[dest={mtk.167},level=1]{161. มูลาทิเภสชฺชกถา}
\csromanheader{1}{161. \textbf{มูลา}\-\textbf{ทิ}\-\textbf{เภสชฺช}\-\textbf{กถา}}
\proseitem{263}{เตน โข ปน สมเยน คิลานานํ ภิกฺขูนํ มูเลหิ เภสชฺเชหิ อตฺโถ โหติ. ภควโต เอตมตฺถํ อาโรเจสุํ. อนุชานามิ ภิกฺขเว มูลานิ เภสชฺชานิ หลิทฺทิํ สิงฺคิเวรํ วจํ วจตฺตํ\csromansharedfootnote{04526b2f6dbd5a40}{วจตฺถํ (สี, สฺยา)} อติวิสํ กฏุกโรหิณิํ อุสีรํ ภทฺทมุตฺตกํ, ยานิ วา ปนญฺญานิปิ อตฺถิ มูลานิ เภสชฺชานิ เนว ขาทนีเย ขาทนียตฺถํ ผรนฺติ, น โภชนีเย โภชนียตฺถํ ผรนฺติ, ตานิ ปฏิคฺคเหตฺวา ยาวชีวํ ปริหริตุํ, สติ ปจฺจเย ปริภุญฺชิตุํ, อสติ ปจฺจเย ปริภุญฺช\-นฺตสฺส อาปตฺติ ทุกฺกฏสฺสาติ.}

\prose{เตน โข ปน สมเยน คิลานานํ ภิกฺขูนํ มูเลหิ เภสชฺเชหิ ปิฏฺเฐหิ อตฺโถ โหติ. ภควโต เอตมตฺถํ อาโรเจสุํ. อนุชานามิ ภิกฺขเว นิสทํ นิสท\-โปตกนฺติ.}

\prose{เตน โข ปน สมเยน คิลานานํ ภิกฺขูนํ กสาเวหิ เภสชฺเชหิ อตฺโถ โหติ. ภควโต เอตมตฺถํ อาโรเจสุํ. อนุชานามิ ภิกฺขเว กสาวานิ เภสชฺชานิ นิมฺพกสาวํ กุฏชกสาวํ ปโฏลกสาวํ ผคฺคว\-กสาวํ นตฺตมาล\-กสาวํ, ยานิ วา ปนญฺญานิปิ อตฺถิ กสาวานิ\csromansharedfootnote{b1caac7f5f1cccb3}{กสาวเภสชฺชานิ (ก)} เภสชฺชานิ เนว ขาทนีเย ขาทนียตฺถํ ผรนฺติ, น โภชนีเย โภชนียตฺถํ ผรนฺติ, ตานิ ปฏิคฺคเหตฺวา ยาวชีวํ ปริหริตุํ, สติ ปจฺจเย ปริภุญฺชิตุํ, อสติ ปจฺจเย ปริภุญฺช\-นฺตสฺส อาปตฺติ ทุกฺกฏสฺสาติ.}

\prose{เตน โข ปน สมเยน คิลานานํ ภิกฺขูนํ ปณฺเณหิ เภสชฺเชหิ อตฺโถ โหติ. ภควโต เอตมตฺถํ อาโรเจสุํ. อนุชานามิ \csromanfolio{293}ภิกฺขเว ปณฺณานิ เภสชฺชานิ นิมฺพปณฺณํ กุฏชปณฺณํ ปโฏลปณฺณํ สุลสิปณฺณํ กปฺปาสปณฺณํ, ยานิ วา ปนญฺญานิปิ อตฺถิ ปณฺณานิ เภสชฺชานิ เนว ขาทนีเย ขาทนียตฺถํ ผรนฺติ, น โภชนีเย โภชนียตฺถํ ผรนฺติ ฯเปฯ.}

\prose{เตน โข ปน สมเยน คิลานานํ ภิกฺขูนํ ผเลหิ เภสชฺเชหิ อตฺโถ โหติ. ภควโต เอตมตฺถํ อาโรเจสุํ. อนุชานามิ ภิกฺขเว ผลานิ เภสชฺชานิ พิลงฺคํ ปิปฺปลิํ มริจํ หรีตกํ วิภีตกํ อามลกํ โคฏฺฐผลํ\csromansharedfootnote{71484194dc47058f}{โคฐผลํ (สฺยา), โกฏฺฐผลํ (ก)}, ยานิ วา ปนญฺญานิปิ อตฺถิ ผลานิ เภสชฺชานิ เนว ขาทนีเย ขาทนียตฺถํ ผรนฺติ, น โภชนีเย โภชนียตฺถํ ผรนฺติ ฯเปฯ.}

\prose{เตน โข ปน สมเยน คิลานานํ ภิกฺขูนํ ชตูหิ เภสชฺเชหิ อตฺโถ โหติ. ภควโต เอตมตฺถํ อาโรเจสุํ. อนุชานามิ ภิกฺขเว ชตูนิ เภสชฺชานิ หิงฺคุํ หิงฺคุชตุํ หิงฺคุสิปาฏิกํ ตกํ ตกปตฺติํ ตกปณฺณิํ สชฺชุลสํ, ยานิ วา ปนญฺญานิปิ อตฺถิ ชตูนิ เภสชฺชานิ เนว ขาทนีเย ขาทนียตฺถํ ผรนฺติ ฯเปฯ.}

\prose{เตน โข ปน สมเยน คิลานานํ ภิกฺขูนํ โลเณหิ เภสชฺเชหิ อตฺโถ โหติ. ภควโต เอตมตฺถํ อาโรเจสุํ. อนุชานามิ ภิกฺขเว โลณานิ เภสชฺชานิ สามุทฺทํ กาฬโลณํ สินฺธวํ อุพฺภิทํ\csromansharedfootnote{6c3d9354280d3d3c}{อุพฺภิรํ (ก)} พิลํ\csromansharedfootnote{3c82aaa4db794a14}{พิฬาลํ (สี)}, ยานิ วา ปนญฺญานิปิ อตฺถิ โลณานิ เภสชฺชานิ เนว ขาทนีเย ขาทนียตฺถํ ผรนฺติ, น โภชนีเย โภชนียตฺถํ ผรนฺติ, ตานิ ปฏิคฺคเหตฺวา ยาวชีวํ ปริหริตุํ, สติ ปจฺจเย ปริภุญฺชิตุํ, อสติ ปจฺจเย ปริภุญฺช\-นฺตสฺส อาปตฺติ ทุกฺกฏสฺสาติ.}

\proseitem{264}{เตน โข ปน สมเยน อายสฺมโต อานนฺทสฺส อุปชฺฌายสฺส อายสฺมโต เพลฏฺฐ\-สีสสฺส ถุลฺลกจฺฉา\-พาโธ โหติ, ตสฺส ลสิกาย จีวรานิ กาเย ลคฺคนฺติ, ตานิ ภิกฺขู อุทเกน เตเมตฺวา เตเมตฺวา อปกฑฺฒนฺติ. อทฺทสา โข ภควา เสนาสนจาริกํ อาหิณฺฑนฺโต เต ภิกฺขู ตานิ จีวรานิ อุทเกน เตเมตฺวา เตเมตฺวา อปกฑฺฒนฺเต, ทิสฺวาน เยน เต ภิกฺขู เตนุปสงฺกมิ, อุปสงฺกมิตฺวา เต ภิกฺขู เอตทโวจ “กิํ อิมสฺส ภิกฺขเว ภิกฺขุโน อาพาโธ”ติ. อิมสฺส ภนฺเต อายสฺมโต ถุลฺลกจฺฉา\-พาโธ, ลสิกาย จีวรานิ กาเย ลคฺคนฺติ, ตานิ มยํ อุทเกน เตเมตฺวา เตเมตฺวา \csromanfolio{294}อปกฑฺฒามาติ. อถ โข ภควา เอตสฺมิํ นิทาเน เอตสฺมิํ ปกรเณ ธมฺมิํ กถํ กตฺวา ภิกฺขู อามนฺเตสิ “อนุชานามิ ภิกฺขเว ยสฺส กณฺฑุ วา ปิฬกา วา อสฺสาโว วา ถุลฺลกจฺฉุ วา อาพาโธ กาโย วา ทุคฺคนฺโธ จุณฺณานิ เภสชฺชานิ, อคิลานสฺส ฉกณํ มตฺติกํ รชน\-นิปฺปกฺกํ, อนุชานามิ ภิกฺขเว อุทุกฺขลํ มุสลนฺ”ติ.}

\prose{เตน โข ปน สมเยน คิลานานํ ภิกฺขูนํ จุณฺเณหิ เภสชฺเชหิ จาลิเตหิ อตฺโถ โหติ. ภควโต เอตมตฺถํ อาโรเจสุํ. อนุชานามิ ภิกฺขเว จุณฺณจาลินินฺติ. สเณฺหหิ อตฺโถ โหติ. ภควโต เอตมตฺถํ อาโรเจสุํ. อนุชานามิ ภิกฺขเว ทุสฺสจาลินินฺติ.}

\prose{เตน โข ปน สมเยน อญฺญตรสฺส ภิกฺขุโน อมนุสฺสิกาพาโธ โหติ, ตํ อาจริยุ\-ปชฺฌายา อุปฏฺฐหนฺตา นาสกฺขิํสุ อโรคํ กาตุํ. โส สูกรสูนํ คนฺตฺวา อามกมํสํ ขาทิ, อามกโลหิตํ ปิวิ, ตสฺส โส อมนุสฺสิกาพาโธ ปฏิปฺปสฺสมฺภิ. ภควโต เอตมตฺถํ อาโรเจสุํ. อนุชานามิ ภิกฺขเว อมนุสฺสิกาพาเธ อามกมํสํ อามกโลหิตนฺติ.}

\proseitem{265}{เตน โข ปน สมเยน อญฺญตรสฺส ภิกฺขุโน จกฺขุ\-โรคา\-พาโธ โหติ, ตํ ภิกฺขู ปริคฺคเหตฺวา อุจฺจารมฺปิ ปสฺสาวมฺปิ นิกฺขาเมนฺติ.}

\prose{อทฺทสา โข ภควา เสนาสนจาริกํ อาหิณฺฑนฺโต เต ภิกฺขู ตํ ภิกฺขุํ ปริคฺคเหตฺวา อุจฺจารมฺปิ ปสฺสาวมฺปิ นิกฺขาเมนฺเต, ทิสฺวาน เยน เต ภิกฺขู เตนุปสงฺกมิ, อุปสงฺกมิตฺวา เต ภิกฺขู เอตทโวจ “กิํ อิมสฺส ภิกฺขเว ภิกฺขุโน อาพาโธ”ติ.}

\prose{อิมสฺส ภนฺเต อายสฺมโต จกฺขุ\-โรคา\-พาโธ อิมํ มยํ ปริคฺคเหตฺวา อุจฺจารมฺปิ ปสฺสาวมฺปิ นิกฺขาเมมาติ. อถ โข ภควา เอตสฺมิํ นิทาเน เอตสฺมิํ ปกรเณ ธมฺมิํ กถํ กตฺวา ภิกฺขู อามนฺเตสิ “อนุชานามิ ภิกฺขเว อญฺชนํ กาฬญฺชนํ รสญฺชนํ โสตญฺชนํ เครุกํ กปลฺลนฺ”ติ. อญฺชนู\-ปปิสเนหิ อตฺโถ โหติ. ภควโต เอตมตฺถํ อาโรเจสุํ. อนุชานามิ ภิกฺขเว จนฺทนํ ตครํ กาฬานุสาริยํ ตาลีสํ ภทฺท\-มุตฺตกนฺติ. เตน โข ปน สมเยน ภิกฺขู ปิฏฺฐานิ อญฺชนานิ จรุเกสุปิ\csromansharedfootnote{338dd9d64fd7ee1e}{[มิสฺสินฺคฺ โนเต 0]} สราวเกสุปิ นิกฺขิปนฺติ, ติณจุณฺเณหิปิ ปํสุเกหิปิ โอกิริยนฺติ. ภควโต เอตมตฺถํ อาโรเจสุํ. อนุชานามิ ภิกฺขเว อญฺชนินฺติ.}

\prose{\csromanfolio{295}เตน โข ปน สมเยน ฉพฺพคฺคิยา ภิกฺขู อุจฺจาวจา อญฺชนิโย ธาเรนฺติ โสวณฺณมยํ รูปิยมยํ. มนุสฺสา อุชฺฌายนฺติ ขิยฺยนฺติ วิปาเจนฺติ “เสยฺยถาปิ คิหี กามโภคิโน”ติ. ภควโต เอตมตฺถํ อาโรเจสุํ. น ภิกฺขเว อุจฺจาวจา อญฺชนี ธาเรตพฺพา, โย ธาเรยฺย, อาปตฺติ ทุกฺกฏสฺส. อนุชานามิ ภิกฺขเว อฏฺฐิมยํ ทนฺตมยํ วิสาณมยํ นฬมยํ เวฬุมยํ กฏฺฐมยํ ชตุมยํ ผลมยํ โลหมยํ สงฺข\-นาภิ\-มยนฺติ.}

\prose{เตน โข ปน สมเยน อญฺชนิโย อปารุตา โหนฺติ, ติณจุณฺเณหิปิ ปํสุเกหิปิ โอกิริยนฺติ.}

\prose{ภควโต เอตมตฺถํ อาโรเจสุํ. อนุชานามิ ภิกฺขเว อปิธานนฺติ. อปิธานํ นิปตติ. ภควโต เอตมตฺถํ อาโรเจสุํ. อนุชานามิ ภิกฺขเว สุตฺตเกน พนฺธิตฺวา อญฺชนิยา พนฺธิตุนฺติ. อญฺชนี ผลติ\csromansharedfootnote{338dd9d64fd7ee1e}{[มิสฺสินฺคฺ โนเต 0]}. ภควโต เอตมตฺถํ อาโรเจสุํ. อนุชานามิ ภิกฺขเว สุตฺตเกน สิพฺเพตุนฺติ.}

\prose{เตน โข ปน สมเยน ภิกฺขู องฺคุลิยา อญฺชนฺติ, อกฺขีนิ ทุกฺขานิ โหนฺติ.}

\setlength{\csromangathapagewidth}{0pt}
\setlength{\csromangathawakwidth}{0pt}
\csromangathameasuregroup{}{ภควโต เอตมตฺถํ อาโรเจสุํ. อนุชานามิ ภิกฺขเว อญฺชนิสลากนฺติ.}
\csromangathameasurewak{ภควโต เอตมตฺถํ อาโรเจสุํ. อนุชานามิ ภิกฺขเว อญฺชนิสลากนฺติ.}
\setlength{\csromangathaleftcol}{0pt}
\csromangathagroup{\csromangathastanza{\csromangathawak{ภควโต เอตมตฺถํ อาโรเจสุํ. อนุชานามิ ภิกฺขเว อญฺชนิ\-สลากนฺติ.}}}

\prose{เตน โข ปน สมเยน ฉพฺพคฺคิยา ภิกฺขู อุจฺจาวจา อญฺชนิสลากาโย ธาเรนฺติ โสวณฺณมยํ รูปิยมยํ. มนุสฺสา อุชฺฌายนฺติ ขิยฺยนฺติ วิปาเจนฺติ “เสยฺยถาปิ คิหี กามโภคิโน”ติ. ภควโต เอตมตฺถํ อาโรเจสุํ. น ภิกฺขเว อุจฺจาวจา อญฺชนิสลากา ธาเรตพฺพา, โย ธาเรยฺย, อาปตฺติ ทุกฺกฏสฺส. อนุชานามิ ภิกฺขเว อฏฺฐิมยํ ฯเปฯ สงฺข\-นาภิ\-มยนฺติ.}

\proseitem{265}{เตน โข ปน สมเยน อญฺชนิสลากา ภูมิยํ ปติตา ผรุสา โหติ. ภควโต เอตมตฺถํ อาโรเจสุํ. อนุชานามิ ภิกฺขเว สลาก\-ฐานิยนฺติ\csromansharedfootnote{6900ff20e39c0a8d}{สลาโกธานิยนฺติ (สี, สฺยา)}.}

\prose{เตน โข ปน สมเยน ภิกฺขู อญฺชนิมฺปิ อญฺชนิ\-สลากมฺปิ หตฺเถน ปริหรนฺติ. ภควโต เอตมตฺถํ อาโรเจสุํ. อนุชานามิ ภิกฺขเว อญฺชนิ\-ตฺถวิกนฺติ. อํสพทฺธโก น โหติ. ภควโต เอตมตฺถํ อาโรเจสุํ. อนุชานามิ ภิกฺขเว อํสพทฺธกํ พนฺธน\-สุตฺตกนฺติ.}

\proseitem{266}{\csromanfolio{296}เตน โข ปน สมเยน อายสฺมโต ปิลินฺท\-วจฺฉสฺส สีสาภิตาโป โหติ. ภควโต เอตมตฺถํ อาโรเจสุํ. อนุชานามิ ภิกฺขเว มุทฺธนิ เตลกนฺติ. นกฺขมนิโย โหติ. ภควโต เอตมตฺถํ อาโรเจสุํ. อนุชานามิ ภิกฺขเว นตฺถุกมฺมนฺติ. นตฺถุ คลติ. ภควโต เอตมตฺถํ อาโรเจสุํ. อนุชานามิ ภิกฺขเว นตฺถุกรณินฺติ.}

\prose{เตน โข ปน สมเยน ฉพฺพคฺคิยา ภิกฺขู อุจฺจาวจา นตฺถุกรณิโย ธาเรนฺติ โสวณฺณมยํ รูปิยมยํ. มนุสฺสา อุชฺฌายนฺติ ขิยฺยนฺติ วิปาเจนฺติ “เสยฺยถาปิ คิหี กามโภคิโน”ติ. ภควโต เอตมตฺถํ อาโรเจสุํ. น ภิกฺขเว อุจฺจาวจา นตฺถุกรณี ธาเรตพฺพา, โย ธาเรยฺย, อาปตฺติ ทุกฺกฏสฺส. อนุชานามิ ภิกฺขเว อฏฺฐิมยํ ฯเปฯ สงฺข\-นาภิ\-มยนฺติ. นตฺถุํ วิสมํ อาสิญฺจนฺติ\csromansharedfootnote{4679f0c05e24104f}{นตฺถุ วิสมํ อาสิญฺชิยติ (สี, สฺยา)}. ภควโต เอตมตฺถํ อาโรเจสุํ. อนุชานามิ ภิกฺขเว ยมก\-นตฺถุ\-กรณินฺติ. นกฺขมนิโย โหติ. ภควโต เอตมตฺถํ อาโรเจสุํ. อนุชานามิ ภิกฺขเว ธูมํ ปาตุนฺติ. ตญฺเญว วฏฺฏิํ อาลิมฺเปตฺวา ปิวนฺติ, กณฺโฐ ทหติ. ภควโต เอตมตฺถํ อาโรเจสุํ. อนุชานามิ ภิกฺขเว ธูมเนตฺตนฺติ.}

\prose{เตน โข สมเยน ฉพฺพคฺคิยา ภิกฺขู อุจฺจาวจานิ ธูมเนตฺตานิ ธาเรนฺติ โสวณฺณมยํ รูปิยมยํ. มนุสฺสา อุชฺฌายนฺติ ขิยฺยนฺติ วิปาเจนฺติ “เสยฺยถาปิ คิหี กามโภคิโน”ติ. ภควโต เอตมตฺถํ อาโรเจสุํ. น ภิกฺขเว อุจฺจาวจานิ ธูมเนตฺตานิ ธาเรตพฺพานิ, โย ธาเรยฺย, อาปตฺติ ทุกฺกฏสฺส. อนุชานามิ ภิกฺขเว อฏฺฐิมยํ ฯเปฯ สงฺข\-นาภิ\-มยนฺติ.}

\setlength{\csromangathapagewidth}{0pt}
\setlength{\csromangathawakwidth}{0pt}
\csromangathameasuregroup{}{เตน โข ปน สมเยน ธูมเนตฺตานิ อปารุตานิ โหนฺติ, \\ ปาณกา ปวิสนฺติ. \\ ภควโต เอตมตฺถํ อาโรเจสุํ. อนุชานามิ ภิกฺขเว อปิธานนฺติ.}
\setlength{\csromangathaleftcol}{0pt}
\csromangathagroup{\csromangathastanza{เตน โข ปน สมเยน ธูมเนตฺตานิ อปารุตานิ โหนฺติ, \\ ปาณกา ปวิสนฺติ. \\ ภควโต เอตมตฺถํ อาโรเจสุํ. อนุชานามิ ภิกฺขเว อปิธานนฺติ.}}

\proseitem{266}{เตน โข ปน สมเยน ภิกฺขู ธูมเนตฺตานิ หตฺเถน ปริหรนฺติ. ภควโต เอตมตฺถํ อาโรเจสุํ. อนุชานามิ ภิกฺขเว ธูมเนตฺตถวิกนฺติ. เอกโต สํสิยนฺติ. ภควโต เอตมตฺถํ อาโรเจสุํ. อนุชานามิ ภิกฺขเว ยมก\-ถวิกนฺติ. อํสพทฺธโก น โหติ. ภควโต เอตมตฺถํ อาโรเจสุํ. อนุชานามิ ภิกฺขเว อํสพทฺธกํ พนฺธน\-สุตฺตกนฺติ.}

\proseitem{267}{\csromanfolio{297}เตน โข ปน สมเยน อายสฺมโต ปิลินฺท\-วจฺฉสฺส วาตาพาโธ โหติ. เวชฺชา เอวมาหํสุ “เตลํ ปจิตพฺพนฺ”ติ. ภควโต เอตมตฺถํ อาโรเจสุํ. อนุชานามิ ภิกฺขเว เตลปากนฺติ. ตสฺมิํ โข ปน เตลปาเก มชฺชํ ปกฺขิปิตพฺพํ โหติ. ภควโต เอตมตฺถํ อาโรเจสุํ. อนุชานามิ ภิกฺขเว เตลปาเก มชฺชํ ปกฺขิปิตุนฺติ.}

\prose{เตน โข ปน สมเยน ฉพฺพคฺคิยา ภิกฺขู อติปกฺขิตฺต\-มชฺชานิ\csromansharedfootnote{68b80144f7c44e86}{อติขิตฺตมชฺชานิ (ก)} เตลานิ ปจนฺติ, ตานิ ปิวิตฺวา มชฺชนฺติ.}

\prose{ภควโต เอตมตฺถํ อาโรเจสุํ. น ภิกฺขเว อติปกฺขิตฺต\-มชฺชํ เตลํ ปาตพฺพํ, โย ปิเวยฺย, ยถาธมฺโม กาเรตพฺโพ. อนุชานามิ ภิกฺขเว ยสฺมิํ เตลปาเก มชฺชสฺส น วณฺโณ น คนฺโธ น รโส ปญฺญายติ, เอวรูปํ มชฺช\-ปกฺขิตฺตํ เตลํ ปาตุนฺติ.}

\prose{เตน โข ปน สมเยน ภิกฺขูนํ พหุํ อติปกฺขิตฺต\-มชฺชํ เตลํ ปกฺกํ โหติ, อถ โข ภิกฺขูนํ เอตทโหสิ “กถํ นุ โข อติปกฺขิตฺต\-มชฺเช เตเล ปฏิปชฺชิตพฺพนฺ”ติ.}

\prose{ภควโต เอตมตฺถํ อาโรเจสุํ. อนุชานามิ ภิกฺขเว อพฺภญฺชนํ อธิฏฺฐาตุนฺติ.}

\prose{เตน โข ปน สมเยน อายสฺมโต ปิลินฺท\-วจฺฉสฺส พหุตรํ เตลํ ปกฺกํ โหติ, เตลภาชนํ น วิชฺชติ.}

\prose{ภควโต เอตมตฺถํ อาโรเจสุํ. อนุชานามิ ภิกฺขเว ตีณิ ตุมฺพานิ โลหตุมฺพํ กฏฺฐตุมฺพํ ผลตุมฺพนฺติ.}

\proseitem{267}{เตน โข ปน สมเยน อายสฺมโต ปิลินฺท\-วจฺฉสฺส องฺควาโต โหติ. ภควโต เอตมตฺถํ อาโรเจสุํ. อนุชานามิ ภิกฺขเว เสทกมฺมนฺติ. นกฺขมนิโย โหติ. ภควโต เอตมตฺถํ อาโรเจสุํ. อนุชานามิ ภิกฺขเว สมฺภาร\-เสทนฺติ. นกฺขมนิโย โหติ. ภควโต เอตมตฺถํ อาโรเจสุํ. อนุชานามิ ภิกฺขเว มหาเสทนฺติ. นกฺขมนิโย โหติ. ภควโต เอตมตฺถํ อาโรเจสุํ. อนุชานามิ ภิกฺขเว ภงฺโคทกนฺติ. นกฺขมนิโย โหติ. ภควโต เอตมตฺถํ อาโรเจสุํ. อนุชานามิ ภิกฺขเว อุทกโกฏฺฐกนฺติ.}

\prose{เตน โข ปน สมเยน อายสฺมโต ปิลินฺท\-วจฺฉสฺส ปพฺพวาโต โหติ. ภควโต เอตมตฺถํ อาโรเจสุํ. อนุชานามิ ภิกฺขเว \csromanfolio{298}โลหิตํ โมเจตุนฺติ. นกฺขมนิโย โหติ. ภควโต เอตมตฺถํ อาโรเจสุํ. อนุชานามิ ภิกฺขเว โลหิตํ โมเจตฺวา วิสาเณน คาเหตุนฺติ\csromansharedfootnote{3c8b7637823fce87}{คเหตุนฺติ (สี, สฺยา)}.}

\prose{เตน โข ปน สมเยน อายสฺมโต ปิลินฺท\-วจฺฉสฺส ปาทา ผลิตา\csromansharedfootnote{2f4703b8b0c23122}{ผาลิตา (ก)} โหนฺติ. ภควโต เอตมตฺถํ อาโรเจสุํ. อนุชานามิ ภิกฺขเว ปาท\-พฺภญฺชนนฺติ. นกฺขมนิโย โหติ. ภควโต เอตมตฺถํ อาโรเจสุํ. อนุชานามิ ภิกฺขเว ปชฺชํ อภิสงฺขริตุนฺติ.}

\prose{เตน โข ปน สมเยน อญฺญตรสฺส ภิกฺขุโน คณฺฑาพาโธ โหติ. ภควโต เอตมตฺถํ อาโรเจสุํ. อนุชานามิ ภิกฺขเว สตฺถกมฺมนฺติ. กสาโวทเกน อตฺโถ โหติ. ภควโต เอตมตฺถํ อาโรเจสุํ. อนุชานามิ ภิกฺขเว กสาโวทกนฺติ. ติลกกฺเกน อตฺโถ โหติ. ภควโต เอตมตฺถํ อาโรเจสุํ. อนุชานามิ ภิกฺขเว ติลกกฺกนฺติ. กพฬิกาย อตฺโถ โหติ. ภควโต เอตมตฺถํ อาโรเจสุํ. อนุชานามิ ภิกฺขเว กพฬิกนฺติ. วณพนฺธน\-โจเฬน อตฺโถ โหติ. ภควโต เอตมตฺถํ อาโรเจสุํ. อนุชานามิ ภิกฺขเว วณพนฺธน\-โจฬนฺติ. วโณ กณฺฑุวติ. ภควโต เอตมตฺถํ อาโรเจสุํ. อนุชานามิ ภิกฺขเว สาสปกุฏฺเฏน\csromansharedfootnote{00213d2fb2d1aacb}{สาสปกุฑฺเฑน (สี, สฺยา)} โผสิตุนฺติ. วโณ กิลิชฺชิตฺถ. ภควโต เอตมตฺถํ อาโรเจสุํ. อนุชานามิ ภิกฺขเว ธูมํ กาตุนฺติ. วฑฺฒมํสํ วุฏฺฐาติ. ภควโต เอตมตฺถํ อาโรเจสุํ. อนุชานามิ ภิกฺขเว โลณ\-สกฺขริกาย ฉินฺทิตุนฺติ. วโณ น รุหติ. ภควโต เอตมตฺถํ อาโรเจสุํ. อนุชานามิ ภิกฺขเว วณเตลนฺติ. เตลํ คลติ. ภควโต เอตมตฺถํ อาโรเจสุํ. อนุชานามิ ภิกฺขเว วิกาสิกํ สพฺพํ วณปฏิกมฺมนฺติ.}

\proseitem{268}{เตน โข ปน สมเยน อญฺญตโร ภิกฺขุ อหินา ทฏฺโฐ โหติ. ภควโต เอตมตฺถํ อาโรเจสุํ. อนุชานามิ ภิกฺขเว จตฺตาริ มหาวิกฏานิ ทาตุํ คูถํ มุตฺตํ ฉาริกํ มตฺติกนฺติ. อถ โข ภิกฺขูนํ เอตทโหสิ “อปฺปฏิคฺคหิตานิ นุ โข อุทาหุ ปฏิคฺคเหตพฺพานี”ติ. ภควโต เอตมตฺถํ อาโรเจสุํ. อนุชานามิ ภิกฺขเว สติ กปฺปิยการเก ปฏิคฺคหาเปตุํ, อสติ กปฺปิยการเก สามํ คเหตฺวา ปริภุญฺชิ\-ตุนฺติ.}

\prose{\csromanfolio{299}เตน โข ปน สมเยน อญฺญตเรน ภิกฺขุนา วิสํ ปีตํ โหติ. ภควโต เอตมตฺถํ อาโรเจสุํ. อนุชานามิ ภิกฺขเว คูถํ ปาเยตุนฺติ. อถ โข ภิกฺขูนํ เอตทโหสิ “อปฺปฏิคฺคหิตํ นุ โข อุทาหุ ปฏิคฺคเหตพฺโพ”ติ. ภควโต เอตมตฺถํ อาโรเจสุํ. อนุชานามิ ภิกฺขเว ยํ กโรนฺโต ปฏิคฺคณฺหาติ, เสฺวว ปฏิคฺคโห กโต, น ปุน\csromansharedfootnote{6c0e579925becd48}{กโต ปน (?)} ปฏิคฺคเหตพฺโพติ.}

\proseitem{269}{เตน โข ปน สมเยน อญฺญตรสฺส ภิกฺขุโน ฆรทินฺนกาพาโธ โหติ. ภควโต เอตมตฺถํ อาโรเจสุํ. อนุชานามิ ภิกฺขเว สีตาโลฬิ ปาเยตุนฺติ.}

\prose{เตน โข ปน สมเยน อญฺญตโร ภิกฺขุ ทุฏฺฐคหณิโก โหติ, ภควโต เอตมตฺถํ อาโรเจสุํ. อนุชานามิ ภิกฺขเว อามิสขารํ ปาเยตุนฺติ.}

\prose{เตน โข ปน สมเยน อญฺญตรสฺส ภิกฺขุโน ปณฺฑุโรคาพาโธ โหติ. ภควโต เอตมตฺถํ อโรเจสุํ. อนุชานามิ ภิกฺขเว มุตฺตหรีตกํ ปาเยตุนฺติ.}

\prose{เตน โข ปน สมเยน อญฺญตรสฺส ภิกฺขุโน ฉวิโทสาพาโธ โหติ. ภควโต เอตมตฺถํ อาโรเจสุํ. อนุชานามิ ภิกฺขเว คนฺธาเลปํ กาตุนฺติ.}

\prose{เตน โข ปน สมเยน อญฺญตโร ภิกฺขุ อภิสนฺนกาโย โหติ. ภควโต เอตมตฺถํ อาโรเจสุํ. อนุชานามิ ภิกฺขเว วิเรจนํ ปาตุนฺติ. อจฺฉกญฺชิยา อตฺโถ โหติ. ภควโต เอตมตฺถํ อาโรเจสุํ. อนุชานามิ ภิกฺขเว อจฺฉกญฺชินฺติ. อกฏยูเสน อตฺโถ โหติ. ภควโต เอตมตฺถํ อาโรเจสุํ. อนุชานามิ ภิกฺขเว อกฏยูสนฺติ. กฏากเฏน อตฺโถ โหติ. ภควโต เอตมตฺถํ อาโรเจสุํ. อนุชานามิ ภิกฺขเว กฏากฏนฺติ. ปฏิจฺฉาทนีเยน อตฺโถ โหติ. ภควโต เอตมตฺถํ อาโรเจสุํ. อนุชานามิ ภิกฺขเว ปฏิจฺฉาทนียนฺติ.}

\csromanmatikaanchor{mtk.168}
\csromantocmarkref{section}{162. ปิลินฺทวจฺฉวตฺถุ}{mtk.168}
\bookmark[dest={mtk.168},level=1]{162. ปิลินฺทวจฺฉวตฺถุ}
\csromanheader{1}{162. \textbf{ปิลินฺทวจฺฉ}\-\textbf{วตฺถุ}}
\proseitem{270}{\csromansymbolfootnote{*}{อิทํ วตฺถุ วิ 1. 360 ปิฏฺฐาทีสุปิ อาคตํ.}เตน โข ปน สมเยน อายสฺมา ปิลินฺทวจฺโฉ ราชคเห ปพฺภารํ โสธาเปติ เลณํ กตฺตุกาโม. อถ โข ราชา มาคโธ \csromanfolio{300}เสนิโย พิมฺพิสาโร เยนายสฺมา ปิลินฺทวจฺโฉ เตนุปสงฺกมิ, อุปสงฺกมิตฺวา อายสฺมนฺตํ ปิลินฺทวจฺฉํ อภิวาเทตฺวา เอกมนฺตํ นิสีทิ, เอกมนฺตํ นิสินฺโน โข ราชา มาคโธ เสนิโย พิมฺพิสาโร อายสฺมนฺตํ ปิลินฺทวจฺฉํ เอตทโวจ “กิํ ภนฺเต เถโร การาเปตี”ติ. ปพฺภารํ มหาราช โสธาเปมิ เลณํ กตฺตุกาโมติ. อตฺโถ ภนฺเต อยฺยสฺส อารามิเกนาติ. น โข มหาราช ภควตา อารามิโก อนุญฺญาโตติ. เตน หิ ภนฺเต ภควนฺตํ ปฏิปุจฺฉิตฺวา มม อาโรเจยฺยาถาติ. เอวํ มหาราชาติ โข อายสฺมา ปิลินฺทวจฺโฉ รญฺโญ มาคธสฺส เสนิยสฺส พิมฺพิสารสฺส ปจฺจสฺโสสิ. อถ โข อายสฺมา ปิลินฺทวจฺโฉ ราชานํ มาคธํ เสนิยํ พิมฺพิสารํ ธมฺมิยา กถาย สนฺทสฺเสสิ สมาทเปสิ สมุตฺเตเชสิ สมฺปหํเสสิ. อถ โข ราชา มาคโธ เสนิโย พิมฺพิสาโร อายสฺมตา ปิลินฺท\-วจฺเฉน ธมฺมิยา กถาย สนฺทสฺสิโต สมาทปิโต สมุตฺเตชิโต สมฺปหํสิโต อุฏฺฐายาสนา อายสฺมนฺตํ ปิลินฺทวจฺฉํ อภิวาเทตฺวา ปทกฺขิณํ กตฺวา ปกฺกามิ.}

\prose{อถ โข อายสฺมา ปิลินฺทวจฺโฉ ภควโต สนฺติเก ทูตํ ปาเหสิ “ราชา ภนฺเต มาคโธ เสนิโย พิมฺพิสาโร อารามิกํ ทาตุกาโม, กถํ นุ โข ภนฺเต มยา ปฏิปชฺชิตพฺพนฺ”ติ. อถ โข ภควา เอตสฺมิํ นิทาเน เอตสฺมิํ ปกรเณ ธมฺมิํ กถํ กตฺวา ภิกฺขู อามนฺเตสิ “อนุชานามิ ภิกฺขเว อารามิกนฺ”ติ. ทุติยมฺปิ โข ราชา มาคโธ เสนิโย พิมฺพิสาโร เยนายสฺมา ปิลินฺทวจฺโฉ เตนุปสงฺกมิ, อุปสงฺกมิตฺวา อายสฺมนฺตํ ปิลินฺทวจฺฉํ อภิวาเทตฺวา เอกมนฺตํ นิสีทิ, เอกมนฺตํ นิสินฺโน โข ราชา มาคโธ เสนิโย พิมฺพิสาโร อายสฺมนฺตํ ปิลินฺทวจฺฉํ เอตทโวจ “อนุญฺญาโต ภนฺเต ภควตา อารามิโก”ติ. เอวํ มหาราชาติ. เตน หิ ภนฺเต อยฺยสฺส อารามิกํ ทมฺมีติ. อถ โข ราชา มาคโธ เสนิโย พิมฺพิสาโร อายสฺมโต ปิลินฺท\-วจฺฉสฺส อารามิกํ ปฏิสฺสุตฺวา วิสฺสริตฺวา จิเรน สติํ ปฏิลภิตฺวา อญฺญตรํ สพฺพตฺถกํ มหามตฺตํ อามนฺเตสิ “โย มยา ภเณ อยฺยสฺส อารามิโก ปฏิสฺสุโต, ทินฺโน โส อารามิโกติ. น โข เทว อยฺยสฺส อารามิโก ทินฺโนติ. กีว จิรํ นุ โข ภเณ อิโต\csromansharedfootnote{73c507dfb25a699a}{อิโต รตฺติ (สฺยา)} หิ ตํ โหตี”ติ. อถ โข \csromanfolio{301}โส มหามตฺโต รตฺติโย คเณตฺวา\csromansharedfootnote{20529ea0b81957b3}{วิคเณตฺวา (สี)} ราชานํ มาคธํ เสนิยํ พิมฺพิสารํ เอตทโวจ “ปญฺจ เทว รตฺติสตานี”ติ. เตน หิ ภเณ อยฺยสฺส ปญฺจ อารามิกสตานิ เทหีติ. เอวํ เทวาติ โข โส มหามตฺโต รญฺโญ มาคธสฺส เสนิยสฺส พิมฺพิสารสฺส ปฏิสฺสุตฺวา อายสฺมโต ปิลินฺท\-วจฺฉสฺส ปญฺจ อารามิกสตานิ ปาทาสิ, ปาฏิเยกฺโก คาโม นิวิสิ, อารามิกคามโกติปิ นํ อาหํสุ, ปิลินฺทคามโกติปิ นํ อาหํสุ.}

\proseitem{271}{เตน โข ปน สมเยน อายสฺมา ปิลินฺทวจฺโฉ ตสฺมิํ คามเก กุลูปโก โหติ. อถ โข อายสฺมา ปิลินฺทวจฺโฉ ปุพฺพณฺหสมยํ นิวาเสตฺวา ปตฺต\-จีวรมา\-ทาย ปิลินฺทคามํ ปิณฺฑาย ปาวิสิ. เตน โข ปน สมเยน ตสฺมิํ คามเก อุสฺสโว โหติ. ทารกา อลงฺกตา มาลากิตา กีฬนฺติ. อถ โข อายสฺมา ปิลินฺทวจฺโฉ ปิลินฺทคามเก สปทานํ ปิณฺฑาย จรมาโน เยน อญฺญตรสฺส อารามิกสฺส นิเวสนํ เตนุปสงฺกมิ, อุปสงฺกมิตฺวา ปญฺญตฺเต อาสเน นิสีทิ. เตน โข ปน สมเยน ตสฺสา อารามิกินิยา ธีตา อญฺเญ ทารเก อลงฺกเต มาลากิเต ปสฺสิตฺวา โรทติ “มาลํ เม เทถ, อลงฺการํ เม เทถา”ติ. อถ โข อายสฺมา ปิลินฺทวจฺโฉ ตํ อารามิกินิํ เอตทโวจ “กิสฺสายํ ทาริกา โรทตี”ติ. อยํ ภนฺเต ทาริกา อญฺเญ ทารเก อลงฺกเต มาลากิเต ปสฺสิตฺวา โรทติ “มาลํ เม เทถ, อลงฺการํ เม เทถา”ติ. กุโต อมฺหากํ ทุคฺคตานํ มาลา, กุโต อลงฺกาโรติ. อถ โข อายสฺมา ปิลินฺทวจฺโฉ อญฺญตรํ ติณณฺฑุปกํ คเหตฺวา ตํ อารามิกินิํ เอตทโวจ “หนฺทิมํ ติณณฺฑุปกํ ตสฺสา ทาริกาย สีเส ปฏิมุญฺจา”ติ. อถ โข สา อารามิกินี ตํ ติณณฺฑุปกํ คเหตฺวา ตสฺสา ทาริกาย สีเส ปฏิมุญฺชิ. สา อโหสิ สุวณฺณมาลา อภิรูปา ทสฺสนียา ปาสาทิกา. นตฺถิ ตาทิสา รญฺโญปิ อนฺเตปุเร สุวณฺณมาลา. มนุสฺสา รญฺโญ มาคธสฺส เสนิยสฺส พิมฺพิสารสฺส อาโรเจสุํ “อมุกสฺส เทว อารามิกสฺส ฆเร สุวณฺณมาลา อภิรูปา ทสฺสนียา ปาสาทิกา, นตฺถิ ตาทิสา เทวสฺสปิ อนฺเตปุเร สุวณฺณมาลา, กุโต ตสฺส ทุคฺคตสฺส, นิสฺสํสยํ โจริกาย อาภตา”ติ.}

\prose{\csromanfolio{302}อถ โข ราชา มาคโธ เสนิโย พิมฺพิสาโร ตํ อารามิกกุลํ พนฺธาเปสิ. ทุติยมฺปิ โข อายสฺมา ปิลินฺทวจฺโฉ ปุพฺพณฺหสมยํ นิวาเสตฺวา ปตฺต\-จีวรมา\-ทาย ปิลินฺทคามํ ปิณฺฑาย ปาวิสิ. ปิลินฺทคามเก สปทานํ ปิณฺฑาย จรมาโน เยน ตสฺส อารามิกสฺส นิเวสนํ เตนุปสงฺกมิ, อุปสงฺกมิตฺวา ปฏิวิสฺสเก ปุจฺฉิ “กหํ อิมํ อารามิกกุลํ คตนฺ”ติ. เอติสฺสา ภนฺเต สุวณฺณมาลาย การณา รญฺญา พนฺธาปิตนฺติ. อถ โข อายสฺมา ปิลินฺทวจฺโฉ เยน รญฺโญ มาคธสฺส เสนิยสฺส พิมฺพิสารสฺส นิเวสนํ เตนุปสงฺกมิ, อุปสงฺกมิตฺวา ปญฺญตฺเต อาสเน นิสีทิ. อถ โข ราชา มาคโธ เสนิโย พิมฺพิสาโร เยนายสฺมา ปิลินฺทวจฺโฉ เตนุปสงฺกมิ, อุปสงฺกมิตฺวา อายสฺมนฺตํ ปิลินฺทวจฺฉํ อภิวาเทตฺวา เอกมนฺตํ นิสีทิ. เอกมนฺตํ นิสินฺนํ โข ราชานํ มาคธํ เสนิยํ พิมฺพิสารํ อายสฺมา ปิลินฺทวจฺโฉ เอตทโวจ “กิสฺส มหาราช อารามิกกุลํ พนฺธาปิตนฺ”ติ. “ตสฺส ภนฺเต อารามิกสฺส ฆเร สุวณฺณมาลา อภิรูปา ทสฺสนียา ปาสาทิกา, นตฺถิ ตาทิสา อมฺหากมฺปิ อนฺเตปุเร สุวณฺณมาลา, กุโต ตสฺส ทุคฺคตสฺส, นิสฺสํสยํ โจริกาย อาภตา”ติ. อถ โข อายสฺมา ปิลินฺทวจฺโฉ รญฺโญ มาคธสฺส เสนิยสฺส พิมฺพิสารสฺส ปาสาทํ สุวณฺณนฺติ อธิมุจฺจิ. โส อโหสิ สพฺพ\-โสวณฺณมโย. อิทํ ปน เต มหาราช ตาว พหุํ สุวณฺณํ กุโตติ. “อญฺญาตํ ภนฺเต, อยฺยสฺเสเวโส อิทฺธานุภาโว”ติ ตํ อารามิกกุลํ มุญฺจาเปสิ.}

\proseitemwithcloserruled{271}{มนุสฺสา “อยฺเยน กิร ปิลินฺท\-วจฺเฉน สราชิกาย ปริสาย อุตฺตริ\-มนุสฺสธมฺมํ อิทฺธิ\-ปาฏิหาริยํ ทสฺสิตนฺ”ติ อตฺตมนา อภิปฺปสนฺนา อายสฺมโต ปิลินฺท\-วจฺฉสฺส ปญฺจ เภสชฺชานิ อภิหริํสุ. เสยฺยถิทํ, สปฺปิํ นวนีตํ เตลํ มธุํ\csromansharedfootnote{7f866a856eff4235}{สปฺปิ นวนียํ เตลํ มธุ (ก)} ผาณิตํ. ปกติยาปิ จ อายสฺมา ปิลินฺทวจฺโฉ ลาภี โหติ ปญฺจนฺนํ เภสชฺชานํ, ลทฺธํ ลทฺธํ ปริสาย วิสฺสชฺเชติ. ปริสา จสฺส โหติ พาหุลฺลิกา, ลทฺธํ ลทฺธํ โกลมฺเพปิ\csromansharedfootnote{9c6a8dbb8f94044f}{โกฬุมฺเพปิ (ก)} ฆเฏปิ ปูเรตฺวา ปฏิสาเมติ, ปริสฺสาว\-นานิปิ ถวิกาโยปิ ปูเรตฺวา วาตปาเนสุ ลคฺเคติ, ตานิ โอลีนวิลีนานิ ติฏฺฐนฺติ. อุนฺทูเรหิปิ วิหารา โอกิณฺณวิกิณฺณา โหนฺติ. มนุสฺสา วิหารจาริกํ อาหิณฺฑนฺตา ปสฺสิตฺวา อุชฺฌายนฺติ ขิยฺยนฺติ วิปาเจนฺติ “อนฺโตโกฏฺฐาคาริกา อิเม สมณา \csromanfolio{303}สกฺยปุตฺติยา เสยฺยถาปิ ราชา มาคโธ เสนิโย พิมฺพิสาโร”ติ. อสฺโสสุํ โข ภิกฺขู เตสํ มนุสฺสานํ อุชฺฌายนฺตานํ ขิยฺยนฺตานํ วิปาเจนฺตานํ. เย เต ภิกฺขู อปฺปิจฺฉา, เต อุชฺฌายนฺติ ขิยฺยนฺติ วิปาเจนฺติ “กถํ หิ นาม ภิกฺขู เอวรูปาย พาหุลฺลาย เจเตสฺสนฺตี”ติ. อถ โข เต ภิกฺขู เต อเนก\-ปริยาเยน วิครหิตฺวา ภควโต เอตมตฺถํ อาโรเจสุํ ฯเปฯ. สจฺจํ กิร ภิกฺขเว ภิกฺขู เอวรูปาย พาหุลฺลาย เจเตนฺตีติ. สจฺจํ ภควาติ - ป- วิครหิตฺวา ธมฺมิํ กถํ กตฺวา ภิกฺขู อามนฺเตสิ “ยานิ โข ปน ตานิ คิลานานํ ภิกฺขูนํ ปฏิสายนียานิ เภสชฺชานิ, เสยฺยถิทํ, สปฺปิ นวนีตํ เตลํ มธุ ผาณิตํ, ตานิ ปฏิคฺคเหตฺวา สตฺตาหปรมํ สนฺนิธิการกํ ปริภุญฺชิ\-ตพฺพานิ, ตํ อติกฺกามยโต ยถาธมฺโม กาเรตพฺโพ”ติ.}{เภสชฺชา\-นุญฺญาต\-ภาณวาโร นิฏฺฐิโต ปฐโม.}
\csromanmatikaanchor{mtk.169}
\csromantocmarkref{section}{163. คุฬาทิอนุชานนา}{mtk.169}
\bookmark[dest={mtk.169},level=1]{163. คุฬาทิอนุชานนา}
\csromanheader{1}{163. คุฬา\-ทิ\-อนุชานนา}
\proseitem{272}{อถ โข ภควา สาวตฺถิยํ ยถาภิรนฺตํ วิหริตฺวา เยน ราชคหํ เตน จาริกํ ปกฺกามิ. อทฺทสา โข อายสฺมา กงฺขาเรวโต อนฺตรามคฺเค คุฬกรณํ โอกฺกมิตฺวา คุเฬ ปิฏฺฐมฺปิ ฉาริกมฺปิ ปกฺขิปนฺเต, ทิสฺวาน “อกปฺปิโย คุโฬ สามิโส, น กปฺปติ คุโฬ วิกาเล ปริภุญฺชิตุนฺ”ติ กุกฺกุจฺจายนฺโต สปริโส คุฬํ น ปริภุญฺชติ. เยปิสฺส โสตพฺพํ มญฺญนฺติ, เตปิ คุฬํ น ปริภุญฺชนฺติ. ภควโต เอตมตฺถํ อาโรเจสุํ. กิมตฺถาย\csromansharedfootnote{859c15cbbeb2a275}{กิมตฺถิยา (ก)} ภิกฺขเว คุเฬ ปิฏฺฐมฺปิ ฉาริกมฺปิ ปกฺขิปนฺตีติ. ถทฺธตฺถาย\csromansharedfootnote{f750b7de68495f91}{พนฺธนตฺถาย (สี, สฺยา)} ภควาติ. สเจ ภิกฺขเว ถทฺธตฺถาย คุเฬ ปิฏฺฐมฺปิ ฉาริกมฺปิ ปกฺขิปนฺติ, โส จ คุโฬเตฺวว สงฺขํ คจฺฉติ. อนุชานามิ ภิกฺขเว ยถาสุขํ คุฬํ ปริภุญฺชิ\-ตุนฺติ.}

\prose{อทฺทสา โข อายสฺมา กงฺขาเรวโต อนฺตรามคฺเค วจฺเจ มุคฺคํ ชาตํ, ปสฺสิตฺวา “อกปฺปิยา มุคฺคา, ปกฺกาปิ มุคฺคา ชายนฺตี”ติ กุกฺกุจฺจายนฺโต สปริโส มุคฺคํ น ปริภุญฺชติ.}

\setlength{\csromangathapagewidth}{0pt}
\setlength{\csromangathawakwidth}{0pt}
\csromangathameasuregroup{เยปิสฺส โสตพฺพํ มญฺญนฺติ,}{\csromangathabat{เยปิสฺส โสตพฺพํ มญฺญนฺติ,}{เตปิ มุคฺคํ น ปริภุญฺชนฺติ.}}
\csromangathasetleft{เยปิสฺส โสตพฺพํ มญฺญนฺติ,}
\csromangathagroup{\csromangathastanza{\csromangathabat{เยปิสฺส โสตพฺพํ มญฺญนฺติ,}{เตปิ มุคฺคํ น ปริภุญฺชนฺติ.}}}

\prose{ภควโต เอตมตฺถํ อาโรเจสุํ. สเจ\csromansharedfootnote{338dd9d64fd7ee1e}{[มิสฺสินฺคฺ โนเต 0]} ภิกฺขเว ปกฺกาปิ มุคฺคา ชายนฺติ, อนุชานามิ ภิกฺขเว ยถาสุขํ มุคฺคํ ปริภุญฺชิ\-ตุนฺติ.}

\proseitem{273}{\csromanfolio{304}เตน โข ปน สมเยน อญฺญตรสฺส ภิกฺขุโน อุทรวาตาพาโธ โหติ, โส โลณโสวีรกํ อปายิ, ตสฺส โส อุทรวาตาพาโธ ปฏิปฺปสฺสมฺภิ. ภควโต เอตมตฺถํ อาโรเจสุํ. อนุชานามิ ภิกฺขเว คิลานสฺส โลณโสวีรกํ, อคิลานสฺส อุทก\-สมฺภินฺนํ ปาน\-ปริโภเคน ปริภุญฺชิ\-ตุนฺติ.}

\csromanmatikaanchor{mtk.170}
\csromantocmarkref{section}{164. อนฺโตวุฏฺฐาทิปฏิกฺเขปกถา}{mtk.170}
\bookmark[dest={mtk.170},level=1]{164. อนฺโตวุฏฺฐาทิปฏิกฺเขปกถา}
\csromanheader{1}{164. \textbf{อนฺโต}\-\textbf{วุฏฺฐา}\-\textbf{ทิ}\-\textbf{ปฏิกฺเขป}\-\textbf{กถา}}
\proseitem{274}{อถ โข ภควา อนุปุพฺเพน จาริกํ จรมาโน เยน ราชคหํ ตทวสริ, ตตฺร สุทํ ภควา ราชคเห วิหรติ เวฬุวเน กลนฺทก\-นิวาเป. เตน โข ปน สมเยน ภควโต อุทรวาตาพาโธ โหติ. อถ โข อายสฺมา อานนฺโท “ปุพฺเพปิ ภควโต อุทรวาตาพาโธ เตกฏุลยาคุยา ผาสุ โหตี”ติ สามํ ติลมฺปิ ตณฺฑุลมฺปิ มุคฺคมฺปิ วิญฺญาเปตฺวา อนฺโต วาเสตฺวา อนฺโต สามํ ปจิตฺวา ภควโต อุปนาเมสิ “ปิวตุ ภควา เตกฏุลยาคุนฺ”ติ. ชานนฺตาปิ ตถาคตา ปุจฺฉนฺติ, ชานนฺตาปิ น ปุจฺฉนฺติ. กาลํ วิทิตฺวา ปุจฺฉนฺติ, กาลํ วิทิตฺวา น ปุจฺฉนฺติ. อตฺถสํหิตํ ตถาคตา ปุจฺฉนฺติ, โน อนตฺถ\-สํหิตํ, อนตฺถสํหิเต เสตุฆาโต ตถาคตานํ. ทฺวีหิ อากาเรหิ พุทฺธา ภควนฺโต ภิกฺขู ปฏิปุจฺฉนฺติ “ธมฺมํ วา เทเสสฺสาม สาวกานํ วา สิกฺขาปทํ ปญฺญเปสฺสามา”ติ. อถ โข ภควา อายสฺมนฺตํ อานนฺทํ อามนฺเตสิ “กุตายํ อานนฺท ยาคู”ติ. อถ โข อายสฺมา อานนฺโท ภควโต เอตมตฺถํ อาโรเจสิ. วิครหิ พุทฺโธ ภควา อนนุจฺฉวิกํ อานนฺท อนนุโลมิกํ อปฺปติรูปํ อสฺสามณกํ อกปฺปิยํ อกรณียํ. กถํ หิ นาม ตฺวํ อานนฺท เอวรูปาย พาหุลฺลาย เจเตสฺสสิ. ยทปิ อานนฺท อนฺโต วุฏฺฐํ\csromansharedfootnote{c3e63967816f811c}{วุตฺถํ (สี, สฺยา, ก)}, ตทปิ อกปฺปิยํ. ยทปิ อนฺโต ปกฺกํ, ตทปิ อกปฺปิยํ. ยทปิ สามํ ปกฺกํ, ตทปิ อกปฺปิยํ. เนตํ อานนฺท อปฺปสนฺนานํ วา ปสาทาย ฯเปฯ วิครหิตฺวา ธมฺมิํ กถํ กตฺวา ภิกฺขู อามนฺเตสิ “น ภิกฺขเว อนฺโต วุฏฺฐํ อนฺโต ปกฺกํ สามํ ปกฺกํ ปริภุญฺชิ\-ตพฺพํ, โย ปริภุญฺเชยฺย, อาปตฺติ ทุกฺกฏสฺส. อนฺโต เจ ภิกฺขเว วุฏฺฐํ อนฺโต ปกฺกํ สามํ ปกฺกํ, ตญฺเจ ปริภุญฺเชยฺย, อาปตฺติ ติณฺณํ ทุกฺกฏานํ. อนฺโต เจ ภิกฺขเว วุฏฺฐํ อนฺโต ปกฺกํ อญฺเญหิ ปกฺกํ, ตญฺเจ ปริภุญฺเชยฺย, อาปตฺติ ทฺวินฺนํ ทุกฺกฏานํ. \csromanfolio{305}อนฺโต เจ ภิกฺขเว วุฏฺฐํ พหิ ปกฺกํ สามํ ปกฺกํ, ตญฺเจ ปริภุญฺเชยฺย, อาปตฺติ ทฺวินฺนํ ทุกฺกฏานํ. พหิ เจ ภิกฺขเว วุฏฺฐํ อนฺโต ปกฺกํ สามํ ปกฺกํ, ตญฺเจ ปริภุญฺเชยฺย, อาปตฺติ ทฺวินฺนํ ทุกฺกฏานํ. อนฺโต เจ ภิกฺขเว วุฏฺฐํ พหิ ปกฺกํ อญฺเญหิ ปกฺกํ, ตญฺเจ ปริภุญฺเชยฺย, อาปตฺติ ทุกฺกฏสฺส. พหิ เจ ภิกฺขเว วุฏฺฐํ อนฺโต ปกฺกํ อญฺเญหิ ปกฺกํ, ตญฺเจ ปริภุญฺเชยฺย, อาปตฺติ ทุกฺกฏสฺส. พหิ เจ ภิกฺขเว วุฏฺฐํ พหิ ปกฺกํ สามํ ปกฺกํ, ตญฺเจ ปริภุญฺเชยฺย, อาปตฺติ ทุกฺกฏสฺส. พหิ เจ ภิกฺขเว วุฏฺฐํ พหิ ปกฺกํ อญฺเญหิ ปกฺกํ, ตญฺเจ ปริภุญฺเชยฺย, อนาปตฺตี”ติ.}

\prose{เตน โข ปน สมเยน ภิกฺขู “ภควตา สามํปาโก ปฏิกฺขิตฺโต”ติ ปุน ปาเก กุกฺกุจฺจายนฺติ. ภควโต เอตมตฺถํ อาโรเจสุํ. อนุชานามิ ภิกฺขเว ปุน ปากํ ปจิตุนฺติ.}

\prose{เตน โข ปน สมเยน ราชคหํ ทุพฺภิกฺขํ โหติ, มนุสฺสา โลณมฺปิ เตลมฺปิ ตณฺฑุลมฺปิ ขาทนียมฺปิ อารามํ อาหรนฺติ, ตานิ ภิกฺขู พหิ วาเสนฺติ, อุกฺกปิณฺฑกาปิ ขาทนฺติ, โจราปิ หรนฺติ. ภควโต เอตมตฺถํ อาโรเจสุํ. อนุชานามิ ภิกฺขเว อนฺโต วาเสตุนฺติ. อนฺโต วาเสตฺวา พหิ ปาเจนฺติ, ทมกา ปริวาเรนฺติ, ภิกฺขู อวิสฺสฏฺฐา ปริภุญฺชนฺติ. ภควโต เอตมตฺถํ อาโรเจสุํ. อนุชานามิ ภิกฺขเว อนฺโต ปจิตุนฺติ. ทุพฺภิกฺเข กปฺปิยการกา พหุตรํ หรนฺติ, อปฺปตรํ ภิกฺขูนํ เทนฺติ. ภควโต เอตมตฺถํ อโรเจสุํ. อนุชานามิ ภิกฺขเว สามํ ปจิตุํ, อนุชานามิ ภิกฺขเว อนฺโต วุฏฺฐํ อนฺโต ปกฺกํ สามํ ปกฺกนฺติ.}

\csromanmatikaanchor{mtk.171}
\csromantocmarkref{section}{165. อุคฺคหิตปฏิคฺคหณา}{mtk.171}
\bookmark[dest={mtk.171},level=1]{165. อุคฺคหิตปฏิคฺคหณา}
\csromanheader{1}{165. \textbf{อุคฺคหิต}\-\textbf{ปฏิคฺคหณา}}
\proseitem{275}{เตน โข ปน สมเยน สมฺพหุลา ภิกฺขู กาสีสุ วสฺสํวุฏฺฐา ราชคหํ คจฺฉนฺตา ภควนฺตํ ทสฺสนาย อนฺตรามคฺเค น ลภิํสุ ลูขสฺส วา ปณีตสฺส วา โภชนสฺส ยาวทตฺถํ ปาริปูริํ, พหุญฺจ ผลขาทนียํ อโหสิ, กปฺปิยการโก จ น อโหสิ. อถ โข เต ภิกฺขู กิลนฺตรูปา เยน ราชคหํ เวฬุวนํ กลนฺทก\-นิวาโป, เยน ภควา เตนุ\-ปสงฺกมิํสุ, อุปสงฺกมิตฺวา ภควนฺตํ อภิวาเทตฺวา เอกมนฺตํ นิสีทิํสุ. อาจิณฺณํ โข ปเนตํ พุทฺธานํ ภควนฺตานํ อาคนฺตุเกหิ ภิกฺขูหิ สทฺธิํ ปฏิสมฺโมทิตุํ. อถ โข ภควา เต ภิกฺขู เอตทโวจ \csromanfolio{306}“กจฺจิ ภิกฺขเว ขมนียํ, กจฺจิ ยาปนียํ, กจฺจิตฺถ อปฺป\-กิลมเถน อทฺธานํ อาคตา, กุโต จ ตุเมฺห ภิกฺขเว อาคจฺฉถา”ติ. ขมนียํ ภควา, ยาปนียํ ภควา, อิธ มยํ ภนฺเต กาสีสุ วสฺสํวุฏฺฐา ราชคหํ อาคจฺฉนฺตา ภควนฺตํ ทสฺสนาย อนฺตรามคฺเค น ลภิมฺหา ลูขสฺส วา ปณีตสฺส วา โภชนสฺส ยาวทตฺถํ ปาริปูริํ, พหุญฺจ ผลขาทนียํ อโหสิ, กปฺปิยการโก จ น อโหสิ, เตน มยํ กิลนฺตรูปา อทฺธานํ อาคตาติ. อถ โข ภควา เอตสฺมิํ นิทาเน เอตสฺมิํ ปกรเณ ธมฺมิํ กถํ กตฺวา ภิกฺขู อามนฺเตสิ “อนุชานามิ ภิกฺขเว ยตฺถ ผลขาทนียํ ปสฺสติ, กปฺปิยการโก จ น โหติ, สามํ คเหตฺวา หริตฺวา กปฺปิยการเก ปสฺสิตฺวา ภูมิยํ นิกฺขิปิตฺวา ปฏิคฺคหาเปตฺวา ปริภุญฺชิตุํ, อนุชานามิ ภิกฺขเว อุคฺคหิตํ ปฏิคฺคหิตุนฺ”ติ.}

\proseitem{276}{เตน โข ปน สมเยน อญฺญตรสฺส พฺราหฺมณสฺส นวา จ ติลา นวญฺจ มธุ อุปฺปนฺนา โหนฺติ. อถ โข ตสฺส พฺราหฺมณสฺส เอตทโหสิ “ยนฺนูนาหํ นเว จ ติเล นวญฺจ มธุํ พุทฺธ\-ปฺปมุขสฺส ภิกฺขุ\-สํฆสฺส ทเทยฺยนฺ”ติ. อถ โข โส พฺราหฺมโณ เยน ภควา เตนุปสงฺกมิ, อุปสงฺกมิตฺวา ภควตา สทฺธิํ ปฏิสมฺโมทิ, สมฺโมทนียํ กถํ สารณียํ วีติสาเรตฺวา เอกมนฺตํ อฏฺฐาสิ. เอกมนฺตํ ฐิโต โข โส พฺราหฺมโณ ภควนฺตํ เอตทโวจ “อธิวาเสตุ เม ภวํ โคตโม สฺวาตนาย ภตฺตํ สทฺธิํ ภิกฺขุ\-สํเฆนา”ติ. อธิวาเสสิ ภควา ตุณฺหีภาเวน. อถ โข โส พฺราหฺมโณ ภควโต อธิวาสนํ วิทิตฺวา ปกฺกามิ. อถ โข โส พฺราหฺมโณ ตสฺสา รตฺติยา อจฺจเยน ปณีตํ ขาทนียํ โภชนียํ ปฏิยาทาเปตฺวา ภควโต กาลํ อาโรจาเปสิ “กาโล โภ โคตม นิฏฺฐิตํ ภตฺตนฺ”ติ. อถ โข ภควา ปุพฺพณฺหสมยํ นิวาเสตฺวา ปตฺต\-จีวรมา\-ทาย เยน ตสฺส พฺราหฺมณสฺส นิเวสนํ เตนุปสงฺกมิ, อุปสงฺกมิตฺวา ปญฺญตฺเต อาสเน นิสีทิ สทฺธิํ ภิกฺขุ\-สํเฆน. อถ โข โส พฺราหฺมโณ พุทฺธ\-ปฺปมุขํ ภิกฺขุสํฆํ ปณีเตน ขาทนีเยน โภชนีเยน สหตฺถา สนฺตปฺเปตฺวา สมฺปวาเรตฺวา ภควนฺตํ ภุตฺตาวิํ โอนีต\-ปตฺต\-ปาณิํ เอกมนฺตํ นิสีทิ. เอกมนฺตํ นิสินฺนํ โข ตํ พฺราหฺมณํ ภควา ธมฺมิยา กถาย สนฺทสฺเสตฺวา สมาทเปตฺวา สมุตฺเตเชตฺวา สมฺปหํเสตฺวา อุฏฺฐายาสนา ปกฺกามิ.}

\prose{\csromanfolio{307}อถ โข ตสฺส พฺราหฺมณสฺส อจิร\-ปกฺกนฺตสฺส ภควโต เอตทโหสิ “เยสํ โข มยา อตฺถาย พุทฺธปฺปมุโข ภิกฺขุสํโฆ นิมนฺติโต ‘นเว จ ติเล นวญฺจ มธุํ ทสฺสามี’ติ, เต มยา ปมุฏฺฐา ทาตุํ, ยนฺนูนาหํ นเว จ ติเล นวญฺจ มธุํ โกลมฺเพหิ จ ฆเฏหิ จ อารามํ หราเปยฺยนฺ”ติ.}

\prose{อถ โข โส พฺราหฺมโณ นเว จ ติเล นวญฺจ มธุํ โกลมฺเพหิ จ ฆเฏหิ จ อารามํ หราเปตฺวา เยน ภควา เตนุปสงฺกมิ, อุปสงฺกมิตฺวา เอกมนฺตํ อฏฺฐาสิ, เอกมนฺตํ ฐิโต โข โส พฺราหฺมโณ ภควนฺตํ เอตทโวจ “เยสํ โข มยา โภ โคตม อตฺถาย พุทฺธปฺปมุโข ภิกฺขุสํโฆ นิมนฺติโต ‘นเว จ ติเล นวญฺจ มธุํ ทสฺสามี’ติ, เต มยา ปมุฏฺฐา ทาตุํ, ปฏิคฺคณฺหาตุ เม ภวํ โคตโม นเว จ ติเล นวญฺจ มธุนฺ”ติ.}

\prose{เตน หิ พฺราหฺมณ ภิกฺขูนํ เทหีติ. เตน โข ปน สมเยน ภิกฺขู ทุพฺภิกฺเข อปฺปมตฺตเกปิ ปวาเรนฺติ, ปฏิสงฺขาปิ ปฏิกฺขิปนฺติ, สพฺโพ จ สํโฆ ปวาริโต โหติ, ภิกฺขู กุกฺกุจฺจายนฺตา น ปฏิคฺคณฺหนฺติ. ปฏิคฺคณฺหถ ภิกฺขเว ปริภุญฺชถ, อนุชานามิ ภิกฺขเว ตโต นีหฏํ ภุตฺตาวินา ปวาริเตน อนติริตฺตํ ปริภุญฺชิ\-ตุนฺติ.}

\csromanmatikaanchor{mtk.172}
\csromantocmarkref{section}{166. ปฏิคฺคหิตาทิอนุชานนา}{mtk.172}
\bookmark[dest={mtk.172},level=1]{166. ปฏิคฺคหิตาทิอนุชานนา}
\csromanheader{1}{166. ปฏิคฺคหิตา\-ทิ\-อนุชานนา}
\proseitem{277}{\csromansymbolfootnote{*}{วิ 2. 132 ปิฏฺเฐปิ.}เตน โข ปน สมเยน อายสฺมโต อุปนนฺทสฺส สกฺยปุตฺตสฺส อุปฏฺฐากกุลํ สํฆสฺสตฺถาย ขาทนียํ ปาเหสิ “อยฺยสฺส อุปนนฺทสฺส ทสฺเสตฺวา สํฆสฺส ทาตพฺพนฺ”ติ. เตน โข ปน สมเยน อายสฺมา อุปนนฺโท สกฺยปุตฺโต คามํ ปิณฺฑาย ปวิฏฺโฐ โหติ. อถ โข เต มนุสฺสา อารามํ คนฺตฺวา ภิกฺขู ปุจฺฉิํสุ “กหํ ภนฺเต อยฺโย อุปนนฺโท”ติ. เอสาวุโส อายสฺมา อุปนนฺโท สกฺยปุตฺโต คามํ ปิณฺฑาย ปวิฏฺโฐติ. อิทํ ภนฺเต ขาทนียํ อยฺยสฺส อุปนนฺทสฺส ทสฺเสตฺวา สํฆสฺส ทาตพฺพนฺติ. ภควโต เอตมตฺถํ อาโรเจสุํ. เตน หิ ภิกฺขเว ปฏิคฺคเหตฺวา นิกฺขิปถ ยาว อุปนนฺโท อาคจฺฉตีติ. อถ โข อายสฺมา อุปนนฺโท สกฺยปุตฺโต ปุเรภตฺตํ กุลานิ ปยิรุปาสิตฺวา ทิวา อาคจฺฉติ. เตน โข ปน สมเยน ภิกฺขู ทุพฺภิกฺเข อปฺปมตฺตเกปิ ปวาเรนฺติ, ปฏิสงฺขาปิ ปฏิกฺขิปนฺติ, สพฺโพ จ สํโฆ ปวาริโต โหติ, ภิกฺขู กุกฺกุจฺจายนฺตา น ปฏิคฺคณฺหนฺติ, ปฏิคฺคณฺหถ ภิกฺขเว ปริภุญฺชถ, อนุชานามิ ภิกฺขเว ปุเรภตฺตํ ปฏิคฺคหิตํ ภุตฺตาวินา ปวาริเตน อนติริตฺตํ ปริภุญฺชิ\-ตุนฺติ.}

\proseitem{278}{\csromanfolio{308}อถ โข ภควา ราชคเห ยถาภิรนฺตํ วิหริตฺวา เยน สาวตฺถิ เตน จาริกํ ปกฺกามิ. อนุปุพฺเพน จาริกํ จรมาโน เยน สาวตฺถิ ตทวสริ, ตตฺร สุทํ ภควา สาวตฺถิยํ วิหรติ เชตวเน อนาถ\-ปิณฺฑิกสฺส อาราเม. เตน โข ปน สมเยน อายสฺมโต สาริปุตฺตสฺส กายฑาหาพาโธ โหติ. อถ โข อายสฺมา มหาโมคฺคลฺลาโน เยนายสฺมา สาริปุตฺโต เตนุปสงฺกมิ, อุปสงฺกมิตฺวา อายสฺมนฺตํ สาริปุตฺตํ เอตทโวจ “ปุพฺเพ เต อาวุโส สาริปุตฺต กายฑาหาพาโธ เกน ผาสุ โหตี”ติ. ภิเสหิ จ เม อาวุโส มุฬาลิกาหิ จาติ. อถ โข อายสฺมา มหาโมคฺคลฺลาโน เสยฺยถาปิ นาม พลวา ปุริโส สมิญฺชิตํ วา พาหํ ปสาเรยฺย, ปสาริตํ วา พาหํ สมิญฺเชยฺย, เอวเมว เชตวเน อนฺตรหิโต มนฺทากินิยา โปกฺขรณิยา ตีเร ปาตุรโหสิ. อทฺทสา โข อญฺญตโร นาโค อายสฺมนฺตํ มหา\-โมคฺคลฺลานํ ทูรโตว อาคจฺฉนฺตํ, ทิสฺวาน อายสฺมนฺตํ มหา\-โมคฺคลฺลานํ เอตทโวจ “เอตุ โข ภนฺเต อยฺโย มหาโมคฺคลฺลาโน, สฺวาคตํ ภนฺเต อยฺยสฺส มหา\-โมคฺคลฺลานสฺส, เกน ภนฺเต อยฺยสฺส อตฺโถ, กิํ ทมฺมี”ติ. ภิเสหิ จ เม อาวุโส อตฺโถ มุฬาลิกาหิ จาติ. อถ โข โส นาโค อญฺญตรํ นาคํ อาณาเปสิ “เตน หิ ภเณ อยฺยสฺส ภิเส จ มุฬาลิกาโย จ ยาวทตฺถํ เทหี”ติ. อถ โข โส นาโค มนฺทากินิํ โปกฺขรณิํ โอคาเหตฺวา โสณฺฑาย ภิสญฺจ มุฬาลิกญฺจ อพฺพาหิตฺวา สุวิกฺขาลิตํ วิกฺขาเลตฺวา ภณฺฑิกํ พนฺธิตฺวา เยนายสฺมา มหาโมคฺคลฺลาโน เตนุปสงฺกมิ. อถ โข อายสฺมา มหาโมคฺคลฺลาโน เสยฺยถาปิ นาม พลวา ปุริโส สมิญฺชิตํ วา พาหํ ปสาเรยฺย, ปสาริตํ วา พาหํ สมิญฺเชยฺย, เอวเมว มนฺทากินิยา โปกฺขรณิยา ตีเร อนฺตรหิโต เชตวเน ปาตุรโหสิ. โสปิ โข นาโค มนฺทากินิยา โปกฺขรณิยา ตีเร อนฺตรหิโต เชตวเน ปาตุรโหสิ. อถ โข โส นาโค อายสฺมโต มหา\-โมคฺคลฺลานสฺส ภิเส จ มุฬาลิกาโย จ ปฏิคฺคหาเปตฺวา เชตวเน อนฺตรหิโต มนฺทากินิยา โปกฺขรณิยา ตีเร ปาตุรโหสิ. อถ โข อายสฺมา มหาโมคฺคลฺลาโน อายสฺมโต สาริปุตฺตสฺส ภิเส จ มุฬาลิกาโย จ อุปนาเมสิ. อถ โข อายสฺมโต สาริปุตฺตสฺส ภิเส จ มุฬาลิกาโย จ ภุตฺตสฺส กายฑาหาพาโธ ปฏิปฺปสฺสมฺภิ, พหู ภิสา \csromanfolio{309}จ มุฬาลิกาโย จ อวสิฏฺฐา โหนฺติ. เตน โข ปน สมเยน ภิกฺขู ทุพฺภิกฺเข อปฺปมตฺตเกปิ ปวาเรนฺติ, ปฏิสงฺขาปิ ปฏิกฺขิปนฺติ, สพฺโพ จ สํโฆ ปวาริโต โหติ, ภิกฺขู กุกฺกุจฺจายนฺตา น ปฏิคฺคณฺหนฺติ. ปฏิคฺคณฺหถ ภิกฺขเว ปริภุญฺชถ, อนุชานามิ ภิกฺขเว วนฏฺฐํ โปกฺขรฏฺฐํ ภุตฺตาวินา ปวาริเตน อนติริตฺตํ ปริภุญฺชิ\-ตุนฺติ.}

\prose{เตน โข ปน สมเยน สาวตฺถิยํ พหุํ ผลขาทนียํ อุปฺปนฺนํ โหติ, กปฺปิยการโก จ น โหติ, ภิกฺขู กุกฺกุจฺจายนฺตา ผลํ น ปริภุญฺชนฺติ. ภควโต เอตมตฺถํ อาโรเจสุํ. อนุชานามิ ภิกฺขเว อพีชํ นิพฺพตฺตพีชํ\csromansharedfootnote{700411318ee7dbe0}{นิพฺพฏฺฏพีชํ (สี), นิพฺพฏพีชํ (สฺยา), นิปฺปฏฺฏพีชํ (ก)} อกตกปฺปํ ผลํ ปริภุญฺชิ\-ตุนฺติ.}

\csromanmatikaanchor{mtk.173}
\csromantocmarkref{section}{167. สตฺถกมฺมปฏิกฺเขปกถา}{mtk.173}
\bookmark[dest={mtk.173},level=1]{167. สตฺถกมฺมปฏิกฺเขปกถา}
\csromanheader{1}{167. \textbf{สตฺถกมฺม}\-\textbf{ปฏิกฺเขป}\-\textbf{กถา}}
\proseitem{279}{อถ โข ภควา สาวตฺถิยํ ยถาภิรนฺตํ วิหริตฺวา เยน ราชคหํ เตน จาริกํ ปกฺกามิ. อนุปุพฺเพน จาริกํ จรมาโน เยน ราชคหํ ตทวสริ, ตตฺร สุทํ ภควา ราชคเห วิหรติ เวฬุวเน กลนฺทก\-นิวาเป. เตน โข ปน สมเยน อญฺญตรสฺส ภิกฺขุโน ภคนฺทลา\-พาโธ โหติ, อากาสโคตฺโต เวชฺโช สตฺถกมฺมํ กโรติ. อถ โข ภควา เสนาสนจาริกํ อาหิณฺฑนฺโต เยน ตสฺส ภิกฺขุโน วิหาโร เตนุปสงฺกมิ. อทฺทสา โข อากาสโคตฺโต เวชฺโช ภควนฺตํ ทูรโตว อาคจฺฉนฺตํ, ทิสฺวาน ภควนฺตํ เอตทโวจ “อาคจฺฉตุ ภวํ โคตโม อิมสฺส ภิกฺขุโน วจฺจมคฺคํ ปสฺสตุ, เสยฺยถาปิ โคธามุขนฺ”ติ. อถ โข ภควา “โส มํ ขฺวายํ โมฆปุริโส อุปฺปณฺเฑตี”ติ ตโตว ปฏินิวตฺติตฺวา เอตสฺมิํ นิทาเน เอตสฺมิํ ปกรเณ ภิกฺขุสํฆํ สนฺนิปาตาเปตฺวา ภิกฺขู ปฏิปุจฺฉิ “อตฺถิ กิร ภิกฺขเว อมุกสฺมิํ วิหาเร ภิกฺขุ คิลาโน”ติ. อตฺถิ ภควาติ. กิํ ตสฺส ภิกฺขเว ภิกฺขุโน อาพาโธติ. ตสฺส ภนฺเต อายสฺมโต ภคนฺทลา\-พาโธ, อากาสโคตฺโต เวชฺโช สตฺถกมฺมํ กโรตีติ. วิครหิ พุทฺโธ ภควา อนนุจฺฉวิกํ ภิกฺขเว ตสฺส โมฆปุริสสฺส อนนุโลมิกํ อปฺปติรูปํ อสฺสามณกํ อกปฺปิยํ อกรณียํ. กถํ หิ นาม โส ภิกฺขเว โมฆปุริโส สมฺพาเธ สตฺถกมฺมํ การาเปสฺสติ. สมฺพาเธ ภิกฺขเว สุขุมา ฉวิ ทุโรปโย \csromanfolio{310}วโณ ทุปฺปริหารํ สตฺถํ. เนตํ ภิกฺขเว อปฺปสนฺนานํ วา ปสาทาย - ป- วิครหิตฺวา ฯเปฯ ธมฺมิํ กถํ กตฺวา ภิกฺขู อามนฺเตสิ “น ภิกฺขเว สมฺพาเธ สตฺถกมฺมํ การาเปตพฺพํ, โย การาเปยฺย, อาปตฺติ ถุลฺลจฺจยสฺสา”ติ.}

\prose{เตน โข ปน สมเยน ฉพฺพคฺคิยา ภิกฺขู “ภควตา สตฺถกมฺมํ ปฏิกฺขิตฺตนฺ”ติ วตฺถิกมฺมํ การาเปนฺติ. เย เต ภิกฺขู อปฺปิจฺฉา, เต อุชฺฌายนฺติ ขิยฺยนฺติ วิปาเจนฺติ “กถํ หิ นาม ฉพฺพคฺคิยา ภิกฺขู วตฺถิกมฺมํ การาเปสฺสนฺตี”ติ. อถ โข เต ภิกฺขู ภควโต เอตมตฺถํ อาโรเจสุํ. สจฺจํ กิร ภิกฺขเว ฉพฺพคฺคิยา ภิกฺขู วตฺถิกมฺมํ การาเปนฺตีติ. สจฺจํ ภควาติ ฯเปฯ วิครหิตฺวา ธมฺมิํ กถํ กตฺวา ภิกฺขู อามนฺเตสิ “น ภิกฺขเว สมฺพาธสฺส สามนฺตา ทฺวงฺคุลา สตฺถกมฺมํ วา วตฺถิกมฺมํ วา การาเปตพฺพํ, โย การาเปยฺย, อาปตฺติ ถุลฺลจฺจยสฺสา”ติ.}

\csromanmatikaanchor{mtk.174}
\csromantocmarkref{section}{168. มนุสฺสมํสปฏิกฺเขปกถา}{mtk.174}
\bookmark[dest={mtk.174},level=1]{168. มนุสฺสมํสปฏิกฺเขปกถา}
\csromanheader{1}{168. \textbf{มนุสฺสมํส}\-\textbf{ปฏิกฺเขป}\-\textbf{กถา}}
\proseitem{280}{อถ โข ภควา ราชคเห ยถาภิรนฺตํ วิหริตฺวา เยน พาราณสี เตน จาริกํ ปกฺกามิ. อนุปุพฺเพน จาริกํ จรมาโน เยน พาราณสี ตทวสริ, ตตฺร สุทํ ภควา พาราณสิยํ วิหรติ อิสิปตเน มิคทาเย. เตน โข ปน สมเยน พาราณสิยํ สุปฺปิโย จ อุปาสโก สุปฺปิยา จ อุปาสิกา อุภโตปสนฺนา โหนฺติ ทายกา การกา สํฆุปฏฺฐากา. อถ โข สุปฺปิยา อุปาสิกา อารามํ คนฺตฺวา วิหาเรน วิหารํ ปริเวเณน ปริเวณํ อุปสงฺกมิตฺวา ภิกฺขู ปุจฺฉติ “โก ภนฺเต คิลาโน, กสฺส กิํ อาหริยตู”ติ. เตน โข ปน สมเยน อญฺญตเรน ภิกฺขุนา วิเรจนํ ปีตํ โหติ, อถ โข โส ภิกฺขุ สุปฺปิยํ อุปาสิกํ เอตทโวจ “มยา โข ภคินิ วิเรจนํ ปีตํ, อตฺโถ เม ปฏิจฺฉาทนีเยนา”ติ. สุฏฺฐุ อยฺย อาหริยิสฺสตีติ ฆรํ คนฺตฺวา อนฺเตวาสิํ อาณาเปสิ “คจฺฉ ภเณ ปวตฺตมํสํ ชานาหี”ติ. เอวํ อยฺเยติ โข โส ปุริโส สุปฺปิยาย อุปาสิกาย ปฏิสฺสุณิตฺวา เกวลกปฺปํ พาราณสิํ อาหิณฺฑนฺโต น อทฺทส ปวตฺตมํสํ. อถ โข โส ปุริโส เยน สุปฺปิยา อุปาสิกา เตนุปสงฺกมิ, อุปสงฺกมิตฺวา สุปฺปิยํ อุปาสิกํ เอตทโวจ “นตฺถยฺเย ปวตฺตมํสํ, มาฆาโต อชฺชา”ติ. อถ โข สุปฺปิยาย อุปาสิกาย เอตทโหสิ “ตสฺส โข คิลานสฺส ภิกฺขุโน ปฏิจฺฉาทนียํ อลภนฺตสฺส อาพาโธ วา อภิวฑฺฒิสฺสติ, \csromanfolio{311}กาลงฺกิริยา วา ภวิสฺสติ, น โข เมตํ ปติรูปํ ยาหํ ปฏิสฺสุณิตฺวา น หราเปยฺยนฺ”ติ, โปตฺถนิกํ คเหตฺวา อูรุมํสํ อุกฺกนฺติตฺวา ทาสิยา อทาสิ “หนฺท เช อิมํ มํสํ สมฺปาเทตฺวา อมุกสฺมิํ วิหาเร ภิกฺขุ คิลาโน, ตสฺส ทชฺชาหิ, โย จ มํ ปุจฺฉติ, คิลานาติ ปฏิเวเทหี”ติ อุตฺตราสงฺเคน อูรุํ เวเฐตฺวา โอวรกํ ปวิสิตฺวา มญฺจเก นิปชฺชิ. อถ โข สุปฺปิโย อุปาสโก ฆรํ คนฺตฺวา ทาสิํ ปุจฺฉิ “กหํ สุปฺปิยา”ติ. เอสายฺย โอวรเก นิปนฺนาติ. อถ โข สุปฺปิโย อุปาสโก เยน สุปฺปิยา อุปาสิกา เตนุปสงฺกมิ, อุปสงฺกมิตฺวา สุปฺปิยํ อุปาสิกํ เอตทโวจ “กิสฺส นิปนฺนาสี”ติ. คิลานามฺหีติ. กิํ เต อาพาโธติ. อถ โข สุปฺปิยา อุปาสิกา สุปฺปิยสฺส อุปาสกสฺส เอตมตฺถํ อาโรเจสิ. อถ โข สุปฺปิโย อุปาสโก “อจฺฉริยํ วต โภ, อพฺภุตํ วต โภ, ยาว สทฺธายํ สุปฺปิยา ปสนฺนา, ยตฺร หิ นาม อตฺตโนปิ มํสานิ ปริจฺจตฺตานิ, กิมฺปิมาย\csromansharedfootnote{70f6b85decbc7e2b}{กิํ ปนิมาย (สี, สฺยา)} อญฺญํ กิญฺจิ อเทยฺยํ ภวิสฺสตี”ติ หฏฺโฐ อุทคฺโค เยน ภควา เตนุปสงฺกมิ, อุปสงฺกมิตฺวา ภควนฺตํ อภิวาเทตฺวา เอกมนฺตํ นิสีทิ. เอกมนฺตํ นิสินฺโน โข สุปฺปิโย อุปาสโก ภควนฺตํ เอตทโวจ “อธิวาเสตุ เม ภนฺเต ภควา สฺวาตนาย ภตฺตํ สทฺธิํ ภิกฺขุ\-สํเฆนา”ติ. อธิวาเสสิ ภควา ตุณฺหีภาเวน. อถ โข สุปฺปิโย อุปาสโก ภควโต อธิวาสนํ วิทิตฺวา อุฏฺฐายาสนา ภควนฺตํ อภิวาเทตฺวา ปทกฺขิณํ กตฺวา ปกฺกามิ. อถ โข สุปฺปิโย อุปาสโก ตสฺสา รตฺติยา อจฺจเยน ปณีตํ ขาทนียํ โภชนียํ ปฏิยาทาเปตฺวา ภควโต กาลํ อาโรจาเปสิ “กาโล ภนฺเต นิฏฺฐิตํ ภตฺตนฺ”ติ. อถ โข ภควา ปุพฺพณฺหสมยํ นิวาเสตฺวา ปตฺต\-จีวรมา\-ทาย เยน สุปฺปิยสฺส อุปาสกสฺส นิเวสนํ เตนุปสงฺกมิ, อุปสงฺกมิตฺวา ปญฺญตฺเต อาสเน นิสีทิ สทฺธิํ ภิกฺขุ\-สํเฆน. อถ โข \textbf{สุปฺปิโย อุปาสโก เยน ภควา เตนุปสงฺกมิ, อุปสงฺกมิตฺวา} \textbf{ภควนฺตํ อภิวาเทตฺวา เอกมนฺตํ อฏฺฐาสิ. เอกมนฺตํ ฐิตํ โข} สุปฺปิยํ อุปาสกํ ภควา เอตทโวจ “กหํ สุปฺปิยา”ติ. คิลานา ภควาติ. เตน หิ อาคจฺฉตูติ. น ภควา อุสฺสหตีติ. เตน หิ ปริคฺคเหตฺวาปิ อาเนถาติ. อถ โข สุปฺปิโย อุปาสโก สุปฺปิยํ อุปาสิกํ \csromanfolio{312}ปริคฺคเหตฺวา อาเนสิ. ตสฺสา สห ทสฺสเนน ภควโต ตาว มหาวโณ รุโฬฺห อโหสิ สุจฺฉวิ\-โลม\-ชาโต. อถ โข สุปฺปิโย จ อุปาสโก สุปฺปิยา จ อุปาสิกา “อจฺฉริยํ วต โภ อพฺภุตํ วต โภ ตถาคตสฺส มหิทฺธิกตา มหานุภาวตา, ยตฺร หิ นาม สห ทสฺสเนน ภควโต ตาว มหาวโณ รุโฬฺห ภวิสฺสติ สุจฺฉวิ\-โลม\-ชาโต”ติ หฏฺฐา อุทคฺคา พุทฺธ\-ปฺปมุขํ ภิกฺขุสํฆํ ปณีเตน ขาทนีเยน โภชนีเยน สหตฺถา สนฺตปฺเปตฺวา สมฺปวาเรตฺวา ภควนฺตํ ภุตฺตาวิํ โอนีต\-ปตฺต\-ปาณิํ เอกมนฺตํ นิสีทิํสุ. อถ โข ภควา สุปฺปิยญฺจ อุปาสกํ สุปฺปิยญฺจ อุปาสิกํ ธมฺมิยา กถาย สนฺทสฺเสตฺวา สมาทเปตฺวา สมุตฺเตเชตฺวา สมฺปหํเสตฺวา อุฏฺฐายาสนา ปกฺกามิ.}

\prose{อถ โข ภควา เอตสฺมิํ นิทาเน เอตสฺมิํ ปกรเณ ภิกฺขุสํฆํ สนฺนิปาตาเปตฺวา ภิกฺขู ปฏิปุจฺฉิ “โก ภิกฺขเว สุปฺปิยํ อุปาสิกํ มํสํ วิญฺญาเปสี”ติ. เอวํ วุตฺเต โส ภิกฺขุ ภควนฺตํ เอตทโวจ “อหํ โข ภนฺเต สุปฺปิยํ อุปาสิกํ มํสํ วิญฺญาเปสินฺ”ติ. อาหริยิตฺถ ภิกฺขูติ. อาหริยิตฺถ ภควาติ. ปริภุญฺชิ ตฺวํ ภิกฺขูติ. ปริภุญฺชามหํ ภควาติ. ปฏิเวกฺขิ ตฺวํ ภิกฺขูติ. นาหํ ภควา ปฏิเวกฺขินฺติ. วิครหิ พุทฺโธ ภควา ฯเปฯ. กถํ หิ นาม ตฺวํ โมฆปุริส อปฺปฏิเวกฺขิตฺวา มํสํ ปริภุญฺชิสฺสสิ. มนุสฺสมํสํ โข ตยา โมฆปุริส ปริภุตฺตํ. เนตํ โมฆปุริส อปฺปสนฺนานํ วา ปสาทาย ฯเปฯ วิครหิตฺวา ธมฺมิํ กถํ กตฺวา ภิกฺขู อามนฺเตสิ “สนฺติ ภิกฺขเว มนุสฺสา สทฺธา ปสนฺนา, เตหิ อตฺตโนปิ มํสานิ ปริจฺจตฺตานิ. น ภิกฺขเว มนุสฺสมํสํ ปริภุญฺชิ\-ตพฺพํ, โย ปริภุญฺเชยฺย, อาปตฺติ ถุลฺลจฺจยสฺส. น จ ภิกฺขเว อปฺปฏิเวกฺขิตฺวา มํสํ ปริภุญฺชิ\-ตพฺพํ, โย ปริภุญฺเชยฺย, อาปตฺติ ทุกฺกฏสฺสา”ติ.}

\csromanmatikaanchor{mtk.175}
\csromantocmarkref{section}{169. หตฺถิมํสาทิปฏิกฺเขปกถา}{mtk.175}
\bookmark[dest={mtk.175},level=1]{169. หตฺถิมํสาทิปฏิกฺเขปกถา}
\csromanheader{1}{169. \textbf{หตฺถิ}\-\textbf{มํสา}\-\textbf{ทิ}\-\textbf{ปฏิกฺเขป}\-\textbf{กถา}}
\proseitem{281}{เตน โข ปน สมเยน รญฺโญ หตฺถี มรนฺติ, มนุสฺสา ทุพฺภิกฺเข หตฺถิมํสํ ปริภุญฺชนฺติ, ภิกฺขูนํ ปิณฺฑาย จรนฺตานํ หตฺถิมํสํ เทนฺติ, ภิกฺขู หตฺถิมํสํ ปริภุญฺชนฺติ. มนุสฺสา อุชฺฌายนฺติ ขิยฺยนฺติ วิปาเจนฺติ “กถํ หิ นาม สมณา สกฺยปุตฺติยา หตฺถิมํสํ ปริภุญฺชิสฺสนฺติ, ราชงฺคํ หตฺถี, สเจ ราชา ชาเนยฺย, น เนสํ อตฺตมโน อสฺสา”ติ. ภควโต \csromanfolio{313}เอตมตฺถํ อาโรเจสุํ. น ภิกฺขเว หตฺถิมํสํ ปริภุญฺชิ\-ตพฺพํ, โย ปริภุญฺเชยฺย, อาปตฺติ ทุกฺกฏสฺสาติ.}

\prose{เตน โข ปน สมเยน รญฺโญ อสฺสา มรนฺติ, มนุสฺสา ทุพฺภิกฺเข อสฺสมํสํ ปริภุญฺชนฺติ, ภิกฺขูนํ ปิณฺฑาย จรนฺตานํ อสฺสมํสํ เทนฺติ, ภิกฺขู อสฺสมํสํ ปริภุญฺชนฺติ.}

\prose{มนุสฺสา อุชฺฌายนฺติ ขิยฺยนฺติ วิปาเจนฺติ “กถํ หิ นาม สมณา สกฺยปุตฺติยา อสฺสมํสํ ปริภุญฺชิสฺสนฺติ, ราชงฺคํ อสฺสา, สเจ ราชา ชาเนยฺย, น เนสํ อตฺตมโน อสฺสา”ติ.}

\prose{ภควโต เอตมตฺถํ อาโรเจสุํ. น ภิกฺขเว อสฺสมํสํ ปริภุญฺชิ\-ตพฺพํ, โย ปริภุญฺเชยฺย, อาปตฺติ ทุกฺกฏสฺสาติ.}

\prose{เตน โข ปน สมเยน มนุสฺสา ทุพฺภิกฺเข สุนขมํสํ ปริภุญฺชนฺติ, ภิกฺขูนํ ปิณฺฑาย จรนฺตานํ สุนขมํสํ เทนฺติ, ภิกฺขู สุนขมํสํ ปริภุญฺชนฺติ, มนุสฺสา อุชฺฌายฺยนฺติ ขิยฺยนฺติ วิปาเจนฺติ “กถํ หิ นาม สมณา สกฺยปุตฺติยา สุนขมํสํ ปริภุญฺชิสฺสนฺติ, เชคุจฺโฉ สุนโข ปฏิกูโล”ติ.}

\prose{ภควโต เอตมตฺถํ อาโรเจสุํ. น ภิกฺขเว สุนขมํสํ ปริภุญฺชิ\-ตพฺพํ, โย ปริภุญฺเชยฺย, อาปตฺติ ทุกฺกฏสฺสาติ.}

\prose{เตน โข ปน สมเยน มนุสฺสา ทุพฺภิกฺเข อหิมํสํ ปริภุญฺชนฺติ, ภิกฺขูนํ ปิณฺฑาย จรนฺตานํ อหิมํสํ เทนฺติ, ภิกฺขู อหิมํสํ ปริภุญฺชนฺติ, มนุสฺสา อุชฺฌายนฺติ ขิยฺยนฺติ วิปาเจนฺติ “กถํ หิ นาม สมณา สกฺยปุตฺติยา อหิมํสํ ปริภุญฺชิสฺสนฺติ, เชคุจฺโฉ อหิ ปฏิกูโล”ติ.}

\prose{สุปสฺโสปิ\csromansharedfootnote{338dd9d64fd7ee1e}{[มิสฺสินฺคฺ โนเต 0]} นาคราชา เยน ภควา เตนุปสงฺกมิ, อุปสงฺกมิตฺวา ภควนฺตํ อภิวาเทตฺวา เอกมนฺตํ อฏฺฐาสิ, เอกมนฺตํ ฐิโต โข สุปสฺโส นาคราชา ภควนฺตํ เอตทโวจ “สนฺติ ภนฺเต นาคา อสฺสทฺธา อปฺปสนฺนา เต อปฺป\-มตฺต\-เกหิปิ ภิกฺขู วิเหเฐยฺยุํ, สาธุ ภนฺเต อยฺยา อหิมํสํ น ปริภุญฺเชยฺยุนฺ”ติ.}

\prose{อถ โข ภควา สุปสฺสํ นาคราชานํ ธมฺมิยา กถาย สนฺทสฺเสสิ ฯเปฯ ปทกฺขิณํ กตฺวา ปกฺกามิ. อถ โข ภควา เอตสฺมิํ นิทาเน เอตสฺมิํ ปกรเณ ธมฺมิํ กถํ กตฺวา ภิกฺขู อามนฺเตสิ “น ภิกฺขเว อหิมํสํ ปริภุญฺชิ\-ตพฺพํ, โย ปริภุญฺเชยฺย, อาปตฺติ ทุกฺกฏสฺสา”ติ.}

\prose{\csromanfolio{314}เตน โข ปน สมเยน ลุทฺทกา สีหํ หนฺตฺวา สีหมํสํ\csromansharedfootnote{d288e243bba4773e}{มํสํ (ก)2. อธิกานิ (สี, สฺยา)} ปริภุญฺชนฺติ, ภิกฺขูนํ ปิณฺฑาย จรนฺตานํ สีหมํสํ เทนฺติ, ภิกฺขู สีหมํสํ ปริภุญฺชิตฺวา อรญฺเญ วิหรนฺติ, สีหา สีหมํส\-คนฺเธน ภิกฺขู ปริปาเตนฺติ.}

\prose{ภควโต เอตมตฺถํ อาโรเจสุํ. น ภิกฺขเว สีหมํสํ ปริภุญฺชิ\-ตพฺพํ, โย ปริภุญฺเชยฺย, อาปตฺติ ทุกฺกฏสฺสาติ.}

\prose{เตน โข ปน สมเยน ลุทฺทกา พฺยคฺฆํ หนฺตฺวา ฯเปฯ ทีปิํ หนฺตฺวา - ป- อจฺฉํ หนฺตฺวา ฯเปฯ ตรจฺฉํ หนฺตฺวา ตรจฺฉมํสํ ปริภุญฺชนฺติ, ภิกฺขูนํ ปิณฺฑาย จรนฺตานํ ตรจฺฉมํสํ เทนฺติ, ภิกฺขู ตรจฺฉมํสํ ปริภุญฺชิตฺวา อรญฺเญ วิหรนฺติ, ตรจฺฉา ตรจฺฉมํส\-คนฺเธน ภิกฺขู ปริปาเตนฺติ.}

\prosewithcloserruled{ภควโต เอตมตฺถํ อาโรเจสุํ. น ภิกฺขเว ตรจฺฉมํสํ ปริภุญฺชิ\-ตพฺพํ, โย ปริภุญฺเชยฺย, อาปตฺติ ทุกฺกฏสฺสาติ.}{สุปฺปิย\-ภาณวาโร นิฏฺฐิโต ทุติโย.}
\csromanmatikaanchor{mtk.176}
\csromantocmarkref{section}{170. ยาคุมธุโคฬกานุชานนา}{mtk.176}
\bookmark[dest={mtk.176},level=1]{170. ยาคุมธุโคฬกานุชานนา}
\csromanheader{1}{170. \textbf{ยาคุ}\-\textbf{มธุโคฬกา}\-\textbf{นุชานนา}}
\proseitem{282}{อถ โข ภควา พาราณสิยํ ยถาภิรนฺตํ วิหริตฺวา เยน อนฺธกวินฺทํ เตน จาริกํ ปกฺกามิ มหตา ภิกฺขุ\-สํเฆน สทฺธิํ อฑฺฒเตลเสหิ ภิกฺขุสเตหิ. เตน โข ปน สมเยน ชานปทา มนุสฺสา พหุํ โลณมฺปิ เตลมฺปิ ตณฺฑุลมฺปิ ขาทนียมฺปิ สกเฏสุ อาโรเปตฺวา พุทฺธ\-ปฺปมุขสฺส ภิกฺขุ\-สํฆสฺส ปิฏฺฐิโต ปิฏฺฐิโต อนุพนฺธา โหนฺติ “ยทา ปฏิปาฏิํ ลภิสฺสาม, ตทา ภตฺตํ กริสฺสามา”ติ. ปญฺจมตฺตานิ จ วิฆาสาท\-สตานิ. อถ โข ภควา อนุปุพฺเพน จาริกํ จรมาโน เยน อนฺธกวินฺทํ ตทวสริ. อถ โข อญฺญตรสฺส พฺราหฺมณสฺส ปฏิปาฏิํ อลภนฺตสฺส เอตทโหสิ “อตีตานิ2 โข เม เทฺว มาสานิ พุทฺธ\-ปฺปมุขํ ภิกฺขุสํฆํ อนุพนฺธนฺตสฺส ‘ยทา ปฏิปาฏิํ ลภิสฺสามิ, ตทา ภตฺตํ กริสฺสามี’ติ, น จ เม ปฏิปาฏิ ลพฺภติ, อหญฺจมฺหิ เอกตฺตโก\csromansharedfootnote{b3e16a44b48ab736}{เอกโก (สี, สฺยา)}, พหุ จ เม ฆราวาสตฺโถ หายติ, ยนฺนูนาหํ ภตฺตคฺคํ โอโลเกยฺยํ, ยํ ภตฺตคฺเค นาสฺส, ตํ ปฏิยาเทยฺยนฺ”ติ. อถ โข โส พฺราหฺมโณ ภตฺตคฺคํ โอโลเกนฺโต เทฺว นาทฺทส \csromanfolio{315}ยาคุญฺจ มธุโคฬกญฺจ. อถ โข โส พฺราหฺมโณ เยนายสฺมา อานนฺโท เตนุปสงฺกมิ, อุปสงฺกมิตฺวา อายสฺมนฺตํ อานนฺทํ เอตทโวจ “อิธ เม โภ อานนฺท ปฏิปาฏิํ อลภนฺตสฺส เอตทโหสิ ‘อตีตานิ โข เม เทฺว มาสานิ พุทฺธ\-ปฺปมุขํ ภิกฺขุสํฆํ อนุพนฺธนฺตสฺส ‘ยทา ปฏิปาฏิํ ลภิสฺสามิ, ตทา ภตฺตํ กริสฺสามี’ติ, น จ เม ปฏิปาฏิ ลพฺภติ, อหญฺจมฺหิ เอกตฺตโก, พหุ จ เม ฆราวาสตฺโถ หายติ, ยนฺนูนาหํ ภตฺตคฺคํ โอโลเกยฺยํ, ยํ ภตฺตคฺเค นาสฺส, ตํ ปฏิยาเทยฺยนฺ’ติ. โส โข อหํ โภ อานนฺท ภตฺตคฺคํ โอโลเกนฺโต เทฺว นาทฺทสํ ยาคุญฺจ มธุโคฬกญฺจ. สจาหํ โภ อานนฺท ปฏิยาเทยฺยํ ยาคุญฺจ มธุโคฬกญฺจ, ปฏิคฺคเณฺหยฺย เม ภวํ โคตโม”ติ. เตน หิ พฺราหฺมณ ภควนฺตํ ปฏิปุจฺฉิสฺสามีติ. อถ โข อายสฺมา อานนฺโท ภควโต เอตมตฺถํ อาโรเจสิ. เตน หานนฺท ปฏิยาเทตูติ. เตน หิ พฺราหฺมณ ปฏิยาเทหีติ. อถ โข โส พฺราหฺมโณ ตสฺสา รตฺติยา อจฺจเยน ปหูตํ ยาคุญฺจ มธุโคฬกญฺจ \textbf{ปฏิยาทาเปตฺวา ภควโต อุปนาเมสิ “ปฏิคฺคณฺหาตุ เม ภวํ โคตโม} ยาคุญฺจ มธุโคฬกญฺจา”ติ. เตน หิ พฺราหฺมณ ภิกฺขูนํ เทหีติ. ภิกฺขู กุกฺกุจฺจายนฺตา น ปฏิคฺคณฺหนฺติ. ปฏิคฺคณฺหถ ภิกฺขเว ปริภุญฺช\-ถาติ. อถ โข โส พฺราหฺมโณ พุทฺธ\-ปฺปมุขํ ภิกฺขุสํฆํ ปหูตาย ยาคุยา จ มธุโคฬเกน จ สหตฺถา สนฺตปฺเปตฺวา สมฺปวาเรตฺวา ภควนฺตํ โธตหตฺถํ โอนีต\-ปตฺต\-ปาณิํ เอกมนฺตํ นิสีทิ. เอกมนฺตํ นิสินฺนํ โข ตํ พฺราหฺมณํ ภควา เอตทโวจ\csromandash{}}

\prose{“ทสยิเม พฺราหฺมณ อานิสํสา ยาคุยา. กตเม ทส, ยาคุํ เทนฺโต อายุํ เทติ, วณฺณํ เทติ, สุขํ เทติ, พลํ เทติ, ปฏิภานํ เทติ, ยาคุ ปีตา ขุทฺทํ\csromansharedfootnote{e7bfcdfe0913e277}{ขุทํ (สี, สฺยา)} ปฏิหนติ, ปิปาสํ วิเนติ, วาตํ อนุโลเมติ, วตฺถิํ โสเธติ, อามาวเสสํ ปาเจติ, อิเม โข พฺราหฺมณ ทสานิสํสา ยาคุยา”ติ *.}

\setlength{\csromangathapagewidth}{0pt}
\setlength{\csromangathawakwidth}{0pt}
\csromangathameasuregroup{}{\textsuperscript{+}โย สญฺญตานํ ปรทตฺตโภชินํ, \\ กาเลน สกฺกจฺจ ททาติ ยาคุํ. \\ ทสสฺส ฐานานิ อนุปฺปเวจฺฉติ, \\ อายุญฺช วณฺณญฺจ สุขํ พลญฺจ. \\ ปฏิภานมสฺส อุปชายเต ตโต, \\ ขุทฺทํ ปิปาสญฺจ พฺยปเนติ วาตํ. \\ โสเธติ วตฺถิํ ปริณาเมติ ภตฺตํ, \\ เภสชฺชเมตํ สุคเตน วณฺณิตํ. \\ ตสฺมา หิ ยาคุํ อลเมว ทาตุํ, \\ นิจฺจํ มนุสฺเสน สุขตฺถิเกน. \\ ทิพฺพานิ วา ปตฺถยตา สุขานิ, \\ มนุสฺสโสภคฺยตมิจฺฉตา วาติ.}
\csromangathameasurewak{\textsuperscript{+}โย สญฺญตานํ ปรทตฺตโภชินํ, \\ กาเลน สกฺกจฺจ ททาติ ยาคุํ. \\ ทสสฺส ฐานานิ อนุปฺปเวจฺฉติ, \\ อายุญฺช วณฺณญฺจ สุขํ พลญฺจ. \\ ปฏิภานมสฺส อุปชายเต ตโต, \\ ขุทฺทํ ปิปาสญฺจ พฺยปเนติ วาตํ. \\ โสเธติ วตฺถิํ ปริณาเมติ ภตฺตํ, \\ เภสชฺชเมตํ สุคเตน วณฺณิตํ. \\ ตสฺมา หิ ยาคุํ อลเมว ทาตุํ, \\ นิจฺจํ มนุสฺเสน สุขตฺถิเกน. \\ ทิพฺพานิ วา ปตฺถยตา สุขานิ, \\ มนุสฺสโสภคฺยตมิจฺฉตา วาติ.}
\setlength{\csromangathaleftcol}{0pt}
\csromangathagroup{\csromangathastanza{\csromangathawak{\csromansymbolfootnote{+}{อํ 1. 375 ปิฏฺเฐ โถกํ วิสทิสํ.}โย สญฺญตานํ ปรทตฺตโภชินํ, \\ กาเลน สกฺกจฺจ ททาติ ยาคุํ. \\ ทสสฺส ฐานานิ อนุปฺปเวจฺฉติ, \\ อายุญฺช วณฺณญฺจ สุขํ พลญฺจ.}}\csromangathastanza{\csromangathawak{\csromanfolio{316}ปฏิภานมสฺส อุปชายเต ตโต, \\ ขุทฺทํ ปิปาสญฺจ พฺยปเนติ วาตํ. \\ โสเธติ วตฺถิํ ปริณาเมติ ภตฺตํ, \\ เภสชฺชเมตํ สุคเตน วณฺณิตํ.}}\csromangathastanza{\csromangathawak{ตสฺมา หิ ยาคุํ อลเมว ทาตุํ, \\ นิจฺจํ มนุสฺเสน สุขตฺถิเกน. \\ ทิพฺพานิ วา ปตฺถยตา สุขานิ, \\ มนุสฺสโสภคฺยตมิจฺฉตา วาติ.}}}

\prose{อถ โข ภควา ตํ พฺราหฺมณํ อิมาหิ คาถาหิ อนุโมทิตฺวา อุฏฺฐายาสนา ปกฺกามิ. อถ โข ภควา เอตสฺมิํ นิทาเน เอตสฺมิํ ปกรเณ ธมฺมิํ กถํ กตฺวา ภิกฺขู อามนฺเตสิ “อนุชานามิ ภิกฺขเว ยาคุญฺจ มธุโคฬกญฺจา”ติ.}

\csromanmatikaanchor{mtk.177}
\csromantocmarkref{section}{171. ตรุณปสนฺนมหามตฺตวตฺถุ}{mtk.177}
\bookmark[dest={mtk.177},level=1]{171. ตรุณปสนฺนมหามตฺตวตฺถุ}
\csromanheader{1}{171. \textbf{ตรุณปสนฺน}\-\textbf{มหามตฺต}\-\textbf{วตฺถุ}}
\proseitem{283}{อสฺโสสุํ โข มนุสฺสา ภควตา กิร ยาคุ อนุญฺญาตา มธุโคฬก\-ญฺจาติ, เต กาลสฺเสว โภชฺชยาคุํ ปฏิยาเทนฺติ มธุโคฬกญฺจ. ภิกฺขู กาลสฺเสว โภชฺชยาคุยา ธาตา มธุโคฬเกน จ ภตฺตคฺเค น จิตฺตรูปํ ปริภุญฺชนฺติ. เตน โข ปน สมเยน อญฺญตเรน ตรุณ\-ปสนฺเนน มหามตฺเตน สฺวาตนาย พุทฺธปฺปมุโข ภิกฺขุสํโฆ นิมนฺติโต โหติ. อถ โข ตสฺส ตรุณ\-ปสนฺนสฺส มหามตฺตสฺส เอตทโหสิ “ยนฺนูนาหํ อฑฺฒ\-เตลสนฺนํ ภิกฺขุสตานํ อฑฺฒเตลสานิ มํสปาติ\-สตานิ ปฏิยาเทยฺยํ, เอกเมกสฺส ภิกฺขุโน เอกเมกํ มํสปาติํ อุปนาเมยฺยนฺ”ติ. อถ โข โส ตรุณปสนฺโน มหามตฺโต ตสฺสา รตฺติยา อจฺจเยน ปณีตํ ขาทนียํ โภชนียํ ปฏิยาทาเปตฺวา อฑฺฒเตลสานิ จ มํสปาติ\-สตานิ ภควโต กาลํ อาโรจาเปสิ “กาโล ภนฺเต นิฏฺฐิตํ ภตฺตนฺ”ติ. อถ โข ภควา ปุพฺพณฺหสมยํ นิวาเสตฺวา ปตฺต\-จีวรมา\-ทาย เยน ตสฺส ตรุณ\-ปสนฺนสฺส มหามตฺตสฺส นิเวสนํ เตนุปสงฺกมิ, อุปสงฺกมิตฺวา ปญฺญตฺเต อาสเน นิสีทิ สทฺธิํ ภิกฺขุ\-สํเฆน. อถ โข โส ตรุณปสนฺโน มหามตฺโต ภตฺตคฺเค ภิกฺขู ปริวิสติ. ภิกฺขู เอวมาหํสุ “โถกํ อาวุโส เทหิ, โถกํ อาวุโส เทหี”ติ. มา โข ตุเมฺห ภนฺเต “อยํ ตรุณปสนฺโน \csromanfolio{317}มหามตฺโต”ติ โถกํ โถกํ ปฏิคฺคณฺหถ, พหุํ เม ขาทนียํ โภชนียํ ปฏิยตฺตํ อฑฺฒเตลสานิ จ มํสปาติ\-สตานิ “เอกเมกสฺส ภิกฺขุโน เอกเมกํ มํสปาติํ อุปนาเมสฺสามี”ติ, ปฏิคฺคณฺหถ ภนฺเต ยาวทตฺถนฺติ. น โข มยํ อาวุโส เอตํการณา โถกํ โถกํ ปฏิคฺคณฺหาม, อปิ จ มยํ กาลสฺเสว โภชฺชยาคุยา ธาตา มธุโคฬเกน จ, เตน มยํ โถกํ โถกํ ปฏิคฺคณฺหามาติ. อถ โข โส ตรุณปสนฺโน มหามตฺโต อุชฺฌายติ ขิยฺยติ วิปาเจติ “กถํ หิ นาม ภทนฺตา มยา นิมนฺติตา อญฺญสฺส โภชฺชยาคุํ ปริภุญฺชิสฺสนฺติ, น จาหํ ปฏิพโล ยาวทตฺถํ ทาตุนฺ”ติ, กุปิโต อนตฺตมโน อาสาทนาเปกฺโข ภิกฺขูนํ ปตฺเต ปูเรนฺโต อคมาสิ “ภุญฺชถ วา หรถ วา”ติ. อถ โข โส ตรุณปสนฺโน มหามตฺโต พุทฺธ\-ปฺปมุขํ ภิกฺขุสํฆํ ปณีเตน ขาทนีเยน โภชนีเยน สหตฺถา สนฺตปฺเปตฺวา สมฺปวาเรตฺวา ภควนฺตํ ภุตฺตาวิํ โอนีต\-ปตฺต\-ปาณิํ เอกมนฺตํ นิสีทิ. เอกมนฺตํ นิสินฺนํ โข ตํ ตรุณปสนฺนํ มหามตฺตํ ภควา ธมฺมิยา กถาย สนฺทสฺเสตฺวา สมาทเปตฺวา สมุตฺเตเชตฺวา สมฺปหํเสตฺวา อุฏฺฐายาสนา ปกฺกามิ.}

\prose{อถ โข ตสฺส ตรุณ\-ปสนฺนสฺส มหามตฺตสฺส อจิร\-ปกฺกนฺตสฺส ภควโต อหุเทว กุกฺกุจฺจํ, อหุ วิปฺปฏิสาโร, “อลาภา วต เม, น วต เม ลาภา, ทุลฺลทฺธํ วต เม, น วต เม สุลทฺธํ, โยหํ กุปิโต อนตฺตมโน อาสาทนาเปกฺโข ภิกฺขูนํ ปตฺเต ปูเรนฺโต อคมาสิํ ‘ภุญฺชถ วา หรถ วา’ติ, กิํ นุ โข มยา พหุํ ปสุตํ ปุญฺญํ วา อปุญฺญํ วา”ติ. อถ โข โส ตรุณปสนฺโน มหามตฺโต เยน ภควา เตนุปสงฺกมิ, อุปสงฺกมิตฺวา ภควนฺตํ อภิวาเทตฺวา เอกมนฺตํ นิสีทิ, เอกมนฺตํ นิสินฺโน โข โส ตรุณปสนฺโน มหามตฺโต ภควนฺตํ เอตทโวจ “อิธ มยฺหํ ภนฺเต อจิร\-ปกฺกนฺตสฺส ภควโต อหุเทว กุกฺกุจฺจํ, อหุ วิปฺปฏิสาโร, ‘อลาภา วต เม, น วต เม ลาภา, ทุลฺลทฺธํ วต เม, น วต เม สุลทฺธํ, โยหํ กุปิโต อนตฺตมโน อาสาทนาเปกฺโข ภิกฺขูนํ ปตฺเต ปูเรนฺโต อคมาสิํ ภุญฺชถ วา หรถ วาติ, กิํ นุ โข มยา พหุํ ปสุตํ ปุญฺญํ วา อปุญฺญํ วา’ติ. กิํ นุ โข มยา ภนฺเต พหุํ ปสุตํ ปุญฺญํ วา อปุญฺญํ วา”ติ. ยทคฺเคน ตยา อาวุโส สฺวาตนาย พุทฺธปฺปมุโข ภิกฺขุสํโฆ นิมนฺติโต, ตทคฺเคน เต พหุํ ปุญฺญํ ปสุตํ. ยทคฺเคน \csromanfolio{318}เต เอกเมเกน ภิกฺขุนา เอกเมกํ สิตฺถํ ปฏิคฺคหิตํ, ตทคฺเคน เต พหุํ ปุญฺญํ ปสุตํ, สคฺคา เต อารทฺธาติ. อถ โข โส ตรุณปสนฺโน มหามตฺโต “ลาภา กิร เม, สุลทฺธํ กิร เม, พหุํ กิร มยา ปุญฺญํ ปสุตํ, สคฺคา กิร เม อารทฺธา”ติ หฏฺโฐ อุทคฺโค อุฏฺฐายาสนา ภควนฺตํ อภิวาเทตฺวา ปทกฺขิณํ กตฺวา ปกฺกามิ. อถ โข ภควา เอตสฺมิํ นิทาเน เอตสฺมิํ ปกรเณ ภิกฺขุสํฆํ สนฺนิปาตาเปตฺวา ภิกฺขู ปฏิปุจฺฉิ “สจฺจํ กิร ภิกฺขเว ภิกฺขู อญฺญตฺร นิมนฺติตา อญฺญสฺส โภชฺชยาคุํ ปริภุญฺชนฺตี”ติ. สจฺจํ ภควาติ. วิครหิ พุทฺโธ ภควา ฯเปฯ. กถํ หิ นาม เต ภิกฺขเว โมฆปุริสา อญฺญตฺร นิมนฺติตา อญฺญสฺส โภชฺชยาคุํ ปริภุญฺชิสฺสนฺติ. เนตํ ภิกฺขเว อปฺปสนฺนานํ วา ปสาทาย ฯเปฯ วิครหิตฺวา ธมฺมิํ กถํ กตฺวา ภิกฺขู อามนฺเตสิ “น ภิกฺขเว อญฺญตฺร นิมนฺติเตน อญฺญสฺส โภชฺชยาคุ ปริภุญฺชิ\-ตพฺพา, โย ปริภุญฺเชยฺย, ยถาธมฺโม กาเรตพฺโพ”ติ.}

\csromanmatikaanchor{mtk.178}
\csromantocmarkref{section}{172. เพลฏฺฐกจฺจานวตฺถุ}{mtk.178}
\bookmark[dest={mtk.178},level=1]{172. เพลฏฺฐกจฺจานวตฺถุ}
\csromanheader{1}{172. \textbf{เพลฏฺฐ}\-\textbf{กจฺจาน}\-\textbf{วตฺถุ}}
\proseitem{284}{อถ โข ภควา อนฺธกวินฺเท ยถาภิรนฺตํ วิหริตฺวา เยน ราชคหํ, เตน จาริกํ ปกฺกามิ มหตา ภิกฺขุ\-สํเฆน สทฺธิํ อฑฺฒเตลเสหิ ภิกฺขุสเตหิ. เตน โข ปน สมเยน เพลฏฺโฐ กจฺจาโน ราชคหา อนฺธกวินฺทํ อทฺธาน\-มคฺค\-ปฺปฏิปนฺโน โหติ ปญฺจมตฺเตหิ สกฏสเตหิ สพฺเพเหว คุฬกุมฺภ\-ปูเรหิ. อทฺทสา โข ภควา เพลฏฺฐํ กจฺจานํ ทูรโตว อาคจฺฉนฺตํ, ทิสฺวาน มคฺคา โอกฺกมฺม อญฺญตรสฺมิํ รุกฺขมูเล นิสีทิ. อถ โข เพลฏฺโฐ กจฺจาโน เยน ภควา เตนุปสงฺกมิ, อุปสงฺกมิตฺวา ภควนฺตํ อภิวาเทตฺวา เอกมนฺตํ อฏฺฐาสิ, เอกมนฺตํ ฐิโต โข เพลฏฺโฐ กจฺจาโน ภควนฺตํ เอตทโวจ “อิจฺฉามหํ ภนฺเต เอกเมกสฺส ภิกฺขุโน เอกเมกํ คุฬกุมฺภํ ทาตุนฺ”ติ. เตน หิ ตฺวํ กจฺจาน เอกํเยว คุฬกุมฺภํ อาหราติ. “เอวํ ภนฺเต”ติ โข เพลฏฺโฐ กจฺจาโน ภควโต ปฏิสฺสุณิตฺวา เอกํเยว คุฬกุมฺภํ อาทาย เยน ภควา เตนุปสงฺกมิ, อุปสงฺกมิตฺวา ภควนฺตํ เอตทโวจ “อาภโต\csromansharedfootnote{74791bc6d2813ab6}{อาหโฏ (สี, สฺยา, ก)} ภนฺเต คุฬกุมฺโภ, กถาหํ ภนฺเต ปฏิปชฺชามี”ติ. เตน หิ ตฺวํ กจฺจาน ภิกฺขูนํ คุฬํ เทหีติ. “เอวํ ภนฺเต”ติ โข เพลฏฺโฐ กจฺจาโน ภควโต \csromanfolio{319}ปฏิสฺสุณิตฺวา ภิกฺขูนํ คุฬํ ทตฺวา ภควนฺตํ เอตทโวจ “ทินฺโน ภนฺเต ภิกฺขูนํ คุโฬ, พหุ จายํ คุโฬ อวสิฏฺโฐ, กถาหํ ภนฺเต ปฏิปชฺชามี”ติ. เตน หิ ตฺวํ กจฺจาน ภิกฺขูนํ คุฬํ ยาวทตฺถํ เทหีติ. “เอวํ ภนฺเต”ติ โข เพลฏฺโฐ กจฺจาโน ภควโต ปฏิสฺสุณิตฺวา ภิกฺขูนํ คุฬํ ยาวทตฺถํ ทตฺวา ภควนฺตํ เอตทโวจ “ทินฺโน ภนฺเต ภิกฺขูนํ คุโฬ ยาวทตฺโถ, พหุ จายํ คุโฬ อวสิฏฺโฐ, กถาหํ ภนฺเต ปฏิปชฺชามี”ติ. เตน หิ ตฺวํ กจฺจาน ภิกฺขู คุเฬหิ สนฺตปฺเปหีติ. “เอวํ ภนฺเต”ติ โข เพลฏฺโฐ กจฺจาโน ภควโต ปฏิสฺสุณิตฺวา ภิกฺขู คุเฬหิ สนฺตปฺเปสิ. เอกจฺเจ ภิกฺขู ปตฺเตปิ ปูเรสุํ, ปริสฺสาว\-นานิปิ ถวิกาโยปิ ปูเรสุํ. อถ โข เพลฏฺโฐ กจฺจาโน ภิกฺขู คุเฬหิ สนฺตปฺเปตฺวา ภควนฺตํ เอตทโวจ “สนฺตปฺปิตา ภนฺเต ภิกฺขู คุเฬหิ, พหุ จายํ คุโฬ อวสิฏฺโฐ, กถาหํ ภนฺเต ปฏิปชฺชามี”ติ. เตน หิ ตฺวํ กจฺจาน วิฆาสาทานํ คุฬํ เทหีติ. “เอวํ ภนฺเต”ติ โข เพลฏฺโฐ กจฺจาโน ภควโต ปฏิสฺสุณิตฺวา วิฆาสาทานํ คุฬํ ทตฺวา ภควนฺตํ เอตทโวจ “ทินฺโน ภนฺเต วิฆาสาทานํ คุโฬ, พหุ จายํ คุโฬ อวสิฏฺโฐ, กถาหํ ภนฺเต ปฏิปชฺชามี”ติ. เตน หิ ตฺวํ กจฺจาน วิฆาสาทานํ คุฬํ ยาวทตฺถํ เทหีติ. “เอวํ ภนฺเต”ติ โข เพลฏฺโฐ กจฺจาโน ภควโต ปฏิสฺสุณิตฺวา วิฆาสาทานํ คุฬํ ยาวทตฺถํ ทตฺวา ภควนฺตํ เอตทโวจ “ทินฺโน ภนฺเต วิฆาสาทานํ คุโฬ ยาวทตฺโถ, พหุ จายํ คุโฬ อวสิฏฺโฐ, กถาหํ ภนฺเต ปฏิปชฺชามี”ติ. เตน หิ ตฺวํ กจฺจาน วิฆาสาเท คุเฬหิ สนฺตปฺเปหีติ. “เอวํ ภนฺเต”ติ โข เพลฏฺโฐ กจฺจาโน ภควโต ปฏิสฺสุณิตฺวา วิฆาสาเท คุเฬหิ สนฺตปฺเปสิ. เอกจฺเจ วิฆาสาทา โกลมฺเพปิ ฆเฏปิ ปูเรสุํ, ปิฏกานิปิ อุจฺฉงฺเคปิ ปูเรสุํ. อถ โข เพลฏฺโฐ กจฺจาโน วิฆาสาเท คุเฬหิ สนฺตปฺเปตฺวา ภควนฺตํ เอตทโวจ “สนฺตปฺปิตา ภนฺเต วิฆาสาทา คุเฬหิ, พหุ จายํ คุโฬ อวสิฏฺโฐ, กถาหํ ภนฺเต ปฏิปชฺชามี”ติ. นาหํ ตํ กจฺจาน ปสฺสามิ สเทวเก โลเก สมารเก สพฺรหฺมเก สสฺสมณ\-พฺราหฺมณิยา ปชาย สเทว\-มนุสฺสาย ยสฺส โส คุโฬ ปริภุตฺโต สมฺมา ปริณามํ คจฺเฉยฺย อญฺญตฺร ตถาคตสฺส วา ตถาคต\-สาวกสฺส วา. เตน หิ ตฺวํ กจฺจาน ตํ คุฬํ อปฺปหริเต วา ฉฑฺเฑหิ, อปฺปาณเก วา อุทเก โอปิลาเปหีติ. “เอวํ ภนฺเต”ติ โข เพลฏฺโฐ กจฺจาโน ภควโต ปฏิสฺสุณิตฺวา ตํ คุฬํ อปฺปาณเก \csromanfolio{320}อุทเก โอปิลาเปติ. อถ โข โส คุโฬ อุทเก ปกฺขิตฺโต จิจฺจิฏายติ จิฏิจิฏายติ ปธูปายติ\csromansharedfootnote{9d980332cf97d3cf}{สนฺธูปายติ (สี, สฺยา)} สมฺปธูปายติ. เสยฺยถาปิ นาม ผาโล ทิวสํ\-สนฺตตฺโต อุทเก ปกฺขิตฺโต จิจฺจิฏายติ จิฏิจิฏายติ ปธูปายติ สมฺปธูปายติ, เอวเมว โส คุโฬ อุทเก ปกฺขิตฺโต จิจฺจิฏายติ จิฏิจิฏายติ ปธูปายติ สมฺปธูปายติ.}

\prose{อถ โข เพลฏฺโฐ กจฺจาโน สํวิคฺโค โลมหฏฺฐชาโต เยน ภควา เตนุปสงฺกมิ, อุปสงฺกมิตฺวา ภควนฺตํ อภิวาเทตฺวา เอกมนฺตํ นิสีทิ. เอกมนฺตํ นิสินฺนสฺส โข เพลฏฺฐสฺส กจฺจานสฺส ภควา อนุปุพฺพิํ กถํ กเถสิ, เสยฺยถิทํ, ทานกถํ สีลกถํ สคฺคกถํ กามานํ อาทีนวํ โอการํ สํกิเลสํ เนกฺขมฺเม อานิสํสํ ปกาเสสิ. ยทา ภควา อญฺญาสิ เพลฏฺฐํ กจฺจานํ กลฺลจิตฺตํ มุทุจิตฺตํ วินีวรณ\-จิตฺตํ อุทคฺคจิตฺตํ ปสนฺนจิตฺตํ, อถ ยา พุทฺธานํ สามุกฺกํสิกา ธมฺมเทสนา, ตํ ปกาเสสิ ฯเปฯ เอวเมว เพลฏฺฐสฺส กจฺจานสฺส ตสฺมิํเยว อาสเน วิรชํ วีตมลํ ธมฺมจกฺขุํ อุทปาทิ “ยํ กิญฺจิ สมุทย\-ธมฺมํ สพฺพํ ตํ นิโรธธมฺมนฺ”ติ.}

\prose{อถ โข เพลฏฺโฐ กจฺจาโน ทิฏฺฐธมฺโม ปตฺตธมฺโม วิทิตธมฺโม ปริโยคาฬฺห\-ธมฺโม ติณฺณ\-วิจิกิจฺโฉ วิคต\-กถํกโถ เวสารชฺช\-ปฺปตฺโต อปรปฺปจฺจโย สตฺถุสาสเน ภควนฺตํ เอตทโวจ “อภิกฺกนฺตํ ภนฺเต, อภิกฺกนฺตํ ภนฺเต, เสยฺยถาปิ ภนฺเต นิกุชฺชิตํ วา อุกฺกุชฺเชยฺย ฯเปฯ เอวเมวํ โข ภควตา อเนก\-ปริยาเยน ธมฺโม ปกาสิโต.}

\prose{เอสาหํ ภนฺเต ภควนฺตํ สรณํ คจฺฉามิ, ธมฺมญฺจ ภิกฺขุสํฆญฺจ, อุปาสกํ มํ ภควา ธาเรตุ อชฺชตคฺเค ปาณุเปตํ สรณคตนฺ”ติ.}

\prose{อถ โข ภควา อนุปุพฺเพน จาริกํ จรมาโน เยน ราชคหํ ตทวสริ, ตตฺร สุทํ ภควา ราชคเห วิหรติ เวฬุวเน กลนฺทก\-นิวาเป.}

\prose{เตน โข ปน สมเยน ราชคเห คุโฬ อุสฺสนฺโน โหติ. ภิกฺขู “คิลานสฺเสว ภควตา คุโฬ อนุญฺญาโต, โน อคิลานสฺสา”ติ กุกฺกุจฺจายนฺตา คุฬํ น ภุญฺชนฺติ. ภควโต เอตมตฺถํ อาโรเจสุํ. อนุชานามิ ภิกฺขเว คิลานสฺส คุฬํ, อคิลานสฺส คุโฬทกนฺติ.}

\csromanmatikaanchor{mtk.179}
\csromantocmarkref{section}{173. ปาฏลิคามวตฺถุ}{mtk.179}
\bookmark[dest={mtk.179},level=1]{173. ปาฏลิคามวตฺถุ}
\csromanheader{1}{173. \textbf{ปาฏลิคาม}\-\textbf{วตฺถุ}}
\proseitem{285}{\csromanfolio{321}\csromansymbolfootnote{*}{อิโต ปรํ ยาว 325 ปิฏฺเฐ ‘ติณฺณา เมธาวิโน ชนา’ติ ปาโฐ ที 2. 71 ปิฏฺฐาทีสุ จ ขุ 1. 183 ปิฏฺฐาทีสุปิ จ อาคโต.}อถ โข ภควา ราชคเห ยถาภิรนฺตํ วิหริตฺวา เยน ปาฏลิคาโม เตน จาริกํ ปกฺกามิ มหตา ภิกฺขุ\-สํเฆน สทฺธิํ อฑฺฒเตลเสหิ ภิกฺขูสเตหิ. อถ โข ภควา อนุปุพฺเพน จาริกํ จรมาโน เยน ปาฏลิคาโม ตทวสริ. อสฺโสสุํ โข ปาฏลิคามิกา อุปาสกา “ภควา กิร ปาฏลิคามํ อนุปฺปตฺโต”ติ. อถ โข ปาฏลิคามิกา อุปาสกา เยน ภควา เตนุ\-ปสงฺกมิํสุ, อุปสงฺกมิตฺวา ภควนฺตํ อภิวาเทตฺวา เอกมนฺตํ นิสีทิํสุ. เอกมนฺตํ นิสินฺเน โข ปาฏลิคามิเก อุปาสเก ภควา ธมฺมิยา กถาย สนฺทสฺเสสิ สมาทเปสิ สมุตฺเตเชสิ สมฺปหํเสสิ. อถ โข ปาฏลิคามิกา อุปาสกา ภควตา ธมฺมิยา กถาย สนฺทสฺสิตา สมาทปิตา สมุตฺเตชิตา สมฺปหํสิตา ภควนฺตํ เอตทโวจุํ “อธิวาเสตุ โน ภนฺเต ภควา อาวสถาคารํ สทฺธิํ ภิกฺขุ\-สํเฆนา”ติ. อธิวาเสสิ ภควา ตุณฺหีภาเวน. อถ โข ปาฏลิคามิกา อุปาสกา ภควโต อธิวาสนํ วิทิตฺวา อุฏฺฐายาสนา ภควนฺตํ อภิวาเทตฺวา ปทกฺขิณํ กตฺวา เยน อาวสถาคารํ เตนุ\-ปสงฺกมิํสุ, อุปสงฺกมิตฺวา สพฺพสนฺถริํ อาวสถาคารํ สนฺถริตฺวา อาสนานิ ปญฺญเปตฺวา อุทกมณิกํ ปติฏฺฐาเปตฺวา เตลปทีปํ อาโรเปตฺวา เยน ภควา เตนุ\-ปสงฺกมิํสุ, อุปสงฺกมิตฺวา ภควนฺตํ อภิวาเทตฺวา เอกมนฺตํ อฏฺฐํสุ, เอกมนฺตํ ฐิตา โข ปาฏลิคามิกา อุปาสกา ภควนฺตํ เอตทโวจุํ “สพฺพสนฺถริ\-สนฺถตํ ภนฺเต อาวสถาคารํ, อาสนานิ ปญฺญตฺตานิ, อุทกมณิโก ปติฏฺฐาปิโต, เตลปทีโป อาโรปิโต, ยสฺสทานิ ภนฺเต ภควา กาลํ มญฺญตี”ติ.}

\prose{อถ โข ภควา นิวาเสตฺวา ปตฺต\-จีวรมา\-ทาย สทฺธิํ ภิกฺขุ\-สํเฆน เยน อาวสถาคารํ เตนุปสงฺกมิ, อุปสงฺกมิตฺวา ปาเท ปกฺขาเลตฺวา อาวสถาคารํ ปวิสิตฺวา มชฺฌิมํ ถมฺภํ นิสฺสาย ปุรตฺถา\-ภิมุโข นิสีทิ. ภิกฺขุ\-สํโฆปิ โข ปาเท ปกฺขาเลตฺวา อาวสถาคารํ ปวิสิตฺวา ปจฺฉิมํ ภิตฺติํ นิสฺสาย ปุรตฺถา\-ภิมุโข นิสีทิ ภควนฺตํเยว ปุรกฺขตฺวา. ปาฏลิคามิกาปิ โข อุปาสกา ปาเท ปกฺขาเลตฺวา อาวสถาคารํ ปวิสิตฺวา ปุรตฺถิมํ ภิตฺติํ นิสฺสาย ปจฺฉิมา\-ภิมุขา นิสีทิํสุ ภควนฺตํเยว ปุรกฺขตฺวา. อถ โข ภควา ปาฏลิคามิเก อุปาสเก อามนฺเตสิ\csromandash{}}

\prose{\csromanfolio{322}\csromansymbolfootnote{*}{ที 3. 196; อํ 2. 221 ปิฏฺเฐสุปิ.}ปญฺจิเม คหปตโย อาทีนวา ทุสฺสีลสฺส สีลวิปตฺติยา. กตเม ปญฺจ, อิธ คหปตโย ทุสฺสีโล สีลวิปนฺโน ปมาทาธิกรณํ มหติํ โภคชานิํ นิคจฺฉติ, อยํ ปฐโม อาทีนโว ทุสฺสีลสฺส สีลวิปตฺติยา. ปุน จปรํ คหปตโย ทุสฺสีลสฺส สีลวิปนฺนสฺส ปาปโก กิตฺติสทฺโท อพฺภุคฺคจฺฉติ, อยํ ทุติโย อาทีนโว ทุสฺสีลสฺส สีลวิปตฺติยา. ปุน จปรํ คหปตโย ทุสฺสีโล สีลวิปนฺโน ยญฺญเทว ปริสํ อุปสงฺกมติ ยทิ ขตฺติย\-ปริสํ ยทิ พฺราหฺมณ\-ปริสํ ยทิ คหปติ\-ปริสํ ยทิ สมณปริสํ, อวิสารโท อุปสงฺกมติ มงฺกุภูโต, อยํ ตติโย อาทีนโว ทุสฺสีลสฺส สีลวิปตฺติยา. ปุน จปรํ คหปตโย ทุสฺสีโล สีลวิปนฺโน สมฺมูโฬฺห กาลํ กโรติ, อยํ จตุตฺโถ อาทีนโว ทุสฺสีลสฺส สีลวิปตฺติยา. ปุน จปรํ คหปตโย ทุสฺสีโล สีลวิปนฺโน กายสฺส เภทา ปรํ มรณา อปายํ ทุคฺคติํ วินิปาตํ นิรยํ อุปปชฺชติ, อยํ ปญฺจโม อาทีนโว ทุสฺสีลสฺส สีลวิปตฺติยา. อิเม โข คหปตโย ปญฺจ อาทีนวา ทุสฺสีลสฺส สีลวิปตฺติยา.}

\prose{\csromansymbolfootnote{+}{ที 3. 197; อํ 2. 221 ปิฏฺเฐสุปิ.}ปญฺจิเม คหปตโย อานิสํสา สีลวโต สีลสมฺปทาย. กตเม ปญฺจ, อิธ คหปตโย สีลวา สีลสมฺปนฺโน อปฺปมาทาธิกรณํ มหนฺตํ โภคกฺขนฺธํ อธิคจฺฉติ, อยํ ปฐโม อานิสํโส สีลวโต สีลสมฺปทาย. ปุน จปรํ คหปตโย สีลวโต สีลสมฺปนฺนสฺส กลฺยาโณ กิตฺติสทฺโท อพฺภุคฺคจฺฉติ, อยํ ทุติโย อานิสํโส สีลวโต ลสมฺปทาย. ปุน จปรํ คหปตโย สีลวา สีลสมฺปนฺโน ยญฺญเทว ปริสํ อุปสงฺกมติ ยทิ ขตฺติย\-ปริสํ ยทิ พฺราหฺมณ\-ปริสํ ยทิ คหปติ\-ปริสํ ยทิ สมณปริสํ, วิสารโท อุปสงฺกมติ อมงฺกุภูโต, อยํ ตติโย อานิสํโส สีลวโต สีลสมฺปทาย. ปุน จปรํ คหปตโย สีลวา สีลสมฺปนฺโน อสมฺมูโฬฺห กาลํ กโรติ, อยํ จตุตฺโถ อานิสํโส สีลวโต สีลสมฺปทาย. ปุน จปรํ คหปตโย สีลวา สีลสมฺปนฺโน กายสฺส เภทา ปรํ มรณา สุคติํ สคฺคํ โลกํ อุปปชฺชติ, อยํ ปญฺจโม อานิสํโส สีลวโต สีลสมฺปทาย, อิเม โข คหปตโย ปญฺจ อานิสํสา สีลวโต สีลสมฺปทายาติ.}

\proseitem{285}{อถ โข ภควา ปาฏลิคามิเก อุปาสเก พหุเทว รตฺติํ ธมฺมิยา กถาย สนฺทสฺเสตฺวา สมาทเปตฺวา สมุตฺเตเชตฺวา สมฺปหํเสตฺวา \csromanfolio{323}อุยฺโยเชสิ “อภิกฺกนฺตา โข คหปตโย รตฺติ, ยสฺสทานิ ตุเมฺห กาลํ มญฺญถา”ติ. “เอวํ ภนฺเต”ติ โข ปาฏลิคามิกา อุปาสกา ภควโต ปฏิสฺสุณิตฺวา อุฏฺฐายาสนา ภควนฺตํ อภิวาเทตฺวา ปทกฺขิณํ กตฺวา ปกฺกมิํสุ. อถ โข ภควา อจิร\-ปกฺกนฺเตสุ ปาฏลิคามิเกสุ อุปาสเกสุ สุญฺญาคารํ ปาวิสิ.}

\csromanmatikaanchor{mtk.180}
\csromantocmarkref{section}{174. สุนิธวสฺสการวตฺถุ}{mtk.180}
\bookmark[dest={mtk.180},level=1]{174. สุนิธวสฺสการวตฺถุ}
\csromanheader{1}{174. \textbf{สุนิธ}\-\textbf{วสฺสการ}\-\textbf{วตฺถุ}}
\proseitem{286}{เตน โข ปน สมเยน สุนิธ\-วสฺสการา มคธ\-มหามตฺตา ปาฏลิคาเม นครํ มาเปนฺติ วชฺชีนํ ปฏิพาหาย. อทฺทสา โข ภควา รตฺติยา ปจฺจูสสมยํ ปจฺจุฏฺฐาย ทิพฺเพน จกฺขุนา วิสุทฺเธน อติกฺกนฺต\-มานุสเกน สมฺพหุลา เทวตาโย ปาฏลิคาเม วตฺถูนิ ปริคฺคณฺห\-นฺติโย. ยสฺมิํ ปเทเส มเหสกฺขา เทวตา วตฺถูนิ ปริคฺคณฺหนฺติ, มเหสกฺขานํ ตตฺถ ราชูนํ ราช\-มหามตฺตานํ จิตฺตานิ นมนฺติ นิเวสนานิ มาเปตุํ. ยสฺมิํ ปเทเส มชฺฌิมา เทวตา วตฺถูนิ ปริคฺคณฺหนฺติ, มชฺฌิมานํ ตตฺถ ราชูนํ ราช\-มหามตฺตานํ จิตฺตานิ นมนฺติ นิเวสนานิ มาเปตุํ. ยสฺมิํ ปเทเส นีจา เทวตา วตฺถูนิ ปริคฺคณฺหนฺติ, นีจานํ ตตฺถ ราชูนํ ราช\-มหามตฺตานํ จิตฺตานิ นมนฺติ นิเวสนานิ มาเปตุํ. อถ โข ภควา อายสฺมนฺตํ อานนฺทํ อามนฺเตสิ “เก นุ โข เต อานนฺท ปาฏลิคาเม นครํ มาเปนฺตี”ติ. สุนิธ\-วสฺสการา ภนฺเต มคธ\-มหามตฺตา ปาฏลิคาเม นครํ มาเปนฺติ วชฺชีนํ ปฏิพาหายาติ. เสยฺยถาปิ อานนฺท เทเวหิ ตาวติํเสหิ สทฺธิํ มนฺเตตฺวา, เอวเมว โข อานนฺท สุนิธ\-วสฺสการา มคธ\-มหามตฺตา ปาฏลิคาเม นครํ มาเปนฺติ วชฺชีนํ ปฏิพาหาย. อิธาหํ อานนฺท รตฺติยา ปจฺจูสสมยํ ปจฺจุฏฺฐาย อทฺทสํ ทิพฺเพน จกฺขุนา วิสุทฺเธน อติกฺกนฺต\-มานุสเกน สมฺพหุลา เทวตาโย ปาฏลิคาเม วตฺถูนิ ปริคฺคณฺห\-นฺติโย. ยสฺมิํ ปเทเส มเหสกฺขา เทวตา วตฺถูนิ ปริคฺคณฺหนฺติ, มเหสกฺขานํ ตตฺถ ราชูนํ ราช\-มหามตฺตานํ จิตฺตานิ นมนฺติ นิเวสนานิ มาเปตุํ. ยสฺมิํ ปเทเส มชฺฌิมา เทวตา วตฺถูนิ ปริคฺคณฺหนฺติ, มชฺฌิมานํ ตตฺถ ราชูนํ ราช\-มหามตฺตานํ จิตฺตานิ นมนฺติ นิเวสนานิ มาเปตุํ. ยสฺมิํ ปเทเส นีจา เทวตา วตฺถูนิ ปริคฺคณฺหนฺติ, นีจานํ ตตฺถ ราชูนํ ราช\-มหามตฺตานํ จิตฺตานิ นมนฺติ นิเวสนานิ มาเปตุํ. ยาวตา อานนฺท อริยํ อายตนํ ยาวตา วณิปฺปโถ, อิทํ อคฺคนครํ ภวิสฺสติ ปาฏลิปุตฺตํ ปุฏเภทนํ. ปาฏลิปุตฺตสฺส โข อานนฺท ตโย \csromanfolio{324}อนฺตรายา ภวิสฺสนฺติ อคฺคิโต วา อุทกโต วา อพฺภนฺตรโต วา มิถุเภทาติ.}

\prose{อถ โข สุนิธ\-วสฺสการา มคธ\-มหามตฺตา เยน ภควา เตนุ\-ปสงฺกมิํสุ, อุปสงฺกมิตฺวา ภควตา สทฺธิํ สมฺโมทิํสุ, สมฺโมทนียํ กถํ สารณียํ วีติสาเรตฺวา เอกมนฺตํ อฏฺฐํสุ, เอกมนฺตํ ฐิตา โข สุนิธ\-วสฺสการา มคธ\-มหามตฺตา ภควนฺตํ เอตทโวจุํ “อธิวาเสตุ โน ภวํ โคตโม อชฺชตนาย ภตฺตํ สทฺธิํ ภิกฺขุ\-สํเฆนา”ติ. อธิวาเสสิ ภควา ตุณฺหีภาเวน. อถ โข สุนิธ\-วสฺสการา มคธ\-มหามตฺตา ภควโต อธิวาสนํ วิทิตฺวา ปกฺกมิํสุ. อถ โข สุนิธ\-วสฺสการา มคธ\-มหามตฺตา ปณีตํ ขาทนียํ โภชนียํ ปฏิยาทาเปตฺวา ภควโต กาลํ อาโรจาเปสุํ “กาโล โภ โคตม นิฏฺฐิตํ ภตฺตนฺ”ติ. อถ โข ภควา ปุพฺพณฺหสมยํ นิวาเสตฺวา ปตฺต\-จีวรมา\-ทาย เยน สุนิธ\-วสฺสการานํ มคธ\-มหามตฺตานํ ปริเวสนา เตนุปสงฺกมิ, อุปสงฺกมิตฺวา ปญฺญตฺเต อาสเน นิสีทิ สทฺธิํ ภิกฺขุ\-สํเฆน. อถ โข สุนิธ\-วสฺสการา มคธ\-มหามตฺตา พุทฺธ\-ปฺปมุขํ ภิกฺขุสํฆํ ปณีเตน ขาทนีเยน โภชนีเยน สหตฺถา สนฺตปฺเปตฺวา สมฺปวาเรตฺวา ภควนฺตํ ภุตฺตาวิํ โอนีต\-ปตฺต\-ปาณิํ เอกมนฺตํ นิสีทิํสุ. เอกมนฺตํ นิสินฺเน โข สุนิธวสฺสกาเร มคธ\-มหามตฺเต ภควา อิมาหิ คาถาหิ อนุโมทิ\csromandash{}}

\setlength{\csromangathapagewidth}{0pt}
\setlength{\csromangathawakwidth}{0pt}
\csromangathameasuregroup{“ยสฺมิํ ปเทเส กปฺเปติ, \\ สีลวนฺเตตฺถ โภเชตฺวา, \\ ยา ตตฺถ เทวตา อาสุํ, \\ ตา ปูชิตา ปูชยนฺติ, \\ ตโต นํ อนุกมฺปนฺติ, \\ เทวตานุกมฺปิโต โปโส,}{\csromangathabat{“ยสฺมิํ ปเทเส กปฺเปติ,}{วาสํ ปณฺฑิตชาติโย.} \\ \csromangathabat{สีลวนฺเตตฺถ โภเชตฺวา,}{สญฺญเต พฺรหฺมจารโย\textsuperscript{1}.} \\ \csromangathabat{ยา ตตฺถ เทวตา อาสุํ,}{ตาสํ ทกฺขิณมาทิเส.} \\ \csromangathabat{ตา ปูชิตา ปูชยนฺติ,}{มานิตา มานยนฺติ นํ.} \\ \csromangathabat{ตโต นํ อนุกมฺปนฺติ,}{มาตา ปุตฺตํว โอรสํ.} \\ \csromangathabat{เทวตานุกมฺปิโต โปโส,}{สทา ภทฺรานิ ปสฺสตี”ติ.}}
\csromangathasetleft{“ยสฺมิํ ปเทเส กปฺเปติ, \\ สีลวนฺเตตฺถ โภเชตฺวา, \\ ยา ตตฺถ เทวตา อาสุํ, \\ ตา ปูชิตา ปูชยนฺติ, \\ ตโต นํ อนุกมฺปนฺติ, \\ เทวตานุกมฺปิโต โปโส,}
\csromangathagroup{\csromangathastanza{\csromangathabat{“ยสฺมิํ ปเทเส กปฺเปติ,}{วาสํ ปณฺฑิตชาติโย.} \\ \csromangathabat{สีลวนฺเตตฺถ โภเชตฺวา,}{สญฺญเต พฺรหฺมจารโย\csromansharedfootnote{3f99d9f630010c44}{พฺรหฺมจาริโน (สฺยา)}.}}\csromangathastanza{\csromangathabat{ยา ตตฺถ เทวตา อาสุํ,}{ตาสํ ทกฺขิณมาทิเส.} \\ \csromangathabat{ตา ปูชิตา ปูชยนฺติ,}{มานิตา มานยนฺติ นํ.}}\csromangathastanza{\csromangathabat{ตโต นํ อนุกมฺปนฺติ,}{มาตา ปุตฺตํว โอรสํ.} \\ \csromangathabat{เทวตานุกมฺปิโต โปโส,}{สทา ภทฺรานิ ปสฺสตี”ติ.}}}

\prose{อถ โข ภควา สุนิธวสฺสกาเร มคธ\-มหามตฺเต อิมาหิ คาถาหิ อนุโมทิตฺวา อุฏฺฐายาสนา ปกฺกามิ. เตน โข ปน สมเยน สุนิธ\-วสฺสการา มคธ\-มหามตฺตา ภควนฺตํ ปิฏฺฐิโต ปิฏฺฐิโต อนุพนฺธา โหนฺติ “เยนชฺช สมโณ โคตโม ทฺวาเรน นิกฺขมิสฺสติ, ตํ โคตมทฺวารํ นาม \csromanfolio{325}ภวิสฺสติ, เยน ติตฺเถน คงฺคํ นทิํ อุตฺตริสฺสติ, ตํ โคตมติตฺถํ นาม ภวิสฺสตี”ติ. อถ โข ภควา เยน ทฺวาเรน นิกฺขมิ, ตํ โคตมทฺวารํ นาม อโหสิ. อถ โข ภควา เยน คงฺคา นที เตนุปสงฺกมิ, เตน โข ปน สมเยน คงฺคา นที ปูรา โหติ สมติตฺติกา กากเปยฺยา, มนุสฺสา อญฺเญ นาวํ ปริเยสนฺติ อญฺเญ อุฬุมฺปํ ปริเยสนฺติ อญฺเญ กุลฺลํ พนฺธนฺติ โอรา ปารํ คนฺตุกามา. อทฺทสา โข ภควา เต มนุสฺเส อญฺเญ นาวํ ปริเยสนฺเต อญฺเญ อุฬุมฺปํ ปริเยสนฺเต อญฺเญ กุลฺลํ พนฺธนฺเต โอรา ปารํ คนฺตุกาเม, ทิสฺวาน เสยฺยถาปิ นาม พลวา ปุริโส สมิญฺชิตํ วา พาหํ ปสาเรยฺย, ปสาริตํ วา พาหํ สมิญฺเชยฺย, เอวเมว โข คงฺคาย นทิยา โอริมตีเร อนฺตรหิโต ปาริมตีเร ปจฺจุฏฺฐาสิ สทฺธิํ ภิกฺขุ\-สํเฆน. อถ โข ภควา เอตมตฺถํ วิทิตฺวา ตายํ เวลายํ อิมํ อุทานํ อุทาเนสิ\csromandash{}}

\setlength{\csromangathapagewidth}{0pt}
\setlength{\csromangathawakwidth}{0pt}
\csromangathameasuregroup{}{“เย ตรนฺติ อณฺณวํ สรํ, \\ เสตุํ กตฺวาน วิสชฺช ปลฺลลานิ. \\ กุลฺลํ หิ ชโน พนฺธติ, \\ ติณฺณา เมธาวิโน ชนา”ติ.}
\csromangathameasurewak{“เย ตรนฺติ อณฺณวํ สรํ, \\ เสตุํ กตฺวาน วิสชฺช ปลฺลลานิ. \\ กุลฺลํ หิ ชโน พนฺธติ, \\ ติณฺณา เมธาวิโน ชนา”ติ.}
\setlength{\csromangathaleftcol}{0pt}
\csromangathagroup{\csromangathastanza{\csromangathawak{“เย ตรนฺติ อณฺณวํ สรํ, \\ เสตุํ กตฺวาน วิสชฺช ปลฺลลานิ. \\ กุลฺลํ หิ ชโน พนฺธติ, \\ ติณฺณา เมธาวิโน ชนา”ติ.}}}

\csromanmatikaanchor{mtk.181}
\csromantocmarkref{section}{175. โกฏิคาเม สจฺจกถา}{mtk.181}
\bookmark[dest={mtk.181},level=1]{175. โกฏิคาเม สจฺจกถา}
\csromanheader{1}{175. \textbf{โกฏิคาเม สจฺจกถา}}
\proseitem{287}{อถ โข ภควา เยน โกฏิคาโม เตนุปสงฺกมิ, ตตฺร สุทํ ภควา โกฏิคาเม วิหรติ. ตตฺร โข ภควา ภิกฺขู อามนฺเตสิ\csromandash{}}

\prose{\csromansymbolfootnote{*}{ที 2. 76 ปิฏฺเฐปิ.}“จตุนฺนํ ภิกฺขเว อริยสจฺจานํ อนนุโพธา อปฺปฏิเวธา เอวมิทํ ทีฆมทฺธานํ สนฺธาวิตํ สํสริตํ มมญฺเจว ตุมฺหากญฺจ. กตเมสํ จตุนฺนํ, ทุกฺขสฺส ภิกฺขเว อริยสจฺจสฺส อนนุโพธา อปฺปฏิเวธา เอวมิทํ ทีฆมทฺธานํ สนฺธาวิตํ สํสริตํ มมญฺเจว ตุมฺหากญฺจ. ทุกฺข\-สมุทยสฺส อริยสจฺจสฺส - ป-. ทุกฺข\-นิโรธสฺส อริยสจฺจสฺส ฯเปฯ. ทุกฺขนิโรธ\-คามินิยา ปฏิปทาย อริยสจฺจสฺส อนนุโพธา อปฺปฏิเวธา เอวมิทํ ทีฆมทฺธานํ สนฺธาวิตํ สํสริตํ มมญฺเจว ตุมฺหากญฺจ. ตยิทํ ภิกฺขเว ทุกฺขํ อริยสจฺจํ อนุพุทฺธํ ปฏิวิทฺธํ, ทุกฺข\-สมุทยํ\csromansharedfootnote{add3418b328ea08c}{ทุกฺขสมุทโย (สฺยา) ทุกฺขนิโรโธ (สฺยา)} อริยสจฺจํ อนุพุทฺธํ ปฏิวิทฺธํ, ทุกฺขนิโรธํ2 อริยสจฺจํ อนุพุทฺธํ ปฏิวิทฺธํ, ทุกฺขนิโรธ\-คามินี ปฏิปทา อริยสจฺจํ อนุพุทฺธํ ปฏิวิทฺธํ, อุจฺฉินฺนา ภวตณฺหา, ขีณา ภวเนตฺติ, นตฺถิทานิ ปุนพฺภโว”ติ.}

\setlength{\csromangathapagewidth}{0pt}
\setlength{\csromangathawakwidth}{0pt}
\csromangathameasuregroup{จตุนฺนํ อริยสจฺจานํ, \\ สํสิตํ ทีฆมทฺธานํ, \\ ตานิ เอตานิ ทิฏฺฐานิ, \\ อุจฺฉินฺนํ มูลํ ทุกฺขสฺส,}{\csromangathabat{จตุนฺนํ อริยสจฺจานํ,}{ยถาภูตํ อทสฺสนา.} \\ \csromangathabat{สํสิตํ ทีฆมทฺธานํ,}{ตาสุ ตาเสฺวว ชาติสุ.} \\ \csromangathabat{ตานิ เอตานิ ทิฏฺฐานิ,}{ภวเนตฺติ สมูหตา.} \\ \csromangathabat{อุจฺฉินฺนํ มูลํ ทุกฺขสฺส,}{นตฺถิทานิ ปุนพฺภโวติ.}}
\csromangathasetleft{จตุนฺนํ อริยสจฺจานํ, \\ สํสิตํ ทีฆมทฺธานํ, \\ ตานิ เอตานิ ทิฏฺฐานิ, \\ อุจฺฉินฺนํ มูลํ ทุกฺขสฺส,}
\csromangathagroup{\csromangathastanza{\csromanfolio{326}\csromangathabat{จตุนฺนํ อริยสจฺจานํ,}{ยถาภูตํ อทสฺสนา.} \\ \csromangathabat{สํสิตํ ทีฆมทฺธานํ,}{ตาสุ ตาเสฺวว ชาติสุ.}}\csromangathastanza{\csromangathabat{ตานิ เอตานิ ทิฏฺฐานิ,}{ภวเนตฺติ สมูหตา.} \\ \csromangathabat{อุจฺฉินฺนํ มูลํ ทุกฺขสฺส,}{นตฺถิทานิ ปุนพฺภโวติ.}}}

\csromanmatikaanchor{mtk.182}
\csromantocmarkref{section}{176. อมฺพปาลีวตฺถุ}{mtk.182}
\bookmark[dest={mtk.182},level=1]{176. อมฺพปาลีวตฺถุ}
\csromanheader{1}{176. \textbf{อมฺพปาลีวตฺถุ}}
\proseitem{288}{\csromansymbolfootnote{*}{ที 2. 81 ปิฏฺฐาทีสุปิ.}อสฺโสสิ โข อมฺพปาลี คณิกา “ภควา กิร โกฏิคามํ อนุปฺปตฺโต”ติ, อถ โข อมฺพปาลี คณิกา ภทฺรานิ ภทฺรานิ ยานานิ โยชาเปตฺวา ภทฺรํ ภทฺรํ ยานํ อภิรุหิตฺวา ภเทฺรหิ ภเทฺรหิ ยาเนหิ เวสาลิยา นิยฺยาสิ ภควนฺตํ ทสฺสนาย, ยาวติกา ยานสฺส ภูมิ, ยาเนน คนฺตฺวา ยานา ปจฺโจโรหิตฺวา ปตฺติกาว เยน ภควา เตนุปสงฺกมิ, อุปสงฺกมิตฺวา ภควนฺตํ อภิวาเทตฺวา เอกมนฺตํ นิสีทิ. เอกมนฺตํ นิสินฺนํ โข อมฺพปาลิํ คณิกํ ภควา ธมฺมิยา กถาย สนฺทสฺเสสิ สมาทเปสิ สมุตฺเตเชสิ สมฺปหํเสสิ. อถ โข อมฺพปาลี คณิกา ภควตา ธมฺมิยา กถาย สนฺทสฺสิตา สมาทปิตา สมุตฺเตชิตา สมฺปหํสิตา ภควนฺตํ เอตทโวจ “อธิวาเสตุ เม ภนฺเต ภควา สฺวาตนาย ภตฺตํ สทฺธิํ ภิกฺขุสํเฆนา”ติ. อธิวาเสสิ ภควา ตุณฺหีภาเวน. อถ โข อมฺพปาลี คณิกา ภควโต อธิวาสนํ วิทิตฺวา อุฏฺฐายาสนา ภควนฺตํ อภิวาเทตฺวา ปทกฺขิณํ กตฺวา ปกฺกามิ.}

\csromanmatikaanchor{mtk.183}
\csromantocmarkref{section}{177. ลิจฺฉวีวตฺถุ}{mtk.183}
\bookmark[dest={mtk.183},level=1]{177. ลิจฺฉวีวตฺถุ}
\csromanheader{1}{177. \textbf{ลิจฺฉวีวตฺถุ}}
\proseitem{289}{\csromansymbolmark{*}อสฺโสสุํ โข เวสาลิกา ลิจฺฉวี “ภควา กิร โกฏิคามํ อนุปฺปตฺโต”ติ, อถ โข เวสาลิกา ลิจฺฉวี ภทฺรานิ ภทฺรานิ ยานานิ โยชาเปตฺวา ภทฺรํ ภทฺรํ ยานํ อภิรุหิตฺวา ภเทฺรหิ ภเทฺรหิ ยาเนหิ เวสาลิยา นิยฺยาสุํ ภควนฺตํ ทสฺสนาย, อปฺเปกจฺเจ ลิจฺฉวี นีลา โหนฺติ นีลวณฺณา นีลวตฺถา นีลาลงฺการา, อปฺเปกจฺเจ ลิจฺฉวี ปีตา โหนฺติ ปีตวณฺณา ปีตวตฺถา ปีตาลงฺการา, อปฺเปกจฺเจ ลิจฺฉวี โลหิตา โหนฺติ โลหิตวณฺณา โลหิตวตฺถา โลหิตาลงฺการา, อปฺเปกจฺเจ ลิจฺฉวี โอทาตา โหนฺติ โอทาตวณฺณา โอทาตวตฺถา โอทาตาลงฺการา. อถ โข อมฺพปาลี คณิกา ทหรานํ ทหรานํ ลิจฺฉวีนํ อีสาย อีสํ ยุเคน ยุคํ จกฺเกน จกฺกํ อกฺเขน อกฺขํ ปฏิวฏฺเฏสิ\csromansharedfootnote{71fcf68482342da8}{ปฏิวตฺเตสิ (ก)}. \csromanfolio{327}อถ โข เต ลิจฺฉวี อมฺพปาลิํ คณิกํ เอตทโวจุํ “กิสฺส เช อมฺพปาลิ ทหารานํ ทหรานํ1 ลิจฺฉวีนํ อีสาย อีสํ ยุเคน ยุคํ จกฺเกน จกฺกํ อกฺเขน อกฺขํ ปฏิวฏฺเฏสี”ติ. ตถา หิ ปน มยา อยฺยปุตฺตา สฺวาตนาย พุทฺธปฺปมุโข ภิกฺขุสํโฆ นิมนฺติโตติ. เทหิ เช อมฺพปาลิ อมฺหากํ เอตํ ภตฺตํ สตสหสฺเสนาติ. สเจปิ เม อยฺยปุตฺตา เวสาลิํ สาหารํ ทชฺเชยฺยาถ, เนว ทชฺชาหํ ตํ ภตฺตนฺติ. อถ โข เต ลิจฺฉวี องฺคุลิํ โผเฏสุํ “ชิตมฺห วต โภ อมฺพกาย, ปราชิตมฺห วต โภ อมฺพกายา”ติ. อถ โข เต ลิจฺฉวี เยน ภควา เตนุ\-ปสงฺกมิํสุ. อทฺทสา โข ภควา เต ลิจฺฉวี ทูรโตว อาคจฺฉนฺเต, ทิสฺวาน ภิกฺขู อามนฺเตสิ “เยหิ ภิกฺขเว ภิกฺขูหิ เทวา ตาวติํสา อทิฏฺฐปุพฺพา, โอโลเกถ ภิกฺขเว ลิจฺฉวี\-ปริสํ, อปโลเกถ ภิกฺขเว ลิจฺฉวี\-ปริสํ, อุปสํหรถ ภิกฺขเว ลิจฺฉวี\-ปริสํ ตาวติํสปริสนฺ”ติ. อถ โข เต ลิจฺฉวี ยาวติกา ยานสฺส ภูมิ, ยาเนน คนฺตฺวา ยานา ปจฺโจโรหิตฺวา ปตฺติกาว เยน ภควา เตนุ\-ปสงฺกมิํสุ, อุปสงฺกมิตฺวา ภควนฺตํ อภิวาเทตฺวา เอกมนฺตํ นิสีทิํสุ. เอกมนฺตํ นิสินฺเน โข เต ลิจฺฉวี ภควา ธมฺมิยา กถาย สนฺทสฺเสสิ สมาทเปสิ สมุตฺเตเชสิ สมฺปหํเสสิ. อถ โข เต ลิจฺฉวี ภควตา ธมฺมิยา กถาย สนฺทสฺสิตา สมาทปิตา สมุตฺเตชิตา สมฺปหํสิตา ภควนฺตํ เอตทโวจุํ “อธิวาเสตุ โน ภนฺเต ภควา สฺวาตนาย ภตฺตํ สทฺธิํ ภิกฺขุสํเฆนา”ติ. อธิวุฏฺโฐมฺหิ ลิจฺฉวี สฺวาตนาย อมฺพปาลิยา คณิกาย ภตฺตนฺติ. อถ โข เต ลิจฺฉวี องฺคุลิํ โผเฏสุํ “ชิตมฺห วต โภ อมฺพกาย, ปราชิตมฺห วต โภ อมฺพกายา”ติ. อถ โข เต ลิจฺฉวี ภควโต ภาสิตํ อภินนฺทิตฺวา อนุโมทิตฺวา อุฏฺฐายาสนา ภควนฺตํ อภิวาเทตฺวา ปทกฺขิณํ กตฺวา ปกฺกมิํสุ.}

\prose{อถ โข ภควา โกฏิคาเม ยถาภิรนฺตํ วิหริตฺวา2 เยน นาติกา3 เตนุปสงฺกมิ, ตตฺร สุทํ ภควา นาติเก วิหรติ คิญฺชกาวสเถ. อถ โข อมฺพปาลี คณิกา ตสฺสา รตฺติยา อจฺจเยน สเก อาราเม ปณีตํ ขาทนียํ โภชนียํ ปฏิยาทาเปตฺวา ภควโต กาลํ อาโรจาเปสิ “กาโล ภนฺเต นิฏฺฐิตํ ภตฺตนฺ”ติ. อถ โข ภควา}

\proseitem{1}{อมฺหากํ ทหรานํ ทหรานํ (สี, สฺยา)}

\proseitem{2}{มหา\-ปรินิพฺพาน\-สุตฺเต อนุสนฺธิ อญฺญถา อาคโต.}

\proseitemwithcloserruled{3}{นาทิกา (สี, สฺยา) \csromanfolio{328}ปุพฺพณฺหสมยํ นิวาเสตฺวา ปตฺต\-จีวรมา\-ทาย เยน อมฺพปาลิยา คณิกาย ปริเวสนา เตนุปสงฺกมิ, อุปสงฺกมิตฺวา ปญฺญตฺเต อาสเน นิสีทิ สทฺธิํ ภิกฺขุ\-สํเฆน. อถ โข อมฺพปาลี คณิกา พุทฺธ\-ปฺปมุขํ ภิกฺขุสํฆํ ปณีเตน ขาทนีเยน โภชนีเยน สหตฺถา สนฺตปฺเปตฺวา สมฺปวาเรตฺวา ภควนฺตํ ภุตฺตาวิํ โอนีต\-ปตฺต\-ปาณิํ เอกมนฺตํ นิสีทิ, เอกมนฺตํ นิสินฺนา โข อมฺพปาลี คณิกา ภควนฺตํ เอตทโวจ “อิมาหํ ภนฺเต อมฺพวนํ พุทฺธ\-ปฺปมุขสฺส ภิกฺขุ\-สํฆสฺส ทมฺมี”ติ. ปฏิคฺคเหสิ ภควา อารามํ. อถ โข ภควา อมฺพปาลิํ คณิกํ ธมฺมิยา กถาย สนฺทสฺเสตฺวา สมาทเปตฺวา สมุตฺเตเชตฺวา สมฺปหํเสตฺวา อุฏฺฐายาสนา เยน มหาวนํ เตนุปสงฺกมิ, ตตฺร สุทํ ภควา เวสาลิยํ วิหรติ มหาวเน กูฏาคาร\-สาลายํ.}{ลิจฺฉวิ\-ภาณวาโร นิฏฺฐิโต ตติโย.}
\csromanmatikaanchor{mtk.184}
\csromantocmarkref{section}{178. สีหเสนาปติวตฺถุ}{mtk.184}
\bookmark[dest={mtk.184},level=1]{178. สีหเสนาปติวตฺถุ}
\csromanheader{1}{178. \textbf{สีหเสนาปติ}\-\textbf{วตฺถุ}}
\proseitem{290}{\csromansymbolfootnote{*}{อํ 3. 24 ปิฏฺฐาทีสุปิ.}เตน โข ปน สมเยน อภิญฺญาตา อภิญฺญาตา ลิจฺฉวี สนฺธาคาเร\csromansharedfootnote{7913fc871d18a9a9}{สนฺถาคาเร (สี, สฺยา)} สนฺนิสินฺนา สนฺนิปติตา อเนก\-ปริยาเยน พุทฺธสฺส วณฺณํ ภาสนฺติ, ธมฺมสฺส วณฺณํ ภาสนฺติ, สํฆสฺส วณฺณํ ภาสนฺติ. เตน โข ปน สมเยน สีโห เสนาปติ นิคณฺฐสาวโก ตสฺสํ ปริสายํ นิสินฺโน โหติ. อถ โข สีหสฺส เสนาปติสฺส เอตทโหสิ “นิสฺสํสยํ โข โส ภควา อรหํ สมฺมาสมฺพุทฺโธ ภวิสฺสติ, ตถา หิเม อภิญฺญาตา อภิญฺญาตา ลิจฺฉวี สนฺธาคาเร สนฺนิสินฺนา สนฺนิปติตา อเนก\-ปริยาเยน พุทฺธสฺส วณฺณํ ภาสนฺติ, ธมฺมสฺส วณฺณํ ภาสนฺติ, สํฆสฺส วณฺณํ ภาสนฺติ, ยนฺนูนาหํ ตํ ภควนฺตํ ทสฺสนาย อุปสงฺกเมยฺยํ อรหนฺตํ สมฺมาสมฺพุทฺธนฺ”ติ. อถ โข สีโห เสนาปติ เยน นิคณฺโฐ นาฏปุตฺโต เตนุปสงฺกมิ, อุปสงฺกมิตฺวา นิคณฺฐํ นาฏปุตฺตํ เอตทโวจ “อิจฺฉามหํ ภนฺเต สมณํ โคตมํ ทสฺสนาย อุปสงฺกมิตุนฺ”ติ. กิํ ปน ตฺวํ สีห กิริยวาโท สมาโน อกิริยวาทํ สมณํ โคตมํ ทสฺสนาย อุปสงฺกมิสฺสสิ, สมโณ หิ สีห โคตโม อกิริยวาโท อกิริยาย ธมฺมํ เทเสติ, เตน จ สาวเก วิเนตีติ. อถ โข สีหสฺส เสนาปติสฺส โย อโหสิ คมิกา\-ภิสงฺขาโร ภควนฺตํ ทสฺสนาย, \csromanfolio{329}โส ปฏิปฺปสฺสมฺภิ. ทุติยมฺปิ โข อภิญฺญาตา อภิญฺญาตา ลิจฺฉวี สนฺธาคาเร สนฺนิสินฺนา สนฺนิปติตา อเนก\-ปริยาเยน พุทฺธสฺส วณฺณํ ภาสนฺติ, ธมฺมสฺส วณฺณํ ภาสนฺติ, สํฆสฺส วณฺณํ ภาสนฺติ. ทุติยมฺปิ โข สีหสฺส เสนาปติสฺส เอตทโหสิ “นิสฺสํสยํ โข โส ภควา อรหํ สมฺมาสมฺพุทฺโธ ภวิสฺสติ, ตถา หิเม อภิญฺญาตา อภิญฺญาตา ลิจฺฉวี สนฺธาคาเร สนฺนิสินฺนา สนฺนิปติตา อเนก\-ปริยาเยน พุทฺธสฺส วณฺณํ ภาสนฺติ, ธมฺมสฺส วณฺณํ ภาสนฺติ, สํฆสฺส วณฺณํ ภาสนฺติ, ยนฺนูนาหํ ตํ ภควนฺตํ ทสฺสนาย อุปสงฺกเมยฺยํ อรหนฺตํ สมฺมาสมฺพุทฺธนฺ”ติ. ทุติยมฺปิ โข สีโห เสนาปติ เยน นิคณฺโฐ นาฏปุตฺโต เตนุปสงฺกมิ, อุปสงฺกมิตฺวา นิคณฺฐํ นาฏปุตฺตํ เอตทโวจ “อิจฺฉามหํ ภนฺเต สมณํ โคตมํ ทสฺสนาย อุปสงฺกมิตุนฺ”ติ. กิํ ปน ตฺวํ สีห กิริยวาโท สมาโน อกิริยวาทํ สมณํ โคตมํ ทสฺสนาย อุปสงฺกมิสฺสสิ, สมโณ หิ สีห โคตโม อกิริยวาโท อกิริยาย ธมฺมํ เทเสติ, เตน จ สาวเก วิเนตีติ. ทุติยมฺปิ โข สีหสฺส เสนาปติสฺส โย อโหสิ คมิกา\-ภิสงฺขาโร ภควนฺตํ ทสฺสนาย, โส ปฏิปฺปสฺสมฺภิ. ตติยมฺปิ โข อภิญฺญาตา อภิญฺญาตา ลิจฺฉวี สนฺธาคาเร สนฺนิสินฺนา สนฺนิปติตา อเนก\-ปริยาเยน พุทฺธสฺส วณฺณํ ภาสนฺติ, ธมฺมสฺส วณฺณํ ภาสนฺติ, สํฆสฺส วณฺณํ ภาสนฺติ. ตติยมฺปิ โข สีหสฺส เสนาปติสฺส เอตทโหสิ “นิสฺสํสยํ โข โส ภควา อรหํ สมฺมาสมฺพุทฺโธ ภวิสฺสติ, ตถา หิเม อภิญฺญาตา อภิญฺญาตา ลิจฺฉวี สนฺธาคาเร สนฺนิสินฺนา สนฺนิปติตา อเนก\-ปริยาเยน พุทฺธสฺส วณฺณํ ภาสนฺติ, ธมฺมสฺส วณฺณํ ภาสนฺติ, สํฆสฺส วณฺณํ ภาสนฺติ, กิํ หิ เม กริสฺสนฺติ นิคณฺฐา อปโลกิตา วา อนปโลกิตา วา, ยนฺนูนาหํ อนปโลเกตฺวาว นิคณฺเฐ ตํ ภควนฺตํ ทสฺสนาย อุปสงฺกเมยฺยํ อรหนฺตํ สมฺมาสมฺพุทฺธนฺ”ติ.}

\prose{อถ โข สีโห เสนาปติ ปญฺจหิ รถสเตหิ ทิวา ทิวสฺส เวสาลิยา นิยฺยาสิ ภควนฺตํ ทสฺสนาย, ยาวติกา ยานสฺส ภูมิ, ยาเนน คนฺตฺวา ยานา ปจฺโจโรหิตฺวา ปตฺติโกว เยน ภควา เตนุปสงฺกมิ, อุปสงฺกมิตฺวา ภควนฺตํ อภิวาเทตฺวา เอกมนฺตํ นิสีทิ, เอกมนฺตํ นิสินฺโน โข สีโห เสนาปติ ภควนฺตํ เอตทโวจ “สุตํ เม ตํ ภนฺเต ‘อกิริยวาโท สมโณ โคตโม อกิริยาย ธมฺมํ เทเสติ, เตน จ สาวเก วิเนตี’ติ. เย เต ภนฺเต เอวมาหํสุ ‘อกิริยวาโท \csromanfolio{330}สมโณ โคตโม อกิริยาย ธมฺมํ เทเสติ, เตน จ สาวเก วิเนตี’ติ. กจฺจิ เต ภนฺเต ภควโต วุตฺตวาทิโน, น จ ภควนฺตํ อภูเตน อพฺภาจิกฺขนฺติ, ธมฺมสฺส จานุธมฺมํ พฺยากโรนฺติ, น จ โกจิ สหธมฺมิโก วาทานุวาโท คารยฺหํ ฐานํ อาคจฺฉติ, อนพฺภกฺขาตุ\-กามา หิ มยํ ภนฺเต ภควนฺตนฺ”ติ.}

\proseitem{291}{อตฺถิ สีห ปริยาโย, เยน มํ ปริยาเยน สมฺมา วทมาโน วเทยฺย “อกิริยวาโท สมโณ โคตโม อกิริยาย ธมฺมํ เทเสติ, เตน จ สาวเก วิเนตี”ติ. อตฺถิ สีห ปริยาโย, เยน มํ ปริยาเยน สมฺมา วทมาโน วเทยฺย “กิริยวาโท สมโณ โคตโม กิริยาย ธมฺมํ เทเสติ, เตน จ สาวเก วิเนตี”ติ. อตฺถิ สีห ปริยาโย, เยน มํ ปริยาเยน สมฺมา วทมาโน วเทยฺย “อุจฺเฉทวาโท สมโณ โคตโม อุจฺเฉทาย ธมฺมํ เทเสติ, เตน จ สาวเก วิเนตี”ติ. อตฺถิ สีห ปริยาโย, เยน มํ ปริยาเยน สมฺมา วทมาโน วเทยฺย “เชคุจฺฉี สมโณ โคตโม เชคุจฺฉิตาย ธมฺมํ เทเสติ, เตน จ สาวเก วิเนตี”ติ. อตฺถิ สีห ปริยาโย, เยน มํ ปริยาเยน สมฺมา วทมาโน วเทยฺย “เวนยิโก สมโณ โคตโม วินยาย ธมฺมํ เทเสติ, เตน จ สาวเก วิเนตี”ติ. อตฺถิ สีห ปริยาโย, เยน มํ ปริยาเยน สมฺมา วทมาโน วเทยฺย “ตปสฺสี สมโณ โคตโม ตปสฺสิตาย ธมฺมํ เทเสติ, เตน จ สาวเก วิเนตี”ติ. อตฺถิ สีห ปริยาโย, เยน มํ ปริยาเยน สมฺมา วทมาโน วเทยฺย “อปคพฺโภ สมโณ โคตโม อปคพฺภตาย ธมฺมํ เทเสติ, เตน จ สาวเก วิเนตี”ติ. อตฺถิ สีห ปริยาโย, เยน มํ ปริยาเยน สมฺมา วทมาโน วเทยฺย “อสฺสตฺโถ สมโณ โคตโม อสฺสาสาย ธมฺมํ เทเสติ, เตน จ สาวเก วิเนตี”ติ.}

\proseitem{292}{กตโม จ สีห ปริยาโย, เยน มํ ปริยาเยน สมฺมา วทมาโน วเทยฺย “อกิริยวาโท สมโณ โคตโม อกิริยาย ธมฺมํ เทเสติ, เตน จ สาวเก วิเนตี”ติ. อหํ หิ สีห อกิริยํ วทามิ กาย\-ทุจฺจริตสฺส วจีทุจฺจริตสฺส มโน\-ทุจฺจริตสฺส, อเนกวิหิตานํ ปาปกานํ อกุสลานํ ธมฺมานํ อกิริยํ วทามิ. อยํ โข สีห \csromanfolio{331}ปริยาโย, เยน มํ ปริยาเยน สมฺมา วทมาโน วเทยฺย “อกิริยวาโท สมโณ โคตโม อกิริยาย ธมฺมํ เทเสติ, เตน จ สาวเก วิเนตี”ติ.}

\prose{กตโม จ สีห ปริยาโย, เยน มํ ปริยาเยน สมฺมา วทมาโน วเทยฺย “กิริยวาโท สมโณ โคตโม กิริยาย ธมฺมํ เทเสติ, เตน จ สาวเก วิเนตี”ติ.}

\prose{อหํ หิ สีห กิริยํ วทามิ กาย\-สุจริตสฺส วจีสุจริตสฺส มโน\-สุจริตสฺส, อเนกวิหิตานํ กุสลานํ ธมฺมานํ กิริยํ วทามิ.}

\prose{อยํ โข สีห ปริยาโย, เยน มํ ปริยาเยน สมฺมา วทมาโน วเทยฺย “กิริยวาโท สมโณ โคตโม กิริยาย ธมฺมํ เทเสติ, เตน จ สาวเก วิเนตี”ติ.}

\prose{กตโม จ สีห ปริยาโย, เยน มํ ปริยาเยน สมฺมา วทมาโน วเทยฺย “อุจฺเฉทวาโท สมโณ โคตโม อุจฺเฉทาย ธมฺมํ เทเสติ, เตน จ สาวเก วิเนตี”ติ.}

\prose{อหํ หิ สีห อุจฺเฉทํ วทามิ ราคสฺส โทสสฺส โมหสฺส, อเนกวิหิตานํ ปาปกานํ อกุสลานํ ธมฺมานํ อุจฺเฉทํ วทามิ.}

\prose{อยํ โข สีห ปริยาโย, เยน มํ ปริยาเยน สมฺมา วทมาโน วเทยฺย “อุจฺเฉทวาโท สมโณ โคตโม อุจฺเฉทาย ธมฺมํ เทเสติ, เตน จ สาวเก วิเนตี”ติ.}

\prose{กตโม จ สีห ปริยาโย, เยน มํ ปริยาเยน สมฺมา วทมาโน วเทยฺย “เชคุจฺฉี สมโณ โคตโม เชคุจฺฉิตาย ธมฺมํ เทเสติ, เตน จ สาวเก วิเนตี”ติ.}

\prose{อหํ หิ สีห ชิคุจฺฉามิ กาย\-ทุจฺจริเตน วจีทุจฺจริเตน มโน\-ทุจฺจริเตน, อเนกวิหิตานํ ปาปกานํ อกุสลานํ ธมฺมานํ สมาปตฺติยา ชิคุจฺฉามิ.}

\prose{อยํ โข สีห ปริยาโย, เยน มํ ปริยาเยน สมฺมา วทมาโน วเทยฺย “เชคุจฺฉี สมโณ โคตโม เชคุจฺฉิตาย ธมฺมํ เทเสติ, เตน จ สาวเก วิเนตี”ติ.}

\proseitem{292}{กตโม จ สีห ปริยาโย, เยน มํ ปริยาเยน สมฺมา วทมาโน วเทยฺย “เวนยิโก สมโณ โคตโม, วินยาย ธมฺมํ เทเสติ, เตน จ สาวเก วิเนตี”ติ. อหํ หิ สีห วินยาย ธมฺมํ เทเสมิ ราคสฺส โทสสฺส โมหสฺส, อเนกวิหิตานํ ปาปกานํ \csromanfolio{332}อกุสลานํ ธมฺมานํ วินยาย ธมฺมํ เทเสมิ. อยํ โข สีห ปริยาโย, เยน มํ ปริยาเยน สมฺมา วทมาโน วเทยฺย “เวนยิโก สมโณ โคตโม วินยาย ธมฺมํ เทเสติ, เตน จ สาวเก วิเนตี”ติ.}

\prose{กตโม จ สีห ปริยาโย, เยน มํ ปริยาเยน สมฺมา วทมาโน วเทยฺย “ตปสฺสี สมโณ โคตโม ตปสฺสิตาย ธมฺมํ เทเสติ, เตน จ สาวเก วิเนตี”ติ.}

\prose{ตปนียาหํ สีห ปาปเก อกุสเล ธมฺเม วทามิ กายทุจฺจริตํ วจีทุจฺจริตํ มโนทุจฺจริตํ, ยสฺส โข สีห ตปนียา ปาปกา อกุสลา ธมฺมา ปหีนา อุจฺฉินฺนมูลา ตาลาวตฺถุกตา อนภาวํกตา อายติํ อนุปฺปาทธมฺมา, ตมหํ ตปสฺสีติ วทามิ.}

\prose{ตถาคตสฺส โข สีห ตปนียา ปาปกา อกุสลา ธมฺมา ปหีนา อุจฺฉินฺนมูลา ตาลาวตฺถุกตา อนภาวํกตา อายติํ อนุปฺปาทธมฺมา. อยํ โข สีห ปริยาโย, เยน มํ ปริยาเยน สมฺมา วทมาโน วเทยฺย “ตปสฺสี สมโณ โคตโม ตปสฺสิตาย ธมฺมํ เทเสติ, เตน จ สาวเก วิเนตี”ติ.}

\prose{กตโม จ สีห ปริยาโย, เยน มํ ปริยาเยน สมฺมา วทมาโน วเทยฺย “อปคพฺโภ สมโณ โคตโม อปคพฺภตาย ธมฺมํ เทเสติ, เตน จ สาวเก วิเนตี”ติ.}

\prose{ยสฺส โข สีห อายติํ คพฺภเสยฺยา ปุนพฺภวา\-ภินิพฺพตฺติ ปหีนา อุจฺฉินฺนมูลา ตาลาวตฺถุกตา อนภาวํกตา อายติํ อนุปฺปาทธมฺมา, ตมหํ อปคพฺโภติ วทามิ.}

\prose{ตถาคตสฺส โข สีห อายติํ คพฺภเสยฺยา ปุนพฺภวา\-ภินิพฺพตฺติ ปหีนา อุจฺฉินฺนมูลา ตาลาวตฺถุกตา อนภาวํกตา อายติํ อนุปฺปาทธมฺมา. อยํ โข สีห ปริยาโย, เยน มํ ปริยาเยน สมฺมา วทมาโน วเทยฺย “อปคพฺโภ สมโณ โคตโม อปคพฺภตาย ธมฺมํ เทเสติ, เตน จ สาวเก วิเนตี”ติ.}

\prose{กตโม จ สีห ปริยาโย, เยน มํ ปริยาเยน สมฺมา วทมาโน วเทยฺย “อสฺสตฺโถ สมโณ โคตโม อสฺสาสาย ธมฺมํ เทเสติ, เตน จ สาวเก วิเนตี”ติ.}

\prose{อหํ หิ สีห อสฺสตฺโถ ปรเมน อสฺสาเสน, อสฺสาสาย ธมฺมํ เทเสมิ, เตน จ สาวเก วิเนมิ.}

\prose{อยํ โข สีห ปริยาโย, เยน มํ ปริยาเยน สมฺมา วทมาโน วเทยฺย “อสฺสตฺโถ สมโณ โคตโม อสฺสาสาย ธมฺมํ เทเสติ, เตน จ สาวเก วิเนตี”ติ.}

\proseitem{293}{\csromanfolio{333}เอวํ วุตฺเต สีโห เสนาปติ ภควนฺตํ เอตทโวจ “อภิกฺกนฺตํ ภนฺเต ฯเปฯ อุปาสกํ มํ ภควา ธาเรตุ อชฺชตคฺเค ปาณุเปตํ สรณํ คตนฺ”ติ. อนุวิจฺจการํ\csromansharedfootnote{f4de6f549738237f}{อนุวิชฺชการํ (ก)} โข สีห กโรหิ, อนุวิจฺจกาโร ตุมฺหาทิสานํ ญาตมนุสฺสานํ สาธุ โหตีติ. อิมินาปาหํ ภนฺเต ภควโต ภิยฺโยโส มตฺตาย อตฺตมโน อภิรทฺโธ, ยํ มํ ภควา เอวมาห “อนุวิจฺจการํ โข สีห กโรหิ, อนุวิจฺจกาโร ตุมฺหาทิสานํ ญาตมนุสฺสานํ สาธุ โหตี”ติ. มมํ หิ ภนฺเต อญฺญติตฺถิยา สาวกํ ลภิตฺวา เกวลกปฺปํ เวสาลิํ ปฏากํ ปริหเรยฺยุํ “สีโห โข อมฺหากํ เสนาปติ สาวกตฺตํ อุปคโต”ติ. อถ จ ปน มํ ภควา เอวมาห “อนุวิจฺจการํ โข สีห กโรหิ, อนุวิจฺจกาโร ตุมฺหาทิสานํ ญาตมนุสฺสานํ สาธุ โหตี”ติ. เอสาหํ ภนฺเต ทุติยมฺปิ ภควนฺตํ สรณํ คจฺฉามิ ธมฺมญฺจ ภิกฺขุสํฆญฺจ, อุปาสกํ มํ ภควา ธาเรตุ อชฺชตคฺเค ปาณุเปตํ สรณํ คตนฺติ. ทีฆรตฺตํ โข เต สีห นิคณฺฐานํ โอปานภูตํ กุลํ, เยน เนสํ อุปคตานํ ปิณฺฑกํ ทาตพฺพํ มญฺเญยฺยาสีติ. อิมินาปาหํ ภนฺเต ภควโต ภิยฺโยโส มตฺตาย อตฺตมโน อภิรทฺโธ, ยํ มํ ภควา เอวมาห “ทีฆรตฺตํ โข เต สีห นิคณฺฐานํ โอปานภูตํ กุลํ, เยน เนสํ อุปคตานํ ปิณฺฑกํ ทาตพฺพํ มญฺเญยฺยาสี”ติ. สุตํ เม ตํ ภนฺเต สมโณ โคตโม เอวมาห “มยฺหเมว ทานํ ทาตพฺพํ, น อญฺเญสํ ทานํ ทาตพฺพํ, มยฺหเมว สาวกานํ ทานํ ทาตพฺพํ, น อญฺเญสํ สาวกานํ ทานํ ทาตพฺพํ, มยฺหเมว ทินฺนํ มหปฺผลํ, น อญฺเญสํ ทินฺนํ มหปฺผลํ, มยฺหเมว สาวกานํ ทินฺนํ มหปฺผลํ, น อญฺเญสํ สาวกานํ ทินฺนํ มหปฺผลนฺ”ติ, อถ จ ปน มํ ภควา นิคณฺเฐสุปิ ทาเน สมาทเปติ, อปิ จ ภนฺเต มยเมตฺถ กาลํ ชานิสฺสาม, เอสาหํ ภนฺเต ตติยมฺปิ ภควนฺตํ สรณํ คจฺฉามิ ธมฺมญฺจ ภิกฺขุสํฆญฺจ, อุปาสกํ มํ ภควา ธาเรตุ อชฺชตคฺเค ปาณุเปตํ สรณํ คตนฺติ.}

\prose{อถ โข ภควา สีหสฺส เสนาปติสฺส อนุปุพฺพิํ กถํ กเถสิ. เสยฺยถิทํ, ทานกถํ ฯเปฯ อปรปฺปจฺจโย สตฺถุสาสเน ภควนฺตํ เอตทโวจ “อธิวาเสตุ เม ภนฺเต ภควา สฺวาตนาย ภตฺตํ สทฺธิํ ภิกฺขุสํเฆนา”ติ. อธิวาเสสิ ภควา ตุณฺหีภาเวน. \csromanfolio{334}อถ โข สีโห เสนาปติ ภควโต อธิวาสนํ วิทิตฺวา อุฏฺฐายาสนา ภควนฺตํ อภิวาเทตฺวา ปทกฺขิณํ กตฺวา ปกฺกามิ.}

\proseitem{294}{อถ โข สีโห เสนาปติ อญฺญตรํ ปุริสํ อาณาเปสิ “คจฺฉ ภเณ ปวตฺตมํสํ ชานาหี”ติ. อถ โข สีโห เสนาปติ ตสฺสา รตฺติยา อจฺจเยน ปณีตํ ขาทนียํ โภชนียํ ปฏิยาทาเปตฺวา ภควโต กาลํ อาโรจาเปสิ “กาโล ภนฺเต นิฏฺฐิตํ ภตฺตนฺ”ติ. อถ โข ภควา ปุพฺพณฺหสมยํ นิวาเสตฺวา ปตฺต\-จีวรมา\-ทาย เยน สีหสฺส เสนาปติสฺส นิเวสนํ เตนุปสงฺกมิ, อุปสงฺกมิตฺวา ปญฺญตฺเต อาสเน นิสีทิ สทฺธิํ ภิกฺขุ\-สํเฆน.}

\prose{เตน โข ปน สมเยน สมฺพหุลา นิคณฺฐา เวสาลิยํ รถิกาย รถิกํ สิงฺฆาฏเกน สิงฺฆาฏกํ พาหา ปคฺคยฺห กนฺทนฺติ “อชฺช สีเหน เสนาปตินา ถูลํ ปสุํ วธิตฺวา สมณสฺส โคตมสฺส ภตฺตํ กตํ, ตํ สมโณ โคตโม ชานํ อุทฺทิสฺสกตํ มํสํ ปริภุญฺชติ ปฏิจฺจกมฺมนฺ”ติ. อถ โข อญฺญตโร ปุริโส เยน สีโห เสนาปติ เตนุปสงฺกมิ, อุปสงฺกมิตฺวา สีหสฺส เสนาปติสฺส อุปกณฺณเก อาโรเจสิ “ยคฺเฆ ภนฺเต ชาเนยฺยาสิ, เอเต สมฺพหุลา นิคณฺฐา เวสาลิยํ รถิกาย รถิกํ สิงฺฆาฏเกน สิงฺฆาฏกํ พาหา ปคฺคยฺห กนฺทนฺติ ‘อชฺช สีเหน เสนาปตินา ถูลํ ปสุํ วธิตฺวา สมณสฺส โคตมสฺส ภตฺตํ กตํ, ตํ สมโณ โคตโม ชานํ อุทฺทิสฺสกตํ มํสํ ปริภุญฺชติ ปฏิจฺจกมฺมนฺ’ติ”. อลํ อยฺโย ทีฆรตฺตมฺปิ เต อายสฺมนฺตา อวณฺณกามา พุทฺธสฺส, อวณฺณกามา ธมฺมสฺส, อวณฺณกามา สํฆสฺส, น จ ปน เต อายสฺมนฺตา ชิริทนฺติ ตํ ภควนฺตํ อสตา ตุจฺฉา มุสา อภูเตน อพฺภาจิกฺขนฺตา, น จ มยํ ชีวิตเหตุปิ สญฺจิจฺจ ปาณํ ชีวิตา โวโรเปยฺยามาติ. อถ โข สีโห เสนาปติ พุทฺธ\-ปฺปมุขํ ภิกฺขุสํฆํ ปณีเตน ขาทนีเยน โภชนีเยน สหตฺถา สนฺตปฺเปตฺวา สมฺปวาเรตฺวา ภควนฺตํ ภุตฺตาวิํ โอนีต\-ปตฺต\-ปาณิํ เอกมนฺตํ นิสีทิ. เอกมนฺตํ นิสินฺนํ โข สีหํ เสนาปติํ ภควา ธมฺมิยา กถาย สนฺทสฺเสตฺวา สมาทเปตฺวา สมุตฺเตเชตฺวา สมฺปหํเสตฺวา อุฏฺฐายาสนา ปกฺกามิ. อถ โข ภควา เอตสฺมิํ นิทาเน เอตสฺมิํ ปกรเณ ธมฺมิํ กถํ กตฺวา ภิกฺขู อามนฺเตสิ “น ภิกฺขเว ชานํ อุทฺทิสฺสกตํ มํสํ ปริภุญฺชิ\-ตพฺพํ, \csromanfolio{335}โย ปริภุญฺเชยฺย, อาปตฺติ ทุกฺกฏสฺส. อนุชานามิ ภิกฺขเว ติโกฏิ\-ปริสุทฺธํ มจฺฉมํสํ อทิฏฺฐํ อสฺสุตํ อปริสงฺกิตนฺ”ติ.}

\csromanmatikaanchor{mtk.185}
\csromantocmarkref{section}{179. กปฺปิยภูมิอนุชานนา}{mtk.185}
\bookmark[dest={mtk.185},level=1]{179. กปฺปิยภูมิอนุชานนา}
\csromanheader{1}{179. กปฺปิย\-ภูมิ\-อนุชานนา}
\proseitem{295}{เตน โข ปน สมเยน เวสาลี สุภิกฺขา โหติ สุสสฺสา สุลภปิณฺฑา สุกรา อุญฺเฉน ปคฺคเหน ยาเปตุํ. อถ โข ภควโต รโหคตสฺส ปฏิสลฺลีนสฺส เอวํ เจตโส ปริวิตกฺโก อุทปาทิ “ยานิ ตานิ มยา ภิกฺขูนํ อนุญฺญาตานิ ทุพฺภิกฺเข ทุสฺสสฺเส ทุลฺลภปิณฺเฑ อนฺโต วุฏฺฐํ อนฺโต ปกฺกํ สามํ ปกฺกํ อุคฺคหิตปฏิคฺคหิตกํ ตโต นีหฏํ ปุเรภตฺตํ ปฏิคฺคหิตํ วนฏฺฐํ โปกฺขรฏฺฐํ, อชฺชาปิ นุ โข ตานิ ภิกฺขู ปริภุญฺชนฺตี”ติ. อถ โข ภควา สายนฺหสมยํ ปฏิสลฺลานา วุฏฺฐิโต อายสฺมนฺตํ อานนฺทํ อามนฺเตสิ “ยานิ ตานิ อานนฺท มยา ภิกฺขูนํ อนุญฺญาตานิ ทุพฺภิกฺเข ทุสฺสสฺเส ทุลฺลภปิณฺเฑ อนฺโต วุฏฺฐํ อนฺโต ปกฺกํ สามํ ปกฺกํ อุคฺคหิตปฏิคฺคหิตกํ ตโต นีหฏํ ปุเรภตฺตํ ปฏิคฺคหิตํ วนฏฺฐํ โปกฺขรฏฺฐํ, อชฺชาปิ นุ โข ตานิ ภิกฺขู ปริภุญฺชนฺตี”ติ. ปริภุญฺชนฺติ ภควาติ. อถ โข ภควา เอตสฺมิํ นิทาเน เอตสฺมิํ ปกรเณ ธมฺมิํ กถํ กตฺวา ภิกฺขู อามนฺเตสิ “ยานิ ตานิ ภิกฺขเว มยา ภิกฺขูนํ อนุญฺญาตานิ ทุพฺภิกฺเข ทุสฺสสฺเส ทุลฺลภปิณฺเฑ อนฺโต วุฏฺฐํ อนฺโต ปกฺกํ สามํ ปกฺกํ อุคฺคหิตปฏิคฺคหิตกํ ตโต นีหฏํ ปุเรภตฺตํ ปฏิคฺคหิตํ วนฏฺฐํ โปกฺขรฏฺฐํ, ตานาหํ อชฺชตคฺเค ปฏิกฺขิปามิ, น ภิกฺขเว อนฺโต วุฏฺฐํ อนฺโต ปกฺกํ สามํ ปกฺกํ อุคฺคหิตปฏิคฺคหิตกํ ปริภุญฺชิ\-ตพฺพํ, โย ปริภุญฺเชยฺย, อาปตฺติ ทุกฺกฏสฺส. น จ ภิกฺขเว ตโต นีหฏํ ปุเรภตฺตํ ปฏิคฺคหิตํ วนฏฺฐํ โปกฺขรฏฺฐํ ภุตฺตาวินา ปวาริเตน อนติริตฺตํ ปริภุญฺชิ\-ตพฺพํ, โย ปริภุญฺเชยฺย, ยถาธมฺโม กาเรตพฺโพ”ติ.}

\prose{เตน โข สมเยน ชานปทา มนุสฺสา พหุํ โลณมฺปิ เตลมฺปิ ตณฺฑุลมฺปิ ขาทนียมฺปิ สกเฏสุ อาโรเปตฺวา พหา\-รามโกฏฺฐเก สกฏ\-ปริวฏฺฏํ กริตฺวา อจฺฉนฺติ “ยทา ปฏิปาฏิํ ลภิสฺสาม, ตทา ภตฺตํ กริสฺสามา”ติ, มหา จ เมโฆ อุคฺคโต โหติ. อถ โข เต มนุสฺสา เยนายสฺมา อานนฺโท เตนุ\-ปสงฺกมิํสุ, อุปสงฺกมิตฺวา อายสฺมนฺตํ อานนฺทํ เอตทโวจุํ “อิธ ภนฺเต อานนฺท พหุํ โลณมฺปิ เตลมฺปิ ตณฺฑุลมฺปิ ขาทนียมฺปิ สกเฏสุ อาโรปิตา ติฏฺฐนฺติ, มหา จ เมโฆ \csromanfolio{336}อุคฺคโต, กถํ นุ โข ภนฺเต อานนฺท ปฏิปชฺชิตพฺพนฺ”ติ. อถ โข อายสฺมา อานนฺโท ภควโต เอตมตฺถํ อาโรเจสิ. เตน หานนฺท สํโฆ ปจฺจนฺติมํ วิหารํ กปฺปิยภูมิํ สมฺมนฺนิตฺวา ตตฺถ วาเสตุ, ยํ สํโฆ อากงฺขติ วิหารํ วา อฑฺฒโยคํ วา ปาสาทํ วา หมฺมิยํ วา คุหํ วา. เอวญฺจ ปน ภิกฺขเว สมฺมนฺนิตพฺพา, พฺยตฺเตน ภิกฺขุนา ปฏิพเลน สํโฆ ญาเปตพฺโพ\csromandash{}}

\prose{“สุณาตุ เม ภนฺเต สํโฆ, ยทิ สํฆสฺส ปตฺตกลฺลํ สํโฆ อิตฺถนฺนามํ วิหารํ กปฺปิยภูมิํ สมฺมนฺเนยฺย, เอสา ญตฺติ.}

\prose{สุณาตุ เม ภนฺเต สํโฆ, สํโฆ อิตฺถนฺนามํ วิหารํ กปฺปิยภูมิํ สมฺมนฺนติ, ยสฺสายสฺมโต ขมติ อิตฺถนฺนามสฺส วิหารสฺส กปฺปิยภูมิยา สมฺมุติ, โส ตุณฺหสฺส, ยสฺส นกฺขมติ, โส ภาเสยฺย.}

\prose{สมฺมโต สํเฆน อิตฺถนฺนาโม วิหาโร กปฺปิยภูมิ, ขมติ สํฆสฺส, ตสฺมา ตุณฺหี, เอวเมตํ ธารยามี”ติ.}

\prose{เตน โข ปน สมเยน มนุสฺสา ตตฺเถว สมฺมุติยา\csromansharedfootnote{cfea2dfaacda5c66}{สมฺมติกาย (สฺยา)} กปฺปิยภูมิยา ยาคุโย ปจนฺติ, ภตฺตานิ ปจนฺติ, สูปานิ สมฺปาเทนฺติ, มํสานิ โกฏฺเฏนฺติ, กฏฺฐานิ ผาเลนฺติ.}

\prose{อสฺโสสิ โข ภควา รตฺติยา ปจฺจูสสมยํ ปจฺจุฏฺฐาย อุจฺจาสทฺทํ มหาสทฺทํ กาโกรวสทฺทํ, สุตฺวาน อายสฺมนฺตํ อานนฺทํ อามนฺเตสิ “กิํ นุ โข โส อานนฺท อุจฺจาสทฺโท มหาสทฺโท กาโกรวสทฺโท”ติ.}

\prose{เอตรหิ ภนฺเต มนุสฺสา ตตฺเถว สมฺมุติยา กปฺปิยภูมิยา ยาคุโย ปจนฺติ, ภตฺตานิ ปจนฺติ, สูปานิ สมฺปาเทนฺติ, มํสานิ โกฏฺเฏนฺติ, กฏฺฐานิ ผาเลนฺติ, โส เอโส ภควา อุจฺจาสทฺโท มหาสทฺโท กาโกรวสทฺโทติ.}

\prose{อถ โข ภควา เอตสฺมิํ นิทาเน เอตสฺมิํ ปกรเณ ธมฺมิํ กถํ กตฺวา ภิกฺขู อามนฺเตสิ “น ภิกฺขเว สมฺมุติ\csromansharedfootnote{61785d03c494c2ae}{[มิสฺสินฺคฺ โนเต 1]} กปฺปิยภูมิ ปริภุญฺชิ\-ตพฺพา, โย ปริภุญฺเชยฺย, อาปตฺติ ทุกฺกฏสฺส.}

\prose{อนุชานามิ ภิกฺขเว ติสฺโส กปฺปิยภูมิโย อุสฺสาวนนฺติกํ โคนิสาทิกํ คหปตินฺ”ติ.}

\prose{เตน โข ปน สมเยน อายสฺมา ยโสโช คิลาโน โหติ, ตสฺสตฺถาย เภสชฺชานิ อาหริยนฺติ, ตานิ ภิกฺขู พหิ วาเสนฺติ, อุกฺกปิณฺฑิกาปิ ขาทนฺติ, โจราปิ หรนฺติ.}

\setlength{\csromangathapagewidth}{0pt}
\setlength{\csromangathawakwidth}{0pt}
\csromangathameasuregroup{}{ภควโต เอตมตฺถํ อาโรเจสุํ.}
\csromangathameasurewak{ภควโต เอตมตฺถํ อาโรเจสุํ.}
\setlength{\csromangathaleftcol}{0pt}
\csromangathagroup{\csromangathastanza{\csromangathawak{ภควโต เอตมตฺถํ อาโรเจสุํ.}}}

\prosecontwithcloserruled{\csromanfolio{337}“อนุชานามิ ภิกฺขเว สมฺมุติํ กปฺปิยภูมิํ ปริภุญฺชิตุํ, อนุชานามิ ภิกฺขเว จตสฺโส กปฺปิยภูมิโย อุสฺสาวนนฺติกํ โคนิสาทิกํ คหปติํ สมฺมุตินฺ”ติ.}{สีหภาณวาโร นิฏฺฐิโต จตุตฺโถ.}
\csromanmatikaanchor{mtk.186}
\csromantocmarkref{section}{180. เมณฺฑกคหปติวตฺถุ}{mtk.186}
\bookmark[dest={mtk.186},level=1]{180. เมณฺฑกคหปติวตฺถุ}
\csromanheader{1}{180. \textbf{เมณฺฑก}\-\textbf{คหปติ}\-\textbf{วตฺถุ}}
\proseitem{296}{เตน โข ปน สมเยน ภทฺทิยนคเร เมณฺฑโก คหปติ ปฏิวสติ, ตสฺส เอวรูโป อิทฺธานุภาโว โหติ, สีสํ นหายิตฺวา ธญฺญาคารํ สมฺมชฺชาเปตฺวา พหิทฺวาเร นิสีทติ, อนฺตลิกฺขา ธญฺญสฺส ธารา โอปติตฺวา ธญฺญาคารํ ปูเรติ. ภริยาย เอวรูโป อิทฺธานุภาโว โหติ, เอกํเยว อาฬฺหกถาลิกํ อุปนิสีทิตฺวา เอกญฺจ สูปภิญฺชนกํ\csromansharedfootnote{e05b4221e6808826}{สูปภิญฺชรกํ (สี)} ทาส\-กมฺมกร\-โปริสํ ภตฺเตน ปริวิสติ, น ตาว ตํ ขิยฺยติ\csromansharedfootnote{f8838eb27d6dbe61}{ขียติ (สี, สฺยา)}, ยาว สา น วุฏฺฐาติ. ปุตฺตสฺส เอวรูโป อิทฺธานุภาโว โหติ, เอกํเยว สหสฺส\-ถวิกํ คเหตฺวา ทาส\-กมฺมกร\-โปริสสฺส ฉมาสิกํ เวตนํ เทติ, น ตาว ตํ ขิยฺยติ, ยาวสฺส หตฺถคตา. สุณิสาย เอวรูโป อิทฺธานุภาโว โหติ, เอกํเยว จตุโทณิกํ ปิฏกํ อุปนิสีทิตฺวา ทาส\-กมฺมกร\-โปริสสฺส ฉมาสิกํ ภตฺตํ เทติ, น ตาว ตํ ขิยฺยติ, ยาว สา น วุฏฺฐาติ. ทาสสฺส เอวรูโป อิทฺธานุภาโว โหติ, เอเกน นงฺคเลน กสนฺตสฺส สตฺต สีตาโย คจฺฉนฺติ.}

\prose{อสฺโสสิ โข ราชา มาคโธ เสนิโย พิมฺพิสาโร “อมฺหากํ กิร วิชิเต ภทฺทิยนคเร เมณฺฑโก คหปติ ปฏิวสติ, ตสฺส เอวรูโป อิทฺธานุภาโว, สีสํ นหายิตฺวา ธญฺญาคารํ สมฺมชฺชาเปตฺวา พหิทฺวาเร นิสีทติ, อนฺตลิกฺขา ธญฺญสฺส ธารา โอปติตฺวา ธญฺญาคารํ ปูเรติ. ภริยาย เอวรูโป อิทฺธานุภาโว, เอกํเยว อาฬฺหกถาลิกํ อุปนิสีทิตฺวา เอกญฺจ สูปภิญฺชนกํ ทาส\-กมฺมกร\-โปริสํ ภตฺเตน ปริวิสติ, น ตาว ตํ ขิยฺยติ, ยาว สา น วุฏฺฐาติ. ปุตฺตสฺส เอวรูโป อิทฺธานุภาโว, เอกํเยว สหสฺส\-ถวิกํ คเหตฺวา ทาส\-กมฺมกร\-โปริสสฺส ฉมาสิกํ เวตนํ เทติ, น ตาว ตํ ขิยฺยติ, ยาวสฺส หตฺถคตา. สุณิสาย \csromanfolio{338}เอวรูโป อิทฺธานุภาโว, เอกํเยว จตุโทณิกํ ปิฏกํ อุปนิสีทิตฺวา ทาส\-กมฺมกร\-โปริสสฺส ฉมาสิกํ ภตฺตํ เทติ, น ตาว ตํ ขิยฺยติ, ยาว สา น วุฏฺฐาติ. ทาสสฺส เอวรูโป อิทฺธานุภาโว, เอเกน นงฺคเลน กสนฺตสฺส สตฺต สีตาโย คจฺฉนฺตี”ติ. อถ โข ราชา มาคโธ เสนิโย พิมฺพิสาโร อญฺญตรํ สพฺพตฺถกํ มหามตฺตํ อามนฺเตสิ “อมฺหากํ กิร ภเณ วิชิเต ภทฺทิยนคเร เมณฺฑโก คหปติ ปฏิวสติ, ตสฺส เอวรูโป อิทฺธานุภาโว, สีสํ นหายิตฺวา ธญฺญาคารํ สมฺมชฺชาเปตฺวา พหิทฺวาเร นิสีทติ, อนฺตลิกฺขา ธญฺญสฺส ธารา โอปติตฺวา ธญฺญาคารํ ปูเรติ. ภริยาย ฯเปฯ. ปุตฺตสฺส. สุณิสาย. ทาสสฺส เอวรูโป อิทฺธานุภาโว, เอเกน นงฺคเลน กสนฺตสฺส สตฺต สีตาโย คจฺฉนฺตีติ. คจฺฉ ภเณ ชานาหิ, ยถา มยา สามํ ทิฏฺโฐ, เอวํ ตว ทิฏฺโฐ ภวิสฺสตี”ติ.}

\proseitem{297}{เอวํ เทวาติ โข โส มหามตฺโต รญฺโญ มาคธสฺส เสนิยสฺส พิมฺพิสารสฺส ปฏิสฺสุณิตฺวา จตุรงฺคินิยา เสนาย เยน ภทฺทิยํ เตน ปายาสิ. อนุปุพฺเพน เยน ภทฺทิยํ เยน เมณฺฑโก คหปติ เตนุปสงฺกมิ, อุปสงฺกมิตฺวา เมณฺฑกํ คหปติํ เอตทโวจ “อหํ หิ คหปติ รญฺญา อาณตฺโต ‘อมฺหากํ กิร ภเณ วิชิเต ภทฺทิยนคเร เมณฺฑโก คหปติ ปฏิวสติ, ตสฺส เอวรูโป อิทฺธานุภาโว, สีสํ นหายิตฺวา ฯเปฯ. ภริยาย. ปุตฺตสฺส. สุณิสาย. ทาสสฺส เอวรูโป อิทฺธานุภาโว, เอเกน นงฺคเลน กสนฺตสฺส สตฺต สีตาโย คจฺฉนฺตี’ติ, คจฺฉ ภเณ ชานาหิ, ยถา มยา สามํ ทิฏฺโฐ, เอวํ ตว ทิฏฺโฐ ภวิสฺสตี’ติ. ปสฺสาม เต คหปติ อิทฺธานุภาวนฺ”ติ. อถ โข เมณฺฑโก คหปติ สีสํ นหายิตฺวา ธญฺญาคารํ สมฺมชฺชาเปตฺวา พหิทฺวาเร นิสีทิ, อนฺตลิกฺขา ธญฺญสฺส ธารา โอปติตฺวา ธญฺญาคารํ ปูเรสิ. ทิฏฺโฐ เต คหปติ อิทฺธานุภาโว, ภริยาย เต อิทฺธานุภาวํ ปสฺสิสฺสามาติ. อถ โข เมณฺฑโก คหปติ ภริยํ อาณาเปสิ “เตน หิ จตุรงฺคินิํ เสนํ ภตฺเตน ปริวิสา”ติ. อถ โข เมณฺฑกสฺส คหปติสฺส ภริยา เอกํเยว อาฬฺหกถาลิกํ อุปนิสีทิตฺวา เอกญฺจ สูปภิญฺชนกํ จตุรงฺคินิํ เสนํ ภตฺเตน ปริวิสิ, น ตาว ตํ ขิยฺยติ, ยาว สา น วุฏฺฐาติ. ทิฏฺโฐ เต คหปติ ภริยายปิ อิทฺธานุภาโว, ปุตฺตสฺส เต อิทฺธานุภาวํ ปสฺสิสฺสามาติ. อถ โข เมณฺฑโก คหปติ ปุตฺตํ อาณาเปสิ “เตน หิ จตุรงฺคินิยา เสนาย ฉมาสิกํ เวตนํ \csromanfolio{339}เทหี”ติ. อถ โข เมณฺฑกสฺส คหปติสฺส ปุตฺโต เอกํเยว สหสฺส\-ถวิกํ คเหตฺวา จตุรงฺคินิยา เสนาย ฉมาสิกํ เวตนํ อทาสิ, น ตาว ตํ ขิยฺยติ, ยาวสฺส หตฺถคตา. ทิฏฺโฐ เต คหปติ ปุตฺตสฺสปิ อิทฺธานุภาโว, สุณิสาย เต อิทฺธานุภาวํ ปสฺสิสฺสามาติ. อถ โข เมณฺฑโก คหปติ สุณิสํ อาณาเปสิ “เตน หิ จตุรงฺคินิยา เสนาย ฉมาสิกํ ภตฺตํ เทหี”ติ. อถ โข เมณฺฑกสฺส คหปติสฺส สุณิสา เอกํเยว จตุโทณิกํ ปิฏกํ อุปนิสีทิตฺวา จตุรงฺคินิยา เสนาย ฉมาสิกํ ภตฺตํ อทาสิ, น ตาว ตํ ขิยฺยติ, ยาว สา น วุฏฺฐาติ. ทิฏฺโฐ เต คหปติ สุณิสายปิ อิทฺธานุภาโว, ทาสสฺส เต อิทฺธานุภาวํ ปสฺสิสฺสามาติ. มยฺหํ โข สามิ ทาสสฺส อิทฺธานุภาโว เขตฺเต ปสฺสิตพฺโพติ. อลํ คหปติ ทิฏฺโฐ เต ทาสสฺสปิ อิทฺธานุภาโวติ. อถ โข โส มหามตฺโต จตุรงฺคินิยา เสนาย ปุนเทว ราชคหํ ปจฺจาคญฺฉิ. เยน ราชา มาคโธ เสนิโย พิมฺพิสาโร เตนุปสงฺกมิ, อุปสงฺกมิตฺวา รญฺโญ มาคธสฺส เสนิยสฺส พิมฺพิสารสฺส เอตมตฺถํ อาโรเจสิ.}

\proseitem{298}{อถ โข ภควา เวสาลิยํ ยถาภิรนฺตํ วิหริตฺวา เยน ภทฺทิยํ เตน จาริกํ ปกฺกามิ มหตา ภิกฺขุ\-สํเฆน สทฺธิํ อฑฺฒเตลเสหิ ภิกฺขุสเตหิ. อถ โข ภควา อนุปุพฺเพน จาริกํ จรมาโน เยน ภทฺทิยํ ตทวสริ, ตตฺร สุทํ ภควา ภทฺทิเย วิหรติ ชาติยา วเน. อสฺโสสิ โข เมณฺฑโก คหปติ “สมโณ ขลุ โภ โคตโม สกฺยปุตฺโต สกฺยกุลา ปพฺพชิโต ภทฺทิยํ อนุปฺปตฺโต ภทฺทิเย วิหรติ ชาติยา วเน, ตํ โข ปน ภควนฺตํ โคตมํ เอวํ กลฺยาโณ กิตฺติสทฺโท อพฺภุคฺคโต ‘อิติปิ โส ภควา อรหํ สมฺมาสมฺพุทฺโธ วิชฺชา\-จรณ\-สมฺปนฺโน สุคโต โลกวิทู อนุตฺตโร ปุริส\-ทมฺม\-สารถิ สตฺถา เทวมนุสฺสานํ พุทฺโธ ภควา’\csromansharedfootnote{475b0b93102c8b46}{ภควาติ (ก)}. โส อิมํ โลกํ สเทวกํ สมารกํ สพฺรหฺมกํ สสฺสมณ\-พฺราหฺมณิํ ปชํ สเทวมนุสฺสํ สยํ อภิญฺญา สจฺฉิกตฺวา ปเวเทติ, โส ธมฺมํ เทเสติ อาทิกลฺยาณํ มชฺเฌกลฺยาณํ ปริโยสาน\-กลฺยาณํ สาตฺถํ สพฺยญฺชนํ เกวล\-ปริปุณฺณํ ปริสุทฺธํ พฺรหฺมจริยํ ปกาเสติ, สาธุ โข ปน ตถารูปานํ อรหตํ ทสฺสนํ โหตี”ติ. อถ โข เมณฺฑโก คหปติ ภทฺรานิ ภทฺรานิ ยานานิ โยชาเปตฺวา ภทฺรํ ภทฺรํ ยานํ อภิรุหิตฺวา ภเทฺรหิ ภเทฺรหิ \csromanfolio{340}ยาเนหิ ภทฺทิยา นิยฺยาสิ ภควนฺตํ ทสฺสนาย. อทฺทสํสุ โข สมฺพหุลา ติตฺถิยา เมณฺฑกํ คหปติํ ทูรโตว อาคจฺฉนฺตํ, ทิสฺวาน เมณฺฑกํ คหปติํ เอตทโวจุํ “กหํ ตฺวํ คหปติ คจฺฉสี”ติ. คจฺฉามหํ ภนฺเต ภควนฺตํ\csromansharedfootnote{0aaa8df640f92eb8}{อิทํ ปทํ สีสฺยาโปตฺถเกสุ นตฺถิ.} สมณํ โคตมํ ทสฺสนายาติ. กิํ ปน ตฺวํ คหปติ กิริยวาโท สมาโน อกิริยวาทํ สมณํ โคตมํ ทสฺสนาย อุปสงฺกมิสฺสสิ, สมโณ หิ คหปติ โคตโม อกิริยวาโท อกิริยาย ธมฺมํ เทเสติ, เตน จ สาวเก วิเนตีติ. อถ โข เมณฺฑกสฺส คหปติสฺส เอตทโหสิ “นิสฺสํสยํ โข โส ภควา อรหํ สมฺมาสมฺพุทฺโธ ภวิสฺสติ, ยถยิเม ติตฺถิยา อุสูยนฺตี”ติ, ยาวติกา ยานสฺส ภูมิ, ยาเนน คนฺตฺวา ยานา ปจฺโจโรหิตฺวา ปตฺติโกว เยน ภควา เตนุปสงฺกมิ, อุปสงฺกมิตฺวา ภควนฺตํ อภิวาเทตฺวา เอกมนฺตํ นิสีทิ. เอกมนฺตํ นิสินฺนสฺส โข เมณฺฑกสฺส คหปติสฺส ภควา อนุปุพฺพิํ กถํ กเถสิ. เสยฺยถิทํ, ทานกถํ ฯเปฯ อปรปฺปจฺจโย สตฺถุสาสเน ภควนฺตํ เอตทโวจ “อภิกฺกนฺตํ ภนฺเต ฯเปฯ อุปาสกํ มํ ภควา ธาเรตุ อชฺชตคฺเค ปาณุเปตํ สรณํ คตํ, อธิวาเสตุ จ เม ภนฺเต ภควา สฺวาตนาย ภตฺตํ สทฺธิํ ภิกฺขุสํเฆนา”ติ. อธิวาเสสิ ภควา ตุณฺหีภาเวน. อถ โข เมณฺฑโก คหปติ ภควโต อธิวาสนํ วิทิตฺวา อุฏฺฐายาสนา ภควนฺตํ อภิวาเทตฺวา ปทกฺขิณํ กตฺวา ปกฺกามิ.}

\prose{อถ โข เมณฺฑโก คหปติ ตสฺสา รตฺติยา อจฺจเยน ปณีตํ ขาทนียํ โภชนียํ ปฏิยาทาเปตฺวา ภควโต กาลํ อาโรจาเปสิ “กาโล ภนฺเต นิฏฺฐิตํ ภตฺตนฺ”ติ. อถ โข ภควา ปุพฺพณฺหสมยํ นิวาเสตฺวา ปตฺต\-จีวรมา\-ทาย เยน เมณฺฑกสฺส คหปติสฺส นิเวสนํ เตนุปสงฺกมิ, อุปสงฺกมิตฺวา ปญฺญตฺเต อาสเน นิสีทิ สทฺธิํ ภิกฺขุ\-สํเฆน. อถ โข เมณฺฑกสฺส คหปติสฺส ภริยา จ ปุตฺโต จ สุณิสา จ ทาโส จ เยน ภควา เตนุ\-ปสงฺกมิํสุ, อุปสงฺกมิตฺวา ภควนฺตํ อภิวาเทตฺวา เอกมนฺตํ นิสีทิํสุ. เตสํ ภควา อนุปุพฺพิํ กถํ กเถสิ. เสยฺยถิทํ, ทานกถํ ฯเปฯ อปรปฺปจฺจยา สตฺถุสาสเน ภควนฺตํ เอตทโวจุํ “อภิกฺกนฺตํ ภนฺเต ฯเปฯ เอเต มยํ ภนฺเต ภควนฺตํ สรณํ คจฺฉาม ธมฺมญฺจ ภิกฺขุสํฆญฺจ, อุปาสเก โน ภควา ธาเรตุ อชฺชตคฺเค ปาณุเปเต \csromanfolio{341}สรณํ คเต”ติ. อถ โข เมณฺฑโก คหปติ พุทฺธ\-ปฺปมุขํ ภิกฺขุสํฆํ ปณีเตน ขาทนีเยน โภชนีเยน สหตฺถา สนฺตปฺเปตฺวา สมฺปวาเรตฺวา ภควนฺตํ ภุตฺตาวิํ โอนีต\-ปตฺต\-ปาณิํ เอกมนฺตํ นิสีทิ, เอกมนฺตํ นิสินฺโน โข เมณฺฑโก คหปติ ภควนฺตํ เอตทโวจ “ยาว ภนฺเต ภควา ภทฺทิเย วิหรติ, ตาว อหํ พุทฺธ\-ปฺปมุขสฺส ภิกฺขุ\-สํฆสฺส ธุวภตฺเตนา”ติ. อถ โข ภควา เมณฺฑกํ คหปติํ ธมฺมิยา กถาย สนฺทสฺเสตฺวา สมาทเปตฺวา สมุตฺเตเชตฺวา สมฺปหํเสตฺวา อุฏฺฐายาสนา ปกฺกามิ.}

\csromanmatikaanchor{mtk.187}
\csromantocmarkref{section}{181. ปญฺจโครสาทิอนุชานนา}{mtk.187}
\bookmark[dest={mtk.187},level=1]{181. ปญฺจโครสาทิอนุชานนา}
\csromanheader{1}{181. ปญฺจโครสา\-ทิ\-อนุชานนา}
\proseitem{299}{อถ โข ภควา ภทฺทิเย ยถาภิรนฺตํ วิหริตฺวา เมณฺฑกํ คหปติํ อนาปุจฺฉา เยน องฺคุตฺตราโป เตน จาริกํ ปกฺกามิ มหตา ภิกฺขุ\-สํเฆน สทฺธิํ อฑฺฒเตลเสหิ ภิกฺขุสเตหิ. อสฺโสสิ โข เมณฺฑโก คหปติ “ภควา กิร เยน องฺคุตฺตราโป เตน จาริกํ ปกฺกนฺโต มหตา ภิกฺขุ\-สํเฆน สทฺธิํ อฑฺฒเตลเสหิ ภิกฺขุสเตหี”ติ. อถ โข เมณฺฑโก คหปติ ทาเส จ กมฺมกเร จ อาณาเปสิ “เตน หิ ภเณ พหุํ โลณมฺปิ เตลมฺปิ ตณฺฑุลมฺปิ ขาทนียมฺปิ สกเฏสุ อาโรเปตฺวา อาคจฺฉถ, อฑฺฒเตลสานิ จ โคปาลกสตานิ อฑฺฒเตลสานิ จ เธนุสตานิ อาทาย อาคจฺฉนฺตุ, ยตฺถ ภควนฺตํ ปสฺสิสฺสาม, ตตฺถ ตรุเณน\csromansharedfootnote{fe828f33089055b5}{ธารุเณฺหน (สี, สฺยา)} ขีเรน โภเชสฺสามา”ติ. อถ โข เมณฺฑโก คหปติ ภควนฺตํ อนฺตรามคฺเค กนฺตาเร สมฺภาเวสิ. อถ โข เมณฺฑโก คหปติ เยน ภควา เตนุปสงฺกมิ, อุปสงฺกมิตฺวา ภควนฺตํ อภิวาเทตฺวา เอกมนฺตํ อฏฺฐาสิ. เอกมนฺตํ ฐิโต โข เมณฺฑโก คหปติ ภควนฺตํ เอตทโวจ “อธิวาเสตุ เม ภนฺเต ภควา สฺวาตนาย ภตฺตํ สทฺธิํ ภิกฺขุ\-สํเฆนา”ติ. อธิวาเสสิ ภควา ตุณฺหีภาเวน. อถ โข เมณฺฑโก คหปติ ภควโต อธิวาสนํ วิทิตฺวา ภควนฺตํ อภิวาเทตฺวา ปทกฺขิณํ กตฺวา ปกฺกามิ.}

\prose{อถ โข เมณฺฑโก คหปติ ตสฺสา รตฺติยา อจฺจเยน ปณีตํ ขาทนียํ โภชนียํ ปฏิยาทาเปตฺวา ภควโต กาลํ อาโรจาเปสิ “กาโล ภนฺเต นิฏฺฐิตํ ภตฺตนฺ”ติ. อถ โข ภควา ปุพฺพณฺหสมยํ นิวาเสตฺวา ปตฺต\-จีวรมา\-ทาย เยน เมณฺฑกสฺส คหปติสฺส ปริเวสนา \csromanfolio{342}เตนุปสงฺกมิ, อุปสงฺกมิตฺวา ปญฺญตฺเต อาสเน นิสีทิ สทฺธิํ ภิกฺขุ\-สํเฆน. อถ โข เมณฺฑโก คหปติ อฑฺฒเตลสานิ โคปาลกสตานิ อาณาเปสิ “เตนหิ\csromansharedfootnote{cb56481398405b41}{เตน หิ เทฺว ปทานิ \csromandash{} ม.พ.ป.} ภเณ เอกเมกํ เธนุํ คเหตฺวา เอกเมกสฺส ภิกฺขุโน อุปติฏฺฐถ ตรุเณน ขีเรน โภเชสฺสามา”ติ. อถ โข เมณฺฑโก คหปติ พุทฺธ\-ปฺปมุขํ ภิกฺขุสํฆํ ปณีเตน ขาทนีเยน โภชนีเยน สหตฺถา สนฺตปฺเปสิ สมฺปวาเรสิ, ตรุเณน จ ขีเรน. ภิกฺขู กุกฺกุจฺจายนฺตา ขีรํ น ปฏิคฺคณฺหนฺติ. ปฏิคฺคณฺหถ ภิกฺขเว ปริภุญฺช\-ถาติ. อถ โข เมณฺฑโก คหปติ พุทฺธ\-ปฺปมุขํ ภิกฺขุสํฆํ ปณีเตน ขาทนีเยน โภชนีเยน สหตฺถา สนฺตปฺเปตฺวา สมฺปวาเรตฺวา ตรุเณน จ ขีเรน ภควนฺตํ ภุตฺตาวิํ โอนีต\-ปตฺต\-ปาณิํ เอกมนฺตํ นิสีทิ, เอกมนฺตํ นิสินฺโน โข เมณฺฑโก คหปติ ภควนฺตํ เอตทโวจ “สนฺติ ภนฺเต มคฺคา กนฺตารา อปฺโปทกา อปฺปภกฺขา, น สุกรา อปาเถยฺเยน คนฺตุํ, สาธุ ภนฺเต ภควา ภิกฺขูนํ ปาเถยฺยํ อนุชานาตู”ติ. อถ โข ภควา เมณฺฑกํ คหปติํ ธมฺมิยา กถาย สนฺทสฺเสตฺวา สมาทเปตฺวา สมุตฺเตเชตฺวา สมฺปหํเสตฺวา อุฏฺฐายาสนา ปกฺกามิ. อถ โข ภควา เอตสฺมิํ นิทาเน เอตสฺมิํ ปกรเณ ธมฺมิํ กถํ กตฺวา ภิกฺขู อามนฺเตสิ “อนุชานามิ ภิกฺขเว ปญฺจ โครเส ขีรํ ทธิํ ตกฺกํ นวนีตํ สปฺปิํ. สนฺติ ภิกฺขเว มคฺคา กนฺตารา อปฺโปทกา อปฺปภกฺขา, น สุกรา อปาเถยฺเยน คนฺตุํ, อนุชานามิ ภิกฺขเว ปาเถยฺยํ ปริเยสิตุํ ตณฺฑุโล ตณฺฑุล\-ตฺถิเกน มุคฺโค มุคฺคตฺถิเกน, มาโส มาสตฺถิเกน, โลณํ โลณตฺถิเกน, คุโฬ คุฬตฺถิเกน, เตลํ เตลตฺถิเกน, สปฺปิ สปฺปิตฺถิเกน. สนฺติ ภิกฺขเว มนุสฺสา สทฺธา ปสนฺนา, เต กปฺปิย\-การกานํ หตฺเถ หิรญฺญํ อุปนิกฺขิปนฺติ ‘อิมินา อยฺยสฺส ยํ กปฺปิยํ ตํ เทถา’ติ, อนุชานามิ ภิกฺขเว ยํ ตโต กปฺปิยํ ตํ สาทิตุํ, น เตฺววาหํ ภิกฺขเว ‘เกนจิ ปริยาเยน ชาตรูป\-รชตํ สาทิตพฺพํ ปริเยสิตพฺพนฺ’ติ วทามี”ติ.}

\csromanmatikaanchor{mtk.188}
\csromantocmarkref{section}{182. เกณิยชฏิลวตฺถุ}{mtk.188}
\bookmark[dest={mtk.188},level=1]{182. เกณิยชฏิลวตฺถุ}
\csromanheader{1}{182. \textbf{เกณิย}\-\textbf{ชฏิล}\-\textbf{วตฺถุ}}
\proseitem{300}{\csromansymbolfootnote{*}{ม 2. 347 ปิฏฺฐาทีสุ; ขุ 1. 363 ปิฏฺฐาทีสุ จ อาคตํ เสลสุตฺตมฺปิ ปสฺสิตพฺพํ.}อถ โข ภควา อนุปุพฺเพน จาริกํ จรมาโน เยน อาปณํ ตทวสริ, อสฺโสสิ โข เกณิโย ชฏิโล “สมโณ ขลุ โภ โคตโม สกฺยปุตฺโต สกฺยกุลา ปพฺพชิโต อาปณํ อนุปฺปตฺโต, ตํ โข ปน ภวนฺตํ โคตมํ เอวํ กลฺยาโณ กิตฺติสทฺโท อพฺภุคฺคโต ฯเปฯ สาธุ โข ปน ตถารูปานํ อรหตํ ทสฺสนํ โหตี”ติ. อถ โข \csromanfolio{343}เกณิยสฺส ชฏิลสฺส เอตทโหสิ “กิํ นุ โข อหํ สมณสฺส โคตมสฺส หราเปยฺยนฺ”ติ. อถ โข เกณิยสฺส ชฏิลสฺส เอตทโหสิ “เยปิ โข เต\csromansymbolfootnote{*}{* อยํ ปาโฐ ที 1. 97, 224, 228; ม 2. 381; อํ 2. 197, 202 ปิฏฺเฐสุปิ อาคโต.}พฺราหฺมณานํ ปุพฺพกา อิสโย มนฺตานํ กตฺตาโร มนฺตานํ ปวตฺตาโร, เยสมิทํ เอตรหิ พฺราหฺมณา โปราณํ มนฺตปทํ คีตํ ปวุตฺตํ สมิหิตํ, ตทนุคายนฺติ, ตทนุภาสนฺติ, ภาสิตมนุภาสนฺติ, วาจิตมนุวาเจนฺติ. เสยฺยถิทํ, อฏฺฐโก วามโก วามเทโว เวสฺสามิตฺโต ยมตคฺคิ\csromansharedfootnote{9873a359e3814ad3}{ยมทคฺคิ (ก)} องฺคีรโส ภารทฺวาโช วาเสฏฺโฐ กสฺสโป ภคุ *, รตฺตูปรตา วิรตา วิกาลโภชนา, เต เอวรูปานิ ปานานิ สาทิยิํสุ. สมโณปิ โคตโม รตฺตูปรโต วิรโต วิกาลโภชนา, อรหติ สมโณปิ โคตโม เอวรูปานิ ปานานิ สาทิยิตุนฺ”ติ ปหูตํ ปานํ ปฏิยาทาเปตฺวา กาเชหิ คาหาเปตฺวา เยน ภควา เตนุปสงฺกมิ, อุปสงฺกมิตฺวา ภควตา สทฺธิํ สมฺโมทิ, สมฺโมทนียํ กถํ สารณียํ วีติสาเรตฺวา เอกมนฺตํ อฏฺฐาสิ, เอกมนฺตํ ฐิโต โข เกณิโย ชฏิโล ภควนฺตํ เอตทโวจ “ปฏิคฺคณฺหาตุ เม ภวํ โคตโม ปานนฺ”ติ. เตน หิ เกณิย ภิกฺขูนํ เทหีติ. อถ โข เกณิโย ชฏิโล ภิกฺขูนํ เทติ. ภิกฺขู กุกฺกุจฺจายนฺตา น ปฏิคฺคณฺหนฺติ. ปฏิคฺคณฺหถ ภิกฺขเว ปริภุญฺช\-ถาติ. อถ โข เกณิโย ชฏิโล พุทฺธ\-ปฺปมุขํ ภิกฺขุสํฆํ ปหูเตหิ ปาเนหิ สหตฺถา สนฺตปฺเปตฺวา สมฺปวาเรตฺวา ภควนฺตํ โธตหตฺถํ โอนีต\-ปตฺต\-ปาณิํ เอกมนฺตํ นิสีทิ. เอกมนฺตํ นิสินฺนํ โข เกณิยํ ชฏิลํ ภควา ธมฺมิยา กถาย สนฺทสฺเสสิ สมาทเปสิ สมุตฺเตเชสิ สมฺปหํเสสิ. อถ โข เกณิโย ชฏิโล ภควตา ธมฺมิยา กถาย สนฺทสฺสิโต สมาทปิโต สมุตฺเตชิโต สมฺปหํสิโต ภควนฺตํ เอตทโวจ “อธิวาเสตุ เม ภวํ โคตโม สฺวาตนาย ภตฺตํ สทฺธิํ ภิกฺขุสํเฆนา”ติ. มหา โข เกณิย ภิกฺขุสํโฆ อฑฺฒเตลสานิ ภิกฺขุสตานิ, ตฺวญฺจ พฺราหฺมเณสุ อภิปฺปสนฺโนติ. ทุติยมฺปิ โข เกณิโย ชฏิโล ภควนฺตํ เอตทโวจ “กิญฺจาปิ โข โภ โคตม มหา ภิกฺขุสํโฆ อฑฺฒเตลสานิ ภิกฺขุสตานิ, อหญฺจ พฺราหฺมเณสุ อภิปฺปสนฺโน, อธิวาเสตุ เม ภวํ โคตโม สฺวาตนาย ภตฺตํ สทฺธิํ ภิกฺขุสํเฆนา”ติ. มหา โข เกณิย ภิกฺขุสํโฆ อฑฺฒเตลสานิ ภิกฺขุสตานิ, ตฺวญฺจ พฺราหฺมเณสุ \csromanfolio{344}อภิปฺปสนฺโนติ. ตติยมฺปิ โข เกณิโย ชฏิโล ภควนฺตํ เอตทโวจ “กิญฺจาปิ โข โภ โคตม มหา ภิกฺขุสํโฆ อฑฺฒเตลสานิ ภิกฺขุสตานิ, อหญฺจ พฺราหฺมเณสุ อภิปฺปสนฺโน, อธิวาเสตุ ภวํ โคตโม สฺวาตนาย ภตฺตํ สทฺธิํ ภิกฺขุสํเฆนา”ติ. อธิวาเสสิ ภควา ตุณฺหีภาเวน. อถ โข เกณิโย ชฏิโล ภควโต อธิวาสนํ วิทิตฺวา อุฏฺฐายาสนา ปกฺกามิ. อถ โข ภควา เอตสฺมิํ นิทาเน เอตสฺมิํ ปกรเณ ธมฺมิํ กถํ กตฺวา ภิกฺขู อามนฺเตสิ “อนุชานามิ ภิกฺขเว อฏฺฐ ปานานิ อมฺพปานํ ชมฺพุปานํ โจจปานํ โมจปานํ มธูกปานํ\csromansharedfootnote{baa2221917b667a7}{มธุปานํ (สี, สฺยา)} มุทฺทิกปานํ สาลูกปานํ ผารุสกปานํ. อนุชานามิ ภิกฺขเว สพฺพํ ผลรสํ ฐเปตฺวา ธญฺญ\-ผลรสํ, อนุชานามิ ภิกฺขเว สพฺพํ ปตฺตรสํ ฐเปตฺวา ฑากรสํ, อนุชานามิ ภิกฺขเว สพฺพํ ปุปฺผรสํ ฐเปตฺวา มธูก\-ปุปฺผรสํ, อนุชานามิ ภิกฺขเว อุจฺฉุรสนฺ”ติ.}

\prose{อถ โข เกณิโย ชฏิโล ตสฺสา รตฺติยา อจฺจเยน สเก อสฺสเม ปณีตํ ขาทนียํ โภชนียํ ปฏิยาทาเปตฺวา ภควโต กาลํ อาโรจาเปสิ “กาโล โภ โคตม นิฏฺฐิตํ ภตฺตนฺ”ติ. อถ โข ภควา ปุพฺพณฺหสมยํ นิวาเสตฺวา ปตฺต\-จีวรมา\-ทาย เยน เกณิยสฺส ชฏิลสฺส อสฺสโม เตนุปสงฺกมิ, อุปสงฺกมิตฺวา ปญฺญตฺเต อาสเน นิสีทิ สทฺธิํ ภิกฺขุ\-สํเฆน. อถ โข เกณิโย ชฏิโล พุทฺธ\-ปฺปมุขํ ภิกฺขุสํฆํ ปณีเตน ขาทนีเยน โภชนีเยน สหตฺถา สนฺตปฺเปตฺวา สมฺปวาเรตฺวา ภควนฺตํ ภุตฺตาวิํ โอนีต\-ปตฺต\-ปาณิํ เอกมนฺตํ นิสีทิ. เอกมนฺตํ นิสินฺนํ โข เกณิยํ ชฏิลํ ภควา อิมาหิ คาถาหิ อนุโมทิ\csromandash{}}

\setlength{\csromangathapagewidth}{0pt}
\setlength{\csromangathawakwidth}{0pt}
\csromangathameasuregroup{อคฺคิหุตฺตมุขา ยญฺญา, \\ ราชา มุขํ มนุสฺสานํ, \\ นกฺขตฺตานํ มุขํ จนฺโท, \\ ปุญฺญํ อากงฺขมานานํ,}{\csromangathabat{อคฺคิหุตฺตมุขา ยญฺญา,}{สาวิตฺตี ฉนฺทโส มุขํ.} \\ \csromangathabat{ราชา มุขํ มนุสฺสานํ,}{นทีนํ สาคโร มุขํ.} \\ \csromangathabat{นกฺขตฺตานํ มุขํ จนฺโท,}{อาทิจฺโจ ตปตํ มุขํ.} \\ \csromangathabat{ปุญฺญํ อากงฺขมานานํ,}{สํโฆ เว ยชตํ มุขนฺติ.}}
\csromangathasetleft{อคฺคิหุตฺตมุขา ยญฺญา, \\ ราชา มุขํ มนุสฺสานํ, \\ นกฺขตฺตานํ มุขํ จนฺโท, \\ ปุญฺญํ อากงฺขมานานํ,}
\csromangathagroup{\csromangathastanza{\csromangathabat{อคฺคิหุตฺตมุขา ยญฺญา,}{สาวิตฺตี ฉนฺทโส มุขํ.} \\ \csromangathabat{ราชา มุขํ มนุสฺสานํ,}{นทีนํ สาคโร มุขํ.}}\csromangathastanza{\csromangathabat{นกฺขตฺตานํ มุขํ จนฺโท,}{อาทิจฺโจ ตปตํ มุขํ.} \\ \csromangathabat{ปุญฺญํ อากงฺขมานานํ,}{สํโฆ เว ยชตํ มุขนฺติ.}}}

\prose{อถ โข ภควา เกณิยํ ชฏิลํ อิมาหิ คาถาหิ อนุโมทิตฺวา อุฏฺฐายาสนา ปกฺกามิ.}

\csromanmatikaanchor{mtk.189}
\csromantocmarkref{section}{183. โรชมลฺลวตฺถุ}{mtk.189}
\bookmark[dest={mtk.189},level=1]{183. โรชมลฺลวตฺถุ}
\csromanheader{1}{183. \textbf{โรช}\-\textbf{มลฺล}\-\textbf{วตฺถุ}}
\proseitem{301}{\csromanfolio{345}อถ โข ภควา อาปเณ ยถาภิรนฺตํ วิหริตฺวา เยน กุสินารา เตน จาริกํ ปกฺกามิ มหตา ภิกฺขุ\-สํเฆน สทฺธิํ อฑฺฒเตลเสหิ ภิกฺขุสเตหิ. อสฺโสสุํ โข โกสินารกา มลฺลา “ภควา กิร กุสินารํ อาคจฺฉติ มหตา ภิกฺขุ\-สํเฆน สทฺธิํ อฑฺฒเตลเสหิ ภิกฺขุสเตหี”ติ. เต สงฺครํ\csromansharedfootnote{cb2a250d03e84819}{สงฺกรํ (ก)} อกํสุ “โย ภควโต ปจฺจุคฺคมนํ น กริสฺสติ, ปญฺจสตานิสฺส ทณฺโฑ”ติ. เตน โข ปน สมเยน โรโช มลฺโล อายสฺมโต อานนฺทสฺส สหาโย โหติ. อถ โข ภควา อนุปุพฺเพน จาริกํ จรมาโน เยน กุสินารา ตทวสริ. อถ โข โกสินารกา มลฺลา ภควโต ปจฺจุคฺคมนํ อกํสุ. อถ โข โรโช มลฺโล ภควโต ปจฺจุคฺคมนํ กริตฺวา เยนายสฺมา อานนฺโท เตนุปสงฺกมิ, อุปสงฺกมิตฺวา อายสฺมนฺตํ อานนฺทํ อภิวาเทตฺวา เอกมนฺตํ อฏฺฐาสิ. เอกมนฺตํ ฐิตํ โข โรชํ มลฺลํ อายสฺมา อานนฺโท เอตทโวจ “อุฬารํ โข เต อิทํ อาวุโส โรช, ยํ ตฺวํ ภควโต ปจฺจุคฺคมนํ อกาสี”ติ. นาหํ ภนฺเต อานนฺท พหุกโต พุทฺเธ วา ธมฺเม วา สํเฆ วา, อปิ จ ญาตีหิ สงฺคโร กโต “โย ภควโต ปจฺจุคฺคมนํ น กริสฺสติ, ปญฺจสตานิสฺส ทณฺโฑ”ติ, โส โข อหํ ภนฺเต อานนฺท ญาตีนํ ทณฺฑภยา เอวาหํ ภควโต ปจฺจุคฺคมนํ อกาสินฺติ. อถ โข อายสฺมา อานนฺโท อนตฺตมโน อโหสิ “กถํ หิ นาม โรโช มลฺโล เอวํ วกฺขตี”ติ. อถ โข อายสฺมา อานนฺโท เยน ภควา เตนุปสงฺกมิ, อุปสงฺกมิตฺวา ภควนฺตํ อภิวาเทตฺวา เอกมนฺตํ นิสีทิ, เอกมนฺตํ นิสินฺโน โข อายสฺมา อานนฺโท ภควนฺตํ เอตทโวจ “อยํ ภนฺเต โรโช มลฺโล อภิญฺญาโต ญาตมนุสฺโส, มหตฺถิโก โข ปน เอวรูปานํ ญาตมนุสฺสานํ อิมสฺมิํ ธมฺมวินเย ปสาโท, สาธุ ภนฺเต ภควา ตถา กโรตุ, ยถา โรโช มลฺโล อิมสฺมิํ ธมฺมวินเย ปสีเทยฺยา”ติ. น โข ตํ อานนฺท ทุกฺกรํ ตถาคเตน, ยถา โรโช มลฺโล อิมสฺมิํ ธมฺมวินเย ปสีเทยฺยาติ.}

\prose{อถ โข ภควา โรชํ มลฺลํ เมตฺเตน จิตฺเตน ผริตฺวา อุฏฺฐายาสนา วิหารํ ปาวิสิ. อถ โข โรโช มลฺโล ภควโต เมตฺเตน จิตฺเตน ผุฏฺโฐ เสยฺยถาปิ นาม คาวิํ ตรุณวจฺโฉ, เอวเมว \csromanfolio{346}วิหาเรน วิหารํ ปริเวเณน ปริเวณํ อุปสงฺกมิตฺวา ภิกฺขู ปุจฺฉติ “กหํ นุ โข ภนฺเต เอตรหิ โส ภควา วิหรติ อรหํ สมฺมาสมฺพุทฺโธ, ทสฺสนกามา หิ มยํ ตํ ภควนฺตํ อรหนฺตํ สมฺมาสมฺพุทฺธนฺ”ติ. เอสาวุโส โรช วิหาโร สํวุตทฺวาโร, เตน อปฺปสทฺโท อุปสงฺกมิตฺวา อตรมาโน อาฬินฺทํ ปวิสิตฺวา อุกฺกาสิตฺวา อคฺคฬํ อาโกเฏหิ, วิวริสฺสติ เต ภควา ทฺวารนฺติ. อถ โข โรโช มลฺโล เยน โส วิหาโร สํวุตทฺวาโร, เตน อปฺปสทฺโท อุปสงฺกมิตฺวา อตรมาโน อาฬินฺทํ ปวิสิตฺวา อุกฺกาสิตฺวา อคฺคฬํ อาโกเฏสิ. วิวริ ภควา ทฺวารํ. อถ โข โรโช มลฺโล วิหารํ ปวิสิตฺวา ภควนฺตํ อภิวาเทตฺวา เอกมนฺตํ นิสีทิ. เอกมนฺตํ นิสินฺนสฺส โข โรชสฺส มลฺลสฺส ภควา อนุปุพฺพิํ กถํ กเถสิ. เสยฺยถิทํ, ทานกถํ ฯเปฯ อปรปฺปจฺจโย สตฺถุสาสเน ภควนฺตํ เอตทโวจ “สาธุ ภนฺเต อยฺยา มมญฺเญว ปฏิคฺคเณฺหยฺยุํ จีวร\-ปิณฺฑปาต\-เสนาสน\-คิลานปฺปจฺจย\-เภสชฺช\-ปริกฺขารํ, โน อญฺเญสนฺ”ติ. เยสํ โข โรช เสกฺเขน ญาเณน เสกฺเขน ทสฺสเนน ธมฺโม ทิฏฺโฐ เสยฺยถาปิ ตยา, เตสมฺปิ เอวํ โหติ “อโห นูน อยฺยา อมฺหากญฺเญว ปฏิคฺคเณฺหยฺยุํ จีวร\-ปิณฺฑปาต\-เสนาสน\-คิลานปฺปจฺจย\-เภสชฺช\-ปริกฺขารํ, โน อญฺเญสนฺ”ติ. เตน หิ โรช ตว เจว ปฏิคฺคหิสฺสนฺติ อญฺเญสญฺจาติ.}

\proseitem{302}{เตน โข ปน สมเยน กุสินารายํ ปณีตานํ ภตฺตานํ ภตฺตปฏิปาฏิ อฏฺฐิตา โหติ. อถ โข โรชสฺส มลฺลสฺส ปฏิปาฏิํ อลภนฺตสฺส เอตทโหสิ “ยนฺนูนาหํ ภตฺตคฺคํ โอโลเกยฺยํ, ยํ ภตฺตคฺเค นาสฺส, ตํ ปฏิยาเทยฺยนฺ”ติ. อถ โข โรโช มลฺโล ภตฺตคฺคํ โอโลเกนฺโต เทฺว นาทฺทส ฑากญฺจ ปิฏฺฐขาทนียญฺจ. อถ โข โรโช มลฺโล เยนายสฺมา อานนฺโท เตนุปสงฺกมิ, อุปสงฺกมิตฺวา อายสฺมนฺตํ อานนฺทํ เอตทโวจ “อิธ เม ภนฺเต อานนฺท ปฏิปาฏิํ อลภนฺตสฺส เอตทโหสิ ‘ยนฺนูนาหํ ภตฺตคฺคํ โอโลเกยฺยํ, ยํ ภตฺตคฺเค นาสฺส, ตํ ปฏิยาเทยฺยนฺ’ติ, โส โข อหํ ภนฺเต อานนฺท ภตฺตคฺคํ โอโลเกนฺโต เทฺว นาทฺทสํ ฑากญฺจ ปิฏฺฐขาทนียญฺจ, สจาหํ ภนฺเต อานนฺท ปฏิยาเทยฺยํ ฑากญฺจ ปิฏฺฐขาทนียญฺจ, ปฏิคฺคเณฺหยฺย เม ภควา”ติ. เตน หิ โรช ภควนฺตํ ปฏิปุจฺฉิสฺสามีติ. อถ โข อายสฺมา อานนฺโท ภควโต เอตมตฺถํ อาโรเจสิ. เตน หานนฺท ปฏิยาเทตูติ. เตน หิ โรช ปฏิยาเทหีติ. อถ โข โรโช มลฺโล ตสฺสา รตฺติยา \csromanfolio{347}อจฺจเยน ปหูตํ ฑากญฺจ ปิฏฺฐขาทนียญฺจ ปฏิยาทาเปตฺวา ภควโต อุปนาเมสิ “ปฏิคฺคณฺหาตุ เม ภนฺเต ภควา ฑากญฺจ ปิฏฺฐขาทนียญฺจา”ติ. เตน หิ โรช ภิกฺขูนํ เทหีติ. อถ โข โรโช มลฺโล ภิกฺขูนํ เทติ. ภิกฺขู กุกฺกุจฺจายนฺตา น ปฏิคฺคณฺหนฺติ. “ปฏิคฺคณฺหถ ภิกฺขเว ปริภุญฺชถา”ติ. อถ โข โรโช มลฺโล พุทฺธ\-ปฺปมุขํ ภิกฺขุสํฆํ ปหูเตหิ ฑาเกหิ จ ปิฏฺฐ\-ขาทนีเยหิ จ สหตฺถา สนฺตปฺเปตฺวา สมฺปวาเรตฺวา ภควนฺตํ โธตหตฺถํ โอนีต\-ปตฺต\-ปาณิํ เอกมนฺตํ นิสีทิ. เอกมนฺตํ นิสินฺนํ โข โรชํ มลฺลํ ภควา ธมฺมิยา กถาย สนฺทสฺเสตฺวา สมาทเปตฺวา สมุตฺเตเชตฺวา สมฺปหํเสตฺวา อุฏฺฐายาสนา ปกฺกามิ. อถ โข ภควา เอตสฺมิํ นิทาเน เอตสฺมิํ ปกรเณ ธมฺมิํ กถํ กตฺวา ภิกฺขู อามนฺเตสิ “อนุชานามิ ภิกฺขเว สพฺพญฺจ ฑากํ สพฺพญฺจ ปิฏฺฐขาทนียนฺ”ติ.}

\csromanmatikaanchor{mtk.190}
\csromantocmarkref{section}{184. วุฑฺฒปพฺพชิตวตฺถุ}{mtk.190}
\bookmark[dest={mtk.190},level=1]{184. วุฑฺฒปพฺพชิตวตฺถุ}
\csromanheader{1}{184. \textbf{วุฑฺฒ}\-\textbf{ปพฺพชิต}\-\textbf{วตฺถุ}}
\proseitem{303}{อถ โข ภควา กุสินารายํ ยถาภิรนฺตํ วิหริตฺวา เยน อาตุมา เตน จาริกํ ปกฺกามิ มหตา ภิกฺขุ\-สํเฆน สทฺธิํ อฑฺฒเตลเสหิ ภิกฺขุสเตหิ. เตน โข ปน สมเยน อญฺญตโร วุฑฺฒ\-ปพฺพชิโต อาตุมายํ ปฏิวสติ นหาปิตปุพฺโพ, ตสฺส เทฺว ทารกา โหนฺติ มญฺชุกา ปฏิภาเนยฺยกา ทกฺขา ปริโยทาต\-สิปฺปา สเก อาจริยเก นหาปิตกมฺเม. อสฺโสสิ โข โส วุฑฺฒ\-ปพฺพชิโต “ภควา กิร อาตุมํ อาคจฺฉติ มหตา ภิกฺขุ\-สํเฆน สทฺธิํ อฑฺฒเตลเสหิ ภิกฺขุสเตหี”ติ. อถ โข โส วุฑฺฒ\-ปพฺพชิโต เต ทารเก เอตทโวจ “ภควา กิร ตาตา อาตุมํ อาคจฺฉติ มหตา ภิกฺขุ\-สํเฆน สทฺธิํ อฑฺฒเตลเสหิ ภิกฺขุสเตหิ, คจฺฉถ ตุเมฺห ตาตา ขุรภณฺฑํ อาทาย นาฬิยาวาปเกน อนุฆรกํ อนุฆรกํ อาหิณฺฑถ, โลณมฺปิ เตลมฺปิ ตณฺฑุลมฺปิ ขาทนียมฺปิ สํหรถ, ภควโต อาคตสฺส ยาคุปานํ กริสฺสามา”ติ. “เอวํ ตาตา”ติ โข เต ทารกา ตสฺส วุฑฺฒ\-ปพฺพชิตสฺส ปฏิสฺสุณิตฺวา ขุรภณฺฑํ อาทาย นาฬิยาวาปเกน อนุฆรกํ อนุฆรกํ อาหิณฺฑนฺติ โลณมฺปิ เตลมฺปิ ตณฺฑุลมฺปิ ขาทนียมฺปิ สํหรนฺตา. มนุสฺสา เต ทารเก มญฺชุเก ปฏิภาเนยฺยเก ปสฺสิตฺวา เยปิ น การาเปตุกามา, เตปิ การาเปนฺติ, การาเปตฺวาปิ พหุํ เทนฺติ. อถ โข เต ทารกา พหุํ โลณมฺปิ เตลมฺปิ ตณฺฑุลมฺปิ ขาทนียมฺปิ สํหริํสุ.}

\prose{\csromanfolio{348}อถ โข ภควา อนุปุพฺเพน จาริกํ จรมาโน เยน อาตุมา ตทวสริ, ตตฺร สุทํ ภควา อาตุมายํ วิหรติ ภุสาคาเร.}

\prose{อถ โข โส วุฑฺฒ\-ปพฺพชิโต ตสฺสา รตฺติยา อจฺจเยน ปหูตํ ยาคุํ ปฏิยาทาเปตฺวา ภควโต อุปนาเมสิ “ปฏิคฺคณฺหาตุ เม ภนฺเต ภควา ยาคุนฺ”ติ. ชานนฺตาปิ ตถาคตา ปุจฺฉนฺติ ฯเปฯ สาวกานํ วา สิกฺขาปทํ ปญฺญเปสฺสามาติ. อถ โข ภควา ตํ วุฑฺฒ\-ปพฺพชิตํ เอตทโวจ “กุตายํ ภิกฺขุ ยาคู”ติ. อถ โข โส วุฑฺฒ\-ปพฺพชิโต ภควโต เอตมตฺถํ อาโรเจสิ. วิครหิ พุทฺโธ ภควา อนนุจฺฉวิกํ โมฆปุริส อนนุโลมิกํ อปฺปติรูปํ อสฺสามณกํ อกปฺปิยํ อกรณียํ, กถํ หิ นาม ตฺวํ โมฆปุริส ปพฺพชิโต อกปฺปิเย สมาทเปสฺสสิ\csromansharedfootnote{338dd9d64fd7ee1e}{[มิสฺสินฺคฺ โนเต 0]}. เนตํ โมฆปุริส อปฺปสนฺนานํ วา ปสาทาย ฯเปฯ วิครหิตฺวา ธมฺมิํ กถํ กตฺวา ภิกฺขู อามนฺเตสิ “น ภิกฺขเว ปพฺพชิเตน อกปฺปิเย สมาทเปตพฺพํ, โย สมาทเปยฺย, อาปตฺติ ทุกฺกฏสฺส. น จ ภิกฺขเว นหาปิต\-ปุพฺเพน ขุรภณฺฑํ ปริหริตพฺพํ, โย ปริหเรยฺย, อาปตฺติ ทุกฺกฏสฺสา”ติ.}

\prose{อถ โข ภควา อาตุมายํ ยถาภิรนฺตํ วิหริตฺวา เยน สาวตฺถิ เตน จาริกํ ปกฺกามิ, อนุปุพฺเพน จาริกํ จรมาโน เยน สาวตฺถิ ตทวสริ, ตตฺร สุทํ ภควา สาวตฺถิยํ วิหรติ เชตวเน อนาถ\-ปิณฺฑิกสฺส อาราเม.}

\prose{เตน โข ปน สมเยน สาวตฺถิยํ พหุํ ผลขาทนียํ อุปฺปนฺนํ โหติ. อถ โข ภิกฺขูนํ เอตทโหสิ “กิํ นุ โข ภควตา ผลขาทนียํ อนุญฺญาตํ, กิํ อนนุญฺญาตนฺ”ติ. ภควโต เอตมตฺถํ อาโรเจสุํ. อนุชานามิ ภิกฺขเว สพฺพํ ผลขาทนียนฺติ.}

\proseitem{304}{เตน โข ปน สมเยน สํฆิกานิ พีชานิ ปุคฺคลิกาย ภูมิยา โรปิยนฺติ, ปุคฺคลิกานิ พีชานิ สํฆิกาย ภูมิยา โรปิยนฺติ. ภควโต เอตมตฺถํ อาโรเจสุํ. สํฆิกานิ ภิกฺขเว พีชานิ ปุคฺคลิกาย ภูมิยา โรปิตานิ ภาคํ ทตฺวา ปริภุญฺชิ\-ตพฺพานิ. ปุคฺคลิกานิ พีชานิ สํฆิกาย ภูมิยา โรปิตานิ ภาคํ ทตฺวา ปริภุญฺชิ\-ตพฺพานีติ.}

\csromanmatikaanchor{mtk.191}
\csromantocmarkref{section}{185. จตุมหาปเทสกถา}{mtk.191}
\bookmark[dest={mtk.191},level=1]{185. จตุมหาปเทสกถา}
\csromanheader{1}{185. \textbf{จตุมหาปเทส}\-\textbf{กถา}}
\csromanmatikaanchor{mtk.192}
\csromantocmarkref{section}{186. อุทฺทานคาถา}{mtk.192}
\bookmark[dest={mtk.192},level=1]{186. อุทฺทานคาถา}
\proseitem{305}{เตน โข ปน สมเยน ภิกฺขูนํ กิสฺมิญฺจิ กิสฺมิญฺจิ ฐาเน กุกฺกุจฺจํ อุปฺปชฺชติ “กิํ นุ โข ภควตา อนุญฺญาตํ, กิํ อนนุญฺญาตนฺ”ติ. ภควโต เอตมตฺถํ อาโรเจสุํ. ยํ ภิกฺขเว มยา “อิทํ น กปฺปตี”ติ อปฺปฏิกฺขิตฺตํ, \csromanfolio{349}ตญฺเจ อกปฺปิยํ อนุโลเมติ, กปฺปิยํ ปฏิพาหติ, ตํ โว น กปฺปติ. ยํ ภิกฺขเว มยา “อิทํ น กปฺปตี”ติ อปฺปฏิกฺขิตฺตํ. ตญฺเจ กปฺปิยํ อนุโลเมติ, อกปฺปิยํ ปฏิพาหติ, ตํ โว กปฺปติ. ยํ ภิกฺขเว มยา “อิทํ กปฺปตี”ติ อนนุญฺญาตํ, ตญฺเจ อกปฺปิยํ อนุโลเมติ, กปฺปิยํ ปฏิพาหติ, ตํ โว น กปฺปติ. ยํ ภิกฺขเว มยา “อิทํ กปฺปตี”ติ อนนุญฺญาตํ, ตญฺเจ กปฺปิยํ อนุโลเมติ, อกปฺปิยํ ปฏิพาหติ, ตํ โว กปฺปตี”ติ.}

\prose{อถ โข ภิกฺขูนํ เอตทโหสิ “กปฺปติ นุ โข ยาวกาลิเกน ยามกาลิกํ, น นุ โข กปฺปติ. กปฺปติ นุ โข ยาวกาลิเกน สตฺตาหกาลิกํ, น นุ โข กปฺปติ. กปฺปติ นุ โข ยาวกาลิเกน ยาวชีวิกํ, น นุ โข กปฺปติ. กปฺปติ นุ โข ยามกาลิเกน สตฺตาหกาลิกํ, น นุ โข กปฺปติ. กปฺปติ นุ โข ยามกาลิเกน ยาวชีวิกํ, น นุ โข กปฺปติ. กปฺปติ นุ โข สตฺตา\-ห\-กาลิ\-เกน ยาวชีวิกํ, น นุ โข กปฺปตี”ติ. ภควโต เอตมตฺถํ อาโรเจสุํ. ยาวกาลิเกน ภิกฺขเว ยามกาลิกํ ตทหุ ปฏิคฺคหิตํ กาเล กปฺปติ, วิกาเล น กปฺปติ. ยาวกาลิเกน ภิกฺขเว สตฺตาหกาลิกํ ตทหุ ปฏิคฺคหิตํ กาเล กปฺปติ, วิกาเล น กปฺปติ. ยาวกาลิเกน ภิกฺขเว ยาวชีวิกํ ตทหุ ปฏิคฺคหิตํ กาเล กปฺปติ, วิกาเล น กปฺปติ. ยามกาลิเกน ภิกฺขเว สตฺตาหกาลิกํ ตทหุ ปฏิคฺคหิตํ ยาเม กปฺปติ, ยามาติกฺกนฺเต น กปฺปติ. ยามกาลิเกน ภิกฺขเว ยาวชีวิกํ ตทหุ ปฏิคฺคหิตํ ยาเม กปฺปติ, ยามาติกฺกนฺเต น กปฺปติ. สตฺตา\-ห\-กาลิ\-เกน ภิกฺขเว ยาวชีวิกํ ปฏิคฺคหิตํ สตฺตาหํ กปฺปติ, สตฺตาหา\-ติกฺกนฺเต น กปฺปตี”ติ.}

\vspace{\tipitakaparskip}
\csromancenter{เภสชฺช\-กฺขนฺธโก ฉฏฺโฐ.}
\csromancenter{186. \textbf{ตสฺสุทฺทานํ}}
\setlength{\csromangathapagewidth}{0pt}
\setlength{\csromangathawakwidth}{0pt}
\csromangathameasuregroup{สารทิเก วิกาเลปิ, \\ กสาเวหิ ปณฺณํ ผลํ, \\ จุณฺณํ จาลินิ มํสญฺจ, \\ อญฺชนี อุจฺจาปารุตา, \\ ถวิกํสพทฺธกํ สุตฺตํ, \\ นตฺถุกรณี ธูมญฺจ, \\ เตลปาเกสุ มชฺชญฺจ, \\ ตุมฺพํ เสทํ สมฺภารญฺจ, \\ ทกโกฏฺฐํ โลหิตญฺจ, \\ ปชฺชํ สตฺถํ กสาวญฺจ, \\ โจฬํ สาสปกุฏฺฏญฺจ, \\ วณเตลํ วิกาสิกํ, \\ คูถํ กโรนฺโต โลฬิญฺจ, \\ คนฺธา วิเรจนญฺเจว, \\ ปฏิจฺฉาทนิ ปพฺภารา, \\ คุฬํ มุคฺคํ โสวีรญฺจ, \\ ปุนานุญฺญาสิ ทุพฺภิกฺเข, \\ ปุเรภตฺตํ กายฑาโห, \\ วตฺถิกมฺมญฺจ สุปฺปิญฺจ, \\ หตฺถิอสฺสา สุนโข จ, \\ อจฺฉตรจฺฉมํสญฺจ, \\ ตรุณํ อญฺญตฺร คุฬํ, \\ คงฺคา โกฏิสจฺจกถา, \\ อุทฺทิสฺส กตํ สุภิกฺขํ, \\ เมโฆ ยโส เมณฺฑโก จ, \\ เกณิ อมฺโพ ชมฺพุ โจจ, \\ ผารุสกา ฑากปิฏฺฐํ, \\ สาวตฺถิยํ ผลํ พีชํ,}{\csromangathabat{สารทิเก วิกาเลปิ,}{วสํ มูเล ปิฏฺเฐหิ จ.} \\ \csromangathabat{กสาเวหิ ปณฺณํ ผลํ,}{ชตุ โลณํ ฉกณํ จ.} \\ \csromangathabat{จุณฺณํ จาลินิ มํสญฺจ,}{อญฺชนํ อุปปิสนี\textsuperscript{1}.} \\ \csromangathabat{อญฺชนี อุจฺจาปารุตา,}{สลากา สลากฐานิํ\textsuperscript{1}.} \\ \csromangathabat{ถวิกํสพทฺธกํ สุตฺตํ,}{มุทฺธนิเตลนตฺถุ จ.} \\ \csromangathabat{นตฺถุกรณี ธูมญฺจ,}{เนตฺตญฺจาปิธนตฺถวิ.} \\ \csromangathabat{เตลปาเกสุ มชฺชญฺจ,}{อติกฺขิตฺตํ อพฺภญฺชนํ.} \\ \csromangathabat{ตุมฺพํ เสทํ สมฺภารญฺจ,}{มหา ภงฺโคทกํ ตถา.} \\ \csromangathabat{ทกโกฏฺฐํ โลหิตญฺจ,}{วิสาณํ ปาทพฺภญฺชนํ.} \\ \csromangathabat{ปชฺชํ สตฺถํ กสาวญฺจ,}{ติลกกฺกํ กพฬิกํ.} \\ \csromangathabat{โจฬํ สาสปกุฏฺฏญฺจ,}{ธูม สกฺขริกาย จ.} \\ \csromangathabat{วณเตลํ วิกาสิกํ,}{วิกฏญฺจ ปฏิคฺคหํ.} \\ \csromangathabat{คูถํ กโรนฺโต โลฬิญฺจ,}{ขารํ มุตฺตหรีตกํ.} \\ \csromangathabat{คนฺธา วิเรจนญฺเจว,}{อจฺฉากฏํ กฏากฏํ.} \\ \csromangathabat{ปฏิจฺฉาทนิ ปพฺภารา,}{อาราม สตฺตาเหน จ.} \\ \csromangathabat{คุฬํ มุคฺคํ โสวีรญฺจ,}{สามํปากา ปุนาปเจ.} \\ \csromangathabat{ปุนานุญฺญาสิ ทุพฺภิกฺเข,}{ผลญฺจ ติลขาทนี.} \\ \csromangathabat{ปุเรภตฺตํ กายฑาโห,}{นิพฺพตฺตญฺจ ภคนฺทลํ.} \\ \csromangathabat{วตฺถิกมฺมญฺจ สุปฺปิญฺจ,}{มนุสฺสมํสเมว จ.} \\ \csromangathabat{หตฺถิอสฺสา สุนโข จ,}{อหิ สีหญฺจ ทีปิกํ\textsuperscript{1}.} \\ \csromangathabat{อจฺฉตรจฺฉมํสญฺจ,}{ปฏิปาฏิ จ ยาคุ จ.} \\ \csromangathabat{ตรุณํ อญฺญตฺร คุฬํ,}{สุนิธาวสถาคารํ.} \\ \csromangathabat{คงฺคา โกฏิสจฺจกถา,}{อมฺพปาลี จ ลิจฺฉวี.} \\ \csromangathabat{อุทฺทิสฺส กตํ สุภิกฺขํ,}{ปุนเทว ปฏิกฺขิปิ.} \\ \csromangathabat{เมโฆ ยโส เมณฺฑโก จ,}{โครสํ ปาเถยฺยเกน จ.} \\ \csromangathabat{เกณิ อมฺโพ ชมฺพุ โจจ,}{โมจมธุมุทฺทิกสาลุกํ.} \\ \csromangathabat{ผารุสกา ฑากปิฏฺฐํ,}{อาตุมายํ นหาปิโต.} \\ \csromangathabat{สาวตฺถิยํ ผลํ พีชํ,}{กิสฺมิํ ฐาเน จ กาลิเกติ.} \\ อิมมฺหิ ขนฺธเก วตฺถู เอกสตํ ฉวตฺถุ.}
\csromangathameasurewak{อิมมฺหิ ขนฺธเก วตฺถู เอกสตํ ฉวตฺถุ.}
\csromangathasetleft{สารทิเก วิกาเลปิ, \\ กสาเวหิ ปณฺณํ ผลํ, \\ จุณฺณํ จาลินิ มํสญฺจ, \\ อญฺชนี อุจฺจาปารุตา, \\ ถวิกํสพทฺธกํ สุตฺตํ, \\ นตฺถุกรณี ธูมญฺจ, \\ เตลปาเกสุ มชฺชญฺจ, \\ ตุมฺพํ เสทํ สมฺภารญฺจ, \\ ทกโกฏฺฐํ โลหิตญฺจ, \\ ปชฺชํ สตฺถํ กสาวญฺจ, \\ โจฬํ สาสปกุฏฺฏญฺจ, \\ วณเตลํ วิกาสิกํ, \\ คูถํ กโรนฺโต โลฬิญฺจ, \\ คนฺธา วิเรจนญฺเจว, \\ ปฏิจฺฉาทนิ ปพฺภารา, \\ คุฬํ มุคฺคํ โสวีรญฺจ, \\ ปุนานุญฺญาสิ ทุพฺภิกฺเข, \\ ปุเรภตฺตํ กายฑาโห, \\ วตฺถิกมฺมญฺจ สุปฺปิญฺจ, \\ หตฺถิอสฺสา สุนโข จ, \\ อจฺฉตรจฺฉมํสญฺจ, \\ ตรุณํ อญฺญตฺร คุฬํ, \\ คงฺคา โกฏิสจฺจกถา, \\ อุทฺทิสฺส กตํ สุภิกฺขํ, \\ เมโฆ ยโส เมณฺฑโก จ, \\ เกณิ อมฺโพ ชมฺพุ โจจ, \\ ผารุสกา ฑากปิฏฺฐํ, \\ สาวตฺถิยํ ผลํ พีชํ,}
\csromangathagroupwithclosermajor{\csromangathastanza{\csromangathabat{สารทิเก วิกาเลปิ,}{วสํ มูเล ปิฏฺเฐหิ จ.} \\ \csromangathabat{กสาเวหิ ปณฺณํ ผลํ,}{ชตุ โลณํ ฉกณํ จ.}}\csromangathastanza{\csromangathabat{จุณฺณํ จาลินิ มํสญฺจ,}{อญฺชนํ อุปปิสนี\csromansharedfootnote{f824749d53076b6d}{อุปปิํสนี (สี), อุปปิํสนํ (สฺยา)}.} \\ \csromangathabat{อญฺชนี อุจฺจาปารุตา,}{สลากา สลากฐานิํ\csromansharedfootnote{fb2945d669729fd2}{สลาโกธนี (สี, สฺยา)}.}}\csromangathastanza{\csromanfolio{350}\csromangathabat{ถวิกํสพทฺธกํ สุตฺตํ,}{มุทฺธนิเตลนตฺถุ จ.} \\ \csromangathabat{นตฺถุกรณี ธูมญฺจ,}{เนตฺตญฺจาปิธนตฺถวิ.}}\csromangathastanza{\csromangathabat{เตลปาเกสุ มชฺชญฺจ,}{อติกฺขิตฺตํ อพฺภญฺชนํ.} \\ \csromangathabat{ตุมฺพํ เสทํ สมฺภารญฺจ,}{มหา ภงฺโคทกํ ตถา.}}\csromangathastanza{\csromangathabat{ทกโกฏฺฐํ โลหิตญฺจ,}{วิสาณํ ปาทพฺภญฺชนํ.} \\ \csromangathabat{ปชฺชํ สตฺถํ กสาวญฺจ,}{ติลกกฺกํ กพฬิกํ.}}\csromangathastanza{\csromangathabat{โจฬํ สาสปกุฏฺฏญฺจ,}{ธูม สกฺขริกาย จ.} \\ \csromangathabat{วณเตลํ วิกาสิกํ,}{วิกฏญฺจ ปฏิคฺคหํ.}}\csromangathastanza{\csromangathabat{คูถํ กโรนฺโต โลฬิญฺจ,}{ขารํ มุตฺตหรีตกํ.} \\ \csromangathabat{คนฺธา วิเรจนญฺเจว,}{อจฺฉากฏํ กฏากฏํ.}}\csromangathastanza{\csromangathabat{ปฏิจฺฉาทนิ ปพฺภารา,}{อาราม สตฺตาเหน จ.} \\ \csromangathabat{คุฬํ มุคฺคํ โสวีรญฺจ,}{สามํปากา ปุนาปเจ.}}\csromangathastanza{\csromangathabat{ปุนานุญฺญาสิ ทุพฺภิกฺเข,}{ผลญฺจ ติลขาทนี.} \\ \csromangathabat{ปุเรภตฺตํ กายฑาโห,}{นิพฺพตฺตญฺจ ภคนฺทลํ.}}\csromangathastanza{\csromangathabat{วตฺถิกมฺมญฺจ สุปฺปิญฺจ,}{มนุสฺสมํสเมว จ.} \\ \csromangathabat{หตฺถิอสฺสา สุนโข จ,}{อหิ สีหญฺจ ทีปิกํ\csromansharedfootnote{03f8e660f96c470c}{หตฺถิอสฺสสุนขาหิ, สีหพฺยคฺฆญฺจ ทีปิกํ (สี)}.}}\csromangathastanza{\csromangathabat{อจฺฉตรจฺฉมํสญฺจ,}{ปฏิปาฏิ จ ยาคุ จ.} \\ \csromangathabat{ตรุณํ อญฺญตฺร คุฬํ,}{สุนิธาวสถาคารํ.}}\csromangathastanza{\csromangathabat{คงฺคา โกฏิสจฺจกถา,}{อมฺพปาลี จ ลิจฺฉวี.} \\ \csromangathabat{อุทฺทิสฺส กตํ สุภิกฺขํ,}{ปุนเทว ปฏิกฺขิปิ.}}\csromangathastanza{\csromangathabat{เมโฆ ยโส เมณฺฑโก จ,}{โครสํ ปาเถยฺยเกน จ.} \\ \csromangathabat{เกณิ อมฺโพ ชมฺพุ โจจ,}{โมจมธุมุทฺทิกสาลุกํ.}}\csromangathastanza{\csromangathabat{ผารุสกา ฑากปิฏฺฐํ,}{อาตุมายํ นหาปิโต.} \\ \csromangathabat{สาวตฺถิยํ ผลํ พีชํ,}{กิสฺมิํ ฐาเน จ กาลิเกติ.}}\csromangathastanza{\csromangathawak{อิมมฺหิ ขนฺธเก วตฺถู เอกสตํ ฉวตฺถุ.}}}{\textbf{เภสชฺช}\-\textbf{กฺขนฺธโก นิฏฺฐิโต.}}
\csromanmatikaanchor{mtk.193}
\csromantocmarknopage{chapter}{7. กถินกฺขนฺธก}{mtk.193}
\bookmark[dest={mtk.193},level=0]{7. กถินกฺขนฺธก}
\chapterheadpageread{7. \textbf{กถิน}\-\textbf{กฺขนฺธก}}
\csromanmatikaanchor{mtk.194}
\csromantocmarkref{section}{187. กถินานุชานนา}{mtk.194}
\bookmark[dest={mtk.194},level=1]{187. กถินานุชานนา}
\csromanheader{1}{187. \textbf{กถินา}\-\textbf{นุชานนา}}
\proseitem{306}{\csromanfolio{351}เตน สมเยน พุทฺโธ ภควา สาวตฺถิยํ วิหรติ เชตวเน อนาถ\-ปิณฺฑิกสฺส อาราเม. เตน โข ปน สมเยน ติํสมตฺตา ปาเวยฺยกา\csromansharedfootnote{aef9850a47077043}{ปาเฐยฺยกา (สี, สฺยา)} ภิกฺขู สพฺเพ อารญฺญิกา สพฺเพ ปิณฺฑปาติกา สพฺเพ ปํสุกูลิกา สพฺเพ เตจีวริกา สาวตฺถิํ อาคจฺฉนฺตา ภควนฺตํ ทสฺสนาย อุปกฏฺฐาย วสฺสู\-ปนายิกาย นาสกฺขิํสุ สาวตฺถิยํ วสฺสูปนายิกํ สมฺภาเวตุํ, อนฺตรามคฺเค สาเกเต วสฺสํ อุปคจฺฉิํสุ, เต อุกฺกณฺฐิตรูปา วสฺสํ วสิํสุ “อาสนฺเนว โน ภควา วิหรติ อิโต ฉสุ โยชเนสุ, น จ มยํ ลภาม ภควนฺตํ ทสฺสนายา”ติ. อถ โข เต ภิกฺขู วสฺสํวุฏฺฐา เตมาสจฺจเยน กตาย ปวารณาย เทเว วสฺสนฺเต อุทกสงฺคเห อุทกจิกฺขลฺเล โอกปุณฺเณหิ จีวเรหิ กิลนฺตรูปา เยน สาวตฺถิ เชตวนํ อนาถ\-ปิณฺฑิกสฺส อาราโม, เยน ภควา เตนุ\-ปสงฺกมิํสุ, อุปสงฺกมิตฺวา ภควนฺตํ อภิวาเทตฺวา เอกมนฺตํ นิสีทิํสุ. อาจิณฺณํ โข ปเนตํ พุทฺธานํ ภควนฺตานํ อาคนฺตุเกหิ ภิกฺขูหิ สทฺธิํ ปฏิสมฺโมทิตุํ. อถ โข ภควา เต ภิกฺขู เอตทโวจ “กจฺจิ ภิกฺขเว ขมนียํ, กจฺจิ ยาปนียํ, กจฺจิ สมคฺคา สมฺโมทมานา อวิวทมานา ผาสุกํ วสฺสํ วสิตฺถ, น จ ปิณฺฑเกน กิลมิตฺถา”ติ. ขมนียํ ภควา, ยาปนียํ ภควา, สมคฺคา จ มยํ ภนฺเต สมฺโมทมานา อวิวทมานา วสฺสํ วสิมฺหา, น จ ปิณฺฑเกน กิลมิมฺหา, อิธ มยํ ภนฺเต ติํสมตฺตา ปาเวยฺยกา ภิกฺขู สาวตฺถิํ อาคจฺฉนฺตา ภควนฺตํ ทสฺสนาย อุปกฏฺฐาย วสฺสู\-ปนายิกาย นาสกฺขิมฺหา สาวตฺถิยํ วสฺสูปนายิกํ สมฺภาเวตุํ, อนฺตรามคฺเค สาเกเต วสฺสํ อุปคจฺฉิมฺหา, เต มยํ ภนฺเต อุกฺกณฺฐิตรูปา วสฺสํ วสิมฺหา “อาสนฺเนว โน ภควา วิหรติ อิโต ฉสุ โยชเนสุ, น จ มยํ ลภาม ภควนฺตํ ทสฺสนายา”ติ. อถ โข มยํ ภนฺเต วสฺสํวุฏฺฐา เตมาสจฺจเยน กตาย ปวารณาย เทเว วสฺสนฺเต อุทกสงฺคเห อุทกจิกฺขลฺเล โอกปุณฺเณหิ จีวเรหิ กิลนฺตรูปา อทฺธานํ อาคตาติ. อถ โข ภควา เอตสฺมิํ นิทาเน เอตสฺมิํ ปกรเณ ธมฺมิํ กถํ กตฺวา ภิกฺขู อามนฺเตสิ\csromandash{}อนุชานามิ ภิกฺขเว วสฺสํ\-วุฏฺฐานํ \csromanfolio{352}ภิกฺขูนํ กถินํ\csromansharedfootnote{a916aaa72abe7fa5}{กฐินํ (สี, สฺยา)} อตฺถริตุํ, อตฺถต\-กถินานํ โว ภิกฺขเว ปญฺจ กปฺปิสฺสนฺติ, อนามนฺตจาโร อสมาทานจาโร คณโภชนํ ยาวทตฺถ\-จีวรํ โย จ ตตฺถ จีวรุปฺปาโท โส เนสํ ภวิสฺสตีติ, อตฺถต\-กถินานํ โว ภิกฺขเว อิมานิ ปญฺจ กปฺปิสฺสนฺติ. เอวญฺจ ปน ภิกฺขเว กถินํ อตฺถริตพฺพํ, พฺยตฺเตน ภิกฺขุนา ปฏิพเลน สํโฆ ญาเปตพฺโพ\csromandash{}}

\proseitem{307}{“สุณาตุ เม ภนฺเต สํโฆ, อิทํ สํฆสฺส กถินทุสฺสํ อุปฺปนฺนํ, ยทิ สํฆสฺส ปตฺตกลฺลํ, สํโฆ อิมํ กถินทุสฺสํ อิตฺถนฺนามสฺส ภิกฺขุโน ทเทยฺย กถินํ อตฺตริตุํ, เอสา ญตฺติ.}

\prose{สุณาตุ เม ภนฺเต สํโฆ, อิทํ สํฆสฺส กถินทุสฺสํ อุปฺปนฺนํ, สํโฆ อิมํ กถินทุสฺสํ อิตฺถนฺนามสฺส ภิกฺขุโน เทติ กถินํ อตฺถริตุํ, ยสฺสายสฺมโต ขมติ อิมสฺส กถิน\-ทุสฺสสฺส อิตฺถนฺนามสฺส ภิกฺขุโน ทานํ กถินํ อตฺถริตุํ, โส ตุณฺหสฺส, ยสฺส นกฺขมติ, โส ภาเสยฺย.}

\prose{ทินฺนํ อิทํ สํเฆน กถินทุสฺสํ อิตฺถนฺนามสฺส ภิกฺขุโน กถินํ อตฺถริตุํ, ขมติ สํฆสฺส ตสฺมา ตุณฺหี, เอวเมตํ ธารยามี”ติ.}

\proseitem{308}{เอวํ โข ภิกฺขเว อตฺถตํ โหติ กถินํ, เอวํ อนตฺถตํ. กถญฺจ ปน ภิกฺขเว อนตฺถตํ โหติ กถินํ, น อุลฺลิขิต\-มตฺเตน อตฺถตํ โหติ กถินํ, น โธวนมตฺเตน อตฺถตํ โหติ กถินํ, น จีวรวิจารณ\-มตฺเตน อตฺถตํ โหติ กถินํ, น เฉทนมตฺเตน อตฺถตํ โหติ กถินํ, น พนฺธน\-มตฺเตน อตฺถตํ โหติ กถินํ, น โอวฏฺฏิย\-กรณ\-มตฺเตน อตฺถตํ โหติ กถินํ, น กณฺฑุส\-กรณ\-มตฺเตน\csromansharedfootnote{45513aeaaa2087d1}{น คณฺฑุสกรณมตฺเตน (ก)} อตฺถตํ โหติ กถินํ, น ทฬฺหีกมฺม\-กรณ\-มตฺเตน อตฺถตํ โหติ กถินํ, น อนุวาต\-กรณ\-มตฺเตน อตฺถตํ โหติ กถินํ, น ปริภณฺฑ\-กรณ\-มตฺเตน อตฺถตํ โหติ กถินํ, น โอวทฺเธยฺย\-กรณ\-มตฺเตน\csromansharedfootnote{ae8546808d609fe6}{น โอวฏฺเฏยฺยกรณมตฺเตน (สี, สฺยา), น โอวเทยฺยกรณมตฺเตน (ก)} อตฺถตํ โหติ กถินํ, น กมฺพล\-มทฺทน\-มตฺเตน อตฺถตํ โหติ กถินํ, น นิมิตฺตกเตน อตฺถตํ โหติ กถินํ, น ปริกถา\-กเตน อตฺถตํ โหติ กถินํ, น กุกฺกุกเตน อตฺถตํ โหติ กถินํ, น สนฺนิธิกเตน อตฺถตํ โหติ กถินํ, น นิสฺสคฺคิเยน อตฺถตํ โหติ กถินํ, น อกปฺปกเตน \csromanfolio{353}อตฺถตํ โหติ กถินํ, น อญฺญตฺร สงฺฆาฏิยา อตฺถตํ โหติ กถินํ, น อญฺญตฺร อุตฺตราสงฺเคน อตฺถตํ โหติ กถินํ, น อญฺญตฺร อนฺตรวาสเกน อตฺถตํ โหติ กถินํ, น อญฺญตฺร ปญฺจเกน วา อติเรก\-ปญฺจเกน วา ตทเหว สญฺฉินฺเนน สมณฺฑลีกเตน อตฺถตํ โหติ กถินํ, น อญฺญตฺร ปุคฺคลสฺส อตฺถารา อตฺถตํ โหติ กถินํ, สมฺมา เจ อตฺถตํ โหติ กถินํ, ตญฺเจ นิสฺสีมฏฺโฐ อนุโมทติ, เอวมฺปิ อนตฺถตํ โหติ กถินํ. เอวํ โข ภิกฺขเว อนตฺถตํ โหติ กถินํ.}

\proseitem{309}{กถญฺจ ภิกฺขเว อตฺถตํ โหติ กถินํ, อหเตน อตฺถตํ โหติ กถินํ, อหตกปฺเปน อตฺถตํ โหติ กถินํ, ปิโลติกาย อตฺถตํ โหติ กถินํ, ปํสุกูเลน อตฺถตํ โหติ กถินํ, ปาปณิเกน อตฺถตํ โหติ กถินํ, อนิมิตฺตกเตน อตฺถตํ โหติ กถินํ, อปริกถากเตน อตฺถตํ โหติ กถินํ, อกุกฺกุกเตน อตฺถตํ โหติ กถินํ, อสนฺนิธิกเตน อตฺถตํ โหติ กถินํ, อนิสฺสคฺคิเยน อตฺถตํ โหติ กถินํ, กปฺปกเตน อตฺถตํ โหติ กถินํ, สงฺฆาฏิยา อตฺถตํ โหติ กถินํ, อุตฺตราสงฺเคน อตฺถตํ โหติ กถินํ, อนฺตรวาสเกน อตฺถตํ โหติ กถินํ, ปญฺจเกน วา อติเรก\-ปญฺจเกน วา ตทเหว สญฺฉินฺเนน สมณฺฑลีกเตน อตฺถตํ โหติ กถินํ, ปุคฺคลสฺส อตฺถารา อตฺถตํ โหติ กถินํ, สมฺมา เจ อตฺถตํ โหติ กถินํ, ตญฺเจ สีมฏฺโฐ อนุโมทติ, เอวมฺปิ อตฺถตํ โหติ กถินํ. เอวํ โข ภิกฺขเว อตฺถตํ โหติ กถินํ.}

\proseitem{310}{กถญฺจ ภิกฺขเว อุพฺภตํ โหติ กถินํ, อฏฺฐิมา ภิกฺขเว มาติกา กถินสฺส อุพฺภาราย, ปกฺกมนนฺติกา นิฏฺฐานนฺติกา สนฺนิฏฺฐานนฺติกา นาสนนฺติกา สวนนฺติกา อาสาวจฺเฉทิกา สีมาติกฺกนฺติกา สหุพฺภาราติ\csromansharedfootnote{b9ef5007b32efd3c}{สอุพฺภาราติ (ก)}.}
\csromansectionrule

\csromanmatikaanchor{mtk.195}
\csromantocmarkref{section}{188. อาทายสตฺตก}{mtk.195}
\bookmark[dest={mtk.195},level=1]{188. อาทายสตฺตก}
\csromanheader{1}{188. \textbf{อาทายสตฺตก}}
\proseitem{311}{ภิกฺขุ อตฺถตกถิโน กตจีวรํ อาทาย ปกฺกมติ “น ปจฺเจสฺสนฺ”ติ, ตสฺส ภิกฺขุโน ปกฺกมนนฺติโก กถินุทฺธาโร. (1)}

\prose{\csromanfolio{354}ภิกฺขุ อตฺถตกถิโน จีวรํ อาทาย ปกฺกมติ, ตสฺส พหิสีม\-คตสฺส เอวํ โหติ “อิเธวิมํ จีวรํ กาเรสฺสํ, น ปจฺเจสฺสนฺ”ติ, โส ตํ จีวรํ กาเรติ, ตสฺส ภิกฺขุโน นิฏฺฐานนฺติโก กถินุทฺธาโร. (2)}

\prose{ภิกฺขุ อตฺถตกถิโน จีวรํ อาทาย ปกฺกมติ, ตสฺส พหิสีม\-คตสฺส เอวํ โหติ “เนวิมํ จีวรํ กาเรสฺสํ, น ปจฺเจสฺสนฺ”ติ, ตสฺส ภิกฺขุโน สนฺนิฏฺฐานนฺติโก กถินุทฺธาโร. (3)}

\prose{ภิกฺขุ อตฺถตกถิโน จีวรํ อาทาย ปกฺกมติ, ตสฺส พหิสีม\-คตสฺส เอวํ โหติ “อิเธวิมํ จีวรํ กาเรสฺสํ, น ปจฺเจสฺสนฺ”ติ, โส ตํ จีวรํ กาเรติ, ตสฺส ตํ จีวรํ กยิรมานํ นสฺสติ, ตสฺส ภิกฺขุโน นาสนนฺติโก กถินุทฺธาโร. (4)}

\prose{ภิกฺขุ อตฺถตกถิโน จีวรํ อาทาย ปกฺกมติ “ปจฺเจสฺสนฺ”ติ, โส พหิสีมคโต ตํ จีวรํ กาเรติ, โส กตจีวโร สุณาติ “อุพฺภตํ กิร ตสฺมิํ อาวาเส กถินนฺ”ติ, ตสฺส ภิกฺขุโน สวนนฺติโก กถินุทฺธาโร. (5)}

\prose{ภิกฺขุ อตฺถตกถิโน จีวรํ อาทาย ปกฺกมติ “ปจฺเจสฺสนฺ”ติ, โส พหิสีมคโต ตํ จีวรํ กาเรติ, โส กตจีวโร “ปจฺเจสฺสํ ปจฺเจสฺสนฺ”ติ พหิทฺธา กถินุทฺธารํ วีตินาเมติ, ตสฺส ภิกฺขุโน สีมาติกฺกนฺติโก กถินุทฺธาโร. (6)}

\prosewithcloser{ภิกฺขุ อตฺถตกถิโน จีวรํ อาทาย ปกฺกมติ “ปจฺเจสฺสนฺ”ติ, โส พหิสีมคโต ตํ จีวรํ กาเรติ, โส กตจีวโร “ปจฺเจสฺสํ ปจฺเจสฺสนฺ”ติ สมฺภุณาติ กถินุทฺธารํ, ตสฺส ภิกฺขุโน สห ภิกฺขูหิ กถินุทฺธาโร.}{(7) อาทายสตฺตกํ นิฏฺฐิตํ.}
\csromanmatikaanchor{mtk.196}
\csromantocmarkref{section}{189. สมาทายสตฺตก}{mtk.196}
\bookmark[dest={mtk.196},level=1]{189. สมาทายสตฺตก}
\csromanheader{1}{189. \textbf{สมาทาย}\-\textbf{สตฺตก}}
\proseitem{312}{ภิกฺขุ อตฺถตกถิโน กตจีวรํ สมาทาย ปกฺกมติ “น ปจฺเจสฺสนฺ”ติ, ตสฺส ภิกฺขุโน ปกฺกมนนฺติโก กถินุทฺธาโร. (1)}

\prose{\csromanfolio{355}ภิกฺขุ อตฺถตกถิโน จีวรํ สมาทาย ปกฺกมติ, ตสฺส พหิสีม\-คตสฺส เอวํ โหติ “อิเธวิมํ จีวรํ กาเรสฺสํ, น ปจฺเจสฺสนฺ”ติ, โส ตํ จีวรํ กาเรติ, ตสฺส ภิกฺขุโน นิฏฺฐานนฺติโก กถินุทฺธาโร. (2)}

\prose{ภิกฺขุ อตฺตตกถิโน จีวรํ สมาทาย ปกฺกมติ, ตสฺส พหิสีม\-คตสฺส เอวํ โหติ “เนวิมํ จีวรํ กาเรสฺสํ, น ปจฺเจสฺสนฺ”ติ, ตสฺส ภิกฺขุโน สนฺนิฏฺฐานนฺติโก กถินุทฺธาโร. (3)}

\prose{ภิกฺขุ อตฺถตกถิโน จีวรํ สมาทาย ปกฺกมติ, ตสฺส พหิสีม\-คตสฺส เอวํ โหติ “อิเธวิมํ จีวรํ กาเรสฺสํ, น ปจฺเจสฺสนฺ”ติ, โส ตํ จีวรํ กาเรติ, ตสฺส ตํ จีวรํ กยิรมานํ นสฺสติ, ตสฺส ภิกฺขุโน นาสนนฺติโก กถินุทฺธาโร. (4)}

\prose{ภิกฺขุ อตฺถตกถิโน จีวรํ สมาทาย ปกฺกมติ “ปจฺเจสฺสนฺ”ติ, โส พหิสีมคโต ตํ จีวรํ กาเรติ, โส กตจีวโร สุณาติ “อุพฺภตํ กิร ตสฺมิํ อาวาเส กถินนฺ”ติ, ตสฺส ภิกฺขุโน สวนนฺติโก กถินุทฺธาโร. (5)}

\prose{ภิกฺขุ อตฺถตกถิโน จีวรํ สมาทาย ปกฺกมติ “ปจฺเจสฺสนฺ”ติ, โส พหิสีมคโต ตํ จีวรํ กาเรติ, โส กตจีวโร “ปจฺเจสฺสํ ปจฺเจสฺสนฺ”ติ พหิทฺธา กถินุทฺธารํ วีตินาเมติ, ตสฺส ภิกฺขุโน สีมาติกฺกนฺติโก กถินุทฺธาโร. (6)}

\prosewithcloserruled{ภิกฺขุ อตฺถตกถิโน จีวรํ สมาทาย ปกฺกมติ “ปจฺเจสฺสนฺ”ติ, โส พหิสีมคโต ตํ จีวรํ กาเรติ, โส กตจีวโร “ปจฺเจสฺสํ ปจฺเจสฺสนฺ”ติ สมฺภุณาติ กถินุทฺธารํ, ตสฺส ภิกฺขุโน สห ภิกฺขูหิ กถินุทฺธาโร. (7)}{สมาทาย\-สตฺตกํ นิฏฺฐิตํ\csromansharedfootnote{aff45f2628162078}{...ทุติยํ นิฏฺฐิตํ (ก)}.}
\csromanmatikaanchor{mtk.197}
\csromantocmarkref{section}{190. อาทายฉกฺก}{mtk.197}
\bookmark[dest={mtk.197},level=1]{190. อาทายฉกฺก}
\csromanheader{1}{190. \textbf{อาทายฉกฺก}}
\proseitem{313}{ภิกฺขุ อตฺถตกถิโน วิปฺปกต\-จีวรํ อาทาย ปกฺกมติ, ตสฺส พหิสีม\-คตสฺส เอวํ โหติ “อิเธวิมํ จีวรํ กาเรสฺสํ, น ปจฺเจสฺสนฺ”ติ, โส ตํ จีวรํ กาเรติ, ตสฺส ภิกฺขุโน นิฏฺฐานนฺติโก กถินุทฺธาโร. (1)}

\prose{\csromanfolio{356}ภิกฺขุ อตฺถตกถิโน วิปฺปกต\-จีวรํ อาทาย ปกฺกมติ, ตสฺส พหิสีม\-คตสฺส เอวํ โหติ “เนวิมํ จีวรํ กาเรสฺสํ, น ปจฺเจสฺสนฺ”ติ, ตสฺส ภิกฺขุโน สนฺนิฏฺฐานนฺติโก กถินุทฺธาโร. (2)}

\prose{ภิกฺขุ อตฺถตกถิโน วิปฺปกต\-จีวรํ อาทาย ปกฺกมติ, ตสฺส พหิสีม\-คตสฺส เอวํ โหติ “อิเธวิมํ จีวรํ กาเรสฺสํ, น ปจฺเจสฺสนฺ”ติ, โส ตํ จีวรํ กาเรติ, ตสฺส ตํ จีวรํ กยิรมานํ นสฺสติ, ตสฺส ภิกฺขุโน นาสนนฺติโก กถินุทฺธาโร. (3)}

\prose{ภิกฺขุ อตฺถตกถิโน วิปฺปกต\-จีวรํ อาทาย ปกฺกมติ “ปจฺเจสฺสนฺ”ติ, โส พหิสีมคโต ตํ จีวรํ กาเรติ, โส กตจีวโร สุณาติ “อุพฺภตํ กิร ตสฺมิํ อาวาเส กถินนฺ”ติ, ตสฺส ภิกฺขุโน สวนนฺติโก กถินุทฺธาโร. (4)}

\prose{ภิกฺขุ อตฺถตกถิโน วิปฺปกต\-จีวรํ อาทาย ปกฺกมติ “ปจฺเจสฺสนฺ”ติ, โส พหิสีมคโต ตํ จีวรํ กาเรติ, โส กตจีวโร “ปจฺเจสฺสํ ปจฺเจสฺสนฺ”ติ พหิทฺธา กถินุทฺธารํ วีตินาเมติ, ตสฺส ภิกฺขุโน สีมาติกฺกนฺติโก กถินุทฺธาโร. (5)}

\prosewithcloserruled{ภิกฺขุ อตฺถตกถิโน วิปฺปกต\-จีวรํ อาทาย ปกฺกมติ “ปจฺเจสฺสนฺ”ติ, โส พหิสีมคโต ตํ จีวรํ กาเรติ, โส กตจีวโร “ปจฺเจสฺสํ ปจฺเจสฺสนฺ”ติ สมฺภุณาติ กถินุทฺธารํ, ตสฺส ภิกฺขุโน สห ภิกฺขูหิ กถินุทฺธาโร. (6)}{อาทายฉกฺกํ นิฏฺฐิตํ.}
\csromanmatikaanchor{mtk.198}
\csromantocmarkref{section}{191. สมาทายฉกฺก}{mtk.198}
\bookmark[dest={mtk.198},level=1]{191. สมาทายฉกฺก}
\csromanheader{1}{191. \textbf{สมาทายฉกฺก}}
\proseitem{314}{ภิกฺขุ อตฺถตกถิโน วิปฺปกต\-จีวรํ สมาทาย ปกฺกมติ, ตสฺส พหิสีม\-คตสฺส เอวํ โหติ “อิเธวิมํ จีวรํ กาเรสฺสํ, น ปจฺเจสฺสนฺ”ติ, โส ตํ จีวรํ กาเรติ, ตสฺส ภิกฺขุโน นิฏฺฐานนฺติโก กถินุทฺธาโร. (1)}

\prose{ภิกฺขุ อตฺถตกถิโน วิปฺปกต\-จีวรํ สมาทาย ปกฺกมติ, ตสฺส พหิสีม\-คตสฺส เอวํ โหติ “เนวิมํ จีวรํ กาเรสฺสํ, น ปจฺเจสฺสนฺ”ติ, ตสฺส ภิกฺขุโน สนฺนิฏฺฐานนฺติโก กถินุทฺธาโร. (2)}

\prose{\csromanfolio{357}ภิกฺขุ อตฺถตกถิโน วิปฺปกต\-จีวรํ สมาทาย ปกฺกมติ, ตสฺส พหิสีม\-คตสฺส เอวํ โหติ “อิเธวิมํ จีวรํ กาเรสฺสํ, น ปจฺเจสฺสนฺ”ติ, โส ตํ จีวรํ กาเรติ, ตสฺส ตํ จีวรํ กยิรมานํ นสฺสติ, ตสฺส ภิกฺขุโน นาสนนฺติโก กถินุทฺธาโร. (3)}

\prose{ภิกฺขุ อตฺถตกถิโน วิปฺปกต\-จีวรํ สมาทาย ปกฺกมติ “ปจฺเจสฺสนฺ”ติ, โส พหิสีมคโต ตํ จีวรํ กาเรติ, โส กตจีวโร สุณาติ “อุพฺภตํ กิร ตสฺมิํ อาวาเส กถินนฺ”ติ, ตสฺส ภิกฺขุโน สวนนฺติโก กถินุทฺธาโร. (4)}

\prose{ภิกฺขุ อตฺถตกถิโน วิปฺปกต\-จีวรํ สมาทาย ปกฺกมติ “ปจฺเจสฺสนฺ”ติ, โส พหิสีมคโต ตํ จีวรํ กาเรติ, โส กตจีวโร “ปจฺเจสฺสํ ปจฺเจสฺสนฺ”ติ พหิทฺธา กถินุทฺธารํ วีตินาเมติ, ตสฺส ภิกฺขุโน สีมาติกฺกนฺติโก กถินุทฺธาโร. (5)}

\prosewithcloserruled{ภิกฺขุ อตฺถตกถิโน วิปฺปกต\-จีวรํ สมาทาย ปกฺกมติ “ปจฺเจสฺสนฺ”ติ, โส พหิสีมคโต ตํ จีวรํ กาเรติ, โส กตจีวโร “ปจฺเจสฺสํ ปจฺเจสฺสนฺ”ติ สมฺภุณาติ กถินุทฺธารํ, ตสฺส ภิกฺขุโน สห ภิกฺขูหิ กถินุทฺธาโร. (6)}{สมาทาย\-ฉกฺกํ.}
\csromanmatikaanchor{mtk.199}
\csromantocmarkref{section}{192. อาทายปนฺนรสก}{mtk.199}
\bookmark[dest={mtk.199},level=1]{192. อาทายปนฺนรสก}
\csromanheader{1}{192. \textbf{อาทาย}\-\textbf{ปนฺนรสก}}
\proseitem{315}{ภิกฺขุ อตฺถตกถิโน จีวรํ อาทาย ปกฺกมติ, ตสฺส พหิสีม\-คตสฺส เอวํ โหติ “อิเธวิมํ จีวรํ กาเรสฺสํ, น ปจฺเจสฺสนฺ”ติ, โส ตํ จีวรํ กาเรติ, ตสฺส ภิกฺขุโน นิฏฺฐานนฺติโก กถินุทฺธาโร. (1)}

\prose{ภิกฺขุ อตฺถตกถิโน จีวรํ อาทาย ปกฺกมติ, ตสฺส พหิสีม\-คตสฺส เอวํ โหติ “เนวิมํ จีวรํ กาเรสฺสํ, น ปจฺเจสฺสนฺ”ติ, ตสฺส ภิกฺขุโน สนฺนิฏฺฐานนฺติโก กถินุทฺธาโร. (2)}

\prose{ภิกฺขุ อตฺถตกถิโน จีวรํ อาทาย ปกฺกมติ, ตสฺส พหิสีม\-คตสฺส เอวํ โหติ “อิเธวิมํ จีวรํ กาเรสฺสํ, น ปจฺเจสฺสนฺ”ติ, โส \csromanfolio{358}ตํ จีวรํ กาเรติ, ตสฺส ตํ จีวรํ กยิรมานํ นสฺสติ, ตสฺส ภิกฺขุโน นาสนนฺติโก กถินุทฺธาโร. (3) ติกํ.}

\prose{ภิกฺขุ อตฺถตกถิโน จีวรํ อาทาย ปกฺกมติ “น ปจฺเจสฺสนฺ”ติ, ตสฺส พหิสีม\-คตสฺส เอวํ โหติ “อิเธวิมํ จีวรํ กาเรสฺสนฺ”ติ, โส ตํ จีวรํ กาเรติ, ตสฺส ภิกฺขุโน นิฏฺฐานนฺติโก กถินุทฺธาโร. (4)}

\prose{ภิกฺขุ อตฺถตกถิโน จีวรํ อาทาย ปกฺกมติ “น ปจฺเจสฺสนฺ”ติ, ตสฺส พหิสีม\-คตสฺส เอวํ โหติ “เนวิมํ จีวรํ กาเรสฺสนฺ”ติ, ตสฺส ภิกฺขุโน สนฺนิฏฺฐานนฺติโก กถินุทฺธาโร. (5)}

\prose{ภิกฺขุ อตฺถตกถิโน จีวรํ อาทาย ปกฺกมติ “น ปจฺเจสฺสนฺ”ติ, ตสฺส พหิสีม\-คตสฺส เอวํ โหติ “อิเธวิมํ จีวรํ กาเรสฺสนฺ”ติ, โส ตํ จีวรํ กาเรติ, ตสฺส ตํ จีวรํ กยิรมานํ นสฺสติ, ตสฺส ภิกฺขุโน นาสนนฺติโก กถินุทฺธาโร. (6) ติกํ.}

\prose{ภิกฺขุ อตฺถตกถิโน จีวรํ อาทาย ปกฺกมติ อนธิฏฺฐิเตน, เนวสฺส โหติ “ปจฺเจสฺสนฺ”ติ, น ปนสฺส โหติ “น ปจฺเจสฺสนฺ”ติ, ตสฺส พหิสีม\-คตสฺส เอวํ โหติ “อิเธวิมํ จีวรํ กาเรสฺสํ, น ปจฺเจสฺสนฺ”ติ, โส ตํ จีวรํ กาเรติ, ตสฺส ภิกฺขุโน นิฏฺฐานนฺติโก กถินุทฺธาโร. (7)}

\prose{ภิกฺขุ อตฺถตกถิโน จีวรํ อาทาย ปกฺกมติ อนธิฏฺฐิเตน, เนวสฺส โหติ “ปจฺเจสฺสนฺ”ติ, น ปนสฺส โหติ “น ปจฺเจสฺสนฺ”ติ, ตสฺส พหิสีม\-คตสฺส เอวํ โหติ “เนวิมํ จีวรํ กาเรสฺสํ, น ปจฺเจสฺสนฺ”ติ, ตสฺส ภิกฺขุโน สนฺนิฏฺฐานนฺติโก กถินุทฺธาโร. (8)}

\prose{ภิกฺขุ อตฺถตกถิโน จีวรํ อาทาย ปกฺกมติ อนธิฏฺฐิเตน, เนวสฺส โหติ “ปจฺเจสฺสนฺ”ติ, น ปนสฺส โหติ “น ปจฺเจสฺสนฺ”ติ, ตสฺส พหิสีม\-คตสฺส เอวํ โหติ “อิเธวิมํ จีวรํ กาเรสฺสํ, น ปจฺเจสฺสนฺ”ติ, โส ตํ จีวรํ กาเรติ, ตสฺส ตํ จีวรํ กยิรมานํ นสฺสติ, ตสฺส ภิกฺขุโน นาสนนฺติโก กถินุทฺธาโร. (9) ติกํ. \csromanfolio{359}ภิกฺขุ อตฺถตกถิโน จีวรํ อาทาย ปกฺกมติ “ปจฺเจสฺสนฺ”ติ, ตสฺส พหิสีม\-คตสฺส เอวํ โหติ “อิเธวิมํ จีวรํ กาเรสฺสํ, น ปจฺเจสฺสนฺ”ติ, โส ตํ จีวรํ กาเรติ, ตสฺส ภิกฺขุโน นิฏฺฐานนฺติโก กถินุทฺธาโร. (10) ภิกฺขุ อตฺถตกถิโน จีวรํ อาทาย ปกฺกมติ “ปจฺเจสฺสนฺ”ติ, ตสฺส พหิสีม\-คตสฺส เอวํ โหติ “เนวิมํ จีวรํ กาเรสฺสํ, น ปจฺเจสฺสนฺ”ติ, ตสฺส ภิกฺขุโน สนฺนิฏฺฐานนฺติโก กถินุทฺธาโร. (11) ภิกฺขุ อตฺถตกถิโน จีวรํ อาทาย ปกฺกมติ “ปจฺเจสฺสนฺ”ติ, ตสฺส พหิสีม\-คตสฺส เอวํ โหติ “อิเธวิมํ จีวรํ กาเรสฺสํ, น ปจฺเจสฺสนฺ”ติ, โส ตํ จีวรํ กาเรติ, ตสฺส ตํ จีวรํ กยิรมานํ นสฺสติ, ตสฺส ภิกฺขุโน นาสนนฺติโก กถินุทฺธาโร. (12)}

\prose{ภิกฺขุ อตฺถตกถิโน จีวรํ อาทาย ปกฺกมติ “ปจฺเจสฺสนฺ”ติ, โส พหิสีมคโต ตํ จีวรํ กาเรติ, โส กตจีวโร สุณาติ “อุพฺภตํ กิร ตสฺมิํ อาวาเส กถินนฺ”ติ, ตสฺส ภิกฺขุโน สวนนฺติโก กถินุทฺธาโร. (13)}

\prose{ภิกฺขุ อตฺถตกถิโน จีวรํ อาทาย ปกฺกมติ “ปจฺเจสฺสนฺ”ติ, โส พหิสีมคโต ตํ จีวรํ กาเรติ, โส กตจีวโร “ปจฺเจสฺสํ ปจฺเจสฺสนฺ”ติ พหิทฺธา กถินุทฺธารํ วีตินาเมติ, ตสฺส ภิกฺขุโน สีมาติกฺกนฺติโก กถินุทฺธาโร. (14)}

\prose{ภิกฺขุ อตฺถตกถิโน จีวรํ อาทาย ปกฺกมติ “ปจฺเจสฺสนฺ”ติ, โส พหิสีมคโต ตํ จีวรํ กาเรติ, โส กตจีวโร “ปจฺเจสฺสํ ปจฺเจสฺสนฺ”ติ สมฺภุณาติ กถินุทฺธารํ, ตสฺส ภิกฺขุโน สห ภิกฺขูหิ กถินุทฺธาโร. (15) ฉกฺกํ.}

\vspace{\tipitakaparskip}
\csromancenter{อาทาย\-ปนฺนรสกํ.}
\csromanmatikaanchor{mtk.200}
\csromantocmarkref{section}{193. สมาทายปนฺนรสกาทิ}{mtk.200}
\bookmark[dest={mtk.200},level=1]{193. สมาทายปนฺนรสกาทิ}
\csromanheader{1}{193. \textbf{สมาทาย}\-\textbf{ปนฺนรสกาทิ}}
\proseitem{316}{\csromanfolio{360}ภิกฺขุ อตฺถตกถิโน จีวรํ สมาทาย ปกฺกมติ ฯเปฯ.}

\prose{(อาทาย\-วาร\-สทิสํ เอวํ วิตฺถาเรตพฺพํ.)}

\prose{ภิกฺขุ อตฺถตกถิโน วิปฺปกต\-จีวรํ อาทาย ปกฺกมติ, ตสฺส พหิสีม\-คตสฺส เอวํ โหติ “อิเธวิมํ จีวรํ กาเรสฺสํ, น ปจฺเจสฺสนฺ”ติ, โส ตํ จีวรํ กาเรติ, ตสฺส ภิกฺขุโน นิฏฺฐานนฺติโก กถินุทฺธาโร ฯเปฯ.}

\vspace{\tipitakaparskip}
\csromancenter{(สมาทาย\-วาร\-สทิสํ เอวํ วิตฺถาเรตพฺพํ)}
\csromansectionrule
\csromanmatikaanchor{mtk.201}
\csromantocmarkref{section}{194. วิปฺปกตสมาทายปนฺนรสก}{mtk.201}
\bookmark[dest={mtk.201},level=1]{194. วิปฺปกตสมาทายปนฺนรสก}
\csromanheader{1}{194. \textbf{วิปฺปกต}\-\textbf{สมาทาย}\-\textbf{ปนฺนรสก}}
\proseitem{317}{ภิกฺขุ อตฺถตกถิโน วิปฺปกต\-จีวรํ สมาทาย ปกฺกมติ, ตสฺส พหิสีม\-คตสฺส เอวํ โหติ “อิเธวิมํ จีวรํ กาเรสฺสํ, น ปจฺเจสฺสนฺ”ติ, โส ตํ จีวรํ กาเรติ, ตสฺส ภิกฺขุโน นิฏฺฐานนฺติโก กถินุทฺธาโร. (1)}

\prose{ภิกฺขุ อตฺถตกถิโน วิปฺปกต\-จีวรํ สมาทาย ปกฺกมติ, ตสฺส พหิสีม\-คตสฺส เอวํ โหติ “เนวิมํ จีวรํ กาเรสฺสํ, น ปจฺเจสฺสนฺ”ติ, ตสฺส ภิกฺขุโน สนฺนิฏฺฐานนฺติโก กถินุทฺธาโร. (2)}

\prose{ภิกฺขุ อตฺถตกถิโน วิปฺปกต\-จีวรํ สมาทาย ปกฺกมติ, ตสฺส พหิสีม\-คตสฺส เอวํ โหติ “อิเธวิมํ จีวรํ กาเรสฺสํ, น ปจฺเจสฺสนฺ”ติ, โส ตํ จีวรํ กาเรติ, ตสฺส ตํ จีวรํ กยิรมานํ นสฺสติ, ตสฺส ภิกฺขุโน นาสนนฺติโก กถินุทฺธาโร. (3) ติกํ.}

\prose{ภิกฺขุ อตฺถตกถิโน วิปฺปกต\-จีวรํ สมาทาย ปกฺกมติ “น ปจฺเจสฺสนฺ”ติ, ตสฺส พหิสีม\-คตสฺส เอวํ โหติ “อิเธวิมํ จีวรํ กาเรสฺสนฺ”ติ, โส ตํ จีวรํ กาเรติ, ตสฺส ภิกฺขุโน นิฏฺฐานนฺติโก กถินุทฺธาโร. (4)}

\prose{ภิกฺขุ อตฺถตกถิโน วิปฺปกต\-จีวรํ สมาทาย ปกฺกมติ “น ปจฺเจสฺสนฺ”ติ, ตสฺส พหิสีม\-คตสฺส เอวํ โหติ “เนวิมํ จีวรํ กาเรสฺสนฺ”ติ, ตสฺส ภิกฺขุโน สนฺนิฏฺฐานนฺติโก กถินุทฺธาโร. (5)}

\prose{\csromanfolio{361}ภิกฺขุ อตฺถตกถิโน วิปฺปกต\-จีวรํ สมาทาย ปกฺกมติ “น ปจฺเจสฺสนฺ”ติ, ตสฺส พหิสีม\-คตสฺส เอวํ โหติ “อิเธวิมํ จีวรํ กาเรสฺสนฺ”ติ, โส ตํ จีวรํ กาเรติ, ตสฺส ตํ จีวรํ กยิรมานํ นสฺสติ, ตสฺส ภิกฺขุโน นาสนนฺติโก กถินุทฺธาโร. (6) ติกํ.}

\prose{ภิกฺขุ อตฺถตกถิโน วิปฺปกต\-จีวรํ สมาทาย ปกฺกมติ อนธิฏฺฐิเตน, เนวสฺส โหติ “ปจฺเจสฺสนฺ”ติ, น ปนสฺส โหติ “น ปจฺเจสฺสนฺ”ติ, ตสฺส พหิสีม\-คตสฺส เอวํ โหติ “อิเธวิมํ จีวรํ กาเรสฺสํ, น ปจฺเจสฺสนฺ”ติ, โส ตํ จีวรํ กาเรติ, ตสฺส ภิกฺขุโน นิฏฺฐานนฺติโก กถินุทฺธาโร. (7)}

\prose{ภิกฺขุ อตฺถตกถิโน วิปฺปกต\-จีวรํ สมาทาย ปกฺกมติ อนธิฏฺฐิเตน, เนวสฺส โหติ “ปจฺเจสฺสนฺ”ติ, น ปนสฺส โหติ “น ปจฺเจสฺสนฺ”ติ, ตสฺส พหิสีม\-คตสฺส เอวํ โหติ “เนวิมํ จีวรํ กาเรสฺสํ, น ปจฺเจสฺสนฺ”ติ, ตสฺส ภิกฺขุโน สนฺนิฏฺฐานนฺติโก กถินุทฺธาโร. (8)}

\prose{ภิกฺขุ อตฺถตกถิโน วิปฺปกต\-จีวรํ สมาทาย ปกฺกมติ อนธิฏฺฐิเตน, เนวสฺส โหติ “ปจฺเจสฺสนฺ”ติ, น ปนสฺส โหติ “น ปจฺเจสฺสนฺ”ติ, ตสฺส พหิสีม\-คตสฺส เอวํ โหติ “อิเธวิมํ จีวรํ กาเรสฺสํ, น ปจฺเจสฺสนฺ”ติ, โส ตํ จีวรํ กาเรติ, ตสฺส ตํ จีวรํ กยิรมานํ นสฺสติ, ตสฺส ภิกฺขุโน นาสนนฺติโก กถินุทฺธาโร. (9) ติกํ.}

\prose{ภิกฺขุ อตฺถตกถิโน วิปฺปกต\-จีวรํ สมาทาย ปกฺกมติ “ปจฺเจสฺสนฺ”ติ, ตสฺส พหิสีม\-คตสฺส เอวํ โหติ “อิเธวิมํ จีวรํ กาเรสฺสํ, น ปจฺเจสฺสนฺ”ติ, โส ตํ จีวรํ กาเรติ, ตสฺส ภิกฺขุโน นิฏฺฐานนฺติโก กถินุทฺธาโร. (10)}

\prose{ภิกฺขุ อตฺถตกถิโน วิปฺปกต\-จีวรํ สมาทาย ปกฺกมติ “ปจฺเจสฺสนฺ”ติ, ตสฺส พหิสีม\-คตสฺส เอวํ โหติ “เนวิมํ จีวรํ กาเรสฺสํ, น ปจฺเจสฺสนฺ”ติ, ตสฺส ภิกฺขุโน สนฺนิฏฺฐานนฺติโก กถินุทฺธาโร. (11)}

\csromanmatikaanchor{mtk.202}
\csromantocmarkref{section}{195. อานาสาโทฬสก}{mtk.202}
\bookmark[dest={mtk.202},level=1]{195. อานาสาโทฬสก}
\prose{ภิกฺขุ อตฺถตกถิโน วิปฺปกต\-จีวรํ สมาทาย ปกฺกมติ “ปจฺเจสฺสนฺ”ติ, ตสฺส พหิสีม\-คตสฺส เอวํ โหติ “อิเธวิมํ จีวรํ กาเรสฺสํ, น \csromanfolio{362}ปจฺเจสฺสนฺ”ติ, โส ตํ จีวรํ กาเรติ, ตสฺส ตํ จีวรํ กยิรมานํ นสฺสติ, ตสฺส ภิกฺขุโน นาสนนฺติโก กถินุทฺธาโร. (12)}

\prose{ภิกฺขุ อตฺถตกถิโน วิปฺปกต\-จีวรํ สมาทาย ปกฺกมติ “ปจฺเจสฺสนฺ”ติ, โส พหิสีมคโต ตํ จีวรํ กาเรติ, โส กตจีวโร สุณาติ “อุพฺภตํ กิร ตสฺมิํ อาวาเส กถินนฺ”ติ, ตสฺส ภิกฺขุโน สวนนฺติโก กถินุทฺธาโร. (13)}

\prose{ภิกฺขุ อตฺถตกถิโน วิปฺปกต\-จีวรํ สมาทาย ปกฺกมติ “ปจฺเจสฺสนฺ”ติ, โส พหิสีมคโต ตํ จีวรํ กาเรติ, โส กตจีวโร “ปจฺเจสฺสํ ปจฺเจสฺสนฺ”ติ พหิทฺธา กถินุทฺธารํ วีตินาเมติ, ตสฺส ภิกฺขุโน สีมาติกฺกนฺติโก กถินุทฺธาโร. (14)}

\prose{ภิกฺขุ อตฺถตกถิโน วิปฺปกต\-จีวรํ สมาทาย ปกฺกมติ “ปจฺเจสฺสนฺ”ติ, โส พหิสีมคโต ตํ จีวรํ กาเรติ, โส กตจีวโร “ปจฺเจสฺสํ ปจฺเจสฺสนฺ”ติ สมฺภุณาติ กถินุทฺธารํ, ตสฺส ภิกฺขุโน สห ภิกฺขูหิ กถินุทฺธาโร. (15) ฉกฺกํ.}

\vspace{\tipitakaparskip}
\csromancenter{สมาทาย\-ปนฺนรสกํ.}
\nitthitamruled{อาทายภาณวาโร.}
\csromanheader{2}{195. \textbf{อนาสาโทฬสก}}
\proseitem{318}{ภิกฺขุ อตฺถตกถิโน จีวราสาย ปกฺกมติ, โส พหิสีมคโต ตํ จีวราสํ ปยิรุปาสติ, อนาสาย ลภติ, อาสาย น ลภติ, ตสฺส เอวํ โหติ “อิเธวิมํ จีวรํ กาเรสฺสํ, น ปจฺเจสฺสนฺ”ติ, โส ตํ จีวรํ กาเรติ, ตสฺส ภิกฺขุโน นิฏฺฐานนฺติโก กถินุทฺธาโร. (1)}

\prose{ภิกฺขุ อตฺถตกถิโน จีวราสาย ปกฺกมติ, โส พหิสีมคโต ตํ จีวราสํ ปยิรุปาสติ, อนาสาย ลภติ, อาสาย น ลภติ, ตสฺส เอวํ โหติ “เนวิมํ จีวรํ กาเรสฺสํ, น ปจฺเจสฺสนฺ”ติ, ตสฺส ภิกฺขุโน สนฺนิฏฺฐานนฺติโก กถินุทฺธาโร. (2)}

\prose{\csromanfolio{363}ภิกฺขุ อตฺถตกถิโน จีวราสาย ปกฺกมติ, โส พหิสีมคโต ตํ จีวราสํ ปยิรุปาสติ, อนาสาย ลภติ, อาสาย น ลภติ, ตสฺส เอวํ โหติ “อิเธวิมํ จีวรํ กาเรสฺสํ, น ปจฺเจสฺสนฺ”ติ, โส ตํ จีวรํ กาเรติ, ตสฺส ตํ จีวรํ กยิรมานํ นสฺสติ, ตสฺส ภิกฺขุโน นาสนนฺติโก กถินุทฺธาโร. (3)}

\prose{ภิกฺขุ อตฺถตกถิโน จีวราสาย ปกฺกมติ, ตสฺส พหิสีม\-คตสฺส เอวํ โหติ “อิเธวิมํ จีวราสํ ปยิรุปาสิสฺสํ, น ปจฺเจสฺสนฺ”ติ, โส ตํ จีวราสํ ปยิรุปาสติ, ตสฺส สา จีวราสา อุปจฺฉิชฺชติ, ตสฺส ภิกฺขุโน อาสาวจฺเฉทิโก กถินุทฺธาโร. (4)}

\prose{ภิกฺขุ อตฺถตกถิโน จีวราสาย ปกฺกมติ “น ปจฺเจสฺสนฺ”ติ, โส พหิสีมคโต ตํ จีวราสํ ปยิรุปาสติ, อนาสาย ลภติ, อาสาย น ลภติ, ตสฺส เอวํ โหติ “อิเธวิมํ จีวรํ กาเรสฺสนฺ”ติ, โส ตํ จีวรํ กาเรติ, ตสฺส ภิกฺขุโน นิฏฺฐานนฺติโก กถินุทฺธาโร. (5)}

\prose{ภิกฺขุ อตฺถตกถิโน จีวราสาย ปกฺกมติ “น ปจฺเจสฺสนฺ”ติ, โส พหิสีมคโต ตํ จีวราสํ ปยิรุปาสติ, อนาสาย ลภติ, อาสาย น ลภติ, ตสฺส เอวํ โหติ “เนวิมํ จีวรํ กาเรสฺสนฺ”ติ, ตสฺส ภิกฺขุโน สนฺนิฏฺฐานนฺติโก กถินุทฺธาโร. (6)}

\prose{ภิกฺขุ อตฺถตกถิโน จีวราสาย ปกฺกมติ “น ปจฺเจสฺสนฺ”ติ, โส พหิสีมคโต ตํ จีวราสํ ปยิรุปาสติ, อนาสาย ลภติ, อาสาย น ลภติ, ตสฺส เอวํ โหติ “อิเธวิมํ จีวรํ กาเรสฺสนฺ”ติ, โส ตํ จีวรํ กาเรติ, ตสฺส ตํ จีวรํ กยิรมานํ นสฺสติ, ตสฺส ภิกฺขุโน นาสนนฺติโก กถินุทฺธาโร. (7)}

\prose{ภิกฺขุ อตฺถตกถิโน จีวราสาย ปกฺกมติ “น ปจฺเจสฺสนฺ”ติ, ตสฺส พหิสีม\-คตสฺส เอวํ โหติ “อิเธวิมํ จีวราสํ ปยิรุปาสิสฺสนฺ”ติ, โส ตํ จีวราสํ ปยิรุปาสติ, ตสฺส สา จีวราสา อุปจฺฉิชฺชติ, ตสฺส ภิกฺขุโน อาสาวจฺเฉทิโก กถินุทฺธาโร. (8)}

\prose{ภิกฺขุ อตฺถตกถิโน จีวราสาย ปกฺกมติ อนธิฏฺฐิเตน, เนวสฺส โหติ “ปจฺเจสฺสนฺ”ติ, น ปนสฺส โหติ “น ปจฺเจสฺสนฺ”ติ, โส พหิสีมคโต ตํ จีวราสํ ปยิรุปาสติ, อนาสาย ลภติ, อาสาย \csromanfolio{364}น ลภติ, ตสฺส เอวํ โหติ “อิเธวิมํ จีวรํ กาเรสฺสํ, น ปจฺเจสฺสนฺ”ติ, โส ตํ จีวรํ กาเรติ, ตสฺส ภิกฺขุโน นิฏฺฐานนฺติโก กถินุทฺธาโร. (9)}

\prose{ภิกฺขุ อตฺถตกถิโน จีวราสาย ปกฺกมติ อนธิฏฺฐิเตน, เนวสฺส โหติ “ปจฺเจสฺสนฺ”ติ, น ปนสฺส โหติ “น ปจฺเจสฺสนฺ”ติ, โส พหิสีมคโต ตํ จีวราสํ ปยิรุปาสติ, อนาสาย ลภติ, อาสาย น ลภติ, ตสฺส เอวํ โหติ “เนวิมํ จีวรํ กาเรสฺสํ, น ปจฺเจสฺสนฺ”ติ, ตสฺส ภิกฺขุโน สนฺนิฏฺฐานนฺติโก กถินุทฺธาโร. (10)}

\prose{ภิกฺขุ อตฺถตกถิโน จีวราสาย ปกฺกมติ อนธิฏฺฐิเตน, เนวสฺส โหติ “ปจฺเจสฺสนฺ”ติ, น ปนสฺส โหติ “น ปจฺเจสฺสนฺ”ติ, โส พหิสีมคโต ตํ จีวราสํ ปยิรุปาสติ, อนาสาย ลภติ, อาสาย น ลภติ, ตสฺส เอวํ โหติ “อิเธวิมํ จีวรํ กาเรสฺสํ, น ปจฺเจสฺสนฺ”ติ, โส ตํ จีวรํ กาเรติ, ตสฺส ตํ จีวรํ กยิรมานํ นสฺสติ, ตสฺส ภิกฺขุโน นาสนนฺติโก กถินุทฺธาโร. (11)}

\prosewithcloser{ภิกฺขุ อตฺถตกถิโน จีวราสาย ปกฺกมติ อนธิฏฺฐิเตน, เนวสฺส โหติ “ปจฺเจสฺสนฺ”ติ, น ปนสฺส โหติ “น ปจฺเจสฺสนฺ”ติ, ตสฺส พหิสีม\-คตสฺส เอวํ โหติ “อิเธวิมํ จีวราสํ ปยิรุปาสิสฺสํ, น ปจฺเจสฺสนฺ”ติ, โส ตํ จีวราสํ ปยิรุปาสติ, ตสฺส สา จีวราสา อุปจฺฉิชฺชติ, ตสฺส ภิกฺขุโน อาสาวจฺเฉทิโก กถินุทฺธาโร.}{(12) อนาสาโทฬสกํ\csromansharedfootnote{338dd9d64fd7ee1e}{[มิสฺสินฺคฺ โนเต 0]} นิฏฺฐิตํ.}
\csromanmatikaanchor{mtk.203}
\csromantocmarkref{section}{196. อาสาโทฬสก}{mtk.203}
\bookmark[dest={mtk.203},level=1]{196. อาสาโทฬสก}
\csromanheader{1}{196. \textbf{อาสาโทฬสก}}
\proseitem{319}{ภิกฺขุ อตฺถตกถิโน จีวราสาย ปกฺกมติ “ปจฺเจสฺสนฺ”ติ, โส พหิสีมคโต ตํ จีวราสํ ปยิรุปาสติ, อาสาย ลภติ, อนาสาย น ลภติ, ตสฺส เอวํ โหติ “อิเธวิมํ จีวรํ กาเรสฺสํ, น ปจฺเจสฺสนฺ”ติ, โส ตํ จีวรํ กาเรติ, ตสฺส ภิกฺขุโน นิฏฺฐานนฺติโก กถินุทฺธาโร. (1)}

\prose{\csromanfolio{365}ภิกฺขุ อตฺถตกถิโน จีวราสาย ปกฺกมติ “ปจฺเจสฺสนฺ”ติ, โส พหิสีมคโต ตํ จีวราสํ ปยิรุปาสติ, อาสาย ลภติ, อนาสาย น ลภติ, ตสฺส เอวํ โหติ “เนวิมํ จีวรํ กาเรสฺสํ, น ปจฺเจสฺสนฺ”ติ, ตสฺส ภิกฺขุโน สนฺนิฏฺฐานนฺติโก กถินุทฺธาโร. (2)}

\prose{ภิกฺขุ อตฺถตกถิโน จีวราสาย ปกฺกมติ “ปจฺเจสฺสนฺ”ติ, โส พหิสีมคโต ตํ จีวราสํ ปยิรุปาสติ, อาสาย ลภติ, อนาสาย น ลภติ, ตสฺส เอวํ โหติ “อิเธวิมํ จีวรํ กาเรสฺสํ, น ปจฺเจสฺสนฺ”ติ, โส ตํ จีวรํ กาเรติ, ตสฺส ตํ จีวรํ กยิรมานํ นสฺสติ, ตสฺส ภิกฺขุโน นาสนนฺติโก กถินุทฺธาโร. (3)}

\prose{ภิกฺขุ อตฺถตกถิโน จีวราสาย ปกฺกมติ “ปจฺเจสฺสนฺ”ติ, ตสฺส พหิสีม\-คตสฺส เอวํ โหติ “อิเธวิมํ จีวราสํ ปยิรุปาสิสฺสํ, น ปจฺเจสฺสนฺ”ติ, โส ตํ จีวราสํ ปยิรุปาสติ, ตสฺส สา จีวราสา อุปจฺฉิชฺชติ, ตสฺส ภิกฺขุโน อาสาวจฺเฉทิโก กถินุทฺธาโร. (4)}

\prose{ภิกฺขุ อตฺถตกถิโน จีวราสาย ปกฺกมติ “ปจฺเจสฺสนฺ”ติ, โส พหิสีมคโต สุณาติ “อุพฺภตํ กิร ตสฺมิํ อาวาเส กถินนฺ”ติ, ตสฺส เอวํ โหติ “ยโต ตสฺมิํ อาวาเส อุพฺภตํ กถินํ, อิเธวิมํ จีวราสํ ปยิรุปาสิสฺสนฺ”ติ, โส ตํ จีวราสํ ปยิรุปาสติ, อาสาย ลภติ, อนาสาย น ลภติ, ตสฺส เอวํ โหติ “อิเธวิมํ จีวรํ กาเรสฺสํ, น ปจฺเจสฺสนฺ”ติ, โส ตํ จีวรํ กาเรติ, ตสฺส ภิกฺขุโน นิฏฺฐานนฺติโก กถินุทฺธาโร. (5)}

\prose{ภิกฺขุ อตฺถตกถิโน จีวราสาย ปกฺกมติ “ปจฺเจสฺสนฺ”ติ, โส พหิสีมคโต สุณาติ “อุพฺภตํ กิร ตสฺมิํ อาวาเส กถินนฺ”ติ, ตสฺส เอวํ โหติ “ยโต ตสฺมิํ อาวาเส อุพฺภตํ กถินํ, อิเธวิมํ จีวราสํ ปยิรุปาสิสฺสนฺ”ติ, โส ตํ จีวราสํ ปยิรุปาสติ, อาสาย ลภติ, อนาสาย น ลภติ, ตสฺส เอวํ โหติ “เนวิมํ จีวรํ กาเรสฺสํ, น ปจฺเจสฺสนฺ”ติ, ตสฺส ภิกฺขุโน สนฺนิฏฺฐานนฺติโก กถินุทฺธาโร. (6)}

\prose{ภิกฺขุ อตฺถตกถิโน จีวราสาย ปกฺกมติ “ปจฺเจสฺสนฺ”ติ, โส พหิสีมคโต สุณาติ “อุพฺภตํ กิร ตสฺมิํ อาวาเส กถินนฺ”ติ, ตสฺส เอวํ โหติ “ยโต ตสฺมิํ อาวาเส อุพฺภตํ กถินํ, อิเธวิมํ \csromanfolio{366}จีวราสํ ปยิรุปาสิสฺสนฺ”ติ, โส ตํ จีวราสํ ปยิรุปาสติ, อาสาย ลภติ, อนาสาย น ลภติ, ตสฺส เอวํ โหติ “อิเธวิมํ จีวรํ กาเรสฺสํ, น ปจฺเจสฺสนฺ”ติ, โส ตํ จีวรํ กาเรติ, ตสฺส ตํ จีวรํ กยิรมานํ นสฺสติ, ตสฺส ภิกฺขุโน นาสนนฺติโก กถินุทฺธาโร. (7) ภิกฺขุ อตฺถตกถิโน จีวราสาย ปกฺกมติ “ปจฺเจสฺสนฺ”ติ, โส พหิสีมคโต สุณาติ “อุพฺภตํ กิร ตสฺมิํ อาวาเส กถินนฺ”ติ, ตสฺส เอวํ โหติ “ยโต ตสฺมิํ อาวาเส อุพฺภตํ กถินํ, อิเธวิมํ จีวราสํ ปยิรุปาสิสฺสํ, น ปจฺเจสฺสนฺ”ติ, โส ตํ จีวราสํ ปยิรุปาสติ, ตสฺส สา จีวราสา อุปจฺฉิชฺชติ, ตสฺส ภิกฺขุโน อาสาวจฺเฉทิโก กถินุทฺธาโร. (8)}

\prose{ภิกฺขุ อตฺถตกถิโน จีวราสาย ปกฺกมติ “ปจฺเจสฺสนฺ”ติ, โส พหิสีมคโต ตํ จีวราสํ ปยิรุปาสติ, อาสาย ลภติ, อนาสาย น ลภติ, โส ตํ จีวรํ กาเรติ, โส กตจีวโร สุณาติ “อุพฺภตํ กิร ตสฺมิํ อาวาเส กถินนฺ”ติ, ตสฺส ภิกฺขุโน สวนนฺติโก กถินุทฺธาโร. (9)}

\prose{ภิกฺขุ อตฺถตกถิโน จีวราสาย ปกฺกมติ “ปจฺเจสฺสนฺ”ติ, ตสฺส พหิสีม\-คตสฺส เอวํ โหติ “อิเธวิมํ จีวราสํ ปยิรุปาสิสฺสํ, น ปจฺเจสฺสนฺ”ติ, โส ตํ จีวราสํ ปยิรุปาสติ, ตสฺส สา จีวราสา อุปจฺฉิชฺชติ, ตสฺส ภิกฺขุโน อาสาวจฺเฉทิโก กถินุทฺธาโร. (10)}

\prose{ภิกฺขุ อตฺถตกถิโน จีวราสาย ปกฺกมติ “ปจฺเจสฺสนฺ”ติ, โส พหิสีมคโต ตํ จีวราสํ ปยิรุปาสติ, อาสาย ลภติ, อนาสาย น ลภติ, โส ตํ จีวรํ กาเรติ, โส กตจีวโร “ปจฺเจสฺสํ ปจฺเจสฺสนฺ”ติ พหิทฺธา กถินุทฺธารํ วีตินาเมติ, ตสฺส ภิกฺขุโน สีมาติกฺกนฺติโก กถินุทฺธาโร. (11)}

\prosewithcloser{ภิกฺขุ อตฺถตกถิโน จีวราสาย ปกฺกมติ “ปจฺเจสฺสนฺ”ติ, โส พหิสีมคโต ตํ จีวราสํ ปยิรุปาสติ, อาสาย ลภติ, อนาสาย น ลภติ, โส ตํ จีวรํ กาเรติ, โส กตจีวโร “ปจฺเจสฺสํ ปจฺเจสฺสนฺ”ติ สมฺภุณาติ กถินุทฺธารํ, ตสฺส ภิกฺขุโน สห ภิกฺขูหิ กถินุทฺธาโร. (12)}{อาสาโทฬสกํ นิฏฺฐิตํ.}
\csromanmatikaanchor{mtk.204}
\csromantocmarkref{section}{197. กรณียโทฬสก}{mtk.204}
\bookmark[dest={mtk.204},level=1]{197. กรณียโทฬสก}
\csromanheader{1}{197. \textbf{กรณีย}\-\textbf{โทฬสก}}
\proseitem{320}{\csromanfolio{367}ภิกฺขุ อตฺถตกถิโน เกนจิเทว กรณีเยน ปกฺกมติ, ตสฺส พหิสีม\-คตสฺส จีวราสา อุปฺปชฺชติ, โส ตํ จีวราสํ ปยิรุปาสติ, อนาสาย ลภติ, อาสาย น ลภติ, ตสฺส เอวํ โหติ “อิเธวิมํ จีวรํ กาเรสฺสํ, น ปจฺเจสฺสนฺ”ติ, โส ตํ จีวรํ กาเรติ, ตสฺส ภิกฺขุโน นิฏฺฐานนฺติโก กถินุทฺธาโร. (1)}

\prose{ภิกฺขุ อตฺถตกถิโน เกนจิเทว กรณีเยน ปกฺกมติ, ตสฺส พหิสีม\-คตสฺส จีวราสา อุปฺปชฺชติ, โส ตํ จีวราสํ ปยิรุปาสติ, อนาสาย ลภติ, อาสาย น ลภติ, ตสฺส เอวํ โหติ “เนวิมํ จีวรํ กาเรสฺสํ, น ปจฺเจสฺสนฺ”ติ, ตสฺส ภิกฺขุโน สนฺนิฏฺฐานนฺติโก กถินุทฺธาโร. (2)}

\prose{ภิกฺขุ อตฺถตกถิโน เกนจิเทว กรณีเยน ปกฺกมติ, ตสฺส พหิสีม\-คตสฺส จีวราสา อุปฺปชฺชติ, โส ตํ จีวราสํ ปยิรุปาสติ, อนาสาย ลภติ, อาสาย น ลภติ, ตสฺส เอวํ โหติ “อิเธวิมํ จีวรํ กาเรสฺสํ, น ปจฺเจสฺสนฺ”ติ, โส ตํ จีวรํ กาเรติ, ตสฺส ตํ จีวรํ กยิรมานํ นสฺสติ, ตสฺส ภิกฺขุโน นาสนนฺติโก กถินุทฺธาโร. (3)}

\prose{ภิกฺขุ อตฺถตกถิโน เกนจิเทว กรณีเยน ปกฺกมติ, ตสฺส พหิสีม\-คตสฺส จีวราสา อุปฺปชฺชติ, ตสฺส เอวํ โหติ “อิเธวิมํ จีวราสํ ปยิรุปาสิสฺสํ, น ปจฺเจสฺสนฺ”ติ, โส ตํ จีวราสํ ปยิรุปาสติ, ตสฺส สา จีวราสา อุปจฺฉิชฺชติ, ตสฺส ภิกฺขุโน อาสาวจฺเฉทิโก กถินุทฺธาโร. (4)}

\prose{ภิกฺขุ อตฺถตกถิโน เกนจิเทว กรณีเยน ปกฺกมติ “น ปจฺเจสฺสนฺ”ติ, ตสฺส พหิสีม\-คตสฺส จีวราสา อุปฺปชฺชติ, โส ตํ จีวราสํ ปยิรุปาสติ, อนาสาย ลภติ, อาสาย น ลภติ, ตสฺส เอวํ โหติ “อิเธวิมํ จีวรํ กาเรสฺสนฺ”ติ, โส ตํ จีวรํ กาเรติ, ตสฺส ภิกฺขุโน นิฏฺฐานนฺติโก กถินุทฺธาโร. (5)}

\prose{ภิกฺขุ อตฺถตกถิโน เกนจิเทว กรณีเยน ปกฺกมติ “น ปจฺเจสฺสนฺ”ติ, ตสฺส พหิสีม\-คตสฺส จีวราสา อุปฺปชฺชติ, โส ตํ จีวราสํ ปยิรุปาสติ, อนาสาย ลภติ, อาสาย น ลภติ, ตสฺส เอวํ \csromanfolio{368}โหติ “เนวิมํ จีวรํ กาเรสฺสนฺ”ติ, ตสฺส ภิกฺขุโน สนฺนิฏฺฐานนฺติโก กถินุทฺธาโร. (6)}

\prose{ภิกฺขุ อตฺถตกถิโน เกนจิเทว กรณีเยน ปกฺกมติ “น ปจฺเจสฺสนฺ”ติ, ตสฺส พหิสีม\-คตสฺส จีวราสา อุปฺปชฺชติ, โส ตํ จีวราสํ ปยิรุปาสติ, อนาสาย ลภติ, อาสาย น ลภติ, ตสฺส เอวํ โหติ “อิเธวิมํ จีวรํ กาเรสฺสนฺ”ติ, โส ตํ จีวรํ กาเรติ, ตสฺส ตํ จีวรํ กยิรมานํ นสฺสติ, ตสฺส ภิกฺขุโน นาสนนฺติโก กถินุทฺธาโร. (7)}

\prose{ภิกฺขุ อตฺถตกถิโน เกนจิเทว กรณีเยน ปกฺกมติ “น ปจฺเจสฺสนฺ”ติ, ตสฺส พหิสีม\-คตสฺส จีวราสา อุปฺปชฺชติ, ตสฺส เอวํ โหติ “อิเธวิมํ จีวราสํ ปยิรุปาสิสฺสนฺ”ติ, โส ตํ จีวราสํ ปยิรุปาสติ, ตสฺส สา จีวราสา อุปจฺฉิชฺชติ, ตสฺส ภิกฺขุโน อาสาวจฺเฉทิโก กถินุทฺธาโร. (8)}

\prose{ภิกฺขุ อตฺถตกถิโน เกนจิเทว กรณีเยน ปกฺกมติ อนธิฏฺฐิเตน, เนวสฺส โหติ “ปจฺเจสฺสนฺ”ติ, น ปนสฺส โหติ “น ปจฺเจสฺสนฺ”ติ, ตสฺส พหิสีม\-คตสฺส จีวราสา อุปฺปชฺชติ, โส ตํ จีวราสํ ปยิรุปาสติ, อนาสาย ลภติ, อาสาย น ลภติ, ตสฺส เอวํ โหติ “อิเธวิมํ จีวรํ กาเรสฺสํ, น ปจฺเจสฺสนฺ”ติ, โส ตํ จีวรํ กาเรติ, ตสฺส ภิกฺขุโน นิฏฺฐานนฺติโก กถินุทฺธาโร. (9)}

\prose{ภิกฺขุ อตฺถตกถิโน เกนจิเทว กรณีเยน ปกฺกมติ อนธิฏฺฐิเตน, เนวสฺส โหติ “ปจฺเจสฺสนฺ”ติ, น ปนสฺส โหติ “น ปจฺเจสฺสนฺ”ติ, ตสฺส พหิสีม\-คตสฺส จีวราสา อุปฺปชฺชติ, โส ตํ จีวราสํ ปยิรุปาสติ, อนาสาย ลภติ, อาสาย น ลภติ, ตสฺส เอวํ โหติ “เนวิมํ จีวรํ กาเรสฺสํ, น ปจฺเจสฺสนฺ”ติ, ตสฺส ภิกฺขุโน สนฺนิฏฺฐานนฺติโก กถินุทฺธาโร. (10)}

\prose{ภิกฺขุ อตฺถตกถิโน เกนจิเทว กรณีเยน ปกฺกมติ อนธิฏฺฐิเตน, เนวสฺส โหติ “ปจฺเจสฺสนฺ”ติ, น ปนสฺส โหติ “น ปจฺเจสฺสนฺ”ติ, ตสฺส พหิสีม\-คตสฺส จีวราสา อุปฺปชฺชติ, โส ตํ จีวราสํ ปยิรุปาสติ, อนาสาย ลภติ, อาสาย น ลภติ, ตสฺส เอวํ โหติ “อิเธวิมํ จีวรํ กาเรสฺสํ, น ปจฺเจสฺสนฺ”ติ, โส ตํ จีวรํ กาเรติ, ตสฺส \csromanfolio{369}ตํ จีวรํ กยิรมานํ นสฺสติ, ตสฺส ภิกฺขุโน นาสนนฺติโก กถินุทฺธาโร. (11)}

\prosewithcloserruled{ภิกฺขุ อตฺถตกถิโน เกนจิเทว กรณีเยน ปกฺกมติ อนธิฏฺฐิเตน, เนวสฺส โหติ “ปจฺเจสฺสนฺ”ติ, น ปนสฺส โหติ “น ปจฺเจสฺสนฺ”ติ, ตสฺส พหิสีม\-คตสฺส จีวราสา อุปฺปชฺชติ, ตสฺส เอวํ โหติ “อิเธวิมํ จีวราสํ ปยิรุปาสิสฺสํ, น ปจฺเจสฺสนฺ”ติ, โส ตํ จีวราสํ ปยิรุปาสติ, ตสฺส สา จีวราสา อุปจฺฉิชฺชติ, ตสฺส ภิกฺขุโน อาสาวจฺเฉทิโก กถินุทฺธาโร. (12)}{กรณีย\-โทฬสกํ นิฏฺฐิตํ.}
\csromanmatikaanchor{mtk.205}
\csromantocmarkref{section}{198. อปวิลายนนวก}{mtk.205}
\bookmark[dest={mtk.205},level=1]{198. อปวิลายนนวก}
\csromanheader{1}{198. \textbf{อปวิลายน}\-\textbf{นวก}}
\proseitem{321}{ภิกฺขุ อตฺถตกถิโน ทิสํคมิโก ปกฺกมติ จีวร\-ปฏิวีสํ อปวิลายมาโน\csromansharedfootnote{d0eb18ee193c3a6c}{อปวินยมาโน (สี), อปจินยมาโน (ก)}, ตเมนํ ทิสํคตํ ภิกฺขู ปุจฺฉนฺติ “กหํ ตฺวํ อาวุโส วสฺสํวุฏฺโฐ, กตฺถ จ เต จีวรปฏิวีโส”ติ. โส เอวํ วเทติ “อมุกสฺมิํ อาวาเส วสฺสํวุฏฺโฐมฺหิ, ตตฺถ จ เม จีวรปฏิวีโส”ติ. เต เอวํ วทนฺติ “คจฺฉาวุโส ตํ จีวรํ อาหร, มยํ เต อิธ จีวรํ กริสฺสามา”ติ. โส ตํ อาวาสํ คนฺตฺวา ภิกฺขู ปุจฺฉติ “กหํ เม อาวุโส จีวรปฏิวีโส”ติ. เต เอวํ วทนฺติ “อยํ เต อาวุโส จีวรปฏิวีโส, กหํ คมิสฺสสี”ติ. โส เอวํ วเทติ “อมุกํ นาม\csromansharedfootnote{8a85c58a4a652678}{อมุกญฺจ (ก)} อาวาสํ คมิสฺสามิ, ตตฺถ เม ภิกฺขู จีวรํ กริสฺสนฺตี”ติ. เต เอวํ วทนฺติ “อลํ อาวุโส มา อคมาสิ, มยํ เต อิธ จีวรํ กริสฺสามา”ติ. ตสฺส เอวํ โหติ “อิเธวิมํ จีวรํ กาเรสฺสํ, น ปจฺเจสฺสนฺ”ติ, โส ตํ จีวรํ กาเรติ, ตสฺส ภิกฺขุโน นิฏฺฐานนฺติโก กถินุทฺธาโร. (1)}

\prose{ภิกฺขุ อตฺถตกถิโน ทิสํคมิโก ปกฺกมติ ฯเปฯ “เนวิมํ จีวรํ กาเรสฺสํ, น ปจฺเจสฺสนฺ”ติ, ตสฺส ภิกฺขุโน สนฺนิฏฺฐานนฺติโก กถินุทฺธาโร. (2)}

\prose{\csromanfolio{370}ภิกฺขุ อตฺถตกถิโน ทิสํคมิโก ปกฺกมติ ฯเปฯ “อิเธวิมํ จีวรํ กาเรสฺสํ, น ปจฺเจสฺสนฺ”ติ, โส ตํ จีวรํ กาเรติ, ตสฺส ตํ จีวรํ กยิรมานํ นสฺสติ, ตสฺส ภิกฺขุโน นาสนนฺติโก กถินุทฺธาโร. (3)}

\proseitem{322}{ภิกฺขุ อตฺถตกถิโน ทิสํคมิโก ปกฺกมติ จีวร\-ปฏิวีสํ อปวิลายมาโน, ตเมนํ ทิสํคตํ ภิกฺขู ปุจฺฉนฺติ “กหํ ตฺวํ อาวุโส วสฺสํวุฏฺโฐ, กตฺถ จ เต จีวรปฏิวีโส”ติ. โส เอวํ วเทติ “อมุกสฺมิํ อาวาเส วสฺสํวุฏฺโฐมฺหิ, ตตฺถ จ เม จีวรปฏิวีโส”ติ. เต เอวํ วทนฺติ “คจฺฉาวุโส ตํ จีวรํ อาหร, มยํ เต อิธ จีวรํ กริสฺสามา”ติ. โส ตํ อาวาสํ คนฺตฺวา ภิกฺขู ปุจฺฉติ “กหํ เม อาวุโส จีวรปฏิวีโส”ติ. เต เอวํ วทนฺติ “อยํ เต อาวุโส จีวรปฏิวีโส”ติ. โส ตํ จีวรํ อาทาย ตํ อาวาสํ คจฺฉติ, ตเมนํ อนฺตรามคฺเค ภิกฺขู ปุจฺฉนฺติ “อาวุโส กหํ คมิสฺสสี”ติ. โส เอวํ วเทติ “อมุกํ นาม อาวาสํ คมิสฺสามิ, ตตฺถ เม ภิกฺขู จีวรํ กริสฺสนฺตี”ติ. เต เอวํ วทนฺติ “อลํ อาวุโส มา อคมาสิ, มยํ เต อิธ จีวรํ กริสฺสามา”ติ. ตสฺส เอวํ โหติ “อิเธวิมํ จีวรํ กาเรสฺสํ, น ปจฺเจสฺสนฺ”ติ, โส ตํ จีวรํ กาเรติ, ตสฺส ภิกฺขุโน นิฏฺฐานนฺติโก กถินุทฺธาโร. (4)}

\prose{ภิกฺขุ อตฺถตกถิโน ทิสํคมิโก ปกฺกมติ จีวร\-ปฏิวีสํ อปวิลายมาโน, ตเมนํ ทิสํคตํ ภิกฺขู ปุจฺฉนฺติ “กหํ ตฺวํ อาวุโส วสฺสํวุฏฺโฐ, กตฺถ จ เต จีวรปฏิวีโส”ติ. โส เอวํ วเทติ “อมุกสฺมิํ อาวาเส วสฺสํวุฏฺโฐมฺหิ, ตตฺถ จ เม จีวรปฏิวีโส”ติ. เต เอวํ วทนฺติ “คจฺฉาวุโส ตํ จีวรํ อาหร, มยํ เต อิธ จีวรํ กริสฺสามา”ติ. โส ตํ อาวาสํ คนฺตฺวา ภิกฺขู ปุจฺฉติ “กหํ เม อาวุโส จีวรปฏิวีโส”ติ. เต เอวํ วทนฺติ “อยํ เต อาวุโส จีวรปฏิวีโส”ติ. โส ตํ จีวรํ อาทาย ตํ อาวาสํ คจฺฉติ, ตเมนํ อนฺตรามคฺเค ภิกฺขู ปุจฺฉนฺติ “อาวุโส กหํ คมิสฺสสี”ติ. โส เอวํ วเทติ “อมุกํ นาม อาวาสํ คมิสฺสามิ, ตตฺถ เม ภิกฺขู จีวรํ กริสฺสนฺตี”ติ. เต เอวํ วทนฺติ “อลํ อาวุโส มา อคมาสิ, มยํ เต อิธ จีวรํ กริสฺสามา”ติ. ตสฺส เอวํ โหติ “เนวิมํ จีวรํ กาเรสฺสํ, น ปจฺเจสฺสนฺ”ติ, ตสฺส ภิกฺขุโน สนฺนิฏฺฐานนฺติโก กถินุทฺธาโร. (5)}

\prose{\csromanfolio{371}ภิกฺขุ อตฺถตกถิโน ทิสํคมิโก ปกฺกมติ ฯเปฯ “อิเธวิมํ จีวรํ กาเรสฺสํ, น ปจฺเจสฺสนฺ”ติ, โส ตํ จีวรํ กาเรติ, ตสฺส ตํ จีวรํ กยิรมานํ นสฺสติ, ตสฺส ภิกฺขุโน นาสนนฺติโก กถินุทฺธาโร. (6)}

\proseitem{323}{ภิกฺขุ อตฺถตกถิโน ทิสํคมิโก ปกฺกมติ จีวร\-ปฏิวีสํ อปวิลายมาโน, ตเมนํ ทิสํคตํ ภิกฺขู ปุจฺฉนฺติ “กหํ ตฺวํ อาวุโส วสฺสํวุฏฺโฐ, กตฺถ จ เต จีวรปฏิวีโส”ติ. โส เอวํ วเทติ “อมุกสฺมิํ อาวาเส วสฺสํวุฏฺโฐมฺหิ, ตตฺถ จ เม จีวรปฏิวีโส”ติ. เต เอวํ วทนฺติ “คจฺฉาวุโส ตํ จีวรํ อาหร, มยํ เต อิธ จีวรํ กริสฺสามา”ติ. โส ตํ อาวาสํ คนฺตฺวา ภิกฺขู ปุจฺฉติ “กหํ เม อาวุโส จีวรปฏิวีโส”ติ. เต เอวํ วทนฺติ “อยํ เต อาวุโส จีวรปฏิวีโส”ติ. โส ตํ จีวรํ อาทาย ตํ อาวาสํ คจฺฉติ, ตสฺส ตํ อาวาสํ คตสฺส เอวํ โหติ “อิเธวิมํ จีวรํ กาเรสฺสํ, น ปจฺเจสฺสนฺ”ติ, โส ตํ จีวรํ กาเรติ, ตสฺส ภิกฺขุโน นิฏฺฐานนฺติโก กถินุทฺธาโร. (7)}

\prose{ภิกฺขุ อตฺถตกถิโน ทิสํคมิโก ปกฺกมติ ฯเปฯ “เนวิมํ จีวรํ กาเรสฺสํ, น ปจฺเจสฺสนฺ”ติ, ตสฺส ภิกฺขุโน สนฺนิฏฺฐานนฺติโก กถินุทฺธาโร. (8)}

\prosewithcloserruled{ภิกฺขุ อตฺถตกถิโน ทิสํคมิโก ปกฺกมติ ฯเปฯ “อิเธวิมํ จีวรํ กาเรสฺสํ, น ปจฺเจสฺสนฺ”ติ, โส ตํ จีวรํ กาเรติ, ตสฺส ตํ จีวรํ กยิรมานํ นสฺสติ, ตสฺส ภิกฺขุโน นาสนนฺติโก กถินุทฺธาโร. (9)}{อปวิลายน\-นวกํ นิฏฺฐิตํ.}
\csromanmatikaanchor{mtk.206}
\csromantocmarkref{section}{199. ผาสุวิหารปญฺจก}{mtk.206}
\bookmark[dest={mtk.206},level=1]{199. ผาสุวิหารปญฺจก}
\csromanheader{1}{199. \textbf{ผาสุวิหาร}\-\textbf{ปญฺจก}}
\proseitem{324}{ภิกฺขุ อตฺถตกถิโน ผาสุวิหาริโก จีวรํ อาทาย ปกฺกมติ “อมุกํ นาม อาวาสํ คมิสฺสามิ, ตตฺถ เม ผาสุ ภวิสฺสติ วสิสฺสามิ, โน เจ เม ผาสุ ภวิสฺสติ, อมุกํ นาม อาวาสํ คมิสฺสามิ, ตตฺถ เม ผาสุ ภวิสฺสติ วสิสฺสามิ, โน เจ เม ผาสุ ภวิสฺสติ, อมุกํ นาม อาวาสํ คมิสฺสามิ, ตตฺถ เม ผาสุ \csromanfolio{372}ภวิสฺสติ วสิสฺสามิ, โน เจ เม ผาสุ ภวิสฺสติ ปจฺเจสฺสนฺ”ติ, ตสฺส พหิสีม\-คตสฺส เอวํ โหติ “อิเธวิมํ จีวรํ กาเรสฺสํ, น ปจฺเจสฺสนฺ”ติ, โส ตํ จีวรํ กาเรติ, ตสฺส ภิกฺขุโน นิฏฺฐานนฺติโก กถินุทฺธาโร. (1)}

\prose{ภิกฺขุ อตฺถตกถิโน ผาสุวิหาริโก จีวรํ อาทาย ปกฺกมติ “อมุกํ นาม อาวาสํ คมิสฺสามิ, ตตฺถ เม ผาสุ ภวิสฺสติ วสิสฺสามิ, โน เจ เม ผาสุ ภวิสฺสติ, อมุกํ นาม อาวาสํ คมิสฺสามิ, ตตฺถ เม ผาสุ ภวิสฺสติ วสิสฺสามิ, โน เจ เม ผาสุ ภวิสฺสติ, อมุกํ นาม อาวาสํ คมิสฺสามิ, ตตฺถ เม ผาสุ ภวิสฺสติ วสิสฺสามิ, โน เจ เม ผาสุ ภวิสฺสติ ปจฺเจสฺสนฺ”ติ, ตสฺส พหิสีม\-คตสฺส เอวํ โหติ “เนวิมํ จีวรํ กาเรสฺสํ, น ปจฺเจสฺสนฺ”ติ, ตสฺส ภิกฺขุโน สนฺนิฏฺฐานนฺติโก กถินุทฺธาโร. (2)}

\prose{ภิกฺขุ อตฺถตกถิโน ผาสุวิหาริโก จีวรํ อาทาย ปกฺกมติ “อมุกํ นาม อาวาสํ คมิสฺสามิ, ตตฺถ เม ผาสุ ภวิสฺสติ วสิสฺสามิ, โน เจ เม ผาสุ ภวิสฺสติ, อมุกํ นาม อาวาสํ คมิสฺสามิ, ตตฺถ เม ผาสุ ภวิสฺสติ วสิสฺสามิ, โน เจ เม ผาสุ ภวิสฺสติ, อมุกํ นาม อาวาสํ คมิสฺสามิ, ตตฺถ เม ผาสุ ภวิสฺสติ วสิสฺสามิ, โน เจ เม ผาสุ ภวิสฺสติ ปจฺเจสฺสนฺ”ติ, ตสฺส พหิสีม\-คตสฺส เอวํ โหติ “อิเธวิมํ จีวรํ กาเรสฺสํ, น ปจฺเจสฺสนฺ”ติ, โส ตํ จีวรํ กาเรติ, ตสฺส ตํ จีวรํ กยิรมานํ นสฺสติ, ตสฺส ภิกฺขุโน นาสนนฺติโก กถินุทฺธาโร. (3)}

\prose{ภิกฺขุ อตฺถตกถิโน ผาสุวิหาริโก จีวรํ อาทาย ปกฺกมติ “อมุกํ นาม อาวาสํ คมิสฺสามิ, ตตฺถ เม ผาสุ ภวิสฺสติ วสิสฺสามิ, โน เจ เม ผาสุ ภวิสฺสติ, อมุกํ นาม อาวาสํ คมิสฺสามิ, ตตฺถ เม ผาสุ ภวิสฺสติ วสิสฺสามิ, โน เจ เม ผาสุ ภวิสฺสติ, อมุกํ นาม อาวาสํ คมิสฺสามิ, ตตฺถ เม ผาสุ ภวิสฺสติ วสิสฺสามิ, โน เจ เม ผาสุ ภวิสฺสติ ปจฺเจสฺสนฺ”ติ, โส พหิสีมคโต ตํ จีวรํ กาเรติ, โส กตจีวโร “ปจฺเจสฺสํ ปจฺเจสฺสนฺ”ติ พหิทฺธา กถินุทฺธารํ วีตินาเมติ, ตสฺส ภิกฺขุโน สีมาติกฺกนฺติโก กถินุทฺธาโร. (4)}

\prosewithcloserruled{\csromanfolio{373}ภิกฺขุ อตฺถตกถิโน ผาสุวิหาริโก จีวรํ อาทาย ปกฺกมติ “อมุกํ นาม อาวาสํ คมิสฺสามิ, ตตฺถ เม ผาสุ ภวิสฺสติ วสิสฺสามิ, โน เจ เม ผาสุ ภวิสฺสติ, อมุกํ นาม อาวาสํ คมิสฺสามิ, ตตฺถ เม ผาสุ ภวิสฺสติ วสิสฺสามิ, โน เจ เม ผาสุ ภวิสฺสติ, อมุกํ นาม อาวาสํ คมิสฺสามิ, ตตฺถ เม ผาสุ ภวิสฺสติ วสิสฺสามิ, โน เจ เม ผาสุ ภวิสฺสติ ปจฺเจสฺสนฺ”ติ, โส พหิสีมคโต ตํ จีวรํ กาเรติ, โส กตจีวโร “ปจฺเจสฺสํ ปจฺเจสฺสนฺ”ติ สมฺภุณาติ กถินุทฺธารํ, ตสฺส ภิกฺขุโน สห ภิกฺขูหิ กถินุทฺธาโร. (5)}{ผาสุวิหาร\-ปญฺจกํ นิฏฺฐิตํ.}
\csromanmatikaanchor{mtk.207}
\csromantocmarkref{section}{200. ปลิโพธาปลิโพธกถา}{mtk.207}
\bookmark[dest={mtk.207},level=1]{200. ปลิโพธาปลิโพธกถา}
\csromanheader{1}{200. \textbf{ปลิโพธา}\-\textbf{ปลิโพธ}\-\textbf{กถา}}
\proseitem{325}{เทฺวเม ภิกฺขเว กถินสฺส ปลิโพธา, เทฺว อปลิโพธา. กตเม จ ภิกฺขเว เทฺว กถินสฺส ปลิโพธา, อาวาสปลิโพโธ จ จีวร\-ปลิโพโธ จ. กถญฺจ ภิกฺขเว อาวาสปลิโพโธ โหติ, อิธ ภิกฺขเว ภิกฺขุ วสติ วา ตสฺมิํ อาวาเส, สาเปกฺโข วา ปกฺกมติ “ปจฺเจสฺสนฺ”ติ, เอวํ โข ภิกฺขเว อาวาสปลิโพโธ โหติ. กถญฺจ ภิกฺขเว จีวร\-ปลิโพโธ โหติ, อิธ ภิกฺขเว ภิกฺขุโน จีวรํ อกตํ วา โหติ วิปฺปกตํ วา, จีวราสา วา อนุปจฺฉินฺนา, เอวํ โข ภิกฺขเว จีวร\-ปลิโพโธ โหติ. อิเม โข ภิกฺขเว เทฺว กถินสฺส ปลิโพธา.}

\csromanmatikaanchor{mtk.208}
\csromantocmarkref{section}{201. อุทฺทานคาถา}{mtk.208}
\bookmark[dest={mtk.208},level=1]{201. อุทฺทานคาถา}
\prose{กตเม จ ภิกฺขเว เทฺว กถินสฺส อปลิโพธา, อาวาส\-อปลิโพโธ จ จีวร\-อปลิโพโธ จ. กถญฺจ ภิกฺขเว อาวาส\-อปลิโพโธ โหติ, อิธ ภิกฺขเว ภิกฺขุ ปกฺกมติ ตมฺหา อาวาสา จตฺเตน วนฺเตน มุตฺเตน อนเปกฺโข\csromansharedfootnote{1d7099367daa8cdf}{อนเปกฺเขน (ก)} “น ปจฺเจสฺสนฺ”ติ, เอวํ โข ภิกฺขเว อาวาส\-อปลิโพโธ โหติ. กถญฺจ ภิกฺขเว จีวร\-อปลิโพโธ โหติ, อิธ ภิกฺขเว ภิกฺขุโน จีวรํ กตํ วา โหติ นฏฺฐํ วา วินฏฺฐํ วา \csromanfolio{374}ทฑฺฒํ วา, จีวราสา วา อุปจฺฉินฺนา, เอวํ โข ภิกฺขเว จีวร\-อปลิโพโธ โหติ. อิเม โข ภิกฺขเว เทฺว กถินสฺส อปลิโพธาติ.}

\vspace{\tipitakaparskip}
\csromancenter{กถิน\-กฺขนฺธโก นิฏฺฐิโต สตฺตโม.}
\csromancenter{201. \textbf{ตสฺสุทฺทานํ}}
\setlength{\csromangathapagewidth}{0pt}
\setlength{\csromangathawakwidth}{0pt}
\csromangathameasuregroup{ติํส ปาเวยฺยกา ภิกฺขู, \\ วสฺสํวุฏฺโฐกปุณฺเณหิ, \\ อิทํ วตฺถุ กถินสฺส, \\ อนามนฺตา อสมาจารา, \\ ยาวทตฺถญฺจ อุปฺปาโท, \\ ญตฺติ เอวตฺถตญฺเจว, \\ อุลฺลิขิ โธวนา เจว, \\ พนฺธโน วฏฺฏิ กณฺฑุส, \\ ปริภณฺฑํ โอวทฺเธยฺยํ, \\ กุกฺกุ สนฺนิธิ นิสฺสคฺคิ, \\ อญฺญตฺร ปญฺจาติเรเก, \\ นาญฺญตฺร ปุคฺคลา สมฺมา, \\ กถินานตฺถตํ โหติ, \\ อหตากปฺปปิโลติ, \\ อนิมิตฺตาปริกถา, \\ อนิสฺสคฺคิ กปฺปกเต, \\ ปญฺจเก วาติเรเก วา, \\ ปุคฺคลสฺสตฺถารา สมฺมา,}{\csromangathabat{ติํส ปาเวยฺยกา ภิกฺขู,}{สาเกตุกฺกณฺฐิตา วสุํ.} \\ \csromangathabat{วสฺสํวุฏฺโฐกปุณฺเณหิ,}{อคมุํ ชินทสฺสนํ.} \\ \csromangathabat{อิทํ วตฺถุ กถินสฺส,}{กปฺปิสฺสนฺติ จ ปญฺจกา.} \\ \csromangathabat{อนามนฺตา อสมาจารา,}{ตเถว คณโภชนํ.} \\ \csromangathabat{ยาวทตฺถญฺจ อุปฺปาโท,}{อตฺถตานํ ภวิสฺสติ.} \\ \csromangathabat{ญตฺติ เอวตฺถตญฺเจว,}{เอวญฺเจว อนตฺถตํ.} \\ \csromangathabat{อุลฺลิขิ โธวนา เจว,}{วิจารณญฺจ เฉทนํ.} \\ \csromangathabat{พนฺธโน วฏฺฏิ กณฺฑุส,}{ทฬฺหีกมฺมานุวาติกา.} \\ \csromangathabat{ปริภณฺฑํ โอวทฺเธยฺยํ,}{มทฺทนา นิมิตฺตํ กถา.} \\ \csromangathabat{กุกฺกุ สนฺนิธิ นิสฺสคฺคิ,}{น กปฺปญฺญตฺร เต ตโย.} \\ \csromangathabat{อญฺญตฺร ปญฺจาติเรเก,}{สญฺฉินฺเนน สมณฺฑลี.} \\ \csromangathabat{นาญฺญตฺร ปุคฺคลา สมฺมา,}{นิสฺสีมฏฺโฐนุโมทติ.} \\ \csromangathabat{กถินานตฺถตํ โหติ,}{เอวํ พุทฺเธน เทสิตํ.} \\ \csromangathabat{อหตากปฺปปิโลติ,}{ปํสุ ปาปณิกาย จ.} \\ \csromangathabat{อนิมิตฺตาปริกถา,}{อกุกฺกุ จ อสนฺนิธิ.} \\ \csromangathabat{อนิสฺสคฺคิ กปฺปกเต,}{ตถา ติจีวเรน จ.} \\ \csromangathabat{ปญฺจเก วาติเรเก วา,}{ฉินฺเน สมณฺฑลีกเต.} \\ \csromangathabat{ปุคฺคลสฺสตฺถารา สมฺมา,}{สีมฏฺโฐ อนุโมทติ.}}
\csromangathasetleft{ติํส ปาเวยฺยกา ภิกฺขู, \\ วสฺสํวุฏฺโฐกปุณฺเณหิ, \\ อิทํ วตฺถุ กถินสฺส, \\ อนามนฺตา อสมาจารา, \\ ยาวทตฺถญฺจ อุปฺปาโท, \\ ญตฺติ เอวตฺถตญฺเจว, \\ อุลฺลิขิ โธวนา เจว, \\ พนฺธโน วฏฺฏิ กณฺฑุส, \\ ปริภณฺฑํ โอวทฺเธยฺยํ, \\ กุกฺกุ สนฺนิธิ นิสฺสคฺคิ, \\ อญฺญตฺร ปญฺจาติเรเก, \\ นาญฺญตฺร ปุคฺคลา สมฺมา, \\ กถินานตฺถตํ โหติ, \\ อหตากปฺปปิโลติ, \\ อนิมิตฺตาปริกถา, \\ อนิสฺสคฺคิ กปฺปกเต, \\ ปญฺจเก วาติเรเก วา, \\ ปุคฺคลสฺสตฺถารา สมฺมา,}
\csromangathagroup{\csromangathastanza{\csromangathabat{ติํส ปาเวยฺยกา ภิกฺขู,}{สาเกตุกฺกณฺฐิตา วสุํ.} \\ \csromangathabat{วสฺสํวุฏฺโฐกปุณฺเณหิ,}{อคมุํ ชินทสฺสนํ.}}\csromangathastanza{\csromangathabat{อิทํ วตฺถุ กถินสฺส,}{กปฺปิสฺสนฺติ จ ปญฺจกา.} \\ \csromangathabat{อนามนฺตา อสมาจารา,}{ตเถว คณโภชนํ.}}\csromangathastanza{\csromangathabat{ยาวทตฺถญฺจ อุปฺปาโท,}{อตฺถตานํ ภวิสฺสติ.} \\ \csromangathabat{ญตฺติ เอวตฺถตญฺเจว,}{เอวญฺเจว อนตฺถตํ.}}\csromangathastanza{\csromangathabat{อุลฺลิขิ โธวนา เจว,}{วิจารณญฺจ เฉทนํ.} \\ \csromangathabat{พนฺธโน วฏฺฏิ กณฺฑุส,}{ทฬฺหีกมฺมานุวาติกา.}}\csromangathastanza{\csromangathabat{ปริภณฺฑํ โอวทฺเธยฺยํ,}{มทฺทนา นิมิตฺตํ กถา.} \\ \csromangathabat{กุกฺกุ สนฺนิธิ นิสฺสคฺคิ,}{น กปฺปญฺญตฺร เต ตโย.}}\csromangathastanza{\csromangathabat{อญฺญตฺร ปญฺจาติเรเก,}{สญฺฉินฺเนน สมณฺฑลี.} \\ \csromangathabat{นาญฺญตฺร ปุคฺคลา สมฺมา,}{นิสฺสีมฏฺโฐนุโมทติ.}}\csromangathastanza{\csromangathabat{กถินานตฺถตํ โหติ,}{เอวํ พุทฺเธน เทสิตํ.} \\ \csromangathabat{อหตากปฺปปิโลติ,}{ปํสุ ปาปณิกาย จ.}}\csromangathastanza{\csromangathabat{อนิมิตฺตาปริกถา,}{อกุกฺกุ จ อสนฺนิธิ.} \\ \csromangathabat{อนิสฺสคฺคิ กปฺปกเต,}{ตถา ติจีวเรน จ.}}\csromangathastanza{\csromangathabat{ปญฺจเก วาติเรเก วา,}{ฉินฺเน สมณฺฑลีกเต.} \\ \csromangathabat{ปุคฺคลสฺสตฺถารา สมฺมา,}{สีมฏฺโฐ อนุโมทติ.}}}

\csromanheader{2}{เอวํ กถิน\-ตฺถรณํ, อุพฺภารสฺส\-ฏฺฐ\-มาติกา.}
\setlength{\csromangathapagewidth}{0pt}
\setlength{\csromangathawakwidth}{0pt}
\csromangathameasuregroup{ปกฺกมนนฺติ นิฏฺฐานํ, \\ สวนํ อาสาวจฺเฉทิ,}{\csromangathabat{ปกฺกมนนฺติ นิฏฺฐานํ,}{สนฺนิฏฺฐานญฺจ นาสนํ.} \\ \csromangathabat{สวนํ อาสาวจฺเฉทิ,}{สีมา สหุพฺภารฏฺฐมี.}}
\csromangathasetleft{ปกฺกมนนฺติ นิฏฺฐานํ, \\ สวนํ อาสาวจฺเฉทิ,}
\csromangathagroup{\csromangathastanza{\csromangathabat{ปกฺกมนนฺติ นิฏฺฐานํ,}{สนฺนิฏฺฐานญฺจ นาสนํ.} \\ \csromangathabat{สวนํ อาสาวจฺเฉทิ,}{สีมา สหุพฺภารฏฺฐมี.}}}

\vspace{\tipitakaparskip}
\csromancenter{กตจีวรมาทาย, “น ปจฺเจสฺสนฺ”ติ คจฺฉติ.}
\setlength{\csromangathapagewidth}{0pt}
\setlength{\csromangathawakwidth}{0pt}
\csromangathameasuregroup{ตสฺส ตํ กถินุทฺธาโร, \\ อาทาย จีวรํ ยาติ,}{\csromangathabat{ตสฺส ตํ กถินุทฺธาโร,}{โหติ ปกฺกมนนฺติโก.} \\ \csromangathabat{อาทาย จีวรํ ยาติ,}{นิสฺสีเม อิทํ จินฺตยิ.}}
\csromangathasetleft{ตสฺส ตํ กถินุทฺธาโร, \\ อาทาย จีวรํ ยาติ,}
\csromangathagroup{\csromangathastanza{\csromanfolio{375}\csromangathabat{ตสฺส ตํ กถินุทฺธาโร,}{โหติ ปกฺกมนนฺติโก.} \\ \csromangathabat{อาทาย จีวรํ ยาติ,}{นิสฺสีเม อิทํ จินฺตยิ.}}}

\vspace{\tipitakaparskip}
\csromancenter{กาเรสฺสํ “น ปจฺเจสฺสนฺ”ติ, นิฏฺฐาเน กถินุทฺธาโร.}
\csromancenter{อาทาย นิสฺสีมํ เนว, “น ปจฺเจสฺสนฺ”ติ มานโส.}
\setlength{\csromangathapagewidth}{0pt}
\setlength{\csromangathawakwidth}{0pt}
\csromangathameasuregroup{ตสฺส ตํ กถินุทฺธาโร, \\ อาทาย จีวรํ ยาติ,}{\csromangathabat{ตสฺส ตํ กถินุทฺธาโร,}{สนฺนิฏฺฐานนฺติโก ภเว.} \\ \csromangathabat{อาทาย จีวรํ ยาติ,}{นิสฺสีเม อิทํ จินฺตยิ.}}
\csromangathasetleft{ตสฺส ตํ กถินุทฺธาโร, \\ อาทาย จีวรํ ยาติ,}
\csromangathagroup{\csromangathastanza{\csromangathabat{ตสฺส ตํ กถินุทฺธาโร,}{สนฺนิฏฺฐานนฺติโก ภเว.} \\ \csromangathabat{อาทาย จีวรํ ยาติ,}{นิสฺสีเม อิทํ จินฺตยิ.}}}

\vspace{\tipitakaparskip}
\csromancenter{กาเรสฺสํ “น ปจฺเจสฺสนฺ”ติ, กยิรํ ตสฺส นสฺสติ.}
\setlength{\csromangathapagewidth}{0pt}
\setlength{\csromangathawakwidth}{0pt}
\csromangathameasuregroup{ตสฺส ตํ กถินุทฺธาโร, \\ อาทาย ยาติ ปจฺเจสฺสํ, \\ กตจีวโร สุณาติ, \\ ตสฺส ตํ กถินุทฺธาโร, \\ อาทาย ยาติ ปจฺเจสฺสํ, \\ กตจีวโร พหิทฺธา, \\ ตสฺส ตํ กถินุทฺธาโร, \\ อาทาย ยาติ ปจฺเจสฺสํ, \\ กตจีวโร ปจฺเจสฺสํ, \\ ตสฺส ตํ กถินุทฺธาโร, \\ อาทาย จ สมาทาย, \\ ปกฺกมนนฺติกา นตฺถิ, \\ อาทาย นิสฺสีมคตํ, \\ นิฏฺฐานํ สนฺนิฏฺฐานญฺจ,}{\csromangathabat{ตสฺส ตํ กถินุทฺธาโร,}{ภวติ นาสนนฺติโก.} \\ \csromangathabat{อาทาย ยาติ ปจฺเจสฺสํ,}{พหิ กาเรติ จีวรํ.} \\ \csromangathabat{กตจีวโร สุณาติ,}{อุพฺภตํ กถินํ ตหิํ.} \\ \csromangathabat{ตสฺส ตํ กถินุทฺธาโร,}{ภวติ สวนนฺติโก.} \\ \csromangathabat{อาทาย ยาติ ปจฺเจสฺสํ,}{พหิ กาเรติ จีวรํ.} \\ \csromangathabat{กตจีวโร พหิทฺธา,}{นาเมติ กถินุทฺธารํ.} \\ \csromangathabat{ตสฺส ตํ กถินุทฺธาโร,}{สีมาติกฺกนฺติโก ภเว.} \\ \csromangathabat{อาทาย ยาติ ปจฺเจสฺสํ,}{พหิ กาเรติ จีวรํ.} \\ \csromangathabat{กตจีวโร ปจฺเจสฺสํ,}{สมฺโภติ กถินุทฺธารํ.} \\ \csromangathabat{ตสฺส ตํ กถินุทฺธาโร,}{สห ภิกฺขูหิ ชายติ.} \\ \csromangathabat{อาทาย จ สมาทาย,}{สตฺตสตฺตวิธา คติ.} \\ \csromangathabat{ปกฺกมนนฺติกา นตฺถิ,}{ฉกฺเก วิปฺปกเต\textsuperscript{1} คติ.} \\ \csromangathabat{อาทาย นิสฺสีมคตํ,}{กาเรสฺสํ อิติ ชายติ.} \\ \csromangathabat{นิฏฺฐานํ สนฺนิฏฺฐานญฺจ,}{นาสนญฺจ อิเม ตโย.}}
\csromangathasetleft{ตสฺส ตํ กถินุทฺธาโร, \\ อาทาย ยาติ ปจฺเจสฺสํ, \\ กตจีวโร สุณาติ, \\ ตสฺส ตํ กถินุทฺธาโร, \\ อาทาย ยาติ ปจฺเจสฺสํ, \\ กตจีวโร พหิทฺธา, \\ ตสฺส ตํ กถินุทฺธาโร, \\ อาทาย ยาติ ปจฺเจสฺสํ, \\ กตจีวโร ปจฺเจสฺสํ, \\ ตสฺส ตํ กถินุทฺธาโร, \\ อาทาย จ สมาทาย, \\ ปกฺกมนนฺติกา นตฺถิ, \\ อาทาย นิสฺสีมคตํ, \\ นิฏฺฐานํ สนฺนิฏฺฐานญฺจ,}
\csromangathagroup{\csromangathastanza{\csromangathabat{ตสฺส ตํ กถินุทฺธาโร,}{ภวติ นาสนนฺติโก.} \\ \csromangathabat{อาทาย ยาติ ปจฺเจสฺสํ,}{พหิ กาเรติ จีวรํ.}}\csromangathastanza{\csromangathabat{กตจีวโร สุณาติ,}{อุพฺภตํ กถินํ ตหิํ.} \\ \csromangathabat{ตสฺส ตํ กถินุทฺธาโร,}{ภวติ สวนนฺติโก.}}\csromangathastanza{\csromangathabat{อาทาย ยาติ ปจฺเจสฺสํ,}{พหิ กาเรติ จีวรํ.} \\ \csromangathabat{กตจีวโร พหิทฺธา,}{นาเมติ กถินุทฺธารํ.}}\csromangathastanza{\csromangathabat{ตสฺส ตํ กถินุทฺธาโร,}{สีมาติกฺกนฺติโก ภเว.} \\ \csromangathabat{อาทาย ยาติ ปจฺเจสฺสํ,}{พหิ กาเรติ จีวรํ.}}\csromangathastanza{\csromangathabat{กตจีวโร ปจฺเจสฺสํ,}{สมฺโภติ กถินุทฺธารํ.} \\ \csromangathabat{ตสฺส ตํ กถินุทฺธาโร,}{สห ภิกฺขูหิ ชายติ.}}\csromangathastanza{\csromangathabat{อาทาย จ สมาทาย,}{สตฺตสตฺตวิธา คติ.} \\ \csromangathabat{ปกฺกมนนฺติกา นตฺถิ,}{ฉกฺเก วิปฺปกเต\csromansharedfootnote{74ad1ef70bfd81b2}{ฉฏฺเฐ วิปฺปกตา (สี), ฉจฺจา วิปฺปกตา (ก)} คติ.}}\csromangathastanza{\csromangathabat{อาทาย นิสฺสีมคตํ,}{กาเรสฺสํ อิติ ชายติ.} \\ \csromangathabat{นิฏฺฐานํ สนฺนิฏฺฐานญฺจ,}{นาสนญฺจ อิเม ตโย.}}}

\prose{อาทาย “น ปจฺเจสฺสนฺ”ติ, พหิสีเม กโรมิติ.}

\setlength{\csromangathapagewidth}{0pt}
\setlength{\csromangathawakwidth}{0pt}
\csromangathameasuregroup{นิฏฺฐานํ สนฺนิฏฺฐานมฺปิ, \\ อนธิฏฺฐิเตน เนวสฺส, \\ อาทาย ยาติ ปจฺเจสฺสํ,}{\csromangathabat{นิฏฺฐานํ สนฺนิฏฺฐานมฺปิ,}{นาสนมฺปิ อิทํ ตโย.} \\ \csromangathabat{อนธิฏฺฐิเตน เนวสฺส,}{เหฏฺฐา ตีณิ นยาวิธิ.} \\ \csromangathabat{อาทาย ยาติ ปจฺเจสฺสํ,}{พหิสีเม กโรมิติ.}}
\csromangathasetleft{นิฏฺฐานํ สนฺนิฏฺฐานมฺปิ, \\ อนธิฏฺฐิเตน เนวสฺส, \\ อาทาย ยาติ ปจฺเจสฺสํ,}
\csromangathagroup{\csromangathastanza{\csromangathabat{นิฏฺฐานํ สนฺนิฏฺฐานมฺปิ,}{นาสนมฺปิ อิทํ ตโย.} \\ \csromangathabat{อนธิฏฺฐิเตน เนวสฺส,}{เหฏฺฐา ตีณิ นยาวิธิ.} \\ \csromangathabat{อาทาย ยาติ ปจฺเจสฺสํ,}{พหิสีเม กโรมิติ.}}}

\vspace{\tipitakaparskip}
\csromancenter{“น ปจฺเจสฺสนฺ”ติ กาเรติ, นิฏฺฐาเน กถินุทฺธาโร.}
\setlength{\csromangathapagewidth}{0pt}
\setlength{\csromangathawakwidth}{0pt}
\csromangathameasuregroup{สนฺนิฏฺฐานํ นาสนญฺจ, \\ สห ภิกฺขูหิ ชาเยถ, \\ สมาทาย วิปฺปกตา, \\ อิเม เต จตุโร วารา, \\ อนาสาย จ อาสาย, \\ นยโต ตํ วิชาเนยฺย, \\ อปวิลานา นเวตฺถ\textsuperscript{1}, \\ ปลิโพธาปลิโพธา,}{\csromangathabat{สนฺนิฏฺฐานํ นาสนญฺจ,}{สวนสีมาติกฺกมา.} \\ \csromangathabat{สห ภิกฺขูหิ ชาเยถ,}{เอวํ ปนฺนรสํ คติ.} \\ \csromangathabat{สมาทาย วิปฺปกตา,}{สมาทาย ปุนา ตถา.} \\ \csromangathabat{อิเม เต จตุโร วารา,}{สพฺเพ ปนฺนรสวิธิ.} \\ \csromangathabat{อนาสาย จ อาสาย,}{กรณีโย จ เต ตโย.} \\ \csromangathabat{นยโต ตํ วิชาเนยฺย,}{ตโย ทฺวาทส ทฺวาทส.} \\ \csromangathabat{อปวิลานา นเวตฺถ\textsuperscript{1},}{ผาสุ ปญฺจวิธา ตหิํ.} \\ \csromangathabat{ปลิโพธาปลิโพธา,}{อุทฺทานํ นยโต กตนฺติ.}}
\csromangathasetleft{สนฺนิฏฺฐานํ นาสนญฺจ, \\ สห ภิกฺขูหิ ชาเยถ, \\ สมาทาย วิปฺปกตา, \\ อิเม เต จตุโร วารา, \\ อนาสาย จ อาสาย, \\ นยโต ตํ วิชาเนยฺย, \\ อปวิลานา นเวตฺถ\textsuperscript{1}, \\ ปลิโพธาปลิโพธา,}
\csromangathagroup{\csromangathastanza{\csromanfolio{376}\csromangathabat{สนฺนิฏฺฐานํ นาสนญฺจ,}{สวนสีมาติกฺกมา.} \\ \csromangathabat{สห ภิกฺขูหิ ชาเยถ,}{เอวํ ปนฺนรสํ คติ.}}\csromangathastanza{\csromangathabat{สมาทาย วิปฺปกตา,}{สมาทาย ปุนา ตถา.} \\ \csromangathabat{อิเม เต จตุโร วารา,}{สพฺเพ ปนฺนรสวิธิ.}}\csromangathastanza{\csromangathabat{อนาสาย จ อาสาย,}{กรณีโย จ เต ตโย.} \\ \csromangathabat{นยโต ตํ วิชาเนยฺย,}{ตโย ทฺวาทส ทฺวาทส.}}\csromangathastanza{\csromangathabat{อปวิลานา นเวตฺถ\csromansharedfootnote{9161532896772f64}{อปวิลายมาเนว (สฺยา), อปวินา นว เจตฺถ (สี)},}{ผาสุ ปญฺจวิธา ตหิํ.} \\ \csromangathabat{ปลิโพธาปลิโพธา,}{อุทฺทานํ นยโต กตนฺติ.}}}

\prosewithclosermajor{อิมมฺหิ ขนฺธเก วตฺถู โทฬสก\-เปยฺยาลมุขานิ เอกสตํ อฏฺฐารส.}{\textbf{กถิน}\-\textbf{กฺขนฺธโก นิฏฺฐิโต.}}
\csromanmatikaanchor{mtk.209}
\csromantocmarknopage{chapter}{8. จีวรกฺขนฺธก}{mtk.209}
\bookmark[dest={mtk.209},level=0]{8. จีวรกฺขนฺธก}
\chapterheadpageread{8. \textbf{จีวร}\-\textbf{กฺขนฺธก}}
\csromanmatikaanchor{mtk.210}
\csromantocmarkref{section}{202. ชีวกวตฺถุ}{mtk.210}
\bookmark[dest={mtk.210},level=1]{202. ชีวกวตฺถุ}
\csromanheader{1}{202. \textbf{ชีวกวตฺถุ}}
\proseitem{326}{\csromanfolio{377}เตน สมเยน พุทฺโธ ภควา ราชคเห วิหรติ เวฬุวเน กลนฺทก\-นิวาเป. เตน โข ปน สมเยน เวสาลี อิทฺธา เจว โหติ ผิตา\csromansharedfootnote{b6821f306cd9b779}{ผีตา (พหูสุ)} จ พหุชนา จ อากิณฺณมนุสฺสา จ สุภิกฺขา จ, สตฺต จ ปาสาท\-สหสฺสานิ สตฺต จ ปาสาทสตานิ สตฺต จ ปาสาทา, สตฺต จ กูฏาคาร\-สหสฺสานิ สตฺต จ กูฏาคารสตานิ สตฺต จ กูฏาคารานิ, สตฺต จ อารามสหสฺสานิ สตฺต จ อารามสตานิ สตฺต จ อารามา, สตฺต จ โปกฺขรณี\-สหสฺสานิ สตฺต จ โปกฺขรณี\-สตานิ สตฺต จ โปกฺขรณิโย, อมฺพปาลี จ คณิกา อภิรูปา โหติ ทสฺสนียา ปาสาทิกา ปรมาย วณฺณ\-โปกฺขร\-ตาย สมนฺนาคตา, ปทกฺขิณา\csromansharedfootnote{11f1080d9c3c639a}{ปทกฺขา (สฺยา)} นจฺเจ จ คีเต จ วาทิเต จ, อภิสฏา อตฺถิกานํ อตฺถิกานํ มนุสฺสานํ ปญฺญาสาย จ รตฺติํ คจฺฉติ, ตาย จ เวสาลี ภิยฺโยโส มตฺตาย อุปโสภติ. อถ โข ราชคหโก เนคโม เวสาลิํ อคมาสิ เกนจิเทว กรณีเยน. อทฺทสา โข ราชคหโก เนคโม เวสาลิํ อิทฺธญฺเจว ผิตญฺจ พหุชนญฺจ อากิณฺณมนุสฺสญฺจ สุภิกฺขญฺจ, สตฺต จ ปาสาท\-สหสฺสานิ สตฺต จ ปาสาทสตานิ สตฺต จ ปาสาเท, สตฺต จ กูฏาคาร\-สหสฺสานิ สตฺต จ กูฏาคารสตานิ สตฺต จ กูฏาคารานิ, สตฺต จ อารามสหสฺสานิ สตฺต จ อารามสตานิ สตฺต จ อาราเม, สตฺต จ โปกฺขรณี\-สหสฺสานิ สตฺต จ โปกฺขรณี\-สตานิ สตฺต จ โปกฺขรณิโย, อมฺพปาลิญฺจ คณิกํ อภิรูปํ ทสฺสนียํ ปาสาทิกํ ปรมาย วณฺณ\-โปกฺขร\-ตาย สมนฺนาคตํ, ปทกฺขิณํ\csromansharedfootnote{19772863603a87a4}{ปทกฺขํ (สฺยา)} นจฺเจ จ คีเต จ วาทิเต จ, อภิสฏํ อตฺถิกานํ อตฺถิกานํ มนุสฺสานํ ปญฺญาสาย จ รตฺติํ คจฺฉนฺติํ, ตาย จ เวสาลิํ ภิยฺโยโส มตฺตาย อุปโสภนฺติํ.}

\proseitem{327}{อถ โข ราชคหโก เนคโม เวสาลิยํ ตํ กรณียํ ตีเรตฺวา ปุนเทว ราชคหํ ปจฺจาคญฺฉิ, เยน ราชา มาคโธ เสนิโย พิมฺพิสาโร เตนุปสงฺกมิ, อุปสงฺกมิตฺวา ราชานํ มาคธํ เสนิยํ \csromanfolio{378}พิมฺพิสารํ เอตทโวจ “เวสาลี เทว อิทฺธา เจว ผิตา จ พหุชนา จ อากิณฺณมนุสฺสา จ สุภิกฺขา จ, สตฺต จ ปาสาท\-สหสฺสานิ ฯเปฯ ตาย จ เวสาลี ภิยฺโยโส มตฺตาย อุปโสภติ, สาธุ เทว มยมฺปิ คณิกํ วุฏฺฐาเปสฺสามา”ติ\csromansharedfootnote{db7a95b30e2a1348}{วุฏฺฐาเปยฺยาม (ก)}. เตน หิ ภเณ ตาทิสิํ กุมาริํ ชานาถ, ยํ ตุเมฺห คณิกํ วุฏฺฐาเปยฺยาถาติ. เตน โข ปน สมเยน ราชคเห สาลวตี นาม กูมารี อภิรูปา โหติ ทสฺสนียา ปาสาทิกา ปรมาย วณฺณ\-โปกฺขร\-ตาย สมนฺนาคตา. อถ โข ราชคหโก เนคโม สาลวติํ กุมาริํ คณิกํ วุฏฺฐาเปสิ. อถ โข สาลวตี คณิกา นจิรสฺเสว ปทกฺขิณา อโหสิ นจฺเจ จ คีเต จ วาทิเต จ, อภิสฏา อตฺถิกานํ อตฺถิกานํ มนุสฺสานํ ปฏิสเตน จ รตฺติํ คจฺฉติ. อถ โข สาลวตี คณิกา นจิรสฺเสว คพฺภินี อโหสิ. อถ โข สาลวติยา คณิกาย เอตทโหสิ “อิตฺถี โข คพฺภินี ปุริสานํ อมนาปา, สเจ มํ โกจิ ชานิสฺสติ ‘สาลวตี คณิกา คพฺภินี’ติ, สพฺโพ เม สกฺกาโร ภญฺชิสฺสติ\csromansharedfootnote{a768b434f324e6de}{ปริหายิสฺสติ (สี, สฺยา)}, ยนฺนูนาหํ คิลานํ ปฏิเวเทยฺยนฺ”ติ. อถ โข สาลวตี คณิกา โทวาริกํ อาณาเปสิ “มา ภเณ โทวาริก โกจิ ปุริโส ปาวิสิ, โย จ มํ ปุจฺฉติ, ‘คิลานา’ติ ปฏิเวเทหี”ติ. “เอวํ อยฺเย”ติ โข โส โทวาริโก สาลวติยา คณิกาย ปจฺจสฺโสสิ. อถ โข สาลวตี คณิกา ตสฺส คพฺภสฺส ปริปากม\-นฺวาย ปุตฺตํ วิชายิ. อถ โข สาลวตี คณิกา ทาสิํ อาณาเปสิ “หนฺท เช อิมํ ทารกํ กตฺตรสุปฺเป ปกฺขิปิตฺวา นีหริตฺวา สงฺการกูเฏ ฉฑฺเฑหี”ติ. “เอวํ อยฺเย”ติ โข สา ทาสี สาลวติยา คณิกาย ปฏิสฺสุตฺวา ตํ ทารกํ กตฺตรสุปฺเป ปกฺขิปิตฺวา นีหริตฺวา สงฺการกูเฏ ฉฑฺเฑสิ.}

\proseitem{328}{เตน โข ปน สมเยน อภโย นาม ราชกุมาโร กาลสฺเสว ราชุปฏฺฐานํ คจฺฉนฺโต อทฺทส ตํ ทารกํ กาเกหิ สมฺปริกิณฺณํ, ทิสฺวาน มนุสฺเส ปุจฺฉิ “กิํ เอตํ ภเณ กาเกหิ สมฺปริกิณฺณนฺ”ติ. ทารโก เทวาติ. ชีวติ ภเณติ. ชีวติ เทวาติ. เตน หิ ภเณ ตํ ทารกํ อมฺหากํ อนฺเตปุรํ เนตฺวา ธาตีนํ เทถ โปเสตุนฺติ. “เอวํ เทวา”ติ โข เต มนุสฺสา อภยสฺส ราชกุมารสฺส ปฏิสฺสุตฺวา \csromanfolio{379}ตํ ทารกํ อภยสฺส ราชกุมารสฺส อนฺเตปุรํ เนตฺวา ธาตีนํ อทํสุ “โปเสถา”ติ. ตสฺส ชีวตีติ “ชีวโก”ติ นามํ อกํสุ, กุมาเรน โปสาปิโตติ “โกมารภจฺโจ”ติ นามํ อกํสุ. อถ โข ชีวโก โกมารภจฺโจ นจิรสฺเสว วิญฺญุตํ ปาปุณิ. อถ โข ชีวโก โกมารภจฺโจ เยน อภโย ราชกุมาโร เตนุปสงฺกมิ, อุปสงฺกมิตฺวา อภยํ ราชกุมารํ เอตทโวจ “กา เม เทว มาตา, โก ปิตา”ติ. อหมฺปิ โข เต ภเณ ชีวก มาตรํ น ชานามิ, อปิ จาหํ เต ปิตา, มยาสิ\csromansharedfootnote{68a935ff99655484}{มยาปิ (ก)} โปสาปิโตติ. อถ โข ชีวกสฺส โกมาร\-ภจฺจสฺส เอตทโหสิ “อิมานิ โข ราชกุลานิ น สุกรานิ อสิปฺเปน อุปชีวิตุํ, ยนฺนูนาหํ สิปฺปํ สิกฺเขยฺยนฺ”ติ.}

\proseitem{329}{เตน โข ปน สมเยน ตกฺกสิลายํ\csromansharedfootnote{6f097363d9673989}{ตกฺกสีลายํ (ก)} ทิสาปาโมกฺโข เวชฺโช ปฏิวสติ. อถ โข ชีวโก โกมารภจฺโจ อภยํ ราชกุมารํ อนาปุจฺฉา เยน ตกฺกสิลา เตน ปกฺกามิ. อนุปุพฺเพน เยน ตกฺกสิลา, เยน เวชฺโช เตนุปสงฺกมิ, อุปสงฺกมิตฺวา ตํ เวชฺชํ เอตทโวจ “อิจฺฉามหํ อาจริย สิปฺปํ สิกฺขิตุนฺ”ติ. เตน หิ ภเณ ชีวก สิกฺขสฺสูติ. อถ โข ชีวโก โกมารภจฺโจ พหุญฺจ คณฺหาติ, ลหุญฺจ คณฺหาติ, สุฏฺฐุ จ อุปธาเรติ, คหิตญฺจสฺส น สมฺมุสฺสติ\csromansharedfootnote{662fe239ad2b949e}{น ปมุสฺสติ (สี, สฺยา)}. อถ โข ชีวกสฺส โกมาร\-ภจฺจสฺส สตฺตนฺนํ วสฺสานํ อจฺจเยน เอตทโหสิ “อหํ โข พหุญฺจ คณฺหามิ, ลหุญฺจ คณฺหามิ, สุฏฺฐุ จ อุปธาเรมิ, คหิตญฺจ เม น สมฺมุสฺสติ, สตฺต จ เม วสฺสานิ อธียนฺตสฺส นยิมสฺส สิปฺปสฺส อนฺโต ปญฺญายติ, กทา อิมสฺส สิปฺปสฺส อนฺโต ปญฺญายิสฺสตี”ติ. อถ โข ชีวโก โกมารภจฺโจ เยน โส เวชฺโช เตนุปสงฺกมิ, อุปสงฺกมิตฺวา ตํ เวชฺชํ เอตทโวจ “อหํ โข อาจริย พหุญฺจ คณฺหามิ, ลหุญฺจ คณฺหามิ, สุฏฺฐุ จ อุปธาเรมิ, คหิตญฺจ เม น สมฺมุสฺสติ, สตฺต จ เม วสฺสานิ อธียนฺตสฺส นยิมสฺส สิปฺปสฺส อนฺโต ปญฺญายติ, กทา อิมสฺส สิปฺปสฺส อนฺโต ปญฺญายิสฺสตี”ติ. เตน หิ ภเณ ชีวก ขณิตฺติํ อาทาย ตกฺกสิลาย สมนฺตา โยชนํ อาหิณฺฑิตฺวา ยํ กิญฺจิ อเภสชฺชํ ปสฺเสยฺยาสิ, ตํ อาหราติ. “เอวํ อาจริยา”ติ โข ชีวโก โกมารภจฺโจ ตสฺส เวชฺชสฺส ปฏิสฺสุตฺวา ขณิตฺติํ อาทาย \csromanfolio{380}ตกฺกสิลาย สมนฺตา โยชนํ อาหิณฺฑนฺโต น กิญฺจิ อเภสชฺชํ อทฺทส. อถ โข ชีวโก โกมารภจฺโจ เยน โส เวชฺโช เตนุปสงฺกมิ, อุปสงฺกมิตฺวา ตํ เวชฺชํ เอตทโวจ “อาหิณฺฑนฺโตมฺหิ อาจริย ตกฺกสิลาย สมนฺตา โยชนํ, น กิญฺจิ\csromansharedfootnote{c23e0fe914ab30f5}{อาหิณฺฑนฺโต น กิญฺจิ (ก)} อเภสชฺชํ อทฺทสนฺ”ติ. “สุสิกฺขิโตสิ ภเณ ชีวก, อลํ เต เอตฺตกํ ชีวิกายา”ติ ชีวกสฺส โกมาร\-ภจฺจสฺส ปริตฺตํ ปาเถยฺยํ ปาทาสิ. อถ โข ชีวโก โกมารภจฺโจ ตํ ปริตฺตํ ปาเถยฺยํ อาทาย เยน ราชคหํ เตน ปกฺกามิ. อถ โข ชีวกสฺส โกมาร\-ภจฺจสฺส ตํ ปริตฺตํ ปาเถยฺยํ อนฺตรามคฺเค สาเกเต ปริกฺขยํ อคมาสิ. อถ โข ชีวกสฺส โกมาร\-ภจฺจสฺส เอตทโหสิ “อิเม โข มคฺคา กนฺตารา อปฺโปทกา อปฺปภกฺขา น สุกรา อปาเถยฺเยน คนฺตุํ, ยนฺนูนาหํ ปาเถยฺยํ ปริเยเสยฺยนฺ”ติ.}

\csromanmatikaanchor{mtk.211}
\csromantocmarkref{section}{203. เสฏฺฐิภริยาวตฺถุ}{mtk.211}
\bookmark[dest={mtk.211},level=1]{203. เสฏฺฐิภริยาวตฺถุ}
\csromanheader{1}{203. \textbf{เสฏฺฐิภริยา}\-\textbf{วตฺถุ}}
\proseitem{330}{เตน โข ปน สมเยน สาเกเต เสฏฺฐิ\-ภริยาย สตฺตวสฺสิโก สีสาพาโธ โหติ, พหู มหนฺตา มหนฺตา ทิสาปาโมกฺขา เวชฺชา อาคนฺตฺวา นาสกฺขิํสุ อโรคํ กาตุํ, พหุํ หิรญฺญํ อาทาย อคมํสุ. อถ โข ชีวโก โกมารภจฺโจ สาเกตํ ปวิสิตฺวา มนุสฺเส ปุจฺฉิ “โก ภเณ คิลาโน, กํ ติกิจฺฉามี”ติ. เอติสฺสา อาจริย เสฏฺฐิ\-ภริยาย สตฺตวสฺสิโก สีสาพาโธ, คจฺฉ อาจริย เสฏฺฐิภริยํ ติกิจฺฉาหีติ. อถ โข ชีวโก โกมารภจฺโจ เยน เสฏฺฐิสฺส คหปติสฺส นิเวสนํ เตนุปสงฺกมิ, อุปสงฺกมิตฺวา โทวาริกํ อาณาเปสิ คจฺฉ ภเณ โทวาริก เสฏฺฐิ\-ภริยาย ปาวท “เวชฺโช อยฺเย อาคโต, โส ตํ ทฏฺฐุกาโม”ติ. “เอวํ อาจริยา”ติ โข โส โทวาริโก ชีวกสฺส โกมาร\-ภจฺจสฺส ปฏิสฺสุตฺวา เยน เสฏฺฐิภริยา เตนุปสงฺกมิ, อุปสงฺกมิตฺวา เสฏฺฐิภริยํ เอตทโวจ “เวชฺโช อยฺเย อาคโต, โส ตํ ทฏฺฐุกาโม”ติ. กีทิโส ภเณ โทวาริก เวชฺโชติ. ทหรโก อยฺเยติ. อลํ ภเณ โทวาริก กิํ เม ทหรโก เวชฺโช กริสฺสติ, พหู มหนฺตา มหนฺตา ทิสาปาโมกฺขา เวชฺชา \csromanfolio{381}อาคนฺตฺวา นาสกฺขิํสุ อโรคํ กาตุํ, พหุํ หิรญฺญํ อาทาย อคมํสูติ. อถ โข โส โทวาริโก เยน ชีวโก โกมารภจฺโจ เตนุปสงฺกมิ, อุปสงฺกมิตฺวา ชีวกํ โกมารภจฺจํ เอตทโวจ “เสฏฺฐิภริยา อาจริย เอวมาห อลํ ภเณ โทวาริก กิํ เม ทหรโก เวชฺโช กริสฺสติ, พหู มหนฺตา มหนฺตา ทิสาปาโมกฺขา เวชฺชา อาคนฺตฺวา นาสกฺขิํสุ อโรคํ กาตุํ, พหุํ หิรญฺญํ อาทาย อคมํสู”ติ. คจฺฉ ภเณ โทวาริก เสฏฺฐิ\-ภริยาย ปาวท “เวชฺโช อยฺเย เอวมาห มา กิร อยฺเย ปุเร กิญฺจิ อทาสิ, ยทา อโรคา อโหสิ, ตทา ยํ อิจฺเฉยฺยาสิ, ตํ ทชฺเชยฺยาสี”ติ. “เอวํ อาจริยา”ติ โข โส โทวาริโก ชีวกสฺส โกมาร\-ภจฺจสฺส ปฏิสฺสุตฺวา เยน เสฏฺฐิภริยา เตนุปสงฺกมิ, อุปสงฺกมิตฺวา เสฏฺฐิภริยํ เอตทโวจ “เวชฺโช อยฺเย เอวมาห มา กิร อยฺเย ปุเร กิญฺจิ อทาสิ, ยทา อโรคา อโหสิ, ตทา ยํ อิจฺเฉยฺยาสิ, ตํ ทชฺเชยฺยาสี”ติ. เตน หิ ภเณ โทวาริก เวชฺโช อาคจฺฉตูติ. “เอวํ อยฺเย”ติ โข โส โทวาริโก เสฏฺฐิ\-ภริยาย ปฏิสฺสุตฺวา เยน ชีวโก โกมารภจฺโจ เตนุปสงฺกมิ, อุปสงฺกมิตฺวา ชีวกํ โกมารภจฺจํ เอตทโวจ “เสฏฺฐิภริยา ตํ อาจริย ปกฺโกสตี”ติ.}

\prose{อถ โข ชีวโก โกมารภจฺโจ เยน เสฏฺฐิภริยา เตนุปสงฺกมิ, อุปสงฺกมิตฺวา เสฏฺฐิ\-ภริยาย วิการํ สลฺลกฺเขตฺวา เสฏฺฐิภริยํ เอตทโวจ “ปสเตน อยฺเย สปฺปินา อตฺโถ”ติ. อถ โข เสฏฺฐิภริยา ชีวกสฺส โกมาร\-ภจฺจสฺส ปสตํ สปฺปิํ ทาเปสิ. อถ โข ชีวโก โกมารภจฺโจ ตํ ปสตํ สปฺปิํ นานาเภสชฺเชหิ นิปฺปจิตฺวา เสฏฺฐิภริยํ มญฺจเก อุตฺตานํ นิปาเตตฺวา\csromansharedfootnote{7003fc0f21583c19}{นิปชฺชาเปตฺวา (สี, สฺยา)} นตฺถุโต อทาสิ. อถ โข ตํ สปฺปิ นตฺถุโต ทินฺนํ มุขโต อุคฺคญฺฉิ. อถ โข เสฏฺฐิภริยา ปฏิคฺคเห นิฏฺฐุภิตฺวา ทาสิํ อาณาเปสิ “หนฺท เช อิมํ สปฺปิํ ปิจุนา คณฺหาหี”ติ. อถ โข ชีวกสฺส โกมาร\-ภจฺจสฺส เอตทโหสิ “อจฺฉริยํ\csromansharedfootnote{ac9fc063912310fd}{อจฺฉริยํ วต โภ (สฺยา)} ยาว ลูขายํ ฆรณี, ยตฺร หิ นาม อิมํ ฉฑฺฑนีย\-ธมฺมํ สปฺปิํ ปิจุนา คาหาเปสฺสติ, พหุกานิ จ เม มหคฺฆานิ\csromansharedfootnote{5befb1a38a316269}{มหคฺฆานิ มหคฺฆานิ (สี, สฺยา)} เภสชฺชานิ อุปคตานิ, กิมฺปิ มา’ยํ กิญฺจิ\csromansharedfootnote{2fc077889885ec43}{กญฺจิ (สฺยา)} เทยฺยธมฺมํ ทสฺสตี”ติ. อถ โข เสฏฺฐิภริยา \csromanfolio{382}ชีวกสฺส โกมาร\-ภจฺจสฺส วิการํ สลฺลกฺเขตฺวา ชีวกํ โกมารภจฺจํ เอตทโวจ “กิสฺส ตฺวํ อาจริย วิมโนสี”ติ. อิธ เม เอตทโหสิ “อจฺฉริยํ ยาว ลูขายํ ฆรณี, ยตฺร หิ นาม อิมํ ฉฑฺฑนีย\-ธมฺมํ สปฺปิํ ปิจุนา คาหาเปสฺสติ, พหุกานิ จ เม มหคฺฆานิ เภสชฺชานิ อุปคตานิ, กิมฺปิ มา’ยํ กิญฺจิ เทยฺยธมฺมํ ทสฺสตี”ติ. มยํ โข อาจริย อาคาริกา นาม อุปชานาเมตสฺส สํยมสฺส, วรเมตํ สปฺปิ ทาสานํ วา กมฺมกรานํ วา ปาทพฺภญฺชนํ วา ปทีปกรเณ วา อาสิตฺตํ, มา โข ตฺวํ อาจริย วิมโน อโหสิ, น เต เทยฺยธมฺโม หายิสฺสตีติ. อถ โข ชีวโก โกมารภจฺโจ เสฏฺฐิ\-ภริยาย สตฺตวสฺสิกํ สีสาพาธํ เอเกเนว นตฺถุกมฺเมน อปกฑฺฒิ. อถ โข เสฏฺฐิภริยา อโรคา สมานา ชีวกสฺส โกมาร\-ภจฺจสฺส จตฺตาริ สหสฺสานิ ปาทาสิ. ปุตฺโต “มาตา เม อโรคา ฐิตา”ติ จตฺตาริ สหสฺสานิ ปาทาสิ. สุณิสา “สสฺสุ เม อโรคา ฐิตา”ติ จตฺตาริ สหสฺสานิ ปาทาสิ. เสฏฺฐิ คหปติ “ภริยา เม อโรคา ฐิตา”ติ จตฺตาริ สหสฺสานิ ปาทาสิ ทาสญฺจ ทาสิญฺจ อสฺสรถญฺจ.}

\prose{อถ โข ชีวโก โกมารภจฺโจ ตานิ โสฬส\-สหสฺสานิ อาทาย ทาสญฺจ ทาสิญฺจ อสฺสรถญฺจ เยน ราชคหํ เตน ปกฺกามิ. อนุปุพฺเพน เยน ราชคหํ, เยน อภโย ราชกุมาโร เตนุปสงฺกมิ, อุปสงฺกมิตฺวา อภยํ ราชกุมารํ เอตทโวจ “อิทํ เม เทว ปฐมกมฺมํ โสฬส\-สหสฺสานิ ทาโส จ ทาสี จ อสฺสรโถ จ, ปฏิคฺคณฺหาตุ เม เทโว โปสาวนิกนฺ”ติ. อลํ ภเณ ชีวก ตุยฺหเมว โหตุ, อมฺหากญฺเญว อนฺเตปุเร นิเวสนํ มาเปหีติ. “เอวํ เทวา”ติ โข ชีวโก โกมารภจฺโจ อภยสฺส ราชกุมารสฺส ปฏิสฺสุตฺวา อภยสฺส ราชกุมารสฺส อนฺเตปุเร นิเวสนํ มาเปสิ.}

\csromanmatikaanchor{mtk.212}
\csromantocmarkref{section}{204. พิมฺพิสารราชวตฺถุ}{mtk.212}
\bookmark[dest={mtk.212},level=1]{204. พิมฺพิสารราชวตฺถุ}
\csromanheader{1}{204. \textbf{พิมฺพิสาร}\-\textbf{ราช}\-\textbf{วตฺถุ}}
\proseitem{331}{เตน โข ปน สมเยน รญฺโญ มาคธสฺส เสนิยสฺส พิมฺพิสารสฺส ภคนฺทลา\-พาโธ โหติ, สาฏกา โลหิเตน มกฺขิยนฺติ. เทวิโย ทิสฺวา อุปฺปณฺเฑนฺติ “อุตุนี ทานิ เทโว, ปุปฺผํ เทวสฺส อุปฺปนฺนํ, น จิรํ\csromansharedfootnote{c391897f2b644151}{นจิรสฺเสว (สฺยา)} เทโว วิชายิสฺสตี”ติ. เตน ราชา มงฺกุ \csromanfolio{383}โหติ. อถ โข ราชา มาคโธ เสนิโย พิมฺพิสาโร อภยํ ราชกุมารํ เอตทโวจ “มยฺหํ โข ภเณ อภย ตาทิโส อาพาโธ, สาฏกา โลหิเตน มกฺขิยนฺติ, เทวิโย มํ ทิสฺวา อุปฺปณฺเฑนฺติ ‘อุตุนี ทานิ เทโว, ปุปฺผํ เทวสฺส อุปฺปนฺนํ, น จิรํ เทโว วิชายิสฺสตี’ติ, อิงฺฆ ภเณ อภย ตาทิสํ เวชฺชํ ชานาหิ, โย มํ ติกิจฺเฉยฺยา”ติ. อยํ เทว อมฺหากํ ชีวโก เวชฺโช ตรุโณ ภทฺรโก, โส เทวํ ติกิจฺฉิสฺสตีติ. เตน หิ ภเณ อภย ชีวกํ เวชฺชํ อาณาเปหิ, โส มํ ติกิจฺฉิสฺสตีติ. อถ โข อภโย ราชกุมาโร ชีวกํ โกมารภจฺจํ อาณาเปสิ “คจฺฉ ภเณ ชีวก ราชานํ ติกิจฺฉาหีติ. “เอวํ เทวา”ติ โข ชีวโก โกมารภจฺโจ อภยสฺส ราชกุมารสฺส ปฏิสฺสุตฺวา นเขน เภสชฺชํ อาทาย เยน ราชา มาคโธ เสนิโย พิมฺพิสาโร เตนุปสงฺกมิ, อุปสงฺกมิตฺวา ราชานํ มาคธํ เสนิยํ พิมฺพิสารํ เอตทโวจ “อาพาธํ เต เทว ปสฺสามา”ติ\csromansharedfootnote{aa0d7cb7f288d0e7}{ปสฺสามีติ (สฺยา)}. อถ โข ชีวโก โกมารภจฺโจ รญฺโญ มาคธสฺส เสนิยสฺส พิมฺพิสารสฺส ภคนฺทลา\-พาธํ เอเกเนว อาเลเปน อปกฑฺฒิ. อถ โข ราชา มาคโธ เสนิโย พิมฺพิสาโร อโรโค สมาโน ปญฺจ อิตฺถิสตานิ สพฺพาลงฺการํ ภูสาเปตฺวา โอมุญฺจาเปตฺวา ปุญฺชํ การาเปตฺวา ชีวกํ โกมารภจฺจํ เอตทโวจ “เอตํ ภเณ ชีวก ปญฺจนฺนํ อิตฺถิสตานํ สพฺพาลงฺการํ ตุยฺหํ โหตู”ติ. อลํ เทว อธิการํ เม เทโว สรตูติ. เตน หิ ภเณ ชีวก มํ อุปฏฺฐห อิตฺถาคารญฺจ พุทฺธปฺปมุขญฺจ ภิกฺขุ\-สํฆนฺติ. “เอวํ เทวา”ติ โข ชีวโก โกมารภจฺโจ รญฺโญ มาคธสฺส เสนิยสฺส พิมฺพิสารสฺส ปจฺจสฺโสสิ.}

\csromanmatikaanchor{mtk.213}
\csromantocmarkref{section}{205. ราชคหเสฏฺฐิวตฺถุ}{mtk.213}
\bookmark[dest={mtk.213},level=1]{205. ราชคหเสฏฺฐิวตฺถุ}
\csromanheader{1}{205. \textbf{ราชคห}\-\textbf{เสฏฺฐิ}\-\textbf{วตฺถุ}}
\proseitem{332}{เตน โข ปน สมเยน ราชคหกสฺส เสฏฺฐิสฺส สตฺตวสฺสิโก สีสาพาโธ โหติ, พหู มหนฺตา มหนฺตา ทิสาปาโมกฺขา เวชฺชา อาคนฺตฺวา นาสกฺขิํสุ อโรคํ กาตุํ, พหุํ หิรญฺญํ อาทาย อคมํสุ, อปิ จ เวชฺเชหิ ปจฺจกฺขาโต โหติ, เอกจฺเจ เวชฺชา เอวมาหํสุ “ปญฺจมํ ทิวสํ เสฏฺฐิ คหปติ กาลํ กริสฺสตี”ติ, เอกจฺเจ เวชฺชา เอวมาหํสุ “สตฺตมํ ทิวสํ เสฏฺฐิ คหปติ กาลํ กริสฺสตี”ติ. \csromanfolio{384}อถ โข ราชคหกสฺส เนคมสฺส เอตทโหสิ “อยํ โข เสฏฺฐิ คหปติ พหูปกาโร รญฺโญ เจว เนคมสฺส จ, อปิ จ เวชฺเชหิ ปจฺจกฺขาโต, เอกจฺเจ เวชฺชา เอวมาหํสุ ‘ปญฺจมํ ทิวสํ เสฏฺฐิ คหปติ กาลํ กริสฺสตี’ติ, เอกจฺเจ เวชฺชา เอวมาหํสุ ‘สตฺตมํ ทิวสํ เสฏฺฐิ คหปติ กาลํ กริสฺสตี'ติ, อยญฺจ รญฺโญ ชีวโก เวชฺโช ตรุโณ ภทฺรโก, ยนฺนูน มยํ ราชานํ ชีวกํ เวชฺชํ ยาเจยฺยาม เสฏฺฐิํ คหปติํ ติกิจฺฉิตุนฺ”ติ. อถ โข ราชคหโก เนคโม เยน ราชา มาคโธ เสนิโย พิมฺพิสาโร เตนุปสงฺกมิ, อุปสงฺกมิตฺวา ราชานํ มาคธํ เสนิยํ พิมฺพิสารํ เอตทโวจ “อยํ เทว เสฏฺฐิ คหปติ พหูปกาโร เทวสฺส เจว เนคมสฺส จ, อปิ จ เวชฺเชหิ ปจฺจกฺขาโต, เอกจฺเจ เวชฺชา เอวมาหํสุ ‘ปญฺจมํ ทิวสํ เสฏฺฐิ คหปติ กาลํ กริสฺสตี’ติ, เอกจฺเจ เวชฺชา เอวมาหํสุ ‘สตฺตมํ ทิวสํ เสฏฺฐิ คหปติ กาลํ กริสฺสตี’ติ, สาธุ เทโว ชีวกํ เวชฺชํ อาณาเปตุ เสฏฺฐิํ คหปติํ ติกิจฺฉิตุนฺ”ติ.}

\prose{อถ โข ราชา มาคโธ เสนิโย พิมฺพิสาโร ชีวกํ โกมารภจฺจํ อาณาเปสิ “คจฺฉ ภเณ ชีวก เสฏฺฐิํ คหปติํ ติกิจฺฉาหี”ติ. “เอวํ เทวา”ติ โข ชีวโก โกมารภจฺโจ รญฺโญ มาคธสฺส เสนิยสฺส พิมฺพิสารสฺส ปฏิสฺสุตฺวา เยน เสฏฺฐิ คหปติ เตนุปสงฺกมิ, อุปสงฺกมิตฺวา เสฏฺฐิสฺส คหปติสฺส วิการํ สลฺลกฺเขตฺวา เสฏฺฐิํ คหปติํ เอตทโวจ “สเจ ตฺวํ คหปติ อโรโค ภเวยฺยาสิ\csromansharedfootnote{7411365e10dfd504}{สจาหํ ตํ คหปติ อโรคาเปยฺยํ (สี), สจาหํ ตํ คหปติ อโรคํ กเรยฺยํ (สฺยา)}, กิํ เม อสฺส เทยฺยธมฺโม”ติ. สพฺพํ สาปเตยฺยญฺจ เต อาจริย โหตุ, อหญฺจ เต ทาโสติ. สกฺขิสฺสสิ ปน ตฺวํ คหปติ เอเกน ปสฺเสน สตฺตมาเส นิปชฺชิตุนฺติ. สกฺโกมหํ อาจริย เอเกน ปสฺเสน สตฺตมาเส นิปชฺชิตุนฺติ. สกฺขิสฺสสิ ปน ตฺวํ คหปติ ทุติเยน ปสฺเสน สตฺตมาเส นิปชฺชิตุนฺติ. สกฺโกมหํ อาจริย ทุติเยน ปสฺเสน สตฺตมาเส นิปชฺชิตุนฺติ. สกฺขิสฺสสิ ปน ตฺวํ คหปติ อุตฺตาโน สตฺตมาเส นิปชฺชิตุนฺติ. สกฺโกมหํ อาจริย อุตฺตาโน สตฺตมาเส นิปชฺชิตุนฺติ.}

\prose{อถ โข ชีวโก โกมารภจฺโจ เสฏฺฐิํ คหปติํ มญฺจเก นิปาเตตฺวา\csromansharedfootnote{7003fc0f21583c19}{นิปชฺชาเปตฺวา (สี, สฺยา)} มญฺจเก\csromansharedfootnote{9f883d1ff96d2833}{มญฺจเกน (สี)} สมฺพนฺธิตฺวา สีสจฺฉวิํ อุปฺปาเฏตฺวา\csromansharedfootnote{c3ac742cbe18f19f}{ผาเลตฺวา (สี)} สิพฺพินิํ วินาเมตฺวา \csromanfolio{385}เทฺว ปาณเก นีหริตฺวา มหาชนสฺส ทสฺเสสิ “ปสฺสถยฺเย\csromansharedfootnote{8d9af6395f64b18e}{ปสฺเสสฺยาถ (สี), ปสฺสถ (สฺยา), ปสฺสถยฺโย (ก)} อิเม เทฺว ปาณเก เอกํ ขุทฺทกํ เอกํ มหลฺลกํ. เย เต อาจริยา เอวมาหํสุ ‘ปญฺจมํ ทิวสํ เสฏฺฐิ คหปติ กาลํ กริสฺสตี’ติ, เตหายํ มหลฺลโก ปาณโก ทิฏฺโฐ, ปญฺจมํ ทิวสํ เสฏฺฐิสฺส คหปติสฺส มตฺถลุงฺคํ ปริยาทิยิสฺสติ, มตฺถลุงฺคสฺส ปริยาทานา เสฏฺฐิ คหปติ กาลํ กริสฺสติ, สุทิฏฺโฐ เตหิ อาจริเยหิ. เย เต อาจริยา เอวมาหํสุ ‘สตฺตมํ ทิวสํ เสฏฺฐิ คหปติ กาลํ กริสฺสตี’ติ, เตหายํ ขุทฺทโก ปาณโก ทิฏฺโฐ, สตฺตมํ ทิวสํ เสฏฺฐิสฺส คหปติสฺส มตฺถลุงฺคํ ปริยาทิยิสฺสติ, มตฺถลุงฺคสฺส ปริยาทานา เสฏฺฐิ คหปติ กาลํ กริสฺสติ, สุทิฏฺโฐ เตหิ อาจริเยหี”ติ. สิพฺพินิํ สมฺปฏิปาเฏตฺวา สีสจฺฉวิํ สิพฺพิตฺวา อาเลปํ อทาสิ. อถ โข เสฏฺฐิ คหปติ สตฺตาหสฺส อจฺจเยน ชีวกํ โกมารภจฺจํ เอตทโวจ “นาหํ อาจริย สกฺโกมิ เอเกน ปสฺเสน สตฺตมาเส นิปชฺชิตุนฺ”ติ. นนุ เม ตฺวํ คหปติ ปฏิสฺสุณิ “สกฺโกมหํ อาจริย เอเกน ปสฺเสน สตฺตมาเส นิปชฺชิตุนฺ”ติ. สจฺจาหํ อาจริย ปฏิสฺสุณิํ, อปาหํ มริสฺสามิ, นาหํ สกฺโกมิ เอเกน ปสฺเสน สตฺตมาเส นิปชฺชิตุนฺติ. เตน หิ ตฺวํ คหปติ ทุติเยน ปสฺเสน สตฺตมาเส นิปชฺชาหีติ. อถ โข เสฏฺฐิ คหปติ สตฺตาหสฺส อจฺจเยน ชีวกํ โกมารภจฺจํ เอตทโวจ “นาหํ อาจริย สกฺโกมิ ทุติเยน ปสฺเสน สตฺตมาเส นิปชฺชิตุนฺ”ติ. นนุ เม ตฺวํ คหปติ ปฏิสฺสุณิ “สกฺโกมหํ อาจริย ทุติเยน ปสฺเสน สตฺตมาเส นิปชฺชิตุนฺ”ติ. สจฺจาหํ อาจริย ปฏิสฺสุณิํ, อปาหํ มริสฺสามิ, นาหํ อาจริย สกฺโกมิ ทุติเยน ปสฺเสน สตฺตมาเส นิปชฺชิตุนฺติ. เตน หิ ตฺวํ คหปติ อุตฺตาโน สตฺตมาเส นิปชฺชาหีติ. อถ โข เสฏฺฐิ คหปติ สตฺตาหสฺส อจฺจเยน ชีวกํ โกมารภจฺจํ เอตทโวจ “นาหํ อาจริย สกฺโกมิ อุตฺตาโน สตฺตมาเส นิปชฺชิตุนฺ”ติ. นนุ เม ตฺวํ คหปติ ปฏิสฺสุณิ “สกฺโกมหํ อาจริย อุตฺตาโน สตฺตมาเส นิปชฺชิตุนฺ”ติ. สจฺจาหํ อาจริย ปฏิสฺสุณิํ, อปาหํ มริสฺสามิ, นาหํ สกฺโกมิ อุตฺตาโน สตฺตมาเส นิปชฺชิตุนฺติ. อหํ เจ ตํ คหปติ น วเทยฺยํ, เอตฺตกมฺปิ ตฺวํ น นิปชฺเชยฺยาสิ, อปิ จ ปฏิกจฺเจว มยา ญาโต “ตีหิ สตฺตาเหหิ เสฏฺฐิ คหปติ อโรโค ภวิสฺสตี”ติ. \csromanfolio{386}อุฏฺเฐหิ คหปติ อโรโคสิ, ชานาสิ กิํ เม เทยฺยธมฺโมติ. สพฺพํ สาปเตยฺยญฺจ เต อาจริย โหตุ, อหญฺจ เต ทาโสติ. อลํ คหปติ มา เม ตฺวํ สพฺพํ สาปเตยฺยํ อทาสิ, มา จ เม ทาโส, รญฺโญ สตสหสฺสํ เทหิ, มยฺหํ สตสหสฺสนฺติ. อถ โข เสฏฺฐิ คหปติ อโรโค สมาโน รญฺโญ สตสหสฺสํ อทาสิ, ชีวกสฺส โกมาร\-ภจฺจสฺส สตสหสฺสํ.}

\csromanmatikaanchor{mtk.214}
\csromantocmarkref{section}{206. เสฏฺฐิปุตฺตวตฺถุ}{mtk.214}
\bookmark[dest={mtk.214},level=1]{206. เสฏฺฐิปุตฺตวตฺถุ}
\csromanheader{1}{206. \textbf{เสฏฺฐิปุตฺต}\-\textbf{วตฺถุ}}
\proseitem{333}{เตน โข ปน สมเยน พาราณเส\-ยฺย\-กสฺส เสฏฺฐิปุตฺตสฺส โมกฺขจิกาย กีฬนฺตสฺส อนฺตคณฺฐา\-พาโธ โหติ, เยน ยาคุปิ ปีตา น สมฺมา ปริณามํ คจฺฉติ, ภตฺตมฺปิ ภุตฺตํ น สมฺมา ปริณามํ คจฺฉติ, อุจฺจาโรปิ ปสฺสาโวปิ น ปคุโณ, โส เตน กิโส โหติ ลูโข ทุพฺพณฺโณ อุปฺปณฺฑุ\-ปฺปณฺฑุก\-ชาโต ธมนิสนฺถต\-คตฺโต. อถ โข พาราณเส\-ยฺย\-กสฺส เสฏฺฐิสฺส เอตทโหสิ “มยฺหํ โข ปุตฺตสฺส ตาทิโส อาพาโธ, เยน ยาคุปิ ปีตา น สมฺมา ปริณามํ คจฺฉติ, ภตฺตมฺปิ ภุตฺตํ น สมฺมา ปริณามํ คจฺฉติ, อุจฺจาโรปิ ปสฺสาโวปิ น ปคุโณ, โส เตน กิโส ลูโข ทุพฺพณฺโณ อุปฺปณฺฑุ\-ปฺปณฺฑุก\-ชาโต ธมนิสนฺถต\-คตฺโต, ยนฺนูนาหํ ราชคหํ คนฺตฺวา ราชานํ ชีวกํ เวชฺชํ ยาเจยฺยํ ปุตฺตํ เม ติกิจฺฉิตุนฺ”ติ. อถ โข พาราณเสยฺยโก เสฏฺฐิ ราชคหํ คนฺตฺวา เยน ราชา มาคโธ เสนิโย พิมฺพิสาโร เตนุปสงฺกมิ, อุปสงฺกมิตฺวา ราชานํ มาคธํ เสนิยํ พิมฺพิสารํ เอตทโวจ “มยฺหํ โข เทว ปุตฺตสฺส ตาทิโส อาพาโธ, เยน ยาคุปิ ปีตา น สมฺมา ปริณามํ คจฺฉติ, ภตฺตมฺปิ ภุตฺตํ น สมฺมา ปริณามํ คจฺฉติ, อุจฺจาโรปิ ปสฺสาโวปิ น ปคุโณ, โส เตน กิโส ลูโข ทุพฺพณฺโณ อุปฺปณฺฑุ\-ปฺปณฺฑุก\-ชาโต ธมนิสนฺถต\-คตฺโต, สาธุ เทโว ชีวกํ เวชฺชํ อาณาเปตุ ปุตฺตํ เม ติกิจฺฉิตุนฺ”ติ.}

\prose{อถ โข ราชา มาคโธ เสนิโย พิมฺพิสาโร ชีวกํ โกมารภจฺจํ อาณาเปสิ “คจฺฉ ภเณ ชีวก พาราณสิํ คนฺตฺวา พาราณเสยฺยกํ เสฏฺฐิปุตฺตํ ติกิจฺฉาหี”ติ. “เอวํ เทวา”ติ โข ชีวโก โกมารภจฺโจ รญฺโญ มาคธสฺส เสนิยสฺส พิมฺพิสารสฺส ปฏิสฺสุตฺวา พาราณสิํ คนฺตฺวา เยน พาราณเสยฺยโก เสฏฺฐิปุตฺโต เตนุปสงฺกมิ, อุปสงฺกมิตฺวา พาราณเส\-ยฺย\-กสฺส เสฏฺฐิปุตฺตสฺส วิการํ สลฺลกฺเขตฺวา ชนํ \csromanfolio{387}อุสฺสาเรตฺวา ติโรกรณิยํ ปริกฺขิปิตฺวา ถมฺเภ อุพฺพนฺธิตฺวา\csromansharedfootnote{fb1c4e5c7ce9d218}{อุปนิพนฺธิตฺวา (สี, สฺยา)} ภริยํ ปุรโต ฐเปตฺวา อุทรจฺฉวิํ อุปฺปาเฏตฺวา อนฺตคณฺฐิํ นีหริตฺวา ภริยาย ทสฺเสสิ “ปสฺส เต สามิกสฺส อาพาธํ, อิมินา ยาคุปิ ปีตา น สมฺมา ปริณามํ คจฺฉติ, ภตฺตมฺปิ ภุตฺตํ น สมฺมา ปริณามํ คจฺฉติ, อุจฺจาโรปิ ปสฺสาโวปิ น ปคุโณ, อิมินายํ กิโส ลูโข ทุพฺพณฺโณ อุปฺปณฺฑุ\-ปฺปณฺฑุก\-ชาโต ธมนิสนฺถต\-คตฺโต”ติ. อนฺตคณฺฐิํ วินิเวเฐตฺวา อนฺตานิ ปฏิปเวเสตฺวา อุทรจฺฉวิํ สิพฺพิตฺวา อาเลปํ อทาสิ. อถ โข พาราณเสยฺยโก เสฏฺฐิปุตฺโต น จิรสฺเสว อโรโค อโหติ. อถ โข พารณเสยฺยโก เสฏฺฐิ “ปุตฺโต เม อโรโค ฐิโต”ติ\csromansharedfootnote{32c44b2308ddab8f}{อโรคาปิโตติ (สี)} ชีวกสฺส โกมาร\-ภจฺจสฺส โสฬส\-สหสฺสานิ ปาทาสิ. อถ โข ชีวโก โกมารภจฺโจ ตานิ โสฬส\-สหสฺสานิ อาทาย ปุนเทว ราชคหํ ปจฺจาคญฺฉิ.}

\csromanmatikaanchor{mtk.215}
\csromantocmarkref{section}{207. ปชฺโชตราชวตฺถุ}{mtk.215}
\bookmark[dest={mtk.215},level=1]{207. ปชฺโชตราชวตฺถุ}
\csromanheader{1}{207. \textbf{ปชฺโชต}\-\textbf{ราช}\-\textbf{วตฺถุ}}
\proseitem{334}{เตน โข ปน สมเยน รญฺโญ\csromansharedfootnote{ae2274917af00741}{อุชฺเชนิยํ รญฺโญ (สฺยา)} ปชฺโชตสฺส ปณฺฑุโรคาพาโธ โหติ, พหู มหนฺตา มหนฺตา ทิสาปาโมกฺขา เวชฺชา อาคนฺตฺวา นาสกฺขิํสุ อโรคํ กาตุํ, พหุํ หิรญฺญํ อาทาย อคมํสุ. อถ โข ราชา ปชฺโชโต รญฺโญ มาคธสฺส เสนิยสฺส พิมฺพิสารสฺส สนฺติเก ทูตํ ปาเหสิ “มยฺหํ โข ตาทิโส อาพาโธ, สาธุ เทโว ชีวกํ เวชฺชํ อาณาเปตุ, โส มํ ติกิจฺฉิสฺสตี”ติ. อถ โข ราชา มาคโธ เสนิโย พิมฺพิสาโร ชีวกํ โกมารภจฺจํ อาณาเปสิ “คจฺฉ ภเณ ชีวก อุชฺเชนิํ คนฺตฺวา ราชานํ ปชฺโชตํ ติกิจฺฉาหี”ติ. “เอวํ เทวา”ติ โข ชีวโก โกมารภจฺโจ รญฺโญ มาคธสฺส เสนิยสฺส พิมฺพิสารสฺส ปฏิสฺสุตฺวา อุชฺเชนิํ คนฺตฺวา เยน ราชา ปชฺโชโต เตนุปสงฺกมิ, อุปสงฺกมิตฺวา รญฺโญ ปชฺโชตสฺส วิการํ สลฺลกฺเขตฺวา ราชานํ ปชฺโชตํ เอตทโวจ “สปฺปิํ เทหิ\csromansharedfootnote{805459fff1190d69}{อิทํ ปททฺวยํ สีสฺยาโปตฺถเกสุ นตฺถิ.}, สปฺปิํ เทว นิปฺปจิสฺสามิ, ตํ เทโว ปิวิสฺสตี”ติ. อลํ ภเณ ชีวก, ยํ เต สกฺกา วินา สปฺปินา อโรคํ กาตุํ, ตํ กโรหิ, เชคุจฺฉํ เม สปฺปิ ปฏิกูลนฺติ. อถ โข ชีวกสฺส โกมาร\-ภจฺจสฺส เอตทโหสิ “อิมสฺส โข รญฺโญ ตาทิโส \csromanfolio{388}อาพาโธ น สกฺกา วินา สปฺปินา อโรคํ กาตุํ, ยนฺนูนาหํ สปฺปิํ นิปฺปเจยฺยํ กสาววณฺณํ กสาวคนฺธํ กสาวรสนฺ”ติ. อถ โข ชีวโก โกมารภจฺโจ นานาเภสชฺเชหิ สปฺปิํ นิปฺปจิ กสาววณฺณํ กสาวคนฺธํ กสาวรสํ. อถ โข ชีวกสฺส โกมาร\-ภจฺจสฺส เอตทโหสิ “อิมสฺส โข รญฺโญ สปฺปิ ปีตํ ปริณาเมนฺตํ อุทฺเทกํ ทสฺสติ, จณฺโฑยํ ราชา ฆาตาเปยฺยาปิ มํ, ยนฺนูนาหํ ปฏิกจฺเจว อาปุจฺเฉยฺยนฺ”ติ. อถ โข ชีวโก โกมารภจฺโจ เยน ราชา ปชฺโชโต เตนุปสงฺกมิ, อุปสงฺกมิตฺวา ราชานํ ปชฺโชตํ เอตทโวจ “มยํ โข เทว เวชฺชา นาม ตาทิเสน มุหุตฺเตน มูลานิ อุทฺธราม, เภสชฺชานิ สํหราม, สาธุ เทโว วาหนาคาเรสุ จ ทฺวาเรสุ จ อาณาเปตุ ‘เยน วาหเนน ชีวโก อิจฺฉติ เตน วาหเนน คจฺฉตุ, เยน ทฺวาเรน อิจฺฉติ เตน ทฺวาเรน คจฺฉตุ, ยํ กาลํ อิจฺฉติ ตํ กาลํ คจฺฉตุ, ยํ กาลํ อิจฺฉติ ตํ กาลํ ปวิสตู’ติ”. อถ โข ราชา ปชฺโชโต วาหนาคาเรสุ จ ทฺวาเรสุ จ อาณาเปสิ “เยน วาหเนน ชีวโก อิจฺฉติ เตน วาหเนน คจฺฉตุ, เยน ทฺวาเรน อิจฺฉติ เตน ทฺวาเรน คจฺฉตุ, ยํ กาลํ อิจฺฉติ ตํ กาลํ คจฺฉตุ, ยํ กาลํ อิจฺฉติ ตํ กาลํ ปวิสตู”ติ.}

\prose{เตน โข ปน สมเยน รญฺโญ ปชฺโชตสฺส ภทฺทวติกา นาม หตฺถินิกา ปญฺญาส\-โยชนิกา โหติ. อถ โข ชีวโก โกมารภจฺโจ รญฺโญ ปชฺโชตสฺส สปฺปิํ\csromansharedfootnote{776ab296fca809ec}{ตํ สปฺปิํ (สฺยา)} อุปนาเมสิ “กสาวํ เทโว ปิวตู”ติ. อถ โข ชีวโก โกมารภจฺโจ ราชานํ ปชฺโชตํ สปฺปิํ ปาเยตฺวา หตฺถิสาลํ คนฺตฺวา ภทฺทวติกาย หตฺถินิกาย นครมฺหา นิปฺปติ.}

\prose{อถ โข รญฺโญ ปชฺโชตสฺส ตํ สปฺปิ ปีตํ ปริณาเมนฺตํ อุทฺเทกํ อทาสิ. อถ โข ราชา ปชฺโชโต มนุสฺเส เอตทโวจ “ทุฏฺเฐน ภเณ ชีวเกน สปฺปิํ ปายิโตมฺหิ, เตน หิ ภเณ ชีวกํ เวชฺชํ วิจินถา”ติ. ภทฺทวติกาย เทว หตฺถินิกาย นครมฺหา นิปฺปติโตติ. เตน โข ปน สมเยน รญฺโญ ปชฺโชตสฺส กาโก นาม ทาโส สฏฺฐิโยชนิโก โหติ, อมนุสฺเสน ปฏิจฺจ ชาโต. อถ โข ราชา ปชฺโชโต กากํ ทาสํ อาณาเปสิ “คจฺฉ ภเณ กาก ชีวกํ \csromanfolio{389}เวชฺชํ นิวตฺเตหิ ‘ราชา ตํ อาจริย นิวตฺตาเปตี’ติ, เอเต โข ภเณ กาก เวชฺชา นาม พหุมายา, มา จสฺส กิญฺจิ ปฏิคฺคเหสี”ติ.}

\prose{อถ โข กาโก ทาโส ชีวกํ โกมารภจฺจํ อนฺตรามคฺเค โกสมฺพิยํ สมฺภาเวสิ ปาตราสํ กโรนฺตํ. อถ โข กาโก ทาโส ชีวกํ กามารภจฺจํ เอตทโวจ “ราชา ตํ อาจริย นิวตฺตาเปตี”ติ. อาคเมหิ ภเณ กาก ยาว ภุญฺชาม\csromansharedfootnote{913d876eedf5a57c}{ภุญฺชามิ (สี, สฺยา)}, หนฺท ภเณ กาก ภุญฺชสฺสูติ. อลํ อาจริย รญฺญามฺหิ อาณตฺโต “เอเต โข ภเณ กาก เวชฺชา นาม พหุมายา, มา จสฺส กิญฺจิ ปฏิคฺคเหสี”ติ. เตน โข ปน สมเยน ชีวโก โกมารภจฺโจ นเขน เภสชฺชํ โอลุมฺเปตฺวา อามลกญฺจ ขาทติ ปานียญฺจ ปิวติ. อถ โข ชีวโก โกมารภจฺโจ กากํ ทาสํ เอตทโวจ “หนฺท ภเณ กาก อามลกญฺจ ขาท ปานียญฺจ ปิวสฺสู”ติ. อถ โข กาโก ทาโส “อยํ โข เวชฺโช อามลกญฺจ ขาทติ ปานียญฺจ ปิวติ, น อรหติ กิญฺจิ ปาปกํ โหตุนฺ”ติ อุปฑฺฒามลกญฺจ ขาทิ ปานียญฺจ อปายิ, ตสฺส ตํ อุปฑฺฒามลกํ ขาทิตํ ตตฺเถว นิจฺฉาเรสิ. อถ โข กาโก ทาโส ชีวกํ โกมารภจฺจํ เอตทโวจ “อตฺถิ เม อาจริย ชีวิตนฺ”ติ. “มา ภเณ กาก ภายิ, ตฺวํ เจว อโรโค ภวิสฺสสิ ราชา จ, จณฺโฑ โส ราชา ฆาตาเปยฺยาปิ มํ, เตนาหํ น นิวตฺตามี”ติ ภทฺทวติกํ หตฺถินิกํ กากสฺส นิยฺยาเทตฺวา เยน ราชคหํ เตน ปกฺกามิ. อนุปุพฺเพน เยน ราชา มาคโธ เสนิโย พิมฺพิสาโร เตนุปสงฺกมิ, อุปสงฺกมิตฺวา รญฺโญ มาคธสฺส เสนิยสฺส พิมฺพิสารสฺส เอตมตฺถํ อาโรเจสิ. สุฏฺฐุ ภเณ ชีวก อกาสิ ยมฺปิ น นิวตฺโต, จณฺโฑ โส ราชา ฆาตาเปยฺยาปิ ตนฺติ. อถ โข ราชา ปชฺโชโต อโรโค สมาโน ชีวกสฺส โกมาร\-ภจฺจสฺส สนฺติเก ทูตํ ปาเหสิ “อาคจฺฉตุ ชีวโก, วรํ ทสฺสามี”ติ. อลํ อยฺโย\csromansharedfootnote{5c7596c023c947a9}{เทว (สฺยา)} อธิการํ เม เทโว สรตูติ.}

\csromanmatikaanchor{mtk.216}
\csromantocmarkref{section}{208. สิเวยฺยกทุสฺสยุคกถา}{mtk.216}
\bookmark[dest={mtk.216},level=1]{208. สิเวยฺยกทุสฺสยุคกถา}
\csromanheader{1}{208. \textbf{สิเวยฺยก}\-\textbf{ทุสฺส}\-\textbf{ยุค}\-\textbf{กถา}}
\csromanmatikaanchor{mtk.217}
\csromantocmarkref{section}{209. สมติํสวิเรจนกถา}{mtk.217}
\bookmark[dest={mtk.217},level=1]{209. สมติํสวิเรจนกถา}
\proseitem{335}{เตน โข ปน สมเยน รญฺโญ ปชฺโชตสฺส สิเวยฺยกํ ทุสฺสยุคํ อุปฺปนฺนํ โหติ พหูนํ\csromansharedfootnote{9793b89f39b16a06}{พหุนฺนํ (สี, สฺยา)} ทุสฺสานํ พหูนํ ทุสฺสยุคานํ พหูนํ ทุสฺสยุค\-สตานํ พหูนํ ทุสฺส\-ยุค\-สหสฺสานํ พหูนํ ทุสฺสยุค\-สตสหสฺสานํ \csromanfolio{390}อคฺคญฺจ เสฏฺฐญฺจ โมกฺขญฺจ อุตฺตมญฺจ ปวรญฺจ. อถ โข ราชา ปชฺโชโต ตํ สิเวยฺยกํ ทุสฺสยุคํ ชีวกสฺส โกมาร\-ภจฺจสฺส ปาเหสิ. อถ โข ชีวกสฺส โกมาร\-ภจฺจสฺส เอตทโหสิ “อิทํ โข เม สิเวยฺยกํ ทุสฺสยุคํ รญฺญา ปชฺโชเตน ปหิตํ พหูนํ ทุสฺสานํ พหูนํ ทุสฺสยุคานํ พหูนํ ทุสฺสยุค\-สตานํ พหูนํ ทุสฺส\-ยุค\-สหสฺสานํ พหูนํ ทุสฺสยุค\-สตสหสฺสานํ อคฺคญฺจ เสฏฺฐญฺจ โมกฺขญฺจ อุตฺตมญฺจ ปวรญฺจ, นยิทํ อญฺโญ โกจิ ปจฺจารหติ อญฺญตฺร เตน ภควตา อรหตา สมฺมา\-สมฺพุทฺเธน รญฺญา วา มาคเธน เสนิเยน พิมฺพิสาเรนา”ติ.}

\csromanheader{2}{209. \textbf{สมตฺติํส}\-\textbf{วิเรจน}\-\textbf{กถา}}
\proseitem{336}{เตน โข ปน สมเยน ภควโต กาโย โทสาภิสนฺโน โหติ. อถ โข ภควา อายสฺมนฺตํ อานนฺทํ อามนฺเตสิ “โทสาภิสนฺโน โข อานนฺท ตถาคตสฺส กาโย, อิจฺฉติ ตถาคโต วิเรจนํ ปาตุนฺ”ติ. อถ โข อายสฺมา อานนฺโท เยน ชีวโก โกมารภจฺโจ เตนุปสงฺกมิ, อุปสงฺกมิตฺวา ชีวกํ โกมารภจฺจํ เอตทโวจ “โทสาภิสนฺโน โข อาวุโส ชีวก ตถาคตสฺส กาโย, อิจฺฉติ ตถาคโต วิเรจนํ ปาตุนฺ”ติ. เตน หิ ภนฺเต อานนฺท ภควโต กายํ กติปาหํ สิเนเหถาติ. อถ โข อายสฺมา อานนฺโท ภควโต กายํ กติปาหํ สิเนเหตฺวา เยน ชีวโก โกมารภจฺโจ เตนุปสงฺกมิ, อุปสงฺกมิตฺวา ชีวกํ โกมารภจฺจํ เอตทโวจ “สินิทฺโธ โข อาวุโส ชีวก ตถาคตสฺส กาโย, ยสฺส ทานิ กาลํ มญฺญสี”ติ. อถ โข ชีวกสฺส โกมาร\-ภจฺจสฺส เอตทโหสิ “น โข เมตํ ปติรูปํ โยหํ ภควโต โอฬาริกํ วิเรจนํ ทเทยฺยนฺ”ติ, ตีณิ อุปฺปลหตฺถานิ นานาเภสชฺเชหิ ปริภาเวตฺวา เยน ภควา เตนุปสงฺกมิ, อุปสงฺกมิตฺวา เอกํ อุปฺปลหตฺถํ ภควโต อุปนาเมสิ “อิมํ ภนฺเต ภควา ปฐมํ อุปฺปลหตฺถํ อุปสิงฺฆตุ, อิทํ ภควนฺตํ ทสกฺขตฺตุํ วิเรเจสฺสตี”ติ, ทุติยํ อุปฺปลหตฺถํ ภควโต อุปนาเมสิ “อิมํ ภนฺเต ภควา ทุติยํ อุปฺปลหตฺถํ อุปสิงฺฆตุ, อิทํ ภควนฺตํ ทสกฺขตฺตุํ วิเรเจสฺสตี”ติ, ตติยํ อุปฺปลหตฺถํ ภควโต อุปนาเมสิ “อิมํ ภนฺเต ภควา ตติยํ อุปฺปลหตฺถํ อุปสิงฺฆตุ, อิทํ ภควนฺตํ ทสกฺขตฺตุํ \csromanfolio{391}วิเรเจสฺสตี”ติ, เอวํ ภควโต สมตฺติํสาย วิเรจนํ ภวิสฺสตีติ. อถ โข ชีวโก โกมารภจฺโจ ภควโต สมตฺติํสาย วิเรจนํ ทตฺวา ภควนฺตํ อภิวาเทตฺวา ปทกฺขิณํ กตฺวา ปกฺกามิ. อถ โข ชีวกสฺส โกมาร\-ภจฺจสฺส พหิ ทฺวารโกฏฺฐกา นิกฺขนฺตสฺส เอตทโหสิ “มยา โข ภควโต สมตฺติํสาย วิเรจนํ ทินฺนํ, โทสาภิสนฺโน ตถาคตสฺส กาโย, น ภควนฺตํ สมตฺติํสกฺขตฺตุํ วิเรเจสฺสติ, เอกูน\-ตฺติํสกฺขตฺตุํ ภควนฺตํ วิเรเจสฺสติ, อปิ จ ภควา วิริตฺโต นหายิสฺสติ, นหาตํ ภควนฺตํ สกิํ วิเรเจสฺสติ, เอวํ ภควโต สมตฺติํสาย วิเรจนํ ภวิสฺสตี”ติ.}

\prose{อถ โข ภควา ชีวกสฺส โกมาร\-ภจฺจสฺส เจตสา เจโต\-ปริวิตกฺกม\-ญฺญาย อายสฺมนฺตํ อานนฺทํ อามนฺเตสิ “อิธานนฺท ชีวกสฺส โกมาร\-ภจฺจสฺส พหิ ทฺวารโกฏฺฐกา นิกฺขนฺตสฺส เอตทโหสิ ‘มยา โข ภควโต สมตฺติํสาย วิเรจนํ ทินฺนํ, โทสาภิสนฺโน ตถาคตสฺส กาโย, น ภควนฺตํ สมตฺติํสกฺขตฺตุํ วิเรเจสฺสติ, เอกูน\-ตฺติํสกฺขตฺตุํ ภควนฺตํ วิเรเจสฺสติ, อปิ จ ภควา วิริตฺโต นหายิสฺสติ, นหาตํ ภควนฺตํ สกิํ วิเรเจสฺสติ, เอวํ ภควโต สมตฺติํสาย วิเรจนํ ภวิสฺสตี’ติ, เตน หานนฺท อุโณฺหทกํ ปฏิยาเทหี”ติ. “เอวํ ภนฺเต”ติ โข อายสฺมา อานนฺโท ภควโต ปฏิสฺสุณิตฺวา อุโณฺหทกํ ปฏิยาเทสิ.}

\prose{อถ โข ชีวโก โกมารภจฺโจ เยน ภควา เตนุปสงฺกมิ, อุปสงฺกมิตฺวา ภควนฺตํ อภิวาเทตฺวา เอกมนฺตํ นิสีทิ, เอกมนฺตํ นิสินฺโน โข ชีวโก โกมารภจฺโจ ภควนฺตํ เอตทโวจ “วิริตฺโต ภนฺเต ภควา”ติ. วิริตฺโตมฺหิ ชีวกาติ. อิธ มยฺหํ ภนฺเต พหิ ทฺวารโกฏฺฐกา นิกฺขนฺตสฺส เอตทโหสิ “มยา โข ภควโต สมตฺติํสาย วิเรจนํ ทินฺนํ, โทสาภิสนฺโน ตถาคตสฺส กาโย, น ภควนฺตํ สมตฺติํสกฺขตฺตุํ วิเรเจสฺสติ, เอกูน\-ตฺติํสกฺขตฺตุํ ภควนฺตํ วิเรเจสฺสติ, อปิ จ ภควา วิริตฺโต นหายิสฺสติ, นหาตํ ภควนฺตํ สกิํ วิเรเจสฺสติ, เอวํ ภควโต สมตฺติํสาย วิเรจนํ ภวิสฺสตี”ติ. นหายตุ ภนฺเต ภควา, นหายตุ สุคโตติ. อถ โข ภควา อุโณฺหทกํ นหายิ, นหาตํ ภควนฺตํ สกิํ วิเรเจสิ, เอวํ ภควโต สมตฺติํสาย วิเรจนํ อโหสิ. อถ โข ชีวโก โกมารภจฺโจ \csromanfolio{392}ภควนฺตํ เอตทโวจ “ยาว ภนฺเต ภควโต กาโย ปกตตฺโต โหติ, อลํ\csromansharedfootnote{278837b608442c94}{อหํ ตาว ยูสปิณฺฑปาเตนาติ (สี), อลํ ยูสปิณฺฑเกนาติ (สฺยา)} ยูสปิณฺฑปาเตนา”ติ.}

\csromanmatikaanchor{mtk.218}
\csromantocmarkref{section}{210. วรยาจนากถา}{mtk.218}
\bookmark[dest={mtk.218},level=1]{210. วรยาจนากถา}
\csromanheader{1}{210. \textbf{วรยาจนา}\-\textbf{กถา}}
\proseitem{337}{อถ โข ภควโต กาโย น จิรสฺเสว ปกตตฺโต อโหสิ. อถ โข ชีวโก โกมารภจฺโจ ตํ สิเวยฺยกํ ทุสฺสยุคํ อาทาย เยน ภควา เตนุปสงฺกมิ, อุปสงฺกมิตฺวา ภควนฺตํ อภิวาเทตฺวา เอกมนฺตํ นิสีทิ, เอกมนฺตํ นิสินฺโน โข ชีวโก โกมารภจฺโจ ภควนฺตํ เอตทโวจ “เอกาหํ ภนฺเต ภควนฺตํ วรํ ยาจามี”ติ. อติกฺกนฺตวรา โข ชีวก ตถาคตาติ. ยญฺจ ภนฺเต กปฺปติ ยญฺจ อนวชฺชนฺติ. วเทหิ ชีวกาติ. ภควา ภนฺเต ปํสุกูลิโก ภิกฺขุสํโฆ จ, อิทํ เม ภนฺเต สิเวยฺยกํ ทุสฺสยุคํ รญฺญา ปชฺโชเตน ปหิตํ พหูนํ ทุสฺสานํ พหูนํ ทุสฺสยุคานํ พหูนํ ทุสฺสยุค\-สตานํ พหูนํ ทุสฺส\-ยุค\-สหสฺสานํ พหูนํ ทุสฺสยุค\-สตสหสฺสานํ อคฺคญฺจ เสฏฺฐญฺจ โมกฺขญฺจ อุตฺตมญฺจ ปวรญฺจ, ปฏิคฺคณฺหาตุ เม ภนฺเต ภควา สิเวยฺยกํ ทุสฺสยุคํ, ภิกฺขุ\-สํฆสฺส จ คหปติ\-จีวรํ อนุชานาตูติ. ปฏิคฺคเหสิ ภควา สิเวยฺยกํ ทุสฺสยุคํ. อถ โข ภควา ชีวกํ โกมารภจฺจํ ธมฺมิยา กถาย สนฺทสฺเสสิ สมาทเปสิ สมุตฺเตเชสิ สมฺปหํเสสิ. อถ โข ชีวโก โกมารภจฺโจ ภควตา ธมฺมิยา กถาย สนฺทสฺสิโต สมาทปิโต สมุตฺเตชิโต สมฺปหํสิโต อุฏฺฐายาสนา ภควนฺตํ อภิวาเทตฺวา ปทกฺขิณํ กตฺวา ปกฺกามิ. อถ โข ภควา เอตสฺมิํ นิทาเน เอตสฺมิํ ปกรเณ ธมฺมิํ กถํ กตฺวา ภิกฺขู อามนฺเตสิ “อนุชานามิ ภิกฺขเว คหปติ\-จีวรํ, โย อิจฺฉติ, ปํสุกูลิโก โหตุ, โย อิจฺฉติ, คหปติ\-จีวรํ สาทิยตุ, อิตรีตเรน\-ปาหํ\csromansharedfootnote{1b62138175829d7c}{ปหํ (สี), จาหํ (สฺยา)} ภิกฺขเว สนฺตุฏฺฐิํ วณฺเณมี”ติ.}

\prose{อสฺโสสุํ โข ราชคเห มนุสฺสา “ภควตา กิร ภิกฺขูนํ คหปติ\-จีวรํ อนุญฺญาตนฺ”ติ, เต จ มนุสฺสา หฏฺฐา อเหสุํ อุทคฺคา “อิทานิ โข มยํ ทานานิ ทสฺสาม ปุญฺญานิ กริสฺสาม, ยโต ภควตา ภิกฺขูนํ คหปติ\-จีวรํ อนุญฺญาตนฺ”ติ, เอกาเหเนว ราชคเห พหูนิ จีวร\-สหสฺสานิ อุปฺปชฺชิํสุ.}

\prose{\csromanfolio{393}อสฺโสสุํ โข ชานปทา มนุสฺสา “ภควตา กิร ภิกฺขูนํ คหปติ\-จีวรํ อนุญฺญาตนฺ”ติ, เต จ มนุสฺสา หฏฺฐา อเหสุํ อุทคฺคา “อิทานิ โข มยํ ทานานิ ทสฺสาม ปุญฺญานิ กริสฺสาม, ยโต ภควตา ภิกฺขูนํ คหปติ\-จีวรํ อนุญฺญาตนฺ”ติ, ชนปเทปิ เอกาเหเนว พหูนิ จีวร\-สหสฺสานิ อุปฺปชฺชิํสุ.}

\proseitem{337}{เตน โข ปน สมเยน สํฆสฺส ปาวาโร อุปฺปนฺโน โหติ. ภควโต เอตมตฺถํ อาโรเจสุํ. อนุชานามิ ภิกฺขเว ปาวารนฺติ.}

\prose{โกเสยฺยปาวาโร อุปฺปนฺโน โหติ. ภควโต เอตมตฺถํ อาโรเจสุํ. อนุชานามิ ภิกฺขเว โกเสยฺยปาวารนฺติ.}

\prosewithcloserruled{โกชวํ อุปฺปนฺนํ โหติ. ภควโต เอตมตฺถํ อาโรเจสุํ. อนุชานามิ ภิกฺขเว โกชวนฺติ.}{ปฐม\-ภาณวาโร นิฏฺฐิโต.}
\csromanmatikaanchor{mtk.219}
\csromantocmarkref{section}{211. กมฺพลานุชานนาทิกถา}{mtk.219}
\bookmark[dest={mtk.219},level=1]{211. กมฺพลานุชานนาทิกถา}
\csromanheader{1}{211. \textbf{กมฺพลา}\-\textbf{นุชานนา}\-\textbf{ทิ}\-\textbf{กถา}}
\proseitem{338}{เตน โข ปน สมเยน กาสิราชา ชีวกสฺส โกมาร\-ภจฺจสฺส อฑฺฒกาสิกํ กมฺพลํ ปาเหสิ อุปฑฺฒกาสินํ ขมมานํ. อถ โข ชีวโก โกมารภจฺโจ ตํ อฑฺฒกาสิกํ กมฺพลํ อาทาย เยน ภควา เตนุปสงฺกมิ, อุปสงฺกมิตฺวา ภควนฺตํ อภิวาเทตฺวา เอกมนฺตํ นิสีทิ, เอกมนฺตํ นิสินฺโน โข ชีวโก โกมารภจฺโจ ภควนฺตํ เอตทโวจ “อยํ เม ภนฺเต อฑฺฒกาสิโก กมฺพโล กาสิรญฺญา ปหิโต อุปฑฺฒกาสินํ ขมมาโน, ปฏิคฺคณฺหาตุ เม ภนฺเต ภควา กมฺพลํ, ยํ มมสฺส ทีฆรตฺตํ หิตาย สุขายา”ติ. ปฏิคฺคเหสิ ภควา กมฺพลํ. อถ โข ภควา ชีวกํ โกมารภจฺจํ ธมฺมิยา กถาย สนฺทสฺเสสิ ฯเปฯ ปทกฺขิณํ กตฺวา ปกฺกามิ. อถ โข ภควา เอตสฺมิํ นิทาเน เอตสฺมิํ ปกรเณ ธมฺมิํ กถํ กตฺวา ภิกฺขู อามนฺเตสิ “อนุชานามิ ภิกฺขเว กมฺพลนฺ”ติ.}

\proseitem{339}{เตน โข ปน สมเยน สํฆสฺส อุจฺจาวจานิ จีวรานิ อุปฺปนฺนานิ โหนฺติ. อถ โข ภิกฺขูนํ เอตทโหสิ “กิํ นุ โข ภควตา จีวรํ \csromanfolio{394}อนุญฺญาตํ กิํ อนนุญฺญาตนฺ”ติ. ภควโต เอตมตฺถํ อาโรเจสุํ. อนุชานามิ ภิกฺขเว ฉ จีวรานิ โขมํ กปฺปาสิกํ โกเสยฺยํ กมฺพลํ สาณํ ภงฺคนฺติ.}

\proseitem{340}{เตน โข ปน สมเยน เย เต ภิกฺขู คหปติ\-จีวรํ สาทิยนฺติ, เต กุกฺกุจฺจายนฺตา ปํสุกูลํ น สาทิยนฺติ “เอกํเยว ภควตา จีวรํ อนุญฺญาตํ, น เทฺว”ติ. ภควโต เอตมตฺถํ อาโรเจสุํ. อนุชานามิ ภิกฺขเว คหปติ\-จีวรํ สาทิยนฺเตน ปํสุกูลมฺปิ สาทิยิตุํ, ตทุภเยน\-ปาหํ ภิกฺขเว สนฺตุฏฺฐิํ วณฺเณมีติ.}

\csromanmatikaanchor{mtk.220}
\csromantocmarkref{section}{212. ปํสุกูลปริเยสนกถา}{mtk.220}
\bookmark[dest={mtk.220},level=1]{212. ปํสุกูลปริเยสนกถา}
\csromanheader{1}{212. \textbf{ปํสุกูล}\-\textbf{ปริเยสน}\-\textbf{กถา}}
\proseitem{341}{เตน โข ปน สมเยน สมฺพหุลา ภิกฺขู โกสเลสุ ชนปเท อทฺธาน\-มคฺค\-ปฺปฏิปนฺนา โหนฺติ. เอกจฺเจ ภิกฺขู สุสานํ โอกฺกมิํสุ ปํสุกูลาย, เอกจฺเจ ภิกฺขู นาคเมสุํ. เย เต ภิกฺขู สุสานํ โอกฺกมิํสุ ปํสุกูลาย, เต ปํสุกูลานิ ลภิํสุ. เย เต ภิกฺขู นาคเมสุํ, เก เอวมาหํสุ “อมฺหากมฺปิ อาวุโส ภาคํ เทถา”ติ. เต เอวมาหํสุ “น มยํ อาวุโส ตุมฺหากํ ภาคํ ทสฺสาม, กิสฺส ตุเมฺห นาคมิตฺถา”ติ. ภควโต เอตมตฺถํ อาโรเจสุํ. อนุชานามิ ภิกฺขเว นาคเมนฺตานํ นากามา ภาคํ ทาตุนฺติ.}

\prose{เตน โข ปน สมเยน สมฺพหุลา ภิกฺขู โกสเลสุ ชนปเท อทฺธาน\-มคฺค\-ปฺปฏิปนฺนา โหนฺติ. เอกจฺเจ ภิกฺขู สุสานํ โอกฺกมิํสุ ปํสุกูลาย, เอกจฺเจ ภิกฺขู อาคเมสุํ. เย เต ภิกฺขู สุสานํ โอกฺกมิํสุ ปํสุกูลาย, เต ปํสุกูลานิ ลภิํสุ. เย เต ภิกฺขู อาคเมสุํ, เต เอวมาหํสุ “อมฺหากมฺปิ อาวุโส ภาคํ เทถา”ติ. เต เอวมาหํสุ “น มยํ อาวุโส ตุมฺหากํ ภาคํ ทสฺสาม, กิสฺส ตุเมฺห น โอกฺกมิตฺถา”ติ. ภควโต เอตมตฺถํ อาโรเจสุํ. อนุชานามิ ภิกฺขเว อาคเมนฺตานํ อกามา ภาคํ ทาตุนฺติ.}

\proseitem{341}{เตน โข ปน สมเยน สมฺพหุลา ภิกฺขู โกสเลสุ ชนปเท อทฺธาน\-มคฺค\-ปฺปฏิปนฺนา โหนฺติ. เอกจฺเจ ภิกฺขู ปฐมํ สุสานํ โอกฺกมิํสุ ปํสุกูลาย, เอกจฺเจ ภิกฺขู ปจฺฉา โอกฺกมิํสุ. เย เต ภิกฺขู ปฐมํ สุสานํ โอกฺกมิํสุ ปํสุกูลาย, เต ปํสุกูลานิ ลภิํสุ. เย เต \csromanfolio{395}ภิกฺขู ปจฺฉา โอกฺกมิํสุ, เต น ลภิํสุ. เต เอวมาหํสุ “อมฺหากมฺปิ อาวุโส ภาคํ เทถา”ติ. เต เอวมาหํสุ “น มยํ อาวุโส ตุมฺหากํ ภาคํ ทสฺสาม, กิสฺส ตุเมฺห ปจฺฉา โอกฺกมิตฺถา”ติ. ภควโต เอตมตฺถํ อาโรเจสุํ. อนุชานามิ ภิกฺขเว ปจฺฉา โอกฺกนฺตานํ นากามา ภาคํ ทาตุนฺติ.}

\prose{เตน โข ปน สมเยน สมฺพหุลา ภิกฺขู โกสเลสุ ชนปเท อทฺธาน\-มคฺค\-ปฺปฏิปนฺนา โหนฺติ. เต สทิสา สุสานํ โอกฺกมิํสุ ปํสุกูลาย, เอกจฺเจ ภิกฺขู ปํสุกูลานิ ลภิํสุ, เอกจฺเจ ภิกฺขู น ลภิํสุ. เย เต ภิกฺขู น ลภิํสุ, เต เอวมาหํสุ “อมฺหากมฺปิ อาวุโส ภาคํ เทถา”ติ. เต เอวมาหํสุ “น มยํ อาวุโส ตุมฺหากํ ภาคํ ทสฺสาม, กิสฺส ตุเมฺห น ลภิตฺถา”ติ. ภควโต เอตมตฺถํ อาโรเจสุํ. อนุชานามิ ภิกฺขเว สทิสานํ โอกฺกนฺตานํ อกามา ภาคํ ทาตุนฺติ.}

\prose{เตน โข ปน สมเยน สมฺพหุลา ภิกฺขู โกสเลสุ ชนปเท อทฺธาน\-มคฺค\-ปฺปฏิปนฺนา โหนฺติ. เต กติกํ กตฺวา สุสานํ โอกฺกมิํสุ ปํสุกูลาย, เอกจฺเจ ภิกฺขู ปํสุกูลานิ ลภิํสุ, เอกจฺเจ ภิกฺขู น ลภิํสุ. เย เต ภิกฺขู น ลภิํสุ, เต เอวมาหํสุ “อมฺหากมฺปิ อาวุโส ภาคํ เทถา”ติ. เต เอวมาหํสุ “น มยํ อาวุโส ตุมฺหากํ ภาคํ ทสฺสาม, กิสฺส ตุเมฺห น ลภิตฺถา”ติ. ภควโต เอตมตฺถํ อาโรเจสุํ. อนุชานามิ ภิกฺขเว กติกํ กตฺวา โอกฺกนฺตานํ อกามา ภาคํ ทาตุนฺติ.}

\csromanmatikaanchor{mtk.221}
\csromantocmarkref{section}{213. จีวรปฏิคฺคาหกสมฺมุติกถา}{mtk.221}
\bookmark[dest={mtk.221},level=1]{213. จีวรปฏิคฺคาหกสมฺมุติกถา}
\csromanheader{1}{213. \textbf{จีวรปฏิคฺคาหก}\-\textbf{สมฺมุติ}\-\textbf{กถา}}
\proseitem{342}{เตน โข ปน สมเยน มนุสฺสา จีวรํ อาทาย อารามํ อาคจฺฉนฺติ, เต ปฏิคฺคาหกํ อลภมานา ปฏิหรนฺติ, จีวรํ ปริตฺตํ อุปฺปชฺชติ. ภควโต เอตมตฺถํ อาโรเจสุํ. อนุชานามิ ภิกฺขเว ปญฺจหงฺเคหิ สมนฺนาคตํ ภิกฺขุํ จีวรปฏิคฺคาหกํ สมฺมนฺนิตุํ, โย น ฉนฺทาคติํ คจฺเฉยฺย, น โทสาคติํ คจฺเฉยฺย, น โมหาคติํ คจฺเฉยฺย, น ภยาคติํ คจฺเฉยฺย, คหิตาคหิตญฺจ ชาเนยฺย. เอวญฺจ ปน ภิกฺขเว สมฺมนฺนิตพฺโพ, ปฐมํ ภิกฺขุ ยาจิตพฺโพ, ยาจิตฺวา พฺยตฺเตน ภิกฺขุนา ปฏิพเลน สํโฆ ญาเปตพฺโพ\csromandash{}}

\prose{“สุณาตุ เม ภนฺเต สํโฆ, ยทิ สํฆสฺส ปตฺตกลฺลํ, สํโฆ อิตฺถนฺนามํ ภิกฺขุํ จีวรปฏิคฺคาหกํ สมฺมนฺเนยฺย, เอสา ญตฺติ.}

\prose{\csromanfolio{396}สุณาตุ เม ภนฺเต สํโฆ, สํโฆ อิตฺถนฺนามํ ภิกฺขุํ จีวรปฏิคฺคาหกํ สมฺมนฺนติ, ยสฺสายสฺมโต ขมติ อิตฺถนฺนามสฺส ภิกฺขุโน จีวรปฏิคฺคาหกสฺส สมฺมุติ, โส ตุณฺหสฺส, ยสฺส นกฺขมติ, โส ภาเสยฺย.}

\prose{สมฺมโต สํเฆน อิตฺถนฺนาโม ภิกฺขุ จีวรปฏิคฺคาหโก, ขมติ สํฆสฺส, ตสฺมา ตุณฺหี, เอวเมตํ ธารยามี”ติ.}

\proseitem{342}{เตน โข ปน สมเยน จีวรปฏิคฺคาหกา ภิกฺขู จีวรํ ปฏิคฺคเหตฺวา ตตฺเถว อุชฺฌิตฺวา ปกฺกมนฺติ, จีวรํ นสฺสติ. ภควโต เอตมตฺถํ อาโรเจสุํ. อนุชานามิ ภิกฺขเว ปญฺจหงฺเคหิ สมนฺนาคตํ ภิกฺขุํ จีวรนิทหกํ สมฺมนฺนิตุํ, โย น ฉนฺทาคติํ คจฺเฉยฺย, น โทสาคติํ คจฺเฉยฺย, น โมหาคติํ คจฺเฉยฺย, น ภยาคติํ คจฺเฉยฺย, นิหิตา\-นิหิตญฺจ ชาเนยฺย. เอวญฺจ ปน ภิกฺขเว สมฺมนฺนิตพฺโพ, ปฐมํ ภิกฺขุ ยาจิตพฺโพ, ยาจิตฺวา พฺยตฺเตน ภิกฺขุนา ปฏิพเลน สํโฆ ญาเปตพฺโพ\csromandash{}}

\prose{“สุณาตุ เม ภนฺเต สํโฆ, ยทิ สํฆสฺส ปตฺตกลฺลํ, สํโฆ อิตฺถนฺนามํ ภิกฺขุํ จีวรนิทหกํ สมฺมนฺเนยฺย, เอสา ญตฺติ. สุณาตุ เม ภนฺเต สํโฆ, สํโฆ อิตฺถนฺนามํ ภิกฺขุํ จีวรนิทหกํ สมฺมนฺนติ, ยสฺสายสฺมโต ขมติ อิตฺถนฺนามสฺส ภิกฺขุโน จีวรนิทหกสฺส สมฺมุติ, โส ตุณฺหสฺส, ยสฺส นกฺขมติ, โส ภาเสยฺย. สมฺมโต สํเฆน อิตฺถนฺนาโม ภิกฺขุ จีวรนิทหโก, ขมติ สํฆสฺส, ตสฺมา ตุณฺหี, เอวเมตํ ธารยามี”ติ.}

\csromanmatikaanchor{mtk.222}
\csromantocmarkref{section}{214. ภณฺฑาคารสมฺมุติอาทิกถา}{mtk.222}
\bookmark[dest={mtk.222},level=1]{214. ภณฺฑาคารสมฺมุติอาทิกถา}
\csromanheader{1}{214. ภณฺฑาคาร\-สมฺมุติ\-อาทิ\-กถา}
\proseitem{343}{เตน โข ปน สมเยน จีวรนิทหโก ภิกฺขุ มณฺฑเปปิ รุกฺขมูเลปิ นิพฺพโกเสปิ จีวรํ นิทหติ, อุนฺทูเรหิปิ อุปจิกาหิปิ ขชฺชนฺติ. ภควโต เอตมตฺถํ อาโรเจสุํ. อนุชานามิ ภิกฺขเว ภณฺฑาคารํ สมฺมนฺนิตุํ, ยํ สํโฆ อากงฺขติ วิหารํ วา อฑฺฒโยคํ วา ปาสาทํ วา หมฺมิยํ วา คุหํ วา. เอวญฺจ ปน ภิกฺขเว สมฺมนฺนิตพฺโพ, พฺยตฺเตน ภิกฺขุนา ปฏิพเลน สํโฆ ญาเปตพฺโพ\csromandash{}}

\prose{“สุณาตุ เม ภนฺเต สํโฆ, ยทิ สํฆสฺส ปตฺตกลฺลํ, สํโฆ อิตฺถนฺนามํ วิหารํ ภณฺฑาคารํ สมฺมนฺเนยฺย, เอสา ญตฺติ. สุณาตุ เม ภนฺเต สํโฆ, สํโฆ อิตฺถนฺนามํ วิหารํ ภณฺฑาคารํ สมฺมนฺนติ, ยสฺสายสฺมโต ขมติ อิตฺถนฺนามสฺส วิหารสฺส ภณฺฑาคารสฺส สมฺมุติ, โส ตุณฺหสฺส, \csromanfolio{397}ยสฺส นกฺขมติ, โส ภาเสยฺย. สมฺมโต สํเฆน อิตฺถนฺนาโม วิหาโร ภณฺฑาคารํ, ขมติ สํฆสฺส, ตสฺมา ตุณฺหี, เอวเมตํ ธารยามี”ติ.}

\prose{เตน โข ปน สมเยน สํฆสฺส ภณฺฑาคาเร จีวรํ อคุตฺตํ โหติ. ภควโต เอตมตฺถํ อาโรเจสุํ. อนุชานามิ ภิกฺขเว ปญฺจหงฺเคหิ สมนฺนาคตํ ภิกฺขุํ ภณฺฑาคาริกํ สมฺมนฺนิตุํ, โย น ฉนฺทาคติํ คจฺเฉยฺย, น โทสาคติํ คจฺเฉยฺย, น โมหาคติํ คจฺเฉยฺย, น ภยาคติํ คจฺเฉยฺย, คุตฺตาคุตฺตญฺจ ชาเนยฺย. เอวญฺจ ปน ภิกฺขเว สมฺมนฺนิตพฺโพ. ปฐมํ ภิกฺขุ ยาจิตพฺโพ, ยาจิตฺวา พฺยตฺเตน ภิกฺขุนา ปฏิพเลน สํโฆ ญาเปตพฺโพ\csromandash{}}

\prose{“สุณาตุ เม ภนฺเต สํโฆ, ยทิ สํฆสฺส ปตฺตกลฺลํ สํโฆ อิตฺถนฺนามํ ภิกฺขุํ ภณฺฑาคาริกํ สมฺมนฺเนยฺย, เอสา ญตฺติ.}

\prose{. สุณาตุ เม ภนฺเต สํโฆ, สํโฆ อิตฺถนฺนามํ ภิกฺขุํ ภณฺฑาคาริกํ สมฺมนฺนติ, ยสฺสายสฺมโต ขมติ อิตฺถนฺนามสฺส ภิกฺขุโน ภณฺฑา\-คาริ\-กสฺส สมฺมุติ, โส ตุณฺหสฺส, ยสฺส นกฺขมติ, โส ภาเสยฺย.}

\prose{. สมฺมโต สํเฆน อิตฺถนฺนาโม ภิกฺขุ ภณฺฑาคาริโก, ขมติ สํฆสฺส, ตสฺมา ตุณฺหี, เอวเมตํ ธารยามี”ติ.}

\prose{เตน โข ปน สมเยน ฉพฺพคฺคิยา ภิกฺขู ภณฺฑาคาริกํ วุฏฺฐาเปนฺติ. ภควโต เอตมตฺถํ อาโรเจสุํ. น ภิกฺขเว ภณฺฑาคาริโก วุฏฺฐาเปตพฺโพ, โย วุฏฺฐาเปยฺย, อาปตฺติ ทุกฺกฏสฺสาติ.}

\proseitem{343}{เตน โข ปน สมเยน สํฆสฺส ภณฺฑาคาเร จีวรํ อุสฺสนฺนํ โหติ. ภควโต เอตมตฺถํ อาโรเจสุํ. อนุชานามิ ภิกฺขเว สมฺมุขีภูเตน สํเฆน ภาเชตุนฺติ.}

\prose{เตน โข ปน สมเยน สํโฆ จีวรํ ภาเชนฺโต โกลาหลํ อกาสิ. ภควโต เอตมตฺถํ อาโรเจสุํ. อนุชานามิ ภิกฺขเว ปญฺจหงฺเคหิ สมนฺนาคตํ ภิกฺขุํ จีวรภาชกํ สมฺมนฺนิตุํ, โย น ฉนฺทาคติํ คจฺเฉยฺย, น โทสาคติํ คจฺเฉยฺย, น โมหาคติํ คจฺเฉยฺย, น ภยาคติํ คจฺเฉยฺย, ภาชิตาภาชิตญฺจ ชาเนยฺย. เอวญฺจ ปน ภิกฺขเว สมฺมนฺนิตพฺโพ. ปฐมํ ภิกฺขุ ยาจิตพฺโพ, ยาจิตฺวา พฺยตฺเตน ภิกฺขุนา ปฏิพเลน สํโฆ ญาเปตพฺโพ\csromandash{}}

\prose{\csromanfolio{398}“สุณาตุ เม ภนฺเต สํโฆ, ยทิ สํฆสฺส ปตฺตกลฺลํ, สํโฆ อิตฺถนฺนามํ ภิกฺขุํ จีวรภาชกํ สมฺมนฺเนยฺย, เอสา ญตฺติ. สุณาตุ เม ภนฺเต สํโฆ, สํโฆ อิตฺถนฺนามํ ภิกฺขุํ จีวรภาชกํ สมฺมนฺนติ, ยสฺสายสฺมโต ขมติ อิตฺถนฺนามสฺส ภิกฺขุโน จีวรภาชกสฺส สมฺมุติ, โส ตุณฺหสฺส, ยสฺส นกฺขมติ, โส ภาเสยฺย. สมฺมโต สํเฆน อิตฺถนฺนาโม ภิกฺขุ จีวรภาชโก, ขมติ สํฆสฺส, ตสฺมา ตุณฺหี, เอวเมตํ ธารยามี”ติ.}

\prose{อถ โข จีวรภาชกานํ ภิกฺขูนํ เอตทโหสิ “กถํ นุ โข จีวรํ ภาเชตพฺพนฺ”ติ. ภควโต เอตมตฺถํ อาโรเจสุํ. อนุชานามิ ภิกฺขเว ปฐมํ อุจฺจินิตฺวา ตุลยิตฺวา วณฺณาวณฺณํ กตฺวา ภิกฺขู คเณตฺวา วคฺคํ พนฺธิตฺวา จีวร\-ปฏิวีสํ ฐเปตุนฺติ.}

\prose{อถ โข จีวรภาชกานํ ภิกฺขูนํ เอตทโหสิ “กถํ นุ โข สามเณรานํ จีวรปฏิวีโส ทาตพฺโพ”ติ. ภควโต เอตมตฺถํ อาโรเจสุํ. อนุชานามิ ภิกฺขเว สามเณรานํ อุปฑฺฒ\-ปฏิวีสํ ทาตุนฺติ.}

\prose{เตน โข ปน สมเยน อญฺญตโร ภิกฺขุ สเกน ภาเคน อุตฺตริตุกาโม โหติ. ภควโต เอตมตฺถํ อาโรเจสุํ. อนุชานามิ ภิกฺขเว อุตฺตรนฺตสฺส สกํ ภาคํ ทาตุนฺติ.}

\prose{เตน โข ปน สมเยน อญฺญตโร ภิกฺขุ อติเรกภาเคน อุตฺตริตุกาโม โหติ. ภควโต เอตมตฺถํ อาโรเจสุํ. อนุชานามิ ภิกฺขเว อนุกฺเขเป ทินฺเน อติเรกภาคํ ทาตุนฺติ.}

\prose{อถ โข จีวรภาชกานํ ภิกฺขูนํ เอตทโหสิ “กถํ นุ โข จีวรปฏิวีโส ทาตพฺโพ, อาคต\-ปฏิปาฏิยา\csromansharedfootnote{8af8e6247721dc7a}{อาคตาคตปฏิปาฏิยา (ก)} นุ โข อุทาหุ ยถาวุฑฺฒนฺ”ติ. ภควโต เอตมตฺถํ อาโรเจสุํ. อนุชานามิ ภิกฺขเว วิกลเก โตเสตฺวา กุสปาตํ กาตุนฺติ.}

\csromanmatikaanchor{mtk.223}
\csromantocmarkref{section}{215. จีวรรชนกถา}{mtk.223}
\bookmark[dest={mtk.223},level=1]{215. จีวรรชนกถา}
\csromanheader{1}{215. \textbf{จีวร}\-\textbf{รชน}\-\textbf{กถา}}
\proseitem{344}{เตน โข ปน สมเยน ภิกฺขู ฉกเณนปิ ปณฺฑุมตฺติกายปิ จีวรํ รชนฺติ, จีวรํ ทุพฺพณฺณํ โหติ. ภควโต เอตมตฺถํ อาโรเจสุํ. \csromanfolio{399}อนุชานามิ ภิกฺขเว ฉ รชนานิ มูลรชนํ ขนฺธรชนํ ตจรชนํ ปตฺตรชนํ ปุปฺผรชนํ ผลรชนนฺติ.}

\prose{เตน โข ปน สมเยน ภิกฺขู สีตุทกาย\csromansharedfootnote{3c9099ffcc563e22}{สีตุนฺทิกาย (สี), สีตูทกาย (สฺยา)} จีวรํ รชนฺติ, จีวรํ ทุคฺคนฺธํ โหติ. ภควโต เอตมตฺถํ อาโรเจสุํ. อนุชานามิ ภิกฺขเว รชนํ ปจิตุํ จุลฺลํ รชน\-กุมฺภินฺติ. รชนํ อุตฺตริยติ. ภควโต เอตมตฺถํ อาโรเจสุํ. อนุชานามิ ภิกฺขเว อุตฺตราฬุมฺปํ\csromansharedfootnote{559f9d078826691d}{อุตฺตราฬุปํ (โยชนา), อุตฺตราฬุวํ (สฺยา)} พนฺธิตุนฺติ.}

\prose{เตน โข ปน สมเยน ภิกฺขู น ชานนฺติ รชนํ ปกฺกํ วา อปกฺกํ วา. ภควโต เอตมตฺถํ อาโรเจสุํ. อนุชานามิ ภิกฺขเว อุทเก วา นขปิฏฺฐิกาย วา เถวกํ ทาตุนฺติ.}

\prose{เตน โข ปน สมเยน ภิกฺขู รชนํ โอโรเปนฺตา กุมฺภิํ อาวิญฺฉนฺติ\csromansharedfootnote{cc34785f20f384c4}{อาวิญฺชนฺติ (สี), อาวฏฺฏนฺติ (สฺยา)}, กุมฺภี ภิชฺชติ. ภควโต เอตมตฺถํ อาโรเจสุํ. อนุชานามิ ภิกฺขเว รชนุฬุงฺกํ\csromansharedfootnote{ffe051b5e39c50a8}{รชนาฬุงฺกํ (โยชนา)} ทณฺฑ\-กถาลกนฺติ.}

\prose{เตน โข ปน สมเยน ภิกฺขูนํ รชนภาชนํ น สํวิชฺชติ. ภควโต เอตมตฺถํ อาโรเจสุํ. อนุชานามิ ภิกฺขเว รชนโกลมฺพํ รชนฆฏนฺติ.}

\prose{เตน โข ปน สมเยน ภิกฺขู ปาติยาปิ ปตฺเตปิ จีวรํ โอมทฺทนฺติ, จีวรํ ปริภิชฺชติ.}

\setlength{\csromangathapagewidth}{0pt}
\setlength{\csromangathawakwidth}{0pt}
\csromangathameasuregroup{}{ภควโต เอตมตฺถํ อาโรเจสุํ. อนุชานามิ ภิกฺขเว รชนโทณิกนฺติ.}
\csromangathameasurewak{ภควโต เอตมตฺถํ อาโรเจสุํ. อนุชานามิ ภิกฺขเว รชนโทณิกนฺติ.}
\setlength{\csromangathaleftcol}{0pt}
\csromangathagroup{\csromangathastanza{\csromangathawak{\csromangathaitemline{344}{ภควโต เอตมตฺถํ อาโรเจสุํ. อนุชานามิ ภิกฺขเว รชน\-โทณิกนฺติ.}}}}

\prose{เตน โข ปน สมเยน ภิกฺขู ฉมาย จีวรํ ปตฺถรนฺติ, จีวรํ ปํสุกิตํ โหติ.}

\setlength{\csromangathapagewidth}{0pt}
\setlength{\csromangathawakwidth}{0pt}
\csromangathameasuregroup{}{ภควโต เอตมตฺถํ อาโรเจสุํ. อนุชานามิ ภิกฺขเว ติณสนฺถารกนฺติ.}
\csromangathameasurewak{ภควโต เอตมตฺถํ อาโรเจสุํ. อนุชานามิ ภิกฺขเว ติณสนฺถารกนฺติ.}
\setlength{\csromangathaleftcol}{0pt}
\csromangathagroup{\csromangathastanza{\csromangathawak{ภควโต เอตมตฺถํ อาโรเจสุํ. อนุชานามิ ภิกฺขเว ติณสนฺถารกนฺติ.}}}

\proseitem{344}{ติณสนฺถารโก อุปจิกาหิ ขชฺชติ. ภควโต เอตมตฺถํ อาโรเจสุํ. อนุชานามิ ภิกฺขเว จีวรวํสํ จีวรรชฺชุนฺติ.}

\prose{มชฺเฌน ลคฺเคนฺติ, รชนํ อุภโต คลติ. ภควโต เอตมตฺถํ อาโรเจสุํ. อนุชานามิ ภิกฺขเว กณฺเณ พนฺธิตุนฺติ.}

\prose{\csromanfolio{400}กณฺโณ ชีรติ. ภควโต เอตมตฺถํ อาโรเจสุํ. อนุชานามิ ภิกฺขเว กณฺณ\-สุตฺตกนฺติ.}

\prose{รชนํ เอกโต คลติ. ภควโต เอตมตฺถํ อาโรเจสุํ. อนุชานามิ ภิกฺขเว สมฺปริวตฺตกํ สมฺปริวตฺตกํ รเชตุํ, น จ อจฺฉินฺเน เถเว ปกฺกมิตุนฺติ.}

\prose{เตน โข ปน สมเยน จีวรํ ปตฺถินฺนํ โหติ. ภควโต เอตมตฺถํ อาโรเจสุํ. อนุชานามิ ภิกฺขเว อุทเก โอสาเรตุนฺติ.}

\prose{เตน โข ปน สมเยน จีวรํ ผรุสํ โหติ. ภควโต เอตมตฺถํ อาโรเจสุํ. อนุชานามิ ภิกฺขเว ปาณินา อาโกเฏตุนฺติ.}

\prose{เตน โข ปน สมเยน ภิกฺขู อจฺฉินฺนกานิ จีวรานิ ธาเรนฺติ ทนฺตกาสาวานิ. มนุสฺสา อุชฺฌายนฺติ ขิยฺยนฺติ วิปาเจนฺติ “เสยฺยถาปิ นาม\csromansharedfootnote{fa9322a172102b9f}{เสยฺยถาปิ (?)} คิหี กามโภคิโน”ติ. ภควโต เอตมตฺถํ อาโรเจสุํ. น ภิกฺขเว อจฺฉินฺนกานิ จีวรานิ ธาเรตพฺพานิ, โย ธาเรยฺย, อาปตฺติ ทุกฺกฏสฺสาติ.}

\csromanmatikaanchor{mtk.224}
\csromantocmarkref{section}{216. ฉินฺนกจีวรานุชานนา}{mtk.224}
\bookmark[dest={mtk.224},level=1]{216. ฉินฺนกจีวรานุชานนา}
\csromanheader{1}{216. \textbf{ฉินฺนก}\-\textbf{จีวรา}\-\textbf{นุชานนา}}
\proseitem{345}{อถ โข ภควา ราชคเห ยถาภิรนฺตํ วิหริตฺวา เยน ทกฺขิณาคิริ เตน จาริกํ ปกฺกามิ. อทฺทสา โข ภควา มคธเขตฺตํ อจฺฉิพทฺธํ\csromansharedfootnote{2a4800af603f43ef}{อจฺจิพทฺธํ (สี, สฺยา), อจฺฉิพนฺธํ (ก)} ปาฬิพทฺธํ มริยาท\-พทฺธํ สิงฺฆาฏก\-พทฺธํ, ทิสฺวาน อายสฺมนฺตํ อานนฺทํ อามนฺเตสิ “ปสฺสสิ โน ตฺวํ อานนฺท มคธเขตฺตํ อจฺฉิพทฺธํ ปาฬิพทฺธํ มริยาท\-พทฺธํ สิงฺฆาฏกพทฺธนฺ”ติ. “เอวํ ภนฺเต”ติ. อุสฺสหสิ ตฺวํ อานนฺท ภิกฺขูนํ เอวรูปานิ จีวรานิ สํวิทหิตุนฺติ. อุสฺสหามิ ภควาติ. อถ โข ภควา ทกฺขิณา\-คิริสฺมิํ ยถาภิรนฺตํ วิหริตฺวา ปุนเทว ราชคหํ ปจฺจาคญฺฉิ. อถ โข อายสฺมา อานนฺโท สมฺพหุลานํ ภิกฺขูนํ จีวรานิ สํวิทหิตฺวา เยน ภควา เตนุปสงฺกมิ, อุปสงฺกมิตฺวา ภควนฺตํ เอตทโวจ “ปสฺสตุ เม\csromansharedfootnote{dd34feab0b33cc97}{ปสฺสถ ตุเมฺห (ก)} ภนฺเต ภควา จีวรานิ สํวิทหิตานี”ติ. อถ โข ภควา เอตสฺมิํ นิทาเน เอตสฺมิํ ปกรเณ ธมฺมิํ กถํ กตฺวา ภิกฺขู อามนฺเตสิ “ปณฺฑิโต ภิกฺขเว อานนฺโท, มหาปญฺโญ ภิกฺขเว อานนฺโท, ยตฺร หิ นาม มยา สํขิตฺเตน ภาสิตสฺส วิตฺถาเรน อตฺถํ อาชานิสฺสติ, กุสิมฺปิ นาม กริสฺสติ, อฑฺฒกุสิมฺปิ นาม กริสฺสติ, มณฺฑลมฺปิ นาม \csromanfolio{401}กริสฺสติ, อฑฺฒมณฺฑลมฺปิ นาม กริสฺสติ, วิวฏฺฏมฺปิ นาม กริสฺสติ, อนุวิวฏฺฏมฺปิ นาม กริสฺสติ, คีเวยฺยกมฺปิ นาม กริสฺสติ, ชงฺเฆยฺยกมฺปิ นาม กริสฺสติ, พาหนฺตมฺปิ นาม กริสฺสติ, ฉินฺนกํ ภวิสฺสติ สตฺถลูขํ สมณสารุปฺปํ ปจฺจตฺถิกานํ จ อนภิจฺฉิตํ. อนุชานามิ ภิกฺขเว ฉินฺนกํ สํฆาฏิํ ฉินฺนกํ อุตฺตราสงฺคํ ฉินฺนกํ อนฺตรวาสกนฺ”ติ.}

\csromanmatikaanchor{mtk.225}
\csromantocmarkref{section}{217. ติจีวรานุชานนา}{mtk.225}
\bookmark[dest={mtk.225},level=1]{217. ติจีวรานุชานนา}
\csromanheader{1}{217. \textbf{ติจีวรา}\-\textbf{นุชานนา}}
\proseitem{346}{อถ โข ภควา ราชคเห ยถาภิรนฺตํ วิหริตฺวา เยน เวสาลี เตน จาริกํ ปกฺกามิ. อทฺทส ภควา อนฺตรา จ ราชคหํ อนฺตรา จ เวสาลิํ อทฺธาน\-มคฺค\-ปฺปฏิปนฺโน สมฺพหุเล ภิกฺขู จีวเรหิ อุพฺภณฺฑิเต\csromansharedfootnote{7ed01e6d9c900a90}{อุพฺภณฺฑีกเต (สฺยา)} สีเสปิ จีวรภิสิํ กริตฺวา ขนฺเธปิ จีวรภิสิํ กริตฺวา กฏิยาปิ จีวรภิสิํ กริตฺวา อาคจฺฉนฺเต, ทิสฺวาน ภควโต เอตทโหสิ “อติลหุํ โข อิเม โมฆปุริสา จีวเร พาหุลฺลาย อาวตฺตา, ยนฺนูนาหํ ภิกฺขูนํ จีวเร สีมํ พนฺเธยฺยํ, มริยาทํ ฐเปยฺยนฺ”ติ. อถ โข ภควา อนุปุพฺเพน จาริกํ จรมาโน เยน เวสาลี ตทวสริ, ตตฺร สุทํ ภควา เวสาลิยํ วิหรติ โคตมเก เจติเย. เตน โข ปน สมเยน ภควา สีตาสุ เหมนฺติกาสุ รตฺตีสุ อนฺตรฏฺฐกาสุ หิมปาตสมเย รตฺติํ อชฺโฌกาเส เอกจีวโร นิสีทิ, น ภควนฺตํ สีตํ อโหสิ. นิกฺขนฺเต ปฐเม ยาเม สีตํ ภควนฺตํ อโหสิ, ทุติยํ ภควา จีวรํ ปารุปิ, น ภควนฺตํ สีตํ อโหสิ. นิกฺขนฺเต มชฺฌิเม ยาเม สีตํ ภควนฺตํ อโหสิ, ตติยํ ภควา จีวรํ ปารุปิ, น ภควนฺตํ สีตํ อโหสิ. นิกฺขนฺเต ปจฺฉิเม ยาเม อุทฺธสฺเต อรุเณ นนฺทิมุขิยา รตฺติยา สีตํ ภควนฺตํ อโหสิ, จตุตฺถํ ภควา จีวรํ ปารุปิ, น ภควนฺตํ สีตํ อโหสิ. อถ โข ภควโต เอตทโหสิ “เยปิ โข เต กุลปุตฺตา อิมสฺมิํ ธมฺมวินเย สีตาลุกา สีตภีรุกา, เตปิ สกฺโกนฺติ ติจีวเรน ยาเปตุํ, ยนฺนูนาหํ ภิกฺขูนํ จีวเร สีมํ พนฺเธยฺยํ มริยาทํ ฐเปยฺยํ ติจีวรํ อนุชาเนยฺยนฺ”ติ. อถ โข ภควา เอตสฺมิํ นิทาเน เอตสฺมิํ ปกรเณ ธมฺมิํ กถํ กตฺวา ภิกฺขู อามนฺเตสิ “อิธาหํ ภิกฺขเว อนฺตรา จ ราชคหํ อนฺตรา จ เวสาลิํ อทฺธาน\-มคฺค\-ปฺปฏิปนฺโน อทฺทสํ สมฺพหุลา \csromanfolio{402}ภิกฺขู จีวเรหิ อุพฺภณฺฑิเต สีเสปิ จีวรภิสิํ กริตฺวา ขนฺเธปิ จีวรภิสิํ กริตฺวา กฏิยาปิ จีวรภิสิํ กริตฺวา อาคจฺฉนฺเต, ทิสฺวาน เม เอตทโหสิ ‘อติลหุํ โข อิเม โมฆปุริสา จีวเร พาหุลฺลาย อาวตฺตา, ยนฺนูนาหํ ภิกฺขูนํ จีวเร สีมํ พนฺเธยฺยํ มริยาทํ ฐเปยฺยนฺ’ติ, อิธาหํ ภิกฺขเว สีตาสุ เหมนฺติกาสุ รตฺตีสุ อนฺตรฏฺฐกาสุ หิมปาตสมเย รตฺติํ อชฺโฌกาเส เอกจีวโร นิสีทิํ, น มํ สีตํ อโหสิ. นิกฺขนฺเต ปฐเม ยาเม สีตํ มํ อโหสิ, ทุติยาหํ จีวรํ ปารุปิํ, น มํ สีตํ อโหสิ. นิกฺขนฺเต มชฺฌิเม ยาเม สีตํ มํ อโหสิ, ตติยาหํ จีวรํ ปารุปิํ, น มํ สีตํ อโหสิ. นิกฺขนฺเต ปจฺฉิเม ยาเม อุทฺธสฺเต อรุเณ นนฺทิมุขิยา รตฺติยา สีตํ มํ อโหสิ, จตุตฺถาหํ จีวรํ ปารุปิํ, น มํ สีตํ อโหสิ. ตสฺส มยฺหํ ภิกฺขเว เอตทโหสิ ‘เยปิ โข เต กุลปุตฺตา อิมสฺมิํ ธมฺมวินเย สีตาลุกา สีตภีรุกา, เตปิ สกฺโกนฺติ ติจีวเรน ยาเปตุํ, ยนฺนูนาหํ ภิกฺขูนํ จีวเร สีมํ พนฺเธยฺยํ มริยาทํ ฐเปยฺยํ ติจีวรํ อนุชาเนยฺยนฺ’ติ. อนุชานามิ ภิกฺขเว ติจีวรํ ทิคุณํ สํฆาฏิํ เอกจฺจิยํ อุตฺตราสงฺคํ เอกจฺจิยํ อนฺตรวาสกนฺ”ติ.}

\csromanmatikaanchor{mtk.226}
\csromantocmarkref{section}{218. อติเรกจีวรกถา}{mtk.226}
\bookmark[dest={mtk.226},level=1]{218. อติเรกจีวรกถา}
\csromanheader{1}{218. \textbf{อติเรก}\-\textbf{จีวร}\-\textbf{กถา}}
\proseitem{347}{\csromansymbolfootnote{*}{วิ 1. 294 ปิฏฺเฐปิ.}เตน โข ปน สมเยน ฉพฺพคฺคิยา ภิกฺขู “ภควตา ติจีวรํ อนุญฺญาตนฺ”ติ อญฺเญเนว ติจีวเรน คามํ ปวิสนฺติ, อญฺเญน ติจีวเรน อาราเม อจฺฉนฺติ, อญฺเญน ติจีวเรน นหานํ โอตรนฺติ. เย เต ภิกฺขู อปฺปิจฺฉา, เต อุชฺฌายนฺติ ขิยฺยนฺติ วิปาเจนฺติ “กถํ หิ นาม ฉพฺพคฺคิยา ภิกฺขู อติเรกจีวรํ ธาเรสฺสนฺตี”ติ. อถ โข เต ภิกฺขู ภควโต เอตมตฺถํ อาโรเจสุํ. อถ โข ภควา เอตสฺมิํ นิทาเน เอตสฺมิํ ปกรเณ ธมฺมิํ กถํ กตฺวา ภิกฺขู อามนฺเตสิ “น ภิกฺขเว อติเรกจีวรํ ธาเรตพฺพํ, โย ธาเรยฺย, ยถาธมฺโม กาเรตพฺโพ”ติ.}

\prose{\csromansymbolmark{*}เตน โข ปน สมเยน อายสฺมโต อานนฺทสฺส อติเรกจีวรํ อุปฺปนฺนํ โหติ, อายสฺมา จ อานนฺโท ตํ จีวรํ อายสฺมโต สาริปุตฺตสฺส ทาตุกาโม โหติ, อายสฺมา จ สาริปุตฺโต สาเกเต วิหรติ. อถ โข อายสฺมโต อานนฺทสฺส เอตทโหสิ “ภควตา สิกฺขาปทํ ปญฺญตฺตํ ‘น อติเรกจีวรํ ธาเรตพฺพนฺ’ติ, อิทญฺจ เม อติเรกจีวรํ \csromanfolio{403}อุปฺปนฺนํ, อหญฺจิมํ จีวรํ อายสฺมโต สาริปุตฺตสฺส ทาตุกาโม, อายสฺมา จ สาริปุตฺโต สาเกเต วิหรติ, กถํ นุ โข มยา ปฏิปชฺชิตพฺพนฺ”ติ. ภควโต เอตมตฺถํ อาโรเจสิ. กีวจิรํ ปนานนฺท สาริปุตฺโต อาคจฺฉิสฺสตีติ. นวมํ วา ภควา ทิวสํ ทสมํ วาติ. อถ โข ภควา เอตสฺมิํ นิทาเน เอตสฺมิํ ปกรเณ ธมฺมิํ กถํ กตฺวา ภิกฺขู อามนฺเตสิ “อนุชานามิ ภิกฺขเว ทสาหปรมํ อติเรกจีวรํ ธาเรตุนฺ”ติ.}

\prose{เตน โข ปน สมเยน ภิกฺขูนํ อติเรกจีวรํ อุปฺปนฺนํ โหติ. อถ โข ภิกฺขูนํ เอตทโหสิ “กถํ นุ โข อเมฺหหิ อติเรกจีวเร ปฏิปชฺชิตพฺพนฺ”ติ. ภควโต เอตมตฺถํ อาโรเจสุํ. อนุชานามิ ภิกฺขเว อติเรกจีวรํ วิกปฺเปตุนฺติ.}

\proseitem{348}{อถ โข ภควา เวสาลิยํ ยถาภิรนฺตํ วิหริตฺวา เยน พาราณสี เตน จาริกํ ปกฺกามิ. อนุปุพฺเพน จาริกํ จรมาโน เยน พาราณสี ตทวสริ, ตตฺร สุทํ ภควา พาราณสิยํ วิหรติ อิสิปตเน มิคทาเย. เตน โข ปน สมเยน อญฺญตรสฺส ภิกฺขุโน อนฺตรวาสโก ฉิทฺโท โหติ. อถโข\csromansharedfootnote{b8dfa686fd155ee3}{อถ โข เทฺว นิปาตา \csromandash{} ม.พ.ป.} ตสฺส ภิกฺขุโน เอตทโหสิ “ภควตา ติจีวรํ อนุญฺญาตํ ทิคุณา สํฆาฏิ เอกจฺจิโย อุตฺตราสงฺโค เอกจฺจิโย อนฺตรวาสโก, อยญฺจ เม อนฺตรวาสโก ฉิทฺโท, ยนฺนูนาหํ อคฺคฬํ อจฺฉุเปยฺยํ, สมนฺตโต ทุปฏฺฏํ ภวิสฺสติ, มชฺเฌ เอกจฺจิยนฺ”ติ. อถ โข โส ภิกฺขุ อคฺคฬํ อจฺฉุเปสิ. อทฺทสา โข ภควา เสนาสนจาริกํ อาหิณฺฑนฺโต ตํ ภิกฺขุํ อคฺคฬํ อจฺฉุเปนฺตํ\csromansharedfootnote{0397fb00a3192b76}{อจฺฉุปนฺตํ (ก)}, ทิสฺวาน เยน โส ภิกฺขุ เตนุปสงฺกมิ, อุปสงฺกมิตฺวา ตํ ภิกฺขุํ เอตทโวจ “กิํ ตฺวํ ภิกฺขุ กโรสี”ติ. อคฺคฬํ ภควา อจฺฉุเปมีติ. สาธุ สาธุ ภิกฺขุ, สาธุ โข ตฺวํ ภิกฺขุ อคฺคฬํ อจฺฉุเปสีติ. อถ โข ภควา เอตสฺมิํ นิทาเน เอตสฺมิํ ปกรเณ ธมฺมิํ กถํ กตฺวา ภิกฺขู อามนฺเตสิ “อนุชานามิ ภิกฺขเว อหตานํ ทุสฺสานํ อหตกปฺปานํ ทิคุณํ สํฆาฏิํ เอกจฺจิยํ อุตฺตราสงฺคํ เอกจฺจิยํ อนฺตรวาสกํ, อุตุทฺธฏานํ ทุสฺสานํ จตุคฺคุณํ สํฆาฏิํ ทิคุณํ อุตฺตราสงฺคํ \csromanfolio{404}ทิคุณํ อนฺตรวาสกํ, ปํสุกูเล ยาวทตฺถํ, ปาปณิเก อุสฺสาโห กรณีโย. อนุชานามิ ภิกฺขเว อคฺคฬํ ตุนฺนํ โอวฏฺฏิกํ กณฺฑุสกํ ทฬฺหีกมฺมนฺ”ติ.}

\csromanmatikaanchor{mtk.227}
\csromantocmarkref{section}{219. วิสาขาวตฺถุ}{mtk.227}
\bookmark[dest={mtk.227},level=1]{219. วิสาขาวตฺถุ}
\csromanheader{1}{219. \textbf{วิสาขาวตฺถุ}}
\proseitem{349}{อถ โข ภควา พาราณสิยํ ยถาภิรนฺตํ วิหริตฺวา เยน สาวตฺถิ เตน จาริกํ ปกฺกามิ. อนุปุพฺเพน จาริกํ จรมาโน เยน สาวตฺถิ ตทวสริ, ตตฺร สุทํ ภควา สาวตฺถิยํ วิหรติ เชตวเน อนาถ\-ปิณฺฑิกสฺส อาราเม. อถ โข วิสาขา มิคารมาตา เยน ภควา เตนุปสงฺกมิ, อุปสงฺกมิตฺวา ภควนฺตํ อภิวาเทตฺวา เอกมนฺตํ นิสีทิ. เอกมนฺตํ นิสินฺนํ โข วิสาขํ มิคารมาตรํ ภควา ธมฺมิยา กถาย สนฺทสฺเสสิ สมาทเปสิ สมุตฺเตเชสิ สมฺปหํเสสิ. อถ โข วิสาขา มิคารมาตา ภควตา ธมฺมิยา กถาย สนฺทสฺสิตา สมาทปิตา สมุตฺเตชิตา สมฺปหํสิตา ภควนฺตํ เอตทโวจ “อธิวาเสตุ เม ภนฺเต ภควา สฺวาตนาย ภตฺตํ สทฺธิํ ภิกฺขุสํเฆนา”ติ. อธิวาเสสิ ภควา ตุณฺหีภาเวน. อถ โข วิสาขา มิคารมาตา ภควโต อธิวาสนํ วิทิตฺวา อุฏฺฐายาสนา ภควนฺตํ อภิวาเทตฺวา ปทกฺขิณํ กตฺวา ปกฺกามิ.}

\prose{เตน โข ปน สมเยน ตสฺสา รตฺติยา อจฺจเยน จาตุทฺทีปิโก มหาเมโฆ ปาวสฺสิ. อถ โข ภควา ภิกฺขู อามนฺเตสิ “ยถา ภิกฺขเว เชตเวเน วสฺสติ, เอวํ จตูสุ ทีเปสุ วสฺสติ, โอวสฺสาเปถ ภิกฺขเว กายํ, อยํ ปจฺฉิมโก จาตุทฺทีปิโก มหาเมโฆ”ติ. “เอวํ ภนฺเต”ติ โข เต ภิกฺขู ภควโต ปฏิสฺสุณิตฺวา นิกฺขิตฺต\-จีวรา กายํ โอวสฺสาเปนฺติ. อถ โข วิสาขา มิคารมาตา ปณีตํ ขาทนียํ โภชนียํ ปฏิยาทาเปตฺวา ทาสิํ อาณาเปสิ “คจฺฉ เช อารามํ คนฺตฺวา กาลํ อาโรเจหิ กาโล ภนฺเต นิฏฺฐิตํ ภตฺตนฺ”ติ. “เอวํ อยฺเย”ติ โข สา ทาสี วิสาขาย มิคารมาตุยา ปฏิสฺสุณิตฺวา อารามํ คนฺตฺวา อทฺทส ภิกฺขู นิกฺขิตฺต\-จีวเร กายํ โอวสฺสาเปนฺเต, ทิสฺวาน “นตฺถิ อาราเม ภิกฺขู, อาชีวกา กายํ โอวสฺสาเปนฺตี”ติ เยน วิสาขา มิคารมาตา เตนุปสงฺกมิ, อุปสงฺกมิตฺวา วิสาขํ มิคารมาตรํ เอตทโวจ “นตฺถยฺเย อาราเม ภิกฺขู, อาชีวกา กายํ โอวสฺสาเปนฺตี”ติ. อถ โข วิสาขาย มิคารมาตุยา ปณฺฑิตาย วิยตฺตาย เมธาวินิยา เอตทโหสิ “นิสฺสํสยํ โข อยฺยา นิกฺขิตฺต\-จีวรา กายํ โอวสฺสาเปนฺติ, \csromanfolio{405}สายํ พาลา มญฺญิตฺถ นตฺถิ อาราเม ภิกฺขู, อาชีวกา กายํ โอวสฺสาเปนฺตี”ติ ปุน ทาสิํ อาณาเปสิ “คจฺฉ เช อารามํ คนฺตฺวา กาลํ อาโรเจหิ กาโล ภนฺเต นิฏฺฐิตํ ภตฺตนฺ”ติ, อถ โข เต ภิกฺขู คตฺตานิ สีติํ กริตฺวา\csromansharedfootnote{e237581161838c32}{สีตีกริตฺวา (สฺยา)} กลฺลกายา จีวรานิ คเหตฺวา ยถาวิหารํ ปวิสิํสุ. อถ โข สา ทาสี อารามํ คนฺตฺวา ภิกฺขู อปสฺสนฺตี “นตฺถิ อาราเม ภิกฺขู, สุญฺโญ อาราโม”ติ เยน วิสาขา มิคารมาตา เตนุปสงฺกมิ, อุปสงฺกมิตฺวา วิสาขํ มิคารมาตรํ เอตทโวจ “นตฺถยฺเย อาราเม ภิกฺขู, สุญฺโญ อาราโม”ติ. อถ โข วิสาขาย มิคารมาตุยา ปณฺฑิตาย วิยตฺตาย เมธาวินิยา เอตทโหสิ “นิสฺสํสยํ โข อยฺยา คตฺตานิ สีติํ กริตฺวา กลฺลกายา จีวรานิ คเหตฺวา ยถาวิหารํ ปวิฏฺฐา, สายํ พาลา มญฺญิตฺถ นตฺถิ อาราเม ภิกฺขู, สุญฺโญ อาราโม”ติ ปุน ทาสิํ อาณาเปสิ “คจฺฉ เช อารามํ คนฺตฺวา กาลํ อาโรเจหิ กาโล ภนฺเต นิฏฺฐิตํ ภตฺตนฺ”ติ.}

\proseitem{350}{อถ โข ภควา ภิกฺขู อามนฺเตสิ “สนฺทหถ\csromansharedfootnote{50051156954bf79f}{สนฺนหถ (สี, สฺยา)} ภิกฺขเว ปตฺตจีวรํ, กาโล ภตฺตสฺสา”ติ. “เอวํ ภนฺเต”ติ โข เต ภิกฺขู ภควโต ปจฺจสฺโสสุํ. อถ โข ภควา ปุพฺพณฺหสมยํ นิวาเสตฺวา ปตฺต\-จีวรมา\-ทาย เสยฺยถาปิ นาม พลวา ปุริโส สมิญฺชิตํ วา พาหํ ปสาเรยฺย, ปสาริตํ วา พาหํ สมิญฺเชยฺย, เอวเมว เชตวเน อนฺตรหิโต วิสาขาย มิคารมาตุยา โกฏฺฐเก ปาตุรโหสิ. นิสีทิ ภควา ปญฺญตฺเต อาสเน สทฺธิํ ภิกฺขุ\-สํเฆน. อถ โข วิสาขา มิคารมาตา “อจฺฉริยํ วต โภ อพฺภุตํ วต โภ ตถาคตสฺส มหิทฺธิกตา มหานุภาวตา, ยตฺร หิ นาม ชณฺณุก\-มตฺเตสุปิ โอเฆสุ ปวตฺตมาเนสุ กฏิมตฺเตสุปิ โอเฆสุ ปวตฺตมาเนสุ น หิ นาม เอกภิกฺขุสฺสปิ\csromansharedfootnote{d6fb6bba8453dce9}{ปวตฺตมาเนสุ น เอกภิกฺขุสฺสปิ (?)} ปาทา วา จีวรานิ วา อลฺลานิ ภวิสฺสนฺตี”ติ หฏฺฐา อุทคฺคา พุทฺธ\-ปฺปมุขํ ภิกฺขุสํฆํ ปณีเตน ขาทนีเยน โภชนีเยน สหตฺถา สนฺตปฺเปตฺวา สมฺปวาเรตฺวา ภควนฺตํ ภุตฺตาวิํ โอนีต\-ปตฺต\-ปาณิํ เอกมนฺตํ นิสีทิ, เอกมนฺตํ นิสินฺนา โข วิสาขา มิคารมาตา ภควนฺตํ เอตทโวจ “อฏฺฐาหํ ภนฺเต ภควนฺตํ วรานิ ยาจามี”ติ. อติกฺกนฺตวรา โข วิสาเข ตถาคตาติ. ยานิ จ ภนฺเต \csromanfolio{406}กปฺปิยานิ ยานิ จ อนวชฺชานีติ. วเทหิ วิสาเขติ. อิจฺฉามหํ ภนฺเต สํฆสฺส ยาวชีวํ วสฺสิกสาฏิกํ ทาตุํ, อาคนฺตุก\-ภตฺตํ ทาตุํ, คมิกภตฺตํ ทาตุํ, คิลานภตฺตํ ทาตุํ, คิลานุปฏฺฐาก\-ภตฺตํ ทาตุํ, คิลาน\-เภสชฺชํ ทาตุํ, ธุวยาคุํ ทาตุํ, ภิกฺขุนิ\-สํฆสฺส อุทกสาฏิกํ ทาตุนฺติ. กิํ ปน ตฺวํ วิสาเข อตฺถวสํ สมฺปสฺสมานา ตถาคตํ อฏฺฐ วรานิ ยาจสีติ.}

\prose{อิธาหํ ภนฺเต ทาสิํ อาณาเปสิํ “คจฺฉ เช อารามํ คนฺตฺวา กาลํ อาโรเจหิ กาโล ภนฺเต นิฏฺฐิตํ ภตฺตนฺ”ติ, อถ โข สา ภนฺเต ทาสี อารามํ คนฺตฺวา อทฺทส ภิกฺขู นิกฺขิตฺต\-จีวเร กายํ โอวสฺสาเปนฺเต, ทิสฺวาน “นตฺถิ อาราเม ภิกฺขู, อาชีวกา กายํ โอวสฺสาเปนฺตี”ติ เยนาหํ เตนุปสงฺกมิ, อุปสงฺกมิตฺวา มํ เอตทโวจ “นตฺถยฺเย อาราเม ภิกฺขู, อาชีวกา กายํ โอวสฺสาเปนฺตี”ติ.}

\prose{อสุจิ ภนฺเต นคฺคิยํ เชคุจฺฉํ ปฏิกูลํ, อิมาหํ ภนฺเต อตฺถวสํ สมฺปสฺสมานา อิจฺฉามิ สํฆสฺส ยาวชีวํ วสฺสิกสาฏิกํ ทาตุํ.}

\prose{ปุน จ ปรํ ภนฺเต อาคนฺตุโก ภิกฺขุ น วีถิกุสโล น โคจรกุสโล กิลนฺโต ปิณฺฑาย จรติ. โส เม อาคนฺตุก\-ภตฺตํ ภุญฺชิตฺวา วีถิกุสโล โคจรกุสโล อกิลนฺโต ปิณฺฑาย จริสฺสติ, อิมาหํ ภนฺเต อตฺถวสํ สมฺปสฺสมานา อิจฺฉามิ สํฆสฺส ยาวชีวํ อาคนฺตุก\-ภตฺตํ ทาตุํ.}

\prose{ปุน จปรํ ภนฺเต คมิโก ภิกฺขุ อตฺตโน ภตฺตํ ปริเยสมาโน สตฺถา วา วิหายิสฺสติ, ยตฺต วา วาสํ คนฺตุกาโม ภวิสฺสติ ตตฺถ วิกาเล อุปคจฺฉิสฺสติ, กิลนฺโต อทฺธานํ คมิสฺสติ.}

\prose{โส เม คมิกภตฺตํ ภุญฺชิตฺวา สตฺถา น วิหายิสฺสติ, ยตฺถ วาสํ คนฺตุกาโม ภวิสฺสติ ตตฺถ กาเล อุปคจฺฉิสฺสติ, อกิลนฺโต อทฺธานํ คมิสฺสติ, อิมาหํ ภนฺเต อตฺถวสํ สมฺปสฺสมานา อิจฺฉามิ สํฆสฺส ยาวชีวํ คมิกภตฺตํ ทาตุํ.}

\prose{ปุน จปรํ ภนฺเต คิลานสฺส ภิกฺขุโน สปฺปายานิ โภชนานิ อลภนฺตสฺส อาพาโธ วา อภิวฑฺฒิสฺสติ, กาลํกิริยา วา ภวิสฺสติ.}

\prose{ตสฺส เม คิลานภตฺตํ ภุตฺตสฺส อาพาโธ น อภิวฑฺฒิสฺสติ, กาลํกิริยา น ภวิสฺสติ, อิมาหํ ภนฺเต อตฺถวสํ สมฺปสฺสมานา อิจฺฉามิ สํฆสฺส ยาวชีวํ คิลานภตฺตํ ทาตุํ.}

\prose{\csromanfolio{407}ปุน จปรํ ภนฺเต คิลานุปฏฺฐาโก ภิกฺขุ อตฺตโน ภตฺตํ ปริเยสมาโน คิลานสฺส อุสฺสูเร ภตฺตํ นีหริสฺสติ, ภตฺตจฺเฉทํ กริสฺสติ.}

\prose{โส เม คิลานุปฏฺฐาก\-ภตฺตํ ภุญฺชิตฺวา คิลานสฺส กาเลน ภตฺตํ นีหริสฺสติ, ภตฺตจฺเฉทํ น กริสฺสติ, อิมาหํ ภนฺเต อตฺถวสํ สมฺปสฺสมานา อิจฺฉามิ สํฆสฺส ยาวชีวํ คิลานุปฏฺฐาก\-ภตฺตํ ทาตุํ.}

\prose{ปุน จปรํ ภนฺเต คิลานสฺส ภิกฺขุโน สปฺปายานิ เภสชฺชานิ อลภนฺตสฺส อาพาโธ วา อภิวฑฺฒิสฺสติ, กาลํกิริยา วา ภวิสฺสติ.}

\prose{ตสฺส เม คิลาน\-เภสชฺชํ ปริภุตฺตสฺส อาพาโธ น อภิวฑฺฒิสฺสติ, กาลํกิริยา น ภวิสฺสติ, อิมาหํ ภนฺเต อตฺถวสํ สมฺปสฺสมานา อิจฺฉามิ สํฆสฺส ยาวชีวํ คิลาน\-เภสชฺชํ ทาตุํ.}

\prose{ปุน จปรํ ภนฺเต ภควตา อนฺธกวินฺเท ทสานิสํเส สมฺปสฺสมาเนน ยาคุ อนุญฺญาตา, ตฺยาหํ ภนฺเต อานิสํเส สมฺปสฺสมานา อิจฺฉามิ สํฆสฺส ยาวชีวํ ธุวยาคุํ ทาตุํ.}

\prose{อิธ ภนฺเต ภิกฺขุนิโย อจิรวติยา นทิยา เวสิยาหิ สทฺธิํ นคฺคา เอกติตฺเถ นหายนฺติ, ตา ภนฺเต เวสิยา ภิกฺขุนิโย อุปฺปณฺเฑสุํ “กิํ นุ โข นาม ตุมฺหากํ อยฺเย ทหรานํ\csromansharedfootnote{a409ef3938a90ec1}{ทหรานํ ทหรานํ (สี)} พฺรหฺมจริยํ จิณฺเณน, นนุ นาม กามา ปริภุญฺชิ\-ตพฺพา, ยทา ชิณฺณา ภวิสฺสถ, ตทา พฺรหฺมจริยํ จริสฺสถ, เอวํ ตุมฺหากํ อุโภ อตฺถา ปริคฺคหิตา ภวิสฺสนฺตี”ติ.}

\prose{ตา ภนฺเต ภิกฺขุนิโย เวสิยาหิ อุปฺปณฺฑิยมานา มงฺกู อเหสุํ. อสุจิ ภนฺเต มาตุคามสฺส นคฺคิยํ เชคุจฺฉํ ปฏิกูลํ, อิมาหํ ภนฺเต อตฺถวสํ สมฺปสฺสมานา อิจฺฉามิ ภิกฺขุนิ\-สํฆสฺส ยาวชีวํ อุทกสาฏิกํ ทาตุนฺติ.}

\proseitem{351}{กิํ ปน ตฺวํ วิสาเข อานิสํสํ สมฺปสฺสมานา ตถาคตํ อฏฺฐ วรานิ ยาจสีติ. อิธ ภนฺเต ทิสาสุ วสฺสํวุฏฺฐา ภิกฺขู สาวตฺถิํ อาคจฺฉิสฺสนฺติ ภควนฺตํ ทสฺสนาย, เต ภควนฺตํ อุปสงฺกมิตฺวา ปุจฺฉิสฺสนฺติ “อิตฺถนฺนาโม ภนฺเต ภิกฺขุ กาลงฺกโต, ตสฺส กา คติ โก อภิสมฺปราโย”ติ. ตํ ภควา พฺยากริสฺสติ โสตาปตฺติผเล วา สกทาคามิ\-ผเล วา อนาคามิผเล วา อรหตฺเต วา. ตฺยาหํ อุปสงฺกมิตฺวา ปุจฺฉิสฺสามิ “อาคตปุพฺพา นุ โข ภนฺเต เตน อยฺเยน สาวตฺถี”ติ. สเจ \csromanfolio{408}เม วกฺขนฺติ “อาคตปุพฺพา เตน ภิกฺขุนา สาวตฺถี”ติ. นิฏฺฐเมตฺถ คจฺฉิสฺสามิ “นิสฺสํสยํ เม ปริภุตฺตํ เตน อยฺเยน วสฺสิกสาฏิกา วา อาคนฺตุก\-ภตฺตํ วา คมิกภตฺตํ วา คิลานาภตฺตํ วา คิลานุปฏฺฐาก\-ภตฺตํ วา คิลาน\-เภสชฺชํ วา ธุวยาคุ วา”ติ. ตสฺสา เม ตทนุสฺสรนฺติยา ปามุชฺชํ ชายิสฺสติ, ปมุทิตาย ปีติ ชายิสฺสติ, ปีติมนาย กาโย ปสฺสมฺภิสฺสติ, ปสฺสทฺธกายา สุขํ เวทิยิสฺสามิ, สุขินิยา จิตฺตํ สมาธิยิสฺสติ, สา เม ภวิสฺสติ อินฺทฺริยภาวนา พลภาวนา โพชฺฌงฺค\-ภาวนา, อิมาหํ ภนฺเต อานิสํสํ สมฺปสฺสมานา ตถาคตํ อฏฺฐ วรานิ ยาจามีติ. สาธุ สาธุ วิสาเข, สาธุ โข ตฺวํ วิสาเข อิมํ อานิสํสํ สมฺปสฺสมานา ตถาคตํ อฏฺฐ วรานิ ยาจสิ, อนุชานามิ เต วิสาเข อฏฺฐ วรานีติ. อถ โข ภควา วิสาขํ มิคารมาตรํ อิมาหิ คาถาหิ อนุโมทิ\csromandash{}}

\setlength{\csromangathapagewidth}{0pt}
\setlength{\csromangathawakwidth}{0pt}
\csromangathameasuregroup{}{“ยา อนฺนปานํ ททติปฺปโมทิตา, \\ สีลูปปนฺนา สุคตสฺส สาวิกา. \\ ททาติ ทานํ อภิภุยฺย มจฺฉรํ, \\ โสวคฺคิกํ โสกนุทํ สุขาวหํ. \\ ทิพฺพํ สา ลภเต อายุํ\textsuperscript{1}, \\ อาคมฺม มคฺคํ วิรชํ อนงฺคณํ. \\ สา ปุญฺญกามา สุขินี อนามยา, \\ สคฺคมฺหิ กายมฺหิ จิรํ ปโมทตี’ติ.}
\csromangathameasurewak{“ยา อนฺนปานํ ททติปฺปโมทิตา, \\ สีลูปปนฺนา สุคตสฺส สาวิกา. \\ ททาติ ทานํ อภิภุยฺย มจฺฉรํ, \\ โสวคฺคิกํ โสกนุทํ สุขาวหํ. \\ ทิพฺพํ สา ลภเต อายุํ\textsuperscript{1}, \\ อาคมฺม มคฺคํ วิรชํ อนงฺคณํ. \\ สา ปุญฺญกามา สุขินี อนามยา, \\ สคฺคมฺหิ กายมฺหิ จิรํ ปโมทตี’ติ.}
\setlength{\csromangathaleftcol}{0pt}
\csromangathagroup{\csromangathastanza{\csromangathawak{“ยา อนฺนปานํ ททติปฺปโมทิตา, \\ สีลูปปนฺนา สุคตสฺส สาวิกา. \\ ททาติ ทานํ อภิภุยฺย มจฺฉรํ, \\ โสวคฺคิกํ โสกนุทํ สุขาวหํ.}}\csromangathastanza{\csromangathawak{ทิพฺพํ สา ลภเต อายุํ\csromansharedfootnote{26c7db1ff930b74c}{ทิพฺพํ พลํ สา ลภเต จ อายุํ (สี, สฺยา)}, \\ อาคมฺม มคฺคํ วิรชํ อนงฺคณํ. \\ สา ปุญฺญกามา สุขินี อนามยา, \\ สคฺคมฺหิ กายมฺหิ จิรํ ปโมทตี’ติ.}}}

\proseitemwithcloser{352}{อถ โข ภควา วิสาขํ มิคารมาตรํ อิมาหิ คาถาหิ อนุโมทิตฺวา อุฏฺฐายาสนา ปกฺกามิ. อถ โข ภควา เอตสฺมิํ นิทาเน เอตสฺมิํ ปกรเณ ธมฺมิํ กถํ กตฺวา ภิกฺขู อามนฺเตสิ “อนุชานามิ ภิกฺขเว วสฺสิกสาฏิกํ อาคนฺตุก\-ภตฺตํ คมิกภตฺตํ คิลานภตฺตํ คิลานุปฏฺฐาก\-ภตฺตํ คิลาน\-เภสชฺชํ ธุวยาคุํ ภิกฺขุนิ\-สํฆสฺส อุทกสาฏิกนฺ”ติ.}{วิสาขา\-ภาณวาโร นิฏฺฐิโต.}
\csromanmatikaanchor{mtk.228}
\csromantocmarkref{section}{220. นิสีทนาทิอนุชานนา}{mtk.228}
\bookmark[dest={mtk.228},level=1]{220. นิสีทนาทิอนุชานนา}
\csromanheader{1}{220. นิสีทนา\-ทิ\-อนุชานนา}
\proseitem{353}{\csromanfolio{409}เตน โข ปน สมเยน ภิกฺขู ปณีตานิ โภชนานิ ภุญฺชิตฺวา มุฏฺฐสฺสตี อสมฺปชานา นิทฺทํ โอกฺกมนฺติ, เตสํ มุฏฺฐสฺสตีนํ อสมฺปชานานํ นิทฺทํ โอกฺกมนฺตานํ สุปินนฺเตน อสุจิ มุจฺจติ, เสนาสนํ อสุจินา มกฺขิยติ.}

\prose{อถ โข ภควา อายสฺมตา อานนฺเทน ปจฺฉาสมเณน เสนาสนจาริกํ อาหิณฺฑนฺโต อทฺทส เสนาสนํ อสุจินา มกฺขิตํ, ทิสฺวาน อายสฺมนฺตํ อานนฺทํ อามนฺเตสิ “กิํ เอตํ อานนฺท เสนาสนํ มกฺขิตนฺ”ติ.}

\prose{เอตรหิ ภนฺเต ภิกฺขู ปณีตานิ โภชนานิ ภุญฺชิตฺวา มุฏฺฐสฺสตี อสมฺปชานา นิทฺทํ โอกฺกมนฺติ, เตสํ มุฏฺฐสฺสตีนํ อสมฺปชานานํ นิทฺทํ โอกฺกมนฺตานํ สุปินนฺเตน อสุจิ มุจฺจติ, ตยิทํ ภควา เสนาสนํ อสุจินา มกฺขิตนฺติ.}

\prose{เอวเมตํ อานนฺท เอวเมตํ อานนฺท, มุจฺจติ หิ อานนฺท มุฏฺฐสฺสตีนํ อสมฺปชานานํ นิทฺทํ โอกฺกมนฺตานํ สุปินนฺเตน อสุจิ.}

\prose{เย เต อานนฺท ภิกฺขู อุปฏฺฐิตสฺสตี สมฺปชานา นิทฺทํ โอกฺกมนฺติ, เตสํ อสุจิ น มุจฺจติ.}

\setlength{\csromangathapagewidth}{0pt}
\setlength{\csromangathawakwidth}{0pt}
\csromangathameasuregroup{}{เยปิ เต อานนฺท ปุถุชฺชนา กาเมสุ วีตราคา, \\ เตสมฺปิ อสุจิ น มุจฺจติ.}
\setlength{\csromangathaleftcol}{0pt}
\csromangathagroup{\csromangathastanza{\csromangathaitemline{353}{เยปิ เต อานนฺท ปุถุชฺชนา กาเมสุ วีตราคา,} \\ \csromangathacontline{เตสมฺปิ อสุจิ น มุจฺจติ.}}}

\prose{อฏฺฐานเมตํ อานนฺท อนวกาโส ยํ อรหโต อสุจิ มุจฺเจยฺยาติ. อถ โข ภควา เอตสฺมิํ นิทาเน เอตสฺมิํ ปกรเณ ธมฺมิํ กถํ กตฺวา ภิกฺขู อามนฺเตสิ\csromandash{}อิธาหํ ภิกฺขเว อานนฺเทน ปจฺฉาสมเณน เสนาสนจาริกํ อาหิณฺฑนฺโต อทฺทสํ เสนาสนํ อสุจินา มกฺขิตํ, ทิสฺวาน อานนฺทํ อามนฺเตสิํ “กิํ เอตํ อานนฺท เสนาสนํ มกฺขิตนฺ”ติ. เอตรหิ ภนฺเต ภิกฺขู ปณีตานิ โภชนานิ ภุญฺชิตฺวา มุฏฺฐสฺสตี อสมฺปชานา นิทฺทํ โอกฺกมนฺติ, เตสํ มุฏฺฐสฺสตีนํ อสมฺปชานานํ นิทฺทํ โอกฺกมนฺตานํ สุปินนฺเตน อสุจิ มุจฺจติ, ตยิทํ ภควา เสนาสนํ อสุจินา มกฺขิตนฺติ. เอวเมตํ อานนฺท เอวเมตํ อานนฺท, มุจฺจติ หิ อานนฺท มุฏฺฐสฺสตีนํ อสมฺปชานานํ นิทฺทํ โอกฺกมนฺตานํ สุปินนฺเตน อสุจิ, เย เต อานนฺท ภิกฺขู อุปฏฺฐิตสฺสตี สมฺปชานา นิทฺทํ โอกฺกมนฺติ, เตสํ อสุจิ น มุจฺจติ. เยปิ เต อานนฺท ปุถุชฺชนา กาเมสุ วีตราคา, เตสมฺปิ อสุจิ น มุจฺจติ. อฏฺฐานเมตํ อานนฺท อนวกาโส ยํ อรหโต อสุจิ มุจฺเจยฺยาติ.}

\prose{ปญฺจิเม ภิกฺขเว อาทีนวา มุฏฺฐสฺสติสฺส อสมฺปชานสฺส นิทฺทํ โอกฺกมโต. ทุกฺขํ สุปติ, ทุกฺขํ ปฏิพุชฺฌติ, ปาปกํ สุปินํ ปสฺสติ, เทวตา น รกฺขนฺติ, อสุจิ มุจฺจติ. อิเม โข ภิกฺขเว ปญฺจ อาทีนวา มุฏฺฐสฺสติสฺส อสมฺปชานสฺส นิทฺทํ โอกฺกมโต.}

\prose{\csromanfolio{410}ปญฺจิเม ภิกฺขเว อานิสํสา อุปฏฺฐิต\-สฺสติสฺส สมฺปชานสฺส นิทฺทํ โอกฺกมโต. สุขํ สุปติ, สุขํ ปฏิพุชฺฌติ, น ปาปกํ สุปินํ ปสฺสติ, เทวตา รกฺขนฺติ, อสุจิ น มุจฺจติ. อิเม โข ภิกฺขเว ปญฺจ อานิสํสา อุปฏฺฐิต\-สฺสติสฺส สมฺปชานสฺส นิทฺทํ โอกฺกมโต.}

\proseitem{353}{อนุชานามิ ภิกฺขเว กายคุตฺติยา จีวรคุตฺติยา เสนาสน\-คุตฺติยา นิสีทนนฺติ.}

\prose{เตน โข ปน สมเยน อติขุทฺทกํ นิสีทนํ น สพฺพํ เสนาสนํ สํโคเปติ. ภควโต เอตมตฺถํ อาโรเจสุํ. อนุชานามิ ภิกฺขเว ยาวมหนฺตํ ปจฺจตฺถรณํ อากงฺขติ, ตาวมหนฺตํ ปจฺจตฺถรณํ กาตุนฺติ.}

\proseitem{354}{เตน โข ปน สมเยน อายสฺมโต อานนฺทสฺส อุปชฺฌายสฺส อายสฺมโต เพลฏฺฐ\-สีสสฺส ถุลฺลกจฺฉา\-พาโธ โหติ, ตสฺส ลสิกาย จีวรานิ กาเย ลคฺคนฺติ, ตานิ ภิกฺขู อุทเกน เตเมตฺวา เตเมตฺวา อปกฑฺฒนฺติ. อทฺทสา โข ภควา เสนาสนจาริกํ อาหิณฺฑนฺโต เต ภิกฺขู ตานิ จีวรานิ อุทเกน เตเมตฺวา เตเมตฺวา อปกฑฺฒนฺเต, ทิสฺวาน เยน เต ภิกฺขู เตนุปสงฺกมิ, อุปสงฺกมิตฺวา เต ภิกฺขู เอตทโวจ “กิํ อิมสฺส ภิกฺขเว ภิกฺขุโน อาพาโธ”ติ. อิมสฺส ภนฺเต อายสฺมโต ถุลฺลกจฺฉา\-พาโธ, ลสิกาย จีวรานิ กาเย ลคฺคนฺติ, ตานิ มยํ อุทเกน เตเมตฺวา เตเมตฺวา อปกฑฺฒามาติ. อถ โข ภควา เอตสฺมิํ นิทาเน เอตสฺมิํ ปกรเณ ธมฺมิํ กถํ กตฺวา ภิกฺขู อามนฺเตสิ “อนุชานามิ ภิกฺขเว ‘ยสฺส กณฺฑุ วา ปิฬกา วา อสฺสาโว วา ถุลฺลกจฺฉุ วา อาพาโธ’ กณฺฑุปฺปฏิจฺฉาทินฺ”ติ.}

\proseitem{355}{อถ โข วิสาขา มิคารมาตา มุขปุญฺฉน\-โจฬํ\csromansharedfootnote{e154fa91ba948628}{มุขปุญฺชนโจฬํ (ก)} อาทาย เยน ภควา เตนุปสงฺกมิ, อุปสงฺกมิตฺวา ภควนฺตํ อภิวาเทตฺวา เอกมนฺตํ นิสีทิ, เอกมนฺตํ นิสินฺนา โข วิสาขา มิคารมาตา ภควนฺตํ เอตทโวจ “ปฏิคฺคณฺหาตุ เม ภนฺเต ภควา มุขปุญฺฉน\-โจฬํ, ยํ มมสฺส ทีฆรตฺตํ หิตาย สุขายา”ติ. ปฏิคฺคเหสิ ภควา มุขปุญฺฉน\-โจฬํ. อถ โข ภควา วิสาขํ มิคารมาตรํ ธมฺมิยา กถาย สนฺทสฺเสสิ สมาทเปสิ สมุตฺเตเชสิ สมฺปหํเสสิ. อถ โข วิสาขา มิคารมาตา ภควตา \csromanfolio{411}ธมฺมิยา กถาย สนฺทสฺสิตา สมาทปิตา สมุตฺเตชิตา สมฺปหํสิตา อุฏฺฐายาสนา ภควนฺตํ อภิวาเทตฺวา ปทกฺขิณํ กตฺวา ปกฺกามิ. อถ โข ภควา เอตสฺมิํ นิทาเน เอตสฺมิํ ปกรเณ ธมฺมิํ กถํ กตฺวา ภิกฺขู อามนฺเตสิ “อนุชานามิ ภิกฺขเว มุขปุญฺฉนโจฬกนฺ”ติ\csromansharedfootnote{cd073845119825d0}{มุขปุญฺฉนโจลนฺติ (สฺยา)}.}

\proseitem{356}{เตน โข ปน สมเยน โรโช มลฺโล อายสฺมโต อานนฺทสฺส สหาโย โหติ, โรชสฺส มลฺลสฺส โขมปิโลติกา อายสฺมโต อานนฺทสฺส หตฺเถ นิกฺขิตฺตา โหติ, อายสฺมโต จ อานนฺทสฺส โขม\-ปิโลติกาย อตฺโถ โหติ. ภควโต เอตมตฺถํ อาโรเจสุํ. อนุชานามิ ภิกฺขเว ปญฺจหงฺเคหิ สมนฺนาคตสฺส วิสฺสาสํ คเหตุํ. สนฺทิฏฺโฐ จ โหติ, สมฺภตฺโต จ, อาลปิโต จ, ชีวติ จ, ชานาติ จ “คหิเต เม อตฺตมโน ภวิสฺสตี”ติ. อนุชานามิ ภิกฺขเว อิเมหิ ปญฺจหงฺเคหิ สมนฺนาคตสฺส วิสฺสาสํ คเหตุนฺติ.}

\proseitem{357}{เตน โข ปน สมเยน ภิกฺขูนํ ปริปุณฺณํ โหติ ติจีวรํ, อตฺโถ จ โหติ ปริสฺสาว\-เนหิปิ ถวิกาหิปิ. ภควโต เอตมตฺถํ อาโรเจสุํ. อนุชานามิ ภิกฺขเว ปริกฺขาร\-โจฬกนฺติ.}

\csromanmatikaanchor{mtk.229}
\csromantocmarkref{section}{221. ปจฺฉิมวิกปฺปนุปคจีวราทิกถา}{mtk.229}
\bookmark[dest={mtk.229},level=1]{221. ปจฺฉิมวิกปฺปนุปคจีวราทิกถา}
\csromanheader{1}{221. \textbf{ปจฺฉิม}\-\textbf{วิกปฺปนุปค}\-\textbf{จีวรา}\-\textbf{ทิ}\-\textbf{กถา}}
\proseitem{358}{อถ โข ภิกฺขูนํ เอตทโหสิ “ยานิ ตานิ ภควตา อนุญฺญาตานิ ติจีวรนฺติ วา วสฺสิกสาฏิกาติ วา นิสีทนนฺติ วา ปจฺจตฺถรณนฺติ วา กณฺฑุปฺปฏิจฺฉาทีติ วา มุขปุญฺฉน\-โจฬนฺติ วา ปริกฺขาร\-โจฬนฺติ วา, สพฺพานิ ตานิ อธิฏฺฐาตพฺพานิ นุ โข, อุทาหุ วิกปฺเปตพฺพานี”ติ. ภควโต เอตมตฺถํ อาโรเจสุํ. อนุชานามิ ภิกฺขเว ติจีวรํ อธิฏฺฐาตุํ น วิกปฺเปตุํ, วสฺสิกสาฏิกํ วสฺสานํ จาตุมาสํ อธิฏฺฐาตุํ ตโต ปรํ วิกปฺเปตุํ, นิสีทนํ อธิฏฺฐาตุํ น วิกปฺเปตุํ, ปจฺจตฺถรณํ อธิฏฺฐาตุํ น วิกปฺเปตุํ, กณฺฑุ\-ปฺปฏิจฺฉาทิํ ยาวอาพาธา อธิฏฺฐาตุํ ตโต ปรํ วิกปฺเปตุํ, มุขปุญฺฉน\-โจฬํ อธิฏฺฐาตุํ น วิกปฺเปตุํ, ปริกฺขาร\-โจฬํ อธิฏฺฐาตุํ น วิกปฺเปตุนฺติ.}

\prose{\csromanfolio{412}อถ โข ภิกฺขูนํ เอตทโหสิ “กิตฺตกํ ปจฺฉิมํ นุ โข จีวรํ วิกปฺเปตพฺพนฺ”ติ. ภควโต เอตมตฺถํ อาโรเจสุํ. อนุชานามิ ภิกฺขเว อายาเมน อฏฺฐงฺคุลํ สุคตงฺคุเลน จตุรงฺคุล\-วิตฺถตํ ปจฺฉิมํ จีวรํ วิกปฺเปตุนฺติ.}

\proseitem{359}{เตน โข ปน สมเยน อายสฺมโต มหากสฺสปสฺส ปํสุกูลกโต ครุโก โหติ. ภควโต เอตมตฺถํ อาโรเจสุํ. อนุชานามิ ภิกฺขเว สุตฺตลูขํ กาตุนฺติ. วิกณฺโณ โหติ. ภควโต เอตมตฺถํ อาโรเจสุํ. อนุชานามิ ภิกฺขเว วิกณฺณํ อุทฺธริตุนฺติ. สุตฺตา โอกิริยนฺติ. ภควโต เอตมตฺถํ อาโรเจสุํ. อนุชานามิ ภิกฺขเว อนุวาตํ ปริภณฺฑํ อาโรเปตุนฺติ.}

\prose{เตน โข ปน สมเยน สํฆาฏิยา ปตฺตา ลุชฺชนฺติ. ภควโต เอตมตฺถํ อาโรเจสุํ. อนุชานามิ ภิกฺขเว อฏฺฐปทกํ กาตุนฺติ.}

\proseitem{360}{เตน โข ปน สมเยน อญฺญตรสฺส ภิกฺขุโน ติจีวเร กยิรมาเน สพฺพํ ฉินฺนกํ นปฺปโหติ. ภควโต เอตมตฺถํ อาโรเจสุํ. อนุชานามิ ภิกฺขเว เทฺว ฉินฺนกานิ เอกํ อจฺฉินฺนกนฺติ.}

\prose{เทฺว ฉินฺนกานิ เอกํ อจฺฉินฺนกํ นปฺปโหติ. ภควโต เอตมตฺถํ อาโรเจสุํ. อนุชานามิ ภิกฺขเว เทฺว อจฺฉินฺนกานิ เอกํ ฉินฺนกนฺติ.}

\prose{เทฺว อจฺฉินฺนกานิ เอกํ ฉินฺนกํ นปฺปโหติ. ภควโต เอตมตฺถํ อาโรเจสุํ. อนุชานามิ ภิกฺขเว อนฺวาธิกมฺปิ อาโรเปตุํ, น จ ภิกฺขเว สพฺพํ อจฺฉินฺนกํ ธาเรตพฺพํ, โย ธาเรยฺย, อาปตฺติ ทุกฺกฏสฺสาติ.}

\proseitem{361}{เตน โข ปน สมเยน อญฺญตรสฺส ภิกฺขุโน พหุํ จีวรํ อุปฺปนฺนํ โหติ, โส จ ตํ จีวรํ มาตาปิตูนํ ทาตุกาโม โหติ. ภควโต เอตมตฺถํ อาโรเจสุํ. มาตาปิตโรติ\csromansharedfootnote{b951989ef18fdf06}{มาตาปิตูนํ โข (สี)} โข ภิกฺขเว ททมาเน\csromansharedfootnote{84f9ab86bd27dbbb}{วทมาโน (ก), วทมาเน (?)} กิํ วเทยฺยาม, อนุชานามิ ภิกฺขเว มาตาปิตูนํ ทาตุํ, น จ ภิกฺขเว สทฺธาเทยฺยํ วินิปาเตตพฺพํ, โย วินิปาเตยฺย, อาปตฺติ ทุกฺกฏสฺสาติ.}

\proseitem{362}{\csromanfolio{413}เตน โข ปน สมเยน อญฺญตโร ภิกฺขุ อนฺธวเน จีวรํ นิกฺขิปิตฺวา สนฺตรุตฺตเรน คามํ ปิณฺฑาย ปาวิสิ, โจรา ตํ จีวรํ อวหริํสุ, โส ภิกฺขุ ทุจฺโจโฬ โหติ ลูขจีวโร.}

\prose{ภิกฺขู เอวมาหํสุ “กิสฺส ตฺวํ อาวุโส ทุจฺโจโฬ ลูขจีวโรสี”ติ. อิธาหํ\csromansharedfootnote{338dd9d64fd7ee1e}{[มิสฺสินฺคฺ โนเต 0]} อาวุโส อนฺธวเน จีวรํ นิกฺขิปิตฺวา สนฺตรุตฺตเรน คามํ ปิณฺฑาย ปาวิสิํ, โจรา ตํ จีวรํ อวหริํสุ, เตนาหํ ทุจฺโจโฬ ลูขจีวโรติ. ภควโต เอตมตฺถํ อาโรเจสุํ. น ภิกฺขเว สนฺตรุตฺตเรน คาโม ปวิสิตพฺโพ, โย ปวิเสยฺย, อาปตฺติ ทุกฺกฏสฺสาติ.}

\prose{เตน โข ปน สมเยน อายสฺมา อานนฺโท อสฺสติยา สนฺตรุตฺตเรน คามํ ปิณฺฑาย ปาวิสิ. ภิกฺขู อายสฺมนฺตํ อานนฺทํ เอตทโวจุํ “นนุ อาวุโส อานนฺท ภควตา ปญฺญตฺตํ ‘น สนฺตรุตฺตเรน คาโม ปวิสิตพฺโพ’ติ, กิสฺส ตฺวํ อาวุโส อานนฺท สนฺตรุตฺตเรน คามํ ปวิฏฺโฐ”ติ. สจฺจํ อาวุโส ภควตา ปญฺญตฺตํ “น สนฺตรุตฺตเรน คาโม ปวิสิตพฺโพ”ติ, อปิ จาหํ อสฺสติยา ปวิฏฺโฐติ. ภควโต เอตมตฺถํ อาโรเจสุํ.}

\prose{ปญฺจิเม ภิกฺขเว ปจฺจยา สํฆาฏิยา นิกฺเขปาย. คิลาโน วา โหติ, วสฺสิก\-สงฺเกตํ วา โหติ, นทีปารํ คนฺตุํ วา โหติ, อคฺคฬ\-คุตฺติ\-วิหาโร วา โหติ, อตฺถต\-กถินํ วา โหติ. อิเม โข ภิกฺขเว ปญฺจ ปจฺจยา สํฆาฏิยา นิกฺเขปาย.}

\prose{ปญฺจิเม ภิกฺขเว ปจฺจยา อุตฺตราสงฺคสฺส นิกฺเขปาย ฯเปฯ อนฺตร\-วาสกสฺส นิกฺเขปาย. คิลาโน วา โหติ, วสฺสิก\-สงฺเกตํ วา โหติ, นทีปารํ คนฺตุํ วา โหติ, อคฺคฬ\-คุตฺติ\-วิหาโร วา โหติ, อตฺถต\-กถินํ วา โหติ. อิเม โข ภิกฺขเว ปญฺจ ปจฺจยา อุตฺตราสงฺคสฺส อนฺตร\-วาสกสฺส นิกฺเขปาย.}

\prose{ปญฺจิเม ภิกฺขเว ปจฺจยา วสฺสิก\-สาฏิกาย นิกฺเขปาย. คิลาโน วา โหติ, นิสฺสีมํ คนฺตุํ วา โหติ, นทีปารํ คนฺตุํ วา โหติ, อคฺคฬ\-คุตฺติ\-วิหาโร วา โหติ. วสฺสิกสาฏิกา อกตา วา โหติ วิปฺปกตา วา. อิเม โข ภิกฺขเว ปญฺจ ปจฺจยา วสฺสิก\-สาฏิกาย นิกฺเขปายาติ.}

\csromanmatikaanchor{mtk.230}
\csromantocmarkref{section}{222. สํฆิกจีวรุปฺปาทกถา}{mtk.230}
\bookmark[dest={mtk.230},level=1]{222. สํฆิกจีวรุปฺปาทกถา}
\csromanheader{1}{222. \textbf{สํฆิกจีวรุปฺปาทกถา}}
\proseitem{363}{\csromanfolio{414}เตน โข ปน สมเยน อญฺญตโร ภิกฺขุ เอโก วสฺสํ วสิ. ตตฺถ มนุสฺสา “สํฆสฺส เทมา”ติ จีวรานิ อทํสุ. อถ โข ตสฺส ภิกฺขุโน เอตทโหสิ “ภควตา ปญฺญตฺตํ ‘จตุวคฺโค ปจฺฉิโม สํโฆ’ติ, อหญฺจมฺหิ เอกโก, อิเม จ มนุสฺสา ‘สํฆสฺส เทมา’ติ จีวรานิ อทํสุ, ยนฺนูนาหํ อิมานิ สํฆิกานิ จีวรานิ สาวตฺถิํ หเรยฺยนฺ”ติ. อถ โข โส ภิกฺขุ ตานิ จีวรานิ อาทาย สาวตฺถิํ คนฺตฺวา ภควโต เอตมตฺถํ อาโรเจสิ. ตุเยฺหว ภิกฺขุ ตานิ จีวรานิ ยาว กถินสฺส อุพฺภารายาติ. อิธ ปน ภิกฺขเว ภิกฺขุ เอโก วสฺสํ วสติ, ตตฺถ มนุสฺสา “สํฆสฺส เทมา”ติ จีวรานิ เทนฺติ, อนุชานามิ ภิกฺขเว ตสฺเสว ตานิ จีวรานิ ยาว กถินสฺส อุพฺภารายาติ.}

\prose{เตน โข ปน สมเยน อญฺญตโร ภิกฺขุ อุตุกาลํ เอโก วสิ. ตตฺถ มนุสฺสา “สํฆสฺส เทมา”ติ จีวรานิ อทํสุ. อถ โข ตสฺส ภิกฺขุโน เอตทโหสิ “ภควตา ปญฺญตฺตํ ‘จตุวคฺโค ปจฺฉิโม สํโฆ’ติ, อหญฺจมฺหิ เอกโก, อิเม จ มนุสฺสา ‘สํฆสฺส เทมา’ติ จีวรานิ อทํสุ, ยนฺนูนาหํ อิมานิ สํฆิกานิ จีวรานิ สาวตฺถิํ หเรยฺยนฺ”ติ. อถ โข โส ภิกฺขุ ตานิ จีวรานิ อาทาย สาวตฺถิํ คนฺตฺวา ภิกฺขูนํ เอตมตฺถํ อาโรเจสิ. ภิกฺขู ภควโต เอตมตฺถํ อาโรเจสุํ. อนุชานามิ ภิกฺขเว สมฺมุขีภูเตน สํเฆน ภาเชตุํ, อิธ ปน ภิกฺขเว ภิกฺขุ อุตุกาลํ เอโก วสติ, ตตฺถ มนุสฺสา “สํฆสฺส เทมา”ติ จีวรานิ เทนฺติ, อนุชานามิ ภิกฺขเว เตน ภิกฺขุนา ตานิ จีวรานิ อธิฏฺฐาตุํ “มยฺหิมานิ จีวรานี”ติ. ตสฺส เจ ภิกฺขเว ภิกฺขุโน ตํ จีวรํ อนธิฏฺฐิเต อญฺโญ ภิกฺขุ อาคจฺฉติ, สมโก ทาตพฺโพ ภาโค. เตหิ เจ ภิกฺขเว ภิกฺขูหิ ตํ จีวรํ ภาชิยมาเน อปาติเต กุเส อญฺโญ ภิกฺขุ อาคจฺฉติ, สมโก ทาตพฺโพ ภาโค. เตหิ เจ ภิกฺขเว ภิกฺขูหิ ตํ จีวรํ ภาชิยมาเน ปาติเต กุเส อญฺโญ ภิกฺขุ อาคจฺฉติ, นากามา ทาตพฺโพ ภาโคติ.}

\proseitem{363}{เตน โข ปน สมเยน เทฺว ภาติกา เถรา อายสฺมา จ อิสิทาโส อายสฺมา จ อิสิภโฏ สาวตฺถิยํ วสฺสํวุฏฺฐา อญฺญตรํ คามกาวาสํ อคมํสุ. มนุสฺสา จิรสฺสาปิ เถรา อาคตาติ สจีวรานิ \csromanfolio{415}ภตฺตานิ อทํสุ. อาวาสิกา ภิกฺขู เถเร ปุจฺฉิํสุ “อิมานิ ภนฺเต สํฆิกานิ จีวรานิ เถเร อาคมฺม อุปฺปนฺนานิ, สาทิยิสฺสนฺติ เถรา ภาคนฺ”ติ. เถรา เอวมาหํสุ “ยถา โข มยํ อาวุโส ภควตา ธมฺมํ เทสิตํ อาชานาม, ตุมฺหากํเยว ตานิ จีวรานิ ยาว กถินสฺส อุพฺภารายา”ติ.}

\prose{เตน โข ปน สมเยน ตโย ภิกฺขู ราชคเห วสฺสํ วสนฺติ. ตตฺถ มนุสฺสา “สํฆสฺส เทมา”ติ จีวรานิ เทนฺติ. อถ โข เตสํ ภิกฺขูนํ เอตทโหสิ “ภควตา ปญฺญตฺตํ ‘จตุวคฺโค ปจฺฉิโม สํโฆ’ติ, มยญฺจมฺหา ตโย ชนา, อิเม จ มนุสฺสา ‘สํฆสฺส เทมา’ติ จีวรานิ เทนฺติ, กถํ นุ โข อเมฺหหิ ปฏิปชฺชิตพฺพนฺ”ติ. เตน โข ปน สมเยน สมฺพหุลา เถรา อายสฺมา จ นิลวาสี อายสฺมา จ สาณวาสี อายสฺมา จ โคตโก อายสฺมา จ ภคุ อายสฺมา จ ผฬิกสนฺตาโน ปาฏลิปุตฺเต วิหรนฺติ กุกฺกุฏาราเม. อถ โข เต ภิกฺขู ปาฏลิปุตฺตํ คนฺตฺวา เถเร ปุจฺฉิํสุ. เถรา เอตมาหํสุ “ยถา โข มยํ อาวุโส ภควตา ธมฺมํ เทสิตํ อาชานาม, ตุมฺหากํเยว ตานิ จีวรานิ ยาว กถินสฺส อุพฺภารายา”ติ.}

\csromanmatikaanchor{mtk.231}
\csromantocmarkref{section}{223. อุปนนฺทสกฺยปุตฺตวตฺถุ}{mtk.231}
\bookmark[dest={mtk.231},level=1]{223. อุปนนฺทสกฺยปุตฺตวตฺถุ}
\csromanheader{1}{223. \textbf{อุปนนฺท}\-\textbf{สกฺยปุตฺต}\-\textbf{วตฺถุ}}
\proseitem{364}{เตน โข ปน สมเยน อายสฺมา อุปนนฺโท สกฺยปุตฺโต สาวตฺถิยํ วสฺสํวุฏฺโฐ อญฺญตรํ คามกาวาสํ อคมาสิ, ตตฺถ จ ภิกฺขู จีวรํ ภาเชตุกามา สนฺนิปติํสุ, เต เอวมาหํสุ “อิมานิ โข อาวุโส สํฆิกานิ จีวรานิ ภาชิยิสฺสนฺติ, สาทิยิสฺสสิ ภาคนฺ”ติ. อามาวุโส สาทิยิสฺสามีติ, ตโต จีวรภาคํ คเหตฺวา อญฺญํ อาวาสํ อคมาสิ, ตตฺถปิ ภิกฺขู จีวรํ ภาเชตุกามา สนฺนิปติํสุ, เตปิ เอวมาหํสุ “อิมานิ โข อาวุโส สํฆิกานิ จีวรานิ ภาชิยิสฺสนฺติ, สาทิยิสฺสสิ ภาคนฺ”ติ. อามาวุโส สาทิยิสฺสามีติ, ตโตปิ จีวรภาคํ คเหตฺวา อญฺญํ อาวาสํ อคมาสิ, ตตฺถปิ ภิกฺขู จีวรํ ภาเชตุกามา สนฺนิปติํสุ, เตปิ เอวมาหํสุ “อิมานิ โข อาวุโส สํฆิกานิ จีวรานิ ภาชิยิสฺสนฺติ, สาทิยิสฺสสิ ภาคนฺ”ติ. อามาวุโส สาทิยิสฺสามีติ, ตโตปิ จีวรภาคํ คเหตฺวา มหนฺตํ จีวร\-ภณฺฑิกํ อาทาย ปุนเทว สาวตฺถิํ ปจฺจาคญฺฉิ. ภิกฺขู เอวมาหํสุ “มหาปุญฺโญสิ ตฺวํ อาวุโส อุปนนฺท พหุํ เต จีวรํ อุปฺปนฺนนฺ”ติ. กุโต \csromanfolio{416}เม อาวุโส ปุญฺญํ, อิธาหํ อาวุโส สาวตฺถิยํ วสฺสํวุฏฺโฐ อญฺญตรํ คามกาวสํ อคมาสิํ, ตตฺถ ภิกฺขู จีวรํ ภาเชตุกามา สนฺนิปติํสุ, เต มํ เอวมาหํสุ “อิมานิ โข อาวุโส สํฆิกานิ จีวรานิ ภาชิยิสฺสนฺติ, สาทิยิสฺสสิ ภาคนฺ”ติ. อามาวุโส สาทิยิสฺสามีติ, ตโต จีวราภาคํ คเหตฺวา อญฺญํ อาวาสํ อคมาสิํ, ตตฺถปิ ภิกฺขู จีวรํ ภาเชตุกามา สนฺนิปติํสุ, เตปิ มํ เอวมาหํสุ “อิมานิ โข อาวุโส สํฆิกานิ จีวรานิ ภาชิยิสฺสนฺติ, สาทิยิสฺสสิ ภาคนฺ”ติ. อามาวุโส สาทิยิสฺสามีติ, ตโตปิ จีวรภาคํ คเหตฺวา อญฺญํ อาวาสํ อคมาสิํ, ตตฺถปิ ภิกฺขู จีวรํ ภาเชตุกามา สนฺนิปติํสุ, เตปิ มํ เอวมาหํสุ “อิมานิ โข อาวุโส สํฆิกานิ จีวรานิ ภาชิยิสฺสนฺติ, สาทิยิสฺสสิ ภาคนฺ”ติ. อามาวุโส สาทิยิสฺสามีติ, ตโตปิ จีวรภาคํ อคฺคเหสิํ, เอวํ เม พหุํ จีวรํ อุปฺปนฺนนฺติ. กิํ ปน ตฺวํ อาวุโส อุปนนฺท อญฺญตฺร วสฺสํวุฏฺโฐ อญฺญตฺร จีวรภาคํ สาทิยีติ. เอวมาวุโสติ. เย เต ภิกฺขู อปฺปิจฺฉา ฯเปฯ เต อุชฺฌายนฺติ ขิยฺยนฺติ วิปาเจนฺติ “กถํ หิ นาม อายสฺมา อุปนนฺโท สกฺยปุตฺโต อญฺญตฺร วสฺสํวุฏฺโฐ อญฺญตฺร จีวรภาคํ สาทิยิสฺสตี”ติ. ภควโต เอตมตฺถํ อาโรเจสุํ ฯเปฯ. สจฺจํ กิร ตฺวํ อุปนนฺท อญฺญตฺร วสฺสํวุฏฺโฐ อญฺญตฺร จีวรภาคํ สาทิยีติ. สจฺจํ ภควาติ. วิครหิ พุทฺโธ ภควา ฯเปฯ. กถํ หิ นาม ตฺวํ โมฆปุริส อญฺญตฺร วสฺสํวุฏฺโฐ อญฺญตฺร จีวรภาคํ สาทิยิสฺสสิ. เนตํ โมฆปุริส อปฺปสนฺนานํ วา ปสาทาย ฯเปฯ วิครหิตฺวา ฯเปฯ ธมฺมิํ กถํ กตฺวา ภิกฺขู อามนฺเตสิ “น ภิกฺขเว อญฺญตฺร วสฺสํวุฏฺเฐน อญฺญตฺร จีวรภาโค สาทิตพฺโพ, โย สาทิเยยฺย, อาปตฺติ ทุกฺกฏสฺสา”ติ.}

\prose{เตน โข ปน สมเยน อายสฺมา อุปนนฺโท สกฺยปุตฺโต เอโก ทฺวีสุ อาวาเสสุ วสฺสํ วสิ “เอวํ เม พหุํ จีวรํ อุปฺปชฺชิสฺสตี”ติ. อถ โข เตสํ ภิกฺขูนํ เอตทโหสิ “กถํ นุ โข อายสฺมโต อุปนนฺทสฺส สกฺยปุตฺตสฺส จีวรปฏิวีโส ทาตพฺโพ”ติ. ภควโต เอตมตฺถํ อาโรเจสุํ. เทถ ภิกฺขเว โมฆปุริสสฺส เอกาธิปฺปายํ. อิธ ปน ภิกฺขเว ภิกฺขุ เอโก ทฺวีสุ อาวาเสสุ วสฺสํ วสติ “เอวํ เม พหุํ จีวรํ อุปฺปชฺชิสฺสตี”ติ, สเจ อมุตฺร อุปฑฺฒํ อมุตฺร อุปฑฺฒํ วสติ, อมุตฺร อุปฑฺโฒ อมุตฺร อุปฑฺโฒ จีวรปฏิวีโส ทาตพฺโพ. ยตฺถ วา ปน พหุตรํ วสติ, ตโต จีวรปฏิวีโส ทาตพฺโพติ.}

\csromanmatikaanchor{mtk.232}
\csromantocmarkref{section}{224. คิลานวตฺถุกถา}{mtk.232}
\bookmark[dest={mtk.232},level=1]{224. คิลานวตฺถุกถา}
\csromanheader{1}{224. \textbf{คิลาน}\-\textbf{วตฺถุ}\-\textbf{กถา}}
\proseitem{365}{\csromanfolio{417}เตน โข ปน สมเยน อญฺญตรสฺส ภิกฺขุโน กุจฺฉิวิการา\-พาโธ โหติ, โส สเก มุตฺตกรีเส ปลิปนฺโน เสติ. อถ โข ภควา อายสฺมตา อานนฺเทน ปจฺฉาสมเณน เสนาสนจาริกํ อาหิณฺฑนฺโต เยน ตสฺส ภิกฺขุโน วิหาโร เตนุปสงฺกมิ. อทฺทสา โข ภควา ตํ ภิกฺขุํ สเก มุตฺตกรีเส ปลิปนฺนํ สยมานํ, ทิสฺวาน เยน โส ภิกฺขุ เตนุปสงฺกมิ, อุปสงฺกมิตฺวา ตํ ภิกฺขุํ เอตทโวจ “กิํ เต ภิกฺขุ อาพาโธ”ติ. กุจฺฉิวิกาโร เม ภควาติ. อตฺถิ ปน เต ภิกฺขุ อุปฏฺฐาโกติ. นตฺถิ ภควาติ. กิสฺส ตํ ภิกฺขู น อุปฏฺเฐนฺตีติ. อหํ โข ภนฺเต ภิกฺขูนํ อการโก, เตน มํ ภิกฺขู น อุปฏฺเฐนฺตีติ. อถ โข ภควา อายสฺมนฺตํ อานนฺทํ อามนฺเตสิ “คจฺฉานนฺท อุทกํ อาหร, อิมํ ภิกฺขุํ นหาเปสฺสามา”ติ. “เอวํ ภนฺเต”ติ โข อายสฺมา อานนฺโท ภควโต ปฏิสฺสุณิตฺวา อุทกํ อาหริ. ภควา อุทกํ อาสิญฺจิ. อายสฺมา อานนฺโท ปริโธวิ. ภควา สีสโต อคฺคเหสิ. อายสฺมา อานนฺโท ปาทโต อุจฺจาเรตฺวา มญฺจเก นิปาเตสุํ. อถ โข ภควา เอตสฺมิํ นิทาเน เอตสฺมิํ ปกรเณ ภิกฺขุสํฆํ สนฺนิปาตาเปตฺวา ภิกฺขู ปฏิปุจฺฉิ “อตฺถิ ภิกฺขเว อมุกสฺมิํ วิหาเร ภิกฺขุ คิลาโน”ติ. อตฺถิ ภควาติ. กิํ ตสฺส ภิกฺขเว ภิกฺขุโน อาพาโธติ. ตสฺส ภนฺเต อายสฺมโต กุจฺฉิวิการา\-พาโธติ. อตฺถิ ปน ภิกฺขเว ตสฺส ภิกฺขุโน อุปฏฺฐาโกติ. นตฺถิ ภควาติ. กิสฺส ตํ ภิกฺขู น อุปฏฺเฐนฺตีติ. เอโส ภนฺเต ภิกฺขุ ภิกฺขูนํ อการโก, เตน ตํ ภิกฺขู น อุปฏฺเฐนฺตีติ. นตฺถิ โว ภิกฺขเว มาตา นตฺถิ ปิตา, เย โว อุปฏฺฐเหยฺยุํ, ตุเมฺห เจ ภิกฺขเว อญฺญมญฺญํ น อุปฏฺฐหิสฺสถ, อถ โก จรหิ อุปฏฺฐหิสฺสติ. โย ภิกฺขเว มํ อุปฏฺฐเหยฺย, โส คิลานํ อุปฏฺฐเหยฺย. สเจ อุปชฺฌาโย โหติ, อุปชฺฌาเยน ยาวชีวํ อุปฏฺฐาตพฺโพ, วุฏฺฐานมสฺส อาคเมตพฺพํ. สเจ อาจริโย โหติ, อาจริเยน ยาวชีวํ อุปฏฺฐาตพฺโพ, วุฏฺฐานมสฺส อาคเมตพฺพํ. สเจ สทฺธิวิหาริโก โหติ, สทฺธิ\-วิหาริเกน ยาวชีวํ อุปฏฺฐาตพฺโพ, วุฏฺฐานมสฺส อาคเมตพฺพํ. สเจ อนฺเตวาสิโก โหติ, อนฺเตวาสิเกน ยาวชีวํ อุปฏฺฐาตพฺโพ, วุฏฺฐานมสฺส อาคเมตพฺพํ. สเจ สมานุ\-ปชฺฌายโก โหติ, สมานุ\-ปชฺฌาย\-เกน ยาวชีวํ \csromanfolio{418}อุปฏฺฐาตพฺโพ, วุฏฺฐานมสฺส อาคเมตพฺพํ. สเจ สมานาจริยโก โหติ, สมานา\-จริย\-เกน ยาวชีวํ อุปฏฺฐาตพฺโพ, วุฏฺฐานมสฺส อาคเมตพฺพํ. สเจ น โหติ อุปชฺฌาโย วา อาจริโย วา สทฺธิวิหาริโก วา อนฺเตวาสิโก วา สมานุ\-ปชฺฌายโก วา สมานาจริยโก วา สํเฆน อุปฏฺฐาตพฺโพ, โน เจ อุปฏฺฐเหยฺย, อาปตฺติ ทุกฺกฏสฺส.}

\proseitem{336}{ปญฺจหิ ภิกฺขเว องฺเคหิ สมนฺนาคโต คิลาโน ทูปฏฺโฐ โหติ. อสปฺปายการี โหติ, สปฺปาเย มตฺตํ น ชานาติ, เภสชฺชํ น ปฏิเสวิตา โหติ, อตฺถกามสฺส คิลานุ\-ปฏฺฐากสฺส ยถาภูตํ อาพาธํ นาวิกตฺตา โหติ “อภิกฺกมนฺตํ วา อภิกฺกมตีติ, ปฏิกฺกมนฺตํ วา ปฏิกฺกมตีติ, ฐิตํ วา ฐิโต”ติ, อุปฺปนฺนานํ สารีริกานํ เวทนานํ ทุกฺขานํ ติพฺพานํ ขรานํ กฏุกานํ อสาตานํ อมนาปานํ ปาณหรานํ อนธิวาสก\-ชาติโก โหติ. อิเมหิ โข ภิกฺขเว ปญฺจหงฺเคหิ สมนฺนาคโต คิลาโน ทูปฏฺโฐ โหติ.}

\prose{ปญฺจหิ ภิกฺขเว องฺเคหิ สมนฺนาคโต คิลาโน สูปฏฺโฐ โหติ. สปฺปายการี โหติ, สปฺปาเย มตฺตํ ชานาติ, เภสชฺชํ ปฏิเสวิตา โหติ, อตฺถกามสฺส คิลานุ\-ปฏฺฐากสฺส ยถาภูตํ อาพาธํ อาวิกตฺตา โหติ “อภิกฺกมนฺตํ วา อภิกฺกมตีติ, ปฏิกฺกมนฺตํ วา ปฏิกฺกมตีติ, ฐิตํ วา ฐิโต”ติ, อุปฺปนฺนานํ สารีริกานํ เวทนานํ ทุกฺขานํ ติพฺพานํ ขรานํ กฏุกานํ อสาตานํ อมนาปานํ ปาณหรานํ อธิวาสก\-ชาติโก โหติ. อิเมหิ โข ภิกฺขเว ปญฺจหงฺเคหิ สมนฺนาคโต คิลาโน สูปฏฺโฐ โหติ.}

\prose{ปญฺจหิ ภิกฺขเว องฺเคหิ สมนฺนาคโต คิลานุปฏฺฐาโก นาลํ คิลานํ อุปฏฺฐาตุํ. น ปฏิพโล โหติ เภสชฺชํ สํวิธาตุํ, สปฺปายา\-สปฺปายํ น ชานาติ, อสปฺปายํ อุปนาเมติ สปฺปายํ อปนาเมติ, อามิสนฺตโร คิลานํ อุปฏฺฐาติ โน เมตฺตจิตฺโต, เชคุจฺฉี โหติ อุจฺจารํ วา ปสฺสาวํ วา เขฬํ วา วนฺตํ วา นีหาตุํ, น ปฏิพโล โหติ คิลานํ กาเลน กาลํ ธมฺมิยา กถาย สนฺทสฺเสตุํ สมาทเปตุํ สมุตฺเตเชตุํ สมฺปหํเสตุํ. อิเมหิ โข ภิกฺขเว ปญฺจหงฺเคหิ สมนฺนาคโต คิลานุปฏฺฐาโก นาลํ คิลานํ อุปฏฺฐาตุํ.}

\prose{\csromanfolio{419}ปญฺจหิ ภิกฺขเว องฺเคหิ สมนฺนาคโต คิลานุปฏฺฐาโก อลํ คิลานํ อุปฏฺฐาตุํ. ปฏิพโล โหติ เภสชฺชํ สํวิธาตุํ, สปฺปายา\-สปฺปายํ ชานาติ, อสปฺปายํ อปนาเมติ สปฺปายํ อุปนาเมติ, เมตฺตจิตฺโต คิลานํ อุปฏฺฐาติ โน อามิสนฺตโร, อเชคุจฺฉี โหติ อุจฺจารํ วา ปสฺสาวํ วา เขฬํ วา วนฺตํ วา นีหาตุํ, ปฏิพโล โหติ คิลานํ กาเลน กาลํ ธมฺมิยา กถาย สนฺทสฺเสตุํ สมาทเปตุํ สมุตฺเตเชตุํ สมฺปหํเสตุํ. อิเมหิ โข ภิกฺขเว ปญฺจหงฺเคหิ สมนฺนาคโต คิลานุปฏฺฐาโก อลํ คิลานํ อุปฏฺฐาตุนฺติ.}

\csromanmatikaanchor{mtk.233}
\csromantocmarkref{section}{225. มตสนฺตกกถา}{mtk.233}
\bookmark[dest={mtk.233},level=1]{225. มตสนฺตกกถา}
\csromanheader{1}{225. \textbf{มตสนฺตก}\-\textbf{กถา}}
\proseitem{367}{เตน โข ปน สมเยน เทฺว ภิกฺขู โกสเลสุ ชนปเท อทฺธาน\-มคฺค\-ปฺปฏิปนฺนา โหนฺติ, เต อญฺญตรํ อาวาสํ อุปคจฺฉิํสุ. ตตฺถ อญฺญตโร ภิกฺขุ คิลาโน โหติ. อถ โข เตสํ ภิกฺขูนํ เอตทโหสิ “ภควตา โข อาวุโส คิลานุ\-ปฏฺฐานํ วณฺณิตํ, หนฺท มยํ อาวุโส อิมํ ภิกฺขุํ อุปฏฺฐเหมา”ติ, เต ตํ อุปฏฺฐหิํสุ. โส เตหิ อุปฏฺฐหิย\-มาโน กาลมกาสิ. อถ โข เต ภิกฺขู ตสฺส ภิกฺขุโน ปตฺต\-จีวรมา\-ทาย สาวตฺถิํ คนฺตฺวา ภควโต เอตมตฺถํ อาโรเจสุํ. ภิกฺขุสฺส ภิกฺขเว กาลงฺกเต สํโฆ สามี ปตฺตจีวเร, อปิ จ คิลานุปฏฺฐากา พหูปการา. อนุชานามิ ภิกฺขเว สํเฆน ติจีวรญฺจ ปตฺตญฺจ คิลานุ\-ปฏฺฐากานํ ทาตุํ. เอวญฺจ ปน ภิกฺขเว ทาตพฺพํ, เตน คิลานุ\-ปฏฺฐาเกน ภิกฺขุนา สํฆํ อุปสงฺกมิตฺวา เอวมสฺส วจนีโย “อิตฺถนฺนาโม ภนฺเต ภิกฺขุ กาลงฺกโต, อิทํ ตสฺส ติจีวรญฺจ ปตฺโต จา”ติ. พฺยตฺเตน ภิกฺขุนา ปฏิพเลน สํโฆ ญาเปตพฺโพ\csromandash{}}

\prose{“สุณาตุ เม ภนฺเต สํโฆ, อิตฺถนฺนาโม ภิกฺขุ กาลงฺกโต, อิทํ ตสฺส ติจีวรญฺจ ปตฺโต จ, ยทิ สํฆสฺส ปตฺตกลฺลํ, สํโฆ อิมํ ติจีวรญฺจ ปตฺตญฺจ คิลานุ\-ปฏฺฐากานํ ทเทยฺย, เอสา ญตฺติ.}

\prose{สุณาตุ เม ภนฺเต สํโฆ, อิตฺถนฺนาโม ภิกฺขุ กาลงฺกโต, อิทํ ตสฺส ติจีวรญฺจ ปตฺโต จ, สํโฆ อิมํ ติจีวรญฺจ ปตฺตญฺจ คิลานุ\-ปฏฺฐากานํ เทติ, ยสฺสายสฺมโต ขมติ อิมสฺส ติจีวรสฺส จ ปตฺตสฺส จ คิลานุ\-ปฏฺฐากานํ ทานํ, โส ตุณฺหสฺส, ยสฺส นกฺขมติ, โส ภาเสยฺย.}

\prose{\csromanfolio{420}ทินฺนํ อิทํ สํเฆน ติจีวรญฺจ ปตฺโต จ คิลานุปฏฺฐากาํนํ, ขมติ สํฆสฺส, ตสฺมา ตุณฺหี, เอวเมตํ ธารยามี”ติ.}

\proseitem{368}{เตน โข ปน สมเยน อญฺญตโร สามเณโร กาลงฺกโต โหติ. ภควโต เอตมตฺถํ อาโรเจสุํ. สามเณรสฺส ภิกฺขเว กาลงฺกเต สํโฆ สามี ปตฺตจีวเร, อปิ จ คิลานุปฏฺฐากา พหูปการา, อนุชานามิ ภิกฺขเว สํเฆน จีวรญฺจ ปตฺตญฺจ คิลานุ\-ปฏฺฐากานํ ทาตุํ. เอวญฺจ ปน ภิกฺขเว ทาตพฺพํ, เตน คิลานุ\-ปฏฺฐาเกน ภิกฺขุนา สํฆํ อุปสงฺกมิตฺวา เอวมสฺส วจนีโย “อิตฺถนฺนาโม ภนฺเต สามเณโร กาลงฺกโต, อิทํ ตสฺส จีวรญฺจ ปตฺโต จา”ติ. พฺยตฺเตน ภิกฺขุนา ปฏิพเลน สํโฆ ญาเปตพฺโพ\csromandash{}}

\prose{“สุณาตุ เม ภนฺเต สํโฆ, อิตฺถนฺนาโม สามเณโร กาลงฺกโต, อิทํ ตสฺส จีวรญฺจ ปตฺโต จ, ยทิ สํฆสฺส ปตฺตกลฺลํ, สํโฆ อิมํ จีวรญฺจ ปตฺตญฺจ คิลานุ\-ปฏฺฐากานํ ทเทยฺย, เอสา ญตฺติ.}

\prose{สุณาตุ เม ภนฺเต สํโฆ, อิตฺถนฺนาโม สามเณโร กาลงฺกโต, อิทํ ตสฺส จีวรญฺจ ปตฺโต จ, สํโฆ อิมํ จีวรญฺจ ปตฺตญฺจ คิลานุ\-ปฏฺฐากานํ เทติ, ยสฺสายสฺมโต ขมติ อิมสฺส จีวรสฺส จ ปตฺตสฺส จ คิลานุ\-ปฏฺฐากานํ ทานํ, โส ตุณฺหสฺส, ยสฺส นกฺขมติ, โส ภาเสยฺย.}

\prose{ทินฺนํ อิทํ สํเฆน จีวรญฺจ ปตฺโต จ คิลานุ\-ปฏฺฐากานํ, ขมติ สํฆสฺส, ตสฺมา ตุณฺหี, เอวเมตํ ธารยามี”ติ.}

\proseitem{369}{เตน โข ปน สมเยน อญฺญตโร ภิกฺขุ จ สามเณโร จ คิลานํ อุปฏฺฐหิํสุ, โส เตหิ อุปฏฺฐหิย\-มาโน กาลมกาสิ. อถ โข ตสฺส คิลานุ\-ปฏฺฐากสฺส ภิกฺขุโน เอตทโหสิ “กถํ นุ โข คิลานุ\-ปฏฺฐากสฺส สามเณรสฺส จีวรปฏิวีโส ทาตพฺโพ”ติ. ภควโต เอตมตฺถํ อาโรเจสุํ. อนุชานามิ ภิกฺขเว คิลานุ\-ปฏฺฐากสฺส สามเณรสฺส สมกํ ปฏิวีสํ ทาตุนฺติ.}

\prose{เตน โข ปน สมเยน อญฺญตโร ภิกฺขุ พหุภณฺโฑ พหุปริกฺขาโร กาลงฺกโต โหติ. ภควโต เอตมตฺถํ อาโรเจสุํ. ภิกฺขุสฺส ภิกฺขเว กาลงฺกเต สํโฆ สามี ปตฺตจีวเร, อปิ จ คิลานุปฏฺฐากา พหูปการา. อนุชานามิ ภิกฺขเว สํเฆน ติจีวรญฺจ \csromanfolio{421}ปตฺตญฺจ คิลานุ\-ปฏฺฐากานํ ทาตุํ, ยํ ตตฺถ ลหุภณฺฑํ ลหุปริกฺขารํ, ตํ สมฺมุขีภูเตน สํเฆน ภาเชตุํ, ยํ ตตฺถ ครุภณฺฑํ ครุปริกฺขารํ, ตํ อาคตานาคตสฺส จาตุทฺทิสสฺส สํฆสฺส อวิสฺสชฺชิกํ อเวภงฺคิกนฺติ.}

\csromanmatikaanchor{mtk.234}
\csromantocmarkref{section}{226. นคฺคิยปฏิกฺเขปกถา}{mtk.234}
\bookmark[dest={mtk.234},level=1]{226. นคฺคิยปฏิกฺเขปกถา}
\csromanheader{1}{226. \textbf{นคฺคิย}\-\textbf{ปฏิกฺเขป}\-\textbf{กถา}}
\proseitem{370}{เตน โข ปน สมเยน อญฺญตโร ภิกฺขุ นคฺโค หุตฺวา เยน ภควา เตนุปสงฺกมิ, อุปสงฺกมิตฺวา ภควนฺตํ เอตทโวจ “ภควา ภนฺเต อเนก\-ปริยาเยน อปฺปิจฺฉสฺส สนฺตุฏฺฐสฺส สลฺเลขสฺส ธุตสฺส ปาสาทิกสฺส อปจยสฺส วีริยา\-รมฺภสฺส วณฺณวาที, อิทํ ภนฺเต นคฺคิยํ อเนก\-ปริยาเยน อปฺปิจฺฉตาย สนฺตุฏฺฐิตาย สลฺเลขาย ธุตตาย\csromansharedfootnote{56b265b29ce7fd38}{ธุตตฺตาย (ก)} ปาสาทิกตาย อปจยาย วีริยารมฺภาย สํวตฺตติ, สาธุ ภนฺเต ภควา ภิกฺขูนํ นคฺคิยํ อนุชานาตู”ติ. วิครหิ พุทฺโธ ภควา อนนุจฺฉวิกํ โมฆปุริส อนนุโลมิกํ อปฺปติรูปํ อสฺสามณกํ อกปฺปิยํ อกรณียํ, กถํ หิ นาม ตฺวํ โมฆปุริส นคฺคิยํ ติตฺถิย\-สมาทานํ สมาทิยิสฺสสิ. เนตํ โมฆปุริส อปฺปสนฺนานํ วา ปสาทาย ฯเปฯ วิครหิตฺวา ฯเปฯ ธมฺมิํ กถํ กตฺวา ภิกฺขู อามนฺเตสิ “น ภิกฺขเว นคฺคิยํ ติตฺถิย\-สมาทานํ สมาทิยิตพฺพํ, โย สมาทิเยยฺย, อาปตฺติ ถุลฺลจฺจยสฺสา”ติ.}

\csromanmatikaanchor{mtk.235}
\csromantocmarkref{section}{227. กุสจีราทิปฏิกฺเขปกถา}{mtk.235}
\bookmark[dest={mtk.235},level=1]{227. กุสจีราทิปฏิกฺเขปกถา}
\csromanheader{1}{227. \textbf{กุสจีรา}\-\textbf{ทิ}\-\textbf{ปฏิกฺเขป}\-\textbf{กถา}}
\proseitem{371}{เตน โข ปน สมเยน อญฺญตโร ภิกฺขุ กุสจีรํ นิวาเสตฺวา - ป- วากจีรํ นิวาเสตฺวา ฯเปฯ ผลกจีรํ นิวาเสตฺวา ฯเปฯ เกสกมฺพลํ นิวาเสตฺวา ฯเปฯ วาฬกมฺพลํ นิวาเสตฺวา ฯเปฯ อุลูกปกฺขํ นิวาเสตฺวา ฯเปฯ อชินกฺขิปํ นิวาเสตฺวา เยน ภควา เตนุปสงฺกมิ, อุปสงฺกมิตฺวา ภควนฺตํ เอตทโวจ “ภควา ภนฺเต อเนก\-ปริยาเยน อปฺปิจฺฉสฺส สนฺตุฏฺฐสฺส สลฺเลขสฺส ธุตสฺส ปาสาทิกสฺส อปจยสฺส วีริยา\-รมฺภสฺส วณฺณวาที, อิทํ ภนฺเต อชินกฺขิปํ อเนก\-ปริยาเยน อปฺปิจฺฉตาย สนฺตุฏฺฐิตาย สลฺเลขาย ธุตตาย ปาสาทิกตาย อปจยาย วีริยารมฺภาย สํวตฺตติ, สาธุ ภนฺเต ภควา ภิกฺขูนํ อชินกฺขิปํ อนุชานาตู”ติ. วิครหิ พุทฺโธ ภควา อนนุจฺฉวิกํ โมฆปุริส อนนุโลมิกํ อปฺปติรูปํ \csromanfolio{422}อสฺสามณกํ อกปฺปิยํ อกรณียํ, กถํ หิ นาม ตฺวํ โมฆปุริส อชินกฺขิปํ ติตฺถิยธชํ ธาเรสฺสสิ. เนตํ โมฆปุริส อปฺปสนฺนานํ วา ปสาทาย ฯเปฯ วิครหิตฺวา ฯเปฯ ธมฺมิํ กถํ กตฺวา ภิกฺขู อามนฺเตสิ “น ภิกฺขเว อชินกฺขิปํ ติตฺถิยธชํ ธาเรตพฺพํ, โย ธาเรยฺย, อาปตฺติ ถุลฺลจฺจยสฺสา”ติ.}

\prose{เตน โข ปน สมเยน อญฺญตโร ภิกฺขุ อกฺกนาฬํ นิวาเสตฺวา ฯเปฯ โปตฺถกํ นิวาเสตฺวา เยน ภควา เตนุปสงฺกมิ, อุปสงฺกมิตฺวา ภควนฺตํ เอตทโวจ “ภควา ภนฺเต อเนก\-ปริยาเยน อปฺปิจฺฉสฺส สนฺตุฏฺฐสฺส สลฺเลขสฺส ธุตสฺส ปาสาทิกสฺส อปจยสฺส วีริยา\-รมฺภสฺส วณฺณวาที, อยํ ภนฺเต โปตฺถโก อเนก\-ปริยาเยน อปฺปิจฺฉตาย สนฺตุฏฺฐิตาย สลฺเลขาย ธุตตาย ปาสาทิกตาย อปจยาย วีริยารมฺภาย สํวตฺตติ, สาธุ ภนฺเต ภควา ภิกฺขูนํ โปตฺถกํ อนุชานาตู”ติ. วิครหิ พุทฺโธ ภควา อนนุจฺฉวิกํ โมฆปุริส อนนุโลมิกํ อปฺปติรูปํ อสฺสามณกํ อกปฺปิยํ อกรณียํ, กถํ หิ นาม ตฺวํ โมฆปุริส โปตฺถกํ นิวาเสสฺสสิ. เนตํ โมฆปุริส อปฺปสนฺนานํ วา ปสาทาย ฯเปฯ วิครหิตฺวา ฯเปฯ ธมฺมิํ กถํ กตฺวา ภิกฺขู อามนฺเตสิ “น ภิกฺขเว โปตฺถโก นิวาเสตพฺโพ, โย นิวาเสยฺย, อาปตฺติ ทุกฺกฏสฺสา”ติ.}

\csromanmatikaanchor{mtk.236}
\csromantocmarkref{section}{228. สพฺพนีลกาทิปฏิกฺเขปกถา}{mtk.236}
\bookmark[dest={mtk.236},level=1]{228. สพฺพนีลกาทิปฏิกฺเขปกถา}
\csromanheader{1}{228. \textbf{สพฺพนีลกา}\-\textbf{ทิ}\-\textbf{ปฏิกฺเขป}\-\textbf{กถา}}
\proseitem{372}{เตน โข ปน สมเยน ฉพฺพคฺคิยา ภิกฺขู สพฺพนีลกานิ จีวรานิ ธาเรนฺติ ฯเปฯ สพฺพปีตกานิ จีวรานิ ธาเรนฺติ ฯเปฯ สพฺพ\-โลหิตกานิ จีวรานิ ธาเรนฺติ ฯเปฯ สพฺพ\-มญฺชิฏฺฐกานิ\csromansharedfootnote{bbb69ae8bff0324a}{สพฺพมญฺเชฏฺฐกานิ (สี, สฺยา)} จีวรานิ ธาเรนฺติ ฯเปฯ สพฺพกณฺหานิ จีวรานิ ธาเรนฺติ ฯเปฯ สพฺพ\-มหารงฺค\-รตฺตานิ จีวรานิ ธาเรนฺติ ฯเปฯ สพฺพ\-มหา\-นาม\-รตฺตานิ จีวรานิ ธาเรนฺติ ฯเปฯ อจฺฉินฺนทสานิ จีวรานิ ธาเรนฺติ ฯเปฯ ทีฆทสานิ จีวรานิ ธาเรนฺติ ฯเปฯ ปุปฺผทสานิ จีวรานิ ธาเรนฺติ ฯเปฯ ผณทสานิ\csromansharedfootnote{1be235d56250fee4}{ผลทสานิ (ก)} จีวรานิ ธาเรนฺติ ฯเปฯ กญฺจุกํ ธาเรนฺติ ฯเปฯ ติรีฏกํ ธาเรนฺติ ฯเปฯ เวฐนํ ธาเรนฺติ. มนุสฺสา อุชฺฌายนฺติ ขิยฺยนฺติ วิปาเจนฺติ “กถํ หิ นาม สมณา สกฺยปุตฺติยา เวฐนํ ธาเรสฺสนฺติ. เสยฺยถาปิ คิหี กามโภคิโน”ติ. ภควโต เอตมตฺถํ \csromanfolio{423}อาโรเจสุํ. น ภิกฺขเว สพฺพนีลกานิ จีวรานิ ธาเรตพฺพานิ. น สพฺพปีตกานิ จีวรานิ ธาเรตพฺพานิ. น สพฺพ\-โลหิตกานิ จีวรานิ ธาเรตพฺพานิ. น สพฺพ\-มญฺชิฏฺฐกานิ จีวรานิ ธาเรตพฺพานิ. น สพฺพกณฺหานิ จีวรานิ ธาเรตพฺพานิ. น สพฺพ\-มหารงฺค\-รตฺตานิ จีวรานิ ธาเรตพฺพานิ. น สพฺพ\-มหา\-นาม\-รตฺตานิ จีวรานิ ธาเรตพฺพานิ. น อจฺฉินฺนทสานิ จีวรานิ ธาเรตพฺพานิ. น ทีฆทสานิ จีวรานิ ธาเรตพฺพานิ. น ปุปฺผทสานิ จีวรานิ ธาเรตพฺพานิ. น ผณทสานิ จีวรานิ ธาเรตพฺพานิ. น กญฺจุกํ ธาเรตพฺพํ. น ติรีฏกํ ธาเรตพฺพํ. น เวฐนํ ธาเรตพฺพํ, โย ธาเรยฺย, อาปตฺติ ทุกฺกฏสฺสาติ.}

\csromanmatikaanchor{mtk.237}
\csromantocmarkref{section}{229. วสฺสํวุฏฺฐานํ อนุปฺปนฺนจีวรกถา ...}{mtk.237}
\bookmark[dest={mtk.237},level=1]{229. วสฺสํวุฏฺฐานํ อนุปฺปนฺนจีวรกถา ...}
\csromanheader{1}{229. \textbf{วสฺสํ}\-\textbf{วุฏฺฐานํ อนุปฺปนฺน}\-\textbf{จีวร}\-\textbf{กถา}}
\proseitem{373}{เตน โข ปน สมเยน วสฺสํวุฏฺฐา ภิกฺขู อนุปฺปนฺเน จีวเร ปกฺกมนฺติปิ, วิพฺภมนฺติปิ, กาลมฺปิ กโรนฺติ, สามเณราปิ ปฏิชานนฺติ, สิกฺขํ ปจฺจกฺขาต\-กาปิ ปฏิชานนฺติ, อนฺติมวตฺถุํ อชฺฌาปนฺนกาปิ ปฏิชานนฺติ, อุมฺมตฺตกาปิ ปฏิชานนฺติ, ขิตฺตจิตฺตาปิ ปฏิชานนฺติ, เวทนาฏฺฏาปิ ปฏิชานนฺติ, อาปตฺติยา อทสฺสเน อุกฺขิตฺตกาปิ ปฏิชานนฺติ, อาปตฺติยา อปฺปฏิกมฺเม อุกฺขิตฺตกาปิ ปฏิชานนฺติ, ปาปิกาย ทิฏฺฐิยา อปฺปฏินิสฺสคฺเค อุกฺขิตฺตกาปิ ปฏิชานนฺติ, ปณฺฑกาปิ ปฏิชานนฺติ, เถยฺยสํวาส\-กาปิ ปฏิชานนฺติ, ติตฺถิยปกฺกนฺตกาปิ ปฏิชานนฺติ, ติรจฺฉาน\-คตาปิ ปฏิชานนฺติ, มาตุฆาตกาปิ ปฏิชานนฺติ, ปิตุฆาตกาปิ ปฏิชานนฺติ, อรหนฺต\-ฆาตกาปิ ปฏิชานนฺติ, ภิกฺขุนิทูสกาปิ ปฏิชานนฺติ, สํฆเภทกาปิ ปฏิชานนฺติ, โลหิตุปฺปาท\-กาปิ ปฏิชานนฺติ, อุภโตพฺยญฺชน\-กาปิ ปฏิชานนฺติ. ภควโต เอตมตฺถํ อาโรเจสุํ.}

\proseitem{374}{อิธ ปน ภิกฺขเว วสฺสํวุฏฺโฐ ภิกฺขุ อนุปฺปนฺเน จีวเร ปกฺกมติ, สนฺเต ปติรูเป คาหเก ทาตพฺพํ.}

\prose{อิธ ปน ภิกฺขเว วสฺสํวุฏฺโฐ ภิกฺขุ อนุปฺปนฺเน จีวเร วิพฺภมติ. กาลํ กโรติ. สามเณโร ปฏิชานาติ. สิกฺขํ ปจฺจกฺขาตโก ปฏิชานาติ. อนฺติมวตฺถุํ อชฺฌาปนฺนโก ปฏิชานาติ, สํโฆ สามี.}

\proseitem{374}{อิธ ปน ภิกฺขเว วสฺสํวุฏฺโฐ ภิกฺขุ อนุปฺปนฺเน จีวเร อุมฺมตฺตโก ปฏิชานาติ. ขิตฺตจิตฺโต ปฏิชานาติ. เวทนาฏฺโฏ ปฏิชานาติ. \csromanfolio{424}\textbf{อาปตฺติยา อทสฺสเน อุกฺขิตฺตโก ปฏิชานาติ. อาปตฺติยา อปฺปฏิกมฺเม อุกฺขิตฺตโก} ปฏิชานาติ. ปาปิกาย ทิฏฺฐิยา อปฺปฏินิสฺสคฺเค อุกฺขิตฺตโก ปฏิชานาติ, สนฺเต ปติรูเป คาหเก ทาตพฺพํ.}

\prose{อิธ ปน ภิกฺขเว วสฺสํวุฏฺโฐ ภิกฺขุ อนุปฺปนฺเน จีวเร ปณฺฑโก ปฏิชานาติ. เถยฺย\-สํวาสโก ปฏิชานาติ. ติตฺถิย\-ปกฺกนฺตโก ปฏิชานาติ. ติรจฺฉานคโต ปฏิชานาติ. มาตุฆาตโก ปฏิชานาติ. ปิตุฆาตโก ปฏิชานาติ. อรหนฺต\-ฆาตโก ปฏิชานาติ. ภิกฺขุนิ\-ทูสโก ปฏิชานาติ. สํฆเภทโก ปฏิชานาติ. โลหิตุปฺปาทโก ปฏิชานาติ. อุภโต\-พฺยญฺชนโก ปฏิชานาติ, สํโฆ สามี.}

\proseitem{375}{อิธ ปน ภิกฺขเว วสฺสํวุฏฺโฐ ภิกฺขุ อุปฺปนฺเน จีวเร อภาชิเต ปกฺกมติ, สนฺเต ปติรูเป คาหเก ทาตพฺพํ.}

\prose{อิธ ปน ภิกฺขเว วสฺสํวุฏฺโฐ ภิกฺขุ อุปฺปนฺเน จีวเร อภาชิเต วิพฺภมติ. กาลํ กโรติ. สามเณโร ปฏิชานาติ. สิกฺขํ ปจฺจกฺขาตโก ปฏิชานาติ. อนฺติมวตฺถุํ อชฺฌาปนฺนโก ปฏิชานาติ, สํโฆ สามี.}

\prose{อิธ ปน ภิกฺขเว วสฺสํวุฏฺโฐ ภิกฺขุ อุปฺปนฺเน จีวเร อภาชิเต อุมฺมตฺตโก ปฏิชานาติ. ขิตฺตจิตฺโต ปฏิชานาติ. เวทนาฏฺโฏ ปฏิชานาติ. อาปตฺติยา อทสฺสเน อุกฺขิตฺตโก ปฏิชานาติ. อาปตฺติยา อปฺปฏิกมฺเม อุกฺขิตฺตโก ปฏิชานาติ. ปาปิกาย ทิฏฺฐิยา อปฺปฏินิสฺสคฺเค อุกฺขิตฺตโก ปฏิชานาติ, สนฺเต ปติรูเป คาหเก ทาตพฺพํ.}

\prose{อิธ ปน ภิกฺขเว วสฺสํวุฏฺโฐ ภิกฺขุ อุปฺปนฺเน จีวเร อภาชิเต ปณฺฑโก ปฏิชานาติ. เถยฺย\-สํวาสโก ปฏิชานาติ. ติตฺถิย\-ปกฺกนฺตโก ปฏิชานาติ. ติรจฺฉานคโต ปฏิชานาติ. มาตุฆาตโก ปฏิชานาติ. ปิตุฆาตโก ปฏิชานาติ. อรหนฺต\-ฆาตโก ปฏิชานาติ. ภิกฺขุนิ\-ทูสโก ปฏิชานาติ. สํฆเภทโก ปฏิชานาติ. โลหิตุปฺปาทโก ปฏิชานาติ. อุภโต\-พฺยญฺชนโก ปฏิชานาติ, สํโฆ สามี.}

\csromanmatikaanchor{mtk.238}
\csromantocmarkref{section}{230. สํเฆ ภินฺเน จีวรุปฺปาทกถา}{mtk.238}
\bookmark[dest={mtk.238},level=1]{230. สํเฆ ภินฺเน จีวรุปฺปาทกถา}
\csromanheader{1}{230. \textbf{สํเฆ ภินฺเน จีวรุ}\-\textbf{ปฺปาท}\-\textbf{กถา}}
\proseitem{376}{\csromanfolio{425}อิธ ปน ภิกฺขเว วสฺสํ\-วุฏฺฐานํ ภิกฺขูนํ อนุปฺปนฺเน จีวเร สํโฆ ภิชฺชติ, ตตฺถ มนุสฺสา เอกสฺมิํ ปกฺเข อุทกํ เทนฺติ, เอกสฺมิํ ปกฺเข จีวรํ เทนฺติ “สํฆสฺส เทมา”ติ, สํฆสฺเสเวตํ.}

\prose{อิธ ปน ภิกฺขเว วสฺสํ\-วุฏฺฐานํ ภิกฺขูนํ อนุปฺปนฺเน จีวเร สํโฆ ภิชฺชติ, ตตฺถ มนุสฺสา เอกสฺมิํ ปกฺเข อุทกํ เทนฺติ, ตสฺมิํ เยว ปกฺเข จีวรํ เทนฺติ “สํฆสฺส เทมา”ติ, สํฆสฺเสเวตํ.}

\prose{อิธ ปน ภิกฺขเว วสฺสํ\-วุฏฺฐานํ ภิกฺขูนํ อนุปฺปนฺเน จีวเร สํโฆ ภิชฺชติ, ตตฺถ มนุสฺสา เอกสฺมิํ ปกฺเข อุทกํ เทนฺติ, เอกสฺมิํ ปกฺเข จีวรํ เทนฺติ “ปกฺขสฺส เทมา”ติ, ปกฺขสฺเสเวตํ.}

\prose{อิธ ปน ภิกฺขเว วสฺสํ\-วุฏฺฐานํ ภิกฺขูนํ อนุปฺปนฺเน จีวเร สํโฆ ภิชฺชติ, ตตฺถ มนุสฺสา เอกสฺมิํ ปกฺเข อุทกํ เทนฺติ, ตสฺมิํเยว ปกฺเข จีวรํ เทนฺติ “ปกฺขสฺส เทมา”ติ, ปกฺขสฺเสเวตํ.}

\prose{อิธ ปน ภิกฺขเว วสฺสํ\-วุฏฺฐานํ ภิกฺขูนํ อุปฺปนฺเน จีวเร อภาชิเต สํโฆ ภิชฺชติ, สพฺเพสํ สมกํ ภาเชตพฺพนฺติ.}

\csromanmatikaanchor{mtk.239}
\csromantocmarkref{section}{231. ทุคฺคหิตสุคฺคหิตาทิกถา}{mtk.239}
\bookmark[dest={mtk.239},level=1]{231. ทุคฺคหิตสุคฺคหิตาทิกถา}
\csromanheader{1}{231. \textbf{ทุคฺคหิต}\-\textbf{สุคฺคหิตา}\-\textbf{ทิ}\-\textbf{กถา}}
\proseitem{377}{เตน โข ปน สมเยน อายสฺมา เรวโต อญฺญตรสฺส ภิกฺขุโน หตฺเถ อายสฺมโต สาริปุตฺตสฺส จีวรํ ปาเหสิ “อิมํ จีวรํ เถรสฺส เทหี”ติ. อถ โข โส ภิกฺขุ อนฺตรามคฺเค อายสฺมโต เรวตสฺส วิสฺสาสา ตํ จีวรํ อคฺคเหสิ. อถ โข อายสฺมา เรวโต อายสฺมตา สาริปุตฺเตน สมาคนฺตฺวา ปุจฺฉิ “อหํ ภนฺเต เถรสฺส จีวรํ ปาเหสิํ, สมฺปตฺตํ ตํ จีวรนฺ”ติ. นาหํ ตํ อาวุโส จีวรํ ปสฺสามีติ. อถ โข อายสฺมา เรวโต ตํ ภิกฺขุํ เอตทโวจ “อหํ อาวุโส อายสฺมโต หตฺเถ เถรสฺส จีวรํ ปาเหสิํ, กหํ ตํ จีวรนฺ”ติ. อหํ ภนฺเต อายสฺมโต วิสฺสาสา ตํ จีวรํ อคฺคเหสินฺติ. ภควโต เอตมตฺถํ อาโรเจสุํ.}

\proseitem{378}{\csromanfolio{426}อิธ ปน ภิกฺขเว ภิกฺขุ ภิกฺขุสฺส หตฺเถ จีวรํ ปหิณติ “อิมํ จีวรํ อิตฺถนฺนามสฺส เทหี”ติ. โส อนฺตรามคฺเค โย ปหิณติ, ตสฺส วิสฺสาสา คณฺหาติ, สุคฺคหิตํ. ยสฺส ปหิยฺยติ, ตสฺส วิสฺสาสา คณฺหาติ, ทุคฺคหิตํ.}

\prose{อิธ ปน ภิกฺขเว ภิกฺขุ ภิกฺขุสฺส หตฺเถ จีวรํ ปหิณติ “อิมํ จีวรํ อิตฺถนฺนามสฺส เทหี”ติ. โส อนฺตรามคฺเค ยสฺส ปหิยฺยติ, ตสฺส วิสฺสาสา คณฺหาติ, ทุคฺคหิตํ. โย ปหิณติ, ตสฺส วิสฺสาสา คณฺหาติ, สุคฺคหิตํ.}

\prose{อิธ ปน ภิกฺขเว ภิกฺขุ ภิกฺขุสฺส หตฺเถ จีวรํ ปหิณติ “อิมํ จีวรํ อิตฺถนฺนามสฺส เทหี”ติ. โส อนฺตรามคฺเค สุณาติ “โย ปหิณติ, โส กาลงฺกโต”ติ, ตสฺส มตกจีวรํ อธิฏฺฐาติ, สฺวาธิฏฺฐิตํ. ยสฺส ปหิยฺยติ, ตสฺส วิสฺสาสา คณฺหาติ, ทุคฺคหิตํ.}

\prose{อิธ ปน ภิกฺขเว ภิกฺขุ ภิกฺขุสฺส หตฺเถ จีวรํ ปหิณติ “อิมํ จีวรํ อิตฺถนฺนามสฺส เทหี”ติ. โส อนฺตรามคฺเค สุณาติ “ยสฺส ปหิยฺยติ, โส กาลงฺกโต”ติ, ตสฺส มตกจีวรํ อธิฏฺฐาติ, ทฺวาธิฏฺฐิตํ. โย ปหิณติ, ตสฺส วิสฺสาสา คณฺหาติ, สุคฺคหิตํ.}

\prose{อิธ ปน ภิกฺขเว ภิกฺขุ ภิกฺขุสฺส หตฺเถ จีวรํ ปหิณติ “อิมํ จีวรํ อิตฺถนฺนามสฺส เทหี”ติ. โส อนฺตรามคฺเค สุณาติ “อุโภ กาลงฺกตา”ติ, โย ปหิณติ, ตสฺส มตกจีวรํ อธิฏฺฐาติ, สฺวาธิฏฺฐิตํ. ยสฺส ปหิยฺยติ, ตสฺส มตกจีวรํ อธิฏฺฐาติ, ทฺวาธิฏฺฐิตํ.}

\prose{อิธ ปน ภิกฺขเว ภิกฺขุ ภิกฺขุสฺส หตฺเถ จีวรํ ปหิณติ “อิมํ จีวรํ อิตฺถนฺนามสฺส ทมฺมี”ติ. โส อนฺตรามคฺเค โย คหิณติ, ตสฺส วิสฺสาสา คณฺหาติ, ทุคฺคหิตํ. ยสฺส ปหิยฺยติ, ตสฺส วิสฺสาสา คณฺหาติ, สุคฺคหิตํ.}

\prose{อิธ ปน ภิกฺขเว ภิกฺขุ ภิกฺขุสฺส หตฺเถ จีวรํ ปหิณติ “อิมํ จีวรํ อิตฺถนฺนามสฺส ทมฺมี”ติ. โส อนฺตรามคฺเค ยสฺส ปหิยฺยติ, ตสฺส วิสฺสาสา คณฺหาติ, สุคฺคหิตํ. โย ปหิณติ, ตสฺส วิสฺสาสา คณฺหาติ, ทุคฺคหิตํ.}

\proseitem{378}{อิธ ปน ภิกฺขเว ภิกฺขุ ภิกฺขุสฺส หตฺเถ จีวรํ ปหิณติ “อิมํ จีวรํ อิตฺถนฺนามสฺส ทมฺมี”ติ. โส อนฺตรามคฺเค สุณาติ “โย ปหิณติ, \csromanfolio{427}โส กาลงฺกโต”ติ, ตสฺส มตกจีวรํ อธิฏฺฐาติ, ทฺวาธิฏฺฐิตํ. ยสฺส ปหิยฺยติ, ตสฺส วิสฺสาสา คณฺหาติ, สุคฺคหิตํ.}

\prose{อิธ ปน ภิกฺขเว ภิกฺขุ ภิกฺขุสฺส หตฺเถ จีวรํ ปหิณติ “อิมํ จีวรํ อิตฺถนฺนามสฺส ทมฺมี”ติ. โส อนฺตรามคฺเค สุณาติ “ยสฺส ปหิยฺยติ, โส กาลงฺกโต”ติ, ตสฺส มตกจีวรํ อธิฏฺฐาติ, สฺวาธิฏฺฐิตํ. โย ปหิณติ, ตสฺส วิสฺสาสา คณฺหาติ, ทุคฺคหิตํ.}

\prose{อิธ ปน ภิกฺขเว ภิกฺขุ ภิกฺขุสฺส หตฺเถ จีวรํ ปหิณติ “อิมํ จีวรํ อิตฺถนฺนามสฺส ทมฺมี”ติ. โส อนฺตรามคฺเค สุณาติ “อุโภ กาลงฺกตา”ติ, โย ปหิณติ, ตสฺส มตกจีวรํ อธิฏฺฐาติ, ทฺวาธิฏฺฐิตํ. ยสฺส ปหิยฺยติ, ตสฺส มตกจีวรํ อธิฏฺฐาติ, สฺวาธิฏฺฐิตํ.}

\csromanmatikaanchor{mtk.240}
\csromantocmarkref{section}{232. อฏฺฐจีวรมาติกา}{mtk.240}
\bookmark[dest={mtk.240},level=1]{232. อฏฺฐจีวรมาติกา}
\csromanheader{1}{232. \textbf{อฏฺฐ}\-\textbf{จีวร}\-\textbf{มาติกา}}
\proseitem{379}{อฏฺฐิมา ภิกฺขเว มาติกา จีวรสฺส อุปฺปาทาย. สีมาย เทติ, กติกาย เทติ, ภิกฺขา\-ปญฺญตฺติยา เทติ, สํฆสฺส เทติ, อุภโต\-สํฆสฺส เทติ, วสฺสํวุฏฺฐสํฆสฺส เทติ, อาทิสฺส เทติ, ปุคฺคลสฺส เทติ.}

\setlength{\csromangathapagewidth}{0pt}
\setlength{\csromangathawakwidth}{0pt}
\csromangathameasuregroup{สีมาย เทติ,}{\csromangathabat{สีมาย เทติ,}{ยาวติกา ภิกฺขู อนฺโตสีมคตา, เตหิ ภาเชตพฺพํ.}}
\csromangathasetleft{สีมาย เทติ,}
\csromangathagroup{\csromangathastanza{\csromangathabat{สีมาย เทติ,}{ยาวติกา ภิกฺขู อนฺโตสีมคตา, เตหิ ภาเชตพฺพํ.}}}

\prose{กติกาย เทติ, สมฺพหุลา อาวาสา สมานลาภา โหนฺติ, เอกสฺมิํ อาวาเส ทินฺเน สพฺพตฺถ ทินฺนํ โหติ.}

\prose{ภิกฺขา\-ปญฺญตฺติยา เทติ, ยตฺถ สํฆสฺส ธุวการา กริยฺยนฺติ, ตตฺถ เทติ.}

\setlength{\csromangathapagewidth}{0pt}
\setlength{\csromangathawakwidth}{0pt}
\csromangathameasuregroup{สํฆสฺส เทติ,}{\csromangathabat{สํฆสฺส เทติ,}{สมฺมุขีภูเตน สํเฆน ภาเชตพฺพํ.}}
\csromangathasetleft{สํฆสฺส เทติ,}
\csromangathagroup{\csromangathastanza{\csromangathabat{สํฆสฺส เทติ,}{สมฺมุขีภูเตน สํเฆน ภาเชตพฺพํ.}}}

\prose{อุภโต\-สํฆสฺส เทติ, พหุกาปิ ภิกฺขู โหนฺติ, เอกา ภิกฺขุนี โหติ, อุปฑฺฒํ ทาตพฺพํ, พหุกาปิ ภิกฺขุนิโย โหนฺติ, เอโก ภิกฺขุ โหติ, อุปฑฺฒํ ทาตพฺพํ.}

\prose{วสฺสํวุฏฺฐสํฆสฺส เทติ, ยาวติกา ภิกฺขู ตสฺมิํ อาวาเส วสฺสํวุฏฺฐา, เตหิ ภาเชตพฺพํ.}

\prose{อาทิสฺส เทติ, ยาคุยา วา ภตฺเต วา ขาทนีเย วา จีวเร วา เสนาสเน วา เภสชฺเช วา.}

\prose{ปุคฺคลสฺส เทติ, “อิมํ จีวรํ อิตฺถนฺนามสฺส ทมฺมี”ติ.}

\vspace{\tipitakaparskip}
\csromancenter{จีวร\-กฺขนฺธโก อฏฺฐโม.}
\csromancenter{233. \textbf{ตสฺสุทฺทานํ}}
\csromanmatikaanchor{mtk.241}
\csromantocmarkref{section}{233. อุทฺทานคาถา}{mtk.241}
\bookmark[dest={mtk.241},level=1]{233. อุทฺทานคาถา}
\setlength{\csromangathapagewidth}{0pt}
\setlength{\csromangathawakwidth}{0pt}
\csromangathameasuregroup{ราชคหโก เนคโม, \\ ปุน ราชคหํ คนฺตฺวา, \\ ปุตฺโต สาลวติกาย, \\ ชีวตีติ กุมาเรน, \\ โส หิ ตกฺกสีลํ คนฺตฺวา, \\ สตฺตวสฺสิกอาพาธํ, \\ รญฺโญ ภคนฺทลาพาธํ, \\ มมญฺจ อิตฺถาคารญฺจ, \\ ราชคหโก จ เสฏฺฐิ, \\ ปชฺโชตสฺส มหาโรคํ, \\ อธิการญฺจ สิเวยฺยํ, \\ ตีหิ อุปฺปลหตฺเถหิ, \\ ปกตตฺตํ วรํ ยาจิ, \\ จีวรญฺจ คิหิทานํ, \\ ราชคเห ชนปเท, \\ ปาวาโร โกสิยญฺเจว, \\ อุจฺจาวจา จ สนฺตุฏฺฐิ, \\ ปฐมํ ปจฺฉา สทิสา, \\ ภณฺฑาคารํ อคุตฺตญฺจ, \\ อุสฺสนฺนํ โกลาหลญฺจ, \\ สกาติเรกภาเคน, \\ ฉกเณน สีตุทกา\textsuperscript{1}, \\ โอโรเปนฺตา ภาชนญฺจ, \\ อุปจิกามชฺเฌ ชีรนฺติ, \\ ผรุสาจฺฉินฺนจฺฉิพนฺธา, \\ วีมํสิตฺวา สกฺยมุนิ, \\ อญฺเญน อติเรเกน, \\ จาตุทฺทีโป วรํ ยาจิ, \\ อาคนฺตุคมิคิลานํ, \\ ธุวํ อุทกสาฏิญฺจ, \\ ถุลฺลกจฺฉุมุขํ โขมํ, \\ ปจฺฉิมํ กโต ครุโก, \\ ลุชฺชนฺติ นปฺปโหนฺติ จ, \\ อนฺธวเน อสฺสติยา, \\ เทฺว ภาตุกา ราชคเห, \\ กุจฺฉิวิกาโร คิลาโน, \\ นคฺคา กุสา วากจีรํ, \\ วาฬอุลูกปกฺขญฺจ, \\ โปตฺถกํ นีลปีตญฺจ, \\ กณฺหา มหารงฺคนาม, \\ ทีฆปุปฺผผณทสา, \\ อนุปฺปนฺเน ปกฺกมติ, \\ ปกฺเข ททนฺติ สํฆสฺส, \\ วิสฺสาสคาหาธิฏฺฐาติ,}{\csromangathabat{ราชคหโก เนคโม,}{ทิสฺวา เวสาลิยํ คณิํ.} \\ \csromangathabat{ปุน ราชคหํ คนฺตฺวา,}{รญฺโญ ตํ ปฏิเวทยิ.} \\ \csromangathabat{ปุตฺโต สาลวติกาย,}{อภยสฺส หิ อตฺรโช.} \\ \csromangathabat{ชีวตีติ กุมาเรน,}{สงฺขาโต ชีวโก อิติ.} \\ \csromangathabat{โส หิ ตกฺกสีลํ คนฺตฺวา,}{อุคฺคเหตฺวา มหาภิโส.} \\ \csromangathabat{สตฺตวสฺสิกอาพาธํ,}{นตฺถุกมฺเมน นาสยิ.} \\ \csromangathabat{รญฺโญ ภคนฺทลาพาธํ,}{อาเลเปน อปากฑฺฒิ.} \\ \csromangathabat{มมญฺจ อิตฺถาคารญฺจ,}{พุทฺธสํฆํ จุปฏฺฐหิ.} \\ \csromangathabat{ราชคหโก จ เสฏฺฐิ,}{อนฺตคณฺฐิ ติกิจฺฉิตํ.} \\ \csromangathabat{ปชฺโชตสฺส มหาโรคํ,}{ฆตปาเนน นาสยิ.} \\ \csromangathabat{อธิการญฺจ สิเวยฺยํ,}{อภิสนฺนํ สิเนหติ.} \\ \csromangathabat{ตีหิ อุปฺปลหตฺเถหิ,}{สมตฺติํสวิเรจนํ.} \\ \csromangathabat{ปกตตฺตํ วรํ ยาจิ,}{สิเวยฺยญฺจ ปฏิคฺคหิ.} \\ \csromangathabat{จีวรญฺจ คิหิทานํ,}{อนุญฺญาสิ ตถาคโต.} \\ \csromangathabat{ราชคเห ชนปเท,}{พหุํ อุปฺปชฺชิ จีวรํ.} \\ \csromangathabat{ปาวาโร โกสิยญฺเจว,}{โกชโว อฑฺฒกาสิกํ.} \\ \csromangathabat{อุจฺจาวจา จ สนฺตุฏฺฐิ,}{นาคเมสาคเมสุํ จ.} \\ \csromangathabat{ปฐมํ ปจฺฉา สทิสา,}{กติกา จ ปฏิหรุํ.} \\ \csromangathabat{ภณฺฑาคารํ อคุตฺตญฺจ,}{วุฏฺฐาเปนฺติ ตเถว จ.} \\ \csromangathabat{อุสฺสนฺนํ โกลาหลญฺจ,}{กถํ ภาเช กถํ ทเท.} \\ \csromangathabat{สกาติเรกภาเคน,}{ปฏิวีโส กถํ ทเท.} \\ \csromangathabat{ฉกเณน สีตุทกา\textsuperscript{1},}{อุตฺตริตุ น ชานเร.} \\ \csromangathabat{โอโรเปนฺตา ภาชนญฺจ,}{ปาติยา จ ฉมาย จ.} \\ \csromangathabat{อุปจิกามชฺเฌ ชีรนฺติ,}{เอกโต ปตฺถินฺเนน จ.} \\ \csromangathabat{ผรุสาจฺฉินฺนจฺฉิพนฺธา,}{อทฺทสาสิ อุพฺภณฺฑิเต.} \\ \csromangathabat{วีมํสิตฺวา สกฺยมุนิ,}{อนุญฺญาสิ ติจีวรํ.} \\ \csromangathabat{อญฺเญน อติเรเกน,}{อุปฺปชฺชิ ฉิทฺทเมว จ.} \\ \csromangathabat{จาตุทฺทีโป วรํ ยาจิ,}{ทาตุํ วสฺสิกสาฏิกํ.} \\ \csromangathabat{อาคนฺตุคมิคิลานํ,}{อุปฏฺฐากญฺจ เภสชฺชํ.} \\ \csromangathabat{ธุวํ อุทกสาฏิญฺจ,}{ปณีตํ อติขุทฺทกํ.} \\ \csromangathabat{ถุลฺลกจฺฉุมุขํ โขมํ,}{ปริปุณฺณํ อธิฏฺฐานํ.} \\ \csromangathabat{ปจฺฉิมํ กโต ครุโก,}{วิกณฺโณ สุตฺตโมกิริ.} \\ \csromangathabat{ลุชฺชนฺติ นปฺปโหนฺติ จ,}{อนฺวาธิกํ พหูนิ จ.} \\ \csromangathabat{อนฺธวเน อสฺสติยา,}{เอโก วสฺสํ อุตุมฺหิ จ.} \\ \csromangathabat{เทฺว ภาตุกา ราชคเห,}{อุปนนฺโท ปุน ทฺวิสุ.} \\ \csromangathabat{กุจฺฉิวิกาโร คิลาโน,}{อุโภ เจว คิลานกา\textsuperscript{1}.} \\ \csromangathabat{นคฺคา กุสา วากจีรํ,}{ผลโก เกสกมฺพลํ.} \\ \csromangathabat{วาฬอุลูกปกฺขญฺจ,}{อชินํ อกฺกนาฬกํ.} \\ \csromangathabat{โปตฺถกํ นีลปีตญฺจ,}{โลหิตํ มญฺชิฏฺเฐน จ.} \\ \csromangathabat{กณฺหา มหารงฺคนาม,}{อจฺฉินฺนทสิกา ตถา.} \\ \csromangathabat{ทีฆปุปฺผผณทสา,}{กญฺจุติรีฏเวฐนํ.} \\ \csromangathabat{อนุปฺปนฺเน ปกฺกมติ,}{สํโฆ ภิชฺชติ ตาวเท.} \\ \csromangathabat{ปกฺเข ททนฺติ สํฆสฺส,}{อายสฺมา เรวโต ปหิ.} \\ \csromangathabat{วิสฺสาสคาหาธิฏฺฐาติ,}{อฏฺฐ จีวรมาติกาติ.} \\ อิมมฺหิ ขนฺธเก วตฺถู ฉนฺนวุติ.}
\csromangathameasurewak{อิมมฺหิ ขนฺธเก วตฺถู ฉนฺนวุติ.}
\csromangathasetleft{ราชคหโก เนคโม, \\ ปุน ราชคหํ คนฺตฺวา, \\ ปุตฺโต สาลวติกาย, \\ ชีวตีติ กุมาเรน, \\ โส หิ ตกฺกสีลํ คนฺตฺวา, \\ สตฺตวสฺสิกอาพาธํ, \\ รญฺโญ ภคนฺทลาพาธํ, \\ มมญฺจ อิตฺถาคารญฺจ, \\ ราชคหโก จ เสฏฺฐิ, \\ ปชฺโชตสฺส มหาโรคํ, \\ อธิการญฺจ สิเวยฺยํ, \\ ตีหิ อุปฺปลหตฺเถหิ, \\ ปกตตฺตํ วรํ ยาจิ, \\ จีวรญฺจ คิหิทานํ, \\ ราชคเห ชนปเท, \\ ปาวาโร โกสิยญฺเจว, \\ อุจฺจาวจา จ สนฺตุฏฺฐิ, \\ ปฐมํ ปจฺฉา สทิสา, \\ ภณฺฑาคารํ อคุตฺตญฺจ, \\ อุสฺสนฺนํ โกลาหลญฺจ, \\ สกาติเรกภาเคน, \\ ฉกเณน สีตุทกา\textsuperscript{1}, \\ โอโรเปนฺตา ภาชนญฺจ, \\ อุปจิกามชฺเฌ ชีรนฺติ, \\ ผรุสาจฺฉินฺนจฺฉิพนฺธา, \\ วีมํสิตฺวา สกฺยมุนิ, \\ อญฺเญน อติเรเกน, \\ จาตุทฺทีโป วรํ ยาจิ, \\ อาคนฺตุคมิคิลานํ, \\ ธุวํ อุทกสาฏิญฺจ, \\ ถุลฺลกจฺฉุมุขํ โขมํ, \\ ปจฺฉิมํ กโต ครุโก, \\ ลุชฺชนฺติ นปฺปโหนฺติ จ, \\ อนฺธวเน อสฺสติยา, \\ เทฺว ภาตุกา ราชคเห, \\ กุจฺฉิวิกาโร คิลาโน, \\ นคฺคา กุสา วากจีรํ, \\ วาฬอุลูกปกฺขญฺจ, \\ โปตฺถกํ นีลปีตญฺจ, \\ กณฺหา มหารงฺคนาม, \\ ทีฆปุปฺผผณทสา, \\ อนุปฺปนฺเน ปกฺกมติ, \\ ปกฺเข ททนฺติ สํฆสฺส, \\ วิสฺสาสคาหาธิฏฺฐาติ,}
\csromangathagroupwithclosermajor{\csromangathastanza{\csromanfolio{428}\csromangathabat{ราชคหโก เนคโม,}{ทิสฺวา เวสาลิยํ คณิํ.} \\ \csromangathabat{ปุน ราชคหํ คนฺตฺวา,}{รญฺโญ ตํ ปฏิเวทยิ.}}\csromangathastanza{\csromangathabat{ปุตฺโต สาลวติกาย,}{อภยสฺส หิ อตฺรโช.} \\ \csromangathabat{ชีวตีติ กุมาเรน,}{สงฺขาโต ชีวโก อิติ.}}\csromangathastanza{\csromangathabat{โส หิ ตกฺกสีลํ คนฺตฺวา,}{อุคฺคเหตฺวา มหาภิโส.} \\ \csromangathabat{สตฺตวสฺสิกอาพาธํ,}{นตฺถุกมฺเมน นาสยิ.}}\csromangathastanza{\csromangathabat{รญฺโญ ภคนฺทลา\-พาธํ,}{อาเลเปน อปากฑฺฒิ.} \\ \csromangathabat{มมญฺจ อิตฺถาคารญฺจ,}{พุทฺธสํฆํ จุปฏฺฐหิ.}}\csromangathastanza{\csromangathabat{ราชคหโก จ เสฏฺฐิ,}{อนฺตคณฺฐิ ติกิจฺฉิตํ.} \\ \csromangathabat{ปชฺโชตสฺส มหาโรคํ,}{ฆตปาเนน นาสยิ.}}\csromangathastanza{\csromangathabat{อธิการญฺจ สิเวยฺยํ,}{อภิสนฺนํ สิเนหติ.} \\ \csromangathabat{ตีหิ อุปฺปลหตฺเถหิ,}{สมตฺติํสวิเรจนํ.}}\csromangathastanza{\csromangathabat{ปกตตฺตํ วรํ ยาจิ,}{สิเวยฺยญฺจ ปฏิคฺคหิ.} \\ \csromangathabat{จีวรญฺจ คิหิทานํ,}{อนุญฺญาสิ ตถาคโต.}}\csromangathastanza{\csromangathabat{ราชคเห ชนปเท,}{พหุํ อุปฺปชฺชิ จีวรํ.} \\ \csromangathabat{ปาวาโร โกสิยญฺเจว,}{โกชโว อฑฺฒกาสิกํ.}}\csromangathastanza{\csromangathabat{อุจฺจาวจา จ สนฺตุฏฺฐิ,}{นาคเมสาคเมสุํ จ.} \\ \csromangathabat{ปฐมํ ปจฺฉา สทิสา,}{กติกา จ ปฏิหรุํ.}}\csromangathastanza{\csromangathabat{ภณฺฑาคารํ อคุตฺตญฺจ,}{วุฏฺฐาเปนฺติ ตเถว จ.} \\ \csromangathabat{อุสฺสนฺนํ โกลาหลญฺจ,}{กถํ ภาเช กถํ ทเท.}}\csromangathastanza{\csromangathabat{สกาติเรกภาเคน,}{ปฏิวีโส กถํ ทเท.} \\ \csromangathabat{ฉกเณน สีตุทกา\csromansharedfootnote{ef39131bd15eae7e}{สีตุนฺที จ (สี), สีตุณฺหิ จ (กตฺถจิ)},}{อุตฺตริตุ น ชานเร.}}\csromangathastanza{\csromangathabat{โอโรเปนฺตา ภาชนญฺจ,}{ปาติยา จ ฉมาย จ.} \\ \csromangathabat{อุปจิกามชฺเฌ ชีรนฺติ,}{เอกโต ปตฺถินฺเนน จ.}}\csromangathastanza{\csromanfolio{429}\csromangathabat{ผรุสาจฺฉินฺนจฺฉิพนฺธา,}{อทฺทสาสิ อุพฺภณฺฑิเต.} \\ \csromangathabat{วีมํสิตฺวา สกฺยมุนิ,}{อนุญฺญาสิ ติจีวรํ.}}\csromangathastanza{\csromangathabat{อญฺเญน อติเรเกน,}{อุปฺปชฺชิ ฉิทฺทเมว จ.} \\ \csromangathabat{จาตุทฺทีโป วรํ ยาจิ,}{ทาตุํ วสฺสิกสาฏิกํ.}}\csromangathastanza{\csromangathabat{อาคนฺตุคมิคิลานํ,}{อุปฏฺฐากญฺจ เภสชฺชํ.} \\ \csromangathabat{ธุวํ อุทกสาฏิญฺจ,}{ปณีตํ อติขุทฺทกํ.}}\csromangathastanza{\csromangathabat{ถุลฺลกจฺฉุมุขํ โขมํ,}{ปริปุณฺณํ อธิฏฺฐานํ.} \\ \csromangathabat{ปจฺฉิมํ กโต ครุโก,}{วิกณฺโณ สุตฺตโมกิริ.}}\csromangathastanza{\csromangathabat{ลุชฺชนฺติ นปฺปโหนฺติ จ,}{อนฺวาธิกํ พหูนิ จ.} \\ \csromangathabat{อนฺธวเน อสฺสติยา,}{เอโก วสฺสํ อุตุมฺหิ จ.}}\csromangathastanza{\csromangathabat{เทฺว ภาตุกา ราชคเห,}{อุปนนฺโท ปุน ทฺวิสุ.} \\ \csromangathabat{กุจฺฉิวิกาโร คิลาโน,}{อุโภ เจว คิลานกา\csromansharedfootnote{47da1535d85e0092}{คิลายนา (ก)}.}}\csromangathastanza{\csromangathabat{นคฺคา กุสา วากจีรํ,}{ผลโก เกสกมฺพลํ.} \\ \csromangathabat{วาฬอุลูกปกฺขญฺจ,}{อชินํ อกฺกนาฬกํ.}}\csromangathastanza{\csromangathabat{โปตฺถกํ นีลปีตญฺจ,}{โลหิตํ มญฺชิฏฺเฐน จ.} \\ \csromangathabat{กณฺหา มหารงฺคนาม,}{อจฺฉินฺนทสิกา ตถา.}}\csromangathastanza{\csromangathabat{ทีฆปุปฺผผณทสา,}{กญฺจุติรีฏเวฐนํ.} \\ \csromangathabat{อนุปฺปนฺเน ปกฺกมติ,}{สํโฆ ภิชฺชติ ตาวเท.}}\csromangathastanza{\csromangathabat{ปกฺเข ททนฺติ สํฆสฺส,}{อายสฺมา เรวโต ปหิ.} \\ \csromangathabat{วิสฺสาสคาหาธิฏฺฐาติ,}{อฏฺฐ จีวรมาติกาติ.}}\csromangathastanza{\csromangathawak{อิมมฺหิ ขนฺธเก วตฺถู ฉนฺนวุติ.}}}{\textbf{จีวร}\-\textbf{กฺขนฺธโก นิฏฺฐิโต.}}
\csromanmatikaanchor{mtk.242}
\csromantocmarknopage{chapter}{9. จมฺเปยฺยกฺขนฺธก}{mtk.242}
\bookmark[dest={mtk.242},level=0]{9. จมฺเปยฺยกฺขนฺธก}
\chapterheadpageread{9. \textbf{จมฺเปยฺย}\-\textbf{กฺขนฺธก}}
\csromanmatikaanchor{mtk.243}
\csromantocmarkref{section}{234. กสฺสปโคตฺตภิกฺขุวตฺถุ}{mtk.243}
\bookmark[dest={mtk.243},level=1]{234. กสฺสปโคตฺตภิกฺขุวตฺถุ}
\csromanheader{1}{234. \textbf{กสฺสป}\-\textbf{โคตฺต}\-\textbf{ภิกฺขุ}\-\textbf{วตฺถุ}}
\proseitem{380}{\csromanfolio{430}เตน สมเยน พุทฺโธ ภควา จมฺปายํ วิหรติ คคฺคราย โปกฺขรณิยา ตีเร. เตน โข ปน สมเยน กาสีสุ ชนปเท วาสภคาโม นาม โหติ, ตตฺถ กสฺสปโคตฺโต นาม ภิกฺขุ อาวาสิโก โหติ ตนฺติพทฺโธ อุสฺสุกฺกํ อาปนฺโน “กินฺติ อนาคตา จ เปสลา ภิกฺขู อาคจฺเฉยฺยุํ, อาคตา จ เปสลา ภิกฺขู ผาสุ วิหเรยฺยุํ, อยญฺจ อาวาโส วุทฺธิํ วิรุฬฺหิํ เวปุลฺลํ อาปชฺเชยฺยา”ติ. เตน โข ปน สมเยน สมฺพหุลา ภิกฺขู กาสีสุ จาริกํ จรมานา เยน วาสภคาโม ตทวสรุํ. อทฺทสา โข กสฺสปโคตฺโต ภิกฺขุ เต ภิกฺขู ทูรโตว อาคจฺฉนฺเต, ทิสฺวาน อาสนํ ปญฺญเปสิ, ปาโททกํ ปาทปีฐํ ปาทกถลิกํ อุปนิกฺขิปิ, ปจฺจุคฺคนฺตฺวา ปตฺตจีวรํ ปฏิคฺคเหสิ, ปานีเยน อาปุจฺฉิ, นหาเน อุสฺสุกฺกํ อกาสิ, อุสฺสุกฺกมฺปิ อกาสิ ยาคุยา ขาทนีเย ภตฺตสฺมิํ. อถ โข เตสํ อาคนฺตุกานํ ภิกฺขูนํ เอตทโหสิ “ภทฺทโก โข อยํ อาวุโส อาวาสิโก ภิกฺขุ นหาเน อุสฺสุกฺกํ กโรติ, อุสฺสุกฺกมฺปิ กโรติ ยาคุยา ขาทนีเย ภตฺตสฺมิํ, หนฺท มยํ อาวุโส อิเธว วาสภคาเม นิวาสํ กปฺเปมา”ติ. อถ โข เต อาคนฺตุกา ภิกฺขู ตตฺเถว วาสภคาเม นิวาสํ กปฺเปสุํ.}

\prose{อถ โข กสฺสป\-โคตฺตสฺส ภิกฺขุโน เอตทโหสิ “โย โข อิเมสํ อาคนฺตุกานํ ภิกฺขูนํ อาคนฺตุก\-กิลมโถ, โส ปฏิปฺปสฺสทฺโธ, เยปิเม โคจเร อปฺปกตญฺญุโน, เตทานิเม โคจเร ปกตญฺญุโน, ทุกฺกรํ โข ปน ปรกุเลสุ ยาวชีวํ อุสฺสุกฺกํ กาตุํ, วิญฺญตฺติ จ มนุสฺสานํ อมนาปา ยนฺนูนาหํ น อุสฺสุกฺกํ กเรยฺยํ ยาคุยา ขาทนีเย ภตฺตสฺมินฺ”ติ. โส น อุสฺสุกฺกํ อกาสิ ยาคุยา ขาทนีเย ภตฺตสฺมิํ. อถ โข เตสํ อาคนฺตุกานํ ภิกฺขูนํ เอตทโหสิ “ปุพฺเพ ขฺวายํ อาวุโส อาวาสิโก ภิกฺขุ นหาเน อุสฺสุกฺกํ อกาสิ, อุสฺสุกฺกมฺปิ อกาสิ ยาคุยา ขาทนีเย ภตฺตสฺมิํ, โสทานายํ น อุสฺสุกฺกํ กโรติ ยาคุยา ขาทนีเย ภตฺตสฺมิํ, ทุฏฺโฐทานายํ อาวุโส อาวาสิโก ภิกฺขุ, หนฺท มยํ อาวุโส อาวาสิกํ\csromansharedfootnote{96cb8ced9603db47}{อิมํ อาวาสิกํ (สฺยา)} ภิกฺขุํ อุกฺขิปามา”ติ. อถ โข \csromanfolio{431}เต อาคนฺตุกา ภิกฺขู สนฺนิปติตฺวา กสฺสปโคตฺตํ ภิกฺขุํ เอตทโวจุํ “ปุพฺเพ โข ตฺวํ อาวุโส นหาเน อุสฺสุกฺกํ กโรสิ, อุสฺสุกฺกมฺปิ กโรสิ ยาคุยา ขาทนีเย ภตฺตสฺมิํ, โส ทานิ ตฺวํ น อุสฺสุกฺกํ กโรสิ ยาคุยา ขาทนีเย ภตฺตสฺมิํ, อาปตฺติํ ตฺวํ อาวุโส อาปนฺโน, ปสฺสเสตํ อาปตฺตินฺ”ติ. นตฺถิ เม อาวุโส อาปตฺติ, ยมหํ ปสฺเสยฺยนฺติ. อถ โข เต อาคนฺตุกา ภิกฺขู กสฺสปโคตฺตํ ภิกฺขุํ อาปตฺติยา อทสฺสเน อุกฺขิปิํสุ.}

\prose{อถ โข กสฺสป\-โคตฺตสฺส ภิกฺขุโน เอตทโหสิ “อหํ โข เอตํ น ชานามิ 'อาปตฺติ วา เอสา อนาปตฺติ วา, อาปนฺโน จมฺหิ\csromansharedfootnote{cf6fd4732a14de9e}{วมฺหิ (?)} อนาปนฺโน วา, อุกฺขิตฺโต จมฺหิ อนุกฺขิตฺโต วา, ธมฺมิเกน วา อธมฺมิเกน วา, กุปฺเปน วา อกุปฺเปน วา, ฐานารเหน วา อฏฺฐานารเหน วา’, ยนฺนูนาหํ จมฺปํ คนฺตฺวา ภควนฺตํ เอตมตฺถํ ปุจฺเฉยฺยนฺ”ติ. อถ โข กสฺสปโคตฺโต ภิกฺขุ เสนาสนํ สํสาเมตฺวา ปตฺต\-จีวรมา\-ทาย เยน จมฺปา เตน ปกฺกามิ. อนุปุพฺเพน เยน จมฺปา, เยน ภควา เตนุปสงฺกมิ, อุปสงฺกมิตฺวา ภควนฺตํ อภิวาเทตฺวา เอกมนฺตํ นิสีทิ. อาจิณฺณํ โข ปเนตํ พุทฺธานํ ภควนฺตานํ อาคนฺตุเกหิ ภิกฺขูหิ สทฺธิํ ปฏิสมฺโมทิตุํ. อถ โข ภควา กสฺสปโคตฺตํ ภิกฺขุํ เอตทโวจ “กจฺจิ ภิกฺขุ ขมนียํ, กจฺจิ ยาปนียํ, กจฺจิ อปฺป\-กิลมเถน อทฺธานํ อาคโต, กุโต จ ตฺวํ ภิกฺขุ อาคจฺฉสี”ติ. ขมนียํ ภควา, ยาปนียํ ภควา, อปฺป\-กิลมเถน จาหํ ภนฺเต อทฺธานํ อาคโต, อตฺถิ ภนฺเต กาสีสุ ชนปเท วาสภคาโม นาม, ตตฺถาหํ ภควา อาวาสิโก ตนฺติพทฺโธ อุสฺสุกฺกํ อาปนฺโน “กินฺติ อนาคตา จ เปสลา ภิกฺขู อาคจฺเฉยฺยุํ, อาคตา จ เปสลา ภิกฺขู ผาสุ วิหเรยฺยุํ, อยญฺจ อาวาโส วุทฺธิํ วิรุฬฺหิํ เวปุลฺลํ อาปชฺเชยฺยา”ติ. อถ โข ภนฺเต สมฺพหุลา ภิกฺขู กาสีสุ จาริกํ จรมานา เยน วาสภคาโม ตทวสรุํ. อทฺทสํ โข อหํ ภนฺเต เต ภิกฺขู ทูรโตว อาคจฺฉนฺเต, ทิสฺวาน อาสนํ ปญฺญเปสิํ, ปาโททกํ ปาทปีฐํ ปาทกถลิกํ อุปนิกฺขิปิํ, ปจฺจุคฺคนฺตฺวา ปตฺตจีวรํ ปฏิคฺคเหสิํ, ปานีเยน อาปุจฺฉิํ, นหาเน อุสฺสุกฺกํ อกาสิํ, อุสฺสุกฺกมฺปิ อกาสิํ ยาคุยา ขาทนีเย ภตฺตสฺมิํ. อถ โข เตสํ ภนฺเต อาคนฺตุกานํ ภิกฺขูนํ เอตทโหสิ “ภทฺทโก \csromanfolio{432}โข อยํ อาวุโส อาวาสิโก ภิกฺขุ นหาเน อุสฺสุกฺกํ กโรติ, อุสฺสุกฺกมฺปิ กโรติ ยาคุยา ขาทนีเย ภตฺตสฺมิํ, หนฺท มยํ อาวุโส อิเธว วาสภคาเม นิวาสํ กปฺเปมา”ติ. อถ โข เต ภนฺเต อาคนฺตุกา ภิกฺขู ตตฺเถว วาสภคาเม นิวาสํ กปฺเปสุํ. ตสฺส มยฺหํ ภนฺเต เอตทโหสิ “โย โข อิเมสํ อาคนฺตุกานํ ภิกฺขูนํ อาคนฺตุก\-กิลมโถ, โส ปฏิปฺปสฺสทฺโธ, เยปิเม โคจเร อปฺปกตญฺญุโน, เตทานิเม โคจเร ปกตญฺญุโน, ทุกฺกรํ โข ปน ปรกุเลสุ ยาวชีวํ อุสฺสุกฺกํ กาตุํ, วิญฺญตฺติ จ มนุสฺสานํ อมนาปา, ยนฺนูนาหํ น อุสฺสุกฺกํ กเรยฺยํ ยาคุยา ขาทนีเย ภตฺตสฺมินฺ”ติ. โส โข อหํ ภนฺเต น อุสฺสุกฺกํ อกาสิํ ยาคุยา ขาทนีเย ภตฺตสฺมิํ. อถ โข เตสํ ภนฺเต อาคนฺตุกานํ ภิกฺขูนํ เอตทโหสิ “ปุพฺเพ ขฺวายํ อาวุโส อาวาสิโก ภิกฺขุ นหาเน อุสฺสุกฺกํ กโรติ, อุสฺสุกฺกมฺปิ กโรติ ยาคุยา ขาทนีเย ภตฺตสฺมิํ, โสทานายํ น อุสฺสุกฺกํ กโรติ ยาคุยา ขาทนีเย ภตฺตสฺมิํ, ทุฏฺโฐทานายํ อาวุโส อาวาสิโก ภิกฺขุ, หนฺท มยํ อาวุโส อาวาสิกํ ภิกฺขุํ อุกฺขิปามา”ติ. อถ โข เต ภนฺเต อาคนฺตุกา ภิกฺขู สนฺนิปติตฺวา มํ เอตทโวจุํ “ปุพฺเพ โข ตฺวํ อาวุโส นหาเน อุสฺสุกฺกํ กโรสิ, อุสฺสุกฺกมฺปิ กโรสิ ยาคุยา ขาทนีเย ภตฺตสฺมิํ, โสทานิ ตฺวํ น อุสฺสุกฺกมฺ กโรสิ ยาคุยา ขาทนีเย ภตฺตสฺมิํ, อาปตฺติํ ตฺวํ อาวุโส อาปนฺโน ปสฺสเสตํ อาปตฺตินฺ”ติ. นตฺถิ เม อาวุโส อาปตฺติ, ยมหํ ปสฺเสยฺยนฺติ. อถ โข เต ภนฺเต อาคนฺตุกา ภิกฺขู มํ อาปตฺติยา อทสฺสเน อุกฺขิปิํสุ. ตสฺส มยฺหํ ภนฺเต เอตทโหสิ “อหํ โข เอตํ น ชานามิ 'อาปตฺติ วา เอสา อนาปตฺติ วา, อาปนฺโน จมฺหิ อนาปนฺโน วา, อุกฺขิตฺโต จมฺหิ อนุกฺขิตฺโต วา, ธมฺมิเกน วา อธมฺมิเกน วา, กุปฺเปน วา อกุปฺเปน วา, ฐานารเหน วา อฏฺฐานารเหน วา’, ยนฺนูนาหํ จมฺปํ คนฺตฺวา ภควนฺตํ เอตมตฺถํ ปุจฺเฉยฺยนฺ”ติ. ตโต อหํ ภควา อาคจฺฉามีติ. อนาปตฺติ เอสา ภิกฺขุ, เนสา อาปตฺติ, อนาปนฺโนสิ, นสิ อาปนฺโน, อนุกฺขิตฺโตสิ, นสิ อุกฺขิตฺโต, อธมฺมิเกนาสิ กมฺเมน อุกฺขิตฺโต กุปฺเปน อฏฺฐานารเหน. คจฺฉ ตฺวํ ภิกฺขุ ตตฺเถว วาสภคาเม นิวาสํ กปฺเปหีติ. “เอวํ ภนฺเต”ติ โข กสฺสปโคตฺโต ภิกฺขุ ภควโต ปฏิสฺสุณิตฺวา อุฏฺฐายาสนา ภควนฺตํ อภิวาเทตฺวา ปทกฺขิณํ กตฺวา เยน วาสภคาโม เตน ปกฺกามิ.}

\proseitem{381}{\csromanfolio{433}อถ โข เตสํ อาคนฺตุกานํ ภิกฺขูนํ อหุเทว กุกฺกุจฺจํ, อหุ วิปฺปฏิสาโร “อลาภา วต โน, น วต โน ลาภา, ทุลฺลทฺธํ วต โน, น วต โน สุลทฺธํ, เย มยํ สุทฺธํ ภิกฺขุํ อนาปตฺติกํ อวตฺถุสฺมิํ อการเณ อุกฺขิปิมฺหา, หนฺท มยํ อาวุโส จมฺปํ คนฺตฺวา ภควโต สนฺติเก อจฺจยํ อจฺจยโต เทเสมา”ติ.}

\prose{อถ โข เต อาคนฺตุกา ภิกฺขู เสนาสนํ สํสาเมตฺวา ปตฺต\-จีวรมา\-ทาย เยน จมฺปา เตน ปกฺกมิํสุ. อนุปุพฺเพน เยน จมฺปา, เยน ภควา เตนุ\-ปสงฺกมิํสุ, อุปสงฺกมิตฺวา ภควนฺตํ อภิวาเทตฺวา เอกมนฺตํ นิสีทิํสุ. อาจิณฺณํ โข ปเนตํ พุทฺธานํ ภควนฺตานํ อาคนฺตุเกหิ ภิกฺขูหิ สทฺธิํ ปฏิสมฺโมทิตุํ. อถ โข ภควา เต ภิกฺขู เอตทโวจ “กจฺจิ ภิกฺขเว ขมนียํ, กจฺจิ ยาปนียํ, กจฺจิตฺถ อปฺป\-กิลมเถน อทฺธานํ อาคตา, กุโต จ ตุเมฺห ภิกฺขเว อาคจฺฉถา”ติ. ขมนียํ ภควา, ยาปนียํ ภควา, อปฺป\-กิลมเถน จ มยํ ภนฺเต อทฺธานํ อาคตา, อตฺถิ ภนฺเต กาสีสุ ชนปเท วาสภคาโม นาม, ตโต มยํ ภควา อาคจฺฉามาติ. ตุเมฺห ภิกฺขเว อาวาสิกํ ภิกฺขุํ อุกฺขิปิตฺถาติ. เอวํ ภนฺเตติ. กิสฺมิํ ภิกฺขเว วตฺถุสฺมิํ การเณติ. อวตฺถุสฺมิํ ภควา อการเณติ. วิครหิ พุทฺโธ ภควา อนนุจฺฉวิกํ โมฆปุริสา อนนุโลมิกํ อปฺปติรูปํ อสฺสามณกํ อกปฺปิยํ อกรณียํ, กถํ หิ นาม ตุเมฺห โมฆปุริสา สุทฺธํ ภิกฺขุํ อนาปตฺติกํ อวตฺถุสฺมิํ อการเณ อุกฺขิปิสฺสถ, เนตํ โมฆปุริสา อปฺปสนฺนานํ วา ปสาทาย ฯเปฯ วิครหิตฺวา ธมฺมิํ กถํ กตฺวา ภิกฺขู อามนฺเตสิ “น ภิกฺขเว สุทฺโธ ภิกฺขุ อนาปตฺติโก อวตฺถุสฺมิํ อการเณ อุกฺขิปิตพฺโพ, โย อุกฺขิเปยฺย, อาปตฺติ ทุกฺกฏสฺสา”ติ.}

\prose{อถ โข เต ภิกฺขู อุฏฺฐายาสนา เอกํสํ อุตฺตราสงฺคํ กริตฺวา ภควโต ปาเทสุ สิรสา นิปติตฺวา ภควนฺตํ เอตทโวจุํ “อจฺจโย โน ภนฺเต อจฺจคมา ยถาพาเล ยถามูเฬฺห ยถาอกุสเล, เย มยํ สุทฺธํ ภิกฺขุํ อนาปตฺติกํ อวตฺถุสฺมิํ อการเณ อุกฺขิปิมฺหา, เตสํ โน ภนฺเต ภควา อจฺจยํ อจฺจยโต ปฏิคฺคณฺหาตุ อายติํ สํวรายา”ติ.}

\prose{ตคฺฆ ตุเมฺห ภิกฺขเว อจฺจโย อจฺจคมา ยถาพาเล ยถามูเฬฺห ยถาอกุสเล, เย ตุเมฺห สุทฺธํ ภิกฺขุํ อนาปตฺติกํ อวตฺถุสฺมิํ อการเณ อุกฺขิปิตฺถ, ยโต จ โข ตุเมฺห ภิกฺขเว อจฺจยํ อจฺจยโต ทิสฺวา ยถาธมฺมํ ปฏิกโรถ, ตํ โว มยํ ปฏิคฺคณฺหาม. \csromanfolio{434}วุทฺธิเหสา ภิกฺขเว อริยสฺส วินเย โย อจฺจยํ อจฺจยโต ทิสฺวา ยถาธมฺมํ ปฏิกโรติ, อายติํ\csromansharedfootnote{b4b93413b13273d9}{อายติํ จ (สี)} สํวรํ อาปชฺชตีติ.}

\csromanmatikaanchor{mtk.244}
\csromantocmarkref{section}{235. อธมฺเมน วคฺคาทิกมฺมกถา}{mtk.244}
\bookmark[dest={mtk.244},level=1]{235. อธมฺเมน วคฺคาทิกมฺมกถา}
\csromanheader{1}{235. \textbf{อธมฺเมน วคฺคา}\-\textbf{ทิ}\-\textbf{กมฺม}\-\textbf{กถา}}
\proseitem{382}{เตน โข ปน สมเยน จมฺปายํ ภิกฺขู เอวรูปานิ กมฺมานิ กโรนฺติ, อธมฺเมน วคฺคกมฺมํ กโรนฺติ, อธมฺเมน สมคฺคกมฺมํ กโรนฺติ, ธมฺเมน วคฺคกมฺมํ กโรนฺติ, ธมฺม\-ปติรูปเกน วคฺคกมฺมํ กโรนฺติ, ธมฺม\-ปติรูปเกน สมคฺคกมฺมํ กโรนฺติ, เอโกปิ เอกํ อุกฺขิปติ, เอโกปิ เทฺว อุกฺขิปติ, เอโกปิ สมฺพหุเล อุกฺขิปติ, เอโกปิ สํฆํ อุกฺขิปติ, เทฺวปิ เอกํ อุกฺขิปนฺติ, เทฺวปิ เทฺว อุกฺขิปนฺติ, เทฺวปิ สมฺพหุเล อุกฺขิปนฺติ, เทฺวปิ สํฆํ อุกฺขิปนฺติ, สมฺพหุลาปิ เอกํ อุกฺขิปนฺติ, สมฺพหุลาปิ เทฺว อุกฺขิปนฺติ, สมฺพหุลาปิ สมฺพหุเล อุกฺขิปนฺติ, สมฺพหุลาปิ สํฆํ อุกฺขิปนฺติ, สํโฆปิ สํฆํ อุกฺขิปติ. เย เต ภิกฺขู อปฺปิจฺฉา ฯเปฯ เต อุชฺฌายนฺติ ขิยฺยนฺติ วิปาเจนฺติ “กถํ หิ นาม จมฺปายํ ภิกฺขู เอวรูปานิ กมฺมานิ กริสฺสนฺติ, อธมฺเมน วคฺคกมฺมํ กริสฺสนฺติ, อธมฺเมน สมคฺคกมฺมํ กริสฺสนฺติ, ธมฺเมน วคฺคกมฺมํ กริสฺสนฺติ, ธมฺม\-ปติรูปเกน วคฺคกมฺมํ กริสฺสนฺติ, ธมฺม\-ปติรูปเกน สมคฺคกมฺมํ กริสฺสนฺติ, เอโกปิ เอกํ อุกฺขิปิสฺสติ, เอโกปิ เทฺว อุกฺขิปิสฺสติ, เอโกปิ สมฺพหุเล อุกฺขิปิสฺสติ, เอโกปิ สํฆํ อุกฺขิปิสฺสติ, เทฺวปิ เอกํ อุกฺขิปิสฺสนฺติ, เทฺวปิ เทฺว อุกฺขิปิสฺสนฺติ, เทฺวปิ สมฺพหุเล อุกฺขิปิสฺสนฺติ, เทฺวปิ สํฆํ อุกฺขิปิสฺสนฺติ, สมฺพหุลาปิ เอกํ อุกฺขิปิสฺสนฺติ, สมฺพหุลาปิ เทฺว อุกฺขิปิสฺสนฺติ, สมฺพหุลาปิ สมฺพหุเล อุกฺขิปิสฺสนฺติ, สมฺพหุลาปิ สํฆํ อุกฺขิปิสฺสนฺติ, สํโฆปิ สํฆํ อุกฺขิปิสฺสตี”ติ. อถ โข เต ภิกฺขู ภควโต เอตมตฺถํ อาโรเจสุํ ฯเปฯ. สจฺจํ กิร ภิกฺขเว จมฺปายํ ภิกฺขู เอวรูปานิ กมฺมานิ กโรนฺติ, อธมฺเมน วคฺคกมฺมํ กโรนฺติ ฯเปฯ สํโฆปิ สํฆํ อุกฺขิปตีติ. สจฺจํ ภควาติ. วิครหิ พุทฺโธ ภควา อนนุจฺฉวิกํ ภิกฺขเว เตสํ โมฆปุริสานํ อนนุโลมิกํ อปฺปติรูปํ อสฺสามณกํ อกปฺปิยํ อกรณียํ, กถํ หิ นาม เต ภิกฺขเว โมฆปุริสา เอวรูปานิ กมฺมานิ กริสฺสนฺติ, อธมฺเมน วคฺคกมฺมํ กริสฺสนฺติ ฯเปฯ สํโฆปิ สํฆํ อุกฺขิปิสฺสติ, เนตํ ภิกฺขเว อปฺปสนฺนานํ วา ปสาทาย ฯเปฯ วิครหิตฺวา ฯเปฯ ธมฺมิํ กถํ กตฺวา ภิกฺขู อามนฺเตสิ\csromandash{}}

\proseitem{383}{\csromanfolio{435}อธมฺเมน เจ ภิกฺขเว วคฺคกมฺมํ อกมฺมํ น จ กรณียํ, อธมฺเมน\csromansharedfootnote{e0e1029e4298eadc}{อธมฺเมน เจ ภิกฺขเว (สฺยา)} สมคฺคกมฺมํ อกมฺมํ น จ กรณียํ, ธมฺเมน วคฺคกมฺมํ อกมฺมํ น จ กรณียํ, ธมฺม\-ปติรูปเกน วคฺคกมฺมํ อกมฺมํ น จ กรณียํ, ธมฺม\-ปติรูปเกน สมคฺคกมฺมํ อกมฺมํ น จ กรณียํ, เอโกปิ เอกํ อุกฺขิปติ อกมฺมํ น จ กรณียํ, เอโกปิ เทฺว อุกฺขิปติ อกมฺมํ น จ กรณียํ, เอโกปิ สมฺพหุเล อุกฺขิปติ อกมฺมํ น จ กรณียํ, เอโกปิ สํฆํ อุกฺขิปติ อกมฺมํ น จ กรณียํ, เทฺวปิ เอกํ อุกฺขิปนฺติ อกมฺมํ น จ กรณียํ, เทฺวปิ เทฺว อุกฺขิปนฺติ อกมฺมํ น จ กรณียํ, เทฺวปิ สมฺพหุเล อุกฺขิปนฺติ อกมฺมํ น จ กรณียํ, เทฺวปิ สํฆํ อุกฺขิปนฺติ อกมฺมํ น จ กรณียํ, สมฺพหุลาปิ เอกํ อุกฺขิปนฺติ อกมฺมํ น จ กรณียํ, สมฺพหุลาปิ เทฺว อุกฺขิปนฺติ อกมฺมํ น จ กรณียํ, สมฺพหุลาปิ สมฺพหุเล อุกฺขิปนฺติ อกมฺมํ น จ กรณียํ, สมฺพหุลาปิ สํฆํ อุกฺขิปนฺติ อกมฺมํ น จ กรณียํ, สํโฆปิ สํฆํ อุกฺขิปติ อกมฺมํ น จ กรณียํ.}

\proseitem{384}{จตฺตาริมานิ ภิกฺขเว กมฺมานิ, อธมฺเมน วคฺคกมฺมํ, อธมฺเมน สมคฺคกมฺมํ, ธมฺเมน วคฺคกมฺมํ, ธมฺเมน สมคฺคกมฺมํ. ตตฺร ภิกฺขเว ยทิทํ อธมฺเมน วคฺคกมฺมํ, อิทํ ภิกฺขเว กมฺมํ อธมฺมตฺตา วคฺคตฺตา กุปฺปํ อฏฺฐานารหํ, น ภิกฺขเว เอวรูปํ กมฺมํ กาตพฺพํ, น จ มยา เอวรูปํ กมฺมํ อนุญฺญาตํ. ตตฺร ภิกฺขเว ยทิทํ อธมฺเมน สมคฺคกมฺมํ, อิทํ ภิกฺขเว กมฺมํ อธมฺมตฺตา กุปฺปํ อฏฺฐานารหํ, น ภิกฺขเว เอวรูปํ กมฺมํ กาตพฺพํ, น จ มยา เอวรูปํ กมฺมํ อนุญฺญาตํ. ตตฺร ภิกฺขเว ยทิทํ ธมฺเมน วคฺคกมฺมํ, อิทํ ภิกฺขเว กมฺมํ วคฺคตฺตา กุปฺปํ อฏฺฐานารหํ, น ภิกฺขเว เอวรูปํ กมฺมํ กาตพฺพํ, น จ มยา เอวรูปํ กมฺมํ อนุญฺญาตํ. ตตฺร ภิกฺขเว ยทิทํ ธมฺเมน สมคฺคกมฺมํ, อิทํ ภิกฺขเว กมฺมํ ธมฺมตฺตา สมคฺคตฺตา อกุปฺปํ ฐานารหํ, เอวรูปํ ภิกฺขเว กมฺมํ กาตพฺพํ, เอวรูปญฺจ มยา กมฺมํ อนุญฺญาตํ. ตสฺมาติห ภิกฺขเว เอวรูปํ กมฺมํ กริสฺสาม ยทิทํ ธมฺเมน สมคฺคนฺติ, เอวํ หิ โว ภิกฺขเว สิกฺขิตพฺพนฺติ.}

\csromanmatikaanchor{mtk.245}
\csromantocmarkref{section}{236. ญตฺติวิปนฺนกมฺมาทิกถา}{mtk.245}
\bookmark[dest={mtk.245},level=1]{236. ญตฺติวิปนฺนกมฺมาทิกถา}
\csromanheader{1}{236. \textbf{ญตฺติวิปนฺน}\-\textbf{กมฺมา}\-\textbf{ทิ}\-\textbf{กถา}}
\proseitem{385}{เตน โข ปน สมเยน ฉพฺพคฺคิยา ภิกฺขู เอวรูปานิ กมฺมานิ กโรนฺติ, อธมฺเมน วคฺคกมฺมํ กโรนฺติ, อธมฺเมน สมคฺคกมฺมํ กโรนฺติ, ธมฺเมน \csromanfolio{436}วคฺคกมฺมํ กโรนฺติ, ธมฺม\-ปติรูปเกน วคฺคกมฺมํ กโรนฺติ, ธมฺม\-ปติรูปเกน สมคฺคกมฺมํ กโรนฺติ, ญตฺติ\-วิปนฺนมฺปิ กมฺมํ กโรนฺติ อนุสฺสาวน\-สมฺปนฺนํ, อนุสฺสาวน\-วิปนฺนมฺปิ กมฺมํ กโรนฺติ ญตฺติสมฺปนฺนํ, ญตฺติ\-วิปนฺนมฺปิ อนุสฺสาวน\-วิปนฺนมฺปิ กมฺมํ กโรนฺติ, อญฺญตฺราปิ ธมฺมา กมฺมํ กโรนฺติ, อญฺญตฺราปิ วินยา กมฺมํ กโรนฺติ, อญฺญตฺราปิ สตฺถุสาสนา กมฺมํ กโรนฺติ, ปฏิกุฏฺฐ\-กตมฺปิ กมฺมํ กโรนฺติ อธมฺมิกํ กุปฺปํ อฏฺฐานารหํ. เย เต ภิกฺขู อปฺปิจฺฉา, เต อุชฺฌายนฺติ ขิยฺยนฺติ วิปาเจนฺติ “กถํ หิ นาม ฉพฺพคฺคิยา ภิกฺขู เอวรูปานิ กมฺมานิ กริสฺสนฺติ, อธมฺเมน วคฺคกมฺมํ กริสฺสนฺติ, อธมฺเมน สมคฺคกมฺมํ กริสฺสนฺติ, ธมฺเมน วคฺคกมฺมํ กริสฺสนฺติ, ธมฺม\-ปติรูปเกน วคฺคกมฺมํ กริสฺสนฺติ, ธมฺม\-ปติรูปเกน สมคฺคกมฺมํ กริสฺสนฺติ, ญตฺติ\-วิปนฺนมฺปิ กมฺมํ กริสฺสนฺติ อนุสฺสาวน\-สมฺปนฺนํ, อนุสฺสาวน\-วิปนฺนมฺปิ กมฺมํ กริสฺสนฺติ ญตฺติสมฺปนฺนํ, ญตฺติ\-วิปนฺนมฺปิ อนุสฺสาวน\-วิปนฺนมฺปิ กมฺมํ กริสฺสนฺติ, อญฺญตฺราปิ ธมฺมา กมฺมํ กริสฺสนฺติ, อญฺญตฺราปิ วินยา กมฺมํ กริสฺสนฺติ, อญฺญตฺราปิ สตฺถุสาสนา กมฺมํ กริสฺสนฺติ, ปฏิกุฏฺฐ\-กตมฺปิ กมฺมํ กริสฺสนฺติ อธมฺมิกํ กุปฺปํ อฏฺฐานารหนฺ”ติ. อถ โข เต ภิกฺขู ภควโต เอตมตฺถํ อาโรเจสุํ ฯเปฯ. สจฺจํ กิร ภิกฺขเว ฉพฺพคฺคิยา ภิกฺขู เอวรูปานิ กมฺมานิ กโรนฺติ, อธมฺเมน วคฺคกมฺมํ กโรนฺติ ฯเปฯ ปฏิกุฏฺฐ\-กตมฺปิ กมฺมํ กโรนฺติ อธมฺมิกํ กุปฺปํ อฏฺฐานารหนฺติ. สจฺจํ ภควาติ. วิครหิ พุทฺโธ ภควา ฯเปฯ วิครหิตฺวา ธมฺมิํ กถํ กตฺวา ภิกฺขู อามนฺเตสิ\csromandash{}}

\proseitem{386}{อธมฺเมน เจ ภิกฺขเว วคฺคกมฺมํ อกมฺมํ น จ กรณียํ, อธมฺเมน สมคฺคกมฺมํ อกมฺมํ น จ กรณียํ, ธมฺเมน วคฺคกมฺมํ อกมฺมํ น จ กรณียํ, ธมฺม\-ปติรูปเกน วคฺคกมฺมํ อกมฺมํ น จ กรณียํ, ธมฺม\-ปติรูปเกน สมคฺคกมฺมํ อกมฺมํ น จ กรณียํ, ญตฺติวิปนฺน\-ญฺเจ ภิกฺขเว กมฺมํ อนุสฺสาวน\-สมฺปนฺนํ อกมฺมํ น จ กรณียํ, อนุสฺสาวนวิปนฺน\-ญฺเจ ภิกฺขเว กมฺมํ ญตฺติสมฺปนฺนํ อกมฺมํ น จ กรณียํ, ญตฺติวิปนฺน\-ญฺเจ ภิกฺขเว กมฺมํ อนุสฺสาวน\-วิปนฺนํ อกมฺมํ น จ กรณียํ, อญฺญตฺราปิ ธมฺมา กมฺมํ อกมฺมํ น จ กรณียํ, อญฺญตฺราปิ วินยา กมฺมํ อกมฺมํ น จ กรณียํ, อญฺญตฺราปิ สตฺถุสาสนา กมฺมํ อกมฺมํ น จ กรณียํ, ปฏิกุฏฺฐกตํ เจ ภิกฺขเว กมฺมํ อธมฺมิกํ กุปฺปํ อฏฺฐานารหํ อกมฺมํ น จ กรณียํ.}

\proseitem{387}{ฉยิมานิ ภิกฺขเว กมฺมานิ อธมฺมกมฺมํ วคฺคกมฺมํ สมคฺคกมฺมํ ธมฺม\-ปติรูปเกน วคฺคกมฺมํ ธมฺม\-ปติรูปเกน สมคฺคกมฺมํ ธมฺเมน สมคฺคกมฺมํ.}

\prose{\csromanfolio{437}กตมญฺจ ภิกฺขเว อธมฺมกมฺมํ, ญตฺติทุติเย เจ ภิกฺขเว กมฺเม เอกาย ญตฺติยา กมฺมํ กโรติ, น จ กมฺมวาจํ อนุสฺสาเวติ, อธมฺมกมฺมํ. ญตฺติทุติเย เจ ภิกฺขเว กมฺเม ทฺวีหิ ญตฺตีหิ กมฺมํ กโรติ, น จ กมฺมวาจํ อนุสฺสาเวติ, อธมฺมกมฺมํ. ญตฺติทุติเย เจ ภิกฺขเว กมฺเม เอกาย กมฺมวาจาย กมฺมํ กโรติ, น จ ญตฺติํ ฐเปติ, อธมฺมกมฺมํ. ญตฺติทุติเย เจ ภิกฺขเว กมฺเม ทฺวีหิ กมฺมวาจาหิ กมฺมํ กโรติ, น จ ญตฺติํ ฐเปติ, อธมฺมกมฺมํ. ญตฺติจตุตฺเถ เจ ภิกฺขเว กมฺเม เอกาย ญตฺติยา กมฺมํ กโรติ, น จ กมฺมวาจํ อนุสฺสาเวติ, อธมฺมกมฺมํ. ญตฺติจตุตฺเถ เจ ภิกฺขเว กมฺเม ทฺวีหิ ญตฺตีหิ กมฺมํ กโรติ, น จ กมฺมวาจํ อนุสฺสาเวติ, อธมฺมกมฺมํ. ญตฺติจตุตฺเถ เจ ภิกฺขเว กมฺเม ตีหิ ญตฺตีหิ กมฺมํ กโรติ, น จ กมฺมวาจํ อนุสฺสาเวติ, อธมฺมกมฺมํ. ญตฺติจตุตฺเถ เจ ภิกฺขเว กมฺเม จตูหิ ญตฺตีหิ กมฺมํ กโรติ, น จ กมฺมวาจํ อนุสฺสาเวติ, อธมฺมกมฺมํ. ญตฺติจตุตฺเถ เจ ภิกฺขเว กมฺเม เอกาย กมฺมวาจาย กมฺมํ กโรติ, น จ ญตฺติํ ฐเปติ, อธมฺมกมฺมํ. ญตฺติจตุตฺเถ เจ ภิกฺขเว กมฺเม ทฺวีหิ กมฺมวาจาหิ กมฺมํ กโรติ, น จ ญตฺติํ ฐเปติ, อธมฺมกมฺมํ. ญตฺติจตุตฺเถ เจ ภิกฺขเว กมฺเม ตีหิ กมฺมวาจาหิ กมฺมํ กโรติ, น จ ญตฺติํ ฐเปติ, อธมฺมกมฺมํ. ญตฺติจตุตฺเถ เจ ภิกฺขเว กมฺเม จตูหิ กมฺมวาจาหิ กมฺมํ กโรติ, น จ ญตฺติํ ฐเปติ, อธมฺมกมฺมํ. อิทํ วุจฺจติ ภิกฺขเว อธมฺมกมฺมํ. (1)}

\prose{กตมญฺจ ภิกฺขเว วคฺคกมฺมํ, ญตฺติทุติเย เจ ภิกฺขเว กมฺเม ยาวติกา ภิกฺขู กมฺมปฺปตฺตา, เต อนาคตา โหนฺติ, ฉนฺทารหานํ ฉนฺโท อนาหโฏ โหติ, สมฺมุขีภูตา ปฏิกฺโกสนฺติ, วคฺคกมฺมํ. ญตฺติทุติเย เจ ภิกฺขเว กมฺเม ยาวติกา ภิกฺขู กมฺมปฺปตฺตา, เต อาคตา โหนฺติ, ฉนฺทารหานํ ฉนฺโท อนาหโฏ โหติ, สมฺมุขีภูตา ปฏิกฺโกสนฺติ, วคฺคกมฺมํ. ญตฺติทุติเย เจ ภิกฺขเว กมฺเม ยาวติกา ภิกฺขู กมฺมปฺปตฺตา, เต อาคตา โหนฺติ, ฉนฺทารหานํ ฉนฺโท อาหโฏ โหติ, สมฺมุขีภูตา ปฏิกฺโกสนฺติ, วคฺคกมฺมํ. ญตฺติจตุตฺเถ เจ ภิกฺขเว กมฺเม ยาวติกา ภิกฺขู กมฺมปฺปตฺตา, เต อนาคตา โหนฺติ, ฉนฺทารหานํ ฉนฺโท อนาหโฏ โหติ, สมฺมุขีภูตา ปฏิกฺโกสนฺติ, วคฺคกมฺมํ. ญตฺติจตุตฺเถ เจ ภิกฺขเว กมฺเม ยาวติกา ภิกฺขู กมฺมปฺปตฺตา, เต อาคตา โหนฺติ, ฉนฺทารหานํ ฉนฺโท อนาหโฏ โหติ, \csromanfolio{438}สมฺมุขีภูตา ปฏิกฺโกสนฺติ, วคฺคกมฺมํ. ญตฺติจตุตฺเถ เจ ภิกฺขเว กมฺเม ยาวติกา ภิกฺขู กมฺมปฺปตฺตา, เต อาคตา โหนฺติ, ฉนฺทารหานํ ฉนฺโท อาหโฏ โหติ, สมฺมุขีภูตา ปฏิกฺโกสนฺติ, วคฺคกมฺมํ. อิทํ วุจฺจติ ภิกฺขเว วคฺคกมฺมํ. (2)}

\prose{กตมญฺจ ภิกฺขเว สมคฺคกมฺมํ, ญตฺติทุติเย เจ ภิกฺขเว กมฺเม ยาวติกา ภิกฺขู กมฺมปฺปตฺตา, เต อาคตา โหนฺติ, ฉนฺทารหานํ ฉนฺโท อาหโฏ โหติ, สมฺมุขีภูตา น ปฏิกฺโกสนฺติ, สมคฺคกมฺมํ. ญตฺติจตุตฺเถ เจ ภิกฺขเว กมฺเม ยาวติกา ภิกฺขู กมฺมปฺปตฺตา, เต อาคตา โหนฺติ, ฉนฺทารหานํ ฉนฺโท อาหโฏ โหติ, สมฺมุขีภูตา น ปฏิกฺโกสนฺติ, สมคฺคกมฺมํ. อิทํ วุจฺจติ ภิกฺขเว สมคฺคกมฺมํ. (3)}

\prose{กตมญฺจ ภิกฺขเว ธมฺม\-ปติรูปเกน วคฺคกมฺมํ, ญตฺติทุติเย เจ ภิกฺขเว กมฺเม ปฐมํ กมฺมวาจํ อนุสฺสาเวติ, ปจฺฉา ญตฺติํ ฐเปติ, ยาวติกา ภิกฺขู กมฺมปฺปตฺตา, เต อนาคตา โหนฺติ, ฉนฺทารหานํ ฉนฺโท อนาหโฏ โหติ, สมฺมุขีภูตา ปฏิกฺโกสนฺติ, ธมฺม\-ปติรูปเกน วคฺคกมฺมํ. ญตฺติทุติเย เจ ภิกฺขเว กมฺเม ปฐมฺมํ กมฺมวาจํ อนุสฺสาเวติ, ปจฺฉา ญตฺติํ ฐเปติ, ยาวติกา ภิกฺขู กมฺมปฺปตฺตา, เต อาคตา โหนฺติ, ฉนฺทารหานํ ฉนฺโท อนาหโฏ โหติ, สมฺมุขีภูตา ปฏิกฺโกสนฺติ, ธมฺม\-ปติรูปเกน วคฺคกมฺมํ. ญตฺติทุติเย เจ ภิกฺขเว กมฺเม ปฐมํ กมฺมวาจํ อนุสฺสาเวติ, ปจฺฉา ญตฺติํ ฐเปติ, ยาวติกา ภิกฺขู กมฺมปฺปตฺตา, เต อาคตา โหนฺติ, ฉนฺทารหานํ ฉนฺโท อาหโฏ โหติ, สมฺมุขีภูตา ปฏิกฺโกสนฺติ, ธมฺม\-ปติรูปเกน วคฺคกมฺมํ. ญตฺติจตุตฺเถ เจ ภิกฺขเว กมฺเม ปฐมํ กมฺมวาจํ อนุสฺสาเวติ, ปจฺฉา ญตฺติํ ฐเปติ, ยาวติกา ภิกฺขู กมฺมปฺปตฺตา, เต อนาคตา โหนฺติ, ฉนฺทารหานํ ฉนฺโท อนาหโฏ โหติ, สมฺมุขีภูตา ปฏิกฺโกสนฺติ, ธมฺม\-ปติรูปเกน วคฺคกมฺมํ. ญตฺติจตุตฺเถ เจ ภิกฺขเว กมฺเม ปฐมํ กมฺมวาจํ อนุสฺสาเวติ, ปจฺฉา ญตฺติํ ฐเปติ, ยาวติกา ภิกฺขู กมฺมปฺปตฺตา, เต อาคตา โหนฺติ, ฉนฺทารหานํ ฉนฺโท อนาหโฏ โหติ, สมฺมุขีภูตา ปฏิกฺโกสนฺติ, ธมฺม\-ปติรูปเกน วคฺคกมฺมํ. ญตฺติจตุตฺเถ เจ ภิกฺขเว กมฺเม ปฐมํ กมฺมวาจํ อนุสฺสาเวติ, ปจฺฉา ญตฺติํ ฐเปติ, ยาวติกา ภิกฺขู กมฺมปฺปตฺตา, เต อาคตา โหนฺติ, ฉนฺทารหานํ ฉนฺโท อาหโฏ \csromanfolio{439}โหติ, สมฺมุขีภูตา ปฏิกฺโกสนฺติ, ธมฺม\-ปติรูปเกน วคฺคกมฺมํ. อิทํ วุจฺจติ ภิกฺขเว ธมฺม\-ปติรูปเกน วคฺคกมฺมํ. (4)}

\prose{กตมญฺจ ภิกฺขเว ธมฺม\-ปติรูปเกน สมคฺคกมฺมํ, ญตฺติทุติเย เจ ภิกฺขเว กมฺเม ปฐมํ กมฺมวาจํ อนุสฺสาเวติ, ปจฺฉา ญตฺติํ ฐเปติ, ยาวติกา ภิกฺขู กมฺมปฺปตฺตา, เต อาคตา โหนฺติ, ฉนฺทารหานํ ฉนฺโท อาหโฏ โหติ, สมฺมุขีภูตา น ปฏิกฺโกสนฺติ, ธมฺม\-ปติรูปเกน สมคฺคกมฺมํ. ญตฺติจตุตฺเถ เจ ภิกฺขเว กมฺเม ปฐมํ กมฺมวาจํ อนุสฺสาเวติ, ปจฺฉา ญตฺติํ ฐเปติ, ยาวติกา ภิกฺขู กมฺมปฺปตฺตา, เต อาคตา โหนฺติ, ฉนฺทารหานํ ฉนฺโท อาหโฏ โหติ, สมฺมุขีภูตา น ปฏิกฺโกสนฺติ, ธมฺม\-ปติรูปเกน สมคฺคกมฺมํ. อิทํ วุจฺจติ ภิกฺขเว ธมฺม\-ปติรูปเกน สมคฺคกมฺมํ. (5)}

\prose{กตมญฺจ ภิกฺขเว ธมฺเมน สมคฺคกมฺมํ, ญตฺติทุติเย เจ ภิกฺขเว กมฺเม ปฐมํ ญตฺติํ ฐเปติ, ปจฺฉา เอกาย กมฺมวาจาย กมฺมํ กโรติ, ยาวติกา ภิกฺขู กมฺมปฺปตฺตา, เต อาคตา โหนฺติ, ฉนฺทารหานํ ฉนฺโท อาหโฏ โหติ, สมฺมุขีภูตา น ปฏิกฺโกสนฺติ, ธมฺเมน สมคฺคกมฺมํ. ญตฺติจตุตฺเถ เจ ภิกฺขเว กมฺเม ปฐมํ ญตฺติํ ฐเปติ, ปจฺฉา ตีหิ กมฺมวาจาหิ กมฺมํ กโรติ, ยาวติกา ภิกฺขู กมฺมปฺปตฺตา, เต อาคตา โหนฺติ, ฉนฺทารหานํ ฉนฺโท อาหโฏ โหติ, สมฺมุขีภูตา น ปฏิกฺโกสนฺติ, ธมฺเมน สมคฺคกมฺมํ. อิทํ วุจฺจติ ภิกฺขเว ธมฺเมน สมคฺคกมฺมํ. (6)}

\csromanmatikaanchor{mtk.246}
\csromantocmarkref{section}{237. จตุวคฺคกรณาทิกถา}{mtk.246}
\bookmark[dest={mtk.246},level=1]{237. จตุวคฺคกรณาทิกถา}
\csromanheader{1}{237. \textbf{จตุวคฺค}\-\textbf{กรณา}\-\textbf{ทิ}\-\textbf{กถา}}
\proseitem{388}{ปญฺจ สํฆา, จตุวคฺโค ภิกฺขุสํโฆ, ปญฺจวคฺโค ภิกฺขุสํโฆ, ทสวคฺโค ภิกฺขุสํโฆ, วีสติวคฺโค ภิกฺขุสํโฆ, อติเรก\-วีสติ\-วคฺโค ภิกฺขุสํโฆ. ตตฺร ภิกฺขเว ยฺวายํ จตุวคฺโค ภิกฺขุสํโฆ ฐเปตฺวา ตีณิ กมฺมานิ อุปสมฺปทํ ปวารณํ อพฺภานํ ธมฺเมน สมคฺโค สพฺพกมฺเมสุ กมฺมปฺปตฺโต. ตตฺร ภิกฺขเว ยฺวายํ ปญฺจวคฺโค ภิกฺขุสํโฆ ฐเปตฺวา เทฺว กมฺมานิ มชฺฌิเมสุ ชนปเทสุ อุปสมฺปทํ อพฺภานํ ธมฺเมน สมคฺโค สพฺพกมฺเมสุ กมฺมปฺปตฺโต. ตตฺร ภิกฺขเว ยฺวายํ ทสวคฺโค ภิกฺขุสํโฆ ฐเปตฺวา เอกํ กมฺมํ อพฺภานํ ธมฺเมน สมคฺโค สพฺพกมฺเมสุ กมฺมปฺปตฺโต. ตตฺร ภิกฺขเว ยฺวายํ วีสติวคฺโค ภิกฺขุสํโฆ ธมฺเมน สมคฺโค \csromanfolio{440}สพฺพกมฺเมสุ กมฺมปฺปตฺโต. ตตฺร ภิกฺขเว ยฺวายํ อติเรก\-วีสติ\-วคฺโค ภิกฺขุสํโฆ ธมฺเมน สมคฺโค สพฺพกมฺเมสุ กมฺมปฺปตฺโต.}

\proseitem{389}{จตุวคฺคกรณ\-ญฺเจ ภิกฺขเว กมฺมํ ภิกฺขุนิ\-จตุตฺโถ กมฺมํ กเรยฺย, อกมฺมํ น จ กรณียํ. จตุวคฺคกรณ\-ญฺเจ ภิกฺขเว กมฺมํ สิกฺขมาน\-จตุตฺโถ กมฺมํ กเรยฺย, อกมฺมํ น จ กรณียํ ฯเปฯ สามเณร\-จตุตฺโถ กมฺมํ กเรยฺย, อกมฺมํ น จ กรณียํ. สามเณริ\-จตุตฺโถ กมฺมํ กเรยฺย, อกมฺมํ น จ กรณียํ. สิกฺขํ ปจฺจกฺขาตก\-จตุตฺโถ กมฺมํ กเรยฺย, อกมฺมํ น จ กรณียํ. อนฺติมวตฺถุํ อชฺฌาปนฺนก\-จตุตฺโถ กมฺมํ กเรยฺย, อกมฺมํ น จ กรณียํ. อาปตฺติยา อทสฺสเน อุกฺขิตฺตก\-จตุตฺโถกมฺมํ กเรยฺย, อกมฺมํ น จ กรณียํ. อาปตฺติยา อปฺปฏิกมฺเม อุกฺขิตฺตก\-จตุตฺโถ กมฺมํ กเรยฺย, อกมฺมํ น จ กรณียํ. ปาปิกาย ทิฏฺฐิยา อปฺปฏินิสฺสคฺเค อุกฺขิตฺตก\-จตุตฺโถ กมฺมํ กเรยฺย, อกมฺมํ น จ กรณียํ. ปณฺฑก\-จตุตฺโถ กมฺมํ กเรยฺย, อกมฺมํ น จ กรณียํ. เถยฺยสํวาสก\-จตุตฺโถ กมฺมํ กเรยฺย, อกมฺมํ น จ กรณียํ. ติตฺถิยปกฺกนฺตก\-จตุตฺโถ กมฺมํ กเรยฺย, อกมฺมํ น จ กรณียํ. ติรจฺฉานคต\-จตุตฺโถ กมฺมํ กเรยฺย, อกมฺมํ น จ กรณียํ. มาตุฆาตก\-จตุตฺโถ กมฺมํ กเรยฺย, อกมฺมํ น จ กริณียํ. ปิตุฆาตก\-จตุตฺโถ กมฺมํ กเรยฺย, อกมฺมํ น จ กรณียํ. อรหนฺตฆาตก\-จตุตฺโถ กมฺมํ กเรยฺย, อกมฺมํ น จ กรณียํ. ภิกฺขุนิทูสก\-จตุตฺโถ กมฺมํ กเรยฺย, อกมฺมํ น จ กรณียํ. สํฆเภทกจตุตฺโถ กมฺมํ กเรยฺย, อกมฺมํ น จ กรณียํ. โลหิตุปฺปาทก\-จตุตฺโถ กมฺมํ กเรยฺย, อกมฺมํ น จ กรณียํ. อุภโตพฺยญฺชนก\-จตุตฺโถ กมฺมํ กเรยฺย, อกมฺมํ น จ กรณียํ. นานาสํวาสก\-จตุตฺโถ กมฺมํ กเรยฺย, อกมฺมํ น จ กรณียํ. นานาสีมาย ฐิตจตุตฺโถ กมฺมํ กเรยฺย, อกมฺมํ น จ กรณียํ. อิทฺธิยา เวหาเส ฐิตจตุตฺโถ กมฺมํ กเรยฺย, อกมฺมํ น จ กรณียํ. ยสฺส สํโฆ กมฺมํ กโรติ, ตํจตุตฺโถ กมฺมํ กเรยฺย, อกมฺมํ น จ กรณียํ.}

\vspace{\tipitakaparskip}
\csromancenter{จตุวคฺค\-กรณํ.}
\proseitem{390}{ปญฺจวคฺคกรณ\-ญฺเจ ภิกฺขเว กมฺมํ ภิกฺขุนิ\-ปญฺจโม กมฺมํ กเรยฺย, อกมฺมํ น จ กรณียํ. ปญฺจวคฺคกรณ\-ญฺเจ ภิกฺขเว กมฺมํ สิกฺขมาน\-ปญฺจโม กมฺมํ กเรยฺย ฯเปฯ สามเณร\-ปญฺจโม กมฺมํ กเรยฺย. สามเณริ\-ปญฺจโม กมฺมํ \csromanfolio{441}กเรยฺย. สิกฺขํ ปจฺจกฺขาตก\-ปญฺจโม กมฺมํ กเรยฺย. อนฺติมวตฺถุํ อชฺฌาปนฺนก\-ปญฺจโม กมฺมํ กเรยฺย. อาปตฺติยา อทสฺสเน อุกฺขิตฺตก\-ปญฺจโม กมฺมํ กเรยฺย. อาปตฺติยา อปฺปฏิกมฺเม อุกฺขิตฺตก\-ปญฺจโม กมฺมํ กเรยฺย. ปาปิกาย ทิฏฺฐิยา อปฺปฏินิสฺสคฺเค อุกฺขิตฺตก\-ปญฺจโม กมฺมํ กเรยฺย. ปณฺฑกปญฺจโม กมฺมํ กเรยฺย. เถยฺยสํวาสก\-ปญฺจโม กมฺมํ กเรยฺย. ติตฺถิยปกฺกนฺตก\-ปญฺจโม กมฺมํ กเรยฺย. ติรจฺฉานคต\-ปญฺจโม กมฺมํ กเรยฺย. มาตุฆาตก\-ปญฺจโม กมฺมํ กเรยฺย. ปิตุฆาตก\-ปญฺจโม กมฺมํ กเรยฺย. อรหนฺตฆาตก\-ปญฺจโม กมฺมํ กเรยฺย. ภิกฺขุนิทูสก\-ปญฺจโม กมฺมํ กเรยฺย. สํฆเภทกปญฺจโม กมฺมํ กเรยฺย. โลหิตุปฺปาทก\-ปญฺจโม กมฺมํ กเรยฺย. อุภโตพฺยญฺชนก\-ปญฺจโม กมฺมํ กเรยฺย. นานาสํวาสก\-ปญฺจโม กมฺมํ กเรยฺย. นานาสีมาย ฐิตปญฺจโม กมฺมํ กเรยฺย. อิทฺธิยา เวหาเส ฐิตปญฺจโม กมฺมํ กเรยฺย. ยสฺส สํโฆ กมฺมํ กโรติ, ตํปญฺจโม กมฺมํ กเรยฺย, อกมฺมํ น จ กรณียํ.}

\vspace{\tipitakaparskip}
\csromancenter{ปญฺจวคฺค\-กรณํ.}
\proseitem{391}{ทสวคฺคกรณ\-ญฺเจ ภิกฺขเว กมฺมํ ภิกฺขุนิ\-ทสโม กมฺมํ กเรยฺย, อกมฺมํ น จ กรณียํ. ทสวคฺคกรณ\-ญฺเจ ภิกฺขเว กมฺมํ สิกฺขมาน\-ทสโม กมฺมํ กเรยฺย, อกมฺมํ น จ กรณียํ ฯเปฯ. ทสวคฺคกรณ\-ญฺเจ ภิกฺขเว กมฺมํ ยสฺส สํโฆ กมฺมํ กโรติ, ตํทสโม กมฺมํ กเรยฺย, อกมฺมํ น จ กรณียํ.}

\vspace{\tipitakaparskip}
\csromancenter{ทสวคฺค\-กรณํ.}
\proseitemwithcloserruled{392}{วีสติวคฺคกรณ\-ญฺเจ ภิกฺขเว กมฺมํ ภิกฺขุนิวีโส กมฺมํ กเรยฺย, อกมฺมํ น จ กรณียํ. วีสติวคฺคกรณ\-ญฺเจ ภิกฺขเว กมฺมํ สิกฺขมานวีโส กมฺมํ กเรยฺย ฯเปฯ สามเณรวีโส กมฺมํ กเรยฺย. สามเณริวีโส กมฺมํ กเรยฺย. สิกฺขํ ปจฺจกฺขาตก\-วีโส กมฺมํ กเรยฺย. อนฺติมวตฺถุํ อชฺฌาปนฺนก\-วีโส กมฺมํ กเรยฺย. อาปตฺติยา อทสฺสเน อุกฺขิตฺตกวีโส กมฺมํ กเรยฺย. อาปตฺติยา อปฺปฏิกมฺเม อุกฺขิตฺตกวีโส กมฺมํ กเรยฺย. ปาปิกาย ทิฏฺฐิยา อปฺปฏินิสฺสคฺเค อุกฺขิตฺตกวีโส กมฺมํ กเรยฺย. ปณฺฑกวีโส กมฺมํ กเรยฺย. เถยฺยสํวาสก\-วีโส กมฺมํ กเรยฺย. ติตฺถิยปกฺกนฺตก\-วีโส กมฺมํ กเรยฺย. ติรจฺฉานคต\-วีโส กมฺมํ กเรยฺย. มาตุฆาตก\-วีโส กมฺมํ กเรยฺย. ปิตุฆาตก\-วีโส กมฺมํ กเรยฺย. \csromanfolio{442}อรหนฺตฆาตก\-วีโส กมฺมํ กเรยฺย. ภิกฺขุนิทูสก\-วีโส กมฺมํ กเรยฺย. สํฆเภทกวีโส กมฺมํ กเรยฺย. โลหิตุปฺปาทก\-วีโส กมฺมํ กเรยฺย. อุภโตพฺยญฺชนก\-วีโส กมฺมํ กเรยฺย. นานาสํวาสก\-วีโส กมฺมํ กเรยฺย. นานาสีมาย ฐิตวีโส กมฺมํ กเรยฺย. อิทฺธิยา เวหาเส ฐิตวีโส กมฺมํ กเรยฺย. ยสฺส สํโฆ กมฺมํ กโรติ, ตํวีโส กมฺมํ กเรยฺย, อกมฺมํ น จ กรณียํ.}{วีสติวคฺค\-กรณํ.}
\csromanmatikaanchor{mtk.247}
\csromantocmarkref{section}{238. ปาริวาสิกาทิกถา}{mtk.247}
\bookmark[dest={mtk.247},level=1]{238. ปาริวาสิกาทิกถา}
\csromanheader{1}{238. \textbf{ปาริวาสิกา}\-\textbf{ทิ}\-\textbf{กถา}}
\proseitem{393}{ปาริวาสิก\-จตุตฺโถ เจ ภิกฺขเว ปริวาสํ ทเทยฺย, มูลาย ปฏิกสฺเสยฺย, มานตฺตํ ทเทยฺย, ตํวีโส อพฺเภยฺย, อกมฺมํ น จ กรณียํ. มูลาย ปฏิกสฺสนารห\-จตุตฺโถ เจ ภิกฺขเว ปริวาสํ ทเทยฺย, มูลาย ปฏิกสฺเสยฺย, มานตฺตํ ทเทยฺย, ตํวีโส อพฺเภยฺย, อกมฺมํ น จ กรณียํ. มานตฺตารห\-จตุตฺโถ เจ ภิกฺขเว ปริวาสํ ทเทยฺย, มูลาย ปฏิกสฺเสยฺย, มานตฺตํ ทเทยฺย, ตํวีโส อพฺเภยฺย, อกมฺมํ น จ กรณียํ. มานตฺตจาริก\-จตุตฺโถ เจ ภิกฺขว ปริวาสํ ทเทยฺย, มูลาย ปฏิกสฺเสยฺย, มานตฺตํ ทเทยฺย, ตํวีโส อพฺเภยฺย, อกมฺมํ น จ กรณียํ. อพฺภานารห\-จตุตฺโถ เจ ภิกฺขเว ปริวาสํ ทเทยฺย, มูลาย ปฏิกสฺเสยฺย, มานตฺตํ ทเทยฺย, ตํวีโส อพฺเภยฺย, อกมฺมํ น จ กรณียํ.}

\proseitem{394}{เอกจฺจสฺส ภิกฺขเว สํฆมชฺเฌ ปฏิกฺโกสนา รุหติ, เอกจฺจสฺส น รุหติ. กสฺส จ ภิกฺขเว สํฆมชฺเฌ ปฏิกฺโกสนา น รุหติ, ภิกฺขุนิยา ภิกฺขเว สํฆมชฺเฌ ปฏิกฺโกสนา น รุหติ. สิกฺขมานาย ภิกฺขเว ฯเปฯ สามเณรสฺส ภิกฺขเว. สามเณริยา ภิกฺขเว. สิกฺขา\-ปจฺจกฺขาตกสฺส ภิกฺขเว. อนฺติมวตฺถุํ อชฺฌาปนฺนกสฺส ภิกฺขเว. อุมฺมตฺตกสฺส ภิกฺขเว. ขิตฺตจิตฺตสฺส ภิกฺขเว. เวทนาฏฺฏสฺส ภิกฺขเว. อาปตฺติยา อทสฺสเน อุกฺขิตฺตกสฺส ภิกฺขเว. อาปตฺติยา อปฺปฏิกมฺเม อุกฺขิตฺตกสฺส ภิกฺขเว. ปาปิกาย ทิฏฺฐิยา อปฺปฏินิสฺสคฺเค อุกฺขิตฺตกสฺส ภิกฺขเว. ปณฺฑกสฺส ภิกฺขเว. เถยฺย\-สํวาสกสฺส ภิกฺขเว. ติตฺถิย\-ปกฺกนฺตกสฺส \csromanfolio{443}ภิกฺขเว. ติรจฺฉาน\-คตสฺส ภิกฺขเว. มาตุฆาตกสฺส ภิกฺขเว. ปิตุฆาตกสฺส ภิกฺขเว. อรหนฺต\-ฆาตกสฺส ภิกฺขเว. ภิกฺขุนิ\-ทูสกสฺส ภิกฺขเว. สํฆเภทกสฺส ภิกฺขเว. โลหิตุ\-ปฺปาทกสฺส ภิกฺขเว. อุภโต\-พฺยญฺชนกสฺส ภิกฺขเว. นานา\-สํวาสกสฺส ภิกฺขเว. นานาสีมาย ฐิตสฺส ภิกฺขเว. อิทฺธิยา เวหาเส ฐิตสฺส ภิกฺขเว. ยสฺส สํโฆ กมฺมํ กโรติ, ตสฺส จ\csromansharedfootnote{0e9f8360f7a3e995}{ตสฺส (สฺยา)} ภิกฺขเว สํฆมชฺเฌ ปฏิกฺโกสนา น รุหติ. อิเมสํ โข ภิกฺขเว สํฆมชฺเฌ ปฏิกฺโกสนา น รุหติ.}

\prose{กสฺส จ ภิกฺขเว สํฆมชฺเฌ ปฏิกฺโกสนา รุหติ. ภิกฺขุสฺส ภิกฺขเว ปกตตฺตสฺส สมาน\-สํวาส\-กสฺส สมานสีมาย ฐิตสฺส อนฺตมโส อานนฺตริกสฺสาปิ\csromansharedfootnote{1a60586af9864017}{อนนฺตริกสฺสาปิ (สฺยา)} ภิกฺขุโน วิญฺญาเปนฺตสฺส สํฆมชฺเฌ ปฏิกฺโกสนา รุหติ. อิมสฺส ภิกฺขเว สํฆมชฺเฌ ปฏิกฺโกสนา รุหติ.}

\csromanmatikaanchor{mtk.248}
\csromantocmarkref{section}{239. เทฺวนิสฺสารณาทิกถา}{mtk.248}
\bookmark[dest={mtk.248},level=1]{239. เทฺวนิสฺสารณาทิกถา}
\csromanheader{1}{239. \textbf{เทฺว}\-\textbf{นิสฺสารณา}\-\textbf{ทิ}\-\textbf{กถา}}
\proseitem{395}{เทฺวมา ภิกฺขเว นิสฺสารณา, อตฺถิ ภิกฺขเว ปุคฺคโล อปฺปตฺโต นิสฺสารณํ, ตญฺเจ สํโฆ นิสฺสาเรติ, เอกจฺโจ สุนิสฺสาริโต, เอกจฺโจ ทุนฺนิสฺสาริโต. กตโม จ ภิกฺขเว ปุคฺคโล อปฺปตฺโต นิสฺสารณํ, ตญฺเจ สํโฆ นิสฺสาเรติ, ทุนฺนิสฺสาริโต. อิธ ปน ภิกฺขเว ภิกฺขุ สุทฺโธ โหติ อนาปตฺติโก, ตญฺเจ สํโฆ นิสฺสาเรติ, ทุนฺนิสฺสาริโต. อยํ วุจฺจติ ภิกฺขเว ปุคฺคโล อปฺปตฺโต นิสฺสารณํ, ตญฺเจ สํโฆ นิสฺสาเรติ, ทุนฺนิสฺสาริโต.}

\prose{กตโม จ ภิกฺขเว ปุคฺคโล อปฺปตฺโต นิสฺสารณํ, ตญฺเจ สํโฆ นิสฺสาเรติ, สุนิสฺสาริโต. อิธ ปน ภิกฺขเว ภิกฺขุ พาโล โหติ อพฺยตฺโต อาปตฺติพหุโล อนปทาโน, คิหิสํสฏฺโฐ วิหรติ อนนุโลมิเกหิ คิหิสํสคฺเคหิ, ตญฺเจ สํโฆ นิสฺสาเรติ, สุนิสฺสาริโต. อยํ วุจฺจติ ภิกฺขเว ปุคฺคโล อปฺปตฺโต นิสฺสารณํ, ตญฺเจ สํโฆ นิสฺสาเรติ, สุนิสฺสาริโต.}

\proseitem{396}{เทฺวมา ภิกฺขเว โอสารณา, อตฺถิ ภิกฺขเว ปุคฺคโล อปฺปตฺโต โอสารณํ, ตญฺเจ สํโฆ โอสาเรติ, เอกจฺโจ โสสาริโต \csromanfolio{444}เอกจฺโจ โทสาริโต. กตโม จ ภิกฺขเว ปุคฺคโล อปฺปตฺโต โอสารณํ, ตญฺเจ สํโฆ โอสาเรติ, โทสาริโต. ปณฺฑโก ภิกฺขเว อปฺปตฺโต โอสารณํ, ตญฺเจ สํโฆ โอสาเรติ, โทสาริโต. เถยฺย\-สํวาสโก ภิกฺขเว อปฺปตฺโต โอสารณํ, ตญฺเจ สํโฆ โอสาเรติ, โทสาริโต. ติตฺถิย\-ปกฺกนฺตโก ภิกฺขเว ฯเปฯ. ติรจฺฉานคโต ภิกฺขเว. มาตุฆาตโก ภิกฺขเว. ปิตุฆาตโก ภิกฺขเว. อรหนฺต\-ฆาตโก ภิกฺขเว. ภิกฺขุนิ\-ทูสโก ภิกฺขเว. สํฆเภทโก ภิกฺขเว. โลหิตุปฺปาทโก ภิกฺขเว. อุภโต\-พฺยญฺชนโก ภิกฺขเว อปฺปตฺโต โอสารณํ, ตญฺเจ สํโฆ โอสาเรติ, โทสาริโต. อยํ วุจฺจติ ภิกฺขเว ปุคฺคโล อปฺปตฺโต โอสารณํ, ตญฺเจ สํโฆ โอสาเรติ, โทสาริโต. อิเม วุจฺจนฺติ ภิกฺขเว ปุคฺคลา อปฺปตฺตา โอสารณํ, เต เจ สํโฆ โอสาเรติ, โทสาริตา.}

\prose{กตโม จ ภิกฺขเว ปุคฺคโล อปฺปตฺโต โอสารณํ, ตญฺเจ สํโฆ โอสาเรติ, โสสาริโต.}

\prose{หตฺถจฺฉินฺโน ภิกฺขเว อปฺปตฺโต โอสารณํ, ตญฺเจ สํโฆ โอสาเรติ, โสสาริโต.}

\prose{ปาทจฺฉินฺโน ภิกฺขเว ฯเปฯ. หตฺถ\-ปาท\-จฺฉินฺโน ภิกฺขเว. กณฺณจฺฉินฺโน ภิกฺขเว. นาสจฺฉินฺโน ภิกฺขเว. กณฺณ\-นาส\-จฺฉินฺโน ภิกฺขเว. องฺคุลิจฺฉินฺโน ภิกฺขเว. อฬจฺฉินฺโน ภิกฺขเว. กณฺฑรจฺฉินฺโน ภิกฺขเว. ผณหตฺถโก ภิกฺขเว. ขุชฺโช ภิกฺขเว. วามโน ภิกฺขเว. คลคณฺฑี ภิกฺขเว. ลกฺขณาหโต ภิกฺขเว. กสาหโต ภิกฺขเว. ลิขิตโก ภิกฺขเว. สีปทิโก ภิกฺขเว. ปาปโรคี ภิกฺขเว. ปริสทูสโก ภิกฺขเว. กาโณ ภิกฺขเว. กุณี ภิกฺขเว. ขญฺโช ภิกฺขเว. ปกฺขหโต ภิกฺขเว. ฉินฺนิริยาปโถ ภิกฺขเว. ชราทุพฺพโล ภิกฺขเว. อนฺโธ ภิกฺขเว. มูโค ภิกฺขเว. พธิโร ภิกฺขเว. อนฺธมูโค ภิกฺขเว. อนฺธพธิโร ภิกฺขเว. มูคพธิโร ภิกฺขเว. อนฺธ\-มูค\-พธิโร ภิกฺขเว อปฺปตฺโต โอสารณํ, ตญฺเจ สํโฆ โอสาเรติ, โสสาริโต. อยํ วุจฺจติ ภิกฺขเว ปุคฺคโล อปฺปตฺโต โอสารณํ, ตญฺเจ สํโฆ โอสาเรติ, โสสาริโต. อิเม วุจฺจนฺติ ภิกฺขเว ปุคฺคลา อปฺปตฺตา โอสารณํ, เต เจ สํโฆ โอสาเรติ, โสสาริตา.}

\setlength{\csromangathapagewidth}{0pt}
\setlength{\csromangathawakwidth}{0pt}
\csromangathameasuregroup{}{วาสภคามภาณวาโร นิฏฺฐิโต ปฐโม.}
\csromangathameasurewak{วาสภคามภาณวาโร นิฏฺฐิโต ปฐโม.}
\setlength{\csromangathaleftcol}{0pt}
\csromangathagroup{\csromangathastanza{\csromangathawak{วาสภคาม\-ภาณวาโร นิฏฺฐิโต ปฐโม.}}}

\csromanmatikaanchor{mtk.249}
\csromantocmarkref{section}{240. อธมฺมกมฺมาทิกถา}{mtk.249}
\bookmark[dest={mtk.249},level=1]{240. อธมฺมกมฺมาทิกถา}
\csromanheader{1}{240. \textbf{อธมฺมกมฺมา}\-\textbf{ทิ}\-\textbf{กถา}}
\proseitem{397}{\csromanfolio{445}อิธ ปน ภิกฺขเว ภิกฺขุสฺส น โหติ อาปตฺติ ทฏฺฐพฺพา, ตเมนํ โจเทติ สํโฆ วา สมฺพหุลา วา เอกปุคฺคโล วา “อาปตฺติํ ตฺวํ อาวุโส อาปนฺโน, ปสฺสเสตํ อาปตฺตินฺ”ติ, โส เอวํ วเทติ “นตฺถิ เม อาวุโส อาปตฺติ, ยมหํ ปสฺเสยฺยนฺ”ติ, ตํ สํโฆ อาปตฺติยา อทสฺสเน อุกฺขิปติ, อธมฺมกมฺมํ.}

\prose{อิธ ปน ภิกฺขเว ภิกฺขุสฺส น โหติ อาปตฺติ ปฏิกาตพฺพา, ตเมนํ โจเทติ สํโฆ วา สมฺพหุลา วา เอกปุคฺคโล วา “อาปตฺติํ ตฺวํ อาวุโส อาปนฺโน, ปฏิกโรหิ ตํ อาปตฺตินฺ”ติ, โส เอวํ วเทติ “นตฺถิ เม อาวุโส อาปตฺติ, ยมยํ ปฏิกเรยฺยนฺ”ติ, ตํ สํโฆ อาปตฺติยา อปฺปฏิกมฺเม อุกฺขิปติ, อธมฺมกมฺมํ.}

\prose{อิธ ปน ภิกฺขเว ภิกฺขุสฺส น โหติ ปาปิกา ทิฏฺฐิ ปฏินิสฺสชฺเชตา, ตเมนํ โจเทติ สํโฆ วา สมฺพหุลา วา เอกปุคฺคโล วา “ปาปิกา เต อาวุโส ทิฏฺฐิ, ปฏินิสฺสชฺเชตํ ปาปิกํ ทิฏฺฐินฺ”ติ, โส เอวํ วเทติ “นตฺถิ เม อาวุโส ปาปิกา ทิฏฺฐิ, ยมหํ ปฏินิสฺสชฺเชยฺยนฺ”ติ, ตํ สํโฆ ปาปิกาย ทิฏฺฐิยา อปฺปฏินิสฺสคฺเค อุกฺขิปติ, อธมฺมกมฺมํ.}

\prose{อิธ ปน ภิกฺขเว ภิกฺขุสฺส น โหติ อาปตฺติ ทฏฺฐพฺพา, น โหติ อาปตฺติ ปฏิกาตพฺพา, ตเมนํ โจเทติ สํโฆ วา สมฺพหุลา วา เอกปุคฺคโล วา “อาปตฺติํ ตฺวํ อาวุโส อาปนฺโน, ปสฺสเสตํ อาปตฺติํ, ปฏิกโรหิ ตํ อาปตฺตินฺ”ติ, โส เอวํ วเทติ “นตฺถิ เม อาวุโส อาปตฺติ, ยมหํ ปสฺเสยฺยํ, นตฺถิ เม อาวุโส อาปตฺติ, ยมหํ ปฏิกเรยฺยนฺ”ติ, ตํ สํโฆ อทสฺสเน วา อปฺปฏิกมฺเม วา อุกฺขิปติ, อธมฺมกมฺมํ.}

\prose{อิธ ปน ภิกฺขเว ภิกฺขุสฺส น โหติ อาปตฺติ ทฏฺฐพฺพา, น โหติ ปาปิกา ทิฏฺฐิ ปฏินิสฺสชฺเชตา, ตเมนํ โจเทติ สํโฆ วา สมฺพหุลา วา เอกปุคฺคโล วา “อาปตฺติํ ตฺวํ อาวุโส อาปนฺโน, ปสฺสเสตํ อาปตฺติํ, ปาปิกา เต ทิฏฺฐิ, ปฏินิสฺสชฺเชตํ ปาปิกํ ทิฏฺฐินฺ”ติ, โส เอวํ วเทติ “นตฺถิ เม อาวุโส อาปตฺติ, ยมหํ ปสฺเสยฺยํ, นตฺถิ เม อาวุโส ปาปิกา ทิฏฺฐิ, ยมหํ ปฏินิสฺสชฺเชยฺยนฺ”ติ, ตํ สํโฆ อทสฺสเน วา อปฺปฏินิสฺสคฺเค วา อุกฺขิปติ, อธมฺมกมฺมํ.}

\prose{\csromanfolio{446}อิธ ปน ภิกฺขเว ภิกฺขุสฺส น โหติ อาปตฺติ ปฏิกาตพฺพา, น โหติ ปาปิกา ทิฏฺฐิ ปฏินิสฺสชฺเชตา, ตเมนํ โจเทติ สํโฆ วา สมฺพหุลา วา เอกปุคฺคโล วา “อาปตฺติํ ตฺวํ อาวุโส อาปนฺโน, ปฏิกโรหิ ตํ อาปตฺติํ, ปาปิกา เต ทิฏฺฐิ, ปฏินิสฺสชฺเชตํ ปาปิกํ ทิฏฺฐินฺ”ติ, โส เอวํ วเทติ “นตฺถิ เม อาวุโส อาปตฺติ, ยมหํ ปฏิกเรยฺยํ, นตฺถิ เม อาวุโส ปาปิกา ทิฏฺฐิ, ยมหํ ปฏินิสฺสชฺเชยฺยนฺ”ติ, ตํ สํโฆ อปฺปฏิกมฺเม วา อปฺปฏินิสฺสคฺเค วา อุกฺขิปติ, อธมฺมกมฺมํ.}

\prose{อิธ ปน ภิกฺขเว ภิกฺขุสฺส น โหติ อาปตฺติ ทฏฺฐพฺพา, น โหติ อาปตฺติ ปฏิกาตพฺพา, น โหติ ปาปิกา ทิฏฺฐิ ปฏินิสฺสชฺเชตา, ตเมนํ โจเทติ สํโฆ วา สมฺพหุลา วา เอกปุคฺคโล วา “อาปตฺติํ ตฺวํ อาวุโส อาปนฺโน, ปสฺสเสตํ อาปตฺติํ, ปฏิกโรหิ ตํ อาปตฺติํ, ปาปิกา เต ทิฏฺฐิ, ปฏินิสฺสชฺเชตํ ปาปิกํ ทิฏฺฐินฺ”ติ, โส เอวํ วเทติ “นตฺถิ เม อาวุโส อาปตฺติ, ยมหํ ปสฺเสยฺยํ, นตฺถิ เม อาวุโส อาปตฺติ, ยมหํ ปฏิกเรยฺยํ, นตฺถิ เม อาวุโส ปาปิกา ทิฏฺฐิ, ยมหํ ปฏินิสฺสชฺเชยฺยนฺ”ติ, ตํ สํโฆ อทสฺสเน วา อปฺปฏิกมฺเม วา อปฺปฏินิสฺสคฺเค วา อุกฺขิปติ, อธมฺมกมฺมํ.}

\proseitem{398}{อิธ ปน ภิกฺขเว ภิกฺขุสฺส โหติ อาปตฺติ ทฏฺฐพฺพา, ตเมนํ โจเทติ สํโฆ วา สมฺพหุลา วา เอกปุคฺคโล วา “อาปตฺติํ ตฺวํ อาวุโส อาปนฺโน, ปสฺสเสตํ อาปตฺตินฺ”ติ, โส เอวํ วเทติ “อามาวุโส ปสฺสามี”ติ, ตํ สํโฆ อาปตฺติยา อทสฺสเน อุกฺขิปติ, อธมฺมกมฺมํ.}

\prose{อิธ ปน ภิกฺขเว ภิกฺขุสฺส โหติ อาปตฺติ ปฏิกาตพฺพา, ตเมนํ โจเทติ สํโฆ วา สมฺพหุลา วา เอกปุคฺคโล วา “อาปตฺติํ ตฺวํ อาวุโส อาปนฺโน, ปฏิกโรหิ ตํ อาปตฺตินฺ”ติ, โส เอวํ วเทติ “อามาวุโส ปฏิกริสฺสามี”ติ, ตํ สํโฆ อาปตฺติยา อปฺปฏิกมฺเม อุกฺขิปติ, อธมฺมกมฺมํ.}

\proseitem{398}{อิธ ปน ภิกฺขเว ภิกฺขุสฺส โหติ ปาปิกา ทิฏฺฐิ ปฏินิสฺสชฺเชตา, ตเมนํ โจเทติ สํโฆ วา สมฺพหุลา วา เอกปุคฺคโล วา “ปาปิกา เต อาวุโส ทิฏฺฐิ, ปฏินิสฺสชฺเชตํ ปาปิกํ ทิฏฺฐินฺ”ติ, โส เอวํ วเทติ \csromanfolio{447}“อามาวุโส ปฏินิสฺสชฺชิสฺสามี”ติ, ตํ สํโฆ ปาปิกาย ทิฏฺฐิยา อปฺปฏินิสฺสคฺเค อุกฺขิปติ, อธมฺมกมฺมํ.}

\prose{อิธ ปน ภิกฺขเว ภิกฺขุสฺส โหติ อาปตฺติ ทฏฺฐพฺพา, โหติ อาปตฺติ ปฏิกาตพฺพา ฯเปฯ โหติ อาปตฺติ ทฏฺฐพฺพา, โหติ ปาปิกา ทิฏฺฐิ ปฏินิสฺสชฺเชตา ฯเปฯ โหติ อาปตฺติ ปฏิกาตพฺพา, โหติ ปาปิกา ทิฏฺฐิ ปฏินิสฺสชฺเชตา ฯเปฯ โหติ อาปตฺติ ทฏฺฐพฺพา, โหติ อาปตฺติ ปฏิกาตพฺพา, โหติ ปาปิกา ทิฏฺฐิ ปฏินิสฺสชฺเชตา, ตเมนํ โจเทติ สํโฆ วา สมฺพหุลา วา เอกปุคฺคโล วา “อาปตฺติํ ตฺวํ อาวุโส อาปนฺโน, ปสฺสเสตํ อาปตฺติํ, ปฏิกโรหิ ตํ อาปตฺติํ, ปาปิกา เต ทิฏฺฐิ, ปฏินิสฺสชฺเชตํ ปาปิกํ ทิฏฺฐินฺ”ติ, โส เอวํ วเทติ “อามาวุโส ปสฺสามิ, อาม ปฏิกริสฺสามิ, อาม ปฏินิสฺสชฺชิสฺสามี”ติ, ตํ สํโฆ อทสฺสเน วา อปฺปฏิกมฺเม วา อปฺปฏินิสฺสคฺเค วา อุกฺขิปติ, อธมฺมกมฺมํ.}

\proseitem{399}{อิธ ปน ภิกฺขเว ภิกฺขุสฺส โหติ อาปตฺติ ทฏฺฐพฺพา, ตเมนํ โจเทติ สํโฆ วา สมฺพหุลา วา เอกปุคฺคโล วา “อาปตฺติํ ตฺวํ อาวุโส อาปนฺโน, ปสฺสเสตํ อาปตฺตินฺ”ติ, โส เอวํ วเทติ “นตฺถิ เม อาวุโส อาปตฺติ, ยมหํ ปสฺเสยฺยนฺ”ติ, ตํ สํโฆ อาปตฺติยา อทสฺสเน อุกฺขิปติ, ธมฺมกมฺมํ.}

\prose{อิธ ปน ภิกฺขเว ภิกฺขุสฺส โหติ อาปตฺติ ปฏิกาตพฺพา, ตเมนํ โจเทติ สํโฆ วา สมฺพหุลา วา เอกปุคฺคโล วา “อาปตฺติํ ตฺวํ อาวุโส อาปนฺโน, ปฏิกโรหิ ตํ อาปตฺตินฺ”ติ, โส เอวํ วเทติ “นตฺถิ เม อาวุโส อาปตฺติ, ยมหํ ปฏิกเรยฺยนฺ”ติ, ตํ สํโฆ อาปตฺติยา อปฺปฏิกมฺเม อุกฺขิปติ, ธมฺมกมฺมํ.}

\prose{อิธ ปน ภิกฺขเว ภิกฺขุสฺส โหติ ปาปิกา ทิฏฺฐิ ปฏินิสฺสชฺเชตา, ตเมนํ โจเทติ สํโฆ วา สมฺพหุลา วา เอกปุคฺคโล วา “ปาปิกา เต อาวุโส ทิฏฺฐิ, ปฏินิสฺสชฺเชตํ ปาปิกํ ทิฏฺฐินฺ”ติ, โส เอวํ วเทติ “นตฺถิ เม อาวุโส ปาปิกา ทิฏฺฐิ, ยมหํ ปฏินิสฺสชฺเชยฺยนฺ”ติ, ตํ สํโฆ ปาปิกาย ทิฏฺฐิยา อปฺปฏินิสฺสคฺเค อุกฺขิปติ, ธมฺมกมฺมํ.}

\proseitem{399}{อิธ ปน ภิกฺขเว ภิกฺขุสฺส โหติ อาปตฺติ ทฏฺฐพฺพา, โหติ อาปตฺติ ปฏิกาตพฺพา ฯเปฯ โหติ อาปตฺติ ทฏฺฐพฺพา, โหติ ปาปิกา ทิฏฺฐิ ปฏินิสฺสชฺเชตา ฯเปฯ โหติ อาปตฺติ ปฏิกาตพฺพา, โหติ ปาปิกา ทิฏฺฐิ \csromanfolio{448}ปฏินิสฺสชฺเชตา ฯเปฯ โหติ อาปตฺติ ทฏฺฐพฺพา, โหติ อาปตฺติ ปฏิกาตพฺพา, โหติ ปาปิกา ทิฏฺฐิ ปฏินิสฺสชฺเชตา, ตเมนํ โจเทติ สํโฆ วา สมฺพหุลา วา เอกปุคฺคโล วา “อาปตฺติํ ตฺวํ อาวุโส อาปนฺโน, ปสฺสเสตํ อาปตฺติํ, ปฏิกโรหิ ตํ อาปตฺติํ, ปาปิกา เต ทิฏฺฐิ, ปฏินิสฺสชฺเชตํ ปาปิกํ ทิฏฺฐินฺ”ติ, โส เอวํ วเทติ “นตฺถิเม อาวุโส อาปตฺติ, ยมหํ ปสฺเสยฺยํ, นตฺถิ เม อาวุโส อาปตฺติ, ยมหํ ปฏิกเรยฺยํ, นตฺถิ เม อาวุโส ปาปิกา ทิฏฺฐิ, ยมหํ ปฏินิสฺสชฺเชยฺยนฺ”ติ, ตํ สํโฆ อทสฺสเน วา อปฺปฏิกมฺเม วา อปฺปฏินิสฺสคฺเค วา อุกฺขิปติ, ธมฺมกมฺมนฺติ.}

\csromanmatikaanchor{mtk.250}
\csromantocmarkref{section}{241. อุปาลิปุจฺฉากถา}{mtk.250}
\bookmark[dest={mtk.250},level=1]{241. อุปาลิปุจฺฉากถา}
\csromanheader{1}{241. \textbf{อุปาลิ}\-\textbf{ปุจฺฉา}\-\textbf{กถา}}
\proseitem{400}{อถ โข อายสฺมา อุปาลิ เยน ภควา เตนุปสงฺกมิ, อุปสงฺกมิตฺวา ภควนฺตํ อภิวาเทตฺวา เอกมนฺตํ นิสีทิ, เอกมนฺตํ นิสินฺโน โข อายสฺมา อุปาลิ ภควนฺตํ เอตทโวจ “โย นุ โข ภนฺเต สมคฺโค สํโฆ สมฺมุขา\-กรณียํ กมฺมํ อสมฺมุขา กโรติ, ธมฺมกมฺมํ นุ โข ตํ ภนฺเต วินยกมฺมนฺ”ติ. อธมฺมกมฺมํ ตํ อุปาลิ อวินย\-กมฺมนฺติ. โย นุ โข ภนฺเต สมคฺโค สํโฆ ปฏิปุจฺฉา\-กรณียํ กมฺมํ อปฺปฏิปุจฺฉา กโรติ ฯเปฯ ปฏิญฺญาย\-กรณียํ กมฺมํ อปฏิญฺญาย กโรติ. สติวินยา\-รหสฺส อมูฬฺหวินยํ เทติ. อมูฬฺหวินยา\-รหสฺส ตสฺสปาปิยสิกา\-กมฺมํ กโรติ. ตสฺสปาปิยสิกา\-กมฺมา\-รหสฺส ตชฺชนีย\-กมฺมํ กโรติ. ตชฺชนียกมฺมา\-รหสฺส นิยสฺสกมฺมํ กโรติ. นิยสฺสกมฺมา\-รหสฺส ปพฺพาชนีย\-กมฺมํ กโรติ. ปพฺพาชนียกมฺมา\-รหสฺส ปฏิสารณีย\-กมฺมํ กโรติ. ปฏิสารณียกมฺมา\-รหสฺส อุกฺเขปนีย\-กมฺมํ กโรติ. อุกฺเขปนียกมฺมา\-รหสฺส ปริวาสํ เทติ. ปริวาสารหํ มูลาย ปฏิกสฺสติ. มูลาย\-ปฏิกสฺสนา\-รหสฺส มานตฺตํ เทติ. มานตฺตารหํ อพฺเภติ. อพฺภานารหํ อุปสมฺปาเทติ. ธมฺมกมฺมํ นุ โข ตํ ภนฺเต วินยกมฺมนฺติ. อธมฺมกมฺมํ ตํ อุปาลิ อวินยกมฺมํ.}

\prose{โย โข อุปาลิ สมคฺโค สํโฆ สมฺมุขา\-กรณียํ กมฺมํ อสมฺมุขา กโรติ, เอวํ โข อุปาลิ อธมฺมกมฺมํ โหติ อวินยกมฺมํ, เอวญฺจ ปน สํโฆ สาติสาโร โหติ. โย โข อุปาลิ สมคฺโค สํโฆ ปฏิปุจฺฉา\-กรณียํ กมฺมํ อปฺปฏิปุจฺฉา กโรติ ฯเปฯ ปฏิญฺญาย\-กรณียํ \csromanfolio{449}กมฺมํ อปฏิญฺญาย กโรติ. สติวินยา\-รหสฺส อมูฬฺหวินยํ เทติ. อมูฬฺหวินยา\-รหสฺส ตสฺสปาปิยสิกา\-กมฺมํ กโรติ. ตสฺสปาปิยสิกา\-กมฺมา\-รหสฺส ตชฺชนีย\-กมฺมํ กโรติ. ตชฺชนียกมฺมา\-รหสฺส นิยสฺสกมฺมํ กโรติ. นิยสฺสกมฺมา\-รหสฺส ปพฺพาชนีย\-กมฺมํ กโรติ. ปพฺพาชนียกมฺมา\-รหสฺส ปฏิสารณีย\-กมฺมํ กโรติ. ปฏิสารณียกมฺมา\-รหสฺส อุกฺเขปนีย\-กมฺมํ กโรติ. อุกฺเขปนียกมฺมา\-รหสฺส ปริวาสํ เทติ. ปริวาสารหํ มูลาย ปฏิกสฺสติ. มูลาย\-ปฏิกสฺสนา\-รหสฺส มานตฺตํ เทติ. มานตฺตารหํ อพฺเภติ. อพฺภานารหํ อุปสมฺปาเทติ, เอวํ โข อุปาลิ อธมฺมกมฺมํ โหติ อวินยกมฺมํ, เอวญฺจ ปน สํโฆ สาติสาโร โหตีติ.}

\proseitem{401}{โย นุ โข ภนฺเต สมคฺโค สํโฆ สมฺมุขา\-กรณียํ กมฺมํ สมฺมุขา กโรติ, ธมฺมกมฺมํ นุ โข ตํ ภนฺเต วินยกมฺมนฺติ. ธมฺมกมฺมํ ตํ อุปาลิ วินยกมฺมนฺติ. โย นุ โข ภนฺเต สมคฺโค สํโฆ ปฏิปุจฺฉา\-กรณียํ กมฺมํ ปฏิปุจฺฉา กโรติ ฯเปฯ ปฏิญฺญาย\-กรณียํ กมฺมํ ปฏิญฺญาย กโรติ. สติวินยา\-รหสฺส สติวินยํ เทติ. อมูฬฺหวินยา\-รหสฺส อมูฬฺหวินยํ เทติ. ตสฺสปาปิยสิกา\-กมฺมา\-รหสฺส ตสฺสปาปิยสิกา\-กมฺมํ กโรติ. ตชฺชนียกมฺมา\-รหสฺส ตชฺชนีย\-กมฺมํ กโรติ. นิยสฺสกมฺมา\-รหสฺส นิยสฺสกมฺมํ กโรติ. ปพฺพาชนียกมฺมา\-รหสฺส ปพฺพาชนีย\-กมฺมํ กโรติ. ปฏิสารณียกมฺมา\-รหสฺส ปฏิสารณีย\-กมฺมํ กโรติ. อุกฺเขปนียกมฺมา\-รหสฺส อุกฺเขปนีย\-กมฺมํ กโรติ. ปริวาสา\-รหสฺส ปริวาสํ เทติ. มูลาย ปฏิกสฺสนารหํ มูลาย ปฏิกสฺสติ. มานตฺตารหสฺส มานตฺตํ เทติ. อพฺภานารหํ อพฺเภติ. อุปสมฺปทารหํ อุปสมฺปาเทติ, ธมฺมกมฺมํ นุ โข ตํ ภนฺเต วินยกมฺมนฺติ. ธมฺมกมฺมํ ตํ อุปาลิ วินยกมฺมํ.}

\prose{โย โข อุปาลิ สมคฺโค สํโฆ สมฺมุขา\-กรณียํ กมฺมํ สมฺมุขา กโรติ, เอวํ โข อุปาลิ ธมฺมกมฺมํ โหติ วินยกมฺมํ, เอวญฺจ ปน สํโฆ อนติสาโร โหติ. โย โข อุปาลิ สมคฺโค สํโฆ ปฏิปุจฺฉา\-กรณียํ กมฺมํ ปฏิปุจฺฉา กโรติ. ปฏิญฺญาย\-กรณียํ กมฺมํ ปฏิญฺญาย กโรติ. สติวินยา\-รหสฺส สติวินยํ เทติ. อมูฬฺหวินยา\-รหสฺส อมูฬฺหวินยํ เทติ. ตสฺสปาปิยสิกา\-กมฺมา\-รหสฺส ตสฺสปาปิยสิกา\-กมฺมํ กโรติ. ตชฺชนียกมฺมา\-รหสฺส ตชฺชนีย\-กมฺมํ กโรติ. นิยสฺสกมฺมา\-รหสฺส นิยสฺสกมฺมํ กโรติ. ปพฺพาชนียกมฺมา\-รหสฺส ปพฺพาชนีย\-กมฺมํ \csromanfolio{450}กโรติ. ปฏิสารณียกมฺมา\-รหสฺส ปฏิสารณีย\-กมฺมํ กโรติ. อุกฺเขปนียกมฺมา\-รหสฺส อุกฺเขปนีย\-กมฺมํ กโรติ. ปริวาสา\-รหสฺส ปริวาสํ เทติ. มูลาย\-ปฏิกสฺสนารหํ มูลาย ปฏิกสฺสติ. มานตฺตารหสฺส มานตฺตํ เทติ. อพฺภานารหํ อพฺเภติ. อุปสมฺปทารหํ อุปสมฺปาเทติ, เอวํ โข อุปาลิ ธมฺมกมฺมํ โหติ วินยกมฺมํ, เอวญฺจ ปน สํโฆ อนติสาโร โหตีติ.}

\proseitem{402}{โย นุ โข ภนฺเต สมคฺโค สํโฆ สติวินยา\-รหสฺส อมูฬฺหวินยํ เทติ, อมูฬฺหวินยา\-รหสฺส สติวินยํ เทติ, ธมฺมกมฺมํ นุ โข ตํ ภนฺเต วินยกมฺมนฺติ. อธมฺมกมฺมํ ตํ อุปาลิ อวินย\-กมฺมนฺติ. โย นุ โข ภนฺเต สมคฺโค สํโฆ อมูฬฺหวินยา\-รหสฺส ตสฺสปาปิยสิกา\-กมฺมํ กโรติ, ตสฺสปาปิยสิกา\-กมฺมา\-รหสฺส อมูฬฺหวินยํ เทติ ฯเปฯ ตสฺสปาปิยสิกา\-กมฺมา\-รหสฺส ตชฺชนีย\-กมฺมํ กโรติ, ตชฺชนียกมฺมา\-รหสฺส ตสฺสปาปิยสิกา\-กมฺมํ กโรติ. ตชฺชนียกมฺมา\-รหสฺส นิยสฺสกมฺมํ กโรติ, นิยสฺสกมฺมา\-รหสฺส ตชฺชนีย\-กมฺมํ กโรติ. นิยสฺสกมฺมา\-รหสฺส ปพฺพาชนีย\-กมฺมํ กโรติ, ปพฺพาชนียกมฺมา\-รหสฺส นิยสฺสกมฺมํ กโรติ. ปพฺพาชนียกมฺมา\-รหสฺส ปฏิสารณีย\-กมฺมํ กโรติ, ปฏิสารณียกมฺมา\-รหสฺส ปพฺพาชนีย\-กมฺมํ กโรติ. ปฏิสารณียกมฺมา\-รหสฺส อุกฺเขปนีย\-กมฺมํ กโรติ, อุกฺเขปนียกมฺมา\-รหสฺส ปฏิสารณีย\-กมฺมํ กโรติ. อุกฺเขปนียกมฺมา\-รหสฺส ปริวาสํ เทติ, ปริวาสา\-รหสฺส อุกฺเขปนีย\-กมฺมํ กโรติ. ปริวาสารหํ มูลาย ปฏิกสฺสติ, มูลาย ปฏิกสฺสนา\-รหสฺส ปริวาสํ เทติ. มูลาย ปฏิกสฺสนา\-รหสฺส มานตฺตํ เทติ, มานตฺตารหํ มูลาย ปฏิกสฺสติ. มานตฺตารหํ อพฺเภติ, อพฺภานารหสฺส มานตฺตํ เทติ. อพฺภานารหํ อุปสมฺปาเทติ, อุปสมฺปทารหํ อพฺเภติ, ธมฺมกมฺมํ นุ โข ตํ ภนฺเต วินยกมฺมนฺติ. อธมฺมกมฺมํ ตํ อุปาลิ อวินย\-กมฺมนฺติ.}

\prose{โย โข อุปาลิ สมคฺโค สํโฆ สติวินยา\-รหสฺส อมูฬฺหวินยํ เทติ, อมูฬฺหวินยา\-รหสฺส สติวินยํ เทติ, เอวํ โข อุปาลิ อธมฺมกมฺมํ โหติ อวินยกมฺมํ, เอวญฺจ ปน สํโฆ สาติสาโร โหติ. โย โข อุปาลิ สมคฺโค สํโฆ อมูฬฺหวินยา\-รหสฺส ตสฺสปาปิยสิกา\-กมฺมํ กโรติ, ตสฺสปาปิยสิกา\-กมฺมา\-รหสฺส อมูฬฺหวินยํ เทติ ฯเปฯ ตสฺสปาปิยสิกา\-กมฺมา\-รหสฺส ตชฺชนีย\-กมฺมํ กโรติ, ตชฺชนียกมฺมา\-รหสฺส ตสฺสปาวิยสิกากมฺมํ กโรติ. ตชฺชนียกมฺมา\-รหสฺส นิยสฺสกมฺมํ \csromanfolio{451}กโรติ, นิยสฺสกมฺมา\-รหสฺส ตชฺชนีย\-กมฺมํ กโรติ. นิยสฺสกมฺมา\-รหสฺส ปพฺพาชนีย\-กมฺมํ กโรติ, ปพฺพาชนียกมฺมา\-รหสฺส นิยสฺสกมฺมํ กโรติ. ปพฺพาชนียกมฺมา\-รหสฺส ปฏิสารณีย\-กมฺมํ กโรติ, ปฏิสารณียกมฺมา\-รหสฺส ปพฺพาชนีย\-กมฺมํ กโรติ. ปฏิสารณียกมฺมา\-รหสฺส อุกฺเขปนีย\-กมฺมํ กโรติ, อุกฺเขปนียกมฺมา\-รหสฺส ปฏิสารณีย\-กมฺมํ กโรติ. อุกฺเขปนียกมฺมา\-รหสฺส ปริวาสํ เทติ, ปริวาสา\-รหสฺส อุกฺเขปนีย\-กมฺมํ กโรติ. ปริวาสารหํ มูลาย ปฏิกสฺสติ, มูลาย\-ปฏิกสฺสนา\-รหสฺส ปริวาสํ เทติ. มูลาย\-ปฏิกสฺสนา\-รหสฺส มานตฺตํ เทติ, มานตฺตารหํ มูลาย ปฏิกสฺสติ. มานตฺตารหํ อพฺเภติ, อพฺภานารหสฺส มานตฺตํ เทติ. อพฺภานารหํ อุปสมฺปาเทติ, อุปสมฺปทารหํ อพฺเภติ, เอวํ โข อุปาลิ อธมฺมกมฺมํ โหติ อวินยกมฺมํ, เอวญฺจ ปน สํโฆ สาติสาโร โหตีติ.}

\proseitem{403}{โย นุ โข ภนฺเต สมคฺโค สํโฆ สติวินยา\-รหสฺส สติวินยํ เทติ, อมูฬฺหวินยา\-รหสฺส อมูฬฺหวินยํ เทติ, ธมฺมกมฺมํ นุ โข ตํ ภนฺเต วินยกมฺมนฺติ. ธมฺมกมฺมํ ตํ อุปาลิ วินยกมฺมนฺติ. โย นุ โข ภนฺเต สมคฺโค สํโฆ อมูฬฺหวินยา\-รหสฺส อมูฬฺหวินยํ เทติ ฯเปฯ ตสฺสปาปิยสิกา\-กมฺมา\-รหสฺส ตสฺสปาปิยสิกา\-กมฺมํ กโรติ ฯเปฯ ตชฺชนียกมฺมา\-รหสฺส ตชฺชนีย\-กมฺมํ กโรติ ฯเปฯ นิยสฺสกมฺมา\-รหสฺส นิยสฺสกมฺมํ กโรติ ฯเปฯ ปพฺพาชนียกมฺมา\-รหสฺส ปพฺพาชนีย\-กมฺมํ กโรติ ฯเปฯ ปฏิสารณียกมฺมา\-รหสฺส ปฏิสารณีย\-กมฺมํ กโรติ ฯเปฯ อุกฺเขปนียกมฺมา\-รหสฺส อุกฺเขปนีย\-กมฺมํ กโรติ ฯเปฯ ปริวาสา\-รหสฺส ปริวาสํ เทติ ฯเปฯ มูลาย\-ปฏิกสฺสนารหํ มูลาย ปฏิกสฺสติ ฯเปฯ มานตฺตารหสฺส มานตฺตํ เทติ ฯเปฯ อพฺภานารหํ อพฺเภติ, อุปสมฺปทารหํ อุปสมฺปาเทติ, ธมฺมกมฺมํ นุ โข ตํ ภนฺเต วินยกมฺมนฺติ. ธมฺมกมฺมํ ตํ อุปาลิ วินยกมฺมํ.}

\prose{โย โข อุปาลิ สมคฺโค สํโฆ สติวินยา\-รหสฺส สติวินยํ เทติ, อมูฬฺหวินยา\-รหสฺส อมูฬฺหวินยํ เทติ, เอวํ โข อุปาลิ ธมฺมกมฺมํ โหติ วินยกมฺมํ, เอวญฺจ ปน สํโฆ อนติสาโร โหติ. โย โข อุปาลิ สมคฺโค สํโฆ อมูฬฺหวินยา\-รหสฺส อมูฬฺหวินยํ เทติ ฯเปฯ ตสฺสปาปิยสิกา\-กมฺมา\-รหสฺส ตสฺสปาปิยสิกา\-กมฺมํ กโรติ ฯเปฯ ตชฺชนียกมฺมา\-รหสฺส ตชฺชนีย\-กมฺมํ กโรติ ฯเปฯ นิยสฺสกมฺมา\-รหสฺส นิยสฺสกมฺมํ กโรติ ฯเปฯ ปพฺพาชนียกมฺมา\-รหสฺส ปพฺพาชนีย\-กมฺมํ กโรติ ฯเปฯ \csromanfolio{452}ปฏิสารณียกมฺมา\-รหสฺส ปฏิสารณีย\-กมฺมํ กโรติ ฯเปฯ อุกฺเขปนียกมฺมา\-รหสฺส อุกฺเขปนีย\-กมฺมํ กโรติ ฯเปฯ ปริวาสา\-รหสฺส ปริวาสํ เทติ ฯเปฯ มูลาย ปฏิกสฺสนารหํ มูลาย ปฏิกสฺสติ ฯเปฯ มานตฺตารหสฺส มานตฺตํ เทติ ฯเปฯ อพฺภานารหํ อพฺเภติ, อุปสมฺปทารหํ อุปสมฺปาเทติ, เอวํ โข อุปาลิ ธมฺมกมฺมํ โหติ วินยกมฺมํ, เอวญฺจ ปน สํโฆ อนติสาโร โหตีติ.}

\proseitem{404}{อถ โข ภควา ภิกฺขู อามนฺเตสิ “โย โข ภิกฺขเว สมคฺโค สํโฆ สติวินยา\-รหสฺส อมูฬฺหวินยํ เทติ, เอวํ โข ภิกฺขเว อธมฺมกมฺมํ โหติ อวินยกมฺมํ, เอวญฺจ ปน สํโฆ สาติสาโร โหติ. โย โข ภิกฺขเว สมคฺโค สํโฆ สติวินยา\-รหสฺส ตสฺสปาปิยสิกา\-กมฺมํ กโรติ ฯเปฯ สติวินยา\-รหสฺส ตชฺชนีย\-กมฺมํ กโรติ. สติวินยา\-รหสฺส นิยสฺสกมฺมํ กโรติ. สติวินยา\-รหสฺส ปพฺพาชนีย\-กมฺมํ กโรติ. สติวินยา\-รหสฺส ปฏิสารณีย\-กมฺมํ กโรติ. สติวินยา\-รหสฺส อุกฺเขปนีย\-กมฺมํ กโรติ. สติวินยา\-รหสฺส ปริวาสํ เทติ. สติวินยารหํ มูลาย ปฏิกสฺสติ. สติวินยา\-รหสฺส มานตฺตํ เทติ. สติวินยารหํ อพฺเภติ. สติวินยารหํ อุปสมฺปาเทติ, เอวํ โข ภิกฺขเว อธมฺมกมฺมํ โหติ อวินยกมฺมํ, เอวญฺจ ปน สํโฆ สาติสาโร โหติ.}

\proseitem{405}{โย โข ภิกฺขเว สมคฺโค สํโฆ อมูฬฺหวินยา\-รหสฺส ตสฺสปาปิยสิกา\-กมฺมํ กโรติ, เอวํ โข ภิกฺขเว อธมฺมกมฺมํ โหติ อวินยกมฺมํ, เอวญฺจ ปน สํโฆ สาติสาโร โหติ. โย โข ภิกฺขเว สมคฺโค สํโฆ อมูฬฺหวินยา\-รหสฺส ตชฺชนีย\-กมฺมํ กโรติ ฯเปฯ อมูฬฺหวินยา\-รหสฺส นิยสฺสกมฺมํ กโรติ. อมูฬฺหวินยา\-รหสฺส ปพฺพาชนีย\-กมฺมํ กโรติ. อมูฬฺหวินยา\-รหสฺส ปฏิสารณีย\-กมฺมํ กโรติ. อมูฬฺหวินยา\-รหสฺส อุกฺเขปนีย\-กมฺมํ กโรติ. อมูฬฺหวินยา\-รหสฺส ปริวาสํ เทติ. อมูฬฺหวินยารหํ มูลาย ปฏิกสฺสติ. อมูฬฺหวินยา\-รหสฺส มานตฺตํ เทติ. อมูฬฺหวินยารหํ อพฺเภติ. อมูฬฺหวินยารหํ อุปสมฺปาเทติ. อมูฬฺหวินยา\-รหสฺส สติวินยํ เทติ, เอวํ โข ภิกฺขเว อธมฺมกมฺมํ โหติ อวินยกมฺมํ, เอวญฺจ ปน สํโฆ สาติสาโร โหติ.}

\proseitemwithcloserruled{406}{\csromanfolio{453}โย โข ภิกฺขเว สมคฺโค สํโฆ ตสฺสปาปิยสิกา\-กมฺมา\-รหสฺส ตชฺชนีย\-กมฺมํ กโรติ ฯเปฯ ตชฺชนียกมฺมา\-รหสฺส. นิยสฺสกมฺมา\-รหสฺส. ปพฺพาชนียกมฺมา\-รหสฺส. ปฏิสารณียกมฺมา\-รหสฺส. อุกฺเขปนียกมฺมา\-รหสฺส. ปริวาสารหํ. มูลาย\-ปฏิกสฺสนา\-รหสฺส. มานตฺตารหํ. อพฺภานารหํ. อุปสมฺปทา\-รหสฺส สติวินยํ เทติ, เอวํ โข ภิกฺขเว อธมฺมกมฺมํ โหติ อวินยกมฺมํ, เอวญฺจ ปน สํโฆ สาติสาโร โหติ. โย โข ภิกฺขเว สมคฺโค สํโฆ อุปสมฺปทา\-รหสฺส อมูฬฺหวินยํ เทติ ฯเปฯ อุปสมฺปทา\-รหสฺส ตสฺสปาปิยสิกา\-กมฺมํ กโรติ. อุปสมฺปทา\-รหสฺส ตชฺชนีย\-กมฺมํ กโรติ. อุปสมฺปทา\-รหสฺส นิยสฺสกมฺมํ กโรติ. อุปสมฺปทา\-รหสฺส ปพฺพาชนีย\-กมฺมํ กโรติ. อุปสมฺปทา\-รหสฺส ปฏิสารณีย\-กมฺมํ กโรติ. อุปสมฺปทา\-รหสฺส อุกฺเขปนีย\-กมฺมํ กโรติ. อุปสมฺปทา\-รหสฺส ปริวาสํ เทติ. อุปสมฺปทารหํ มูลาย ปฏิกสฺสติ. อุปสมฺปทา\-รหสฺส มานตฺตํ เทติ. อุปสมฺปทารหํ อพฺเภติ, เอวํ โข ภิกฺขเว อธมฺมกมฺมํ โหติ อวินยกมฺมํ, เอวญฺจ ปน สํโฆ สาติสาโร โหตีติ.}{อุปาลิ\-ปุจฺฉา\-ภาณวาโร นิฏฺฐิโต ทุติโย.}
\csromanmatikaanchor{mtk.251}
\csromantocmarkref{section}{242. ตชฺชนียกมฺมกถา}{mtk.251}
\bookmark[dest={mtk.251},level=1]{242. ตชฺชนียกมฺมกถา}
\csromanheader{1}{242. \textbf{ตชฺชนียกมฺม}\-\textbf{กถา}}
\proseitem{407}{อิธ ปน ภิกฺขเว ภิกฺขุ ภณฺฑนการโก โหติ กลหการโก วิวาทการโก ภสฺสการโก สํเฆ อธิกรณ\-การโก. ตตฺร เจ ภิกฺขูนํ เอวํ โหติ “อยํ โข อาวุโส ภิกฺขุ ภณฺฑนการโก กลหการโก วิวาทการโก ภสฺสการโก สํเฆ อธิกรณ\-การโก, หนฺทสฺส มยํ ตชฺชนีย\-กมฺมํ กโรมา”ติ, เต ตสฺส ตชฺชนีย\-กมฺมํ กโรนฺติ อธมฺเมน วคฺคา. โส ตมฺหา อาวาสา อญฺญํ อาวาสํ คจฺฉติ, ตตฺถปิ ภิกฺขูนํ เอวํ โหติ “อยํ โข อาวุโส ภิกฺขุ สํเฆน ตชฺชนียกมฺมกโต อธมฺเมน วคฺเคหิ, หนฺทสฺส มยํ ตชฺชนีย\-กมฺมํ กโรมา”ติ, เต ตสฺส ตชฺชนีย\-กมฺมํ กโรนฺติ อธมฺเมน สมคฺคา. โส ตมฺหาปิ อาวาสา อญฺญํ อาวาสํ คจฺฉติ, ตตฺถปิ ภิกฺขูนํ เอวํ โหติ “อยํ โข อาวุโส ภิกฺขุ สํเฆน ตชฺชนียกมฺมกโต อธมฺเมน สมคฺเคหิ, หนฺทสฺส มยํ ตชฺชนีย\-กมฺมํ กโรมา”ติ, เต ตสฺส ตชฺชนีย\-กมฺมํ กโรนฺติ ธมฺเมน วคฺคา. \csromanfolio{454}โส ตมฺหาปิ อาวาสา อญฺญํ อาวาสํ คจฺฉติ, ตตฺถปิ ภิกฺขูนํ เอวํ โหติ “อยํ โข อาวุโส ภิกฺขุ สํเฆน ตชฺชนียกมฺมกโต ธมฺเมน วคฺเคหิ, หนฺทสฺส มยํ ตชฺชนีย\-กมฺมํ กโรมา”ติ, เต ตสฺส ตชฺชนีย\-กมฺมํ กโรนฺติ ธมฺม\-ปติรูปเกน วคฺคา. โส ตมฺหาปิ อาวาสา อญฺญํ อาวาสํ คจฺฉติ, ตตฺถปิ ภิกฺขูนํ เอวํ โหติ “อยํ โข อาวุโส ภิกฺขุ สํเฆน ตชฺชนียกมฺมกโต ธมฺม\-ปติรูปเกน วคฺเคหิ, หนฺทสฺส มยํ ตชฺชนีย\-กมฺมํ กโรมา”ติ, เต ตสฺส ตชฺชนีย\-กมฺมํ กโรนฺติ ธมฺม\-ปติรูปเกน สมคฺคา.}

\proseitem{408}{อิธ ปน ภิกฺขเว ภิกฺขุ ภณฺฑนการโก โหติ กลหการโก วิวาทการโก ภสฺสการโก สํเฆ อธิกรณ\-การโก. ตตฺร เจ ภิกฺขูนํ เอวํ โหติ “อยํ โข อาวุโส ภิกฺขุ ภณฺฑนการโก กลหการโก วิวาทการโก ภสฺสการโก สํเฆ อธิกรณ\-การโก, หนฺทสฺส มยํ ตชฺชนีย\-กมฺมํ กโรมา”ติ, เต ตสฺส ตชฺชนีย\-กมฺมํ กโรนฺติ อธมฺเมน สมคฺคา. โส ตมฺหา อาวาสา อญฺญํ อาวาสํ คจฺฉติ, ตตฺถปิ ภิกฺขูนํ เอวํ โหติ “อยํ โข อาวุโส ภิกฺขุ สํเฆน ตชฺชนียกมฺมกโต อธมฺเมน สมคฺเคหิ, หนฺทสฺส มยํ ตชฺชนีย\-กมฺมํ กโรมา”ติ, เต ตสฺส ตชฺชนีย\-กมฺมํ กโรนฺติ ธมฺเมน วคฺคา. โส ตมฺหาปิ อาวาสา อญฺญํ อาวาสํ คจฺฉติ, ตตฺถปิ ภิกฺขูนํ เอวํ โหติ “อยํ โข อาวุโส ภิกฺขุ สํเฆน ตชฺชนียกมฺมกโต ธมฺเมน วคฺเคหิ, หนฺทสฺส มยํ ตชฺชนีย\-กมฺมํ กโรมา”ติ, เต ตสฺส ตชฺชนีย\-กมฺมํ กโรนฺติ ธมฺม\-ปติรูปเกน วคฺคา. โส ตมฺหาปิ อาวาสา อญฺญํ อาวาสํ คจฺฉติ, ตตฺถปิ ภิกฺขูนํ เอวํ โหติ “อยํ โข อาวุโส ภิกฺขุ สํเฆน ตชฺชนียกมฺมกโต ธมฺม\-ปติรูปเกน วคฺเคหิ, หนฺทสฺส มยํ ตชฺชนีย\-กมฺมํ กโรมา”ติ, เต ตสฺส ตชฺชนีย\-กมฺมํ กโรนฺติ ธมฺม\-ปติรูปเกน สมคฺคา. โส ตมฺหาปิ อาวาสา อญฺญํ อาวาสํ คจฺฉติ, ตตฺถปิ ภิกฺขูนํ เอวํ โหติ “อยํ โข อาวุโส ภิกฺขุ สํเฆน ตชฺชนียกมฺมกโต ธมฺม\-ปติรูปเกน สมคฺเคหิ, หนฺทสฺส มยํ ตชฺชนีย\-กมฺมํ กโรมา”ติ, เต ตสฺส ตชฺชนีย\-กมฺมํ กโรนฺติ อธมฺเมน วคฺคา.}

\proseitem{409}{อิธ ปน ภิกฺขเว ภิกฺขุ ภณฺฑนการโก โหติ กลหการโก วิวาทการโก ภสฺสการโก สํเฆ อธิกรณ\-การโก. ตตฺร เจ ภิกฺขูนํ เอวํ โหติ “อยํ โข อาวุโส ภิกฺขุ \csromanfolio{455}ภณฺฑนการโก กลหการโก วิวาทการโก ภสฺสการโก สํเฆ อธิกรณ\-การโก, หนฺทสฺส มยํ ตชฺชนีย\-กมฺมํ กโรมา”ติ, เต ตสฺส ตชฺชนีย\-กมฺมํ กโรนฺติ ธมฺเมน วคฺคา. โส ตมฺหา อาวาสา อญฺญํ อาวาสํ คจฺฉติ, ตตฺถปิ ภิกฺขูนํ เอวํ โหติ “อยํ โข อาวุโส ภิกฺขุ สํเฆน ตชฺชนียกมฺมกโต ธมฺเมน วคฺเคหิ, หนฺทสฺส มยํ ตชฺชนีย\-กมฺมํ กโรมา”ติ, เต ตสฺส ตชฺชนีย\-กมฺมํ กโรนฺติ ธมฺม\-ปติรูปเกน วคฺคา. โส ตมฺหาปิ อาวาสา อญฺญํ อาวาสํ คจฺฉติ, ตตฺถปิ ภิกฺขูนํ เอวํ โหติ “อยํ โข อาวุโส ภิกฺขุ สํเฆน ตชฺชนียกมฺมกโต ธมฺม\-ปติรูปเกน วคฺเคหิ, หนฺทสฺส มยํ ตชฺชนีย\-กมฺมํ กโรมา”ติ, เต ตสฺส ตชฺชนีย\-กมฺมํ กโรนฺติ ธมฺม\-ปติรูปเกน สมคฺคา. โส ตมฺหาปิ อาวาสา อญฺญํ อาวาสํ คจฺฉติ, ตตฺถปิ ภิกฺขูนํ เอวํ โหติ “อยํ โข อาวุโส ภิกฺขุ\-สํเฆน ตชฺชนียกมฺมกโต ธมฺม\-ปติรูปเกน สมคฺเคหิ, หนฺทสฺส มยํ ตชฺชนีย\-กมฺมํ กโรมา”ติ, เต ตสฺส ตชฺชนีย\-กมฺมํ กโรนฺติ อธมฺเมน วคฺคา. โส ตมฺหาปิ อาวาสา อญฺญํ อาวาสํ คจฺฉติ, ตตฺถปิ ภิกฺขูนํ เอวํ โหติ “อยํ โข อาวุโส ภิกฺขุ สํเฆน ตชฺชนียกมฺมกโต อธมฺเมน วคฺเคหิ, หนฺทสฺส มยํ ตชฺชนีย\-กมฺมํ กโรมา”ติ, เต ตสฺส ตชฺชนีย\-กมฺมํ กโรนฺติ อธมฺเมน สมคฺคา.}

\proseitem{410}{อิธ ปน ภิกฺขเว ภิกฺขุ ภณฺฑนการโก โหติ กลหการโก วิวาทการโก ภสฺสการโก สํเฆ อธิกรณ\-การโก. ตตฺร เจ ภิกฺขูนํ เอวํ โหติ “อยํ โข อาวุโส ภิกฺขุ ภณฺฑนการโก กลหการโก วิวาทการโก ภสฺสการโก สํเฆ อธิกรณ\-การโก, หนฺทสฺส มยํ ตชฺชนีย\-กมฺมํ กโรมา”ติ, เต ตสฺส ตชฺชนีย\-กมฺมํ กโรนฺติ ธมฺม\-ปติรูปเกน วคฺคา. โส ตมฺหา อาวาสา อญฺญํ อาวาสํ คจฺฉติ, ตตฺถปิ ภิกฺขูนํ เอวํ โหติ “อยํ โข อาวุโส ภิกฺขุ สํเฆน ตชฺชนียกมฺมกโต ธมฺม\-ปติรูปเกน วคฺเคหิ, หนฺทสฺส มยํ ตชฺชนีย\-กมฺมํ กโรมา”ติ, เต ตสฺส ตชฺชนีย\-กมฺมํ กโรนฺติ ธมฺม\-ปติรูปเกน สมคฺคา. โส ตมฺหาปิ อาวาสา อญฺญํ อาวาสํ คจฺฉฺติ, ตตฺถปิ ภิกฺขูนํ เอวํ โหติ “อยํ โข อาวุโส ภิกฺขุ สํเฆน ตชฺชนียกมฺมกโต ธมฺม\-ปติรูปเกน สมคฺเคหิ, หนฺทสฺส มยํ ตชฺชนีย\-กมฺมํ กโรมา”ติ, เต ตสฺส ตชฺชนีย\-กมฺมํ กโรนฺติ อธมฺเมน วคฺคา. โส ตมฺหาปิ อาวาสา อญฺญํ อาวาสํ คจฺฉติ, ตตฺถปิ \csromanfolio{456}ภิกฺขูนํ เอวํ โหติ “อยํ โข อาวุโส ภิกฺขุ สํเฆน ตชฺชนียกมฺมกโต อธมฺเมน วคฺเคหิ, หนฺทสฺส มยํ ตชฺชนีย\-กมฺมํ กโรมา”ติ, เต ตสฺส ตชฺชนีย\-กมฺมํ กโรนฺติ อธมฺเมน สมคฺคา. โส ตมฺหาปิ อาวาสา อญฺญํ อาวาสํ คจฺฉติ, ตตฺถปิ ภิกฺขูนํ เอวํ โหติ “อยํ โข อาวุโส ภิกฺขุ สํเฆน ตชฺชนียกมฺมกโต อธมฺเมน สมคฺเคหิ, หนฺทสฺส มยํ ตชฺชนีย\-กมฺมํ กโรมา”ติ, เต ตสฺส ตชฺชนีย\-กมฺมํ กโรนฺติ ธมฺเมน วคฺคา.}

\proseitem{411}{อิธ ปน ภิกฺขเว ภิกฺขุ ภณฺฑนการโก โหติ กลหการโก วิวาทการโก ภสฺสการโก สํเฆ อธิกรณ\-การโก. ตตฺร เจ ภิกฺขูนํ เอวํ โหติ “อยํ โข อาวุโส ภิกฺขุ ภณฺฑนการโก กลหการโก วิวาทการโก ภสฺสการโก สํเฆ อธิกรณ\-การโก, หนฺทสฺส มยํ ตชฺชนีย\-กมฺมํ กโรมา”ติ, เต ตสฺส ตชฺชนีย\-กมฺมํ กโรนฺติ ธมฺม\-ปติรูปเกน สมคฺคา. โส ตมฺหา อาวาสา อญฺญํ อาวาสํ คจฺฉติ, ตตฺถปิ ภิกฺขูนํ เอวํ โหติ “อยํ โข อาวุโส ภิกฺขุ สํเฆน ตชฺชนียกมฺมกโต ธมฺม\-ปติรูปเกน สมคฺเคหิ, หนฺทสฺส มยํ ตชฺชนีย\-กมฺมํ กโรมา”ติ, เต ตสฺส ตชฺชนีย\-กมฺมํ กโรนฺติ อธมฺเมน วคฺคา. โส ตมฺหาปิ อาวาสา อญฺญํ อาวาสํ คจฺฉติ, ตตฺถปิ ภิกฺขูนํ เอวํ โหติ “อยํ โข อาวุโส ภิกฺขุ สํเฆน ตชฺชนียกมฺมกโต อธมฺเมน วคฺเคหิ, หนฺทสฺส มยํ ตชฺชนีย\-กมฺมํ กโรมา”ติ, เต ตสฺส ตชฺชนีย\-กมฺมํ กโรนฺติ อธมฺเมน สมคฺคา. โส ตมฺหาปิ อาวาสา อญฺญํ อาวาสํ คจฺฉติ, ตตฺถปิ ภิกฺขูนํ เอวํ โหติ “อยํ โข อาวุโส ภิกฺขุ สํเฆน ตชฺชนียกมฺมกโต อธมฺเมน สมคฺเคหิ, หนฺทสฺส มยํ ตชฺชนีย\-กมฺมํ กโรมา”ติ, เต ตสฺส ตชฺชนีย\-กมฺมํ กโรนฺติ ธมฺเมน วคฺคา. โส ตมฺหาปิ อาวาสา อญฺญํ อาวาสํ คจฺฉติ, ตตฺถปิ ภิกฺขูนํ เอวํ โหติ “อยํ โข อาวุโส ภิกฺขุ สํเฆน ตชฺชนียกมฺมกโต ธมฺเมน วคฺเคหิ, หนฺทสฺส มยํ ตชฺชนีย\-กมฺมํ กโรมา”ติ, เต ตสฺส ตชฺชนีย\-กมฺมํ กโรนฺติ ธมฺม\-ปติรูปเกน วคฺคา.}

\csromanmatikaanchor{mtk.252}
\csromantocmarkref{section}{243. นิยสฺสกมฺมกถา}{mtk.252}
\bookmark[dest={mtk.252},level=1]{243. นิยสฺสกมฺมกถา}
\csromanheader{1}{243. \textbf{นิยสฺสกมฺม}\-\textbf{กถา}}
\proseitem{412}{อิธ ปน ภิกฺขเว ภิกฺขุ พาโล โหติ อพฺยตฺโต อาปตฺติพหุโล อนปทาโน, คิหิสํสฏฺโฐ วิหรติ อนนุโลมิเกหิ \csromanfolio{457}คิหิสํสคฺเคหิ. ตตฺร เจ ภิกฺขูนํ เอวํ โหติ “อยํ โข อาวุโส ภิกฺขุ พาโล อพฺยตฺโต อาปตฺติพหุโล อนปทาโน, คิหิสํสฏฺโฐ วิหรติ อนนุโลมิเกหิ คิหิสํสคฺเคหิ, หนฺทสฺส มยํ นิยสฺสกมฺมํ กโรมา”ติ, เต ตสฺส นิยสฺสกมฺมํ กโรนฺติ อธมฺเมน วคฺคา. โส ตมฺหา อาวาสา อญฺญํ อาวาสํ คจฺฉติ, ตตฺถปิ ภิกฺขูนํ เอวํ โหติ “อยํ โข อาวุโส ภิกฺขุ \textbf{สํเฆน นิยสฺสกมฺมกโต อธมฺเมน วคฺเคหิ, หนฺทสฺส มยํ} \textbf{นิยสฺสกมฺมํ กโรมา”ติ, เต ตสฺส นิยสฺสกมฺมํ กโรนฺติ อธมฺเมน} สมคฺคา ฯเปฯ ธมฺเมน วคฺคา. ธมฺม\-ปติรูปเกน วคฺคา. ธมฺม\-ปติรูปเกน สมคฺคา ฯเปฯ.}

\vspace{\tipitakaparskip}
\csromancenter{ยถา เหฏฺฐา, ตถา จกฺกํ กาตพฺพํ.}
\csromanmatikaanchor{mtk.253}
\csromantocmarkref{section}{244. ปพฺพาชนียกมฺมกถา}{mtk.253}
\bookmark[dest={mtk.253},level=1]{244. ปพฺพาชนียกมฺมกถา}
\csromanheader{1}{244. \textbf{ปพฺพาชนีย}\-\textbf{กมฺมกถา}}
\proseitem{413}{อิธ ปน ภิกฺขเว ภิกฺขุ กุลทูสโก โหติ ปาปสมาจาโร. ตตฺร เจ ภิกฺขูนํ เอวํ โหติ “อยํ โข อาวุโส ภิกฺขุ กุลทูสโก ปาปสมาจาโร, หนฺทสฺส มยํ ปพฺพาชนีย\-กมฺมํ กโรมา”ติ, เต ตสฺส ปพฺพาชนีย\-กมฺมํ กโรนฺติ อธมฺเมน วคฺคา. โส ตมฺหา อาวาสา อญฺญํ อาวาสํ คจฺฉติ, ตตฺถปิ ภิกฺขูนํ เอวํ โหติ “อยํ โข อาวุโส ภิกฺขุ สํเฆน ปพฺพาชนีย\-กมฺมกโต อธมฺเมน วคฺเคหิ, หนฺทสฺส มยํ ปพฺพาชนีย\-กมฺมํ กโรมา”ติ. เต ตสฺส ปพฺพาชนีย\-กมฺมํ กโรนฺติ อธมฺเมน สมคฺคา ฯเปฯ ธมฺเมน วคฺคา. ธมฺม\-ปติรูปเกน วคฺคา. ธมฺม\-ปติรูปเกน สมคฺคา ฯเปฯ.}

\vspace{\tipitakaparskip}
\csromancenter{จกฺกํ กาตพฺพํ.}
\csromanmatikaanchor{mtk.254}
\csromantocmarkref{section}{245. ปฏิสารณียกมฺมกถา}{mtk.254}
\bookmark[dest={mtk.254},level=1]{245. ปฏิสารณียกมฺมกถา}
\csromanheader{1}{245. \textbf{ปฏิสารณีย}\-\textbf{กมฺมกถา}}
\proseitem{414}{อิธ ปน ภิกฺขเว ภิกฺขุ คิหี อกฺโกสติ ปริภาสติ. ตตฺร เจ ภิกฺขูนํ เอวํ โหติ “อยํ โข อาวุโส ภิกฺขุ คิหี อกฺโกสติ ปริภาสติ, หนฺทสฺส มยํ ปฏิสารณีย\-กมฺมํ กโรมา”ติ, เต ตสฺส ปฏิสารณีย\-กมฺมํ กโรนฺติ อธมฺเมน วคฺคา. โส ตมฺหา อาวาสา อญฺญํ อาวาสํ คจฺฉติ, ตตฺถปิ ภิกฺขูนํ เอวํ โหติ “อยํ โข อาวุโส ภิกฺขุ สํเฆน ปฏิสารณีย\-กมฺมกโต อธมฺเมน วคฺเคหิ, หนฺทสฺส มยํ \csromanfolio{458}\textbf{ปฏิสารณีย}\-\textbf{กมฺมํ กโรมา”ติ, เต ตสฺส ปฏิสารณีย}\-\textbf{กมฺมํ กโรนฺติ} \textbf{อธมฺเมน สมคฺคา ฯเปฯ ธมฺเมน วคฺคา. ธมฺม}\-\textbf{ปติรูปเกน วคฺคา.} ธมฺม\-ปติรูปเกน สมคฺคา ฯเปฯ.}

\vspace{\tipitakaparskip}
\csromancenter{จกฺกํ กาตพฺพํ.}
\csromanmatikaanchor{mtk.255}
\csromantocmarkref{section}{246. อทสฺสเน อุกฺเขปนียกมฺมกถา}{mtk.255}
\bookmark[dest={mtk.255},level=1]{246. อทสฺสเน อุกฺเขปนียกมฺมกถา}
\csromanheader{1}{246. \textbf{อทสฺสเน อุกฺเขปนีย}\-\textbf{กมฺมกถา}}
\proseitem{415}{อิธ ปน ภิกฺขเว ภิกฺขุ อาปตฺติํ อาปชฺชิตฺวา น อิจฺฉติ อาปตฺติํ ปสฺสิตุํ. ตตฺร เจ ภิกฺขูนํ เอวํ โหติ “อยํ โข อาวุโส ภิกฺขุ อาปตฺติํ อาปชฺชิตฺวา น อิจฺฉติ อาปตฺติํ ปสฺสิตุํ, หนฺทสฺส มยํ อาปตฺติยา อทสฺสเน อุกฺเขปนีย\-กมฺมํ กโรมา”ติ, เต ตสฺส อาปตฺติยา อทสฺสเน อุกฺเขปนีย\-กมฺมํ กโรนฺติ อธมฺเมน วคฺคา. โส ตมฺหา อาวาสา อญฺญํ อาวาสํ คจฺฉติ, ตตฺถปิ ภิกฺขูนํ เอวํ โหติ “อยํ โข อาวุโส ภิกฺขุ สํเฆน อาปตฺติยา อทสฺสเน อุกฺเขปนีย\-กมฺมกโต อธมฺเมน วคฺเคหิ, หนฺทสฺส มยํ อาปตฺติยา อทสฺสเน อุกฺเขปนีย\-กมฺมํ กโรมา”ติ, เต ตสฺส อาปตฺติยา อทสฺสเน อุกฺเขปนีย\-กมฺมํ กโรนฺติ อธมฺเมน สมคฺคา ฯเปฯ ธมฺเมน วคฺคา. ธมฺม\-ปติรูปเกน วคฺคา. ธมฺม\-ปติรูปเกน สมคฺคา ฯเปฯ.}

\vspace{\tipitakaparskip}
\csromancenter{จกฺกํ กาตพฺพํ.}
\csromanmatikaanchor{mtk.256}
\csromantocmarkref{section}{247. อปฺปฏิกมฺเม อุกฺเขปนียกมฺมกถา}{mtk.256}
\bookmark[dest={mtk.256},level=1]{247. อปฺปฏิกมฺเม อุกฺเขปนียกมฺมกถา}
\csromanheader{1}{247. \textbf{อปฺปฏิกมฺเม อุกฺเขปนีย}\-\textbf{กมฺมกถา}}
\proseitem{416}{อิธ ปน ภิกฺขเว ภิกฺขุ อาปตฺติํ อาปชฺชิตฺวา น อิจฺฉติ อาปตฺติํ ปฏิกาตุํ. ตตฺร เจ ภิกฺขูนํ เอวํ โหติ “อยํ โข อาวุโส ภิกฺขุ อาปตฺติํ อาปชฺชิตฺวา น อิจฺฉติ อาปตฺติํ ปฏิกาตุํ, หนฺทสฺส มยํ อาปตฺติยา อปฺปฏิกมฺเม อุกฺเขปนีย\-กมฺมํ กโรมา”ติ, เต ตสฺส อาปตฺติยา อปฺปฏิกมฺเม อุกฺเขปนีย\-กมฺมํ กโรนฺติ อธมฺเมน วคฺคา. โส ตมฺหา อาวาสา อญฺญํ อาวาสํ คจฺฉติ, ตตฺถปิ ภิกฺขูนํ เอวํ โหติ “อยํ โข อาวุโส ภิกฺขุ สํเฆน อาปตฺติยา อปฺปฏิกมฺเม อุกฺเขปนีย\-กมฺมกโต อธมฺเมน วคฺเคหิ, หนฺทสฺส มยํ อาปตฺติยา อปฺปฏิกมฺเม อุกฺเขปนีย\-กมฺมํ กโรมา”ติ, เต ตสฺส อาปตฺติยา อปฺปฏิกมฺเม อุกฺเขปนีย\-กมฺมํ กโรนฺติ อธมฺเมน สมคฺคา ฯเปฯ ธมฺเมน วคฺคา. ธมฺม\-ปติรูปเกน วคฺคา. ธมฺม\-ปติรูปเกน สมคฺคา ฯเปฯ.}

\vspace{\tipitakaparskip}
\csromancenter{จกฺกํ กาตพฺพํ.}
\csromanmatikaanchor{mtk.257}
\csromantocmarkref{section}{248. อปฺปฏินิสฺสคฺเค อุกฺเขปนียกมฺมกถา ...}{mtk.257}
\bookmark[dest={mtk.257},level=1]{248. อปฺปฏินิสฺสคฺเค อุกฺเขปนียกมฺมกถา ...}
\csromanheader{1}{248. \textbf{อปฺปฏินิสฺสคฺเค อุกฺเขปนีย}\-\textbf{กมฺมกถา}}
\proseitem{417}{\csromanfolio{459}อิธ ปน ภิกฺขเว ภิกฺขุ น อิจฺฉติ ปาปิกํ ทิฏฺฐิํ ปฏินิสฺสชฺชิตุํ. ตตฺร เจ ภิกฺขูนํ เอวํ โหติ “อยํ โข อาวุโส ภิกฺขุ น อิจฺฉติ ปาปิกํ ทิฏฺฐิํ ปฏินิสฺสชฺชิตุํ, หนฺทสฺส มยํ ปาปิกาย ทิฏฺฐิยา อปฺปฏินิสฺสคฺเค อุกฺเขปนีย\-กมฺมํ กโรมา”ติ, เต ตสฺส ปาปิกาย ทิฏฺฐิยา อปฺปฏินิสฺสคฺเค อุกฺเขปนีย\-กมฺมํ กโรนฺติ อธมฺเมน วคฺคา. โส ตมฺหา อาวาสา อญฺญํ อาวาสํ คจฺฉติ, ตตฺถปิ ภิกฺขูนํ เอวํ โหติ “อยํ โข อาวุโส ภิกฺขุ สํเฆน ปาปิกาย ทิฏฺฐิยา อปฺปฏินิสฺสคฺเค อุกฺเขปนีย\-กมฺมกโต อธมฺเมน วคฺเคหิ, หนฺทสฺส มยํ ปาปิกาย ทิฏฺฐิยา อปฺปฏินิสฺสคฺเค อุกฺเขปนีย\-กมฺมํ กโรมา”ติ, เต ตสฺส ปาปิกาย ทิฏฺฐิยา อปฺปฏินิสฺสคฺเค อุกฺเขปนีย\-กมฺมํ กโรนฺติ อธมฺเมน สมคฺคา - ป- ธมฺเมน วคฺคา. ธมฺม\-ปติรูปเกน วคฺคา. ธมฺม\-ปติรูปเกน สมคฺคา ฯเปฯ.}

\vspace{\tipitakaparskip}
\csromancenter{จกฺกํ กาตพฺพํ.}
\csromanmatikaanchor{mtk.258}
\csromantocmarkref{section}{249. ตชฺชนียกมฺมปฏิปฺปสฺสทฺธิกถา}{mtk.258}
\bookmark[dest={mtk.258},level=1]{249. ตชฺชนียกมฺมปฏิปฺปสฺสทฺธิกถา}
\csromanheader{1}{249. \textbf{ตชฺชนียกมฺม}\-\textbf{ปฏิปฺปสฺสทฺธิ}\-\textbf{กถา}}
\proseitem{418}{อิธ ปน ภิกฺขเว ภิกฺขุ สํเฆน ตชฺชนียกมฺมกโต สมฺมา วตฺตติ, โลมํ ปาเตติ, เนตฺถารํ วตฺตติ, ตชฺชนียสฺส กมฺมสฺส ปฏิปฺปสฺสทฺธิํ ยาจติ. ตตฺร เจ ภิกฺขูนํ เอวํ โหติ “อยํ โข อาวุโส ภิกฺขุ สํเฆน ตชฺชนียกมฺมกโต สมฺมา วตฺตติ, โลมํ ปาเตติ, เนตฺถารํ วตฺตติ, ตชฺชนียสฺส กมฺมสฺส ปฏิปฺปสฺสทฺธิํ ยาจติ, หนฺทสฺส มยํ ตชฺชนีย\-กมฺมํ ปฏิปฺปสฺสมฺเภมา”ติ, เต ตสฺส ตชฺชนีย\-กมฺมํ ปฏิปฺปสฺสมฺเภนฺติ อธมฺเมน วคฺคา. โส ตมฺหา อาวาสา อญฺญํ อาวาสํ คจฺฉติ, ตตฺถปิ ภิกฺขูนํ เอวํ โหติ “อิมสฺส โข อาวุโส ภิกฺขุโน สํเฆน ตชฺชนีย\-กมฺมํ ปฏิปฺปสฺสทฺธํ อธมฺเมน วคฺเคหิ, หนฺทสฺส มยํ ตชฺชนีย\-กมฺมํ ปฏิปฺปสฺสมฺเภมา”ติ, เต ตสฺส ตชฺชนีย\-กมฺมํ ปฏิปฺปสฺสมฺเภนฺติ อธมฺเมน สมคฺคา. โส ตมฺหาปิ อาวาสา อญฺญํ อาวาสํ คจฺฉติ, ตตฺถปิ ภิกฺขูนํ เอวํ โหติ “อิมสฺส โข อาวุโส ภิกฺขุโน สํเฆน ตชฺชนีย\-กมฺมํ ปฏิปฺปสฺสทฺธํ อธมฺเมน สมคฺเคหิ, หนฺทสฺส มยํ ตชฺชนีย\-กมฺมํ ปฏิปฺปสฺสมฺเภมา”ติ, เต ตสฺส ตชฺชนีย\-กมฺมํ ปฏิปฺปสฺสมฺเภนฺติ ธมฺเมน วคฺคา. โส ตมฺหาปิ อาวาสา อญฺญํ อาวาสํ คจฺฉติ, ตตฺถปิ ภิกฺขูนํ เอวํ โหติ “อิมสฺส โข อาวุโส ภิกฺขุโน สํเฆน ตชฺชนีย\-กมฺมํ ปฏิปฺปสฺสทฺธํ ธมฺเมน วคฺเคหิ, \csromanfolio{460}หนฺทสฺส มยํ ตชฺชนีย\-กมฺมํ ปฏิปฺปสฺสมฺเภมา”ติ, เต ตสฺส ตชฺชนีย\-กมฺมํ ปฏิปฺปสฺสมฺเภนฺติ ธมฺม\-ปติรูปเกน วคฺคา. โส ตมฺหาปิ อาวาสา อญฺญํ อาวาสํ คจฺฉติ, ตตฺถปิ ภิกฺขูนํ เอวํ โหติ “อิมสฺส โข อาวุโส ภิกฺขุโน สํเฆน ตชฺชนีย\-กมฺมํ ปฏิปฺปสฺสทฺธํ ธมฺม\-ปติรูปเกน วคฺเคหิ, หนฺทสฺส มยํ ตชฺชนีย\-กมฺมํ ปฏิปฺปสฺสมฺเภมา”ติ, เต ตสฺส ตชฺชนีย\-กมฺมํ ปฏิปฺปสฺสมฺเภนฺติ ธมฺม\-ปติรูปเกน สมคฺคา.}

\proseitem{419}{อิธ ปน ภิกฺขเว ภิกฺขุ สํเฆน ตชฺชนียกมฺมกโต สมฺมา วตฺตติ, โลมํ ปาเตติ, เนตฺถารํ วตฺตติ, ตชฺชนียสฺส กมฺมสฺส ปฏิปฺปสฺสทฺธิํ ยาจติ. ตตฺร เจ ภิกฺขูนํ เอวํ โหติ “อยํ โข อาวุโส ภิกฺขุ สํเฆน ตชฺชนียกมฺมกโต สมฺมา วตฺตติ, โลมํ ปาเตติ, เนตฺถารํ วตฺตติ, ตชฺชนียสฺส กมฺมสฺส ปฏิปฺปสฺสทฺธิํ ยาจติ, หนฺทสฺส มยํ ตชฺชนีย\-กมฺมํ ปฏิปฺปสฺสมฺเภมา”ติ, เต ตสฺส ตชฺชนีย\-กมฺมํ ปฏิปฺปสฺสมฺเภนฺติ อธมฺเมน สมคฺคา. โส ตมฺหา อาวาสา อญฺญํ อาวาสํ คจฺฉติ, ตตฺถปิ ภิกฺขูนํ เอวํ โหติ “อิมสฺส โข อาวุโส ภิกฺขุโน สํเฆน ตชฺชนีย\-กมฺมํ ปฏิปฺปสฺสทฺธํ อธมฺเมน สมคฺเคหิ, หนฺทสฺส มยํ ตชฺชนีย\-กมฺมํ ปฏิปฺปสฺสมฺเภมา”ติ, เต ตสฺส ตชฺชนีย\-กมฺมํ ปฏิปฺปสฺสมฺเภนฺติ ธมฺเมน วคฺคา. โส ตมฺหาปิ อาวาสา อญฺญํ อาวาสํ คจฺฉติ, ตตฺถปิ ภิกฺขูนํ เอวํ โหติ “อิมสฺส โข อาวุโส ภิกฺขุโน สํเฆน ตชฺชนีย\-กมฺมํ ปฏิปฺปสฺสทฺธํ ธมฺเมน วคฺเคหิ, หนฺทสฺส มยํ ตชฺชนีย\-กมฺมํ ปฏิปฺปสฺสมฺเภมา”ติ, เต ตสฺส ตชฺชนีย\-กมฺมํ ปฏิปฺปสฺสมฺเภนฺติ ธมฺม\-ปติรูปเกน วคฺคา. โส ตมฺหาปิ อาวาสา อญฺญํ อาวาสํ คจฺฉติ, ตตฺถปิ ภิกฺขูนํ เอวํ โหติ “อิมสฺส โข อาวุโส ภิกฺขุโน สํเฆน ตชฺชนีย\-กมฺมํ ปฏิปฺปสฺสทฺธํ ธมฺม\-ปติรูปเกน วคฺเคหิ, หนฺทสฺส มยํ ตชฺชนีย\-กมฺมํ ปฏิปฺปสฺสมฺเภมา”ติ, เต ตสฺส ตชฺชนีย\-กมฺมํ ปฏิปฺปสฺสมฺเภนฺติ ธมฺม\-ปติรูปเกน สมคฺคา. โส ตมฺหาปิ อาวาสา อญฺญํ อาวาสํ คจฺฉติ, ตตฺถปิ ภิกฺขูนํ เอวํ โหติ “อิมสฺส โข อาวุโส ภิกฺขุโน สํเฆน ตชฺชนีย\-กมฺมํ ปฏิปฺปสฺสทฺธํ ธมปติรูปเกน สมคฺเคหิ, หนฺทสฺส มยํ ตชฺชนีย\-กมฺมํ ปฏิปฺปสฺสมฺเภมา”ติ, เต ตสฺส ตชฺชนีย\-กมฺมํ ปฏิปฺปสฺสมฺเภนฺติ อธมฺเมน วคฺคา.}

\proseitem{420}{อิธ ปน ภิกฺขเว ภิกฺขุ สํเฆน ตชฺชนียกมฺมกโต สมฺมา วตฺตติ, โลมํ ปาเตติ, เนตฺถารํ วตฺตติ, ตชฺชนียสฺส กมฺมสฺส ปฏิปฺปสฺสทฺธิํ \csromanfolio{461}ยาจติ. ตตฺร เจ ภิกฺขูนํ เอวํ โหติ “อยํ โข อาวุโส ภิกฺขุ สํเฆน ตชฺชนียกมฺมกโต สมฺมา วตฺตติ, โลมํ ปาเตติ, เนตฺถารํ วตฺตติ, ตชฺชนียสฺส กมฺมสฺส ปฏิปฺปสฺสทฺธิํ ยาจติ, หนฺทสฺส มยํ ตชฺชนีย\-กมฺมํ ปฏิปฺปสฺสมฺเภมา”ติ, เต ตสฺส ตชฺชนีย\-กมฺมํ ปฏิปฺปสฺสมฺเภนฺติ ธมฺเมน วคฺคา. โส ตมฺหา อาวาสา อญฺญํ อาวาสํ คจฺฉติ, ตตฺถปิ ภิกฺขูนํ เอวํ โหติ “อิมสฺส โข อาวุโส ภิกฺขุโน สํเฆน ตชฺชนีย\-กมฺมํ ปฏิปฺปสฺสทฺธํ ธมฺเมน วคฺเคหิ, หนฺทสฺส มยํ ตชฺชนีย\-กมฺมํ ปฏิปฺปสฺสมฺเภมา”ติ, เต ตสฺส ตชฺชนีย\-กมฺมํ ปฏิปฺปสฺสมฺเภนฺติ ธมฺม\-ปติรูปเกน วคฺคา. โส ตมฺหาปิ อาวาสา อญฺญํ อาวาสํ คจฺฉติ, ตตฺถปิ ภิกฺขูนํ เอวํ โหติ “อิมสฺส โข อาวุโส ภิกฺขุโน สํเฆน ตชฺชนีย\-กมฺมํ ปฏิปฺปสฺสทฺธํ ธมฺม\-ปติรูปเกน วคฺเคหิ, หนฺทสฺส มยํ ตชฺชนีย\-กมฺมํ ปฏิปฺปสฺสมฺเภมา”ติ เต ตสฺส ตชฺชนีย\-กมฺมํ ปฏิปฺปสฺสมฺเภนฺติ ธมฺม\-ปติรูปเกน สมคฺคา. โส ตมฺหาปิ อาวาสา อญฺญํ อาวาสํ คจฺฉติ, ตตฺถปิ ภิกฺขูนํ เอวํ โหติ “อิมสฺส โข อาวุโส ภิกฺขุโน สํเฆน ตชฺชนีย\-กมฺมํ ปฏิปฺปสฺสทฺธํ ธมฺม\-ปติรูปเกน สมคฺเคหิ, หนฺทสฺส มยํ ตชฺชนีย\-กมฺมํ ปฏิปฺปสฺสมฺเภมา”ติ, เต ตสฺส ตชฺชนีย\-กมฺมํ ปฏิปฺปสฺสมฺเภนฺติ อธมฺเมน วคฺคา. โส ตมฺหาปิ อาวาสา อญฺญํ อาวาสํ คจฺฉติ, ตตฺถปิ ภิกฺขูนํ เอวํ โหติ “อิมสฺส โข อาวุโส ภิกฺขุโน สํเฆน ตชฺชนีย\-กมฺมํ ปฏิปฺปสฺสทฺธํ อธมฺเมน วคฺเคหิ, หนฺทสฺส มยํ ตชฺชนีย\-กมฺมํ ปฏิปฺปสฺสมฺเภมา”ติ, เต ตสฺส ตชฺชนีย\-กมฺมํ ปฏิปฺปสฺสมฺเภนฺติ อธมฺเมน สมคฺคา.}

\proseitem{421}{อิธ ปน ภิกฺขเว ภิกฺขุ สํเฆน ตชฺชนียกมฺมกโต สมฺมา วตฺตติ, โลมํ ปาเตติ, เนตฺถารํ วตฺตติ, ตชฺชนียสฺส กมฺมสฺส ปฏิปฺปสฺสทฺธิํ ยาจติ. ตตฺร เจ ภิกฺขูนํ เอวํ โหติ “อยํ โข อาวุโส ภิกฺขุ สํเฆน ตชฺชนียกมฺมกโต สมฺมา วตฺตติ, โลมํ ปาเตติ, เนตฺถารํ วตฺตติ, ตชฺชนียสฺส กมฺมสฺส ปฏิปฺปสฺสทฺธิํ ยาจติ, หนฺทสฺส มยํ ตชฺชนีย\-กมฺมํ ปฏิปฺปสฺสมฺเภมา”ติ, เต ตสฺส ตชฺชนีย\-กมฺมํ ปฏิปฺปสฺสมฺเภนฺติ ธมฺม\-ปติรูปเกน วคฺคา. โส ตมฺหา อาวาสา อญฺญํ อาวาสํ คจฺฉติ, ตตฺถปิ ภิกฺขูนํ เอวํ โหติ “อิมสฺส โข อาวุโส ภิกฺขุโน สํเฆน ตชฺชนีย\-กมฺมํ ปฏิปฺปสฺสทฺธํ ธมฺม\-ปติรูปเกน วคฺเคหิ, หนฺทสฺส มยํ ตชฺชนีย\-กมฺมํ ปฏิปฺปสฺสมฺเภมา”ติ, เต ตสฺส ตชฺชนีย\-กมฺมํ ปฏิปฺปสฺสมฺเภนฺติ ธมฺม\-ปติรูปเกน สมคฺคา. โส ตมฺหาปิ อาวาสา อญฺญํ อาวาสํ คจฺฉติ, ตตฺถปิ ภิกฺขูนํ \csromanfolio{462}เอวํ โหติ “อิมสฺส โข อาวุโส ภิกฺขุโน สํเฆน ตชฺชนีย\-กมฺมํ ปฏิปฺปสฺสทฺธํ ธมฺม\-ปติรูปเกน สมคฺเคหิ, หนฺทสฺส มยํ ตชฺชนีย\-กมฺมํ ปฏิปฺปสฺสมฺเภมา”ติ, เต ตสฺส ตชฺชนีย\-กมฺมํ \textbf{ปฏิปฺปสฺสมฺเภนฺติ อธมฺเมน วคฺคา. โส ตมฺหาปิ อาวาสา อญฺญํ อาวาสํ} \textbf{คจฺฉติ, ตตฺถปิ ภิกฺขูนํ เอวํ โหติ “อิมสฺส โข อาวุโส ภิกฺขุโน} สํเฆน ตชฺชนีย\-กมฺมํ ปฏิปฺปสฺสทฺธํ อธมฺเมน วคฺเคหิ, หนฺทสฺส มยํ ตชฺชนีย\-กมฺมํ ปฏิปฺปสฺสมฺเภมา”ติ, เต ตสฺส ตชฺชนีย\-กมฺมํ ปฏิปฺปสฺสมฺเภนฺติ อธมฺเมน สมคฺคา. โส ตมฺหาปิ อาวาสา อญฺญํ อาวาสํ คจฺฉติ, ตตฺถปิ ภิกฺขูนํ เอวํ โหติ “อิมสฺส โข อาวุโส ภิกฺขุโน สํเฆน ตชฺชนีย\-กมฺมํ ปฏิปฺปสฺสทฺธํ อธมฺเมน สมคฺเคหิ, หนฺทสฺส มยํ ตชฺชนีย\-กมฺมํ ปฏิปฺปสฺสมฺเภมา”ติ, เต ตสฺส ตชฺชนีย\-กมฺมํ ปฏิปฺปสฺสมฺเภนฺติ ธมฺเมน วคฺคา.}

\proseitem{422}{อิธ ปน ภิกฺขเว ภิกฺขุ สํเฆน ตชฺชนียกมฺมกโต สมฺมา วตฺตติ, โลมํ ปาเตติ, เนตฺถารํ วตฺตติ, ตชฺชนียสฺส กมฺมสฺส ปฏิปฺปสฺสทฺธิํ ยาจติ. ตตฺร เจ ภิกฺขูนํ เอวํ โหติ “อยํ โข อาวุโส ภิกฺขุ สํเฆน ตชฺชนียกมฺมกโต สมฺมา วตฺตติ, โลมํ ปาเตติ, เนตฺถารํ วตฺตติ, ตชฺชนียสฺส กมฺมสฺส ปฏิปฺปสฺสทฺธิํ ยาจติ, หนฺทสฺส มยํ ตชฺชนีย\-กมฺมํ ปฏิปฺปสฺสมฺเภมา”ติ, เต ตสฺส ตชฺชนีย\-กมฺมํ ปฏิปฺปสฺสมฺเภนฺติ ธมฺม\-ปติรูปเกน สมคฺคา. โส ตมฺหา อาวาสา อญฺญํ อาวาสํ คจฺฉติ, ตตฺถปิ ภิกฺขูนํ เอวํ โหติ “อิมสฺส โข อาวุโส ภิกฺขุโน สํเฆน ตชฺชนีย\-กมฺมํ ปฏิปฺปสฺสทฺธํ ธมฺม\-ปติรูปเกน สมคฺเคหิ, หนฺทสฺส มยํ ตชฺชนีย\-กมฺมํ ปฏิปฺปสฺสมฺเภมา”ติ, เต ตสฺส ตชฺชนีย\-กมฺมํ ปฏิปฺปสฺสมฺเภนฺติ อธมฺเมน วคฺคา. โส ตมฺหาปิ อาวาสา อญฺญํ อาวาสํ คจฺฉติ, ตตฺถปิ ภิกฺขูนํ เอวํ โหติ “อิมสฺส โข อาวุโส ภิกฺขุโน สํเฆน ตชฺชนีย\-กมฺมํ ปฏิปฺปสฺสทฺธํ อธมฺเมน วคฺเคหิ, หนฺทสฺส มยํ ตชฺชนีย\-กมฺมํ ปฏิปฺปสฺสมฺเภมา”ติ, เต ตสฺส ตชฺชนีย\-กมฺมํ ปฏิปฺปสฺสมฺเภนฺติ อธมฺเมน สมคฺคา. โส ตมฺหาปิ อาวาสา อญฺญํ อาวาสํ คจฺฉติ, ตตฺถปิ ภิกฺขูนํ เอวํ โหติ “อิมสฺส โข อาวุโส ภิกฺขุโน สํเฆน ตชฺชนีย\-กมฺมํ ปฏิปฺปสฺสทฺธํ อธมฺเมน สมคฺเคหิ, หนฺทสฺส มยํ ตชฺชนีย\-กมฺมํ ปฏิปฺปสฺสมฺเภมา”ติ, เต ตสฺส ตชฺชนีย\-กมฺมํ ปฏิปฺปสฺสมฺเภนฺติ ธมฺเมน วคฺคา. โส ตมฺหาปิ อาวาสา อญฺญํ อาวาสํ คจฺฉติ, ตตฺถปิ ภิกฺขูนํ เอวํ โหติ “อิมสฺส โข อาวุโส ภิกฺขุโน สํเฆน ตชฺชนีย\-กมฺมํ \csromanfolio{463}\textbf{ปฏิปฺปสฺสทฺธํ ธมฺเมน วคฺเคหิ, หนฺทสฺส มยํ ตชฺชนีย}\-\textbf{กมฺมํ} \textbf{ปฏิปฺปสฺสมฺเภมา”ติ, เต ตสฺส ตชฺชนีย}\-\textbf{กมฺมํ ปฏิปฺปสฺสมฺเภนฺติ} ธมฺม\-ปติรูปเกน วคฺคา.}

\csromanmatikaanchor{mtk.259}
\csromantocmarkref{section}{250. นิยสฺสกมฺมปฏิปฺปสฺสทฺธิกถา}{mtk.259}
\bookmark[dest={mtk.259},level=1]{250. นิยสฺสกมฺมปฏิปฺปสฺสทฺธิกถา}
\csromanheader{1}{250. \textbf{นิยสฺสกมฺม}\-\textbf{ปฏิปฺปสฺสทฺธิ}\-\textbf{กถา}}
\proseitem{423}{อิธ ปน ภิกฺขเว ภิกฺขุ สํเฆน นิยสฺสกมฺมกโต สมฺมา วตฺตติ, โลมํ ปาเตติ, เนตฺถารํ วตฺตติ, นิยสฺสสฺส กมฺมสฺส ปฏิปฺปสฺสทฺธิํ ยาจติ. ตตฺร เจ ภิกฺขูนํ เอวํ โหติ “อยํ โข อาวุโส ภิกฺขุ สํเฆน นิยสฺสกมฺมกโต สมฺมา วตฺตติ, โลมํ ปาเตติ, เนตฺถารํ วตฺตติ, นิยสฺสสฺส กมฺมสฺส ปฏิปฺปสฺสทฺธิํ ยาจติ, หนฺทสฺส มยํ นิยสฺสกมฺมํ ปฏิปฺปสฺสมฺเภมา”ติ, เต ตสฺส นิยสฺสกมฺมํ ปฏิปฺปสฺสมฺเภนฺติ อธมฺเมน วคฺคา. โส ตมฺหา อาวาสา อญฺญํ อาวาสํ คจฺฉติ, ตตฺถปิ ภิกฺขูนํ เอวํ โหติ “อิมสฺส โข อาวุโส ภิกฺขุโน สํเฆน นิยสฺสกมฺมํ ปฏิปฺปสฺสทฺธํ อธมฺเมน วคฺเคหิ, หนฺทสฺส มยํ นิยสฺสกมฺมํ ปฏิปฺปสฺสมฺเภมา”ติ, เต ตสฺส นิยสฺสกมฺมํ ปฏิปฺปสฺสมฺเภนฺติ อธมฺเมน สมคฺคา ฯเปฯ ธมฺเมน วคฺคา. ธมฺม\-ปติรูปเกน วคฺคา. ธมฺม\-ปติรูปเกน สมคฺคา ฯเปฯ.}

\vspace{\tipitakaparskip}
\csromancenter{จกฺกํ กาตพฺพํ.}
\csromanmatikaanchor{mtk.260}
\csromantocmarkref{section}{251. ปพฺพาชนียกมฺมปฏิปฺปสฺสทฺธิกถา}{mtk.260}
\bookmark[dest={mtk.260},level=1]{251. ปพฺพาชนียกมฺมปฏิปฺปสฺสทฺธิกถา}
\csromanheader{1}{251. \textbf{ปพฺพาชนียกมฺม}\-\textbf{ปฏิปฺปสฺสทฺธิ}\-\textbf{กถา}}
\proseitem{424}{อิธ ปน ภิกฺขเว ภิกฺขุ สํเฆน ปพฺพาชนีย\-กมฺมกโต สมฺมา วตฺตติ, โลมํ ปาเตติ, เนตฺถารํ วตฺตติ, ปพฺพาชนียสฺส กมฺมสฺส ปฏิปฺปสฺสทฺธิํ ยาจติ. ตตฺร เจ ภิกฺขูนํ เอวํ โหติ “อยํ โข อาวุโส ภิกฺขุ สํเฆน ปพฺพาชนีย\-กมฺมกโต สมฺมา วตฺตติ, โลมํ ปาเตติ, เนตฺถารํ วตฺตติ, ปพฺพาชนียสฺส กมฺมสฺส ปฏิปฺปสฺสทฺธิํ ยาจติ, หนฺทสฺส มยํ ปพฺพาชนีย\-กมฺมํ ปฏิปฺปสฺสมฺเภมา”ติ, เต ตสฺส ปพฺพาชนีย\-กมฺมํ ปฏิปฺปสฺสมฺเภนฺติ อธมฺเมน วคฺคา. โส ตมฺหา อาวาสา อญฺญํ อาวาสํ คจฺฉติ, ตตฺถปิ ภิกฺขูนํ เอวํ โหติ “อิมสฺส โข อาวุโส ภิกฺขุโน สํเฆน ปพฺพาชนีย\-กมฺมํ ปฏิปฺปสฺสทฺธํ อธมฺเมน วคฺเคหิ, หนฺทสฺส มยํ ปพฺพาชนีย\-กมฺมํ ปฏิปฺปสฺสมฺเภมา”ติ, เต ตสฺส ปพฺพาชนีย\-กมฺมํ ปฏิปฺปสฺสมฺเภนฺติ อธมฺเมน สมคฺคา ฯเปฯ ธมฺเมน วคฺคา. ธมฺม\-ปติรูปเกน วคฺคา. ธมฺม\-ปติรูปเกน สมคฺคา ฯเปฯ.}

\vspace{\tipitakaparskip}
\csromancenter{จกฺกํ กาตพฺพํ.}
\csromanheader{2}{252. \textbf{ปฏิสารณียกมฺมฏิปฺปสฺสทฺธิกถา}}
\csromanmatikaanchor{mtk.261}
\csromanmatikaanchor{mtk.262}
\csromantocmarkref{section}{252. ปฏิสารณียกมฺมปฏิปฺปสฺสทฺธิกถา}{mtk.261}
\bookmark[dest={mtk.261},level=1]{252. ปฏิสารณียกมฺมปฏิปฺปสฺสทฺธิกถา}
\csromantocmarkref{section}{253. อทสฺสเน อุกฺเขปนียกมฺมปฏิปฺปสฺสทฺธิกถา}{mtk.262}
\bookmark[dest={mtk.262},level=1]{253. อทสฺสเน อุกฺเขปนียกมฺมปฏิปฺปสฺสทฺธิกถา}
\proseitem{425}{\csromanfolio{464}อิธ ปน ภิกฺขเว ภิกฺขุ สํเฆน ปฏิสารณีย\-กมฺมกโต สมฺมา วตฺตติ, โลมํ ปาเตติ, เนตฺถารํ วตฺตติ, ปฏิสารณียสฺส กมฺมสฺส ปฏิปฺปสฺสทฺธิํ ยาจติ. ตตฺร เจ ภิกฺขูนํ เอวํ โหติ “อยํ โข อาวุโส ภิกฺขุ สํเฆน ปฏิสารณีย\-กมฺมกโต สมฺมา วตฺตติ, โลมํ ปาเตติ, เนตฺถารํ วตฺตติ, ปฏิสารณียสฺส กมฺมสฺส ปฏิปฺปสฺสทฺธิํ ยาจติ, หนฺทสฺส มยํ ปฏิสารณีย\-กมฺมํ ปฏิปฺปสฺสมฺเภมา”ติ, เต ตสฺส ปฏิสารณีย\-กมฺมํ ปฏิปฺปสฺสมฺเภนฺติ อธมฺเมน วคฺคา. โส ตมฺหา อาวาสา อญฺญํ อาวาสํ คจฺฉติ, ตตฺถปิ ภิกฺขูนํ เอวํ โหติ “อิมสฺส โข อาวุโส ภิกฺขุโน สํเฆน ปฏิสารณีย\-กมฺมํ ปฏิปฺปสฺสทฺธํ อธมฺเมน วคฺเคหิ, หนฺทสฺส มยํ ปฏิสารณีย\-กมฺมํ ปฏิปฺปสฺสมฺเภมา”ติ, เต ตสฺส ปฏิสารณีย\-กมฺมํ ปฏิปฺปสฺสมฺเภนฺติ อธมฺเมน สมคฺคา ฯเปฯ ธมฺเมน วคฺคา. ธมฺม\-ปติรูปเกน วคฺคา. ธมฺม\-ปติรูปเกน สมคฺคา ฯเปฯ.}

\vspace{\tipitakaparskip}
\csromancenter{จกฺกํ กาตพฺพํ.}
\prose{253. \textbf{อทสฺสเน อุกฺเขปนียกมฺม}\-\textbf{ปฏิปฺปสฺสทฺธิ}\-\textbf{กถา}}

\csromanmatikaanchor{mtk.263}
\csromanmatikaanchor{mtk.264}
\csromantocmarkref{section}{254. อปฺปฏิกมฺเม อุกฺเขปนียกมฺมปฏิปฺปสฺสทฺธิกถา}{mtk.263}
\bookmark[dest={mtk.263},level=1]{254. อปฺปฏิกมฺเม อุกฺเขปนียกมฺมปฏิปฺปสฺสทฺธิกถา}
\csromantocmarkref{section}{255. อปฺปฏินิสฺสคฺเค อุกฺเขปนียกมฺมปฏิปฺปสฺสทฺธิกถา}{mtk.264}
\bookmark[dest={mtk.264},level=1]{255. อปฺปฏินิสฺสคฺเค อุกฺเขปนียกมฺมปฏิปฺปสฺสทฺธิกถา}
\proseitem{426}{อิธ ปน ภิกฺขเว ภิกฺขุ สํเฆน อาปตฺติยา อทสฺสเน อุกฺเขปนีย\-กมฺมกโต สมฺมา วตฺตติ, โลมํ ปาเตติ, เนตฺถารํ วตฺตติ, อาปตฺติยา อทสฺสเน อุกฺเขปนียสฺส กมฺมสฺส ปฏิปฺปสฺสทฺธิํ ยาจติ. ตตฺร เจ ภิกฺขูนํ เอวํ โหติ “อยํ โข อาวุโส ภิกฺขุ สํเฆน อาปตฺติยา อทสฺสเน อุกฺเขปนีย\-กมฺมกโต สมฺมา วตฺตติ, โลมํ ปาเตติ, เนตฺถารํ วตฺตติ, อาปตฺติยา อทสฺสเน อุกฺเขปนียสฺส กมฺมสฺส ปฏิปฺปสฺสทฺธิํ ยาจติ, หนฺทสฺส มยํ อาปตฺติยา อทสฺสเน อุกฺเขปนีย\-กมฺมํ ปฏิปฺปสฺสมฺเภมา”ติ, เต ตสฺส อาปตฺติยา อทสฺสเน อุกฺเขปนีย\-กมฺมํ ปฏิปฺปสฺสมฺเภนฺติ อธมฺเมน วคฺคา. โส ตมฺหา อาวาสา อญฺญํ อาวาสํ คจฺฉติ, ตตฺถปิ ภิกฺขุนํ เอวํ โหติ “อิมสฺส โข อาวุโส ภิกฺขุโน สํเฆน อาปตฺติยา อทสฺสเน อุกฺเขปนีย\-กมฺมํ ปฏิปฺปสฺสทฺธํ อธมฺเมน วคฺเคหิ, หนฺทสฺส มยํ อาปตฺติยา อทสฺสเน อุกฺเขปนีย\-กมฺมํ ปฏิปฺปสฺสมฺเภมา”ติ, เต ตสฺส อาปตฺติยา อทสฺสเน \csromanfolio{465}อุกฺเขปนีย\-กมฺมํ ปฏิปฺปสฺสมฺเภนฺติ อธมฺเมน สมคฺคา ฯเปฯ ธมฺเมน วคฺคา. ธมฺม\-ปติรูปเกน วคฺคา. ธมฺม\-ปติรูปเกน สมคฺคา ฯเปฯ.}

\vspace{\tipitakaparskip}
\csromancenter{จกฺกํ กาตพฺพํ.}
\prose{254. \textbf{อปฺปฏิกมฺเม อุกฺเขปนียกมฺม}\-\textbf{ปฏิปฺปสฺสทฺธิ}\-\textbf{กถา}}

\proseitem{427}{อิธ ปน ภิกฺขเว ภิกฺขุ สํเฆน อาปตฺติยา อปฺปฏิกมฺเม อุกฺเขปนีย\-กมฺมกโต สมฺมา วตฺตติ, โลมํ ปาเตติ, เนตฺถารํ วตฺตติ, อาปตฺติยา อปฺปฏิกมฺเม อุกฺเขปนียสฺส กมฺมสฺส ปฏิปฺปสฺสทฺธิํ ยาจติ. ตตฺร เจ ภิกฺขูนํ เอวํ โหติ “อยํ โข อาวุโส ภิกฺขุ สํเฆน อาปตฺติยา อปฺปฏิกมฺเม อุกฺเขปนีย\-กมฺมกโต สมฺมา วตฺตติ, โลมํ ปาเตติ, เนตฺถารํ วตฺตติ, อาปตฺติยา อปฺปฏิกมฺเม อุกฺเขปนียสฺส กมฺมสฺส ปฏิปฺปสฺสทฺธิํ ยาจติ, หนฺทสฺส มยํ อาปตฺติยา อปฺปฏิกมฺเม อุกฺเขปนีย\-กมฺมํ ปฏิปฺปสฺสมฺเภมา”ติ, เต ตสฺส อาปตฺติยา อปฺปฏิกมฺเม อุกฺเขปนีย\-กมฺมํ ปฏิปฺปสฺสมฺเภนฺติ อธมฺเมน วคฺคา. โส ตมฺหา อาวาสา อญฺญํ อาวาสํ คจฺฉติ, ตตฺถปิ ภิกฺขูนํ เอวํ โหติ “อิมสฺส โข อาวุโส ภิกฺขุโน สํเฆน อาปตฺติยา อปฺปฏิกมฺเม อุกฺเขปนีย\-กมฺมํ ปฏิปฺปสฺสทฺธํ อธมฺเมน วคฺเคหิ, หนฺทสฺส มยํ อาปตฺติยา อปฺปฏิกมฺเม อุกฺเขปนีย\-กมฺมํ ปฏิปฺปสฺสมฺเภมา”ติ, เต ตสฺส อาปตฺติยา อปฺปฏิกมฺเม อุกฺเขปนีย\-กมฺมํ ปฏิปฺปสฺสมฺเภนฺติ อธมฺเมน สมคฺคา ฯเปฯ ธมฺเมน วคฺคา. ธมฺม\-ปติรูปเกน วคฺคา. ธมฺม\-ปติรูปเกน สมคฺคา ฯเปฯ.}

\vspace{\tipitakaparskip}
\csromancenter{จกฺกํ กาตพฺพํ.}
\prose{255. \textbf{อปฺปฏินิสฺสคฺเค อุกฺเขปนียกมฺม}\-\textbf{ปฏิปฺปสฺสทฺธิ}\-\textbf{กถา}}

\proseitem{428}{อิธ ปน ภิกฺขเว ภิกฺขุ สํเฆน ปาปิกาย ทิฏฺฐิยา อปฺปฏินิสฺสคฺเค อุกฺเขปนีย\-กมฺมกโต สมฺมา วตฺตติ, โลมํ ปาเตติ, เนตฺถารํ วตฺตติ, ปาปิกาย ทิฏฺฐิยา อปฺปฏินิสฺสคฺเค อุกฺเขปนียสฺส กมฺมสฺส ปฏิปฺปสฺสทฺธิํ ยาจติ. ตตฺร เจ ภิกฺขูนํ เอวํ โหติ “อยํ โข อาวุโส ภิกฺขุ สํเฆน ปาปิกาย ทิฏฺฐิยา อปฺปฏินิสฺสคฺเค อุกฺเขปนีย\-กมฺมกโต สมฺมา วตฺตติ, โลมํ ปาเตติ, เนตฺถารํ วตฺตติ, ปาปิกาย ทิฏฺฐิยา อปฺปฏินิสฺสคฺเค อุกฺเขปนียสฺส กมฺมสฺส ปฏิปฺปสฺสทฺธิํ ยาจติ, หนฺทสฺส มยํ ปาปิกาย ทิฏฺฐิยา อปฺปฏินิสฺสคฺเค อุกฺเขปนีย\-กมฺมํ ปฏิปฺปสฺสมฺเภมา”ติ, เต \csromanfolio{466}ตสฺส ปาปิกาย ทิฏฺฐิยา อปฺปฏินิสฺสคฺเค อุกฺเขปนีย\-กมฺมํ ปฏิปฺปสฺสมฺเภนฺติ อธมฺเมน วคฺคา. โส ตมฺหา อาวาสา อญฺญํ อาวาสํ คจฺฉติ, ตตฺถปิ ภิกฺขูนํ เอวํ โหติ “อิมสฺส โข อาวุโส ภิกฺขุโน สํเฆน ปาปิกาย ทิฏฺฐิยา อปฺปฏินิสฺสคฺเค อุกฺเขปนีย\-กมฺมํ ปฏิปฺปสฺสทฺธํ อธมฺเมน วคฺเคหิ, หนฺทสฺส มยํ ปาปิกาย ทิฏฺฐิยา อปฺปฏินิสฺสคฺเค อุกฺเขปนีย\-กมฺมํ ปฏิปฺปสฺสมฺเภมา”ติ, เต ตสฺส ปาปิกาย ทิฏฺฐิยา อปฺปฏินิสฺสคฺเค อุกฺเขปนีย\-กมฺมํ ปฏิปฺปสฺสมฺเภนฺติ อธมฺเมน สมคฺคา ฯเปฯ ธมฺเมน วคฺคา. ธมฺม\-ปติรูปเกน วคฺคา. ธมฺม\-ปติรูปเกน สมคฺคา ฯเปฯ.}

\vspace{\tipitakaparskip}
\csromancenter{จกฺกํ กาตพฺพํ.}
\csromanmatikaanchor{mtk.265}
\csromantocmarkref{section}{256. ตชฺชนียกมฺมวิวาทกถา}{mtk.265}
\bookmark[dest={mtk.265},level=1]{256. ตชฺชนียกมฺมวิวาทกถา}
\csromanheader{1}{256. \textbf{ตชฺชนียกมฺม}\-\textbf{วิวาท}\-\textbf{กถา}}
\proseitem{429}{อิธ ปน ภิกฺขเว ภิกฺขุ ภณฺฑนการโก โหติ กลหการโก วิวาทการโก ภสฺสการโก สํเฆ อธิกรณ\-การโก. ตตฺร เจ ภิกฺขูนํ เอวํ โหติ “อยํ โข อาวุโส ภิกฺขุ ภณฺฑนการโก ฯเปฯ สํเฆ อธิกรณ\-การโก, หนฺทสฺส มยํ ตชฺชนีย\-กมฺมํ กโรมา”ติ, เต ตสฺส ตชฺชนีย\-กมฺมํ กโรนฺติ อธมฺเมน วคฺคา. ตตฺรฏฺโฐ สํโฆ วิวทติ “อธมฺเมน วคฺคกมฺมํ, อธมฺเมน สมคฺคกมฺมํ, ธมฺเมน วคฺคกมฺมํ, ธมฺม\-ปติรูปเกน วคฺคกมฺมํ, ธมฺม\-ปติรูปเกน สมคฺคกมฺมํ, อกตํ กมฺมํ ทุกฺกฏํ กมฺมํ ปุน กาตพฺพํ กมฺมนฺ”ติ. ตตฺร ภิกฺขเว เย เต ภิกฺขู เอวมาหํสุ “อธมฺเมน วคฺคกมฺมนฺ”ติ, เย จ เต ภิกฺขู เอวมาหํสุ “อกตํ กมฺมํ ทุกฺกฏํ กมฺมํ ปุน กาตพฺพํ กมฺมนฺ”ติ, อิเม ตตฺถ ภิกฺขู ธมฺมวาทิโน.}

\proseitem{430}{อิธ ปน ภิกฺขเว ภิกฺขุ ภณฺฑนการโก โหติ ฯเปฯ สํเฆ อธิกรณ\-การโก. ตตฺร เจ ภิกฺขูนํ เอวํ โหติ “อยํ โข อาวุโส ภิกฺขุ ภณฺฑนการโก ฯเปฯ สํเฆ อธิกรณ\-การโก, หนฺทสฺส มยํ ตชฺชนีย\-กมฺมํ กโรมา”ติ, เต ตสฺส ตชฺชนีย\-กมฺมํ กโรนฺติ อธมฺเมน สมคฺคา. ตตฺรฏฺโฐ สํโฆ วิวทติ “อธมฺเมน วคฺคกมฺมํ, อธมฺเมน สมคฺคกมฺมํ, ธมฺเมน วคฺคกมฺมํ, ธมฺม\-ปติรูปเกน วคฺคกมฺมํ, ธมฺม\-ปติรูปเกน สมคฺคกมฺมํ, อกตํ กมฺมํ ทุกฺกฏํ กมฺมํ ปุน กาตพฺพํ กมฺมนฺ”ติ. ตตฺร ภิกฺขเว เย เต ภิกฺขู เอวมาหํสุ “อธมฺเมน สมคฺคกมฺมนฺ”ติ, เย จ เต ภิกฺขู เอวมาหํสุ “อกตํ กมฺมํ ทุกฺกฏํ กมฺมํ ปุน กาตพฺพํ กมฺมนฺ”ติ, อิเม ตตฺถ ภิกฺขู ธมฺมวาทิโน.}

\proseitem{431}{\csromanfolio{467}อิธ ปน ภิกฺขเว ภิกฺขุ ภณฺฑนการโก โหติ ฯเปฯ สํเฆ อธิกรณ\-การโก. ตตฺร เจ ภิกฺขูนํ เอวํ โหติ “อยํ โข อาวุโส ภิกฺขุ ภณฺฑนการโก ฯเปฯ สํเฆ อธิกรณ\-การโก, หนฺทสฺส มยํ ตชฺชนีย\-กมฺมํ กาโรมา”ติ, เต ตสฺส ตชฺชนีย\-กมฺมํ กโรนฺติ ธมฺเมน วคฺคา. ตตฺรฏฺโฐ สํโฆ วิวทติ “อธมฺเมน วคฺคกมฺมํ, อธมฺเมน สมคฺคกมฺมํ, ธมฺเมน วคฺคกมฺมํ, ธมฺม\-ปติรูปเกน วคฺคกมฺมํ, ธมฺม\-ปติรูปเกน สมคฺคกมฺมํ, อกตํ กมฺมํ ทุกฺกฏํ กมฺมํ ปุน กาตพฺพํ กมฺมนฺ”ติ. ตตฺร ภิกฺขเว เย เต ภิกฺขู เอวมาหํสุ “ธมฺเมน วคฺคกมฺมนฺ”ติ, เย จ เต ภิกฺขู เอวมาหํสุ “อกตํ กมฺมํ ทุกฺกฏํ กมฺมํ ปุน กาตพฺพํ กมฺมนฺ”ติ, อิเม ตตฺถ ภิกฺขู ธมฺมวาทิโน.}

\proseitem{432}{อิธ ปน ภิกฺขเว ภิกฺขุ ภณฺฑนการโก โหติ ฯเปฯ สํเฆ อธิกรณ\-การโก. ตตฺร เจ ภิกฺขูนํ เอวํ โหติ “อยํ โข อาวุโส ภิกฺขุ ภณฺฑนการโก ฯเปฯ สํเฆ อธิกรณ\-การโก, หนฺทสฺส มยํ ตชฺชนีย\-กมฺมํ กโรมา”ติ, เต ตสฺส ตชฺชนีย\-กมฺมํ กโรนฺติ ธมฺม\-ปติรูปเกน วคฺคา. ตตฺรฏฺโฐ สํโฆ วิวทติ “อธมฺเมน วคฺคกมฺมํ, อธมฺเมน สมคฺคกมฺมํ, ธมฺเมน วคฺคกมฺมํ, ธมฺม\-ปติรูปเกน วคฺคกมฺมํ, ธมฺม\-ปติรูปเกน สมคฺคกมฺมํ, อกตํ กมฺมํ ทุกฺกฏํ กมฺมํ ปุน กาตพฺพํ กมฺมนฺ”ติ. ตตฺร ภิกฺขเว เย เต ภิกฺขู เอวมาหํสุ “ธมฺม\-ปติรูปเกน สมคฺคกมฺมนฺ”ติ, เย จ เต ภิกฺขู เอวมาหํสุ “อกตํ กมฺมํ ทุกฺกฏํ กมฺมํ ปุน กาตพฺพํ กมฺมนฺ”ติ, อิเม ตตฺถ ภิกฺขู ธมฺมวาทิโน.}

\proseitem{433}{อิธ ปน ภิกฺขเว ภิกฺขุ ภณฺฑนการโก โหติ ฯเปฯ สํเฆ อธิกรณ\-การโก. ตตฺร เจ ภิกฺขูนํ เอวํ โหติ “อยํ โข อาวุโส ภิกฺขุ ภณฺฑนกรโก โหติ ฯเปฯ สํเฆ อธิกรณ\-การโก, หนฺทสฺส มยํ ตชฺชนีย\-กมฺมํ กโรมา”ติ, เต ตสฺส ตชฺชนีย\-กมฺมํ กโรนฺติ ธมฺม\-ปติรูปเกน สมคฺคา. ตตฺรฏฺโฐ สํโฆ วิวทติ “อธมฺเมน วคฺคกมฺมํ, อธมฺเมน สมคฺคกมฺมํ, ธมฺเมน วคฺคกมฺมํ, ธมฺม\-ปติรูปเกน วคฺคกมฺมํ, ธมฺม\-ปติรูปเกน สมคฺคกมฺมํ, อกตํ กมฺมํ ทุกฺกฏํ กมฺมํ ปุน กาตพฺพํ กมฺมนฺ”ติ. ตตฺร ภิกฺขเว เย เต ภิกฺขู เอวมาหํสุ “ธมฺม\-ปติรูปเกน สมคฺคกมฺมนฺ”ติ, เย จ เต ภิกฺขู เอวมาหํสุ “อกตํ กมฺมํ ทุกฺกฏํ กมฺมํ ปุน กาตพฺพํ กมฺมนฺ”ติ, อิเม ตตฺถ ภิกฺขู ธมฺมวาทิโน.}

\csromanmatikaanchor{mtk.266}
\csromantocmarkref{section}{257. นิยสฺสกมฺมวิวาทกถา}{mtk.266}
\bookmark[dest={mtk.266},level=1]{257. นิยสฺสกมฺมวิวาทกถา}
\csromanheader{1}{257. \textbf{นิยสฺสกมฺม}\-\textbf{วิวาท}\-\textbf{กถา}}
\proseitem{434}{\csromanfolio{468}อิธ ปน ภิกฺขเว ภิกฺขุ พาโล โหติ อพฺยตฺโต อาปตฺติพหุโล อนปทาโน, คิหิสํสฏฺโฐ วิหรติ อนนุโลมิเกหิ คิหิสํสคฺเคหิ. ตตฺร เจ ภิกฺขูนํ เอวํ โหติ “อยํ โข อาวุโส ภิกฺขุ พาโล อพฺยตฺโต อาปตฺติพหุโล อนปทาโน, คิหิสํสฏฺโฐ วิหรติ อนนุโลมิเกหิ คิหิสํสคฺเคหิ, หนฺทสฺส มยํ นิยสฺสกมฺมํ กโรมา”ติ, เต ตสฺส นิยสฺสกมฺมํ กโรนฺติ อธมฺเมน วคฺคา ฯเปฯ อธมฺเมน สมคฺคา. ธมฺเมน วคฺคา. ธมฺม\-ปติรูปเกน วคฺคา. ธมฺม\-ปติรูปเกน สมคฺคา. ตตฺรฏฺโฐ สํโฆ วิวทติ “อธมฺเมน วคฺคกมฺมํ, อธมฺเมน สมคฺคกมฺมํ, ธมฺเมน วคฺคกมฺมํ, ธมฺม\-ปติรูปเกน วคฺคกมฺมํ, ธมฺม\-ปติรูปเกน สมคฺคกมฺมํ, อกตํ กมฺมํ ทุกฺกฏํ กมฺมํ ปุน กาตพฺพํ กมฺมนฺ”ติ. ตตฺร ภิกฺขเว เย เต ภิกฺขู เอวมาหํสุ “ธมฺม\-ปติรูปเกน สมคฺคกมฺมนฺ”ติ, เย จ เต ภิกฺขู เอวมาหํสุ “อกตํ กมฺมํ ทุกฺกฏํ กมฺมํ ปุน กาตพฺพํ กมฺมนฺ”ติ, อิเม ตตฺถ ภิกฺขู ธมฺมวาทิโน. อิเม ปญฺจ วารา สํขิตฺตา.}

\csromanmatikaanchor{mtk.267}
\csromantocmarkref{section}{258. ปพฺพาชนียกมฺมวิวาทกถา}{mtk.267}
\bookmark[dest={mtk.267},level=1]{258. ปพฺพาชนียกมฺมวิวาทกถา}
\csromanheader{1}{258. \textbf{ปพฺพาชนียกมฺม}\-\textbf{วิวาท}\-\textbf{กถา}}
\proseitem{435}{อิธ ปน ภิกฺขเว ภิกฺขุ กุลทูสโก โหติ ปาปสมาจาโร. ตตฺร เจ ภิกฺขูนํ เอวํ โหติ “อยํ โข อาวุโส ภิกฺขุ กุลทูสโก ปาปสมาจาโร, หนฺทสฺส มยํ ปพฺพาชนีย\-กมฺมํ กโรมา”ติ, เต ตสฺส ปพฺพาชนีย\-กมฺมํ กโรนฺติ อธมฺเมน วคฺคา ฯเปฯ อธมฺเมน สมคฺคา. ธมฺเมน วคฺคา. ธมฺม\-ปติรูปเกน วคฺคา. ธมฺม\-ปติรูปเกน สมคฺคา. ตตฺรฏฺโฐ สํโฆ วิวทติ “อธมฺเมน วคฺคกมฺมํ, อธมฺเมน สมคฺคกมฺมํ, ธมฺเมน วคฺคกมฺมํ, ธมฺม\-ปติรูปเกน วคฺคกมฺมํ, ธมฺม\-ปติรูปเกน สมคฺคกมฺมํ, อกตํ กมฺมํ ทุกฺกฏํ กมฺมํ ปุน กาตพฺพํ กมฺมนฺ”ติ. ตตฺร ภิกฺขเว เย เต ภิกฺขู เอวมาหํสุ “ธมฺม\-ปติรูปเกน สมคฺคกมฺมนฺ”ติ, เย จ เต ภิกฺขู เอวมาหํสุ “อกตํ กมฺมํ ทุกฺกฏํ กมฺมํ ปุน กาตพฺพํ กมฺมนฺ”ติ, อิเม ตตฺถ ภิกฺขู ธมฺมวาทิโน. อิเม ปญฺจ วารา สํขิตฺตา.}

\csromanmatikaanchor{mtk.268}
\csromantocmarkref{section}{259. ปฏิสารณียกมฺมวิวาทกถา}{mtk.268}
\bookmark[dest={mtk.268},level=1]{259. ปฏิสารณียกมฺมวิวาทกถา}
\csromanheader{1}{259. \textbf{ปฏิสารณียกมฺม}\-\textbf{วิวาท}\-\textbf{กถา}}
\proseitem{436}{อิธ ปน ภิกฺขเว ภิกฺขุ คิหี อกฺโกสติ ปริภาสติ. ตตฺร เจ ภิกฺขูนํ เอวํ โหติ “อยํ โข อาวุโส ภิกฺขุ คิหี \csromanfolio{469}อกฺโกสติ ปริภาสติ, หนฺทสฺส มยํ ปฏิสารณีย\-กมฺมํ กโรมา”ติ, เต ตสฺส ปฏิสารณีย\-กมฺมํ กโรนฺติ อธมฺเมน วคฺคา ฯเปฯ อธมฺเมน สมคฺคา. ธมฺเมน วคฺคา. ธมฺม\-ปติรูปเกน วคฺคา. ธมฺม\-ปติรูปเกน สมคฺคา. ตตฺรฏฺโฐ สํโฆ วิวทติ “อธมฺเมน วคฺคกมฺมํ, อธมฺเมน สมคฺคกมฺมํ, ธมฺเมน วคฺคกมฺมํ, ธมฺม\-ปติรูปเกน วคฺคกมฺมํ, ธมฺม\-ปติรูปเกน สมคฺคกมฺมํ, อกตํ กมฺมํ ทุกฺกฏํ กมฺมํ ปุน กาตพฺพํ \textbf{กมฺมนฺ”ติ. ตตฺร ภิกฺขเว เย เต ภิกฺขู เอวมาหํสุ} “\textbf{ธมฺม}\-\textbf{ปติรูปเกน สมคฺคกมฺมนฺ”ติ, เย จ เต ภิกฺขู เอวมาหํสุ} “อกตํ กมฺมํ ทุกฺกฏํ กมฺมํ ปุน กาตพฺพํ กมฺมนฺ”ติ, อิเม ตตฺถ ภิกฺขู ธมฺมวาทิโน. อิเม ปญฺจ วารา สํขิตฺตา.}

\csromanmatikaanchor{mtk.269}
\csromantocmarkref{section}{260. อทสฺสเน อุกฺเขปนียกมฺมวิวาทกถา ...}{mtk.269}
\bookmark[dest={mtk.269},level=1]{260. อทสฺสเน อุกฺเขปนียกมฺมวิวาทกถา ...}
\csromanheader{1}{260. \textbf{อทสฺสเน อุกฺเขปนียกมฺม}\-\textbf{วิวาท}\-\textbf{กถา}}
\proseitem{437}{อิธ ปน ภิกฺขเว ภิกฺขุ อาปตฺติํ อาปชฺชิตฺวา น อิจฺฉติ อาปตฺติํ ปสฺสิตุํ. ตตฺร เจ ภิกฺขูนํ เอวํ โหติ “อยํ โข อาวุโส ภิกฺขุ อาปตฺติํ อาปชฺชิตฺวา น อิจฺฉติ อาปตฺติํ ปสฺสิตุํ, หนฺทสฺส มยํ อาปตฺติยา อทสฺสเน อุกฺเขปนีย\-กมฺมํ กโรมา”ติ, เต ตสฺส อาปตฺติยา อทสฺสเน อุกฺเขปนีย\-กมฺมํ กโรนฺติ อธมฺเมน วคฺคา ฯเปฯ อธมฺเมน สมคฺคา. ธมฺเมน วคฺคา. ธมฺม\-ปติรูปเกน วคฺคา. ธมฺม\-ปติรูปเกน สมคฺคา. ตตฺรฏฺโฐ สํโฆ วิวทติ “อธมฺเมน วคฺคกมฺมํ, อธมฺเมน สมคฺคกมฺมํ, ธมฺเมน วคฺคกมฺมํ, ธมฺม\-ปติรูปเกน วคฺคกมฺมํ, ธมฺม\-ปติรูปเกน สมคฺคกมฺมํ, อกตํ กมฺมํ ทุกฺกฏํ กมฺมํ ปุน กาตพฺพํ กมฺมนฺ”ติ. ตตฺร ภิกฺขเว เย เต ภิกฺขู เอวมาหํสุ “ธมฺม\-ปติรูปเกน สมคฺคกมฺมนฺ”ติ, เย จ เต ภิกฺขู เอวมาหํสุ “อกตํ กมฺมํ ทุกฺกฏํ กมฺมํ ปุน กาตพฺพํ กมฺมนฺ”ติ, อิเม ตตฺถ ภิกฺขู ธมฺมวาทิโน. อิเม ปญฺจ วารา สํขิตฺตา.}

\csromanmatikaanchor{mtk.270}
\csromantocmarkref{section}{261. อปฺปฏิกมฺเม อุกฺเขปนียกมฺมวิวาทกถา}{mtk.270}
\bookmark[dest={mtk.270},level=1]{261. อปฺปฏิกมฺเม อุกฺเขปนียกมฺมวิวาทกถา}
\csromanheader{1}{261. \textbf{อปฺปฏิกมฺเม อุกฺเขปนียกมฺม}\-\textbf{วิวาท}\-\textbf{กถา}}
\proseitem{438}{อิธ ปน ภิกฺขเว ภิกฺขุ อาปตฺติํ อาปชฺชิตฺวา น อิจฺฉติ อาปตฺติํ ปฏิกาตุํ. ตตฺร เจ ภิกฺขูนํ เอวํ โหติ “อยํ โข อาวุโส ภิกฺขุ อาปตฺติํ อาปชฺชิตฺวา น อิจฺฉติ อาปตฺติํ ปฏิกาตุํ, หนฺทสฺส มยํ อาปตฺติยา อปฺปฏิกมฺเม อุกฺเขปนีย\-กมฺมํ กโรมา”ติ, เต ตสฺส อาปตฺติยา อปฺปฏิกมฺเม อุกฺเขปนีย\-กมฺมํ กโรนฺติ อธมฺเมน วคฺคา ฯเปฯ อธมฺเมน สมคฺคา. ธมฺเมน วคฺคา. ธมฺม\-ปติรูปเกน วคฺคา. ธมฺม\-ปติรูปเกน สมคฺคา. ตตฺรฏฺโฐ สํโฆ วิวทติ “อธมฺเมน วคฺคกมฺมํ, อธมฺเมน สมคฺคกมฺมํ, \csromanfolio{470}ธมฺเมน วคฺคกมฺมํ, ธมฺม\-ปติรูปเกน วคฺคกมฺมํ, ธมฺม\-ปติรูปเกน สมคฺคกมฺมํ, อกตํ กมฺมํ ทุกฺกฏํ กมฺมํ ปุน กาตพฺพํ กมฺมนฺ”ติ, ตตฺร ภิกฺขเว เย เต ภิกฺขู เอวมาหํสุ “ธมฺม\-ปติรูปเกน สมคฺคกมฺมนฺ”ติ, เย จ เต ภิกฺขู เอวมาหํสุ “อกตํ กมฺมํ ทุกฺกฏํ กมฺมํ ปุน กาตพฺพํ กมฺมนฺ”ติ, อิเม ตตฺถ ภิกฺขู ธมฺมวาทิโน. อิเม ปญฺจ วารา สํขิตฺตา.}

\csromanmatikaanchor{mtk.271}
\csromantocmarkref{section}{262. อปฺปฏินิสฺสคฺเค อุกฺเขปนียกมฺมวิวาทกถา}{mtk.271}
\bookmark[dest={mtk.271},level=1]{262. อปฺปฏินิสฺสคฺเค อุกฺเขปนียกมฺมวิวาทกถา}
\csromanheader{1}{262. \textbf{อปฺปฏินิสฺสคฺเค อุกฺเขปนียกมฺม}\-\textbf{วิวาท}\-\textbf{กถา}}
\proseitem{439}{อิธ ปน ภิกฺขเว ภิกฺขุ น อิจฺฉติ ปาปิกํ ทิฏฺฐิํ ปฏินิสฺสชฺชิตุํ. ตตฺร เจ ภิกฺขูนํ เอวํ โหติ “อยํ โข อาวุโส ภิกฺขุ น อิจฺฉติ ปาปิกํ ทิฏฺฐิํ ปฏินิสฺสชฺชิตุํ, หนฺทสฺส มยํ ปาปิกาย ทิฏฺฐิยา อปฺปฏินิสฺสคฺเค อุกฺเขปนีย\-กมฺมํ กโรมา”ติ, เต ตสฺส ปาปิกาย ทิฏฺฐิยา อปฺปฏินิสฺสคฺเค อุกฺเขปนีย\-กมฺมํ กโรนฺติ อธมฺเมน วคฺคา ฯเปฯ อธมฺเมน สมคฺคา. ธมฺเมน วคฺคา. ธมฺม\-ปติรูปเกน วคฺคา. ธมฺม\-ปติรูปเกน สมคฺคา. ตตฺรฏฺโฐ สํโฆ วิวทติ “อธมฺเมน วคฺคกมฺมํ, อธมฺเมน สมคฺคกมฺมํ, ธมฺเมน วคฺคกมฺมํ, ธมฺม\-ปติรูปเกน วคฺคกมฺมํ, ธมฺม\-ปติรูปเกน สมคฺคกมฺมํ, อกตํ กมฺมํ ทุกฺกฏํ กมฺมํ ปุน กาตพฺพํ กมฺมนฺ”ติ. ตตฺร ภิกฺขเว เย เต ภิกฺขู เอวมาหํสุ “ธมฺม\-ปติรูปเกน สมคฺคกมฺมนฺ”ติ, เย จ เต ภิกฺขู เอวมาหํสุ “อกตํ กมฺมํ ทุกฺกฏํ กมฺมํ ปุน กาตพฺพํ กมฺมนฺ”ติ, อิเม ตตฺถ ภิกฺขู ธมฺมวาทิโน. อิเม ปญฺจ วารา สํขิตฺตา.}

\csromanmatikaanchor{mtk.272}
\csromantocmarkref{section}{263. ตชฺชนียกมฺมปฏิปฺปสฺสทฺธิกถา}{mtk.272}
\bookmark[dest={mtk.272},level=1]{263. ตชฺชนียกมฺมปฏิปฺปสฺสทฺธิกถา}
\csromanheader{1}{263. \textbf{ตชฺชนียกมฺม}\-\textbf{ปฏิปฺปสฺสทฺธิ}\-\textbf{กถา}}
\proseitem{440}{อิธ ปน ภิกฺขเว ภิกฺขุ สํเฆน ตชฺชนียกมฺมกโต สมฺมา วตฺตติ, โลมํ ปาเตติ, เนตฺถารํ วตฺตติ, ตชฺชนียสฺส กมฺมสฺส ปฏิปฺปสฺสทฺธิํ ยาจติ. ตตฺร เจ ภิกฺขูนํ เอวํ โหติ “อยํ โข อาวุโส ภิกฺขุ สํเฆน ตชฺชนียกมฺมกโต สมฺมา วตฺตติ, โลมํ ปาเตติ, เนตฺถารํ วตฺตติ, ตชฺชนียสฺส กมฺมสฺส ปฏิปฺปสฺสทฺธิํ ยาจติ, หนฺทสฺส มยํ ตชฺชนีย\-กมฺมํ ปฏิปฺปสฺสมฺเภมา”ติ, เต ตสฺส ตชฺชนีย\-กมฺมํ ปฏิปฺปสฺสมฺเภนฺติ อธมฺเมน วคฺคา. ตตฺรฏฺโฐ สํโฆ วิวทติ “อธมฺเมน วคฺคกมฺมํ, อธมฺเมน สมคฺคกมฺมํ, ธมฺเมน วคฺคกมฺมํ, ธมฺม\-ปติรูปเกน วคฺคกมฺมํ, ธมฺม\-ปติรูปเกน สมคฺคกมฺมํ, อกตํ กมฺมํ ทุกฺกฏํ กมฺมํ ปุน กาตพฺพํ กมฺมนฺ”ติ. ตตฺร ภิกฺขเว เย เต ภิกฺขู เอวมาหํสุ “อธมฺเมน วคฺคกมฺมนฺ”ติ, เย จ เต ภิกฺขู \csromanfolio{471}เอวมาหํสุ “อกตํ กมฺมํ ทุกฺกฏํ กมฺมํ ปุน กาตพฺพํ กมฺมนฺ”ติ, อิเม ตตฺถ ภิกฺขู ธมฺมวาทิโน.}

\proseitem{441}{อิธ ปน ภิกฺขเว ภิกฺขุ สํเฆน ตชฺชนียกมฺมกโต สมฺมา วตฺตติ, โลมํ ปาเตติ, เนตฺถารํ วตฺตติ, ตชฺชนียสฺส กมฺมสฺส ปฏิปฺปสฺสทฺธิํ ยาจติ. ตตฺร เจ ภิกฺขูนํ เอวํ โหติ “อยํ โข อาวุโส ภิกฺขุ สํเฆน ตชฺชนียกมฺมกโต สมฺมา วตฺตติ, โลมํ ปาเตติ, เนตฺถารํ วตฺตติ, ตชฺชนียสฺส กมฺมสฺส ปฏิปฺปสฺสทฺธิํ ยาจติ, หนฺทสฺส มยํ ตชฺชนีย\-กมฺมํ ปฏิปฺปสฺสมฺเภมา”ติ, เต ตสฺส ตชฺชนีย\-กมฺมํ ปฏิปฺปสฺสมฺเภนฺติ อธมฺเมน สมคฺคา. ตตฺรฏฺโฐ สํโฆ วิวทติ “อธมฺเมน วคฺคกมฺมํ, อธมฺเมน สมคฺคกมฺมํ, ธมฺเมน วคฺคกมฺมํ, ธมฺม\-ปติรูปเกน วคฺคกมฺมํ, ธมฺม\-ปติรูปเกน สมคฺคกมฺมํ, อกตํ กมฺมํ ทุกฺกฏํ กมฺมํ ปุน กาตพฺพํ กมฺมนฺ”ติ. ตตฺร ภิกฺขเว เย เต ภิกฺขู เอวมาหํสุ “อธมฺเมน สมคฺคกมฺมนฺ”ติ, เย จ เต ภิกฺขู เอวมาหํสุ “อกตํ กมฺมํ ทุกฺกฏํ กมฺมํ ปุน กาตพฺพํ กมฺมนฺ”ติ, อิเม ตตฺถ ภิกฺขู ธมฺมวาทิโน.}

\proseitem{442}{อิธ ปน ภิกฺขเว ภิกฺขุ สํเฆน ตชฺชนียกมฺมกโต สมฺมา วตฺตติ, โลมํ ปาเตติ, เนตฺถารํ วตฺตติ, ตชฺชนียสฺส กมฺมสฺส ปฏิปฺปสฺสทฺธิํ ยาจติ. ตตฺร เจ ภิกฺขูนํ เอวํ โหติ “อยํ โข อาวุโส ภิกฺขุ สํเฆน ตชฺชนียกมฺมกโต สมฺมา วตฺตติ, โลมํ ปาเตติ, เนตฺถารํ วตฺตติ, ตชฺชนียสฺส กมฺมสฺส ปฏิปฺปสฺสทฺธิํ ยาจติ, หนฺทสฺส มยํ ตชฺชนีย\-กมฺมํ ปฏิปฺปสฺสมฺเภมา”ติ, เต ตสฺส ตชฺชนีย\-กมฺมํ ปฏิปฺปสฺสมฺเภนฺติ ธมฺเมน วคฺคา. ตตฺรฏฺโฐ สํโฆ วิวทติ “อธมฺเมน วคฺคกมฺมํ, อธมฺเมน สมคฺคกมฺมํ, ธมฺเมน วคฺคกมฺมํ, ธมฺม\-ปติรูปเกน วคฺคกมฺมํ, ธมฺม\-ปติรูปเกน สมคฺคกมฺมํ, อกตํ กมฺมํ ทุกฺกฏํ กมฺมํ ปุน กาตพฺพํ กมฺมนฺ”ติ. ตตฺร ภิกฺขเว เย เต ภิกฺขู เอวมาหํสุ “ธมฺเมน วคฺคกมฺมนฺ”ติ, เย จ เต ภิกฺขู เอวมาหํสุ “อกตํ กมฺมํ ทุกฺกฏํ กมฺมํ ปุน กาตพฺพํ กมฺมนฺ”ติ, อิเม ตตฺถ ภิกฺขู ธมฺมวาทิโน.}

\proseitem{443}{อิธ ปน ภิกฺขเว ภิกฺขุ สํเฆน ตชฺชนียกมฺมกโต สมฺมา วตฺตติ, โลมํ ปาเตติ, เนตฺถารํ วตฺตติ, ตชฺชนียสฺส กมฺมสฺส ปฏิปฺปสฺสทฺธิํ ยาจติ. ตตฺร เจ ภิกฺขูนํ เอวํ โหติ “อยํ โข อาวุโส ภิกฺขุ สํเฆน ตชฺชนียกมฺมกโต สมฺมา วตฺตติ, โลมํ ปาเตติ, เนตฺถารํ วตฺตติ, ตชฺชนียสฺส กมฺมสฺส ปฏิปฺปสฺสทฺธิํ ยาจติ, หนฺทสฺส มยํ \csromanfolio{472}ตชฺชนีย\-กมฺมํ ปฏิปฺปสฺสมฺเภมา”ติ, เต ตสฺส ตชฺชนีย\-กมฺมํ ปฏิปฺปสฺสมฺเภนฺติ ธมฺม\-ปติรูปเกน วคฺคา. ตตฺรฏฺโฐ สํโฆ วิวทติ “อธมฺเมน วคฺคกมฺมํ อธมฺเมน สมคฺคกมฺมํ, ธมฺเมน วคฺคกมฺมํ, ธมฺม\-ปติรูปเกน วคฺคกมฺมํ, ธมฺม\-ปติรูปเกน สมคฺคกมฺมํ, อกตํ กมฺมํ ทุกฺกฏํ กมฺมํ ปุน กาตพฺพํ \textbf{กมฺมนฺ”ติ. ตตฺร ภิกฺขเว เย เต ภิกฺขู เอวมาหํสุ} “\textbf{ธมฺม}\-\textbf{ปติรูปเกน วคฺคกมฺมนฺ”ติ, เย จ เต ภิกฺขู เอวมาหํสุ} “อกตํ กมฺมํ ทุกฺกฏํ กมฺมํ ปุน กาตพฺพํ กมฺมนฺ”ติ, อิเม ตตฺถ ภิกฺขู ธมฺมวาทิโน.}

\proseitem{444}{อิธ ปน ภิกฺขเว ภิกฺขุ สํเฆน ตชฺชนียกมฺมกโต สมฺมา วตฺตติ, โลมํ ปาเตติ, เนตฺถารํ วตฺตติ, ตชฺชนียสฺส กมฺมสฺส ปฏิปฺปสฺสทฺธิํ ยาจติ. ตตฺร เจ ภิกฺขูนํ เอวํ โหติ “อยํ โข อาวุโส ภิกฺขุ สํเฆน ตชฺชนียกมฺมกโต สมฺมา วตฺตติ, โลมํ ปาเตติ, เนตฺถารํ วตฺตติ, ตชฺชนียสฺส กมฺมสฺส ปฏิปฺปสฺสทฺธิํ ยาจติ, หนฺทสฺส มยํ ตชฺชนีย\-กมฺมํ ปฏิปฺปสฺสมฺเภมา”ติ, เต ตสฺส ตชฺชนีย\-กมฺมํ ปฏิปฺปสฺสมฺเภนฺติ ธมฺม\-ปติรูปเกน สมคฺคา. ตตฺรฏฺโฐ สํโฆ วิวทติ “อธมฺเมน วคฺคกมฺมํ, อธมฺเมน สมคฺคกมฺมํ, ธมฺเมน วคฺคกมฺมํ, ธมฺม\-ปติรูปเกน วคฺคกมฺมํ, ธมฺม\-ปติรูปเกน สมคฺคกมฺมํ, อกตํ กมฺมํ ทุกฺกฏํ กมฺมํ ปุน กาตพฺพํ กมฺมนฺ”ติ. ตตฺร ภิกฺขเว เย เต ภิกฺขู เอวมาหํสุ “ธมฺม\-ปติรูปเกน สมคฺคกมฺมนฺ”ติ, เย จ เต ภิกฺขู เอวมาหํสุ “อกตํ กมฺมํ ทุกฺกฏํ กมฺมํ ปุน กาตพฺพํ กมฺมนฺ”ติ, อิเม ตตฺถ ภิกฺขู ธมฺมวาทิโน.}

\csromanmatikaanchor{mtk.273}
\csromantocmarkref{section}{264. นิยสฺสกมฺมปฏิปฺปสฺสทฺธิกถา}{mtk.273}
\bookmark[dest={mtk.273},level=1]{264. นิยสฺสกมฺมปฏิปฺปสฺสทฺธิกถา}
\csromanheader{1}{264. \textbf{นิยสฺสกมฺม}\-\textbf{ปฏิปฺปสฺสทฺธิ}\-\textbf{กถา}}
\proseitem{445}{อิธ ปน ภิกฺขเว ภิกฺขุ สํเฆน นิยสฺสกมฺมกโต สมฺมา วตฺตติ, โลมํ ปาเตติ, เนตฺถารํ วตฺตติ, นิยสฺสสฺส กมฺมสฺส ปฏิปฺปสฺสทฺธิํ ยาจติ. ตตฺร เจ ภิกฺขูนํ เอวํ โหติ “อยํ โข อาวุโส ภิกฺขุ สํเฆน นิยสฺสกมฺมกโต สมฺมา วตฺตติ, โลมํ ปาเตติ, เนตฺถารํ วตฺตติ, นิยสฺสสฺส กมฺมสฺส ปฏิปฺปสฺสทฺธิํ ยาจติ, หนฺทสฺส มยํ นิยสฺสกมฺมํ ปฏิปฺปสฺสมฺเภมา”ติ, เต ตสฺส นิยสฺสกมฺมํ ปฏิปฺปสฺสมฺเภนฺติ อธมฺเมน วคฺคา ฯเปฯ อธมฺเมน สมคฺคา. ธมฺเมน วคฺคา. ธมฺม\-ปติรูปเกน วคฺคา. ธมฺม\-ปติรูปเกน สมคฺคา. ตตฺรฏฺโฐ สํโฆ วิวทติ “อธมฺเมน วคฺคกมฺมํ, อธมฺเมน สมคฺคกมฺมํ, ธมฺเมน วคฺคกมฺมํ, ธมฺม\-ปติรูปเกน วคฺคกมฺมํ, ธมฺม\-ปติรูปเกน สมคฺคกมฺมํ, อกตํ กมฺมํ ทุกฺกฏํ กมฺมํ ปุน กาตพฺพํ กมฺมนฺ”ติ. ตตฺร \csromanfolio{473}ภิกฺขเว เย เต ภิกฺขู เอวมาหํสุ “ธมฺม\-ปติรูปเกน สมคฺคกมฺมนฺ”ติ, เย จ เต ภิกฺขู เอวมาหํสุ “อกตํ กมฺมํ ทุกฺกฏํ กมฺมํ ปุน กาตพฺพํ กมฺมนฺ”ติ, อิเม ตตฺถ ภิกฺขู ธมฺมวาทิโน. อิเมปิ\csromansharedfootnote{60c6a223736e91d9}{อิเม (สี, สฺยา, เอวมุปริปิ)} ปญฺจ วารา สํขิตฺตา.}

\csromanmatikaanchor{mtk.274}
\csromantocmarkref{section}{265. ปพฺพาชนียกมฺมปฏิปฺปสฺสทฺธิกถา}{mtk.274}
\bookmark[dest={mtk.274},level=1]{265. ปพฺพาชนียกมฺมปฏิปฺปสฺสทฺธิกถา}
\csromanheader{1}{265. \textbf{ปพฺพาชนียกมฺม}\-\textbf{ปฏิปฺปสฺสทฺธิ}\-\textbf{กถา}}
\proseitem{446}{อิธ ปน ภิกฺขเว ภิกฺขุ สํเฆน ปพฺพาชนีย\-กมฺมกโต สมฺมา วตฺตติ, โลมํ ปาเตติ, เนตฺถารํ วตฺตติ, ปพฺพาชนียสฺส กมฺมสฺส ปฏิปฺปสฺสทฺธิํ ยาจติ. ตตฺร เจ ภิกฺขูนํ เอวํ โหติ “อยํ โข อาวุโส ภิกฺขุ สํเฆน ปพฺพาชนีย\-กมฺมกโต สมฺมา วตฺตติ, โลมํ ปาเตติ, เนตฺถารํ วตฺตติ, ปพฺพาชนียสฺส กมฺมสฺส ปฏิปฺปสฺสทฺธิํ ยาจติ, หนฺทสฺส มยํ ปพฺพาชนีย\-กมฺมํ ปฏิปฺปสฺสมฺเภมา”ติ, เต ตสฺส ปพฺพาชนีย\-กมฺมํ ปฏิปฺปสฺสมฺเภนฺติ อธมฺเมน วคฺคา ฯเปฯ อธมฺเมน สมคฺคา. ธมฺเมน วคฺคา. ธมฺม\-ปติรูปเกน วคฺคา. ธมฺม\-ปติรูปเกน สมคฺคา. ตตฺรฏฺโฐ สํโฆ วิวทติ “อธมฺเมน วคฺคกมฺมํ, อธมฺเมน สมคฺคกมฺมํ, ธมฺเมน วคฺคกมฺมํ, ธมฺม\-ปติรูปเกน วคฺคกมฺมํ, ธมฺม\-ปติรูปเกน สมคฺคกมฺมํ, อกตํ กมฺมํ ทุกฺกฏํ กมฺมํ ปุน กาตพฺพํ กมฺมนฺ”ติ. ตตฺร ภิกฺขเว เย เต ภิกฺขู เอวมาหํสุ “ธมฺม\-ปติรูปเกน สมคฺคกมฺมนฺ”ติ, เย จ เต ภิกฺขู เอวมาหํสุ “อกตํ กมฺมํ ทุกฺกฏํ กมฺมํ ปุน กาตพฺพํ กมฺมนฺ”ติ, อิเม ตตฺถ ภิกฺขู ธมฺมวาทิโน. อิเมปิ ปญฺจ วารา สํขิตฺตา.}

\csromanmatikaanchor{mtk.275}
\csromantocmarkref{section}{266. ปฏิสารณียกมฺมปฏิปฺปสฺสทฺธิกถา}{mtk.275}
\bookmark[dest={mtk.275},level=1]{266. ปฏิสารณียกมฺมปฏิปฺปสฺสทฺธิกถา}
\csromanheader{1}{266. \textbf{ปฏิสารณียกมฺม}\-\textbf{ปฏิปฺปสฺสทฺธิ}\-\textbf{กถา}}
\csromanmatikaanchor{mtk.276}
\csromanmatikaanchor{mtk.277}
\csromantocmarkref{section}{267. อทสฺสเน อุกฺเขปนียกมฺมปฏิปฺปสฺสทฺธิกถา}{mtk.276}
\bookmark[dest={mtk.276},level=1]{267. อทสฺสเน อุกฺเขปนียกมฺมปฏิปฺปสฺสทฺธิกถา}
\csromantocmarkref{section}{268. อปฺปฏิกมฺเม อุกฺเขปนียกมฺมปฏิปฺปสฺสทฺธิกถา}{mtk.277}
\bookmark[dest={mtk.277},level=1]{268. อปฺปฏิกมฺเม อุกฺเขปนียกมฺมปฏิปฺปสฺสทฺธิกถา}
\proseitem{447}{อิธ ปน ภิกฺขเว ภิกฺขุ สํเฆน ปฏิสารณีย\-กมฺมกโต สมฺมา วตฺตติ, โลมํ ปาเตติ, เนตฺถารํ วตฺตติ, ปฏิสารณียสฺส กมฺมสฺส ปฏิปฺปสฺสทฺธิํ ยาจติ. ตตฺร เจ ภิกฺขูนํ เอวํ โหติ “อยํ โข อาวุโส ภิกฺขุ สํเฆน ปฏิสารณีย\-กมฺมกโต สมฺมา วตฺตติ, โลมํ ปาเตติ, เนตฺถารํ วตฺตติ, ปฏิสารณียสฺส กมฺมสฺส ปฏิปฺปสฺสทฺธิํ ยาจติ, หนฺทสฺส มยํ ปฏิสารณีย\-กมฺมํ ปฏิปฺปสฺสมฺเภมา”ติ, เต ตสฺส ปฏิสารณีย\-กมฺมํ ปฏิปฺปสฺสมฺเภนฺติ อธมฺเมน วคฺคา ฯเปฯ อธมฺเมน สมคฺคา. ธมฺเมน วคฺคา. ธมฺม\-ปติรูปเกน วคฺคา. ธมฺม\-ปติรูปเกน สมคฺคา. ตตฺรฏฺโฐ สํโฆ วิวทติ “อธมฺเมน วคฺคกมฺมํ, อธมฺเมน สมคฺคกมฺมํ, ธมฺเมน \csromanfolio{474}วคฺคกมฺมํ, ธมฺม\-ปติรูปเกน วคฺคกมฺมํ, ธมฺม\-ปติรูปเกน สมคฺคกมฺมํ, อกตํ กมฺมํ ทุกฺกฏํ กมฺมํ ปุน กาตพฺพํ กมฺมนฺ”ติ. ตตฺร ภิกฺขเว เย เต ภิกฺขู เอวมาหํสุ “ธมฺม\-ปติรูปเกน สมคฺคกมฺมนฺ”ติ, เย จ เต ภิกฺขู เอวมาหํสุ “อกตํ กมฺมํ ทุกฺกฏํ กมฺมํ ปุน กาตพฺพํ กมฺมนฺ”ติ, อิเม ตตฺถ ภิกฺขู ธมฺมวาทิโน. อิเมปิ ปญฺจ วารา สํขิตฺตา.}

\prose{267. \textbf{อทสฺสเน อุกฺเขปนียกมฺม}\-\textbf{ปฏิปฺปสฺสทฺธิ}\-\textbf{กถา}}

\proseitem{448}{อิธ ปน ภิกฺขเว ภิกฺขุ สํเฆน อาปตฺติยา อทสฺสเน อุกฺเขปนีย\-กมฺมกโต สมฺมา วตฺตติ, โลมํ ปาเตติ, เนตฺถารํ วตฺตติ, อาปตฺติยา อทสฺสเน อุกฺเขปนียสฺส กมฺมสฺส ปฏิปฺปสฺสทฺธิํ ยาจติ. ตตฺร เจ ภิกฺขูนํ เอวํ โหติ “อยํ โข อาวุโส ภิกฺขุ สํเฆน อาปตฺติยา อทสฺสเน อุกฺเขปนีย\-กมฺมกโต สมฺมา วตฺตติ, โลมํ ปาเตติ, เนตฺถารํ วตฺตติ, อาปตฺติยา อทสฺสเน อุกฺเขปนียสฺส กมฺมสฺส ปฏิปฺปสฺสทฺธิํ ยาจติ, หนฺทสฺส มยํ อาปตฺติยา อทสฺสเน อุกฺเขปนีย\-กมฺมํ ปฏิปฺปสฺสมฺเภมา”ติ, เต ตสฺส อาปตฺติยา อทสฺสเน อุกฺเขปนีย\-กมฺมํ ปฏิปฺปสฺสมฺเภนฺติ อธมฺเมน วคฺคา ฯเปฯ อธมฺเมน สมคฺคา. ธมฺเมน วคฺคา. ธมฺม\-ปติรูปเกน วคฺคา. ธมฺม\-ปติรูปเกน สมคฺคา. ตตฺรฏฺโฐ สํโฆ วิวทติ “อธมฺเมน วคฺคกมฺมํ, อธมฺเมน สมคฺคกมฺมํ, ธมฺเมน วคฺคกมฺมํ, ธมฺม\-ปติรูปเกน วคฺคกมฺมํ, ธมฺม\-ปติรูปเกน สมคฺคกมฺมํ, อกตํ กมฺมํ ทุกฺกฏํ กมฺมํ ปุน กาตพฺพํ กมฺมนฺ”ติ. ตตฺร ภิกฺขเว เย เต ภิกฺขู เอวมาหํสุ “ธมฺม\-ปติรูปเกน สมคฺคกมฺมนฺ”ติ, เย จ เต ภิกฺขู เอวมาหํสุ “อกตํ กมฺมํ ทุกฺกฏํ กมฺมํ ปุน กาตพฺพํ กมฺมนฺ”ติ, อิเม ตตฺถ ภิกฺขู ธมฺมวาทิโน. อิเมปิ ปญฺจ วารา สํขิตฺตา.}

\prose{268. \textbf{อปฺปฏิกมฺเม อุกฺเขปนียกมฺม}\-\textbf{ปฏิปฺปสฺสทฺธิ}\-\textbf{กถา}}

\csromanmatikaanchor{mtk.278}
\csromantocmarkref{section}{269. อปฺปฏินิสฺสคฺเค อุกฺเขปนียกมฺมปฏิปฺปสฺสทฺธิกถา}{mtk.278}
\bookmark[dest={mtk.278},level=1]{269. อปฺปฏินิสฺสคฺเค อุกฺเขปนียกมฺมปฏิปฺปสฺสทฺธิกถา}
\proseitem{449}{อิธ ปน ภิกฺขเว ภิกฺขุ สํเฆน อาปตฺติยา อปฺปฏิกมฺเม อุกฺเขปนีย\-กมฺมกโต สมฺมา วตฺตติ, โลมํ ปาเตติ, เนตฺถารํ วตฺตติ, อาปตฺติยา อปฺปฏิกมฺเม อุกฺเขปนียสฺส กมฺมสฺส ปฏิปฺปสฺสทฺธิํ ยาจติ. ตตฺร เจ ภิกฺขูนํ เอวํ โหติ “อยํ โข อาวุโส ภิกฺขุ สํเฆน อาปตฺติยา อปฺปฏิกมฺเม อุกฺเขปนีย\-กมฺมกโต สมฺมา วตฺตติ, โลมํ ปาเตติ, เนตฺถารํ วตฺตติ, อาปตฺติยา อปฺปฏิกมฺเม อุกฺเขปนียสฺส กมฺมสฺส ปฏิปฺปสฺสทฺธิํ ยาจติ, หนฺทสฺส มยํ อาปตฺติยา อปฺปฏิกมฺเม \csromanfolio{475}อุกฺเขปนีย\-กมฺมํ ปฏิปฺปสฺสมฺเภมา”ติ, เต ตสฺส อาปตฺติยา อปฺปฏิกมฺเม อุกฺเขปนีย\-กมฺมํ ปฏิปฺปสฺสมฺเภนฺติ อธมฺเมน วคฺคา ฯเปฯ อธมฺเมน สมคฺคา. ธมฺเมน วคฺคา. ธมฺม\-ปติรูปเกน วคฺคา. ธมฺม\-ปติรูปเกน สมคฺคา. ตตฺรฏฺโฐ สํโฆ วิวทติ “อธมฺเมน วคฺคกมฺมํ, อธมฺเมน สมคฺคกมฺมํ, ธมฺเมน วคฺคกมฺมํ, ธมฺม\-ปติรูปเกน วคฺคกมฺมํ, ธมฺม\-ปติรูปเกน สมคฺคกมฺมํ, อกตํ กมฺมํ ทุกฺกฏํ กมฺมํ ปุน กาตพฺพํ กมฺมนฺ”ติ. ตตฺร ภิกฺขเว เย เต ภิกฺขู เอวมาหํสุ “ธมฺม\-ปติรูปเกน สมคฺคกมฺมนฺ”ติ, เย จ เต ภิกฺขู เอวมาหํสุ “อกตํ กมฺมํ ทุกฺกฏํ กมฺมํ ปุน กาตพฺพํ กมฺมนฺ”ติ, อิเม ตตฺถ ภิกฺขู ธมฺมวาทิโน. อิเมปิ ปญฺจ วารา สํขิตฺตา.}

\prose{269. \textbf{อปฺปฏินิสฺสคฺเค อุกฺเขปนียกมฺม}\-\textbf{ปฏิปฺปสฺสทฺธิ}\-\textbf{กถา}}

\proseitem{450}{อิธ ปน ภิกฺขเว ภิกฺขุ สํเฆน ปาปิกาย ทิฏฺฐิยา อปฺปฏินิสฺสคฺเค อุกฺเขปนีย\-กมฺมกโต สมฺมา วตฺตติ, โลมํ ปาเตติ, เนตฺถารํ วตฺตติ, ปาปิกาย ทิฏฺฐิยา อปฺปฏินิสฺสคฺเค อุกฺเขปนียสฺส กมฺมสฺส ปฏิปฺปสฺสทฺธิํ ยาจติ. ตตฺร เจ ภิกฺขูนํ เอวํ โหติ “อยํ โข อาวุโส ภิกฺขุ สํเฆน ปาปิกาย ทิฏฺฐิยา อปฺปฏินิสฺสคฺเค อุกฺเขปนีย\-กมฺมกโต สมฺมา วตฺตติ, โลมํ ปาเตติ, เนตฺถารํ วตฺตติ, ปาปิกาย ทิฏฺฐิยา อปฺปฏินิสฺสคฺเค อุกฺเขปนียสฺส กมฺมสฺส ปฏิปฺปสฺสทฺธิํ ยาจติ, หนฺทสฺส มยํ ปาปิกาย ทิฏฺฐิยา อปฺปฏินิสฺสคฺเค อุกฺเขปนีย\-กมฺมํ ปฏิปฺปสฺสมฺเภมา”ติ, เต ตสฺส ปาปิกาย ทิฏฺฐิยา อปฺปฏินิสฺสคฺเค อุกฺเขปนีย\-กมฺมํ ปฏิปฺปสฺสมฺเภนฺติ อธมฺเมน วคฺคา ฯเปฯ อธมฺเมน สมคฺคา. ธมฺเมน วคฺคา. ธมฺม\-ปติรูปเกน วคฺคา. ธมฺม\-ปติรูปเกน สมคฺคา. ตตฺรฏฺโฐ สํโฆ วิวทติ “อธมฺเมน วคฺคกมฺมํ, อธมฺเมน สมคฺคกมฺมํ, ธมฺเมน วคฺคกมฺมํ, ธมฺม\-ปติรูปเกน วคฺคกมฺมํ, ธมฺม\-ปติรูปเกน สมคฺคกมฺมํ, อกตํ กมฺมํ ทุกฺกฏํ กมฺมํ ปุน กาตพฺพํ กมฺมนฺ”ติ. ตตฺร ภิกฺขเว เย เต ภิกฺขู เอวมาหํสุ “ธมฺม\-ปติรูปเกน สมคฺคกมฺมนฺ”ติ, เย จ เต ภิกฺขู เอวมาหํสุ “อกตํ กมฺมํ ทุกฺกฏํ กมฺมํ ปุน กาตพฺพํ กมฺมนฺ”ติ, อิเม ตตฺถ ภิกฺขู ธมฺมวาทิโน. อิเมปิ ปญฺจ วารา สํขิตฺตา.}

\vspace{\tipitakaparskip}
\csromancenter{จมฺเปยฺย\-กฺขนฺธโก นวโม.}
\csromancenter{270. \textbf{ตสฺสุทฺทานํ}}
\csromanmatikaanchor{mtk.279}
\csromantocmarkref{section}{270. อุทฺทานคาถา}{mtk.279}
\bookmark[dest={mtk.279},level=1]{270. อุทฺทานคาถา}
\setlength{\csromangathapagewidth}{0pt}
\setlength{\csromangathawakwidth}{0pt}
\csromangathameasuregroup{จมฺปายํ ภควา อาสิ, \\ อาคนฺตุกานมุสฺสุกฺกํ, \\ ปกตญฺญุโนติ ญตฺวา, \\ อุกฺขิตฺโต น กโรตีติ, \\ อธมฺเมน วคฺคกมฺมํ, \\ ธมฺเมน วคฺคกมฺมญฺจ, \\ ปติรูปเกน สมคฺคํ, \\ เอโก จ เทฺว สมฺพหุเล, \\ ทุเวปิ สมฺพหุลาปิ,}{\csromangathabat{จมฺปายํ ภควา อาสิ,}{วตฺถุ วาสภคามเก.} \\ \csromangathabat{อาคนฺตุกานมุสฺสุกฺกํ,}{อกาสิ อิจฺฉิตพฺพเก\textsuperscript{1}.} \\ \csromangathabat{ปกตญฺญุโนติ ญตฺวา,}{อุสฺสุกฺกํ น กรี ตทา.} \\ \csromangathabat{อุกฺขิตฺโต น กโรตีติ,}{สาคมา ชินสนฺติเก.} \\ \csromangathabat{อธมฺเมน วคฺคกมฺมํ,}{สมคฺคํ อธมฺเมน จ.} \\ \csromangathabat{ธมฺเมน วคฺคกมฺมญฺจ,}{ปติรูปเกน วคฺคิกํ.} \\ \csromangathabat{ปติรูปเกน สมคฺคํ,}{เอโก อุกฺขิปเตกกํ.} \\ \csromangathabat{เอโก จ เทฺว สมฺพหุเล,}{สํฆํ อุกฺขิปเตกโก.} \\ \csromangathabat{ทุเวปิ สมฺพหุลาปิ,}{สํโฆ สํฆญฺจ อุกฺขิปิ.}}
\csromangathasetleft{จมฺปายํ ภควา อาสิ, \\ อาคนฺตุกานมุสฺสุกฺกํ, \\ ปกตญฺญุโนติ ญตฺวา, \\ อุกฺขิตฺโต น กโรตีติ, \\ อธมฺเมน วคฺคกมฺมํ, \\ ธมฺเมน วคฺคกมฺมญฺจ, \\ ปติรูปเกน สมคฺคํ, \\ เอโก จ เทฺว สมฺพหุเล, \\ ทุเวปิ สมฺพหุลาปิ,}
\csromangathagroup{\csromangathastanza{\csromanfolio{476}\csromangathabat{จมฺปายํ ภควา อาสิ,}{วตฺถุ วาสภคามเก.} \\ \csromangathabat{อาคนฺตุกานมุสฺสุกฺกํ,}{อกาสิ อิจฺฉิตพฺพเก\csromansharedfootnote{a4d1beee6767665b}{อิจฺฉิตพฺพโก (ก)}.}}\csromangathastanza{\csromangathabat{ปกตญฺญุโนติ ญตฺวา,}{อุสฺสุกฺกํ น กรี ตทา.} \\ \csromangathabat{อุกฺขิตฺโต น กโรตีติ,}{สาคมา ชินสนฺติเก.}}\csromangathastanza{\csromangathabat{อธมฺเมน วคฺคกมฺมํ,}{สมคฺคํ อธมฺเมน จ.} \\ \csromangathabat{ธมฺเมน วคฺคกมฺมญฺจ,}{ปติรูปเกน วคฺคิกํ.}}\csromangathastanza{\csromangathabat{ปติรูปเกน สมคฺคํ,}{เอโก อุกฺขิปเตกกํ.} \\ \csromangathabat{เอโก จ เทฺว สมฺพหุเล,}{สํฆํ อุกฺขิปเตกโก.} \\ \csromangathabat{ทุเวปิ สมฺพหุลาปิ,}{สํโฆ สํฆญฺจ อุกฺขิปิ.}}}

\vspace{\tipitakaparskip}
\csromancenter{สพฺพญฺญุปวโร สุตฺวา, “อธมฺมนฺ”ติ ปฏิกฺขิปิ.}
\setlength{\csromangathapagewidth}{0pt}
\setlength{\csromangathawakwidth}{0pt}
\csromangathameasuregroup{ญตฺติวิปนฺนํ ยํ กมฺมํ, \\ อนุสฺสาวนวิปนฺนํ, \\ อุภเยน วิปนฺนญฺจ, \\ วินยา สตฺถุ ปฏิกุฏฺฐํ, \\ อธมฺมวคฺคํ สมคฺคํ, \\ ธมฺเมเนว จ สามคฺคิํ, \\ จตุวคฺโค ปญฺจวคฺโค, \\ ปโรวีสติวคฺโค จ\textsuperscript{1}, \\ ฐเปตฺวา อุปสมฺปทํ, \\ อพฺภานกมฺเมน สห, \\ ทุเว กมฺเม ฐเปตฺวาน, \\ อพฺภานํ ปญฺจวคฺคิโก, \\ อพฺภาเนกํ ฐเปตฺวาน, \\ สพฺพกมฺมกโร สํโฆ, \\ ภิกฺขุนี สิกฺขมานา จ, \\ ปจฺจกฺขาตนฺติมวตฺถู, \\ อปฺปฏิกมฺเม ทิฏฺฐิยา, \\ ติตฺถิยา ติรจฺฉานคตํ, \\ อรหํ ภิกฺขุนีทูสิ, \\ พฺยญฺชนํ นานาสํวาสํ, \\ ยสฺส สํโฆ กเร กมฺมํ, \\ สมฺพุทฺเธน ปฏิกฺขิตฺตา, \\ ปาริวาสิกจตุตฺโถ, \\ มูลา มานตฺตมพฺเภยฺย, \\ มูลา อรหมานตฺตา, \\ น กมฺมการกา ปญฺจ, \\ ภิกฺขุนี สิกฺขมานา จ, \\ ปจฺจกฺขนฺติมอุมฺมตฺตา, \\ อปฺปฏิกมฺเม ทิฏฺฐิยา, \\ นานาสํวาสกา สีมา, \\ อฏฺฐารสนฺนเมเตสํ, \\ ภิกฺขุสฺส ปกตตฺตสฺส, \\ สุทฺธสฺส ทุนฺนิสาริโต, \\ ปณฺฑโก เถยฺยสํวาโส, \\ มาตุ ปิตุ อรหนฺต, \\ โลหิตุปฺปาทโก เจว, \\ เอกาทสนฺนํ เอเตสํ, \\ หตฺถปาทํ ตทุภยํ, \\ องฺคุลิ อฬกณฺฑรํ, \\ คณฺฑี ลกฺขณกสา จ, \\ ปาปา ปริสกาโณ จ, \\ อิริยาปถทุพฺพโล, \\ อนฺธมูคนฺธพธิโร, \\ อนฺธมูคพธิโร จ, \\ เตสํ โอสารณํ โหติ, \\ ทฏฺฐพฺพา ปฏิกาตพฺพา, \\ ตสฺส อุกฺเขปนา กมฺมา, \\ อาปนฺนํ อนุวตฺตนฺตํ, \\ อาปนฺนํ นานุวตฺตนฺตํ, \\ สมฺมุขา ปฏิปุจฺฉา จ, \\ สติ อมูฬฺหปาปิกา, \\ ปพฺพาชนีย ปฏิสาโร, \\ มูลา มานตฺตอพฺภานา, \\ อญฺญํ กเรยฺย อญฺญสฺส, \\ ตํ ตํ กเรยฺย ตํ ตสฺส, \\ ปจฺจาโรเปยฺย อญฺญญฺญํ, \\ เทฺว เทฺว ตมฺมูลกํ ตสฺส\textsuperscript{1},}{\csromangathabat{ญตฺติวิปนฺนํ ยํ กมฺมํ,}{สมฺปนฺนํ อนุสาวนํ.} \\ \csromangathabat{อนุสฺสาวนวิปนฺนํ,}{สมฺปนฺนํ ญตฺติยา จ ยํ.} \\ \csromangathabat{อุภเยน วิปนฺนญฺจ,}{อญฺญตฺร ธมฺมเมว จ.} \\ \csromangathabat{วินยา สตฺถุ ปฏิกุฏฺฐํ,}{กุปฺปํ อฏฺฐานารหิกํ.} \\ \csromangathabat{อธมฺมวคฺคํ สมคฺคํ,}{ธมฺม ปติรูปานิ เย ทุเว.} \\ \csromangathabat{ธมฺเมเนว จ สามคฺคิํ,}{อนุญฺญาสิ ตถาคโต.} \\ \csromangathabat{จตุวคฺโค ปญฺจวคฺโค,}{ทสวคฺโค จ วีสติ.} \\ \csromangathabat{ปโรวีสติวคฺโค จ\textsuperscript{1},}{สํโฆ ปญฺจวิโธ ตถา.} \\ \csromangathabat{ฐเปตฺวา อุปสมฺปทํ,}{ยญฺจ กมฺมํ ปวารณํ.} \\ \csromangathabat{อพฺภานกมฺเมน สห,}{จตุวคฺเคหิ กมฺมิโก.} \\ \csromangathabat{ทุเว กมฺเม ฐเปตฺวาน,}{มชฺฌเทสูปสมฺปทํ.} \\ \csromangathabat{อพฺภานํ ปญฺจวคฺคิโก,}{สพฺพกมฺเมสุ กมฺมิโก.} \\ \csromangathabat{อพฺภาเนกํ ฐเปตฺวาน,}{เย ภิกฺขู ทสวคฺคิกา.} \\ \csromangathabat{สพฺพกมฺมกโร สํโฆ,}{วีโส สพฺพตฺถ กมฺมิโก.} \\ \csromangathabat{ภิกฺขุนี สิกฺขมานา จ,}{สามเณโร สามเณรี.} \\ \csromangathabat{ปจฺจกฺขาตนฺติมวตฺถู,}{อุกฺขิตฺตาปตฺติทสฺสเน.} \\ \csromangathabat{อปฺปฏิกมฺเม ทิฏฺฐิยา,}{ปณฺฑโก เถยฺยสํวาสกํ.} \\ \csromangathabat{ติตฺถิยา ติรจฺฉานคตํ,}{มาตุ ปิตุ จ ฆาตกํ.} \\ \csromangathabat{อรหํ ภิกฺขุนีทูสิ,}{เภทกํ โลหิตุปฺปาทํ.} \\ \csromangathabat{พฺยญฺชนํ นานาสํวาสํ,}{นานาสีมาย อิทฺธิยา.} \\ \csromangathabat{ยสฺส สํโฆ กเร กมฺมํ,}{โหนฺเตเต จตุวีสติ.} \\ \csromangathabat{สมฺพุทฺเธน ปฏิกฺขิตฺตา,}{น เหเต คณปูรกา.} \\ \csromangathabat{ปาริวาสิกจตุตฺโถ,}{ปริวาสํ ทเทยฺย วา.} \\ \csromangathabat{มูลา มานตฺตมพฺเภยฺย,}{อกมฺมํ น จ กรณํ.} \\ \csromangathabat{มูลา อรหมานตฺตา,}{อพฺภานารหเมว จ.} \\ \csromangathabat{น กมฺมการกา ปญฺจ,}{สมฺพุทฺเธน ปกาสิตา.} \\ \csromangathabat{ภิกฺขุนี สิกฺขมานา จ,}{สามเณโร สามเณริกา.} \\ \csromangathabat{ปจฺจกฺขนฺติมอุมฺมตฺตา,}{ขิตฺตาเวทนทสฺสเน.} \\ \csromangathabat{อปฺปฏิกมฺเม ทิฏฺฐิยา,}{ปณฺฑกาปิ จ พฺยญฺชนา\textsuperscript{1}.} \\ \csromangathabat{นานาสํวาสกา สีมา,}{เวหาสํ ยสฺส กมฺม จ.} \\ \csromangathabat{อฏฺฐารสนฺนเมเตสํ,}{ปฏิกฺโกสํ น รุหติ.} \\ \csromangathabat{ภิกฺขุสฺส ปกตตฺตสฺส,}{รุหติ ปฏิกฺโกสนา.} \\ \csromangathabat{สุทฺธสฺส ทุนฺนิสาริโต,}{พาโล หิ สุนิสฺสาริโต.} \\ \csromangathabat{ปณฺฑโก เถยฺยสํวาโส,}{ปกฺกนฺโต ติรจฺฉานคโต.} \\ \csromangathabat{มาตุ ปิตุ อรหนฺต,}{ทูสโก สํฆเภทโก.} \\ \csromangathabat{โลหิตุปฺปาทโก เจว,}{อุภโตพฺยญฺชโน จ โย.} \\ \csromangathabat{เอกาทสนฺนํ เอเตสํ,}{โอสารณํ น ยุชฺชติ.} \\ \csromangathabat{หตฺถปาทํ ตทุภยํ,}{กณฺณนาสํ ตทูภยํ.} \\ \csromangathabat{องฺคุลิ อฬกณฺฑรํ,}{ผณํ ขุชฺโช จ วามโน.} \\ \csromangathabat{คณฺฑี ลกฺขณกสา จ,}{ลิขิตโก จ สีปที.} \\ \csromangathabat{ปาปา ปริสกาโณ จ,}{กุณี ขญฺโช หโตปิ จ.} \\ \csromangathabat{อิริยาปถทุพฺพโล,}{อนฺโธ มูโค จ พธิโร.} \\ \csromangathabat{อนฺธมูคนฺธพธิโร,}{มูคพธิรเมว จ.} \\ \csromangathabat{อนฺธมูคพธิโร จ,}{ทฺวตฺติํเสเต อนูนกา.} \\ \csromangathabat{เตสํ โอสารณํ โหติ,}{สมฺพุทฺเธน ปกาสิตํ.} \\ \csromangathabat{ทฏฺฐพฺพา ปฏิกาตพฺพา,}{นิสฺสชฺเชตา น วิชฺชติ.} \\ \csromangathabat{ตสฺส อุกฺเขปนา กมฺมา,}{สตฺต โหนฺติ อธมฺมิกา.} \\ \csromangathabat{อาปนฺนํ อนุวตฺตนฺตํ,}{สตฺต เตปิ อธมฺมิกา.} \\ \csromangathabat{อาปนฺนํ นานุวตฺตนฺตํ,}{สตฺต กมฺมา สุธมฺมิกา.} \\ \csromangathabat{สมฺมุขา ปฏิปุจฺฉา จ,}{ปฏิญฺญาย จ การณา.} \\ \csromangathabat{สติ อมูฬฺหปาปิกา,}{ตชฺชนีนิยสฺเสน จ.} \\ \csromangathabat{ปพฺพาชนีย ปฏิสาโร,}{อุกฺเขปปริวาส จ.} \\ \csromangathabat{มูลา มานตฺตอพฺภานา,}{ตเถว อุปสมฺปทา.} \\ \csromangathabat{อญฺญํ กเรยฺย อญฺญสฺส,}{โสฬเสเต อธมฺมิกา.} \\ \csromangathabat{ตํ ตํ กเรยฺย ตํ ตสฺส,}{โสฬเสเต สุธมฺมิกา.} \\ \csromangathabat{ปจฺจาโรเปยฺย อญฺญญฺญํ,}{โสฬเสเต อธมฺมิกา.} \\ \csromangathabat{เทฺว เทฺว ตมฺมูลกํ ตสฺส\textsuperscript{1},}{เตปิ โสฬส ธมฺมิกา.}}
\csromangathasetleft{ญตฺติวิปนฺนํ ยํ กมฺมํ, \\ อนุสฺสาวนวิปนฺนํ, \\ อุภเยน วิปนฺนญฺจ, \\ วินยา สตฺถุ ปฏิกุฏฺฐํ, \\ อธมฺมวคฺคํ สมคฺคํ, \\ ธมฺเมเนว จ สามคฺคิํ, \\ จตุวคฺโค ปญฺจวคฺโค, \\ ปโรวีสติวคฺโค จ\textsuperscript{1}, \\ ฐเปตฺวา อุปสมฺปทํ, \\ อพฺภานกมฺเมน สห, \\ ทุเว กมฺเม ฐเปตฺวาน, \\ อพฺภานํ ปญฺจวคฺคิโก, \\ อพฺภาเนกํ ฐเปตฺวาน, \\ สพฺพกมฺมกโร สํโฆ, \\ ภิกฺขุนี สิกฺขมานา จ, \\ ปจฺจกฺขาตนฺติมวตฺถู, \\ อปฺปฏิกมฺเม ทิฏฺฐิยา, \\ ติตฺถิยา ติรจฺฉานคตํ, \\ อรหํ ภิกฺขุนีทูสิ, \\ พฺยญฺชนํ นานาสํวาสํ, \\ ยสฺส สํโฆ กเร กมฺมํ, \\ สมฺพุทฺเธน ปฏิกฺขิตฺตา, \\ ปาริวาสิกจตุตฺโถ, \\ มูลา มานตฺตมพฺเภยฺย, \\ มูลา อรหมานตฺตา, \\ น กมฺมการกา ปญฺจ, \\ ภิกฺขุนี สิกฺขมานา จ, \\ ปจฺจกฺขนฺติมอุมฺมตฺตา, \\ อปฺปฏิกมฺเม ทิฏฺฐิยา, \\ นานาสํวาสกา สีมา, \\ อฏฺฐารสนฺนเมเตสํ, \\ ภิกฺขุสฺส ปกตตฺตสฺส, \\ สุทฺธสฺส ทุนฺนิสาริโต, \\ ปณฺฑโก เถยฺยสํวาโส, \\ มาตุ ปิตุ อรหนฺต, \\ โลหิตุปฺปาทโก เจว, \\ เอกาทสนฺนํ เอเตสํ, \\ หตฺถปาทํ ตทุภยํ, \\ องฺคุลิ อฬกณฺฑรํ, \\ คณฺฑี ลกฺขณกสา จ, \\ ปาปา ปริสกาโณ จ, \\ อิริยาปถทุพฺพโล, \\ อนฺธมูคนฺธพธิโร, \\ อนฺธมูคพธิโร จ, \\ เตสํ โอสารณํ โหติ, \\ ทฏฺฐพฺพา ปฏิกาตพฺพา, \\ ตสฺส อุกฺเขปนา กมฺมา, \\ อาปนฺนํ อนุวตฺตนฺตํ, \\ อาปนฺนํ นานุวตฺตนฺตํ, \\ สมฺมุขา ปฏิปุจฺฉา จ, \\ สติ อมูฬฺหปาปิกา, \\ ปพฺพาชนีย ปฏิสาโร, \\ มูลา มานตฺตอพฺภานา, \\ อญฺญํ กเรยฺย อญฺญสฺส, \\ ตํ ตํ กเรยฺย ตํ ตสฺส, \\ ปจฺจาโรเปยฺย อญฺญญฺญํ, \\ เทฺว เทฺว ตมฺมูลกํ ตสฺส\textsuperscript{1},}
\csromangathagroup{\csromangathastanza{\csromangathabat{ญตฺติวิปนฺนํ ยํ กมฺมํ,}{สมฺปนฺนํ อนุสาวนํ.} \\ \csromangathabat{อนุสฺสาวน\-วิปนฺนํ,}{สมฺปนฺนํ ญตฺติยา จ ยํ.}}\csromangathastanza{\csromangathabat{อุภเยน วิปนฺนญฺจ,}{อญฺญตฺร ธมฺมเมว จ.} \\ \csromangathabat{วินยา สตฺถุ ปฏิกุฏฺฐํ,}{กุปฺปํ อฏฺฐานารหิกํ.}}\csromangathastanza{\csromangathabat{อธมฺมวคฺคํ สมคฺคํ,}{ธมฺม ปติรูปานิ เย ทุเว.} \\ \csromangathabat{ธมฺเมเนว จ สามคฺคิํ,}{อนุญฺญาสิ ตถาคโต.}}\csromangathastanza{\csromangathabat{จตุวคฺโค ปญฺจวคฺโค,}{ทสวคฺโค จ วีสติ.} \\ \csromangathabat{ปโรวีสติวคฺโค จ\csromansharedfootnote{e71d60070c2b1044}{อติเรกวีสติวคฺโค (สฺยา)},}{สํโฆ ปญฺจวิโธ ตถา.}}\csromangathastanza{\csromangathabat{ฐเปตฺวา อุปสมฺปทํ,}{ยญฺจ กมฺมํ ปวารณํ.} \\ \csromangathabat{อพฺภานกมฺเมน สห,}{จตุวคฺเคหิ กมฺมิโก.}}\csromangathastanza{\csromangathabat{ทุเว กมฺเม ฐเปตฺวาน,}{มชฺฌเทสูปสมฺปทํ.} \\ \csromangathabat{อพฺภานํ ปญฺจวคฺคิโก,}{สพฺพกมฺเมสุ กมฺมิโก.}}\csromangathastanza{\csromangathabat{อพฺภาเนกํ ฐเปตฺวาน,}{เย ภิกฺขู ทสวคฺคิกา.} \\ \csromangathabat{สพฺพกมฺมกโร สํโฆ,}{วีโส สพฺพตฺถ กมฺมิโก.}}\csromangathastanza{\csromanfolio{477}\csromangathabat{ภิกฺขุนี สิกฺขมานา จ,}{สามเณโร สามเณรี.} \\ \csromangathabat{ปจฺจกฺขาตนฺติมวตฺถู,}{อุกฺขิตฺตาปตฺติทสฺสเน.}}\csromangathastanza{\csromangathabat{อปฺปฏิกมฺเม ทิฏฺฐิยา,}{ปณฺฑโก เถยฺย\-สํวาสกํ.} \\ \csromangathabat{ติตฺถิยา ติรจฺฉานคตํ,}{มาตุ ปิตุ จ ฆาตกํ.}}\csromangathastanza{\csromangathabat{อรหํ ภิกฺขุนีทูสิ,}{เภทกํ โลหิตุปฺปาทํ.} \\ \csromangathabat{พฺยญฺชนํ นานาสํวาสํ,}{นานาสีมาย อิทฺธิยา.}}\csromangathastanza{\csromangathabat{ยสฺส สํโฆ กเร กมฺมํ,}{โหนฺเตเต จตุวีสติ.} \\ \csromangathabat{สมฺพุทฺเธน ปฏิกฺขิตฺตา,}{น เหเต คณปูรกา.}}\csromangathastanza{\csromangathabat{ปาริวาสิก\-จตุตฺโถ,}{ปริวาสํ ทเทยฺย วา.} \\ \csromangathabat{มูลา มานตฺตมพฺเภยฺย,}{อกมฺมํ น จ กรณํ.}}\csromangathastanza{\csromangathabat{มูลา อรหมานตฺตา,}{อพฺภานารหเมว จ.} \\ \csromangathabat{น กมฺมการกา ปญฺจ,}{สมฺพุทฺเธน ปกาสิตา.}}\csromangathastanza{\csromangathabat{ภิกฺขุนี สิกฺขมานา จ,}{สามเณโร สามเณริกา.} \\ \csromangathabat{ปจฺจกฺขนฺติมอุมฺมตฺตา,}{ขิตฺตาเวทนทสฺสเน.}}\csromangathastanza{\csromangathabat{อปฺปฏิกมฺเม ทิฏฺฐิยา,}{ปณฺฑกาปิ จ พฺยญฺชนา\csromansharedfootnote{ef204b3c7003b9ba}{อิโต ปรํ สฺยามมูเล ทิยฑฺฒคาถาหิ อภพฺพปุคฺคลา สมตฺตํ ทสฺสิตา.}.} \\ \csromangathabat{นานาสํวาสกา สีมา,}{เวหาสํ ยสฺส กมฺม จ.}}\csromangathastanza{\csromangathabat{อฏฺฐารสนฺนเมเตสํ,}{ปฏิกฺโกสํ น รุหติ.} \\ \csromangathabat{ภิกฺขุสฺส ปกตตฺตสฺส,}{รุหติ ปฏิกฺโกสนา.}}\csromangathastanza{\csromangathabat{สุทฺธสฺส ทุนฺนิสาริโต,}{พาโล หิ สุนิสฺสาริโต.} \\ \csromangathabat{ปณฺฑโก เถยฺยสํวาโส,}{ปกฺกนฺโต ติรจฺฉานคโต.}}\csromangathastanza{\csromangathabat{มาตุ ปิตุ อรหนฺต,}{ทูสโก สํฆเภทโก.} \\ \csromangathabat{โลหิตุปฺปาทโก เจว,}{อุภโตพฺยญฺชโน จ โย.}}\csromangathastanza{\csromangathabat{เอกาทสนฺนํ เอเตสํ,}{โอสารณํ น ยุชฺชติ.} \\ \csromangathabat{หตฺถปาทํ ตทุภยํ,}{กณฺณนาสํ ตทูภยํ.}}\csromangathastanza{\csromangathabat{องฺคุลิ อฬกณฺฑรํ,}{ผณํ ขุชฺโช จ วามโน.} \\ \csromangathabat{คณฺฑี ลกฺขณกสา จ,}{ลิขิตโก จ สีปที.}}\csromangathastanza{\csromanfolio{478}\csromangathabat{ปาปา ปริสกาโณ จ,}{กุณี ขญฺโช หโตปิ จ.} \\ \csromangathabat{อิริยาปถทุพฺพโล,}{อนฺโธ มูโค จ พธิโร.}}\csromangathastanza{\csromangathabat{อนฺธมูคนฺธพธิโร,}{มูคพธิรเมว จ.} \\ \csromangathabat{อนฺธ\-มูค\-พธิโร จ,}{ทฺวตฺติํเสเต อนูนกา.}}\csromangathastanza{\csromangathabat{เตสํ โอสารณํ โหติ,}{สมฺพุทฺเธน ปกาสิตํ.} \\ \csromangathabat{ทฏฺฐพฺพา ปฏิกาตพฺพา,}{นิสฺสชฺเชตา น วิชฺชติ.}}\csromangathastanza{\csromangathabat{ตสฺส อุกฺเขปนา กมฺมา,}{สตฺต โหนฺติ อธมฺมิกา.} \\ \csromangathabat{อาปนฺนํ อนุวตฺตนฺตํ,}{สตฺต เตปิ อธมฺมิกา.}}\csromangathastanza{\csromangathabat{อาปนฺนํ นานุวตฺตนฺตํ,}{สตฺต กมฺมา สุธมฺมิกา.} \\ \csromangathabat{สมฺมุขา ปฏิปุจฺฉา จ,}{ปฏิญฺญาย จ การณา.}}\csromangathastanza{\csromangathabat{สติ อมูฬฺหปาปิกา,}{ตชฺชนีนิยสฺเสน จ.} \\ \csromangathabat{ปพฺพาชนีย ปฏิสาโร,}{อุกฺเขปปริวาส จ.}}\csromangathastanza{\csromangathabat{มูลา มานตฺตอพฺภานา,}{ตเถว อุปสมฺปทา.} \\ \csromangathabat{อญฺญํ กเรยฺย อญฺญสฺส,}{โสฬเสเต อธมฺมิกา.}}\csromangathastanza{\csromangathabat{ตํ ตํ กเรยฺย ตํ ตสฺส,}{โสฬเสเต สุธมฺมิกา.} \\ \csromangathabat{ปจฺจาโรเปยฺย อญฺญญฺญํ,}{โสฬเสเต อธมฺมิกา.} \\ \csromangathabat{เทฺว เทฺว ตมฺมูลกํ ตสฺส\csromansharedfootnote{fbe17010294ed701}{เทฺว เทฺว มูลา กตา ตสฺส (สฺยา), โทโทตมูลกนฺตสฺส (ก)},}{เตปิ โสฬส ธมฺมิกา.}}}

\vspace{\tipitakaparskip}
\csromancenter{เอเกกมูลกํ จกฺกํ, “อธมฺมนฺ”ติ ชิโนพฺรวิ.}
\setlength{\csromangathapagewidth}{0pt}
\setlength{\csromangathawakwidth}{0pt}
\csromangathameasuregroup{อกาสิ ตชฺชนียํ กมฺมํ, \\ อธมฺเมน วคฺคกมฺมํ, \\ ตตฺถาธมฺเมน สมคฺคา, \\ อญฺญตฺถ วคฺคาธมฺเมน, \\ ปติรูเปน วคฺคาปิ, \\ อธมฺเมน สมคฺคา จ, \\ ปติรูปเกน วคฺคา จ, \\ เอเกกมูลกํ กตฺวา, \\ พาลา พฺยตฺตสฺส นิยสฺสํ, \\ ปฏิสารณียํ กมฺมํ, \\ อทสฺสนาปฺปฏิกมฺเม, \\ เตสํ อุกฺเขปนียกมฺมํ, \\ อุปริ นยกมฺมานํ\textsuperscript{1}, \\ เตสํเยว อนุโลมํ, \\ ปสฺสทฺธิ เตสํ กมฺมานํ, \\ ตสฺมิํ ตสฺมิํ ตุ กมฺเมสุ, \\ อกตํ ทุกฺกฏํ เจว, \\ กมฺเม ปสฺสทฺธิยา จาปิ, \\ วิปตฺติพฺยาธิเต ทิสฺวา, \\ ปฏิปฺปสฺสทฺธิมกฺขาสิ,}{\csromangathabat{อกาสิ ตชฺชนียํ กมฺมํ,}{สํโฆ ภณฺฑนการโก.} \\ \csromangathabat{อธมฺเมน วคฺคกมฺมํ,}{อญฺญํ อาวาสํ คจฺฉิ โส.} \\ \csromangathabat{ตตฺถาธมฺเมน สมคฺคา,}{ตสฺส ตชฺชนียํ กรุํ.} \\ \csromangathabat{อญฺญตฺถ วคฺคาธมฺเมน,}{ตสฺส ตชฺชนียํ กรุํ.} \\ \csromangathabat{ปติรูเปน วคฺคาปิ,}{สมคฺคาปิ ตถา กรุํ.} \\ \csromangathabat{อธมฺเมน สมคฺคา จ,}{ธมฺเมน วคฺคเมว จ.} \\ \csromangathabat{ปติรูปเกน วคฺคา จ,}{สมคฺคา จ อิเม ปทา.} \\ \csromangathabat{เอเกกมูลกํ กตฺวา,}{จกฺกํ พนฺเธ วิจกฺขโณ.} \\ \csromangathabat{พาลา พฺยตฺตสฺส นิยสฺสํ,}{ปพฺพาเช กุลทูสกํ.} \\ \csromangathabat{ปฏิสารณียํ กมฺมํ,}{กเร อกฺโกสกสฺส จ.} \\ \csromangathabat{อทสฺสนาปฺปฏิกมฺเม,}{โย จ ทิฏฺฐิํ น นิสฺสชฺเช.} \\ \csromangathabat{เตสํ อุกฺเขปนียกมฺมํ,}{สตฺถวาเหน ภาสิตํ.} \\ \csromangathabat{อุปริ นยกมฺมานํ\textsuperscript{1},}{ปญฺโญ ตชฺชนียํ นเย.} \\ \csromangathabat{เตสํเยว อนุโลมํ,}{สมฺมา วตฺตติ ยาจิเต\textsuperscript{1}.} \\ \csromangathabat{ปสฺสทฺธิ เตสํ กมฺมานํ,}{เหฏฺฐา กมฺมนเยน จ.} \\ \csromangathabat{ตสฺมิํ ตสฺมิํ ตุ กมฺเมสุ,}{ตตฺรฏฺโฐ จ วิวทติ.} \\ \csromangathabat{อกตํ ทุกฺกฏํ เจว,}{ปุนกาตพฺพกนฺติ จ.} \\ \csromangathabat{กมฺเม ปสฺสทฺธิยา จาปิ,}{เต ภิกฺขู ธมฺมวาทิโน.} \\ \csromangathabat{วิปตฺติพฺยาธิเต ทิสฺวา,}{กมฺมปฺปตฺเต มหามุนิ.} \\ \csromangathabat{ปฏิปฺปสฺสทฺธิมกฺขาสิ,}{สลฺลกตฺโตว โอสธนฺติ.} \\ อิมมฺหิ ขนฺธเก วตฺถูนิ ฉตฺติํสาติ.}
\csromangathameasurewak{อิมมฺหิ ขนฺธเก วตฺถูนิ ฉตฺติํสาติ.}
\csromangathasetleft{อกาสิ ตชฺชนียํ กมฺมํ, \\ อธมฺเมน วคฺคกมฺมํ, \\ ตตฺถาธมฺเมน สมคฺคา, \\ อญฺญตฺถ วคฺคาธมฺเมน, \\ ปติรูเปน วคฺคาปิ, \\ อธมฺเมน สมคฺคา จ, \\ ปติรูปเกน วคฺคา จ, \\ เอเกกมูลกํ กตฺวา, \\ พาลา พฺยตฺตสฺส นิยสฺสํ, \\ ปฏิสารณียํ กมฺมํ, \\ อทสฺสนาปฺปฏิกมฺเม, \\ เตสํ อุกฺเขปนียกมฺมํ, \\ อุปริ นยกมฺมานํ\textsuperscript{1}, \\ เตสํเยว อนุโลมํ, \\ ปสฺสทฺธิ เตสํ กมฺมานํ, \\ ตสฺมิํ ตสฺมิํ ตุ กมฺเมสุ, \\ อกตํ ทุกฺกฏํ เจว, \\ กมฺเม ปสฺสทฺธิยา จาปิ, \\ วิปตฺติพฺยาธิเต ทิสฺวา, \\ ปฏิปฺปสฺสทฺธิมกฺขาสิ,}
\csromangathagroupwithclosermajor{\csromangathastanza{\csromangathabat{อกาสิ ตชฺชนียํ กมฺมํ,}{สํโฆ ภณฺฑนการโก.} \\ \csromangathabat{อธมฺเมน วคฺคกมฺมํ,}{อญฺญํ อาวาสํ คจฺฉิ โส.}}\csromangathastanza{\csromangathabat{ตตฺถาธมฺเมน สมคฺคา,}{ตสฺส ตชฺชนียํ กรุํ.} \\ \csromangathabat{อญฺญตฺถ วคฺคาธมฺเมน,}{ตสฺส ตชฺชนียํ กรุํ.}}\csromangathastanza{\csromangathabat{ปติรูเปน วคฺคาปิ,}{สมคฺคาปิ ตถา กรุํ.} \\ \csromangathabat{อธมฺเมน สมคฺคา จ,}{ธมฺเมน วคฺคเมว จ.}}\csromangathastanza{\csromangathabat{ปติรูปเกน วคฺคา จ,}{สมคฺคา จ อิเม ปทา.} \\ \csromangathabat{เอเกกมูลกํ กตฺวา,}{จกฺกํ พนฺเธ วิจกฺขโณ.}}\csromangathastanza{\csromanfolio{479}\csromangathabat{พาลา พฺยตฺตสฺส นิยสฺสํ,}{ปพฺพาเช กุลทูสกํ.} \\ \csromangathabat{ปฏิสารณียํ กมฺมํ,}{กเร อกฺโกสกสฺส จ.}}\csromangathastanza{\csromangathabat{อทสฺสนาปฺปฏิกมฺเม,}{โย จ ทิฏฺฐิํ น นิสฺสชฺเช.} \\ \csromangathabat{เตสํ อุกฺเขปนีย\-กมฺมํ,}{สตฺถวาเหน ภาสิตํ.}}\csromangathastanza{\csromangathabat{อุปริ นยกมฺมานํ\csromansharedfootnote{3c3b12d18e2d0474}{อุปวินยกมฺมานํ (สฺยา), อุกฺเขปนียกมฺมานํ (ก)},}{ปญฺโญ ตชฺชนียํ นเย.} \\ \csromangathabat{เตสํเยว อนุโลมํ,}{สมฺมา วตฺตติ ยาจิเต\csromansharedfootnote{5bdc27a95981df01}{ยาจติ (สฺยา), ยาจิโต (สี)}.}}\csromangathastanza{\csromangathabat{ปสฺสทฺธิ เตสํ กมฺมานํ,}{เหฏฺฐา กมฺมนเยน จ.} \\ \csromangathabat{ตสฺมิํ ตสฺมิํ ตุ กมฺเมสุ,}{ตตฺรฏฺโฐ จ วิวทติ.}}\csromangathastanza{\csromangathabat{อกตํ ทุกฺกฏํ เจว,}{ปุนกาตพฺพกนฺติ จ.} \\ \csromangathabat{กมฺเม ปสฺสทฺธิยา จาปิ,}{เต ภิกฺขู ธมฺมวาทิโน.}}\csromangathastanza{\csromangathabat{วิปตฺติพฺยาธิเต ทิสฺวา,}{กมฺมปฺปตฺเต มหามุนิ.} \\ \csromangathabat{ปฏิปฺปสฺสทฺธิมกฺขาสิ,}{สลฺลกตฺโตว โอสธนฺติ.}}\csromangathastanza{\csromangathawak{อิมมฺหิ ขนฺธเก วตฺถูนิ ฉตฺติํสาติ.}}}{\textbf{จมฺเปยฺย}\-\textbf{กฺขนฺธโก นิฏฺฐิโต.}}
\csromanmatikaanchor{mtk.280}
\csromantocmarknopage{chapter}{10. โกสมฺพกกฺขนฺธก}{mtk.280}
\bookmark[dest={mtk.280},level=0]{10. โกสมฺพกกฺขนฺธก}
\chapterheadpageread{10. \textbf{โกสมฺพก}\-\textbf{กฺขนฺธก}}
\csromanmatikaanchor{mtk.281}
\csromantocmarkref{section}{271. โกสมฺพกวิวาทกถา}{mtk.281}
\bookmark[dest={mtk.281},level=1]{271. โกสมฺพกวิวาทกถา}
\csromanheader{1}{271. \textbf{โกสมฺพก}\-\textbf{วิวาท}\-\textbf{กถา}}
\proseitem{451}{\csromanfolio{480}เตน สมเยน พุทฺโธ ภควา โกสมฺพิยํ วิหรติ โฆสิตาราเม. เตน โข ปน สมเยน อญฺญตโร ภิกฺขุ อาปตฺติํ อาปนฺโน โหติ, โส ตสฺสา อาปตฺติยา อาปตฺติทิฏฺฐิ\csromansharedfootnote{7752a56add6798d6}{อาปตฺติทิฏฺฐี (สี)} โหติ, อญฺเญ ภิกฺขู ตสฺสา อาปตฺติยา อนาปตฺติ\-ทิฏฺฐิโน โหนฺติ, โส อปเรน สมเยน ตสฺสา อาปตฺติยา อนาปตฺติทิฏฺฐิ โหติ, อญฺเญ ภิกฺขู ตสฺสา อาปตฺติยา อาปตฺติทิฏฺฐิโน โหนฺติ. อถ โข เต ภิกฺขู ตํ ภิกฺขุํ เอตทโวจุํ “อาปตฺติํ ตฺวํ อาวุโส อาปนฺโน, ปสฺสเสตํ อาปตฺตินฺ”ติ. นตฺถิ เม อาวุโส อาปตฺติ, ยมหํ ปสฺเสยฺยนฺติ. อถ โข เต ภิกฺขู สามคฺคิํ ลภิตฺวา ตํ ภิกฺขุํ อาปตฺติยา อทสฺสเน อุกฺขิปิํสุ. โส จ ภิกฺขุ พหุสฺสุโต โหติ อาคตาคโม ธมฺมธโร วินยธโร มาติกาธโร ปณฺฑิโต พฺยตฺโต เมธาวี ลชฺชี กุกฺกุจฺจโก สิกฺขากาโม. อถ โข โส ภิกฺขุ สนฺทิฏฺเฐ สมฺภตฺเต ภิกฺขู อุปสงฺกมิตฺวา เอตทโวจ “อนาปตฺติ เอสา อาวุโส, เนสา อาปตฺติ, อนาปนฺโนมฺหิ, นมฺหิ อาปนฺโน, อนุกฺขิตฺโตมฺหิ, นมฺหิ อุกฺขิตฺโต, อธมฺมิเกนมฺหิ กมฺเมน อุกฺขิตฺโต กุปฺเปน อฏฺฐานารเหน, โหถ เม อายสฺมนฺโต ธมฺมโต วินยโต ปกฺขา”ติ. อลภิ โข โส ภิกฺขุ สนฺทิฏฺเฐ สมฺภตฺเต ภิกฺขู ปกฺเข. ชานปทานมฺปิ สนฺทิฏฺฐานํ สมฺภตฺตานํ ภิกฺขูนํ สนฺติเก ทูตํ ปาเหสิ “อนาปตฺติ เอสา อาวุโส, เนสา อาปตฺติ, อนาปนฺโนมฺหิ, นมฺหิ อาปนฺโน, อนุกฺขิตฺโตมฺหิ, นมฺหิ อุกฺขิตฺโต, อธมฺมิเกนมฺหิ กมฺเมน อุกฺขิตฺโต กุปฺเปน อฏฺฐานารเหน, โหนฺตุ เม อายสฺมนฺโต ธมฺมโต วินยโต ปกฺขา”ติ. อลภิ โข โส ภิกฺขุ ชานปเทปิ สนฺทิฏฺเฐ สมฺภตฺเต ภิกฺขู ปกฺเข. อถ โข เต อุกฺขิตฺตา\-นุวตฺตกา ภิกฺขู เยน อุกฺเขปกา ภิกฺขู เตนุ\-ปสงฺกมิํสุ, อุปสงฺกมิตฺวา อุกฺเขปเก ภิกฺขู เอตทโวจุํ “อนาปตฺติ เอสา อาวุโส, เนสา อาปตฺติ, อนาปนฺโน เอโส ภิกฺขุ, เนโส ภิกฺขุ อาปนฺโน, อนุกฺขิตฺโต เอโส ภิกฺขุ, เนโส ภิกฺขุ อุกฺขิตฺโต, อธมฺมิเกน กมฺเมน อุกฺขิตฺโต \csromanfolio{481}กุปฺเปน อฏฺฐานารเหนา”ติ. เอวํ วุตฺเต อุกฺเขปกา ภิกฺขู อุกฺขิตฺตา\-นุวตฺตเก ภิกฺขู เอตทโวจุํ “อาปตฺติ เอสา อาวุโส, เนสา อนาปตฺติ, อาปนฺโน เอโส ภิกฺขุ, เนโส ภิกฺขุ อนาปนฺโน, อุกฺขิตฺโต เอโส ภิกฺขุ, เนโส ภิกฺขุ อนุกฺขิตฺโต, ธมฺมิเกน กมฺเมน อุกฺขิตฺโต อกุปฺเปน ฐานารเหน, มา โข ตุเมฺห อายสฺมนฺโต เอตํ อุกฺขิตฺตกํ ภิกฺขุํ อนุวตฺติตฺถ อนุปริวาเรถา”ติ. เอวมฺปิ โข เต อุกฺขิตฺตา\-นุวตฺตกา ภิกฺขู อุกฺเขปเกหิ ภิกฺขูหิ วุจฺจมานา ตเถว ตํ อุกฺขิตฺตกํ ภิกฺขุํ อนุวตฺติํสุ อนุปริวาเรสุํ.}

\proseitem{452}{อถ โข อญฺญตโร ภิกฺขุ เยน ภควา เตนุปสงฺกมิ, อุปสงฺกมิตฺวา ภควนฺตํ อภิวาเทตฺวา เอกมนฺตํ นิสีทิ, เอกมนฺตํ นิสินฺโน โข โส ภิกฺขุ ภควนฺตํ เอตทโวจ\csromandash{}อิธ ภนฺเต อญฺญตโร ภิกฺขุ อาปตฺติํ อาปนฺโน อโหสิ, โส ตสฺสา อาปตฺติยา อาปตฺติทิฏฺฐิ อโหสิ, อญฺเญ ภิกฺขู ตสฺสา อาปตฺติยา อนาปตฺติ\-ทิฏฺฐิโน อเหสุํ, โส อปเรน สมเยน ตสฺสา อาปตฺติยา อนาปตฺติทิฏฺฐิ อโหสิ, อญฺเญ ภิกฺขู ตสฺสา อาปตฺติยา อาปตฺติทิฏฺฐิโน อเหสุํ. อถ โข เต ภนฺเต ภิกฺขู ตํ ภิกฺขุํ เอตทโวจุํ “อาปตฺติํ ตฺวํ อาวุโส อาปนฺโน, ปสฺสเสตํ อาปตฺตินฺ”ติ. นตฺถิ เม อาวุโส อาปตฺติ, ยมหํ ปสฺเสยฺยนฺติ. อถ โข เต ภนฺเต ภิกฺขู สามคฺคิํ ลภิตฺวา ตํ ภิกฺขุํ อาปตฺติยา อทสฺสเน อุกฺขิปิํสุ, โส จ ภนฺเต ภิกฺขุ พหุสฺสุโต อาคตาคโม ธมฺมธโร วินยธโร มาติกาธโร ปณฺฑิโต พฺยตฺโต เมธาวี ลชฺชี กุกฺกุจฺจโก สิกฺขากาโม. อถ โข โส ภนฺเต ภิกฺขุ สนฺทิฏฺเฐ สมฺภตฺเต ภิกฺขู อุปสงฺกมิตฺวา เอตทโวจ “อนาปตฺติ เอสา อาวุโส, เนสา อาปตฺติ, อนาปนฺโนมฺหิ, นมฺหิ อาปนฺโน, อนุกฺขิตฺโตมฺหิ, นมฺหิ อุกฺขิตฺโต, อธมฺมิเกนมฺหิ กมฺเมน อุกฺขิตฺโต กุปฺเปน อฏฺฐานารเหน, โหถ เม อายสฺมนฺโต ธมฺมโต วินยโต ปกฺขา”ติ. อลภิ โข โส ภนฺเต ภิกฺขุ สนฺทิฏฺเฐ สมฺภตฺเต ภิกฺขู ปกฺเข. ชานปทานมฺปิ สนฺทิฏฺฐานํ สมฺภตฺตานํ ภิกฺขูนํ สนฺติเก ทูตํ ปาเหสิ “อนาปตฺติ เอสา อาวุโส, เนสา อาปตฺติ, อนาปนฺโนมฺหิ, นมฺหิ อาปนฺโน, อนุกฺขิตฺโตมฺหิ, นมฺหิ อุกฺขิตฺโต, อธมฺมิเกนมฺหิ กมฺเมน อุกฺขิตฺโต กุปฺเปน อฏฺฐานารเหน, โหนฺตุ เม อายสฺมนฺโต ธมฺมโต วินยโต ปกฺขา”ติ. อลภิ โข โส ภนฺเต ภิกฺขุ ชานปเทปิ สนฺทิฏฺเฐ \csromanfolio{482}สมฺภตฺเต ภิกฺขู ปกฺเข. อถ โข เต ภนฺเต อุกฺขิตฺตา\-นุวตฺตกา ภิกฺขู เยน อุกฺเขปกา ภิกฺขู เตนุ\-ปสงฺกมิํสุ, อุปสงฺกมิตฺวา อุกฺเขปเก ภิกฺขู เอตทโวจุํ “อนาปตฺติ เอสา อาวุโส, เนสา อาปตฺติ, อนาปนฺโน เอโส ภิกฺขุ, เนโส ภิกฺขุ อาปนฺโน, อนุกฺขิตฺโต เอโส ภิกฺขุ, เนโส ภิกฺขุ อุกฺขิตฺโต, อธมฺมิเกน กมฺเมน อุกฺขิตฺโต กุปฺเปน อฏฺฐานารเหนา”ติ. เอวํ วุตฺเต เต ภนฺเต อุกฺเขปกา ภิกฺขู อุกฺขิตฺตา\-นุวตฺตเก ภิกฺขู เอตทโวจุํ “อาปตฺติ เอสา อาวุโส, เนสา อนาปตฺติ, อาปนฺโน เอโส ภิกฺขุ, เนโส ภิกฺขุ อนาปนฺโน, อุกฺขิตฺโต เอโส ภิกฺขุ, เนโส ภิกฺขุ อนุกฺขิตฺโต, ธมฺมิเกน กมฺเมน อุกฺขิตฺโต อกุปฺเปน ฐานารเหน, มา โข ตุเมฺห อายสฺมนฺโต เอตํ อุกฺขิตฺตกํ ภิกฺขุํ อนุวตฺติตฺถ อนุปริวาเรถา”ติ. เอวมฺปิ โข เต ภนฺเต อุกฺขิตฺตา\-นุวตฺตกา ภิกฺขู อุกฺเขปเกหิ ภิกฺขูหิ วุจฺจมานา ตเถว ตํ อุกฺขิตฺตกํ ภิกฺขุํ อนุวตฺตนฺติ อนุปริ\-วาเร\-นฺตีติ.}

\proseitem{453}{อถ โข ภควา “ภินฺโน ภิกฺขุสํโฆ, ภินฺโน ภิกฺขุสํโฆ”ติ อุฏฺฐายาสนา เยน อุกฺเขปกา ภิกฺขู เตนุปสงฺกมิ, อุปสงฺกมิตฺวา ปญฺญตฺเต อาสเน นิสีทิ, นิสชฺช โข ภควา อุกฺเขปเก ภิกฺขู เอตทโวจ\csromandash{}มา โข ตุเมฺห ภิกฺขเว “ปฏิภาติ โน, ปฏิภาติ โน”ติ ยสฺมิํ วา ตสฺมิํ วา ภิกฺขุํ อุกฺขิปิตพฺพํ มญฺญิตฺถ.}

\prose{อิธ ปน ภิกฺขเว ภิกฺขุ อาปตฺติํ อาปนฺโน โหติ, โส ตสฺสา อาปตฺติยา อนาปตฺติทิฏฺฐิ โหติ, อญฺเญ ภิกฺขู ตสฺสา อาปตฺติยา อาปตฺติทิฏฺฐิโน โหนฺติ, เต เจ ภิกฺขเว ภิกฺขู ตํ ภิกฺขุํ เอวํ ชานนฺติ “อยํ โข อายสฺมา พหุสฺสุโต อาคตาคโม ธมฺมธโร วินยธโร มาติกาธโร ปณฺฑิโต พฺยตฺโต เมธาวี ลชฺชี กุกฺกุจฺจโก สิกฺขากาโม, สเจ มยํ อิมํ ภิกฺขุํ อาปตฺติยา อทสฺสเน อุกฺขิปิสฺสาม, น มยํ อิมินา ภิกฺขุนา สทฺธิํ อุโปสถํ กริสฺสาม, วินา อิมินา ภิกฺขุนา อุโปสถํ กริสฺสาม, ภวิสฺสติ สํฆสฺส ตโตนิทานํ ภณฺฑนํ กลโห วิคฺคโห วิวาโท สํฆเภโท สํฆราชิ สํฆววตฺถานํ สํฆนานากรณนฺ”ติ, เภทครุเกหิ ภิกฺขเว ภิกฺขูหิ น โส ภิกฺขุ อาปตฺติยา อทสฺสเน อุกฺขิปิตพฺโพ.}

\proseitem{453}{อิธ ปน ภิกฺขเว ภิกฺขุ อาปตฺติํ อาปนฺโน โหติ, โส ตสฺสา อาปตฺติยา อนาปตฺติทิฏฺฐิ โหติ, อญฺเญ ภิกฺขู ตสฺสา อาปตฺติยา \csromanfolio{483}อาปตฺติทิฏฺฐิโน โหนฺติ, เต เจ ภิกฺขเว ภิกฺขู ตํ ภิกฺขุํ เอวํ ชานนฺติ “อยํ โข อายสฺมา พหุสฺสุโต อาคตาคโม ธมฺมธโร วินยธโร มาติกาธโร ปณฺฑิโต พฺยตฺโต เมธาวี ลชฺชี กุกฺกุจฺจโก สิกฺขากาโม, สเจ มยํ อิมํ ภิกฺขุํ อาปตฺติยา อทสฺสเน อุกฺขิปิสฺสาม, น มยํ อิมินา ภิกฺขุนา สทฺธิํ ปวาเรสฺสาม, วินา อิมินา ภิกฺขุนา ปวาเรสฺสาม. น มยํ อิมินา ภิกฺขุนา สทฺธิํ สํฆกมฺมํ กริสฺสาม, วินา อิมินา ภิกฺขุนา สํฆกมฺมํ กริสฺสาม. น มยํ อิมินา ภิกฺขุนา สทฺธิํ อาสเน นิสีทิสฺสาม, วินา อิมินา ภิกฺขุนา อาสเน นิสีทิสฺสาม. น มยํ อิมินา ภิกฺขุนา สทฺธิํ ยาคุปาเน นิสีทิสฺสาม, วินา อิมินา ภิกฺขุนา ยาคุปาเน นิสีทิสฺสาม. น มยํ อิมินา ภิกฺขุนา สทฺธิํ ภตฺตคฺเค นิสีทิสฺสาม, วินา อิมินา ภิกฺขุนา ภตฺเตคฺเค นิสีทิสฺสาม. น มยํ อิมินา ภิกฺขุนา สทฺธิํ เอกจฺฉนฺเน วสิสฺสาม. วินา อิมินา ภิกฺขุนา เอกจฺฉนฺเน วสิสฺสาม. น มยํ อิมินา ภิกฺขุนา สทฺธิํ ยถาวุฑฺฒํ อภิวาทนํ ปจฺจุฏฺฐานํ อญฺชลิกมฺมํ สามีจิกมฺมํ กริสฺสาม. วินา อิมินา ภิกฺขุนา ยถาวุฑฺฒํ อภิวาทนํ ปจฺจุฏฺฐานํ อญฺชลิกมฺมํ สามีจิกมฺมํ กริสฺสาม. ภวิสฺสติ สํฆสฺส ตโตนิทานํ ภณฺฑนํ กลโห วิคฺคโห วิวาโท สํฆเภโท สํฆราชิ สํฆววตฺถานํ สํฆนานากรณนฺติ. เภทครุเกหิ ภิกฺขเว ภิกฺขูหิ น โส ภิกฺขุ อาปตฺติยา อทสฺสเน อุกฺขิปิตพฺโพติ.}

\proseitem{454}{อถ โข ภควา อุกฺเขปกานํ ภิกฺขูนํ เอตมตฺถํ ภาสิตฺวา อุฏฺฐายาสนา เยน อุกฺขิตฺตา\-นุวตฺตกา ภิกฺขู เตนุปสงฺกมิ, อุปสงฺกมิตฺวา ปญฺญตฺเต อาสเน นิสีทิ, นิสชฺช โข ภควา อุกฺขิตฺตา\-นุวตฺตเก ภิกฺขู เอตทโวจ\csromandash{}มา โข ตุเมฺห ภิกฺขเว อาปตฺติํ อาปชฺชิตฺวา “นามฺห อาปนฺนา นามฺห อาปนฺนา”ติ อาปตฺติํ น ปฏิกาตพฺพํ มญฺญิตฺถ.}

\prose{อิธ ปน ภิกฺขเว ภิกฺขุ อาปตฺติํ อาปนฺโน โหติ, โส ตสฺสา อาปตฺติยา อนาปตฺติทิฏฺฐิ โหติ, อญฺเญ ภิกฺขู ตสฺสา อาปตฺติยา อาปตฺติทิฏฺฐิโน โหนฺติ, โส เจ ภิกฺขเว ภิกฺขุ เต ภิกฺขู เอวํ ชานาติ “อิเม โข อายสฺมนฺโต\csromansharedfootnote{72ff93293253703f}{อายสฺมนฺตา (ก)} พหุสฺสุตา อาคตาคมา ธมฺมธรา วินยธรา มาติกาธรา ปณฺฑิตา พฺยตฺตา เมธาวิโน ลชฺชิโน กุกฺกุจฺจกา สิกฺขากามา, นาลํ มมํ วา การณา อญฺเญสํ วา การณา ฉนฺทา \csromanfolio{484}โทสา โมหา ภยา อคติํ คนฺตุํ, สเจ มํ อิเม ภิกฺขู อาปตฺติยา อทสฺสเน อุกฺขิปิสฺสนฺติ, น มยา สทฺธิํ อุโปสถํ กริสฺสนฺติ, วินา มยา อุโปสถํ กริสฺสนฺติ, ภวิสฺสติ สํฆสฺส ตโตนิทานํ ภณฺฑนํ กลโห วิคฺคโห วิวาโท สํฆเภโท สํฆราชิ สํฆววตฺถานํ สํฆนานากรณนฺ”ติ, เภทครุเกน ภิกฺขเว ภิกฺขุนา ปเรสมฺปิ สทฺธาย อาปตฺติ เทเสตพฺพา.}

\prose{อิธ ปน ภิกฺขเว ภิกฺขุ อาปตฺติํ อาปนฺโน โหติ, โส ตสฺสา อาปตฺติยา อนาปตฺติทิฏฺฐิ โหติ, อญฺเญ ภิกฺขู ตสฺสา อาปตฺติยา อาปตฺติทิฏฺฐิโน โหนฺติ, โส เจ ภิกฺขเว ภิกฺขุ เต ภิกฺขู เอวํ ชานาติ “อิเม โข อายสฺมนฺโต พหุสฺสุตา อาคตาคมา ธมฺมธรา วินยธรา มาติกาธรา ปณฺฑิตา พฺยตฺตา เมธาวิโน ลชฺชิโน กุกฺกุจฺจกา สิกฺขากามา, นาลํ มมํ วา การณา อญฺเญสํ วา การณา ฉนฺทา โทสา โมหา ภยา อคติํ คนฺตุํ, สเจ มํ อิเม ภิกฺขู อาปตฺติยา อทสฺสเน อุกฺขิปิสฺสนฺติ, น มยา สทฺธิํ ปวาเรสฺสนฺติ, วินา มยา ปวาเรสฺสนฺติ. น มยา สทฺธิํ สํฆกมฺมํ กริสฺสนฺติ, วินา มยา สํฆกมฺมํ กริสฺสนฺติ. น มยา สทฺธิํ อาสเน นิสีทิสฺสนฺติ, วินา มยา อาสเน นิสีทิสฺสนฺติ. น มยา สทฺธิํ ยาคุปาเน นิสีทิสฺสนฺติ, วินา มยา ยาคุปาเน นิสีทิสฺสนฺติ. น มยา สทฺธิํ ภตฺตคฺเค นิสีทิสฺสนฺติ, วินา มยา ภตฺตคฺเค นิสีทิสฺสนฺติ. น มยา สทฺธิํ เอกจฺฉนฺเน วสิสฺสนฺติ, วินา มยา เอกจฺฉนฺเน วสิสฺสนฺติ. น มยา สทฺธิํ ยถาวุฑฺฒํ อภิวาทนํ ปจฺจุฏฺฐานํ อญฺชลิกมฺมํ สามีจิกมฺมํ กริสฺสนฺติ, วินา มยา ยถาวุฑฺฒํ อภิวาทนํ ปจฺจุฏฺฐานํ อญฺชลิกมฺมํ สามีจิกมฺมํ กริสฺสนฺติ, ภวิสฺสติ สํฆสฺส ตโตนิทานํ ภณฺฑนํ กลโห วิคฺคโห วิวาโท สํฆเภโท สํฆราชิ สํฆววตฺถานํ สํฆนานากรณนฺ”ติ. เภทครุเกน ภิกฺขเว ภิกฺขุนา ปเรสมฺปิ สทฺธาย อาปตฺติ เทเสตพฺพาติ. อถ โข ภควา อุกฺขิตฺตา\-นุวตฺตกานํ ภิกฺขูนํ เอตมตฺถํ ภาสิตฺวา อุฏฺฐายาสนา ปกฺกามิ.}

\proseitem{455}{เตน โข ปน สมเยน อุกฺขิตฺตา\-นุวตฺตกา ภิกฺขู ตตฺเถว อนฺโตสีมาย อุโปสถํ กโรนฺติ, สํฆกมฺมํ กโรนฺติ. อุกฺเขปกา ปน ภิกฺขู นิสฺสีมํ คนฺตฺวา อุโปสถํ กโรนฺติ, สํฆกมฺมํ กโรนฺติ. อถ โข อญฺญตโร อุกฺเขปโก ภิกฺขุ เยน ภควา เตนุปสงฺกมิ, อุปสงฺกมิตฺวา ภควนฺตํ อภิวาเทตฺวา เอกมนฺตํ นิสีทิ, เอกมนฺตํ นิสินฺโน \csromanfolio{485}โข โส ภิกฺขุ ภควนฺตํ เอตทโวจ “เต ภนฺเต อุกฺขิตฺตา\-นุวตฺตกา ภิกฺขู ตตฺเถว อนฺโตสีมาย อุโปสถํ กโรนฺติ, สํฆกมฺมํ กโรนฺติ, มยํ ปน อุกฺเขปกา ภิกฺขู นิสฺสีมํ คนฺตฺวา อุโปสถํ กโรม, สํฆกมฺมํ กโรมา”ติ. เต เจ ภิกฺขุ อุกฺขิตฺตา\-นุวตฺตกา ภิกฺขู ตตฺเถว อนฺโตสีมาย อุโปสถํ กริสฺสนฺติ, สํฆกมฺมํ กริสฺสนฺติ, ยถา มยา ญตฺติ จ อนุสฺสาวนา จ ปญฺญตฺตา, เตสํ ตานิ กมฺมานิ ธมฺมิกานิ กมฺมานิ ภวิสฺสนฺติ อกุปฺปานิ ฐานารหานิ. ตุเมฺห เจ ภิกฺขุ อุกฺเขปกา ภิกฺขู ตตฺเถว อนฺโตสีมาย อุโปสถํ กริสฺสถ, สํฆกมฺมํ กริสฺสถ, ยถา มยา ญตฺติ จ อนุสฺสาวนา จ ปญฺญตฺตา, ตุมฺหากมฺปิ ตานิ กมฺมานิ ธมฺมิกานิ กมฺมานิ ภวิสฺสนฺติ อกุปฺปานิ ฐานารหานิ. ตํ กิสฺส เหตุ, นานาสํวาสกา เอเต\csromansharedfootnote{ce0751f9b818cc3b}{เต (สฺยา)} ภิกฺขู\csromansharedfootnote{48e9438c3279fb4d}{ภิกฺขุ (สี, สฺยา)} ตุเมฺหหิ, ตุเมฺห จ เตหิ นานาสํวาสกา.}

\prose{เทฺวมา ภิกฺขุ นานาสํวาสก\-ภูมิโย. อตฺตนา วา อตฺตานํ นานาสํวาสกํ กโรติ, สมคฺโค วา นํ สํโฆ อุกฺขิปติ อทสฺสเน วา อปฺปฏิกมฺเม วา อปฺปฏินิสฺสคฺเค วา. อิมา โข ภิกฺขุ เทฺว นานาสํวาสก\-ภูมิโย. เทฺวมา ภิกฺขุ สมานสํวาสก\-ภูมิโย. อตฺตนา วา อตฺตานํ สมานสํวาสํ กโรติ, สมคฺโค วา นํ สํโฆ อุกฺขิตฺตํ โอสาเรติ อทสฺสเน วา อปฺปฏิกมฺเม วา อปฺปฏินิสฺสคฺเค วา. อิมา โข ภิกฺขุ เทฺว สมานสํวาสก\-ภูมิโยติ.}

\proseitem{456}{เตน โข ปน สมเยน ภิกฺขู ภตฺตคฺเค อนฺตรฆเร ภณฺฑนชาตา กลหชาตา วิวาทาปนฺนา อญฺญมญฺญํ อนนุโลมิกํ กายกมฺมํ วจีกมฺมํ อุปทํเสนฺติ, หตฺถ\-ปรามาสํ กโรนฺติ. มนุสฺสา อุชฺฌายนฺติ ขิยฺยนฺติ วิปาเจนฺติ “กถํ หิ นาม สมณา สกฺยปุตฺติยา ภตฺตคฺเค อนฺตรฆเร ภณฺฑนชาตา กลหชาตา วิวาทาปนฺนา อญฺญมญฺญํ อนนุโลมิกํ กายกมฺมํ วจีกมฺมํ อุปทํเสสฺสนฺติ, หตฺถ\-ปรามาสํ กริสฺสนฺตี”ติ. อสฺโสสุํ โข ภิกฺขู เตสํ มนุสฺสานํ อุชฺฌายนฺตานํ ขิยฺยนฺตานํ วิปาเจนฺตานํ. เย เต ภิกฺขู อปฺปิจฺฉา, เต อุชฺฌายนฺติ ขิยฺยนฺติ วิปาเจนฺติ “กถํ หิ นาม ภิกฺขู ภตฺตคฺเค อนฺตรฆเร ภณฺฑนชาตา กลหชาตา วิวาทาปนฺนา อญฺญมญฺญํ อนนุโลมิกํ กายกมฺมํ วจีกมฺมํ \csromanfolio{486}อุปทํเสสฺสนฺติ, หตฺถ\-ปรามาสํ กริสฺสนฺตี”ติ. อถ โข เต ภิกฺขู ภควโต เอตมตฺถํ อาโรเจสุํ ฯเปฯ. สจฺจํ กิร ภิกฺขเว ภิกฺขู ภตฺตคฺเค อนฺตรฆเร ภณฺฑนชาตา ฯเปฯ หตฺถ\-ปรามาสํ กโรนฺตีติ. สจฺจํ ภควาติ. วิครหิ พุทฺโธ ภควา ฯเปฯ วิครหิตฺวา ธมฺมิํ กถํ กตฺวา ภิกฺขู อามนฺเตสิ “ภินฺเน ภิกฺขเว สํเฆ อธมฺมิยายมาเน\csromansharedfootnote{07133c1168c55f88}{อธมฺมิยมาเน (สี, สฺยา, กตฺถจิ) อธมฺมียมาเน (ก)} อสมฺโมทิกาย วตฺตมานาย\csromansharedfootnote{5eecfda1dd2eb586}{“อสมฺโมทิกาวตฺตมานาย” อิติ อฏฺฐกถายํ สํวณฺเณตพฺพปาโฐ.} ‘เอตฺตาวตา น อญฺญมญฺญํ อนนุโลมิกํ กายกมฺมํ วจีกมฺมํ อุปทํเสสฺสาม, หตฺถ\-ปรามาสํ กริสฺสามา’ติ อาสเน นิสีทิตพฺพํ, ภินฺเน ภิกฺขเว สํเฆ ธมฺมิยายมาเน สมฺโมทิกาย วตฺตมานาย อาสนนฺตริกาย นิสีทิตพฺพนฺ”ติ.}

\proseitem{457}{\csromansymbolfootnote{*}{ม 3. 191 ปิฏฺเฐ โถกํ วิสทิสํ.}เตน โข ปน สมเยน ภิกฺขู สํฆมชฺเฌ ภณฺฑนชาตา กลหชาตา วิวาทาปนฺนา อญฺญมญฺญํ มุขสตฺตีหิ วิตุทนฺตา วิหรนฺติ, เต น สกฺโกนฺติ ตํ อธิกรณํ วูปสเมตุํ. อถ โข อญฺญตโร ภิกฺขุ เยน ภควา เตนุปสงฺกมิ, อุปสงฺกมิตฺวา ภควนฺตํ อภิวาเทตฺวา เอกมนฺตํ อฏฺฐาสิ, เอกมนฺตํ ฐิโต โข โส ภิกฺขุ ภควนฺตํ เอตทโวจ “อิธ ภนฺเต ภิกฺขู สํฆมชฺเฌ ภณฺฑนชาตา กลหชาตา วิวาทาปนฺนา อญฺญมญฺญํ มุขสตฺตีหิ วิตุทนฺตา วิหรนฺติ, เต น สกฺโกนฺติ ตํ อธิกรณํ วูปสเมตุํ, สาธุ ภนฺเต ภควา เยน เต ภิกฺขู เตนุ\-ปสงฺกมตุ อนุกมฺปํ อุปาทายา”ติ. อธิวาเสสิ ภควา ตุณฺหีภาเวน. อถ โข ภควา เยน เต ภิกฺขู เตนุปสงฺกมิ, อุปสงฺกมิตฺวา ปญฺญตฺเต อาสเน นิสีทิ, นิสชฺช โข ภควา เต ภิกฺขู เอตทโวจ “อลํ ภิกฺขเว มา ภณฺฑนํ มา กลหํ มา วิคฺคหํ มา วิวาทนฺ”ติ. เอวํ วุตฺเต อญฺญตโร อธมฺมวาที ภิกฺขุ ภควนฺตํ เอตทโวจ “อาคเมตุ ภนฺเต ภควา ธมฺมสฺสามี, อปฺโปสฺสุกฺโก ภนฺเต ภควา ทิฏฺฐ\-ธมฺม\-สุข\-วิหารม\-นุยุตฺโต วิหรตุ, มยเมเตน ภณฺฑเนน กลเหน วิคฺคเหน วิวาเทน ปญฺญายิสฺสามา”ติ. ทุติยมฺปิ โข ภควา เต ภิกฺขู เอตทโวจ “อลํ ภิกฺขเว มา ภณฺฑนํ มา กลหํ มา วิคฺคหํ มา วิวาทนฺ”ติ. ทุติยมฺปิ โข โส อธมฺมวาที ภิกฺขุ ภควนฺตํ เอตทโวจ “อาคเมตุ ภนฺเต ภควา ธมฺมสฺสามี, อปฺโปสฺสุกฺโก ภนฺเต ภควา ทิฏฺฐ\-ธมฺม\-สุข\-วิหารม\-นุยุตฺโต \csromanfolio{487}วิหรตุ, มยเมเตน ภณฺฑเนน กลเหน วิคฺคเหน วิวาเทน ปญฺญายิสฺสามา”ติ.}

\csromanmatikaanchor{mtk.282}
\csromantocmarkref{section}{272. ทีฆาวุวตฺถุ}{mtk.282}
\bookmark[dest={mtk.282},level=1]{272. ทีฆาวุวตฺถุ}
\csromanheader{1}{272. \textbf{ทีฆาวุวตฺถุ}}
\proseitem{458}{อถ โข ภควา ภิกฺขู อามนฺเตสิ “ภูตปุพฺพํ ภิกฺขเว พาราณสิยํ\csromansharedfootnote{60f01dee4d664c73}{(วชิรพุทฺธิฏีกา โอโลเกตพฺพา)} พฺรหฺมทตฺโต นาม กาสิราชา อโหสิ อฑฺโฒ มหทฺธโน มหาโภโค มหพฺพโล มหาวาหโน มหาวิชิโต ปริปุณฺณ\-โกส\-โกฏฺฐาคาโร. ทีฆีติ นาม โกสลราชา อโหสิ ทลิทฺโท อปฺปธโน อปฺปโภโค อปฺปพโล อปฺปวาหโน อปฺปวิชิโต อปริปุณฺณ\-โกสโกฏฺฐาคาโร. อถ โข ภิกฺขเว พฺรหฺมทตฺโต กาสิราชา จตุรงฺคินิํ เสนํ สนฺนยฺหิตฺวา ทีฆีติํ โกสลราชานํ อพฺภุยฺยาสิ. อสฺโสสิ โข ภิกฺขเว ทีฆีติ โกสลราชา “พฺรหฺมทตฺโต กิร กาสิราชา จตุรงฺคินิํ เสนํ สนฺนยฺหิตฺวา มมํ อพฺภุยฺยาโต”ติ. อถ โข ภิกฺขเว ทีฆีติสฺส โกสลรญฺโญ เอตทโหสิ “พฺรหฺมทตฺโต โข กาสิราชา อฑฺโฒ มหทฺธโน มหาโภโค มหพฺพโล มหาวาหโน มหาวิชิโต ปริปุณฺณ\-โกส\-โกฏฺฐาคาโร, อหํ ปนมฺหิ ทลิทฺโท อปฺปธโน อปฺปโภโค อปฺปพโล อปฺปวาหโน อปฺปวิชิโต อปริปุณฺณ\-โกสโกฏฺฐาคาโร, นาหํ ปฏิพโล พฺรหฺมทตฺเตน กาสิรญฺญา เอกสํฆาตมฺปิ สหิตุํ, ยนฺนูนาหํ ปฏิกจฺเจว นครมฺหา นิปฺปเตยฺยนฺ”ติ.}

\prose{อถ โข ภิกฺขเว ทีฆีติ โกสลราชา มเหสิํ อาทาย ปฏิกจฺเจว นครมฺหา นิปฺปติ. อถ โข ภิกฺขเว พฺรหฺมทตฺโต กาสิราชา ทีฆีติสฺส โกสลรญฺโญ พลญฺจ วาหนญฺจ ชนปทญฺจ โกสญฺจ โกฏฺฐาคารญฺจ อภิวิชิย อชฺฌาวสติ. อถ โข ภิกฺขเว ทีฆีติ โกสลราชา สปชาปติโก เยน พาราณสี เตน ปกฺกามิ. อนุปุพฺเพน เยน พาราณสี ตทวสริ, ตตฺร สุทํ ภิกฺขเว ทีฆีติ โกสลราชา สปชาปติโก พาราณสิยํ อญฺญตรสฺมิํ ปจฺจนฺติเม โอกาเส กุมฺภการนิเวสเน อญฺญาตกเวเสน ปริพฺพาชก\-จฺฉนฺเนน ปฏิวสติ. อถ โข ภิกฺขเว ทีฆีติสฺส โกสลรญฺโญ มเหสี นจิรสฺเสว คพฺภินี อโหสิ, ตสฺสา เอวรูโป โทหโฬ อุปฺปนฺโน โหติ “อิจฺฉติ สูริยสฺส อุคฺคมนกาเล จตุรงฺคินิํ เสนํ สนฺนทฺธํ วมฺมิกํ\csromansharedfootnote{f4100ef6efdec38c}{วมฺมิตํ (สี)} สุภูเม ฐิตํ \csromanfolio{488}ปสฺสิตุํ, ขคฺคานญฺจ โธวนํ ปาตุํ”. อถ โข ภิกฺขเว ทีฆีติสฺส โกสลรญฺโญ มเหสี ทีฆีติํ โกสลราชานํ เอตทโวจ คพฺภินีมฺหิ เทว, ตสฺสา เม เอวรูโป โทหโฬ อุปฺปนฺโน “อิจฺฉามิ สูริยสฺส อุคฺคมนกาเล จตุรงฺคินิํ เสนํ สนฺนทฺธํ วมฺมิกํ สุภูเม ฐิตํ ปสฺสิตุํ, ขคฺคานญฺจ โธวนํ ปาตุนฺ”ติ. กุโต เทวิ อมฺหากํ ทุคฺคตานํ จตุรงฺคินี เสนา สนฺนทฺธา วมฺมิกา สุภูเม ฐิตา ขคฺคานญฺจ โธวนํ ปาตุนฺติ\csromansharedfootnote{32b074449eeb64a1}{ขคฺคานญฺจ โธวนนฺติ (สี, สฺยา)}. สจาหํ เทว น ลภิสฺสามิ มริสฺสามีติ.}

\proseitem{459}{เตน โข ปน สมเยน พฺรหฺมทตฺตสฺส กาสิรญฺโญ ปุโรหิโต พฺราหฺมโณ ทีฆีติสฺส โกสลรญฺโญ สหาโย โหติ. อถ โข ภิกฺขเว ทีฆีติ โกสลราชา เยน พฺรหฺมทตฺตสฺส กาสิรญฺโญ ปุโรหิโต พฺราหฺมโณ เตนุปสงฺกมิ, อุปสงฺกมิตฺวา พฺรหฺมทตฺตสฺส กาสิรญฺโญ ปุโรหิตํ พฺราหฺมณํ เอตทโวจ สขี เต สมฺม คพฺภินี, ตสฺสา เอวรูโป โทหโฬ อุปฺปนฺโน “อิจฺฉติ สูริยสฺส อุคฺคมนกาเล จตุรงฺคินิํ เสนํ สนฺนทฺธํ วมฺมิกํ สุภูเม ฐิตํ ปสฺสิตุํ, ขคฺคานญฺจ โธวนํ ปาตุนฺ”ติ. เตน หิ เทว มยมฺปิ เทวิํ ปสฺสามาติ. อถ โข ภิกฺขเว ทีฆีติสฺส โกสลรญฺโญ มเหสี เยน พฺรหฺมทตฺตสฺส กาสิรญฺโญ ปุโรหิโต พฺราหฺมโณ เตนุปสงฺกมิ. อทฺทสา โข ภิกฺขเว พฺรหฺมทตฺตสฺส กาสิรญฺโญ ปุโรหิโต พฺราหฺมโณ ทีฆีติสฺส โกสลรญฺโญ มเหสิํ ทูรโตว อาคจฺฉนฺติํ, ทิสฺวาน อุฏฺฐายาสนา เอกํสํ อุตฺตราสงฺคํ กริตฺวา เยน ทีฆีติสฺส โกสลรญฺโญ มเหสี เตนญฺชลิํ ปณาเมตฺวา ติกฺขตฺตุํ อุทานํ อุทาเนสิ “โกสลราชา วต โภ กุจฺฉิคโต, โกสลราชา วต โภ กุจฺฉิคโต”ติ. อตฺตมนา\csromansharedfootnote{8f4b0bfe722555b8}{อวิมนา (สี, สฺยา, กตฺถจิ)} เทวิ โหหิ, ลจฺฉสิ สูริยสฺส อุคฺคมนกาเล จตุรงฺคินิํ เสนํ สนฺนทฺธํ วมฺมิกํ สุภูเม ฐิตํ ปสฺสิตุํ, ขคฺคานญฺจ โธวนํ ปาตุนฺติ.}

\prose{อถ โข ภิกฺขเว พฺรหฺมทตฺตสฺส กาสิรญฺโญ ปุโรหิโต พฺราหฺมโณ เยน พฺรหฺมทตฺโต กาสิราชา เตนุปสงฺกมิ, อุปสงฺกมิตฺวา พฺรหฺมทตฺตํ กาสิราชานํ เอตทโวจ “ตถา เทว นิมิตฺตานิ ทิสฺสนฺติ, เสฺว สูริยุคฺคมนกาเล จตุรงฺคินี เสนา สนฺนทฺธา วมฺมิกา สุภูเม ติฏฺฐตุ ขคฺคา จ โธวิยนฺตู”ติ. อถ โข ภิกฺขเว พฺรหฺมทตฺโต กาสิราชา มนุสฺเส \csromanfolio{489}อาณาเปสิ “ยถา ภเณ ปุโรหิโต พฺราหฺมโณ อาห, ตถา กโรถา”ติ. อลภิ โข ภิกฺขเว ทีฆีติสฺส โกสลรญฺโญ มเหสี สูริยสฺส อุคฺคมนกาเล จตุรงฺคินิํ เสนํ สนฺนทฺธํ วมฺมิกํ สุภูเม ฐิตํ ปสฺสิตุํ, ขคฺคานญฺจ โธวนํ ปาตุํ. อถ โข ภิกฺขเว ทีฆีติสฺส โกสลรญฺโญ มเหสี ตสฺส คพฺภสฺส ปริปากม\-นฺวาย ปุตฺตํ วิชายิ, ตสฺส ทีฆาวูติ นามํ อกํสุ. อถ โข ภิกฺขเว ทีฆาวุ กุมาโร นจิรสฺเสว วิญฺญุตํ ปาปุณิ. อถ โข ภิกฺขเว ทีฆีติสฺส โกสลรญฺโญ เอตทโหสิ “อยํ โข พฺรหฺมทตฺโต กาสิราชา พหุโน อมฺหากํ อนตฺถสฺส การโก, อิมินา อมฺหากํ พลญฺจ วาหนญฺจ ชนปโท จ โกโส จ โกฏฺฐาคารญฺจ อจฺฉินฺนํ, สจายํ อเมฺห ชานิสฺสติ, สพฺเพว ตโย ฆาตาเปสฺสติ, ยนฺนูนาหํ ทีฆาวุํ กุมารํ พหินคเร วาเสยฺยนฺ”ติ. อถ โข ภิกฺขเว ทีฆีติ โกสลราชา ทีฆาวุํ กุมารํ พหินคเร วาเสสิ. อถ โข ภิกฺขเว ทีฆาวุ กุมาโร พหินคเร ปฏิวสนฺโต นจิรสฺเสว สพฺพสิปฺปานิ สิกฺขิ.}

\proseitem{460}{เตน โข ปน สมเยน ทีฆีติสฺส โกสลรญฺโญ กปฺปโก พฺรหฺมทตฺเต กาสิรญฺเญ ปฏิวสติ. อทฺทสา โข ภิกฺขเว ทีฆีติสฺส โกสลรญฺโญ กปฺปโก ทีฆีติํ โกสลราชานํ สปชาปติกํ พาราณสิยํ อญฺญตรสฺมิํ ปจฺจนฺติเม โอกาเส กุมฺภการนิเวสเน อญฺญาตกเวเสน ปริพฺพาชก\-จฺฉนฺเนน ปฏิวสนฺตํ, ทิสฺวาน เยน พฺรหฺมทตฺโต กาสิราชา เตนุปสงฺกมิ, อุปสงฺกมิตฺวา พฺรหฺมทตฺตํ กาสิราชานํ เอตทโวจ “ทีฆีติ เทว โกสลราชา สปชาปติโก พาราณสิยํ อญฺญตรสฺมิํ ปจฺจนฺติเม โอกาเส กุมฺภการนิเวสเน อญฺญาตกเวเสน ปริพฺพาชก\-จฺฉนฺเนน ปฏิวสตี”ติ. อถ โข ภิกฺขเว พฺรหฺมทตฺโต กาสิราชา มนุสฺเส อาณาเปสิ “เตน หิ ภเณ ทีฆีติํ โกสลราชานํ สปชาปติกํ อาเนถา”ติ. “เอวํ เทวา”ติ โข ภิกฺขเว เต มนุสฺสา พฺรหฺมทตฺตสฺส กาสิรญฺโญ ปฏิสฺสุตฺวา ทีฆีติํ โกสลราชานํ สปชาปติกํ อาเนสุํ. อถ โข ภิกฺขเว พฺรหฺมทตฺโต กาสิราชา มนุสฺเส อาณาเปสิ “เตน หิ ภเณ ทีฆีติํ โกสลราชานํ สปชาปติกํ ทฬฺหาย รชฺชุยา ปจฺฉาพาหํ คาฬฺหพนฺธนํ พนฺธิตฺวา ขุรมุณฺฑํ กริตฺวา ขรสฺสเรน ปณเวน รถิกาย รถิกํ สิงฺฆาฏเกน สิงฺฆาฏกํ ปริเนตฺวา ทกฺขิเณน ทฺวาเรน นิกฺขาเมตฺวา \csromanfolio{490}ทกฺขิณโต นครสฺส จตุธา ฉินฺทิตฺวา จตุทฺทิสา พิลานิ นิกฺขิปถา”ติ. “เอวํ เทวา”ติ โข ภิกฺขเว เต มนุสฺสา พฺรหฺมทตฺตสฺส กาสิรญฺโญ ปฏิสฺสุตฺวา ทีฆีติํ โกสลราชานํ สปชาปติกํ ทฬฺหาย รชฺชุยา ปจฺฉาพาหํ คาฬฺหพนฺธนํ พนฺธิตฺวา ขุรมุณฺฑํ กริตฺวา ขรสฺสเรน ปณเวน รถิกาย รถิกํ สิงฺฆาฏเกน สิงฺฆาฏกํ ปริเนนฺติ.}

\prose{อถ โข ภิกฺขเว ทีฆาวุสฺส กุมารสฺส เอตทโหสิ “จิรํทิฏฺฐา โข เม มาตาปิตโร, ยนฺนูนาหํ มาตาปิตโร ปสฺเสยฺยนฺ”ติ. อถ โข ภิกฺขเว ทีฆาวุ กุมาโร พาราณสิํ ปวิสิตฺวา อทฺทส มาตาปิตโร ทฬฺหาย รชฺชุยา ปจฺฉาพาหํ คาฬฺหพนฺธนํ พนฺธิตฺวา ขุรมุณฺฑํ กริตฺวา ขรสฺสเรน ปณเวน รถิกาย รถิกํ สิงฺฆาฏเกน สิงฺฆาฏกํ ปริเนนฺเต, ทิสฺวาน เยน มาตาปิตโร เตนุปสงฺกมิ. อทฺทสา โข ภิกฺขเว ทีฆีติ โกสลราชา ทีฆาวุํ กุมารํ ทูรโตว อาคจฺฉนฺตํ, ทิสฺวาน ทีฆาวุํ กุมารํ เอตทโวจ “มา โข ตฺวํ ตาต ทีฆาวุ ทีฆํ ปสฺส, มา รสฺสํ, น หิ ตาต ทีฆาวุ เวเรน เวรา สมฺมนฺติ, อเวเรน หิ ตาต ทีฆาวุ เวรา สมฺมนฺตี”ติ. เอวํ วุตฺเต ภิกฺขเว เต มนุสฺสา ทีฆีติํ โกสลราชานํ เอตทโวจุํ “อุมฺมตฺตโก อยํ ทีฆีติ โกสลราชา วิปฺปลปติ, โก อิมสฺส ทีฆาวุ, กํ อยํ เอวมาห มา โข ตฺวํ ตาต ทีฆาวุ ทีฆํ ปสฺส, มา รสฺสํ, น หิ ตาต ทีฆาวุ เวเรน เวรา สมฺมนฺติ, อเวเรน หิ ตาต ทีฆาวุ เวรา สมฺมนฺตี”ติ. นาหํ ภเณ อุมฺมตฺตโก วิปฺปลปามิ, อปิ จ โย วิญฺญู โส วิภาเวสฺสตีติ. ทุติยมฺปิ โข ภิกฺขเว ฯเปฯ. ตติยมฺปิ โข ภิกฺขเว ทีฆีติ โกสลราชา ทีฆาวุํ กุมารํ เอตทโวจ “มา โข ตฺวํ ตาต ทีฆาวุ ทีฆํ ปสฺส, มา รสฺสํ, น หิ ตาต ทีฆาวุ เวเรน เวรา สมฺมนฺติ, อเวเรน หิ ตาต ทีฆาวุ เวรา สมฺมนฺตี”ติ. ตติยมฺปิ โข ภิกฺขเว เต มนุสฺสา ทีฆีติํ โกสลราชานํ เอตทโวจุํ “อุมฺมตฺตโก อยํ ทีฆีติ โกสลราชา วิปฺปลปติ, โก อิมสฺส ทีฆาวุ, กํ อยํ เอวมาห มา โข ตฺวํ ตาต ทีฆาวุ ทีฆํ ปสฺส, มา รสฺสํ, น หิ ตาต ทีฆาวุ เวเรน เวรา สมฺมนฺติ, อเวเรน หิ ตาต ทีฆาวุ เวรา สมฺมนฺตี”ติ. นาหํ ภเณ อุมฺมตฺตโก วิปฺปลปามิ, อปิ จ โย วิญฺญู โส วิภาเวสฺสตีติ. อถ โข ภิกฺขเว เต มนุสฺสา ทีฆีติํ โกสลราชานํ สปชาปติกํ รถิกาย รถิกํ สิงฺฆาฏเกน สิงฺฆาฏกํ \csromanfolio{491}ปริเนตฺวา ทกฺขิเณน ทฺวาเรน นิกฺขาเมตฺวา ทกฺขิณโต นครสฺส จตุธา ฉินฺทิตฺวา จตุทฺทิสา พิลานิ นิกฺขิปิตฺวา คุมฺพํ ฐเปตฺวา ปกฺกมิํสุ. อถ โข ภิกฺขเว ทีฆาวุ กุมาโร พาราณสิํ ปวิสิตฺวา สุรํ นีหริตฺวา คุมฺพิเย ปาเยสิ. ยทา เต มตฺตา อเหสุํ ปติตา, อถ กฏฺฐานิ สํกฑฺฒิตฺวา จิตกํ กริตฺวา มาตาปิตูนํ สรีรํ จิตกํ อาโรเปตฺวา อคฺคิํ ทตฺวา ปญฺชลิโก ติกฺขตฺตุํ จิตกํ ปทกฺขิณํ อกาสิ.}

\proseitem{461}{เตน โข ปน สมเยน พฺรหฺมทตฺโต กาสิราชา อุปริปาสาท\-วรคโต โหติ. อทฺทสา โข ภิกฺขเว พฺรหฺมทตฺโต กาสิราชา ทีฆาวุํ กุมารํ ปญฺชลิกํ ติกฺขตฺตุํ จิตกํ ปทกฺขิณํ กโรนฺตํ, ทิสฺวานสฺส เอตทโหสิ “นิสฺสํสยํ โข โส มนุสฺโส ทีฆีติสฺส โกสลรญฺโญ ญาติ วา สาโลหิโต วา, อโห เม อนตฺถโต, น หิ นาม เม โกจิ อาโรเจสฺสตี”ติ. อถ โข ภิกฺขเว ทีฆาวุ กุมาโร อรญฺญํ คนฺตฺวา ยาวทตฺถํ กนฺทิตฺวา โรทิตฺวา ขปฺปํ\csromansharedfootnote{d74b9ad9476a0cfa}{พปฺปํ (สี, สฺยา)} ปุญฺฉิตฺวา พาราณสิํ ปวิสิตฺวา อนฺเตปุรสฺส สามนฺตา หตฺถิสาลํ คนฺตฺวา หตฺถาจริยํ เอตทโวจ “อิจฺฉามหํ อาจริย สิปฺปํ สิกฺขิตุนฺ”ติ. เตน หิ ภเณ มาณวกสิกฺขสฺสูติ. อถ โข ภิกฺขเว ทีฆาวุ กุมาโร รตฺติยา ปจฺจูสสมยํ ปจฺจุฏฺฐาย หตฺถิสาลายํ มญฺชุนา สเรน คายิ, วีณญฺจ วาเทสิ. อสฺโสสิ โข ภิกฺขเว พฺรหฺมทตฺโต กาสิราชา รตฺติยา ปจฺจูสสมยํ ปจฺจุฏฺฐาย หตฺถิสาลายํ มญฺชุนา สเรน คีตํ วีณญฺจ วาทิตํ, สุตฺวาน มนุสฺเส ปุจฺฉิ “โก ภเณ รตฺติยา ปจฺจูสสมยํ ปจฺจุฏฺฐาย หตฺถิสาลายํ มญฺชุนา สเรน คายิ, วีณญฺจ วาเทสีติ. อมุกสฺส เทว หตฺถา\-จริยสฺส อนฺเตวาสี มาณวโก รตฺติยา ปจฺจูสสมยํ ปจฺจุฏฺฐาย หตฺถิสาลายํ มญฺจุนา สเรน คายิ, วีณญฺจ วาเทสีติ. เตน หิ ภเณ ตํ มาณวกํ อาเนถาติ. “เอวํ เทวา”ติ โข ภิกฺขเว เต มนุสฺสา พฺรหฺมทตฺตสฺส กาสิรญฺโญ ปฏิสฺสุตฺวา ทีฆาวุํ กุมารํ อาเนสุํ. ตฺวํ ภเณ มาณวก รตฺติยา ปจฺจูสสมยํ ปจฺจุฏฺฐาย หตฺถิสาลายํ มญฺชุนา สเรน คายิ, วีณญฺจ วาเทสีติ. เอวํ เทวาติ. เตน หิ ตฺวํ ภเณ มาณวก คายสฺสุ, วีณญฺจ วาเทหีติ. “เอวํ เทวา”ติ โข ภิกฺขเว ทีฆาวุ กุมาโร พฺรหฺมทตฺตสฺส กาสิรญฺโญ ปฏิสฺสุตฺวา อาราธาเปกฺโข มญฺชุนา สเรน \csromanfolio{492}คายิ, วีณญฺจ วาเทสิ. ตฺวํ ภเณ มาณวก มํ อุปฏฺฐหาติ. “เอวํ เทวา”ติ โข ภิกฺขเว ทีฆาวุ กุมาโร พฺรหฺมทตฺตสฺส กาสิรญฺโญ ปจฺจสฺโสสิ. อถ โข ภิกฺขเว ทีฆาวุ กุมาโร พฺรหฺมทตฺตสฺส กาสิรญฺโญ ปุพฺพุฏฺฐายี อโหสิ ปจฺฉานิปาตี กิํการ\-ปฏิสฺสาวี มนาปจารี ปิยวาที. อถ โข ภิกฺขเว พฺรหฺมทตฺโต กาสิราชา ทีฆาวุํ กุมารํ นจิรสฺเสว อพฺภนฺตริเม วิสฺสาสิก\-ฏฺฐาเน ฐเปสิ.}

\proseitem{462}{อถ โข ภิกฺขเว พฺรหฺมทตฺโต กาสิราชา ทีฆาวุํ กุมารํ เอตทโวจ “เตน หิ ภเณ มาณวก รถํ โยเชหิ มิควํ คมิสฺสามา”ติ. “เอวํ เทวา”ติ โข ภิกฺขเว ทีฆาวุ กุมาโร พฺรหฺมทตฺตสฺส กาสิรญฺโญ ปฏิสฺสุตฺวา รถํ โยเชตฺวา พฺรหฺมทตฺตํ กาสิราชานํ เอตทโวจ “ยุตฺโต โข เต เทว รโถ, ยสฺส ทานิ กาลํ มญฺญสี”ติ. อถ โข ภิกฺขเว พฺรหฺมทตฺโต กาสิราชา รถํ อภิรุหิ, ทีฆาวุกุมาโร รถํ เปเสสิ, ตถา ตถา รถํ เปเสสิ, ยถา ยถา อญฺเญเนว เสนา อคมาสิ อญฺเญเนว รโถ. อถ โข ภิกฺขเว พฺรหฺมทตฺโต กาสิราชา ทูรํ คนฺตฺวา ทีฆาวุํ กุมารํ เอตทโวจ “เตน หิ ภเณ มาณวก รถํ มุญฺจสฺสุ, กิลนฺโตมฺหิ นิปชฺชิสฺสามี”ติ. “เอวํ เทวา”ติ โข ภิกฺขเว ทีฆาวุ กุมาโร พฺรหฺมทตฺตสฺส กาสิรญฺโญ ปฏิสฺสุตฺวา รถํ มุญฺจิตฺวา ปถวิยํ ปลฺลงฺเกน นิสีทิ. อถ โข ภิกฺขเว พฺรหฺมทตฺโต กาสิราชา ทีฆาวุสฺส กุมารสฺส อุจฺฉงฺเค สีสํ กตฺวา เสยฺยํ กปฺเปสิ, ตสฺส กิลนฺตสฺส มุหุตฺตเกเนว นิทฺทา โอกฺกมิ. อถ โข ภิกฺขเว ทีฆาวุสฺส กุมารสฺส เอตทโหสิ “อยํ โข พฺรหฺมทตฺโต กาสิราชา พหุโน อมฺหากํ อนตฺถสฺส การโก, อิมินา อมฺหากํ พลญฺจ วาหนญฺจ ชนปโท จ โกโส จ โกฏฺฐาคารญฺจ อจฺฉินฺนํ, อิมินา จ เม มาตาปิตโร หตา, อยํ ขฺวสฺส กาโล โยหํ เวรํ อปฺเปยฺยนฺ”ติ โกสิยา ขคฺคํ นิพฺพาหิ. อถ โข ภิกฺขเว ทีฆาวุสฺส กุมารสฺส เอตทโหสิ “ปิตา โข มํ มรณกาเล อวจ ‘มา โข ตฺวํ ตาต ทีฆาวุ ทีฆํ ปสฺส, มา รสฺสํ, น หิ ตาต ทีฆาวุ เวเรน เวรา สมฺมนฺติ, อเวเรน หิ ตาต ทีฆาวุ เวรา สมฺมนฺตี’ติ, น โข เมตํ ปติรูปํ, ยฺวาหํ ปิตุวจนํ อติกฺกเมยฺยนฺ”ติ โกสิยา ขคฺคํ ปเวเสสิ. ทุติยมฺปิ โข ภิกฺขเว ทีฆาวุสฺส กุมารสฺส เอตทโหสิ “อยํ โข พฺรหฺมทตฺโต กาสิราชา พหุโน อมฺหากํ อนตฺถสฺส การโก, อิมินา อมฺหากํ \csromanfolio{493}พลญฺจ วาหนญฺจ ชนปโท จ โกโส จ โกฏฺฐาคารญฺจ อจฺฉินฺนํ, อิมินา จ เม มาตาปิตโร หตา, อยํ ขฺวสฺส กาโล โยหํ เวรํ อปฺเปยฺยนฺ”ติ โกสิยา ขคฺคํ นิพฺพาหิ. ทุติยมฺปิ โข ภิกฺขเว ทีฆาวุสฺส กุมารสฺส เอตทโหสิ “ปิตา โข มํ มรณกาเล อวจ ‘มา โข ตฺวํ ตาต ทีฆาวุ ทีฆํ ปสฺส, มา รสฺสํ, น หิ ตาต ทีฆาวุ เวเรน เวรา สมฺมนฺติ, อเวเรน หิ ตาต ทีฆาวุ เวรา สมฺมนฺตี’ติ, น โข เมตํ ปติรูปํ, ยฺวาหํ ปิตุวจนํ อติกฺกเมยฺยนฺ”ติ ปุนเทว โกสิยา ขคฺคํ ปเวเสสิ. ตติยมฺปิ โข ภิกฺขเว ทีฆาวุสฺส กุมารสฺส เอตทโหสิ “อยํ โข พฺรหฺมทตฺโต กาสิราชา พหุโน อมฺหากํ อนตฺถสฺส การโก, อิมินา อมฺหากํ พลญฺจ วาหนญฺจ ชนปโท จ โกโส จ โกฏฺฐาคารญฺจ อจฺฉินฺนํ, อิมินา จ เม มาตาปิตโร หตา, อยํ ขฺวสฺส กาโล โยหํ เวรํ อปฺเปยฺยนฺ”ติ โกสิยา ขคฺคํ นิพฺพาหิ. ตติยมฺปิ โข ภิกฺขเว ทีฆาวุสฺส กุมารสฺส เอตทโหสิ “ปิตา โข มํ มรณกาเล อวจ ‘มา โข ตฺวํ ตาต ทีฆาวุ ทีฆํ ปสฺส, มา รสฺสํ, น หิ ตาต ทีฆาวุ เวเรน เวรา สมฺมนฺติ, อเวเรน หิ ตาต ทีฆาวุ เวรา สมฺมนฺตี’ติ, น โข เมตํ ปติรูปํ, ยฺวาหํ ปิตุวจนํ อติกฺกเมยฺยนฺ”ติ ปุนเทว โกสิยา ขคฺคํ ปเวเสสิ. อถ โข ภิกฺขเว พฺรหฺมทตฺโต กาสิราชา ภีโต อุพฺพิคฺโค อุสฺสงฺกี อุตฺรสฺโต สหสา วุฏฺฐาสิ. อถ โข ภิกฺขเว ทีฆาวุ กุมาโร พฺรหฺมทตฺตํ กาสิราชานํ เอตทโวจ “กิสฺส ตฺวํ เทว ภีโต อุพฺพิคฺโค อุสฺสงฺกี อุตฺรสฺโต สหสา วุฏฺฐาสี”ติ. อิธ มํ ภเณ มาณวก ทีฆีติสฺส โกสลรญฺโญ ปุตฺโต ทีฆาวุ กุมาโร สุปินนฺเตน ขคฺเคน ปริปาเตสิ, เตนาหํ ภีโต อุพฺพิคฺโค อุสฺสงฺกี อุตฺรสฺโต สหสา วุฏฺฐาสินฺติ. อถ โข ภิกฺขเว ทีฆาวุ กุมาโร วาเมน หตฺเถน พฺรหฺมทตฺตสฺส กาสิรญฺโญ สีสํ ปรามสิตฺวา ทกฺขิเณน หตฺเถน ขคฺคํ นิพฺพาเหตฺวา พฺรหฺมทตฺตํ กาสิราชานํ เอตทโวจ “อหํ โข โส เทว ทีฆีติสฺส โกสลรญฺโญ ปุตฺโต ทีฆาวุ กุมาโร, พหุโน ตฺวํ อมฺหากํ อนตฺถสฺส การโก, ตยา อมฺหากํ พลญฺจ วาหนญฺจ ชนปโท จ โกโส จ โกฏฺฐาคารญฺจ อจฺฉินฺนํ, ตยา จ เม มาตาปิตโร หตา, อยํ ขฺวสฺส กาโล ยฺวาหํ เวรํ อปฺเปยฺยนฺ”ติ. อถ โข ภิกฺขเว พฺรหฺมทตฺโต กาสิราชา ทีฆาวุสฺส กุมารสฺส ปาเทสุ สิรสา นิปติตฺวา ทีฆาวุํ กุมารํ เอตทโวจ “ชีวิตํ เม ตาต ทีฆาวุ เทหิ, ชีวิตํ เม ตาต ทีฆาวุ เทหี”ติ. กฺยาหํ อุสฺสหามิ เทวสฺส ชีวิตํ \csromanfolio{494}ทาตุํ, เทโว โข เม ชีวิตํ ทเทยฺยาติ. เตน หิ ตาต ทีฆาวุ ตฺวญฺเจว เม ชีวิตํ เทหิ, อหญฺจ เต ชีวิตํ ทมฺมีติ. อถ โข ภิกฺขเว พฺรหฺมทตฺโต จ กาสิราชา ทีฆาวุ จ กุมาโร อญฺญมญฺญสฺส ชีวิตํ อทํสุ, ปาณิญฺจ อคฺคเหสุํ, สปถญฺจ อกํสุ อทฺทูภาย\csromansharedfootnote{ed99292ac9d37547}{อทฺรูภาย, อทุพฺภาย (ก)}.}

\prose{อถ โข ภิกฺขเว พฺรหฺมทตฺโต กาสิราชา ทีฆาวุํ กุมารํ เอตทโวจ “เตน หิ ตาต ทีฆาวุ รถํ โยเชหิ คมิสฺสามา”ติ. “เอวํ เทวา”ติ โข ภิกฺขเว ทีฆาวุ กุมาโร พฺรหฺมทตฺตสฺส กาสิรญฺโญ ปฏิสฺสุตฺวา รถํ โยเชตฺวา พฺรหฺมทตฺตํ กาสิราชานํ เอตทโวจ “ยุตฺโต โข เต เทว รโถ, ยสฺส ทานิ กาลํ มญฺญสี”ติ. อถ โข ภิกฺขเว พฺรหฺมทตฺโต กาสิราชา รถํ อภิรุหิ, ทีฆาวุ กุมาโร รถํ เปเสสิ, ตถา ตถา รถํ เปเสสิ, ยถา ยถา นจิรสฺเสว เสนาย สมาคญฺฉิ. อถ โข ภิกฺขเว พฺรหฺมทตฺโต กาสิราชา พาราณสิํ ปวิสิตฺวา อมจฺเจ ปาริสชฺเช สนฺนิปาตาเปตฺวา เอตทโวจ “สเจ ภเณ ทีฆีติสฺส โกสลรญฺโญ ปุตฺตํ ทีฆาวุํ กุมารํ ปสฺเสยฺยาถ, กินฺติ นํ กเรยฺยาถา”ติ. เอกจฺเจ เอวมาหํสุ “มยํ เทว หตฺเถ ฉินฺเทยฺยาม. มยํ เทว ปาเท ฉินฺเทยฺยาม. มยํ เทว หตฺถปาเท ฉินฺเทยฺยาม. มยํ เทว กณฺเณ ฉินฺเทยฺยาม. มยํ เทว นาสํ ฉินฺเทยฺยาม. มยํ เทว กณฺณนาสํ ฉินฺเทยฺยาม. มยํ เทว สีสํ ฉินฺเทยฺยามา”ติ. อยํ โข ภเณ ทีฆีติสฺส โกสลรญฺโญ ปุตฺโต ทีฆาวุ กุมาโร, นายํ ลพฺภา กิญฺจิ กาตุํ, อิมินา จ เม ชีวิตํ ทินฺนํ, มยา จ อิมสฺส ชีวิตํ ทินฺนนฺติ.}

\proseitemwithcloserruled{463}{อถ โข ภิกฺขเว พฺรหฺมทตฺโต กาสิราชา ทีฆาวุํ กุมารํ เอตทโวจ “ยํ โข เต ตาต ทีฆาวุ ปิตา มรณกาเล อวจ ‘มา โข ตฺวํ ตาต ทีฆาวุ ทีฆํ ปสฺส, มา รสฺสํ, น หิ ตาต ทีฆาวุ เวเรน เวรา สมฺมนฺติ, อเวเรน หิ ตาต ทีฆาวุ เวรา สมฺมนฺตี’ติ, กิํ เต ปิตา สนฺธาย อวจา”ติ. ยํ โข เม เทว ปิตา มรณกาเล อวจ “มา ทีฆนฺ”ติ, “มา จิรํ เวรํ อกาสี”ติ อิมํ โข เม เทว ปิตา มรณกาเล อวจ “มา ทีฆนฺ”ติ. ยํ โข เม เทว ปิตา มรณกาเล อวจ “มา รสฺสนฺ”ติ, “มา ขิปฺปํ มิตฺเตหิ ภิชฺชิตฺถา”ติ อิมํ โข เม เทว ปิตา มรณกาเล \csromanfolio{495}อวจ “มา รสฺสนฺ”ติ. ยํ โข เม เทว ปิตา มรณกาเล อวจ “น หิ ตาต ทีฆาวุ เวเรน เวรา สมฺมนฺติ, อเวเรน หิ ตาต ทีฆาวุ เวรา สมฺมนฺตี”ติ, “เทเวน เม มาตาปิตโร หตา”ติ สจาหํ เทวํ ชีวิตา โวโรเปยฺยํ, เย เทวสฺส อตฺถกามา, เต มํ ชีวิตา โวโรเปยฺยุํ, เย เม อตฺถกามา, เต เต ชีวิตา โวโรเปยฺยุํ, เอวํ ตํ เวรํ เวเรน น วูปสเมยฺย, อิทานิ จ ปน เม เทเวน ชีวิตํ ทินฺนํ, มยา จ เทวสฺส ชีวิตํ ทินฺนํ, เอวํ ตํ เวรํ อเวเรน วูปสนฺตํ, อิมํ โข เม เทว ปิตา มรณกาเล อวจ “น หิ ตาต ทีฆาวุ เวเรน เวรา สมฺมนฺติ, อเวเรน หิ ตาต ทีฆาวุ เวรา สมฺมนฺตี”ติ. อถ โข ภิกฺขเว พฺรหฺมทตฺโต กาสิราชา “อจฺฉริยํ วต โภ, อพฺภุตํ วต โภ, ยาว ปณฺฑิโต อยํ ทีฆาวุ กุมาโร, ยตฺร หิ นาม ปิตุโน สํขิตฺเตน ภาสิตสฺส วิตฺถาเรน อตฺถํ อาชานิสฺสตี”ติ เปตฺติกํ พลญฺจ วาหนญฺจ ชนปทญฺจ โกสญฺจ โกฏฺฐาคารญฺจ ปฏิปาเทสิ, ธีตรญฺจ อทาสิ. เตสํ หิ นาม ภิกฺขเว ราชูนํ อาทินฺนทณฺฑานํ อาทินฺน\-สตฺถานํ เอวรูปํ ขนฺติโสรจฺจํ ภวิสฺสติ. อิธ โข ปน ตํ ภิกฺขเว โสเภถ, ยํ ตุเมฺห เอวํ สฺวากฺขาเต ธมฺมวินเย ปพฺพชิตา สมานา ขมา จ ภเวยฺยาถ โสรตา จาติ. ตติยมฺปิ โข ภควา เต ภิกฺขู เอตทโวจ “อลํ ภิกฺขเว มา ภณฺฑนํ มา กลหํ มา วิคฺคหํ มา วิวาทนฺ”ติ. ตติยมฺปิ โข โส อธมฺมวาที ภิกฺขุ ภควนฺตํ เอตทโวจ “อาคเมตุ ภนฺเต ภควา ธมฺมสฺสามี, อปฺโปสฺสุกฺโก ภนฺเต ภควา ทิฏฺฐ\-ธมฺม\-สุข\-วิหารม\-นุยุตฺโต วิหรตุ, มยเมเตน ภณฺฑเนน กลเหน วิคฺคเหน วิวาเทน ปญฺญายิสฺสามา”ติ. อถ โข ภควา “ปริยาทินฺน\-รูปา โข อิเม โมฆปุริสา, นยิเม สุกรา สญฺญาเปตุนฺ”ติ อุฏฺฐายาสนา ปกฺกามิ.}{ทีฆาวุ\-ภาณวาโร นิฏฺฐิโต ปฐโม.}
\proseitem{464}{\csromansymbolfootnote{*}{ม 3. 192 ปิฏฺเฐปิ.}อถ โข ภควา ปุพฺพณฺหสมยํ นิวาเสตฺวา ปตฺต\-จีวรมา\-ทาย โกสมฺพิํ ปิณฺฑาย ปาวิสิ, โกสมฺพิยํ ปิณฺฑาย จริตฺวา ปจฺฉาภตฺตํ ปิณฺฑปาต\-ปฏิกฺกนฺโต เสนาสนํ สํสาเมตฺวา ปตฺต\-จีวรมา\-ทาย สํฆมชฺเฌ ฐิตโกว อิมา คาถาโย อภาสิ\csromandash{}}

\setlength{\csromangathapagewidth}{0pt}
\setlength{\csromangathawakwidth}{0pt}
\csromangathameasuregroup{\textsuperscript{*}ปุถุสทฺโท สมชโน, \\ สํฆสฺมิํ ภิชฺชมานสฺมิํ, \\ \textsuperscript{*}ปริมุฏฺฐา ปณฺฑิตาภาสา, \\ ยาวิจฺฉนฺติ มุขายามํ, \\ \textsuperscript{*}อกฺโกจฺฉิ มํ อวธิ มํ, \\ เย จ ตํ อุปนยฺหนฺติ, \\ \textsuperscript{*}อกฺโกจฺฉิ มํ อวธิ มํ, \\ เย จ ตํ นุปนยฺหนฺติ, \\ \textsuperscript{*}น หิ เวเรน เวรานิ, \\ อเวเรน จ สมฺมนฺติ, \\ \textsuperscript{*}ปเร จ น วิชานนฺติ, \\ เย จ ตตฺถ วิชานนฺติ, \\ \textsuperscript{*}อฏฺฐิจฺฉินฺนา ปาณหรา, \\ รฏฺฐํ วิลุมฺปมานานํ,}{\csromangathabat{\textsuperscript{*}ปุถุสทฺโท สมชโน,}{น พาโล โกจิ มญฺญถ.} \\ \csromangathabat{สํฆสฺมิํ ภิชฺชมานสฺมิํ,}{นาญฺญํ ภิยฺโย อมญฺญรุํ.} \\ \csromangathabat{\textsuperscript{*}ปริมุฏฺฐา ปณฺฑิตาภาสา,}{วาจาโคจรภาณิโน.} \\ \csromangathabat{ยาวิจฺฉนฺติ มุขายามํ,}{เยน นีตา น ตํ วิทู.} \\ \csromangathabat{\textsuperscript{*}อกฺโกจฺฉิ มํ อวธิ มํ,}{อชินิ มํ อหาสิ เม.} \\ \csromangathabat{เย จ ตํ อุปนยฺหนฺติ,}{เวรํ เตสํ น สมฺมติ.} \\ \csromangathabat{\textsuperscript{*}อกฺโกจฺฉิ มํ อวธิ มํ,}{อชินิ มํ อหาสิ เม.} \\ \csromangathabat{เย จ ตํ นุปนยฺหนฺติ,}{เวรํ เตสูปสมฺมติ.} \\ \csromangathabat{\textsuperscript{*}น หิ เวเรน เวรานิ,}{สมฺมนฺตีธ กุทาจนํ.} \\ \csromangathabat{อเวเรน จ สมฺมนฺติ,}{เอสธมฺโม สนนฺตโน.} \\ \csromangathabat{\textsuperscript{*}ปเร จ น วิชานนฺติ,}{มยเมตฺถ ยมามเส.} \\ \csromangathabat{เย จ ตตฺถ วิชานนฺติ,}{ตโต สมฺมนฺติ เมธคา.} \\ \csromangathabat{\textsuperscript{*}อฏฺฐิจฺฉินฺนา ปาณหรา,}{ควาสฺสธนหาริโน.} \\ \csromangathabat{รฏฺฐํ วิลุมฺปมานานํ,}{เตสมฺปิ โหติ สงฺคติ.} \\ กสฺมา ตุมฺหาก โน สิยา. \\ \textsuperscript{*}สเจ ลเภถ นิปกํ สหายํ, \\ สทฺธิํ จรํ สาธุวิหาริ ธีรํ. \\ อภิภุยฺย สพฺพานิ ปริสฺสยานิ, \\ จเรยฺย เตนตฺตมโน สตีมา. \\ \textsuperscript{+}โน เจ ลเภถ นิปกํ สหายํ, \\ สทฺธิํ จรํ สาธุวิหาริ ธีรํ. \\ ราชาว รฏฺฐํ วิชิตํ ปหาย, \\ เอโก จเร มาตงฺครญฺเญว นาโค. \\ \textsuperscript{+}เอกสฺส จริตํ เสยฺโย, \\ นตฺถิ พาเล สหายตา. \\ เอโก จเร น จ ปาปานิ กยิรา, \\ อปฺโปสฺสุกฺโก มาตงฺครญฺเญว นาโคติ.}
\csromangathameasurewak{กสฺมา ตุมฺหาก โน สิยา. \\ \textsuperscript{*}สเจ ลเภถ นิปกํ สหายํ, \\ สทฺธิํ จรํ สาธุวิหาริ ธีรํ. \\ อภิภุยฺย สพฺพานิ ปริสฺสยานิ, \\ จเรยฺย เตนตฺตมโน สตีมา. \\ \textsuperscript{+}โน เจ ลเภถ นิปกํ สหายํ, \\ สทฺธิํ จรํ สาธุวิหาริ ธีรํ. \\ ราชาว รฏฺฐํ วิชิตํ ปหาย, \\ เอโก จเร มาตงฺครญฺเญว นาโค. \\ \textsuperscript{+}เอกสฺส จริตํ เสยฺโย, \\ นตฺถิ พาเล สหายตา. \\ เอโก จเร น จ ปาปานิ กยิรา, \\ อปฺโปสฺสุกฺโก มาตงฺครญฺเญว นาโคติ.}
\csromangathasetleft{\textsuperscript{*}ปุถุสทฺโท สมชโน, \\ สํฆสฺมิํ ภิชฺชมานสฺมิํ, \\ \textsuperscript{*}ปริมุฏฺฐา ปณฺฑิตาภาสา, \\ ยาวิจฺฉนฺติ มุขายามํ, \\ \textsuperscript{*}อกฺโกจฺฉิ มํ อวธิ มํ, \\ เย จ ตํ อุปนยฺหนฺติ, \\ \textsuperscript{*}อกฺโกจฺฉิ มํ อวธิ มํ, \\ เย จ ตํ นุปนยฺหนฺติ, \\ \textsuperscript{*}น หิ เวเรน เวรานิ, \\ อเวเรน จ สมฺมนฺติ, \\ \textsuperscript{*}ปเร จ น วิชานนฺติ, \\ เย จ ตตฺถ วิชานนฺติ, \\ \textsuperscript{*}อฏฺฐิจฺฉินฺนา ปาณหรา, \\ รฏฺฐํ วิลุมฺปมานานํ,}
\csromangathagroup{\csromangathastanza{\csromanfolio{496}\csromangathabat{\csromansymbolfootnote{*}{ม 3. 192 ปิฏฺเฐปิ.}ปุถุสทฺโท สมชโน,}{น พาโล โกจิ มญฺญถ.} \\ \csromangathabat{สํฆสฺมิํ ภิชฺชมานสฺมิํ,}{นาญฺญํ ภิยฺโย อมญฺญรุํ.}}\csromangathastanza{\csromangathabat{\csromansymbolmark{*}ปริมุฏฺฐา ปณฺฑิตาภาสา,}{วาจา\-โคจร\-ภาณิโน.} \\ \csromangathabat{ยาวิจฺฉนฺติ มุขายามํ,}{เยน นีตา น ตํ วิทู.}}\csromangathastanza{\csromangathabat{\csromansymbolmark{*}อกฺโกจฺฉิ มํ อวธิ มํ,}{อชินิ มํ อหาสิ เม.} \\ \csromangathabat{เย จ ตํ อุปนยฺหนฺติ,}{เวรํ เตสํ น สมฺมติ.}}\csromangathastanza{\csromangathabat{\csromansymbolmark{*}อกฺโกจฺฉิ มํ อวธิ มํ,}{อชินิ มํ อหาสิ เม.} \\ \csromangathabat{เย จ ตํ นุปนยฺหนฺติ,}{เวรํ เตสูปสมฺมติ.}}\csromangathastanza{\csromangathabat{\csromansymbolmark{*}น หิ เวเรน เวรานิ,}{สมฺมนฺตีธ กุทาจนํ.} \\ \csromangathabat{อเวเรน จ สมฺมนฺติ,}{เอสธมฺโม สนนฺตโน.}}\csromangathastanza{\csromangathabat{\csromansymbolmark{*}ปเร จ น วิชานนฺติ,}{มยเมตฺถ ยมามเส.} \\ \csromangathabat{เย จ ตตฺถ วิชานนฺติ,}{ตโต สมฺมนฺติ เมธคา.}}\csromangathastanza{\csromangathabat{\csromansymbolmark{*}อฏฺฐิจฺฉินฺนา ปาณหรา,}{ควาสฺสธนหาริโน.} \\ \csromangathabat{รฏฺฐํ วิลุมฺปมานานํ,}{เตสมฺปิ โหติ สงฺคติ.}}\csromangathastanza{\csromangathawak{กสฺมา ตุมฺหาก โน สิยา. \\ \csromansymbolmark{*}สเจ ลเภถ นิปกํ สหายํ, \\ สทฺธิํ จรํ สาธุวิหาริ ธีรํ. \\ อภิภุยฺย สพฺพานิ ปริสฺสยานิ,}}\csromangathastanza{\csromangathawak{จเรยฺย เตนตฺตมโน สตีมา. \\ \csromansymbolfootnote{+}{ม 3. 193 ปิฏฺเฐปิ.}โน เจ ลเภถ นิปกํ สหายํ, \\ สทฺธิํ จรํ สาธุวิหาริ ธีรํ. \\ ราชาว รฏฺฐํ วิชิตํ ปหาย,}}\csromangathastanza{\csromangathawak{เอโก จเร มาตงฺครญฺเญว นาโค. \\ \csromansymbolmark{+}เอกสฺส จริตํ เสยฺโย, \\ นตฺถิ พาเล สหายตา. \\ เอโก จเร น จ ปาปานิ กยิรา,}}\csromangathastanza{\csromangathawak{อปฺโปสฺสุกฺโก มาตงฺครญฺเญว นาโคติ.}}}

\csromanmatikaanchor{mtk.283}
\csromantocmarkref{section}{273. พาลกโลณกคมนกถา}{mtk.283}
\bookmark[dest={mtk.283},level=1]{273. พาลกโลณกคมนกถา}
\csromanheader{1}{273. \textbf{พาลกโลณกคมนกถา}}
\proseitem{465}{\csromanfolio{497}อถ โข ภควา สํฆมชฺเฌ ฐิตโกว อิมา คาถาโย ภาสิตฺวา เยน พาลก\-โลณก\-คาโม\csromansharedfootnote{5c0f5d90e44bf408}{พาลกโลณการคาโม (สี, สฺยา)} เตนุปสงฺกมิ. เตน โข ปน สมเยน อายสฺมา ภคุ พาลก\-โลณก\-คาเม วิหรติ. อทฺทสา โข อายสฺมา ภคุ ภควนฺตํ ทูรโตว อาคจฺฉนฺตํ, ทิสฺวาน อาสนํ ปญฺญเปสิ, ปาโททกํ ปาทปีฐํ ปาทกถลิกํ อุปนิกฺขิปิ, ปจฺจุคฺคนฺตฺวา ปตฺตจีวรํ ปฏิคฺคเหสิ. นิสีทิ ภควา ปญฺญตฺเต อาสเน, นิสชฺช โข ภควา ปาเท ปกฺขาเลสิ. อายสฺมาปิ โข ภคุ ภควนฺตํ อภิวาเทตฺวา เอกมนฺตํ นิสีทิ. เอกมนฺตํ นิสินฺนํ โข อายสฺมนฺตํ ภคุํ ภควา เอตทโวจ “กจฺจิ ภิกฺขุ ขมนียํ, กจฺจิ ยาปนียํ, กจฺจิ ปิณฺฑเกน น กิลมสี”ติ. ขมนียํ ภควา, ยาปนียํ ภควา, น จาหํ ภนฺเต ปิณฺฑเกน กิลมามีติ. อถ โข ภควา อายสฺมนฺตํ ภคุํ ธมฺมิยา กถาย สนฺทสฺเสตฺวา สมาทเปตฺวา สมุตฺเตเชตฺวา สมฺปหํเสตฺวา อุฏฺฐายาสนา เยน ปาจินวํสทาโย เตนุปสงฺกมิ.}

\csromanmatikaanchor{mtk.284}
\csromantocmarkref{section}{274. ปาจินวํสทายคมนกถา}{mtk.284}
\bookmark[dest={mtk.284},level=1]{274. ปาจินวํสทายคมนกถา}
\csromanheader{1}{274. \textbf{ปาจินวํสทาย}\-\textbf{คมน}\-\textbf{กถา}}
\proseitem{466}{\csromansymbolfootnote{*}{ม 1. 266-8 ปิฏฺเฐสุปิ ปสฺสิตพฺพํ.}เตน โข ปน สมเยน อายสฺมา จ อนุรุทฺโธ อายสฺมา จ นนฺทิโย อายสฺมา จ กิมิโล\csromansharedfootnote{d6640201af0a7bb8}{กิมฺพิโล (สี, สฺยา)} ปาจินวํสทาเย วิหรนฺติ. อทฺทสา โข ทายปาโล ภควนฺตํ ทูรโตว อาคจฺฉนฺตํ, ทิสฺวาน ภควนฺตํ เอตทโวจ “มา สมณ เอตํ ทายํ ปาวิสิ, สนฺเตตฺถ ตโย กุลปุตฺตา อตฺตกามรูปา วิหรนฺติ, มา เตสํ อผาสุมกาสี”ติ. อสฺโสสิ โข อายสฺมา อนุรุทฺโธ ทายปาลสฺส ภควตา สทฺธิํ มนฺตยมานสฺส, สุตฺวาน ทายปาลํ เอตทโวจ “มาวุโส ทายปาล ภควนฺตํ วาเรสิ, สตฺถา โน ภควา อนุปฺปตฺโต”ติ. อถ โข อายสฺมา อนุรุทฺโธ เยนายสฺมา จ นนฺทิโย อายสฺมา จ กิมิโล เตนุปสงฺกมิ, อุปสงฺกมิตฺวา อายสฺมนฺตญฺจ นนฺทิยํ อายสฺมนฺตญฺจ กิมิลํ เอตทโวจ “อภิกฺกมถา\-ยสฺมนฺโต อภิกฺกมถา\-ยสฺมนฺโต, สตฺถา โน ภควา อนุปฺปตฺโต”ติ. อถ โข อายสฺมา จ อนุรุทฺโธ อายสฺมา จ นนฺทิโย อายสฺมา จ กิมิโล ภควนฺตํ ปจฺจุคฺคนฺตฺวา เอโก ภควโต ปตฺตจีวรํ ปฏิคฺคเหสิ, เอโก อาสนํ ปญฺญเปสิ, เอโก ปาโททกํ ปาทปีฐํ ปาทกถลิกํ อุปนิกฺขิปิ. นิสีทิ ภควา ปญฺญตฺเต \csromanfolio{498}อาสเน, นิสชฺช โข ภควา ปาเท ปกฺขาเลสิ. เตปิ โข อายสฺมนฺโต ภควนฺตํ อภิวาเทตฺวา เอกมนฺตํ นิสีทิํสุ. เอกมนฺตํ นิสินฺนํ โข อายสฺมนฺตํ อนุรุทฺธํ ภควา เอตทโวจ “กจฺจิ โว อนุรุทฺธา ขมนียํ, กจฺจิ ยาปนียํ, กจฺจิ ปิณฺฑเกน น กิลมถา”ติ. ขมนียํ ภควา, ยาปนียํ ภควา, น จ มยํ ภนฺเต ปิณฺฑเกน กิลมามาติ.}

\prose{กจฺจิ ปน โว อนุรุทฺธา สมคฺคา สมฺโมทมานา อวิวทมานา ขีโรทกีภูตา อญฺญมญฺญํ ปิยจกฺขูหิ สมฺปสฺสนฺตา วิหรถาติ. ตคฺฆ มยํ ภนฺเต สมคฺคา สมฺโมทมานา อวิวทมานา ขีโรทกีภูตา อญฺญมญฺญํ ปิยจกฺขูหิ สมฺปสฺสนฺตา วิหรามาติ. ยถา กถํ ปน ตุเมฺห อนุรุทฺธา สมคฺคา สมฺโมทมานา อวิวทมานา ขีโรทกีภูตา อญฺญมญฺญํ ปิยจกฺขูหิ สมฺปสฺสนฺตา วิหรถาติ. อิธ มยฺหํ ภนฺเต เอวํ โหติ “ลาภา วต เม, สุลทฺธํ วต เม, โยหํ เอวรูเปหิ สพฺรหฺมจารีหิ สทฺธิํ วิหรามี”ติ, ตสฺส มยฺหํ ภนฺเต อิเมสุ อายสฺมนฺเตสุ เมตฺตํ กายกมฺมํ ปจฺจุปฏฺฐิตํ อาวิ เจว รโห จ, เมตฺตํ วจีกมฺมํ, เมตฺตํ มโนกมฺมํ ปจฺจุปฏฺฐิตํ อาวิ เจว รโห จ, ตสฺส มยฺหํ ภนฺเต เอวํ โหติ “ยนฺนูนาหํ สกํ จิตฺตํ นิกฺขิปิตฺวา อิเมสํเยว อายสฺมนฺตานํ จิตฺตสฺส วเสน วตฺเตยฺยนฺ”ติ, โส โข อหํ ภนฺเต สกํ จิตฺตํ นิกฺขิปิตฺวา อิเมสํเยว อายสฺมนฺตานํ จิตฺตสฺส วเสน วตฺตามิ, นานา หิ โข โน ภนฺเต กายา, เอกญฺจ ปน มญฺเญ จิตฺตนฺติ.}

\prose{อายสฺมาปิ โข นนฺทิโย ฯเปฯ. อายสฺมาปิ โข กิมิโล ภควนฺตํ เอตทโวจ “มยฺหมฺปิ โข ภนฺเต เอวํ โหติ ‘ลาภา วต เม, สุลทฺธํ วต เม, โยหํ เอวรูเปหิ สพฺรหฺมจารีหิ สทฺธิํ วิหรามี’ติ, ตสฺส มยฺหํ ภนฺเต อิเมสุ อายสฺมนฺเตสุ เมตฺตํ กายกมฺมํ ปจฺจุปฏฺฐิตํ อาวิ เจว รโห จ, เมตฺตํ วจีกมฺมํ, เมตฺตํ มโนกมฺมํ ปจฺจุปฏฺฐิตํ อาวิ เจว รโห จ, ตสฺส มยฺหํ ภนฺเต เอวํ โหติ ‘ยนฺนูนาหํ สกํ จิตฺตํ นิกฺขิปิตฺวา อิเมสํเยว อายสฺมนฺตานํ จิตฺตสฺส วเสน วตฺเตยฺยนฺ’ติ, โส โข อหํ ภนฺเต สกํ จิตฺตํ นิกฺขิปิตฺวา อิเมสํเยว อายสฺมนฺตานํ จิตฺตสฺส วเสน วตฺตามิ, นานา หิ โข โน ภนฺเต กายา, เอกญฺจ ปน มญฺเญ จิตฺตนฺติ. เอวํ โข มยํ ภนฺเต สมคฺคา สมฺโมทมานา อวิวทมานา ขีโรทกีภูตา อญฺญมญฺญํ ปิยจกฺขูหิ สมฺปสฺสนฺตา วิหรามา”ติ.}

\prose{\csromanfolio{499}กจฺจิ ปน โว อนุรุทฺธา อปฺปมตฺตา อาตาปิโน ปหิตตฺตา วิหรถาติ. ตคฺฆ มยํ ภนฺเต อปฺปมตฺตา อาตาปิโน ปหิตตฺตา วิหรามาติ. ยถา กถํ ปน ตุเมฺห อนุรุทฺธา อปฺปมตฺตา อาตาปิโน ปหิตตฺตา วิหรถาติ. อิธ ภนฺเต อมฺหากํ โย ปฐมํ คามโต ปิณฺฑาย ปฏิกฺกมติ, โส อาสนํ ปญฺญเปติ, ปาโททกํ ปาทปีฐํ ปาทกถลิกํ อุปนิกฺขิปติ, อวกฺการปาติํ โธวิตฺวา อุปฏฺฐาเปติ, ปานียํ ปริโภชนียํ อุปฏฺฐาเปติ. โย ปจฺฉา คามโต ปิณฺฑาย ปฏิกฺกมติ, สเจ โหติ ภุตฺตาวเสโส, สเจ อากงฺขติ ภุญฺชติ, โน เจ อากงฺขติ อปฺปหริเต วา ฉฑฺเฑติ, อปฺปาณเก วา อุทเก โอปิลาเปติ, โส อาสนํ อุทฺธรติ, ปาโททกํ ปาทปีฐํ ปาทกถลิกํ ปฏิสาเมติ, อวกฺการปาติํ โธวิตฺวา ปฏิสาเมติ, ปานียํ ปริโภชนียํ ปฏิสาเมติ, ภตฺตคฺคํ สมฺมชฺชติ. โย ปสฺสติ ปานียฆฏํ วา ปริโภชนียฆฏํ วา วจฺจฆฏํ วา ริตฺตํ ตุจฺฉํ, โส อุปฏฺฐาเปติ. สจสฺส โหติ อวิสยฺหํ, หตฺถวิกาเรน ทุติยํ อามนฺเตตฺวา หตฺถ\-วิลงฺฆเกน อุปฏฺฐาเปม, น เตฺวว มยํ ภนฺเต ตปฺปจฺจยา วาจํ ภินฺทาม. ปญฺจาหิกํ โข ปน มยํ ภนฺเต สพฺพรตฺติํ ธมฺมิยา กถาย สนฺนิสีทาม. เอวํ โข มยํ ภนฺเต อปฺปมตฺตา อาตาปิโน ปหิตตฺตา วิหรามาติ.}

\csromanmatikaanchor{mtk.285}
\csromantocmarkref{section}{275. ปาลิเลยฺยกคมนกถา}{mtk.285}
\bookmark[dest={mtk.285},level=1]{275. ปาลิเลยฺยกคมนกถา}
\csromanheader{1}{275. \textbf{ปาลิเลยฺยก}\-\textbf{คมน}\-\textbf{กถา}}
\proseitem{467}{\csromansymbolfootnote{*}{ขุ 1. 125-6 ปิฏฺเฐสุ โถกํ วิสทิสํ.}อถ โข ภควา อายสฺมนฺตญฺจ อนุรุทฺธํ อายสฺมนฺตญฺจ นนฺทิยํ อายสฺมนฺตญฺจ กิมิลํ ธมฺมิยา กถาย สนฺทสฺเสตฺวา สมาทเปตฺวา สมุตฺเตเชตฺวา สมฺปหํเสตฺวา อุฏฺฐายาสนา เยน ปาลิเลยฺยกํ\csromansharedfootnote{7518e63a06df090b}{ปาริเลยฺยกํ (สี, สฺยา)} เตน จาริกํ ปกฺกามิ. อนุปุพฺเพน จาริกํ จรมาโน เยน ปาลิเลยฺยกํ ตทวสริ, ตตฺร สุทํ ภควา ปาลิเลยฺยเก วิหรติ รกฺขิต\-วนสณฺเฑ ภทฺทสาลมูเล. อถ โข ภควโต รโหคตสฺส ปฏิสลฺลีนสฺส เอวํ เจตโส ปริวิตกฺโก อุทปาทิ “อหํ โข ปุพฺเพ อากิณฺโณ น ผาสุ วิหาสิํ เตหิ โกสมฺพเกหิ\csromansharedfootnote{234946681376f40c}{โกสมฺพิเกหิ (สฺยา)} ภิกฺขูหิ ภณฺฑน\-การเกหิ กลหการเกหิ วิวาทการเกหิ ภสฺสการเกหิ สํเฆ อธิกรณ\-การเกหิ, โสมฺหิ เอตรหิ เอโก อทุติโย สุขํ ผาสุ วิหรามิ อญฺญเตฺรว เตหิ โกสมฺพเกหิ ภิกฺขูหิ ภณฺฑน\-การเกหิ กลหการเกหิ วิวาทการเกหิ ภสฺสการเกหิ สํเฆ อธิกรณการเกหี”ติ.}

\prose{\csromanfolio{500}อญฺญตโรปิ โข หตฺถินาโค อากิณฺโณ วิหรติ หตฺถีหิ หตฺถินีหิ หตฺถิ\-กฬเภหิ หตฺถิจฺฉาเปหิ, ฉินฺนคฺคานิ เจว ติณานิ ขาทติ, โอภคฺโคภคฺค\-ญฺจสฺส สาขาภงฺคํ ขาทนฺติ, อาวิลานิ จ ปานียานิ ปิวติ, โอคาหา จสฺส โอติณฺณสฺส\csromansharedfootnote{e6d2c5d8234cecbf}{โอคาหญฺจสฺส โอติณฺณสฺส (สฺยา), โอคาหา จสฺส อุตฺติณฺณสฺส (สี)} หตฺถินิโย กายํ อุปนิฆํสนฺติโย คจฺฉนฺติ.}

\prose{อถ โข ตสฺส หตฺถินาคสฺส เอตทโหสิ “อหํ โข อากิณฺโณ วิหรามิ หตฺถีหิ หตฺถินีหิ หตฺถิ\-กฬเภหิ หตฺถิจฺฉาเปหิ, ฉินฺนคฺคานิ เจว ติณานิ ขาทามิ, โอภคฺโคภคฺคญฺจ เม สาขาภงฺคํ ขาทนฺติ, อาวิลานิ จ ปานียานิ ปิวามิ, โอคาหา จ เม โอติณฺณสฺส หตฺถินิโย กายํ อุปนิฆํสนฺติโย คจฺฉนฺติ, ยนฺนูนาหํ เอโกว คณสฺมา วูปกฏฺโฐ วิหเรยฺยนฺ”ติ.}

\prose{อถ โข โส หตฺถินาโค ยูถา อปกฺกมฺม เยน ปาลิเลยฺยกํ รกฺขิต\-วนสณฺโฑ ภทฺท\-สาลมูลํ เยน ภควา เตนุปสงฺกมิ, อุปสงฺกมิตฺวา โสณฺฑาย ภควโต ปานียํ ปริโภชนียํ อุปฏฺฐาเปติ, อปฺปหริตญฺจ กโรติ.}

\prose{อถ โข ตสฺส หตฺถินาคสฺส เอตทโหสิ “อหํ โข ปุพฺเพ อากิณฺโณ น ผาสุ วิหาสิํ หตฺถีหิ หตฺถินีหิ หตฺถิ\-กฬเภหิ หตฺถิจฺฉาเปหิ, ฉินฺนคฺคานิ เจว ติณานิ ขาทิํ, โอภคฺโคภคฺคญฺจ เม สาขาภงฺคํ ขาทิํสุ, อาวิลานิ จ ปานียานิ อปายิํ, โอคาหา จ เม โอติณฺณสฺส หตฺถินิโย กายํ อุปนิฆํสนฺติโย อคมํสุ, โสมฺหิ เอตรหิ เอโก อทุติโย สุขํ ผาสุ วิหรามิ อญฺญเตฺรว หตฺถีหิ หตฺถินีหิ หตฺถิ\-กฬเภหิ หตฺถิจฺฉาเปหี”ติ.}

\proseitem{467}{อถ โข ภควา อตฺตโน จ ปวิเวกํ วิทิตฺวา ตสฺส จ หตฺถินาคสฺส เจตสา เจโต\-ปริวิตกฺกม\-ญฺญาย ตายํ เวลายํ อิมํ อุทานํ อุทาเนสิ\csromandash{}}

\setlength{\csromangathapagewidth}{0pt}
\setlength{\csromangathawakwidth}{0pt}
\csromangathameasuregroup{\textsuperscript{*}“เอตํ\textsuperscript{1} นาคสฺส นาเคน, \\ สเมติ จิตฺตํ จิตฺเตน,}{\csromangathabat{\textsuperscript{*}“เอตํ\textsuperscript{1} นาคสฺส นาเคน,}{อีสาทนฺตสฺส หตฺถิโน.} \\ \csromangathabat{สเมติ จิตฺตํ จิตฺเตน,}{ยเทโก รมตี วเน”ติ.}}
\csromangathasetleft{\textsuperscript{*}“เอตํ\textsuperscript{1} นาคสฺส นาเคน, \\ สเมติ จิตฺตํ จิตฺเตน,}
\csromangathagroup{\csromangathastanza{\csromangathabat{\csromansymbolfootnote{*}{ขุ 1. 127 ปิฏฺเฐปิ.}“เอตํ\csromansharedfootnote{865fdde227c93b9a}{เอวํ (ก)} นาคสฺส นาเคน,}{อีสาทนฺตสฺส หตฺถิโน.} \\ \csromangathabat{สเมติ จิตฺตํ จิตฺเตน,}{ยเทโก รมตี วเน”ติ.}}}

\prose{อถ โข ภควา ปาลิเลยฺยเก ยถาภิรนฺตํ วิหริตฺวา เยน สาวตฺถิ เตน จาริกํ ปกฺกามิ. อนุปุพฺเพน จาริกํ จรมาโน เยน สาวตฺถิ ตทวสริ, ตตฺร สุทํ ภควา สาวตฺถิยํ วิหรติ เชตวเน อนาถ\-ปิณฺฑิกสฺส อาราเม. อถ โข โกสมฺพโก อุปาสกา “อิเม โข อยฺยา โกสมฺพกา ภิกฺขู พหุโน อมฺหากํ อนตฺถสฺส \csromanfolio{501}การกา, อิเมหิ อุพฺพาโฬฺห ภควา ปกฺกนฺโต, หนฺท มยํ อยฺเย โกสมฺพเก ภิกฺขู เนว อภิวาเทยฺยาม, น ปจฺจุฏฺเฐยฺยาม, น อญฺชลิกมฺมํ สามีจิกมฺมํ กเรยฺยาม, น สกฺกเรยฺยาม, น ครุํ กเรยฺยาม, น มาเนยฺยาม, น ภเชยฺยาม, น ปูเชยฺยาม, อุปคตานมฺปิ ปิณฺฑกํ น ทชฺเชยฺยาม, เอวํ อิเม อเมฺหหิ อสกฺกริยมานา อครุกริยมานา อมานิยมานา อภชิยมานา อปูชิยมานา อสกฺการปกตา ปกฺกมิสฺสนฺติ วา วิพฺภมิสฺสนฺติ วา ภควนฺตํ วา ปสาเทสฺสนฺตี”ติ. อถ โข โกสมฺพกา อุปาสกา โกสมฺพเก ภิกฺขู เนว อภิวาเทสุํ, น ปจฺจุฏฺเฐสุํ, น อญฺชลิกมฺมํ สามีจิกมฺมํ อกํสุ, น สกฺกริํสุ, น ครุํ กริํสุ, น มาเนสุํ, น ภเชสุํ, น ปูเชสุํ, อุปคตานมฺปิ ปิณฺฑกํ น อทํสุ. อถ โข โกสมฺพกา ภิกฺขู โกสมฺพเกหิ อุปาสเกหิ อสกฺกริยมานา อครุกริยมานา อมานิยมานา อภชิยมานา อปูชิยมานา อสกฺการปกตา เอวมาหํสุ “หนฺท มยํ อาวุโส สาวตฺถิํ คนฺตฺวา ภควโต สนฺติเก อิมํ อธิกรณํ วูปสเมยฺยามา”ติ.}

\csromanmatikaanchor{mtk.286}
\csromantocmarkref{section}{276. อฏฺฐารสวตฺถุกถา}{mtk.286}
\bookmark[dest={mtk.286},level=1]{276. อฏฺฐารสวตฺถุกถา}
\csromanheader{1}{276. \textbf{อฏฺฐารส}\-\textbf{วตฺถุ}\-\textbf{กถา}}
\proseitem{468}{อถ โข โกสมฺพกา ภิกฺขู เสนาสนํ สํสาเมตฺวา ปตฺต\-จีวรมา\-ทาย เยน สาวตฺถิ เตนุ\-ปสงฺกมิํสุ. อสฺโสสิ โข อายสฺมา สาริปุตฺโต “เต กิร โกสมฺพกา ภิกฺขู ภณฺฑนการกา กลหการกา วิวาทการกา ภสฺสการกา สํเฆ อธิกรณ\-การกา สาวตฺถิํ อาคจฺฉนฺตี”ติ. อถ โข อายสฺมา สาริปุตฺโต เยน ภควา เตนุปสงฺกมิ, อุปสงฺกมิตฺวา ภควนฺตํ อภิวาเทตฺวา เอกมนฺตํ นิสีทิ, เอกมนฺตํ นิสินฺโน โข อายสฺมา สาริปุตฺโต ภควนฺตํ เอตทโวจ “เต กิร ภนฺเต โกสมฺพกา ภิกฺขู ภณฺฑนการกา กลหการกา วิวาทการกา ภสฺสการกา สํเฆ อธิกรณ\-การกา สาวตฺถิํ อาคจฺฉนฺติ, กถาหํ ภนฺเต เตสุ ภิกฺขูสุ ปฏิปชฺชามี”ติ. เตน หิ ตฺวํ สาริปุตฺต ยถา ธมฺโม, ตถา ติฏฺฐาหีติ. กถาหํ ภนฺเต ชาเนยฺยํ ธมฺมํ วา อธมฺมํ วาติ.}

\prose{อฏฺฐารสหิ โข สาริปุตฺต วตฺถูหิ อธมฺมวาที ชานิตพฺโพ. อิธ สาริปุตฺต ภิกฺขุ อธมฺมํ ธมฺโมติ ทีเปติ, ธมฺมํ อธมฺโมติ ทีเปติ, อวินยํ วินโยติ ทีเปติ, วินยํ อวินโยติ ทีเปติ, อภาสิตํ อลปิตํ \csromanfolio{502}ตถาคเตน ภาสิตํ ลปิตํ ตถาคเตนาติ ทีเปติ, ภาสิตํ ลปิตํ ตถาคเตน อภาสิตํ อลปิตํ ตถาคเตนาติ ทีเปติ, อนาจิณฺณํ ตถาคเตน อาจิณฺณํ ตถาคเตนาติ ทีเปติ, อาจิณฺณํ ตถาคเตน อนาจิณฺณํ ตถาคเตนาติ ทีเปติ, อปญฺญตฺตํ ตถาคเตน ปญฺญตฺตํ ตถาคเตนาติ ทีเปติ, ปญฺญตฺตํ ตถาคเตน อปญฺญตฺตํ ตถาคเตนาติ ทีเปติ, อนาปตฺติํ อาปตฺตีติ ทีเปติ, อาปตฺติํ อนาปตฺตีติ ทีเปติ, ลหุกํ อาปตฺติํ ครุกา อาปตฺตีติ ทีเปติ, ครุกํ อาปตฺติํ ลหุกา อาปตฺตีติ ทีเปติ, สาวเสสํ อาปตฺติํ อนวเสสา อาปตฺตีติ ทีเปติ, อนวเสสํ อาปตฺติํ สาวเสสา อาปตฺตีติ ทีเปติ, ทุฏฺฐลฺลํ อาปตฺติํ อทุฏฺฐุลฺลา อาปตฺตีติ ทีเปติ, อทุฏฺฐุลฺลํ อาปตฺติํ ทุฏฺฐุลฺลา อาปตฺตีติ ทีเปติ. อิเมหิ โข สาริปุตฺต อฏฺฐารสหิ วตฺถูหิ อธมฺมวาที ชานิตพฺโพ.}

\prose{อฏฺฐารสหิ จ โข สาริปุตฺต วตฺถูหิ ธมฺมวาที ชานิตพฺโพ. อิธ สาริปุตฺต ภิกฺขุ อธมฺมํ อธมฺโมติ ทีเปติ, ธมฺมํ ธมฺโมติ ทีเปติ, อวินยํ อวินโยติ ทีเปติ, วินยํ วินโยติ ทีเปติ, อภาสิตํ อลปิตํ ตถาคเตน อภาสิตํ อลปิตํ ตถาคเตนาติ ทีเปติ, ภาสิตํ ลปิตํ ตถาคเตน ภาสิตํ ลปิตํ ตถาคเตนาติ ทีเปติ, อนาจิณฺณํ ตถาคเตน อนาจิณฺณํ ตถาคเตนาติ ทีเปติ, อาจิณฺณํ ตถาคเตน อาจิณฺณํ ตถาคเตนาติ ทีเปติ, อปญฺญตฺตํ ตถาคเตน อปญฺญตฺตํ ตถาคเตนาติ ทีเปติ, ปญฺญตฺตํ ตถาคเตน ปญฺญตฺตํ ตถาคเตนาติ ทีเปติ, อนาปตฺติํ อนาปตฺตีติ ทีเปติ, อาปตฺติํ อาปตฺตีติ ทีเปติ, ลหุกํ อาปตฺติํ ลหุกา อาปตฺตีติ ทีเปติ, ครุกํ อาปตฺติํ ครุกา อาปตฺตีติ ทีเปติ, สาวเสสํ อาปตฺติํ สาวเสสา อาปตฺตีติ ทีเปติ, อนวเสสํ อาปตฺติํ อนวเสสา อาปตฺตีติ ทีเปติ, ทุฏฺฐุลฺลํ อาปตฺติํ ทุฏฺฐุลฺลา อาปตฺตีติ ทีเปติ, อทุฏฺฐุลฺลํ อาปตฺติํ อทุฏฺฐุลฺลา อาปตฺตีติ ทีเปติ. อิเมหิ โข สาริปุตฺต อฏฺฐารสหิ วตฺถูหิ ธมฺมวาที ชานิตพฺโพติ.}

\proseitem{469}{อสฺโสสิ โข อายสฺมา มหาโมคฺคลฺลาโน ฯเปฯ. อสฺโสสิ โข อายสฺมา มหากสฺสโป. อสฺโสสิ โข อายสฺมา มหากจฺจาโน. อสฺโสสิ โข อายสฺมา มหาโกฏฺฐิโก. อสฺโสสิ โข อายสฺมา มหากปฺปิโน. อสฺโสสิ โข อายสฺมา มหาจุนฺโท. อสฺโสสิ โข อายสฺมา อนุรุทฺโธ. อสฺโสสิ โข อายสฺมา \csromanfolio{503}เรวโต. อสฺโสสิ โข อายสฺมา อุปาลิ. อสฺโสสิ โข อายสฺมา อานนฺโท. อสฺโสสิ โข อายสฺมา ราหุโล “เต กิร โกสมฺพกา ภิกฺขู ภณฺฑนการกา กลหการกา วิวาทการกา ภสฺสการกา สํเฆ อธิกรณ\-การกา สาวตฺถิํ อาคจฺฉนฺตี”ติ. อถ โข อายสฺมา ราหุโล เยน ภควา เตนุปสงฺกมิ, อุปสงฺกมิตฺวา ภควนฺตํ อภิวาเทตฺวา เอกมนฺตํ นิสีทิ, เอกมนฺตํ นิสินฺโน โข อายสฺมา ราหุโล ภควนฺตํ เอตทโวจ “เต กิร ภนฺเต โกสมฺพกา ภิกฺขู ภณฺฑนการกา กลหการกา วิวาทการกา ภสฺสการกา สํเฆ อธิกรณ\-การกา สาวตฺถิํ อาคจฺฉนฺติ, กถาหํ ภนฺเต เตสุ ภิกฺขูสุ ปฏิปชฺชามี”ติ. เตน หิ ตฺวํ ราหุล ยถา ธมฺโม, ตถา ติฏฺฐาหีติ. กถาหํ ภนฺเต ชาเนยฺยํ ธมฺมํ วา อธมฺมํ วาติ.}

\prose{อฏฺฐารสหิ โข ราหุล วตฺถูหิ อธมฺมวาที ชานิตพฺโพ. อิธ ราหุล ภิกฺขุ อธมฺมํ ธมฺโมติ ทีเปติ, ธมฺมํ อธมฺโมติ ทีเปติ, อวินยํ วินโยติ ทีเปติ, วินยํ อวินโยติ ทีเปติ, อภาสิตํ อลปิตํ ตถาคเตน ภาสิตํ ลปิตํ ตถาคเตนาติ ทีเปติ, ภาสิตํ ลปิตํ ตถาคเตน อภาสิตํ อลปิตํ ตถาคเตนาติ ทีเปติ, อนาจิณฺณํ ตถาคเตน อาจิณฺณํ ตถาคเตนาติ ทีเปติ, อาจิณฺณํ ตถาคเตน อนาจิณฺณํ ตถาคเตนาติ ทีเปติ, อปญฺญตฺตํ ตถาคเตน ปญฺญตฺตํ ตถาคเตนาติ ทีเปติ, ปญฺญตฺตํ ตถาคเตน อปญฺญตฺตํ ตถาคเตนาติ ทีเปติ, อนาปตฺติํ อาปตฺตีติ ทีเปติ, อาปตฺติํ อนาปตฺตีติ ทีเปติ, ลหุกํ อาปตฺติํ ครุกา อาปตฺตีติ ทีเปติ, ครุกํ อาปตฺติํ ลหุกา อาปตฺตีติ ทีเปติ, สาวเสสํ อาปตฺติํ อนวเสสา อาปตฺตีติ ทีเปติ, อนวเสสํ อาปตฺติํ สาวเสสา อาปตฺตีติ ทีเปติ, ทุฏฺฐุลฺลํ อาปตฺติํ อทุฏฺฐุลฺลา อาปตฺตีติ ทีเปติ, อทุฏฺฐุลฺลํ อาปตฺติํ ทุฏฺฐุลฺลา อาปตฺตีติ ทีเปติ. อิเมหิ โข ราหุล อฏฺฐารสหิ วตฺถูหิ อธมฺมวาที ชานิตพฺโพ.}

\prose{อฏฺฐารสหิ จ โข ราหุล วตฺถูหิ ธมฺมวาที ชานิตพฺโพ. อิธ ราหุล ภิกฺขุ อธมฺมํ อธมฺโมติ ทีเปติ, ธมฺมํ ธมฺโมติ ทีเปติ, อวินยํ อวินโยติ ทีเปติ, วินยํ วินโยติ ทีเปติ, อภาสิตํ อลปิตํ ตถาคเตน อภาสิตํ อลปิตํ ตถาคเตนาติ ทีเปติ, ภาสิตํ ลปิตํ ตถาคเตน ภาสิตํ ลปิตํ ตถาคเตนาติ ทีเปติ, \csromanfolio{504}อนาจิณฺณํ ตถาคเตน อนาจิณฺณํ ตถาคเตนาติ ทีเปติ, อาจิณฺณํ ตถาคเตน อาจิณฺณํ ตถาคเตนาติ ทิเปติ, อปญฺญตฺตํ ตถาคเตน อปญฺญตฺตํ ตถาคเตนาติ ทีเปติ, ปญฺญตฺตํ ตถาคเตน ปญฺญตฺตํ ตถาคเตนาติ ทีเปติ, อนาปตฺติํ อนาปตฺตีติ ทีเปติ, อาปตฺติํ อาปตฺตีติ ทีเปติ, ลหุกํ อาปตฺติํ ลหุกา อาปตฺตีติ ทีเปติ, ครุกํ อาปตฺติํ ครุกา อาปตฺตีติ ทีเปติ, สาวเสสํ อาปตฺติํ สาวเสสา อาปตฺตีติ ทีเปติ, อนวเสสํ อาปตฺติํ อนวเสสา อาปตฺตีติ ทีเปติ, ทุฏฺฐุลฺลํ อาปตฺติํ ทุฏฺฐุลฺลา อาปตฺตีติ ทีเปติ, อทุฏฺฐุลฺลํ อาปตฺติํ อทุฏฺฐุลฺลา อาปตฺตีติ ทีเปติ. อิเมหิ โข ราหุล อฏฺฐารสหิ วตฺถูหิ ธมฺมวาที ชานิตพฺโพติ.}

\proseitem{470}{อสฺโสสิ โข มหาปชาปติ\csromansharedfootnote{3d7073d5ee5a87cc}{มหาปชาปตี (สี, สฺยา)} โคตมี “เต กิร โกสมฺพกา ภิกฺขู ภณฺฑนการกา กลหการกา วิวาทการกา ภสฺสการกา สํเฆ อธิกรณ\-การกา สาวตฺถิํ อาคจฺฉนฺตี”ติ. อถ โข มหาปชาปติ โคตมี เยน ภควา เตนุปสงฺกมิ, อุปสงฺกมิตฺวา ภควนฺตํ อภิวาเทตฺวา เอกมนฺตํ อฏฺฐาสิ, เอกมนฺตํ ฐิตา โข มหาปชาปติ โคตมี ภควนฺตํ เอตทโวจ “เต กิร ภนฺเต โกสมฺพกา ภิกฺขู ภณฺฑนการกา กลหการกา วิวาทการกา ภสฺสการกา สํเฆ อธิกรณ\-การกา สาวตฺถิํ อาคจฺฉนฺติ, กถาหํ ภนฺเต เตสุ ภิกฺขูสุ ปฏิปชฺชามี”ติ. เตน หิ ตฺวํ โคตมิ อุภยตฺถ ธมฺมํ สุณ, อุภยตฺถ ธมฺมํ สุตฺวา เย ตตฺถ ภิกฺขู ธมฺมวาทิโน, เตสํ ทิฏฺฐิญฺจ ขนฺติญฺจ รุจิญฺจ อาทายญฺจ โรเจหิ, ยญฺจ กิญฺจิ ภิกฺขุนิ\-สํเฆน ภิกฺขุ\-สํฆโต ปจฺจาสีสิตพฺพํ, สพฺพํ ตํ ธมฺมวาทิโตว ปจฺจาสีสิตพฺพนฺติ.}

\proseitem{471}{อสฺโสสิ โข อนาถปิณฺฑิโก คหปติ “เต กิร โกสมฺพกา ภิกฺขู ภณฺฑนการกา กลหการกา วิวาทการกา ภสฺสการกา สํเฆ อธิกรณ\-การกา สาวตฺถิํ อาคจฺฉนฺตี”ติ. อถ โข อนาถปิณฺฑิโก คหปติ เยน ภควา เตนุปสงฺกมิ, อุปสงฺกมิตฺวา ภควนฺตํ อภิวาเทตฺวา เอกมนฺตํ นิสีทิ, เอกมนฺตํ นิสินฺโน โข อนาถปิณฺฑิโก คหปติ ภควนฺตํ เอตทโวจ “เต กิร ภนฺเต โกสมฺพกา ภิกฺขู ภณฺฑนการกา กลหการกา วิวาทการกา \csromanfolio{505}ภสฺสการกา สํเฆ อธิกรณ\-การกา สาวตฺถิํ อาคจฺฉนฺติ, กถาหํ ภนฺเต เตสุ ภิกฺขูสุ ปฏิปชฺชามี”ติ. เตน หิ ตฺวํ คหปติ อุภยตฺถ ทานํ เทหิ, อุภยตฺถ ทานํ ทตฺวา อุภยตฺถ ธมฺมํ สุณ, อุภยตฺถ ธมฺมํ สุตฺวา เย ตตฺถ ภิกฺขู ธมฺมวาทิโน, เตสํ ทิฏฺฐิญฺจ ขนฺติญฺจ รุจิญฺจ อาทายญฺจ โรเจหีติ.}

\proseitem{472}{อสฺโสสิ โข วิสาขา มิคารมาตา “เต กิร โกสมฺพกา ภิกฺขู ภณฺฑนการกา กลหการกา วิวาทการกา ภสฺสการกา สํเฆ อธิกรณ\-การกา สาวตฺถิํ อาคจฺฉนฺตี”ติ. อถ โข วิสาขา มิคารมาตา เยน ภควา เตนุปสงฺกมิ, อุปสงฺกมิตฺวา ภควนฺตํ อภิวาเทตฺวา เอกมนฺตํ นิสีทิ, เอกมนฺตํ นิสินฺนา โข วิสาขา มิคารมาตา ภควนฺตํ เอตทโวจ “เต กิร ภนฺเต โกสมฺพกา ภิกฺขู ภณฺฑนการกา กลหการกา วิวาทการกา ภสฺสการกา สํเฆ อธิกรณ\-การกา สาวตฺถิํ อาคจฺฉนฺติ, กถาหํ ภนฺเต เตสุ ภิกฺขูสุ ปฏิปชฺชามี”ติ. เตน หิ ตฺวํ วิสาเข อุภยตฺถ ทานํ เทหิ, อุภยตฺถ ทานํ ทตฺวา อุภยตฺถ ธมฺมํ สุณ, อุภยตฺถ ธมฺมํ สุตฺวา เย ตตฺถ ภิกฺขู ธมฺมวาทิโน, เตสํ ทิฏฺฐิญฺจ ขนฺติญฺจ รุจิญฺจ อาทายญฺจ โรเจหีติ.}

\proseitem{473}{อถ โข โกสมฺพกา ภิกฺขู อนุปุพฺเพน เยน สาวตฺถิ ตทวสรุํ. อถ โข อายสฺมา สาริปุตฺโต เยน ภควา เตนุปสงฺกมิ, อุปสงฺกมิตฺวา ภควนฺตํ อภิวาเทตฺวา เอกมนฺตํ นิสีทิ, เอกมนฺตํ นิสินฺโน โข อายสฺมา สาริปุตฺโต ภควนฺตํ เอตทโวจ “เต กิร ภนฺเต โกสมฺพกา ภิกฺขู ภณฺฑนการกา กลหการกา วิวาทการกา ภสฺสการกา สํเฆ อธิกรณ\-การกา สาวตฺถิํ อนุปฺปตฺตา, กถํ นุ โข ภนฺเต เตสุ ภิกฺขูสุ เสนาสเน\csromansharedfootnote{824d20f80611c5e8}{เสนาสเนสุ (ก), เสนาสนํ (สฺยา)} ปฏิปชฺชิตพฺพนฺ”ติ. เตน หิ สาริปุตฺต วิวิตฺตํ เสนาสนํ ทาตพฺพนฺติ. สเจ ปน ภนฺเต วิวิตฺตํ น โหติ, กถํ ปฏิปชฺชิตพฺพนฺติ. เตน หิ สาริปุตฺต วิวิตฺตํ กตฺวาปิ ทาตพฺพํ, น เตฺววาหํ สาริปุตฺต เกนจิ ปริยาเยน วุฑฺฒตรสฺส ภิกฺขุโน เสนาสนํ ปฏิพาหิตพฺพนฺติ วทามิ, โย ปฏิพาเหยฺย, อาปตฺติ ทุกฺกฏสฺสาติ.}

\prose{อามิเส ปน ภนฺเต กถํ ปฏิปชฺชิตพฺพนฺติ. อามิสํ โข สาริปุตฺต สพฺเพสํ สมกํ ภาเชตพฺพนฺติ.}

\csromanmatikaanchor{mtk.287}
\csromantocmarkref{section}{277. โอสารณานุชานนา}{mtk.287}
\bookmark[dest={mtk.287},level=1]{277. โอสารณานุชานนา}
\csromanheader{1}{277. \textbf{โอสารณา}\-\textbf{นุชานนา}}
\proseitem{474}{\csromanfolio{506}อถ โข ตสฺส อุกฺขิตฺตกสฺส ภิกฺขุโน ธมฺมญฺจ วินยญฺจ ปจฺจเวกฺขนฺตสฺส เอตทโหสิ “อาปตฺติ เอสา เนสา อนาปตฺติ, อาปนฺโนมฺหิ นมฺหิ อนาปนฺโน, อุกฺขิตฺโตมฺหิ นมฺหิ อนุกฺขิตฺโต, ธมฺมิเกนมฺหิ กมฺเมน อุกฺขิตฺโต อกุปฺเปน ฐานารเหนา”ติ. อถ โข โส อุกฺขิตฺตโก ภิกฺขุ เยน อุกฺขิตฺตา\-นุวตฺตกา ภิกฺขู เตนุปสงฺกมิ, อุปสงฺกมิตฺวา อุกฺขิตฺตา\-นุวตฺตเก ภิกฺขู เอตทโวจ “อาปตฺติ เอสา อาวุโส เนสา อนาปตฺติ, อาปนฺโนมฺหิ นมฺหิ อนาปนฺโน, อุกฺขิตฺโตมฺหิ นมฺหิ อนุกฺขิตฺโต, ธมฺมิเกนมฺหิ กมฺเมน อุกฺขิตฺโต อกุปฺเปน ฐานารเหน, เอถ มํ อายสฺมนฺโต โอสาเรถา”ติ. อถ โข เต อุกฺขิตฺตา\-นุวตฺตกา ภิกฺขู ตํ อุกฺขิตฺตกํ ภิกฺขุํ อาทาย เยน ภควา เตนุ\-ปสงฺกมิํสุ, อุปสงฺกมิตฺวา ภควนฺตํ อภิวาเทตฺวา เอกมนฺตํ นิสีทิํสุ, เอกมนฺตํ นิสินฺนา โข เต ภิกฺขู ภควนฺตํ เอตทโวจุํ “อยํ ภนฺเต อุกฺขิตฺตโก ภิกฺขุ เอวมาห 'อาปตฺติ เอสา อาวุโส เนสา อนาปตฺติ, อาปนฺโนมฺหิ นมฺหิ อนาปนฺโน, อุกฺขิตฺโตมฺหิ นมฺหิ อนุกฺขิตฺโต, ธมฺมิเกนมฺหิ กมฺเมน อุกฺขิตฺโต อกุปฺเปน ฐานารเหน, เอถ มํ อายสฺมนฺโต โอสาเรถา’ติ, กถํ นุ โข ภนฺเต ปฏิปชฺชิตพฺพนฺ”ติ. อาปตฺติ เอสา ภิกฺขเว เนสา อนาปตฺติ, อาปนฺโน เอโส ภิกฺขุ เนโส ภิกฺขุ อนาปนฺโน, อุกฺขิตฺโต เอโส ภิกฺขุ เนโส ภิกฺขุ อนุกฺขิตฺโต, ธมฺมิเกน กมฺเมน อุกฺขิตฺโต อกุปฺเปน ฐานารเหน, ยโต จ โข โส ภิกฺขเว ภิกฺขุ อาปนฺโน จ อุกฺขิตฺโต จ ปสฺสติ จ, เตน หิ ภิกฺขเว ตํ ภิกฺขุํ โอสาเรถาติ.}

\csromanmatikaanchor{mtk.288}
\csromantocmarkref{section}{278. สํฆสามคฺคีกถา}{mtk.288}
\bookmark[dest={mtk.288},level=1]{278. สํฆสามคฺคีกถา}
\csromanheader{1}{278. \textbf{สํฆสามคฺคีกถา}}
\proseitem{475}{อถ โข เต อุกฺขิตฺตา\-นุวตฺตกา ภิกฺขู ตํ อุกฺขิตฺตกํ ภิกฺขุํ โอสาเรตฺวา เยน อุกฺเขปกา ภิกฺขู เตนุ\-ปสงฺกมิํสุ, อุปสงฺกมิตฺวา อุกฺเขปเก ภิกฺขู เอตทโวจุํ “ยสฺมิํ อาวุโส วตฺถุสฺมิํ อโหสิ สํฆสฺส ภณฺฑนํ กลโห วิคฺคโห วิวาโท สํฆเภโท สํฆราชิ สํฆววตฺถานํ สํฆนานากรณํ, โส เอโส ภิกฺขุ อาปนฺโน จ อุกฺขิตฺโต จ ปสฺสิ\csromansharedfootnote{f28daa37034b1ac5}{ปสฺสี (อิติปิ)} จ โอสาริโต จ, หนฺท มยํ อาวุโส ตสฺส วตฺถุสฺส วูปสมาย สํฆสามคฺคิํ กโรมา”ติ.}

\prose{\csromanfolio{507}อถ โข เต อุกฺเขปกา ภิกฺขู เยน ภควา เตนุ\-ปสงฺกมิํสุ, อุปสงฺกมิตฺวา ภควนฺตํ อภิวาเทตฺวา เอกมนฺตํ นิสีทิํสุ, เอกมนฺตํ นิสินฺนา โข เต ภิกฺขู ภควนฺตํ เอตทโวจุํ “เต ภนฺเต อุกฺขิตฺตา\-นุวตฺตกา ภิกฺขู เอวมาหํสุ ‘ยสฺมิํ อาวุโส วตฺถุสฺมิํ อโหสิ สํฆสฺส ภณฺฑนํ กลโห วิคฺคโห วิวาโท สํฆเภโท สํฆราชิ สํฆววตฺถานํ สํฆนานากรณํ, โส เอโส ภิกฺขุ อาปนฺโน จ อุกฺขิตฺโต จ ปสฺสิ จ โอสาริโต จ, หนฺท มยํ อาวุโส ตสฺส วตฺถุสฺส วูปสมาย สํฆสามคฺคิํ กโรมา’ติ, กถํ นุ โข ภนฺเต ปฏิปชฺชิตพฺพนฺ”ติ. ยโต จ โข โส ภิกฺขเว ภิกฺขุ อาปนฺโน จ อุกฺขิตฺโต จ ปสฺสิ จ โอสาริโต จ, เตน หิ ภิกฺขเว สํโฆ ตสฺส วตฺถุสฺส วูปสมาย สํฆสามคฺคิํ กโรตุ. เอวญฺจ ปน ภิกฺขเว กาตพฺพา, สพฺเพเหว เอกชฺฌํ สนฺนิปติตพฺพํ คิลาเนหิ จ อคิลาเนหิ จ, น เกหิจิ ฉนฺโท ทาตพฺโพ, สนฺนิปติตฺวา พฺยตฺเตน ภิกฺขุนา ปฏิพเลน สํโฆ ญาเปตพฺโพ\csromandash{}}

\prose{“สุณาตุ เม ภนฺเต สํโฆ, ยสฺมิํ วตฺถุสฺมิํ อโหสิ สํฆสฺส ภณฺฑนํ กลโห วิคฺคโห วิวาโท สํฆเภโท สํฆราชิ สํฆววตฺถานํ สํฆนานากรณํ, โส เอโส ภิกฺขุ อาปนฺโน จ อุกฺขิตฺโต จ ปสฺสิ จ โอสาริโต จ, ยทิ สํฆสฺส ปตฺตกลฺลํ, สํโฆ ตสฺส วตฺถุสฺส วูปสมาย สํฆสามคฺคิํ กเรยฺย, เอสา ญตฺติ.}

\prose{สุณาตุ เม ภนฺเต สํโฆ, ยสฺมิํ วตฺถุสฺมิํ อโหสิ สํฆสฺส ภณฺฑนํ กลโห วิคฺคโห วิวาโท สํฆเภโท สํฆราชิ สํฆววตฺถานํ สํฆนานากรณํ, โส เอโส ภิกฺขุ อาปนฺโน จ อุกฺขิตฺโต จ ปสฺสิ จ โอสาริโต จ, สํโฆ ตสฺส วตฺถุสฺส วูปสมาย สํฆสามคฺคิํ กโรติ, ยสฺสายสฺมโต ขมติ ตสฺส วตฺถุสฺส วูปสมาย สํฆสามคฺคิยา กรณํ, โส ตุณฺหสฺส, ยสฺส นกฺขมติ, โส ภาเสยฺย.}

\prose{กตา สํเฆน ตสฺส วตฺถุสฺส วูปสมาย สํฆสามคฺคี, นิหโต สํฆเภโท นิหตา สํฆราชิ นิหตํ สํฆววตฺถานํ นิหตํ สํฆนานากรณํ, ขมติ สํฆสฺส, ตสฺมา ตุณฺหี, เอวเมตํ ธารยามี”ติ.}

\setlength{\csromangathapagewidth}{0pt}
\setlength{\csromangathawakwidth}{0pt}
\csromangathameasuregroup{ตาวเทว อุโปสโถ กาตพฺโพ,}{\csromangathabat{ตาวเทว อุโปสโถ กาตพฺโพ,}{ปาติโมกฺขํ อุทฺทิสิตพฺพนฺติ.}}
\csromangathasetleft{ตาวเทว อุโปสโถ กาตพฺโพ,}
\csromangathagroup{\csromangathastanza{\csromangathabat{ตาวเทว อุโปสโถ กาตพฺโพ,}{ปาติโมกฺขํ อุทฺทิสิตพฺพนฺติ.}}}

\csromanmatikaanchor{mtk.289}
\csromantocmarkref{section}{279. อุปาลิสํฆสามคฺคีปุจฺฉา}{mtk.289}
\bookmark[dest={mtk.289},level=1]{279. อุปาลิสํฆสามคฺคีปุจฺฉา}
\csromanheader{1}{279. \textbf{อุปาลิสํฆสามคฺคีปุจฺฉา}}
\proseitem{476}{\csromanfolio{508}อถ โข อายสฺมา อุปาลิ เยน ภควา เตนุปสงฺกมิ, อุปสงฺกมิตฺวา ภควนฺตํ อภิวาเทตฺวา เอกมนฺตํ นิสีทิ, เอกมนฺตํ นิสินฺโน โข อายสฺมา อุปาลิ ภควนฺตํ เอตทโวจ “ยสฺมิํ ภนฺเต วตฺถุสฺมิํ โหติ สํฆสฺส ภณฺฑนํ กลโห วิคฺคโห วิวาโท สํฆเภโท สํฆราชิ สํฆววตฺถานํ สํฆนานากรณํ, สํโฆ ตํ วตฺถุํ อวินิจฺฉินิตฺวา อมูลา มูลํ คนฺตฺวา สํฆสามคฺคิํ กโรติ, ธมฺมิกา นุ โข สา ภนฺเต สํฆสามคฺคี”ติ.}

\prose{ยสฺมิํ อุปาลิ วตฺถุสฺมิํ โหติ สํฆสฺส ภณฺฑนํ กลโห วิคฺคโห วิวาโท สํฆเภโท สํฆราชิ สํฆววตฺถานํ สํฆนานากรณํ, สํโฆ ตํ วตฺถุํ อวินิจฺฉินิตฺวา อมูลา มูลํ คนฺตฺวา สํฆสามคฺคิํ กโรติ, อธมฺมิกา สา อุปาลิ สํฆสามคฺคีติ.}

\prose{ยสฺมิํ ปน ภนฺเต วตฺถุสฺมิํ โหติ สํฆสฺส ภณฺฑนํ กลโห วิคฺคโห วิวาโท สํฆเภโท สํฆราชิ สํฆววตฺถานํ สํฆนานากรณํ, สํโฆ ตํ วตฺถุํ วินิจฺฉินิตฺวา มูลา มูลํ คนฺตฺวา สํฆสามคฺคิํ กโรติ, ธมฺมิกา นุ โข สา ภนฺเต สํฆสามคฺคีติ.}

\prose{ยสฺมิํ อุปาลิ วตฺถุสฺมิํ โหติ สํฆสฺส ภณฺฑนํ กลโห วิคฺคโห วิวาโท สํฆเภโท สํฆราชิ สํฆววตฺถานํ สํฆนานากรณํ, สํโฆ ตํ วตฺถุํ วินิจฺฉินิตฺวา มูลา มูลํ คนฺตฺวา สํฆสามคฺคิํ กโรติ, ธมฺมิกา สา อุปาลิ สํฆสามคฺคีติ.}

\prose{กติ นุ โข ภนฺเต สํฆสามคฺคิโยติ. เทฺวมา อุปาลิ สํฆสามคฺคิโย, อตฺถุปาลิ สํฆสามคฺคี อตฺถาเปตา พฺยญฺชนุเปตา, อตฺถุปาลิ สํฆสามคฺคี อตฺถุเปตา จ พฺยญฺชนุเปตา จ. กตมา จ อุปาลิ สํฆสามคฺคี อตฺถาเปตา พฺยญฺชนุเปตา. ยสฺมิํ อุปาลิ วตฺถุสฺมิํ โหติ สํฆสฺส ภณฺฑนํ กลโห วิคฺคโห วิวาโท สํฆเภโท สํฆราชิ สํฆววตฺถานํ สํฆนานากรณํ, สํโฆ ตํ วตฺถุํ อวินิจฺฉินิตฺวา อมูลา มูลํ คนฺตฺวา สํฆสามคฺคิํ กโรติ, อยํ วุจฺจติ อุปาลิ สํฆสามคฺคี อตฺถาเปตา พฺยญฺชนุเปตา. กตมา จ อุปาลิ สํฆสามคฺคี อตฺถุเปตา จ พฺยญฺชนุเปตา จ. ยสฺมิํ อุปาลิ วตฺถุสฺมิํ โหติ สํฆสฺส ภณฺฑนํ กลโห วิคฺคโห วิวาโท สํฆเภโท สํฆราชิ สํฆววตฺถานํ สํฆนานากรณํ, สํโฆ ตํ วตฺถุํ วินิจฺฉินิตฺวา มูลา มูลํ คนฺตฺวา สํฆสามคฺคิํ กโรติ, อยํ วุจฺจติ อุปาลิ สํฆสามคฺคี อตฺถุเปตา จ พฺยญฺชนุเปตา จ. อิมา โข อุปาลิ เทฺว สํฆสามคฺคิโยติ.}

\proseitem{477}{\csromanfolio{509}อถ โข อายสฺมา อุปาลิ อุฏฺฐายาสนา เอกํสํ อุตฺตราสงฺคํ กริตฺวา เยน ภควา เตนญฺชลิํ ปณาเมตฺวา ภควนฺตํ คาถาย อชฺฌภาสิ.}

\csromanmatikaanchor{mtk.290}
\csromantocmarkref{section}{280. อุทฺทานคาถา}{mtk.290}
\bookmark[dest={mtk.290},level=1]{280. อุทฺทานคาถา}
\setlength{\csromangathapagewidth}{0pt}
\setlength{\csromangathawakwidth}{0pt}
\csromangathameasuregroup{}{สํฆสฺส กิจฺเจสุ จ มนฺตนาสุ จ, \\ อตฺเถสุ ชาเตสุ วินิจฺฉเยสุ จ. \\ กถํปกาโรธ นโร มหตฺถิโก, \\ ภิกฺขุ กถํ โหติธ ปคฺคหารโหติ. \\ อนานุวชฺโช ปฐเมน สีลโต, \\ อเวกฺขิตาจาโร สุสํวุตินฺทฺริโย. \\ ปจฺจตฺถิกา นูปวทนฺติ ธมฺมโต, \\ น หิสฺส ตํ โหติ วเทยฺยุ เยน นํ. \\ โส ตาทิโส สีลวิสุทฺธิยา ฐิโต, \\ วิสารโท โหติ วิสยฺห ภาสติ. \\ นจฺฉมฺภติ ปริสคโต น เวธติ, \\ อตฺถํ น หาเปติ อนุยฺยุตํ ภณํ. \\ ตเถว ปญฺหํ ปริสาสุ ปุจฺฉิโต, \\ น เจว ปชฺฌายติ น มงฺกุ โหติ. \\ โส กาลาคตํ พฺยากรณารหํ วโจ, \\ รญฺเชติ วิญฺญูปริสํ วิจกฺขโณ. \\ สคารโว วุฑฺฒตเรสุ ภิกฺขุสุ, \\ อาเจรกมฺหิ จ สเก วิสารโท. \\ อลํ ปเมตุํ ปคุโณ กเถตเว, \\ ปจฺจตฺถิกานญฺจ วิรทฺธิโกวิโท. \\ ปจฺจตฺถิกา เยน วชนฺติ นิคฺคหํ, \\ มหาชโน สญฺญปนญฺจ คจฺฉติ. \\ สกญฺจ อาทายมยํ น ริญฺจติ, \\ วิยากรํ\textsuperscript{1} ปญฺหมนุปฆาติกํ. \\ ทูเตยฺยกมฺเมสุ อลํ สมุคฺคโห, \\ สํฆสฺส กิจฺเจสุ จ อาหุ นํ ยถา. \\ กรํ วโจ ภิกฺขุคเณน เปสิโต, \\ “อหํ กโรมี”ติ น เตน มญฺญติ. \\ อาปชฺชติ ยาวตเกสุ วตฺถุสุ, \\ อาปตฺติยา โหติ ยถา จ วุฏฺฐิติ. \\ เอเต วิภงฺคา อุภยสฺส สฺวาคตา, \\ อาปตฺติ วุฏฺฐานปทสฺส โกวิโท. \\ นิสฺสารณํ คจฺฉติ ยานิ จาจรํ, \\ นิสฺสาริโต โหติ ยถา จ วตฺตนา\textsuperscript{1}. \\ โอสารณํ ตํวุสิตสฺส ชนฺตุโน, \\ เอตมฺปิ ชานาติ วิภงฺคโกวิโท. \\ สคารโว วุฑฺฒตเรสุ ภิกฺขุสุ, \\ นเวสุ เถเรสุ จ มชฺฌิเมสุ จ. \\ มหาชนสฺสตฺถจโรธ ปณฺฑิโต, \\ โส ตาทิโส ภิกฺขุ อิธ ปคฺคหารโหติ.}
\csromangathameasurewak{สํฆสฺส กิจฺเจสุ จ มนฺตนาสุ จ, \\ อตฺเถสุ ชาเตสุ วินิจฺฉเยสุ จ. \\ กถํปกาโรธ นโร มหตฺถิโก, \\ ภิกฺขุ กถํ โหติธ ปคฺคหารโหติ. \\ อนานุวชฺโช ปฐเมน สีลโต, \\ อเวกฺขิตาจาโร สุสํวุตินฺทฺริโย. \\ ปจฺจตฺถิกา นูปวทนฺติ ธมฺมโต, \\ น หิสฺส ตํ โหติ วเทยฺยุ เยน นํ. \\ โส ตาทิโส สีลวิสุทฺธิยา ฐิโต, \\ วิสารโท โหติ วิสยฺห ภาสติ. \\ นจฺฉมฺภติ ปริสคโต น เวธติ, \\ อตฺถํ น หาเปติ อนุยฺยุตํ ภณํ. \\ ตเถว ปญฺหํ ปริสาสุ ปุจฺฉิโต, \\ น เจว ปชฺฌายติ น มงฺกุ โหติ. \\ โส กาลาคตํ พฺยากรณารหํ วโจ, \\ รญฺเชติ วิญฺญูปริสํ วิจกฺขโณ. \\ สคารโว วุฑฺฒตเรสุ ภิกฺขุสุ, \\ อาเจรกมฺหิ จ สเก วิสารโท. \\ อลํ ปเมตุํ ปคุโณ กเถตเว, \\ ปจฺจตฺถิกานญฺจ วิรทฺธิโกวิโท. \\ ปจฺจตฺถิกา เยน วชนฺติ นิคฺคหํ, \\ มหาชโน สญฺญปนญฺจ คจฺฉติ. \\ สกญฺจ อาทายมยํ น ริญฺจติ, \\ วิยากรํ\textsuperscript{1} ปญฺหมนุปฆาติกํ. \\ ทูเตยฺยกมฺเมสุ อลํ สมุคฺคโห, \\ สํฆสฺส กิจฺเจสุ จ อาหุ นํ ยถา. \\ กรํ วโจ ภิกฺขุคเณน เปสิโต, \\ “อหํ กโรมี”ติ น เตน มญฺญติ. \\ อาปชฺชติ ยาวตเกสุ วตฺถุสุ, \\ อาปตฺติยา โหติ ยถา จ วุฏฺฐิติ. \\ เอเต วิภงฺคา อุภยสฺส สฺวาคตา, \\ อาปตฺติ วุฏฺฐานปทสฺส โกวิโท. \\ นิสฺสารณํ คจฺฉติ ยานิ จาจรํ, \\ นิสฺสาริโต โหติ ยถา จ วตฺตนา\textsuperscript{1}. \\ โอสารณํ ตํวุสิตสฺส ชนฺตุโน, \\ เอตมฺปิ ชานาติ วิภงฺคโกวิโท. \\ สคารโว วุฑฺฒตเรสุ ภิกฺขุสุ, \\ นเวสุ เถเรสุ จ มชฺฌิเมสุ จ. \\ มหาชนสฺสตฺถจโรธ ปณฺฑิโต, \\ โส ตาทิโส ภิกฺขุ อิธ ปคฺคหารโหติ.}
\setlength{\csromangathaleftcol}{0pt}
\csromangathagroup{\csromangathastanza{\csromangathawak{สํฆสฺส กิจฺเจสุ จ มนฺตนาสุ จ, \\ อตฺเถสุ ชาเตสุ วินิจฺฉเยสุ จ. \\ กถํปกาโรธ นโร มหตฺถิโก, \\ ภิกฺขุ กถํ โหติธ ปคฺคหารโหติ.}}\csromangathastanza{\csromangathawak{อนานุวชฺโช ปฐเมน สีลโต, \\ อเวกฺขิตาจาโร สุสํวุตินฺทฺริโย. \\ ปจฺจตฺถิกา นูปวทนฺติ ธมฺมโต, \\ น หิสฺส ตํ โหติ วเทยฺยุ เยน นํ.}}\csromangathastanza{\csromangathawak{โส ตาทิโส สีลวิสุทฺธิยา ฐิโต, \\ วิสารโท โหติ วิสยฺห ภาสติ. \\ นจฺฉมฺภติ ปริสคโต น เวธติ, \\ อตฺถํ น หาเปติ อนุยฺยุตํ ภณํ.}}\csromangathastanza{\csromangathawak{ตเถว ปญฺหํ ปริสาสุ ปุจฺฉิโต, \\ น เจว ปชฺฌายติ น มงฺกุ โหติ. \\ โส กาลาคตํ พฺยากรณารหํ วโจ, \\ รญฺเชติ วิญฺญูปริสํ วิจกฺขโณ.}}\csromangathastanza{\csromangathawak{สคารโว วุฑฺฒตเรสุ ภิกฺขุสุ, \\ อาเจรกมฺหิ จ สเก วิสารโท. \\ อลํ ปเมตุํ ปคุโณ กเถตเว, \\ ปจฺจตฺถิกานญฺจ วิรทฺธิโกวิโท.}}\csromangathastanza{\csromangathawak{ปจฺจตฺถิกา เยน วชนฺติ นิคฺคหํ, \\ มหาชโน สญฺญปนญฺจ คจฺฉติ. \\ สกญฺจ อาทายมยํ น ริญฺจติ, \\ วิยากรํ\csromansharedfootnote{170eea6ce79b8767}{โส พฺยากรํ (สี), เวยฺยากรํ (สฺยา)} ปญฺหมนุปฆาติกํ.}}\csromangathastanza{\csromangathawak{\csromanfolio{510}ทูเตยฺยกมฺเมสุ อลํ สมุคฺคโห, \\ สํฆสฺส กิจฺเจสุ จ อาหุ นํ ยถา. \\ กรํ วโจ ภิกฺขุคเณน เปสิโต, \\ “อหํ กโรมี”ติ น เตน มญฺญติ.}}\csromangathastanza{\csromangathawak{อาปชฺชติ ยาวตเกสุ วตฺถุสุ, \\ อาปตฺติยา โหติ ยถา จ วุฏฺฐิติ. \\ เอเต วิภงฺคา อุภยสฺส สฺวาคตา, \\ อาปตฺติ วุฏฺฐานปทสฺส โกวิโท.}}\csromangathastanza{\csromangathawak{นิสฺสารณํ คจฺฉติ ยานิ จาจรํ, \\ นิสฺสาริโต โหติ ยถา จ วตฺตนา\csromansharedfootnote{faf7b918297d05e0}{วตฺถุนา (สี, สฺยา)}. \\ โอสารณํ ตํวุสิตสฺส ชนฺตุโน, \\ เอตมฺปิ ชานาติ วิภงฺคโกวิโท.}}\csromangathastanza{\csromangathawak{สคารโว วุฑฺฒตเรสุ ภิกฺขุสุ, \\ นเวสุ เถเรสุ จ มชฺฌิเมสุ จ. \\ มหาชนสฺสตฺถจโรธ ปณฺฑิโต, \\ โส ตาทิโส ภิกฺขุ อิธ ปคฺคหารโหติ.}}}

\vspace{\tipitakaparskip}
\csromancenter{โกสมฺพก\-กฺขนฺธโก ทสโม.}
\csromancenter{280. \textbf{ตสฺสุทฺทานํ}}
\setlength{\csromangathapagewidth}{0pt}
\setlength{\csromangathawakwidth}{0pt}
\csromangathameasuregroup{โกสมฺพิยํ ชินวโร, \\ นุกฺขิเปยฺย ยสฺมิํ ตสฺมิํ, \\ อนฺโตสีมายํ ตตฺเถว, \\ ปาลิเลยฺยา จ สาวตฺถิ, \\ มหากสฺสปกจฺจานา, \\ มหาจุนฺโท อนุรุทฺโธ, \\ อานนฺโท ราหุโล เจว, \\ เสนาสนํ วิวิตฺตญฺจ, \\ น เกหิ ฉนฺโท ทาตพฺโพ, \\ อนานุวชฺโช สีเลน,}{\csromangathabat{โกสมฺพิยํ ชินวโร,}{วิวาทาปตฺติทสฺสเน.} \\ \csromangathabat{นุกฺขิเปยฺย ยสฺมิํ ตสฺมิํ,}{สทฺธายาปตฺติ เทสเย.} \\ \csromangathabat{อนฺโตสีมายํ ตตฺเถว,}{พาลกญฺเจว วํสทา.} \\ \csromangathabat{ปาลิเลยฺยา จ สาวตฺถิ,}{สาริปุตฺโต จ โกลิโต.} \\ \csromangathabat{มหากสฺสปกจฺจานา,}{โกฏฺฐิโก กปฺปิเนน จ.} \\ \csromangathabat{มหาจุนฺโท อนุรุทฺโธ,}{เรวโต อุปาลิ จุโภ.} \\ \csromangathabat{อานนฺโท ราหุโล เจว,}{โคตมีนาถปิณฺฑิโก.} \\ \csromangathabat{เสนาสนํ วิวิตฺตญฺจ,}{อามิสํ สมกมฺปิ จ.} \\ \csromangathabat{น เกหิ ฉนฺโท ทาตพฺโพ,}{อุปาลิปริปุจฺฉิโต.} \\ \csromangathabat{อนานุวชฺโช สีเลน,}{สามคฺคี ชินสาสเนติ.}}
\csromangathasetleft{โกสมฺพิยํ ชินวโร, \\ นุกฺขิเปยฺย ยสฺมิํ ตสฺมิํ, \\ อนฺโตสีมายํ ตตฺเถว, \\ ปาลิเลยฺยา จ สาวตฺถิ, \\ มหากสฺสปกจฺจานา, \\ มหาจุนฺโท อนุรุทฺโธ, \\ อานนฺโท ราหุโล เจว, \\ เสนาสนํ วิวิตฺตญฺจ, \\ น เกหิ ฉนฺโท ทาตพฺโพ, \\ อนานุวชฺโช สีเลน,}
\csromangathagroupwithclosermajor{\csromangathastanza{\csromangathabat{โกสมฺพิยํ ชินวโร,}{วิวาทาปตฺติทสฺสเน.} \\ \csromangathabat{นุกฺขิเปยฺย ยสฺมิํ ตสฺมิํ,}{สทฺธายาปตฺติ เทสเย.}}\csromangathastanza{\csromangathabat{อนฺโตสีมายํ ตตฺเถว,}{พาลกญฺเจว วํสทา.} \\ \csromangathabat{ปาลิเลยฺยา จ สาวตฺถิ,}{สาริปุตฺโต จ โกลิโต.}}\csromangathastanza{\csromangathabat{มหากสฺสปกจฺจานา,}{โกฏฺฐิโก กปฺปิเนน จ.} \\ \csromangathabat{มหาจุนฺโท อนุรุทฺโธ,}{เรวโต อุปาลิ จุโภ.}}\csromangathastanza{\csromanfolio{511}\csromangathabat{อานนฺโท ราหุโล เจว,}{โคตมีนาถปิณฺฑิโก.} \\ \csromangathabat{เสนาสนํ วิวิตฺตญฺจ,}{อามิสํ สมกมฺปิ จ.}}\csromangathastanza{\csromangathabat{น เกหิ ฉนฺโท ทาตพฺโพ,}{อุปาลิปริปุจฺฉิโต.} \\ \csromangathabat{อนานุวชฺโช สีเลน,}{สามคฺคี ชินสาสเนติ.}}}{โกสมฺพก\-กฺขนฺธโก นิฏฺฐิโต.}
\nitthitammajor{\textbf{มหาวคฺคปาฬิ นิฏฺฐิตา.}}
\orphannote{(เอวํ วุตฺเต อายสฺมา มหากจฺจายโน โสณํ อุปาสกํ กุฏิกณฺณํ เอตทโวจ) (สฺยา, ขุ 1. 45)}

\orphannote{ปจฺฉิมา ปญฺจ อานิสํสา อํ 2. 218 ปิฏฺเฐปิ อาคตา.}

